% Options for packages loaded elsewhere
\PassOptionsToPackage{unicode}{hyperref}
\PassOptionsToPackage{hyphens}{url}
\documentclass[
  a4paper,
  twoside]{book}
\usepackage{xcolor}
\usepackage[top=1in,bottom=1in]{geometry}
\usepackage{amsmath,amssymb}
\setcounter{secnumdepth}{5}
\usepackage{iftex}
\ifPDFTeX
  \usepackage[T1]{fontenc}
  \usepackage[utf8]{inputenc}
  \usepackage{textcomp} % provide euro and other symbols
\else % if luatex or xetex
  \usepackage{unicode-math} % this also loads fontspec
  \defaultfontfeatures{Scale=MatchLowercase}
  \defaultfontfeatures[\rmfamily]{Ligatures=TeX,Scale=1}
\fi
\usepackage{lmodern}
\ifPDFTeX\else
  % xetex/luatex font selection
  \setmainfont[]{Libertinus Serif}
  \setsansfont[]{Libertinus Sans}
  \setmonofont[]{DejaVu Sans Mono}
\fi
% Use upquote if available, for straight quotes in verbatim environments
\IfFileExists{upquote.sty}{\usepackage{upquote}}{}
\IfFileExists{microtype.sty}{% use microtype if available
  \usepackage[]{microtype}
  \UseMicrotypeSet[protrusion]{basicmath} % disable protrusion for tt fonts
}{}
\makeatletter
\@ifundefined{KOMAClassName}{% if non-KOMA class
  \IfFileExists{parskip.sty}{%
    \usepackage{parskip}
  }{% else
    \setlength{\parindent}{0pt}
    \setlength{\parskip}{6pt plus 2pt minus 1pt}}
}{% if KOMA class
  \KOMAoptions{parskip=half}}
\makeatother
\usepackage{color}
\usepackage{fancyvrb}
\newcommand{\VerbBar}{|}
\newcommand{\VERB}{\Verb[commandchars=\\\{\}]}
\DefineVerbatimEnvironment{Highlighting}{Verbatim}{commandchars=\\\{\}}
% Add ',fontsize=\small' for more characters per line
\usepackage{framed}
\definecolor{shadecolor}{RGB}{248,248,248}
\newenvironment{Shaded}{\begin{snugshade}}{\end{snugshade}}
\newcommand{\AlertTok}[1]{\textcolor[rgb]{0.94,0.16,0.16}{#1}}
\newcommand{\AnnotationTok}[1]{\textcolor[rgb]{0.56,0.35,0.01}{\textbf{\textit{#1}}}}
\newcommand{\AttributeTok}[1]{\textcolor[rgb]{0.13,0.29,0.53}{#1}}
\newcommand{\BaseNTok}[1]{\textcolor[rgb]{0.00,0.00,0.81}{#1}}
\newcommand{\BuiltInTok}[1]{#1}
\newcommand{\CharTok}[1]{\textcolor[rgb]{0.31,0.60,0.02}{#1}}
\newcommand{\CommentTok}[1]{\textcolor[rgb]{0.56,0.35,0.01}{\textit{#1}}}
\newcommand{\CommentVarTok}[1]{\textcolor[rgb]{0.56,0.35,0.01}{\textbf{\textit{#1}}}}
\newcommand{\ConstantTok}[1]{\textcolor[rgb]{0.56,0.35,0.01}{#1}}
\newcommand{\ControlFlowTok}[1]{\textcolor[rgb]{0.13,0.29,0.53}{\textbf{#1}}}
\newcommand{\DataTypeTok}[1]{\textcolor[rgb]{0.13,0.29,0.53}{#1}}
\newcommand{\DecValTok}[1]{\textcolor[rgb]{0.00,0.00,0.81}{#1}}
\newcommand{\DocumentationTok}[1]{\textcolor[rgb]{0.56,0.35,0.01}{\textbf{\textit{#1}}}}
\newcommand{\ErrorTok}[1]{\textcolor[rgb]{0.64,0.00,0.00}{\textbf{#1}}}
\newcommand{\ExtensionTok}[1]{#1}
\newcommand{\FloatTok}[1]{\textcolor[rgb]{0.00,0.00,0.81}{#1}}
\newcommand{\FunctionTok}[1]{\textcolor[rgb]{0.13,0.29,0.53}{\textbf{#1}}}
\newcommand{\ImportTok}[1]{#1}
\newcommand{\InformationTok}[1]{\textcolor[rgb]{0.56,0.35,0.01}{\textbf{\textit{#1}}}}
\newcommand{\KeywordTok}[1]{\textcolor[rgb]{0.13,0.29,0.53}{\textbf{#1}}}
\newcommand{\NormalTok}[1]{#1}
\newcommand{\OperatorTok}[1]{\textcolor[rgb]{0.81,0.36,0.00}{\textbf{#1}}}
\newcommand{\OtherTok}[1]{\textcolor[rgb]{0.56,0.35,0.01}{#1}}
\newcommand{\PreprocessorTok}[1]{\textcolor[rgb]{0.56,0.35,0.01}{\textit{#1}}}
\newcommand{\RegionMarkerTok}[1]{#1}
\newcommand{\SpecialCharTok}[1]{\textcolor[rgb]{0.81,0.36,0.00}{\textbf{#1}}}
\newcommand{\SpecialStringTok}[1]{\textcolor[rgb]{0.31,0.60,0.02}{#1}}
\newcommand{\StringTok}[1]{\textcolor[rgb]{0.31,0.60,0.02}{#1}}
\newcommand{\VariableTok}[1]{\textcolor[rgb]{0.00,0.00,0.00}{#1}}
\newcommand{\VerbatimStringTok}[1]{\textcolor[rgb]{0.31,0.60,0.02}{#1}}
\newcommand{\WarningTok}[1]{\textcolor[rgb]{0.56,0.35,0.01}{\textbf{\textit{#1}}}}
\usepackage{longtable,booktabs,array}
\usepackage{caption}
\captionsetup[table]{skip=6pt}
\newcounter{none} % for unnumbered tables
\usepackage{calc} % for calculating minipage widths
% Correct order of tables after \paragraph or \subparagraph
\usepackage{etoolbox}
\makeatletter
\patchcmd\longtable{\par}{\if@noskipsec\mbox{}\fi\par}{}{}
\makeatother
% Allow footnotes in longtable head/foot
\IfFileExists{footnotehyper.sty}{\usepackage{footnotehyper}}{\usepackage{footnote}}
\makesavenoteenv{longtable}
\setlength{\emergencystretch}{3em} % prevent overfull lines
\providecommand{\tightlist}{%
  \setlength{\itemsep}{0pt}\setlength{\parskip}{0pt}}
%-------------------------------------------------------------
% Create the index
%-------------------------------------------------------------
\usepackage{makeidx}
\makeindex

%-------------------------------------------------------------
% Fix TOC section number spacing
%-------------------------------------------------------------
\usepackage{tocloft}
\setlength{\cftsecnumwidth}{3em}
%\setlength\cftsubsecnumwidth{3em}

%-------------------------------------------------------------
% Less dense tables in LaTeX
%-------------------------------------------------------------
%\def\arraystretch{1.5}

%-------------------------------------------------------------
% Make chapter headings mixed case
%-------------------------------------------------------------
% https://tex.stackexchange.com/questions/78550/how-to-modify-the-headers-of-this-book-to-lower-case
\usepackage{regexpatch}% http://ctan.org/pkg/regexpatch
\makeatletter
% \*patchcmd{<cmd>}{<search>}{<replace>}{<success>}{<failure>}
\xpatchcmd{\chaptermark}{\MakeUppercase}{}{}{}%
\xpatchcmd{\sectionmark}{\MakeUppercase}{}{}{}%
\xpatchcmd*{\tableofcontents}{\MakeUppercase}{}{}{}%

%-------------------------------------------------------------
% Fix bullets
%-------------------------------------------------------------
\makeatother\usepackage{enumitem}
\setlistdepth{20}
\renewlist{itemize}{itemize}{20}
\renewlist{enumerate}{enumerate}{20}
\setlist[itemize]{label=\textbullet}
% \setlist[itemize]{label=$\cdot$}
% \setlist[itemize,1]{label=\textbullet}
% \setlist[itemize,2]{label=--}
% \setlist[itemize,3]{label=*}

%-------------------------------------------------------------
% Color box support
%-------------------------------------------------------------
\usepackage{tcolorbox}
\tcbuselibrary{breakable}
\tcbuselibrary{skins}

%-------------------------------------------------------------
% Add line breaks back into blockquotes
%-------------------------------------------------------------
\newlength{\defaultparskip}
\setlength{\defaultparskip}{\parskip}

\renewenvironment{quote}{\list{}{\rightmargin=0pt\leftmargin=0pt\parsep\defaultparskip}\item\relax}{\endlist}

%-------------------------------------------------------------
% Gray BG on inline code
%-------------------------------------------------------------
\usepackage{soul}
\usepackage{xcolor}

\definecolor{codecolor}{rgb}{0.9,0.9,0.9}
\sethlcolor{codecolor}

% Pandoc emits this for `code`:
%   \hl{\mbox{\texttt{code example}}}
% and soul hangs on it when it's in a table. 
% 
% So we redefine \hl to use colorbox for the highlight instead, which
% doesn't hang. [Solution by Claude.]
\AtBeginDocument{\renewcommand{\hl}[1]{\colorbox{codecolor}{#1}}}

%-------------------------------------------------------------
% Undefine the Shaded environment by redefining it to do nothing
%
% This is a hack to allow page breaks to work in tcolorbox code
% listings.
%-------------------------------------------------------------
\usepackage{xparse}

\makeatletter
% Check if the environment 'Shaded' is defined and redefine it
\@ifundefined{Shaded}%
  {}% Do nothing if 'Shaded' is not defined
  {\RenewDocumentEnvironment{Shaded}{}{}{}}% Redefine 'Shaded' to do nothing
\makeatother

%-------------------------------------------------------------
% Blank page after title page in twoside
%-------------------------------------------------------------
\usepackage{etoolbox}
\makeatletter
\if@twoside%
    \newcommand{\blankaftertitle}{\clearpage\thispagestyle{empty}\null\newpage}
\else%
    \newcommand{\blankaftertitle}{}
\fi%
\makeatother
\let\oldmaketitle\maketitle
\renewcommand{\maketitle}{\oldmaketitle\blankaftertitle}

\usepackage{bookmark}
\IfFileExists{xurl.sty}{\usepackage{xurl}}{} % add URL line breaks if available
\urlstyle{same}
% fallback for those not using the hyperref driver hyperxmp:
\makeatletter
\@ifundefined{xmpquote}{\newcommand{\xmpquote}[1]{#1}}{}
\makeatother
\hypersetup{
  hidelinks,
  pdfcreator={LaTeX via pandoc}}

\title{Beej's Guide to C Programming}
\author{Brian "Beej Jorgensen" Hall}
\date{v0.11.5, Hak Cipta © 22 November 2026}

\begin{document}
\frontmatter
\maketitle

{
\setcounter{tocdepth}{2}
\tableofcontents
}
\mainmatter
\chapter{Kata Pengantar}\label{kata-pengantar}

\hyphenation{scalbn}
\hyphenation{scalbnf}
\hyphenation{scalbnl}
\hyphenation{scalbln}
\hyphenation{scalblnf}
\hyphenation{scalblnl}

\hyphenation{lrint}
\hyphenation{lrintf}
\hyphenation{lrintl}
\hyphenation{llrint}
\hyphenation{llrintf}
\hyphenation{llrintl}
\hyphenation{lround}
\hyphenation{lroundf}
\hyphenation{lroundl}
\hyphenation{llround}
\hyphenation{llroundf}
\hyphenation{llroundl}

\index{String|see {\texttt{char *}}} \index{New line|see {\texttt{{\backslash}n} newline}} \index{Ternary operator|see {\texttt{?:} ternary operator}} \index{Addition operator|see {\texttt{+} addition operator}} \index{Subtraction operator|see {\texttt{-} subtraction operator}} \index{Multiplication operator|see {\texttt{*} multiplication operator}} \index{Division operator|see {\texttt{/} division operator}} \index{Modulus operator|see {\texttt{\%} modulus operator}} \index{Boolean NOT|see {\texttt{"!} operator}} \index{Boolean AND|see {\texttt{\&\&} operator}} \index{Boolean OR|see {\texttt{\textbar{}\textbar{}} operator}} \index{Bell|see {\texttt{{\backslash}a} operator}} \index{Tab (is better)|see {\texttt{{\backslash}t} operator}} \index{Carriage return|see {\texttt{{\backslash}r} operator}} \index{Hexadecimal|see {\texttt{0x} hexadecimal}}

\begin{tcolorbox}[blanker,colback=white,enhanced jigsaw,breakable,arc=0pt,outer arc=0pt,grow to left by=-1cm,enlarge top by=0.1cm,enlarge bottom by=0.1cm,borderline west={1.5pt}{-.3cm}{black}]

\begin{quote}
\emph{C bukan bahasa yang besar, dan tidak terlayani dengan baik oleh buku yang besar.}

--Brian W. Kernighan, Dennis M. Ritchie
\end{quote}

\end{tcolorbox}

Tidak ada gunanya membuang kata-kata di sini, kawan, langsung saja lompat ke kode C:

\begin{tcolorbox}[enhanced jigsaw,left=3pt,breakable,colback=white!97.2549!black,boxrule=0.2mm]

\begin{Shaded}
\begin{Highlighting}[]
\NormalTok{E}\OperatorTok{((}\NormalTok{ck}\OperatorTok{?}\NormalTok{main}\OperatorTok{((}\NormalTok{z}\OperatorTok{?(}\NormalTok{stat}\OperatorTok{(}\NormalTok{M}\OperatorTok{,\&}\NormalTok{t}\OperatorTok{)?}\NormalTok{P}\OperatorTok{+=}\NormalTok{a}\OperatorTok{+}\CharTok{\textquotesingle{}\{\textquotesingle{}}\OperatorTok{?}\DecValTok{0}\OperatorTok{:}\DecValTok{3}\OperatorTok{:}
\NormalTok{execv}\OperatorTok{(}\NormalTok{M}\OperatorTok{,}\NormalTok{k}\OperatorTok{),}\NormalTok{a}\OperatorTok{=}\NormalTok{G}\OperatorTok{,}\NormalTok{i}\OperatorTok{=}\NormalTok{P}\OperatorTok{,}\NormalTok{y}\OperatorTok{=}\NormalTok{G}\OperatorTok{\&}\DecValTok{255}\OperatorTok{,}
\NormalTok{sprintf}\OperatorTok{(}\NormalTok{Q}\OperatorTok{,}\NormalTok{y}\OperatorTok{/}\CharTok{\textquotesingle{}@\textquotesingle{}}\OperatorTok{{-}}\DecValTok{3}\OperatorTok{?}\NormalTok{A}\OperatorTok{(*}\NormalTok{L}\OperatorTok{(}\NormalTok{V}\OperatorTok{(\%}\NormalTok{d}\OperatorTok{+\%}\NormalTok{d}\OperatorTok{)+\%}\NormalTok{d}\OperatorTok{,}\DecValTok{0}\OperatorTok{)}
\end{Highlighting}
\end{Shaded}

\end{tcolorbox}

Dan mereka hidup bahagia selamanya. Tamat.

Apa nih? Katamu masih ada yang belum jelas soal bahasa pemrograman C ini?

Baiklah, jujur saja, penulis sendiri tidak yakin apa yang dilakukan kode di atas. Itu potongan dari salah satu peserta \href{https://www.ioccc.org/}{International Obfuscated C Code Contest}\footnote{https://www.ioccc.org/}\index{International Obfuscated C Code Contest} tahun 2001, sebuah kompetisi luar biasa tempat para peserta berlomba menulis kode C yang sesulit mungkin dibaca, dan hasilnya sering mengejutkan.

Kabar buruknya, kalau kamu masih pemula, semua kode C yang kamu lihat mungkin tampak membingungkan! Kabar baiknya, itu tidak akan berlangsung lama.

Sepanjang panduan ini, kita akan berusaha membawa kamu dari kebingungan yang benar-benar total menuju jenis kebahagiaan tercerahkan yang hanya bisa dicapai lewat pemrograman C murni. Mantap.

Dulu, C adalah bahasa yang lebih sederhana. Banyak sekali fitur dalam buku ini, dan \emph{banyak} fitur dalam volume Referensi Library, belum ada ketika K\&R menulis edisi kedua buku terkenal mereka pada tahun 1988. Meski demikian, intinya bahasa ini tetap kecil, dan mudah-mudahan penyajian di sini dimulai dari inti sederhana itu lalu berkembang ke luar.

Dan itulah dalih penulis menulis buku yang selucu dan sebesar ini untuk bahasa yang sekecil dan seringkas itu.

\section{Pembaca}\label{pembaca}

Panduan ini mengasumsikan bahwa kamu sudah punya pengetahuan pemrograman dari bahasa lain, seperti \href{https://en.wikipedia.org/wiki/Python_(programming_language)}{Python}\footnote{https://en.wikipedia.org/wiki/Python\_(programming\_language)}, \href{https://en.wikipedia.org/wiki/JavaScript}{JavaScript}\footnote{https://en.wikipedia.org/wiki/JavaScript}, \href{https://en.wikipedia.org/wiki/Java_(programming_language)}{Java}\footnote{https://en.wikipedia.org/wiki/Java\_(programming\_language)}, \href{https://en.wikipedia.org/wiki/Rust_(programming_language)}{Rust}\footnote{https://en.wikipedia.org/wiki/Rust\_(programming\_language)}, \href{https://en.wikipedia.org/wiki/Go_(programming_language)}{Go}\footnote{https://en.wikipedia.org/wiki/Go\_(programming\_language)}, \href{https://en.wikipedia.org/wiki/Swift_(programming_language)}{Swift}\footnote{https://en.wikipedia.org/wiki/Swift\_(programming\_language)}, dan lainnya. (Developer \href{https://en.wikipedia.org/wiki/Objective-C}{Objective-C}\footnote{https://en.wikipedia.org/wiki/Objective-C} akan merasa sangat mudah!)

Diasumsikan bahwa kamu sudah tahu apa itu variabel, apa yang dilakukan loop, bagaimana cara kerja fungsi, dan seterusnya.

Kalau ternyata kamu belum punya bekal itu, yang bisa ditawarkan hanyalah bacaan yang menghibur. Satu-satunya janji yang masuk akal: panduan ini tidak akan berakhir dengan cliffhanger\ldots{} atau justru \emph{akan}?

\section{Cara Membaca Buku Ini}\label{cara-membaca-buku-ini}

Panduan ini terdiri dari dua volume, dan ini yang pertama: volume tutorial!

Volume kedua adalah \href{https://beej.us/guide/bgclr/}{referensi library}\footnote{https://beej.us/guide/bgclr/}, dan isinya jauh lebih banyak referensi daripada tutorial.

Kalau kamu masih baru, bacalah bagian tutorial secara berurutan, kurang lebih. Makin ke belakang, makin tidak penting urutannya.

Dan apa pun tingkat keahlianmu, bagian referensi selalu tersedia, lengkap dengan contoh pemanggilan fungsi standard library untuk menyegarkan ingatan kapan pun dibutuhkan. Cocok dibaca sambil makan sereal atau di waktu-waktu lainnya.

Terakhir, kalau kamu membaca versi cetak, coba intip indeksnya: entri yang berasal dari bagian referensi dicetak miring.

\section{Platform dan Compiler}\label{platform-dan-compiler}

Panduan ini akan berusaha untuk tetap pada \href{https://en.wikipedia.org/wiki/ANSI_C}{ISO-standard C}\footnote{https://en.wikipedia.org/wiki/ANSI\_C} yang Klasik dan Sederhana. Ya, sebagian besar. Di sana-sini mungkin ada bahasan tentang \href{https://en.wikipedia.org/wiki/POSIX}{POSIX}\footnote{https://en.wikipedia.org/wiki/POSIX} atau semacamnya, tapi kita lihat nanti.

Pengguna \textbf{Unix} (misalnya Linux, BSD, dll.) coba jalankan \hl{\mbox{\texttt{cc}}} atau \hl{\mbox{\texttt{gcc}}} dari command line---mungkin compiler-nya sudah terpasang. Kalau belum, cari tahu cara memasang \hl{\mbox{\texttt{gcc}}} atau \hl{\mbox{\texttt{clang}}} di distribusimu.

Pengguna \textbf{Windows} sebaiknya cek \href{https://visualstudio.microsoft.com/vs/community/}{Visual Studio Community}\footnote{https://visualstudio.microsoft.com/vs/community/}. Atau, jika menginginkan pengalaman yang lebih mirip Unix (direkomendasikan!), instal \href{https://docs.microsoft.com/en-us/windows/wsl/install-win10}{WSL}\footnote{https://docs.microsoft.com/en-us/windows/wsl/install-win10} dan \hl{\mbox{\texttt{gcc}}}.

Pengguna \textbf{Mac} perlu memasang \href{https://developer.apple.com/xcode/}{XCode}\footnote{https://developer.apple.com/xcode/}, khususnya command line tools-nya.

Ada banyak compiler di luar sana, dan hampir semuanya bisa dipakai untuk buku ini. Compiler C++ juga bisa mengompilasi banyak (tapi tidak semua!) kode C. Sebisa mungkin, pakailah compiler C yang sebenarnya.

\section{Halaman Utama Resmi}\label{halaman-utama-resmi}

Lokasi resmi dokumen ini ada di \url{https://beej.us/guide/bgc/}\footnote{https://beej.us/guide/bgc/}. Alamat ini mungkin berubah di masa depan, tapi yang lebih mungkin terjadi adalah semua panduan lainnya dipindahkan dari komputer Chico State.

\section{Kebijakan Email}\label{kebijakan-email}

Pada umumnya penulis bersedia membantu lewat email, jadi jangan ragu menulis, tapi respons tidak dijamin. Hidup penulis cukup sibuk, dan ada kalanya pertanyaanmu memang tidak sempat dijawab. Kalau begitu, biasanya pesannya dihapus saja. Jangan diambil hati; penulis hanya tidak akan pernah punya waktu untuk memberi jawaban serinci yang kamu butuhkan.

Sebagai aturan, makin rumit pertanyaannya, makin kecil kemungkinan dijawab. Kalau kamu bisa mempersempit pertanyaan sebelum dikirim, dan menyertakan informasi yang relevan (seperti platform, compiler, pesan error yang muncul, dan hal lain yang mungkin membantu pelacakan masalah), peluang mendapat balasan jauh lebih besar.

Kalau tidak ada balasan, coba lagi, teruslah mencari jawabannya, dan kalau tetap buntu, kirim email lagi berisi hasil temuanmu. Mudah-mudahan itu cukup untuk membantu.

Setelah panjang lebar soal cara menulis (dan tidak menulis) email, perlu diketahui bahwa semua pujian yang diterima panduan ini selama bertahun-tahun sangat dihargai. Itu dorongan moral yang nyata, dan senang rasanya tahu panduan ini dipakai untuk kebaikan! \hl{\mbox{\texttt{:-)}}} Terima kasih!

\section{Mirroring}\label{mirroring}

Kamu dipersilakan me-mirror situs ini, baik secara publik maupun privat. Kalau kamu me-mirror-nya secara publik dan ingin ditautkan dari halaman utama, kirim pesan ke \href{mailto:beej@beej.us}{\hl{\mbox{\texttt{beej@beej.us}}}}.

\section{Catatan untuk Penerjemah}\label{catatan-untuk-penerjemah}

Kalau kamu ingin menerjemahkan panduan ini ke bahasa lain, tulis ke \href{beej@beej.us}{\hl{\mbox{\texttt{beej@beej.us}}}} dan terjemahanmu akan ditautkan dari halaman utama. Silakan cantumkan nama dan informasi kontakmu dalam terjemahan.

Perhatikan pembatasan lisensi di bagian Hak Cipta dan Distribusi di bawah.

\section{Hak Cipta dan Distribusi}\label{hak-cipta-dan-distribusi}

Beej's Guide to C memiliki Hak Cipta © 2021 Brian ``Beej Jorgensen'' Hall.

Dengan pengecualian khusus untuk kode sumber dan terjemahan di bawah, karya ini dilisensikan di bawah Lisensi Creative Commons Atribusi-NonKomersial-TanpaDerivasi 3.0. Untuk melihat salinan lisensi ini, kunjungi \href{https://creativecommons.org/licenses/by-nc-nd/3.0/}{\hl{\mbox{\texttt{https://creativecommons.org/licenses/by-nc-nd/3.0/}}}} atau kirim surat ke Creative Commons, 171 Second Street, Suite 300, San Francisco, California, 94105, USA.

Satu pengecualian khusus untuk bagian ``TanpaDerivasi'' dari lisensi adalah sebagai berikut: panduan ini dapat diterjemahkan secara bebas ke bahasa apapun, selama terjemahannya akurat dan panduan dicetak ulang secara keseluruhan. Pembatasan lisensi yang sama berlaku untuk terjemahan seperti untuk panduan asli. Terjemahan juga dapat mencantumkan nama dan informasi kontak penerjemah.

Kode sumber C yang disajikan dalam dokumen ini diberikan ke domain publik dan sepenuhnya bebas dari segala pembatasan lisensi.

Para pendidik didorong secara bebas untuk merekomendasikan atau menyediakan salinan panduan ini kepada murid-murid mereka.

Hubungi \href{beej@beej.us}{\hl{\mbox{\texttt{beej@beej.us}}}} untuk informasi lebih lanjut.

Font sampul adalah Lazenby Computer oleh Andrew Young.

\section{Dedikasi}\label{dedikasi}

Hal-hal tersulit dalam menulis panduan-panduan ini adalah:

\begin{itemize}
\tightlist
\item
  Mempelajari materinya cukup mendalam sampai bisa menjelaskannya
\item
  Mencari cara terbaik untuk menjelaskannya dengan jernih, sebuah proses iteratif yang rasanya tidak pernah selesai
\item
  Tampil sebagai \emph{otoritas}, padahal sebenarnya hanyalah manusia biasa yang berusaha memahami semuanya, sama seperti orang lain
\item
  Tetap konsisten ketika begitu banyak hal lain menyita perhatian
\end{itemize}

Banyak orang telah membantu sepanjang proses ini, dan penulis ingin mengucapkan terima kasih kepada mereka yang telah mewujudkan buku ini.

\begin{itemize}
\tightlist
\item
  Semua orang di Internet yang memutuskan berbagi pengetahuan mereka dalam satu bentuk atau lainnya. Berbagi informasi pembelajaran secara bebas itulah yang membuat Internet menjadi tempat yang luar biasa.
\item
  Para sukarelawan di \href{https://en.cppreference.com/}{cppreference.com}\footnote{https://en.cppreference.com/} yang membangun jembatan dari spesifikasi menuju dunia nyata.
\item
  Orang-orang yang suka membantu dan berpengetahuan luas di \href{https://groups.google.com/g/comp.lang.c}{comp.lang.c}\footnote{https://groups.google.com/g/comp.lang.c} dan \href{https://www.reddit.com/r/C_Programming/}{r/C\_Programming}\footnote{https://www.reddit.com/r/C\_Programming/}, yang menemani penulis melewati bagian-bagian bahasa yang paling sulit.
\item
  Semua orang yang mengirim koreksi dan pull-request, mulai dari instruksi yang menyesatkan sampai typo.
\end{itemize}

Terima kasih! ♥

\chapter{Hello, World!}\label{hello-world}

\section{Apa yang Diharapkan dari C}\label{apa-yang-diharapkan-dari-c}

\begin{tcolorbox}[blanker,colback=white,enhanced jigsaw,breakable,arc=0pt,outer arc=0pt,grow to left by=-1cm,enlarge top by=0.1cm,enlarge bottom by=0.1cm,borderline west={1.5pt}{-.3cm}{black}]

\begin{quote}
\emph{``Where do these stairs go?''}\\
\emph{``They go up.''}

---Ray Stantz and Peter Venkman, Ghostbusters
\end{quote}

\end{tcolorbox}

C adalah bahasa low-level.

Dulu tidak begitu. Di zaman ketika orang masih mengukir punch card dari granit, C adalah cara luar biasa untuk bebas dari kerepotan bahasa lower-level seperti \href{https://en.wikipedia.org/wiki/Assembly_language}{assembly}\footnote{https://en.wikipedia.org/wiki/Assembly\_language}.

Tapi di era modern ini, bahasa-bahasa generasi sekarang menawarkan berbagai fitur yang belum ada pada 1972, ketika C diciptakan. Artinya, C adalah bahasa yang cukup mendasar dengan fitur yang tidak banyak. C bisa melakukan \emph{apa saja}, tapi kamu harus bekerja keras untuk itu.

Lalu kenapa kita masih memakainya sekarang?

\begin{itemize}
\item
  Sebagai alat belajar: C bukan hanya bagian terhormat dari sejarah komputasi, tapi juga terhubung langsung ke \href{https://en.wikipedia.org/wiki/Bare_machine}{bare metal}\footnote{https://en.wikipedia.org/wiki/Bare\_machine} dengan cara yang tidak dimiliki bahasa masa kini. Saat belajar C, kamu belajar bagaimana software berinteraksi dengan memori komputer di level rendah. Tidak ada sabuk pengaman. Dijamin, kamu akan menulis software yang crash. Dan itu semua bagian dari keseruannya!
\item
  Sebagai alat yang berguna: C masih dipakai untuk aplikasi tertentu, seperti membangun \href{https://en.wikipedia.org/wiki/Operating_system}{sistem operasi}\footnote{https://en.wikipedia.org/wiki/Operating\_system} atau \href{https://en.wikipedia.org/wiki/Embedded_system}{embedded systems}\footnote{https://en.wikipedia.org/wiki/Embedded\_system}. (Walaupun bahasa pemrograman \href{https://en.wikipedia.org/wiki/Rust_(programming_language)}{Rust}\footnote{https://en.wikipedia.org/wiki/Rust\_(programming\_language)} sedang mengincar kedua bidang ini!)
\end{itemize}

Kalau kamu sudah mengenal bahasa lain, banyak hal di C yang terasa mudah. C menginspirasi banyak bahasa lain, jadi kamu akan menemukan potongan-potongannya di Go, Rust, Swift, Python, JavaScript, Java, dan berbagai bahasa lainnya. Bagian-bagian itu akan terasa akrab.

\emph{Pointer} yang bikin orang tersandung di C. Hampir semua hal lain sudah akrab, tapi pointer adalah yang aneh. Konsep di balik pointer kemungkinan besar sudah kamu ketahui, tapi C memaksamu membahasnya secara eksplisit, dengan operator yang mungkin belum pernah kamu lihat.

Bagian ini licik karena begitu kamu \href{https://en.wikipedia.org/wiki/Grok}{\emph{grok}}\footnote{https://en.wikipedia.org/wiki/Grok} pointer, tiba-tiba semuanya jadi mudah. Tapi sebelum saat itu tiba, pointer seperti belut licin yang susah dipegang.

Segala hal lain di C hanyalah menghafal cara lain (kadang cara yang sama!) untuk melakukan sesuatu yang sudah pernah kamu lakukan. Pointer-lah bagian yang aneh. Dan bisa dibilang, pointer pun hanyalah variasi dari tema yang mungkin sudah kamu kenal.

Jadi bersiaplah untuk petualangan seru, sedekat mungkin dengan inti komputer tanpa perlu assembly, dalam bahasa komputer paling berpengaruh sepanjang masa\footnote{Penulis tahu pasti ada yang akan membantah, tapi setidaknya masuk tiga besar, kan?}. Pegang erat!

\section{Hello, World!}\label{hello-world-1}

\index{Hello, world|}Inilah contoh program C yang paling klasik. Semua orang memakainya. (Perhatikan bahwa angka di sebelah kiri hanya sebagai acuan bagi pembaca dan bukan bagian dari kode sumber.)

\begin{tcolorbox}[enhanced jigsaw,left=3pt,breakable,colback=white!97.2549!black,boxrule=0.2mm]

\begin{Shaded}
\begin{Highlighting}[numbers=left,,]
\CommentTok{/* Program Hello world */}

\PreprocessorTok{\#include }\ImportTok{\textless{}stdio.h\textgreater{}}

\DataTypeTok{int}\NormalTok{ main}\OperatorTok{(}\DataTypeTok{void}\OperatorTok{)}
\OperatorTok{\{}
\NormalTok{    printf}\OperatorTok{(}\StringTok{"Hello, World!}\SpecialCharTok{\textbackslash{}n}\StringTok{"}\OperatorTok{);}  \CommentTok{// Di sinilah pekerjaan sebenarnya dilakukan}
\OperatorTok{\}}
\end{Highlighting}
\end{Shaded}

\end{tcolorbox}

Sekarang kita pakai sarung tangan karet lengan panjang yang tangguh, ambil scalpel, lalu bedah program ini untuk melihat apa yang membuatnya bekerja. Bersiaplah, kita mulai. Memotong dengan sangat hati-hati\ldots{}

\index{Komentar|(}

Kita selesaikan yang mudah dulu: apa pun di antara digraf \hl{\mbox{\texttt{/*}}} dan \hl{\mbox{\texttt{*/}}} adalah komentar dan sepenuhnya diabaikan oleh compiler. Begitu juga apa pun di baris yang sama setelah \hl{\mbox{\texttt{//}}}. Dengan komentar, kamu bisa meninggalkan pesan untuk dirimu sendiri dan orang lain, sehingga ketika kamu membaca ulang kodenya di masa depan yang jauh, kamu tahu apa yang dulu coba kamu lakukan. Percayalah, kamu akan lupa; itu hal biasa.

\index{Komentar|)}

\index{C Preprocessor|(} \index{\#include directive@\texttt{\#include} directive|(}

Nah, apa itu \hl{\mbox{\texttt{\#include}}}? MENJIJIKKAN! Yah, intinya ia menyuruh C Preprocessor untuk mengambil isi file lain dan menyisipkannya ke dalam kode tepat \emph{di situ}.

Tunggu---apa itu C Preprocessor? Pertanyaan bagus. Kompilasi punya dua tahap\footnote{Secara teknis sih lebih dari dua, tapi anggap saja dua---tidak tahu itu menyenangkan, kan?}: preprocessor dan compiler. \index{Octothorpe}Apa pun yang diawali tanda pagar, simbol hash, atau ``octothorpe'' (\hl{\mbox{\texttt{\#}}}) akan diproses oleh preprocessor\index{Preprocessor} sebelum compiler mulai bekerja. \emph{Directive} preprocessor yang umum adalah \hl{\mbox{\texttt{\#include}}} dan \hl{\mbox{\texttt{\#define}}}.\index{\#define directive@\texttt{\#define} directive} Soal itu kita bahas nanti.\index{\#include directive@\texttt{\#include} directive|)}\index{C Preprocessor|)}

Sebelum lanjut, kenapa repot-repot menyebut bahwa tanda pagar itu namanya octothorpe? Jawabannya sederhana: kata octothorpe itu terlalu lucu, jadi namanya harus disebarkan setiap ada kesempatan. Octothorpe. Octothorpe, octothorpe, octothorpe.

Nah, \emph{anyway}. Setelah C preprocessor selesai memproses semuanya, hasilnya siap diambil compiler untuk dijadikan \href{https://en.wikipedia.org/wiki/Assembly_language}{kode assembly}\footnote{https://en.wikipedia.org/wiki/Assembly\_language}, \href{https://en.wikipedia.org/wiki/Machine_code}{machine code}\footnote{https://en.wikipedia.org/wiki/Machine\_code}, atau apa pun yang dikerjakannya. Machine code adalah ``bahasa'' yang dipahami CPU, dan CPU memahaminya \emph{dengan sangat cepat}. Ini salah satu alasan program C cenderung cepat.

Untuk sekarang, jangan pusingkan detail teknis kompilasi; cukup ketahui bahwa source-mu dilewatkan ke preprocessor, lalu outputnya dilewatkan ke compiler, yang kemudian menghasilkan executable untuk kamu jalankan.

Bagaimana dengan sisa barisnya? \index{stdio.h header file@\texttt{stdio.h} header file|(}Apa itu \hl{\mbox{\texttt{\textless{}stdio.h\textgreater{}}}}? Itulah yang disebut \emph{header file}. Ekstensi .h di akhir nama file jadi penandanya. Yang satu ini header file ``Standard I/O'' (\hl{\mbox{\texttt{stdio}}}), yang nantinya akan akrab dan kamu sukai. Header ini memberi kita akses ke berbagai fungsi I/O\footnote{Secara teknis, isinya directive preprocessor dan function prototype (nanti dibahas) untuk kebutuhan input dan output yang umum.}. Program demo kita mengeluarkan string ``Hello, World!'', jadi kita perlu akses ke fungsi \index{printf() fungsi@\texttt{printf()} fungsi|(} \hl{\mbox{\texttt{printf()}}}. File \hl{\mbox{\texttt{\textless{}stdio.h\textgreater{}}}} memberi kita akses itu. Intinya, kalau kita memakai \hl{\mbox{\texttt{printf()}}} tanpa \hl{\mbox{\texttt{\#include\ \textless{}stdio.h\textgreater{}}}}, compiler akan protes.

Dari mana tahu bahwa \hl{\mbox{\texttt{printf()}}} butuh \hl{\mbox{\texttt{\#include\ \textless{}stdio.h\textgreater{}}}}?\index{printf() fungsi@\texttt{printf()} fungsi|)} Jawabannya: dokumentasi. Di sistem Unix, ketik \hl{\mbox{\texttt{man\ 3\ printf}}} dan di bagian paling atas halaman man akan tertulis header file apa yang diperlukan. Atau lihat bagian referensi di buku ini. \hl{\mbox{\texttt{:-)}}} \index{stdio.h header file@\texttt{stdio.h} header file|)}

Astaga. Semua itu hanya untuk membahas baris pertama! Tapi harus diakui, baris itu sudah dibedah tuntas. Tidak ada misteri yang tersisa!

Jadi tarik napas dulu\ldots lalu lihat lagi kode contohnya. Tinggal beberapa baris mudah.

Selamat datang kembali! Penulis tahu kamu sebenarnya tidak istirahat; tadi hanya menuruti kemauanmu.

\index{main() fungsi@\texttt{main()} fungsi|(}Baris berikutnya adalah \hl{\mbox{\texttt{main()}}}. Ini definisi fungsi \hl{\mbox{\texttt{main()}}}; semua yang ada di antara kurung kurawal (\hl{\mbox{\texttt{\{}}} dan \hl{\mbox{\texttt{\}}}}) adalah bagian dari definisi fungsi tersebut.

(Lalu bagaimana cara \emph{memanggil} fungsi lain? Jawabannya ada di baris \hl{\mbox{\texttt{printf()}}}, yang sebentar lagi kita bahas.)

Fungsi \hl{\mbox{\texttt{main}}} istimewa dalam banyak hal, tapi satu hal mengalahkan yang lain: fungsi inilah yang dipanggil otomatis ketika programmu mulai berjalan. Tidak ada kodemu yang dipanggil sebelum \hl{\mbox{\texttt{main()}}}. Untuk contoh kita, ini tidak jadi masalah karena kita hanya ingin mencetak satu baris lalu keluar.

Oh ya, satu lagi: begitu eksekusi program melewati ujung \hl{\mbox{\texttt{main()}}}, yaitu di kurung kurawal penutup di bawah sana, program akan keluar dan kamu kembali ke command prompt.

Jadi sekarang kita tahu bahwa program itu mengambil sebuah header file, \hl{\mbox{\texttt{stdio.h}}}\index{stdio.h|}, dan mendeklarasikan fungsi \hl{\mbox{\texttt{main()}}} yang dijalankan saat program dimulai. Apa saja isi menarik di dalam \hl{\mbox{\texttt{main()}}}\index{main() fungsi@\texttt{main()} fungsi|)}?

Senang sekali kamu bertanya. Sungguh! Isinya hanya satu: pemanggilan fungsi \index{printf() fungsi@\texttt{printf()} fungsi|(}\hl{\mbox{\texttt{printf()}}}. Kamu bisa membedakan pemanggilan fungsi dari definisi fungsi dengan beberapa cara, salah satunya adalah tidak adanya kurung kurawal setelahnya. Pemanggilan fungsi diakhiri titik koma supaya compiler tahu ekspresinya sudah selesai. Kamu akan menaruh titik koma di akhir hampir segalanya, seperti yang akan terlihat nanti.

Kamu mengoper satu argumen ke fungsi \hl{\mbox{\texttt{printf()}}}\index{printf() fungsi@\texttt{printf()} fungsi|)}: sebuah string yang akan dicetak saat fungsi dipanggil. Oh, ya---kita memanggil fungsi! Keren! Tunggu, tunggu---jangan sombong. \index{{\backslash}n newline@\texttt{{\backslash}n} newline|(}Apa itu \hl{\mbox{\texttt{\textbackslash{}n}}} yang aneh di ujung string? Begini: sebagian besar karakter dalam string tercetak persis seperti yang tersimpan. Tapi ada karakter tertentu yang sulit ditampilkan di layar, sehingga ditulis sebagai kode backslash dua karakter. Salah satu yang paling populer adalah \hl{\mbox{\texttt{\textbackslash{}n}}} (dibaca ``backslash-N'' atau cukup ``newline''), yang mewakili karakter \emph{newline}. Karakter ini membuat pencetakan berikutnya dilanjutkan di awal baris selanjutnya, bukan di baris yang sama. Mirip menekan Enter di akhir baris.\index{{\backslash}n newline@\texttt{{\backslash}n} newline|)}

Sekarang salin kode itu ke file bernama \hl{\mbox{\texttt{hello.c}}} lalu build. Di platform mirip Unix (misalnya Linux, BSD, Mac, atau WSL), dari command line kamu bisa build dengan perintah seperti ini:

\index{gcc compiler@\texttt{gcc} compiler}

\begin{tcolorbox}[enhanced jigsaw,left=3pt,breakable,colback=white!97.2549!black,boxrule=0.2mm]

\begin{Shaded}
\begin{Highlighting}[]
\FunctionTok{gcc} \AttributeTok{{-}o}\NormalTok{ hello hello.c}
\end{Highlighting}
\end{Shaded}

\end{tcolorbox}

(Artinya: ``kompilasi \hl{\mbox{\texttt{hello.c}}}, dan hasilkan executable bernama \hl{\mbox{\texttt{hello}}}''.)

Setelah selesai, kamu akan punya file bernama \hl{\mbox{\texttt{hello}}} yang bisa dijalankan dengan perintah ini:

\begin{tcolorbox}[enhanced jigsaw,left=3pt,breakable,colback=white!97.2549!black,boxrule=0.2mm]

\begin{Shaded}
\begin{Highlighting}[]
\NormalTok{./hello}
\end{Highlighting}
\end{Shaded}

\end{tcolorbox}

(Awalan \hl{\mbox{\texttt{./}}} memberi tahu shell untuk ``jalankan dari direktori saat ini''.)

Lihat apa yang terjadi:

\begin{tcolorbox}[enhanced jigsaw,left=3pt,breakable,colback=white!97.2549!black,boxrule=0.2mm]

\begin{Shaded}
\begin{Highlighting}[]
\NormalTok{Hello, World! }
\end{Highlighting}
\end{Shaded}

\end{tcolorbox}

Selesai dan sudah teruji! Kirim!\index{Hello, world|)}

\section{Detail Kompilasi}\label{detail-kompilasi}

\index{Kompilasi|(}Mari kita bahas lebih jauh cara membangun program C dan apa yang terjadi di balik layar.

Seperti bahasa lain, C punya \emph{kode sumber}. Tapi tergantung bahasa asalmu, mungkin kamu belum pernah harus \emph{mengompilasi} kode sumber menjadi \emph{executable}.

Kompilasi adalah proses mengubah kode sumber C menjadi program yang bisa dijalankan oleh sistem operasi.

Developer JavaScript dan Python tidak terbiasa dengan langkah kompilasi yang terpisah---padahal di balik layar kompilasi tetap terjadi! Python mengompilasi kode sumbermu menjadi \emph{bytecode} yang dijalankan oleh mesin virtual Python. Developer Java sudah terbiasa dengan kompilasi, tapi hasilnya adalah bytecode untuk Java Virtual Machine.

Kompilasi C menghasilkan \emph{machine code}: deretan 1 dan 0 yang langsung dan cepat dieksekusi oleh CPU.

\begin{tcolorbox}[blanker,colback=white,enhanced jigsaw,breakable,arc=0pt,outer arc=0pt,grow to left by=-1cm,enlarge top by=0.1cm,enlarge bottom by=0.1cm,borderline west={1.5pt}{-.3cm}{black}]

\begin{quote}
Bahasa yang biasanya tidak dikompilasi disebut bahasa \emph{interpreted}. Tapi seperti yang sudah disinggung soal Java dan Python, bahasa-bahasa itu juga punya langkah kompilasi. Dan tidak ada aturan yang melarang C diinterpretasikan. (Interpreter C itu ada!) Singkatnya, semua ini area abu-abu. Secara umum, kompilasi hanyalah mengubah kode sumber menjadi bentuk lain yang lebih mudah dieksekusi.
\end{quote}

\end{tcolorbox}

C compiler adalah program yang melakukan kompilasi.

Seperti sudah disebutkan, \hl{\mbox{\texttt{gcc}}} adalah compiler yang terpasang di banyak \href{https://en.wikipedia.org/wiki/Unix}{sistem operasi mirip Unix}\footnote{https://en.wikipedia.org/wiki/Unix}. Biasanya ia dijalankan dari command line di terminal, tapi tidak selalu. Kamu juga bisa menjalankannya dari IDE.

Jadi, bagaimana cara build lewat command line?

\section{\texorpdfstring{Build dengan \hl{\mbox{\texttt{gcc}}}}{Build dengan gcc}}\label{build-dengan-gcc}

\index{gcc compiler@\texttt{gcc} compiler|(}Kalau kamu punya file sumber bernama \hl{\mbox{\texttt{hello.c}}} di direktori saat ini, kamu bisa membangunnya menjadi program bernama \hl{\mbox{\texttt{hello}}} dengan mengetik perintah ini di terminal:

\begin{tcolorbox}[enhanced jigsaw,left=3pt,breakable,colback=white!97.2549!black,boxrule=0.2mm]

\begin{Shaded}
\begin{Highlighting}[]
\FunctionTok{gcc} \AttributeTok{{-}o}\NormalTok{ hello hello.c}
\end{Highlighting}
\end{Shaded}

\end{tcolorbox}

\hl{\mbox{\texttt{-o}}} berarti ``output ke file ini''\footnote{Kalau nama file output tidak diberikan, hasilnya otomatis ditulis ke file bernama \hl{\mbox{\texttt{a.out}}}---nama ini berakar jauh di sejarah Unix.}. Lalu \hl{\mbox{\texttt{hello.c}}} di akhir adalah nama file yang ingin dikompilasi.

Kalau sumbermu terbagi menjadi beberapa file, kamu bisa mengompilasi semuanya sekaligus (hampir seperti satu file, meski aturan sebenarnya lebih rumit) dengan menuliskan semua file \hl{\mbox{\texttt{.c}}} di command line:

\begin{tcolorbox}[enhanced jigsaw,left=3pt,breakable,colback=white!97.2549!black,boxrule=0.2mm]

\begin{Shaded}
\begin{Highlighting}[]
\FunctionTok{gcc} \AttributeTok{{-}o}\NormalTok{ awesomegame ui.c characters.c npc.c items.c}
\end{Highlighting}
\end{Shaded}

\end{tcolorbox}

\index{gcc compiler@\texttt{gcc} compiler|)}

dan semuanya akan dibangun menjadi satu executable besar.

Itu cukup untuk memulai---nanti kita akan membahas detail soal banyak file sumber, object file, dan segala macam hal menarik lainnya.\index{Kompilasi|)}

\section{\texorpdfstring{Build dengan \hl{\mbox{\texttt{clang}}}}{Build dengan clang}}\label{build-dengan-clang}

Di Mac, compiler bawaannya bukan \hl{\mbox{\texttt{gcc}}}, melainkan \hl{\mbox{\texttt{clang}}}\index{clang compiler@\texttt{clang} compiler}. Tapi ada wrapper yang ikut terpasang, jadi menjalankan \hl{\mbox{\texttt{gcc}}} tetap berfungsi.

Kamu juga bisa memasang compiler \hl{\mbox{\texttt{gcc}}}\index{gcc compiler@\texttt{gcc} compiler} yang asli lewat \href{https://formulae.brew.sh/formula/gcc}{Homebrew}\footnote{https://formulae.brew.sh/formula/gcc} atau cara lain.

\section{Build dari IDE}\label{build-dari-ide}

\index{Integrated Development Environment|(}Kalau kamu memakai \emph{Integrated Development Environment} (IDE), kemungkinan besar kamu tidak perlu build dari command line.

Di Visual Studio, \hl{\mbox{\texttt{CTRL-F7}}} untuk build dan \hl{\mbox{\texttt{CTRL-F5}}} untuk run.

Di VS Code, tekan \hl{\mbox{\texttt{F5}}} untuk run lewat debugger. (Kamu perlu memasang C/C++ Extension.)

Di XCode, build dengan \hl{\mbox{\texttt{COMMAND-B}}} dan run dengan \hl{\mbox{\texttt{COMMAND-R}}}. Untuk mendapatkan command line tools, cari ``XCode command line tools'' dan kamu akan menemukan petunjuk pemasangannya.

Sebagai permulaan, ada baiknya juga mencoba build lewat command line---itu bagian dari sejarahnya!\index{Integrated Development Environment|)}

\section{Versi C}\label{versi-c}

\index{Language versions|(}C sudah berkembang jauh selama bertahun-tahun, dan punya banyak nomor versi untuk menyebut dialek bahasa yang sedang dipakai.

Nomor ini biasanya merujuk ke tahun spesifikasinya.

Versi yang paling terkenal ialah C89, C99, C11, dan C23. Buku ini berfokus pada yang terakhir.

Berikut tabel yang lebih lengkap:

{\def\LTcaptype{none} % do not increment counter
\begin{longtable}[]{@{}
  >{\raggedright\arraybackslash}p{(\linewidth - 2\tabcolsep) * \real{0.2632}}
  >{\raggedright\arraybackslash}p{(\linewidth - 2\tabcolsep) * \real{0.7368}}@{}}
\toprule\noalign{}
\begin{minipage}[b]{\linewidth}\raggedright
Versi
\end{minipage} & \begin{minipage}[b]{\linewidth}\raggedright
Deskripsi
\end{minipage} \\
\midrule\noalign{}
\endhead
\bottomrule\noalign{}
\endlastfoot
K\&R C & 1978, yang asli. Dinamai dari Brian Kernighan dan Dennis Ritchie. Ritchie merancang bahasa dan menulis kodenya, sementara Kernighan ikut menulis bukunya. Kode K\&R asli jarang ditemui sekarang. Kalaupun ada, kodenya akan terlihat aneh, seperti Bahasa Inggris Pertengahan bagi pembaca Inggris modern. \\
\textbf{C89}, ANSI C, C90 & Pada 1989, American National Standards Institute (ANSI) menerbitkan spesifikasi bahasa C yang menentukan corak C hingga hari ini. Setahun kemudian, kendali diserahkan ke International Organization for Standardization (ISO), yang menerbitkan C90 yang identik. \\
C95 & Tambahan pada C89 yang jarang disebut, berisi dukungan wide character. \\
\textbf{C99} & Perombakan besar pertama dengan banyak tambahan bahasa. Yang paling diingat orang adalah hadirnya komentar gaya \hl{\mbox{\texttt{//}}}. Saat tulisan ini dibuat, ini versi C yang paling populer dipakai. \\
\textbf{C11} & Pembaruan besar ini mencakup dukungan Unicode dan multi-threading. Ingat, kalau kamu mulai memakai fitur-fitur ini, portabilitas ke tempat-tempat yang masih terjebak di C99 bisa jadi korbannya. Tapi, jujur saja, 1999 sudah lama sekali. \\
C17, C18 & Pembaruan perbaikan bug untuk C11. C17 tampaknya nama resminya, tapi penerbitannya tertunda sampai 2018. Setahu penulis, keduanya bisa dipertukarkan, dan C17 lebih disukai. \\
C23 & Spesifikasi terbaru. \\
\end{longtable}
}

\index{gcc compiler@\texttt{gcc} compiler|(}Kamu bisa memaksa GCC memakai salah satu standar ini lewat argumen command line \hl{\mbox{\texttt{-std=}}}. Kalau ingin ketat mengikuti standar, tambahkan \hl{\mbox{\texttt{-pedantic}}}.

Misalnya:

\begin{tcolorbox}[enhanced jigsaw,left=3pt,breakable,colback=white!97.2549!black,boxrule=0.2mm]

\begin{Shaded}
\begin{Highlighting}[]
\FunctionTok{gcc} \AttributeTok{{-}std}\OperatorTok{=}\NormalTok{c11 }\AttributeTok{{-}pedantic}\NormalTok{ foo.c}
\end{Highlighting}
\end{Shaded}

\end{tcolorbox}

Untuk buku ini, program dikompilasi sebagai C23 dengan semua peringatan diaktifkan:

\begin{tcolorbox}[enhanced jigsaw,left=3pt,breakable,colback=white!97.2549!black,boxrule=0.2mm]

\begin{Shaded}
\begin{Highlighting}[]
\FunctionTok{gcc} \AttributeTok{{-}Wall} \AttributeTok{{-}Wextra} \AttributeTok{{-}std}\OperatorTok{=}\NormalTok{c23 }\AttributeTok{{-}pedantic}\NormalTok{ foo.c}
\end{Highlighting}
\end{Shaded}

\end{tcolorbox}

\index{gcc compiler@\texttt{gcc} compiler|)}

\chapter{Variabel dan Pernyataan}\label{variabel-dan-pernyataan}

\begin{tcolorbox}[blanker,colback=white,enhanced jigsaw,breakable,arc=0pt,outer arc=0pt,grow to left by=-1cm,enlarge top by=0.1cm,enlarge bottom by=0.1cm,borderline west={1.5pt}{-.3cm}{black}]

\begin{quote}
\emph{``It takes all kinds to make a world, does it not, Padre?''}\\
\emph{``So it does, my son, so it does.''}

---Pirate Captain Thomas Bartholomew Red to the Padre, Pirates
\end{quote}

\end{tcolorbox}

Program C memang berisi banyak hal.

Yap.

Dan karena berbagai alasan, kita semua akan terbantu kalau mengelompokkan jenis-jenis hal yang ada dalam sebuah program, supaya jelas apa yang sedang dibicarakan.

\section{Variabel}\label{variabel}

\index{variabel|(}Orang bilang, ``variabel menyimpan nilai''. Tapi cara lain memandangnya: variabel adalah nama yang mudah dibaca manusia untuk merujuk ke sebuah data di memori.

Kita berhenti sebentar di sini untuk mengintip ke lubang kelinci bernama pointer. Jangan khawatir dulu.

Bayangkan memori sebagai array raksasa berisi byte\footnote{Satu ``byte'' biasanya adalah bilangan biner 8-bit. Bayangkan sebagai integer yang hanya bisa menyimpan nilai 0 sampai 255, inklusif. Secara teknis, C membolehkan byte punya jumlah bit berapa pun, dan kalau ingin merujuk dengan jelas ke bilangan 8-bit, istilah yang tepat adalah \emph{oktet}. Tapi programmer akan menganggap kamu bermaksud 8-bit kalau bilang ``byte'', kecuali kamu menyebutkan yang lain.}. Data disimpan di dalam ``array'' ini\footnote{Cara kerja memori modern di sini disederhanakan sedikit berlebihan. Tapi model mental ini cukup berhasil, jadi mohon dimaklumi.}. Kalau sebuah angka lebih besar dari satu byte, angka itu disimpan dalam beberapa byte. Karena memori mirip array, setiap byte memori bisa dirujuk lewat indeksnya. Indeks ke dalam memori ini juga disebut \emph{alamat}, atau \emph{lokasi}, atau \emph{pointer}.

Ketika kamu punya variabel di C, nilainya berada di memori \emph{di suatu tempat}, pada suatu alamat. Tentu saja. Memangnya di mana lagi? Tapi merujuk ke sebuah nilai lewat alamat numeriknya itu merepotkan, jadi kita memberinya nama, dan itulah yang disebut variabel.

Ada dua alasan kenapa hal ini dibahas:

\begin{enumerate}
\def\labelenumi{\arabic{enumi}.}
\tightlist
\item
  Ini memudahkan pemahaman variabel pointer nanti---variabel yang menyimpan alamat variabel lain!
\item
  Dan, ini juga memudahkan pemahaman pointer nanti.
\end{enumerate}

Jadi, variabel adalah nama untuk data yang disimpan di memori pada suatu alamat.

\subsection{Nama Variabel}\label{nama-variabel}

\index{variabel|(}Nama variabel boleh memakai karakter apa pun dalam rentang 0-9, A-Z, a-z, dan underscore, dengan aturan berikut:

\begin{itemize}
\tightlist
\item
  Nama variabel tidak boleh diawali digit 0-9.
\item
  Nama variabel tidak boleh diawali dua underscore.
\item
  Nama variabel tidak boleh diawali underscore yang diikuti huruf kapital A-Z.
\end{itemize}

Untuk Unicode, coba saja. Spesifikasi di §D.2 punya sejumlah aturan tentang rentang codepoint Unicode mana yang boleh dipakai di bagian mana dari identifier, tapi itu terlalu panjang untuk ditulis di sini dan kemungkinan besar tidak perlu kamu pikirkan.\index{variabel|)}

\subsection{Tipe Variabel}\label{tipe-variabel}

\index{tipe}Tergantung bahasa apa yang ada di gudang senjatamu, kamu mungkin sudah akrab, mungkin juga belum, dengan konsep tipe. Tapi C cukup ketat soal ini, jadi kita perlu mengulasnya sedikit.

Berikut beberapa contoh tipe, yang paling dasar saja:

\index{int tipe@\texttt{int} tipe}\index{float tipe@\texttt{float} tipe}\index{char tipe@\texttt{char} tipe}\index{char * tipe@\texttt{char *} tipe}

{\def\LTcaptype{none} % do not increment counter
\begin{longtable}[]{@{}
  >{\raggedright\arraybackslash}p{(\linewidth - 4\tabcolsep) * \real{0.2353}}
  >{\raggedleft\arraybackslash}p{(\linewidth - 4\tabcolsep) * \real{0.4118}}
  >{\raggedright\arraybackslash}p{(\linewidth - 4\tabcolsep) * \real{0.3529}}@{}}
\toprule\noalign{}
\begin{minipage}[b]{\linewidth}\raggedright
Tipe
\end{minipage} & \begin{minipage}[b]{\linewidth}\raggedleft
Contoh
\end{minipage} & \begin{minipage}[b]{\linewidth}\raggedright
Tipe C
\end{minipage} \\
\midrule\noalign{}
\endhead
\bottomrule\noalign{}
\endlastfoot
Integer & \hl{\mbox{\texttt{3490}}} & \hl{\mbox{\texttt{int}}} \\
Floating point & \hl{\mbox{\texttt{3.14159}}} & \hl{\mbox{\texttt{float}}}\footnote{Ini sedikit bohong. Secara teknis \hl{\mbox{\texttt{3.14159}}} bertipe \hl{\mbox{\texttt{double}}}, tapi kita belum sampai ke sana, dan tujuannya adalah supaya kamu mengasosiasikan \hl{\mbox{\texttt{float}}} dengan ``Floating Point''. C dengan senang hati akan mengonversinya menjadi \hl{\mbox{\texttt{float}}}. Singkatnya, jangan dipusingkan dulu.} \\
Karakter (tunggal) & \hl{\mbox{\texttt{\textquotesingle{}c\textquotesingle{}}}} & \hl{\mbox{\texttt{char}}} \\
String & \hl{\mbox{\texttt{"Hello,\ world!"}}} & \hl{\mbox{\texttt{char\ *}}}\footnote{Bacanya ``pointer ke char'' atau ``char pointer''. ``Char'' kependekan dari karakter. Belum ada studi resmi, tapi secara anekdot kebanyakan orang melafalkannya ``char'' (seperti ``charcoal''), sebagian kecil bilang ``car'', dan segelintir bilang ``care''. Pointer akan dibahas lebih lanjut nanti.} \\
\end{longtable}
}

C berusaha mengonversi secara otomatis antara kebanyakan tipe numerik kalau diminta. Selain itu, semua konversi dilakukan manual, terutama antara string dan numerik.

Hampir semua tipe di C adalah varian dari tipe-tipe ini.

Sebelum memakai variabel, kamu harus \emph{mendeklarasikannya} dan memberi tahu C tipe apa yang akan dipegangnya. Setelah dideklarasikan, tipe variabel tidak bisa diubah lagi saat runtime. Apa yang kamu tetapkan, itulah yang berlaku sampai variabel keluar dari scope dan diserap kembali oleh alam semesta.

Mari kita ambil kode ``Hello, world'' tadi dan tambahkan beberapa variabel:

\begin{tcolorbox}[enhanced jigsaw,left=3pt,breakable,colback=white!97.2549!black,boxrule=0.2mm]

\begin{Shaded}
\begin{Highlighting}[numbers=left,,]
\PreprocessorTok{\#include }\ImportTok{\textless{}stdio.h\textgreater{}}

\DataTypeTok{int}\NormalTok{ main}\OperatorTok{(}\DataTypeTok{void}\OperatorTok{)}
\OperatorTok{\{}
    \DataTypeTok{int}\NormalTok{ i}\OperatorTok{;}    \CommentTok{// Menyimpan integer bertanda, mis. {-}3, {-}2, 0, 1, 10}
    \DataTypeTok{float}\NormalTok{ f}\OperatorTok{;}  \CommentTok{// Menyimpan bilangan floating point bertanda, mis. {-}3.1416}

\NormalTok{    printf}\OperatorTok{(}\StringTok{"Hello, World!}\SpecialCharTok{\textbackslash{}n}\StringTok{"}\OperatorTok{);}  \CommentTok{// Ah, keakraban yang menenangkan}
\OperatorTok{\}}
\end{Highlighting}
\end{Shaded}

\end{tcolorbox}

Nah! Kita sudah mendeklarasikan beberapa variabel. Keduanya belum kita pakai dan belum diinisialisasi. Yang satu menyimpan integer, yang lain menyimpan floating point number (pada dasarnya bilangan real, kalau kamu punya latar belakang matematika).

\index{variabel!belum diinisialisasi}Variabel yang belum diinisialisasi bernilai tak tentu\footnote{Secara informal kita bilang nilainya ``acak'', padahal sebenarnya bukan---bahkan bukan bilangan pseudo-acak.}. Variabel harus diinisialisasi, atau kamu harus menganggap isinya angka sampah.

\begin{tcolorbox}[blanker,colback=white,enhanced jigsaw,breakable,arc=0pt,outer arc=0pt,grow to left by=-1cm,enlarge top by=0.1cm,enlarge bottom by=0.1cm,borderline west={1.5pt}{-.3cm}{black}]

\begin{quote}
Di sinilah C bisa ``menjebak'' kamu. Seringnya, dari pengalaman, nilai yang tak tentu itu nol\ldots{} tapi bisa berubah dari satu run ke run berikutnya! Jangan pernah menganggap nilainya nol, meskipun kamu melihatnya nol. \emph{Selalu} inisialisasi variabel secara eksplisit sebelum dipakai\footnote{Ini tidak sepenuhnya benar. Ketika kita sampai ke pembahasan static storage duration, kamu akan tahu bahwa beberapa variabel otomatis diinisialisasi ke nol. Tapi cara yang aman adalah selalu menginisialisasinya sendiri.}.
\end{quote}

\end{tcolorbox}

Apa nih? Kamu mau menyimpan angka di variabel itu? Gila kamu!

Ayo kita lakukan: \index{= operator penugasan@\texttt{=} operator penugasan}

\begin{tcolorbox}[enhanced jigsaw,left=3pt,breakable,colback=white!97.2549!black,boxrule=0.2mm]

\begin{Shaded}
\begin{Highlighting}[numbers=left,,]
\DataTypeTok{int}\NormalTok{ main}\OperatorTok{(}\DataTypeTok{void}\OperatorTok{)}
\OperatorTok{\{}
    \DataTypeTok{int}\NormalTok{ i}\OperatorTok{;}

\NormalTok{    i }\OperatorTok{=} \DecValTok{2}\OperatorTok{;} \CommentTok{// Menugaskan nilai 2 ke variabel i}

\NormalTok{    printf}\OperatorTok{(}\StringTok{"Hello, World!}\SpecialCharTok{\textbackslash{}n}\StringTok{"}\OperatorTok{);}
\OperatorTok{\}}
\end{Highlighting}
\end{Shaded}

\end{tcolorbox}

Mantap. Kita sudah menyimpan sebuah nilai. Sekarang kita cetak.

\index{printf() fungsi@\texttt{printf()} fungsi|(}Caranya, kita oper \emph{dua} argumen keren ke fungsi \hl{\mbox{\texttt{printf()}}}. Argumen pertama adalah string yang menjelaskan apa yang akan dicetak dan bagaimana mencetaknya (disebut \emph{format string}), dan argumen kedua adalah nilai yang dicetak, yaitu apa pun isi variabel \hl{\mbox{\texttt{i}}}.

\hl{\mbox{\texttt{printf()}}} mencari urutan-urutan khusus dalam format string yang diawali tanda persen (\hl{\mbox{\texttt{\%}}}), yang memberi tahu apa yang harus dicetak. Misalnya, kalau menemukan \hl{\mbox{\texttt{\%d}}}, ia mengambil parameter berikutnya dan mencetaknya sebagai integer. Kalau menemukan \hl{\mbox{\texttt{\%f}}}, nilainya dicetak sebagai float. Kalau menemukan \hl{\mbox{\texttt{\%s}}}, ia mencetak string.

Dengan begitu, kita bisa mencetak nilai berbagai tipe seperti ini:

\begin{tcolorbox}[enhanced jigsaw,left=3pt,breakable,colback=white!97.2549!black,boxrule=0.2mm]

\begin{Shaded}
\begin{Highlighting}[numbers=left,,]
\PreprocessorTok{\#include }\ImportTok{\textless{}stdio.h\textgreater{}}

\DataTypeTok{int}\NormalTok{ main}\OperatorTok{(}\DataTypeTok{void}\OperatorTok{)}
\OperatorTok{\{}
    \DataTypeTok{int}\NormalTok{ i }\OperatorTok{=} \DecValTok{2}\OperatorTok{;}
    \DataTypeTok{float}\NormalTok{ f }\OperatorTok{=} \FloatTok{3.14}\OperatorTok{;}
    \DataTypeTok{char} \OperatorTok{*}\NormalTok{s }\OperatorTok{=} \StringTok{"Hello, world!"}\OperatorTok{;}  \CommentTok{// char * ("char pointer") adalah sebuah string}

\NormalTok{    printf}\OperatorTok{(}\StringTok{"}\SpecialCharTok{\%s}\StringTok{  i = }\SpecialCharTok{\%d}\StringTok{ and f = }\SpecialCharTok{\%f}\StringTok{!}\SpecialCharTok{\textbackslash{}n}\StringTok{"}\OperatorTok{,}\NormalTok{ s}\OperatorTok{,}\NormalTok{ i}\OperatorTok{,}\NormalTok{ f}\OperatorTok{);}
\OperatorTok{\}}
\end{Highlighting}
\end{Shaded}

\end{tcolorbox}

Outputnya akan seperti ini:

\begin{tcolorbox}[enhanced jigsaw,left=3pt,breakable,colback=white!97.2549!black,boxrule=0.2mm]

\begin{Shaded}
\begin{Highlighting}[]
\NormalTok{Hello, world!  i = 2 and f = 3.14!}
\end{Highlighting}
\end{Shaded}

\end{tcolorbox}

Dengan begitu, \hl{\mbox{\texttt{printf()}}} mungkin mirip dengan format string atau parameterized string di bahasa lain yang sudah kamu kenal. \index{printf() fungsi@\texttt{printf()} fungsi|)}

\textbf{\emph{CATATAN:}} Di bagian-bagian berikut, kita anggap kamu sudah mendeklarasikan variabelnya di awal. Kalau sebuah contoh memakai integer \hl{\mbox{\texttt{i}}} dan \hl{\mbox{\texttt{j}}}, bayangkan saja ada baris ini di suatu tempat di atasnya:

\begin{tcolorbox}[enhanced jigsaw,left=3pt,breakable,colback=white!97.2549!black,boxrule=0.2mm]

\begin{Shaded}
\begin{Highlighting}[]
\DataTypeTok{int}\NormalTok{ i}\OperatorTok{,}\NormalTok{ j}\OperatorTok{;}
\end{Highlighting}
\end{Shaded}

\end{tcolorbox}

\subsection{Tipe Boolean}\label{tipe-boolean}

\index{tipe Boolean|(}Apakah C punya tipe Boolean, ya atau tidak?

\hl{\mbox{\texttt{1}}}!

Secara historis, C tidak punya tipe Boolean, dan sebagian orang berpendapat sampai sekarang pun belum.

Di C, \hl{\mbox{\texttt{0}}} berarti ``false'', dan nilai bukan nol berarti ``true''.

Jadi \hl{\mbox{\texttt{1}}} itu true. \hl{\mbox{\texttt{-37}}} juga true. Dan \hl{\mbox{\texttt{0}}} adalah false.

Kamu bisa mendeklarasikan Boolean sebagai \hl{\mbox{\texttt{int}}} biasa:

\begin{tcolorbox}[enhanced jigsaw,left=3pt,breakable,colback=white!97.2549!black,boxrule=0.2mm]

\begin{Shaded}
\begin{Highlighting}[]
\DataTypeTok{int}\NormalTok{ x }\OperatorTok{=} \DecValTok{1}\OperatorTok{;}

\ControlFlowTok{if} \OperatorTok{(}\NormalTok{x}\OperatorTok{)} \OperatorTok{\{}
\NormalTok{    printf}\OperatorTok{(}\StringTok{"x is true!}\SpecialCharTok{\textbackslash{}n}\StringTok{"}\OperatorTok{);}
\OperatorTok{\}}
\end{Highlighting}
\end{Shaded}

\end{tcolorbox}

Di C23, kita mendapat \hl{\mbox{\texttt{bool}}}, \hl{\mbox{\texttt{true}}}, dan \hl{\mbox{\texttt{false}}} yang sesungguhnya. Sebelum itu, kalau versi C-mu cukup modern, kamu bisa \hl{\mbox{\texttt{\#include\ \textless{}stdbool.h\textgreater{}}}}\index{stdbool.h header file@\texttt{stdbool.h} header file} untuk hasil yang sama.

\begin{tcolorbox}[enhanced jigsaw,left=3pt,breakable,colback=white!97.2549!black,boxrule=0.2mm]

\begin{Shaded}
\begin{Highlighting}[numbers=left,,]
\PreprocessorTok{\#include }\ImportTok{\textless{}stdio.h\textgreater{}}
\PreprocessorTok{\#include }\ImportTok{\textless{}stdbool.h\textgreater{}}\PreprocessorTok{  }\CommentTok{// tidak diperlukan di C23}

\DataTypeTok{int}\NormalTok{ main}\OperatorTok{(}\DataTypeTok{void}\OperatorTok{)} \OperatorTok{\{}
    \DataTypeTok{bool}\NormalTok{ x }\OperatorTok{=} \KeywordTok{true}\OperatorTok{;}

    \ControlFlowTok{if} \OperatorTok{(}\NormalTok{x}\OperatorTok{)} \OperatorTok{\{}
\NormalTok{        printf}\OperatorTok{(}\StringTok{"x is true!}\SpecialCharTok{\textbackslash{}n}\StringTok{"}\OperatorTok{);}
    \OperatorTok{\}}
\OperatorTok{\}}
\end{Highlighting}
\end{Shaded}

\end{tcolorbox}

Secara teknis, variabel \hl{\mbox{\texttt{bool}}} seharusnya diisi \hl{\mbox{\texttt{true}}}, \hl{\mbox{\texttt{false}}}, atau hasil ekspresi yang bernilai true atau false. Tapi sebenarnya kamu bisa mengonversi berbagai hal ke \hl{\mbox{\texttt{bool}}}. Ada aturan khusus, tapi yang bernilai nol cenderung menjadi \hl{\mbox{\texttt{false}}}, dan yang bukan nol menjadi true.

Hati-hati kalau keduanya dicampur, karena nilai numerik \hl{\mbox{\texttt{true}}} kemungkinan besar adalah \hl{\mbox{\texttt{1}}}\footnote{Secara teknis hanya satu bit dari sebuah \hl{\mbox{\texttt{char}}} yang dipakai untuk merepresentasikan \hl{\mbox{\texttt{bool}}}, jadi nilainya hanya bisa nol atau satu. Hanya saja, isi bit-bit sisanya (padding) dalam \hl{\mbox{\texttt{char}}} itu tidak ditentukan. Untuk \hl{\mbox{\texttt{false}}}, pasti semuanya nol. Tapi untuk \hl{\mbox{\texttt{true}}}, penulis tidak yakin apakah sisanya harus nol semua.}. Kalau kamu mengandalkan nilai positif lain sebagai true, hasilnya bisa tidak cocok. Misalnya:

\begin{tcolorbox}[enhanced jigsaw,left=3pt,breakable,colback=white!97.2549!black,boxrule=0.2mm]

\begin{Shaded}
\begin{Highlighting}[numbers=left,,]
\NormalTok{printf}\OperatorTok{(}\StringTok{"}\SpecialCharTok{\%d\textbackslash{}n}\StringTok{"}\OperatorTok{,} \KeywordTok{true} \OperatorTok{==} \DecValTok{12}\OperatorTok{);}  \CommentTok{// Mencetak "0", false!}
\end{Highlighting}
\end{Shaded}

\end{tcolorbox}

\index{tipe Boolean|)}

\section{Operator dan Ekspresi}\label{operators}

Operator C seharusnya sudah akrab dari bahasa lain. Kita bahas sekilas beberapa di antaranya di sini.

(Detailnya masih lebih banyak dari ini, tapi bagian ini sudah cukup untuk memulai.)

\subsection{Aritmetika}\label{aritmetika}

\index{Operator Aritmetika|}Semoga ini sudah akrab: \index{+ operator penjumlahan@\texttt{+} operator penjumlahan} \index{- operator pengurangan@\texttt{-} operator pengurangan} \index{* operator perkalian@\texttt{*} operator perkalian} \index{/ operator pembagian@\texttt{/} operator pembagian} \index{\% operator modulus@\texttt{\%} operator modulus}

\begin{tcolorbox}[enhanced jigsaw,left=3pt,breakable,colback=white!97.2549!black,boxrule=0.2mm]

\begin{Shaded}
\begin{Highlighting}[]
\NormalTok{i }\OperatorTok{=}\NormalTok{ i }\OperatorTok{+} \DecValTok{3}\OperatorTok{;}  \CommentTok{// Operator penjumlahan (+) dan penugasan (=), tambahkan 3 ke i}
\NormalTok{i }\OperatorTok{=}\NormalTok{ i }\OperatorTok{{-}} \DecValTok{8}\OperatorTok{;}  \CommentTok{// Pengurangan, kurangi i dengan 8}
\NormalTok{i }\OperatorTok{=}\NormalTok{ i }\OperatorTok{*} \DecValTok{9}\OperatorTok{;}  \CommentTok{// Perkalian}
\NormalTok{i }\OperatorTok{=}\NormalTok{ i }\OperatorTok{/} \DecValTok{2}\OperatorTok{;}  \CommentTok{// Pembagian}
\NormalTok{i }\OperatorTok{=}\NormalTok{ i }\OperatorTok{\%} \DecValTok{5}\OperatorTok{;}  \CommentTok{// Modulo (sisa pembagian)}
\end{Highlighting}
\end{Shaded}

\end{tcolorbox}

Semua itu punya varian singkat. Setiap baris di atas bisa ditulis lebih ringkas sebagai: \index{+= operator penugasan@\texttt{+=} operator penugasan} \index{-= operator penugasan@\texttt{-=} operator penugasan} \index{*= operator penugasan@\texttt{*=} operator penugasan} \index{/= operator penugasan@\texttt{/=} operator penugasan} \index{\%= operator penugasan@\texttt{\%=} operator penugasan}

\begin{tcolorbox}[enhanced jigsaw,left=3pt,breakable,colback=white!97.2549!black,boxrule=0.2mm]

\begin{Shaded}
\begin{Highlighting}[]
\NormalTok{i }\OperatorTok{+=} \DecValTok{3}\OperatorTok{;}  \CommentTok{// Sama dengan "i = i + 3", tambahkan 3 ke i}
\NormalTok{i }\OperatorTok{{-}=} \DecValTok{8}\OperatorTok{;}  \CommentTok{// Sama dengan "i = i {-} 8"}
\NormalTok{i }\OperatorTok{*=} \DecValTok{9}\OperatorTok{;}  \CommentTok{// Sama dengan "i = i * 9"}
\NormalTok{i }\OperatorTok{/=} \DecValTok{2}\OperatorTok{;}  \CommentTok{// Sama dengan "i = i / 2"}
\NormalTok{i }\OperatorTok{\%=} \DecValTok{5}\OperatorTok{;}  \CommentTok{// Sama dengan "i = i \% 5"}
\end{Highlighting}
\end{Shaded}

\end{tcolorbox}

Tidak ada operator pangkat. Kamu harus memakai salah satu varian fungsi \hl{\mbox{\texttt{pow()}}}\index{pow()|} dari \hl{\mbox{\texttt{math.h}}}.

Sekarang kita masuk ke hal-hal yang lebih aneh, yang mungkin tidak ada di bahasa lain!\index{Operator Aritmetika|)}

\subsection{Operator Ternary}\label{operator-ternary}

\index{?: operator ternary@\texttt{?:} operator ternary|(}C juga punya \emph{operator ternary}, yaitu ekspresi yang nilainya bergantung pada hasil sebuah kondisional yang tertanam di dalamnya.

\begin{tcolorbox}[enhanced jigsaw,left=3pt,breakable,colback=white!97.2549!black,boxrule=0.2mm]

\begin{Shaded}
\begin{Highlighting}[]
\CommentTok{// Jika x \textgreater{} 10, tambahkan 17 ke y. Jika tidak, tambahkan 37 ke y.}

\NormalTok{y }\OperatorTok{+=}\NormalTok{ x }\OperatorTok{\textgreater{}} \DecValTok{10}\OperatorTok{?} \DecValTok{17}\OperatorTok{:} \DecValTok{37}\OperatorTok{;}
\end{Highlighting}
\end{Shaded}

\end{tcolorbox}

Berantakan sekali! Tapi kamu akan terbiasa seiring sering membacanya. Supaya lebih jelas, ekspresi di atas bisa ditulis ulang dengan pernyataan \hl{\mbox{\texttt{if}}}:

\begin{tcolorbox}[enhanced jigsaw,left=3pt,breakable,colback=white!97.2549!black,boxrule=0.2mm]

\begin{Shaded}
\begin{Highlighting}[]
\CommentTok{// Ekspresi ini:}

\NormalTok{y }\OperatorTok{+=}\NormalTok{ x }\OperatorTok{\textgreater{}} \DecValTok{10}\OperatorTok{?} \DecValTok{17}\OperatorTok{:} \DecValTok{37}\OperatorTok{;}

\CommentTok{// setara dengan non{-}ekspresi ini:}

\ControlFlowTok{if} \OperatorTok{(}\NormalTok{x }\OperatorTok{\textgreater{}} \DecValTok{10}\OperatorTok{)}
\NormalTok{    y }\OperatorTok{+=} \DecValTok{17}\OperatorTok{;}
\ControlFlowTok{else}
\NormalTok{    y }\OperatorTok{+=} \DecValTok{37}\OperatorTok{;}
\end{Highlighting}
\end{Shaded}

\end{tcolorbox}

Bandingkan keduanya sampai kamu bisa melihat masing-masing komponen operator ternary.

Atau contoh lain, yang mencetak apakah angka di \hl{\mbox{\texttt{x}}} ganjil atau genap:

\begin{tcolorbox}[enhanced jigsaw,left=3pt,breakable,colback=white!97.2549!black,boxrule=0.2mm]

\begin{Shaded}
\begin{Highlighting}[]
\NormalTok{printf}\OperatorTok{(}\StringTok{"The number }\SpecialCharTok{\%d}\StringTok{ is }\SpecialCharTok{\%s}\StringTok{.}\SpecialCharTok{\textbackslash{}n}\StringTok{"}\OperatorTok{,}\NormalTok{ x}\OperatorTok{,}\NormalTok{ x }\OperatorTok{\%} \DecValTok{2} \OperatorTok{==} \DecValTok{0}\OperatorTok{?} \StringTok{"even"}\OperatorTok{:} \StringTok{"odd"}\OperatorTok{);}
\end{Highlighting}
\end{Shaded}

\end{tcolorbox}

Format specifier \hl{\mbox{\texttt{\%s}}} di \hl{\mbox{\texttt{printf()}}}\index{printf()|} berarti ``cetak string''. Kalau ekspresi \hl{\mbox{\texttt{x\ \%\ 2}}} bernilai \hl{\mbox{\texttt{0}}}, nilai seluruh ekspresi ternary adalah string \hl{\mbox{\texttt{"even"}}}. Kalau tidak, nilainya string \hl{\mbox{\texttt{"odd"}}}. Keren sekali!

Perlu diingat bahwa operator ternary bukan kendali alur seperti pernyataan \hl{\mbox{\texttt{if}}}. Ia hanyalah ekspresi yang menghasilkan sebuah nilai. \index{?: operator ternary@\texttt{?:} operator ternary|)}

\subsection{Pre-dan-Post Increment-dan-Decrement}\label{pre-dan-post-increment-dan-decrement}

\index{++ operator increment@\texttt{++} operator increment|(} \index{-- operator decrement@\texttt{--} operator decrement|(} Sekarang kita bahas hal lain yang mungkin belum pernah kamu lihat.

Inilah operator post-increment dan post-decrement yang legendaris:

\begin{tcolorbox}[enhanced jigsaw,left=3pt,breakable,colback=white!97.2549!black,boxrule=0.2mm]

\begin{Shaded}
\begin{Highlighting}[]
\NormalTok{i}\OperatorTok{++;}        \CommentTok{// Tambahkan satu ke i (post{-}increment)}
\NormalTok{i}\OperatorTok{{-}{-};}        \CommentTok{// Kurangi i dengan satu (post{-}decrement)}
\end{Highlighting}
\end{Shaded}

\end{tcolorbox}

Operator ini sangat umum dipakai sebagai versi singkat dari:

\begin{tcolorbox}[enhanced jigsaw,left=3pt,breakable,colback=white!97.2549!black,boxrule=0.2mm]

\begin{Shaded}
\begin{Highlighting}[]
\NormalTok{i }\OperatorTok{+=} \DecValTok{1}\OperatorTok{;}        \CommentTok{// Tambahkan satu ke i}
\NormalTok{i }\OperatorTok{{-}=} \DecValTok{1}\OperatorTok{;}        \CommentTok{// Kurangi i dengan satu}
\end{Highlighting}
\end{Shaded}

\end{tcolorbox}

tapi perbedaannya lebih halus dari itu, dasar penipu cerdas.

Sekarang varian lainnya, pre-increment dan pre-decrement:

\begin{tcolorbox}[enhanced jigsaw,left=3pt,breakable,colback=white!97.2549!black,boxrule=0.2mm]

\begin{Shaded}
\begin{Highlighting}[]
\OperatorTok{++}\NormalTok{i}\OperatorTok{;}        \CommentTok{// Tambahkan satu ke i (pre{-}increment)}
\OperatorTok{{-}{-}}\NormalTok{i}\OperatorTok{;}        \CommentTok{// Kurangi i dengan satu (pre{-}decrement)}
\end{Highlighting}
\end{Shaded}

\end{tcolorbox}

Pada pre-increment dan pre-decrement, nilai variabel di-increment atau di-decrement \emph{sebelum} ekspresi dievaluasi. Setelah itu, ekspresi dievaluasi dengan nilai yang baru.

Pada post-increment dan post-decrement, nilai ekspresi dihitung lebih dulu dengan nilai variabel apa adanya, dan \emph{baru kemudian} variabelnya di-increment atau di-decrement.

Kedua operator ini bisa disematkan dalam ekspresi, seperti ini:

\begin{tcolorbox}[enhanced jigsaw,left=3pt,breakable,colback=white!97.2549!black,boxrule=0.2mm]

\begin{Shaded}
\begin{Highlighting}[]
\NormalTok{i }\OperatorTok{=} \DecValTok{10}\OperatorTok{;}
\NormalTok{j }\OperatorTok{=} \DecValTok{5} \OperatorTok{+}\NormalTok{ i}\OperatorTok{++;}  \CommentTok{// Hitung 5 + i, \_baru kemudian\_ increment i}

\NormalTok{printf}\OperatorTok{(}\StringTok{"}\SpecialCharTok{\%d}\StringTok{, }\SpecialCharTok{\%d\textbackslash{}n}\StringTok{"}\OperatorTok{,}\NormalTok{ i}\OperatorTok{,}\NormalTok{ j}\OperatorTok{);}  \CommentTok{// Mencetak 11, 15}
\end{Highlighting}
\end{Shaded}

\end{tcolorbox}

Bandingkan dengan operator pre-increment:

\begin{tcolorbox}[enhanced jigsaw,left=3pt,breakable,colback=white!97.2549!black,boxrule=0.2mm]

\begin{Shaded}
\begin{Highlighting}[]
\NormalTok{i }\OperatorTok{=} \DecValTok{10}\OperatorTok{;}
\NormalTok{j }\OperatorTok{=} \DecValTok{5} \OperatorTok{+} \OperatorTok{++}\NormalTok{i}\OperatorTok{;}  \CommentTok{// Increment i, \_baru kemudian\_ hitung 5 + i}

\NormalTok{printf}\OperatorTok{(}\StringTok{"}\SpecialCharTok{\%d}\StringTok{, }\SpecialCharTok{\%d\textbackslash{}n}\StringTok{"}\OperatorTok{,}\NormalTok{ i}\OperatorTok{,}\NormalTok{ j}\OperatorTok{);}  \CommentTok{// Mencetak 11, 16}
\end{Highlighting}
\end{Shaded}

\end{tcolorbox}

Teknik ini sering dipakai pada akses dan manipulasi array dan pointer. Dengannya kamu bisa memakai nilai sebuah variabel sekaligus meng-increment atau men-decrement nilai itu, sebelum atau sesudah dipakai.

Tapi sejauh ini, tempat yang paling sering memakainya adalah loop \hl{\mbox{\texttt{for}}}:

\begin{tcolorbox}[enhanced jigsaw,left=3pt,breakable,colback=white!97.2549!black,boxrule=0.2mm]

\begin{Shaded}
\begin{Highlighting}[]
\ControlFlowTok{for} \OperatorTok{(}\NormalTok{i }\OperatorTok{=} \DecValTok{0}\OperatorTok{;}\NormalTok{ i }\OperatorTok{\textless{}} \DecValTok{10}\OperatorTok{;}\NormalTok{ i}\OperatorTok{++)}
\NormalTok{    printf}\OperatorTok{(}\StringTok{"i is }\SpecialCharTok{\%d\textbackslash{}n}\StringTok{"}\OperatorTok{,}\NormalTok{ i}\OperatorTok{);}
\end{Highlighting}
\end{Shaded}

\end{tcolorbox}

Soal itu kita bahas nanti. \index{++ operator increment@\texttt{++} operator increment|)} \index{-- operator decrement@\texttt{--} operator decrement|)}

\subsection{Operator Koma}\label{operator-koma}

\index{, operator koma@\texttt{,} operator koma|(} Ini cara yang jarang dipakai untuk memisahkan ekspresi yang akan dijalankan dari kiri ke kanan:

\begin{tcolorbox}[enhanced jigsaw,left=3pt,breakable,colback=white!97.2549!black,boxrule=0.2mm]

\begin{Shaded}
\begin{Highlighting}[]
\NormalTok{x }\OperatorTok{=} \DecValTok{10}\OperatorTok{,}\NormalTok{ y }\OperatorTok{=} \DecValTok{20}\OperatorTok{;}  \CommentTok{// Mula{-}mula tugaskan 10 ke x, lalu 20 ke y}
\end{Highlighting}
\end{Shaded}

\end{tcolorbox}

Kelihatannya agak konyol, karena koma itu bisa saja diganti titik koma, kan?

\begin{tcolorbox}[enhanced jigsaw,left=3pt,breakable,colback=white!97.2549!black,boxrule=0.2mm]

\begin{Shaded}
\begin{Highlighting}[]
\NormalTok{x }\OperatorTok{=} \DecValTok{10}\OperatorTok{;}\NormalTok{ y }\OperatorTok{=} \DecValTok{20}\OperatorTok{;}  \CommentTok{// Mula{-}mula tugaskan 10 ke x, lalu 20 ke y}
\end{Highlighting}
\end{Shaded}

\end{tcolorbox}

Tapi keduanya sedikit berbeda. Yang terakhir berupa dua ekspresi terpisah, sedangkan yang pertama hanya satu ekspresi tunggal!

Dengan operator koma, nilai ekspresi koma adalah nilai ekspresi yang paling kanan:

\begin{tcolorbox}[enhanced jigsaw,left=3pt,breakable,colback=white!97.2549!black,boxrule=0.2mm]

\begin{Shaded}
\begin{Highlighting}[]
\NormalTok{x }\OperatorTok{=} \OperatorTok{(}\DecValTok{1}\OperatorTok{,} \DecValTok{2}\OperatorTok{,} \DecValTok{3}\OperatorTok{);}

\NormalTok{printf}\OperatorTok{(}\StringTok{"x is }\SpecialCharTok{\%d\textbackslash{}n}\StringTok{"}\OperatorTok{,}\NormalTok{ x}\OperatorTok{);}  \CommentTok{// Mencetak 3, karena 3 adalah yang paling kanan di daftar koma}
\end{Highlighting}
\end{Shaded}

\end{tcolorbox}

Tapi itu pun terasa dipaksakan. Tempat yang umum untuk operator koma adalah loop \hl{\mbox{\texttt{for}}}, untuk melakukan beberapa hal sekaligus di setiap bagiannya:

\begin{tcolorbox}[enhanced jigsaw,left=3pt,breakable,colback=white!97.2549!black,boxrule=0.2mm]

\begin{Shaded}
\begin{Highlighting}[]
\ControlFlowTok{for} \OperatorTok{(}\NormalTok{i }\OperatorTok{=} \DecValTok{0}\OperatorTok{,}\NormalTok{ j }\OperatorTok{=} \DecValTok{10}\OperatorTok{;}\NormalTok{ i }\OperatorTok{\textless{}} \DecValTok{100}\OperatorTok{;}\NormalTok{ i}\OperatorTok{++,}\NormalTok{ j}\OperatorTok{++)}
\NormalTok{    printf}\OperatorTok{(}\StringTok{"}\SpecialCharTok{\%d}\StringTok{, }\SpecialCharTok{\%d\textbackslash{}n}\StringTok{"}\OperatorTok{,}\NormalTok{ i}\OperatorTok{,}\NormalTok{ j}\OperatorTok{);}
\end{Highlighting}
\end{Shaded}

\end{tcolorbox}

Kita akan kembali ke sana nanti. \index{, operator koma@\texttt{,} operator koma|)}

\subsection{Operator Kondisional}\label{operator-kondisional}

\index{Operator Kondisional|(} Untuk nilai Boolean, tersedia banyak operator standar: \index{operator perbandingan} \index{== operator kesamaan@\texttt{==} operator kesamaan} \index{"!= operator ketidaksamaan@\texttt{"!=} operator ketidaksamaan} \index{< operator kurang dari@\texttt{<} operator kurang dari} \index{> operator lebih dari@\texttt{>} operator lebih dari} \index{<= operator kurang dari atau sama dengan@\texttt{<=} operator kurang dari atau sama dengan} \index{>= operator lebih dari atau sama dengan@\texttt{>=} operator lebih dari atau sama dengan}

\begin{tcolorbox}[enhanced jigsaw,left=3pt,breakable,colback=white!97.2549!black,boxrule=0.2mm]

\begin{Shaded}
\begin{Highlighting}[]
\NormalTok{a }\OperatorTok{==}\NormalTok{ b}\OperatorTok{;}  \CommentTok{// True jika a sama dengan b}
\NormalTok{a }\OperatorTok{!=}\NormalTok{ b}\OperatorTok{;}  \CommentTok{// True jika a tidak sama dengan b}
\NormalTok{a }\OperatorTok{\textless{}}\NormalTok{ b}\OperatorTok{;}   \CommentTok{// True jika a kurang dari b}
\NormalTok{a }\OperatorTok{\textgreater{}}\NormalTok{ b}\OperatorTok{;}   \CommentTok{// True jika a lebih dari b}
\NormalTok{a }\OperatorTok{\textless{}=}\NormalTok{ b}\OperatorTok{;}  \CommentTok{// True jika a kurang dari atau sama dengan b}
\NormalTok{a }\OperatorTok{\textgreater{}=}\NormalTok{ b}\OperatorTok{;}  \CommentTok{// True jika a lebih dari atau sama dengan b}
\end{Highlighting}
\end{Shaded}

\end{tcolorbox}

Jangan tertukar antara penugasan (\hl{\mbox{\texttt{=}}}) dan perbandingan (\hl{\mbox{\texttt{==}}})! Pakai dua tanda sama dengan untuk membandingkan, satu untuk penugasan.

Ekspresi perbandingan bisa dipakai dalam pernyataan \hl{\mbox{\texttt{if}}}:

\begin{tcolorbox}[enhanced jigsaw,left=3pt,breakable,colback=white!97.2549!black,boxrule=0.2mm]

\begin{Shaded}
\begin{Highlighting}[]
\ControlFlowTok{if} \OperatorTok{(}\NormalTok{a }\OperatorTok{\textless{}=} \DecValTok{10}\OperatorTok{)}
\NormalTok{    printf}\OperatorTok{(}\StringTok{"Success!}\SpecialCharTok{\textbackslash{}n}\StringTok{"}\OperatorTok{);}
\end{Highlighting}
\end{Shaded}

\end{tcolorbox}

\index{Operator Kondisional|)}

\subsection{Operator Boolean}\label{operator-boolean}

\index{Operator Boolean|(} Ekspresi kondisional bisa dirangkai atau diubah dengan operator Boolean untuk \emph{dan}, \emph{atau}, dan \emph{tidak}. \index{\&\& boolean AND@\texttt{\&\&} boolean AND} \index{"! boolean NOT@\texttt{"!} boolean NOT} \index{\textbar{}\textbar{} boolean OR@\texttt{\textbar{}\textbar{}} boolean OR}

{\def\LTcaptype{none} % do not increment counter
\begin{longtable}[]{@{}cc@{}}
\toprule\noalign{}
Operator & Makna Boolean \\
\midrule\noalign{}
\endhead
\bottomrule\noalign{}
\endlastfoot
\hl{\mbox{\texttt{\&\&}}} & and \\
\hl{\mbox{\texttt{\textbar{}\textbar{}}}} & or \\
\hl{\mbox{\texttt{!}}} & not \\
\end{longtable}
}

Contoh Boolean ``and'':

\begin{tcolorbox}[enhanced jigsaw,left=3pt,breakable,colback=white!97.2549!black,boxrule=0.2mm]

\begin{Shaded}
\begin{Highlighting}[]
\CommentTok{// Lakukan sesuatu jika x kurang dari 10 dan y lebih dari 20:}

\ControlFlowTok{if} \OperatorTok{(}\NormalTok{x }\OperatorTok{\textless{}} \DecValTok{10} \OperatorTok{\&\&}\NormalTok{ y }\OperatorTok{\textgreater{}} \DecValTok{20}\OperatorTok{)}
\NormalTok{    printf}\OperatorTok{(}\StringTok{"Doing something!}\SpecialCharTok{\textbackslash{}n}\StringTok{"}\OperatorTok{);}
\end{Highlighting}
\end{Shaded}

\end{tcolorbox}

Contoh Boolean ``not'':

\begin{tcolorbox}[enhanced jigsaw,left=3pt,breakable,colback=white!97.2549!black,boxrule=0.2mm]

\begin{Shaded}
\begin{Highlighting}[]
\ControlFlowTok{if} \OperatorTok{(!(}\NormalTok{x }\OperatorTok{\textless{}} \DecValTok{12}\OperatorTok{))}
\NormalTok{    printf}\OperatorTok{(}\StringTok{"x is not less than 12}\SpecialCharTok{\textbackslash{}n}\StringTok{"}\OperatorTok{);}
\end{Highlighting}
\end{Shaded}

\end{tcolorbox}

\hl{\mbox{\texttt{!}}} punya prioritas (precedence) lebih tinggi daripada operator Boolean lainnya, jadi dalam kasus ini kita harus memakai tanda kurung.

Tentu saja, itu sama saja dengan:

\begin{tcolorbox}[enhanced jigsaw,left=3pt,breakable,colback=white!97.2549!black,boxrule=0.2mm]

\begin{Shaded}
\begin{Highlighting}[]
\ControlFlowTok{if} \OperatorTok{(}\NormalTok{x }\OperatorTok{\textgreater{}=} \DecValTok{12}\OperatorTok{)}
\NormalTok{    printf}\OperatorTok{(}\StringTok{"x is not less than 12}\SpecialCharTok{\textbackslash{}n}\StringTok{"}\OperatorTok{);}
\end{Highlighting}
\end{Shaded}

\end{tcolorbox}

tapi namanya juga contoh! \index{Operator Boolean|)}

\subsection{\texorpdfstring{Operator \hl{\mbox{\texttt{sizeof}}}}{Operator sizeof}}\label{sizeof-operator}

\index{sizeof operator@\texttt{sizeof} operator|(} Operator ini memberi tahu ukuran (dalam byte) yang dipakai sebuah variabel atau tipe data di memori.

Lebih tepatnya, ia memberi tahu ukuran (dalam byte) yang dipakai oleh \emph{tipe dari suatu ekspresi} (yang bisa saja hanya satu variabel) di memori.

Ukurannya bisa berbeda di tiap sistem, kecuali untuk \hl{\mbox{\texttt{char}}}\index{char tipe@\texttt{char} tipe} dan variannya (yang selalu 1 byte).

Mungkin sekarang ini belum terlihat berguna, tapi kita akan sering merujuknya, jadi layak dibahas.

Karena \hl{\mbox{\texttt{sizeof}}} menghitung jumlah byte yang dibutuhkan untuk menyimpan sebuah tipe, kamu mungkin mengira ia mengembalikan \hl{\mbox{\texttt{int}}}. Atau\ldots{} karena ukuran tidak mungkin negatif, mungkin \hl{\mbox{\texttt{unsigned}}}?

\index{size\_t tipe@\texttt{size\_t} tipe|(} Ternyata C punya tipe khusus untuk menampung nilai kembalian \hl{\mbox{\texttt{sizeof}}}, yaitu \hl{\mbox{\texttt{size\_t}}}, dibaca ``\emph{size tee}''\footnote{\hl{\mbox{\texttt{\_t}}} kependekan dari \hl{\mbox{\texttt{type}}}.}. Yang kita tahu hanyalah bahwa ini tipe integer tak bertanda (unsigned) yang sanggup menyimpan ukuran dalam byte dari apa pun yang bisa diberikan ke \hl{\mbox{\texttt{sizeof}}}.

\hl{\mbox{\texttt{size\_t}}} muncul di banyak tempat yang melibatkan jumlah sesuatu, baik sebagai argumen maupun nilai kembalian. Anggap saja ia nilai yang menyatakan ``banyaknya''. \index{size\_t tipe@\texttt{size\_t} tipe|)}

\hl{\mbox{\texttt{sizeof}}} bisa diterapkan pada variabel atau ekspresi:

\begin{tcolorbox}[enhanced jigsaw,left=3pt,breakable,colback=white!97.2549!black,boxrule=0.2mm]

\begin{Shaded}
\begin{Highlighting}[]
\DataTypeTok{int}\NormalTok{ a }\OperatorTok{=} \DecValTok{999}\OperatorTok{;}

\CommentTok{// \%zu adalah format specifier untuk tipe size\_t}
\CommentTok{// Jika compiler{-}mu menolak bagian "z", buang saja}

\NormalTok{printf}\OperatorTok{(}\StringTok{"}\SpecialCharTok{\%zu\textbackslash{}n}\StringTok{"}\OperatorTok{,} \KeywordTok{sizeof}\NormalTok{ a}\OperatorTok{);}      \CommentTok{// Mencetak 4 di sistem penulis}
\NormalTok{printf}\OperatorTok{(}\StringTok{"}\SpecialCharTok{\%zu\textbackslash{}n}\StringTok{"}\OperatorTok{,} \KeywordTok{sizeof}\OperatorTok{(}\DecValTok{2} \OperatorTok{+} \DecValTok{7}\OperatorTok{));} \CommentTok{// Mencetak 4 di sistem penulis}
\NormalTok{printf}\OperatorTok{(}\StringTok{"}\SpecialCharTok{\%zu\textbackslash{}n}\StringTok{"}\OperatorTok{,} \KeywordTok{sizeof} \FloatTok{3.14}\OperatorTok{);}   \CommentTok{// Mencetak 8 di sistem penulis}

\CommentTok{// Jika perlu mencetak nilai size\_t negatif, pakai \%zd}
\end{Highlighting}
\end{Shaded}

\end{tcolorbox}

Ingat: yang dihasilkan adalah ukuran dalam byte dari \emph{tipe} ekspresi, bukan ukuran ekspresi itu sendiri. Itulah sebabnya ukuran \hl{\mbox{\texttt{2+7}}} sama dengan ukuran \hl{\mbox{\texttt{a}}}---keduanya bertipe \hl{\mbox{\texttt{int}}}. Kita akan kembali ke angka \hl{\mbox{\texttt{4}}} ini di blok kode berikutnya\ldots{}

\ldots yang menunjukkan bahwa \hl{\mbox{\texttt{sizeof}}} juga bisa diterapkan pada sebuah tipe (perhatikan bahwa tanda kurung wajib di sekitar nama tipe, tidak seperti pada ekspresi):

\begin{tcolorbox}[enhanced jigsaw,left=3pt,breakable,colback=white!97.2549!black,boxrule=0.2mm]

\begin{Shaded}
\begin{Highlighting}[]
\NormalTok{printf}\OperatorTok{(}\StringTok{"}\SpecialCharTok{\%zu\textbackslash{}n}\StringTok{"}\OperatorTok{,} \KeywordTok{sizeof}\OperatorTok{(}\DataTypeTok{int}\OperatorTok{));}   \CommentTok{// Mencetak 4 di sistem penulis}
\NormalTok{printf}\OperatorTok{(}\StringTok{"}\SpecialCharTok{\%zu\textbackslash{}n}\StringTok{"}\OperatorTok{,} \KeywordTok{sizeof}\OperatorTok{(}\DataTypeTok{char}\OperatorTok{));}  \CommentTok{// Mencetak 1 di semua sistem}
\end{Highlighting}
\end{Shaded}

\end{tcolorbox}

Perlu dicatat bahwa \hl{\mbox{\texttt{sizeof}}} adalah operasi \emph{compile-time}\footnote{Kecuali untuk variable length array---tapi itu cerita untuk lain waktu.}. Hasilnya ditentukan sepenuhnya saat compile-time, bukan saat runtime.

Kita akan memakainya nanti. \index{sizeof operator@\texttt{sizeof} operator|)}

\section{Kendali Alur}\label{kendali-alur}

\index{Kendali Alur|(} Boolean memang bagus, tapi kita tidak akan ke mana-mana kalau tidak bisa mengendalikan alur program. Mari kita lihat beberapa konstruksi: \hl{\mbox{\texttt{if}}}, \hl{\mbox{\texttt{for}}}, \hl{\mbox{\texttt{while}}}, dan \hl{\mbox{\texttt{do-while}}}.

Pertama, catatan umum tentang pernyataan dan blok pernyataan, dipersembahkan oleh developer C lokal yang ramah ini:

Setelah pernyataan seperti \hl{\mbox{\texttt{if}}} atau \hl{\mbox{\texttt{while}}}, kamu bisa menaruh satu pernyataan tunggal untuk dieksekusi, atau sebuah blok pernyataan yang semuanya dieksekusi berurutan.

\index{if pernyataan@\texttt{if} pernyataan|(} Kita mulai dengan satu pernyataan:

\begin{tcolorbox}[enhanced jigsaw,left=3pt,breakable,colback=white!97.2549!black,boxrule=0.2mm]

\begin{Shaded}
\begin{Highlighting}[]
\ControlFlowTok{if} \OperatorTok{(}\NormalTok{x }\OperatorTok{==} \DecValTok{10}\OperatorTok{)}\NormalTok{ printf}\OperatorTok{(}\StringTok{"x is 10}\SpecialCharTok{\textbackslash{}n}\StringTok{"}\OperatorTok{);}
\end{Highlighting}
\end{Shaded}

\end{tcolorbox}

Ini juga kadang ditulis di baris terpisah. (Whitespace hampir tidak berpengaruh di C---tidak seperti Python.)

\begin{tcolorbox}[enhanced jigsaw,left=3pt,breakable,colback=white!97.2549!black,boxrule=0.2mm]

\begin{Shaded}
\begin{Highlighting}[]
\ControlFlowTok{if} \OperatorTok{(}\NormalTok{x }\OperatorTok{==} \DecValTok{10}\OperatorTok{)}
\NormalTok{    printf}\OperatorTok{(}\StringTok{"x is 10}\SpecialCharTok{\textbackslash{}n}\StringTok{"}\OperatorTok{);}
\end{Highlighting}
\end{Shaded}

\end{tcolorbox}

Tapi bagaimana kalau kamu ingin beberapa hal terjadi sekaligus karena satu kondisi? Pakai kurung kurawal untuk menandai sebuah \emph{blok} atau \emph{pernyataan majemuk}.

\begin{tcolorbox}[enhanced jigsaw,left=3pt,breakable,colback=white!97.2549!black,boxrule=0.2mm]

\begin{Shaded}
\begin{Highlighting}[]
\ControlFlowTok{if} \OperatorTok{(}\NormalTok{x }\OperatorTok{==} \DecValTok{10}\OperatorTok{)} \OperatorTok{\{}
\NormalTok{    printf}\OperatorTok{(}\StringTok{"x is 10}\SpecialCharTok{\textbackslash{}n}\StringTok{"}\OperatorTok{);}
\NormalTok{    printf}\OperatorTok{(}\StringTok{"And also this happens when x is 10}\SpecialCharTok{\textbackslash{}n}\StringTok{"}\OperatorTok{);}
\OperatorTok{\}}
\end{Highlighting}
\end{Shaded}

\end{tcolorbox}

Gaya yang sangat umum adalah \emph{selalu} memakai kurung kurawal, meskipun sebenarnya tidak perlu:

\begin{tcolorbox}[enhanced jigsaw,left=3pt,breakable,colback=white!97.2549!black,boxrule=0.2mm]

\begin{Shaded}
\begin{Highlighting}[]
\ControlFlowTok{if} \OperatorTok{(}\NormalTok{x }\OperatorTok{==} \DecValTok{10}\OperatorTok{)} \OperatorTok{\{}
\NormalTok{    printf}\OperatorTok{(}\StringTok{"x is 10}\SpecialCharTok{\textbackslash{}n}\StringTok{"}\OperatorTok{);}
\OperatorTok{\}}
\end{Highlighting}
\end{Shaded}

\end{tcolorbox}

Sebagian developer merasa kodenya lebih mudah dibaca, dan terhindar dari error seperti ini, ketika baris yang tampak berada di dalam blok \hl{\mbox{\texttt{if}}} ternyata tidak.

\begin{tcolorbox}[enhanced jigsaw,left=3pt,breakable,colback=white!97.2549!black,boxrule=0.2mm]

\begin{Shaded}
\begin{Highlighting}[]
\CommentTok{// CONTOH ERROR BURUK}

\ControlFlowTok{if} \OperatorTok{(}\NormalTok{x }\OperatorTok{==} \DecValTok{10}\OperatorTok{)}
\NormalTok{    printf}\OperatorTok{(}\StringTok{"This happens if x is 10}\SpecialCharTok{\textbackslash{}n}\StringTok{"}\OperatorTok{);}
\NormalTok{    printf}\OperatorTok{(}\StringTok{"This happens ALWAYS}\SpecialCharTok{\textbackslash{}n}\StringTok{"}\OperatorTok{);}  \CommentTok{// Kejutan!! Tanpa kondisi!}
\end{Highlighting}
\end{Shaded}

\end{tcolorbox}

\hl{\mbox{\texttt{while}}}, \hl{\mbox{\texttt{for}}}, dan konstruksi loop lainnya bekerja dengan cara yang sama seperti contoh di atas. Kalau ingin melakukan beberapa hal dalam loop atau setelah \hl{\mbox{\texttt{if}}}, bungkus semuanya dalam kurung kurawal.

Dengan kata lain, \hl{\mbox{\texttt{if}}} menjalankan satu hal yang ada setelahnya. Satu hal itu bisa berupa satu pernyataan atau satu blok pernyataan. \index{if pernyataan@\texttt{if} pernyataan|)}

\subsection{\texorpdfstring{Pernyataan \hl{\mbox{\texttt{if}}}-\hl{\mbox{\texttt{else}}}}{Pernyataan if-else}}\label{ifstat}

\index{if-else pernyataan@\texttt{if}-\texttt{else} pernyataan|(} Kita sudah memakai \hl{\mbox{\texttt{if}}} di beberapa contoh, karena kemungkinan besar kamu sudah pernah melihatnya di bahasa lain, tapi ini satu contoh lagi:

\begin{tcolorbox}[enhanced jigsaw,left=3pt,breakable,colback=white!97.2549!black,boxrule=0.2mm]

\begin{Shaded}
\begin{Highlighting}[]
\DataTypeTok{int}\NormalTok{ i }\OperatorTok{=} \DecValTok{10}\OperatorTok{;}

\ControlFlowTok{if} \OperatorTok{(}\NormalTok{i }\OperatorTok{\textgreater{}} \DecValTok{10}\OperatorTok{)} \OperatorTok{\{}
\NormalTok{    printf}\OperatorTok{(}\StringTok{"Yes, i is greater than 10.}\SpecialCharTok{\textbackslash{}n}\StringTok{"}\OperatorTok{);}
\NormalTok{    printf}\OperatorTok{(}\StringTok{"And this will also print if i is greater than 10.}\SpecialCharTok{\textbackslash{}n}\StringTok{"}\OperatorTok{);}
\OperatorTok{\}}

\ControlFlowTok{if} \OperatorTok{(}\NormalTok{i }\OperatorTok{\textless{}=} \DecValTok{10}\OperatorTok{)}\NormalTok{ printf}\OperatorTok{(}\StringTok{"i is less than or equal to 10.}\SpecialCharTok{\textbackslash{}n}\StringTok{"}\OperatorTok{);}
\end{Highlighting}
\end{Shaded}

\end{tcolorbox}

Pada kode contoh, pesan dicetak kalau \hl{\mbox{\texttt{i}}} lebih besar dari 10, kalau tidak eksekusi lanjut ke baris berikutnya. Perhatikan kurung kurawal setelah pernyataan \hl{\mbox{\texttt{if}}}; kalau kondisinya true, yang dieksekusi adalah pernyataan tunggal tepat setelah \hl{\mbox{\texttt{if}}}, atau sekumpulan kode dalam kurung kurawal setelah \hl{\mbox{\texttt{if}}}. Perilaku \emph{blok kode} seperti ini berlaku untuk semua pernyataan.

Tentu saja, karena C menyenangkan dengan caranya sendiri, kamu juga bisa melakukan sesuatu saat kondisinya false lewat klausa \hl{\mbox{\texttt{else}}} pada \hl{\mbox{\texttt{if}}}:

\begin{tcolorbox}[enhanced jigsaw,left=3pt,breakable,colback=white!97.2549!black,boxrule=0.2mm]

\begin{Shaded}
\begin{Highlighting}[]
\DataTypeTok{int}\NormalTok{ i }\OperatorTok{=} \DecValTok{99}\OperatorTok{;}

\ControlFlowTok{if} \OperatorTok{(}\NormalTok{i }\OperatorTok{==} \DecValTok{10}\OperatorTok{)}
\NormalTok{    printf}\OperatorTok{(}\StringTok{"i is 10!}\SpecialCharTok{\textbackslash{}n}\StringTok{"}\OperatorTok{);}

\ControlFlowTok{else} \OperatorTok{\{}
\NormalTok{    printf}\OperatorTok{(}\StringTok{"i is decidedly not 10.}\SpecialCharTok{\textbackslash{}n}\StringTok{"}\OperatorTok{);}
\NormalTok{    printf}\OperatorTok{(}\StringTok{"Which irritates me a little, frankly.}\SpecialCharTok{\textbackslash{}n}\StringTok{"}\OperatorTok{);}
\OperatorTok{\}}
\end{Highlighting}
\end{Shaded}

\end{tcolorbox}

Kamu bahkan bisa merangkainya secara bertingkat untuk menguji berbagai kondisi, seperti ini:

\begin{tcolorbox}[enhanced jigsaw,left=3pt,breakable,colback=white!97.2549!black,boxrule=0.2mm]

\begin{Shaded}
\begin{Highlighting}[]
\DataTypeTok{int}\NormalTok{ i }\OperatorTok{=} \DecValTok{99}\OperatorTok{;}

\ControlFlowTok{if} \OperatorTok{(}\NormalTok{i }\OperatorTok{==} \DecValTok{10}\OperatorTok{)}
\NormalTok{    printf}\OperatorTok{(}\StringTok{"i is 10!}\SpecialCharTok{\textbackslash{}n}\StringTok{"}\OperatorTok{);}

\ControlFlowTok{else} \ControlFlowTok{if} \OperatorTok{(}\NormalTok{i }\OperatorTok{==} \DecValTok{20}\OperatorTok{)}
\NormalTok{    printf}\OperatorTok{(}\StringTok{"i is 20!}\SpecialCharTok{\textbackslash{}n}\StringTok{"}\OperatorTok{);}

\ControlFlowTok{else} \ControlFlowTok{if} \OperatorTok{(}\NormalTok{i }\OperatorTok{==} \DecValTok{99}\OperatorTok{)} \OperatorTok{\{}
\NormalTok{    printf}\OperatorTok{(}\StringTok{"i is 99! My favorite}\SpecialCharTok{\textbackslash{}n}\StringTok{"}\OperatorTok{);}
\NormalTok{    printf}\OperatorTok{(}\StringTok{"I can\textquotesingle{}t tell you how happy I am.}\SpecialCharTok{\textbackslash{}n}\StringTok{"}\OperatorTok{);}
\NormalTok{    printf}\OperatorTok{(}\StringTok{"Really.}\SpecialCharTok{\textbackslash{}n}\StringTok{"}\OperatorTok{);}
\OperatorTok{\}}
    
\ControlFlowTok{else}
\NormalTok{    printf}\OperatorTok{(}\StringTok{"i is some crazy number I\textquotesingle{}ve never heard of.}\SpecialCharTok{\textbackslash{}n}\StringTok{"}\OperatorTok{);}
\end{Highlighting}
\end{Shaded}

\end{tcolorbox}

Kalau kamu sampai menempuh jalan itu, coba lihat pernyataan \hyperref[switch-statement]{\hl{\mbox{\texttt{switch}}}} yang mungkin solusi lebih baik. Syaratnya, \hl{\mbox{\texttt{switch}}} hanya bisa membandingkan kesamaan dengan angka konstan. Rangkaian \hl{\mbox{\texttt{if}}}-\hl{\mbox{\texttt{else}}} di atas bisa memeriksa ketidaksamaan, rentang, variabel, atau apa pun yang bisa dibuat dalam ekspresi kondisional. \index{if-else pernyataan@\texttt{if}-\texttt{else} pernyataan|)}

\subsection{\texorpdfstring{Pernyataan \hl{\mbox{\texttt{while}}}}{Pernyataan while}}\label{whilestat}

\index{while pernyataan@\texttt{while} pernyataan|(} \hl{\mbox{\texttt{while}}} adalah konstruksi loop biasa yang lazim. Lakukan sesuatu selama ekspresi kondisinya bernilai true.

Yuk, coba!

\begin{tcolorbox}[enhanced jigsaw,left=3pt,breakable,colback=white!97.2549!black,boxrule=0.2mm]

\begin{Shaded}
\begin{Highlighting}[]
\CommentTok{// Cetak output berikut:}
\CommentTok{//}
\CommentTok{//   i is now 0!}
\CommentTok{//   i is now 1!}
\CommentTok{//   [ sama saja untuk 2 sampai 7 ]}
\CommentTok{//   i is now 8!}
\CommentTok{//   i is now 9!}

\DataTypeTok{int}\NormalTok{ i }\OperatorTok{=} \DecValTok{0}\OperatorTok{;}

\ControlFlowTok{while} \OperatorTok{(}\NormalTok{i }\OperatorTok{\textless{}} \DecValTok{10}\OperatorTok{)} \OperatorTok{\{}
\NormalTok{    printf}\OperatorTok{(}\StringTok{"i is now }\SpecialCharTok{\%d}\StringTok{!}\SpecialCharTok{\textbackslash{}n}\StringTok{"}\OperatorTok{,}\NormalTok{ i}\OperatorTok{);}
\NormalTok{    i}\OperatorTok{++;}
\OperatorTok{\}}

\NormalTok{printf}\OperatorTok{(}\StringTok{"All done!}\SpecialCharTok{\textbackslash{}n}\StringTok{"}\OperatorTok{);}
\end{Highlighting}
\end{Shaded}

\end{tcolorbox}

Itulah loop dasar. C juga punya loop \hl{\mbox{\texttt{for}}} yang akan membuat contoh tadi lebih rapi.

Pemakaian \hl{\mbox{\texttt{while}}} yang cukup sering adalah infinite loop, yaitu mengulang selama true:

\begin{tcolorbox}[enhanced jigsaw,left=3pt,breakable,colback=white!97.2549!black,boxrule=0.2mm]

\begin{Shaded}
\begin{Highlighting}[]
\ControlFlowTok{while} \OperatorTok{(}\DecValTok{1}\OperatorTok{)} \OperatorTok{\{}
\NormalTok{    printf}\OperatorTok{(}\StringTok{"1 is always true, so this repeats forever.}\SpecialCharTok{\textbackslash{}n}\StringTok{"}\OperatorTok{);}
\OperatorTok{\}}
\end{Highlighting}
\end{Shaded}

\end{tcolorbox}

\index{while pernyataan@\texttt{while} pernyataan|)}

\subsection{\texorpdfstring{Pernyataan \hl{\mbox{\texttt{do}}}-\hl{\mbox{\texttt{while}}}}{Pernyataan do-while}}\label{dowhilestat}

\index{do-while pernyataan@\texttt{do}-\texttt{while} pernyataan|(} Setelah menguasai pernyataan \hl{\mbox{\texttt{while}}}, mari kita lihat saudara dekatnya, \hl{\mbox{\texttt{do}}}-\hl{\mbox{\texttt{while}}}.

Keduanya pada dasarnya sama, kecuali kalau kondisi loop bernilai false pada pass pertama: \hl{\mbox{\texttt{do}}}-\hl{\mbox{\texttt{while}}} tetap dieksekusi sekali, sedangkan \hl{\mbox{\texttt{while}}} tidak dieksekusi sama sekali. Dengan kata lain, pengujian apakah blok akan dieksekusi terjadi di \emph{akhir} blok pada \hl{\mbox{\texttt{do}}}-\hl{\mbox{\texttt{while}}}, dan di \emph{awal} blok pada \hl{\mbox{\texttt{while}}}.

Kita lihat lewat contoh:

\begin{tcolorbox}[enhanced jigsaw,left=3pt,breakable,colback=white!97.2549!black,boxrule=0.2mm]

\begin{Shaded}
\begin{Highlighting}[]
\CommentTok{// Memakai pernyataan while:}

\NormalTok{i }\OperatorTok{=} \DecValTok{10}\OperatorTok{;}

\CommentTok{// ini tidak dieksekusi karena i tidak kurang dari 10:}
\ControlFlowTok{while}\OperatorTok{(}\NormalTok{i }\OperatorTok{\textless{}} \DecValTok{10}\OperatorTok{)} \OperatorTok{\{}
\NormalTok{    printf}\OperatorTok{(}\StringTok{"while: i is }\SpecialCharTok{\%d\textbackslash{}n}\StringTok{"}\OperatorTok{,}\NormalTok{ i}\OperatorTok{);}
\NormalTok{    i}\OperatorTok{++;}
\OperatorTok{\}}

\CommentTok{// Memakai pernyataan do{-}while:}

\NormalTok{i }\OperatorTok{=} \DecValTok{10}\OperatorTok{;}

\CommentTok{// ini dieksekusi sekali, karena kondisi loop baru diperiksa setelah}
\CommentTok{// badan loop dijalankan:}

\ControlFlowTok{do} \OperatorTok{\{}
\NormalTok{    printf}\OperatorTok{(}\StringTok{"do{-}while: i is }\SpecialCharTok{\%d\textbackslash{}n}\StringTok{"}\OperatorTok{,}\NormalTok{ i}\OperatorTok{);}
\NormalTok{    i}\OperatorTok{++;}
\OperatorTok{\}} \ControlFlowTok{while} \OperatorTok{(}\NormalTok{i }\OperatorTok{\textless{}} \DecValTok{10}\OperatorTok{);}

\NormalTok{printf}\OperatorTok{(}\StringTok{"All done!}\SpecialCharTok{\textbackslash{}n}\StringTok{"}\OperatorTok{);}
\end{Highlighting}
\end{Shaded}

\end{tcolorbox}

Perhatikan bahwa pada kedua kasus, kondisi loop langsung bernilai false. Pada \hl{\mbox{\texttt{while}}}, kondisinya gagal, sehingga blok kodenya tidak pernah dieksekusi. Tapi pada \hl{\mbox{\texttt{do}}}-\hl{\mbox{\texttt{while}}}, kondisi diperiksa \emph{setelah} blok kode dieksekusi, jadi blok itu selalu berjalan setidaknya sekali. Di sini, blok mencetak pesan, meng-increment \hl{\mbox{\texttt{i}}}, lalu kondisinya gagal, dan program lanjut ke output ``All done!''.

Intinya: kalau kamu ingin loop dieksekusi setidaknya sekali, apa pun kondisinya, pakai \hl{\mbox{\texttt{do}}}-\hl{\mbox{\texttt{while}}}.

Semua contoh ini sebenarnya lebih cocok memakai loop \hl{\mbox{\texttt{for}}}. Mari kita coba sesuatu yang kurang deterministik---ulangi sampai muncul bilangan acak tertentu!

\begin{tcolorbox}[enhanced jigsaw,left=3pt,breakable,colback=white!97.2549!black,boxrule=0.2mm]

\begin{Shaded}
\begin{Highlighting}[numbers=left,,]
\PreprocessorTok{\#include }\ImportTok{\textless{}stdio.h\textgreater{}}\PreprocessorTok{   }\CommentTok{// Untuk printf}
\PreprocessorTok{\#include }\ImportTok{\textless{}stdlib.h\textgreater{}}\PreprocessorTok{  }\CommentTok{// Untuk rand}

\DataTypeTok{int}\NormalTok{ main}\OperatorTok{(}\DataTypeTok{void}\OperatorTok{)}
\OperatorTok{\{}
    \DataTypeTok{int}\NormalTok{ r}\OperatorTok{;}

    \ControlFlowTok{do} \OperatorTok{\{}
\NormalTok{        r }\OperatorTok{=}\NormalTok{ rand}\OperatorTok{()} \OperatorTok{\%} \DecValTok{100}\OperatorTok{;} \CommentTok{// Ambil bilangan acak antara 0 dan 99}
\NormalTok{        printf}\OperatorTok{(}\StringTok{"}\SpecialCharTok{\%d\textbackslash{}n}\StringTok{"}\OperatorTok{,}\NormalTok{ r}\OperatorTok{);}
    \OperatorTok{\}} \ControlFlowTok{while} \OperatorTok{(}\NormalTok{r }\OperatorTok{!=} \DecValTok{37}\OperatorTok{);}    \CommentTok{// Ulangi sampai 37 muncul}
\OperatorTok{\}}
\end{Highlighting}
\end{Shaded}

\end{tcolorbox}

Catatan sampingan: pernahkah kamu menjalankannya lebih dari sekali? Kalau ya, apakah kamu melihat urutan angka yang sama muncul lagi. Dan lagi. Dan lagi? Itu karena \hl{\mbox{\texttt{rand()}}} adalah pseudo-random number generator yang harus di-\emph{seed} dengan angka berbeda supaya menghasilkan urutan yang berbeda. Lihat fungsi \href{https://beej.us/guide/bgclr/html/split/stdlib.html\#man-srand}{\hl{\mbox{\texttt{srand()}}}}\footnote{https://beej.us/guide/bgclr/html/split/stdlib.html\#man-srand} untuk detail lebih lanjut. \index{do-while pernyataan@\texttt{do}-\texttt{while} pernyataan|)}

\subsection{\texorpdfstring{Pernyataan \hl{\mbox{\texttt{for}}}}{Pernyataan for}}\label{forstat}

\index{for pernyataan@\texttt{for} pernyataan|(} Selamat datang di salah satu loop paling populer di dunia: loop \hl{\mbox{\texttt{for}}}!

Loop ini bagus kalau kamu sudah tahu jumlah iterasinya di awal.

Hal yang sama bisa dilakukan hanya dengan loop \hl{\mbox{\texttt{while}}}, tapi loop \hl{\mbox{\texttt{for}}} membuat kode lebih rapi.

Berikut dua potongan kode yang setara---perhatikan bahwa loop \hl{\mbox{\texttt{for}}} hanyalah bentuk yang lebih ringkas:

\begin{tcolorbox}[enhanced jigsaw,left=3pt,breakable,colback=white!97.2549!black,boxrule=0.2mm]

\begin{Shaded}
\begin{Highlighting}[]
\CommentTok{// Cetak angka antara 0 dan 9, inklusif...}

\CommentTok{// Memakai pernyataan while:}

\NormalTok{i }\OperatorTok{=} \DecValTok{0}\OperatorTok{;}
\ControlFlowTok{while} \OperatorTok{(}\NormalTok{i }\OperatorTok{\textless{}} \DecValTok{10}\OperatorTok{)} \OperatorTok{\{}
\NormalTok{    printf}\OperatorTok{(}\StringTok{"i is }\SpecialCharTok{\%d\textbackslash{}n}\StringTok{"}\OperatorTok{,}\NormalTok{ i}\OperatorTok{);}
\NormalTok{    i}\OperatorTok{++;}
\OperatorTok{\}}

\CommentTok{// Lakukan hal yang persis sama dengan for{-}loop:}

\ControlFlowTok{for} \OperatorTok{(}\NormalTok{i }\OperatorTok{=} \DecValTok{0}\OperatorTok{;}\NormalTok{ i }\OperatorTok{\textless{}} \DecValTok{10}\OperatorTok{;}\NormalTok{ i}\OperatorTok{++)} \OperatorTok{\{}
\NormalTok{    printf}\OperatorTok{(}\StringTok{"i is }\SpecialCharTok{\%d\textbackslash{}n}\StringTok{"}\OperatorTok{,}\NormalTok{ i}\OperatorTok{);}
\OperatorTok{\}}
\end{Highlighting}
\end{Shaded}

\end{tcolorbox}

Betul, kawan---keduanya melakukan hal yang persis sama. Tapi kamu bisa lihat sendiri bahwa pernyataan \hl{\mbox{\texttt{for}}} sedikit lebih ringkas dan enak dipandang. (Pengguna JavaScript pasti paham betul asal-usul C-nya di titik ini.)

Pernyataan \hl{\mbox{\texttt{for}}} terbagi menjadi tiga bagian yang dipisahkan titik koma. Bagian pertama berisi inisialisasi, bagian kedua kondisi loop, dan bagian ketiga apa yang dijalankan di akhir blok kalau kondisi loop bernilai true. Ketiga bagian ini opsional.

\begin{tcolorbox}[enhanced jigsaw,left=3pt,breakable,colback=white!97.2549!black,boxrule=0.2mm]

\begin{Shaded}
\begin{Highlighting}[]
\ControlFlowTok{for} \OperatorTok{(}\NormalTok{initialize things}\OperatorTok{;}\NormalTok{ loop }\ControlFlowTok{if}\NormalTok{ this is }\KeywordTok{true}\OperatorTok{;} \ControlFlowTok{do}\NormalTok{ this after each loop}\OperatorTok{)}
\end{Highlighting}
\end{Shaded}

\end{tcolorbox}

Perhatikan bahwa loop tidak akan dieksekusi sama sekali kalau kondisi loop sejak awal bernilai false.

\begin{tcolorbox}[blanker,colback=white,enhanced jigsaw,breakable,arc=0pt,outer arc=0pt,grow to left by=-1cm,enlarge top by=0.1cm,enlarge bottom by=0.1cm,borderline west={1.5pt}{-.3cm}{black}]

\begin{quote}
\textbf{Fakta menarik tentang loop \hl{\mbox{\texttt{for}}}!}

Kamu bisa memakai operator koma untuk melakukan beberapa hal di setiap klausa loop \hl{\mbox{\texttt{for}}}!

\begin{tcolorbox}[enhanced jigsaw,left=3pt,breakable,colback=white!97.2549!black,boxrule=0.2mm]

\begin{Shaded}
\begin{Highlighting}[]
\ControlFlowTok{for} \OperatorTok{(}\NormalTok{i }\OperatorTok{=} \DecValTok{0}\OperatorTok{,}\NormalTok{ j }\OperatorTok{=} \DecValTok{999}\OperatorTok{;}\NormalTok{ i }\OperatorTok{\textless{}} \DecValTok{10}\OperatorTok{;}\NormalTok{ i}\OperatorTok{++,}\NormalTok{ j}\OperatorTok{{-}{-})} \OperatorTok{\{}
\NormalTok{    printf}\OperatorTok{(}\StringTok{"}\SpecialCharTok{\%d}\StringTok{, }\SpecialCharTok{\%d\textbackslash{}n}\StringTok{"}\OperatorTok{,}\NormalTok{ i}\OperatorTok{,}\NormalTok{ j}\OperatorTok{);}
\OperatorTok{\}}
\end{Highlighting}
\end{Shaded}

\end{tcolorbox}
\end{quote}

\end{tcolorbox}

\hl{\mbox{\texttt{for}}} yang kosong berjalan selamanya:

\begin{tcolorbox}[enhanced jigsaw,left=3pt,breakable,colback=white!97.2549!black,boxrule=0.2mm]

\begin{Shaded}
\begin{Highlighting}[]
\ControlFlowTok{for}\OperatorTok{(;;)} \OperatorTok{\{}  \CommentTok{// "selamanya"}
\NormalTok{    printf}\OperatorTok{(}\StringTok{"I will print this again and again and again}\SpecialCharTok{\textbackslash{}n}\StringTok{"} \OperatorTok{);}
\NormalTok{    printf}\OperatorTok{(}\StringTok{"for all eternity until the heat{-}death of the universe.}\SpecialCharTok{\textbackslash{}n}\StringTok{"}\OperatorTok{);}

\NormalTok{    printf}\OperatorTok{(}\StringTok{"Or until you hit CTRL{-}C.}\SpecialCharTok{\textbackslash{}n}\StringTok{"}\OperatorTok{);}
\OperatorTok{\}}
\end{Highlighting}
\end{Shaded}

\end{tcolorbox}

\index{for pernyataan@\texttt{for} pernyataan|)}

\subsection{\texorpdfstring{Pernyataan \hl{\mbox{\texttt{switch}}}}{Pernyataan switch}}\label{switch-statement}

\index{switch pernyataan@\texttt{switch} pernyataan|(} Tergantung bahasa asalmu, kamu mungkin sudah atau belum mengenal \hl{\mbox{\texttt{switch}}}, atau versi di C ini mungkin lebih terbatas daripada yang biasa kamu pakai. Pernyataan ini memungkinkan kita mengambil tindakan yang berbeda tergantung nilai sebuah ekspresi integer.

Pada dasarnya, ekspresi dievaluasi menjadi nilai integer, lalu eksekusi melompat ke \index{case pernyataan@\texttt{case} pernyataan|(}\hl{\mbox{\texttt{case}}} yang cocok dengan nilai itu, dan berlanjut dari titik tersebut. Kalau pernyataan \hl{\mbox{\texttt{break}}}\index{break pernyataan@\texttt{break} pernyataan|(} ditemukan, eksekusi melompat keluar dari \hl{\mbox{\texttt{switch}}}.

Berikut contohnya: untuk jumlah kambing tertentu, kita cetak gambaran seberapa banyak jumlah itu.

\begin{tcolorbox}[enhanced jigsaw,left=3pt,breakable,colback=white!97.2549!black,boxrule=0.2mm]

\begin{Shaded}
\begin{Highlighting}[numbers=left,,]
\PreprocessorTok{\#include }\ImportTok{\textless{}stdio.h\textgreater{}}

\DataTypeTok{int}\NormalTok{ main}\OperatorTok{(}\DataTypeTok{void}\OperatorTok{)}
\OperatorTok{\{}
    \DataTypeTok{int}\NormalTok{ goat\_count }\OperatorTok{=} \DecValTok{2}\OperatorTok{;}

    \ControlFlowTok{switch} \OperatorTok{(}\NormalTok{goat\_count}\OperatorTok{)} \OperatorTok{\{}
        \ControlFlowTok{case} \DecValTok{0}\OperatorTok{:}
\NormalTok{            printf}\OperatorTok{(}\StringTok{"You have no goats.}\SpecialCharTok{\textbackslash{}n}\StringTok{"}\OperatorTok{);}
            \ControlFlowTok{break}\OperatorTok{;}

        \ControlFlowTok{case} \DecValTok{1}\OperatorTok{:}
\NormalTok{            printf}\OperatorTok{(}\StringTok{"You have a singular goat.}\SpecialCharTok{\textbackslash{}n}\StringTok{"}\OperatorTok{);}
            \ControlFlowTok{break}\OperatorTok{;}

        \ControlFlowTok{case} \DecValTok{2}\OperatorTok{:}
\NormalTok{            printf}\OperatorTok{(}\StringTok{"You have a brace of goats.}\SpecialCharTok{\textbackslash{}n}\StringTok{"}\OperatorTok{);}
            \ControlFlowTok{break}\OperatorTok{;}

        \ControlFlowTok{default}\OperatorTok{:}
\NormalTok{            printf}\OperatorTok{(}\StringTok{"You have a bona fide plethora of goats!}\SpecialCharTok{\textbackslash{}n}\StringTok{"}\OperatorTok{);}
            \ControlFlowTok{break}\OperatorTok{;}
    \OperatorTok{\}}
\OperatorTok{\}}
\end{Highlighting}
\end{Shaded}

\end{tcolorbox}

Pada contoh itu, \hl{\mbox{\texttt{switch}}} melompat ke \hl{\mbox{\texttt{case\ 2}}} dan mengeksekusi dari sana. Ketika (kalau) menemukan \hl{\mbox{\texttt{break}}}, ia melompat keluar dari \hl{\mbox{\texttt{switch}}}. \index{break pernyataan@\texttt{break} pernyataan|)}

Kamu juga mungkin melihat label \hl{\mbox{\texttt{default}}}\index{default label@\texttt{default} label} di bagian bawah. Inilah yang dijalankan kalau tidak ada case yang cocok.

Setiap \hl{\mbox{\texttt{case}}}, termasuk \hl{\mbox{\texttt{default}}}, bersifat opsional. Urutannya bebas, tapi sangat umum \hl{\mbox{\texttt{default}}}, kalau ada, ditaruh paling akhir.\index{case pernyataan@\texttt{case} pernyataan|)}

Jadi, secara keseluruhan ia bekerja seperti rangkaian \hl{\mbox{\texttt{if}}}-\hl{\mbox{\texttt{else}}}:

\begin{tcolorbox}[enhanced jigsaw,left=3pt,breakable,colback=white!97.2549!black,boxrule=0.2mm]

\begin{Shaded}
\begin{Highlighting}[]
\ControlFlowTok{if} \OperatorTok{(}\NormalTok{goat\_count }\OperatorTok{==} \DecValTok{0}\OperatorTok{)}
\NormalTok{    printf}\OperatorTok{(}\StringTok{"You have no goats.}\SpecialCharTok{\textbackslash{}n}\StringTok{"}\OperatorTok{);}
\ControlFlowTok{else} \ControlFlowTok{if} \OperatorTok{(}\NormalTok{goat\_count }\OperatorTok{==} \DecValTok{1}\OperatorTok{)}
\NormalTok{    printf}\OperatorTok{(}\StringTok{"You have a singular goat.}\SpecialCharTok{\textbackslash{}n}\StringTok{"}\OperatorTok{);}
\ControlFlowTok{else} \ControlFlowTok{if} \OperatorTok{(}\NormalTok{goat\_count }\OperatorTok{==} \DecValTok{2}\OperatorTok{)}
\NormalTok{    printf}\OperatorTok{(}\StringTok{"You have a brace of goats.}\SpecialCharTok{\textbackslash{}n}\StringTok{"}\OperatorTok{);}
\ControlFlowTok{else}
\NormalTok{    printf}\OperatorTok{(}\StringTok{"You have a bona fide plethora of goats!}\SpecialCharTok{\textbackslash{}n}\StringTok{"}\OperatorTok{);}
\end{Highlighting}
\end{Shaded}

\end{tcolorbox}

Dengan beberapa perbedaan utama:

\begin{itemize}
\tightlist
\item
  \hl{\mbox{\texttt{switch}}} sering kali lebih cepat melompat ke kode yang benar (meski spesifikasi tidak menjamin hal itu).
\item
  \hl{\mbox{\texttt{if}}}-\hl{\mbox{\texttt{else}}} bisa memakai kondisional relasional seperti \hl{\mbox{\texttt{\textless{}}}} dan \hl{\mbox{\texttt{\textgreater{}=}}}, serta floating point dan tipe lainnya, sedangkan \hl{\mbox{\texttt{switch}}} tidak bisa.
\end{itemize}

Ada satu hal lagi tentang \hl{\mbox{\texttt{switch}}} yang kadang kamu temui dan cukup menarik: \emph{fall through}.

\index{break pernyataan@\texttt{break} pernyataan|(} Ingat bagaimana \hl{\mbox{\texttt{break}}} membuat kita melompat keluar dari switch?

\index{Fall through|(} Nah, apa yang terjadi kalau kita \emph{tidak} \hl{\mbox{\texttt{break}}}?

Ternyata eksekusi terus masuk ke \hl{\mbox{\texttt{case}}} berikutnya! Demo!

\begin{tcolorbox}[enhanced jigsaw,left=3pt,breakable,colback=white!97.2549!black,boxrule=0.2mm]

\begin{Shaded}
\begin{Highlighting}[]
\ControlFlowTok{switch} \OperatorTok{(}\NormalTok{x}\OperatorTok{)} \OperatorTok{\{}
    \ControlFlowTok{case} \DecValTok{1}\OperatorTok{:}
\NormalTok{        printf}\OperatorTok{(}\StringTok{"1}\SpecialCharTok{\textbackslash{}n}\StringTok{"}\OperatorTok{);}
        \CommentTok{// Fall through!}
    \ControlFlowTok{case} \DecValTok{2}\OperatorTok{:}
\NormalTok{        printf}\OperatorTok{(}\StringTok{"2}\SpecialCharTok{\textbackslash{}n}\StringTok{"}\OperatorTok{);}
        \ControlFlowTok{break}\OperatorTok{;}
    \ControlFlowTok{case} \DecValTok{3}\OperatorTok{:}
\NormalTok{        printf}\OperatorTok{(}\StringTok{"3}\SpecialCharTok{\textbackslash{}n}\StringTok{"}\OperatorTok{);}
        \ControlFlowTok{break}\OperatorTok{;}
\OperatorTok{\}}
\end{Highlighting}
\end{Shaded}

\end{tcolorbox}

Kalau \hl{\mbox{\texttt{x\ ==\ 1}}}, \hl{\mbox{\texttt{switch}}} ini mula-mula mengenai \hl{\mbox{\texttt{case\ 1}}}, mencetak \hl{\mbox{\texttt{1}}}, lalu langsung lanjut ke baris kode berikutnya\ldots{} yang mencetak \hl{\mbox{\texttt{2}}}!

Setelah itu, akhirnya kita bertemu \hl{\mbox{\texttt{break}}} dan melompat keluar dari \hl{\mbox{\texttt{switch}}}.

Kalau \hl{\mbox{\texttt{x\ ==\ 2}}}, kita hanya mengenai \hl{\mbox{\texttt{case\ 2}}}, mencetak \hl{\mbox{\texttt{2}}}, dan \hl{\mbox{\texttt{break}}} seperti biasa.

Ketiadaan \hl{\mbox{\texttt{break}}} inilah yang disebut \emph{fall through}.

Kalau kamu memakai C sebelum C23, \emph{SELALU} taruh komentar di kode pada tempat yang memang kamu maksudkan untuk fall through, seperti di atas. Dengan begitu programmer lain tidak bertanya-tanya apakah kamu sengaja melakukannya. GCC mengenali komentar ini dan menekan peringatan.

Atau kamu bisa memakai atribut GCC non-standar, yang kurang disukai karena jelek dan non-standar, tapi dikenali oleh GCC maupun Clang:

\begin{tcolorbox}[enhanced jigsaw,left=3pt,breakable,colback=white!97.2549!black,boxrule=0.2mm]

\begin{Shaded}
\begin{Highlighting}[]
\ControlFlowTok{switch} \OperatorTok{(}\NormalTok{x}\OperatorTok{)} \OperatorTok{\{}
    \ControlFlowTok{case} \DecValTok{1}\OperatorTok{:}
\NormalTok{        printf}\OperatorTok{(}\StringTok{"1}\SpecialCharTok{\textbackslash{}n}\StringTok{"}\OperatorTok{);}
\NormalTok{        \_\_attribute\_\_}\OperatorTok{((}\NormalTok{fallthrough}\OperatorTok{));}
    \ControlFlowTok{case} \DecValTok{2}\OperatorTok{:}
        \OperatorTok{...}
\end{Highlighting}
\end{Shaded}

\end{tcolorbox}

Tapi sekarang, akhirnya, di C23 kita punya cara standar lewat atribut:

\begin{tcolorbox}[enhanced jigsaw,left=3pt,breakable,colback=white!97.2549!black,boxrule=0.2mm]

\begin{Shaded}
\begin{Highlighting}[]
\ControlFlowTok{switch} \OperatorTok{(}\NormalTok{x}\OperatorTok{)} \OperatorTok{\{}
    \ControlFlowTok{case} \DecValTok{1}\OperatorTok{:}
\NormalTok{        printf}\OperatorTok{(}\StringTok{"1}\SpecialCharTok{\textbackslash{}n}\StringTok{"}\OperatorTok{);}
        \OperatorTok{[[}\AttributeTok{fallthrough}\OperatorTok{]];}
    \ControlFlowTok{case} \DecValTok{2}\OperatorTok{:}
        \OperatorTok{...}
\end{Highlighting}
\end{Shaded}

\end{tcolorbox}

Tapi ingat, ini salah satu tempat paling umum munculnya bug dalam program C: lupa menaruh \hl{\mbox{\texttt{break}}} di \hl{\mbox{\texttt{case}}}. Kamu harus menaruhnya kalau tidak ingin langsung jatuh ke case berikutnya\footnote{Hal ini dianggap sebagai bahaya yang begitu besar, sampai-sampai para perancang bahasa pemrograman Go menjadikan \hl{\mbox{\texttt{break}}} sebagai perilaku bawaan; kamu harus secara eksplisit memakai pernyataan \hl{\mbox{\texttt{fallthrough}}} di Go kalau ingin jatuh ke case berikutnya.}.

\index{Fall through|)} \index{break pernyataan@\texttt{break} pernyataan|)}

Tadi sudah disebut bahwa \hl{\mbox{\texttt{switch}}} bekerja dengan tipe integer---patuhi itu. Jangan pakai floating point atau tipe string di sana. Satu celahnya: tipe karakter boleh dipakai karena diam-diam mereka juga integer. Jadi ini sepenuhnya sah:

\begin{tcolorbox}[enhanced jigsaw,left=3pt,breakable,colback=white!97.2549!black,boxrule=0.2mm]

\begin{Shaded}
\begin{Highlighting}[]
\DataTypeTok{char}\NormalTok{ c }\OperatorTok{=} \CharTok{\textquotesingle{}b\textquotesingle{}}\OperatorTok{;}

\ControlFlowTok{switch} \OperatorTok{(}\NormalTok{c}\OperatorTok{)} \OperatorTok{\{}
    \ControlFlowTok{case} \CharTok{\textquotesingle{}a\textquotesingle{}}\OperatorTok{:}
\NormalTok{        printf}\OperatorTok{(}\StringTok{"It\textquotesingle{}s \textquotesingle{}a\textquotesingle{}!}\SpecialCharTok{\textbackslash{}n}\StringTok{"}\OperatorTok{);}
        \ControlFlowTok{break}\OperatorTok{;}

    \ControlFlowTok{case} \CharTok{\textquotesingle{}b\textquotesingle{}}\OperatorTok{:}
\NormalTok{        printf}\OperatorTok{(}\StringTok{"It\textquotesingle{}s \textquotesingle{}b\textquotesingle{}!}\SpecialCharTok{\textbackslash{}n}\StringTok{"}\OperatorTok{);}
        \ControlFlowTok{break}\OperatorTok{;}

    \ControlFlowTok{case} \CharTok{\textquotesingle{}c\textquotesingle{}}\OperatorTok{:}
\NormalTok{        printf}\OperatorTok{(}\StringTok{"It\textquotesingle{}s \textquotesingle{}c\textquotesingle{}!}\SpecialCharTok{\textbackslash{}n}\StringTok{"}\OperatorTok{);}
        \ControlFlowTok{break}\OperatorTok{;}
\OperatorTok{\}}
\end{Highlighting}
\end{Shaded}

\end{tcolorbox}

Terakhir, kamu bisa memakai \index{enum kata kunci@\texttt{enum} kata kunci}\hl{\mbox{\texttt{enum}}} dalam \hl{\mbox{\texttt{switch}}} karena \hl{\mbox{\texttt{enum}}} juga tipe integer. Soal itu dibahas lebih lanjut di bab \hl{\mbox{\texttt{enum}}}.

\index{switch pernyataan@\texttt{switch} pernyataan|)} \index{break pernyataan@\texttt{break} pernyataan|)} \index{Kendali Alur|)}

\chapter{Fungsi}\label{functions}

\begin{tcolorbox}[blanker,colback=white,enhanced jigsaw,breakable,arc=0pt,outer arc=0pt,grow to left by=-1cm,enlarge top by=0.1cm,enlarge bottom by=0.1cm,borderline west={1.5pt}{-.3cm}{black}]

\begin{quote}
\emph{``Sir, not in an environment such as this. That's why I've also been programmed for over thirty secondary functions that---''}

---C3PO\index{C3PO}, sebelum diinterupsi dengan kasar, melaporkan jumlah fungsi tambahan yang sekarang sudah tidak mengesankan, \emph{Star Wars} script
\end{quote}

\end{tcolorbox}

\index{Fungsi|(} Seperti bahasa lain yang sudah kamu pakai, C punya konsep \emph{fungsi}.

Fungsi bisa menerima berbagai \emph{argumen}\index{argumen Fungsi} dan mengembalikan sebuah nilai. Satu hal penting: tipe argumen dan tipe nilai kembalian harus dideklarasikan lebih dulu---begitulah cara C menyukainya!

Mari kita lihat sebuah fungsi. Fungsi ini menerima sebuah \hl{\mbox{\texttt{int}}} sebagai argumen, dan mengembalikan\index{return pernyataan@\texttt{return} pernyataan} sebuah \hl{\mbox{\texttt{int}}}.

\begin{tcolorbox}[enhanced jigsaw,left=3pt,breakable,colback=white!97.2549!black,boxrule=0.2mm]

\begin{Shaded}
\begin{Highlighting}[numbers=left,,]
\PreprocessorTok{\#include }\ImportTok{\textless{}stdio.h\textgreater{}}

\DataTypeTok{int}\NormalTok{ plus\_one}\OperatorTok{(}\DataTypeTok{int}\NormalTok{ n}\OperatorTok{)}  \CommentTok{// "Definisi"{-}nya}
\OperatorTok{\{}
    \ControlFlowTok{return}\NormalTok{ n }\OperatorTok{+} \DecValTok{1}\OperatorTok{;}
\OperatorTok{\}}
 
\end{Highlighting}
\end{Shaded}

\end{tcolorbox}

\hl{\mbox{\texttt{int}}} sebelum \hl{\mbox{\texttt{plus\_one}}} menunjukkan tipe kembalian.

\hl{\mbox{\texttt{int\ n}}} menunjukkan bahwa fungsi ini menerima satu argumen \hl{\mbox{\texttt{int}}}, yang disimpan dalam \emph{parameter} \hl{\mbox{\texttt{n}}}\index{parameter Fungsi}. Parameter adalah variabel lokal khusus yang menjadi tempat argumen disalin.

Poin ini perlu ditegaskan: argumen disalin ke dalam parameter. Banyak hal di C jadi lebih mudah dipahami kalau kamu ingat bahwa parameter adalah \emph{salinan} argumen, bukan argumen itu sendiri. Soal ini sebentar lagi kita bahas.

Lanjut ke \hl{\mbox{\texttt{main()}}}, di sana ada pemanggilan fungsi, dan nilai kembaliannya kita tugaskan ke variabel lokal \hl{\mbox{\texttt{j}}}:

\begin{tcolorbox}[enhanced jigsaw,left=3pt,breakable,colback=white!97.2549!black,boxrule=0.2mm]

\begin{Shaded}
\begin{Highlighting}[numbers=left,,firstnumber=8,]
\DataTypeTok{int}\NormalTok{ main}\OperatorTok{(}\DataTypeTok{void}\OperatorTok{)}
\OperatorTok{\{}
    \DataTypeTok{int}\NormalTok{ i }\OperatorTok{=} \DecValTok{10}\OperatorTok{,}\NormalTok{ j}\OperatorTok{;}
    
\NormalTok{    j }\OperatorTok{=}\NormalTok{ plus\_one}\OperatorTok{(}\NormalTok{i}\OperatorTok{);}  \CommentTok{// "Pemanggilan"{-}nya}

\NormalTok{    printf}\OperatorTok{(}\StringTok{"i + 1 is }\SpecialCharTok{\%d\textbackslash{}n}\StringTok{"}\OperatorTok{,}\NormalTok{ j}\OperatorTok{);}
\OperatorTok{\}}
\end{Highlighting}
\end{Shaded}

\end{tcolorbox}

\begin{tcolorbox}[blanker,colback=white,enhanced jigsaw,breakable,arc=0pt,outer arc=0pt,grow to left by=-1cm,enlarge top by=0.1cm,enlarge bottom by=0.1cm,borderline west={1.5pt}{-.3cm}{black}]

\begin{quote}
Sebelum lupa: perhatikan bahwa fungsi didefinisikan sebelum dipakai. Kalau tidak, compiler belum mengenalnya ketika mengompilasi \hl{\mbox{\texttt{main()}}} dan akan memberi error pemanggilan fungsi yang tidak dikenal. Ada cara yang lebih tepat untuk menulis kode di atas, yaitu dengan \emph{function prototype}, tapi itu dibahas nanti.
\end{quote}

\end{tcolorbox}

Perhatikan juga bahwa \hl{\mbox{\texttt{main()}}}\index{main() fungsi@\texttt{main()} fungsi} adalah sebuah fungsi!

Fungsi ini mengembalikan \hl{\mbox{\texttt{int}}}.

Lalu, apa itu \hl{\mbox{\texttt{void}}}\index{void tipe@\texttt{void} tipe}? Kata kunci ini menandakan bahwa fungsi tidak menerima argumen.

Kamu juga bisa mengembalikan \hl{\mbox{\texttt{void}}} untuk menandakan bahwa fungsi tidak mengembalikan nilai:

\begin{tcolorbox}[enhanced jigsaw,left=3pt,breakable,colback=white!97.2549!black,boxrule=0.2mm]

\begin{Shaded}
\begin{Highlighting}[numbers=left,,]
\PreprocessorTok{\#include }\ImportTok{\textless{}stdio.h\textgreater{}}

\CommentTok{// Fungsi ini tidak menerima argumen dan tidak mengembalikan nilai:}

\DataTypeTok{void}\NormalTok{ hello}\OperatorTok{(}\DataTypeTok{void}\OperatorTok{)}
\OperatorTok{\{}
\NormalTok{    printf}\OperatorTok{(}\StringTok{"Hello, world!}\SpecialCharTok{\textbackslash{}n}\StringTok{"}\OperatorTok{);}
\OperatorTok{\}}

\DataTypeTok{int}\NormalTok{ main}\OperatorTok{(}\DataTypeTok{void}\OperatorTok{)}
\OperatorTok{\{}
\NormalTok{    hello}\OperatorTok{();}  \CommentTok{// Mencetak "Hello, world!"}
\OperatorTok{\}}
\end{Highlighting}
\end{Shaded}

\end{tcolorbox}

\section{Passing by Value}\label{passvalue}

\index{Pass by value|}Tadi sudah disebutkan bahwa saat sebuah argumen dioper ke fungsi, salinan argumen itu dibuat dan disimpan dalam parameter yang sesuai.

Kalau argumennya berupa variabel, yang disalin ke parameter adalah nilai variabel tersebut.

Lebih umum lagi, seluruh ekspresi argumen dievaluasi dulu untuk mendapatkan nilainya. Nilai itulah yang disalin ke parameter.

Apa pun yang terjadi, nilai dalam parameter itu berdiri sendiri. Ia tidak bergantung pada nilai atau variabel apa pun yang kamu pakai sebagai argumen saat memanggil fungsi.

Mari kita lihat contohnya. Pelajari, dan coba tebak outputnya sebelum dijalankan:

\begin{tcolorbox}[enhanced jigsaw,left=3pt,breakable,colback=white!97.2549!black,boxrule=0.2mm]

\begin{Shaded}
\begin{Highlighting}[numbers=left,,]
\PreprocessorTok{\#include }\ImportTok{\textless{}stdio.h\textgreater{}}

\DataTypeTok{void}\NormalTok{ increment}\OperatorTok{(}\DataTypeTok{int}\NormalTok{ a}\OperatorTok{)}
\OperatorTok{\{}
\NormalTok{    a}\OperatorTok{++;}
\OperatorTok{\}}

\DataTypeTok{int}\NormalTok{ main}\OperatorTok{(}\DataTypeTok{void}\OperatorTok{)}
\OperatorTok{\{}
    \DataTypeTok{int}\NormalTok{ i }\OperatorTok{=} \DecValTok{10}\OperatorTok{;}

\NormalTok{    increment}\OperatorTok{(}\NormalTok{i}\OperatorTok{);}

\NormalTok{    printf}\OperatorTok{(}\StringTok{"i == }\SpecialCharTok{\%d\textbackslash{}n}\StringTok{"}\OperatorTok{,}\NormalTok{ i}\OperatorTok{);}  \CommentTok{// Apa yang dicetak ini?}
\OperatorTok{\}}
\end{Highlighting}
\end{Shaded}

\end{tcolorbox}

Sekilas, \hl{\mbox{\texttt{i}}} bernilai \hl{\mbox{\texttt{10}}}, lalu kita mengopernya ke fungsi \hl{\mbox{\texttt{increment()}}}. Di sana nilainya di-increment, jadi waktu dicetak pasti \hl{\mbox{\texttt{11}}}, kan?

\begin{tcolorbox}[blanker,colback=white,enhanced jigsaw,breakable,arc=0pt,outer arc=0pt,grow to left by=-1cm,enlarge top by=0.1cm,enlarge bottom by=0.1cm,borderline west={1.5pt}{-.3cm}{black}]

\begin{quote}
\emph{``Get used to disappointment.''}

---Dread Pirate Roberts, \emph{The Princess Bride}
\end{quote}

\end{tcolorbox}

Tapi hasilnya bukan \hl{\mbox{\texttt{11}}}---yang tercetak \hl{\mbox{\texttt{10}}}! Kenapa?

Semuanya berpangkal pada kenyataan bahwa ekspresi yang kamu oper ke fungsi \emph{disalin} ke parameter yang sesuai. Parameter itu salinan, bukan yang aslinya.

Jadi \hl{\mbox{\texttt{i}}} bernilai \hl{\mbox{\texttt{10}}} di dalam \hl{\mbox{\texttt{main()}}}. Lalu kita mengopernya ke \hl{\mbox{\texttt{increment()}}}. Parameter yang sesuai di dalam fungsi itu bernama \hl{\mbox{\texttt{a}}}.

Penyalinan pun terjadi, seolah-olah lewat penugasan. Kurang lebih, \hl{\mbox{\texttt{a\ =\ i}}}. Pada titik itu \hl{\mbox{\texttt{a}}} bernilai \hl{\mbox{\texttt{10}}}, dan di dalam \hl{\mbox{\texttt{main()}}}, \hl{\mbox{\texttt{i}}} juga \hl{\mbox{\texttt{10}}}.

Kemudian \hl{\mbox{\texttt{a}}} di-increment menjadi \hl{\mbox{\texttt{11}}}. Tapi \hl{\mbox{\texttt{i}}} sama sekali tidak tersentuh! Nilainya tetap \hl{\mbox{\texttt{10}}}.

Akhirnya fungsi selesai. Semua variabel lokalnya dibuang (selamat tinggal, \hl{\mbox{\texttt{a}}}!) dan kita kembali ke \hl{\mbox{\texttt{main()}}}, tempat \hl{\mbox{\texttt{i}}} masih \hl{\mbox{\texttt{10}}}.

Lalu \hl{\mbox{\texttt{i}}} dicetak, hasilnya \hl{\mbox{\texttt{10}}}, dan selesai.

Itulah sebabnya pada contoh \hl{\mbox{\texttt{plus\_one()}}} tadi, kita \hl{\mbox{\texttt{return}}} nilai yang sudah diubah secara lokal supaya bisa dilihat lagi di \hl{\mbox{\texttt{main()}}}.

Kelihatannya cukup membatasi, ya? Sepertinya dari sebuah fungsi kamu hanya bisa mendapat satu data kembali. Setidaknya itu yang ada di pikiranmu. Tapi ada cara lain untuk mendapat data kembali; orang C menyebutnya \emph{passing by reference}, dan itu cerita untuk lain waktu.

Namun, sehebat apa pun namanya, itu tidak mengubah kenyataan bahwa \emph{SEMUA} yang kamu oper ke fungsi \emph{TANPA PENGECUALIAN} disalin ke parameter yang sesuai, dan fungsi bekerja pada salinan lokal itu, \emph{APA PUN YANG TERJADI}. Ingat itu, bahkan ketika nanti kita membahas passing by reference yang konon hebat itu. \index{Pass by value|)}

\section{Function Prototype}\label{prototypes}

\index{Function prototype|(}Kalau kamu ingat, beberapa bagian lalu (rasanya sudah zaman es) disebutkan bahwa fungsi harus didefinisikan sebelum dipakai, kalau tidak compiler belum mengenalnya dan akan meledak dengan error.

Itu tidak sepenuhnya benar. Kamu bisa memberi tahu compiler lebih dulu bahwa kamu akan memakai fungsi bertipe tertentu dengan daftar parameter tertentu. Dengan begitu, fungsinya bisa didefinisikan di mana saja (bahkan di file lain), asalkan \emph{function prototype} sudah dideklarasikan sebelum fungsinya dipanggil.

Untungnya, function prototype cukup mudah. Ia hanyalah salinan baris pertama definisi fungsi, ditambah titik koma di akhir. Misalnya, kode berikut memanggil fungsi yang didefinisikan belakangan, karena prototype-nya sudah dideklarasikan lebih dulu:

\begin{tcolorbox}[enhanced jigsaw,left=3pt,breakable,colback=white!97.2549!black,boxrule=0.2mm]

\begin{Shaded}
\begin{Highlighting}[numbers=left,,]
\PreprocessorTok{\#include }\ImportTok{\textless{}stdio.h\textgreater{}}

\DataTypeTok{int}\NormalTok{ foo}\OperatorTok{(}\DataTypeTok{void}\OperatorTok{);}  \CommentTok{// Ini prototype{-}nya!}

\DataTypeTok{int}\NormalTok{ main}\OperatorTok{(}\DataTypeTok{void}\OperatorTok{)}
\OperatorTok{\{}
    \DataTypeTok{int}\NormalTok{ i}\OperatorTok{;}
    
    \CommentTok{// Kita bisa memanggil foo() di sini sebelum definisinya karena}
    \CommentTok{// prototype{-}nya sudah dideklarasikan di atas!}

\NormalTok{    i }\OperatorTok{=}\NormalTok{ foo}\OperatorTok{();}
    
\NormalTok{    printf}\OperatorTok{(}\StringTok{"}\SpecialCharTok{\%d\textbackslash{}n}\StringTok{"}\OperatorTok{,}\NormalTok{ i}\OperatorTok{);}  \CommentTok{// 3490}
\OperatorTok{\}}

\DataTypeTok{int}\NormalTok{ foo}\OperatorTok{(}\DataTypeTok{void}\OperatorTok{)}  \CommentTok{// Ini definisinya, persis seperti prototype{-}nya!}
\OperatorTok{\{}
    \ControlFlowTok{return} \DecValTok{3490}\OperatorTok{;}
\OperatorTok{\}}
\end{Highlighting}
\end{Shaded}

\end{tcolorbox}

Kalau fungsi tidak dideklarasikan sebelum dipakai (baik lewat prototype maupun definisinya), yang terjadi disebut \emph{implicit declaration}\index{Implicit declaration}. Ini dulu diizinkan di standar C pertama (C89), yang punya aturan khusus untuknya, tapi sekarang tidak lagi diizinkan. Dan tidak ada alasan sah untuk mengandalkannya di kode baru.

Mungkin kamu menyadari sesuatu pada kode contoh yang sudah kita pakai\ldots{} Kita memakai fungsi \hl{\mbox{\texttt{printf()}}} yang setia itu tanpa mendefinisikannya atau mendeklarasikan prototype-nya! Bagaimana kita bisa lolos dari pelanggaran hukum ini? Sebenarnya tidak lolos. Prototype-nya ada, yaitu di header file \hl{\mbox{\texttt{stdio.h}}} yang kita \hl{\mbox{\texttt{\#include}}}, ingat? Jadi kita tetap sah, Pak Polisi!\index{Function prototype|)}

\section{Daftar Parameter Kosong}\label{daftar-parameter-kosong}

\index{Empty parameter lists} Di C23, daftar parameter kosong dalam definisi atau prototype fungsi artinya persis sama dengan menulis \hl{\mbox{\texttt{void}}}\index{void tipe@\texttt{void} tipe} di situ.

\begin{tcolorbox}[enhanced jigsaw,left=3pt,breakable,colback=white!97.2549!black,boxrule=0.2mm]

\begin{Shaded}
\begin{Highlighting}[]
\DataTypeTok{void}\NormalTok{ foo}\OperatorTok{();}
\DataTypeTok{void}\NormalTok{ foo}\OperatorTok{(}\DataTypeTok{void}\OperatorTok{);} \CommentTok{// Sama saja di C23}
\end{Highlighting}
\end{Shaded}

\end{tcolorbox}

Meski begitu, kode C zaman sekarang masih cenderung memilih \hl{\mbox{\texttt{void}}}\index{void tipe@\texttt{void} tipe} yang eksplisit di daftar parameter, terutama karena sebelum C23, perilaku daftar parameter kosong sangat berbeda. Dan bukan perbedaan yang menyenangkan.

Kita bahas sekarang, tapi singkatnya: kalau kodenya mungkin suatu saat dibangun dengan compiler pra-C23, \textbf{selalu} pakai \hl{\mbox{\texttt{void}}}\index{void tipe@\texttt{void} tipe} untuk menyatakan daftar parameter kosong di semua deklarasi fungsi.

Ada dua konteks yang perlu dipertimbangkan.

\begin{itemize}
\tightlist
\item
  Menghilangkan semua parameter di tempat fungsi didefinisikan
\item
  Menghilangkan semua parameter dalam sebuah prototype
\end{itemize}

Kita lihat dulu definisi fungsinya:

\begin{tcolorbox}[enhanced jigsaw,left=3pt,breakable,colback=white!97.2549!black,boxrule=0.2mm]

\begin{Shaded}
\begin{Highlighting}[]
\DataTypeTok{void}\NormalTok{ foo}\OperatorTok{()}  \CommentTok{// Seharusnya ada \textasciigrave{}void\textasciigrave{} di sini}
\OperatorTok{\{}
\NormalTok{    printf}\OperatorTok{(}\StringTok{"Hello, world!}\SpecialCharTok{\textbackslash{}n}\StringTok{"}\OperatorTok{);}
\OperatorTok{\}}
\end{Highlighting}
\end{Shaded}

\end{tcolorbox}

Spesifikasi memang menyebut bahwa ini berarti fungsi \hl{\mbox{\texttt{foo()}}} tidak menerima parameter, tapi compiler boleh saja tidak memperingatkanmu kalau kamu keliru mengoper argumen. Dan itu undefined behavior. Untuk menghindari semua itu, cukup tulis \hl{\mbox{\texttt{void}}}\index{void tipe@\texttt{void} tipe} di sana.

Pada function prototype, ada perbedaan yang \emph{signifikan} antara memakai \hl{\mbox{\texttt{void}}}\index{void tipe@\texttt{void} tipe!in function prototypes} dan tidak:

\begin{tcolorbox}[enhanced jigsaw,left=3pt,breakable,colback=white!97.2549!black,boxrule=0.2mm]

\begin{Shaded}
\begin{Highlighting}[]
\DataTypeTok{void}\NormalTok{ foo}\OperatorTok{();}
\DataTypeTok{void}\NormalTok{ foo}\OperatorTok{(}\DataTypeTok{void}\OperatorTok{);}  \CommentTok{// Tidak sama!}
\end{Highlighting}
\end{Shaded}

\end{tcolorbox}

Kalau \hl{\mbox{\texttt{void}}} dihilangkan dari prototype, compiler diberi tahu bahwa tidak ada informasi tambahan tentang parameter fungsi. Ini secara efektif mematikan semua pengecekan tipe parameter, sehingga \hl{\mbox{\texttt{void\ foo();}}} pada dasarnya bukan ``prototype'' sama sekali. Ia hanyalah forward declaration.

Jadi, pada prototype \textbf{pasti} pakai \hl{\mbox{\texttt{void}}} kalau daftar parameternya kosong.

Dan, ya, itu praktik yang baik dalam semua kasus. Mungkin suatu hari, ketika semua compiler pra-C23 sudah tidak dipakai lagi, kita bisa meninggalkannya.

\index{Fungsi|)}

\chapter{Pointer---Gemetar Ketakutan!}\label{pointers}

\begin{tcolorbox}[blanker,colback=white,enhanced jigsaw,breakable,arc=0pt,outer arc=0pt,grow to left by=-1cm,enlarge top by=0.1cm,enlarge bottom by=0.1cm,borderline west={1.5pt}{-.3cm}{black}]

\begin{quote}
\emph{``How do you get to Carnegie Hall?''}\\
\emph{``Practice!''}

---20th-century joke of unknown origin
\end{quote}

\end{tcolorbox}

\index{Pointers|(}Pointer adalah salah satu hal yang paling ditakuti di bahasa C. Bahkan, hanya pointer-lah yang membuat bahasa ini terasa menantang. Tapi kenapa?

Karena, jujur saja, pointer bisa mengirim sengatan listrik lewat keyboard dan secara fisik \emph{mengelas} lenganmu di tempat selamanya, mengutukmu menghabiskan hidup di depan keyboard dengan bahasa dari tahun 70-an ini!

Benar begitu? Tentu tidak. Itu cuma usaha untuk menyiapkan kamu supaya berhasil.

Tergantung bahasa asalmu, mungkin kamu sudah paham konsep \emph{referensi}, yaitu variabel yang merujuk ke objek bertipe tertentu.

Pointer sangat mirip dengan itu, hanya saja di C kita harus lebih eksplisit membedakan kapan kita membicarakan referensinya dan kapan hal yang dirujuknya.

\section{Memori dan Variabel}\label{ptmem}

Memori komputer menyimpan segala jenis data, kan? \hl{\mbox{\texttt{float}}}, \hl{\mbox{\texttt{int}}}, atau apa pun yang kamu punya. Agar mudah dikelola, setiap byte memori diberi identitas berupa sebuah integer. Integer-integer ini naik berurutan sepanjang memori\footnote{Biasanya. Pasti ada pengecualian di lorong-lorong gelap sejarah komputasi.}. Bayangkan sekumpulan kotak bernomor, dan setiap kotak menyimpan satu byte\footnote{Byte adalah angka yang terdiri dari tidak lebih dari 8 digit biner, atau singkatnya \emph{bit}. Dalam digit desimal seperti yang biasa dipakai nenek kita, byte bisa menyimpan bilangan tak bertanda antara 0 dan 255, inklusif.} data. Atau bayangkan array besar yang setiap elemennya menyimpan satu byte, kalau kamu berasal dari bahasa yang punya array. Angka yang menandai setiap kotak disebut \emph{alamat}\index{Memory address}-nya.

Nah, tidak semua tipe data hanya memakai satu byte. Misalnya, sebuah \hl{\mbox{\texttt{int}}} sering kali empat byte, begitu juga \hl{\mbox{\texttt{float}}}, tapi itu benar-benar tergantung sistem. Kamu bisa memakai operator \hl{\mbox{\texttt{sizeof}}} untuk mengetahui berapa byte memori yang dipakai sebuah tipe.

\begin{tcolorbox}[enhanced jigsaw,left=3pt,breakable,colback=white!97.2549!black,boxrule=0.2mm]

\begin{Shaded}
\begin{Highlighting}[]
\CommentTok{// \%zu adalah format specifier untuk tipe size\_t}

\NormalTok{printf}\OperatorTok{(}\StringTok{"an int uses }\SpecialCharTok{\%zu}\StringTok{ bytes of memory}\SpecialCharTok{\textbackslash{}n}\StringTok{"}\OperatorTok{,} \KeywordTok{sizeof}\OperatorTok{(}\DataTypeTok{int}\OperatorTok{));}

\CommentTok{// Di sistem penulis ini mencetak "4", tapi bisa berbeda tiap sistem.}
\end{Highlighting}
\end{Shaded}

\end{tcolorbox}

\begin{tcolorbox}[blanker,colback=white,enhanced jigsaw,breakable,arc=0pt,outer arc=0pt,grow to left by=-1cm,enlarge top by=0.1cm,enlarge bottom by=0.1cm,borderline west={1.5pt}{-.3cm}{black}]

\begin{quote}
\textbf{Fakta Menarik tentang Memori}: Kalau sebuah tipe data (seperti \hl{\mbox{\texttt{int}}} biasa) memakai lebih dari satu byte memori, byte-byte penyusunnya selalu berdekatan di memori. Kadang urutannya sesuai harapanmu, kadang tidak\footnote{Urutan byte disebut \emph{endianness} dari angka tersebut. Yang paling sering ditemui adalah \emph{big-endian} (byte paling signifikan di depan) dan \emph{little-endian} (byte paling signifikan di belakang), atau yang kini jarang, \emph{mixed-endian} (byte paling signifikan di tempat lain).}. Walau C tidak menjamin urutan memori tertentu (tergantung platform), pada umumnya kamu tetap bisa menulis kode yang tidak bergantung platform, tanpa harus memikirkan urutan byte yang menjengkelkan ini.
\end{quote}

\end{tcolorbox}

Jadi, \emph{anyway}, kalau sudah siap, putar drum roll dan musik menegangkan untuk definisi pointer: \emph{pointer adalah variabel yang menyimpan sebuah alamat}. Bayangkan musik klasik dari 2001: A Space Odyssey di sini. Ba bum ba bum ba bum BAAAAH!

Ok, mungkin agak berlebihan ya? Sebenarnya pointer tidak punya banyak misteri. Pointer adalah alamat dari data. Sama seperti variabel \hl{\mbox{\texttt{int}}} bisa menyimpan nilai \hl{\mbox{\texttt{12}}}, variabel pointer bisa menyimpan alamat data.

Artinya, semua ini bermakna sama, yaitu sebuah angka yang menyatakan suatu titik di memori:

\begin{itemize}
\tightlist
\item
  Indeks ke dalam memori (kalau kamu membayangkan memori sebagai array besar)
\item
  Alamat
\item
  Lokasi
\end{itemize}

Ketiga istilah ini dipakai bergantian. Dan ya, \emph{lokasi} baru saja ditambahkan karena kata yang bermakna sama itu tidak pernah cukup.

Variabel pointer menyimpan nomor alamat itu, sama seperti variabel \hl{\mbox{\texttt{float}}} bisa menyimpan \hl{\mbox{\texttt{3.14159}}}.

Bayangkan kamu punya setumpuk Post-it® yang bernomor berurutan sebagai alamatnya. (Yang pertama di indeks \hl{\mbox{\texttt{0}}}, berikutnya di indeks \hl{\mbox{\texttt{1}}}, dan seterusnya.)

Selain angka penanda posisinya, kamu juga bisa menulis angka lain di masing-masing Post-it. Misalnya jumlah anjingmu. Atau jumlah bulan di sekitar Mars\ldots{}

\ldots Atau, \emph{bisa juga indeks Post-it yang lain!}

Kalau yang kamu tulis jumlah anjingmu, itu variabel biasa. Tapi kalau yang kamu tulis indeks Post-it lain, \emph{itulah pointer}. Ia menunjuk ke catatan lain!

Analogi lainnya adalah alamat rumah. Sebuah rumah punya ciri-ciri tertentu: halaman, atap metal, panel surya, dan sebagainya. Lalu ada alamat rumah itu. Alamat bukanlah rumahnya. Yang satu adalah rumah lengkap, yang lain hanya beberapa baris teks. Tapi alamat rumah adalah \emph{pointer} ke rumah tersebut. Ia bukan rumahnya, tapi memberi tahu di mana rumahnya.

Hal yang sama bisa kita lakukan di komputer. Ada variabel data yang menyimpan suatu nilai, dan nilai itu berada di memori pada suatu alamat. Lalu ada \emph{variabel pointer} yang berbeda, yang menyimpan alamat variabel data tadi.

Variabel pointer bukan variabel data itu sendiri, tapi seperti alamat rumah, ia memberi tahu kita di mana menemukannya.

Kalau sudah begitu, kita bilang kita punya ``pointer ke'' data tersebut. Dan kita bisa mengikuti pointer itu untuk mengakses datanya.

(Sekarang mungkin belum terlihat berguna, tapi semua ini jadi sangat penting ketika dipakai bersama pemanggilan fungsi. Sabar dulu sampai kita ke sana.)

Jadi, kalau kita punya sebuah \hl{\mbox{\texttt{int}}} dan ingin pointer ke sana, kita butuh cara untuk mendapatkan alamat \hl{\mbox{\texttt{int}}} itu, kan? Sebab pointer hanya menyimpan \emph{alamat} data. Operator apa yang dipakai untuk mencari \emph{alamat} sebuah \hl{\mbox{\texttt{int}}}?

\index{\& operator address-of@\texttt{\&} operator address-of|(}Dengan kejutan yang pasti mengejutkan kamu, pembaca yang baik, kita memakai operator \hl{\mbox{\texttt{address-of}}} (yang ternyata adalah ampersand: ``\hl{\mbox{\texttt{\&}}}'') untuk mencari alamat data. Ampersand.

Sebagai contoh cepat, kita perkenalkan \emph{format specifier} baru untuk \hl{\mbox{\texttt{printf()}}} agar bisa mencetak pointer. Kamu sudah tahu \hl{\mbox{\texttt{\%d}}} mencetak integer desimal, kan? Nah, \hl{\mbox{\texttt{\%p}}}\index{printf() fungsi@\texttt{printf()} fungsi!with pointers} mencetak pointer. Pointer ini akan tampak seperti angka sampah (dan mungkin dicetak dalam heksadesimal\footnote{Yaitu basis 16 dengan digit 0, 1, 2, 3, 4, 5, 6, 7, 8, 9, A, B, C, D, E, dan F.} alih-alih desimal), tapi itu hanyalah indeks ke memori tempat data disimpan. (Atau indeks ke memori tempat byte pertama data disimpan, kalau datanya multi-byte.) Dalam hampir semua keadaan, termasuk yang ini, nilai angka yang tercetak tidak penting bagimu, dan hanya ditampilkan untuk mendemonstrasikan operator \hl{\mbox{\texttt{address-of}}}.

\begin{tcolorbox}[enhanced jigsaw,left=3pt,breakable,colback=white!97.2549!black,boxrule=0.2mm]

\begin{Shaded}
\begin{Highlighting}[numbers=left,,]
\PreprocessorTok{\#include }\ImportTok{\textless{}stdio.h\textgreater{}}

\DataTypeTok{int}\NormalTok{ main}\OperatorTok{(}\DataTypeTok{void}\OperatorTok{)}
\OperatorTok{\{}
    \DataTypeTok{int}\NormalTok{ i }\OperatorTok{=} \DecValTok{10}\OperatorTok{;}

\NormalTok{    printf}\OperatorTok{(}\StringTok{"The value of i is }\SpecialCharTok{\%d\textbackslash{}n}\StringTok{"}\OperatorTok{,}\NormalTok{ i}\OperatorTok{);}
\NormalTok{    printf}\OperatorTok{(}\StringTok{"And its address is }\SpecialCharTok{\%p\textbackslash{}n}\StringTok{"}\OperatorTok{,} \OperatorTok{(}\DataTypeTok{void} \OperatorTok{*)\&}\NormalTok{i}\OperatorTok{);}
\OperatorTok{\}}
\end{Highlighting}
\end{Shaded}

\end{tcolorbox}

\begin{tcolorbox}[blanker,colback=white,enhanced jigsaw,breakable,arc=0pt,outer arc=0pt,grow to left by=-1cm,enlarge top by=0.1cm,enlarge bottom by=0.1cm,borderline west={1.5pt}{-.3cm}{black}]

\begin{quote}
\textbf{Kode di atas mengandung sebuah \emph{cast}}, yang memaksa tipe ekspresi \hl{\mbox{\texttt{\&i}}} menjadi \hl{\mbox{\texttt{void*}}}. Tujuannya mencegah compiler mengeluarkan peringatan. Ini semua belum kita bahas, jadi untuk sekarang abaikan saja \hl{\mbox{\texttt{(void*)}}} di kode di atas dan anggap tidak ada.
\end{quote}

\end{tcolorbox}

Di komputer penulis, ini mencetak:

\begin{tcolorbox}[enhanced jigsaw,left=3pt,breakable,colback=white!97.2549!black,boxrule=0.2mm]

\begin{Shaded}
\begin{Highlighting}[]
\NormalTok{The value of i is 10}
\NormalTok{And its address is 0x7ffddf7072a4}
\end{Highlighting}
\end{Shaded}

\end{tcolorbox}

\index{\& operator address-of@\texttt{\&} operator address-of|)}

Kalau penasaran, angka heksadesimal itu sama dengan 140.727.326.896.068 dalam desimal (basis 10, seperti yang biasa dipakai nenek kita). Itulah indeks ke memori tempat data variabel \hl{\mbox{\texttt{i}}} disimpan. Itulah alamat \hl{\mbox{\texttt{i}}}. Itulah lokasi \hl{\mbox{\texttt{i}}}. Itulah pointer ke \hl{\mbox{\texttt{i}}}.

\begin{tcolorbox}[blanker,colback=white,enhanced jigsaw,breakable,arc=0pt,outer arc=0pt,grow to left by=-1cm,enlarge top by=0.1cm,enlarge bottom by=0.1cm,borderline west={1.5pt}{-.3cm}{black}]

\begin{quote}
\textbf{Tunggu---RAM kamu 140 terabyte?} Ya! Kamu juga kan? Tapi itu bohong besar; tentu saja tidak (sekitar 2024). Komputer modern memakai teknologi ajaib bernama \href{https://en.wikipedia.org/wiki/Virtual_memory}{virtual memory}\footnote{https://en.wikipedia.org/wiki/Virtual\_memory} yang membuat setiap proses mengira seluruh ruang memori komputer adalah miliknya sendiri, berapa pun RAM fisik yang ada di belakangnya. Jadi, meskipun alamatnya angka sebesar itu, oleh sistem virtual memory di CPU alamat itu dipetakan ke alamat memori fisik yang lebih rendah. Komputer penulis ini punya RAM 16 GB (sekali lagi, sekitar 2024, tapi pakai Linux, jadi itu sudah lebih dari cukup). Terabyte RAM? Penulis ini guru, bukan miliarder dotcom. Tidak ada dari semua ini yang perlu kita khawatirkan, kecuali soal penulis yang tidak kaya raya.
\end{quote}

\end{tcolorbox}

Angka ini disebut pointer karena ia memberi tahu di mana \hl{\mbox{\texttt{i}}} berada di memori. Seperti alamat rumah di secarik kertas yang memberi tahu di mana rumah tertentu bisa ditemukan, angka ini menunjukkan di mana di memori nilai \hl{\mbox{\texttt{i}}} bisa ditemukan. Ia menunjuk ke \hl{\mbox{\texttt{i}}}.

Sekali lagi, pada umumnya kita tidak peduli persis berapa nomor alamatnya. Kita hanya peduli bahwa itu pointer ke \hl{\mbox{\texttt{i}}}.

\section{Tipe Pointer}\label{pttypes}

\index{Pointer types|(}Jadi\ldots{} semua ini bagus. Sekarang kamu sudah bisa mengambil alamat variabel dan mencetaknya ke layar. Lumayan buat dipamerkan di resume, kan? Nah, di titik ini kamu menyela dan dengan sopan bertanya apa gunanya pointer.

Pertanyaan yang sangat bagus, dan akan kita jawab setelah pesan dari sponsor kita.

\begin{tcolorbox}[enhanced jigsaw,left=3pt,breakable,colback=white!97.2549!black,boxrule=0.2mm]

\begin{verbatim}
ACME ROBOTIC HOUSING UNIT CLEANING SERVICES. YOUR HOMESTEAD WILL BE
DRAMATICALLY IMPROVED OR YOU WILL BE TERMINATED. MESSAGE ENDS.
\end{verbatim}

\end{tcolorbox}

Selamat datang kembali di edisi lain Beej's Guide. Terakhir kali kita bertemu, kita sedang membahas cara memakai pointer. Sekarang kita akan menyimpan pointer di dalam variabel supaya bisa dipakai nanti. \emph{Tipe pointer} bisa dikenali dari tanda asterisk (\hl{\mbox{\texttt{*}}}) setelah nama tipe dan sebelum nama variabel:

\begin{tcolorbox}[enhanced jigsaw,left=3pt,breakable,colback=white!97.2549!black,boxrule=0.2mm]

\begin{Shaded}
\begin{Highlighting}[numbers=left,,]
\DataTypeTok{int}\NormalTok{ main}\OperatorTok{(}\DataTypeTok{void}\OperatorTok{)}
\OperatorTok{\{}
    \DataTypeTok{int}\NormalTok{ i}\OperatorTok{;}  \CommentTok{// tipe i adalah "int"}
    \DataTypeTok{int} \OperatorTok{*}\NormalTok{p}\OperatorTok{;} \CommentTok{// tipe p adalah "pointer ke int", atau "int{-}pointer"}
\OperatorTok{\}}
\end{Highlighting}
\end{Shaded}

\end{tcolorbox}

Nah, di sini kita punya variabel bertipe pointer, yang bisa menunjuk ke sebuah \hl{\mbox{\texttt{int}}} lain. Artinya, ia bisa menyimpan alamat \hl{\mbox{\texttt{int}}} lain. Kita tahu ia menunjuk ke \hl{\mbox{\texttt{int}}} karena tipenya \hl{\mbox{\texttt{int*}}} (dibaca ``int-pointer'').

Saat melakukan penugasan ke variabel pointer, tipe di sisi kanan harus sama dengan tipe variabel pointer-nya. Untungnya, kalau kita mengambil \hl{\mbox{\texttt{address-of}}} sebuah variabel, tipe hasilnya adalah pointer ke tipe variabel itu, sehingga penugasan seperti berikut ini sempurna:

\begin{tcolorbox}[enhanced jigsaw,left=3pt,breakable,colback=white!97.2549!black,boxrule=0.2mm]

\begin{Shaded}
\begin{Highlighting}[]
\DataTypeTok{int}\NormalTok{ i}\OperatorTok{;}
\DataTypeTok{int} \OperatorTok{*}\NormalTok{p}\OperatorTok{;}  \CommentTok{// p adalah pointer, tapi belum diinisialisasi dan menunjuk ke sampah}

\NormalTok{p }\OperatorTok{=} \OperatorTok{\&}\NormalTok{i}\OperatorTok{;}  \CommentTok{// p diberi alamat i{-}{-}sekarang p "menunjuk ke" i}
\end{Highlighting}
\end{Shaded}

\end{tcolorbox}

Di sisi kiri penugasan, ada variabel bertipe pointer-to-\hl{\mbox{\texttt{int}}} (\hl{\mbox{\texttt{int*}}}), dan di sisi kanan ada ekspresi bertipe pointer-to-\hl{\mbox{\texttt{int}}} karena \hl{\mbox{\texttt{i}}} adalah \hl{\mbox{\texttt{int}}} (address-of sebuah \hl{\mbox{\texttt{int}}} menghasilkan pointer ke \hl{\mbox{\texttt{int}}}). Alamat sesuatu bisa disimpan dalam pointer ke sesuatu itu.

Paham? Memang belum terlalu masuk akal, sebab kamu belum melihat pemakaian nyata variabel pointer, tapi kita melangkah pelan-pelan supaya tidak ada yang tersesat. Sekarang, kita kenalkan operator anti-address-of. Kira-kira ia melakukan kebalikan dari \hl{\mbox{\texttt{address-of}}}, seperti di Dunia Bizarro.\index{Pointer types|)}

\section{Dereferencing}\label{deref}

\index{Dereferencing|(}Variabel pointer bisa dianggap \emph{merujuk} ke variabel lain dengan cara menunjuk ke sana. Jarang sekali orang C membicarakan ``merujuk'' atau ``referensi'', tapi itu disebut di sini supaya nama operator ini lebih masuk akal.

Kalau kamu punya pointer ke sebuah variabel (secara kasar, ``referensi ke variabel''), kamu bisa mengakses variabel aslinya lewat pointer itu dengan \emph{men-dereference}-nya. (Boleh dibilang ini ``de-pointer''-ing pointer, tapi tidak ada yang pernah mengatakan itu.)

Kembali ke analogi kita, ini seperti melihat alamat rumah lalu pergi ke rumah itu.

Lalu apa maksud ``mengakses variabel aslinya''? Begini: kalau kamu punya variabel \hl{\mbox{\texttt{i}}} dan pointer ke \hl{\mbox{\texttt{i}}} bernama \hl{\mbox{\texttt{p}}}, kamu bisa memakai pointer yang sudah di-dereference, \hl{\mbox{\texttt{*p}}}, \emph{persis seperti variabel asli \hl{\mbox{\texttt{i}}}}!

\index{* indirection operator@\texttt{*} indirection operator|(}Pengetahuanmu sudah hampir cukup untuk mengerjakan contoh. Satu informasi terakhir: apa operator dereference-nya? Sebenarnya ia disebut \emph{indirection operator}, karena nilainya diakses secara tidak langsung lewat pointer. Dan operatornya adalah tanda asterisk, lagi: \hl{\mbox{\texttt{*}}}. Jangan tertukar dengan asterisk di deklarasi pointer tadi. Karakternya sama, tapi maknanya berbeda tergantung konteks\footnote{Dan itu belum semua! Asterisk juga dipakai dalam \hl{\mbox{\texttt{/*komentar*/}}}, perkalian, dan function prototype dengan variable length array! Semuanya \hl{\mbox{\texttt{*}}} yang sama, tapi konteksnya memberi makna yang berbeda.}.

Berikut contoh lengkapnya:

\begin{tcolorbox}[enhanced jigsaw,left=3pt,breakable,colback=white!97.2549!black,boxrule=0.2mm]

\begin{Shaded}
\begin{Highlighting}[numbers=left,,]
\PreprocessorTok{\#include }\ImportTok{\textless{}stdio.h\textgreater{}}

\DataTypeTok{int}\NormalTok{ main}\OperatorTok{(}\DataTypeTok{void}\OperatorTok{)}
\OperatorTok{\{}
    \DataTypeTok{int}\NormalTok{ i}\OperatorTok{;}
    \DataTypeTok{int} \OperatorTok{*}\NormalTok{p}\OperatorTok{;}  \CommentTok{// ini BUKAN dereference{-}{-}ini tipe "int*"}

\NormalTok{    p }\OperatorTok{=} \OperatorTok{\&}\NormalTok{i}\OperatorTok{;}  \CommentTok{// p sekarang menunjuk ke i, p menyimpan alamat i}

\NormalTok{    i }\OperatorTok{=} \DecValTok{10}\OperatorTok{;}  \CommentTok{// i sekarang 10}
    \OperatorTok{*}\NormalTok{p }\OperatorTok{=} \DecValTok{20}\OperatorTok{;} \CommentTok{// yang ditunjuk p (yaitu i!) sekarang 20!!}

\NormalTok{    printf}\OperatorTok{(}\StringTok{"i is }\SpecialCharTok{\%d\textbackslash{}n}\StringTok{"}\OperatorTok{,}\NormalTok{ i}\OperatorTok{);}   \CommentTok{// mencetak "20"}
\NormalTok{    printf}\OperatorTok{(}\StringTok{"i is }\SpecialCharTok{\%d\textbackslash{}n}\StringTok{"}\OperatorTok{,} \OperatorTok{*}\NormalTok{p}\OperatorTok{);}  \CommentTok{// "20"! dereference{-}p sama dengan i!}
\OperatorTok{\}}
\end{Highlighting}
\end{Shaded}

\end{tcolorbox}

Ingat bahwa \hl{\mbox{\texttt{p}}} menyimpan alamat \hl{\mbox{\texttt{i}}}, seperti terlihat pada penugasan ke \hl{\mbox{\texttt{p}}} di baris 8. Indirection operator menyuruh komputer untuk \emph{memakai objek yang ditunjuk pointer}, bukan pointer itu sendiri. Dengan begitu, \hl{\mbox{\texttt{*p}}} menjadi semacam alias bagi \hl{\mbox{\texttt{i}}}.\index{Dereferencing|)}\index{* indirection operator@\texttt{*} indirection operator|)}

Bagus, tapi \emph{kenapa}? Kenapa repot-repot semua ini?

\section{Mengoper Pointer sebagai Argumen}\label{ptpass}

\index{Pointers!as arguments|(}Sekarang kamu mungkin merasa sudah punya banyak pengetahuan tentang pointer, tapi nol penerapan, kan? Maksudnya, buat apa \hl{\mbox{\texttt{*p}}} kalau kamu bisa langsung menulis \hl{\mbox{\texttt{i}}}?

Nah, kawan, kekuatan sebenarnya dari pointer muncul ketika kamu mulai mengopernya ke fungsi. Kenapa ini penting? Mungkin kamu ingat bahwa argumen apa pun yang dioper ke fungsi akan dengan tekun disalin ke parameter, lalu kamu bisa mengutak-atik salinan lokal itu di dalam fungsi, dan akhirnya mengembalikan satu nilai.

Bagaimana kalau kamu ingin mendapat lebih dari satu data dari sebuah fungsi? Kamu kan hanya bisa mengembalikan satu hal. Bagaimana kalau pertanyaan itu dijawab dengan pertanyaan lain? \ldots Eh, dua pertanyaan?

Apa yang terjadi kalau sebuah pointer dioper sebagai argumen ke fungsi? Apakah salinan pointer itu ditaruh di parameter yang sesuai? \emph{Berani taruhan apa saja: iya.} Ingat bagaimana tadi kita bahas panjang lebar bahwa \emph{SETIAP ARGUMEN} disalin ke parameter dan fungsi memakai salinannya? Hal yang sama berlaku di sini. Fungsi mendapat salinan pointer.

Tapi, dan inilah bagian cerdasnya: kita sudah mengatur pointer itu lebih dulu agar menunjuk ke suatu variabel\ldots{} sehingga fungsi bisa men-dereference salinan pointer-nya untuk sampai ke variabel aslinya! Fungsi tidak bisa melihat variabel itu secara langsung, tapi ia jelas bisa men-dereference pointer ke variabel itu!

Analoginya, kamu menulis alamat rumah di secarik kertas lalu menyalinnya ke kertas lain. Sekarang ada \emph{dua} pointer ke rumah itu, dan keduanya sama-sama bisa membawamu ke rumah yang sama.

Pada pemanggilan fungsi, satu salinan disimpan di variabel pointer pada calling scope, dan salinan lainnya di variabel pointer yang menjadi parameter fungsi.

Contoh! Kita kembali ke fungsi \hl{\mbox{\texttt{increment()}}} lama, tapi kali ini kita buat ia benar-benar meng-increment nilai di luar caller.

\begin{tcolorbox}[enhanced jigsaw,left=3pt,breakable,colback=white!97.2549!black,boxrule=0.2mm]

\begin{Shaded}
\begin{Highlighting}[numbers=left,,]
\PreprocessorTok{\#include }\ImportTok{\textless{}stdio.h\textgreater{}}

\DataTypeTok{void}\NormalTok{ increment}\OperatorTok{(}\DataTypeTok{int} \OperatorTok{*}\NormalTok{p}\OperatorTok{)}  \CommentTok{// perhatikan bahwa fungsi ini menerima pointer ke int}
\OperatorTok{\{}
    \OperatorTok{*}\NormalTok{p }\OperatorTok{=} \OperatorTok{*}\NormalTok{p }\OperatorTok{+} \DecValTok{1}\OperatorTok{;}        \CommentTok{// tambahkan satu ke yang ditunjuk p}
\OperatorTok{\}}

\DataTypeTok{int}\NormalTok{ main}\OperatorTok{(}\DataTypeTok{void}\OperatorTok{)}
\OperatorTok{\{}
    \DataTypeTok{int}\NormalTok{ i }\OperatorTok{=} \DecValTok{10}\OperatorTok{;}
    \DataTypeTok{int} \OperatorTok{*}\NormalTok{j }\OperatorTok{=} \OperatorTok{\&}\NormalTok{i}\OperatorTok{;}  \CommentTok{// perhatikan address{-}of; mengubahnya menjadi pointer ke i}

\NormalTok{    printf}\OperatorTok{(}\StringTok{"i is }\SpecialCharTok{\%d\textbackslash{}n}\StringTok{"}\OperatorTok{,}\NormalTok{ i}\OperatorTok{);}        \CommentTok{// mencetak "10"}
\NormalTok{    printf}\OperatorTok{(}\StringTok{"i is also }\SpecialCharTok{\%d\textbackslash{}n}\StringTok{"}\OperatorTok{,} \OperatorTok{*}\NormalTok{j}\OperatorTok{);}  \CommentTok{// mencetak "10"}

\NormalTok{    increment}\OperatorTok{(}\NormalTok{j}\OperatorTok{);}                  \CommentTok{// j adalah int*{-}{-}ke i}

\NormalTok{    printf}\OperatorTok{(}\StringTok{"i is }\SpecialCharTok{\%d\textbackslash{}n}\StringTok{"}\OperatorTok{,}\NormalTok{ i}\OperatorTok{);}        \CommentTok{// mencetak "11"!}
\OperatorTok{\}}
\end{Highlighting}
\end{Shaded}

\end{tcolorbox}

Oke! Ada beberapa hal yang perlu diperhatikan di sini\ldots{} terutama bahwa fungsi \hl{\mbox{\texttt{increment()}}} menerima \hl{\mbox{\texttt{int*}}} sebagai argumen. Kita mengoper \hl{\mbox{\texttt{int*}}} itu dalam pemanggilan dengan mengubah variabel \hl{\mbox{\texttt{int\ i}}} menjadi \hl{\mbox{\texttt{int*}}} lewat operator \hl{\mbox{\texttt{address-of}}}. (Ingat, pointer menyimpan alamat, jadi kita membuat pointer ke sebuah variabel dengan memasangkan operator \hl{\mbox{\texttt{address-of}}} padanya.)

Fungsi \hl{\mbox{\texttt{increment()}}} mendapat salinan pointer. Pointer asli \hl{\mbox{\texttt{j}}} (di \hl{\mbox{\texttt{main()}}}) dan salinannya \hl{\mbox{\texttt{p}}} (parameter di \hl{\mbox{\texttt{increment()}}}) sama-sama menunjuk ke alamat yang sama, yaitu tempat nilai \hl{\mbox{\texttt{i}}} tersimpan. (Lagi-lagi secara analogi, seperti dua lembar kertas yang bertuliskan alamat rumah yang sama.) Men-dereference salah satunya memungkinkan kamu mengubah variabel asli \hl{\mbox{\texttt{i}}}! Fungsi bisa mengubah variabel di scope lain! Keren!

Contoh di atas sering ditulis lebih ringkas dengan menaruh address-of langsung di daftar argumen saat pemanggilan:

\begin{tcolorbox}[enhanced jigsaw,left=3pt,breakable,colback=white!97.2549!black,boxrule=0.2mm]

\begin{Shaded}
\begin{Highlighting}[]
\NormalTok{printf}\OperatorTok{(}\StringTok{"i is }\SpecialCharTok{\%d\textbackslash{}n}\StringTok{"}\OperatorTok{,}\NormalTok{ i}\OperatorTok{);}  \CommentTok{// mencetak "10"}
\NormalTok{increment}\OperatorTok{(\&}\NormalTok{i}\OperatorTok{);}
\NormalTok{printf}\OperatorTok{(}\StringTok{"i is }\SpecialCharTok{\%d\textbackslash{}n}\StringTok{"}\OperatorTok{,}\NormalTok{ i}\OperatorTok{);}  \CommentTok{// mencetak "11"!}
\end{Highlighting}
\end{Shaded}

\end{tcolorbox}

Sebagai aturan umum, kalau kamu ingin fungsi mengubah sesuatu yang kamu oper dan hasilnya terlihat dari luar, oper pointer ke sesuatu itu.

\section{\texorpdfstring{Pointer \hl{\mbox{\texttt{NULL}}}}{Pointer NULL}}\label{pointer-null}

\index{NULL pointer@\texttt{NULL} pointer|(}Variabel pointer dari tipe pointer apa pun bisa diisi nilai khusus bernama \hl{\mbox{\texttt{NULL}}}. Nilai ini menandakan bahwa pointer tersebut tidak menunjuk ke apa pun.

\begin{tcolorbox}[enhanced jigsaw,left=3pt,breakable,colback=white!97.2549!black,boxrule=0.2mm]

\begin{Shaded}
\begin{Highlighting}[]
\DataTypeTok{int} \OperatorTok{*}\NormalTok{p}\OperatorTok{;}

\NormalTok{p }\OperatorTok{=}\NormalTok{ NULL}\OperatorTok{;}
\end{Highlighting}
\end{Shaded}

\end{tcolorbox}

Karena tidak menunjuk ke nilai apa pun, men-dereference-nya adalah undefined behavior, dan kemungkinan besar berakhir dengan crash:

\begin{tcolorbox}[enhanced jigsaw,left=3pt,breakable,colback=white!97.2549!black,boxrule=0.2mm]

\begin{Shaded}
\begin{Highlighting}[]
\DataTypeTok{int} \OperatorTok{*}\NormalTok{p }\OperatorTok{=}\NormalTok{ NULL}\OperatorTok{;}

\OperatorTok{*}\NormalTok{p }\OperatorTok{=} \DecValTok{12}\OperatorTok{;}  \CommentTok{// CRASH atau SESUATU YANG BURUK. SEBAIKNYA DIHINDARI.}
\end{Highlighting}
\end{Shaded}

\end{tcolorbox}

Walaupun penciptanya menyebutnya \href{https://en.wikipedia.org/wiki/Null_pointer\#History}{kesalahan satu miliar dolar}\footnote{https://en.wikipedia.org/wiki/Null\_pointer\#History}, pointer \hl{\mbox{\texttt{NULL}}} adalah \href{https://en.wikipedia.org/wiki/Sentinel_value}{sentinel value}\footnote{https://en.wikipedia.org/wiki/Sentinel\_value} yang bagus dan penanda umum bahwa pointer belum diinisialisasi.

(Tentu saja, seperti variabel lain, pointer akan menunjuk ke sampah kecuali kamu secara eksplisit menugaskannya untuk menunjuk ke sebuah alamat atau \hl{\mbox{\texttt{NULL}}}.) \index{NULL pointer@\texttt{NULL} pointer|)}

\section{Catatan tentang Mendeklarasikan Pointer}\label{catatan-tentang-mendeklarasikan-pointer}

\index{Pointers!declarations|(}Sintaks deklarasi pointer bisa terlihat sedikit aneh. Lihat contoh ini:

\begin{tcolorbox}[enhanced jigsaw,left=3pt,breakable,colback=white!97.2549!black,boxrule=0.2mm]

\begin{Shaded}
\begin{Highlighting}[]
\DataTypeTok{int}\NormalTok{ a}\OperatorTok{;}
\DataTypeTok{int}\NormalTok{ b}\OperatorTok{;}
\end{Highlighting}
\end{Shaded}

\end{tcolorbox}

Kita bisa memadatkan itu menjadi satu baris, kan?

\begin{tcolorbox}[enhanced jigsaw,left=3pt,breakable,colback=white!97.2549!black,boxrule=0.2mm]

\begin{Shaded}
\begin{Highlighting}[]
\DataTypeTok{int}\NormalTok{ a}\OperatorTok{,}\NormalTok{ b}\OperatorTok{;}  \CommentTok{// Sama saja}
\end{Highlighting}
\end{Shaded}

\end{tcolorbox}

Jadi \hl{\mbox{\texttt{a}}} dan \hl{\mbox{\texttt{b}}} keduanya adalah \hl{\mbox{\texttt{int}}}. Tidak masalah.

Tapi bagaimana dengan ini?

\begin{tcolorbox}[enhanced jigsaw,left=3pt,breakable,colback=white!97.2549!black,boxrule=0.2mm]

\begin{Shaded}
\begin{Highlighting}[]
\DataTypeTok{int}\NormalTok{ a}\OperatorTok{;}
\DataTypeTok{int} \OperatorTok{*}\NormalTok{p}\OperatorTok{;}
\end{Highlighting}
\end{Shaded}

\end{tcolorbox}

Bisakah kita jadikan itu satu baris? Bisa. Tapi ke mana \hl{\mbox{\texttt{*}}} pergi?

Aturannya: \hl{\mbox{\texttt{*}}} diletakkan di depan setiap variabel yang bertipe pointer. Artinya, \hl{\mbox{\texttt{*}}} \emph{bukan} bagian dari \hl{\mbox{\texttt{int}}} pada contoh ini. Ia bagian dari variabel \hl{\mbox{\texttt{p}}}.

Dengan begitu, kita bisa menulis ini:

\begin{tcolorbox}[enhanced jigsaw,left=3pt,breakable,colback=white!97.2549!black,boxrule=0.2mm]

\begin{Shaded}
\begin{Highlighting}[]
\DataTypeTok{int}\NormalTok{ a}\OperatorTok{,} \OperatorTok{*}\NormalTok{p}\OperatorTok{;}  \CommentTok{// Sama saja}
\end{Highlighting}
\end{Shaded}

\end{tcolorbox}

Perlu dicatat bahwa baris berikut \emph{tidak} mendeklarasikan dua pointer:

\begin{tcolorbox}[enhanced jigsaw,left=3pt,breakable,colback=white!97.2549!black,boxrule=0.2mm]

\begin{Shaded}
\begin{Highlighting}[]
\DataTypeTok{int} \OperatorTok{*}\NormalTok{p}\OperatorTok{,}\NormalTok{ q}\OperatorTok{;}  \CommentTok{// p adalah pointer ke int; q hanya int biasa.}
\end{Highlighting}
\end{Shaded}

\end{tcolorbox}

Ini bisa sangat menjebak kalau programmer menulis baris kode berikut (valid), yang secara fungsional identik dengan yang di atas.

\begin{tcolorbox}[enhanced jigsaw,left=3pt,breakable,colback=white!97.2549!black,boxrule=0.2mm]

\begin{Shaded}
\begin{Highlighting}[]
\DataTypeTok{int}\OperatorTok{*}\NormalTok{ p}\OperatorTok{,}\NormalTok{ q}\OperatorTok{;}  \CommentTok{// p adalah pointer ke int; q hanya int biasa.}
\end{Highlighting}
\end{Shaded}

\end{tcolorbox}

Coba perhatikan ini dan tentukan variabel mana yang pointer dan mana yang bukan:

\begin{tcolorbox}[enhanced jigsaw,left=3pt,breakable,colback=white!97.2549!black,boxrule=0.2mm]

\begin{Shaded}
\begin{Highlighting}[]
\DataTypeTok{int} \OperatorTok{*}\NormalTok{a}\OperatorTok{,}\NormalTok{ b}\OperatorTok{,}\NormalTok{ c}\OperatorTok{,} \OperatorTok{*}\NormalTok{d}\OperatorTok{,}\NormalTok{ e}\OperatorTok{,} \OperatorTok{*}\NormalTok{f}\OperatorTok{,}\NormalTok{ g}\OperatorTok{,}\NormalTok{ h}\OperatorTok{,} \OperatorTok{*}\NormalTok{i}\OperatorTok{;}
\end{Highlighting}
\end{Shaded}

\end{tcolorbox}

Jawabannya ada di footnote\footnote{Variabel bertipe pointer adalah \hl{\mbox{\texttt{a}}}, \hl{\mbox{\texttt{d}}}, \hl{\mbox{\texttt{f}}}, dan \hl{\mbox{\texttt{i}}}, karena hanya merekalah yang punya \hl{\mbox{\texttt{*}}} di depannya.}.\index{Pointers!declarations|)}

\section{\texorpdfstring{\hl{\mbox{\texttt{sizeof}}} dan Pointer}{sizeof dan Pointer}}\label{sizeof-dan-pointer}

\index{Pointers!with sizeof@with \texttt{sizeof}|(}Ada sedikit sintaks di sini yang mungkin membingungkan dan sesekali akan kamu temui.

Ingat bahwa \hl{\mbox{\texttt{sizeof}}} bekerja pada \emph{tipe} dari ekspresi.

\begin{tcolorbox}[enhanced jigsaw,left=3pt,breakable,colback=white!97.2549!black,boxrule=0.2mm]

\begin{Shaded}
\begin{Highlighting}[]
\DataTypeTok{int} \OperatorTok{*}\NormalTok{p}\OperatorTok{;}

\CommentTok{// Mencetak ukuran \textquotesingle{}int\textquotesingle{}}
\NormalTok{printf}\OperatorTok{(}\StringTok{"}\SpecialCharTok{\%zu\textbackslash{}n}\StringTok{"}\OperatorTok{,} \KeywordTok{sizeof}\OperatorTok{(}\DataTypeTok{int}\OperatorTok{));}

\CommentTok{// p bertipe \textquotesingle{}int *\textquotesingle{}, jadi mencetak ukuran \textquotesingle{}int*\textquotesingle{}}
\NormalTok{printf}\OperatorTok{(}\StringTok{"}\SpecialCharTok{\%zu\textbackslash{}n}\StringTok{"}\OperatorTok{,} \KeywordTok{sizeof}\NormalTok{ p}\OperatorTok{);}

\CommentTok{// *p bertipe \textquotesingle{}int\textquotesingle{}, jadi mencetak ukuran \textquotesingle{}int\textquotesingle{}}
\NormalTok{printf}\OperatorTok{(}\StringTok{"}\SpecialCharTok{\%zu\textbackslash{}n}\StringTok{"}\OperatorTok{,} \KeywordTok{sizeof} \OperatorTok{*}\NormalTok{p}\OperatorTok{);}
\end{Highlighting}
\end{Shaded}

\end{tcolorbox}

Kamu mungkin akan menemui kode dengan \hl{\mbox{\texttt{sizeof}}} terakhir itu. Cukup ingat bahwa \hl{\mbox{\texttt{sizeof}}} hanya peduli pada tipe ekspresi, bukan variabel di dalam ekspresi itu sendiri.\index{Pointers!with sizeof@with \texttt{sizeof}|)}\index{Pointers|)}

\chapter{Arrays}\label{arrays}

\begin{tcolorbox}[blanker,colback=white,enhanced jigsaw,breakable,arc=0pt,outer arc=0pt,grow to left by=-1cm,enlarge top by=0.1cm,enlarge bottom by=0.1cm,borderline west={1.5pt}{-.3cm}{black}]

\begin{quote}
\emph{``Should array indices start at 0 or 1? My compromise of 0.5 was rejected without, I thought, proper consideration.''}

---Stan Kelly-Bootle, computer scientist
\end{quote}

\end{tcolorbox}

\index{Arrays|(} Untungnya, C punya array. Memang, C dianggap bahasa low-level\footnote{Setidaknya zaman sekarang.}, tapi konsep array sudah tertanam di dalamnya. Dan karena banyak sekali bahasa yang terinspirasi dari sintaks C, kemungkinan besar kamu sudah akrab dengan \hl{\mbox{\texttt{{[}}}} dan \hl{\mbox{\texttt{{]}}}} untuk mendeklarasikan dan memakai array.

Tapi C hanya \emph{nyaris} punya array! Seperti yang akan kita lihat nanti, array di C hanyalah syntactic sugar---jauh di dalamnya, semuanya adalah pointer dan semacamnya. \emph{Panik!} Tapi untuk sekarang, kita pakai saja sebagai array. \emph{Fiuh}.

\section{Contoh Mudah}\label{contoh-mudah}

Langsung saja kita buat contohnya:

\index{Arrays!indexing|(}

\begin{tcolorbox}[enhanced jigsaw,left=3pt,breakable,colback=white!97.2549!black,boxrule=0.2mm]

\begin{Shaded}
\begin{Highlighting}[numbers=left,,]
\PreprocessorTok{\#include }\ImportTok{\textless{}stdio.h\textgreater{}}

\DataTypeTok{int}\NormalTok{ main}\OperatorTok{(}\DataTypeTok{void}\OperatorTok{)}
\OperatorTok{\{}
    \DataTypeTok{int}\NormalTok{ i}\OperatorTok{;}
    \DataTypeTok{float}\NormalTok{ f}\OperatorTok{[}\DecValTok{4}\OperatorTok{];}  \CommentTok{// Deklarasikan array berisi 4 float}

\NormalTok{    f}\OperatorTok{[}\DecValTok{0}\OperatorTok{]} \OperatorTok{=} \FloatTok{3.14159}\OperatorTok{;}  \CommentTok{// Indeks dimulai dari 0, tentu saja.}
\NormalTok{    f}\OperatorTok{[}\DecValTok{1}\OperatorTok{]} \OperatorTok{=} \FloatTok{1.41421}\OperatorTok{;}
\NormalTok{    f}\OperatorTok{[}\DecValTok{2}\OperatorTok{]} \OperatorTok{=} \FloatTok{1.61803}\OperatorTok{;}
\NormalTok{    f}\OperatorTok{[}\DecValTok{3}\OperatorTok{]} \OperatorTok{=} \FloatTok{2.71828}\OperatorTok{;}

    \CommentTok{// Cetak semuanya:}

    \ControlFlowTok{for} \OperatorTok{(}\NormalTok{i }\OperatorTok{=} \DecValTok{0}\OperatorTok{;}\NormalTok{ i }\OperatorTok{\textless{}} \DecValTok{4}\OperatorTok{;}\NormalTok{ i}\OperatorTok{++)} \OperatorTok{\{}
\NormalTok{        printf}\OperatorTok{(}\StringTok{"}\SpecialCharTok{\%f\textbackslash{}n}\StringTok{"}\OperatorTok{,}\NormalTok{ f}\OperatorTok{[}\NormalTok{i}\OperatorTok{]);}
    \OperatorTok{\}}
\OperatorTok{\}}
\end{Highlighting}
\end{Shaded}

\end{tcolorbox}

Saat mendeklarasikan array, kamu harus menentukan ukurannya. Dan ukuran itu harus tetap\footnote{Sebenarnya tidak juga, tapi variable-length array---yang penulis sendiri kurang suka---adalah cerita untuk lain waktu.}.

Pada contoh di atas, kita membuat array berisi 4 \hl{\mbox{\texttt{float}}}. Angka di dalam kurung siku pada deklarasi menunjukkan hal itu.

Di baris-baris berikutnya, kita mengakses nilai di dalam array, baik untuk mengisi maupun membacanya, lagi-lagi dengan kurung siku. \index{Arrays!indexing|)}

Semoga ini terlihat akrab dari bahasa-bahasa yang sudah kamu kenal!

\section{Mendapatkan Panjang Array}\label{mendapatkan-panjang-array}

\index{Arrays!getting the length|(} Tidak bisa\ldots sebenarnya. C tidak mencatat informasi ini\footnote{Karena di balik layar array hanyalah pointer ke elemen pertama array, tidak ada informasi tambahan yang mencatat panjangnya.}. Kamu harus mengelolanya sendiri di variabel lain.

Kalau dibilang ``tidak bisa'', maksudnya ada beberapa keadaan ketika kamu \emph{bisa}. Ada trik untuk mendapatkan jumlah elemen array di scope tempat array itu dideklarasikan. Tapi pada umumnya, trik ini tidak akan berfungsi seperti yang kamu harapkan kalau array dioper ke fungsi\footnote{Karena ketika sebuah array dioper ke fungsi, yang sebenarnya dioper hanyalah pointer ke elemen pertama array itu, bukan ``seluruh'' array.}.

Mari kita lihat triknya. Idenya, kita ambil \index{sizeof operator@\texttt{sizeof} operator!with arrays|(}\hl{\mbox{\texttt{sizeof}}} array itu, lalu dibagi ukuran tiap elemennya untuk mendapatkan panjangnya. Misalnya, kalau sebuah \hl{\mbox{\texttt{int}}} berukuran 4 byte dan array-nya berukuran 32 byte, berarti di dalamnya ada tempat untuk \(\frac{32}{4}\) atau \(8\) \hl{\mbox{\texttt{int}}}.

\begin{tcolorbox}[enhanced jigsaw,left=3pt,breakable,colback=white!97.2549!black,boxrule=0.2mm]

\begin{Shaded}
\begin{Highlighting}[]
\DataTypeTok{int}\NormalTok{ x}\OperatorTok{[}\DecValTok{12}\OperatorTok{];}  \CommentTok{// 12 int}

\NormalTok{printf}\OperatorTok{(}\StringTok{"}\SpecialCharTok{\%zu\textbackslash{}n}\StringTok{"}\OperatorTok{,} \KeywordTok{sizeof}\NormalTok{ x}\OperatorTok{);}     \CommentTok{// total 48 byte}
\NormalTok{printf}\OperatorTok{(}\StringTok{"}\SpecialCharTok{\%zu\textbackslash{}n}\StringTok{"}\OperatorTok{,} \KeywordTok{sizeof}\OperatorTok{(}\DataTypeTok{int}\OperatorTok{));}  \CommentTok{// 4 byte per int}

\NormalTok{printf}\OperatorTok{(}\StringTok{"}\SpecialCharTok{\%zu\textbackslash{}n}\StringTok{"}\OperatorTok{,} \KeywordTok{sizeof}\NormalTok{ x }\OperatorTok{/} \KeywordTok{sizeof}\OperatorTok{(}\DataTypeTok{int}\OperatorTok{));}  \CommentTok{// 48/4 = 12 int!}
\end{Highlighting}
\end{Shaded}

\end{tcolorbox}

Kalau array-nya berisi \hl{\mbox{\texttt{char}}}, maka \hl{\mbox{\texttt{sizeof}}} array itu \emph{adalah} jumlah elemennya, karena \hl{\mbox{\texttt{sizeof(char)}}} didefinisikan sebagai 1. Untuk tipe lainnya, kamu harus membaginya dengan ukuran tiap elemen.

Tapi trik ini hanya berfungsi di scope tempat array didefinisikan. Kalau array dioper ke fungsi, trik ini gagal. Bahkan kalau ukurannya ditulis ``besar'' di signature fungsi:

\begin{tcolorbox}[enhanced jigsaw,left=3pt,breakable,colback=white!97.2549!black,boxrule=0.2mm]

\begin{Shaded}
\begin{Highlighting}[]
\DataTypeTok{void}\NormalTok{ foo}\OperatorTok{(}\DataTypeTok{int}\NormalTok{ x}\OperatorTok{[}\DecValTok{12}\OperatorTok{])}
\OperatorTok{\{}
\NormalTok{    printf}\OperatorTok{(}\StringTok{"}\SpecialCharTok{\%zu\textbackslash{}n}\StringTok{"}\OperatorTok{,} \KeywordTok{sizeof}\NormalTok{ x}\OperatorTok{);}     \CommentTok{// 8?! Ke mana perginya 48?}
\NormalTok{    printf}\OperatorTok{(}\StringTok{"}\SpecialCharTok{\%zu\textbackslash{}n}\StringTok{"}\OperatorTok{,} \KeywordTok{sizeof}\OperatorTok{(}\DataTypeTok{int}\OperatorTok{));}  \CommentTok{// 4 byte per int}

\NormalTok{    printf}\OperatorTok{(}\StringTok{"}\SpecialCharTok{\%zu\textbackslash{}n}\StringTok{"}\OperatorTok{,} \KeywordTok{sizeof}\NormalTok{ x }\OperatorTok{/} \KeywordTok{sizeof}\OperatorTok{(}\DataTypeTok{int}\OperatorTok{));}  \CommentTok{// 8/4 = 2 int?? SALAH.}
\OperatorTok{\}}
\end{Highlighting}
\end{Shaded}

\end{tcolorbox}

Ini terjadi karena ketika array ``dioper'' ke fungsi, yang dioper hanyalah pointer ke elemen pertama, dan itulah yang diukur \hl{\mbox{\texttt{sizeof}}}. Pembahasan lebih lanjut ada di bagian \hyperref[passing1darrays]{Mengoper Array Satu Dimensi ke Fungsi} di bawah.

Satu hal lagi yang bisa dilakukan dengan \hl{\mbox{\texttt{sizeof}}} dan array: mengetahui ukuran array dengan jumlah elemen tertentu tanpa mendeklarasikan array-nya. Ini seperti mengetahui ukuran \hl{\mbox{\texttt{int}}} lewat \hl{\mbox{\texttt{sizeof(int)}}}.

Misalnya, untuk mengetahui berapa byte yang dibutuhkan array berisi 48 \hl{\mbox{\texttt{double}}}, tulis begini:

\begin{tcolorbox}[enhanced jigsaw,left=3pt,breakable,colback=white!97.2549!black,boxrule=0.2mm]

\begin{Shaded}
\begin{Highlighting}[]
\KeywordTok{sizeof}\OperatorTok{(}\DataTypeTok{double} \OperatorTok{[}\DecValTok{48}\OperatorTok{]);}
\end{Highlighting}
\end{Shaded}

\end{tcolorbox}

\index{sizeof operator@\texttt{sizeof} operator!with arrays|)} \index{Arrays!getting the length|)}

\section{Array Initializer}\label{array-initializer}

\index{Array initializers|(} Array bisa diinisialisasi lebih dulu dengan konstanta:

\begin{tcolorbox}[enhanced jigsaw,left=3pt,breakable,colback=white!97.2549!black,boxrule=0.2mm]

\begin{Shaded}
\begin{Highlighting}[numbers=left,,]
\PreprocessorTok{\#include }\ImportTok{\textless{}stdio.h\textgreater{}}

\DataTypeTok{int}\NormalTok{ main}\OperatorTok{(}\DataTypeTok{void}\OperatorTok{)}
\OperatorTok{\{}
    \DataTypeTok{int}\NormalTok{ i}\OperatorTok{;}
    \DataTypeTok{int}\NormalTok{ a}\OperatorTok{[}\DecValTok{5}\OperatorTok{]} \OperatorTok{=} \OperatorTok{\{}\DecValTok{22}\OperatorTok{,} \DecValTok{37}\OperatorTok{,} \DecValTok{3490}\OperatorTok{,} \DecValTok{18}\OperatorTok{,} \DecValTok{95}\OperatorTok{\};}  \CommentTok{// Inisialisasi dengan nilai{-}nilai ini}

    \ControlFlowTok{for} \OperatorTok{(}\NormalTok{i }\OperatorTok{=} \DecValTok{0}\OperatorTok{;}\NormalTok{ i }\OperatorTok{\textless{}} \DecValTok{5}\OperatorTok{;}\NormalTok{ i}\OperatorTok{++)} \OperatorTok{\{}
\NormalTok{        printf}\OperatorTok{(}\StringTok{"}\SpecialCharTok{\%d\textbackslash{}n}\StringTok{"}\OperatorTok{,}\NormalTok{ a}\OperatorTok{[}\NormalTok{i}\OperatorTok{]);}
    \OperatorTok{\}}
\OperatorTok{\}}
\end{Highlighting}
\end{Shaded}

\end{tcolorbox}

Jangan pernah memasukkan item di initializer lebih banyak daripada tempat yang tersedia di array, kalau tidak compiler akan ngambek:

\begin{tcolorbox}[enhanced jigsaw,left=3pt,breakable,colback=white!97.2549!black,boxrule=0.2mm]

\begin{Shaded}
\begin{Highlighting}[]
\ExtensionTok{foo.c:}\NormalTok{ In function ‘main’:}
\ExtensionTok{foo.c:6:39:}\NormalTok{ warning: excess elements in array initializer}
    \ExtensionTok{6} \KeywordTok{|}     \ExtensionTok{int}\NormalTok{ a}\PreprocessorTok{[}\StringTok{5}\PreprocessorTok{]} \KeywordTok{=}\NormalTok{ \{22, 37, 3490, 18, 95, 999}\ErrorTok{\}}\KeywordTok{;}
      \KeywordTok{|}                                       \ExtensionTok{\^{}\textasciitilde{}\textasciitilde{}}
\ExtensionTok{foo.c:6:39:}\NormalTok{ note: }\PreprocessorTok{(}\SpecialStringTok{near initialization for ‘a’}\PreprocessorTok{)}
\end{Highlighting}
\end{Shaded}

\end{tcolorbox}

Tapi (fakta menarik!) item di initializer boleh \emph{lebih sedikit} daripada tempat yang tersedia di array. Elemen sisanya otomatis diinisialisasi dengan nol. Ini berlaku umum untuk semua jenis array initializer: kalau ada initializer, semua yang tidak diberi nilai secara eksplisit akan bernilai nol.

\begin{tcolorbox}[enhanced jigsaw,left=3pt,breakable,colback=white!97.2549!black,boxrule=0.2mm]

\begin{Shaded}
\begin{Highlighting}[]
\DataTypeTok{int}\NormalTok{ a}\OperatorTok{[}\DecValTok{5}\OperatorTok{]} \OperatorTok{=} \OperatorTok{\{}\DecValTok{22}\OperatorTok{,} \DecValTok{37}\OperatorTok{,} \DecValTok{3490}\OperatorTok{\};}

\CommentTok{// sama dengan:}

\DataTypeTok{int}\NormalTok{ a}\OperatorTok{[}\DecValTok{5}\OperatorTok{]} \OperatorTok{=} \OperatorTok{\{}\DecValTok{22}\OperatorTok{,} \DecValTok{37}\OperatorTok{,} \DecValTok{3490}\OperatorTok{,} \DecValTok{0}\OperatorTok{,} \DecValTok{0}\OperatorTok{\};}
\end{Highlighting}
\end{Shaded}

\end{tcolorbox}

Bentuk pintas yang sering dipakai untuk mengisi seluruh array dengan nol:

\begin{tcolorbox}[enhanced jigsaw,left=3pt,breakable,colback=white!97.2549!black,boxrule=0.2mm]

\begin{Shaded}
\begin{Highlighting}[]
\DataTypeTok{int}\NormalTok{ a}\OperatorTok{[}\DecValTok{100}\OperatorTok{]} \OperatorTok{=} \OperatorTok{\{}\DecValTok{0}\OperatorTok{\};}
\end{Highlighting}
\end{Shaded}

\end{tcolorbox}

Artinya, ``Buat elemen pertama bernilai nol, lalu secara otomatis buat sisanya nol juga.''

Kamu juga bisa mengisi elemen tertentu di initializer dengan menyebutkan indeksnya! Kalau kamu melakukan ini, C dengan senang hati melanjutkan pengisian nilai berikutnya sampai initializer habis, dan mengisi sisanya dengan \hl{\mbox{\texttt{0}}}.

Caranya, taruh indeks di dalam kurung siku diikuti \hl{\mbox{\texttt{=}}}, lalu tulis nilainya.

Berikut contoh pembuatan array:

\begin{tcolorbox}[enhanced jigsaw,left=3pt,breakable,colback=white!97.2549!black,boxrule=0.2mm]

\begin{Shaded}
\begin{Highlighting}[]
\DataTypeTok{int}\NormalTok{ a}\OperatorTok{[}\DecValTok{10}\OperatorTok{]} \OperatorTok{=} \OperatorTok{\{}\DecValTok{0}\OperatorTok{,} \DecValTok{11}\OperatorTok{,} \DecValTok{22}\OperatorTok{,} \OperatorTok{[}\DecValTok{5}\OperatorTok{]=}\DecValTok{55}\OperatorTok{,} \DecValTok{66}\OperatorTok{,} \DecValTok{77}\OperatorTok{\};}
\end{Highlighting}
\end{Shaded}

\end{tcolorbox}

Karena kita menetapkan indeks 5 sebagai tempat awal \hl{\mbox{\texttt{55}}}, data hasil di dalam array adalah:

\begin{tcolorbox}[enhanced jigsaw,left=3pt,breakable,colback=white!97.2549!black,boxrule=0.2mm]

\begin{Shaded}
\begin{Highlighting}[]
\NormalTok{0 11 22 0 0 55 66 77 0 0}
\end{Highlighting}
\end{Shaded}

\end{tcolorbox}

Ekspresi konstan sederhana juga boleh dipakai di sana.

\begin{tcolorbox}[enhanced jigsaw,left=3pt,breakable,colback=white!97.2549!black,boxrule=0.2mm]

\begin{Shaded}
\begin{Highlighting}[]
\PreprocessorTok{\#define COUNT }\DecValTok{5}

\DataTypeTok{int}\NormalTok{ a}\OperatorTok{[}\NormalTok{COUNT}\OperatorTok{]} \OperatorTok{=} \OperatorTok{\{[}\NormalTok{COUNT}\OperatorTok{{-}}\DecValTok{3}\OperatorTok{]=}\DecValTok{3}\OperatorTok{,} \DecValTok{2}\OperatorTok{,} \DecValTok{1}\OperatorTok{\};}
\end{Highlighting}
\end{Shaded}

\end{tcolorbox}

yang menghasilkan:

\begin{tcolorbox}[enhanced jigsaw,left=3pt,breakable,colback=white!97.2549!black,boxrule=0.2mm]

\begin{Shaded}
\begin{Highlighting}[]
\NormalTok{0 0 3 2 1}
\end{Highlighting}
\end{Shaded}

\end{tcolorbox}

Terakhir, kamu juga bisa menyuruh C menghitung ukuran array dari initializer-nya, cukup dengan tidak menulis ukurannya:

\begin{tcolorbox}[enhanced jigsaw,left=3pt,breakable,colback=white!97.2549!black,boxrule=0.2mm]

\begin{Shaded}
\begin{Highlighting}[]
\DataTypeTok{int}\NormalTok{ a}\OperatorTok{[}\DecValTok{3}\OperatorTok{]} \OperatorTok{=} \OperatorTok{\{}\DecValTok{22}\OperatorTok{,} \DecValTok{37}\OperatorTok{,} \DecValTok{3490}\OperatorTok{\};}

\CommentTok{// sama dengan:}

\DataTypeTok{int}\NormalTok{ a}\OperatorTok{[]} \OperatorTok{=} \OperatorTok{\{}\DecValTok{22}\OperatorTok{,} \DecValTok{37}\OperatorTok{,} \DecValTok{3490}\OperatorTok{\};}  \CommentTok{// Ukurannya tidak ditulis!}
\end{Highlighting}
\end{Shaded}

\end{tcolorbox}

\index{Array initializers|)}

\section{Di Luar Batas!}\label{di-luar-batas}

\index{Arrays!out of bounds|(} C tidak melarang kamu mengakses array di luar batasnya. Bahkan mungkin tidak ada peringatan sama sekali.

Kita pinjam contoh di atas dan terus mencetak melewati ujung array. Array ini hanya punya 5 elemen, tapi mari coba cetak 10 dan lihat apa yang terjadi:

\begin{tcolorbox}[enhanced jigsaw,left=3pt,breakable,colback=white!97.2549!black,boxrule=0.2mm]

\begin{Shaded}
\begin{Highlighting}[numbers=left,,]
\PreprocessorTok{\#include }\ImportTok{\textless{}stdio.h\textgreater{}}

\DataTypeTok{int}\NormalTok{ main}\OperatorTok{(}\DataTypeTok{void}\OperatorTok{)}
\OperatorTok{\{}
    \DataTypeTok{int}\NormalTok{ i}\OperatorTok{;}
    \DataTypeTok{int}\NormalTok{ a}\OperatorTok{[}\DecValTok{5}\OperatorTok{]} \OperatorTok{=} \OperatorTok{\{}\DecValTok{22}\OperatorTok{,} \DecValTok{37}\OperatorTok{,} \DecValTok{3490}\OperatorTok{,} \DecValTok{18}\OperatorTok{,} \DecValTok{95}\OperatorTok{\};}

    \ControlFlowTok{for} \OperatorTok{(}\NormalTok{i }\OperatorTok{=} \DecValTok{0}\OperatorTok{;}\NormalTok{ i }\OperatorTok{\textless{}} \DecValTok{10}\OperatorTok{;}\NormalTok{ i}\OperatorTok{++)} \OperatorTok{\{}  \CommentTok{// KABAR BURUK: mencetak terlalu banyak elemen!}
\NormalTok{        printf}\OperatorTok{(}\StringTok{"}\SpecialCharTok{\%d\textbackslash{}n}\StringTok{"}\OperatorTok{,}\NormalTok{ a}\OperatorTok{[}\NormalTok{i}\OperatorTok{]);}
    \OperatorTok{\}}
\OperatorTok{\}}
\end{Highlighting}
\end{Shaded}

\end{tcolorbox}

Di komputer penulis, hasilnya:

\begin{tcolorbox}[enhanced jigsaw,left=3pt,breakable,colback=white!97.2549!black,boxrule=0.2mm]

\begin{Shaded}
\begin{Highlighting}[]
\NormalTok{22}
\NormalTok{37}
\NormalTok{3490}
\NormalTok{18}
\NormalTok{95}
\NormalTok{32765}
\NormalTok{1847052032}
\NormalTok{1780534144}
\NormalTok{{-}56487472}
\NormalTok{21890}
\end{Highlighting}
\end{Shaded}

\end{tcolorbox}

Waduh! Apa itu? Ternyata mencetak melewati ujung array menghasilkan apa yang disebut developer C sebagai \emph{undefined behavior}. Makhluk ini akan kita bahas lebih lanjut nanti, tapi untuk sekarang artinya: ``Kamu sudah melakukan hal buruk, dan apa pun bisa terjadi selama programmu berjalan.''

Dan ``apa pun'' di sini biasanya berarti menemukan angka nol, menemukan angka sampah, atau crash. Tapi sebenarnya spesifikasi C menyatakan bahwa dalam keadaan ini compiler boleh menghasilkan kode yang melakukan \emph{apa saja}\footnote{Pada masa MS-DOS dulu, sebelum ada proteksi memori, penulis pernah menulis kode C yang sangat kasar dan sengaja melakukan segala macam undefined behavior. Tapi penulis tahu apa yang dilakukan, dan semuanya berjalan cukup lancar. Sampai suatu langkah keliru menyebabkan komputer hang dan, seperti yang ketahuan saat reboot, menghapus semua pengaturan BIOS. Seru sekali. (Salam untuk @man atas masa-masa seru itu.)}.

Singkatnya: jangan pernah melakukan apa pun yang menyebabkan undefined behavior. Pernah\footnote{Banyak hal yang menyebabkan undefined behavior, bukan hanya akses array di luar batas. Inilah yang membuat bahasa C begitu \emph{menarik}.}. \index{Arrays!out of bounds|)}

\section{Array Multidimensi}\label{array-multidimensi}

\index{Arrays!multidimensional|(} Kamu boleh menambahkan dimensi sebanyak apa pun pada array.

\begin{tcolorbox}[enhanced jigsaw,left=3pt,breakable,colback=white!97.2549!black,boxrule=0.2mm]

\begin{Shaded}
\begin{Highlighting}[]
\DataTypeTok{int}\NormalTok{ a}\OperatorTok{[}\DecValTok{10}\OperatorTok{];}
\DataTypeTok{int}\NormalTok{ b}\OperatorTok{[}\DecValTok{2}\OperatorTok{][}\DecValTok{7}\OperatorTok{];}
\DataTypeTok{int}\NormalTok{ c}\OperatorTok{[}\DecValTok{4}\OperatorTok{][}\DecValTok{5}\OperatorTok{][}\DecValTok{6}\OperatorTok{];}
\end{Highlighting}
\end{Shaded}

\end{tcolorbox}

Array ini disimpan di memori dengan \href{https://en.wikipedia.org/wiki/Row-_and_column-major_order}{row-major order}\footnote{https://en.wikipedia.org/wiki/Row-\_and\_column-major\_order}. Artinya, pada array 2D, indeks pertama menunjukkan baris, dan indeks kedua menunjukkan kolom.

Array multidimensi juga bisa diberi initializer dengan cara menyarangkannya:

\begin{tcolorbox}[enhanced jigsaw,left=3pt,breakable,colback=white!97.2549!black,boxrule=0.2mm]

\begin{Shaded}
\begin{Highlighting}[numbers=left,,]
\PreprocessorTok{\#include }\ImportTok{\textless{}stdio.h\textgreater{}}

\DataTypeTok{int}\NormalTok{ main}\OperatorTok{(}\DataTypeTok{void}\OperatorTok{)}
\OperatorTok{\{}
    \DataTypeTok{int}\NormalTok{ row}\OperatorTok{,}\NormalTok{ col}\OperatorTok{;}

    \DataTypeTok{int}\NormalTok{ a}\OperatorTok{[}\DecValTok{2}\OperatorTok{][}\DecValTok{5}\OperatorTok{]} \OperatorTok{=} \OperatorTok{\{}      \CommentTok{// Inisialisasi array 2D}
        \OperatorTok{\{}\DecValTok{0}\OperatorTok{,} \DecValTok{1}\OperatorTok{,} \DecValTok{2}\OperatorTok{,} \DecValTok{3}\OperatorTok{,} \DecValTok{4}\OperatorTok{\},}
        \OperatorTok{\{}\DecValTok{5}\OperatorTok{,} \DecValTok{6}\OperatorTok{,} \DecValTok{7}\OperatorTok{,} \DecValTok{8}\OperatorTok{,} \DecValTok{9}\OperatorTok{\}}
    \OperatorTok{\};}

    \ControlFlowTok{for} \OperatorTok{(}\NormalTok{row }\OperatorTok{=} \DecValTok{0}\OperatorTok{;}\NormalTok{ row }\OperatorTok{\textless{}} \DecValTok{2}\OperatorTok{;}\NormalTok{ row}\OperatorTok{++)} \OperatorTok{\{}
        \ControlFlowTok{for} \OperatorTok{(}\NormalTok{col }\OperatorTok{=} \DecValTok{0}\OperatorTok{;}\NormalTok{ col }\OperatorTok{\textless{}} \DecValTok{5}\OperatorTok{;}\NormalTok{ col}\OperatorTok{++)} \OperatorTok{\{}
\NormalTok{            printf}\OperatorTok{(}\StringTok{"(}\SpecialCharTok{\%d}\StringTok{,}\SpecialCharTok{\%d}\StringTok{) = }\SpecialCharTok{\%d\textbackslash{}n}\StringTok{"}\OperatorTok{,}\NormalTok{ row}\OperatorTok{,}\NormalTok{ col}\OperatorTok{,}\NormalTok{ a}\OperatorTok{[}\NormalTok{row}\OperatorTok{][}\NormalTok{col}\OperatorTok{]);}
        \OperatorTok{\}}
    \OperatorTok{\}}
\OperatorTok{\}}
\end{Highlighting}
\end{Shaded}

\end{tcolorbox}

Outputnya:

\begin{tcolorbox}[enhanced jigsaw,left=3pt,breakable,colback=white!97.2549!black,boxrule=0.2mm]

\begin{Shaded}
\begin{Highlighting}[]
\NormalTok{(0,0) = 0}
\NormalTok{(0,1) = 1}
\NormalTok{(0,2) = 2}
\NormalTok{(0,3) = 3}
\NormalTok{(0,4) = 4}
\NormalTok{(1,0) = 5}
\NormalTok{(1,1) = 6}
\NormalTok{(1,2) = 7}
\NormalTok{(1,3) = 8}
\NormalTok{(1,4) = 9}
\end{Highlighting}
\end{Shaded}

\end{tcolorbox}

Dan kamu juga bisa memakai indeks eksplisit:

\begin{tcolorbox}[enhanced jigsaw,left=3pt,breakable,colback=white!97.2549!black,boxrule=0.2mm]

\begin{Shaded}
\begin{Highlighting}[]
\CommentTok{// Buat matriks identitas 3x3}

\DataTypeTok{int}\NormalTok{ a}\OperatorTok{[}\DecValTok{3}\OperatorTok{][}\DecValTok{3}\OperatorTok{]} \OperatorTok{=} \OperatorTok{\{[}\DecValTok{0}\OperatorTok{][}\DecValTok{0}\OperatorTok{]=}\DecValTok{1}\OperatorTok{,} \OperatorTok{[}\DecValTok{1}\OperatorTok{][}\DecValTok{1}\OperatorTok{]=}\DecValTok{1}\OperatorTok{,} \OperatorTok{[}\DecValTok{2}\OperatorTok{][}\DecValTok{2}\OperatorTok{]=}\DecValTok{1}\OperatorTok{\};}
\end{Highlighting}
\end{Shaded}

\end{tcolorbox}

yang membangun array 2D seperti ini:

\begin{tcolorbox}[enhanced jigsaw,left=3pt,breakable,colback=white!97.2549!black,boxrule=0.2mm]

\begin{Shaded}
\begin{Highlighting}[]
\NormalTok{1 0 0}
\NormalTok{0 1 0}
\NormalTok{0 0 1}
\end{Highlighting}
\end{Shaded}

\end{tcolorbox}

\index{Arrays!multidimensional|)}

\section{Array dan Pointer}\label{array-dan-pointer}

\index{Arrays!as pointers|(} {[}\emph{Sambil lalu}{]} Nah\ldots{} tadi di atas mungkin ada disinggung bahwa array itu pointer, jauh di dalamnya? Sekarang kita menyelam sebentar di perairan dangkal supaya semuanya tidak benar-benar membingungkan. Nanti kita bahas hubungan sebenarnya antara array dan pointer, tapi untuk saat ini kita hanya akan melihat cara mengoper array ke fungsi.

\subsection{Mendapatkan Pointer ke Array}\label{mendapatkan-pointer-ke-array}

Ada rahasia yang perlu disampaikan. Secara umum, ketika programmer C berbicara tentang pointer ke array, yang mereka maksud adalah pointer \emph{ke elemen pertama} array\footnote{Ini secara teknis tidak tepat, karena pointer ke array dan pointer ke elemen pertama array punya tipe yang berbeda. Tapi soal itu kita urus nanti saat sudah sampai.}.

Jadi mari kita ambil pointer ke elemen pertama sebuah array.

\begin{tcolorbox}[enhanced jigsaw,left=3pt,breakable,colback=white!97.2549!black,boxrule=0.2mm]

\begin{Shaded}
\begin{Highlighting}[numbers=left,,]
\PreprocessorTok{\#include }\ImportTok{\textless{}stdio.h\textgreater{}}

\DataTypeTok{int}\NormalTok{ main}\OperatorTok{(}\DataTypeTok{void}\OperatorTok{)}
\OperatorTok{\{}
    \DataTypeTok{int}\NormalTok{ a}\OperatorTok{[}\DecValTok{5}\OperatorTok{]} \OperatorTok{=} \OperatorTok{\{}\DecValTok{11}\OperatorTok{,} \DecValTok{22}\OperatorTok{,} \DecValTok{33}\OperatorTok{,} \DecValTok{44}\OperatorTok{,} \DecValTok{55}\OperatorTok{\};}
    \DataTypeTok{int} \OperatorTok{*}\NormalTok{p}\OperatorTok{;}

\NormalTok{    p }\OperatorTok{=} \OperatorTok{\&}\NormalTok{a}\OperatorTok{[}\DecValTok{0}\OperatorTok{];}  \CommentTok{// p menunjuk ke array}
                \CommentTok{// Tepatnya, ke elemen pertama}

\NormalTok{    printf}\OperatorTok{(}\StringTok{"}\SpecialCharTok{\%d\textbackslash{}n}\StringTok{"}\OperatorTok{,} \OperatorTok{*}\NormalTok{p}\OperatorTok{);}  \CommentTok{// Mencetak "11"}
\OperatorTok{\}}
\end{Highlighting}
\end{Shaded}

\end{tcolorbox}

Karena begitu sering dipakai di C, bahasanya menyediakan jalan pintas:

\begin{tcolorbox}[enhanced jigsaw,left=3pt,breakable,colback=white!97.2549!black,boxrule=0.2mm]

\begin{Shaded}
\begin{Highlighting}[numbers=left,,]
\NormalTok{p }\OperatorTok{=} \OperatorTok{\&}\NormalTok{a}\OperatorTok{[}\DecValTok{0}\OperatorTok{];}  \CommentTok{// p menunjuk ke array}

\CommentTok{// sama dengan:}

\NormalTok{p }\OperatorTok{=}\NormalTok{ a}\OperatorTok{;}      \CommentTok{// p menunjuk ke array, tapi tampilannya jauh lebih enak!}
\end{Highlighting}
\end{Shaded}

\end{tcolorbox}

Menyebut nama array saja, tanpa embel-embel, sama dengan mengambil pointer ke elemen pertama array! Ini akan sering kita pakai di contoh-contoh berikutnya.

Tapi tunggu sebentar---bukankah \hl{\mbox{\texttt{p}}} itu \hl{\mbox{\texttt{int*}}}? Dan \hl{\mbox{\texttt{*p}}} menghasilkan \hl{\mbox{\texttt{11}}}, sama seperti \hl{\mbox{\texttt{a{[}0{]}}}}? Yessss. Kamu mulai melihat sekilas hubungan antara array dan pointer di C. (Kita akan membahasnya jauh lebih dalam di bab \hyperref[pointers2]{Pointer II}, di bagian \hyperref[arraypointerequiv]{Kesetaraan Array/Pointer}.) \index{Arrays!as pointers|)}

\subsection{Mengoper Array Satu Dimensi ke Fungsi}\label{passing1darrays}

\index{Arrays!passing to functions|(} Kita coba contoh dengan array satu dimensi. Penulis akan menulis dua fungsi yang bisa menerima array dan melakukan hal berbeda.

Bersiaplah untuk signature fungsi yang bikin takjub!

\begin{tcolorbox}[enhanced jigsaw,left=3pt,breakable,colback=white!97.2549!black,boxrule=0.2mm]

\begin{Shaded}
\begin{Highlighting}[numbers=left,,]
\PreprocessorTok{\#include }\ImportTok{\textless{}stdio.h\textgreater{}}

\CommentTok{// Dioper sebagai pointer ke elemen pertama}
\DataTypeTok{void}\NormalTok{ times2}\OperatorTok{(}\DataTypeTok{int} \OperatorTok{*}\NormalTok{a}\OperatorTok{,} \DataTypeTok{int}\NormalTok{ len}\OperatorTok{)}
\OperatorTok{\{}
    \ControlFlowTok{for} \OperatorTok{(}\DataTypeTok{int}\NormalTok{ i }\OperatorTok{=} \DecValTok{0}\OperatorTok{;}\NormalTok{ i }\OperatorTok{\textless{}}\NormalTok{ len}\OperatorTok{;}\NormalTok{ i}\OperatorTok{++)}
\NormalTok{        printf}\OperatorTok{(}\StringTok{"}\SpecialCharTok{\%d\textbackslash{}n}\StringTok{"}\OperatorTok{,}\NormalTok{ a}\OperatorTok{[}\NormalTok{i}\OperatorTok{]} \OperatorTok{*} \DecValTok{2}\OperatorTok{);}
\OperatorTok{\}}

\CommentTok{// Sama saja, tapi memakai notasi array}
\DataTypeTok{void}\NormalTok{ times3}\OperatorTok{(}\DataTypeTok{int}\NormalTok{ a}\OperatorTok{[],} \DataTypeTok{int}\NormalTok{ len}\OperatorTok{)}
\OperatorTok{\{}
    \ControlFlowTok{for} \OperatorTok{(}\DataTypeTok{int}\NormalTok{ i }\OperatorTok{=} \DecValTok{0}\OperatorTok{;}\NormalTok{ i }\OperatorTok{\textless{}}\NormalTok{ len}\OperatorTok{;}\NormalTok{ i}\OperatorTok{++)}
\NormalTok{        printf}\OperatorTok{(}\StringTok{"}\SpecialCharTok{\%d\textbackslash{}n}\StringTok{"}\OperatorTok{,}\NormalTok{ a}\OperatorTok{[}\NormalTok{i}\OperatorTok{]} \OperatorTok{*} \DecValTok{3}\OperatorTok{);}
\OperatorTok{\}}

\CommentTok{// Sama saja, tapi memakai notasi array dengan ukuran}
\DataTypeTok{void}\NormalTok{ times4}\OperatorTok{(}\DataTypeTok{int}\NormalTok{ a}\OperatorTok{[}\DecValTok{5}\OperatorTok{],} \DataTypeTok{int}\NormalTok{ len}\OperatorTok{)}
\OperatorTok{\{}
    \ControlFlowTok{for} \OperatorTok{(}\DataTypeTok{int}\NormalTok{ i }\OperatorTok{=} \DecValTok{0}\OperatorTok{;}\NormalTok{ i }\OperatorTok{\textless{}}\NormalTok{ len}\OperatorTok{;}\NormalTok{ i}\OperatorTok{++)}
\NormalTok{        printf}\OperatorTok{(}\StringTok{"}\SpecialCharTok{\%d\textbackslash{}n}\StringTok{"}\OperatorTok{,}\NormalTok{ a}\OperatorTok{[}\NormalTok{i}\OperatorTok{]} \OperatorTok{*} \DecValTok{4}\OperatorTok{);}
\OperatorTok{\}}

\DataTypeTok{int}\NormalTok{ main}\OperatorTok{(}\DataTypeTok{void}\OperatorTok{)}
\OperatorTok{\{}
    \DataTypeTok{int}\NormalTok{ x}\OperatorTok{[}\DecValTok{5}\OperatorTok{]} \OperatorTok{=} \OperatorTok{\{}\DecValTok{11}\OperatorTok{,} \DecValTok{22}\OperatorTok{,} \DecValTok{33}\OperatorTok{,} \DecValTok{44}\OperatorTok{,} \DecValTok{55}\OperatorTok{\};}

\NormalTok{    times2}\OperatorTok{(}\NormalTok{x}\OperatorTok{,} \DecValTok{5}\OperatorTok{);}
\NormalTok{    times3}\OperatorTok{(}\NormalTok{x}\OperatorTok{,} \DecValTok{5}\OperatorTok{);}
\NormalTok{    times4}\OperatorTok{(}\NormalTok{x}\OperatorTok{,} \DecValTok{5}\OperatorTok{);}
\OperatorTok{\}}
\end{Highlighting}
\end{Shaded}

\end{tcolorbox}

Semua cara menuliskan array sebagai parameter fungsi tersebut identik.

\begin{tcolorbox}[enhanced jigsaw,left=3pt,breakable,colback=white!97.2549!black,boxrule=0.2mm]

\begin{Shaded}
\begin{Highlighting}[]
\DataTypeTok{void}\NormalTok{ times2}\OperatorTok{(}\DataTypeTok{int} \OperatorTok{*}\NormalTok{a}\OperatorTok{,} \DataTypeTok{int}\NormalTok{ len}\OperatorTok{)}
\DataTypeTok{void}\NormalTok{ times3}\OperatorTok{(}\DataTypeTok{int}\NormalTok{ a}\OperatorTok{[],} \DataTypeTok{int}\NormalTok{ len}\OperatorTok{)}
\DataTypeTok{void}\NormalTok{ times4}\OperatorTok{(}\DataTypeTok{int}\NormalTok{ a}\OperatorTok{[}\DecValTok{5}\OperatorTok{],} \DataTypeTok{int}\NormalTok{ len}\OperatorTok{)}
\end{Highlighting}
\end{Shaded}

\end{tcolorbox}

Di kalangan pemakai C yang sudah berpengalaman, bentuk pertama jauh lebih umum.

Bahkan pada bentuk terakhir, compiler tidak peduli angka berapa yang kamu tulis (selain harus lebih besar dari nol\footnote{C11 §6.7.6.2¶1 mensyaratkan angkanya lebih besar dari nol. Tapi kamu mungkin menemukan kode yang mendeklarasikan array berukuran nol di akhir \hl{\mbox{\texttt{struct}}}, dan GCC cukup longgar soal ini kecuali kamu mengompilasi dengan \hl{\mbox{\texttt{-pedantic}}}. Array berukuran nol ini dulu adalah mekanisme akal-akalan untuk membuat struktur dengan ukuran variabel. Sayangnya, secara teknis mengakses array seperti itu adalah undefined behavior, meskipun di mana-mana praktis berfungsi. C99 menetapkan pengganti yang terdefinisi dengan baik, yang disebut \emph{flexible array member}, dan akan kita bahas nanti.}). Compiler sama sekali tidak memaksakan apa pun.

Setelah mengatakan itu, ukuran array di deklarasi fungsi sebenarnya \emph{penting} ketika kita mengoper array multidimensi ke fungsi, tapi nanti kita kembali ke sana. \index{Arrays!passing to functions|)}

\subsection{Mengubah Array di Dalam Fungsi}\label{mengubah-array-di-dalam-fungsi}

\index{Arrays!modifying within functions|(} Kita sudah bilang bahwa array hanyalah pointer yang menyamar. Artinya, kalau kamu mengoper array ke fungsi, kemungkinan besar yang kamu oper adalah pointer ke elemen pertama array.

Tapi kalau fungsi punya pointer ke data, fungsi itu bisa mengubah data tersebut! Jadi perubahan yang dilakukan fungsi pada array akan terlihat juga di caller.

Berikut contoh ketika kita mengoper pointer ke array ke sebuah fungsi, fungsi itu mengubah nilai-nilai di array, dan perubahannya terlihat di caller.

\begin{tcolorbox}[enhanced jigsaw,left=3pt,breakable,colback=white!97.2549!black,boxrule=0.2mm]

\begin{Shaded}
\begin{Highlighting}[numbers=left,,]
\PreprocessorTok{\#include }\ImportTok{\textless{}stdio.h\textgreater{}}

\DataTypeTok{void}\NormalTok{ double\_array}\OperatorTok{(}\DataTypeTok{int} \OperatorTok{*}\NormalTok{a}\OperatorTok{,} \DataTypeTok{int}\NormalTok{ len}\OperatorTok{)}
\OperatorTok{\{}
    \CommentTok{// Kalikan setiap elemen dengan 2}
    \CommentTok{//}
    \CommentTok{// Ini menggandakan nilai di x pada main() karena x dan a sama{-}sama menunjuk}
    \CommentTok{// ke array yang sama di memori!}

    \ControlFlowTok{for} \OperatorTok{(}\DataTypeTok{int}\NormalTok{ i }\OperatorTok{=} \DecValTok{0}\OperatorTok{;}\NormalTok{ i }\OperatorTok{\textless{}}\NormalTok{ len}\OperatorTok{;}\NormalTok{ i}\OperatorTok{++)}
\NormalTok{        a}\OperatorTok{[}\NormalTok{i}\OperatorTok{]} \OperatorTok{*=} \DecValTok{2}\OperatorTok{;}
\OperatorTok{\}}

\DataTypeTok{int}\NormalTok{ main}\OperatorTok{(}\DataTypeTok{void}\OperatorTok{)}
\OperatorTok{\{}
    \DataTypeTok{int}\NormalTok{ x}\OperatorTok{[}\DecValTok{5}\OperatorTok{]} \OperatorTok{=} \OperatorTok{\{}\DecValTok{1}\OperatorTok{,} \DecValTok{2}\OperatorTok{,} \DecValTok{3}\OperatorTok{,} \DecValTok{4}\OperatorTok{,} \DecValTok{5}\OperatorTok{\};}

\NormalTok{    double\_array}\OperatorTok{(}\NormalTok{x}\OperatorTok{,} \DecValTok{5}\OperatorTok{);}

    \ControlFlowTok{for} \OperatorTok{(}\DataTypeTok{int}\NormalTok{ i }\OperatorTok{=} \DecValTok{0}\OperatorTok{;}\NormalTok{ i }\OperatorTok{\textless{}} \DecValTok{5}\OperatorTok{;}\NormalTok{ i}\OperatorTok{++)}
\NormalTok{        printf}\OperatorTok{(}\StringTok{"}\SpecialCharTok{\%d\textbackslash{}n}\StringTok{"}\OperatorTok{,}\NormalTok{ x}\OperatorTok{[}\NormalTok{i}\OperatorTok{]);}  \CommentTok{// 2, 4, 6, 8, 10!}
\OperatorTok{\}}
\end{Highlighting}
\end{Shaded}

\end{tcolorbox}

Walaupun array dioper sebagai parameter \hl{\mbox{\texttt{a}}} yang bertipe \hl{\mbox{\texttt{int*}}}, lihat bagaimana kita mengaksesnya dengan notasi array \hl{\mbox{\texttt{a{[}i{]}}}}! Haaah. Ini sepenuhnya diperbolehkan.

Nanti ketika kita membahas kesetaraan antara array dan pointer, hal ini akan jauh lebih masuk akal. Untuk sekarang, cukup ketahui bahwa fungsi bisa mengubah array dan perubahannya terlihat di caller. \index{Arrays!modifying within functions|)}

\subsection{Mengoper Array Multidimensi ke Fungsi}\label{mengoper-array-multidimensi-ke-fungsi}

\index{Arrays!passing to functions|(} Ceritanya sedikit berubah ketika kita membicarakan array multidimensi. C perlu tahu semua dimensi (kecuali yang pertama) supaya punya cukup informasi untuk mengetahui di mana di memori nilai tertentu harus dicari.

Berikut contoh ketika kita eksplisit menyebutkan semua dimensi:

\begin{tcolorbox}[enhanced jigsaw,left=3pt,breakable,colback=white!97.2549!black,boxrule=0.2mm]

\begin{Shaded}
\begin{Highlighting}[numbers=left,,]
\PreprocessorTok{\#include }\ImportTok{\textless{}stdio.h\textgreater{}}

\DataTypeTok{void}\NormalTok{ print\_2D\_array}\OperatorTok{(}\DataTypeTok{int}\NormalTok{ a}\OperatorTok{[}\DecValTok{2}\OperatorTok{][}\DecValTok{3}\OperatorTok{])}
\OperatorTok{\{}
    \ControlFlowTok{for} \OperatorTok{(}\DataTypeTok{int}\NormalTok{ row }\OperatorTok{=} \DecValTok{0}\OperatorTok{;}\NormalTok{ row }\OperatorTok{\textless{}} \DecValTok{2}\OperatorTok{;}\NormalTok{ row}\OperatorTok{++)} \OperatorTok{\{}
        \ControlFlowTok{for} \OperatorTok{(}\DataTypeTok{int}\NormalTok{ col }\OperatorTok{=} \DecValTok{0}\OperatorTok{;}\NormalTok{ col }\OperatorTok{\textless{}} \DecValTok{3}\OperatorTok{;}\NormalTok{ col}\OperatorTok{++)}
\NormalTok{            printf}\OperatorTok{(}\StringTok{"}\SpecialCharTok{\%d}\StringTok{ "}\OperatorTok{,}\NormalTok{ a}\OperatorTok{[}\NormalTok{row}\OperatorTok{][}\NormalTok{col}\OperatorTok{]);}
\NormalTok{        printf}\OperatorTok{(}\StringTok{"}\SpecialCharTok{\textbackslash{}n}\StringTok{"}\OperatorTok{);}
    \OperatorTok{\}}
\OperatorTok{\}}

\DataTypeTok{int}\NormalTok{ main}\OperatorTok{(}\DataTypeTok{void}\OperatorTok{)}
\OperatorTok{\{}
    \DataTypeTok{int}\NormalTok{ x}\OperatorTok{[}\DecValTok{2}\OperatorTok{][}\DecValTok{3}\OperatorTok{]} \OperatorTok{=} \OperatorTok{\{}
        \OperatorTok{\{}\DecValTok{1}\OperatorTok{,} \DecValTok{2}\OperatorTok{,} \DecValTok{3}\OperatorTok{\},}
        \OperatorTok{\{}\DecValTok{4}\OperatorTok{,} \DecValTok{5}\OperatorTok{,} \DecValTok{6}\OperatorTok{\}}
    \OperatorTok{\};}

\NormalTok{    print\_2D\_array}\OperatorTok{(}\NormalTok{x}\OperatorTok{);}
\OperatorTok{\}}
\end{Highlighting}
\end{Shaded}

\end{tcolorbox}

Tapi dalam kasus ini, kedua bentuk\footnote{Ini juga setara: \hl{\mbox{\texttt{void\ print\_2D\_array(int\ (*a){[}3{]})}}}, tapi itu sudah melebihi yang ingin dibahas sekarang.} berikut setara:

\begin{tcolorbox}[enhanced jigsaw,left=3pt,breakable,colback=white!97.2549!black,boxrule=0.2mm]

\begin{Shaded}
\begin{Highlighting}[]
\DataTypeTok{void}\NormalTok{ print\_2D\_array}\OperatorTok{(}\DataTypeTok{int}\NormalTok{ a}\OperatorTok{[}\DecValTok{2}\OperatorTok{][}\DecValTok{3}\OperatorTok{])}
\DataTypeTok{void}\NormalTok{ print\_2D\_array}\OperatorTok{(}\DataTypeTok{int}\NormalTok{ a}\OperatorTok{[][}\DecValTok{3}\OperatorTok{])}
\end{Highlighting}
\end{Shaded}

\end{tcolorbox}

Compiler sebenarnya hanya butuh dimensi kedua supaya tahu seberapa jauh harus melompat di memori untuk setiap kenaikan dimensi pertama. Secara umum, compiler harus tahu semua dimensi kecuali yang pertama.

Ingat juga bahwa compiler hanya melakukan pengecekan batas saat kompilasi secara minimal (kalau beruntung), dan C sama sekali tidak mengecek batas saat runtime. Tanpa sabuk pengaman! Jangan sampai crash karena mengakses elemen array di luar batas! \index{Arrays!passing to functions|)} \index{Arrays|)}

\chapter{Strings}\label{strings}

\index{Strings|(}Akhirnya! Strings! Apa yang bisa lebih sederhana dari ini?

Nah, ternyata string di C sebenarnya bukan string. Betul! Mereka itu pointer! Tentu saja mereka pointer!

Seperti array, string di C \emph{nyaris tidak ada}.

Tapi kita lihat saja---sebenarnya tidak seseram itu.

\section{String Literal}\label{string-literal}

\index{String literals|(}Sebelum mulai, mari kita bahas dulu string literal di C. String literal berupa deretan karakter di dalam tanda kutip \emph{ganda} (\hl{\mbox{\texttt{"}}}). (Tanda kutip tunggal mengapit karakter, dan itu makhluk yang sama sekali berbeda.)

Contoh:

\begin{tcolorbox}[enhanced jigsaw,left=3pt,breakable,colback=white!97.2549!black,boxrule=0.2mm]

\begin{Shaded}
\begin{Highlighting}[]
\StringTok{"Hello, world!}\SpecialCharTok{\textbackslash{}n}\StringTok{"}
\StringTok{"This is a test."}
\StringTok{"When asked if this string had quotes in it, she replied, }\SpecialCharTok{\textbackslash{}"}\StringTok{It does.}\SpecialCharTok{\textbackslash{}"}\StringTok{"}
\end{Highlighting}
\end{Shaded}

\end{tcolorbox}

Yang pertama punya newline di ujungnya---pemandangan yang sangat umum.

Yang terakhir punya tanda kutip di dalamnya, dan kamu bisa lihat masing-masing didahului (kita bilang ``di-escape oleh'') sebuah backslash (\hl{\mbox{\texttt{\textbackslash{}}}}) yang menandakan bahwa di titik itu ada tanda kutip literal yang menjadi bagian string. Begitulah cara C compiler membedakan antara mencetak tanda kutip ganda dan tanda kutip ganda penutup string.\index{String literals|)}

\section{Variabel String}\label{variabel-string}

\index{String variables|(}Setelah tahu cara membuat string literal, sekarang kita tugaskan ke sebuah variabel supaya bisa kita pakai.

\begin{tcolorbox}[enhanced jigsaw,left=3pt,breakable,colback=white!97.2549!black,boxrule=0.2mm]

\begin{Shaded}
\begin{Highlighting}[]
\DataTypeTok{char} \OperatorTok{*}\NormalTok{s }\OperatorTok{=} \StringTok{"Hello, world!"}\OperatorTok{;}
\end{Highlighting}
\end{Shaded}

\end{tcolorbox}

Perhatikan tipenya: pointer ke \hl{\mbox{\texttt{char}}}. Variabel string \hl{\mbox{\texttt{s}}} sebenarnya adalah pointer ke karakter pertama string itu, yaitu \hl{\mbox{\texttt{H}}}.

Kita bisa mencetaknya dengan format specifier \hl{\mbox{\texttt{\%s}}} (untuk ``string''):

\begin{tcolorbox}[enhanced jigsaw,left=3pt,breakable,colback=white!97.2549!black,boxrule=0.2mm]

\begin{Shaded}
\begin{Highlighting}[]
\DataTypeTok{char} \OperatorTok{*}\NormalTok{s }\OperatorTok{=} \StringTok{"Hello, world!"}\OperatorTok{;}

\NormalTok{printf}\OperatorTok{(}\StringTok{"}\SpecialCharTok{\%s\textbackslash{}n}\StringTok{"}\OperatorTok{,}\NormalTok{ s}\OperatorTok{);}  \CommentTok{// "Hello, world!"}
\end{Highlighting}
\end{Shaded}

\end{tcolorbox}

\index{String variables|)}

\section{Variabel String sebagai Array}\label{variabel-string-sebagai-array}

\index{String variables!as arrays|(} Pilihan lainnya adalah ini, yang hampir setara dengan pemakaian \hl{\mbox{\texttt{char*}}} di atas:

\begin{tcolorbox}[enhanced jigsaw,left=3pt,breakable,colback=white!97.2549!black,boxrule=0.2mm]

\begin{Shaded}
\begin{Highlighting}[]
\DataTypeTok{char}\NormalTok{ s}\OperatorTok{[}\DecValTok{14}\OperatorTok{]} \OperatorTok{=} \StringTok{"Hello, world!"}\OperatorTok{;}

\CommentTok{// atau, kalau kita malas dengan benar dan menyuruh compiler}
\CommentTok{// menghitung panjangnya untuk kita:}

\DataTypeTok{char}\NormalTok{ s}\OperatorTok{[]} \OperatorTok{=} \StringTok{"Hello, world!"}\OperatorTok{;}
\end{Highlighting}
\end{Shaded}

\end{tcolorbox}

Artinya, kamu bisa memakai notasi array untuk mengakses karakter dalam string. Mari kita lakukan persis itu untuk mencetak semua karakter string dalam satu baris:

\begin{tcolorbox}[enhanced jigsaw,left=3pt,breakable,colback=white!97.2549!black,boxrule=0.2mm]

\begin{Shaded}
\begin{Highlighting}[numbers=left,,]
\PreprocessorTok{\#include }\ImportTok{\textless{}stdio.h\textgreater{}}

\DataTypeTok{int}\NormalTok{ main}\OperatorTok{(}\DataTypeTok{void}\OperatorTok{)}
\OperatorTok{\{}
    \DataTypeTok{char}\NormalTok{ s}\OperatorTok{[]} \OperatorTok{=} \StringTok{"Hello, world!"}\OperatorTok{;}

    \ControlFlowTok{for} \OperatorTok{(}\DataTypeTok{int}\NormalTok{ i }\OperatorTok{=} \DecValTok{0}\OperatorTok{;}\NormalTok{ i }\OperatorTok{\textless{}} \DecValTok{13}\OperatorTok{;}\NormalTok{ i}\OperatorTok{++)}
\NormalTok{        printf}\OperatorTok{(}\StringTok{"}\SpecialCharTok{\%c}\StringTok{"}\OperatorTok{,}\NormalTok{ s}\OperatorTok{[}\NormalTok{i}\OperatorTok{]);}
\NormalTok{    printf}\OperatorTok{(}\StringTok{"}\SpecialCharTok{\textbackslash{}n}\StringTok{"}\OperatorTok{);}
\OperatorTok{\}}
\end{Highlighting}
\end{Shaded}

\end{tcolorbox}

Perhatikan bahwa kita memakai format specifier \hl{\mbox{\texttt{\%c}}} untuk mencetak satu karakter.

Lihat juga ini. Programnya tetap berjalan dengan baik kalau definisi \hl{\mbox{\texttt{s}}} kita ubah menjadi tipe \hl{\mbox{\texttt{char*}}}:

\begin{tcolorbox}[enhanced jigsaw,left=3pt,breakable,colback=white!97.2549!black,boxrule=0.2mm]

\begin{Shaded}
\begin{Highlighting}[numbers=left,,]
\PreprocessorTok{\#include }\ImportTok{\textless{}stdio.h\textgreater{}}

\DataTypeTok{int}\NormalTok{ main}\OperatorTok{(}\DataTypeTok{void}\OperatorTok{)}
\OperatorTok{\{}
    \DataTypeTok{char} \OperatorTok{*}\NormalTok{s }\OperatorTok{=} \StringTok{"Hello, world!"}\OperatorTok{;}   \CommentTok{// char* di sini}

    \ControlFlowTok{for} \OperatorTok{(}\DataTypeTok{int}\NormalTok{ i }\OperatorTok{=} \DecValTok{0}\OperatorTok{;}\NormalTok{ i }\OperatorTok{\textless{}} \DecValTok{13}\OperatorTok{;}\NormalTok{ i}\OperatorTok{++)}
\NormalTok{        printf}\OperatorTok{(}\StringTok{"}\SpecialCharTok{\%c}\StringTok{"}\OperatorTok{,}\NormalTok{ s}\OperatorTok{[}\NormalTok{i}\OperatorTok{]);}    \CommentTok{// Tapi masih memakai array di sini...?}
\NormalTok{    printf}\OperatorTok{(}\StringTok{"}\SpecialCharTok{\textbackslash{}n}\StringTok{"}\OperatorTok{);}
\OperatorTok{\}}
\end{Highlighting}
\end{Shaded}

\end{tcolorbox}

Dan kita masih bisa memakai notasi array untuk mencetaknya! Ini mengejutkan, tapi itu hanya karena kita belum membahas kesetaraan array/pointer. Meski begitu, ini satu petunjuk lagi bahwa array dan pointer itu sama saja, jauh di dalamnya. \index{String variables!as arrays|)}

\section{Inisialisasi String}\label{inisialisasi-string}

\index{Strings!initializers|(} Kita sudah melihat beberapa contoh inisialisasi variabel string dengan string literal:

\begin{tcolorbox}[enhanced jigsaw,left=3pt,breakable,colback=white!97.2549!black,boxrule=0.2mm]

\begin{Shaded}
\begin{Highlighting}[]
\DataTypeTok{char} \OperatorTok{*}\NormalTok{s }\OperatorTok{=} \StringTok{"Hello, world!"}\OperatorTok{;}
\DataTypeTok{char}\NormalTok{ t}\OperatorTok{[]} \OperatorTok{=} \StringTok{"Hello, again!"}\OperatorTok{;}
\end{Highlighting}
\end{Shaded}

\end{tcolorbox}

Tapi kedua inisialisasi ini berbeda secara halus. Sebuah string literal, seperti integer literal, memorinya dikelola otomatis oleh compiler untukmu! Untuk integer, yaitu data berukuran tetap, compiler bisa mengelolanya dengan mudah. Tapi string adalah makhluk berukuran byte yang bervariasi, yang dijinakkan compiler dengan cara dilempar ke sepotong memori, lalu kamu diberi pointer ke sana.

Bentuk ini menunjuk ke mana pun string itu ditaruh. Biasanya tempat itu berada di negeri yang jauh dari memori program lainnya -- yaitu memori read-only -- karena alasan performa dan keamanan.

\begin{tcolorbox}[enhanced jigsaw,left=3pt,breakable,colback=white!97.2549!black,boxrule=0.2mm]

\begin{Shaded}
\begin{Highlighting}[]
\DataTypeTok{char} \OperatorTok{*}\NormalTok{s }\OperatorTok{=} \StringTok{"Hello, world!"}\OperatorTok{;}
\end{Highlighting}
\end{Shaded}

\end{tcolorbox}

Jadi, kalau kamu mencoba mengubah string itu dengan ini:

\begin{tcolorbox}[enhanced jigsaw,left=3pt,breakable,colback=white!97.2549!black,boxrule=0.2mm]

\begin{Shaded}
\begin{Highlighting}[]
\DataTypeTok{char} \OperatorTok{*}\NormalTok{s }\OperatorTok{=} \StringTok{"Hello, world!"}\OperatorTok{;}

\NormalTok{s}\OperatorTok{[}\DecValTok{0}\OperatorTok{]} \OperatorTok{=} \CharTok{\textquotesingle{}z\textquotesingle{}}\OperatorTok{;}  \CommentTok{// KABAR BURUK: mencoba mengubah string literal!}
\end{Highlighting}
\end{Shaded}

\end{tcolorbox}

Perilakunya undefined. Tergantung sistemmu, kemungkinan besar hasilnya crash.

Tapi mendeklarasikannya sebagai array itu berbeda. Compiler tidak menyembunyikan byte-byte itu di bagian kota yang lain; mereka ada di jalan seberang rumahmu. Yang satu ini \emph{salinan} string yang bisa diubah -- bisa kita ubah sesuka hati:

\begin{tcolorbox}[enhanced jigsaw,left=3pt,breakable,colback=white!97.2549!black,boxrule=0.2mm]

\begin{Shaded}
\begin{Highlighting}[]
\DataTypeTok{char}\NormalTok{ t}\OperatorTok{[]} \OperatorTok{=} \StringTok{"Hello, again!"}\OperatorTok{;}  \CommentTok{// t adalah salinan array dari string itu}
\NormalTok{t}\OperatorTok{[}\DecValTok{0}\OperatorTok{]} \OperatorTok{=} \CharTok{\textquotesingle{}z\textquotesingle{}}\OperatorTok{;} \CommentTok{//  Tidak masalah}

\NormalTok{printf}\OperatorTok{(}\StringTok{"}\SpecialCharTok{\%s\textbackslash{}n}\StringTok{"}\OperatorTok{,}\NormalTok{ t}\OperatorTok{);}  \CommentTok{// "zello, again!"}
\end{Highlighting}
\end{Shaded}

\end{tcolorbox}

Jadi ingat: kalau kamu punya pointer ke string literal, jangan coba mengubahnya! Dan kalau kamu memakai string di dalam tanda kutip ganda untuk menginisialisasi array, itu sebenarnya bukan string literal. \index{Strings!initializers|)}

\section{Mendapatkan Panjang String}\label{mendapatkan-panjang-string}

\index{Strings!getting the length|(} Tidak bisa, karena C tidak mencatatnya untukmu. Dan kalau dibilang ``tidak bisa'', maksudnya sebenarnya ``bisa''\footnote{Walaupun memang benar bahwa C tidak mencatat panjang string.}. Ada fungsi di \hl{\mbox{\texttt{\textless{}string.h\textgreater{}}}} bernama \hl{\mbox{\texttt{strlen()}}} yang bisa dipakai untuk menghitung panjang string apa pun dalam byte\footnote{Kalau kamu memakai character set dasar atau character set 8-bit, kamu terbiasa dengan satu karakter sama dengan satu byte. Tapi itu tidak berlaku di semua character encoding.}.

\begin{tcolorbox}[enhanced jigsaw,left=3pt,breakable,colback=white!97.2549!black,boxrule=0.2mm]

\begin{Shaded}
\begin{Highlighting}[numbers=left,,]
\PreprocessorTok{\#include }\ImportTok{\textless{}stdio.h\textgreater{}}
\PreprocessorTok{\#include }\ImportTok{\textless{}string.h\textgreater{}}

\DataTypeTok{int}\NormalTok{ main}\OperatorTok{(}\DataTypeTok{void}\OperatorTok{)}
\OperatorTok{\{}
    \DataTypeTok{char} \OperatorTok{*}\NormalTok{s }\OperatorTok{=} \StringTok{"Hello, world!"}\OperatorTok{;}

\NormalTok{    printf}\OperatorTok{(}\StringTok{"The string is }\SpecialCharTok{\%zu}\StringTok{ bytes long.}\SpecialCharTok{\textbackslash{}n}\StringTok{"}\OperatorTok{,}\NormalTok{ strlen}\OperatorTok{(}\NormalTok{s}\OperatorTok{));}
\OperatorTok{\}}
\end{Highlighting}
\end{Shaded}

\end{tcolorbox}

Fungsi \hl{\mbox{\texttt{strlen()}}} mengembalikan tipe \hl{\mbox{\texttt{size\_t}}}, yaitu tipe integer sehingga bisa dipakai untuk perhitungan integer. Kita mencetak \hl{\mbox{\texttt{size\_t}}} dengan \hl{\mbox{\texttt{\%zu}}}.

Program di atas mencetak:

\begin{tcolorbox}[enhanced jigsaw,left=3pt,breakable,colback=white!97.2549!black,boxrule=0.2mm]

\begin{Shaded}
\begin{Highlighting}[]
\NormalTok{The string is 13 bytes long.}
\end{Highlighting}
\end{Shaded}

\end{tcolorbox}

Bagus! Jadi ternyata \emph{bisa} mendapatkan panjang string! \index{Strings!getting the length|)}

Tapi\ldots{} kalau C tidak mencatat panjang string di mana pun, bagaimana ia tahu seberapa panjang string itu?

\section{Terminasi String}\label{terminasi-string}

\index{Strings!termination|(} C memperlakukan string sedikit berbeda dari banyak bahasa pemrograman, dan sebenarnya berbeda dari hampir semua bahasa pemrograman modern.

Saat membuat bahasa baru, pada dasarnya ada dua pilihan untuk menyimpan string di memori:

\begin{enumerate}
\def\labelenumi{\arabic{enumi}.}
\item
  Simpan byte-byte string beserta sebuah angka yang menyatakan panjang string.
\item
  Simpan byte-byte string, dan tandai akhir string dengan byte khusus yang disebut \emph{terminator}.
\end{enumerate}

Kalau kamu ingin string yang lebih panjang dari 255 karakter, pilihan 1 membutuhkan setidaknya dua byte untuk menyimpan panjangnya. Sedangkan pilihan 2 hanya butuh satu byte untuk mengakhiri string. Jadi ada sedikit penghematan di sana.

Tentu saja, zaman sekarang menghemat satu byte (atau 3---banyak bahasa dengan senang hati membolehkan string sepanjang 4 gigabyte) terasa konyol. Tapi dulu, itu urusan yang lebih besar.

Jadi C memilih pendekatan \#2. Di C, sebuah ``string'' didefinisikan oleh dua ciri dasar:

\begin{itemize}
\tightlist
\item
  Sebuah pointer ke karakter pertama string.
\item
  Sebuah byte bernilai nol (atau karakter \hl{\mbox{\texttt{NUL}}}\footnote{Ini berbeda dari pointer \hl{\mbox{\texttt{NULL}}}, dan penulis menyingkatnya \hl{\mbox{\texttt{NUL}}} kalau membicarakan karakternya dan \hl{\mbox{\texttt{NULL}}} untuk pointer-nya.}) di suatu tempat di memori setelah pointer itu, yang menandai akhir string.
\end{itemize}

Karakter \hl{\mbox{\texttt{NUL}}} bisa ditulis di kode C sebagai \hl{\mbox{\texttt{\textbackslash{}0}}}, meski jarang kamu perlu melakukannya.

Ketika kamu menulis string di dalam tanda kutip ganda di kodemu, karakter \hl{\mbox{\texttt{NUL}}} otomatis dan implisit ikut disertakan.

\begin{tcolorbox}[enhanced jigsaw,left=3pt,breakable,colback=white!97.2549!black,boxrule=0.2mm]

\begin{Shaded}
\begin{Highlighting}[]
\DataTypeTok{char} \OperatorTok{*}\NormalTok{s }\OperatorTok{=} \StringTok{"Hello!"}\OperatorTok{;}  \CommentTok{// Sebenarnya "Hello!\textbackslash{}0" di balik layar}
\end{Highlighting}
\end{Shaded}

\end{tcolorbox}

Dengan ini, mari kita tulis fungsi \hl{\mbox{\texttt{strlen()}}} kita sendiri yang menghitung \hl{\mbox{\texttt{char}}} dalam string sampai menemukan \hl{\mbox{\texttt{NUL}}}.

Caranya, telusuri string untuk mencari satu karakter \hl{\mbox{\texttt{NUL}}}, sambil menghitung\footnote{Nanti kita akan belajar cara yang lebih rapi dengan pointer arithmetic.}:

\begin{tcolorbox}[enhanced jigsaw,left=3pt,breakable,colback=white!97.2549!black,boxrule=0.2mm]

\begin{Shaded}
\begin{Highlighting}[]
\DataTypeTok{int}\NormalTok{ my\_strlen}\OperatorTok{(}\DataTypeTok{char} \OperatorTok{*}\NormalTok{s}\OperatorTok{)}
\OperatorTok{\{}
    \DataTypeTok{int}\NormalTok{ count }\OperatorTok{=} \DecValTok{0}\OperatorTok{;}

    \ControlFlowTok{while} \OperatorTok{(}\NormalTok{s}\OperatorTok{[}\NormalTok{count}\OperatorTok{]} \OperatorTok{!=} \CharTok{\textquotesingle{}}\SpecialCharTok{\textbackslash{}0}\CharTok{\textquotesingle{}}\OperatorTok{)}  \CommentTok{// Tanda kutip tunggal untuk satu karakter}
\NormalTok{        count}\OperatorTok{++;}

    \ControlFlowTok{return}\NormalTok{ count}\OperatorTok{;}
\OperatorTok{\}}
\end{Highlighting}
\end{Shaded}

\end{tcolorbox}

Dan pada dasarnya begitulah cara \hl{\mbox{\texttt{strlen()}}} bawaan menyelesaikan tugasnya. \index{Strings!termination|)}

\section{Menyalin String}\label{menyalin-string}

\index{Strings!copying|(} String tidak bisa disalin lewat operator penugasan (\hl{\mbox{\texttt{=}}}). Yang terjadi hanyalah menyalin pointer ke karakter pertama\ldots{} sehingga akhirnya kamu punya dua pointer ke string yang sama:

\begin{tcolorbox}[enhanced jigsaw,left=3pt,breakable,colback=white!97.2549!black,boxrule=0.2mm]

\begin{Shaded}
\begin{Highlighting}[numbers=left,,]
\PreprocessorTok{\#include }\ImportTok{\textless{}stdio.h\textgreater{}}

\DataTypeTok{int}\NormalTok{ main}\OperatorTok{(}\DataTypeTok{void}\OperatorTok{)}
\OperatorTok{\{}
    \DataTypeTok{char}\NormalTok{ s}\OperatorTok{[]} \OperatorTok{=} \StringTok{"Hello, world!"}\OperatorTok{;}
    \DataTypeTok{char} \OperatorTok{*}\NormalTok{t}\OperatorTok{;}

    \CommentTok{// Ini menyalin pointer{-}nya, bukan menyalin string{-}nya!}
\NormalTok{    t }\OperatorTok{=}\NormalTok{ s}\OperatorTok{;}

    \CommentTok{// Kita ubah t}
\NormalTok{    t}\OperatorTok{[}\DecValTok{0}\OperatorTok{]} \OperatorTok{=} \CharTok{\textquotesingle{}z\textquotesingle{}}\OperatorTok{;}

    \CommentTok{// Tapi mencetak s menampilkan perubahannya!}
    \CommentTok{// Karena t dan s menunjuk ke string yang sama!}

\NormalTok{    printf}\OperatorTok{(}\StringTok{"}\SpecialCharTok{\%s\textbackslash{}n}\StringTok{"}\OperatorTok{,}\NormalTok{ s}\OperatorTok{);}  \CommentTok{// "zello, world!"}
\OperatorTok{\}}
\end{Highlighting}
\end{Shaded}

\end{tcolorbox}

Kalau kamu ingin menyalin string, kamu harus menyalinnya byte demi byte---artinya kamu mengambil byte-byte string dari satu tempat di memori dan menduplikasinya ke tempat lain di memori. Ini dipermudah oleh fungsi \hl{\mbox{\texttt{strcpy()}}}.\footnote{Ada fungsi lain bernama \hl{\mbox{\texttt{strncpy()}}} yang membatasi jumlah byte yang disalin. Sebagian orang bilang kamu harus selalu memakai \hl{\mbox{\texttt{strncpy()}}} karena ada perlindungan dari buffer overrun. Sebagian lain bilang jangan pernah memakai \hl{\mbox{\texttt{strncpy()}}} karena ia belum tentu mengakhiri string-mu, footgun mengerikan lainnya. Kalau mau benar-benar aman, kamu bisa menulis versi \hl{\mbox{\texttt{strncpy()}}} sendiri yang selalu mengakhiri string.}

Sebelum menyalin string, pastikan ada ruang untuk menampung salinannya, yaitu array tujuan yang akan menampung karakter-karakternya harus setidaknya sepanjang string yang disalin.

\begin{tcolorbox}[enhanced jigsaw,left=3pt,breakable,colback=white!97.2549!black,boxrule=0.2mm]

\begin{Shaded}
\begin{Highlighting}[numbers=left,,]
\PreprocessorTok{\#include }\ImportTok{\textless{}stdio.h\textgreater{}}
\PreprocessorTok{\#include }\ImportTok{\textless{}string.h\textgreater{}}

\DataTypeTok{int}\NormalTok{ main}\OperatorTok{(}\DataTypeTok{void}\OperatorTok{)}
\OperatorTok{\{}
    \DataTypeTok{char}\NormalTok{ s}\OperatorTok{[]} \OperatorTok{=} \StringTok{"Hello, world!"}\OperatorTok{;}
    \DataTypeTok{char}\NormalTok{ t}\OperatorTok{[}\DecValTok{100}\OperatorTok{];}  \CommentTok{// Setiap char satu byte, jadi ruangnya lebih dari cukup}

    \CommentTok{// Ini menyalin string{-}nya!}
\NormalTok{    strcpy}\OperatorTok{(}\NormalTok{t}\OperatorTok{,}\NormalTok{ s}\OperatorTok{);}

    \CommentTok{// Kita ubah t}
\NormalTok{    t}\OperatorTok{[}\DecValTok{0}\OperatorTok{]} \OperatorTok{=} \CharTok{\textquotesingle{}z\textquotesingle{}}\OperatorTok{;}

    \CommentTok{// Dan s tidak terpengaruh karena itu string yang berbeda}
\NormalTok{    printf}\OperatorTok{(}\StringTok{"}\SpecialCharTok{\%s\textbackslash{}n}\StringTok{"}\OperatorTok{,}\NormalTok{ s}\OperatorTok{);}  \CommentTok{// "Hello, world!"}

    \CommentTok{// Tapi t sudah berubah}
\NormalTok{    printf}\OperatorTok{(}\StringTok{"}\SpecialCharTok{\%s\textbackslash{}n}\StringTok{"}\OperatorTok{,}\NormalTok{ t}\OperatorTok{);}  \CommentTok{// "zello, world!"}
\OperatorTok{\}}
\end{Highlighting}
\end{Shaded}

\end{tcolorbox}

Perhatikan bahwa pada \hl{\mbox{\texttt{strcpy()}}}, pointer tujuan adalah argumen pertama, dan pointer sumber adalah argumen kedua. Trik untuk mengingatnya: ini urutan yang akan kamu pakai untuk menulis \hl{\mbox{\texttt{t}}} dan \hl{\mbox{\texttt{s}}} seandainya penugasan \hl{\mbox{\texttt{=}}} bekerja untuk string, dengan sumber di kanan dan tujuan di kiri. \index{Strings!copying|)}\index{Strings|)}

\chapter{Structs}\label{structs}

\index{struct keyword@\texttt{struct} keyword|(} Di C ada yang disebut \hl{\mbox{\texttt{struct}}}, yaitu tipe yang bisa didefinisikan sendiri oleh pengguna dan menyimpan beberapa data sekaligus, yang bisa saja bertipe berbeda-beda.

Ini cara praktis untuk membungkus beberapa variabel menjadi satu. Dengan begitu, mengoper variabel ke fungsi jadi lebih ringkas (cukup satu yang dioper, bukan banyak), data lebih rapi, dan kode lebih mudah dibaca.

Kalau kamu berasal dari bahasa lain, mungkin kamu sudah akrab dengan konsep \emph{class} dan \emph{object}. Keduanya tidak ada di C secara bawaan\footnote{Walaupun di C setiap item di memori seperti \hl{\mbox{\texttt{int}}} disebut ``object'', itu bukan object dalam arti pemrograman berorientasi objek.}. Kamu bisa membayangkan \hl{\mbox{\texttt{struct}}} sebagai class yang hanya punya data member, tanpa method.

\section{Mendeklarasikan Struct}\label{mendeklarasikan-struct}

\index{struct keyword@\texttt{struct} keyword!declaring|(} Sebuah \hl{\mbox{\texttt{struct}}} bisa dideklarasikan di kodemu seperti ini:

\begin{tcolorbox}[enhanced jigsaw,left=3pt,breakable,colback=white!97.2549!black,boxrule=0.2mm]

\begin{Shaded}
\begin{Highlighting}[]
\KeywordTok{struct}\NormalTok{ car }\OperatorTok{\{}
    \DataTypeTok{char} \OperatorTok{*}\NormalTok{name}\OperatorTok{;}
    \DataTypeTok{float}\NormalTok{ price}\OperatorTok{;}
    \DataTypeTok{int}\NormalTok{ speed}\OperatorTok{;}
\OperatorTok{\};}
\end{Highlighting}
\end{Shaded}

\end{tcolorbox}

Ini sering dilakukan di scope global, di luar semua fungsi, supaya \hl{\mbox{\texttt{struct}}} itu tersedia secara global.

Dengan begitu, kamu membuat sebuah \emph{tipe} baru. Nama lengkap tipenya adalah \hl{\mbox{\texttt{struct\ car}}}. (Bukan sekadar \hl{\mbox{\texttt{car}}}---itu tidak akan berfungsi.)

Belum ada variabel bertipe itu, tapi kita bisa mendeklarasikannya:

\begin{tcolorbox}[enhanced jigsaw,left=3pt,breakable,colback=white!97.2549!black,boxrule=0.2mm]

\begin{Shaded}
\begin{Highlighting}[]
\KeywordTok{struct}\NormalTok{ car saturn}\OperatorTok{;}  \CommentTok{// Variabel "saturn" bertipe "struct car"}
\end{Highlighting}
\end{Shaded}

\end{tcolorbox}

Sekarang kita punya variabel \hl{\mbox{\texttt{saturn}}}\footnote{Saturn dulu adalah merek mobil ekonomi yang populer di Amerika Serikat, sampai akhirnya gulung tikar akibat krisis 2008, sayangnya bagi kami para penggemarnya.} yang belum diinisialisasi, bertipe \hl{\mbox{\texttt{struct\ car}}}.

Sebaiknya kita inisialisasi! Tapi bagaimana cara mengisi nilai masing-masing field-nya?

Seperti di banyak bahasa lain yang mencurinya dari C, kita memakai operator titik (\hl{\mbox{\texttt{.}}}) untuk mengakses tiap field.

\begin{tcolorbox}[enhanced jigsaw,left=3pt,breakable,colback=white!97.2549!black,boxrule=0.2mm]

\begin{Shaded}
\begin{Highlighting}[]
\NormalTok{saturn}\OperatorTok{.}\NormalTok{name }\OperatorTok{=} \StringTok{"Saturn SL/2"}\OperatorTok{;}
\NormalTok{saturn}\OperatorTok{.}\NormalTok{price }\OperatorTok{=} \FloatTok{15999.99}\OperatorTok{;}
\NormalTok{saturn}\OperatorTok{.}\NormalTok{speed }\OperatorTok{=} \DecValTok{175}\OperatorTok{;}

\NormalTok{printf}\OperatorTok{(}\StringTok{"Name:           }\SpecialCharTok{\%s\textbackslash{}n}\StringTok{"}\OperatorTok{,}\NormalTok{ saturn}\OperatorTok{.}\NormalTok{name}\OperatorTok{);}
\NormalTok{printf}\OperatorTok{(}\StringTok{"Price (USD):    }\SpecialCharTok{\%f\textbackslash{}n}\StringTok{"}\OperatorTok{,}\NormalTok{ saturn}\OperatorTok{.}\NormalTok{price}\OperatorTok{);}
\NormalTok{printf}\OperatorTok{(}\StringTok{"Top Speed (km): }\SpecialCharTok{\%d\textbackslash{}n}\StringTok{"}\OperatorTok{,}\NormalTok{ saturn}\OperatorTok{.}\NormalTok{speed}\OperatorTok{);}
\end{Highlighting}
\end{Shaded}

\end{tcolorbox}

Di baris-baris awal, kita mengisi nilai-nilai di \hl{\mbox{\texttt{struct\ car}}}, lalu di bagian berikutnya kita mencetak nilai-nilai itu. \index{struct keyword@\texttt{struct} keyword!declaring|)}

\section{Struct Initializer}\label{struct-initializers}

\index{struct keyword@\texttt{struct} keyword!initializers|(} Contoh di bagian sebelumnya agak merepotkan. Pasti ada cara yang lebih baik untuk menginisialisasi variabel \hl{\mbox{\texttt{struct}}} itu!

Kamu bisa memakai initializer dengan mengisi nilai untuk field-field-nya \emph{sesuai urutan kemunculannya di \hl{\mbox{\texttt{struct}}}} saat mendefinisikan variabel. (Ini tidak bisa dilakukan setelah variabelnya didefinisikan---harus terjadi saat definisi.)

\begin{tcolorbox}[enhanced jigsaw,left=3pt,breakable,colback=white!97.2549!black,boxrule=0.2mm]

\begin{Shaded}
\begin{Highlighting}[]
\KeywordTok{struct}\NormalTok{ car }\OperatorTok{\{}
    \DataTypeTok{char} \OperatorTok{*}\NormalTok{name}\OperatorTok{;}
    \DataTypeTok{float}\NormalTok{ price}\OperatorTok{;}
    \DataTypeTok{int}\NormalTok{ speed}\OperatorTok{;}
\OperatorTok{\};}

\CommentTok{// Sekarang dengan initializer! Urutan field sama seperti di deklarasi struct:}
\KeywordTok{struct}\NormalTok{ car saturn }\OperatorTok{=} \OperatorTok{\{}\StringTok{"Saturn SL/2"}\OperatorTok{,} \FloatTok{16000.99}\OperatorTok{,} \DecValTok{175}\OperatorTok{\};}

\NormalTok{printf}\OperatorTok{(}\StringTok{"Name:      }\SpecialCharTok{\%s\textbackslash{}n}\StringTok{"}\OperatorTok{,}\NormalTok{ saturn}\OperatorTok{.}\NormalTok{name}\OperatorTok{);}
\NormalTok{printf}\OperatorTok{(}\StringTok{"Price:     }\SpecialCharTok{\%f\textbackslash{}n}\StringTok{"}\OperatorTok{,}\NormalTok{ saturn}\OperatorTok{.}\NormalTok{price}\OperatorTok{);}
\NormalTok{printf}\OperatorTok{(}\StringTok{"Top Speed: }\SpecialCharTok{\%d}\StringTok{ km}\SpecialCharTok{\textbackslash{}n}\StringTok{"}\OperatorTok{,}\NormalTok{ saturn}\OperatorTok{.}\NormalTok{speed}\OperatorTok{);}
\end{Highlighting}
\end{Shaded}

\end{tcolorbox}

Fakta bahwa field di initializer harus berurutan sama itu agak menyeramkan. Kalau ada yang mengubah urutan di \hl{\mbox{\texttt{struct\ car}}}, semua kode lain bisa rusak!

Kita bisa lebih spesifik dengan initializer:

\begin{tcolorbox}[enhanced jigsaw,left=3pt,breakable,colback=white!97.2549!black,boxrule=0.2mm]

\begin{Shaded}
\begin{Highlighting}[]
\KeywordTok{struct}\NormalTok{ car saturn }\OperatorTok{=} \OperatorTok{\{.}\NormalTok{speed}\OperatorTok{=}\DecValTok{175}\OperatorTok{,} \OperatorTok{.}\NormalTok{name}\OperatorTok{=}\StringTok{"Saturn SL/2"}\OperatorTok{\};}
\end{Highlighting}
\end{Shaded}

\end{tcolorbox}

Sekarang initializer tidak bergantung pada urutan di deklarasi \hl{\mbox{\texttt{struct}}}. Kode seperti ini jelas lebih aman.

Seperti pada array initializer, field designator yang tidak disebutkan akan diinisialisasi dengan nol (di sini yaitu \hl{\mbox{\texttt{.price}}}, yang sengaja dihilangkan). \index{struct keyword@\texttt{struct} keyword!initializers|)}

\section{Mengoper Struct ke Fungsi}\label{mengoper-struct-ke-fungsi}

\index{struct keyword@\texttt{struct} keyword!passing and returning|(} Ada dua cara untuk mengoper \hl{\mbox{\texttt{struct}}} ke fungsi.

\begin{enumerate}
\def\labelenumi{\arabic{enumi}.}
\tightlist
\item
  Oper \hl{\mbox{\texttt{struct}}}-nya.
\item
  Oper pointer ke \hl{\mbox{\texttt{struct}}}-nya.
\end{enumerate}

Ingat bahwa ketika sesuatu dioper ke fungsi, sebuah \emph{salinan} dari sesuatu itu dibuat untuk dikerjakan oleh fungsi, baik itu salinan pointer, \hl{\mbox{\texttt{int}}}, \hl{\mbox{\texttt{struct}}}, atau apa pun.

Pada dasarnya ada dua keadaan ketika kamu ingin mengoper pointer ke \hl{\mbox{\texttt{struct}}}:

\begin{enumerate}
\def\labelenumi{\arabic{enumi}.}
\tightlist
\item
  Fungsi perlu bisa mengubah \hl{\mbox{\texttt{struct}}} yang dioper, dan perubahannya harus terlihat di caller.
\item
  \hl{\mbox{\texttt{struct}}}-nya cukup besar sehingga menyalinnya ke stack lebih mahal daripada menyalin sebuah pointer\footnote{Pointer kemungkinan berukuran 8 byte di sistem 64-bit.}.
\end{enumerate}

Karena dua alasan itu, jauh lebih umum mengoper pointer ke \hl{\mbox{\texttt{struct}}} ke fungsi, walaupun mengoper \hl{\mbox{\texttt{struct}}}-nya langsung sama sekali tidak dilarang.

Mari kita coba mengoper pointer, dengan membuat fungsi yang memungkinkan kita mengisi field \hl{\mbox{\texttt{.price}}} dari \hl{\mbox{\texttt{struct\ car}}}:

\begin{tcolorbox}[enhanced jigsaw,left=3pt,breakable,colback=white!97.2549!black,boxrule=0.2mm]

\begin{Shaded}
\begin{Highlighting}[numbers=left,,]
\PreprocessorTok{\#include }\ImportTok{\textless{}stdio.h\textgreater{}}

\KeywordTok{struct}\NormalTok{ car }\OperatorTok{\{}
    \DataTypeTok{char} \OperatorTok{*}\NormalTok{name}\OperatorTok{;}
    \DataTypeTok{float}\NormalTok{ price}\OperatorTok{;}
    \DataTypeTok{int}\NormalTok{ speed}\OperatorTok{;}
\OperatorTok{\};}

\DataTypeTok{int}\NormalTok{ main}\OperatorTok{(}\DataTypeTok{void}\OperatorTok{)}
\OperatorTok{\{}
    \KeywordTok{struct}\NormalTok{ car saturn }\OperatorTok{=} \OperatorTok{\{.}\NormalTok{speed}\OperatorTok{=}\DecValTok{175}\OperatorTok{,} \OperatorTok{.}\NormalTok{name}\OperatorTok{=}\StringTok{"Saturn SL/2"}\OperatorTok{\};}

    \CommentTok{// Oper pointer ke struct car ini, beserta harga baru yang}
    \CommentTok{// lebih realistis:}
\NormalTok{    set\_price}\OperatorTok{(\&}\NormalTok{saturn}\OperatorTok{,} \FloatTok{799.99}\OperatorTok{);}

\NormalTok{    printf}\OperatorTok{(}\StringTok{"Price: }\SpecialCharTok{\%f\textbackslash{}n}\StringTok{"}\OperatorTok{,}\NormalTok{ saturn}\OperatorTok{.}\NormalTok{price}\OperatorTok{);}
\OperatorTok{\}}
\end{Highlighting}
\end{Shaded}

\end{tcolorbox}

Kamu seharusnya bisa menyusun signature \hl{\mbox{\texttt{set\_price()}}} hanya dengan melihat tipe argumen yang kita punya di sana.

\hl{\mbox{\texttt{saturn}}} adalah \hl{\mbox{\texttt{struct\ car}}}, jadi \hl{\mbox{\texttt{\&saturn}}} pasti alamat dari \hl{\mbox{\texttt{struct\ car}}}, alias pointer ke \hl{\mbox{\texttt{struct\ car}}}, yaitu \hl{\mbox{\texttt{struct\ car*}}}.

Dan \hl{\mbox{\texttt{799.99}}} adalah \hl{\mbox{\texttt{float}}}.

Jadi deklarasi fungsinya pasti seperti ini:

\begin{tcolorbox}[enhanced jigsaw,left=3pt,breakable,colback=white!97.2549!black,boxrule=0.2mm]

\begin{Shaded}
\begin{Highlighting}[]
\DataTypeTok{void}\NormalTok{ set\_price}\OperatorTok{(}\KeywordTok{struct}\NormalTok{ car }\OperatorTok{*}\NormalTok{c}\OperatorTok{,} \DataTypeTok{float}\NormalTok{ new\_price}\OperatorTok{)}
\end{Highlighting}
\end{Shaded}

\end{tcolorbox}

Tinggal menulis badannya. Satu percobaan mungkin begini:

\begin{tcolorbox}[enhanced jigsaw,left=3pt,breakable,colback=white!97.2549!black,boxrule=0.2mm]

\begin{Shaded}
\begin{Highlighting}[]
\DataTypeTok{void}\NormalTok{ set\_price}\OperatorTok{(}\KeywordTok{struct}\NormalTok{ car }\OperatorTok{*}\NormalTok{c}\OperatorTok{,} \DataTypeTok{float}\NormalTok{ new\_price}\OperatorTok{)} \OperatorTok{\{}
\NormalTok{    c}\OperatorTok{.}\NormalTok{price }\OperatorTok{=}\NormalTok{ new\_price}\OperatorTok{;}  \CommentTok{// ERROR!!}
\OperatorTok{\}}
\end{Highlighting}
\end{Shaded}

\end{tcolorbox}

Itu tidak akan berhasil karena operator titik hanya bekerja pada \hl{\mbox{\texttt{struct}}}\ldots{} bukan pada \emph{pointer} ke \hl{\mbox{\texttt{struct}}}.

Oke, jadi kita bisa men-dereference variabel \hl{\mbox{\texttt{c}}} untuk ``de-pointer''-nya agar sampai ke \hl{\mbox{\texttt{struct}}}-nya. Men-dereference \hl{\mbox{\texttt{struct\ car*}}} menghasilkan \hl{\mbox{\texttt{struct\ car}}} yang ditunjuk pointer itu, yang seharusnya bisa kita beri operator titik:

\begin{tcolorbox}[enhanced jigsaw,left=3pt,breakable,colback=white!97.2549!black,boxrule=0.2mm]

\begin{Shaded}
\begin{Highlighting}[]
\DataTypeTok{void}\NormalTok{ set\_price}\OperatorTok{(}\KeywordTok{struct}\NormalTok{ car }\OperatorTok{*}\NormalTok{c}\OperatorTok{,} \DataTypeTok{float}\NormalTok{ new\_price}\OperatorTok{)} \OperatorTok{\{}
    \OperatorTok{(*}\NormalTok{c}\OperatorTok{).}\NormalTok{price }\OperatorTok{=}\NormalTok{ new\_price}\OperatorTok{;}  \CommentTok{// Berfungsi, tapi jelek dan tidak idiomatis :(}
\OperatorTok{\}}
\end{Highlighting}
\end{Shaded}

\end{tcolorbox}

Dan itu berhasil! Tapi mengetik semua tanda kurung dan asterisk itu agak merepotkan. C punya syntactic sugar bernama \emph{arrow operator} yang membantu soal ini. \index{struct keyword@\texttt{struct} keyword!passing and returning|)}

\section{Arrow Operator}\label{arrow-operator}

\index{-> arrow operator@\texttt{->} arrow operator|(} Arrow operator membantu merujuk ke field pada pointer ke \hl{\mbox{\texttt{struct}}}.

\begin{tcolorbox}[enhanced jigsaw,left=3pt,breakable,colback=white!97.2549!black,boxrule=0.2mm]

\begin{Shaded}
\begin{Highlighting}[]
\DataTypeTok{void}\NormalTok{ set\_price}\OperatorTok{(}\KeywordTok{struct}\NormalTok{ car }\OperatorTok{*}\NormalTok{c}\OperatorTok{,} \DataTypeTok{float}\NormalTok{ new\_price}\OperatorTok{)} \OperatorTok{\{}
    \CommentTok{// (*c).price = new\_price;  // Berfungsi, tapi tidak idiomatis :(}
    \CommentTok{//}
    \CommentTok{// Baris di atas 100\% setara dengan baris di bawah:}

\NormalTok{    c}\OperatorTok{{-}\textgreater{}}\NormalTok{price }\OperatorTok{=}\NormalTok{ new\_price}\OperatorTok{;}  \CommentTok{// Nah, yang ini!}
\OperatorTok{\}}
\end{Highlighting}
\end{Shaded}

\end{tcolorbox}

Jadi saat mengakses field, kapan memakai titik dan kapan memakai panah?

\begin{itemize}
\tightlist
\item
  Kalau kamu punya \hl{\mbox{\texttt{struct}}}, pakai titik (\hl{\mbox{\texttt{.}}}).
\item
  Kalau kamu punya pointer ke \hl{\mbox{\texttt{struct}}}, pakai panah (\hl{\mbox{\texttt{-\textgreater{}}}}). \index{-> arrow operator@\texttt{->} arrow operator|)}
\end{itemize}

\section{\texorpdfstring{Menyalin dan Mengembalikan \hl{\mbox{\texttt{struct}}}}{Menyalin dan Mengembalikan struct}}\label{menyalin-dan-mengembalikan-struct}

\index{struct keyword@\texttt{struct} keyword!copying|(} Yang ini gampang!

Tinggal tugaskan satu ke yang lain!

\begin{tcolorbox}[enhanced jigsaw,left=3pt,breakable,colback=white!97.2549!black,boxrule=0.2mm]

\begin{Shaded}
\begin{Highlighting}[]
\KeywordTok{struct}\NormalTok{ car a}\OperatorTok{,}\NormalTok{ b}\OperatorTok{;}

\NormalTok{b }\OperatorTok{=}\NormalTok{ a}\OperatorTok{;}  \CommentTok{// Salin struct{-}nya}
\end{Highlighting}
\end{Shaded}

\end{tcolorbox}

Mengembalikan \hl{\mbox{\texttt{struct}}} (bukan pointer ke \hl{\mbox{\texttt{struct}}}) dari sebuah fungsi juga membuat salinan serupa ke variabel penerimanya.

Ini bukan ``deep copy''\footnote{\emph{Deep copy} mengikuti pointer di dalam \hl{\mbox{\texttt{struct}}} dan ikut menyalin data yang ditunjuknya. \emph{Shallow copy} hanya menyalin pointer-nya, bukan hal yang ditunjuk. C tidak punya fungsionalitas deep copy bawaan.}. Semua field disalin apa adanya, termasuk pointer ke sesuatu. \index{struct keyword@\texttt{struct} keyword!copying|)}

\section{\texorpdfstring{Membandingkan \hl{\mbox{\texttt{struct}}}}{Membandingkan struct}}\label{membandingkan-struct}

\index{struct keyword@\texttt{struct} keyword!comparing|(} Hanya ada satu cara yang aman: bandingkan setiap field satu per satu.

Mungkin kamu mengira bisa memakai \href{https://beej.us/guide/bgclr/html/split/stringref.html\#man-strcmp}{\hl{\mbox{\texttt{memcmp()}}}}\footnote{https://beej.us/guide/bgclr/html/split/stringref.html\#man-strcmp}, tapi itu tidak menangani kemungkinan adanya \hyperref[struct-padding-bytes]{padding byte} di dalamnya.

Kalau \hl{\mbox{\texttt{struct}}} dikosongkan dulu dengan nol memakai \href{https://beej.us/guide/bgclr/html/split/stringref.html\#man-memset}{\hl{\mbox{\texttt{memset()}}}}\footnote{https://beej.us/guide/bgclr/html/split/stringref.html\#man-memset}, hasilnya \emph{mungkin} berhasil, tapi masih ada elemen aneh yang \href{https://stackoverflow.com/questions/141720/how-do-you-compare-structs-for-equality-in-c}{mungkin tidak dibandingkan seperti yang kamu harapkan}\footnote{https://stackoverflow.com/questions/141720/how-do-you-compare-structs-for-equality-in-c}. \index{struct keyword@\texttt{struct} keyword!comparing|)} \index{struct keyword@\texttt{struct} keyword|)}

\chapter{Input/Output File}\label{inputoutput-file}

\index{File I/O|(} Kita sudah melihat beberapa contoh I/O dengan \hl{\mbox{\texttt{printf()}}} untuk melakukan I/O di konsol.

Tapi di bab ini kita akan mendorong konsep-konsep itu sedikit lebih jauh.

\section{\texorpdfstring{Tipe Data \hl{\mbox{\texttt{FILE*}}}}{Tipe Data FILE*}}\label{tipe-data-file}

\index{FILE* type@\texttt{FILE*} type|(} Setiap kali kita melakukan I/O apa pun di C, kita melakukannya lewat sepotong data yang kita dapat dalam bentuk tipe \hl{\mbox{\texttt{FILE*}}}. \hl{\mbox{\texttt{FILE*}}} ini menyimpan semua informasi yang dibutuhkan untuk berkomunikasi dengan subsistem I/O tentang file mana yang kamu buka, di mana posisimu dalam file, dan sebagainya.

Spesifikasi menyebut semua ini \emph{stream}, yaitu aliran data dari sebuah file atau dari sumber apa pun. Penulis akan memakai ``file'' dan ``stream'' secara bergantian, tapi sebenarnya kamu sebaiknya menganggap ``file'' sebagai kasus khusus dari ``stream''. Ada cara lain untuk mengalirkan data ke dalam program selain membaca dari file.

Sebentar lagi kita akan melihat cara mendapatkan \hl{\mbox{\texttt{FILE*}}} yang terbuka dari sebuah nama file, tapi pertama-tama perlu disebutkan tiga stream yang sudah terbuka dan siap dipakai.

\index{stdin standard input@\texttt{stdin} standard input|(} \index{stdout standard output@\texttt{stdout} standard output|(} \index{stderr standard error@\texttt{stderr} standard error|(}

{\def\LTcaptype{none} % do not increment counter
\begin{longtable}[]{@{}ll@{}}
\toprule\noalign{}
Nama \hl{\mbox{\texttt{FILE*}}} & Deskripsi \\
\midrule\noalign{}
\endhead
\bottomrule\noalign{}
\endlastfoot
\hl{\mbox{\texttt{stdin}}} & Standard Input, umumnya keyboard secara bawaan \\
\hl{\mbox{\texttt{stdout}}} & Standard Output, umumnya layar secara bawaan \\
\hl{\mbox{\texttt{stderr}}} & Standard Error, umumnya layar juga secara bawaan \\
\end{longtable}
}

Ternyata kita sebenarnya sudah memakai ketiganya secara implisit. Misalnya, dua pemanggilan ini sama:

\begin{tcolorbox}[enhanced jigsaw,left=3pt,breakable,colback=white!97.2549!black,boxrule=0.2mm]

\begin{Shaded}
\begin{Highlighting}[]
\NormalTok{printf}\OperatorTok{(}\StringTok{"Hello, world!}\SpecialCharTok{\textbackslash{}n}\StringTok{"}\OperatorTok{);}
\NormalTok{fprintf}\OperatorTok{(}\NormalTok{stdout}\OperatorTok{,} \StringTok{"Hello, world!}\SpecialCharTok{\textbackslash{}n}\StringTok{"}\OperatorTok{);}  \CommentTok{// printf ke sebuah file}
\end{Highlighting}
\end{Shaded}

\end{tcolorbox}

Tapi soal itu dibahas nanti.

Kamu juga akan memperhatikan bahwa \hl{\mbox{\texttt{stdout}}} dan \hl{\mbox{\texttt{stderr}}} sama-sama menuju layar. Sekilas ini tampak seperti kelalaian atau redundansi, tapi sebenarnya bukan. Sistem operasi pada umumnya memungkinkan kamu \emph{mengalihkan} output salah satunya ke file yang berbeda, dan memisahkan pesan error dari output biasa yang bukan error bisa sangat praktis.

Misalnya, di shell POSIX (seperti sh, ksh, bash, zsh, dan sebagainya) pada sistem mirip Unix, kita bisa menjalankan program dan mengirim hanya output non-error (\hl{\mbox{\texttt{stdout}}}) ke satu file, dan semua output error (\hl{\mbox{\texttt{stderr}}}) ke file lain.

\begin{tcolorbox}[enhanced jigsaw,left=3pt,breakable,colback=white!97.2549!black,boxrule=0.2mm]

\begin{Shaded}
\begin{Highlighting}[]
\ExtensionTok{./foo} \OperatorTok{\textgreater{}}\NormalTok{ output.txt }\DecValTok{2}\OperatorTok{\textgreater{}}\NormalTok{ errors.txt   }\CommentTok{\# Perintah ini khusus Unix}
\end{Highlighting}
\end{Shaded}

\end{tcolorbox}

Karena alasan ini, kamu sebaiknya mengirim pesan error yang serius ke \hl{\mbox{\texttt{stderr}}}, bukan \hl{\mbox{\texttt{stdout}}}. \index{stdin standard input@\texttt{stdin} standard input|)} \index{stdout standard output@\texttt{stdout} standard output|)} \index{stderr standard error@\texttt{stderr} standard error|)}

Cara melakukannya dibahas lebih lanjut nanti. \index{FILE* type@\texttt{FILE*} type|(}

\section{Membaca File Teks}\label{membaca-file-teks}

\index{File I/O!text files, reading|(} Stream sebagian besar dikategorikan dalam dua cara: \emph{teks} dan \emph{biner}.

Stream teks boleh melakukan terjemahan data yang signifikan, terutama penerjemahan newline ke berbagai representasinya\footnote{Dulu ada tiga jenis newline yang berlaku luas: Carriage Return (CR, dipakai Mac lama), Linefeed (LF, dipakai sistem Unix), dan Carriage Return/Linefeed (CRLF, dipakai sistem Windows). Untungnya kehadiran OS X, yang berbasis Unix, mengurangi jumlah ini menjadi dua.}. File teks secara logis adalah deretan \emph{baris} yang dipisahkan newline. Agar portabel, data inputmu sebaiknya selalu diakhiri newline.

Tapi aturan umumnya: kalau kamu bisa menyunting file itu di text editor biasa, ia file teks. Kalau tidak, ia biner. Soal biner dibahas nanti.

Jadi mari kita bekerja---bagaimana cara membuka file untuk dibaca, dan mengambil data darinya?

Kita buat file bernama \hl{\mbox{\texttt{hello.txt}}} yang isinya hanya ini:

\begin{tcolorbox}[enhanced jigsaw,left=3pt,breakable,colback=white!97.2549!black,boxrule=0.2mm]

\begin{Shaded}
\begin{Highlighting}[]
\NormalTok{Hello, world!}
\end{Highlighting}
\end{Shaded}

\end{tcolorbox}

Lalu kita tulis program untuk membuka file itu, membaca satu karakter darinya, dan menutup file saat selesai. Itulah rencana permainannya!

\index{fopen() function@\texttt{fopen()} function|(} \index{fgetc() function@\texttt{fgetc()} function|(} \index{fclose() function@\texttt{fclose()} function|(}

\begin{tcolorbox}[enhanced jigsaw,left=3pt,breakable,colback=white!97.2549!black,boxrule=0.2mm]

\begin{Shaded}
\begin{Highlighting}[numbers=left,,]
\PreprocessorTok{\#include }\ImportTok{\textless{}stdio.h\textgreater{}}

\DataTypeTok{int}\NormalTok{ main}\OperatorTok{(}\DataTypeTok{void}\OperatorTok{)}
\OperatorTok{\{}
    \DataTypeTok{FILE} \OperatorTok{*}\NormalTok{fp}\OperatorTok{;}                      \CommentTok{// Variabel untuk mewakili file yang terbuka}

\NormalTok{    fp }\OperatorTok{=}\NormalTok{ fopen}\OperatorTok{(}\StringTok{"hello.txt"}\OperatorTok{,} \StringTok{"r"}\OperatorTok{);}  \CommentTok{// Buka file untuk dibaca}

    \DataTypeTok{int}\NormalTok{ c }\OperatorTok{=}\NormalTok{ fgetc}\OperatorTok{(}\NormalTok{fp}\OperatorTok{);}             \CommentTok{// Baca satu karakter}
\NormalTok{    printf}\OperatorTok{(}\StringTok{"}\SpecialCharTok{\%c\textbackslash{}n}\StringTok{"}\OperatorTok{,}\NormalTok{ c}\OperatorTok{);}             \CommentTok{// Cetak karakter ke stdout}

\NormalTok{    fclose}\OperatorTok{(}\NormalTok{fp}\OperatorTok{);}                    \CommentTok{// Tutup file saat selesai}
\OperatorTok{\}}
\end{Highlighting}
\end{Shaded}

\end{tcolorbox}

Lihat bagaimana saat kita membuka file dengan \hl{\mbox{\texttt{fopen()}}}, ia mengembalikan \hl{\mbox{\texttt{FILE*}}} kepada kita supaya bisa dipakai nanti.

(Pengecekan ini dihilangkan demi keringkasan, tapi \hl{\mbox{\texttt{fopen()}}} akan mengembalikan \hl{\mbox{\texttt{NULL}}} kalau ada yang salah, misalnya file tidak ditemukan, jadi kamu seharusnya benar-benar melakukan pengecekan error!)

Perhatikan juga \hl{\mbox{\texttt{"r"}}} yang kita oper---ini berarti ``buka stream teks untuk dibaca''. (Ada berbagai string yang bisa kita oper ke \hl{\mbox{\texttt{fopen()}}} dengan makna tambahan, seperti menulis, menambahkan, dan sebagainya.) \index{fopen() function@\texttt{fopen()} function|)}

Setelah itu, kita memakai fungsi \hl{\mbox{\texttt{fgetc()}}} untuk mengambil satu karakter dari stream. Mungkin kamu bertanya-tanya kenapa \hl{\mbox{\texttt{c}}} dibuat \hl{\mbox{\texttt{int}}} dan bukan \hl{\mbox{\texttt{char}}}---tahan dulu pikiran itu! \index{fgetc() function@\texttt{fgetc()} function|)}

Terakhir, kita menutup stream saat selesai memakainya. Semua stream otomatis ditutup saat program keluar, tapi menutup sendiri setiap file yang sudah selesai dipakai itu sopan santun dan kebiasaan rumah tangga yang baik. \index{fclose() function@\texttt{fclose()} function|)}

\index{FILE* type@\texttt{FILE*} type}\hl{\mbox{\texttt{FILE*}}} melacak posisi kita di dalam file. Jadi pemanggilan \hl{\mbox{\texttt{fgetc()}}} berikutnya akan mengambil karakter berikutnya dalam file, lalu berikutnya lagi, sampai akhir.

Tapi itu terdengar merepotkan. Mari kita lihat apakah bisa dipermudah. \index{File I/O!text files, reading|)}

\section{\texorpdfstring{Akhir File: \hl{\mbox{\texttt{EOF}}}}{Akhir File: EOF}}\label{akhir-file-eof}

\index{EOF end of file@\texttt{EOF} end of file|(} Ada karakter khusus yang didefinisikan sebagai makro: \hl{\mbox{\texttt{EOF}}}. Inilah yang dikembalikan \hl{\mbox{\texttt{fgetc()}}} ketika akhir file tercapai dan kamu mencoba membaca karakter lagi.

Saatnya berbagi Fakta Menarik™ itu. Ternyata \hl{\mbox{\texttt{EOF}}} adalah alasan kenapa \hl{\mbox{\texttt{fgetc()}}} dan fungsi sejenisnya mengembalikan \hl{\mbox{\texttt{int}}} dan bukan \hl{\mbox{\texttt{char}}}. \hl{\mbox{\texttt{EOF}}} bukan karakter yang sebenarnya, dan nilainya kemungkinan berada di luar rentang \hl{\mbox{\texttt{char}}}. Karena \hl{\mbox{\texttt{fgetc()}}} harus bisa mengembalikan byte apa pun \textbf{dan} \hl{\mbox{\texttt{EOF}}}, ia harus bertipe lebih lebar yang bisa menampung lebih banyak nilai. Jadi dipilihlah \hl{\mbox{\texttt{int}}}. Tapi kecuali kamu membandingkan nilai kembaliannya dengan \hl{\mbox{\texttt{EOF}}}, kamu boleh tahu, jauh di dalamnya, ia adalah \hl{\mbox{\texttt{char}}}.

Baiklah! Kembali ke dunia nyata! Kita bisa memakai ini untuk membaca seluruh file dalam loop.

\index{fopen() function@\texttt{fopen()} function|(} \index{fgetc() function@\texttt{fgetc()} function|(} \index{fclose() function@\texttt{fclose()} function|(}

\begin{tcolorbox}[enhanced jigsaw,left=3pt,breakable,colback=white!97.2549!black,boxrule=0.2mm]

\begin{Shaded}
\begin{Highlighting}[numbers=left,,]
\PreprocessorTok{\#include }\ImportTok{\textless{}stdio.h\textgreater{}}

\DataTypeTok{int}\NormalTok{ main}\OperatorTok{(}\DataTypeTok{void}\OperatorTok{)}
\OperatorTok{\{}
    \DataTypeTok{FILE} \OperatorTok{*}\NormalTok{fp}\OperatorTok{;}
    \DataTypeTok{int}\NormalTok{ c}\OperatorTok{;}

\NormalTok{    fp }\OperatorTok{=}\NormalTok{ fopen}\OperatorTok{(}\StringTok{"hello.txt"}\OperatorTok{,} \StringTok{"r"}\OperatorTok{);}

    \ControlFlowTok{while} \OperatorTok{((}\NormalTok{c }\OperatorTok{=}\NormalTok{ fgetc}\OperatorTok{(}\NormalTok{fp}\OperatorTok{))} \OperatorTok{!=}\NormalTok{ EOF}\OperatorTok{)}
\NormalTok{        printf}\OperatorTok{(}\StringTok{"}\SpecialCharTok{\%c}\StringTok{"}\OperatorTok{,}\NormalTok{ c}\OperatorTok{);}

\NormalTok{    fclose}\OperatorTok{(}\NormalTok{fp}\OperatorTok{);}
\OperatorTok{\}}
\end{Highlighting}
\end{Shaded}

\end{tcolorbox}

\index{fopen() function@\texttt{fopen()} function|)} \index{fclose() function@\texttt{fclose()} function|)}

(Kalau baris 10 terlalu aneh, uraikan saja mulai dari tanda kurung yang paling dalam. Hal pertama yang dilakukan adalah menugaskan hasil \hl{\mbox{\texttt{fgetc()}}} ke \hl{\mbox{\texttt{c}}}, dan \emph{baru kemudian} membandingkan \emph{nilai itu} dengan \hl{\mbox{\texttt{EOF}}}. Kita hanya menjejalkannya dalam satu baris. Mungkin kelihatan sulit dibaca, tapi pelajari---ini idiomatis di C.) \index{fgetc() function@\texttt{fgetc()} function|)}

Dan saat dijalankan, kita lihat:

\begin{tcolorbox}[enhanced jigsaw,left=3pt,breakable,colback=white!97.2549!black,boxrule=0.2mm]

\begin{Shaded}
\begin{Highlighting}[]
\NormalTok{Hello, world!}
\end{Highlighting}
\end{Shaded}

\end{tcolorbox}

Tapi tetap saja kita bekerja satu karakter demi satu karakter, padahal banyak file teks lebih masuk akal dibaca per baris. Mari beralih ke itu. \index{EOF end of file@\texttt{EOF} end of file|)}

\subsection{Membaca Satu Baris Sekaligus}\label{membaca-satu-baris-sekaligus}

\index{File I/O!line by line|(} Jadi bagaimana mendapat satu baris utuh sekaligus? \index{fgets() function@\texttt{fgets()} function|(}\hl{\mbox{\texttt{fgets()}}} datang menyelamatkan! Sebagai argumen, ia menerima pointer ke buffer \hl{\mbox{\texttt{char}}} untuk menampung byte, jumlah byte maksimum yang dibaca, dan \hl{\mbox{\texttt{FILE*}}} tempat membaca. Ia mengembalikan \hl{\mbox{\texttt{NULL}}} saat end-of-file atau error. \hl{\mbox{\texttt{fgets()}}} bahkan cukup baik hati untuk mengakhiri string dengan NUL saat selesai\footnote{Kalau buffer tidak cukup besar untuk membaca satu baris penuh, ia akan berhenti membaca di tengah baris, dan pemanggilan \hl{\mbox{\texttt{fgets()}}} berikutnya akan melanjutkan sisa barisnya.}.\index{fgets() function@\texttt{fgets()} function|)}

Kita buat loop yang mirip dengan tadi, hanya saja dengan file multibaris yang dibaca satu baris sekaligus.

Berikut file \hl{\mbox{\texttt{quote.txt}}}:

\begin{tcolorbox}[enhanced jigsaw,left=3pt,breakable,colback=white!97.2549!black,boxrule=0.2mm]

\begin{Shaded}
\begin{Highlighting}[]
\NormalTok{A wise man can learn more from}
\NormalTok{a foolish question than a fool}
\NormalTok{can learn from a wise answer.}
\NormalTok{                  {-}{-}Bruce Lee}
\end{Highlighting}
\end{Shaded}

\end{tcolorbox}

Dan berikut kode yang membaca file itu satu baris sekaligus dan mencetak nomor baris sebelum masing-masing baris:

\index{fgets() function@\texttt{fgets()} function|(}

\begin{tcolorbox}[enhanced jigsaw,left=3pt,breakable,colback=white!97.2549!black,boxrule=0.2mm]

\begin{Shaded}
\begin{Highlighting}[numbers=left,,]
\PreprocessorTok{\#include }\ImportTok{\textless{}stdio.h\textgreater{}}

\DataTypeTok{int}\NormalTok{ main}\OperatorTok{(}\DataTypeTok{void}\OperatorTok{)}
\OperatorTok{\{}
    \DataTypeTok{FILE} \OperatorTok{*}\NormalTok{fp}\OperatorTok{;}
    \DataTypeTok{char}\NormalTok{ s}\OperatorTok{[}\DecValTok{1024}\OperatorTok{];}  \CommentTok{// Cukup besar untuk baris apa pun yang ditemui program ini}
    \DataTypeTok{int}\NormalTok{ linecount }\OperatorTok{=} \DecValTok{0}\OperatorTok{;}

\NormalTok{    fp }\OperatorTok{=}\NormalTok{ fopen}\OperatorTok{(}\StringTok{"quote.txt"}\OperatorTok{,} \StringTok{"r"}\OperatorTok{);}

    \ControlFlowTok{while} \OperatorTok{(}\NormalTok{fgets}\OperatorTok{(}\NormalTok{s}\OperatorTok{,} \KeywordTok{sizeof}\NormalTok{ s}\OperatorTok{,}\NormalTok{ fp}\OperatorTok{)} \OperatorTok{!=}\NormalTok{ NULL}\OperatorTok{)} 
\NormalTok{        printf}\OperatorTok{(}\StringTok{"}\SpecialCharTok{\%d}\StringTok{: }\SpecialCharTok{\%s}\StringTok{"}\OperatorTok{,} \OperatorTok{++}\NormalTok{linecount}\OperatorTok{,}\NormalTok{ s}\OperatorTok{);}

\NormalTok{    fclose}\OperatorTok{(}\NormalTok{fp}\OperatorTok{);}
\OperatorTok{\}}
\end{Highlighting}
\end{Shaded}

\end{tcolorbox}

\index{fgets() function@\texttt{fgets()} function|)}

Yang memberi output:

\begin{tcolorbox}[enhanced jigsaw,left=3pt,breakable,colback=white!97.2549!black,boxrule=0.2mm]

\begin{Shaded}
\begin{Highlighting}[]
\NormalTok{1: A wise man can learn more from}
\NormalTok{2: a foolish question than a fool}
\NormalTok{3: can learn from a wise answer.}
\NormalTok{4:                   {-}{-}Bruce Lee}
\end{Highlighting}
\end{Shaded}

\end{tcolorbox}

\index{File I/O!line by line|)}

\section{Input Terformat}\label{input-terformat}

\index{File I/O!formatted input|(} Kamu tahu bagaimana kamu bisa mendapat output terformat dengan \hl{\mbox{\texttt{printf()}}} (dan, dengan demikian, \hl{\mbox{\texttt{fprintf()}}} seperti yang akan kita lihat di bawah)?

\index{fscanf() function@\texttt{fscanf()} function|(} Hal yang sama bisa dilakukan dengan \hl{\mbox{\texttt{fscanf()}}}.

\begin{tcolorbox}[blanker,colback=white,enhanced jigsaw,breakable,arc=0pt,outer arc=0pt,grow to left by=-1cm,enlarge top by=0.1cm,enlarge bottom by=0.1cm,borderline west={1.5pt}{-.3cm}{black}]

\begin{quote}
Sebelum mulai, perlu diingatkan bahwa memakai fungsi gaya \hl{\mbox{\texttt{scanf()}}} bisa berbahaya untuk input yang tidak tepercaya. Kalau kamu tidak menentukan lebar field pada \hl{\mbox{\texttt{\%s}}}, buffer bisa meluap. Lebih buruk lagi, konversi numerik yang tidak valid menghasilkan undefined behavior. Cara aman untuk input yang tidak tepercaya adalah memakai \hl{\mbox{\texttt{\%s}}} dengan lebar field, lalu memakai fungsi seperti \index{strtol function@\texttt{strtol} function} \hl{\mbox{\texttt{strtol()}}} atau \index{strtod function@\texttt{strtod} function} \hl{\mbox{\texttt{strtod()}}} untuk melakukan konversinya.
\end{quote}

\end{tcolorbox}

Kita siapkan file berisi serangkaian record data. Dalam hal ini, paus, dengan nama, panjang dalam meter, dan bobot dalam ton. \hl{\mbox{\texttt{whales.txt}}}:

\begin{tcolorbox}[enhanced jigsaw,left=3pt,breakable,colback=white!97.2549!black,boxrule=0.2mm]

\begin{Shaded}
\begin{Highlighting}[]
\NormalTok{blue 29.9 173}
\NormalTok{right 20.7 135}
\NormalTok{gray 14.9 41}
\NormalTok{humpback 16.0 30}
\end{Highlighting}
\end{Shaded}

\end{tcolorbox}

Ya, kita bisa membacanya dengan \index{fgets() function@\texttt{fgets()} function}\hl{\mbox{\texttt{fgets()}}} lalu mem-parse string-nya dengan \hl{\mbox{\texttt{sscanf()}}} (dan itu lebih tahan terhadap file yang rusak), tapi dalam kasus ini, kita pakai saja \hl{\mbox{\texttt{fscanf()}}} dan mengambilnya langsung.

Fungsi \hl{\mbox{\texttt{fscanf()}}} melewati whitespace di depan saat membaca, dan mengembalikan \hl{\mbox{\texttt{EOF}}} saat end-of-file atau error.

\begin{tcolorbox}[enhanced jigsaw,left=3pt,breakable,colback=white!97.2549!black,boxrule=0.2mm]

\begin{Shaded}
\begin{Highlighting}[numbers=left,,]
\PreprocessorTok{\#include }\ImportTok{\textless{}stdio.h\textgreater{}}

\DataTypeTok{int}\NormalTok{ main}\OperatorTok{(}\DataTypeTok{void}\OperatorTok{)}
\OperatorTok{\{}
    \DataTypeTok{FILE} \OperatorTok{*}\NormalTok{fp}\OperatorTok{;}
    \DataTypeTok{char}\NormalTok{ name}\OperatorTok{[}\DecValTok{1024}\OperatorTok{];}  \CommentTok{// Cukup besar untuk baris apa pun yang ditemui program ini}
    \DataTypeTok{float}\NormalTok{ length}\OperatorTok{;}
    \DataTypeTok{int}\NormalTok{ mass}\OperatorTok{;}

\NormalTok{    fp }\OperatorTok{=}\NormalTok{ fopen}\OperatorTok{(}\StringTok{"whales.txt"}\OperatorTok{,} \StringTok{"r"}\OperatorTok{);}

    \ControlFlowTok{while} \OperatorTok{(}\NormalTok{fscanf}\OperatorTok{(}\NormalTok{fp}\OperatorTok{,} \StringTok{"}\SpecialCharTok{\%s}\StringTok{ }\SpecialCharTok{\%f}\StringTok{ }\SpecialCharTok{\%d}\StringTok{"}\OperatorTok{,}\NormalTok{ name}\OperatorTok{,} \OperatorTok{\&}\NormalTok{length}\OperatorTok{,} \OperatorTok{\&}\NormalTok{mass}\OperatorTok{)} \OperatorTok{!=}\NormalTok{ EOF}\OperatorTok{)}
\NormalTok{        printf}\OperatorTok{(}\StringTok{"}\SpecialCharTok{\%s}\StringTok{ whale, }\SpecialCharTok{\%d}\StringTok{ tonnes, }\SpecialCharTok{\%.1f}\StringTok{ meters}\SpecialCharTok{\textbackslash{}n}\StringTok{"}\OperatorTok{,}\NormalTok{ name}\OperatorTok{,}\NormalTok{ mass}\OperatorTok{,}\NormalTok{ length}\OperatorTok{);}

\NormalTok{    fclose}\OperatorTok{(}\NormalTok{fp}\OperatorTok{);}
\OperatorTok{\}}
\end{Highlighting}
\end{Shaded}

\end{tcolorbox}

\index{fscanf() function@\texttt{fscanf()} function|)}

Yang memberi hasil:

\begin{tcolorbox}[enhanced jigsaw,left=3pt,breakable,colback=white!97.2549!black,boxrule=0.2mm]

\begin{Shaded}
\begin{Highlighting}[]
\NormalTok{blue whale, 173 tonnes, 29.9 meters}
\NormalTok{right whale, 135 tonnes, 20.7 meters}
\NormalTok{gray whale, 41 tonnes, 14.9 meters}
\NormalTok{humpback whale, 30 tonnes, 16.0 meters}
\end{Highlighting}
\end{Shaded}

\end{tcolorbox}

\index{File I/O!formatted input|)}

\section{Menulis File Teks}\label{menulis-file-teks}

\index{File I/O!text files, writing|(} \index{fputc() function@\texttt{fputc()} function|(} \index{fputs() function@\texttt{fputs()} function|(} \index{fprintf() function@\texttt{fprintf()} function|(} Sama seperti kita bisa memakai \hl{\mbox{\texttt{fgetc()}}}, \hl{\mbox{\texttt{fgets()}}}, dan \hl{\mbox{\texttt{fscanf()}}} untuk membaca stream teks, kita bisa memakai \hl{\mbox{\texttt{fputc()}}}, \hl{\mbox{\texttt{fputs()}}}, dan \hl{\mbox{\texttt{fprintf()}}} untuk menulis stream teks.

Untuk itu, kita harus meng-\hl{\mbox{\texttt{fopen()}}} file dalam mode tulis dengan mengoper \hl{\mbox{\texttt{"w"}}} sebagai argumen kedua. Membuka file yang sudah ada dalam mode \hl{\mbox{\texttt{"w"}}} akan langsung memotong file itu menjadi 0 byte untuk ditimpa sepenuhnya.

Kita susun program sederhana yang mengeluarkan file \hl{\mbox{\texttt{output.txt}}} dengan berbagai fungsi output.

\begin{tcolorbox}[enhanced jigsaw,left=3pt,breakable,colback=white!97.2549!black,boxrule=0.2mm]

\begin{Shaded}
\begin{Highlighting}[numbers=left,,]
\PreprocessorTok{\#include }\ImportTok{\textless{}stdio.h\textgreater{}}

\DataTypeTok{int}\NormalTok{ main}\OperatorTok{(}\DataTypeTok{void}\OperatorTok{)}
\OperatorTok{\{}
    \DataTypeTok{FILE} \OperatorTok{*}\NormalTok{fp}\OperatorTok{;}
    \DataTypeTok{int}\NormalTok{ x }\OperatorTok{=} \DecValTok{32}\OperatorTok{;}

\NormalTok{    fp }\OperatorTok{=}\NormalTok{ fopen}\OperatorTok{(}\StringTok{"output.txt"}\OperatorTok{,} \StringTok{"w"}\OperatorTok{);}

\NormalTok{    fputc}\OperatorTok{(}\CharTok{\textquotesingle{}B\textquotesingle{}}\OperatorTok{,}\NormalTok{ fp}\OperatorTok{);}
\NormalTok{    fputc}\OperatorTok{(}\CharTok{\textquotesingle{}}\SpecialCharTok{\textbackslash{}n}\CharTok{\textquotesingle{}}\OperatorTok{,}\NormalTok{ fp}\OperatorTok{);}   \CommentTok{// newline}
\NormalTok{    fprintf}\OperatorTok{(}\NormalTok{fp}\OperatorTok{,} \StringTok{"x = }\SpecialCharTok{\%d\textbackslash{}n}\StringTok{"}\OperatorTok{,}\NormalTok{ x}\OperatorTok{);}
\NormalTok{    fputs}\OperatorTok{(}\StringTok{"Hello, world!}\SpecialCharTok{\textbackslash{}n}\StringTok{"}\OperatorTok{,}\NormalTok{ fp}\OperatorTok{);}

\NormalTok{    fclose}\OperatorTok{(}\NormalTok{fp}\OperatorTok{);}
\OperatorTok{\}}
\end{Highlighting}
\end{Shaded}

\end{tcolorbox}

\index{fputc() function@\texttt{fputc()} function|)} \index{fputs() function@\texttt{fputs()} function|)} \index{fprintf() function@\texttt{fprintf()} function|)}

Dan ini menghasilkan file \hl{\mbox{\texttt{output.txt}}} dengan isi:

\begin{tcolorbox}[enhanced jigsaw,left=3pt,breakable,colback=white!97.2549!black,boxrule=0.2mm]

\begin{Shaded}
\begin{Highlighting}[]
\NormalTok{B}
\NormalTok{x = 32}
\NormalTok{Hello, world!}
\end{Highlighting}
\end{Shaded}

\end{tcolorbox}

Fakta menarik: karena \hl{\mbox{\texttt{stdout}}} adalah sebuah file, kamu bisa mengganti baris 8 dengan:

\begin{tcolorbox}[enhanced jigsaw,left=3pt,breakable,colback=white!97.2549!black,boxrule=0.2mm]

\begin{Shaded}
\begin{Highlighting}[]
\NormalTok{fp }\OperatorTok{=}\NormalTok{ stdout}\OperatorTok{;}
\end{Highlighting}
\end{Shaded}

\end{tcolorbox}

dan program akan mengeluarkan output ke konsol alih-alih ke file. Coba saja! \index{File I/O!text files, writing|)}

\section{I/O File Biner}\label{io-file-biner}

\index{File I/O!binary files|(} Sejauh ini kita hanya membahas file teks. Tapi ada makhluk lain yang sudah disebut di awal, yaitu file \emph{biner}, atau stream biner.

Cara kerjanya sangat mirip file teks, kecuali subsistem I/O tidak melakukan terjemahan apa pun pada data seperti yang mungkin dilakukannya pada file teks. Pada file biner, kamu mendapat aliran byte mentah, dan hanya itu.

Perbedaan besar saat membuka file: kamu harus menambahkan \hl{\mbox{\texttt{"b"}}} ke mode. Artinya, untuk membaca file biner, buka dalam mode \hl{\mbox{\texttt{"rb"}}}. Untuk menulis file, buka dalam mode \hl{\mbox{\texttt{"wb"}}}.

Karena ini aliran byte, dan aliran byte bisa berisi karakter NUL, dan karakter NUL adalah penanda akhir string di C, jarang orang memakai fungsi \hl{\mbox{\texttt{fprintf()}}} dan kawan-kawannya untuk mengolah file biner.

\index{fwrite() function@\texttt{fwrite()} function|(} Sebagai gantinya, fungsi yang paling umum adalah \index{fread() function@\texttt{fread()} function}\hl{\mbox{\texttt{fread()}}} dan \hl{\mbox{\texttt{fwrite()}}}. Fungsi-fungsi ini membaca dan menulis sejumlah byte tertentu ke stream.

Untuk demo, kita tulis dua program. Yang satu menulis serangkaian nilai byte ke disk sekaligus. Dan program kedua membaca satu byte per satu dan mencetaknya\footnote{Biasanya program kedua akan membaca semua byte sekaligus, lalu \emph{baru kemudian} mencetaknya dalam loop. Itu lebih efisien. Tapi di sini kita mengejar nilai demonstrasi.}.

\begin{tcolorbox}[enhanced jigsaw,left=3pt,breakable,colback=white!97.2549!black,boxrule=0.2mm]

\begin{Shaded}
\begin{Highlighting}[numbers=left,,]
\PreprocessorTok{\#include }\ImportTok{\textless{}stdio.h\textgreater{}}

\DataTypeTok{int}\NormalTok{ main}\OperatorTok{(}\DataTypeTok{void}\OperatorTok{)}
\OperatorTok{\{}
    \DataTypeTok{FILE} \OperatorTok{*}\NormalTok{fp}\OperatorTok{;}
    \DataTypeTok{unsigned} \DataTypeTok{char}\NormalTok{ bytes}\OperatorTok{[}\DecValTok{6}\OperatorTok{]} \OperatorTok{=} \OperatorTok{\{}\DecValTok{5}\OperatorTok{,} \DecValTok{37}\OperatorTok{,} \DecValTok{0}\OperatorTok{,} \DecValTok{88}\OperatorTok{,} \DecValTok{255}\OperatorTok{,} \DecValTok{12}\OperatorTok{\};}

\NormalTok{    fp }\OperatorTok{=}\NormalTok{ fopen}\OperatorTok{(}\StringTok{"output.bin"}\OperatorTok{,} \StringTok{"wb"}\OperatorTok{);}  \CommentTok{// mode wb untuk "write binary"!}

    \CommentTok{// Pada pemanggilan fwrite, argumennya adalah:}
    \CommentTok{//}
    \CommentTok{// * Pointer ke data yang akan ditulis}
    \CommentTok{// * Ukuran setiap "potongan" data}
    \CommentTok{// * Jumlah "potongan" data}
    \CommentTok{// * FILE*}

\NormalTok{    fwrite}\OperatorTok{(}\NormalTok{bytes}\OperatorTok{,} \KeywordTok{sizeof}\OperatorTok{(}\DataTypeTok{char}\OperatorTok{),} \DecValTok{6}\OperatorTok{,}\NormalTok{ fp}\OperatorTok{);}

\NormalTok{    fclose}\OperatorTok{(}\NormalTok{fp}\OperatorTok{);}
\OperatorTok{\}}
\end{Highlighting}
\end{Shaded}

\end{tcolorbox}

\index{fwrite() function@\texttt{fwrite()} function|)}

Dua argumen tengah \hl{\mbox{\texttt{fwrite()}}} itu agak aneh. Tapi pada dasarnya yang ingin kita sampaikan kepada fungsi itu: ``Kita punya item berukuran \emph{sebesar ini}, dan kita ingin menulis \emph{sebanyak itu}.'' Ini praktis kalau kamu punya record dengan panjang tetap, dan banyak record dalam sebuah array. Kamu cukup memberi tahu ukuran satu record dan berapa yang ditulis.

Pada contoh di atas, kita memberi tahu bahwa setiap record berukuran sebesar \hl{\mbox{\texttt{char}}}, dan kita punya 6 buah.

Menjalankan program ini menghasilkan file \hl{\mbox{\texttt{output.bin}}}, tapi membukanya di text editor tidak menampilkan apa pun yang ramah! Itu data biner---bukan teks. Dan data biner acak yang baru saja dikarang, pula!

Kalau kita jalankan lewat program \href{https://en.wikipedia.org/wiki/Hex_dump}{hex dump}\footnote{https://en.wikipedia.org/wiki/Hex\_dump}, kita bisa melihat outputnya sebagai byte:

\begin{tcolorbox}[enhanced jigsaw,left=3pt,breakable,colback=white!97.2549!black,boxrule=0.2mm]

\begin{Shaded}
\begin{Highlighting}[]
\NormalTok{05 25 00 58 ff 0c}
\end{Highlighting}
\end{Shaded}

\end{tcolorbox}

\begin{tcolorbox}[blanker,colback=white,enhanced jigsaw,breakable,arc=0pt,outer arc=0pt,grow to left by=-1cm,enlarge top by=0.1cm,enlarge bottom by=0.1cm,borderline west={1.5pt}{-.3cm}{black}]

\begin{quote}
Banyak sistem Unix menyertakan program bernama \hl{\mbox{\texttt{hexdump}}} untuk ini. Kamu bisa memakainya seperti ini dengan switch \hl{\mbox{\texttt{-C}}} (``canonical'') untuk mendapat output yang rapi:

\begin{tcolorbox}[enhanced jigsaw,left=3pt,breakable,colback=white!97.2549!black,boxrule=0.2mm]

\begin{Shaded}
\begin{Highlighting}[]
\NormalTok{$ hexdump {-}C output.bin}
\NormalTok{00000000  05 25 00 58 ff 0c                              |.\%.X..|}
\end{Highlighting}
\end{Shaded}

\end{tcolorbox}

\hl{\mbox{\texttt{00000000}}} adalah offset dalam file tempat baris output ini dimulai. \hl{\mbox{\texttt{05\ 25\ 00\ 58\ ff\ 0c}}} adalah nilai-nilai byte (dan akan lebih panjang (sampai 16 byte per baris) kalau file berisi lebih banyak byte). Dan di sebelah kanan, di antara simbol pipa (\hl{\mbox{\texttt{\textbar{}}}}), adalah upaya terbaik \hl{\mbox{\texttt{hexdump}}} untuk mencetak karakter yang sesuai dengan byte-byte itu. Ia mencetak titik kalau karakternya tidak bisa dicetak. Dalam kasus ini, karena kita hanya mencetak data biner acak, bagian output ini hanya sampah. Tapi kalau kita menulis string ASCII ke file, kita akan melihatnya di sana.
\end{quote}

\end{tcolorbox}

Dan nilai-nilai hex itu memang cocok dengan nilai-nilai (dalam desimal) yang kita tulis.

Tapi sekarang mari kita coba membacanya kembali dengan program lain. Yang ini akan membuka file untuk pembacaan biner (mode \hl{\mbox{\texttt{"rb"}}}) dan membaca byte satu per satu dalam loop.

\index{fread() function@\texttt{fread()} function|(} \hl{\mbox{\texttt{fread()}}} punya fitur menarik: ia mengembalikan jumlah byte yang dibaca, atau \hl{\mbox{\texttt{0}}} saat EOF. Jadi kita bisa me-loop sampai melihat itu, sambil mencetak angka.

\begin{tcolorbox}[enhanced jigsaw,left=3pt,breakable,colback=white!97.2549!black,boxrule=0.2mm]

\begin{Shaded}
\begin{Highlighting}[numbers=left,,]
\PreprocessorTok{\#include }\ImportTok{\textless{}stdio.h\textgreater{}}

\DataTypeTok{int}\NormalTok{ main}\OperatorTok{(}\DataTypeTok{void}\OperatorTok{)}
\OperatorTok{\{}
    \DataTypeTok{FILE} \OperatorTok{*}\NormalTok{fp}\OperatorTok{;}
    \DataTypeTok{unsigned} \DataTypeTok{char}\NormalTok{ c}\OperatorTok{;}

\NormalTok{    fp }\OperatorTok{=}\NormalTok{ fopen}\OperatorTok{(}\StringTok{"output.bin"}\OperatorTok{,} \StringTok{"rb"}\OperatorTok{);} \CommentTok{// rb untuk "read binary"!}

    \ControlFlowTok{while} \OperatorTok{(}\NormalTok{fread}\OperatorTok{(\&}\NormalTok{c}\OperatorTok{,} \KeywordTok{sizeof}\OperatorTok{(}\DataTypeTok{char}\OperatorTok{),} \DecValTok{1}\OperatorTok{,}\NormalTok{ fp}\OperatorTok{)} \OperatorTok{\textgreater{}} \DecValTok{0}\OperatorTok{)}
\NormalTok{        printf}\OperatorTok{(}\StringTok{"}\SpecialCharTok{\%d\textbackslash{}n}\StringTok{"}\OperatorTok{,}\NormalTok{ c}\OperatorTok{);}

\NormalTok{    fclose}\OperatorTok{(}\NormalTok{fp}\OperatorTok{);}
\OperatorTok{\}}
\end{Highlighting}
\end{Shaded}

\end{tcolorbox}

\index{fread() function@\texttt{fread()} function|)}

Dan saat dijalankan, kita melihat angka-angka asli kita!

\begin{tcolorbox}[enhanced jigsaw,left=3pt,breakable,colback=white!97.2549!black,boxrule=0.2mm]

\begin{Shaded}
\begin{Highlighting}[]
\NormalTok{5}
\NormalTok{37}
\NormalTok{0}
\NormalTok{88}
\NormalTok{255}
\NormalTok{12}
\end{Highlighting}
\end{Shaded}

\end{tcolorbox}

Horee! \index{File I/O!binary files|)}

\subsection{\texorpdfstring{Peringatan tentang \hl{\mbox{\texttt{struct}}} dan Angka}{Peringatan tentang struct dan Angka}}\label{peringatan-tentang-struct-dan-angka}

\index{File I/O!with structs@with \texttt{struct}s|(} Seperti yang kita lihat di bagian \hl{\mbox{\texttt{struct}}}, compiler bebas menambahkan padding pada \hl{\mbox{\texttt{struct}}} sesuai kehendaknya. Dan compiler yang berbeda mungkin melakukannya secara berbeda. Dan compiler yang sama pada arsitektur berbeda bisa melakukannya secara berbeda. Dan compiler yang sama pada arsitektur yang sama pun bisa melakukannya secara berbeda.

Intinya: tidak portabel kalau langsung \hl{\mbox{\texttt{fwrite()}}} seluruh \hl{\mbox{\texttt{struct}}} ke file ketika kamu tidak tahu di mana padding akan berakhir. \index{File I/O!with structs@with \texttt{struct}s|)}

Bagaimana cara memperbaikinya? Tahan dulu pikiran itu---kita akan melihat beberapa caranya setelah membahas masalah lain yang terkait.

\index{File I/O!with numeric values|(} Angka!

Ternyata tidak semua arsitektur merepresentasikan angka di memori dengan cara yang sama.

Mari kita lihat \hl{\mbox{\texttt{fwrite()}}} sederhana untuk angka 2-byte. Kita tulis dalam hex supaya setiap byte jelas. Byte paling signifikan bernilai \hl{\mbox{\texttt{0x12}}} dan yang paling tidak signifikan bernilai \hl{\mbox{\texttt{0x34}}}.

\begin{tcolorbox}[enhanced jigsaw,left=3pt,breakable,colback=white!97.2549!black,boxrule=0.2mm]

\begin{Shaded}
\begin{Highlighting}[]
\DataTypeTok{unsigned} \DataTypeTok{short}\NormalTok{ v }\OperatorTok{=} \BaseNTok{0x1234}\OperatorTok{;}  \CommentTok{// Dua byte, 0x12 dan 0x34}

\NormalTok{fwrite}\OperatorTok{(\&}\NormalTok{v}\OperatorTok{,} \KeywordTok{sizeof}\NormalTok{ v}\OperatorTok{,} \DecValTok{1}\OperatorTok{,}\NormalTok{ fp}\OperatorTok{);}
\end{Highlighting}
\end{Shaded}

\end{tcolorbox}

Apa yang berakhir di stream?

Kelihatannya seharusnya \hl{\mbox{\texttt{0x12}}} diikuti \hl{\mbox{\texttt{0x34}}}, kan?

Tapi kalau ini dijalankan di mesin penulis dan hasilnya di-hex dump, hasilnya:

\begin{tcolorbox}[enhanced jigsaw,left=3pt,breakable,colback=white!97.2549!black,boxrule=0.2mm]

\begin{Shaded}
\begin{Highlighting}[]
\NormalTok{34 12}
\end{Highlighting}
\end{Shaded}

\end{tcolorbox}

Terbalik! Kenapa?

Ini ada hubungannya dengan yang disebut \index{Endianess}\href{https://en.wikipedia.org/wiki/Endianess}{\emph{endianess}}\footnote{https://en.wikipedia.org/wiki/Endianess} arsitektur. Sebagian menulis byte paling signifikan lebih dulu, dan sebagian menulis byte paling tidak signifikan lebih dulu.

Artinya, kalau kamu menulis angka multibyte langsung dari memori, kamu tidak bisa melakukannya secara portabel\footnote{Dan inilah alasan penulis dengan cerdik memakai byte individual pada contoh \hl{\mbox{\texttt{fwrite()}}} dan \hl{\mbox{\texttt{fread()}}} di atas.}.

Masalah serupa ada pada floating point. Kebanyakan sistem memakai format yang sama untuk bilangan floating point, tapi sebagian tidak. Tidak ada jaminan!

\index{File I/O!with structs@with \texttt{struct}s|(} Jadi\ldots{} bagaimana kita memperbaiki semua masalah angka dan \hl{\mbox{\texttt{struct}}} ini supaya data kita tertulis secara portabel?

Ringkasnya adalah dengan men-\index{Data serialization}\emph{serialize} data, yaitu istilah umum yang berarti mengambil semua data dan menuliskannya dalam format yang kamu kendalikan, yang terkenal, dan bisa diprogram agar bekerja sama di semua platform.

Seperti yang mungkin kamu duga, ini masalah yang sudah terpecahkan. Ada banyak library serialisasi yang bisa kamu manfaatkan, seperti \href{https://en.wikipedia.org/wiki/Protocol_buffers}{\emph{protocol buffers}}\footnote{https://en.wikipedia.org/wiki/Protocol\_buffers} dari Google, yang sudah tersedia dan siap dipakai. Mereka akan mengurus semua detail rumitnya, bahkan memungkinkan data dari programmu di C berinteroperasi dengan bahasa lain yang mendukung metode serialisasi yang sama.

Lakukan kebaikan untuk dirimu dan semua orang! Serialize data binermu saat menulisnya ke stream! Ini akan menjaga semuanya tetap rapi dan portabel, bahkan kalau kamu memindahkan file data dari satu arsitektur ke arsitektur lain. \index{File I/O!with structs@with \texttt{struct}s|)} \index{File I/O!with numeric values|)} \index{File I/O|)}

\chapter{\texorpdfstring{\hl{\mbox{\texttt{typedef}}}: Membuat Tipe Baru}{typedef: Membuat Tipe Baru}}\label{typedef-membuat-tipe-baru}

\index{typedef keyword@\texttt{typedef} keyword|(} Sebenarnya bukan membuat tipe \emph{baru}, melainkan memberi nama baru pada tipe yang sudah ada. Sekilas kedengarannya sia-sia, tapi ini bisa kita manfaatkan untuk membuat kode lebih bersih.

\section{\texorpdfstring{\hl{\mbox{\texttt{typedef}}} dalam Teori}{typedef dalam Teori}}\label{typedef-dalam-teori}

Pada dasarnya, kamu mengambil tipe yang sudah ada lalu membuat alias untuk tipe itu dengan \hl{\mbox{\texttt{typedef}}}.

Seperti ini:

\begin{tcolorbox}[enhanced jigsaw,left=3pt,breakable,colback=white!97.2549!black,boxrule=0.2mm]

\begin{Shaded}
\begin{Highlighting}[]
\KeywordTok{typedef} \DataTypeTok{int}\NormalTok{ antelope}\OperatorTok{;}  \CommentTok{// Jadikan "antelope" alias untuk "int"}

\NormalTok{antelope x }\OperatorTok{=} \DecValTok{10}\OperatorTok{;}       \CommentTok{// Tipe "antelope" sama dengan tipe "int"}
\end{Highlighting}
\end{Shaded}

\end{tcolorbox}

Kamu bisa melakukannya pada tipe apa pun yang sudah ada. Kamu bahkan bisa membuat beberapa tipe sekaligus dengan daftar yang dipisahkan koma:

\begin{tcolorbox}[enhanced jigsaw,left=3pt,breakable,colback=white!97.2549!black,boxrule=0.2mm]

\begin{Shaded}
\begin{Highlighting}[]
\KeywordTok{typedef} \DataTypeTok{int}\NormalTok{ antelope}\OperatorTok{,}\NormalTok{ bagel}\OperatorTok{,}\NormalTok{ mushroom}\OperatorTok{;}  \CommentTok{// Semuanya "int"}
\end{Highlighting}
\end{Shaded}

\end{tcolorbox}

Itu sangat berguna, kan? Bisa mengetik \hl{\mbox{\texttt{mushroom}}} alih-alih \hl{\mbox{\texttt{int}}}? Kamu pasti \emph{sangat bersemangat} dengan fitur ini!

Oke, Profesor Sarkasme---sebentar lagi kita lihat penerapan yang lebih umum.

\subsection{Scoping}\label{scoping}

\index{typedef keyword@\texttt{typedef} keyword!scoping rules|(} \hl{\mbox{\texttt{typedef}}} mengikuti \hyperref[scope]{aturan scope} biasa.

Karena itu, sangat umum menemukan \hl{\mbox{\texttt{typedef}}} di file scope (``global'') supaya semua fungsi bisa memakai tipe baru itu sesuka hati.

\section{\texorpdfstring{\hl{\mbox{\texttt{typedef}}} dalam Praktik}{typedef dalam Praktik}}\label{typedef-dalam-praktik}

Jadi mengganti nama \hl{\mbox{\texttt{int}}} menjadi sesuatu yang lain memang tidak terlalu menarik. Mari kita lihat di mana \hl{\mbox{\texttt{typedef}}} biasanya muncul.

\subsection{\texorpdfstring{\hl{\mbox{\texttt{typedef}}} dan \hl{\mbox{\texttt{struct}}}}{typedef dan struct}}\label{typedef-struct}

\index{typedef keyword@\texttt{typedef} keyword!with structs@with \texttt{struct}s|(} Kadang sebuah \hl{\mbox{\texttt{struct}}} di-\hl{\mbox{\texttt{typedef}}} menjadi nama baru supaya kita tidak perlu mengetik kata \hl{\mbox{\texttt{struct}}} berulang-ulang.

\begin{tcolorbox}[enhanced jigsaw,left=3pt,breakable,colback=white!97.2549!black,boxrule=0.2mm]

\begin{Shaded}
\begin{Highlighting}[]
\KeywordTok{struct}\NormalTok{ animal }\OperatorTok{\{}
    \DataTypeTok{char} \OperatorTok{*}\NormalTok{name}\OperatorTok{;}
    \DataTypeTok{int}\NormalTok{ leg\_count}\OperatorTok{,}\NormalTok{ speed}\OperatorTok{;}
\OperatorTok{\};}

\CommentTok{//  nama asli          nama baru}
\CommentTok{//            |         |}
\CommentTok{//            v         v}
\CommentTok{//      |{-}{-}{-}{-}{-}{-}{-}{-}{-}{-}{-}| |{-}{-}{-}{-}|}
\KeywordTok{typedef} \KeywordTok{struct}\NormalTok{ animal animal}\OperatorTok{;}

\KeywordTok{struct}\NormalTok{ animal y}\OperatorTok{;}  \CommentTok{// Ini berfungsi}
\NormalTok{animal z}\OperatorTok{;}         \CommentTok{// Ini juga berfungsi karena "animal" adalah alias}
\end{Highlighting}
\end{Shaded}

\end{tcolorbox}

Secara pribadi, penulis tidak menyukai praktik ini. Penulis suka kejelasan yang didapat kode ketika kata \hl{\mbox{\texttt{struct}}} ditambahkan ke tipe; programmer tahu apa yang mereka dapat. Tapi ini sangat umum, jadi tetap dimasukkan di sini.

Sekarang kita jalankan contoh yang persis sama dengan cara yang mungkin sering kamu lihat. Kita akan menaruh \hl{\mbox{\texttt{struct\ animal}}} \emph{di dalam} \hl{\mbox{\texttt{typedef}}}. Semuanya bisa digabung seperti ini:

\begin{tcolorbox}[enhanced jigsaw,left=3pt,breakable,colback=white!97.2549!black,boxrule=0.2mm]

\begin{Shaded}
\begin{Highlighting}[]
\CommentTok{//  nama asli}
\CommentTok{//            |}
\CommentTok{//            v}
\CommentTok{//      |{-}{-}{-}{-}{-}{-}{-}{-}{-}{-}{-}|}
\KeywordTok{typedef} \KeywordTok{struct}\NormalTok{ animal }\OperatorTok{\{}
    \DataTypeTok{char} \OperatorTok{*}\NormalTok{name}\OperatorTok{;}
    \DataTypeTok{int}\NormalTok{ leg\_count}\OperatorTok{,}\NormalTok{ speed}\OperatorTok{;}
\OperatorTok{\}}\NormalTok{ animal}\OperatorTok{;}                         \CommentTok{// \textless{}{-}{-} nama baru}

\KeywordTok{struct}\NormalTok{ animal y}\OperatorTok{;}  \CommentTok{// Ini berfungsi}
\NormalTok{animal z}\OperatorTok{;}         \CommentTok{// Ini juga berfungsi karena "animal" adalah alias}
\end{Highlighting}
\end{Shaded}

\end{tcolorbox}

Itu persis sama dengan contoh sebelumnya, hanya lebih ringkas.

\index{typedef keyword@\texttt{typedef} keyword!with anonymous structs@with anonymous \texttt{struct}s|(} Tapi itu belum semua! Ada jalan pintas umum lain yang mungkin kamu lihat di kode yang memakai apa yang disebut \emph{struktur anonim}\footnote{Soal ini dibahas lebih lanjut nanti.}. Ternyata di berbagai tempat kita tidak perlu memberi nama pada struktur, dan \hl{\mbox{\texttt{typedef}}} adalah salah satunya.

Kita buat contoh yang sama dengan struktur anonim:

\begin{tcolorbox}[enhanced jigsaw,left=3pt,breakable,colback=white!97.2549!black,boxrule=0.2mm]

\begin{Shaded}
\begin{Highlighting}[]
\CommentTok{//  struct anonim! Tidak punya nama!}
\CommentTok{//         |}
\CommentTok{//         v}
\CommentTok{//      |{-}{-}{-}{-}|}
\KeywordTok{typedef} \KeywordTok{struct} \OperatorTok{\{}
    \DataTypeTok{char} \OperatorTok{*}\NormalTok{name}\OperatorTok{;}
    \DataTypeTok{int}\NormalTok{ leg\_count}\OperatorTok{,}\NormalTok{ speed}\OperatorTok{;}
\OperatorTok{\}}\NormalTok{ animal}\OperatorTok{;}                         \CommentTok{// \textless{}{-}{-} nama baru}

\CommentTok{//struct animal y;  // ERROR: ini tidak berfungsi lagi{-}{-}tidak ada struct itu!}
\NormalTok{animal z}\OperatorTok{;}           \CommentTok{// Ini berfungsi karena "animal" adalah alias}
\end{Highlighting}
\end{Shaded}

\end{tcolorbox}

Sebagai contoh lain, kita mungkin menemukan sesuatu seperti ini:

\begin{tcolorbox}[enhanced jigsaw,left=3pt,breakable,colback=white!97.2549!black,boxrule=0.2mm]

\begin{Shaded}
\begin{Highlighting}[]
\KeywordTok{typedef} \KeywordTok{struct} \OperatorTok{\{}
    \DataTypeTok{int}\NormalTok{ x}\OperatorTok{,}\NormalTok{ y}\OperatorTok{;}
\OperatorTok{\}}\NormalTok{ point}\OperatorTok{;}

\NormalTok{point p }\OperatorTok{=} \OperatorTok{\{.}\NormalTok{x}\OperatorTok{=}\DecValTok{20}\OperatorTok{,} \OperatorTok{.}\NormalTok{y}\OperatorTok{=}\DecValTok{40}\OperatorTok{\};}

\NormalTok{printf}\OperatorTok{(}\StringTok{"}\SpecialCharTok{\%d}\StringTok{, }\SpecialCharTok{\%d\textbackslash{}n}\StringTok{"}\OperatorTok{,}\NormalTok{ p}\OperatorTok{.}\NormalTok{x}\OperatorTok{,}\NormalTok{ p}\OperatorTok{.}\NormalTok{y}\OperatorTok{);}  \CommentTok{// 20, 40}
\end{Highlighting}
\end{Shaded}

\end{tcolorbox}

\index{typedef keyword@\texttt{typedef} keyword!with anonymous structs@with anonymous \texttt{struct}s|)} \index{typedef keyword@\texttt{typedef} keyword!with structs@with \texttt{struct}s|)} \index{typedef keyword@\texttt{typedef} keyword!scoping rules|)}

\subsection{\texorpdfstring{\hl{\mbox{\texttt{typedef}}} dan Tipe Lain}{typedef dan Tipe Lain}}\label{typedef-dan-tipe-lain}

Bukan berarti memakai \hl{\mbox{\texttt{typedef}}} pada tipe sederhana seperti \hl{\mbox{\texttt{int}}} sama sekali tidak berguna\ldots{} ia membantu mengabstraksikan tipe agar lebih mudah diubah belakangan.

Misalnya, kalau \hl{\mbox{\texttt{float}}} tersebar di 100 zilyun tempat di kodemu, akan menyakitkan mengubah semuanya menjadi \hl{\mbox{\texttt{double}}} kalau belakangan kamu harus melakukannya karena suatu alasan.

Tapi kalau kamu sudah bersiap sedikit dengan:

\begin{tcolorbox}[enhanced jigsaw,left=3pt,breakable,colback=white!97.2549!black,boxrule=0.2mm]

\begin{Shaded}
\begin{Highlighting}[]
\KeywordTok{typedef} \DataTypeTok{float}\NormalTok{ app\_float}\OperatorTok{;}

\CommentTok{// dan}

\NormalTok{app\_float f1}\OperatorTok{,}\NormalTok{ f2}\OperatorTok{,}\NormalTok{ f3}\OperatorTok{;}
\end{Highlighting}
\end{Shaded}

\end{tcolorbox}

Maka kalau belakangan kamu ingin beralih ke tipe lain, seperti \hl{\mbox{\texttt{long\ double}}}, kamu cukup mengubah \hl{\mbox{\texttt{typedef}}}-nya:

\begin{tcolorbox}[enhanced jigsaw,left=3pt,breakable,colback=white!97.2549!black,boxrule=0.2mm]

\begin{Shaded}
\begin{Highlighting}[]
\CommentTok{//        voila!}
\CommentTok{//      |{-}{-}{-}{-}{-}{-}{-}{-}{-}|}
\KeywordTok{typedef} \DataTypeTok{long} \DataTypeTok{double}\NormalTok{ app\_float}\OperatorTok{;}

\CommentTok{// dan tidak perlu mengubah baris ini:}

\NormalTok{app\_float f1}\OperatorTok{,}\NormalTok{ f2}\OperatorTok{,}\NormalTok{ f3}\OperatorTok{;}  \CommentTok{// Sekarang semuanya long double}
\end{Highlighting}
\end{Shaded}

\end{tcolorbox}

\subsection{\texorpdfstring{\hl{\mbox{\texttt{typedef}}} dan Pointer}{typedef dan Pointer}}\label{typedef-dan-pointer}

\index{typedef keyword@\texttt{typedef} keyword!with pointers|(} Kamu bisa membuat tipe yang berupa pointer.

\begin{tcolorbox}[enhanced jigsaw,left=3pt,breakable,colback=white!97.2549!black,boxrule=0.2mm]

\begin{Shaded}
\begin{Highlighting}[]
\KeywordTok{typedef} \DataTypeTok{int} \OperatorTok{*}\NormalTok{intptr}\OperatorTok{;}

\DataTypeTok{int}\NormalTok{ a }\OperatorTok{=} \DecValTok{10}\OperatorTok{;}
\NormalTok{intptr x }\OperatorTok{=} \OperatorTok{\&}\NormalTok{a}\OperatorTok{;}  \CommentTok{// "intptr" bertipe "int*"}
\end{Highlighting}
\end{Shaded}

\end{tcolorbox}

Penulis benar-benar tidak menyukai praktik ini. Ia menyembunyikan fakta bahwa \hl{\mbox{\texttt{x}}} bertipe pointer karena tidak ada \hl{\mbox{\texttt{*}}} di deklarasinya.

Menurut penulis, lebih baik secara eksplisit menunjukkan bahwa kamu mendeklarasikan tipe pointer supaya developer lain bisa melihatnya dengan jelas dan tidak keliru mengira \hl{\mbox{\texttt{x}}} bukan tipe pointer.

Tapi menurut hitungan terakhir, katakanlah, 832.007 orang punya pendapat berbeda. \index{typedef keyword@\texttt{typedef} keyword!with pointers|)}

\subsection{\texorpdfstring{\hl{\mbox{\texttt{typedef}}} dan Kapitalisasi}{typedef dan Kapitalisasi}}\label{typedef-dan-kapitalisasi}

Penulis pernah melihat segala macam kapitalisasi pada \hl{\mbox{\texttt{typedef}}}.

\begin{tcolorbox}[enhanced jigsaw,left=3pt,breakable,colback=white!97.2549!black,boxrule=0.2mm]

\begin{Shaded}
\begin{Highlighting}[]
\KeywordTok{typedef} \KeywordTok{struct} \OperatorTok{\{}
    \DataTypeTok{int}\NormalTok{ x}\OperatorTok{,}\NormalTok{ y}\OperatorTok{;}
\OperatorTok{\}}\NormalTok{ my\_point}\OperatorTok{;}          \CommentTok{// snake case huruf kecil}

\KeywordTok{typedef} \KeywordTok{struct} \OperatorTok{\{}
    \DataTypeTok{int}\NormalTok{ x}\OperatorTok{,}\NormalTok{ y}\OperatorTok{;}
\OperatorTok{\}}\NormalTok{ MyPoint}\OperatorTok{;}          \CommentTok{// CamelCase}

\KeywordTok{typedef} \KeywordTok{struct} \OperatorTok{\{}
    \DataTypeTok{int}\NormalTok{ x}\OperatorTok{,}\NormalTok{ y}\OperatorTok{;}
\OperatorTok{\}}\NormalTok{ Mypoint}\OperatorTok{;}          \CommentTok{// Huruf besar di awal}

\KeywordTok{typedef} \KeywordTok{struct} \OperatorTok{\{}
    \DataTypeTok{int}\NormalTok{ x}\OperatorTok{,}\NormalTok{ y}\OperatorTok{;}
\OperatorTok{\}}\NormalTok{ MY\_POINT}\OperatorTok{;}          \CommentTok{// SNAKE CASE HURUF BESAR}
\end{Highlighting}
\end{Shaded}

\end{tcolorbox}

Spesifikasi C11 tidak menentukan salah satunya, dan menunjukkan contoh dalam huruf besar semua maupun huruf kecil semua.

K\&R2 terutama memakai huruf besar di awal, tapi menunjukkan beberapa contoh dalam huruf besar dan snake case (dengan \hl{\mbox{\texttt{\_t}}}).

Kalau kamu punya style guide yang dipakai, ikuti saja. Kalau tidak punya, pilih satu dan ikuti.

\section{\texorpdfstring{Array dan \hl{\mbox{\texttt{typedef}}}}{Array dan typedef}}\label{array-dan-typedef}

\index{typedef keyword@\texttt{typedef} keyword!with arrays|(} Sintaksnya agak aneh, dan menurut pengalaman penulis ini jarang terlihat, tapi kamu bisa me-\hl{\mbox{\texttt{typedef}}} array berisi sejumlah item.

\begin{tcolorbox}[enhanced jigsaw,left=3pt,breakable,colback=white!97.2549!black,boxrule=0.2mm]

\begin{Shaded}
\begin{Highlighting}[]
\CommentTok{// Jadikan tipe five\_ints sebagai array berisi 5 int}
\KeywordTok{typedef} \DataTypeTok{int}\NormalTok{ five\_ints}\OperatorTok{[}\DecValTok{5}\OperatorTok{];}

\NormalTok{five\_ints x }\OperatorTok{=} \OperatorTok{\{}\DecValTok{11}\OperatorTok{,} \DecValTok{22}\OperatorTok{,} \DecValTok{33}\OperatorTok{,} \DecValTok{44}\OperatorTok{,} \DecValTok{55}\OperatorTok{\};}
\end{Highlighting}
\end{Shaded}

\end{tcolorbox}

Penulis tidak menyukainya karena ia menyembunyikan sifat array dari variabelnya, tapi itu mungkin dilakukan. \index{typedef keyword@\texttt{typedef} keyword!with arrays|)} \index{typedef keyword@\texttt{typedef} keyword|)}

\chapter{Pointer II: Aritmetika}\label{pointers2}

\index{Pointers!arithmetic|(}

Saatnya menyelam lebih dalam dengan sejumlah topik pointer baru! Kalau kamu belum paham betul soal pointer, \hyperref[pointers]{lihat dulu bagian pertama panduan ini tentang topik tersebut}.

\section{Aritmetika Pointer}\label{aritmetika-pointer}

Ternyata pointer bisa dioperasikan secara matematis, terutama penjumlahan dan pengurangan.

Tapi apa artinya kalau itu dilakukan?

Singkatnya, kalau kamu punya pointer ke suatu tipe, menambahkan satu pada pointer itu akan memindahkannya ke item bertipe sama yang berada tepat setelahnya di memori.

\textbf{Penting} untuk diingat bahwa saat kita menggeser pointer dan melihat tempat-tempat berbeda di memori, kita harus selalu memastikan pointer itu menunjuk ke tempat yang valid sebelum di-dereference. Kalau kita nyasar entah ke mana lalu mencoba melihat isinya, perilakunya undefined, dan crash adalah hasil yang umum.

Ini agak seperti masalah ayam-dan-telur dengan \hyperref[arraypointerequiv]{Kesetaraan Array/Pointer, di bawah}, tapi kita coba saja.

\subsection{Menambah Pointer}\label{menambah-pointer}

Pertama, kita ambil sebuah array berisi angka.

\begin{tcolorbox}[enhanced jigsaw,left=3pt,breakable,colback=white!97.2549!black,boxrule=0.2mm]

\begin{Shaded}
\begin{Highlighting}[]
\DataTypeTok{int}\NormalTok{ a}\OperatorTok{[}\DecValTok{5}\OperatorTok{]} \OperatorTok{=} \OperatorTok{\{}\DecValTok{11}\OperatorTok{,} \DecValTok{22}\OperatorTok{,} \DecValTok{33}\OperatorTok{,} \DecValTok{44}\OperatorTok{,} \DecValTok{55}\OperatorTok{\};}
\end{Highlighting}
\end{Shaded}

\end{tcolorbox}

Lalu kita ambil pointer ke elemen pertama array itu:

\begin{tcolorbox}[enhanced jigsaw,left=3pt,breakable,colback=white!97.2549!black,boxrule=0.2mm]

\begin{Shaded}
\begin{Highlighting}[]
\DataTypeTok{int}\NormalTok{ a}\OperatorTok{[}\DecValTok{5}\OperatorTok{]} \OperatorTok{=} \OperatorTok{\{}\DecValTok{11}\OperatorTok{,} \DecValTok{22}\OperatorTok{,} \DecValTok{33}\OperatorTok{,} \DecValTok{44}\OperatorTok{,} \DecValTok{55}\OperatorTok{\};}

\DataTypeTok{int} \OperatorTok{*}\NormalTok{p }\OperatorTok{=} \OperatorTok{\&}\NormalTok{a}\OperatorTok{[}\DecValTok{0}\OperatorTok{];}  \CommentTok{// Atau "int *p = a;" sama saja hasilnya}
\end{Highlighting}
\end{Shaded}

\end{tcolorbox}

Lalu kita cetak nilai di sana dengan men-dereference pointer-nya:

\begin{tcolorbox}[enhanced jigsaw,left=3pt,breakable,colback=white!97.2549!black,boxrule=0.2mm]

\begin{Shaded}
\begin{Highlighting}[]
\NormalTok{printf}\OperatorTok{(}\StringTok{"}\SpecialCharTok{\%d\textbackslash{}n}\StringTok{"}\OperatorTok{,} \OperatorTok{*}\NormalTok{p}\OperatorTok{);}  \CommentTok{// Mencetak 11}
\end{Highlighting}
\end{Shaded}

\end{tcolorbox}

Sekarang kita pakai aritmetika pointer untuk mencetak elemen berikutnya di array, yaitu yang berindeks 1:

\begin{tcolorbox}[enhanced jigsaw,left=3pt,breakable,colback=white!97.2549!black,boxrule=0.2mm]

\begin{Shaded}
\begin{Highlighting}[]
\NormalTok{printf}\OperatorTok{(}\StringTok{"}\SpecialCharTok{\%d\textbackslash{}n}\StringTok{"}\OperatorTok{,} \OperatorTok{*(}\NormalTok{p }\OperatorTok{+} \DecValTok{1}\OperatorTok{));}  \CommentTok{// Mencetak 22!!}
\end{Highlighting}
\end{Shaded}

\end{tcolorbox}

Apa yang terjadi di situ? C tahu bahwa \hl{\mbox{\texttt{p}}} adalah pointer ke \hl{\mbox{\texttt{int}}}. Jadi ia tahu \hl{\mbox{\texttt{sizeof}}} sebuah \hl{\mbox{\texttt{int}}}\footnote{Ingat, operator \hl{\mbox{\texttt{sizeof}}} memberi tahu ukuran sebuah object di memori dalam byte.} dan tahu harus melompati sebanyak itu byte untuk sampai ke \hl{\mbox{\texttt{int}}} setelah yang pertama!

Bahkan, contoh sebelumnya bisa ditulis dengan dua cara yang setara:

\begin{tcolorbox}[enhanced jigsaw,left=3pt,breakable,colback=white!97.2549!black,boxrule=0.2mm]

\begin{Shaded}
\begin{Highlighting}[]
\NormalTok{printf}\OperatorTok{(}\StringTok{"}\SpecialCharTok{\%d\textbackslash{}n}\StringTok{"}\OperatorTok{,} \OperatorTok{*}\NormalTok{p}\OperatorTok{);}        \CommentTok{// Mencetak 11}
\NormalTok{printf}\OperatorTok{(}\StringTok{"}\SpecialCharTok{\%d\textbackslash{}n}\StringTok{"}\OperatorTok{,} \OperatorTok{*(}\NormalTok{p }\OperatorTok{+} \DecValTok{0}\OperatorTok{));}  \CommentTok{// Mencetak 11}
\end{Highlighting}
\end{Shaded}

\end{tcolorbox}

karena menambahkan \hl{\mbox{\texttt{0}}} pada pointer menghasilkan pointer yang sama.

Mari kita pikirkan manfaatnya. Kita bisa menelusuri elemen array dengan cara ini, tanpa memakai array:

\begin{tcolorbox}[enhanced jigsaw,left=3pt,breakable,colback=white!97.2549!black,boxrule=0.2mm]

\begin{Shaded}
\begin{Highlighting}[]
\DataTypeTok{int}\NormalTok{ a}\OperatorTok{[}\DecValTok{5}\OperatorTok{]} \OperatorTok{=} \OperatorTok{\{}\DecValTok{11}\OperatorTok{,} \DecValTok{22}\OperatorTok{,} \DecValTok{33}\OperatorTok{,} \DecValTok{44}\OperatorTok{,} \DecValTok{55}\OperatorTok{\};}

\DataTypeTok{int} \OperatorTok{*}\NormalTok{p }\OperatorTok{=} \OperatorTok{\&}\NormalTok{a}\OperatorTok{[}\DecValTok{0}\OperatorTok{];}  \CommentTok{// Atau "int *p = a;" sama saja hasilnya}

\ControlFlowTok{for} \OperatorTok{(}\DataTypeTok{int}\NormalTok{ i }\OperatorTok{=} \DecValTok{0}\OperatorTok{;}\NormalTok{ i }\OperatorTok{\textless{}} \DecValTok{5}\OperatorTok{;}\NormalTok{ i}\OperatorTok{++)} \OperatorTok{\{}
\NormalTok{    printf}\OperatorTok{(}\StringTok{"}\SpecialCharTok{\%d\textbackslash{}n}\StringTok{"}\OperatorTok{,} \OperatorTok{*(}\NormalTok{p }\OperatorTok{+}\NormalTok{ i}\OperatorTok{));}  \CommentTok{// Sama dengan p[i]!}
\OperatorTok{\}}
\end{Highlighting}
\end{Shaded}

\end{tcolorbox}

Dan hasilnya sama seperti kalau kita memakai notasi array! Oooo! Makin dekat ke soal kesetaraan array/pointer itu! Nanti dibahas lebih lanjut di bab ini.

Tapi sebenarnya apa yang terjadi di sini? Bagaimana cara kerjanya?

Ingat dari awal tadi bahwa memori itu seperti array besar, tempat sebuah byte disimpan di setiap indeks array?

Dan indeks array ke dalam memori itu punya beberapa nama:

\begin{itemize}
\tightlist
\item
  Indeks ke dalam memori
\item
  Lokasi
\item
  Alamat
\item
  \emph{Pointer!}
\end{itemize}

Jadi pointer adalah indeks ke dalam memori, di suatu tempat.

Sebagai contoh acak, katakanlah angka 3490 disimpan di alamat (``indeks'') 23.237.489.202. Kalau kita punya pointer \hl{\mbox{\texttt{int}}} ke 3490 itu, nilai pointer itu adalah 23.237.489.202\ldots{} karena pointer adalah alamat memori. Kata berbeda untuk hal yang sama.

Sekarang katakanlah ada angka lain, 4096, yang disimpan tepat setelah 3490 di alamat 23.237.489.210 (lebih tinggi 8 dari 3490 karena setiap \hl{\mbox{\texttt{int}}} dalam contoh ini berukuran 8 byte).

Kalau kita tambahkan \hl{\mbox{\texttt{1}}} pada pointer itu, ia sebenarnya melompat maju sejauh \hl{\mbox{\texttt{sizeof(int)}}} byte ke \hl{\mbox{\texttt{int}}} berikutnya. Ia tahu harus melompat sejauh itu karena ia pointer \hl{\mbox{\texttt{int}}}. Kalau ia pointer \hl{\mbox{\texttt{float}}}, ia akan melompat sejauh \hl{\mbox{\texttt{sizeof(float)}}} byte untuk sampai ke float berikutnya!

Jadi kamu bisa melihat \hl{\mbox{\texttt{int}}} berikutnya dengan menambahkan \hl{\mbox{\texttt{1}}} pada pointer, yang sesudahnya dengan menambahkan \hl{\mbox{\texttt{2}}}, dan seterusnya.

\subsection{Mengubah Pointer}\label{mengubah-pointer}

Tadi kita sudah melihat bahwa sebuah integer bisa ditambahkan ke pointer. Sekarang, mari kita \emph{ubah pointer itu sendiri}.

Nilai integer bisa langsung ditambahkan ke (atau dikurangkan dari) pointer mana pun!

Kita ulangi contoh tadi dengan beberapa perubahan. Pertama, kita tambahkan \hl{\mbox{\texttt{999}}} di akhir deretan angka sebagai sentinel value. Ini memberi tahu kita di mana data berakhir.

\begin{tcolorbox}[enhanced jigsaw,left=3pt,breakable,colback=white!97.2549!black,boxrule=0.2mm]

\begin{Shaded}
\begin{Highlighting}[]
\DataTypeTok{int}\NormalTok{ a}\OperatorTok{[]} \OperatorTok{=} \OperatorTok{\{}\DecValTok{11}\OperatorTok{,} \DecValTok{22}\OperatorTok{,} \DecValTok{33}\OperatorTok{,} \DecValTok{44}\OperatorTok{,} \DecValTok{55}\OperatorTok{,} \DecValTok{999}\OperatorTok{\};}  \CommentTok{// Tambahkan 999 di sini sebagai sentinel}

\DataTypeTok{int} \OperatorTok{*}\NormalTok{p }\OperatorTok{=} \OperatorTok{\&}\NormalTok{a}\OperatorTok{[}\DecValTok{0}\OperatorTok{];}  \CommentTok{// p menunjuk ke 11}
\end{Highlighting}
\end{Shaded}

\end{tcolorbox}

Dan \hl{\mbox{\texttt{p}}} menunjuk ke elemen berindeks \hl{\mbox{\texttt{0}}} dari \hl{\mbox{\texttt{a}}}, yaitu \hl{\mbox{\texttt{11}}}, persis seperti tadi.

Sekarang---kita mulai \emph{meng-increment} \hl{\mbox{\texttt{p}}} supaya menunjuk ke elemen-elemen array berikutnya. Kita lakukan ini sampai \hl{\mbox{\texttt{p}}} menunjuk ke \hl{\mbox{\texttt{999}}}; artinya, sampai \hl{\mbox{\texttt{*p\ ==\ 999}}}:

\begin{tcolorbox}[enhanced jigsaw,left=3pt,breakable,colback=white!97.2549!black,boxrule=0.2mm]

\begin{Shaded}
\begin{Highlighting}[]
\ControlFlowTok{while} \OperatorTok{(*}\NormalTok{p }\OperatorTok{!=} \DecValTok{999}\OperatorTok{)} \OperatorTok{\{}       \CommentTok{// Selama yang ditunjuk p bukan 999}
\NormalTok{    printf}\OperatorTok{(}\StringTok{"}\SpecialCharTok{\%d\textbackslash{}n}\StringTok{"}\OperatorTok{,} \OperatorTok{*}\NormalTok{p}\OperatorTok{);}   \CommentTok{// Cetak}
\NormalTok{    p}\OperatorTok{++;}                  \CommentTok{// Geser p supaya menunjuk ke int berikutnya!}
\OperatorTok{\}}
\end{Highlighting}
\end{Shaded}

\end{tcolorbox}

Gila, kan?

Kalau dijalankan, mula-mula \hl{\mbox{\texttt{p}}} menunjuk ke \hl{\mbox{\texttt{11}}}. Lalu \hl{\mbox{\texttt{p}}} di-increment, dan ia menunjuk ke \hl{\mbox{\texttt{22}}}, lalu sekali lagi, ia menunjuk ke \hl{\mbox{\texttt{33}}}. Begitu seterusnya, sampai ia menunjuk ke \hl{\mbox{\texttt{999}}} dan kita berhenti.

\subsection{Mengurangi Pointer}\label{mengurangi-pointer}

\index{Pointers!subtracting|(} Kamu bisa mengurangi pointer dengan sebuah nilai untuk sampai ke alamat yang lebih awal, sama seperti tadi kita menambahkannya.

Tapi kita juga bisa mengurangkan dua pointer untuk mengetahui selisihnya, misalnya untuk menghitung berapa banyak \hl{\mbox{\texttt{int}}} yang ada di antara dua \hl{\mbox{\texttt{int*}}}. Syaratnya, ini hanya berlaku di dalam satu array\footnote{Atau string, yang sebenarnya array berisi \hl{\mbox{\texttt{char}}}. Agak aneh, kamu juga boleh punya pointer yang merujuk ke \emph{satu elemen setelah} ujung array tanpa masalah dan tetap bisa menghitungnya. Hanya saja, pointer itu tidak boleh di-dereference selagi berada di sana.}---kalau pointer-nya menunjuk ke hal lain, hasilnya undefined behavior.

Ingat bahwa string di C adalah \hl{\mbox{\texttt{char*}}}? Mari kita lihat apakah ini bisa dipakai untuk menulis varian lain dari \hl{\mbox{\texttt{strlen()}}} untuk menghitung panjang string dengan pengurangan pointer.

Idenya: kalau kita punya pointer ke awal string, kita bisa mencari pointer ke akhir string dengan menelusuri maju sampai menemukan karakter \hl{\mbox{\texttt{NUL}}}.

Dan kalau kita punya pointer ke awal string, serta sudah menghitung pointer ke akhir string, kita tinggal mengurangkan kedua pointer itu untuk mendapatkan panjangnya!

\begin{tcolorbox}[enhanced jigsaw,left=3pt,breakable,colback=white!97.2549!black,boxrule=0.2mm]

\begin{Shaded}
\begin{Highlighting}[numbers=left,,]
\PreprocessorTok{\#include }\ImportTok{\textless{}stdio.h\textgreater{}}

\DataTypeTok{int}\NormalTok{ my\_strlen}\OperatorTok{(}\DataTypeTok{char} \OperatorTok{*}\NormalTok{s}\OperatorTok{)}
\OperatorTok{\{}
    \CommentTok{// Mulai menelusuri dari awal string}
    \DataTypeTok{char} \OperatorTok{*}\NormalTok{p }\OperatorTok{=}\NormalTok{ s}\OperatorTok{;}

    \CommentTok{// Telusuri sampai menemukan karakter NUL}
    \ControlFlowTok{while} \OperatorTok{(*}\NormalTok{p }\OperatorTok{!=} \CharTok{\textquotesingle{}}\SpecialCharTok{\textbackslash{}0}\CharTok{\textquotesingle{}}\OperatorTok{)}
\NormalTok{        p}\OperatorTok{++;}

    \CommentTok{// Kembalikan selisih pointer}
    \ControlFlowTok{return}\NormalTok{ p }\OperatorTok{{-}}\NormalTok{ s}\OperatorTok{;}
\OperatorTok{\}}

\DataTypeTok{int}\NormalTok{ main}\OperatorTok{(}\DataTypeTok{void}\OperatorTok{)}
\OperatorTok{\{}
\NormalTok{    printf}\OperatorTok{(}\StringTok{"}\SpecialCharTok{\%d\textbackslash{}n}\StringTok{"}\OperatorTok{,}\NormalTok{ my\_strlen}\OperatorTok{(}\StringTok{"Hello, world!"}\OperatorTok{));}  \CommentTok{// Mencetak "13"}
\OperatorTok{\}}
\end{Highlighting}
\end{Shaded}

\end{tcolorbox}

Ingat, pengurangan pointer hanya boleh dilakukan antara dua pointer yang menunjuk ke array yang sama! \index{Pointers!subtracting|)}

\section{Kesetaraan Array/Pointer}\label{arraypointerequiv}

\index{Pointers!array equivalence|(} Akhirnya kita siap membahas ini! Kita sudah melihat banyak contoh tempat kita mencampur notasi array, tapi sekarang kita sampaikan \emph{rumus dasar kesetaraan array/pointer}:

\begin{tcolorbox}[enhanced jigsaw,left=3pt,breakable,colback=white!97.2549!black,boxrule=0.2mm]

\begin{Shaded}
\begin{Highlighting}[]
\NormalTok{a}\OperatorTok{[}\NormalTok{b}\OperatorTok{]} \OperatorTok{==} \OperatorTok{*(}\NormalTok{a }\OperatorTok{+}\NormalTok{ b}\OperatorTok{)}
\end{Highlighting}
\end{Shaded}

\end{tcolorbox}

Pelajari itu! Keduanya setara dan bisa dipakai bergantian!

Ini sedikit disederhanakan, karena pada contoh di atas \hl{\mbox{\texttt{a}}} dan \hl{\mbox{\texttt{b}}} sama-sama bisa berupa ekspresi, dan mungkin kita perlu beberapa tanda kurung tambahan untuk memaksa urutan operasi kalau ekspresinya rumit.

Spesifikasi, seperti biasa, tegas soal ini, menyatakan (di C11 §6.5.2.1¶2):

\begin{tcolorbox}[blanker,colback=white,enhanced jigsaw,breakable,arc=0pt,outer arc=0pt,grow to left by=-1cm,enlarge top by=0.1cm,enlarge bottom by=0.1cm,borderline west={1.5pt}{-.3cm}{black}]

\begin{quote}
\hl{\mbox{\texttt{E1{[}E2{]}}}} is identical to \hl{\mbox{\texttt{(*((E1)+(E2)))}}}
\end{quote}

\end{tcolorbox}

tapi itu agak lebih sulit dicerna. Pokoknya pastikan kamu menambahkan tanda kurung kalau ekspresinya rumit, supaya semua perhitunganmu terjadi dalam urutan yang benar.

Artinya, kita bisa \emph{memutuskan} mau memakai notasi array atau notasi pointer untuk array atau pointer apa pun (dengan asumsi pointer-nya menunjuk ke elemen sebuah array).

Mari kita pakai sebuah array dan pointer dengan notasi array maupun notasi pointer:

\begin{tcolorbox}[enhanced jigsaw,left=3pt,breakable,colback=white!97.2549!black,boxrule=0.2mm]

\begin{Shaded}
\begin{Highlighting}[numbers=left,,]
\PreprocessorTok{\#include }\ImportTok{\textless{}stdio.h\textgreater{}}

\DataTypeTok{int}\NormalTok{ main}\OperatorTok{(}\DataTypeTok{void}\OperatorTok{)}
\OperatorTok{\{}
    \DataTypeTok{int}\NormalTok{ a}\OperatorTok{[]} \OperatorTok{=} \OperatorTok{\{}\DecValTok{11}\OperatorTok{,} \DecValTok{22}\OperatorTok{,} \DecValTok{33}\OperatorTok{,} \DecValTok{44}\OperatorTok{,} \DecValTok{55}\OperatorTok{\};}

    \DataTypeTok{int} \OperatorTok{*}\NormalTok{p }\OperatorTok{=}\NormalTok{ a}\OperatorTok{;}  \CommentTok{// p menunjuk ke elemen pertama a, yaitu 11}

    \CommentTok{// Cetak semua elemen array dengan berbagai cara:}

    \ControlFlowTok{for} \OperatorTok{(}\DataTypeTok{int}\NormalTok{ i }\OperatorTok{=} \DecValTok{0}\OperatorTok{;}\NormalTok{ i }\OperatorTok{\textless{}} \DecValTok{5}\OperatorTok{;}\NormalTok{ i}\OperatorTok{++)}
\NormalTok{        printf}\OperatorTok{(}\StringTok{"}\SpecialCharTok{\%d\textbackslash{}n}\StringTok{"}\OperatorTok{,}\NormalTok{ a}\OperatorTok{[}\NormalTok{i}\OperatorTok{]);}      \CommentTok{// Notasi array dengan a}

    \ControlFlowTok{for} \OperatorTok{(}\DataTypeTok{int}\NormalTok{ i }\OperatorTok{=} \DecValTok{0}\OperatorTok{;}\NormalTok{ i }\OperatorTok{\textless{}} \DecValTok{5}\OperatorTok{;}\NormalTok{ i}\OperatorTok{++)}
\NormalTok{        printf}\OperatorTok{(}\StringTok{"}\SpecialCharTok{\%d\textbackslash{}n}\StringTok{"}\OperatorTok{,}\NormalTok{ p}\OperatorTok{[}\NormalTok{i}\OperatorTok{]);}      \CommentTok{// Notasi array dengan p}

    \ControlFlowTok{for} \OperatorTok{(}\DataTypeTok{int}\NormalTok{ i }\OperatorTok{=} \DecValTok{0}\OperatorTok{;}\NormalTok{ i }\OperatorTok{\textless{}} \DecValTok{5}\OperatorTok{;}\NormalTok{ i}\OperatorTok{++)}
\NormalTok{        printf}\OperatorTok{(}\StringTok{"}\SpecialCharTok{\%d\textbackslash{}n}\StringTok{"}\OperatorTok{,} \OperatorTok{*(}\NormalTok{a }\OperatorTok{+}\NormalTok{ i}\OperatorTok{));}  \CommentTok{// Notasi pointer dengan a}

    \ControlFlowTok{for} \OperatorTok{(}\DataTypeTok{int}\NormalTok{ i }\OperatorTok{=} \DecValTok{0}\OperatorTok{;}\NormalTok{ i }\OperatorTok{\textless{}} \DecValTok{5}\OperatorTok{;}\NormalTok{ i}\OperatorTok{++)}
\NormalTok{        printf}\OperatorTok{(}\StringTok{"}\SpecialCharTok{\%d\textbackslash{}n}\StringTok{"}\OperatorTok{,} \OperatorTok{*(}\NormalTok{p }\OperatorTok{+}\NormalTok{ i}\OperatorTok{));}  \CommentTok{// Notasi pointer dengan p}

    \ControlFlowTok{for} \OperatorTok{(}\DataTypeTok{int}\NormalTok{ i }\OperatorTok{=} \DecValTok{0}\OperatorTok{;}\NormalTok{ i }\OperatorTok{\textless{}} \DecValTok{5}\OperatorTok{;}\NormalTok{ i}\OperatorTok{++)}
\NormalTok{        printf}\OperatorTok{(}\StringTok{"}\SpecialCharTok{\%d\textbackslash{}n}\StringTok{"}\OperatorTok{,} \OperatorTok{*(}\NormalTok{p}\OperatorTok{++));}    \CommentTok{// Menggeser pointer p}
        \CommentTok{//printf("\%d\textbackslash{}n", *(a++));    // Menggeser variabel array a{-}{-}ERROR!}
\OperatorTok{\}}
\end{Highlighting}
\end{Shaded}

\end{tcolorbox}

Jadi kamu bisa lihat bahwa pada umumnya, kalau kamu punya variabel array, kamu bisa memakai notasi pointer maupun notasi array untuk mengakses elemennya. Begitu juga dengan variabel pointer.

Satu perbedaan besarnya: sebuah pointer bisa \emph{diubah} agar menunjuk ke alamat lain, tapi variabel array tidak bisa. Dengan kata lain, kamu sama sekali tidak bisa menugaskan nilai ke variabel array---hanya ke elemen-elemen individual array itu.

Kalau kamu benar-benar ingin menyalin satu array ke array lain, kamu harus memakai fungsi seperti \hl{\mbox{\texttt{memcpy()}}} (atau sebuah loop).

\subsection{Kesetaraan Array/Pointer pada Pemanggilan Fungsi}\label{kesetaraan-arraypointer-pada-pemanggilan-fungsi}

Di sinilah konsep ini paling sering kamu temui, pasti.

Kalau ada fungsi yang menerima argumen pointer, misalnya:

\begin{tcolorbox}[enhanced jigsaw,left=3pt,breakable,colback=white!97.2549!black,boxrule=0.2mm]

\begin{Shaded}
\begin{Highlighting}[]
\DataTypeTok{int}\NormalTok{ my\_strlen}\OperatorTok{(}\DataTypeTok{char} \OperatorTok{*}\NormalTok{s}\OperatorTok{)}
\end{Highlighting}
\end{Shaded}

\end{tcolorbox}

artinya kamu bisa mengoper array maupun pointer ke fungsi ini, dan keduanya berfungsi!

\begin{tcolorbox}[enhanced jigsaw,left=3pt,breakable,colback=white!97.2549!black,boxrule=0.2mm]

\begin{Shaded}
\begin{Highlighting}[]
\DataTypeTok{char}\NormalTok{ s}\OperatorTok{[]} \OperatorTok{=} \StringTok{"Antelopes"}\OperatorTok{;}
\DataTypeTok{char} \OperatorTok{*}\NormalTok{t }\OperatorTok{=} \StringTok{"Wombats"}\OperatorTok{;}

\NormalTok{printf}\OperatorTok{(}\StringTok{"}\SpecialCharTok{\%d\textbackslash{}n}\StringTok{"}\OperatorTok{,}\NormalTok{ my\_strlen}\OperatorTok{(}\NormalTok{s}\OperatorTok{));}  \CommentTok{// Berfungsi!}
\NormalTok{printf}\OperatorTok{(}\StringTok{"}\SpecialCharTok{\%d\textbackslash{}n}\StringTok{"}\OperatorTok{,}\NormalTok{ my\_strlen}\OperatorTok{(}\NormalTok{t}\OperatorTok{));}  \CommentTok{// Berfungsi juga!}
\end{Highlighting}
\end{Shaded}

\end{tcolorbox}

Dan itu juga alasan kenapa dua signature fungsi ini setara:

\begin{tcolorbox}[enhanced jigsaw,left=3pt,breakable,colback=white!97.2549!black,boxrule=0.2mm]

\begin{Shaded}
\begin{Highlighting}[]
\DataTypeTok{int}\NormalTok{ my\_strlen}\OperatorTok{(}\DataTypeTok{char} \OperatorTok{*}\NormalTok{s}\OperatorTok{)}    \CommentTok{// Berfungsi!}
\DataTypeTok{int}\NormalTok{ my\_strlen}\OperatorTok{(}\DataTypeTok{char}\NormalTok{ s}\OperatorTok{[])}   \CommentTok{// Berfungsi juga!}
\end{Highlighting}
\end{Shaded}

\end{tcolorbox}

\index{Pointers!array equivalence|)}

\section{\texorpdfstring{Pointer \hl{\mbox{\texttt{void}}}}{Pointer void}}\label{pointer-void}

\index{void* void pointer@\texttt{void*} void pointer|(}

Kamu sudah melihat kata kunci \hl{\mbox{\texttt{void}}} dipakai pada fungsi untuk menandakan tidak ada parameter atau tidak ada nilai kembalian, tapi ini makhluk yang sama sekali berbeda dan tidak berhubungan.

\hl{\mbox{\texttt{void*}}} jelas adalah pointer ke sebuah \emph{benda} yang ada. Tapi bagian \hl{\mbox{\texttt{void}}}-nya menandakan bahwa kita tidak tahu \emph{tipe} benda itu. Dan kadang, percaya atau tidak, itu sangat berguna. Dengannya kita bisa menulis kode yang sedikit lebih agnostik terhadap tipe, fleksibilitas yang menyenangkan di bahasa bertipe seperti C.

Pada dasarnya ada dua kasus pemakaian---mari kita lihat dan coba mengurangi misterinya.

\index{memcpy() function@\texttt{memcpy()} function|(}

\begin{enumerate}
\def\labelenumi{\arabic{enumi}.}
\item
  Sebuah fungsi akan bekerja pada sesuatu byte demi byte. Misalnya, \hl{\mbox{\texttt{memcpy()}}} menyalin byte memori dari satu pointer ke pointer lain, tapi pointer-pointer itu bisa menunjuk ke tipe apa pun. \hl{\mbox{\texttt{memcpy()}}} memanfaatkan kenyataan bahwa kalau kita menelusuri \hl{\mbox{\texttt{char*}}}, kita menelusuri byte-byte sebuah object apa pun tipenya. Pembahasan lebih lanjut ada di subbagian \hyperref[multibyte-values]{Nilai Multibyte}.
\item
  Fungsi lain memanggil fungsi yang kamu oper kepadanya (sebuah callback), dan ia mengoper data kepadamu. Kamu tahu tipe datanya, tapi fungsi yang memanggilmu tidak. Jadi ia mengoper \hl{\mbox{\texttt{void*}}}---karena tidak tahu tipenya---dan kamu mengonversinya ke tipe yang kamu butuhkan. Fungsi bawaan \href{https://beej.us/guide/bgclr/html/split/stdlib.html\#man-qsort}{\hl{\mbox{\texttt{qsort()}}}}\footnote{https://beej.us/guide/bgclr/html/split/stdlib.html\#man-qsort} dan \href{https://beej.us/guide/bgclr/html/split/stdlib.html\#man-bsearch}{\hl{\mbox{\texttt{bsearch()}}}}\footnote{https://beej.us/guide/bgclr/html/split/stdlib.html\#man-bsearch} memakai teknik ini.
\end{enumerate}

Mari kita lihat contohnya, fungsi bawaan \hl{\mbox{\texttt{memcpy()}}}:

\begin{tcolorbox}[enhanced jigsaw,left=3pt,breakable,colback=white!97.2549!black,boxrule=0.2mm]

\begin{Shaded}
\begin{Highlighting}[]
\DataTypeTok{void} \OperatorTok{*}\NormalTok{memcpy}\OperatorTok{(}\DataTypeTok{void} \OperatorTok{*}\NormalTok{s1}\OperatorTok{,} \DataTypeTok{void} \OperatorTok{*}\NormalTok{s2}\OperatorTok{,} \DataTypeTok{size\_t}\NormalTok{ n}\OperatorTok{);}
\end{Highlighting}
\end{Shaded}

\end{tcolorbox}

Fungsi ini menyalin \hl{\mbox{\texttt{n}}} byte memori mulai dari alamat \hl{\mbox{\texttt{s2}}} ke memori yang dimulai di alamat \hl{\mbox{\texttt{s1}}}.

Tapi lihat! \hl{\mbox{\texttt{s1}}} dan \hl{\mbox{\texttt{s2}}} adalah \hl{\mbox{\texttt{void*}}}! Kenapa? Apa artinya? Mari kita coba contoh-contoh lain.

Misalnya, kita bisa menyalin string dengan \hl{\mbox{\texttt{memcpy()}}} (walaupun \hl{\mbox{\texttt{strcpy()}}} lebih pantas untuk string):

\begin{tcolorbox}[enhanced jigsaw,left=3pt,breakable,colback=white!97.2549!black,boxrule=0.2mm]

\begin{Shaded}
\begin{Highlighting}[numbers=left,,]
\PreprocessorTok{\#include }\ImportTok{\textless{}stdio.h\textgreater{}}
\PreprocessorTok{\#include }\ImportTok{\textless{}string.h\textgreater{}}

\DataTypeTok{int}\NormalTok{ main}\OperatorTok{(}\DataTypeTok{void}\OperatorTok{)}
\OperatorTok{\{}
    \DataTypeTok{char}\NormalTok{ s}\OperatorTok{[]} \OperatorTok{=} \StringTok{"Goats!"}\OperatorTok{;}
    \DataTypeTok{char}\NormalTok{ t}\OperatorTok{[}\DecValTok{100}\OperatorTok{];}

\NormalTok{    memcpy}\OperatorTok{(}\NormalTok{t}\OperatorTok{,}\NormalTok{ s}\OperatorTok{,} \DecValTok{7}\OperatorTok{);}  \CommentTok{// Salin 7 byte{-}{-}termasuk terminator NUL!}

\NormalTok{    printf}\OperatorTok{(}\StringTok{"}\SpecialCharTok{\%s\textbackslash{}n}\StringTok{"}\OperatorTok{,}\NormalTok{ t}\OperatorTok{);}  \CommentTok{// "Goats!"}
\OperatorTok{\}}
\end{Highlighting}
\end{Shaded}

\end{tcolorbox}

Atau kita bisa menyalin beberapa \hl{\mbox{\texttt{int}}}:

\begin{tcolorbox}[enhanced jigsaw,left=3pt,breakable,colback=white!97.2549!black,boxrule=0.2mm]

\begin{Shaded}
\begin{Highlighting}[numbers=left,,]
\PreprocessorTok{\#include }\ImportTok{\textless{}stdio.h\textgreater{}}
\PreprocessorTok{\#include }\ImportTok{\textless{}string.h\textgreater{}}

\DataTypeTok{int}\NormalTok{ main}\OperatorTok{(}\DataTypeTok{void}\OperatorTok{)}
\OperatorTok{\{}
    \DataTypeTok{int}\NormalTok{ a}\OperatorTok{[]} \OperatorTok{=} \OperatorTok{\{}\DecValTok{11}\OperatorTok{,} \DecValTok{22}\OperatorTok{,} \DecValTok{33}\OperatorTok{\};}
    \DataTypeTok{int}\NormalTok{ b}\OperatorTok{[}\DecValTok{3}\OperatorTok{];}

\NormalTok{    memcpy}\OperatorTok{(}\NormalTok{b}\OperatorTok{,}\NormalTok{ a}\OperatorTok{,} \DecValTok{3} \OperatorTok{*} \KeywordTok{sizeof}\OperatorTok{(}\DataTypeTok{int}\OperatorTok{));}  \CommentTok{// Salin data 3 int}

\NormalTok{    printf}\OperatorTok{(}\StringTok{"}\SpecialCharTok{\%d\textbackslash{}n}\StringTok{"}\OperatorTok{,}\NormalTok{ b}\OperatorTok{[}\DecValTok{1}\OperatorTok{]);}  \CommentTok{// 22}
\OperatorTok{\}}
\end{Highlighting}
\end{Shaded}

\end{tcolorbox}

Yang ini agak liar---kamu lihat apa yang kita lakukan dengan \hl{\mbox{\texttt{memcpy()}}} di situ? Kita menyalin data dari \hl{\mbox{\texttt{a}}} ke \hl{\mbox{\texttt{b}}}, tapi kita harus menentukan berapa \emph{byte} yang disalin, dan sebuah \hl{\mbox{\texttt{int}}} lebih dari satu byte.

Oke, kalau begitu---berapa byte yang dipakai sebuah \hl{\mbox{\texttt{int}}}? Jawabannya: tergantung sistem. Tapi kita bisa mengetahui berapa byte yang dipakai tipe apa pun dengan operator \hl{\mbox{\texttt{sizeof}}}.

Jadi inilah jawabannya: sebuah \hl{\mbox{\texttt{int}}} memakai \hl{\mbox{\texttt{sizeof(int)}}} byte memori untuk menyimpannya.

Dan kalau ada 3 \hl{\mbox{\texttt{int}}} di array kita, seperti pada contoh tadi, seluruh ruang yang dipakai 3 \hl{\mbox{\texttt{int}}} itu pasti \hl{\mbox{\texttt{3\ *\ sizeof(int)}}}.

(Pada contoh string tadi, yang lebih akurat secara teknis adalah menyalin \hl{\mbox{\texttt{7\ *\ sizeof(char)}}} byte. Tapi \hl{\mbox{\texttt{char}}} selalu berukuran satu byte menurut definisi, jadi itu hanya menyusut menjadi \hl{\mbox{\texttt{7\ *\ 1}}}.)

Kita bahkan bisa menyalin sebuah \hl{\mbox{\texttt{float}}} atau \hl{\mbox{\texttt{struct}}} dengan \hl{\mbox{\texttt{memcpy()}}}! (Walaupun ini keterlaluan---untuk itu sebaiknya pakai \hl{\mbox{\texttt{=}}} saja):

\begin{tcolorbox}[enhanced jigsaw,left=3pt,breakable,colback=white!97.2549!black,boxrule=0.2mm]

\begin{Shaded}
\begin{Highlighting}[]
\KeywordTok{struct}\NormalTok{ antelope my\_antelope}\OperatorTok{;}
\KeywordTok{struct}\NormalTok{ antelope my\_clone\_antelope}\OperatorTok{;}

\CommentTok{// ...}

\NormalTok{memcpy}\OperatorTok{(\&}\NormalTok{my\_clone\_antelope}\OperatorTok{,} \OperatorTok{\&}\NormalTok{my\_antelope}\OperatorTok{,} \KeywordTok{sizeof}\NormalTok{ my\_antelope}\OperatorTok{);}
\end{Highlighting}
\end{Shaded}

\end{tcolorbox}

Lihat betapa serbaguna \hl{\mbox{\texttt{memcpy()}}}! Kalau kamu punya pointer ke sumber dan pointer ke tujuan, serta jumlah byte yang ingin disalin, kamu bisa menyalin \emph{data bertipe apa pun}.

Bayangkan kalau tidak ada \hl{\mbox{\texttt{void*}}}. Kita harus menulis fungsi \hl{\mbox{\texttt{memcpy()}}} khusus untuk setiap tipe:

\index{memcpy() function@\texttt{memcpy()} function|)}

\begin{tcolorbox}[enhanced jigsaw,left=3pt,breakable,colback=white!97.2549!black,boxrule=0.2mm]

\begin{Shaded}
\begin{Highlighting}[]
\NormalTok{memcpy\_int}\OperatorTok{(}\DataTypeTok{int} \OperatorTok{*}\NormalTok{a}\OperatorTok{,} \DataTypeTok{int} \OperatorTok{*}\NormalTok{b}\OperatorTok{,} \DataTypeTok{int}\NormalTok{ count}\OperatorTok{);}
\NormalTok{memcpy\_float}\OperatorTok{(}\DataTypeTok{float} \OperatorTok{*}\NormalTok{a}\OperatorTok{,} \DataTypeTok{float} \OperatorTok{*}\NormalTok{b}\OperatorTok{,} \DataTypeTok{int}\NormalTok{ count}\OperatorTok{);}
\NormalTok{memcpy\_double}\OperatorTok{(}\DataTypeTok{double} \OperatorTok{*}\NormalTok{a}\OperatorTok{,} \DataTypeTok{double} \OperatorTok{*}\NormalTok{b}\OperatorTok{,} \DataTypeTok{int}\NormalTok{ count}\OperatorTok{);}
\NormalTok{memcpy\_char}\OperatorTok{(}\DataTypeTok{char} \OperatorTok{*}\NormalTok{a}\OperatorTok{,} \DataTypeTok{char} \OperatorTok{*}\NormalTok{b}\OperatorTok{,} \DataTypeTok{int}\NormalTok{ count}\OperatorTok{);}
\NormalTok{memcpy\_unsigned\_char}\OperatorTok{(}\DataTypeTok{unsigned} \DataTypeTok{char} \OperatorTok{*}\NormalTok{a}\OperatorTok{,} \DataTypeTok{unsigned} \DataTypeTok{char} \OperatorTok{*}\NormalTok{b}\OperatorTok{,} \DataTypeTok{int}\NormalTok{ count}\OperatorTok{);}

\CommentTok{// dan seterusnya... iyuh!}
\end{Highlighting}
\end{Shaded}

\end{tcolorbox}

Jauh lebih baik memakai \hl{\mbox{\texttt{void*}}} dan punya satu fungsi yang bisa mengerjakan semuanya.

Itulah kekuatan \hl{\mbox{\texttt{void*}}}. Kamu bisa menulis fungsi yang tidak peduli tipe variabelnya, tapi tetap bisa melakukan sesuatu dengannya.

Tapi dengan kekuatan besar datang tanggung jawab besar. Mungkin tidak sebesar \emph{itu} dalam kasus ini, tapi memang ada beberapa batasan.

\index{void* void pointer@\texttt{void*} void pointer!caveats|(}

\begin{enumerate}
\def\labelenumi{\arabic{enumi}.}
\tightlist
\item
  Kamu tidak bisa melakukan aritmetika pointer pada \hl{\mbox{\texttt{void*}}}.
\item
  Kamu tidak bisa men-dereference \hl{\mbox{\texttt{void*}}}.
\item
  Kamu tidak bisa memakai arrow operator pada \hl{\mbox{\texttt{void*}}}, karena itu juga dereference.
\item
  Kamu tidak bisa memakai notasi array pada \hl{\mbox{\texttt{void*}}}, karena itu juga dereference\footnote{Karena ingat, notasi array hanyalah dereference ditambah sedikit aritmetika pointer, dan \hl{\mbox{\texttt{void*}}} tidak bisa di-dereference!}.
\end{enumerate}

Dan kalau dipikir-pikir, aturan-aturan ini masuk akal. Semua operasi itu bergantung pada pengetahuan tentang \hl{\mbox{\texttt{sizeof}}} tipe data yang ditunjuk, dan pada \hl{\mbox{\texttt{void*}}}, kita tidak tahu ukuran data yang ditunjuk---bisa apa saja!

\index{void* void pointer@\texttt{void*} void pointer!caveats|)}

Tapi tunggu---kalau \hl{\mbox{\texttt{void*}}} tidak bisa di-dereference, apa gunanya?

Seperti pada \hl{\mbox{\texttt{memcpy()}}}, ia membantu menulis fungsi generik yang bisa menangani berbagai tipe data. Tapi rahasianya: di dalamnya, \emph{kamu mengonversi \hl{\mbox{\texttt{void*}}} ke tipe lain sebelum memakainya}!

Dan konversinya mudah: cukup tugaskan ke variabel bertipe yang diinginkan\footnote{Kamu juga bisa melakukan \emph{cast} pada \hl{\mbox{\texttt{void*}}} ke tipe lain, tapi kita belum sampai ke cast.}.

\begin{tcolorbox}[enhanced jigsaw,left=3pt,breakable,colback=white!97.2549!black,boxrule=0.2mm]

\begin{Shaded}
\begin{Highlighting}[]
\DataTypeTok{char}\NormalTok{ a }\OperatorTok{=} \CharTok{\textquotesingle{}X\textquotesingle{}}\OperatorTok{;}  \CommentTok{// Satu char}

\DataTypeTok{void} \OperatorTok{*}\NormalTok{p }\OperatorTok{=} \OperatorTok{\&}\NormalTok{a}\OperatorTok{;}  \CommentTok{// p menunjuk ke \textquotesingle{}X\textquotesingle{}}
\DataTypeTok{char} \OperatorTok{*}\NormalTok{q }\OperatorTok{=}\NormalTok{ p}\OperatorTok{;}   \CommentTok{// q juga menunjuk ke \textquotesingle{}X\textquotesingle{}}

\NormalTok{printf}\OperatorTok{(}\StringTok{"}\SpecialCharTok{\%c\textbackslash{}n}\StringTok{"}\OperatorTok{,} \OperatorTok{*}\NormalTok{p}\OperatorTok{);}  \CommentTok{// ERROR{-}{-}tidak bisa men{-}dereference void*!}
\NormalTok{printf}\OperatorTok{(}\StringTok{"}\SpecialCharTok{\%c\textbackslash{}n}\StringTok{"}\OperatorTok{,} \OperatorTok{*}\NormalTok{q}\OperatorTok{);}  \CommentTok{// Mencetak "X"}
\end{Highlighting}
\end{Shaded}

\end{tcolorbox}

\index{memcpy() function@\texttt{memcpy()} function|(} Mari kita tulis \hl{\mbox{\texttt{memcpy()}}} sendiri untuk mencobanya. Kita bisa menyalin byte (\hl{\mbox{\texttt{char}}}), dan kita tahu jumlah byte-nya karena dioper sebagai argumen.

\begin{tcolorbox}[enhanced jigsaw,left=3pt,breakable,colback=white!97.2549!black,boxrule=0.2mm]

\begin{Shaded}
\begin{Highlighting}[]
\DataTypeTok{void} \OperatorTok{*}\NormalTok{my\_memcpy}\OperatorTok{(}\DataTypeTok{void} \OperatorTok{*}\NormalTok{dest}\OperatorTok{,} \DataTypeTok{void} \OperatorTok{*}\NormalTok{src}\OperatorTok{,} \DataTypeTok{int}\NormalTok{ byte\_count}\OperatorTok{)}
\OperatorTok{\{}
    \CommentTok{// Ubah void* menjadi char*}
    \DataTypeTok{char} \OperatorTok{*}\NormalTok{s }\OperatorTok{=}\NormalTok{ src}\OperatorTok{,} \OperatorTok{*}\NormalTok{d }\OperatorTok{=}\NormalTok{ dest}\OperatorTok{;}

    \CommentTok{// Setelah punya char*, kita bisa men{-}dereference dan menyalinnya}
    \ControlFlowTok{while} \OperatorTok{(}\NormalTok{byte\_count}\OperatorTok{{-}{-})} \OperatorTok{\{}
        \OperatorTok{*}\NormalTok{d}\OperatorTok{++} \OperatorTok{=} \OperatorTok{*}\NormalTok{s}\OperatorTok{++;}
    \OperatorTok{\}}

    \CommentTok{// Kebanyakan fungsi seperti ini mengembalikan tujuan, siapa tahu}
    \CommentTok{// berguna bagi caller.}
    \ControlFlowTok{return}\NormalTok{ dest}\OperatorTok{;}
\OperatorTok{\}}
\end{Highlighting}
\end{Shaded}

\end{tcolorbox}

Tepat di awal, kita menyalin \hl{\mbox{\texttt{void*}}} ke \hl{\mbox{\texttt{char*}}} supaya bisa dipakai sebagai \hl{\mbox{\texttt{char*}}}. Semudah itu.

Lalu ada keseruan di loop while, tempat kita men-decrement \hl{\mbox{\texttt{byte\_count}}} sampai menjadi false (\hl{\mbox{\texttt{0}}}). Ingat bahwa pada post-decrement, nilai ekspresi dihitung lebih dulu (untuk dipakai \hl{\mbox{\texttt{while}}}) dan \emph{baru kemudian} variabelnya di-decrement.

Dan ada keseruan di penyalinan, tempat kita menugaskan \hl{\mbox{\texttt{*d\ =\ *s}}} untuk menyalin byte, tapi memakai post-increment sehingga \hl{\mbox{\texttt{d}}} dan \hl{\mbox{\texttt{s}}} sama-sama bergeser ke byte berikutnya setelah penugasan dilakukan.

Terakhir, kebanyakan fungsi memori dan string mengembalikan salinan pointer ke tujuan, siapa tahu caller ingin memakainya.

Setelah semua itu, penulis ingin menunjukkan sekilas bahwa kita bisa memakai teknik ini untuk menelusuri byte-byte dari object \emph{apa pun} di C, \hl{\mbox{\texttt{float}}}, \hl{\mbox{\texttt{struct}}}, atau apa saja! \index{memcpy() function@\texttt{memcpy()} function|)}

\index{qsort() function@\texttt{qsort()} function|(} \protect\phantomsection\label{qsort-example}{Mari kita} coba satu contoh dunia nyata lagi dengan rutin bawaan \hl{\mbox{\texttt{qsort()}}} yang bisa mengurutkan \emph{apa saja} berkat keajaiban \hl{\mbox{\texttt{void*}}}.

(Pada contoh berikut, abaikan saja kata \hl{\mbox{\texttt{const}}}, yang belum kita bahas.)

\begin{tcolorbox}[enhanced jigsaw,left=3pt,breakable,colback=white!97.2549!black,boxrule=0.2mm]

\begin{Shaded}
\begin{Highlighting}[numbers=left,,]
\PreprocessorTok{\#include }\ImportTok{\textless{}stdio.h\textgreater{}}
\PreprocessorTok{\#include }\ImportTok{\textless{}stdlib.h\textgreater{}}

\CommentTok{// Tipe struktur yang akan kita urutkan}
\KeywordTok{struct}\NormalTok{ animal }\OperatorTok{\{}
    \DataTypeTok{char} \OperatorTok{*}\NormalTok{name}\OperatorTok{;}
    \DataTypeTok{int}\NormalTok{ leg\_count}\OperatorTok{;}
\OperatorTok{\};}

\CommentTok{// Ini fungsi pembanding yang dipanggil oleh qsort() untuk membantunya}
\CommentTok{// menentukan apa persisnya yang dijadikan dasar pengurutan. Kita pakai}
\CommentTok{// untuk mengurutkan array struct animal berdasarkan leg\_count.}
\DataTypeTok{int}\NormalTok{ compar}\OperatorTok{(}\DataTypeTok{const} \DataTypeTok{void} \OperatorTok{*}\NormalTok{elem1}\OperatorTok{,} \DataTypeTok{const} \DataTypeTok{void} \OperatorTok{*}\NormalTok{elem2}\OperatorTok{)}
\OperatorTok{\{}
    \CommentTok{// Kita tahu yang diurutkan adalah struct animal, jadi mari kita}
    \CommentTok{// jadikan kedua argumen pointer ke struct animal}
    \DataTypeTok{const} \KeywordTok{struct}\NormalTok{ animal }\OperatorTok{*}\NormalTok{animal1 }\OperatorTok{=}\NormalTok{ elem1}\OperatorTok{;}
    \DataTypeTok{const} \KeywordTok{struct}\NormalTok{ animal }\OperatorTok{*}\NormalTok{animal2 }\OperatorTok{=}\NormalTok{ elem2}\OperatorTok{;}

    \CommentTok{// Kembalikan \textless{}0, =0, atau \textgreater{}0 tergantung dasar pengurutan yang kita mau.}

    \CommentTok{// Kita urutkan menaik berdasarkan leg\_count, jadi kita kembalikan}
    \CommentTok{// selisih leg\_count}
    \ControlFlowTok{if} \OperatorTok{(}\NormalTok{animal1}\OperatorTok{{-}\textgreater{}}\NormalTok{leg\_count }\OperatorTok{\textgreater{}}\NormalTok{ animal2}\OperatorTok{{-}\textgreater{}}\NormalTok{leg\_count}\OperatorTok{)}
        \ControlFlowTok{return} \DecValTok{1}\OperatorTok{;}
    
    \ControlFlowTok{if} \OperatorTok{(}\NormalTok{animal1}\OperatorTok{{-}\textgreater{}}\NormalTok{leg\_count }\OperatorTok{\textless{}}\NormalTok{ animal2}\OperatorTok{{-}\textgreater{}}\NormalTok{leg\_count}\OperatorTok{)}
        \ControlFlowTok{return} \OperatorTok{{-}}\DecValTok{1}\OperatorTok{;}

    \ControlFlowTok{return} \DecValTok{0}\OperatorTok{;}
\OperatorTok{\}}

\DataTypeTok{int}\NormalTok{ main}\OperatorTok{(}\DataTypeTok{void}\OperatorTok{)}
\OperatorTok{\{}
    \CommentTok{// Kita buat array berisi 4 struct animal dengan karakteristik}
    \CommentTok{// berbeda. Array ini tidak berurutan menurut leg\_count, tapi}
    \CommentTok{// sebentar lagi akan kita urutkan.}
    \KeywordTok{struct}\NormalTok{ animal a}\OperatorTok{[}\DecValTok{4}\OperatorTok{]} \OperatorTok{=} \OperatorTok{\{}
        \OperatorTok{\{.}\NormalTok{name}\OperatorTok{=}\StringTok{"Dog"}\OperatorTok{,} \OperatorTok{.}\NormalTok{leg\_count}\OperatorTok{=}\DecValTok{4}\OperatorTok{\},}
        \OperatorTok{\{.}\NormalTok{name}\OperatorTok{=}\StringTok{"Monkey"}\OperatorTok{,} \OperatorTok{.}\NormalTok{leg\_count}\OperatorTok{=}\DecValTok{2}\OperatorTok{\},}
        \OperatorTok{\{.}\NormalTok{name}\OperatorTok{=}\StringTok{"Antelope"}\OperatorTok{,} \OperatorTok{.}\NormalTok{leg\_count}\OperatorTok{=}\DecValTok{4}\OperatorTok{\},}
        \OperatorTok{\{.}\NormalTok{name}\OperatorTok{=}\StringTok{"Snake"}\OperatorTok{,} \OperatorTok{.}\NormalTok{leg\_count}\OperatorTok{=}\DecValTok{0}\OperatorTok{\}}
    \OperatorTok{\};}

    \CommentTok{// Panggil qsort() untuk mengurutkan array. qsort() perlu diberi tahu}
    \CommentTok{// persis berdasarkan apa data ini diurutkan, dan itu kita lakukan}
    \CommentTok{// di dalam fungsi compar().}
    \CommentTok{//}
    \CommentTok{// Pemanggilan ini berarti: qsort array a, yang punya 4 elemen, dan}
    \CommentTok{// setiap elemen berukuran sizeof(struct animal) byte, dan ini}
    \CommentTok{// fungsi yang akan membandingkan dua elemen mana pun.}
\NormalTok{    qsort}\OperatorTok{(}\NormalTok{a}\OperatorTok{,} \DecValTok{4}\OperatorTok{,} \KeywordTok{sizeof}\OperatorTok{(}\KeywordTok{struct}\NormalTok{ animal}\OperatorTok{),}\NormalTok{ compar}\OperatorTok{);}

    \CommentTok{// Cetak semuanya}
    \ControlFlowTok{for} \OperatorTok{(}\DataTypeTok{int}\NormalTok{ i }\OperatorTok{=} \DecValTok{0}\OperatorTok{;}\NormalTok{ i }\OperatorTok{\textless{}} \DecValTok{4}\OperatorTok{;}\NormalTok{ i}\OperatorTok{++)} \OperatorTok{\{}
\NormalTok{        printf}\OperatorTok{(}\StringTok{"}\SpecialCharTok{\%d}\StringTok{: }\SpecialCharTok{\%s\textbackslash{}n}\StringTok{"}\OperatorTok{,}\NormalTok{ a}\OperatorTok{[}\NormalTok{i}\OperatorTok{].}\NormalTok{leg\_count}\OperatorTok{,}\NormalTok{ a}\OperatorTok{[}\NormalTok{i}\OperatorTok{].}\NormalTok{name}\OperatorTok{);}
    \OperatorTok{\}}
\OperatorTok{\}}
\end{Highlighting}
\end{Shaded}

\end{tcolorbox}

Selama kamu memberi \hl{\mbox{\texttt{qsort()}}} fungsi yang bisa membandingkan dua item dalam array yang akan diurutkan, ia bisa mengurutkan apa saja. Dan itu dilakukan tanpa perlu menanamkan tipe item di mana pun. \hl{\mbox{\texttt{qsort()}}} hanya menata ulang blok-blok byte berdasarkan hasil fungsi \hl{\mbox{\texttt{compar()}}} yang kamu oper. \index{qsort() function@\texttt{qsort()} function|)} \index{void* void pointer@\texttt{void*} void pointer|)} \index{Pointers!arithmetic|)}

\chapter{Alokasi Memori Manual}\label{alokasi-memori-manual}

\index{Manual memory management|(} Inilah salah satu area besar tempat C kemungkinan berbeda dari bahasa yang sudah kamu kenal: \emph{manajemen memori manual}.

Bahasa lain memakai reference counting, garbage collection, atau cara lain untuk menentukan kapan harus mengalokasikan memori baru untuk suatu data---dan kapan harus mendealokasikannya ketika tidak ada lagi variabel yang merujuknya.

Dan itu menyenangkan. Enak rasanya tidak perlu memikirkannya, cukup melepas semua referensi ke suatu item dan percaya bahwa pada suatu saat memori yang terkait dengannya akan dibebaskan.

Tapi C tidak sepenuhnya begitu.

\index{Automatic variables|(} Tentu saja, di C ada variabel yang dialokasikan dan didealokasikan otomatis ketika masuk dan keluar scope. Kita menyebutnya variabel otomatis. Mereka adalah variabel ``lokal'' biasa dengan scope blok. Tidak ada masalah. \index{Automatic variables|)}

Tapi bagaimana kalau kamu ingin sesuatu bertahan lebih lama dari sebuah blok tertentu? Di sinilah manajemen memori manual berperan.

Kamu bisa menyuruh C secara eksplisit untuk mengalokasikan sejumlah byte yang bisa kamu pakai sesukamu. Dan byte-byte itu akan tetap teralokasi sampai kamu membebaskan memori itu secara eksplisit\footnote{Atau sampai program berakhir, dan saat itu semua memori yang dialokasikannya dibebaskan. Tanda bintang: beberapa sistem mengizinkan kamu mengalokasikan memori yang bertahan setelah program berakhir, tapi itu bergantung pada sistem, di luar cakupan panduan ini, dan kamu pasti tidak akan melakukannya secara tidak sengaja.}.

Penting untuk membebaskan memori yang sudah selesai kamu pakai! Kalau tidak, itu disebut \emph{memory leak} dan prosesmu akan terus menahan memori itu sampai berakhir.

\emph{Kalau kamu mengalokasikannya secara manual, kamu harus membebaskannya secara manual saat selesai memakainya.}

Jadi bagaimana caranya? Kita akan belajar beberapa fungsi baru, dan memakai operator \hl{\mbox{\texttt{sizeof}}} untuk membantu mengetahui berapa byte yang harus dialokasikan.

\index{The stack|(} \index{The heap|(} Dalam istilah umum di dunia C, para developer bilang variabel lokal otomatis dialokasikan ``di stack'', dan memori yang dialokasikan manual ada ``di heap''. Spesifikasi tidak membicarakan keduanya, tapi semua developer C akan paham yang kamu maksud kalau kamu menyebutnya. \index{The heap|)} \index{The stack|)}

Semua fungsi yang kita pelajari di bab ini ada di \hl{\mbox{\texttt{\textless{}stdlib.h\textgreater{}}}}.

\section{\texorpdfstring{Mengalokasikan dan Mendealokasikan, \hl{\mbox{\texttt{malloc()}}} dan \hl{\mbox{\texttt{free()}}}}{Mengalokasikan dan Mendealokasikan, malloc() dan free()}}\label{mengalokasikan-dan-mendealokasikan-malloc-dan-free}

\index{malloc() function@\texttt{malloc()} function|(} Fungsi \hl{\mbox{\texttt{malloc()}}} menerima jumlah byte yang akan dialokasikan, dan mengembalikan void pointer ke blok memori yang baru dialokasikan itu.

Karena berupa \hl{\mbox{\texttt{void*}}}, hasilnya bisa ditugaskan ke tipe pointer apa pun yang kamu mau\ldots{} biasanya tipe ini berkaitan dengan jumlah byte yang kamu alokasikan.

\index{sizeof operator@\texttt{sizeof} operator!with malloc()@with \texttt{malloc()}|(} Jadi\ldots{} berapa byte yang harus dialokasikan? \hl{\mbox{\texttt{sizeof}}} bisa membantu. Kalau kita ingin mengalokasikan ruang untuk satu \hl{\mbox{\texttt{int}}}, kita bisa memakai \hl{\mbox{\texttt{sizeof(int)}}} dan mengopernya ke \hl{\mbox{\texttt{malloc()}}}. \index{sizeof operator@\texttt{sizeof} operator!with malloc()@with \texttt{malloc()}|)}

\index{free() function@\texttt{free()} function|(} Setelah selesai memakai memori yang dialokasikan, kita bisa memanggil \hl{\mbox{\texttt{free()}}} untuk menandakan bahwa kita sudah selesai dan memori itu bisa dipakai untuk hal lain. Sebagai argumen, oper pointer yang sama dengan yang kamu dapat dari \hl{\mbox{\texttt{malloc()}}} (atau salinannya). Memakai wilayah memori setelah di-\hl{\mbox{\texttt{free()}}} adalah undefined behavior.

Mari kita coba. Kita alokasikan memori yang cukup untuk satu \hl{\mbox{\texttt{int}}}, lalu simpan sesuatu di sana, lalu cetak.

\begin{tcolorbox}[enhanced jigsaw,left=3pt,breakable,colback=white!97.2549!black,boxrule=0.2mm]

\begin{Shaded}
\begin{Highlighting}[]
\CommentTok{// Alokasikan ruang untuk satu int (sebesar sizeof(int) byte):}

\DataTypeTok{int} \OperatorTok{*}\NormalTok{p }\OperatorTok{=}\NormalTok{ malloc}\OperatorTok{(}\KeywordTok{sizeof}\OperatorTok{(}\DataTypeTok{int}\OperatorTok{));}

\OperatorTok{*}\NormalTok{p }\OperatorTok{=} \DecValTok{12}\OperatorTok{;}  \CommentTok{// Simpan sesuatu di sana}

\NormalTok{printf}\OperatorTok{(}\StringTok{"}\SpecialCharTok{\%d\textbackslash{}n}\StringTok{"}\OperatorTok{,} \OperatorTok{*}\NormalTok{p}\OperatorTok{);}  \CommentTok{// Cetak: 12}

\NormalTok{free}\OperatorTok{(}\NormalTok{p}\OperatorTok{);}  \CommentTok{// Selesai memakai memori itu}

\CommentTok{//*p = 3490;  // ERROR: undefined behavior! Dipakai setelah free()!}
\end{Highlighting}
\end{Shaded}

\end{tcolorbox}

\index{free() function@\texttt{free()} function|)}

Nah, pada contoh yang dibuat-buat itu, sebenarnya tidak ada untungnya. Kita bisa saja memakai \hl{\mbox{\texttt{int}}} otomatis dan hasilnya sama. Tapi kita akan lihat bahwa kemampuan mengalokasikan memori seperti ini punya keunggulan, terutama untuk struktur data yang lebih rumit.

\index{sizeof operator@\texttt{sizeof} operator!with malloc()@with \texttt{malloc()}|(} Satu hal lagi yang sering kamu lihat memanfaatkan fakta bahwa \hl{\mbox{\texttt{sizeof}}} bisa memberi ukuran tipe hasil dari ekspresi konstan apa pun. Jadi kamu juga bisa menaruh nama variabel di sana dan memakainya. Berikut contohnya, sama seperti sebelumnya:

\begin{tcolorbox}[enhanced jigsaw,left=3pt,breakable,colback=white!97.2549!black,boxrule=0.2mm]

\begin{Shaded}
\begin{Highlighting}[]
\DataTypeTok{int} \OperatorTok{*}\NormalTok{p }\OperatorTok{=}\NormalTok{ malloc}\OperatorTok{(}\KeywordTok{sizeof} \OperatorTok{*}\NormalTok{p}\OperatorTok{);}  \CommentTok{// *p adalah int, jadi sama dengan sizeof(int)}
\end{Highlighting}
\end{Shaded}

\end{tcolorbox}

\index{sizeof operator@\texttt{sizeof} operator!with malloc()@with \texttt{malloc()}|)} \index{malloc() function@\texttt{malloc()} function|)}

\section{Pengecekan Error}\label{pengecekan-error}

\index{malloc() function@\texttt{malloc()} function!error checking|(} Semua fungsi alokasi mengembalikan pointer ke bentangan memori yang baru dialokasikan, atau \hl{\mbox{\texttt{NULL}}} kalau memori tidak bisa dialokasikan karena suatu alasan.

Beberapa OS seperti Linux bisa dikonfigurasi sedemikian rupa sehingga \hl{\mbox{\texttt{malloc()}}} tidak pernah mengembalikan \hl{\mbox{\texttt{NULL}}}, bahkan ketika memori habis. Tapi meskipun begitu, kamu tetap harus menulis kodenya dengan perlindungan.

\begin{tcolorbox}[enhanced jigsaw,left=3pt,breakable,colback=white!97.2549!black,boxrule=0.2mm]

\begin{Shaded}
\begin{Highlighting}[]
\DataTypeTok{int} \OperatorTok{*}\NormalTok{x}\OperatorTok{;}

\NormalTok{x }\OperatorTok{=}\NormalTok{ malloc}\OperatorTok{(}\KeywordTok{sizeof}\OperatorTok{(}\DataTypeTok{int}\OperatorTok{)} \OperatorTok{*} \DecValTok{10}\OperatorTok{);}

\ControlFlowTok{if} \OperatorTok{(}\NormalTok{x }\OperatorTok{==}\NormalTok{ NULL}\OperatorTok{)} \OperatorTok{\{}
\NormalTok{    printf}\OperatorTok{(}\StringTok{"Error allocating 10 ints}\SpecialCharTok{\textbackslash{}n}\StringTok{"}\OperatorTok{);}
    \CommentTok{// lakukan sesuatu di sini untuk menanganinya}
\OperatorTok{\}}
\end{Highlighting}
\end{Shaded}

\end{tcolorbox}

Berikut pola umum yang akan sering kamu lihat, ketika penugasan dan kondisi dilakukan di baris yang sama:

\begin{tcolorbox}[enhanced jigsaw,left=3pt,breakable,colback=white!97.2549!black,boxrule=0.2mm]

\begin{Shaded}
\begin{Highlighting}[]
\DataTypeTok{int} \OperatorTok{*}\NormalTok{x}\OperatorTok{;}

\ControlFlowTok{if} \OperatorTok{((}\NormalTok{x }\OperatorTok{=}\NormalTok{ malloc}\OperatorTok{(}\KeywordTok{sizeof}\OperatorTok{(}\DataTypeTok{int}\OperatorTok{)} \OperatorTok{*} \DecValTok{10}\OperatorTok{))} \OperatorTok{==}\NormalTok{ NULL}\OperatorTok{)} \OperatorTok{\{}
\NormalTok{    printf}\OperatorTok{(}\StringTok{"Error allocating 10 ints}\SpecialCharTok{\textbackslash{}n}\StringTok{"}\OperatorTok{);}
    \CommentTok{// lakukan sesuatu di sini untuk menanganinya}
\OperatorTok{\}}
\end{Highlighting}
\end{Shaded}

\end{tcolorbox}

\index{malloc() function@\texttt{malloc()} function!error checking|)}

\section{Mengalokasikan Ruang untuk Array}\label{mengalokasikan-ruang-untuk-array}

\index{malloc() function@\texttt{malloc()} function!and arrays|(} Kita sudah melihat cara mengalokasikan ruang untuk satu benda; bagaimana kalau untuk banyak benda dalam sebuah array?

Di C, array adalah sekumpulan benda yang sama, berurutan dalam satu bentangan memori yang kontigu.

Kita bisa mengalokasikan bentangan memori yang kontigu---kita sudah tahu caranya. Kalau kita ingin 3490 byte memori, tinggal minta saja:

\begin{tcolorbox}[enhanced jigsaw,left=3pt,breakable,colback=white!97.2549!black,boxrule=0.2mm]

\begin{Shaded}
\begin{Highlighting}[]
\DataTypeTok{char} \OperatorTok{*}\NormalTok{p }\OperatorTok{=}\NormalTok{ malloc}\OperatorTok{(}\DecValTok{3490}\OperatorTok{);}  \CommentTok{// Voila}
\end{Highlighting}
\end{Shaded}

\end{tcolorbox}

Dan---memang!---itu array berisi 3490 \hl{\mbox{\texttt{char}}} (alias string!) karena setiap \hl{\mbox{\texttt{char}}} berukuran 1 byte. Dengan kata lain, \hl{\mbox{\texttt{sizeof(char)}}} adalah \hl{\mbox{\texttt{1}}}.

Catatan: tidak ada inisialisasi pada memori yang baru dialokasikan---isinya penuh sampah. Bersihkan dengan \hl{\mbox{\texttt{memset()}}} kalau perlu, atau lihat \hl{\mbox{\texttt{calloc()}}} di bawah.

Tapi kita bisa mengalikan ukuran benda yang diinginkan dengan jumlah elemen yang diinginkan, lalu mengaksesnya dengan notasi pointer maupun array. Contoh!

\index{sizeof operator@\texttt{sizeof} operator!with malloc()@with \texttt{malloc()}|(}

\begin{tcolorbox}[enhanced jigsaw,left=3pt,breakable,colback=white!97.2549!black,boxrule=0.2mm]

\begin{Shaded}
\begin{Highlighting}[numbers=left,,]
\PreprocessorTok{\#include }\ImportTok{\textless{}stdio.h\textgreater{}}
\PreprocessorTok{\#include }\ImportTok{\textless{}stdlib.h\textgreater{}}

\DataTypeTok{int}\NormalTok{ main}\OperatorTok{(}\DataTypeTok{void}\OperatorTok{)}
\OperatorTok{\{}
    \CommentTok{// Alokasikan ruang untuk 10 int}
    \DataTypeTok{int} \OperatorTok{*}\NormalTok{p }\OperatorTok{=}\NormalTok{ malloc}\OperatorTok{(}\KeywordTok{sizeof}\OperatorTok{(}\DataTypeTok{int}\OperatorTok{)} \OperatorTok{*} \DecValTok{10}\OperatorTok{);}

    \CommentTok{// Isi dengan nilai 0{-}45:}
    \ControlFlowTok{for} \OperatorTok{(}\DataTypeTok{int}\NormalTok{ i }\OperatorTok{=} \DecValTok{0}\OperatorTok{;}\NormalTok{ i }\OperatorTok{\textless{}} \DecValTok{10}\OperatorTok{;}\NormalTok{ i}\OperatorTok{++)}
\NormalTok{        p}\OperatorTok{[}\NormalTok{i}\OperatorTok{]} \OperatorTok{=}\NormalTok{ i }\OperatorTok{*} \DecValTok{5}\OperatorTok{;}

    \CommentTok{// Cetak semua nilai 0, 5, 10, 15, ..., 40, 45}
    \ControlFlowTok{for} \OperatorTok{(}\DataTypeTok{int}\NormalTok{ i }\OperatorTok{=} \DecValTok{0}\OperatorTok{;}\NormalTok{ i }\OperatorTok{\textless{}} \DecValTok{10}\OperatorTok{;}\NormalTok{ i}\OperatorTok{++)}
\NormalTok{        printf}\OperatorTok{(}\StringTok{"}\SpecialCharTok{\%d\textbackslash{}n}\StringTok{"}\OperatorTok{,}\NormalTok{ p}\OperatorTok{[}\NormalTok{i}\OperatorTok{]);}

    \CommentTok{// Bebaskan ruangnya}
\NormalTok{    free}\OperatorTok{(}\NormalTok{p}\OperatorTok{);}
\OperatorTok{\}}
\end{Highlighting}
\end{Shaded}

\end{tcolorbox}

Kuncinya ada di baris \hl{\mbox{\texttt{malloc()}}} itu. Kalau kita tahu setiap \hl{\mbox{\texttt{int}}} membutuhkan \hl{\mbox{\texttt{sizeof(int)}}} byte, dan kita ingin 10 buah, kita tinggal mengalokasikan persis sebanyak itu byte dengan:

\begin{tcolorbox}[enhanced jigsaw,left=3pt,breakable,colback=white!97.2549!black,boxrule=0.2mm]

\begin{Shaded}
\begin{Highlighting}[]
\KeywordTok{sizeof}\OperatorTok{(}\DataTypeTok{int}\OperatorTok{)} \OperatorTok{*} \DecValTok{10}
\end{Highlighting}
\end{Shaded}

\end{tcolorbox}

Dan trik ini berlaku untuk semua tipe. Cukup oper ke \hl{\mbox{\texttt{sizeof}}} dan kalikan dengan ukuran array. \index{sizeof operator@\texttt{sizeof} operator!with malloc()@with \texttt{malloc()}|)} \index{malloc() function@\texttt{malloc()} function!and arrays|)}

\section{\texorpdfstring{Alternatif: \hl{\mbox{\texttt{calloc()}}}}{Alternatif: calloc()}}\label{alternatif-calloc}

\index{calloc() function@\texttt{calloc()} function|(} Ini fungsi alokasi lain yang bekerja mirip \hl{\mbox{\texttt{malloc()}}}, dengan dua perbedaan utama:

\begin{itemize}
\tightlist
\item
  Alih-alih satu argumen, kamu mengoper ukuran satu elemen dan jumlah elemen yang ingin dialokasikan. Seolah-olah ia dibuat khusus untuk mengalokasikan array.
\item
  Ia membersihkan memori menjadi nol.
\end{itemize}

Kamu tetap memakai \hl{\mbox{\texttt{free()}}} untuk mendealokasikan memori yang didapat lewat \hl{\mbox{\texttt{calloc()}}}.

Berikut perbandingan \hl{\mbox{\texttt{calloc()}}} dan \hl{\mbox{\texttt{malloc()}}}.

\begin{tcolorbox}[enhanced jigsaw,left=3pt,breakable,colback=white!97.2549!black,boxrule=0.2mm]

\begin{Shaded}
\begin{Highlighting}[]
\CommentTok{// Alokasikan ruang untuk 10 int dengan calloc(), diinisialisasi ke 0:}
\DataTypeTok{int} \OperatorTok{*}\NormalTok{p }\OperatorTok{=}\NormalTok{ calloc}\OperatorTok{(}\DecValTok{10}\OperatorTok{,} \KeywordTok{sizeof}\OperatorTok{(}\DataTypeTok{int}\OperatorTok{));}

\CommentTok{// Alokasikan ruang untuk 10 int dengan malloc(), diinisialisasi ke 0:}
\DataTypeTok{int} \OperatorTok{*}\NormalTok{q }\OperatorTok{=}\NormalTok{ malloc}\OperatorTok{(}\DecValTok{10} \OperatorTok{*} \KeywordTok{sizeof}\OperatorTok{(}\DataTypeTok{int}\OperatorTok{));}
\NormalTok{memset}\OperatorTok{(}\NormalTok{q}\OperatorTok{,} \DecValTok{0}\OperatorTok{,} \DecValTok{10} \OperatorTok{*} \KeywordTok{sizeof}\OperatorTok{(}\DataTypeTok{int}\OperatorTok{));}   \CommentTok{// isi dengan 0}
\end{Highlighting}
\end{Shaded}

\end{tcolorbox}

Sekali lagi, hasilnya sama untuk keduanya, kecuali \hl{\mbox{\texttt{malloc()}}} tidak mengosongkan memori secara bawaan. \index{calloc() function@\texttt{calloc()} function|)}

\section{\texorpdfstring{Mengubah Ukuran Alokasi dengan \hl{\mbox{\texttt{realloc()}}}}{Mengubah Ukuran Alokasi dengan realloc()}}\label{mengubah-ukuran-alokasi-dengan-realloc}

\index{realloc() function@\texttt{realloc()} function|(} Kalau kamu sudah mengalokasikan 10 \hl{\mbox{\texttt{int}}}, tapi kemudian memutuskan butuh 20, apa yang bisa dilakukan?

Satu pilihan adalah mengalokasikan ruang baru, lalu \hl{\mbox{\texttt{memcpy()}}} isi memorinya\ldots{} tapi ternyata kadang tidak ada yang perlu dipindahkan. Dan ada satu fungsi yang cukup pintar untuk melakukan hal yang benar di semua keadaan yang tepat: \hl{\mbox{\texttt{realloc()}}}.

Ia menerima pointer ke memori yang sebelumnya dialokasikan (oleh \hl{\mbox{\texttt{malloc()}}} atau \hl{\mbox{\texttt{calloc()}}}) dan ukuran baru untuk wilayah memori itu.

Ia lalu membesarkan atau mengecilkan memori itu, dan mengembalikan pointer ke sana. Kadang ia mengembalikan pointer yang sama (kalau data tidak perlu disalin ke tempat lain), atau pointer yang berbeda (kalau data memang harus disalin).

Pastikan saat memanggil \hl{\mbox{\texttt{realloc()}}}, kamu menentukan jumlah \emph{byte} yang dialokasikan, bukan sekadar jumlah elemen array! Yaitu:

\begin{tcolorbox}[enhanced jigsaw,left=3pt,breakable,colback=white!97.2549!black,boxrule=0.2mm]

\begin{Shaded}
\begin{Highlighting}[]
\NormalTok{num\_floats }\OperatorTok{*=} \DecValTok{2}\OperatorTok{;}

\NormalTok{np }\OperatorTok{=}\NormalTok{ realloc}\OperatorTok{(}\NormalTok{p}\OperatorTok{,}\NormalTok{ num\_floats}\OperatorTok{);}  \CommentTok{// SALAH: butuh byte, bukan jumlah elemen!}

\NormalTok{np }\OperatorTok{=}\NormalTok{ realloc}\OperatorTok{(}\NormalTok{p}\OperatorTok{,}\NormalTok{ num\_floats }\OperatorTok{*} \KeywordTok{sizeof}\OperatorTok{(}\DataTypeTok{float}\OperatorTok{));}  \CommentTok{// Lebih baik!}
\end{Highlighting}
\end{Shaded}

\end{tcolorbox}

Mari kita alokasikan array berisi 20 \hl{\mbox{\texttt{float}}}, lalu berubah pikiran dan menjadikannya array berisi 40.

Kita akan menugaskan nilai kembalian realloc() ke pointer lain supaya bisa memeriksa apakah ia NULL. Kalau bukan, baru kita tugaskan ulang ke pointer asli. (Kalau nilai kembalian langsung ditugaskan ke pointer asli, kita akan kehilangan pointer itu ketika fungsi mengembalikan \hl{\mbox{\texttt{NULL}}}, dan tidak ada cara untuk mendapatkannya lagi.)

\begin{tcolorbox}[enhanced jigsaw,left=3pt,breakable,colback=white!97.2549!black,boxrule=0.2mm]

\begin{Shaded}
\begin{Highlighting}[numbers=left,,]
\PreprocessorTok{\#include }\ImportTok{\textless{}stdio.h\textgreater{}}
\PreprocessorTok{\#include }\ImportTok{\textless{}stdlib.h\textgreater{}}

\DataTypeTok{int}\NormalTok{ main}\OperatorTok{(}\DataTypeTok{void}\OperatorTok{)}
\OperatorTok{\{}
    \CommentTok{// Alokasikan ruang untuk 20 float}
    \DataTypeTok{float} \OperatorTok{*}\NormalTok{p }\OperatorTok{=}\NormalTok{ malloc}\OperatorTok{(}\KeywordTok{sizeof} \OperatorTok{*}\NormalTok{p }\OperatorTok{*} \DecValTok{20}\OperatorTok{);}  \CommentTok{// sizeof *p sama dengan sizeof(float)}

    \CommentTok{// Isi dengan nilai pecahan 0.0{-}1.0:}
    \ControlFlowTok{for} \OperatorTok{(}\DataTypeTok{int}\NormalTok{ i }\OperatorTok{=} \DecValTok{0}\OperatorTok{;}\NormalTok{ i }\OperatorTok{\textless{}} \DecValTok{20}\OperatorTok{;}\NormalTok{ i}\OperatorTok{++)}
\NormalTok{        p}\OperatorTok{[}\NormalTok{i}\OperatorTok{]} \OperatorTok{=}\NormalTok{ i }\OperatorTok{/} \FloatTok{20.0}\OperatorTok{;}

    \CommentTok{// Tapi tunggu! Ternyata kita ingin array ini berisi 40 elemen}
    \DataTypeTok{float} \OperatorTok{*}\NormalTok{new\_p }\OperatorTok{=}\NormalTok{ realloc}\OperatorTok{(}\NormalTok{p}\OperatorTok{,} \KeywordTok{sizeof} \OperatorTok{*}\NormalTok{p }\OperatorTok{*} \DecValTok{40}\OperatorTok{);}

    \CommentTok{// Periksa apakah realokasi berhasil}
    \ControlFlowTok{if} \OperatorTok{(}\NormalTok{new\_p }\OperatorTok{==}\NormalTok{ NULL}\OperatorTok{)} \OperatorTok{\{}
\NormalTok{        printf}\OperatorTok{(}\StringTok{"Error reallocing}\SpecialCharTok{\textbackslash{}n}\StringTok{"}\OperatorTok{);}
        \ControlFlowTok{return} \DecValTok{1}\OperatorTok{;}
    \OperatorTok{\}}

    \CommentTok{// Kalau berhasil, kita tinggal menugaskan ulang p}
\NormalTok{    p }\OperatorTok{=}\NormalTok{ new\_p}\OperatorTok{;}

    \CommentTok{// Lalu isi elemen{-}elemen baru dengan nilai dalam rentang 1.0{-}2.0}
    \ControlFlowTok{for} \OperatorTok{(}\DataTypeTok{int}\NormalTok{ i }\OperatorTok{=} \DecValTok{20}\OperatorTok{;}\NormalTok{ i }\OperatorTok{\textless{}} \DecValTok{40}\OperatorTok{;}\NormalTok{ i}\OperatorTok{++)}
\NormalTok{        p}\OperatorTok{[}\NormalTok{i}\OperatorTok{]} \OperatorTok{=} \FloatTok{1.0} \OperatorTok{+} \OperatorTok{(}\NormalTok{i }\OperatorTok{{-}} \DecValTok{20}\OperatorTok{)} \OperatorTok{/} \FloatTok{20.0}\OperatorTok{;}

    \CommentTok{// Cetak semua nilai 0.0{-}2.0 dari 40 elemen:}
    \ControlFlowTok{for} \OperatorTok{(}\DataTypeTok{int}\NormalTok{ i }\OperatorTok{=} \DecValTok{0}\OperatorTok{;}\NormalTok{ i }\OperatorTok{\textless{}} \DecValTok{40}\OperatorTok{;}\NormalTok{ i}\OperatorTok{++)}
\NormalTok{        printf}\OperatorTok{(}\StringTok{"}\SpecialCharTok{\%f\textbackslash{}n}\StringTok{"}\OperatorTok{,}\NormalTok{ p}\OperatorTok{[}\NormalTok{i}\OperatorTok{]);}

    \CommentTok{// Bebaskan ruangnya}
\NormalTok{    free}\OperatorTok{(}\NormalTok{p}\OperatorTok{);}
\OperatorTok{\}}
\end{Highlighting}
\end{Shaded}

\end{tcolorbox}

Perhatikan bagaimana kita mengambil nilai kembalian \hl{\mbox{\texttt{realloc()}}} dan menugaskannya ulang ke variabel pointer \hl{\mbox{\texttt{p}}} yang sama dengan yang kita oper. Itu hal yang cukup umum dilakukan.

Kalau baris 7 tampak aneh dengan \hl{\mbox{\texttt{sizeof\ *p}}} di dalamnya, ingat bahwa \hl{\mbox{\texttt{sizeof}}} bekerja pada ukuran tipe dari ekspresi. Dan tipe \hl{\mbox{\texttt{*p}}} adalah \hl{\mbox{\texttt{float}}}, jadi baris itu setara dengan \hl{\mbox{\texttt{sizeof(float)}}}.

Terakhir, mungkin terasa janggal bahwa tidak ada \hl{\mbox{\texttt{free(new\_p)}}} di mana pun, padahal itu pointer yang dikembalikan \hl{\mbox{\texttt{realloc()}}}. Alasannya, \hl{\mbox{\texttt{new\_p}}} disalin ke \hl{\mbox{\texttt{p}}} di baris 23, jadi keduanya bernilai sama; artinya keduanya menunjuk ke gumpalan memori yang sama, dan gumpalannya hanya satu. Jadi saat \hl{\mbox{\texttt{free()}}}, kita bisa membebaskan yang mana saja dengan efek yang sama.

\index{realloc() function@\texttt{realloc()} function|)}

\subsection{Membaca Baris dengan Panjang Sembarang}\label{membaca-baris-dengan-panjang-sembarang}

Dengan contoh lengkap ini, penulis ingin menunjukkan dua hal.

\begin{enumerate}
\def\labelenumi{\arabic{enumi}.}
\tightlist
\item
  Pemakaian \hl{\mbox{\texttt{realloc()}}} untuk membesarkan buffer saat kita membaca lebih banyak data.
\item
  Pemakaian \hl{\mbox{\texttt{realloc()}}} untuk mengecilkan buffer ke ukuran yang pas setelah pembacaan selesai.
\end{enumerate}

Di sini ada loop yang memanggil \hl{\mbox{\texttt{fgetc()}}} berulang kali untuk menambah isi buffer, sampai karakter terakhirnya newline.

Begitu newline ditemukan, buffer dikecilkan ke ukuran yang tepat lalu dikembalikan.

\begin{tcolorbox}[enhanced jigsaw,left=3pt,breakable,colback=white!97.2549!black,boxrule=0.2mm]

\begin{Shaded}
\begin{Highlighting}[numbers=left,,]
\PreprocessorTok{\#include }\ImportTok{\textless{}stdio.h\textgreater{}}
\PreprocessorTok{\#include }\ImportTok{\textless{}stdlib.h\textgreater{}}

\CommentTok{// Baca satu baris berukuran sembarang dari file}
\CommentTok{//}
\CommentTok{// Mengembalikan pointer ke baris tersebut.}
\CommentTok{// Mengembalikan NULL saat EOF atau error.}
\CommentTok{//}
\CommentTok{// Caller bertanggung jawab men{-}free() pointer ini setelah selesai.}
\CommentTok{//}
\CommentTok{// Perhatikan bahwa newline dibuang dari hasilnya. Kalau newline}
\CommentTok{// diperlukan, sebaiknya ubah ini menjadi do{-}while.}

\DataTypeTok{char} \OperatorTok{*}\NormalTok{readline}\OperatorTok{(}\DataTypeTok{FILE} \OperatorTok{*}\NormalTok{fp}\OperatorTok{)}
\OperatorTok{\{}
    \DataTypeTok{int}\NormalTok{ offset }\OperatorTok{=} \DecValTok{0}\OperatorTok{;}   \CommentTok{// Indeks tempat karakter berikutnya masuk ke buffer}
    \DataTypeTok{int}\NormalTok{ bufsize }\OperatorTok{=} \DecValTok{4}\OperatorTok{;}  \CommentTok{// Sebaiknya ukuran awal berupa pangkat 2}
    \DataTypeTok{char} \OperatorTok{*}\NormalTok{buf}\OperatorTok{;}        \CommentTok{// Buffer{-}nya}
    \DataTypeTok{int}\NormalTok{ c}\OperatorTok{;}            \CommentTok{// Karakter yang baru dibaca}

\NormalTok{    buf }\OperatorTok{=}\NormalTok{ malloc}\OperatorTok{(}\NormalTok{bufsize}\OperatorTok{);}  \CommentTok{// Alokasikan buffer awal}

    \ControlFlowTok{if} \OperatorTok{(}\NormalTok{buf }\OperatorTok{==}\NormalTok{ NULL}\OperatorTok{)}   \CommentTok{// Cek error}
        \ControlFlowTok{return}\NormalTok{ NULL}\OperatorTok{;}

    \CommentTok{// Loop utama{-}{-}baca sampai newline atau EOF}
    \ControlFlowTok{while} \OperatorTok{(}\NormalTok{c }\OperatorTok{=}\NormalTok{ fgetc}\OperatorTok{(}\NormalTok{fp}\OperatorTok{),}\NormalTok{ c }\OperatorTok{!=} \CharTok{\textquotesingle{}}\SpecialCharTok{\textbackslash{}n}\CharTok{\textquotesingle{}} \OperatorTok{\&\&}\NormalTok{ c }\OperatorTok{!=}\NormalTok{ EOF}\OperatorTok{)} \OperatorTok{\{}

        \CommentTok{// Periksa apakah ruang di buffer habis, dengan memperhitungkan}
        \CommentTok{// satu byte ekstra untuk terminator NUL}
        \ControlFlowTok{if} \OperatorTok{(}\NormalTok{offset }\OperatorTok{==}\NormalTok{ bufsize }\OperatorTok{{-}} \DecValTok{1}\OperatorTok{)} \OperatorTok{\{}  \CommentTok{// {-}1 untuk terminator NUL}
\NormalTok{            bufsize }\OperatorTok{*=} \DecValTok{2}\OperatorTok{;}  \CommentTok{// Gandakan ruangnya}

            \DataTypeTok{char} \OperatorTok{*}\NormalTok{new\_buf }\OperatorTok{=}\NormalTok{ realloc}\OperatorTok{(}\NormalTok{buf}\OperatorTok{,}\NormalTok{ bufsize}\OperatorTok{);}

            \ControlFlowTok{if} \OperatorTok{(}\NormalTok{new\_buf }\OperatorTok{==}\NormalTok{ NULL}\OperatorTok{)} \OperatorTok{\{}
\NormalTok{                free}\OperatorTok{(}\NormalTok{buf}\OperatorTok{);}   \CommentTok{// Saat error, free lalu keluar}
                \ControlFlowTok{return}\NormalTok{ NULL}\OperatorTok{;}
            \OperatorTok{\}}

\NormalTok{            buf }\OperatorTok{=}\NormalTok{ new\_buf}\OperatorTok{;}  \CommentTok{// Realokasi berhasil}
        \OperatorTok{\}}

\NormalTok{        buf}\OperatorTok{[}\NormalTok{offset}\OperatorTok{++]} \OperatorTok{=}\NormalTok{ c}\OperatorTok{;}  \CommentTok{// Tambahkan byte ke buffer}
    \OperatorTok{\}}

    \CommentTok{// Kita menemui newline atau EOF...}

    \CommentTok{// Kalau di EOF dan tidak ada byte yang terbaca, free buffer dan}
    \CommentTok{// kembalikan NULL untuk menandakan kita di EOF:}
    \ControlFlowTok{if} \OperatorTok{(}\NormalTok{c }\OperatorTok{==}\NormalTok{ EOF }\OperatorTok{\&\&}\NormalTok{ offset }\OperatorTok{==} \DecValTok{0}\OperatorTok{)} \OperatorTok{\{}
\NormalTok{        free}\OperatorTok{(}\NormalTok{buf}\OperatorTok{);}
        \ControlFlowTok{return}\NormalTok{ NULL}\OperatorTok{;}
    \OperatorTok{\}}

    \CommentTok{// Kecilkan sesuai ukuran}
    \ControlFlowTok{if} \OperatorTok{(}\NormalTok{offset }\OperatorTok{\textless{}}\NormalTok{ bufsize }\OperatorTok{{-}} \DecValTok{1}\OperatorTok{)} \OperatorTok{\{}  \CommentTok{// Kalau belum sampai ujung}
        \DataTypeTok{char} \OperatorTok{*}\NormalTok{new\_buf }\OperatorTok{=}\NormalTok{ realloc}\OperatorTok{(}\NormalTok{buf}\OperatorTok{,}\NormalTok{ offset }\OperatorTok{+} \DecValTok{1}\OperatorTok{);} \CommentTok{// +1 untuk terminator NUL}

        \CommentTok{// Kalau berhasil, arahkan buf ke new\_buf;}
        \CommentTok{// kalau tidak, biarkan buf di tempatnya}
        \ControlFlowTok{if} \OperatorTok{(}\NormalTok{new\_buf }\OperatorTok{!=}\NormalTok{ NULL}\OperatorTok{)}
\NormalTok{            buf }\OperatorTok{=}\NormalTok{ new\_buf}\OperatorTok{;}
    \OperatorTok{\}}

    \CommentTok{// Tambahkan terminator NUL}
\NormalTok{    buf}\OperatorTok{[}\NormalTok{offset}\OperatorTok{]} \OperatorTok{=} \CharTok{\textquotesingle{}}\SpecialCharTok{\textbackslash{}0}\CharTok{\textquotesingle{}}\OperatorTok{;}

    \ControlFlowTok{return}\NormalTok{ buf}\OperatorTok{;}
\OperatorTok{\}}

\DataTypeTok{int}\NormalTok{ main}\OperatorTok{(}\DataTypeTok{void}\OperatorTok{)}
\OperatorTok{\{}
    \DataTypeTok{FILE} \OperatorTok{*}\NormalTok{fp }\OperatorTok{=}\NormalTok{ fopen}\OperatorTok{(}\StringTok{"foo.txt"}\OperatorTok{,} \StringTok{"r"}\OperatorTok{);}

    \DataTypeTok{char} \OperatorTok{*}\NormalTok{line}\OperatorTok{;}

    \ControlFlowTok{while} \OperatorTok{((}\NormalTok{line }\OperatorTok{=}\NormalTok{ readline}\OperatorTok{(}\NormalTok{fp}\OperatorTok{))} \OperatorTok{!=}\NormalTok{ NULL}\OperatorTok{)} \OperatorTok{\{}
\NormalTok{        printf}\OperatorTok{(}\StringTok{"}\SpecialCharTok{\%s\textbackslash{}n}\StringTok{"}\OperatorTok{,}\NormalTok{ line}\OperatorTok{);}
\NormalTok{        free}\OperatorTok{(}\NormalTok{line}\OperatorTok{);}
    \OperatorTok{\}}

\NormalTok{    fclose}\OperatorTok{(}\NormalTok{fp}\OperatorTok{);}
\OperatorTok{\}}
\end{Highlighting}
\end{Shaded}

\end{tcolorbox}

Saat membesarkan memori seperti ini, lazimnya (meski bukan aturan baku) ruang yang dibutuhkan digandakan di setiap langkah, supaya jumlah \hl{\mbox{\texttt{realloc()}}} yang terjadi sesedikit mungkin.

Terakhir, kamu mungkin memperhatikan bahwa \hl{\mbox{\texttt{readline()}}} mengembalikan pointer ke buffer hasil \hl{\mbox{\texttt{malloc()}}}. Karena itu, caller bertanggung jawab untuk men-\hl{\mbox{\texttt{free()}}} memori itu secara eksplisit saat selesai memakainya.

\subsection{\texorpdfstring{\hl{\mbox{\texttt{realloc()}}} dengan \hl{\mbox{\texttt{NULL}}}}{realloc() dengan NULL}}\label{realloc-dengan-null}

\index{realloc() function@\texttt{realloc()} function!with NULL argument@with \texttt{NULL} argument|(} Waktunya trivia! Dua baris ini setara:

\begin{tcolorbox}[enhanced jigsaw,left=3pt,breakable,colback=white!97.2549!black,boxrule=0.2mm]

\begin{Shaded}
\begin{Highlighting}[]
\DataTypeTok{char} \OperatorTok{*}\NormalTok{p }\OperatorTok{=}\NormalTok{ malloc}\OperatorTok{(}\DecValTok{3490}\OperatorTok{);}
\DataTypeTok{char} \OperatorTok{*}\NormalTok{p }\OperatorTok{=}\NormalTok{ realloc}\OperatorTok{(}\NormalTok{NULL}\OperatorTok{,} \DecValTok{3490}\OperatorTok{);}
\end{Highlighting}
\end{Shaded}

\end{tcolorbox}

Ini bisa praktis kalau kamu punya semacam loop alokasi dan tidak ingin memperlakukan \hl{\mbox{\texttt{malloc()}}} pertama secara khusus.

\begin{tcolorbox}[enhanced jigsaw,left=3pt,breakable,colback=white!97.2549!black,boxrule=0.2mm]

\begin{Shaded}
\begin{Highlighting}[]
\DataTypeTok{int} \OperatorTok{*}\NormalTok{p }\OperatorTok{=}\NormalTok{ NULL}\OperatorTok{;}
\DataTypeTok{int}\NormalTok{ length }\OperatorTok{=} \DecValTok{0}\OperatorTok{;}

\ControlFlowTok{while} \OperatorTok{(!}\NormalTok{done}\OperatorTok{)} \OperatorTok{\{}
    \CommentTok{// Alokasikan 10 int lagi:}
\NormalTok{    length }\OperatorTok{+=} \DecValTok{10}\OperatorTok{;}
\NormalTok{    p }\OperatorTok{=}\NormalTok{ realloc}\OperatorTok{(}\NormalTok{p}\OperatorTok{,} \KeywordTok{sizeof} \OperatorTok{*}\NormalTok{p }\OperatorTok{*}\NormalTok{ length}\OperatorTok{);}

    \CommentTok{// Lakukan hal{-}hal luar biasa}
    \CommentTok{// ...}
\OperatorTok{\}}
\end{Highlighting}
\end{Shaded}

\end{tcolorbox}

Pada contoh itu, kita tidak butuh \hl{\mbox{\texttt{malloc()}}} awal karena \hl{\mbox{\texttt{p}}} memang \hl{\mbox{\texttt{NULL}}} sejak awal. \index{realloc() function@\texttt{realloc()} function!with NULL argument@with \texttt{NULL} argument|)}

\section{Alokasi dengan Alignment}\label{alokasi-dengan-alignment}

\index{Memory alignment|(} Kemungkinan besar kamu tidak akan memerlukan ini.

Dan penulis tidak ingin terlalu melenceng membahasnya sekarang, tapi ada yang disebut \emph{memory alignment}, yang berkaitan dengan alamat memori (nilai pointer) yang harus merupakan kelipatan dari angka tertentu.

Misalnya, sebuah sistem mungkin mewajibkan nilai 16-bit dimulai di alamat memori kelipatan 2. Atau nilai 64-bit dimulai di alamat memori kelipatan 2, 4, atau 8, misalnya. Itu tergantung CPU.

Sebagian sistem membutuhkan alignment semacam ini agar akses memori cepat, bahkan ada yang membutuhkannya agar akses memori bisa dilakukan sama sekali.

Nah, kalau kamu memakai \hl{\mbox{\texttt{malloc()}}}, \hl{\mbox{\texttt{calloc()}}}, atau \hl{\mbox{\texttt{realloc()}}}, C akan memberimu gumpalan memori yang ter-align dengan baik untuk nilai apa pun, bahkan \hl{\mbox{\texttt{struct}}}. Berlaku di semua kasus.

Tapi mungkin ada saatnya kamu tahu bahwa data tertentu bisa di-align pada batas yang lebih kecil, atau harus di-align pada batas yang lebih besar karena suatu alasan. Penulis membayangkan ini lebih umum pada pemrograman embedded system.

\index{aligned\_alloc() function@\texttt{aligned\_alloc()} function|(} Dalam kasus itu, kamu bisa menentukan alignment dengan \hl{\mbox{\texttt{aligned\_alloc()}}}.

Alignment adalah integer pangkat dua yang lebih besar dari nol, jadi \hl{\mbox{\texttt{2}}}, \hl{\mbox{\texttt{4}}}, \hl{\mbox{\texttt{8}}}, \hl{\mbox{\texttt{16}}}, dan seterusnya, dan itu diberikan ke \hl{\mbox{\texttt{aligned\_alloc()}}} sebelum jumlah byte yang kamu inginkan.

Pembatas lainnya: jumlah byte yang dialokasikan harus merupakan kelipatan alignment. Tapi ini mungkin akan berubah. Lihat \href{http://www.open-std.org/jtc1/sc22/wg14/www/docs/summary.htm\#dr_460}{C Defect Report 460}\footnote{http://www.open-std.org/jtc1/sc22/wg14/www/docs/summary.htm\#dr\_460}

Mari kita coba contoh, mengalokasikan pada batas 64 byte:

\begin{tcolorbox}[enhanced jigsaw,left=3pt,breakable,colback=white!97.2549!black,boxrule=0.2mm]

\begin{Shaded}
\begin{Highlighting}[numbers=left,,]
\PreprocessorTok{\#include }\ImportTok{\textless{}stdio.h\textgreater{}}
\PreprocessorTok{\#include }\ImportTok{\textless{}stdlib.h\textgreater{}}
\PreprocessorTok{\#include }\ImportTok{\textless{}string.h\textgreater{}}

\DataTypeTok{int}\NormalTok{ main}\OperatorTok{(}\DataTypeTok{void}\OperatorTok{)}
\OperatorTok{\{}
    \CommentTok{// Alokasikan 256 byte yang ter{-}align pada batas 64 byte}
    \DataTypeTok{char} \OperatorTok{*}\NormalTok{p }\OperatorTok{=}\NormalTok{ aligned\_alloc}\OperatorTok{(}\DecValTok{64}\OperatorTok{,} \DecValTok{256}\OperatorTok{);}  \CommentTok{// 256 == 64 * 4}

    \CommentTok{// Salin sebuah string ke sana dan cetak}
\NormalTok{    strcpy}\OperatorTok{(}\NormalTok{p}\OperatorTok{,} \StringTok{"Hello, world!"}\OperatorTok{);}
\NormalTok{    printf}\OperatorTok{(}\StringTok{"}\SpecialCharTok{\%s\textbackslash{}n}\StringTok{"}\OperatorTok{,}\NormalTok{ p}\OperatorTok{);}

    \CommentTok{// Bebaskan ruangnya}
\NormalTok{    free}\OperatorTok{(}\NormalTok{p}\OperatorTok{);}
\OperatorTok{\}}
\end{Highlighting}
\end{Shaded}

\end{tcolorbox}

Ada catatan soal \hl{\mbox{\texttt{realloc()}}} dan \hl{\mbox{\texttt{aligned\_alloc()}}}. \hl{\mbox{\texttt{realloc()}}} tidak punya jaminan alignment apa pun, jadi kalau kamu butuh ruang hasil realokasi yang ter-align, kamu harus melakukannya dengan cara susah memakai \hl{\mbox{\texttt{memcpy()}}}. \index{aligned\_alloc() function@\texttt{aligned\_alloc()} function|)}

Berikut fungsi \hl{\mbox{\texttt{aligned\_realloc()}}} non-standar, kalau kamu membutuhkannya:

\begin{tcolorbox}[enhanced jigsaw,left=3pt,breakable,colback=white!97.2549!black,boxrule=0.2mm]

\begin{Shaded}
\begin{Highlighting}[]
\DataTypeTok{void} \OperatorTok{*}\NormalTok{aligned\_realloc}\OperatorTok{(}\DataTypeTok{void} \OperatorTok{*}\NormalTok{ptr}\OperatorTok{,} \DataTypeTok{size\_t}\NormalTok{ old\_size}\OperatorTok{,} \DataTypeTok{size\_t}\NormalTok{ alignment}\OperatorTok{,} \DataTypeTok{size\_t}\NormalTok{ size}\OperatorTok{)}
\OperatorTok{\{}
    \DataTypeTok{char} \OperatorTok{*}\NormalTok{new\_ptr }\OperatorTok{=}\NormalTok{ aligned\_alloc}\OperatorTok{(}\NormalTok{alignment}\OperatorTok{,}\NormalTok{ size}\OperatorTok{);}

    \ControlFlowTok{if} \OperatorTok{(}\NormalTok{new\_ptr }\OperatorTok{==}\NormalTok{ NULL}\OperatorTok{)}
        \ControlFlowTok{return}\NormalTok{ NULL}\OperatorTok{;}

    \DataTypeTok{size\_t}\NormalTok{ copy\_size }\OperatorTok{=}\NormalTok{ old\_size }\OperatorTok{\textless{}}\NormalTok{ size}\OperatorTok{?}\NormalTok{ old\_size}\OperatorTok{:}\NormalTok{ size}\OperatorTok{;}  \CommentTok{// ambil yang terkecil}

    \ControlFlowTok{if} \OperatorTok{(}\NormalTok{ptr }\OperatorTok{!=}\NormalTok{ NULL}\OperatorTok{)}
\NormalTok{        memcpy}\OperatorTok{(}\NormalTok{new\_ptr}\OperatorTok{,}\NormalTok{ ptr}\OperatorTok{,}\NormalTok{ copy\_size}\OperatorTok{);}

\NormalTok{    free}\OperatorTok{(}\NormalTok{ptr}\OperatorTok{);}

    \ControlFlowTok{return}\NormalTok{ new\_ptr}\OperatorTok{;}
\OperatorTok{\}}
\end{Highlighting}
\end{Shaded}

\end{tcolorbox}

Perhatikan bahwa fungsi ini \emph{selalu} menyalin data, yang memakan waktu, sedangkan \hl{\mbox{\texttt{realloc()}}} yang asli akan menghindarinya kalau bisa. Jadi ini jauh dari efisien. Hindari kebutuhan untuk merealokasi data yang alignment-nya kustom. \index{Memory alignment|)} \index{Manual memory management|)}

\chapter{Scope}\label{scope}

\index{Scope|(} Scope adalah soal variabel apa saja yang terlihat di konteks mana.

\section{Block Scope}\label{block-scope}

\index{Scope!block|(} Hampir semua variabel yang didefinisikan developer memakai scope ini, termasuk yang di bahasa lain mungkin disebut ``function scope'', yaitu variabel yang dideklarasikan di dalam fungsi.

Aturan dasarnya: kalau kamu mendeklarasikan variabel di dalam blok yang dibatasi kurung kurawal, scope variabel itu adalah blok tersebut.

Kalau ada blok di dalam blok, variabel yang dideklarasikan di blok \emph{dalam} bersifat lokal untuk blok itu, dan tidak terlihat di scope luar.

Begitu scope sebuah variabel berakhir, variabel itu tidak bisa lagi dirujuk, dan nilainya boleh dianggap sudah lenyap ke \href{https://en.wikipedia.org/wiki/Bit_bucket}{bit bucket}\footnote{https://en.wikipedia.org/wiki/Bit\_bucket} raksasa di langit.

Contoh dengan scope bersarang:

\begin{tcolorbox}[enhanced jigsaw,left=3pt,breakable,colback=white!97.2549!black,boxrule=0.2mm]

\begin{Shaded}
\begin{Highlighting}[numbers=left,,]
\PreprocessorTok{\#include }\ImportTok{\textless{}stdio.h\textgreater{}}

\DataTypeTok{int}\NormalTok{ main}\OperatorTok{(}\DataTypeTok{void}\OperatorTok{)}
\OperatorTok{\{}
    \DataTypeTok{int}\NormalTok{ a }\OperatorTok{=} \DecValTok{12}\OperatorTok{;}         \CommentTok{// Lokal untuk blok luar, tapi terlihat di blok dalam}

    \ControlFlowTok{if}  \OperatorTok{(}\NormalTok{a }\OperatorTok{==} \DecValTok{12}\OperatorTok{)} \OperatorTok{\{}
        \DataTypeTok{int}\NormalTok{ b }\OperatorTok{=} \DecValTok{99}\OperatorTok{;}     \CommentTok{// Lokal untuk blok dalam, tidak terlihat di blok luar}

\NormalTok{        printf}\OperatorTok{(}\StringTok{"}\SpecialCharTok{\%d}\StringTok{ }\SpecialCharTok{\%d\textbackslash{}n}\StringTok{"}\OperatorTok{,}\NormalTok{ a}\OperatorTok{,}\NormalTok{ b}\OperatorTok{);}  \CommentTok{// OK: "12 99"}
    \OperatorTok{\}}

\NormalTok{    printf}\OperatorTok{(}\StringTok{"}\SpecialCharTok{\%d\textbackslash{}n}\StringTok{"}\OperatorTok{,}\NormalTok{ a}\OperatorTok{);}  \CommentTok{// OK, kita masih di scope a}

\NormalTok{    printf}\OperatorTok{(}\StringTok{"}\SpecialCharTok{\%d\textbackslash{}n}\StringTok{"}\OperatorTok{,}\NormalTok{ b}\OperatorTok{);}  \CommentTok{// ILEGAL, di luar scope b}
\OperatorTok{\}}
\end{Highlighting}
\end{Shaded}

\end{tcolorbox}

\subsection{Di Mana Variabel Didefinisikan}\label{di-mana-variabel-didefinisikan}

Fakta menarik lainnya: variabel bisa didefinisikan di mana saja dalam blok, dalam batas kewajaran---scope-nya adalah blok itu, tapi tidak bisa dipakai sebelum didefinisikan.

\begin{tcolorbox}[enhanced jigsaw,left=3pt,breakable,colback=white!97.2549!black,boxrule=0.2mm]

\begin{Shaded}
\begin{Highlighting}[numbers=left,,]
\PreprocessorTok{\#include }\ImportTok{\textless{}stdio.h\textgreater{}}

\DataTypeTok{int}\NormalTok{ main}\OperatorTok{(}\DataTypeTok{void}\OperatorTok{)}
\OperatorTok{\{}
    \DataTypeTok{int}\NormalTok{ i }\OperatorTok{=} \DecValTok{0}\OperatorTok{;}

\NormalTok{    printf}\OperatorTok{(}\StringTok{"}\SpecialCharTok{\%d\textbackslash{}n}\StringTok{"}\OperatorTok{,}\NormalTok{ i}\OperatorTok{);}     \CommentTok{// OK: "0"}

    \CommentTok{//printf("\%d\textbackslash{}n", j);   // ILEGAL{-}{-}j tidak bisa dipakai sebelum didefinisikan}

    \DataTypeTok{int}\NormalTok{ j }\OperatorTok{=} \DecValTok{5}\OperatorTok{;}

\NormalTok{    printf}\OperatorTok{(}\StringTok{"}\SpecialCharTok{\%d}\StringTok{ }\SpecialCharTok{\%d\textbackslash{}n}\StringTok{"}\OperatorTok{,}\NormalTok{ i}\OperatorTok{,}\NormalTok{ j}\OperatorTok{);}   \CommentTok{// OK: "0 5"}
\OperatorTok{\}}
\end{Highlighting}
\end{Shaded}

\end{tcolorbox}

Secara historis, C mewajibkan semua variabel didefinisikan sebelum kode apa pun di dalam blok, tapi sejak standar C99 hal itu tidak berlaku lagi.

\subsection{Variable Hiding}\label{variable-hiding}

\index{Variable hiding|(} Kalau ada variabel bernama sama di scope dalam dan di scope luar, yang di scope dalam diutamakan selama kode berjalan di scope dalam. Artinya, ia \emph{menyembunyikan} variabel di scope luar selama masa hidupnya.

\begin{tcolorbox}[enhanced jigsaw,left=3pt,breakable,colback=white!97.2549!black,boxrule=0.2mm]

\begin{Shaded}
\begin{Highlighting}[numbers=left,,]
\PreprocessorTok{\#include }\ImportTok{\textless{}stdio.h\textgreater{}}

\DataTypeTok{int}\NormalTok{ main}\OperatorTok{(}\DataTypeTok{void}\OperatorTok{)}
\OperatorTok{\{}
    \DataTypeTok{int}\NormalTok{ i }\OperatorTok{=} \DecValTok{10}\OperatorTok{;}

    \OperatorTok{\{}
        \DataTypeTok{int}\NormalTok{ i }\OperatorTok{=} \DecValTok{20}\OperatorTok{;}

\NormalTok{        printf}\OperatorTok{(}\StringTok{"}\SpecialCharTok{\%d\textbackslash{}n}\StringTok{"}\OperatorTok{,}\NormalTok{ i}\OperatorTok{);}  \CommentTok{// i di scope dalam, 20 (i luar tersembunyi)}
    \OperatorTok{\}}

\NormalTok{    printf}\OperatorTok{(}\StringTok{"}\SpecialCharTok{\%d\textbackslash{}n}\StringTok{"}\OperatorTok{,}\NormalTok{ i}\OperatorTok{);}  \CommentTok{// i di scope luar, 10}
\OperatorTok{\}}
\end{Highlighting}
\end{Shaded}

\end{tcolorbox}

Mungkin kamu memperhatikan bahwa pada contoh itu, penulis cuma melempar sebuah blok begitu saja di baris 7, tanpa \hl{\mbox{\texttt{for}}} atau \hl{\mbox{\texttt{if}}} yang mengawalinya! Ini sepenuhnya legal. Kadang developer ingin mengelompokkan beberapa variabel lokal untuk perhitungan singkat dan melakukan ini, tapi jarang terlihat. \index{Variable hiding|)} \index{Scope!block|)}

\section{File Scope}\label{file-scope}

\index{Scope!file|(} Kalau kamu mendefinisikan variabel di luar blok, variabel itu punya \emph{file scope}. Ia terlihat di semua fungsi dalam file yang muncul setelahnya, dan dipakai bersama oleh fungsi-fungsi itu. (Pengecualiannya: kalau sebuah blok mendefinisikan variabel bernama sama, variabel itu akan menyembunyikan yang ada di file scope.)

Ini yang paling mendekati apa yang kamu anggap sebagai scope ``global'' di bahasa lain.

Contohnya:

\begin{tcolorbox}[enhanced jigsaw,left=3pt,breakable,colback=white!97.2549!black,boxrule=0.2mm]

\begin{Shaded}
\begin{Highlighting}[numbers=left,,]
\PreprocessorTok{\#include }\ImportTok{\textless{}stdio.h\textgreater{}}

\DataTypeTok{int}\NormalTok{ shared }\OperatorTok{=} \DecValTok{10}\OperatorTok{;}    \CommentTok{// File scope! Terlihat di seluruh file setelah ini!}

\DataTypeTok{void}\NormalTok{ func1}\OperatorTok{(}\DataTypeTok{void}\OperatorTok{)}
\OperatorTok{\{}
\NormalTok{    shared }\OperatorTok{+=} \DecValTok{100}\OperatorTok{;}  \CommentTok{// Sekarang shared bernilai 110}
\OperatorTok{\}}

\DataTypeTok{void}\NormalTok{ func2}\OperatorTok{(}\DataTypeTok{void}\OperatorTok{)}
\OperatorTok{\{}
\NormalTok{    printf}\OperatorTok{(}\StringTok{"}\SpecialCharTok{\%d\textbackslash{}n}\StringTok{"}\OperatorTok{,}\NormalTok{ shared}\OperatorTok{);}  \CommentTok{// Mencetak "110"}
\OperatorTok{\}}

\DataTypeTok{int}\NormalTok{ main}\OperatorTok{(}\DataTypeTok{void}\OperatorTok{)}
\OperatorTok{\{}
\NormalTok{    func1}\OperatorTok{();}
\NormalTok{    func2}\OperatorTok{();}
\OperatorTok{\}}
\end{Highlighting}
\end{Shaded}

\end{tcolorbox}

Perhatikan bahwa kalau \hl{\mbox{\texttt{shared}}} dideklarasikan di bagian bawah file, kode itu tidak akan dikompilasi. Variabel harus dideklarasikan \emph{sebelum} dipakai oleh fungsi mana pun.

Ada cara untuk memodifikasi item di file scope lebih jauh, yaitu dengan \hyperref[static]{static} dan \hyperref[extern]{extern}, tapi itu kita bahas nanti. \index{Scope!file|)}

\section{\texorpdfstring{Scope pada Loop \hl{\mbox{\texttt{for}}}}{Scope pada Loop for}}\label{scope-pada-loop-for}

\index{Scope!for loop@\texttt{for} loop|(} Penulis sungguh tidak tahu harus menyebut ini apa, karena C11 §6.8.5.3¶1 tidak memberinya nama yang jelas. Kita juga sudah beberapa kali melakukannya di panduan ini. Yaitu ketika kamu mendeklarasikan variabel di klausa pertama sebuah loop \hl{\mbox{\texttt{for}}}:

\begin{tcolorbox}[enhanced jigsaw,left=3pt,breakable,colback=white!97.2549!black,boxrule=0.2mm]

\begin{Shaded}
\begin{Highlighting}[]
\ControlFlowTok{for} \OperatorTok{(}\DataTypeTok{int}\NormalTok{ i }\OperatorTok{=} \DecValTok{0}\OperatorTok{;}\NormalTok{ i }\OperatorTok{\textless{}} \DecValTok{10}\OperatorTok{;}\NormalTok{ i}\OperatorTok{++)}
\NormalTok{    printf}\OperatorTok{(}\StringTok{"}\SpecialCharTok{\%d\textbackslash{}n}\StringTok{"}\OperatorTok{,}\NormalTok{ i}\OperatorTok{);}

\NormalTok{printf}\OperatorTok{(}\StringTok{"}\SpecialCharTok{\%d\textbackslash{}n}\StringTok{"}\OperatorTok{,}\NormalTok{ i}\OperatorTok{);}  \CommentTok{// ILEGAL{-}{-}i hanya ada di scope selama for{-}loop}
\end{Highlighting}
\end{Shaded}

\end{tcolorbox}

Pada contoh itu, masa hidup \hl{\mbox{\texttt{i}}} dimulai saat ia didefinisikan, dan berlanjut selama loop berjalan.

Kalau badan loop dibungkus dalam sebuah blok, variabel yang didefinisikan di \hl{\mbox{\texttt{for}}}-loop terlihat dari scope dalam itu.

Kecuali, tentu saja, kalau scope dalam itu menyembunyikannya. Contoh gila ini mencetak \hl{\mbox{\texttt{999}}} sebanyak lima kali:

\begin{tcolorbox}[enhanced jigsaw,left=3pt,breakable,colback=white!97.2549!black,boxrule=0.2mm]

\begin{Shaded}
\begin{Highlighting}[numbers=left,,]
\PreprocessorTok{\#include }\ImportTok{\textless{}stdio.h\textgreater{}}

\DataTypeTok{int}\NormalTok{ main}\OperatorTok{(}\DataTypeTok{void}\OperatorTok{)}
\OperatorTok{\{}
    \ControlFlowTok{for} \OperatorTok{(}\DataTypeTok{int}\NormalTok{ i }\OperatorTok{=} \DecValTok{0}\OperatorTok{;}\NormalTok{ i }\OperatorTok{\textless{}} \DecValTok{5}\OperatorTok{;}\NormalTok{ i}\OperatorTok{++)} \OperatorTok{\{}
        \DataTypeTok{int}\NormalTok{ i }\OperatorTok{=} \DecValTok{999}\OperatorTok{;}  \CommentTok{// Menyembunyikan i di scope for{-}loop}
\NormalTok{        printf}\OperatorTok{(}\StringTok{"}\SpecialCharTok{\%d\textbackslash{}n}\StringTok{"}\OperatorTok{,}\NormalTok{ i}\OperatorTok{);}
    \OperatorTok{\}}
\OperatorTok{\}}
\end{Highlighting}
\end{Shaded}

\end{tcolorbox}

\index{Scope!for loop@\texttt{for} loop|)}

\section{Catatan tentang Function Scope}\label{catatan-tentang-function-scope}

\index{Scope!function|(} Spesifikasi C memang menyebut \emph{function scope}, tapi istilah itu dipakai khusus untuk \emph{label}, yang belum kita bahas. Nanti di lain hari. \index{Scope!function|)} \index{Scope|)}

\chapter{Tipe II: Jauh Lebih Banyak Tipe!}\label{tipe-ii-jauh-lebih-banyak-tipe}

Kita sudah terbiasa dengan tipe \hl{\mbox{\texttt{char}}}, \hl{\mbox{\texttt{int}}}, dan \hl{\mbox{\texttt{float}}}, tapi sekarang saatnya membawanya ke level berikutnya dan melihat apa lagi yang tersedia di departemen tipe!

\section{Integer Signed dan Unsigned}\label{integer-signed-dan-unsigned}

\index{Types!signed and unsigned|(} Sejauh ini kita memakai \hl{\mbox{\texttt{int}}} sebagai tipe \emph{signed}, yaitu nilai yang bisa negatif maupun positif. Tapi C juga punya tipe integer \emph{unsigned} khusus yang hanya bisa menyimpan angka positif.

Tipe-tipe ini diawali dengan kata kunci \index{unsigned type@\texttt{unsigned} type|(}\hl{\mbox{\texttt{unsigned}}}.

\begin{tcolorbox}[enhanced jigsaw,left=3pt,breakable,colback=white!97.2549!black,boxrule=0.2mm]

\begin{Shaded}
\begin{Highlighting}[]
\DataTypeTok{int}\NormalTok{ a}\OperatorTok{;}           \CommentTok{// signed}
\DataTypeTok{signed} \DataTypeTok{int}\NormalTok{ a}\OperatorTok{;}    \CommentTok{// signed}
\DataTypeTok{signed}\NormalTok{ a}\OperatorTok{;}        \CommentTok{// signed, "singkatan" dari "int" atau "signed int", jarang dipakai}
\DataTypeTok{unsigned} \DataTypeTok{int}\NormalTok{ b}\OperatorTok{;}  \CommentTok{// unsigned}
\DataTypeTok{unsigned}\NormalTok{ c}\OperatorTok{;}      \CommentTok{// unsigned, singkatan dari "unsigned int"}
\end{Highlighting}
\end{Shaded}

\end{tcolorbox}

Kenapa? Kenapa kamu memutuskan hanya ingin menyimpan angka positif?

Jawabannya: variabel unsigned bisa menampung angka yang lebih besar dibanding variabel signed.

Tapi kenapa begitu?

Bayangkan integer direpresentasikan oleh sejumlah \emph{bit}\footnote{``Bit'' kependekan dari \emph{binary digit}. Biner hanyalah cara lain merepresentasikan angka. Alih-alih digit 0-9 yang biasa kita pakai, digitnya hanya 0-1.}. Di komputerku, sebuah \hl{\mbox{\texttt{int}}} direpresentasikan oleh 64 bit.

Dan setiap permutasi bit yang bernilai \hl{\mbox{\texttt{1}}} atau \hl{\mbox{\texttt{0}}} mewakili sebuah angka. Kita yang memutuskan bagaimana membagi angka-angka ini.

Pada angka signed, kita memakai (kurang lebih) separuh permutasi untuk angka negatif, dan separuh lainnya untuk angka positif.

Pada unsigned, kita memakai \emph{semua} permutasi untuk angka positif.

Di komputerku, dengan \hl{\mbox{\texttt{int}}} 64-bit yang memakai \href{https://en.wikipedia.org/wiki/Two\%27s_complement}{two's complement}\footnote{https://en.wikipedia.org/wiki/Two\%27s\_complement} untuk merepresentasikan angka unsigned, rentang integernya adalah sebagai berikut:

{\def\LTcaptype{none} % do not increment counter
\begin{longtable}[]{@{}lrr@{}}
\toprule\noalign{}
Tipe & Minimum & Maksimum \\
\midrule\noalign{}
\endhead
\bottomrule\noalign{}
\endlastfoot
\hl{\mbox{\texttt{int}}} & \hl{\mbox{\texttt{-9,223,372,036,854,775,808}}} & \hl{\mbox{\texttt{9,223,372,036,854,775,807}}} \\
\hl{\mbox{\texttt{unsigned\ int}}} & \hl{\mbox{\texttt{0}}} & \hl{\mbox{\texttt{18,446,744,073,709,551,615}}} \\
\end{longtable}
}

Perhatikan bahwa \hl{\mbox{\texttt{unsigned\ int}}} positif terbesar kira-kira dua kali lebih besar dari \hl{\mbox{\texttt{int}}} positif terbesar. Jadi ada sedikit keleluasaan di sana. \index{unsigned type@\texttt{unsigned} type|)} \index{Types!signed and unsigned|)}

\section{Tipe Karakter}\label{tipe-karakter}

\index{Types!character|(} \index{char type@\texttt{char} type|(} Ingat \hl{\mbox{\texttt{char}}}? Tipe yang bisa kita pakai untuk menyimpan satu karakter?

\begin{tcolorbox}[enhanced jigsaw,left=3pt,breakable,colback=white!97.2549!black,boxrule=0.2mm]

\begin{Shaded}
\begin{Highlighting}[]
\DataTypeTok{char}\NormalTok{ c }\OperatorTok{=} \CharTok{\textquotesingle{}B\textquotesingle{}}\OperatorTok{;}

\NormalTok{printf}\OperatorTok{(}\StringTok{"}\SpecialCharTok{\%c\textbackslash{}n}\StringTok{"}\OperatorTok{,}\NormalTok{ c}\OperatorTok{);}  \CommentTok{// "B"}
\end{Highlighting}
\end{Shaded}

\end{tcolorbox}

Ada kejutan untukmu: sebenarnya ia adalah integer.

\begin{tcolorbox}[enhanced jigsaw,left=3pt,breakable,colback=white!97.2549!black,boxrule=0.2mm]

\begin{Shaded}
\begin{Highlighting}[]
\DataTypeTok{char}\NormalTok{ c }\OperatorTok{=} \CharTok{\textquotesingle{}B\textquotesingle{}}\OperatorTok{;}

\CommentTok{// Ubah ini dari \%c menjadi \%d:}
\NormalTok{printf}\OperatorTok{(}\StringTok{"}\SpecialCharTok{\%d\textbackslash{}n}\StringTok{"}\OperatorTok{,}\NormalTok{ c}\OperatorTok{);}  \CommentTok{// 66 (!!)}
\end{Highlighting}
\end{Shaded}

\end{tcolorbox}

Jauh di dalamnya, \hl{\mbox{\texttt{char}}} hanyalah \hl{\mbox{\texttt{int}}} kecil, yaitu integer yang hanya memakai satu byte ruang, jadi rentangnya terbatas pada\ldots{}

Di sini spesifikasi C jadi sedikit aneh. Ia menjamin bahwa \hl{\mbox{\texttt{char}}} adalah satu byte, yaitu \hl{\mbox{\texttt{sizeof(char)\ ==\ 1}}}. Tapi kemudian di C11 §3.6¶3 ia berusaha keras menyatakan:

\begin{tcolorbox}[blanker,colback=white,enhanced jigsaw,breakable,arc=0pt,outer arc=0pt,grow to left by=-1cm,enlarge top by=0.1cm,enlarge bottom by=0.1cm,borderline west={1.5pt}{-.3cm}{black}]

\begin{quote}
A byte is composed of a contiguous sequence of bits, \emph{the number of which is implementation-defined.}
\end{quote}

\end{tcolorbox}

Tunggu---apa? Mungkin sebagian dari kamu terbiasa dengan anggapan bahwa byte itu 8 bit, kan? Maksudnya, memang begitu, kan? Dan jawabannya: ``Hampir pasti.''\footnote{Istilah industri untuk deretan yang tepat dan pasti 8 bit adalah \emph{oktet}.} Tapi C bahasa yang tua, dan mesin-mesin zaman dulu punya pendapat yang, katakanlah, lebih \emph{santai} soal berapa bit dalam satu byte. Dan selama bertahun-tahun, C mempertahankan fleksibilitas ini.

Tapi dengan asumsi byte di C-mu berukuran 8 bit, seperti hampir semua mesin di dunia yang akan pernah kamu temui, rentang \hl{\mbox{\texttt{char}}} adalah\ldots{}

\index{unsigned char type@\texttt{unsigned char} type|(} \index{signed char type@\texttt{signed char} type|(} ---Nah, sebelum aku bisa menyebutkannya, ternyata \hl{\mbox{\texttt{char}}} bisa signed atau unsigned tergantung compiler-mu. Kecuali kalau kamu menyebutkannya secara eksplisit.

Dalam banyak kasus, cukup memakai \hl{\mbox{\texttt{char}}} karena kamu tidak peduli soal tanda data. Tapi kalau kamu butuh \hl{\mbox{\texttt{char}}} signed atau unsigned, kamu \emph{harus} eksplisit:

\begin{tcolorbox}[enhanced jigsaw,left=3pt,breakable,colback=white!97.2549!black,boxrule=0.2mm]

\begin{Shaded}
\begin{Highlighting}[]
\DataTypeTok{char}\NormalTok{ a}\OperatorTok{;}           \CommentTok{// Bisa signed atau unsigned}
\DataTypeTok{signed} \DataTypeTok{char}\NormalTok{ b}\OperatorTok{;}    \CommentTok{// Pasti signed}
\DataTypeTok{unsigned} \DataTypeTok{char}\NormalTok{ c}\OperatorTok{;}  \CommentTok{// Pasti unsigned}
\end{Highlighting}
\end{Shaded}

\end{tcolorbox}

Oke, sekarang akhirnya kita bisa menentukan rentang angkanya, dengan asumsi \hl{\mbox{\texttt{char}}} berukuran 8 bit dan sistemmu memakai representasi two's complement yang hampir universal untuk signed dan unsigned\footnote{Secara umum, kalau ada bilangan two's complement \(n\) bit, rentang signed-nya adalah \(-2^{n-1}\) sampai \(2^{n-1}-1\). Dan rentang unsigned-nya adalah \(0\) sampai \(2^n-1\).}.

Jadi, dengan asumsi-asumsi itu, akhirnya kita bisa menghitung rentangnya:

{\def\LTcaptype{none} % do not increment counter
\begin{longtable}[]{@{}lrr@{}}
\toprule\noalign{}
Tipe \hl{\mbox{\texttt{char}}} & Minimum & Maksimum \\
\midrule\noalign{}
\endhead
\bottomrule\noalign{}
\endlastfoot
\hl{\mbox{\texttt{signed\ char}}} & \hl{\mbox{\texttt{-128}}} & \hl{\mbox{\texttt{127}}} \\
\hl{\mbox{\texttt{unsigned\ char}}} & \hl{\mbox{\texttt{0}}} & \hl{\mbox{\texttt{255}}} \\
\end{longtable}
}

Dan rentang untuk \hl{\mbox{\texttt{char}}} ditentukan oleh implementasi. \index{unsigned char type@\texttt{unsigned char} type|)} \index{signed char type@\texttt{signed char} type|)}

Biar kuluruskan dulu. \hl{\mbox{\texttt{char}}} sebenarnya angka, berarti bisa dihitung secara matematis?

Bisa! Cuma ingat untuk menjaga nilainya dalam rentang \hl{\mbox{\texttt{char}}}!

\begin{tcolorbox}[enhanced jigsaw,left=3pt,breakable,colback=white!97.2549!black,boxrule=0.2mm]

\begin{Shaded}
\begin{Highlighting}[numbers=left,,]
\PreprocessorTok{\#include }\ImportTok{\textless{}stdio.h\textgreater{}}

\DataTypeTok{int}\NormalTok{ main}\OperatorTok{(}\DataTypeTok{void}\OperatorTok{)}
\OperatorTok{\{}
    \DataTypeTok{char}\NormalTok{ a }\OperatorTok{=} \DecValTok{10}\OperatorTok{,}\NormalTok{ b }\OperatorTok{=} \DecValTok{20}\OperatorTok{;}

\NormalTok{    printf}\OperatorTok{(}\StringTok{"}\SpecialCharTok{\%d\textbackslash{}n}\StringTok{"}\OperatorTok{,}\NormalTok{ a }\OperatorTok{+}\NormalTok{ b}\OperatorTok{);}  \CommentTok{// 30!}
\OperatorTok{\}}
\end{Highlighting}
\end{Shaded}

\end{tcolorbox}

\index{' single quote@\texttt{'} single quote|(} Bagaimana dengan konstanta karakter dalam tanda kutip tunggal, seperti \hl{\mbox{\texttt{\textquotesingle{}B\textquotesingle{}}}}? Bagaimana ia punya nilai numerik?

Spesifikasi juga bersikap mengambang di sini, karena C tidak dirancang untuk berjalan di satu jenis sistem dasar saja.

Tapi anggap saja untuk sementara character set-mu berbasis \href{https://en.wikipedia.org/wiki/ASCII}{ASCII}\footnote{https://en.wikipedia.org/wiki/ASCII} setidaknya untuk 128 karakter pertama. Kalau begitu, konstanta karakter akan dikonversi menjadi \hl{\mbox{\texttt{char}}} yang nilainya sama dengan nilai ASCII karakter tersebut.

Ribet juga, ya. Langsung saja contoh:

\begin{tcolorbox}[enhanced jigsaw,left=3pt,breakable,colback=white!97.2549!black,boxrule=0.2mm]

\begin{Shaded}
\begin{Highlighting}[numbers=left,,]
\PreprocessorTok{\#include }\ImportTok{\textless{}stdio.h\textgreater{}}

\DataTypeTok{int}\NormalTok{ main}\OperatorTok{(}\DataTypeTok{void}\OperatorTok{)}
\OperatorTok{\{}
    \DataTypeTok{char}\NormalTok{ a }\OperatorTok{=} \DecValTok{10}\OperatorTok{;}
    \DataTypeTok{char}\NormalTok{ b }\OperatorTok{=} \CharTok{\textquotesingle{}B\textquotesingle{}}\OperatorTok{;}  \CommentTok{// Nilai ASCII 66}

\NormalTok{    printf}\OperatorTok{(}\StringTok{"}\SpecialCharTok{\%d\textbackslash{}n}\StringTok{"}\OperatorTok{,}\NormalTok{ a }\OperatorTok{+}\NormalTok{ b}\OperatorTok{);}  \CommentTok{// 76!}
\OperatorTok{\}}
\end{Highlighting}
\end{Shaded}

\end{tcolorbox}

Ini bergantung pada lingkungan eksekusimu dan \href{https://en.wikipedia.org/wiki/List_of_information_system_character_sets}{character set yang dipakai}\footnote{https://en.wikipedia.org/wiki/List\_of\_information\_system\_character\_sets}. Salah satu character set paling populer saat ini adalah \href{https://en.wikipedia.org/wiki/Unicode}{Unicode}\footnote{https://en.wikipedia.org/wiki/Unicode} (yang merupakan superset ASCII), jadi untuk 0-9, A-Z, a-z dan tanda baca dasar, hampir pasti kamu mendapat nilai ASCII-nya. \index{' single quote@\texttt{'} single quote|)} \index{char type@\texttt{char} type|)} \index{Types!character|)}

\section{\texorpdfstring{Tipe Integer Lainnya: \hl{\mbox{\texttt{short}}}, \hl{\mbox{\texttt{long}}}, \hl{\mbox{\texttt{long\ long}}}}{Tipe Integer Lainnya: short, long, long long}}\label{tipe-integer-lainnya-short-long-long-long}

Sejauh ini kita hanya memakai dua tipe integer secara umum:

\begin{itemize}
\tightlist
\item
  \hl{\mbox{\texttt{char}}}
\item
  \hl{\mbox{\texttt{int}}}
\end{itemize}

dan baru-baru ini kita belajar varian unsigned dari tipe integer. Kita juga tahu bahwa \hl{\mbox{\texttt{char}}} diam-diam adalah \hl{\mbox{\texttt{int}}} kecil yang menyamar. Jadi kita tahu \hl{\mbox{\texttt{int}}} bisa punya berbagai ukuran bit.

Tapi ada beberapa tipe integer lagi yang perlu kita lihat, beserta rentang minimum yang bisa mereka tampung. (Implementasimu kemungkinan punya rentang lebih lebar daripada yang disyaratkan spesifikasi, tapi rentang di sini adalah yang \textbf{pasti} tersedia secara portabel.)

Header file \hl{\mbox{\texttt{\textless{}limits.h\textgreater{}}}} mendefinisikan makro yang menyimpan rentang berbagai tipe; andalkan itu untuk kepastian, dan \emph{jangan pernah meng-hardcode atau mengasumsikan nilai-nilai ini}.

\index{short type@\texttt{short} type|(} \index{long type@\texttt{long} type|(} \index{long long type@\texttt{long long} type|(} Tipe tambahan ini adalah \hl{\mbox{\texttt{short\ int}}}, \hl{\mbox{\texttt{long\ int}}}, dan \hl{\mbox{\texttt{long\ long\ int}}}. Biasanya, saat memakai tipe-tipe ini, developer C menghilangkan bagian \hl{\mbox{\texttt{int}}} (misalnya \hl{\mbox{\texttt{long\ long}}}), dan compiler tetap senang.

\begin{tcolorbox}[enhanced jigsaw,left=3pt,breakable,colback=white!97.2549!black,boxrule=0.2mm]

\begin{Shaded}
\begin{Highlighting}[]
\CommentTok{// Dua baris ini setara:}
\DataTypeTok{long} \DataTypeTok{long} \DataTypeTok{int}\NormalTok{ x}\OperatorTok{;}
\DataTypeTok{long} \DataTypeTok{long}\NormalTok{ x}\OperatorTok{;}

\CommentTok{// Begitu juga dua baris ini:}
\DataTypeTok{short} \DataTypeTok{int}\NormalTok{ x}\OperatorTok{;}
\DataTypeTok{short}\NormalTok{ x}\OperatorTok{;}
\end{Highlighting}
\end{Shaded}

\end{tcolorbox}

Mari kita lihat tipe data integer dan ukurannya secara menaik, dikelompokkan menurut tanda. Sekali lagi, batas minimum dan maksimum ini adalah yang dijamin secara portabel oleh spesifikasi; sistemmu kemungkinan punya rentang lebih lebar.

{\def\LTcaptype{none} % do not increment counter
\begin{longtable}[]{@{}
  >{\raggedright\arraybackslash}p{(\linewidth - 6\tabcolsep) * \real{0.2500}}
  >{\raggedleft\arraybackslash}p{(\linewidth - 6\tabcolsep) * \real{0.2500}}
  >{\raggedleft\arraybackslash}p{(\linewidth - 6\tabcolsep) * \real{0.2500}}
  >{\raggedleft\arraybackslash}p{(\linewidth - 6\tabcolsep) * \real{0.2500}}@{}}
\toprule\noalign{}
\begin{minipage}[b]{\linewidth}\raggedright
Tipe
\end{minipage} & \begin{minipage}[b]{\linewidth}\raggedleft
Byte Minimum
\end{minipage} & \begin{minipage}[b]{\linewidth}\raggedleft
Nilai Minimum
\end{minipage} & \begin{minipage}[b]{\linewidth}\raggedleft
Nilai Maksimum
\end{minipage} \\
\midrule\noalign{}
\endhead
\bottomrule\noalign{}
\endlastfoot
\hl{\mbox{\texttt{char}}} & 1 & -128 atau 0 & 127 atau 255\footnote{Tergantung apakah \hl{\mbox{\texttt{char}}} secara bawaan berupa \hl{\mbox{\texttt{signed\ char}}} atau \hl{\mbox{\texttt{unsigned\ char}}}} \\
\hl{\mbox{\texttt{signed\ char}}} & 1 & -128 & 127 \\
\hl{\mbox{\texttt{short}}} & 2 & -32768 & 32767 \\
\hl{\mbox{\texttt{int}}} & 2 & -32768 & 32767 \\
\hl{\mbox{\texttt{long}}} & 4 & -2147483648 & 2147483647 \\
\hl{\mbox{\texttt{long\ long}}} & 8 & -9223372036854775808 & 9223372036854775807 \\
\hl{\mbox{\texttt{unsigned\ char}}} & 1 & 0 & 255 \\
\hl{\mbox{\texttt{unsigned\ short}}} & 2 & 0 & 65535 \\
\hl{\mbox{\texttt{unsigned\ int}}} & 2 & 0 & 65535 \\
\hl{\mbox{\texttt{unsigned\ long}}} & 4 & 0 & 4294967295 \\
\hl{\mbox{\texttt{unsigned\ long\ long}}} & 8 & 0 & 18446744073709551615 \\
\end{longtable}
}

Tidak ada tipe \hl{\mbox{\texttt{long\ long\ long}}}. Kamu tidak bisa terus menambah \hl{\mbox{\texttt{long}}} seperti itu. Jangan konyol.

\begin{tcolorbox}[blanker,colback=white,enhanced jigsaw,breakable,arc=0pt,outer arc=0pt,grow to left by=-1cm,enlarge top by=0.1cm,enlarge bottom by=0.1cm,borderline west={1.5pt}{-.3cm}{black}]

\begin{quote}
Penggemar two's complement mungkin menyadari ada yang janggal pada angka-angka itu. Kenapa, misalnya, \hl{\mbox{\texttt{signed\ char}}} berhenti di -127 dan bukan -128? Ingat: itu hanya minimum yang disyaratkan spesifikasi. Beberapa representasi angka (seperti \href{https://en.wikipedia.org/wiki/Signed_number_representations\#Signed_magnitude_representation}{sign and magnitude}\footnote{https://en.wikipedia.org/wiki/Signed\_number\_representations\#Signed\_magnitude\_representation}) mentok di ±127.
\end{quote}

\end{tcolorbox}

Mari kita jalankan tabel yang sama pada sistem 64-bit two's complement milikku dan lihat hasilnya:

{\def\LTcaptype{none} % do not increment counter
\begin{longtable}[]{@{}lrrr@{}}
\toprule\noalign{}
Tipe & Byte di Sistemku & Nilai Minimum & Nilai Maksimum \\
\midrule\noalign{}
\endhead
\bottomrule\noalign{}
\endlastfoot
\hl{\mbox{\texttt{char}}} & 1 & -128 & 127\footnote{\hl{\mbox{\texttt{char}}} di sistemku adalah signed.} \\
\hl{\mbox{\texttt{signed\ char}}} & 1 & -128 & 127 \\
\hl{\mbox{\texttt{short}}} & 2 & -32768 & 32767 \\
\hl{\mbox{\texttt{int}}} & 4 & -2147483648 & 2147483647 \\
\hl{\mbox{\texttt{long}}} & 8 & -9223372036854775808 & 9223372036854775807 \\
\hl{\mbox{\texttt{long\ long}}} & 8 & -9223372036854775808 & 9223372036854775807 \\
\hl{\mbox{\texttt{unsigned\ char}}} & 1 & 0 & 255 \\
\hl{\mbox{\texttt{unsigned\ short}}} & 2 & 0 & 65535 \\
\hl{\mbox{\texttt{unsigned\ int}}} & 4 & 0 & 4294967295 \\
\hl{\mbox{\texttt{unsigned\ long}}} & 8 & 0 & 18446744073709551615 \\
\hl{\mbox{\texttt{unsigned\ long\ long}}} & 8 & 0 & 18446744073709551615 \\
\end{longtable}
}

Ini sedikit lebih masuk akal, tapi kita bisa melihat bahwa sistemku punya batas yang lebih besar daripada minimum di spesifikasi.

\protect\phantomsection\label{limits-macros}{Lalu apa saja makro di \hl{\mbox{\texttt{\textless{}limits.h\textgreater{}}}}?}

{\def\LTcaptype{none} % do not increment counter
\begin{longtable}[]{@{}lll@{}}
\toprule\noalign{}
Tipe & Makro Min & Makro Maks \\
\midrule\noalign{}
\endhead
\bottomrule\noalign{}
\endlastfoot
\hl{\mbox{\texttt{char}}} & \hl{\mbox{\texttt{CHAR\_MIN}}} & \hl{\mbox{\texttt{CHAR\_MAX}}} \\
\hl{\mbox{\texttt{signed\ char}}} & \hl{\mbox{\texttt{SCHAR\_MIN}}} & \hl{\mbox{\texttt{SCHAR\_MAX}}} \\
\hl{\mbox{\texttt{short}}} & \hl{\mbox{\texttt{SHRT\_MIN}}} & \hl{\mbox{\texttt{SHRT\_MAX}}} \\
\hl{\mbox{\texttt{int}}} & \hl{\mbox{\texttt{INT\_MIN}}} & \hl{\mbox{\texttt{INT\_MAX}}} \\
\hl{\mbox{\texttt{long}}} & \hl{\mbox{\texttt{LONG\_MIN}}} & \hl{\mbox{\texttt{LONG\_MAX}}} \\
\hl{\mbox{\texttt{long\ long}}} & \hl{\mbox{\texttt{LLONG\_MIN}}} & \hl{\mbox{\texttt{LLONG\_MAX}}} \\
\hl{\mbox{\texttt{unsigned\ char}}} & \hl{\mbox{\texttt{0}}} & \hl{\mbox{\texttt{UCHAR\_MAX}}} \\
\hl{\mbox{\texttt{unsigned\ short}}} & \hl{\mbox{\texttt{0}}} & \hl{\mbox{\texttt{USHRT\_MAX}}} \\
\hl{\mbox{\texttt{unsigned\ int}}} & \hl{\mbox{\texttt{0}}} & \hl{\mbox{\texttt{UINT\_MAX}}} \\
\hl{\mbox{\texttt{unsigned\ long}}} & \hl{\mbox{\texttt{0}}} & \hl{\mbox{\texttt{ULONG\_MAX}}} \\
\hl{\mbox{\texttt{unsigned\ long\ long}}} & \hl{\mbox{\texttt{0}}} & \hl{\mbox{\texttt{ULLONG\_MAX}}} \\
\end{longtable}
}

Perhatikan ada cara tersembunyi di sana untuk menentukan apakah sebuah sistem memakai \hl{\mbox{\texttt{char}}} signed atau unsigned. Kalau \hl{\mbox{\texttt{CHAR\_MAX\ ==\ UCHAR\_MAX}}}, berarti ia unsigned.

Perhatikan juga bahwa tidak ada makro minimum untuk varian \hl{\mbox{\texttt{unsigned}}}---nilainya cuma \hl{\mbox{\texttt{0}}}. \index{short type@\texttt{short} type|)} \index{long type@\texttt{long} type|)} \index{long long type@\texttt{long long} type|)}

\section{\texorpdfstring{Float Lainnya: \hl{\mbox{\texttt{double}}} dan \hl{\mbox{\texttt{long\ double}}}}{Float Lainnya: double dan long double}}\label{float-lainnya-double-dan-long-double}

Mari kita lihat apa kata spesifikasi tentang bilangan floating point di §5.2.4.2.2¶1-2:

\begin{tcolorbox}[blanker,colback=white,enhanced jigsaw,breakable,arc=0pt,outer arc=0pt,grow to left by=-1cm,enlarge top by=0.1cm,enlarge bottom by=0.1cm,borderline west={1.5pt}{-.3cm}{black}]

\begin{quote}
The following parameters are used to define the model for each floating-point type:

{\def\LTcaptype{none} % do not increment counter
\begin{longtable}[]{@{}
  >{\raggedright\arraybackslash}p{(\linewidth - 2\tabcolsep) * \real{0.5000}}
  >{\raggedright\arraybackslash}p{(\linewidth - 2\tabcolsep) * \real{0.5000}}@{}}
\toprule\noalign{}
\begin{minipage}[b]{\linewidth}\raggedright
Parameter
\end{minipage} & \begin{minipage}[b]{\linewidth}\raggedright
Definition
\end{minipage} \\
\midrule\noalign{}
\endhead
\bottomrule\noalign{}
\endlastfoot
\(s\) & sign (\(\pm1\)) \\
\(b\) & base or radix of exponent representation (an integer \(> 1\)) \\
\(e\) & exponent (an integer between a minimum \(e_{min}\) and a maximum \(e_{max}\)) \\
\(p\) & precision (the number of base-\(b\) digits in the significand) \\
\(f_k\) & nonnegative integers less than \(b\) (the significand digits) \\
\end{longtable}
}

A \emph{floating-point number} (\(x\)) is defined by the following model:

\begin{tcolorbox}[blanker,colback=white,enhanced jigsaw,breakable,arc=0pt,outer arc=0pt,grow to left by=-1cm,enlarge top by=0.1cm,enlarge bottom by=0.1cm,borderline west={1.5pt}{-.3cm}{black}]

\begin{quote}
\(x=sb^e\sum\limits_{k=1}^p f_kb^{-k},\)~~~~\(e_{min}\le e\le e_{max}\)
\end{quote}

\end{tcolorbox}
\end{quote}

\end{tcolorbox}

Kuharap itu langsung bikin semuanya jelas.

Baiklah. Mundur selangkah dan lihat apa yang praktis.

Catatan: kita merujuk ke sejumlah makro di bagian ini. Makro-makro itu ada di header \hl{\mbox{\texttt{\textless{}float.h\textgreater{}}}}.

Bilangan floating point di-encode dalam urutan bit tertentu (\href{https://en.wikipedia.org/wiki/IEEE_754}{format IEEE-754}\footnote{https://en.wikipedia.org/wiki/IEEE\_754} sangat populer) di dalam byte.

Lebih dalam lagi, bilangan itu pada dasarnya direpresentasikan sebagai \emph{significand} (bagian angkanya---digit-digit signifikannya sendiri, kadang juga disebut \emph{mantissa}) dan \emph{exponent}, yaitu pangkat yang dipakai untuk menaikkan digit-digit itu. Ingat bahwa eksponen negatif bisa membuat angka menjadi lebih kecil.

Bayangkan kita memakai \(10\) sebagai bilangan yang dipangkatkan. Kita bisa merepresentasikan angka-angka berikut dengan significand \(12345\) dan eksponen \(-3\), \(4\), dan \(0\) untuk meng-encode nilai floating point berikut:

\(12345\times10^{-3}=12.345\)

\(12345\times10^4=123450000\)

\(12345\times10^0=12345\)

Untuk semua angka itu, significand-nya tetap sama. Satu-satunya perbedaan adalah eksponennya.

Di mesinmu, basis untuk eksponen kemungkinan adalah \(2\), bukan \(10\), karena komputer suka biner. Kamu bisa mengeceknya dengan mencetak makro \hl{\mbox{\texttt{FLT\_RADIX}}}.

Jadi kita punya sebuah angka yang direpresentasikan oleh sejumlah byte, yang di-encode dengan cara tertentu. Karena pola bit terbatas, jumlah bilangan floating point yang bisa direpresentasikan juga terbatas.

Lebih khusus lagi, hanya sejumlah tertentu digit desimal signifikan yang bisa direpresentasikan secara akurat.

Bagaimana cara mendapat lebih banyak? Pakai tipe data yang lebih besar!

Dan kita punya beberapa. Kita sudah kenal \hl{\mbox{\texttt{float}}}, tapi untuk presisi lebih besar ada \hl{\mbox{\texttt{double}}}. Dan untuk presisi yang lebih besar lagi, ada \hl{\mbox{\texttt{long\ double}}} (tidak berhubungan dengan \hl{\mbox{\texttt{long\ int}}} kecuali namanya).

Spesifikasi tidak menentukan berapa byte penyimpanan yang dipakai setiap tipe, tapi di sistemku, kita bisa melihat kenaikan ukuran relatifnya:

\index{double type@\texttt{double} type|(} \index{long double type@\texttt{long double} type|(}

{\def\LTcaptype{none} % do not increment counter
\begin{longtable}[]{@{}lr@{}}
\toprule\noalign{}
Tipe & \hl{\mbox{\texttt{sizeof}}} \\
\midrule\noalign{}
\endhead
\bottomrule\noalign{}
\endlastfoot
\hl{\mbox{\texttt{float}}} & 4 \\
\hl{\mbox{\texttt{double}}} & 8 \\
\hl{\mbox{\texttt{long\ double}}} & 16 \\
\end{longtable}
}

Jadi setiap tipe (di sistemku) memakai bit tambahan itu untuk presisi yang lebih tinggi.

Tapi \emph{seberapa besar} presisinya? Berapa banyak angka desimal yang bisa direpresentasikan oleh nilai-nilai ini?

Nah, C menyediakan sejumlah makro di \hl{\mbox{\texttt{\textless{}float.h\textgreater{}}}} untuk membantu kita mengetahuinya.

Ini jadi agak aneh kalau kamu memakai sistem basis-2 (biner) untuk menyimpan angka (yang dipakai hampir semua orang di planet ini, mungkin termasuk kamu), tapi sabar saja sementara kita mengurainya. \index{double type@\texttt{double} type|)} \index{long double type@\texttt{long double} type|)}

\subsection{Berapa Banyak Digit Desimal?}\label{berapa-banyak-digit-desimal}

\index{Significant digits|(} Pertanyaan sejuta dolarnya adalah, ``Berapa banyak digit desimal signifikan yang bisa disimpan dalam tipe floating point tertentu sehingga ketika dicetak, kita mendapat angka desimal yang sama?''

Jumlah digit desimal yang bisa disimpan dalam tipe floating point dan pasti kembali sama saat dicetak diberikan oleh makro-makro ini:

\index{FLT\_DIG macro@\texttt{FLT\_DIG} macro|(} \index{DBL\_DIG macro@\texttt{DBL\_DIG} macro|(} \index{LDBL\_DIG macro@\texttt{LDBL\_DIG} macro|(}

{\def\LTcaptype{none} % do not increment counter
\begin{longtable}[]{@{}lrr@{}}
\toprule\noalign{}
Tipe & Digit Desimal yang Bisa Disimpan & Minimum \\
\midrule\noalign{}
\endhead
\bottomrule\noalign{}
\endlastfoot
\hl{\mbox{\texttt{float}}} & \hl{\mbox{\texttt{FLT\_DIG}}} & 6 \\
\hl{\mbox{\texttt{double}}} & \hl{\mbox{\texttt{DBL\_DIG}}} & 10 \\
\hl{\mbox{\texttt{long\ double}}} & \hl{\mbox{\texttt{LDBL\_DIG}}} & 10 \\
\end{longtable}
}

\index{DBL\_DIG macro@\texttt{DBL\_DIG} macro|)} \index{LDBL\_DIG macro@\texttt{LDBL\_DIG} macro|)}

Di sistemku, \hl{\mbox{\texttt{FLT\_DIG}}} adalah 6, jadi aku bisa yakin kalau sebuah \hl{\mbox{\texttt{float}}} 6 digit dicetak, hasilnya sama. (Bisa saja lebih dari 6 digit---sebagian angka tetap kembali dengan benar dengan lebih banyak digit. Tapi 6 pasti kembali.)

Misalnya, kalau kita mencetak \hl{\mbox{\texttt{float}}} dengan pola digit yang bertambah seperti ini, kelihatannya kita sampai 8 digit sebelum ada yang salah, tapi setelah itu kembali ke 7 digit yang benar.

\begin{tcolorbox}[enhanced jigsaw,left=3pt,breakable,colback=white!97.2549!black,boxrule=0.2mm]

\begin{Shaded}
\begin{Highlighting}[]
\NormalTok{0.12345}
\NormalTok{0.123456}
\NormalTok{0.1234567}
\NormalTok{0.12345678}
\NormalTok{0.123456791  \textless{}{-}{-} Things start going wrong}
\NormalTok{0.1234567910}
\end{Highlighting}
\end{Shaded}

\end{tcolorbox}

Kita coba demo lain. Pada kode ini ada dua \hl{\mbox{\texttt{float}}} yang masing-masing menyimpan angka dengan \hl{\mbox{\texttt{FLT\_DIG}}} digit desimal signifikan\footnote{Program ini berjalan sesuai komentarnya pada sistem dengan \hl{\mbox{\texttt{FLT\_DIG}}} bernilai \hl{\mbox{\texttt{6}}} yang memakai bilangan floating point IEEE-754 basis-2. Kalau tidak, hasilnya mungkin berbeda.}. Lalu keduanya dijumlahkan, yang seharusnya menghasilkan 12 digit desimal signifikan. Tapi itu lebih dari yang bisa disimpan dalam \hl{\mbox{\texttt{float}}} dan dipulihkan dengan benar sebagai string---jadi ketika dicetak, terlihat ada yang mulai salah setelah digit signifikan ke-7.

\begin{tcolorbox}[enhanced jigsaw,left=3pt,breakable,colback=white!97.2549!black,boxrule=0.2mm]

\begin{Shaded}
\begin{Highlighting}[numbers=left,,]
\PreprocessorTok{\#include }\ImportTok{\textless{}stdio.h\textgreater{}}
\PreprocessorTok{\#include }\ImportTok{\textless{}float.h\textgreater{}}

\DataTypeTok{int}\NormalTok{ main}\OperatorTok{(}\DataTypeTok{void}\OperatorTok{)}
\OperatorTok{\{}
    \CommentTok{// Kedua angka ini punya 6 digit signifikan, jadi bisa}
    \CommentTok{// disimpan secara akurat dalam float:}

    \DataTypeTok{float}\NormalTok{ f }\OperatorTok{=} \FloatTok{3.14159}\BuiltInTok{f}\OperatorTok{;}
    \DataTypeTok{float}\NormalTok{ g }\OperatorTok{=} \FloatTok{0.00000265358}\BuiltInTok{f}\OperatorTok{;}

\NormalTok{    printf}\OperatorTok{(}\StringTok{"}\SpecialCharTok{\%.5f\textbackslash{}n}\StringTok{"}\OperatorTok{,}\NormalTok{ f}\OperatorTok{);}   \CommentTok{// 3.14159       {-}{-} benar!}
\NormalTok{    printf}\OperatorTok{(}\StringTok{"}\SpecialCharTok{\%.11f\textbackslash{}n}\StringTok{"}\OperatorTok{,}\NormalTok{ g}\OperatorTok{);}  \CommentTok{// 0.00000265358 {-}{-} benar!}

    \CommentTok{// Sekarang jumlahkan}
\NormalTok{    f }\OperatorTok{+=}\NormalTok{ g}\OperatorTok{;}                \CommentTok{// 3.14159265358 adalah nilai f yang \_seharusnya\_}

\NormalTok{    printf}\OperatorTok{(}\StringTok{"}\SpecialCharTok{\%.11f\textbackslash{}n}\StringTok{"}\OperatorTok{,}\NormalTok{ f}\OperatorTok{);}  \CommentTok{// 3.14159274101 {-}{-} salah!}
\OperatorTok{\}}
\end{Highlighting}
\end{Shaded}

\end{tcolorbox}

(Kode di atas punya \hl{\mbox{\texttt{f}}} setelah konstanta numeriknya---ini menandakan bahwa konstanta itu bertipe \hl{\mbox{\texttt{float}}}, bukan \hl{\mbox{\texttt{double}}} yang menjadi bawaannya. Soal ini dibahas nanti.)

Ingat bahwa \hl{\mbox{\texttt{FLT\_DIG}}} adalah jumlah digit aman yang bisa disimpan dalam \hl{\mbox{\texttt{float}}} dan diambil kembali dengan benar.

Kadang kamu bisa mendapat satu atau dua digit lebih. Tapi kadang kamu hanya mendapat \hl{\mbox{\texttt{FLT\_DIG}}} digit. Yang pasti: kalau kamu menyimpan angka dengan jumlah digit sampai dengan \hl{\mbox{\texttt{FLT\_DIG}}} dalam \hl{\mbox{\texttt{float}}}, dijamin kembali dengan benar.

Jadi begitulah ceritanya. \hl{\mbox{\texttt{FLT\_DIG}}}. Tamat.

\index{FLT\_DIG macro@\texttt{FLT\_DIG} macro|)}

\ldots Atau belum?

\subsection{Mengonversi ke Desimal dan Kembali}\label{mengonversi-ke-desimal-dan-kembali}

Tapi menyimpan bilangan basis 10 dalam bilangan floating point lalu mengambilnya kembali baru setengah cerita.

Ternyata bilangan floating point bisa meng-encode angka yang butuh lebih banyak angka desimal untuk dicetak sepenuhnya. Hanya saja angka desimal besarmu mungkin tidak cocok dengan salah satu angka itu.

Artinya, kalau kamu melihat bilangan floating point dari satu ke yang berikutnya, ada celah di antaranya. Kalau kamu mencoba meng-encode angka desimal di celah itu, yang dipakai adalah bilangan floating point terdekat. Itulah kenapa kamu hanya bisa meng-encode \hl{\mbox{\texttt{FLT\_DIG}}} pada \hl{\mbox{\texttt{float}}}.

Tapi bagaimana dengan bilangan floating point yang \emph{tidak} ada di celah? Berapa tempat desimal yang dibutuhkan untuk mencetaknya secara akurat?

Cara lain menyatakan pertanyaan ini: untuk bilangan floating point tertentu, berapa digit desimal yang harus dipertahankan kalau kita ingin mengonversi angka desimal itu kembali menjadi bilangan floating point yang identik? Artinya, berapa digit yang harus dicetak di basis 10 untuk memulihkan \textbf{semua} digit basis 2 pada angka aslinya?

Kadang hanya perlu beberapa. Tapi supaya yakin, kamu perlu mengonversi ke desimal dengan jumlah tempat desimal aman tertentu. Angka itu di-encode dalam makro-makro berikut:

\index{FLT\_DECMIAL\_DIG macro@\texttt{FLT\_DECMIAL\_DIG} macro|(} \index{DBL\_DECMIAL\_DIG macro@\texttt{DBL\_DECMIAL\_DIG} macro|(} \index{LDBL\_DECMIAL\_DIG macro@\texttt{LDBL\_DECMIAL\_DIG} macro|(} \index{DECMIAL\_DIG macro@\texttt{DECMIAL\_DIG} macro|(}

{\def\LTcaptype{none} % do not increment counter
\begin{longtable}[]{@{}
  >{\raggedright\arraybackslash}p{(\linewidth - 2\tabcolsep) * \real{0.1774}}
  >{\raggedright\arraybackslash}p{(\linewidth - 2\tabcolsep) * \real{0.8226}}@{}}
\toprule\noalign{}
\begin{minipage}[b]{\linewidth}\raggedright
Makro
\end{minipage} & \begin{minipage}[b]{\linewidth}\raggedright
Deskripsi
\end{minipage} \\
\midrule\noalign{}
\endhead
\bottomrule\noalign{}
\endlastfoot
\hl{\mbox{\texttt{FLT\_DECIMAL\_DIG}}} & Jumlah digit desimal yang di-encode dalam \hl{\mbox{\texttt{float}}}. \\
\hl{\mbox{\texttt{DBL\_DECIMAL\_DIG}}} & Jumlah digit desimal yang di-encode dalam \hl{\mbox{\texttt{double}}}. \\
\hl{\mbox{\texttt{LDBL\_DECIMAL\_DIG}}} & Jumlah digit desimal yang di-encode dalam \hl{\mbox{\texttt{long\ double}}}. \\
\hl{\mbox{\texttt{DECIMAL\_DIG}}} & Sama dengan encoding terlebar, \hl{\mbox{\texttt{LDBL\_DECIMAL\_DIG}}}. \\
\end{longtable}
}

\index{LDBL\_DECMIAL\_DIG macro@\texttt{LDBL\_DECMIAL\_DIG} macro|)} \index{DECMIAL\_DIG macro@\texttt{DECMIAL\_DIG} macro|)}

Mari kita lihat contoh ketika \hl{\mbox{\texttt{DBL\_DIG}}} adalah 15 (jadi hanya segitu yang bisa ada di konstanta), tapi \hl{\mbox{\texttt{DBL\_DECIMAL\_DIG}}} adalah 17 (jadi kita harus mengonversi ke 17 angka desimal untuk mempertahankan semua bit dari \hl{\mbox{\texttt{double}}} aslinya).

Kita tugaskan angka 15 digit signifikan \hl{\mbox{\texttt{0.123456789012345}}} ke \hl{\mbox{\texttt{x}}}, dan angka 1 digit signifikan \hl{\mbox{\texttt{0.0000000000000006}}} ke \hl{\mbox{\texttt{y}}}.

\begin{tcolorbox}[enhanced jigsaw,left=3pt,breakable,colback=white!97.2549!black,boxrule=0.2mm]

\begin{Shaded}
\begin{Highlighting}[]
\NormalTok{x is exact: 0.12345678901234500    Printed to 17 decimal places}
\NormalTok{y is exact: 0.00000000000000060}
\end{Highlighting}
\end{Shaded}

\end{tcolorbox}

Tapi mari kita jumlahkan. Hasilnya seharusnya \hl{\mbox{\texttt{0.1234567890123456}}}, tapi itu lebih dari \hl{\mbox{\texttt{DBL\_DIG}}}, jadi mungkin terjadi hal-hal aneh\ldots{} kita lihat:

\begin{tcolorbox}[enhanced jigsaw,left=3pt,breakable,colback=white!97.2549!black,boxrule=0.2mm]

\begin{Shaded}
\begin{Highlighting}[]
\NormalTok{x + y not quite right: 0.12345678901234559    Should end in 4560!}
\end{Highlighting}
\end{Shaded}

\end{tcolorbox}

Itulah yang kita dapat karena mencetak lebih dari \hl{\mbox{\texttt{DBL\_DIG}}}, kan? Tapi lihat ini\ldots{} angka di atas itu bisa direpresentasikan persis seperti adanya!

Kalau kita tugaskan \hl{\mbox{\texttt{0.12345678901234559}}} (17 digit) ke \hl{\mbox{\texttt{z}}} dan mencetaknya, hasilnya:

\begin{tcolorbox}[enhanced jigsaw,left=3pt,breakable,colback=white!97.2549!black,boxrule=0.2mm]

\begin{Shaded}
\begin{Highlighting}[]
\NormalTok{z is exact: 0.12345678901234559   17 digits correct! More than DBL\_DIG!}
\end{Highlighting}
\end{Shaded}

\end{tcolorbox}

Kalau \hl{\mbox{\texttt{z}}} dipotong menjadi 15 digit, angkanya tidak akan sama. Itulah kenapa untuk mempertahankan semua bit sebuah \hl{\mbox{\texttt{double}}}, kita butuh \hl{\mbox{\texttt{DBL\_DECIMAL\_DIG}}} dan bukan sekadar \hl{\mbox{\texttt{DBL\_DIG}}} yang lebih kecil.

\index{DBL\_DECMIAL\_DIG macro@\texttt{DBL\_DECMIAL\_DIG} macro|)}

Jadi, jelas bahwa saat bermain dengan angka desimal secara umum, tidak aman mencetak lebih dari \index{FLT\_DIG macro@\texttt{FLT\_DIG} macro}\hl{\mbox{\texttt{FLT\_DIG}}}, \index{DBL\_DIG macro@\texttt{DBL\_DIG} macro}\hl{\mbox{\texttt{DBL\_DIG}}}, atau \index{LDBL\_DIG macro@\texttt{LDBL\_DIG} macro}\hl{\mbox{\texttt{LDBL\_DIG}}} digit agar tetap masuk akal terhadap angka basis 10 aslinya dan perhitungan matematika berikutnya.

Tapi saat mengonversi dari \hl{\mbox{\texttt{float}}} ke representasi desimal dan \emph{kembali} ke \hl{\mbox{\texttt{float}}}, pasti pakai \hl{\mbox{\texttt{FLT\_DECIMAL\_DIG}}} agar semua bit terpertahankan persis. \index{FLT\_DECMIAL\_DIG macro@\texttt{FLT\_DECMIAL\_DIG} macro|)}

\index{Significant digits|)}

\section{Tipe Konstanta Numerik}\label{tipe-konstanta-numerik}

Saat kamu menulis konstanta angka, seperti \hl{\mbox{\texttt{1234}}}, ia punya tipe. Tapi tipe apa? Mari kita lihat bagaimana C menentukan tipe konstanta itu, dan cara memaksanya memilih tipe tertentu.

\subsection{Heksadesimal dan Oktal}\label{heksadesimal-dan-oktal}

Selain desimal biasa seperti buatan nenek, C juga mendukung konstanta dengan basis berbeda.

\index{0x hexadecimal@\texttt{0x} hexadecimal|(}

Kalau angka diawali \hl{\mbox{\texttt{0x}}}, ia dibaca sebagai angka hex:

\begin{tcolorbox}[enhanced jigsaw,left=3pt,breakable,colback=white!97.2549!black,boxrule=0.2mm]

\begin{Shaded}
\begin{Highlighting}[]
\DataTypeTok{int}\NormalTok{ a }\OperatorTok{=} \BaseNTok{0x1A2B}\OperatorTok{;}   \CommentTok{// Heksadesimal}
\DataTypeTok{int}\NormalTok{ b }\OperatorTok{=} \BaseNTok{0x1a2b}\OperatorTok{;}   \CommentTok{// Huruf besar/kecil tidak berpengaruh untuk digit hex}

\NormalTok{printf}\OperatorTok{(}\StringTok{"}\SpecialCharTok{\%x}\StringTok{"}\OperatorTok{,}\NormalTok{ a}\OperatorTok{);}  \CommentTok{// Cetak angka hex, "1a2b"}
\end{Highlighting}
\end{Shaded}

\end{tcolorbox}

\index{0x hexadecimal@\texttt{0x} hexadecimal|)}

\index{0 octal@\texttt{0} octal|(} Kalau angka diawali \hl{\mbox{\texttt{0}}}, ia dibaca sebagai angka oktal:

\begin{tcolorbox}[enhanced jigsaw,left=3pt,breakable,colback=white!97.2549!black,boxrule=0.2mm]

\begin{Shaded}
\begin{Highlighting}[]
\DataTypeTok{int}\NormalTok{ a }\OperatorTok{=} \BaseNTok{012}\OperatorTok{;}

\NormalTok{printf}\OperatorTok{(}\StringTok{"}\SpecialCharTok{\%o\textbackslash{}n}\StringTok{"}\OperatorTok{,}\NormalTok{ a}\OperatorTok{);}  \CommentTok{// Cetak angka oktal, "12"}
\end{Highlighting}
\end{Shaded}

\end{tcolorbox}

Ini terutama jadi masalah bagi programmer pemula yang mencoba mengisi angka desimal dengan \hl{\mbox{\texttt{0}}} di kiri agar rapi dan sejajar, dan tanpa sadar mengubah basis angkanya:

\begin{tcolorbox}[enhanced jigsaw,left=3pt,breakable,colback=white!97.2549!black,boxrule=0.2mm]

\begin{Shaded}
\begin{Highlighting}[]
\DataTypeTok{int}\NormalTok{ x }\OperatorTok{=} \DecValTok{11111}\OperatorTok{;}  \CommentTok{// Desimal 11111}
\DataTypeTok{int}\NormalTok{ y }\OperatorTok{=} \BaseNTok{00111}\OperatorTok{;}  \CommentTok{// Desimal 73 (Oktal 111)}
\DataTypeTok{int}\NormalTok{ z }\OperatorTok{=} \BaseNTok{01111}\OperatorTok{;}  \CommentTok{// Desimal 585 (Oktal 1111)}
\end{Highlighting}
\end{Shaded}

\end{tcolorbox}

\index{0 octal@\texttt{0} octal|)}

\subsubsection{Catatan tentang Biner}\label{catatan-tentang-biner}

\index{0b binary@\texttt{0b} binary|(}

Ekstensi tidak resmi\footnote{Aku sungguh heran C belum memasukkan ini ke spesifikasi. Di dokumen C99 Rationale, tertulis, ``A proposal to add binary constants was rejected due to lack of precedent and insufficient utility.'' Padahal itu terdengar agak konyol, mengingat fitur-fitur lain yang dijejalkan ke sana! Aku berani taruhan salah satu rilis berikutnya akan memilikinya.} di banyak compiler C memungkinkan kita menulis bilangan biner dengan awalan \hl{\mbox{\texttt{0b}}}:

\begin{tcolorbox}[enhanced jigsaw,left=3pt,breakable,colback=white!97.2549!black,boxrule=0.2mm]

\begin{Shaded}
\begin{Highlighting}[]
\DataTypeTok{int}\NormalTok{ x }\OperatorTok{=} \BaseNTok{0b101010}\OperatorTok{;}    \CommentTok{// Biner 101010}

\NormalTok{printf}\OperatorTok{(}\StringTok{"}\SpecialCharTok{\%d\textbackslash{}n}\StringTok{"}\OperatorTok{,}\NormalTok{ x}\OperatorTok{);}   \CommentTok{// Mencetak 42 desimal}
\end{Highlighting}
\end{Shaded}

\end{tcolorbox}

Tidak ada format specifier \hl{\mbox{\texttt{printf()}}} untuk mencetak bilangan biner. Kamu harus melakukannya satu karakter demi satu karakter dengan operator bitwise.

\index{0b binary@\texttt{0b} binary|)}

\subsection{Konstanta Integer}\label{konstanta-integer}

\index{Integer constants|(}

Kamu bisa memaksa konstanta integer menjadi tipe tertentu dengan menambahkan sufiks yang menunjukkan tipenya.

Kita akan melakukan beberapa penugasan sebagai demo, tapi paling sering developer tidak menulis sufiks kecuali perlu presisi. Compiler cukup pandai memastikan tipe-tipenya kompatibel.

\index{U unsigned constant@\texttt{U} unsigned constant|(} \index{L long constant@\texttt{L} long constant|(} \index{LL long long constant@\texttt{LL} long long constant|(} \index{UL unsigned long constant@\texttt{UL} unsigned long constant|(} \index{ULL unsigned long long constant@\texttt{ULL} unsigned long long constant|(}

\begin{tcolorbox}[enhanced jigsaw,left=3pt,breakable,colback=white!97.2549!black,boxrule=0.2mm]

\begin{Shaded}
\begin{Highlighting}[]
\DataTypeTok{int}\NormalTok{           x }\OperatorTok{=} \DecValTok{1234}\OperatorTok{;}
\DataTypeTok{long} \DataTypeTok{int}\NormalTok{      x }\OperatorTok{=} \DecValTok{1234}\BuiltInTok{L}\OperatorTok{;}
\DataTypeTok{long} \DataTypeTok{long} \DataTypeTok{int}\NormalTok{ x }\OperatorTok{=} \DecValTok{1234}\BuiltInTok{LL}

\DataTypeTok{unsigned} \DataTypeTok{int}\NormalTok{           x }\OperatorTok{=} \DecValTok{1234}\BuiltInTok{U}\OperatorTok{;}
\DataTypeTok{unsigned} \DataTypeTok{long} \DataTypeTok{int}\NormalTok{      x }\OperatorTok{=} \DecValTok{1234}\BuiltInTok{UL}\OperatorTok{;}
\DataTypeTok{unsigned} \DataTypeTok{long} \DataTypeTok{long} \DataTypeTok{int}\NormalTok{ x }\OperatorTok{=} \DecValTok{1234}\BuiltInTok{ULL}\OperatorTok{;}
\end{Highlighting}
\end{Shaded}

\end{tcolorbox}

Sufiks boleh huruf besar maupun kecil. Dan \hl{\mbox{\texttt{U}}} serta \hl{\mbox{\texttt{L}}} atau \hl{\mbox{\texttt{LL}}} boleh muncul dalam urutan mana pun.

{\def\LTcaptype{none} % do not increment counter
\begin{longtable}[]{@{}ll@{}}
\toprule\noalign{}
Tipe & Sufiks \\
\midrule\noalign{}
\endhead
\bottomrule\noalign{}
\endlastfoot
\hl{\mbox{\texttt{int}}} & Tidak ada \\
\hl{\mbox{\texttt{long\ int}}} & \hl{\mbox{\texttt{L}}} \\
\hl{\mbox{\texttt{long\ long\ int}}} & \hl{\mbox{\texttt{LL}}} \\
\hl{\mbox{\texttt{unsigned\ int}}} & \hl{\mbox{\texttt{U}}} \\
\hl{\mbox{\texttt{unsigned\ long\ int}}} & \hl{\mbox{\texttt{UL}}} \\
\hl{\mbox{\texttt{unsigned\ long\ long\ int}}} & \hl{\mbox{\texttt{ULL}}} \\
\end{longtable}
}

Tadi di tabel kusebut bahwa ``tanpa sufiks'' berarti \hl{\mbox{\texttt{int}}}\ldots{} tapi sebenarnya lebih rumit dari itu.

Jadi apa yang terjadi kalau ada angka tanpa sufiks seperti:

\begin{tcolorbox}[enhanced jigsaw,left=3pt,breakable,colback=white!97.2549!black,boxrule=0.2mm]

\begin{Shaded}
\begin{Highlighting}[]
\DataTypeTok{int}\NormalTok{ x }\OperatorTok{=} \DecValTok{1234}\OperatorTok{;}
\end{Highlighting}
\end{Shaded}

\end{tcolorbox}

Bertipe apa angka itu?

Umumnya C memilih tipe terkecil, mulai dari \hl{\mbox{\texttt{int}}}, yang bisa menampung nilainya.

Tapi lebih spesifik lagi, itu juga bergantung pada basis angkanya (desimal, hex, atau oktal).

Spesifikasi punya tabel bagus yang menunjukkan tipe mana yang dipakai untuk nilai tanpa sufiks. Bahkan, tabel itu akan kusalin utuh di sini.

C11 §6.4.4.1¶5 berbunyi, ``The type of an integer constant is the first of the first of the corresponding list in which its value can be represented.''

Lalu dilanjutkan dengan tabel ini:

{\def\LTcaptype{none} % do not increment counter
\begin{longtable}[]{@{}
  >{\raggedright\arraybackslash}p{(\linewidth - 4\tabcolsep) * \real{0.2361}}
  >{\raggedright\arraybackslash}p{(\linewidth - 4\tabcolsep) * \real{0.3472}}
  >{\raggedright\arraybackslash}p{(\linewidth - 4\tabcolsep) * \real{0.3611}}@{}}
\toprule\noalign{}
\begin{minipage}[b]{\linewidth}\raggedright
Sufiks
\end{minipage} & \begin{minipage}[b]{\linewidth}\raggedright
Konstanta Desimal
\end{minipage} & \begin{minipage}[b]{\linewidth}\raggedright
Konstanta Oktal atau\\
Heksadesimal\strut
\end{minipage} \\
\midrule\noalign{}
\endhead
\bottomrule\noalign{}
\endlastfoot
tidak ada & \begin{minipage}[t]{\linewidth}\raggedright
\hl{\mbox{\texttt{int}}}\\
\hl{\mbox{\texttt{long\ int}}}\strut
\end{minipage} & \begin{minipage}[t]{\linewidth}\raggedright
\hl{\mbox{\texttt{int}}}\\
\hl{\mbox{\texttt{unsigned\ int}}}\\
\hl{\mbox{\texttt{long\ int}}}\\
\hl{\mbox{\texttt{unsigned\ long\ int}}}\\
\hl{\mbox{\texttt{long\ long\ int}}}\\
\hl{\mbox{\texttt{unsigned\ long\ long\ int}}}\\
\strut
\end{minipage} \\
\hl{\mbox{\texttt{u}}} atau \hl{\mbox{\texttt{U}}} & \begin{minipage}[t]{\linewidth}\raggedright
\hl{\mbox{\texttt{unsigned\ int}}}\\
\hl{\mbox{\texttt{unsigned\ long\ int}}}\\
\hl{\mbox{\texttt{unsigned\ long\ long\ int}}}\strut
\end{minipage} & \begin{minipage}[t]{\linewidth}\raggedright
\hl{\mbox{\texttt{unsigned\ int}}}\\
\hl{\mbox{\texttt{unsigned\ long\ int}}}\\
\hl{\mbox{\texttt{unsigned\ long\ long\ int}}}\\
\strut
\end{minipage} \\
\hl{\mbox{\texttt{l}}} atau \hl{\mbox{\texttt{L}}} & \begin{minipage}[t]{\linewidth}\raggedright
\hl{\mbox{\texttt{long\ int}}}\\
\hl{\mbox{\texttt{long\ long\ int}}}\strut
\end{minipage} & \begin{minipage}[t]{\linewidth}\raggedright
\hl{\mbox{\texttt{long\ int}}}\\
\hl{\mbox{\texttt{unsigned\ long\ int}}}\\
\hl{\mbox{\texttt{long\ long\ int}}}\\
\hl{\mbox{\texttt{unsigned\ long\ long\ int}}}\\
\strut
\end{minipage} \\
\begin{minipage}[t]{\linewidth}\raggedright
\hl{\mbox{\texttt{u}}}/\hl{\mbox{\texttt{U}}} dan\\
\hl{\mbox{\texttt{l}}}/\hl{\mbox{\texttt{L}}}\strut
\end{minipage} & \begin{minipage}[t]{\linewidth}\raggedright
\hl{\mbox{\texttt{unsigned\ long\ int}}}\\
\hl{\mbox{\texttt{unsigned\ long\ long\ int}}}\strut
\end{minipage} & \begin{minipage}[t]{\linewidth}\raggedright
\hl{\mbox{\texttt{unsigned\ long\ int}}}\\
\hl{\mbox{\texttt{unsigned\ long\ long\ int}}}\\
\strut
\end{minipage} \\
\hl{\mbox{\texttt{ll}}} atau \hl{\mbox{\texttt{LL}}} & \hl{\mbox{\texttt{long\ long\ int}}} & \begin{minipage}[t]{\linewidth}\raggedright
\hl{\mbox{\texttt{long\ long\ int}}}\\
\hl{\mbox{\texttt{unsigned\ long\ long\ int}}}\\
\strut
\end{minipage} \\
\begin{minipage}[t]{\linewidth}\raggedright
\hl{\mbox{\texttt{u}}}/\hl{\mbox{\texttt{U}}} dan\\
\hl{\mbox{\texttt{ll}}}/\hl{\mbox{\texttt{LL}}}\strut
\end{minipage} & \hl{\mbox{\texttt{unsigned\ long\ long\ int}}} & \hl{\mbox{\texttt{unsigned\ long\ long\ int}}} \\
\end{longtable}
}

\index{L long constant@\texttt{L} long constant|)} \index{LL long long constant@\texttt{LL} long long constant|)} \index{UL unsigned long constant@\texttt{UL} unsigned long constant|)} \index{ULL unsigned long long constant@\texttt{ULL} unsigned long long constant|)}

Artinya, misalnya kalau kamu menulis angka seperti \hl{\mbox{\texttt{123456789U}}}, mula-mula C akan melihat apakah ia bisa menjadi \hl{\mbox{\texttt{unsigned\ int}}}. Kalau tidak muat, dicoba \hl{\mbox{\texttt{unsigned\ long\ int}}}. Lalu \hl{\mbox{\texttt{unsigned\ long\ long\ int}}}. Ia memakai tipe terkecil yang bisa menampung angka itu.

\index{U unsigned constant@\texttt{U} unsigned constant|)}

\index{Integer constants|)}

\subsection{Konstanta Floating Point}\label{konstanta-floating-point}

\index{Floating point constants|(}

Kamu mungkin mengira konstanta floating point seperti \hl{\mbox{\texttt{1.23}}} bertipe bawaan \hl{\mbox{\texttt{float}}}, kan?

Kejutan! Ternyata bilangan floating point tanpa sufiks bertipe \hl{\mbox{\texttt{double}}}! Selamat ulang tahun yang terlambat!

\index{F float constant@\texttt{F} float constant|(} \index{L long double constant@\texttt{L} long double constant|(}

Kamu bisa memaksanya menjadi \hl{\mbox{\texttt{float}}} dengan menambahkan \hl{\mbox{\texttt{f}}} (atau \hl{\mbox{\texttt{F}}}---tidak membedakan huruf besar/kecil). Kamu bisa memaksanya menjadi \hl{\mbox{\texttt{long\ double}}} dengan menambahkan \hl{\mbox{\texttt{l}}} (atau \hl{\mbox{\texttt{L}}}).

{\def\LTcaptype{none} % do not increment counter
\begin{longtable}[]{@{}ll@{}}
\toprule\noalign{}
Tipe & Sufiks \\
\midrule\noalign{}
\endhead
\bottomrule\noalign{}
\endlastfoot
\hl{\mbox{\texttt{float}}} & \hl{\mbox{\texttt{F}}} \\
\hl{\mbox{\texttt{double}}} & Tidak ada \\
\hl{\mbox{\texttt{long\ double}}} & \hl{\mbox{\texttt{L}}} \\
\end{longtable}
}

Misalnya:

\begin{tcolorbox}[enhanced jigsaw,left=3pt,breakable,colback=white!97.2549!black,boxrule=0.2mm]

\begin{Shaded}
\begin{Highlighting}[]
\DataTypeTok{float}\NormalTok{ x       }\OperatorTok{=} \FloatTok{3.14}\BuiltInTok{f}\OperatorTok{;}
\DataTypeTok{double}\NormalTok{ x      }\OperatorTok{=} \FloatTok{3.14}\OperatorTok{;}
\DataTypeTok{long} \DataTypeTok{double}\NormalTok{ x }\OperatorTok{=} \FloatTok{3.14}\BuiltInTok{L}\OperatorTok{;}
\end{Highlighting}
\end{Shaded}

\end{tcolorbox}

\index{F float constant@\texttt{F} float constant|)} \index{L long double constant@\texttt{L} long double constant|)}

Tapi selama ini kita melakukan ini saja, kan?

\begin{tcolorbox}[enhanced jigsaw,left=3pt,breakable,colback=white!97.2549!black,boxrule=0.2mm]

\begin{Shaded}
\begin{Highlighting}[]
\DataTypeTok{float}\NormalTok{ x }\OperatorTok{=} \FloatTok{3.14}\OperatorTok{;}
\end{Highlighting}
\end{Shaded}

\end{tcolorbox}

Bukankah sisi kiri \hl{\mbox{\texttt{float}}} dan sisi kanan \hl{\mbox{\texttt{double}}}? Betul! Tapi C cukup pandai soal konversi numerik otomatis, jadi lebih umum menulis konstanta floating point tanpa sufiks daripada dengan sufiks. Nanti dibahas lagi.

\index{Floating point constants|)}

\subsubsection{Notasi Ilmiah}\label{notasi-ilmiah}

\index{Scientific notation|(}

Ingat tadi saat kita membahas bahwa bilangan floating point bisa direpresentasikan dengan significand, basis, dan eksponen?

Nah, ada cara umum menulis bilangan seperti itu, ditampilkan di sini diikuti padanannya yang lebih mudah dikenali, yaitu hasil perhitungan sebenarnya:

\(1.2345\times10^3 = 1234.5\)

Menulis angka dalam bentuk \(s\times b^e\) disebut \href{https://en.wikipedia.org/wiki/Scientific_notation}{\emph{notasi ilmiah}}\footnote{https://en.wikipedia.org/wiki/Scientific\_notation}. Di C, ini ditulis dengan ``notasi E'', jadi dua bentuk ini setara:

{\def\LTcaptype{none} % do not increment counter
\begin{longtable}[]{@{}ll@{}}
\toprule\noalign{}
Notasi Ilmiah & Notasi E \\
\midrule\noalign{}
\endhead
\bottomrule\noalign{}
\endlastfoot
\(1.2345\times10^{-3}=0.0012345\) & \hl{\mbox{\texttt{1.2345e-3}}} \\
\(1.2345\times10^8=123450000\) & \hl{\mbox{\texttt{1.2345e+8}}} \\
\end{longtable}
}

Kamu bisa mencetak angka dalam notasi ini dengan \hl{\mbox{\texttt{\%e}}}:

\begin{tcolorbox}[enhanced jigsaw,left=3pt,breakable,colback=white!97.2549!black,boxrule=0.2mm]

\begin{Shaded}
\begin{Highlighting}[]
\NormalTok{printf}\OperatorTok{(}\StringTok{"}\SpecialCharTok{\%e\textbackslash{}n}\StringTok{"}\OperatorTok{,} \FloatTok{123456.0}\OperatorTok{);}  \CommentTok{// Mencetak 1.234560e+05}
\end{Highlighting}
\end{Shaded}

\end{tcolorbox}

Beberapa fakta kecil yang menyenangkan tentang notasi ilmiah:

\begin{itemize}
\item
  Angka tidak harus ditulis dengan satu digit di depan titik desimal. Berapa pun jumlah angka boleh ditaruh di depan.

  \begin{tcolorbox}[enhanced jigsaw,left=3pt,breakable,colback=white!97.2549!black,boxrule=0.2mm]

\begin{Shaded}
\begin{Highlighting}[]
\DataTypeTok{double}\NormalTok{ x }\OperatorTok{=} \FloatTok{123.456e+3}\OperatorTok{;}  \CommentTok{// 123456}
\end{Highlighting}
\end{Shaded}

  \end{tcolorbox}

  Tapi saat dicetak, eksponennya akan diubah sehingga hanya ada satu digit di depan titik desimal.
\item
  Tanda plus pada eksponen boleh dihilangkan karena itu bawaannya, tapi dari yang pernah kulihat, ini jarang dipakai.

  \begin{tcolorbox}[enhanced jigsaw,left=3pt,breakable,colback=white!97.2549!black,boxrule=0.2mm]

\begin{Shaded}
\begin{Highlighting}[]
\FloatTok{1.2345e10} \OperatorTok{==} \FloatTok{1.2345e+10}
\end{Highlighting}
\end{Shaded}

  \end{tcolorbox}
\item
  Kamu bisa menambahkan sufiks \hl{\mbox{\texttt{F}}} atau \hl{\mbox{\texttt{L}}} pada konstanta notasi E:

  \begin{tcolorbox}[enhanced jigsaw,left=3pt,breakable,colback=white!97.2549!black,boxrule=0.2mm]

\begin{Shaded}
\begin{Highlighting}[]
\FloatTok{1.2345e10}\BuiltInTok{F}
\FloatTok{1.2345e10}\BuiltInTok{L}
\end{Highlighting}
\end{Shaded}

  \end{tcolorbox}
\end{itemize}

\index{Scientific notation|)}

\subsubsection{Konstanta Floating Point Heksadesimal}\label{konstanta-floating-point-heksadesimal}

\index{Hex floating point constants|(}

Tapi tunggu, masih ada lagi yang mengambang!

Ternyata ada juga konstanta floating point heksadesimal!

Cara kerjanya mirip bilangan floating point desimal, tapi diawali \hl{\mbox{\texttt{0x}}} seperti angka integer.

Syaratnya, kamu \emph{harus} menyebutkan eksponen, dan eksponen ini menghasilkan pangkat 2. Yaitu: \(2^x\).

Lalu kamu memakai \hl{\mbox{\texttt{p}}} alih-alih \hl{\mbox{\texttt{e}}} saat menulis angkanya:

Jadi \hl{\mbox{\texttt{0xa.1p3}}} adalah \(10.0625\times2^3 == 80.5\).

Saat memakai konstanta hex floating point, Kita bisa mencetak notasi ilmiah hex dengan \hl{\mbox{\texttt{\%a}}}:

\begin{tcolorbox}[enhanced jigsaw,left=3pt,breakable,colback=white!97.2549!black,boxrule=0.2mm]

\begin{Shaded}
\begin{Highlighting}[]
\DataTypeTok{double}\NormalTok{ x }\OperatorTok{=} \FloatTok{0xa.1p3}\OperatorTok{;}

\NormalTok{printf}\OperatorTok{(}\StringTok{"}\SpecialCharTok{\%a\textbackslash{}n}\StringTok{"}\OperatorTok{,}\NormalTok{ x}\OperatorTok{);}  \CommentTok{// 0x1.42p+6}
\NormalTok{printf}\OperatorTok{(}\StringTok{"}\SpecialCharTok{\%f\textbackslash{}n}\StringTok{"}\OperatorTok{,}\NormalTok{ x}\OperatorTok{);}  \CommentTok{// 80.500000}
\end{Highlighting}
\end{Shaded}

\end{tcolorbox}

\index{Hex floating point constants|)}

\chapter{Tipe III: Konversi}\label{tipe-iii-konversi}

\index{Type conversions|(}

Di bab ini kita akan membahas konversi dari satu tipe ke tipe lain. C punya berbagai cara untuk melakukannya, dan sebagian mungkin sedikit berbeda dari yang biasa kamu temui di bahasa lain.

Sebelum membahas cara membuat konversi terjadi, mari kita lihat dulu bagaimana konversi bekerja saat ia \emph{memang} terjadi.

\section{Konversi String}\label{konversi-string}

\index{Type conversions!strings|(}

Tidak seperti banyak bahasa, C tidak melakukan konversi string ke angka (dan sebaliknya) semulus konversi antartipe numerik.

Untuk itu, kita harus memanggil fungsi yang mengerjakan pekerjaan kotornya.

\subsection{Nilai Numerik ke String}\label{nilai-numerik-ke-string}

Ketika ingin mengonversi angka menjadi string, kita bisa memakai \hl{\mbox{\texttt{sprintf()}}} (dibaca \emph{SPRINT-f}) atau \hl{\mbox{\texttt{snprintf()}}} (\emph{s-n-print-f})\footnote{Keduanya sama, kecuali \hl{\mbox{\texttt{snprintf()}}} memungkinkan kita menentukan jumlah byte maksimum yang ditulis, sehingga mencegah kelebihan menulis melewati ujung string. Jadi lebih aman.}

Keduanya pada dasarnya bekerja seperti \hl{\mbox{\texttt{printf()}}}, hanya saja outputnya ke string, dan string itu bisa dicetak nanti, atau diapakan saja.

Misalnya, mengubah sebagian nilai π menjadi string:

\begin{tcolorbox}[enhanced jigsaw,left=3pt,breakable,colback=white!97.2549!black,boxrule=0.2mm]

\begin{Shaded}
\begin{Highlighting}[numbers=left,,]
\PreprocessorTok{\#include }\ImportTok{\textless{}stdio.h\textgreater{}}

\DataTypeTok{int}\NormalTok{ main}\OperatorTok{(}\DataTypeTok{void}\OperatorTok{)}
\OperatorTok{\{}
    \DataTypeTok{char}\NormalTok{ s}\OperatorTok{[}\DecValTok{10}\OperatorTok{];}
    \DataTypeTok{float}\NormalTok{ f }\OperatorTok{=} \FloatTok{3.14159}\OperatorTok{;}

    \CommentTok{// Konversi "f" menjadi string, simpan di "s", tulis paling banyak 10}
    \CommentTok{// karakter termasuk terminator NUL}

\NormalTok{    snprintf}\OperatorTok{(}\NormalTok{s}\OperatorTok{,} \DecValTok{10}\OperatorTok{,} \StringTok{"}\SpecialCharTok{\%f}\StringTok{"}\OperatorTok{,}\NormalTok{ f}\OperatorTok{);}

\NormalTok{    printf}\OperatorTok{(}\StringTok{"String value: }\SpecialCharTok{\%s\textbackslash{}n}\StringTok{"}\OperatorTok{,}\NormalTok{ s}\OperatorTok{);}  \CommentTok{// String value: 3.141590}
\OperatorTok{\}}
\end{Highlighting}
\end{Shaded}

\end{tcolorbox}

Jadi kamu bisa memakai \hl{\mbox{\texttt{\%d}}} atau \hl{\mbox{\texttt{\%u}}} seperti biasa untuk integer.

\subsection{String ke Nilai Numerik}\label{string-ke-nilai-numerik}

Ada beberapa keluarga fungsi untuk ini di C. Kita sebut keluarga \hl{\mbox{\texttt{atoi}}} (dibaca \emph{a-to-i}) dan keluarga \hl{\mbox{\texttt{strtol}}} (\emph{stir-to-long}).

Untuk konversi dasar dari string ke angka, coba fungsi \hl{\mbox{\texttt{atoi}}} dari \hl{\mbox{\texttt{\textless{}stdlib.h\textgreater{}}}}. Fungsi-fungsi ini punya penanganan error yang buruk (termasuk undefined behavior kalau kamu mengoper string yang salah), jadi pakailah dengan hati-hati.

{\def\LTcaptype{none} % do not increment counter
\begin{longtable}[]{@{}ll@{}}
\toprule\noalign{}
Fungsi & Deskripsi \\
\midrule\noalign{}
\endhead
\bottomrule\noalign{}
\endlastfoot
\hl{\mbox{\texttt{atoi}}} & String ke \hl{\mbox{\texttt{int}}} \\
\hl{\mbox{\texttt{atof}}} & String ke \hl{\mbox{\texttt{float}}} \\
\hl{\mbox{\texttt{atol}}} & String ke \hl{\mbox{\texttt{long\ int}}} \\
\hl{\mbox{\texttt{atoll}}} & String ke \hl{\mbox{\texttt{long\ long\ int}}} \\
\end{longtable}
}

Walaupun spesifikasi tidak mengakuinya, huruf \hl{\mbox{\texttt{a}}} di awal nama fungsi singkatan dari \href{https://en.wikipedia.org/wiki/ASCII}{ASCII}\footnote{https://en.wikipedia.org/wiki/ASCII}, jadi sebenarnya \hl{\mbox{\texttt{atoi()}}} adalah ``ASCII-to-integer'', tapi mengatakannya hari ini agak terlalu berpusat pada ASCII.

Berikut contoh mengonversi string menjadi \hl{\mbox{\texttt{float}}}:

\begin{tcolorbox}[enhanced jigsaw,left=3pt,breakable,colback=white!97.2549!black,boxrule=0.2mm]

\begin{Shaded}
\begin{Highlighting}[numbers=left,,]
\PreprocessorTok{\#include }\ImportTok{\textless{}stdio.h\textgreater{}}
\PreprocessorTok{\#include }\ImportTok{\textless{}stdlib.h\textgreater{}}

\DataTypeTok{int}\NormalTok{ main}\OperatorTok{(}\DataTypeTok{void}\OperatorTok{)}
\OperatorTok{\{}
    \DataTypeTok{char} \OperatorTok{*}\NormalTok{pi }\OperatorTok{=} \StringTok{"3.14159"}\OperatorTok{;}
    \DataTypeTok{float}\NormalTok{ f}\OperatorTok{;}

\NormalTok{    f }\OperatorTok{=}\NormalTok{ atof}\OperatorTok{(}\NormalTok{pi}\OperatorTok{);}

\NormalTok{    printf}\OperatorTok{(}\StringTok{"}\SpecialCharTok{\%f\textbackslash{}n}\StringTok{"}\OperatorTok{,}\NormalTok{ f}\OperatorTok{);}
\OperatorTok{\}}
\end{Highlighting}
\end{Shaded}

\end{tcolorbox}

Tapi, seperti kubilang tadi, hal-hal aneh seperti ini menghasilkan undefined behavior:

\begin{tcolorbox}[enhanced jigsaw,left=3pt,breakable,colback=white!97.2549!black,boxrule=0.2mm]

\begin{Shaded}
\begin{Highlighting}[]
\DataTypeTok{int}\NormalTok{ x }\OperatorTok{=}\NormalTok{ atoi}\OperatorTok{(}\StringTok{"what"}\OperatorTok{);}  \CommentTok{// "What" bukan angka yang pernah dikenal siapa pun}
\end{Highlighting}
\end{Shaded}

\end{tcolorbox}

(Waktu kujalankan, hasilnya \hl{\mbox{\texttt{0}}}, tapi jangan sekali-kali mengandalkannya. Hasilnya bisa sama sekali berbeda.)

Untuk penanganan error yang lebih baik, mari kita lihat fungsi-fungsi \hl{\mbox{\texttt{strtol}}}, yang juga ada di \hl{\mbox{\texttt{\textless{}stdlib.h\textgreater{}}}}. Bukan hanya itu, mereka juga mengonversi ke lebih banyak tipe dan lebih banyak basis!

{\def\LTcaptype{none} % do not increment counter
\begin{longtable}[]{@{}ll@{}}
\toprule\noalign{}
Fungsi & Deskripsi \\
\midrule\noalign{}
\endhead
\bottomrule\noalign{}
\endlastfoot
\hl{\mbox{\texttt{strtol}}} & String ke \hl{\mbox{\texttt{long\ int}}} \\
\hl{\mbox{\texttt{strtoll}}} & String ke \hl{\mbox{\texttt{long\ long\ int}}} \\
\hl{\mbox{\texttt{strtoul}}} & String ke \hl{\mbox{\texttt{unsigned\ long\ int}}} \\
\hl{\mbox{\texttt{strtoull}}} & String ke \hl{\mbox{\texttt{unsigned\ long\ long\ int}}} \\
\hl{\mbox{\texttt{strtof}}} & String ke \hl{\mbox{\texttt{float}}} \\
\hl{\mbox{\texttt{strtod}}} & String ke \hl{\mbox{\texttt{double}}} \\
\hl{\mbox{\texttt{strtold}}} & String ke \hl{\mbox{\texttt{long\ double}}} \\
\end{longtable}
}

Fungsi-fungsi ini semuanya punya pola pemakaian yang mirip, dan bagi banyak orang ini pengalaman pertama dengan pointer ke pointer! Tapi tenang saja---tidak sesulit kelihatannya.

Mari kita coba contoh ketika string dikonversi menjadi \hl{\mbox{\texttt{unsigned\ long}}}, dengan membuang informasi error (yaitu informasi tentang karakter yang salah di string input):

\begin{tcolorbox}[enhanced jigsaw,left=3pt,breakable,colback=white!97.2549!black,boxrule=0.2mm]

\begin{Shaded}
\begin{Highlighting}[numbers=left,,]
\PreprocessorTok{\#include }\ImportTok{\textless{}stdio.h\textgreater{}}
\PreprocessorTok{\#include }\ImportTok{\textless{}stdlib.h\textgreater{}}

\DataTypeTok{int}\NormalTok{ main}\OperatorTok{(}\DataTypeTok{void}\OperatorTok{)}
\OperatorTok{\{}
    \DataTypeTok{char} \OperatorTok{*}\NormalTok{s }\OperatorTok{=} \StringTok{"3490"}\OperatorTok{;}

    \CommentTok{// Konversi string s, angka basis 10, menjadi unsigned long int.}
    \CommentTok{// NULL berarti kita tidak peduli dengan informasi error apa pun.}

    \DataTypeTok{unsigned} \DataTypeTok{long} \DataTypeTok{int}\NormalTok{ x }\OperatorTok{=}\NormalTok{ strtoul}\OperatorTok{(}\NormalTok{s}\OperatorTok{,}\NormalTok{ NULL}\OperatorTok{,} \DecValTok{10}\OperatorTok{);}

\NormalTok{    printf}\OperatorTok{(}\StringTok{"}\SpecialCharTok{\%lu\textbackslash{}n}\StringTok{"}\OperatorTok{,}\NormalTok{ x}\OperatorTok{);}  \CommentTok{// 3490}
\OperatorTok{\}}
\end{Highlighting}
\end{Shaded}

\end{tcolorbox}

Perhatikan beberapa hal di sana. Meskipun kita tidak sudi menangkap informasi apa pun tentang karakter yang salah di string, \hl{\mbox{\texttt{strtoul()}}} tidak akan menghasilkan undefined behavior; ia hanya mengembalikan \hl{\mbox{\texttt{0}}}.

Perhatikan juga bahwa kita menyebutkan bahwa ini angka desimal (basis 10).

Apakah ini berarti kita bisa mengonversi angka dengan basis berbeda? Tentu! Kita coba biner!

\begin{tcolorbox}[enhanced jigsaw,left=3pt,breakable,colback=white!97.2549!black,boxrule=0.2mm]

\begin{Shaded}
\begin{Highlighting}[numbers=left,,]
\PreprocessorTok{\#include }\ImportTok{\textless{}stdio.h\textgreater{}}
\PreprocessorTok{\#include }\ImportTok{\textless{}stdlib.h\textgreater{}}

\DataTypeTok{int}\NormalTok{ main}\OperatorTok{(}\DataTypeTok{void}\OperatorTok{)}
\OperatorTok{\{}
    \DataTypeTok{char} \OperatorTok{*}\NormalTok{s }\OperatorTok{=} \StringTok{"101010"}\OperatorTok{;}  \CommentTok{// Apa makna angka ini?}

    \CommentTok{// Konversi string s, angka basis 2, menjadi unsigned long int.}

    \DataTypeTok{unsigned} \DataTypeTok{long} \DataTypeTok{int}\NormalTok{ x }\OperatorTok{=}\NormalTok{ strtoul}\OperatorTok{(}\NormalTok{s}\OperatorTok{,}\NormalTok{ NULL}\OperatorTok{,} \DecValTok{2}\OperatorTok{);}

\NormalTok{    printf}\OperatorTok{(}\StringTok{"}\SpecialCharTok{\%lu\textbackslash{}n}\StringTok{"}\OperatorTok{,}\NormalTok{ x}\OperatorTok{);}  \CommentTok{// 42}
\OperatorTok{\}}
\end{Highlighting}
\end{Shaded}

\end{tcolorbox}

Oke, itu semua seru-seruan, tapi \hl{\mbox{\texttt{NULL}}} di situ untuk apa?

Itu membantu kita mengetahui apakah terjadi error saat memproses string. Ia adalah pointer ke pointer ke \hl{\mbox{\texttt{char}}}, yang terdengar menakutkan, tapi tidak lagi begitu setelah kamu paham.

Mari kita coba contoh dengan sengaja memberi angka yang salah, dan kita lihat bagaimana \hl{\mbox{\texttt{strtol()}}} memberi tahu di mana digit tidak valid yang pertama.

\begin{tcolorbox}[enhanced jigsaw,left=3pt,breakable,colback=white!97.2549!black,boxrule=0.2mm]

\begin{Shaded}
\begin{Highlighting}[numbers=left,,]
\PreprocessorTok{\#include }\ImportTok{\textless{}stdio.h\textgreater{}}
\PreprocessorTok{\#include }\ImportTok{\textless{}stdlib.h\textgreater{}}

\DataTypeTok{int}\NormalTok{ main}\OperatorTok{(}\DataTypeTok{void}\OperatorTok{)}
\OperatorTok{\{}
    \DataTypeTok{char} \OperatorTok{*}\NormalTok{s }\OperatorTok{=} \StringTok{"34x90"}\OperatorTok{;}  \CommentTok{// "x" bukan digit yang valid di basis 10!}
    \DataTypeTok{char} \OperatorTok{*}\NormalTok{badchar}\OperatorTok{;}

    \CommentTok{// Konversi string s, angka basis 10, menjadi unsigned long int.}

    \DataTypeTok{unsigned} \DataTypeTok{long} \DataTypeTok{int}\NormalTok{ x }\OperatorTok{=}\NormalTok{ strtoul}\OperatorTok{(}\NormalTok{s}\OperatorTok{,} \OperatorTok{\&}\NormalTok{badchar}\OperatorTok{,} \DecValTok{10}\OperatorTok{);}

    \CommentTok{// Ia mencoba mengonversi sebanyak mungkin, jadi sampai sini:}

\NormalTok{    printf}\OperatorTok{(}\StringTok{"}\SpecialCharTok{\%lu\textbackslash{}n}\StringTok{"}\OperatorTok{,}\NormalTok{ x}\OperatorTok{);}  \CommentTok{// 34}

    \CommentTok{// Tapi kita bisa melihat karakter bermasalah itu karena badchar}
    \CommentTok{// menunjuk ke sana!}

\NormalTok{    printf}\OperatorTok{(}\StringTok{"Invalid character: }\SpecialCharTok{\%c\textbackslash{}n}\StringTok{"}\OperatorTok{,} \OperatorTok{*}\NormalTok{badchar}\OperatorTok{);}  \CommentTok{// "x"}
\OperatorTok{\}}
\end{Highlighting}
\end{Shaded}

\end{tcolorbox}

Jadi di sana \hl{\mbox{\texttt{strtoul()}}} mengubah apa yang ditunjuk \hl{\mbox{\texttt{badchar}}} untuk menunjukkan di mana yang salah\footnote{Kita harus mengoper pointer ke \hl{\mbox{\texttt{badchar}}} ke \hl{\mbox{\texttt{strtoul()}}}, kalau tidak ia tidak akan bisa mengubahnya dengan cara yang terlihat oleh kita, sama seperti alasan kenapa kamu harus mengoper pointer ke \hl{\mbox{\texttt{int}}} ke fungsi kalau ingin fungsi itu bisa mengubah nilai \hl{\mbox{\texttt{int}}} itu.}.

Tapi bagaimana kalau tidak ada yang salah? Dalam kasus itu, \hl{\mbox{\texttt{badchar}}} akan menunjuk ke terminator \hl{\mbox{\texttt{NUL}}} di akhir string. Jadi kita bisa mengujinya:

\begin{tcolorbox}[enhanced jigsaw,left=3pt,breakable,colback=white!97.2549!black,boxrule=0.2mm]

\begin{Shaded}
\begin{Highlighting}[numbers=left,,]
\PreprocessorTok{\#include }\ImportTok{\textless{}stdio.h\textgreater{}}
\PreprocessorTok{\#include }\ImportTok{\textless{}stdlib.h\textgreater{}}

\DataTypeTok{int}\NormalTok{ main}\OperatorTok{(}\DataTypeTok{void}\OperatorTok{)}
\OperatorTok{\{}
    \DataTypeTok{char} \OperatorTok{*}\NormalTok{s }\OperatorTok{=} \StringTok{"3490"}\OperatorTok{;}  \CommentTok{// "x" bukan digit yang valid di basis 10!}
    \DataTypeTok{char} \OperatorTok{*}\NormalTok{badchar}\OperatorTok{;}

    \CommentTok{// Konversi string s, angka basis 10, menjadi unsigned long int.}

    \DataTypeTok{unsigned} \DataTypeTok{long} \DataTypeTok{int}\NormalTok{ x }\OperatorTok{=}\NormalTok{ strtoul}\OperatorTok{(}\NormalTok{s}\OperatorTok{,} \OperatorTok{\&}\NormalTok{badchar}\OperatorTok{,} \DecValTok{10}\OperatorTok{);}

    \CommentTok{// Periksa apakah semuanya berjalan baik}

    \ControlFlowTok{if} \OperatorTok{(*}\NormalTok{badchar }\OperatorTok{==} \CharTok{\textquotesingle{}}\SpecialCharTok{\textbackslash{}0}\CharTok{\textquotesingle{}}\OperatorTok{)} \OperatorTok{\{}
\NormalTok{        printf}\OperatorTok{(}\StringTok{"Success! }\SpecialCharTok{\%lu\textbackslash{}n}\StringTok{"}\OperatorTok{,}\NormalTok{ x}\OperatorTok{);}
    \OperatorTok{\}} \ControlFlowTok{else}  \OperatorTok{\{}
\NormalTok{        printf}\OperatorTok{(}\StringTok{"Partial conversion: }\SpecialCharTok{\%lu\textbackslash{}n}\StringTok{"}\OperatorTok{,}\NormalTok{ x}\OperatorTok{);}
\NormalTok{        printf}\OperatorTok{(}\StringTok{"Invalid character: }\SpecialCharTok{\%c\textbackslash{}n}\StringTok{"}\OperatorTok{,} \OperatorTok{*}\NormalTok{badchar}\OperatorTok{);}
    \OperatorTok{\}}
\OperatorTok{\}}
\end{Highlighting}
\end{Shaded}

\end{tcolorbox}

Jadi begitulah. Fungsi gaya \hl{\mbox{\texttt{atoi()}}} cukup baik untuk keadaan terkendali, tapi fungsi gaya \hl{\mbox{\texttt{strtol()}}} memberimu kendali jauh lebih besar atas penanganan error dan basis input.

\index{Type conversions!strings|)}

\section{\texorpdfstring{Konversi \hl{\mbox{\texttt{char}}}}{Konversi char}}\label{konversi-char}

\index{Type conversions!char@\texttt{char}|(}

Bagaimana kalau ada satu karakter berisi digit, seperti \hl{\mbox{\texttt{\textquotesingle{}5\textquotesingle{}}}}\ldots{} Apakah itu sama dengan nilai \hl{\mbox{\texttt{5}}}?

Kita coba dan lihat.

\begin{tcolorbox}[enhanced jigsaw,left=3pt,breakable,colback=white!97.2549!black,boxrule=0.2mm]

\begin{Shaded}
\begin{Highlighting}[]
\NormalTok{printf}\OperatorTok{(}\StringTok{"}\SpecialCharTok{\%d}\StringTok{ }\SpecialCharTok{\%d\textbackslash{}n}\StringTok{"}\OperatorTok{,} \DecValTok{5}\OperatorTok{,} \CharTok{\textquotesingle{}5\textquotesingle{}}\OperatorTok{);}
\end{Highlighting}
\end{Shaded}

\end{tcolorbox}

Di sistemku yang memakai UTF-8, ini mencetak:

\begin{tcolorbox}[enhanced jigsaw,left=3pt,breakable,colback=white!97.2549!black,boxrule=0.2mm]

\begin{Shaded}
\begin{Highlighting}[]
\NormalTok{5 53}
\end{Highlighting}
\end{Shaded}

\end{tcolorbox}

Jadi\ldots{} bukan. Dan 53? Apa itu? Itu code point UTF-8 (dan ASCII) untuk simbol karakter \hl{\mbox{\texttt{\textquotesingle{}5\textquotesingle{}}}}\footnote{Setiap karakter punya nilai yang terkait dengannya pada skema character encoding mana pun.}

Lalu bagaimana cara mengonversi karakter \hl{\mbox{\texttt{\textquotesingle{}5\textquotesingle{}}}} (yang ternyata bernilai 53) menjadi nilai \hl{\mbox{\texttt{5}}}?

Dengan satu trik cerdik, begitulah caranya!

Standar C menjamin bahwa karakter-karakter ini punya code point yang berurutan dan dalam urutan ini:

\begin{tcolorbox}[enhanced jigsaw,left=3pt,breakable,colback=white!97.2549!black,boxrule=0.2mm]

\begin{Shaded}
\begin{Highlighting}[]
\NormalTok{0  1  2  3  4  5  6  7  8  9}
\end{Highlighting}
\end{Shaded}

\end{tcolorbox}

Renungkan sebentar--bagaimana kita bisa memakainya? Awas, spoiler\ldots{}

Mari kita lihat karakter-karakter itu beserta code point-nya di UTF-8:

\begin{tcolorbox}[enhanced jigsaw,left=3pt,breakable,colback=white!97.2549!black,boxrule=0.2mm]

\begin{Shaded}
\begin{Highlighting}[]
\NormalTok{0  1  2  3  4  5  6  7  8  9}
\NormalTok{48 49 50 51 52 53 54 55 56 57}
\end{Highlighting}
\end{Shaded}

\end{tcolorbox}

Kamu bisa lihat di sana bahwa \hl{\mbox{\texttt{\textquotesingle{}5\textquotesingle{}}}} adalah \hl{\mbox{\texttt{53}}}, persis seperti yang kita dapat tadi. Dan \hl{\mbox{\texttt{\textquotesingle{}0\textquotesingle{}}}} adalah \hl{\mbox{\texttt{48}}}.

Jadi kita bisa mengurangkan \hl{\mbox{\texttt{\textquotesingle{}0\textquotesingle{}}}} dari karakter digit mana pun untuk mendapatkan nilai numeriknya:

\begin{tcolorbox}[enhanced jigsaw,left=3pt,breakable,colback=white!97.2549!black,boxrule=0.2mm]

\begin{Shaded}
\begin{Highlighting}[]
\DataTypeTok{char}\NormalTok{ c }\OperatorTok{=} \CharTok{\textquotesingle{}6\textquotesingle{}}\OperatorTok{;}

\DataTypeTok{int}\NormalTok{ x }\OperatorTok{=}\NormalTok{ c}\OperatorTok{;}  \CommentTok{// x bernilai 54, code point untuk \textquotesingle{}6\textquotesingle{}}

\DataTypeTok{int}\NormalTok{ y }\OperatorTok{=}\NormalTok{ c }\OperatorTok{{-}} \CharTok{\textquotesingle{}0\textquotesingle{}}\OperatorTok{;} \CommentTok{// y bernilai 6, persis seperti yang kita mau}
\end{Highlighting}
\end{Shaded}

\end{tcolorbox}

Dan kita bisa mengonversi ke arah sebaliknya, cukup dengan menambahkan nilainya.

\begin{tcolorbox}[enhanced jigsaw,left=3pt,breakable,colback=white!97.2549!black,boxrule=0.2mm]

\begin{Shaded}
\begin{Highlighting}[]
\DataTypeTok{int}\NormalTok{ x }\OperatorTok{=} \DecValTok{6}\OperatorTok{;}

\DataTypeTok{char}\NormalTok{ c }\OperatorTok{=}\NormalTok{ x }\OperatorTok{+} \CharTok{\textquotesingle{}0\textquotesingle{}}\OperatorTok{;}  \CommentTok{// c bernilai 54}

\NormalTok{printf}\OperatorTok{(}\StringTok{"}\SpecialCharTok{\%d\textbackslash{}n}\StringTok{"}\OperatorTok{,}\NormalTok{ c}\OperatorTok{);}  \CommentTok{// mencetak 54}
\NormalTok{printf}\OperatorTok{(}\StringTok{"}\SpecialCharTok{\%c\textbackslash{}n}\StringTok{"}\OperatorTok{,}\NormalTok{ c}\OperatorTok{);}  \CommentTok{// mencetak 6 dengan \%c}
\end{Highlighting}
\end{Shaded}

\end{tcolorbox}

Mungkin kamu mengira ini cara yang aneh untuk melakukan konversi, dan menurut standar zaman sekarang, memang begitu. Tapi di zaman dahulu kala ketika komputer dibuat dari kayu secara harfiah, inilah metode untuk konversi itu. Dan karena tidak rusak, C tidak pernah memperbaikinya.

\index{Type conversions!char@\texttt{char}|)}

\section{Konversi Numerik}\label{konversi-numerik}

\index{Type conversions!numeric|(}

\subsection{Boolean}\label{boolean}

\index{Type conversions!Boolean}

Kalau nol dikonversi ke \hl{\mbox{\texttt{bool}}}, hasilnya \hl{\mbox{\texttt{0}}}. Selain itu hasilnya \hl{\mbox{\texttt{1}}}.

\subsection{Konversi Integer ke Integer}\label{konversi-integer-ke-integer}

\index{Type conversions!integer|(}

Kalau sebuah tipe integer dikonversi menjadi unsigned dan nilainya tidak muat, hasil unsigned-nya berputar seperti odometer sampai muat di unsigned itu\footnote{Dalam praktiknya, implementasimu kemungkinan cuma membuang bit-bit orde tinggi dari hasil, jadi angka 16-bit \hl{\mbox{\texttt{0x1234}}} yang dikonversi menjadi angka 8-bit berakhir sebagai \hl{\mbox{\texttt{0x0034}}}, atau cukup \hl{\mbox{\texttt{0x34}}}.}.

Kalau sebuah tipe integer dikonversi menjadi angka signed dan tidak muat, hasilnya ditentukan oleh implementasi! Sesuatu yang terdokumentasi akan terjadi, tapi kamu harus mencarinya sendiri\footnote{Lagi-lagi, dalam praktiknya, di sistemmu kemungkinan pola bit angka aslinya dipotong lalu dipakai begitu saja untuk merepresentasikan angka signed, two's complement. Misalnya, sistemku mengambil \hl{\mbox{\texttt{unsigned\ char}}} \hl{\mbox{\texttt{192}}} dan mengonversinya menjadi \hl{\mbox{\texttt{signed\ char}}} \hl{\mbox{\texttt{-64}}}. Dalam two's complement, pola bit kedua angka ini adalah biner \hl{\mbox{\texttt{11000000}}}.}

\subsection{Konversi Integer dan Floating Point}\label{konversi-integer-dan-floating-point}

\index{Type conversions!floating point|(}

Kalau tipe floating point dikonversi menjadi tipe integer, bagian pecahannya dibuang tanpa ampun\footnote{Sebenarnya tidak---cuma dibuang seperti biasa.}.

Tapi---dan ini jebakannya---kalau angkanya terlalu besar untuk muat di integer, hasilnya undefined behavior. Jadi jangan lakukan itu.

Dari integer atau floating point ke floating point, C berusaha sebaik mungkin mencari bilangan floating point terdekat dengan integer itu.

Sekali lagi, kalau nilai aslinya tidak bisa direpresentasikan, hasilnya undefined behavior.

\index{Type conversions!floating point|)} \index{Type conversions!integer|)}

\section{Konversi Implisit}\label{konversi-implisit}

\index{Type conversions!implicit|(}

Konversi ini dilakukan compiler secara otomatis ketika kamu mencampur berbagai tipe.

\subsection{Integer Promotion}\label{integer-promotions}

Di sejumlah tempat, kalau sebuah \hl{\mbox{\texttt{int}}} bisa dipakai untuk merepresentasikan nilai dari \hl{\mbox{\texttt{char}}} atau \hl{\mbox{\texttt{short}}} (signed maupun unsigned), nilai itu \emph{dipromosikan} naik menjadi \hl{\mbox{\texttt{int}}}. Kalau tidak muat di \hl{\mbox{\texttt{int}}}, ia dipromosikan menjadi \hl{\mbox{\texttt{unsigned\ int}}}.

Begitulah cara kita bisa melakukan hal seperti ini:

\begin{tcolorbox}[enhanced jigsaw,left=3pt,breakable,colback=white!97.2549!black,boxrule=0.2mm]

\begin{Shaded}
\begin{Highlighting}[]
\DataTypeTok{char}\NormalTok{ x }\OperatorTok{=} \DecValTok{10}\OperatorTok{,}\NormalTok{ y }\OperatorTok{=} \DecValTok{20}\OperatorTok{;}
\DataTypeTok{int}\NormalTok{ i }\OperatorTok{=}\NormalTok{ x }\OperatorTok{+}\NormalTok{ y}\OperatorTok{;}
\end{Highlighting}
\end{Shaded}

\end{tcolorbox}

Dalam kasus itu, \hl{\mbox{\texttt{x}}} dan \hl{\mbox{\texttt{y}}} dipromosikan oleh C menjadi \hl{\mbox{\texttt{int}}} sebelum perhitungan terjadi.

Integer promotion terjadi saat Usual Arithmetic Conversion, pada fungsi variadic\footnote{Fungsi dengan jumlah argumen yang bervariasi.}, operator unary \hl{\mbox{\texttt{+}}} dan \hl{\mbox{\texttt{-}}}, atau saat mengoper nilai ke fungsi tanpa prototype\footnote{Ini jarang dilakukan karena compiler akan protes dan punya prototype adalah \emph{Hal yang Benar}. Kurasa ini masih berfungsi karena alasan historis, sebelum prototype ada.}.

\subsection{Usual Arithmetic Conversion}\label{usual-arithmetic-conversions}

Konversi otomatis ini dikerjakan C pada operasi numerik yang kamu minta. (Memang itulah namanya, omong-omong, menurut C11 §6.3.1.8.) Perhatikan bahwa di bagian ini kita hanya membicarakan tipe numerik---string dibahas nanti.

Konversi ini menjawab pertanyaan tentang apa yang terjadi ketika tipe dicampur, seperti ini:

\begin{tcolorbox}[enhanced jigsaw,left=3pt,breakable,colback=white!97.2549!black,boxrule=0.2mm]

\begin{Shaded}
\begin{Highlighting}[]
\DataTypeTok{int}\NormalTok{ x }\OperatorTok{=} \DecValTok{3} \OperatorTok{+} \FloatTok{1.2}\OperatorTok{;}   \CommentTok{// Mencampur int dan double}
                   \CommentTok{// 4.2 dikonversi menjadi int}
                   \CommentTok{// 4 disimpan di x}

\DataTypeTok{float}\NormalTok{ y }\OperatorTok{=} \DecValTok{12} \OperatorTok{*} \DecValTok{2}\OperatorTok{;}  \CommentTok{// Mencampur float dan int}
                   \CommentTok{// 24 dikonversi menjadi float}
                   \CommentTok{// 24.0 disimpan di y}
\end{Highlighting}
\end{Shaded}

\end{tcolorbox}

Apakah mereka menjadi \hl{\mbox{\texttt{int}}}? Atau \hl{\mbox{\texttt{float}}}? Bagaimana cara kerjanya?

Berikut langkah-langkahnya, diparafrasekan agar mudah dicerna.

\begin{enumerate}
\def\labelenumi{\arabic{enumi}.}
\item
  Kalau salah satu bagian ekspresi bertipe floating, konversikan bagian lainnya ke tipe floating itu.
\item
  Kalau tidak, dan kedua tipe adalah tipe integer, lakukan integer promotion pada masing-masing, lalu jadikan tipe operand sebesar yang dibutuhkan untuk menampung nilai terbesar yang sama. Kadang ini melibatkan perubahan signed menjadi unsigned.
\end{enumerate}

Kalau ingin tahu detail lengkapnya, lihat C11 §6.3.1.8. Tapi kemungkinan besar kamu tidak perlu.

Pokoknya ingat bahwa tipe int menjadi tipe float kalau ada tipe floating point di mana pun di dalamnya, dan compiler berusaha memastikan tipe integer campuran tidak overflow.

Terakhir, kalau kamu mengonversi dari satu tipe floating point ke tipe floating point lain, compiler akan berusaha membuat konversi yang tepat. Kalau tidak bisa, ia akan membuat perkiraan terbaik yang ia bisa. Kalau angkanya terlalu besar untuk muat di tipe tujuan, \emph{duar}: undefined behavior!

\subsection{\texorpdfstring{\hl{\mbox{\texttt{void*}}}}{void*}}\label{void}

Tipe \hl{\mbox{\texttt{void*}}} menarik karena bisa dikonversi dari atau ke tipe pointer mana pun.

\begin{tcolorbox}[enhanced jigsaw,left=3pt,breakable,colback=white!97.2549!black,boxrule=0.2mm]

\begin{Shaded}
\begin{Highlighting}[]
\DataTypeTok{int}\NormalTok{ x }\OperatorTok{=} \DecValTok{10}\OperatorTok{;}

\DataTypeTok{void} \OperatorTok{*}\NormalTok{p }\OperatorTok{=} \OperatorTok{\&}\NormalTok{x}\OperatorTok{;}  \CommentTok{// \&x bertipe int*, tapi kita simpan di void*}

\DataTypeTok{int} \OperatorTok{*}\NormalTok{q }\OperatorTok{=}\NormalTok{ p}\OperatorTok{;}    \CommentTok{// p adalah void*, tapi kita simpan di int*}
\end{Highlighting}
\end{Shaded}

\end{tcolorbox}

\index{Type conversions!implicit|)}

\section{Konversi Eksplisit}\label{konversi-eksplisit}

\index{Type conversions!explicit|(}

Konversi jenis ini harus kamu minta sendiri; compiler tidak akan melakukannya untukmu.

Kamu bisa mengonversi dari satu tipe ke tipe lain dengan menugaskan satu tipe ke tipe lain memakai \hl{\mbox{\texttt{=}}}.

Kamu juga bisa mengonversi secara eksplisit dengan \emph{cast}.

\index{Type conversions!numeric|)}

\subsection{Casting}\label{casting}

\index{Type conversions!casting|(}

Kamu bisa mengubah tipe sebuah ekspresi secara eksplisit dengan menaruh tipe baru dalam tanda kurung di depannya. Sebagian developer C tidak menyukai praktik ini kecuali benar-benar perlu, tapi kemungkinan besar kamu akan menemui kode C yang memakainya.

Mari kita coba contoh ketika kita ingin mengonversi \hl{\mbox{\texttt{int}}} menjadi \hl{\mbox{\texttt{long}}} supaya bisa disimpan di \hl{\mbox{\texttt{long}}}.

Catatan: contoh ini dibuat-buat dan cast di sini sama sekali tidak perlu karena ekspresi \hl{\mbox{\texttt{x\ +\ 12}}} akan otomatis diubah menjadi \hl{\mbox{\texttt{long\ int}}} agar sesuai dengan tipe \hl{\mbox{\texttt{y}}} yang lebih lebar.

\begin{tcolorbox}[enhanced jigsaw,left=3pt,breakable,colback=white!97.2549!black,boxrule=0.2mm]

\begin{Shaded}
\begin{Highlighting}[]
\DataTypeTok{int}\NormalTok{ x }\OperatorTok{=} \DecValTok{10}\OperatorTok{;}
\DataTypeTok{long} \DataTypeTok{int}\NormalTok{ y }\OperatorTok{=} \OperatorTok{(}\DataTypeTok{long} \DataTypeTok{int}\OperatorTok{)}\NormalTok{x }\OperatorTok{+} \DecValTok{12}\OperatorTok{;}
\end{Highlighting}
\end{Shaded}

\end{tcolorbox}

Pada contoh itu, meskipun \hl{\mbox{\texttt{x}}} tadinya bertipe \hl{\mbox{\texttt{int}}}, ekspresi \hl{\mbox{\texttt{(long\ int)x}}} bertipe \hl{\mbox{\texttt{long\ int}}}. Kita bilang, ``Kita meng-cast \hl{\mbox{\texttt{x}}} menjadi \hl{\mbox{\texttt{long\ int}}}.''

Lebih sering, kamu mungkin melihat cast dipakai untuk mengonversi \hl{\mbox{\texttt{void*}}} menjadi tipe pointer tertentu agar bisa di-dereference.

Callback dari fungsi bawaan \hl{\mbox{\texttt{qsort()}}} bisa menunjukkan perilaku ini karena \hl{\mbox{\texttt{void*}}} dioper kepadanya:

\begin{tcolorbox}[enhanced jigsaw,left=3pt,breakable,colback=white!97.2549!black,boxrule=0.2mm]

\begin{Shaded}
\begin{Highlighting}[]
\DataTypeTok{int}\NormalTok{ compar}\OperatorTok{(}\DataTypeTok{const} \DataTypeTok{void} \OperatorTok{*}\NormalTok{elem1}\OperatorTok{,} \DataTypeTok{const} \DataTypeTok{void} \OperatorTok{*}\NormalTok{elem2}\OperatorTok{)}
\OperatorTok{\{}
    \ControlFlowTok{if} \OperatorTok{(*((}\DataTypeTok{const} \DataTypeTok{int}\OperatorTok{*)}\NormalTok{elem2}\OperatorTok{)} \OperatorTok{\textgreater{}} \OperatorTok{*((}\DataTypeTok{const} \DataTypeTok{int}\OperatorTok{*)}\NormalTok{elem1}\OperatorTok{))} \ControlFlowTok{return} \DecValTok{1}\OperatorTok{;}
    \ControlFlowTok{if} \OperatorTok{(*((}\DataTypeTok{const} \DataTypeTok{int}\OperatorTok{*)}\NormalTok{elem2}\OperatorTok{)} \OperatorTok{\textless{}} \OperatorTok{*((}\DataTypeTok{const} \DataTypeTok{int}\OperatorTok{*)}\NormalTok{elem1}\OperatorTok{))} \ControlFlowTok{return} \OperatorTok{{-}}\DecValTok{1}\OperatorTok{;}
    \ControlFlowTok{return} \DecValTok{0}\OperatorTok{;}
\OperatorTok{\}}
\end{Highlighting}
\end{Shaded}

\end{tcolorbox}

Tapi kamu juga bisa menulisnya dengan jelas lewat penugasan:

\begin{tcolorbox}[enhanced jigsaw,left=3pt,breakable,colback=white!97.2549!black,boxrule=0.2mm]

\begin{Shaded}
\begin{Highlighting}[]
\DataTypeTok{int}\NormalTok{ compar}\OperatorTok{(}\DataTypeTok{const} \DataTypeTok{void} \OperatorTok{*}\NormalTok{elem1}\OperatorTok{,} \DataTypeTok{const} \DataTypeTok{void} \OperatorTok{*}\NormalTok{elem2}\OperatorTok{)}
\OperatorTok{\{}
    \DataTypeTok{const} \DataTypeTok{int} \OperatorTok{*}\NormalTok{e1 }\OperatorTok{=}\NormalTok{ elem1}\OperatorTok{;}
    \DataTypeTok{const} \DataTypeTok{int} \OperatorTok{*}\NormalTok{e2 }\OperatorTok{=}\NormalTok{ elem2}\OperatorTok{;}

    \ControlFlowTok{return} \OperatorTok{*}\NormalTok{e2 }\OperatorTok{{-}} \OperatorTok{*}\NormalTok{e1}\OperatorTok{;}
\OperatorTok{\}}
\end{Highlighting}
\end{Shaded}

\end{tcolorbox}

Salah satu tempat yang lebih sering memakai cast adalah untuk menghindari peringatan saat mencetak nilai pointer dengan \hl{\mbox{\texttt{\%p}}} yang jarang dipakai, yang rewel terhadap apa pun selain \hl{\mbox{\texttt{void*}}}:

\begin{tcolorbox}[enhanced jigsaw,left=3pt,breakable,colback=white!97.2549!black,boxrule=0.2mm]

\begin{Shaded}
\begin{Highlighting}[]
\DataTypeTok{int}\NormalTok{ x }\OperatorTok{=} \DecValTok{3490}\OperatorTok{;}
\DataTypeTok{int} \OperatorTok{*}\NormalTok{p }\OperatorTok{=} \OperatorTok{\&}\NormalTok{x}\OperatorTok{;}

\NormalTok{printf}\OperatorTok{(}\StringTok{"}\SpecialCharTok{\%p\textbackslash{}n}\StringTok{"}\OperatorTok{,}\NormalTok{ p}\OperatorTok{);}
\end{Highlighting}
\end{Shaded}

\end{tcolorbox}

menghasilkan peringatan ini:

\begin{tcolorbox}[enhanced jigsaw,left=3pt,breakable,colback=white!97.2549!black,boxrule=0.2mm]

\begin{Shaded}
\begin{Highlighting}[]
\NormalTok{warning: format ‘\%p’ expects argument of type ‘void *’, but argument}
\NormalTok{         2 has type ‘int *’}
\end{Highlighting}
\end{Shaded}

\end{tcolorbox}

Kamu bisa memperbaikinya dengan cast:

\begin{tcolorbox}[enhanced jigsaw,left=3pt,breakable,colback=white!97.2549!black,boxrule=0.2mm]

\begin{Shaded}
\begin{Highlighting}[]
\NormalTok{printf}\OperatorTok{(}\StringTok{"}\SpecialCharTok{\%p\textbackslash{}n}\StringTok{"}\OperatorTok{,} \OperatorTok{(}\DataTypeTok{void} \OperatorTok{*)}\NormalTok{p}\OperatorTok{);}
\end{Highlighting}
\end{Shaded}

\end{tcolorbox}

Tempat lainnya adalah pada perubahan pointer eksplisit, kalau kamu tidak ingin memakai \hl{\mbox{\texttt{void*}}} sebagai perantara, tapi ini juga cukup jarang:

\begin{tcolorbox}[enhanced jigsaw,left=3pt,breakable,colback=white!97.2549!black,boxrule=0.2mm]

\begin{Shaded}
\begin{Highlighting}[]
\DataTypeTok{long}\NormalTok{ x }\OperatorTok{=} \DecValTok{3490}\OperatorTok{;}
\DataTypeTok{long} \OperatorTok{*}\NormalTok{p }\OperatorTok{=} \OperatorTok{\&}\NormalTok{x}\OperatorTok{;}
\DataTypeTok{unsigned} \DataTypeTok{char} \OperatorTok{*}\NormalTok{c }\OperatorTok{=} \OperatorTok{(}\DataTypeTok{unsigned} \DataTypeTok{char} \OperatorTok{*)}\NormalTok{p}\OperatorTok{;}
\end{Highlighting}
\end{Shaded}

\end{tcolorbox}

Tempat ketiga yang sering mewajibkannya adalah pada fungsi konversi karakter di \href{https://beej.us/guide/bgclr/html/split/ctype.html}{\hl{\mbox{\texttt{\textless{}ctype.h\textgreater{}}}}}\footnote{https://beej.us/guide/bgclr/html/split/ctype.html} yang mengharuskan kamu meng-cast nilai yang tandanya meragukan menjadi \hl{\mbox{\texttt{unsigned\ char}}} untuk menghindari undefined behavior.

Sekali lagi, cast jarang sekali \emph{dibutuhkan} dalam praktik. Kalau kamu mendapati dirimu sering meng-cast, mungkin ada cara lain untuk melakukan hal yang sama, atau mungkin kamu meng-cast secara tidak perlu.

Atau mungkin memang perlu. Secara pribadi, aku berusaha menghindarinya, tapi tidak takut memakainya kalau harus.

\index{Type conversions!explicit|)} \index{Type conversions!casting|)} \index{Type conversions|)}

\chapter{Tipe IV: Qualifier dan Specifier}\label{tipe-iv-qualifier-dan-specifier}

Sekarang setelah kita punya lebih banyak tipe, ternyata tipe-tipe ini bisa diberi atribut tambahan yang mengatur perilakunya. Inilah \emph{type qualifier} dan \emph{storage-class specifier}.

\section{Type Qualifier}\label{type-qualifier}

\index{Type qualifiers|(}

Dengan ini kamu bisa mendeklarasikan nilai konstan, dan juga memberi petunjuk optimasi kepada compiler yang bisa dimanfaatkannya.

\subsection{\texorpdfstring{\hl{\mbox{\texttt{const}}}}{const}}\label{const}

\index{const type qualifier@\texttt{const} type qualifier|(}

Ini type qualifier yang paling umum kamu temui. Artinya variabel itu konstan, dan setiap usaha mengubahnya akan membuat compiler sangat marah.

\begin{tcolorbox}[enhanced jigsaw,left=3pt,breakable,colback=white!97.2549!black,boxrule=0.2mm]

\begin{Shaded}
\begin{Highlighting}[]
\DataTypeTok{const} \DataTypeTok{int}\NormalTok{ x }\OperatorTok{=} \DecValTok{2}\OperatorTok{;}

\NormalTok{x }\OperatorTok{=} \DecValTok{4}\OperatorTok{;}  \CommentTok{// SUARA COMPILER MUNTAH, tidak bisa menugaskan ke konstanta}
\end{Highlighting}
\end{Shaded}

\end{tcolorbox}

Nilai \hl{\mbox{\texttt{const}}} tidak bisa diubah.

Sering kali kamu melihat \hl{\mbox{\texttt{const}}} di daftar parameter fungsi:

\begin{tcolorbox}[enhanced jigsaw,left=3pt,breakable,colback=white!97.2549!black,boxrule=0.2mm]

\begin{Shaded}
\begin{Highlighting}[]
\DataTypeTok{void}\NormalTok{ foo}\OperatorTok{(}\DataTypeTok{const} \DataTypeTok{int}\NormalTok{ x}\OperatorTok{)}
\OperatorTok{\{}
\NormalTok{    printf}\OperatorTok{(}\StringTok{"}\SpecialCharTok{\%d\textbackslash{}n}\StringTok{"}\OperatorTok{,}\NormalTok{ x }\OperatorTok{+} \DecValTok{30}\OperatorTok{);}  \CommentTok{// OK, tidak mengubah "x"}
\OperatorTok{\}}
\end{Highlighting}
\end{Shaded}

\end{tcolorbox}

\subsubsection{\texorpdfstring{\hl{\mbox{\texttt{const}}} dan Pointer}{const dan Pointer}}\label{const-dan-pointer}

\index{const type qualifier@\texttt{const} type qualifier!and pointers|(}

Ini agak membingungkan, karena ada dua pemakaian yang punya dua makna ketika menyangkut pointer.

Pertama, kita bisa membuat agar nilai yang ditunjuk pointer tidak bisa diubah. Caranya dengan menaruh \hl{\mbox{\texttt{const}}} di depan nama tipe (sebelum asterisk) pada deklarasi tipe.

\begin{tcolorbox}[enhanced jigsaw,left=3pt,breakable,colback=white!97.2549!black,boxrule=0.2mm]

\begin{Shaded}
\begin{Highlighting}[]
\DataTypeTok{int}\NormalTok{ x}\OperatorTok{[]} \OperatorTok{=} \OperatorTok{\{}\DecValTok{10}\OperatorTok{,} \DecValTok{20}\OperatorTok{\};}
\DataTypeTok{const} \DataTypeTok{int} \OperatorTok{*}\NormalTok{p }\OperatorTok{=}\NormalTok{ x}\OperatorTok{;} 

\NormalTok{p}\OperatorTok{++;}  \CommentTok{// Kita bisa mengubah p, tidak masalah}

\OperatorTok{*}\NormalTok{p }\OperatorTok{=} \DecValTok{30}\OperatorTok{;} \CommentTok{// Error compiler! Tidak bisa mengubah yang ditunjuknya}
\end{Highlighting}
\end{Shaded}

\end{tcolorbox}

Agak membingungkan, dua hal ini setara:

\begin{tcolorbox}[enhanced jigsaw,left=3pt,breakable,colback=white!97.2549!black,boxrule=0.2mm]

\begin{Shaded}
\begin{Highlighting}[]
\DataTypeTok{const} \DataTypeTok{int} \OperatorTok{*}\NormalTok{p}\OperatorTok{;}  \CommentTok{// Tidak bisa mengubah yang ditunjuk p}
\DataTypeTok{int} \DataTypeTok{const} \OperatorTok{*}\NormalTok{p}\OperatorTok{;}  \CommentTok{// Tidak bisa mengubah yang ditunjuk p, sama seperti di atas}
\end{Highlighting}
\end{Shaded}

\end{tcolorbox}

Bagus, jadi kita tidak bisa mengubah yang ditunjuk pointer, tapi kita bisa mengubah pointer-nya sendiri. Bagaimana kalau kita ingin sebaliknya? Kita ingin bisa mengubah yang ditunjuk pointer, tapi \emph{tidak} bisa mengubah pointer-nya sendiri?

Cukup pindahkan \hl{\mbox{\texttt{const}}} ke setelah asterisk pada deklarasi:

\begin{tcolorbox}[enhanced jigsaw,left=3pt,breakable,colback=white!97.2549!black,boxrule=0.2mm]

\begin{Shaded}
\begin{Highlighting}[]
\DataTypeTok{int} \OperatorTok{*}\DataTypeTok{const}\NormalTok{ p}\OperatorTok{;}   \CommentTok{// Kita tidak bisa mengubah "p" dengan aritmetika pointer}

\NormalTok{p}\OperatorTok{++;}  \CommentTok{// Error compiler!}
\end{Highlighting}
\end{Shaded}

\end{tcolorbox}

Tapi kita bisa mengubah yang ditunjuknya:

\begin{tcolorbox}[enhanced jigsaw,left=3pt,breakable,colback=white!97.2549!black,boxrule=0.2mm]

\begin{Shaded}
\begin{Highlighting}[]
\DataTypeTok{int}\NormalTok{ x }\OperatorTok{=} \DecValTok{10}\OperatorTok{;}
\DataTypeTok{int} \OperatorTok{*}\DataTypeTok{const}\NormalTok{ p }\OperatorTok{=} \OperatorTok{\&}\NormalTok{x}\OperatorTok{;}

\OperatorTok{*}\NormalTok{p }\OperatorTok{=} \DecValTok{20}\OperatorTok{;}   \CommentTok{// Ubah "x" menjadi 20, tidak masalah}
\end{Highlighting}
\end{Shaded}

\end{tcolorbox}

Kamu juga bisa membuat keduanya \hl{\mbox{\texttt{const}}}:

\begin{tcolorbox}[enhanced jigsaw,left=3pt,breakable,colback=white!97.2549!black,boxrule=0.2mm]

\begin{Shaded}
\begin{Highlighting}[]
\DataTypeTok{const} \DataTypeTok{int} \OperatorTok{*}\DataTypeTok{const}\NormalTok{ p}\OperatorTok{;}  \CommentTok{// Tidak bisa mengubah p maupun *p!}
\end{Highlighting}
\end{Shaded}

\end{tcolorbox}

Terakhir, kalau ada beberapa tingkat indirection, kamu harus memberi \hl{\mbox{\texttt{const}}} pada tingkat yang sesuai. Hanya karena sebuah pointer \hl{\mbox{\texttt{const}}}, bukan berarti pointer yang ditunjuknya juga harus \hl{\mbox{\texttt{const}}}. Kamu bisa mengaturnya secara eksplisit seperti pada contoh-contoh berikut:

\begin{tcolorbox}[enhanced jigsaw,left=3pt,breakable,colback=white!97.2549!black,boxrule=0.2mm]

\begin{Shaded}
\begin{Highlighting}[]
\DataTypeTok{char} \OperatorTok{**}\NormalTok{p}\OperatorTok{;}
\NormalTok{p}\OperatorTok{++;}     \CommentTok{// OK!}
\OperatorTok{(*}\NormalTok{p}\OperatorTok{)++;}  \CommentTok{// OK!}

\DataTypeTok{char} \OperatorTok{**}\DataTypeTok{const}\NormalTok{ p}\OperatorTok{;}
\NormalTok{p}\OperatorTok{++;}     \CommentTok{// Error!}
\OperatorTok{(*}\NormalTok{p}\OperatorTok{)++;}  \CommentTok{// OK!}

\DataTypeTok{char} \OperatorTok{*}\DataTypeTok{const} \OperatorTok{*}\NormalTok{p}\OperatorTok{;}
\NormalTok{p}\OperatorTok{++;}     \CommentTok{// OK!}
\OperatorTok{(*}\NormalTok{p}\OperatorTok{)++;}  \CommentTok{// Error!}

\DataTypeTok{char} \OperatorTok{*}\DataTypeTok{const} \OperatorTok{*}\DataTypeTok{const}\NormalTok{ p}\OperatorTok{;}
\NormalTok{p}\OperatorTok{++;}     \CommentTok{// Error!}
\OperatorTok{(*}\NormalTok{p}\OperatorTok{)++;}  \CommentTok{// Error!}
\end{Highlighting}
\end{Shaded}

\end{tcolorbox}

\index{const type qualifier@\texttt{const} type qualifier!and pointers|)}

\subsubsection{\texorpdfstring{\hl{\mbox{\texttt{const}}} Correctness}{const Correctness}}\label{const-correctness}

\index{const type qualifier@\texttt{const} type qualifier!correctness|(}

Ada satu hal lagi yang harus kusebutkan: compiler akan memberi peringatan untuk kode seperti ini:

\begin{tcolorbox}[enhanced jigsaw,left=3pt,breakable,colback=white!97.2549!black,boxrule=0.2mm]

\begin{Shaded}
\begin{Highlighting}[]
\DataTypeTok{const} \DataTypeTok{int}\NormalTok{ x }\OperatorTok{=} \DecValTok{20}\OperatorTok{;}
\DataTypeTok{int} \OperatorTok{*}\NormalTok{p }\OperatorTok{=} \OperatorTok{\&}\NormalTok{x}\OperatorTok{;}
\end{Highlighting}
\end{Shaded}

\end{tcolorbox}

dengan pesan yang kira-kira berbunyi:

\begin{tcolorbox}[enhanced jigsaw,left=3pt,breakable,colback=white!97.2549!black,boxrule=0.2mm]

\begin{Shaded}
\begin{Highlighting}[]
\NormalTok{initialization discards \textquotesingle{}const\textquotesingle{} qualifier from pointer type target}
\end{Highlighting}
\end{Shaded}

\end{tcolorbox}

Apa yang terjadi di sana?

Kita perlu melihat tipe di kedua sisi penugasan:

\begin{tcolorbox}[enhanced jigsaw,left=3pt,breakable,colback=white!97.2549!black,boxrule=0.2mm]

\begin{Shaded}
\begin{Highlighting}[]
    \DataTypeTok{const} \DataTypeTok{int}\NormalTok{ x }\OperatorTok{=} \DecValTok{20}\OperatorTok{;}
    \DataTypeTok{int} \OperatorTok{*}\NormalTok{p }\OperatorTok{=} \OperatorTok{\&}\NormalTok{x}\OperatorTok{;}
\CommentTok{//    \^{}       \^{}}
\CommentTok{//    |       |}
\CommentTok{//  int*    const int*}
\end{Highlighting}
\end{Shaded}

\end{tcolorbox}

Compiler memperingatkan bahwa nilai di sisi kanan penugasan itu \hl{\mbox{\texttt{const}}}, tapi yang di sisi kiri tidak. Compiler memberi tahu bahwa ia membuang sifat ``const'' dari ekspresi di sisi kanan.

Artinya, kita \emph{masih bisa} mencoba hal berikut, tapi itu salah. Compiler akan memperingatkan, dan hasilnya undefined behavior:

\begin{tcolorbox}[enhanced jigsaw,left=3pt,breakable,colback=white!97.2549!black,boxrule=0.2mm]

\begin{Shaded}
\begin{Highlighting}[]
\DataTypeTok{const} \DataTypeTok{int}\NormalTok{ x }\OperatorTok{=} \DecValTok{20}\OperatorTok{;}
\DataTypeTok{int} \OperatorTok{*}\NormalTok{p }\OperatorTok{=} \OperatorTok{\&}\NormalTok{x}\OperatorTok{;}

\OperatorTok{*}\NormalTok{p }\OperatorTok{=} \DecValTok{40}\OperatorTok{;}  \CommentTok{// Undefined behavior{-}{-}mungkin mengubah "x", mungkin tidak!}

\NormalTok{printf}\OperatorTok{(}\StringTok{"}\SpecialCharTok{\%d\textbackslash{}n}\StringTok{"}\OperatorTok{,}\NormalTok{ x}\OperatorTok{);}  \CommentTok{// 40, kalau beruntung}
\end{Highlighting}
\end{Shaded}

\end{tcolorbox}

\index{const type qualifier@\texttt{const} type qualifier!correctness|)} \index{const type qualifier@\texttt{const} type qualifier|)}

\subsection{\texorpdfstring{\hl{\mbox{\texttt{restrict}}}}{restrict}}\label{restrict}

\index{restrict type qualifier@\texttt{restrict} type qualifier|(}

Singkatnya: kamu tidak pernah wajib memakai ini dan boleh mengabaikannya setiap kali melihatnya. Kalau dipakai dengan benar, kemungkinan ada sedikit peningkatan performa. Kalau dipakai dengan salah, kamu akan mendapat undefined behavior.

\hl{\mbox{\texttt{restrict}}} adalah petunjuk bagi compiler bahwa sepotong memori tertentu hanya akan diakses oleh satu pointer dan tidak pernah oleh pointer lain. (Artinya, tidak ada aliasing pada object yang ditunjuk pointer \hl{\mbox{\texttt{restrict}}} itu.) Kalau developer mendeklarasikan pointer sebagai \hl{\mbox{\texttt{restrict}}} lalu mengakses object yang ditunjuknya dengan cara lain (misalnya lewat pointer lain), perilakunya undefined.

Pada dasarnya kamu berkata kepada C, ``Hei---kujamin pointer tunggal ini satu-satunya cara aku mengakses memori ini, dan kalau aku bohong, silakan beri aku undefined behavior.''

Dengan informasi itu, C bisa mengoptimasi hal-hal tertentu. Misalnya, kalau kamu men-dereference pointer \hl{\mbox{\texttt{restrict}}} berulang kali dalam loop, C mungkin memutuskan menyimpan hasilnya di register dan hanya menyimpan hasil akhirnya setelah loop selesai. Kalau ada pointer lain yang merujuk memori yang sama dan mengaksesnya di dalam loop, hasilnya tidak akan akurat.

(Perhatikan bahwa \hl{\mbox{\texttt{restrict}}} tidak berpengaruh kalau object yang ditunjuk tidak pernah ditulis. Semuanya tentang optimasi seputar penulisan ke memori.)

Mari kita tulis fungsi untuk menukar dua variabel, dan kita pakai kata kunci \hl{\mbox{\texttt{restrict}}} untuk meyakinkan C bahwa kita tidak akan pernah mengoper pointer ke hal yang sama. Lalu mari kita sengaja membuat kesalahan dan mencoba mengoper pointer ke hal yang sama.

\begin{tcolorbox}[enhanced jigsaw,left=3pt,breakable,colback=white!97.2549!black,boxrule=0.2mm]

\begin{Shaded}
\begin{Highlighting}[numbers=left,,]
\DataTypeTok{void}\NormalTok{ swap}\OperatorTok{(}\DataTypeTok{int} \OperatorTok{*}\DataTypeTok{restrict}\NormalTok{ a}\OperatorTok{,} \DataTypeTok{int} \OperatorTok{*}\DataTypeTok{restrict}\NormalTok{ b}\OperatorTok{)}
\OperatorTok{\{}
    \DataTypeTok{int}\NormalTok{ t}\OperatorTok{;}

\NormalTok{    t }\OperatorTok{=} \OperatorTok{*}\NormalTok{a}\OperatorTok{;}
    \OperatorTok{*}\NormalTok{a }\OperatorTok{=} \OperatorTok{*}\NormalTok{b}\OperatorTok{;}
    \OperatorTok{*}\NormalTok{b }\OperatorTok{=}\NormalTok{ t}\OperatorTok{;}
\OperatorTok{\}}

\DataTypeTok{int}\NormalTok{ main}\OperatorTok{(}\DataTypeTok{void}\OperatorTok{)}
\OperatorTok{\{}
    \DataTypeTok{int}\NormalTok{ x }\OperatorTok{=} \DecValTok{10}\OperatorTok{,}\NormalTok{ y }\OperatorTok{=} \DecValTok{20}\OperatorTok{;}

\NormalTok{    swap}\OperatorTok{(\&}\NormalTok{x}\OperatorTok{,} \OperatorTok{\&}\NormalTok{y}\OperatorTok{);}  \CommentTok{// OK! "a" dan "b" di atas menunjuk ke hal yang berbeda}

\NormalTok{    swap}\OperatorTok{(\&}\NormalTok{x}\OperatorTok{,} \OperatorTok{\&}\NormalTok{x}\OperatorTok{);}  \CommentTok{// Undefined behavior! "a" dan "b" sama!}
\OperatorTok{\}}
\end{Highlighting}
\end{Shaded}

\end{tcolorbox}

Kalau kata kunci \hl{\mbox{\texttt{restrict}}} di atas dihapus, kedua pemanggilan akan aman. Tapi compiler mungkin tidak bisa mengoptimasi.

\hl{\mbox{\texttt{restrict}}} punya block scope, artinya pembatasannya hanya berlaku selama scope tempat ia dipakai. Kalau ada di daftar parameter fungsi, ia ada di block scope fungsi itu.

Kalau pointer yang dibatasi menunjuk ke array, pembatasan itu hanya berlaku pada object individual di array. Pointer lain boleh membaca dan menulis array itu selama tidak membaca atau menulis elemen yang sama dengan yang dibatasi.

Kalau \hl{\mbox{\texttt{restrict}}} dipakai di file scope, di luar fungsi apa pun, pembatasannya mencakup seluruh program.

Kamu kemungkinan besar akan melihatnya di fungsi library seperti \hl{\mbox{\texttt{printf()}}}:

\begin{tcolorbox}[enhanced jigsaw,left=3pt,breakable,colback=white!97.2549!black,boxrule=0.2mm]

\begin{Shaded}
\begin{Highlighting}[]
\DataTypeTok{int}\NormalTok{ printf}\OperatorTok{(}\DataTypeTok{const} \DataTypeTok{char} \OperatorTok{*} \DataTypeTok{restrict}\NormalTok{ format}\OperatorTok{,} \OperatorTok{...);}
\end{Highlighting}
\end{Shaded}

\end{tcolorbox}

Sekali lagi, itu hanya memberi tahu compiler bahwa di dalam fungsi \hl{\mbox{\texttt{printf()}}}, hanya akan ada satu pointer yang merujuk ke bagian mana pun dari string \hl{\mbox{\texttt{format}}} itu.

Satu catatan terakhir: kalau karena suatu alasan kamu memakai notasi array pada parameter fungsi alih-alih notasi pointer, kamu bisa memakai \hl{\mbox{\texttt{restrict}}} seperti ini:

\begin{tcolorbox}[enhanced jigsaw,left=3pt,breakable,colback=white!97.2549!black,boxrule=0.2mm]

\begin{Shaded}
\begin{Highlighting}[]
\DataTypeTok{void}\NormalTok{ foo}\OperatorTok{(}\DataTypeTok{int}\NormalTok{ p}\OperatorTok{[}\DataTypeTok{restrict}\OperatorTok{])}     \CommentTok{// Tanpa ukuran}

\DataTypeTok{void}\NormalTok{ foo}\OperatorTok{(}\DataTypeTok{int}\NormalTok{ p}\OperatorTok{[}\DataTypeTok{restrict} \DecValTok{10}\OperatorTok{])}  \CommentTok{// Atau dengan ukuran}
\end{Highlighting}
\end{Shaded}

\end{tcolorbox}

Tapi notasi pointer lebih umum.

\index{restrict type qualifier@\texttt{restrict} type qualifier|)}

\subsection{\texorpdfstring{\hl{\mbox{\texttt{volatile}}}}{volatile}}\label{volatile}

\index{volatile type qualifier@\texttt{volatile} type qualifier|(}

Kecil kemungkinan kamu melihat atau membutuhkan ini kecuali kamu berurusan langsung dengan hardware.

\hl{\mbox{\texttt{volatile}}} memberi tahu compiler bahwa sebuah nilai mungkin berubah di luar pengetahuannya dan harus dibaca ulang setiap kali.

Contohnya ketika compiler membaca memori di sebuah alamat yang terus diperbarui di balik layar, misalnya semacam timer hardware.

Kalau compiler memutuskan mengoptimasinya dan menyimpan nilai di register dalam waktu lama, nilai di memori akan berubah dan tidak tercermin di register.

Dengan mendeklarasikan sesuatu \hl{\mbox{\texttt{volatile}}}, kamu berkata kepada compiler, ``Hei, yang ditunjuk ini mungkin berubah kapan saja karena alasan di luar kode program ini.''

\begin{tcolorbox}[enhanced jigsaw,left=3pt,breakable,colback=white!97.2549!black,boxrule=0.2mm]

\begin{Shaded}
\begin{Highlighting}[]
\DataTypeTok{volatile} \DataTypeTok{int} \OperatorTok{*}\NormalTok{p}\OperatorTok{;}
\end{Highlighting}
\end{Shaded}

\end{tcolorbox}

\subsection{\texorpdfstring{\hl{\mbox{\texttt{\_Atomic}}}}{\_Atomic}}\label{atomic}

Ini fitur C opsional yang akan kita bahas di \hyperref[chapter-atomics]{bab Atomics}.

\index{volatile type qualifier@\texttt{volatile} type qualifier|)} \index{Type qualifiers|)}

\section{Storage-Class Specifier}\label{storage-class-specifier}

\index{Storage-Class Specifiers|(}

Storage-class specifier mirip type qualifier. Mereka memberi compiler informasi tambahan tentang tipe sebuah variabel.

\subsection{\texorpdfstring{\hl{\mbox{\texttt{auto}}}}{auto}}\label{auto}

\index{auto storage class@\texttt{auto} storage class|(}

Kata kunci ini hampir tidak pernah terlihat, karena \hl{\mbox{\texttt{auto}}} adalah bawaan untuk variabel block scope. Ia tersirat.

Keduanya sama:

\begin{tcolorbox}[enhanced jigsaw,left=3pt,breakable,colback=white!97.2549!black,boxrule=0.2mm]

\begin{Shaded}
\begin{Highlighting}[]
\OperatorTok{\{}
    \DataTypeTok{int}\NormalTok{ a}\OperatorTok{;}         \CommentTok{// auto adalah bawaannya...}
    \KeywordTok{auto} \DataTypeTok{int}\NormalTok{ a}\OperatorTok{;}    \CommentTok{// Jadi ini mubazir}
\OperatorTok{\}}
\end{Highlighting}
\end{Shaded}

\end{tcolorbox}

Kata kunci \hl{\mbox{\texttt{auto}}} menandakan bahwa object ini punya \emph{automatic storage duration}. Artinya, ia ada di scope tempat ia didefinisikan, dan otomatis didealokasikan saat scope itu ditinggalkan.

Satu jebakan pada variabel otomatis: nilainya tak tentu sampai kamu menginisialisasinya secara eksplisit. Kita bilang isinya ``acak'' atau ``sampah'', walau aku sendiri kurang puas dengan keduanya. Pokoknya, kamu tidak akan tahu isinya kecuali kamu menginisialisasinya.

\index{auto storage class@\texttt{auto} storage class|)}

Selalu inisialisasi semua variabel otomatis sebelum dipakai!

\subsection{\texorpdfstring{\hl{\mbox{\texttt{static}}}}{static}}\label{static}

\index{static storage class@\texttt{static} storage class|(}

Kata kunci ini punya dua makna, tergantung apakah variabelnya file scope atau block scope.

Kita mulai dengan block scope.

\subsubsection{\texorpdfstring{\hl{\mbox{\texttt{static}}} di Block Scope}{static di Block Scope}}\label{static-di-block-scope}

\index{static storage class@\texttt{static} storage class!in block scope|(}

Dalam kasus ini, intinya kita berkata, ``Aku cuma ingin satu instance variabel ini ada, dipakai bersama antarpemanggilan.''

Artinya, nilainya bertahan di antara pemanggilan.

\hl{\mbox{\texttt{static}}} di block scope dengan initializer hanya diinisialisasi satu kali saat program mulai, bukan setiap kali fungsi dipanggil.

Kita coba contoh:

\begin{tcolorbox}[enhanced jigsaw,left=3pt,breakable,colback=white!97.2549!black,boxrule=0.2mm]

\begin{Shaded}
\begin{Highlighting}[numbers=left,,]
\PreprocessorTok{\#include }\ImportTok{\textless{}stdio.h\textgreater{}}

\DataTypeTok{void}\NormalTok{ counter}\OperatorTok{(}\DataTypeTok{void}\OperatorTok{)}
\OperatorTok{\{}
    \DataTypeTok{static} \DataTypeTok{int}\NormalTok{ count }\OperatorTok{=} \DecValTok{1}\OperatorTok{;}  \CommentTok{// Ini diinisialisasi satu kali}

\NormalTok{    printf}\OperatorTok{(}\StringTok{"This has been called }\SpecialCharTok{\%d}\StringTok{ time(s)}\SpecialCharTok{\textbackslash{}n}\StringTok{"}\OperatorTok{,}\NormalTok{ count}\OperatorTok{);}

\NormalTok{    count}\OperatorTok{++;}
\OperatorTok{\}}

\DataTypeTok{int}\NormalTok{ main}\OperatorTok{(}\DataTypeTok{void}\OperatorTok{)}
\OperatorTok{\{}
\NormalTok{    counter}\OperatorTok{();}  \CommentTok{// "This has been called 1 time(s)"}
\NormalTok{    counter}\OperatorTok{();}  \CommentTok{// "This has been called 2 time(s)"}
\NormalTok{    counter}\OperatorTok{();}  \CommentTok{// "This has been called 3 time(s)"}
\NormalTok{    counter}\OperatorTok{();}  \CommentTok{// "This has been called 4 time(s)"}
\OperatorTok{\}}
\end{Highlighting}
\end{Shaded}

\end{tcolorbox}

Lihat bagaimana nilai \hl{\mbox{\texttt{count}}} bertahan di antara pemanggilan?

Satu hal yang perlu dicatat: variabel \hl{\mbox{\texttt{static}}} block scope diinisialisasi ke \hl{\mbox{\texttt{0}}} secara bawaan.

\begin{tcolorbox}[enhanced jigsaw,left=3pt,breakable,colback=white!97.2549!black,boxrule=0.2mm]

\begin{Shaded}
\begin{Highlighting}[]
\DataTypeTok{static} \DataTypeTok{int}\NormalTok{ foo}\OperatorTok{;}      \CommentTok{// Nilai awal bawaannya \textasciigrave{}0\textasciigrave{}...}
\DataTypeTok{static} \DataTypeTok{int}\NormalTok{ foo }\OperatorTok{=} \DecValTok{0}\OperatorTok{;}  \CommentTok{// Jadi penugasan \textasciigrave{}0\textasciigrave{} ini mubazir}
\end{Highlighting}
\end{Shaded}

\end{tcolorbox}

Terakhir, perlu diingat bahwa kalau kamu menulis program multithread, kamu harus memastikan beberapa thread tidak saling menginjak variabel yang sama.

\index{static storage class@\texttt{static} storage class!in block scope|)}

\subsubsection{\texorpdfstring{\hl{\mbox{\texttt{static}}} di File Scope}{static di File Scope}}\label{static-di-file-scope}

\index{static storage class@\texttt{static} storage class!in file scope|(}

Ketika kita sampai di file scope, di luar blok mana pun, maknanya agak berubah.

Variabel di file scope sudah bertahan di antara pemanggilan fungsi, jadi perilaku itu sudah ada.

Sebagai gantinya, dalam konteks ini \hl{\mbox{\texttt{static}}} berarti variabel tersebut tidak terlihat di luar file sumber ini. Mirip ``global'', tapi hanya di file ini.

Lebih lanjut dibahas di bagian tentang membangun dengan banyak file sumber.

\index{static storage class@\texttt{static} storage class!in file scope|)} \index{static storage class@\texttt{static} storage class|)}

\subsection{\texorpdfstring{\hl{\mbox{\texttt{extern}}}}{extern}}\label{extern}

\index{extern storage class@\texttt{extern} storage class|(}

Storage-class specifier \hl{\mbox{\texttt{extern}}} memberi kita cara untuk merujuk ke object di file sumber lain.

Katakanlah, misalnya, file \hl{\mbox{\texttt{bar.c}}} isinya hanya ini:

\begin{tcolorbox}[enhanced jigsaw,left=3pt,breakable,colback=white!97.2549!black,boxrule=0.2mm]

\begin{Shaded}
\begin{Highlighting}[numbers=left,,]
\CommentTok{// bar.c}

\DataTypeTok{int}\NormalTok{ a }\OperatorTok{=} \DecValTok{37}\OperatorTok{;}
\end{Highlighting}
\end{Shaded}

\end{tcolorbox}

Hanya itu. Mendeklarasikan \hl{\mbox{\texttt{int\ a}}} baru di file scope.

Tapi bagaimana kalau ada file sumber lain, \hl{\mbox{\texttt{foo.c}}}, dan kita ingin merujuk ke \hl{\mbox{\texttt{a}}} yang ada di \hl{\mbox{\texttt{bar.c}}}?

Mudah dengan kata kunci \hl{\mbox{\texttt{extern}}}:

\begin{tcolorbox}[enhanced jigsaw,left=3pt,breakable,colback=white!97.2549!black,boxrule=0.2mm]

\begin{Shaded}
\begin{Highlighting}[numbers=left,,]
\CommentTok{// foo.c}

\KeywordTok{extern} \DataTypeTok{int}\NormalTok{ a}\OperatorTok{;}

\DataTypeTok{int}\NormalTok{ main}\OperatorTok{(}\DataTypeTok{void}\OperatorTok{)}
\OperatorTok{\{}
\NormalTok{    printf}\OperatorTok{(}\StringTok{"}\SpecialCharTok{\%d\textbackslash{}n}\StringTok{"}\OperatorTok{,}\NormalTok{ a}\OperatorTok{);}  \CommentTok{// 37, dari bar.c!}

\NormalTok{    a }\OperatorTok{=} \DecValTok{99}\OperatorTok{;}

\NormalTok{    printf}\OperatorTok{(}\StringTok{"}\SpecialCharTok{\%d\textbackslash{}n}\StringTok{"}\OperatorTok{,}\NormalTok{ a}\OperatorTok{);}  \CommentTok{// "a" yang sama dari bar.c, tapi sekarang 99}
\OperatorTok{\}}
\end{Highlighting}
\end{Shaded}

\end{tcolorbox}

Kita juga bisa menaruh \hl{\mbox{\texttt{extern\ int\ a}}} di block scope, dan ia tetap akan merujuk ke \hl{\mbox{\texttt{a}}} di \hl{\mbox{\texttt{bar.c}}}:

\begin{tcolorbox}[enhanced jigsaw,left=3pt,breakable,colback=white!97.2549!black,boxrule=0.2mm]

\begin{Shaded}
\begin{Highlighting}[numbers=left,,]
\CommentTok{// foo.c}

\DataTypeTok{int}\NormalTok{ main}\OperatorTok{(}\DataTypeTok{void}\OperatorTok{)}
\OperatorTok{\{}
    \KeywordTok{extern} \DataTypeTok{int}\NormalTok{ a}\OperatorTok{;}

\NormalTok{    printf}\OperatorTok{(}\StringTok{"}\SpecialCharTok{\%d\textbackslash{}n}\StringTok{"}\OperatorTok{,}\NormalTok{ a}\OperatorTok{);}  \CommentTok{// 37, dari bar.c!}

\NormalTok{    a }\OperatorTok{=} \DecValTok{99}\OperatorTok{;}

\NormalTok{    printf}\OperatorTok{(}\StringTok{"}\SpecialCharTok{\%d\textbackslash{}n}\StringTok{"}\OperatorTok{,}\NormalTok{ a}\OperatorTok{);}  \CommentTok{// "a" yang sama dari bar.c, tapi sekarang 99}
\OperatorTok{\}}
\end{Highlighting}
\end{Shaded}

\end{tcolorbox}

Nah, kalau \hl{\mbox{\texttt{a}}} di \hl{\mbox{\texttt{bar.c}}} ditandai \hl{\mbox{\texttt{static}}}, ini tidak akan berfungsi. Variabel \hl{\mbox{\texttt{static}}} di file scope tidak terlihat di luar file itu.

Catatan terakhir tentang \hl{\mbox{\texttt{extern}}} pada fungsi. Untuk fungsi, \hl{\mbox{\texttt{extern}}} adalah bawaan, jadi mubazir. Kamu bisa mendeklarasikan fungsi \hl{\mbox{\texttt{static}}} kalau hanya ingin fungsi itu terlihat di satu file sumber.

\index{extern storage class@\texttt{extern} storage class|)}

\subsection{\texorpdfstring{\hl{\mbox{\texttt{register}}}}{register}}\label{register}

\index{register storage class@\texttt{register} storage class|(}

Ini kata kunci untuk memberi petunjuk ke compiler bahwa variabel ini sering dipakai, dan harus dibuat secepat mungkin diaksesnya. Compiler tidak wajib menyetujuinya.

Sekarang, optimizer compiler C modern cukup pandai menentukan hal ini sendiri, jadi jarang terlihat belakangan ini. Ada beberapa argumen untuk memakainya, yang kucantumkan di bawah. Meski begitu, jujur saja, aku belum pernah memakainya di kode C mana pun yang pernah kutulis.

Tapi kalau memang harus:

\begin{tcolorbox}[enhanced jigsaw,left=3pt,breakable,colback=white!97.2549!black,boxrule=0.2mm]

\begin{Shaded}
\begin{Highlighting}[numbers=left,,]
\PreprocessorTok{\#include }\ImportTok{\textless{}stdio.h\textgreater{}}

\DataTypeTok{int}\NormalTok{ main}\OperatorTok{(}\DataTypeTok{void}\OperatorTok{)}
\OperatorTok{\{}
    \DataTypeTok{register} \DataTypeTok{int}\NormalTok{ a}\OperatorTok{;}   \CommentTok{// Buat "a" secepat mungkin dipakai.}

    \ControlFlowTok{for} \OperatorTok{(}\NormalTok{a }\OperatorTok{=} \DecValTok{0}\OperatorTok{;}\NormalTok{ a }\OperatorTok{\textless{}} \DecValTok{10}\OperatorTok{;}\NormalTok{ a}\OperatorTok{++)}
\NormalTok{        printf}\OperatorTok{(}\StringTok{"}\SpecialCharTok{\%d\textbackslash{}n}\StringTok{"}\OperatorTok{,}\NormalTok{ a}\OperatorTok{);}
\OperatorTok{\}}
\end{Highlighting}
\end{Shaded}

\end{tcolorbox}

Tapi ada harganya. Kamu tidak bisa mengambil alamat sebuah register:

\begin{tcolorbox}[enhanced jigsaw,left=3pt,breakable,colback=white!97.2549!black,boxrule=0.2mm]

\begin{Shaded}
\begin{Highlighting}[]
\DataTypeTok{register} \DataTypeTok{int}\NormalTok{ a}\OperatorTok{;}
\DataTypeTok{int} \OperatorTok{*}\NormalTok{p }\OperatorTok{=} \OperatorTok{\&}\NormalTok{a}\OperatorTok{;}    \CommentTok{// ERROR COMPILER! Tidak bisa mengambil alamat register}
\end{Highlighting}
\end{Shaded}

\end{tcolorbox}

Hal yang sama berlaku untuk bagian mana pun dari array, karena memperlakukan array \hl{\mbox{\texttt{register}}} sebagai pointer adalah undefined behavior. Kamu boleh mengambil \hl{\mbox{\texttt{sizeof()}}}-nya. Tapi itu bukan keuntungan. Artinya array \hl{\mbox{\texttt{register}}} pada dasarnya tidak berguna.

GCC dan Clang sama-sama error pada semua yang berikut (kecuali GCC, anehnya, tidak error pada \hl{\mbox{\texttt{a{[}2{]}}}}; ia hanya memberi peringatan kalau kamu memakai \hl{\mbox{\texttt{-pedantic}}}):

\begin{tcolorbox}[enhanced jigsaw,left=3pt,breakable,colback=white!97.2549!black,boxrule=0.2mm]

\begin{Shaded}
\begin{Highlighting}[]
\DataTypeTok{register} \DataTypeTok{int}\NormalTok{ a}\OperatorTok{[]} \OperatorTok{=} \OperatorTok{\{}\DecValTok{11}\OperatorTok{,} \DecValTok{22}\OperatorTok{,} \DecValTok{33}\OperatorTok{,} \DecValTok{44}\OperatorTok{,} \DecValTok{55}\OperatorTok{\};}
\DataTypeTok{int} \OperatorTok{*}\NormalTok{p }\OperatorTok{=}\NormalTok{ a}\OperatorTok{;}       \CommentTok{// Undefined behavior!}
\DataTypeTok{int}\NormalTok{ b }\OperatorTok{=} \OperatorTok{*(}\NormalTok{a }\OperatorTok{+} \DecValTok{2}\OperatorTok{);} \CommentTok{// Undefined behavior!}
\DataTypeTok{int}\NormalTok{ c }\OperatorTok{=}\NormalTok{ a}\OperatorTok{[}\DecValTok{2}\OperatorTok{];}     \CommentTok{// Undefined behavior!}
\end{Highlighting}
\end{Shaded}

\end{tcolorbox}

Fakta bahwa alamat variabel register tidak bisa diambil membebaskan compiler untuk melakukan optimasi berdasarkan asumsi itu kalau belum menemukannya sendiri. Selain itu, menambahkan \hl{\mbox{\texttt{register}}} pada variabel \hl{\mbox{\texttt{const}}} mencegah seseorang tanpa sengaja mengoper pointer-nya ke fungsi lain yang dengan sengaja mengabaikan sifat konstannya\footnote{https://gustedt.wordpress.com/2010/08/17/a-common-misconsception-the-register-keyword/}.

Sedikit latar belakang sejarah: jauh di dalam CPU ada ``variabel'' kecil khusus bernama \href{https://en.wikipedia.org/wiki/Processor_register}{\emph{register}}\footnote{https://en.wikipedia.org/wiki/Processor\_register}. Aksesnya sangat cepat dibanding RAM, jadi memakainya memberi dorongan kecepatan. Tapi mereka tidak ada di RAM, jadi tidak punya alamat memori (itulah kenapa kamu tidak bisa mengambil address-of atau pointer ke mereka).

Tapi, seperti kubilang tadi, compiler modern sangat pandai menghasilkan kode optimal, memakai register kapan pun mungkin terlepas dari apakah kamu menulis kata kunci \hl{\mbox{\texttt{register}}} atau tidak. Bukan hanya itu, spesifikasi membolehkan mereka memperlakukannya seolah-olah kamu menulis \hl{\mbox{\texttt{auto}}}, kalau mereka mau. Jadi tidak ada jaminan.

\index{register storage class@\texttt{register} storage class|)}

\subsection{\texorpdfstring{\hl{\mbox{\texttt{\_Thread\_local}}}}{\_Thread\_local}}\label{thread_local}

\index{\_Thread\_local storage class@\texttt{\_Thread\_local} storage class|(}

Saat memakai banyak thread dan ada variabel di global atau block scope \hl{\mbox{\texttt{static}}}, ini cara untuk memastikan setiap thread mendapat salinannya sendiri dari variabel itu. Ini membantu menghindari race condition dan thread yang saling menginjak kaki.

Kalau berada di block scope, kamu harus memakainya bersama \hl{\mbox{\texttt{extern}}} atau \hl{\mbox{\texttt{static}}}.

Selain itu, kalau kamu menyertakan \hl{\mbox{\texttt{\textless{}threads.h\textgreater{}}}}, kamu bisa memakai \hl{\mbox{\texttt{thread\_local}}} yang jauh lebih enak dilihat sebagai alias dari \hl{\mbox{\texttt{\_Thread\_local}}} yang lebih jelek.

Informasi lebih lanjut ada di \hyperref[thread-local]{bagian Threads}.

\index{\_Thread\_local storage class@\texttt{\_Thread\_local} storage class|)} \index{Storage-Class Specifiers|)}

\chapter{Proyek Multifile}\label{proyek-multifile}

\index{Multifile projects|(}

Sejauh ini kita melihat program-program mainan yang sebagian besar muat di satu file. Tapi program C yang kompleks terdiri dari banyak file yang semuanya dikompilasi dan di-link menjadi satu executable.

Di bab ini kita akan melihat beberapa pola dan praktik umum untuk menyusun proyek yang lebih besar.

\section{Include dan Function Prototype}\label{includes-func-protos}

\index{Multifile projects!includes|(} \index{Multifile projects!function prototypes|(}

Situasi yang sangat umum: sebagian fungsimu didefinisikan di satu file, dan kamu ingin memanggilnya dari file lain.

Sebenarnya ini langsung berfungsi dengan sebuah peringatan\ldots{} kita coba dulu, lalu kita lihat cara yang benar untuk memperbaiki peringatannya.

Untuk mengompilasinya, kamu perlu menyebutkan semua file sumber di command line:

\begin{tcolorbox}[enhanced jigsaw,left=3pt,breakable,colback=white!97.2549!black,boxrule=0.2mm]

\begin{Shaded}
\begin{Highlighting}[]
\CommentTok{\# file output   file sumber}
\CommentTok{\#     v            v}
\CommentTok{\#   |{-}{-}{-}{-}| |{-}{-}{-}{-}{-}{-}{-}{-}{-}|}
\FunctionTok{gcc} \AttributeTok{{-}o}\NormalTok{ foo foo.c bar.c}
\end{Highlighting}
\end{Shaded}

\end{tcolorbox}

Pada contoh itu, \hl{\mbox{\texttt{foo.c}}} dan \hl{\mbox{\texttt{bar.c}}} dibangun menjadi executable bernama \hl{\mbox{\texttt{foo}}}.

Pada contoh-contoh ini, nama file kita taruh sebagai komentar pertama di sumbernya. Kita lihat file sumber \hl{\mbox{\texttt{bar.c}}}:

\begin{tcolorbox}[enhanced jigsaw,left=3pt,breakable,colback=white!97.2549!black,boxrule=0.2mm]

\begin{Shaded}
\begin{Highlighting}[numbers=left,,]
\CommentTok{// File bar.c}

\DataTypeTok{int}\NormalTok{ add}\OperatorTok{(}\DataTypeTok{int}\NormalTok{ x}\OperatorTok{,} \DataTypeTok{int}\NormalTok{ y}\OperatorTok{)}
\OperatorTok{\{}
    \ControlFlowTok{return}\NormalTok{ x }\OperatorTok{+}\NormalTok{ y}\OperatorTok{;}
\OperatorTok{\}}
\end{Highlighting}
\end{Shaded}

\end{tcolorbox}

Dan file \hl{\mbox{\texttt{foo.c}}} yang berisi main:

\begin{tcolorbox}[enhanced jigsaw,left=3pt,breakable,colback=white!97.2549!black,boxrule=0.2mm]

\begin{Shaded}
\begin{Highlighting}[numbers=left,,]
\CommentTok{// File foo.c}

\PreprocessorTok{\#include }\ImportTok{\textless{}stdio.h\textgreater{}}

\DataTypeTok{int}\NormalTok{ main}\OperatorTok{(}\DataTypeTok{void}\OperatorTok{)}
\OperatorTok{\{}
\NormalTok{    printf}\OperatorTok{(}\StringTok{"}\SpecialCharTok{\%d\textbackslash{}n}\StringTok{"}\OperatorTok{,}\NormalTok{ add}\OperatorTok{(}\DecValTok{2}\OperatorTok{,} \DecValTok{3}\OperatorTok{));}  \CommentTok{// 5!}
\OperatorTok{\}}
\end{Highlighting}
\end{Shaded}

\end{tcolorbox}

Lihat bagaimana dari \hl{\mbox{\texttt{main()}}} kita memanggil \hl{\mbox{\texttt{add()}}}---tapi \hl{\mbox{\texttt{add()}}} ada di file sumber yang sama sekali berbeda! Ia ada di \hl{\mbox{\texttt{bar.c}}}, sedangkan pemanggilannya ada di \hl{\mbox{\texttt{foo.c}}}!

Kalau kita membangunnya dengan:

\begin{tcolorbox}[enhanced jigsaw,left=3pt,breakable,colback=white!97.2549!black,boxrule=0.2mm]

\begin{Shaded}
\begin{Highlighting}[]
\FunctionTok{gcc} \AttributeTok{{-}o}\NormalTok{ foo foo.c bar.c}
\end{Highlighting}
\end{Shaded}

\end{tcolorbox}

kita mendapat error ini:

\begin{tcolorbox}[enhanced jigsaw,left=3pt,breakable,colback=white!97.2549!black,boxrule=0.2mm]

\begin{Shaded}
\begin{Highlighting}[]
\NormalTok{error: implicit declaration of function \textquotesingle{}add\textquotesingle{} is invalid in C99}
\end{Highlighting}
\end{Shaded}

\end{tcolorbox}

(Atau mungkin kamu mendapat peringatan, yang tidak boleh kamu abaikan. Jangan pernah mengabaikan peringatan di C; tangani semuanya.)

Kalau kamu ingat dari \hyperref[prototypes]{bagian tentang prototype}, implicit declaration dilarang di C modern dan tidak ada alasan sah untuk memasukkannya ke kode baru. Kita harus memperbaikinya.

Maksud \hl{\mbox{\texttt{implicit\ declaration}}} adalah kita memakai sebuah fungsi, yaitu \hl{\mbox{\texttt{add()}}} dalam kasus ini, tanpa memberi tahu C apa pun tentangnya sebelumnya. C ingin tahu apa yang dikembalikannya, tipe argumen apa yang diterimanya, dan hal-hal semacam itu.

Kita sudah melihat cara memperbaikinya dengan \emph{function prototype}. Memang, kalau kita menambahkannya ke \hl{\mbox{\texttt{foo.c}}} sebelum pemanggilan, semuanya berjalan baik:

\begin{tcolorbox}[enhanced jigsaw,left=3pt,breakable,colback=white!97.2549!black,boxrule=0.2mm]

\begin{Shaded}
\begin{Highlighting}[numbers=left,,]
\CommentTok{// File foo.c}

\PreprocessorTok{\#include }\ImportTok{\textless{}stdio.h\textgreater{}}

\DataTypeTok{int}\NormalTok{ add}\OperatorTok{(}\DataTypeTok{int}\OperatorTok{,} \DataTypeTok{int}\OperatorTok{);}  \CommentTok{// Tambahkan prototype{-}nya}

\DataTypeTok{int}\NormalTok{ main}\OperatorTok{(}\DataTypeTok{void}\OperatorTok{)}
\OperatorTok{\{}
\NormalTok{    printf}\OperatorTok{(}\StringTok{"}\SpecialCharTok{\%d\textbackslash{}n}\StringTok{"}\OperatorTok{,}\NormalTok{ add}\OperatorTok{(}\DecValTok{2}\OperatorTok{,} \DecValTok{3}\OperatorTok{));}  \CommentTok{// 5!}
\OperatorTok{\}}
\end{Highlighting}
\end{Shaded}

\end{tcolorbox}

Tidak ada error lagi!

Tapi itu merepotkan---harus mengetik prototype setiap kali ingin memakai sebuah fungsi. Maksudnya, kita memakai \hl{\mbox{\texttt{printf()}}} di situ dan tidak perlu mengetik prototype; kok bisa?

\emph{Sebenarnya kita memang sudah menyertakan prototype \hl{\mbox{\texttt{printf()}}}}! Ada di file \hl{\mbox{\texttt{stdio.h}}}! Dan kita meng-include-nya dengan \hl{\mbox{\texttt{\#include}}}!

Bisakah kita melakukan hal yang sama untuk fungsi \hl{\mbox{\texttt{add()}}}? Membuat prototype-nya dan menaruhnya di header file?

Tentu bisa!

Secara konvensi, header file di C berekstensi \hl{\mbox{\texttt{.h}}}. Dan sering, tapi tidak selalu, namanya sama dengan file \hl{\mbox{\texttt{.c}}} pasangannya. Jadi mari kita buat file \hl{\mbox{\texttt{bar.h}}} untuk file \hl{\mbox{\texttt{bar.c}}}, dan kita taruh prototype-nya di situ:

\begin{tcolorbox}[enhanced jigsaw,left=3pt,breakable,colback=white!97.2549!black,boxrule=0.2mm]

\begin{Shaded}
\begin{Highlighting}[numbers=left,,]
\CommentTok{// File bar.h}

\DataTypeTok{int}\NormalTok{ add}\OperatorTok{(}\DataTypeTok{int}\OperatorTok{,} \DataTypeTok{int}\OperatorTok{);}
\end{Highlighting}
\end{Shaded}

\end{tcolorbox}

Sekarang kita ubah \hl{\mbox{\texttt{foo.c}}} agar meng-include file itu. Dengan asumsi file itu ada di direktori yang sama, kita meng-include-nya di dalam tanda kutip ganda (bukan kurung siku):

\begin{tcolorbox}[enhanced jigsaw,left=3pt,breakable,colback=white!97.2549!black,boxrule=0.2mm]

\begin{Shaded}
\begin{Highlighting}[numbers=left,,]
\CommentTok{// File foo.c}

\PreprocessorTok{\#include }\ImportTok{\textless{}stdio.h\textgreater{}}

\PreprocessorTok{\#include }\ImportTok{"bar.h"}\PreprocessorTok{  }\CommentTok{// Include dari direktori saat ini}

\DataTypeTok{int}\NormalTok{ main}\OperatorTok{(}\DataTypeTok{void}\OperatorTok{)}
\OperatorTok{\{}
\NormalTok{    printf}\OperatorTok{(}\StringTok{"}\SpecialCharTok{\%d\textbackslash{}n}\StringTok{"}\OperatorTok{,}\NormalTok{ add}\OperatorTok{(}\DecValTok{2}\OperatorTok{,} \DecValTok{3}\OperatorTok{));}  \CommentTok{// 5!}
\OperatorTok{\}}
\end{Highlighting}
\end{Shaded}

\end{tcolorbox}

Perhatikan bahwa prototype sudah tidak ada di \hl{\mbox{\texttt{foo.c}}}---kita meng-include-nya dari \hl{\mbox{\texttt{bar.h}}}. Sekarang file \emph{mana pun} yang butuh fungsionalitas \hl{\mbox{\texttt{add()}}} cukup \hl{\mbox{\texttt{\#include\ "bar.h"}}}, dan kamu tidak perlu repot mengetik function prototype.

Seperti yang mungkin sudah kamu duga, \hl{\mbox{\texttt{\#include}}} secara harfiah menyisipkan file yang disebut \emph{tepat di situ} di kode sumbermu, seolah-olah kamu mengetiknya sendiri.

Lalu build dan jalankan:

\begin{tcolorbox}[enhanced jigsaw,left=3pt,breakable,colback=white!97.2549!black,boxrule=0.2mm]

\begin{Shaded}
\begin{Highlighting}[]
\ExtensionTok{./foo}
\ExtensionTok{5}
\end{Highlighting}
\end{Shaded}

\end{tcolorbox}

Memang, kita mendapat hasil \(2+3\)! Hore!

Tapi jangan buka minuman pilihanmu dulu. Kita hampir sampai! Tinggal satu boilerplate lagi yang harus ditambahkan.

\index{Multifile projects!function prototypes|)}

\section{Menangani Include yang Berulang}\label{menangani-include-yang-berulang}

Bukan hal aneh kalau sebuah header file itu sendiri meng-\hl{\mbox{\texttt{\#include}}} header lain yang dibutuhkan untuk fungsionalitas file C pasangannya. Maksudnya, kenapa tidak?

Dan bisa jadi sebuah header di-\hl{\mbox{\texttt{\#include}}} berkali-kali dari tempat yang berbeda. Mungkin itu tidak masalah, tapi mungkin juga menyebabkan error compiler. Dan kita tidak bisa mengendalikan berapa banyak tempat yang meng-\hl{\mbox{\texttt{\#include}}}-nya!

Bahkan lebih buruk, kita bisa terjebak situasi gila ketika header \hl{\mbox{\texttt{a.h}}} meng-include header \hl{\mbox{\texttt{b.h}}}, dan \hl{\mbox{\texttt{b.h}}} meng-include \hl{\mbox{\texttt{a.h}}}! Siklus \hl{\mbox{\texttt{\#include}}} tanpa akhir!

Mencoba membangun yang seperti itu menghasilkan error:

\begin{tcolorbox}[enhanced jigsaw,left=3pt,breakable,colback=white!97.2549!black,boxrule=0.2mm]

\begin{Shaded}
\begin{Highlighting}[]
\NormalTok{error: \#include nested depth 200 exceeds maximum of 200}
\end{Highlighting}
\end{Shaded}

\end{tcolorbox}

Mungkin siklusnya akan terselesaikan pada langkah ke-201\ldots{}

Kita perlu mengaturnya supaya kalau sebuah file sudah di-include satu kali, \hl{\mbox{\texttt{\#include}}} berikutnya untuk file itu diabaikan.

\textbf{Langkah yang akan kita ambil ini sangat umum, jadi lakukan saja secara otomatis setiap kali membuat header file!}

Cara umumnya adalah dengan variabel preprocessor yang kita set saat pertama kali file di-\hl{\mbox{\texttt{\#include}}}. Lalu pada \hl{\mbox{\texttt{\#include}}} berikutnya, kita periksa dulu bahwa variabel itu belum didefinisikan.

Untuk nama variabel itu, sangat umum mengambil nama header file, misalnya \hl{\mbox{\texttt{bar.h}}}, mengubahnya menjadi huruf besar, dan mengganti titik dengan underscore: \hl{\mbox{\texttt{BAR\_H}}}.

Jadi taruh pengecekan di bagian paling, paling atas file untuk melihat apakah file sudah pernah di-include, dan secara efektif mengomentari seluruh isinya kalau sudah.

(Jangan memakai underscore di awal (karena underscore di awal yang diikuti huruf kapital itu dicadangkan) atau dua underscore di awal (karena itu juga dicadangkan).)

\begin{tcolorbox}[enhanced jigsaw,left=3pt,breakable,colback=white!97.2549!black,boxrule=0.2mm]

\begin{Shaded}
\begin{Highlighting}[numbers=left,,]
\PreprocessorTok{\#ifndef BAR\_H   }\CommentTok{// Jika BAR\_H belum didefinisikan...}
\PreprocessorTok{\#define BAR\_H   }\CommentTok{// Definisikan (tanpa nilai tertentu)}

\CommentTok{// File bar.h}

\DataTypeTok{int}\NormalTok{ add}\OperatorTok{(}\DataTypeTok{int}\OperatorTok{,} \DataTypeTok{int}\OperatorTok{);}

\PreprocessorTok{\#endif          }\CommentTok{// Akhir dari \#ifndef BAR\_H}
\end{Highlighting}
\end{Shaded}

\end{tcolorbox}

Ini secara efektif membuat header file hanya di-include satu kali, berapa pun banyaknya tempat yang mencoba meng-\hl{\mbox{\texttt{\#include}}}-nya.

\index{Multifile projects!includes|)}

\section{\texorpdfstring{\hl{\mbox{\texttt{static}}} dan \hl{\mbox{\texttt{extern}}}}{static dan extern}}\label{static-dan-extern}

\index{static storage class@\texttt{static} storage class|(} \index{extern storage class@\texttt{extern} storage class|(} \index{Multifile projects!static storage class@\texttt{static} storage class|(} \index{Multifile projects!extern storage class@\texttt{extern} storage class|(}

Kamu bisa merujuk ke object di file lain dengan \hl{\mbox{\texttt{extern}}}.

Kamu bisa memastikan variabel dan fungsi di file scope \emph{tidak} terlihat dari file sumber lain (bahkan kalau \hl{\mbox{\texttt{extern}}} dipakai) dengan kata kunci \hl{\mbox{\texttt{static}}}.

Untuk info lebih lanjut, lihat bagian di buku ini tentang storage-class specifier \hyperref[static]{\hl{\mbox{\texttt{static}}}} dan \hyperref[extern]{\hl{\mbox{\texttt{extern}}}}.

\index{static storage class@\texttt{static} storage class|)} \index{extern storage class@\texttt{extern} storage class|)} \index{Multifile projects!static storage class@\texttt{static} storage class|)} \index{Multifile projects!extern storage class@\texttt{extern} storage class|)}

\section{Mengompilasi dengan Object File}\label{mengompilasi-dengan-object-file}

\index{Object files|(}

Ini bukan bagian dari spesifikasi, tapi 99,999\% umum di dunia C.

Kamu bisa mengompilasi file C menjadi representasi perantara yang disebut \emph{object file}. Isinya machine code (yaitu 1 dan 0 dari instruksi sebenarnya) yang belum dibangun menjadi executable.

Object file di Windows berekstensi \hl{\mbox{\texttt{.OBJ}}}; di sistem mirip Unix, \hl{\mbox{\texttt{.o}}}.

\index{gcc compiler@\texttt{gcc} compiler|(}

Di gcc, kita bisa membangunnya seperti ini, dengan flag \hl{\mbox{\texttt{-c}}} (hanya kompilasi!):

\begin{tcolorbox}[enhanced jigsaw,left=3pt,breakable,colback=white!97.2549!black,boxrule=0.2mm]

\begin{Shaded}
\begin{Highlighting}[]
\FunctionTok{gcc} \AttributeTok{{-}c}\NormalTok{ foo.c     }\CommentTok{\# menghasilkan foo.o}
\FunctionTok{gcc} \AttributeTok{{-}c}\NormalTok{ bar.c     }\CommentTok{\# menghasilkan bar.o}
\end{Highlighting}
\end{Shaded}

\end{tcolorbox}

Lalu kita bisa me-\emph{link} keduanya menjadi satu executable:

\begin{tcolorbox}[enhanced jigsaw,left=3pt,breakable,colback=white!97.2549!black,boxrule=0.2mm]

\begin{Shaded}
\begin{Highlighting}[]
\FunctionTok{gcc} \AttributeTok{{-}o}\NormalTok{ foo foo.o bar.o}
\end{Highlighting}
\end{Shaded}

\end{tcolorbox}

\emph{Voila}, kita sudah menghasilkan executable \hl{\mbox{\texttt{foo}}} dari dua object file.

Tapi kamu mungkin berpikir, buat apa repot-repot? Bukankah kita bisa langsung:

\begin{tcolorbox}[enhanced jigsaw,left=3pt,breakable,colback=white!97.2549!black,boxrule=0.2mm]

\begin{Shaded}
\begin{Highlighting}[]
\FunctionTok{gcc} \AttributeTok{{-}o}\NormalTok{ foo foo.c bar.c}
\end{Highlighting}
\end{Shaded}

\end{tcolorbox}

dan sekali tepuk membunuh dua \href{https://en.wikipedia.org/wiki/Boids}{boid}\footnote{https://en.wikipedia.org/wiki/Boids}?

\index{gcc compiler@\texttt{gcc} compiler|)}

Untuk program kecil, itu tidak masalah. Aku sendiri sering melakukannya.

Tapi untuk program yang lebih besar, kita bisa memanfaatkan fakta bahwa kompilasi dari sumber ke object file relatif lambat, sedangkan me-link banyak object file relatif cepat.

Ini benar-benar terlihat pada utilitas \hl{\mbox{\texttt{make}}} yang hanya membangun ulang sumber yang lebih baru daripada outputnya.

Misalkan kamu punya seribu file C. Kamu bisa mengompilasi semuanya menjadi object file di awal (lambat) lalu menggabungkan semua object file itu menjadi executable (cepat).

Sekarang katakanlah kamu hanya mengubah satu dari file sumber C itu---ini keajaibannya: \emph{kamu hanya perlu membangun ulang satu object file untuk file sumber itu}! Lalu membangun ulang executable-nya (cepat). Semua file C lainnya tidak perlu disentuh.

Dengan kata lain, dengan hanya membangun ulang object file yang perlu, kita memangkas waktu kompilasi secara drastis. (Kecuali tentu saja kalau kamu melakukan build ``bersih'', yang mengharuskan semua object file dibuat.)

\index{Object files|)} \index{Multifile projects|)}

\chapter{Lingkungan di Luar Program}\label{lingkungan-di-luar-program}

Ketika kamu menjalankan sebuah program, sebenarnya kamu sedang berbicara dengan shell, berkata, ``Hei, tolong jalankan ini.'' Dan shell menjawab, ``Siap,'' lalu memberi tahu sistem operasi, ``Hei, bisa tolong buatkan proses baru dan jalankan ini?'' Dan kalau semuanya lancar, OS menurut dan programmu berjalan.

Tapi ada dunia luas di luar programmu, di dalam shell, yang bisa berinteraksi dengan C. Beberapa di antaranya kita lihat di bab ini.

\section{Argumen Command Line}\label{argumen-command-line}

\index{Command line arguments|(}

Banyak utilitas command line menerima \emph{argumen command line}. Misalnya, kalau kita ingin melihat semua file berakhiran \hl{\mbox{\texttt{.txt}}}, kita bisa mengetik sesuatu seperti ini di sistem mirip Unix:

\begin{tcolorbox}[enhanced jigsaw,left=3pt,breakable,colback=white!97.2549!black,boxrule=0.2mm]

\begin{Shaded}
\begin{Highlighting}[]
\FunctionTok{ls} \PreprocessorTok{*}\NormalTok{.txt}
\end{Highlighting}
\end{Shaded}

\end{tcolorbox}

(atau \hl{\mbox{\texttt{dir}}} alih-alih \hl{\mbox{\texttt{ls}}} di sistem Windows).

Dalam kasus ini, perintahnya adalah \hl{\mbox{\texttt{ls}}}, sedangkan argumennya adalah semua file yang berakhiran \hl{\mbox{\texttt{.txt}}}\footnote{Secara historis, program MS-DOS dan Windows melakukan ini secara berbeda dari Unix. Di Unix, shell akan \emph{memperluas} wildcard menjadi semua file yang cocok sebelum programmu melihatnya, sedangkan varian Microsoft mengoper ekspresi wildcard itu ke program untuk ditangani sendiri. Bagaimanapun juga, ada argumen yang dioper ke program.}.

Jadi bagaimana kita bisa melihat apa yang dioper ke program dari command line?

Katakanlah kita punya program bernama \hl{\mbox{\texttt{add}}} yang menjumlahkan semua angka yang dioper di command line dan mencetak hasilnya:

\begin{tcolorbox}[enhanced jigsaw,left=3pt,breakable,colback=white!97.2549!black,boxrule=0.2mm]

\begin{Shaded}
\begin{Highlighting}[]
\ExtensionTok{./add}\NormalTok{ 10 30 5}
\ExtensionTok{45}
\end{Highlighting}
\end{Shaded}

\end{tcolorbox}

Dijamin ini bisa membayar tagihan!

Tapi serius, ini alat yang bagus untuk melihat cara mengambil argumen dari command line dan memecahnya.

Pertama, mari kita lihat dulu cara mendapatkannya. Untuk itu, kita butuh \hl{\mbox{\texttt{main()}}} yang baru!

Berikut program yang mencetak semua argumen command line. Misalnya, kalau executable-nya kita beri nama \hl{\mbox{\texttt{foo}}}, kita bisa menjalankannya seperti ini:

\begin{tcolorbox}[enhanced jigsaw,left=3pt,breakable,colback=white!97.2549!black,boxrule=0.2mm]

\begin{Shaded}
\begin{Highlighting}[]
\ExtensionTok{./foo}\NormalTok{ i like turtles}
\end{Highlighting}
\end{Shaded}

\end{tcolorbox}

dan kita akan melihat output ini:

\begin{tcolorbox}[enhanced jigsaw,left=3pt,breakable,colback=white!97.2549!black,boxrule=0.2mm]

\begin{Shaded}
\begin{Highlighting}[]
\NormalTok{arg 0: ./foo}
\NormalTok{arg 1: i}
\NormalTok{arg 2: like}
\NormalTok{arg 3: turtles}
\end{Highlighting}
\end{Shaded}

\end{tcolorbox}

Agak aneh, karena argumen ke-nol adalah nama executable itu sendiri. Tapi itu hanya soal membiasakan diri. Argumen sebenarnya menyusul langsung setelahnya.

Sumber:

\index{argc parameter@\texttt{argc} parameter|(} \index{argv parameter@\texttt{argv} parameter|(} \index{main() function@\texttt{main()} function!command line options|(}

\begin{tcolorbox}[enhanced jigsaw,left=3pt,breakable,colback=white!97.2549!black,boxrule=0.2mm]

\begin{Shaded}
\begin{Highlighting}[numbers=left,,]
\PreprocessorTok{\#include }\ImportTok{\textless{}stdio.h\textgreater{}}

\DataTypeTok{int}\NormalTok{ main}\OperatorTok{(}\DataTypeTok{int}\NormalTok{ argc}\OperatorTok{,} \DataTypeTok{char} \OperatorTok{*}\NormalTok{argv}\OperatorTok{[])}
\OperatorTok{\{}
    \ControlFlowTok{for} \OperatorTok{(}\DataTypeTok{int}\NormalTok{ i }\OperatorTok{=} \DecValTok{0}\OperatorTok{;}\NormalTok{ i }\OperatorTok{\textless{}}\NormalTok{ argc}\OperatorTok{;}\NormalTok{ i}\OperatorTok{++)} \OperatorTok{\{}
\NormalTok{        printf}\OperatorTok{(}\StringTok{"arg }\SpecialCharTok{\%d}\StringTok{: }\SpecialCharTok{\%s\textbackslash{}n}\StringTok{"}\OperatorTok{,}\NormalTok{ i}\OperatorTok{,}\NormalTok{ argv}\OperatorTok{[}\NormalTok{i}\OperatorTok{]);}
    \OperatorTok{\}}
\OperatorTok{\}}
\end{Highlighting}
\end{Shaded}

\end{tcolorbox}

Wah! Ada apa dengan signature fungsi \hl{\mbox{\texttt{main()}}} itu? Apa itu \hl{\mbox{\texttt{argc}}} dan \hl{\mbox{\texttt{argv}}}\footnote{Karena keduanya hanya nama parameter biasa, kamu sebenarnya tidak harus menamainya \hl{\mbox{\texttt{argc}}} dan \hl{\mbox{\texttt{argv}}}. Tapi nama-nama itu begitu idiomatis, sehingga kalau kamu kreatif, programmer C lain pasti akan menatapmu dengan curiga!} (dibaca \emph{arg-cee} dan \emph{arg-vee})?

Mulai dari yang mudah: \hl{\mbox{\texttt{argc}}}. Ia singkatan dari \emph{argument count} (jumlah argumen), termasuk nama program itu sendiri. Kalau kamu menganggap semua argumen sebagai array berisi string, yang memang begitulah adanya, maka \hl{\mbox{\texttt{argc}}} bisa dianggap sebagai panjang array itu, yang memang begitulah adanya.

Jadi di loop itu kita menelusuri semua \hl{\mbox{\texttt{argv}}} dan mencetaknya satu per satu, jadi untuk input tertentu:

\begin{tcolorbox}[enhanced jigsaw,left=3pt,breakable,colback=white!97.2549!black,boxrule=0.2mm]

\begin{Shaded}
\begin{Highlighting}[]
\ExtensionTok{./foo}\NormalTok{ i like turtles}
\end{Highlighting}
\end{Shaded}

\end{tcolorbox}

kita mendapat output yang sesuai:

\begin{tcolorbox}[enhanced jigsaw,left=3pt,breakable,colback=white!97.2549!black,boxrule=0.2mm]

\begin{Shaded}
\begin{Highlighting}[]
\NormalTok{arg 0: ./foo}
\NormalTok{arg 1: i}
\NormalTok{arg 2: like}
\NormalTok{arg 3: turtles}
\end{Highlighting}
\end{Shaded}

\end{tcolorbox}

Dengan itu, kita siap membuat program penjumlah.

Rencana kita:

\begin{itemize}
\tightlist
\item
  Lihat semua argumen command line (setelah \hl{\mbox{\texttt{argv{[}0{]}}}}, nama program)
\item
  Konversi menjadi integer
\item
  Tambahkan ke total yang berjalan
\item
  Cetak hasilnya
\end{itemize}

Mari kita kerjakan!

\begin{tcolorbox}[enhanced jigsaw,left=3pt,breakable,colback=white!97.2549!black,boxrule=0.2mm]

\begin{Shaded}
\begin{Highlighting}[numbers=left,,]
\PreprocessorTok{\#include }\ImportTok{\textless{}stdio.h\textgreater{}}
\PreprocessorTok{\#include }\ImportTok{\textless{}stdlib.h\textgreater{}}

\DataTypeTok{int}\NormalTok{ main}\OperatorTok{(}\DataTypeTok{int}\NormalTok{ argc}\OperatorTok{,} \DataTypeTok{char} \OperatorTok{**}\NormalTok{argv}\OperatorTok{)}
\OperatorTok{\{}
    \DataTypeTok{int}\NormalTok{ total }\OperatorTok{=} \DecValTok{0}\OperatorTok{;}

    \ControlFlowTok{for} \OperatorTok{(}\DataTypeTok{int}\NormalTok{ i }\OperatorTok{=} \DecValTok{1}\OperatorTok{;}\NormalTok{ i }\OperatorTok{\textless{}}\NormalTok{ argc}\OperatorTok{;}\NormalTok{ i}\OperatorTok{++)} \OperatorTok{\{}  \CommentTok{// Mulai dari 1, argumen pertama}
        \DataTypeTok{int}\NormalTok{ value }\OperatorTok{=}\NormalTok{ atoi}\OperatorTok{(}\NormalTok{argv}\OperatorTok{[}\NormalTok{i}\OperatorTok{]);}    \CommentTok{// strtol() penanganan error{-}nya lebih baik}

\NormalTok{        total }\OperatorTok{+=}\NormalTok{ value}\OperatorTok{;}
    \OperatorTok{\}}

\NormalTok{    printf}\OperatorTok{(}\StringTok{"}\SpecialCharTok{\%d\textbackslash{}n}\StringTok{"}\OperatorTok{,}\NormalTok{ total}\OperatorTok{);}
\OperatorTok{\}}
\end{Highlighting}
\end{Shaded}

\end{tcolorbox}

\index{main() function@\texttt{main()} function!command line options|)} \index{argc parameter@\texttt{argc} parameter|)}

Contoh eksekusi:

\begin{tcolorbox}[enhanced jigsaw,left=3pt,breakable,colback=white!97.2549!black,boxrule=0.2mm]

\begin{Shaded}
\begin{Highlighting}[]
\FunctionTok{$}\NormalTok{ ./add}
\ExtensionTok{0}
\FunctionTok{$}\NormalTok{ ./add 1}
\ExtensionTok{1}
\FunctionTok{$}\NormalTok{ ./add 1 2}
\ExtensionTok{3}
\FunctionTok{$}\NormalTok{ ./add 1 2 3}
\ExtensionTok{6}
\FunctionTok{$}\NormalTok{ ./add 1 2 3 4}
\ExtensionTok{10}
\end{Highlighting}
\end{Shaded}

\end{tcolorbox}

Tentu saja, program ini mungkin muntah kalau yang dioper bukan integer, tapi membuatnya tahan terhadap hal itu kita serahkan sebagai latihan untuk pembaca.

\subsection{\texorpdfstring{\hl{\mbox{\texttt{argv}}} Terakhir adalah \hl{\mbox{\texttt{NULL}}}}{argv Terakhir adalah NULL}}\label{argv-terakhir-adalah-null}

Satu trivia menyenangkan tentang \hl{\mbox{\texttt{argv}}}: setelah string terakhir ada pointer ke \hl{\mbox{\texttt{NULL}}}.

Artinya:

\begin{tcolorbox}[enhanced jigsaw,left=3pt,breakable,colback=white!97.2549!black,boxrule=0.2mm]

\begin{Shaded}
\begin{Highlighting}[]
\NormalTok{argv}\OperatorTok{[}\NormalTok{argc}\OperatorTok{]} \OperatorTok{==}\NormalTok{ NULL}
\end{Highlighting}
\end{Shaded}

\end{tcolorbox}

selalu bernilai true!

Ini mungkin tampak tidak berguna, tapi ternyata berguna di beberapa tempat; salah satunya kita lihat sekarang.

\subsection{\texorpdfstring{Alternatifnya: \hl{\mbox{\texttt{char\ **argv}}}}{Alternatifnya: char **argv}}\label{alternatifnya-char-argv}

Ingat bahwa saat memanggil fungsi, C tidak membedakan notasi array dan notasi pointer pada signature fungsi.

Artinya, keduanya sama:

\begin{tcolorbox}[enhanced jigsaw,left=3pt,breakable,colback=white!97.2549!black,boxrule=0.2mm]

\begin{Shaded}
\begin{Highlighting}[]
\DataTypeTok{void}\NormalTok{ foo}\OperatorTok{(}\DataTypeTok{char}\NormalTok{ a}\OperatorTok{[])}
\DataTypeTok{void}\NormalTok{ foo}\OperatorTok{(}\DataTypeTok{char} \OperatorTok{*}\NormalTok{a}\OperatorTok{)}
\end{Highlighting}
\end{Shaded}

\end{tcolorbox}

Nah, selama ini kita nyaman menganggap \hl{\mbox{\texttt{argv}}} sebagai array berisi string, yaitu array berisi \hl{\mbox{\texttt{char*}}}, jadi ini masuk akal:

\begin{tcolorbox}[enhanced jigsaw,left=3pt,breakable,colback=white!97.2549!black,boxrule=0.2mm]

\begin{Shaded}
\begin{Highlighting}[]
\DataTypeTok{int}\NormalTok{ main}\OperatorTok{(}\DataTypeTok{int}\NormalTok{ argc}\OperatorTok{,} \DataTypeTok{char} \OperatorTok{*}\NormalTok{argv}\OperatorTok{[])}
\end{Highlighting}
\end{Shaded}

\end{tcolorbox}

tapi karena kesetaraan itu, kamu juga bisa menulis:

\begin{tcolorbox}[enhanced jigsaw,left=3pt,breakable,colback=white!97.2549!black,boxrule=0.2mm]

\begin{Shaded}
\begin{Highlighting}[]
\DataTypeTok{int}\NormalTok{ main}\OperatorTok{(}\DataTypeTok{int}\NormalTok{ argc}\OperatorTok{,} \DataTypeTok{char} \OperatorTok{**}\NormalTok{argv}\OperatorTok{)}
\end{Highlighting}
\end{Shaded}

\end{tcolorbox}

Ya, itu memang pointer ke pointer! Kalau lebih mudah, anggap saja sebagai pointer ke string. Tapi sebenarnya ia adalah pointer ke nilai yang menunjuk ke \hl{\mbox{\texttt{char}}}.

Ingat juga bahwa dua hal ini setara:

\begin{tcolorbox}[enhanced jigsaw,left=3pt,breakable,colback=white!97.2549!black,boxrule=0.2mm]

\begin{Shaded}
\begin{Highlighting}[]
\NormalTok{argv}\OperatorTok{[}\NormalTok{i}\OperatorTok{]}
\OperatorTok{*(}\NormalTok{argv }\OperatorTok{+}\NormalTok{ i}\OperatorTok{)}
\end{Highlighting}
\end{Shaded}

\end{tcolorbox}

yang berarti kamu bisa melakukan aritmetika pointer pada \hl{\mbox{\texttt{argv}}}.

Jadi cara alternatif untuk memproses argumen command line adalah dengan menelusuri array \hl{\mbox{\texttt{argv}}} sambil menaikkan pointer sampai bertemu \hl{\mbox{\texttt{NULL}}} di ujung.

Mari kita ubah penjumlah kita agar bekerja seperti itu:

\begin{tcolorbox}[enhanced jigsaw,left=3pt,breakable,colback=white!97.2549!black,boxrule=0.2mm]

\begin{Shaded}
\begin{Highlighting}[numbers=left,,]
\PreprocessorTok{\#include }\ImportTok{\textless{}stdio.h\textgreater{}}
\PreprocessorTok{\#include }\ImportTok{\textless{}stdlib.h\textgreater{}}

\DataTypeTok{int}\NormalTok{ main}\OperatorTok{(}\DataTypeTok{int}\NormalTok{ argc}\OperatorTok{,} \DataTypeTok{char} \OperatorTok{**}\NormalTok{argv}\OperatorTok{)}
\OperatorTok{\{}
    \DataTypeTok{int}\NormalTok{ total }\OperatorTok{=} \DecValTok{0}\OperatorTok{;}
    
    \CommentTok{// Trik manis supaya compiler berhenti memperingatkan soal variabel}
    \CommentTok{// argc yang tidak terpakai:}
    \OperatorTok{(}\DataTypeTok{void}\OperatorTok{)}\NormalTok{argc}\OperatorTok{;}

    \ControlFlowTok{for} \OperatorTok{(}\DataTypeTok{char} \OperatorTok{**}\NormalTok{p }\OperatorTok{=}\NormalTok{ argv }\OperatorTok{+} \DecValTok{1}\OperatorTok{;} \OperatorTok{*}\NormalTok{p }\OperatorTok{!=}\NormalTok{ NULL}\OperatorTok{;}\NormalTok{ p}\OperatorTok{++)} \OperatorTok{\{}
        \DataTypeTok{int}\NormalTok{ value }\OperatorTok{=}\NormalTok{ atoi}\OperatorTok{(*}\NormalTok{p}\OperatorTok{);}  \CommentTok{// Pakai strtol() untuk penanganan error yang lebih baik}

\NormalTok{        total }\OperatorTok{+=}\NormalTok{ value}\OperatorTok{;}
    \OperatorTok{\}}

\NormalTok{    printf}\OperatorTok{(}\StringTok{"}\SpecialCharTok{\%d\textbackslash{}n}\StringTok{"}\OperatorTok{,}\NormalTok{ total}\OperatorTok{);}
\OperatorTok{\}}
\end{Highlighting}
\end{Shaded}

\end{tcolorbox}

Aku sendiri memakai notasi array untuk mengakses \hl{\mbox{\texttt{argv}}}, tapi gaya ini juga sering kulihat beredar.

\subsection{Fakta Menarik}\label{fakta-menarik}

\index{argc parameter@\texttt{argc} parameter|(}

Beberapa hal lagi tentang \hl{\mbox{\texttt{argc}}} dan \hl{\mbox{\texttt{argv}}}.

\begin{itemize}
\item
  Beberapa lingkungan mungkin tidak mengisi \hl{\mbox{\texttt{argv{[}0{]}}}} dengan nama program. Kalau tidak tersedia, \hl{\mbox{\texttt{argv{[}0{]}}}} akan berupa string kosong. Aku belum pernah melihat ini terjadi.
\item
  Spesifikasi cukup longgar soal apa yang boleh dilakukan implementasi terhadap \hl{\mbox{\texttt{argv}}} dan dari mana nilai-nilai itu berasal. Tapi setiap sistem yang pernah kupakai bekerja dengan cara yang sama, seperti yang kita bahas di bagian ini.
\item
  Kamu bisa mengubah \hl{\mbox{\texttt{argc}}}, \hl{\mbox{\texttt{argv}}}, atau string mana pun yang ditunjuk \hl{\mbox{\texttt{argv}}}. (Hanya saja jangan membuat string-string itu lebih panjang dari sebelumnya!)
\item
  Di beberapa sistem mirip Unix, mengubah string \hl{\mbox{\texttt{argv{[}0{]}}}} membuat output \hl{\mbox{\texttt{ps}}} ikut berubah\footnote{\hl{\mbox{\texttt{ps}}}, Process Status, adalah perintah Unix untuk melihat proses apa saja yang sedang berjalan.}.

  Biasanya, kalau kamu punya program bernama \hl{\mbox{\texttt{foo}}} yang dijalankan dengan \hl{\mbox{\texttt{./foo}}}, kamu mungkin melihat ini pada output \hl{\mbox{\texttt{ps}}}:

  \begin{tcolorbox}[enhanced jigsaw,left=3pt,breakable,colback=white!97.2549!black,boxrule=0.2mm]

\begin{Shaded}
\begin{Highlighting}[]
\NormalTok{ 4078 tty1     S      0:00 ./foo}
\end{Highlighting}
\end{Shaded}

  \end{tcolorbox}

  Tapi kalau kamu mengubah \hl{\mbox{\texttt{argv{[}0{]}}}} seperti ini, dengan hati-hati agar string baru \hl{\mbox{\texttt{"Hi!\ \ "}}} panjangnya sama dengan yang lama \hl{\mbox{\texttt{"./foo"}}}:

  \begin{tcolorbox}[enhanced jigsaw,left=3pt,breakable,colback=white!97.2549!black,boxrule=0.2mm]

\begin{Shaded}
\begin{Highlighting}[]
\NormalTok{strcpy}\OperatorTok{(}\NormalTok{argv}\OperatorTok{[}\DecValTok{0}\OperatorTok{],} \StringTok{"Hi!  "}\OperatorTok{);}
\end{Highlighting}
\end{Shaded}

  \end{tcolorbox}

  lalu menjalankan \hl{\mbox{\texttt{ps}}} selagi program \hl{\mbox{\texttt{./foo}}} masih berjalan, kita akan melihat ini sebagai gantinya:

  \begin{tcolorbox}[enhanced jigsaw,left=3pt,breakable,colback=white!97.2549!black,boxrule=0.2mm]

\begin{Shaded}
\begin{Highlighting}[]
\NormalTok{ 4079 tty1     S      0:00 Hi!  }
\end{Highlighting}
\end{Shaded}

  \end{tcolorbox}

  Perilaku ini tidak ada di spesifikasi dan sangat bergantung pada sistem.
\end{itemize}

\index{argc parameter@\texttt{argc} parameter|)} \index{argv parameter@\texttt{argv} parameter|)} \index{Command line arguments|)}

\section{Exit Status}\label{exit-status}

\index{Exit status|(}

Pernahkah kamu perhatikan bahwa signature fungsi \hl{\mbox{\texttt{main()}}} mengembalikan tipe \hl{\mbox{\texttt{int}}}? Untuk apa itu? Ini berkaitan dengan yang disebut \emph{exit status}, yaitu integer yang bisa dikembalikan ke program yang menjalankan programmu untuk memberi tahu bagaimana hasilnya.

Ada beberapa cara program keluar di C, termasuk \hl{\mbox{\texttt{return}}} dari \hl{\mbox{\texttt{main()}}}, atau memanggil salah satu varian \hl{\mbox{\texttt{exit()}}}.

Semua cara itu menerima sebuah \hl{\mbox{\texttt{int}}} sebagai argumen.

\index{main() function@\texttt{main()} function!returning from|(} Catatan sampingan: sudah lihat belum bahwa di hampir semua contohku, meskipun \hl{\mbox{\texttt{main()}}} seharusnya mengembalikan \hl{\mbox{\texttt{int}}}, aku tidak menulis \hl{\mbox{\texttt{return}}} apa pun? Di fungsi lain, ini ilegal, tapi di C ada kasus khusus: kalau eksekusi mencapai akhir \hl{\mbox{\texttt{main()}}} tanpa menemukan \hl{\mbox{\texttt{return}}}, ia otomatis melakukan \hl{\mbox{\texttt{return\ 0}}}. \index{main() function@\texttt{main()} function!returning from|)}

Tapi apa arti \hl{\mbox{\texttt{0}}} itu? Angka apa lagi yang bisa kita taruh di sana? Dan bagaimana angka-angka itu dipakai?

Spesifikasi jelas sekaligus samar soal ini, seperti biasa. Jelas karena ia merinci apa yang bisa dilakukan, tapi samar karena tidak membatasinya secara khusus.

Tidak ada pilihan lain selain \emph{maju terus} dan mencari tahu!

Mari kita ber-\href{https://en.wikipedia.org/wiki/Inception}{Inception}\footnote{https://en.wikipedia.org/wiki/Inception} sebentar: ternyata ketika kamu menjalankan programmu, \emph{kamu menjalankannya dari program lain}.

Biasanya program lain ini semacam \href{https://en.wikipedia.org/wiki/Shell_(computing)}{shell}\footnote{https://en.wikipedia.org/wiki/Shell\_(computing)} yang tidak banyak melakukan apa-apa sendiri kecuali menjalankan program lain.

Tapi ini proses multitahap, terutama terlihat pada shell command-line:

\begin{enumerate}
\def\labelenumi{\arabic{enumi}.}
\tightlist
\item
  Shell menjalankan programmu
\item
  Shell biasanya tidur (untuk shell command-line)
\item
  Programmu berjalan
\item
  Programmu berhenti
\item
  Shell bangun dan menunggu perintah berikutnya
\end{enumerate}

Nah, ada sedikit komunikasi yang terjadi antara langkah 4 dan 5: program bisa mengembalikan sebuah \emph{status value} yang bisa ditanyakan oleh shell. Biasanya, nilai ini dipakai untuk menandakan sukses atau gagalnya programmu, dan, kalau gagal, jenis kegagalannya.

Nilai inilah yang selama ini kita \hl{\mbox{\texttt{return}}} dari \hl{\mbox{\texttt{main()}}}. Itulah statusnya.

Spesifikasi C mengizinkan dua nilai status berbeda, yang nama makronya didefinisikan di \hl{\mbox{\texttt{\textless{}stdlib.h\textgreater{}}}}:

\index{EXIT\_SUCCESS macro@\texttt{EXIT\_SUCCESS} macro|(} \index{EXIT\_FAILURE macro@\texttt{EXIT\_FAILURE} macro|(}

{\def\LTcaptype{none} % do not increment counter
\begin{longtable}[]{@{}ll@{}}
\toprule\noalign{}
Status & Deskripsi \\
\midrule\noalign{}
\endhead
\bottomrule\noalign{}
\endlastfoot
\hl{\mbox{\texttt{EXIT\_SUCCESS}}} atau \hl{\mbox{\texttt{0}}} & Program berakhir dengan sukses. \\
\hl{\mbox{\texttt{EXIT\_FAILURE}}} & Program berakhir dengan error. \\
\end{longtable}
}

Mari kita tulis program singkat yang mengalikan dua angka dari command line. Kamu wajib menyebutkan tepat dua nilai. Kalau tidak, kita cetak pesan error, dan keluar dengan status error.

\begin{tcolorbox}[enhanced jigsaw,left=3pt,breakable,colback=white!97.2549!black,boxrule=0.2mm]

\begin{Shaded}
\begin{Highlighting}[numbers=left,,]
\PreprocessorTok{\#include }\ImportTok{\textless{}stdio.h\textgreater{}}
\PreprocessorTok{\#include }\ImportTok{\textless{}stdlib.h\textgreater{}}

\DataTypeTok{int}\NormalTok{ main}\OperatorTok{(}\DataTypeTok{int}\NormalTok{ argc}\OperatorTok{,} \DataTypeTok{char} \OperatorTok{**}\NormalTok{argv}\OperatorTok{)}
\OperatorTok{\{}
    \ControlFlowTok{if} \OperatorTok{(}\NormalTok{argc }\OperatorTok{!=} \DecValTok{3}\OperatorTok{)} \OperatorTok{\{}
\NormalTok{        printf}\OperatorTok{(}\StringTok{"usage: mult x y}\SpecialCharTok{\textbackslash{}n}\StringTok{"}\OperatorTok{);}
        \ControlFlowTok{return}\NormalTok{ EXIT\_FAILURE}\OperatorTok{;}   \CommentTok{// Beri tahu shell bahwa ini tidak berhasil}
    \OperatorTok{\}}

\NormalTok{    printf}\OperatorTok{(}\StringTok{"}\SpecialCharTok{\%d\textbackslash{}n}\StringTok{"}\OperatorTok{,}\NormalTok{ atoi}\OperatorTok{(}\NormalTok{argv}\OperatorTok{[}\DecValTok{1}\OperatorTok{])} \OperatorTok{*}\NormalTok{ atoi}\OperatorTok{(}\NormalTok{argv}\OperatorTok{[}\DecValTok{2}\OperatorTok{]));}

    \ControlFlowTok{return} \DecValTok{0}\OperatorTok{;}  \CommentTok{// sama dengan EXIT\_SUCCESS, semuanya baik{-}baik saja.}
\OperatorTok{\}}
\end{Highlighting}
\end{Shaded}

\end{tcolorbox}

\index{EXIT\_SUCCESS macro@\texttt{EXIT\_SUCCESS} macro|)} \index{EXIT\_FAILURE macro@\texttt{EXIT\_FAILURE} macro|)}

Kalau kita coba menjalankannya, kita mendapat efek yang diharapkan sampai kita menyebutkan jumlah argumen command-line yang tepat:

\begin{tcolorbox}[enhanced jigsaw,left=3pt,breakable,colback=white!97.2549!black,boxrule=0.2mm]

\begin{Shaded}
\begin{Highlighting}[]
\FunctionTok{$}\NormalTok{ ./mult}
\ExtensionTok{usage:}\NormalTok{ mult x y}

\FunctionTok{$}\NormalTok{ ./mult 3 4 5}
\ExtensionTok{usage:}\NormalTok{ mult x y}

\FunctionTok{$}\NormalTok{ ./mult 3 4}
\ExtensionTok{12}
\end{Highlighting}
\end{Shaded}

\end{tcolorbox}

\index{Exit status!obtaining from shell|(}

Tapi itu belum benar-benar menunjukkan exit status yang kita kembalikan, kan? Kita bisa menyuruh shell mencetaknya. Dengan asumsi kamu memakai Bash atau shell POSIX lainnya, kamu bisa memakai \hl{\mbox{\texttt{echo\ \$?}}} untuk melihatnya\footnote{Di Windows \hl{\mbox{\texttt{cmd.exe}}}, ketik \hl{\mbox{\texttt{echo\ \%errorlevel\%}}}. Di PowerShell, ketik \hl{\mbox{\texttt{\$LastExitCode}}}.}.

Kita coba:

\begin{tcolorbox}[enhanced jigsaw,left=3pt,breakable,colback=white!97.2549!black,boxrule=0.2mm]

\begin{Shaded}
\begin{Highlighting}[]
\FunctionTok{$}\NormalTok{ ./mult}
\ExtensionTok{usage:}\NormalTok{ mult x y}
\FunctionTok{$}\NormalTok{ echo }\VariableTok{$?}
\ExtensionTok{1}

\FunctionTok{$}\NormalTok{ ./mult 3 4 5}
\ExtensionTok{usage:}\NormalTok{ mult x y}
\FunctionTok{$}\NormalTok{ echo }\VariableTok{$?}
\ExtensionTok{1}

\FunctionTok{$}\NormalTok{ ./mult 3 4}
\ExtensionTok{12}
\FunctionTok{$}\NormalTok{ echo }\VariableTok{$?}
\ExtensionTok{0}
\end{Highlighting}
\end{Shaded}

\end{tcolorbox}

\index{Exit status!obtaining from shell|)}

Menarik! Kita lihat bahwa di sistemku, \index{EXIT\_FAILURE macro@\texttt{EXIT\_FAILURE} macro}\hl{\mbox{\texttt{EXIT\_FAILURE}}} adalah \hl{\mbox{\texttt{1}}}. Spesifikasi tidak menyebutkannya, jadi bisa angka apa saja. Tapi coba saja; kemungkinan di sistemmu juga \hl{\mbox{\texttt{1}}}.

\subsection{Nilai Exit Status Lainnya}\label{nilai-exit-status-lainnya}

Status \hl{\mbox{\texttt{0}}} jelas berarti sukses, tapi bagaimana dengan semua integer lain, bahkan yang negatif?

Di sini kita keluar dari spesifikasi C dan masuk ke wilayah Unix. Secara umum, \hl{\mbox{\texttt{0}}} berarti sukses, sedangkan angka positif bukan nol berarti gagal. Jadi hanya ada satu jenis sukses, dan banyak jenis gagal. Bash menyatakan exit code harus antara 0 dan 255, meski sejumlah kode dicadangkan.

Singkatnya, kalau kamu ingin menandakan exit status error yang berbeda di lingkungan Unix, kamu bisa mulai dari \hl{\mbox{\texttt{1}}} dan naik terus.

Di Linux, kalau kamu mencoba kode di luar rentang 0-255, kodenya akan di-AND secara bitwise dengan \hl{\mbox{\texttt{0xff}}}, yang secara efektif menjepitnya ke rentang itu.

Kamu bisa membuat script shell yang memakai kode status ini untuk memutuskan apa yang dilakukan berikutnya.

\index{Exit status|)}

\section{Environment Variable}\label{env-var}

\index{Environment variables|(}

Sebelum masuk ke bagian ini, aku harus memperingatkan bahwa C tidak menentukan apa itu environment variable. Jadi aku akan menjelaskan sistem environment variable yang berfungsi di semua platform besar yang kuketahui.

Pada dasarnya, environment adalah program yang akan menjalankan programmu, misalnya shell bash. Dan ia mungkin punya beberapa variabel bash yang didefinisikan. Kalau kamu belum tahu, shell bisa membuat variabelnya sendiri. Setiap shell berbeda, tapi di bash kamu bisa mengetik \hl{\mbox{\texttt{set}}} dan ia akan menampilkan semuanya.

Berikut potongan dari 61 variabel yang didefinisikan di bash shell-ku:

\begin{tcolorbox}[enhanced jigsaw,left=3pt,breakable,colback=white!97.2549!black,boxrule=0.2mm]

\begin{Shaded}
\begin{Highlighting}[]
\NormalTok{HISTFILE=/home/beej/.bash\_history}
\NormalTok{HISTFILESIZE=500}
\NormalTok{HISTSIZE=500}
\NormalTok{HOME=/home/beej}
\NormalTok{HOSTNAME=FBILAPTOP}
\NormalTok{HOSTTYPE=x86\_64}
\NormalTok{IFS=$\textquotesingle{} \textbackslash{}t\textbackslash{}n\textquotesingle{}}
\end{Highlighting}
\end{Shaded}

\end{tcolorbox}

Perhatikan bentuknya berupa pasangan key/value. Misalnya, salah satu key-nya adalah \hl{\mbox{\texttt{HOSTTYPE}}} dan value-nya \hl{\mbox{\texttt{x86\_64}}}. Dari sudut pandang C, semua value adalah string, bahkan kalau berupa angka\footnote{Kalau kamu butuh nilai numerik, konversikan string itu dengan sesuatu seperti \hl{\mbox{\texttt{atoi()}}} atau \hl{\mbox{\texttt{strtol()}}}.}.

Jadi, \emph{anyway}! Singkat cerita, nilai-nilai ini bisa diambil dari dalam programmu.

\index{getenv() function@\texttt{getenv()} function|(}

Mari kita tulis program yang memakai fungsi standar \hl{\mbox{\texttt{getenv()}}} untuk mencari nilai yang kamu set di shell.

\hl{\mbox{\texttt{getenv()}}} mengembalikan pointer ke string value, atau \hl{\mbox{\texttt{NULL}}} kalau environment variable-nya tidak ada.

\begin{tcolorbox}[enhanced jigsaw,left=3pt,breakable,colback=white!97.2549!black,boxrule=0.2mm]

\begin{Shaded}
\begin{Highlighting}[numbers=left,,]
\PreprocessorTok{\#include }\ImportTok{\textless{}stdio.h\textgreater{}}
\PreprocessorTok{\#include }\ImportTok{\textless{}stdlib.h\textgreater{}}

\DataTypeTok{int}\NormalTok{ main}\OperatorTok{(}\DataTypeTok{void}\OperatorTok{)}
\OperatorTok{\{}
    \DataTypeTok{char} \OperatorTok{*}\NormalTok{val }\OperatorTok{=}\NormalTok{ getenv}\OperatorTok{(}\StringTok{"FROTZ"}\OperatorTok{);}  \CommentTok{// Coba ambil nilainya}

    \CommentTok{// Pastikan nilainya ada}
    \ControlFlowTok{if} \OperatorTok{(}\NormalTok{val }\OperatorTok{==}\NormalTok{ NULL}\OperatorTok{)} \OperatorTok{\{}
\NormalTok{        printf}\OperatorTok{(}\StringTok{"Cannot find the FROTZ environment variable}\SpecialCharTok{\textbackslash{}n}\StringTok{"}\OperatorTok{);}
        \ControlFlowTok{return}\NormalTok{ EXIT\_FAILURE}\OperatorTok{;}
    \OperatorTok{\}}

\NormalTok{    printf}\OperatorTok{(}\StringTok{"Value: }\SpecialCharTok{\%s\textbackslash{}n}\StringTok{"}\OperatorTok{,}\NormalTok{ val}\OperatorTok{);}
\OperatorTok{\}}
\end{Highlighting}
\end{Shaded}

\end{tcolorbox}

\index{getenv() function@\texttt{getenv()} function|)}

Kalau kujalankan langsung, hasilnya begini:

\begin{tcolorbox}[enhanced jigsaw,left=3pt,breakable,colback=white!97.2549!black,boxrule=0.2mm]

\begin{Shaded}
\begin{Highlighting}[]
\FunctionTok{$}\NormalTok{ ./foo}
\ExtensionTok{Cannot}\NormalTok{ find the FROTZ environment variable}
\end{Highlighting}
\end{Shaded}

\end{tcolorbox}

yang masuk akal, karena belum kuset.

Di bash, aku bisa mengisinya dengan\footnote{Di Windows CMD.EXE, pakai \hl{\mbox{\texttt{set\ FROTZ=value}}}. Di PowerShell, pakai \hl{\mbox{\texttt{\$Env:FROTZ=value}}}.}:

\begin{tcolorbox}[enhanced jigsaw,left=3pt,breakable,colback=white!97.2549!black,boxrule=0.2mm]

\begin{Shaded}
\begin{Highlighting}[]
\FunctionTok{$}\NormalTok{ export FROTZ=}\StringTok{"C is awesome!"}
\end{Highlighting}
\end{Shaded}

\end{tcolorbox}

Lalu kalau kujalankan, hasilnya:

\begin{tcolorbox}[enhanced jigsaw,left=3pt,breakable,colback=white!97.2549!black,boxrule=0.2mm]

\begin{Shaded}
\begin{Highlighting}[]
\FunctionTok{$}\NormalTok{ ./foo}
\ExtensionTok{Value:}\NormalTok{ C is awesome!}
\end{Highlighting}
\end{Shaded}

\end{tcolorbox}

Dengan cara ini, kamu bisa menyiapkan data di environment variable, lalu mengambilnya di kode C-mu dan menyesuaikan perilaku program.

\subsection{Mengatur Environment Variable}\label{mengatur-environment-variable}

Ini bukan standar, tapi banyak sistem menyediakan cara untuk mengatur environment variable.

Kalau di sistem mirip Unix, lihat dokumentasi \index{putenv() function@\texttt{putenv()} function}\hl{\mbox{\texttt{putenv()}}}, \index{setenv()@\texttt{setenv()}}\hl{\mbox{\texttt{setenv()}}}, dan \index{unsetenv() function@\texttt{unsetenv()} function}\hl{\mbox{\texttt{unsetenv()}}}. Di Windows, lihat \index{\_putenv() function@\texttt{\_putenv()} function}\hl{\mbox{\texttt{\_putenv()}}}.

\subsection{Environment Variable Alternatif di Sistem Mirip Unix}\label{environment-variable-alternatif-di-sistem-mirip-unix}

Kalau kamu di sistem mirip Unix, kemungkinan kamu punya beberapa cara lain untuk mengakses environment variable. Perhatikan bahwa walaupun spesifikasi menunjukkannya sebagai ekstensi umum, ini bukan benar-benar bagian dari standar C. Meski demikian, ini bagian dari standar POSIX.

\index{environ variable@\texttt{environ} variable|(}

Salah satunya adalah variabel bernama \hl{\mbox{\texttt{environ}}} yang harus dideklarasikan seperti ini:

\begin{tcolorbox}[enhanced jigsaw,left=3pt,breakable,colback=white!97.2549!black,boxrule=0.2mm]

\begin{Shaded}
\begin{Highlighting}[]
\KeywordTok{extern} \DataTypeTok{char} \OperatorTok{**}\NormalTok{environ}\OperatorTok{;}
\end{Highlighting}
\end{Shaded}

\end{tcolorbox}

Ia adalah array berisi string yang diakhiri pointer \hl{\mbox{\texttt{NULL}}}.

Kamu harus mendeklarasikannya sendiri sebelum memakainya, atau mungkin kamu menemukannya di header file non-standar \hl{\mbox{\texttt{\textless{}unistd.h\textgreater{}}}}.

Setiap string berbentuk \hl{\mbox{\texttt{"key=value"}}} sehingga kamu harus memecah dan mem-parse-nya sendiri kalau ingin mengambil key dan value-nya.

Berikut contoh menelusuri dan mencetak environment variable dengan beberapa cara berbeda:

\begin{tcolorbox}[enhanced jigsaw,left=3pt,breakable,colback=white!97.2549!black,boxrule=0.2mm]

\begin{Shaded}
\begin{Highlighting}[numbers=left,,]
\PreprocessorTok{\#include }\ImportTok{\textless{}stdio.h\textgreater{}}

\KeywordTok{extern} \DataTypeTok{char} \OperatorTok{**}\NormalTok{environ}\OperatorTok{;}  \CommentTok{// HARUS extern DAN bernama "environ"}

\DataTypeTok{int}\NormalTok{ main}\OperatorTok{(}\DataTypeTok{void}\OperatorTok{)}
\OperatorTok{\{}
    \ControlFlowTok{for} \OperatorTok{(}\DataTypeTok{char} \OperatorTok{**}\NormalTok{p }\OperatorTok{=}\NormalTok{ environ}\OperatorTok{;} \OperatorTok{*}\NormalTok{p }\OperatorTok{!=}\NormalTok{ NULL}\OperatorTok{;}\NormalTok{ p}\OperatorTok{++)} \OperatorTok{\{}
\NormalTok{        printf}\OperatorTok{(}\StringTok{"}\SpecialCharTok{\%s\textbackslash{}n}\StringTok{"}\OperatorTok{,} \OperatorTok{*}\NormalTok{p}\OperatorTok{);}
    \OperatorTok{\}}

    \CommentTok{// Atau bisa juga begini:}
    \ControlFlowTok{for} \OperatorTok{(}\DataTypeTok{int}\NormalTok{ i }\OperatorTok{=} \DecValTok{0}\OperatorTok{;}\NormalTok{ environ}\OperatorTok{[}\NormalTok{i}\OperatorTok{]} \OperatorTok{!=}\NormalTok{ NULL}\OperatorTok{;}\NormalTok{ i}\OperatorTok{++)} \OperatorTok{\{}
\NormalTok{        printf}\OperatorTok{(}\StringTok{"}\SpecialCharTok{\%s\textbackslash{}n}\StringTok{"}\OperatorTok{,}\NormalTok{ environ}\OperatorTok{[}\NormalTok{i}\OperatorTok{]);}
    \OperatorTok{\}}
\OperatorTok{\}}
\end{Highlighting}
\end{Shaded}

\end{tcolorbox}

Yang menghasilkan banyak output seperti ini:

\begin{tcolorbox}[enhanced jigsaw,left=3pt,breakable,colback=white!97.2549!black,boxrule=0.2mm]

\begin{Shaded}
\begin{Highlighting}[]
\NormalTok{SHELL=/bin/bash}
\NormalTok{COLORTERM=truecolor}
\NormalTok{TERM\_PROGRAM\_VERSION=1.53.2}
\NormalTok{LOGNAME=beej}
\NormalTok{HOME=/home/beej}
\NormalTok{... etc ...}
\end{Highlighting}
\end{Shaded}

\end{tcolorbox}

Pakai \hl{\mbox{\texttt{getenv()}}} kalau memungkinkan karena lebih portabel. Tapi kalau harus menelusuri environment variable, memakai \hl{\mbox{\texttt{environ}}} mungkin pilihan yang tepat.

\index{environ variable@\texttt{environ} variable|)} \index{env parameter@\texttt{env} parameter|(}

Cara non-standar lain untuk mendapat environment variable adalah sebagai parameter \hl{\mbox{\texttt{main()}}}. Cara kerjanya hampir sama, tapi kamu tidak perlu menambahkan variabel \hl{\mbox{\texttt{extern}}} \hl{\mbox{\texttt{environ}}}. \href{https://pubs.opengroup.org/onlinepubs/9699919799/functions/exec.html}{Bahkan spesifikasi POSIX pun tidak mendukung ini}\footnote{https://pubs.opengroup.org/onlinepubs/9699919799/functions/exec.html} sejauh yang kutahu, tapi ini umum di dunia Unix.

\begin{tcolorbox}[enhanced jigsaw,left=3pt,breakable,colback=white!97.2549!black,boxrule=0.2mm]

\begin{Shaded}
\begin{Highlighting}[numbers=left,,]
\PreprocessorTok{\#include }\ImportTok{\textless{}stdio.h\textgreater{}}

\DataTypeTok{int}\NormalTok{ main}\OperatorTok{(}\DataTypeTok{int}\NormalTok{ argc}\OperatorTok{,} \DataTypeTok{char} \OperatorTok{**}\NormalTok{argv}\OperatorTok{,} \DataTypeTok{char} \OperatorTok{**}\NormalTok{env}\OperatorTok{)}  \CommentTok{// \textless{}{-}{-} env!}
\OperatorTok{\{}
    \OperatorTok{(}\DataTypeTok{void}\OperatorTok{)}\NormalTok{argc}\OperatorTok{;} \OperatorTok{(}\DataTypeTok{void}\OperatorTok{)}\NormalTok{argv}\OperatorTok{;}  \CommentTok{// Tekan peringatan variabel tak terpakai}

    \ControlFlowTok{for} \OperatorTok{(}\DataTypeTok{char} \OperatorTok{**}\NormalTok{p }\OperatorTok{=}\NormalTok{ env}\OperatorTok{;} \OperatorTok{*}\NormalTok{p }\OperatorTok{!=}\NormalTok{ NULL}\OperatorTok{;}\NormalTok{ p}\OperatorTok{++)} \OperatorTok{\{}
\NormalTok{        printf}\OperatorTok{(}\StringTok{"}\SpecialCharTok{\%s\textbackslash{}n}\StringTok{"}\OperatorTok{,} \OperatorTok{*}\NormalTok{p}\OperatorTok{);}
    \OperatorTok{\}}

    \CommentTok{// Atau bisa juga begini:}
    \ControlFlowTok{for} \OperatorTok{(}\DataTypeTok{int}\NormalTok{ i }\OperatorTok{=} \DecValTok{0}\OperatorTok{;}\NormalTok{ env}\OperatorTok{[}\NormalTok{i}\OperatorTok{]} \OperatorTok{!=}\NormalTok{ NULL}\OperatorTok{;}\NormalTok{ i}\OperatorTok{++)} \OperatorTok{\{}
\NormalTok{        printf}\OperatorTok{(}\StringTok{"}\SpecialCharTok{\%s\textbackslash{}n}\StringTok{"}\OperatorTok{,}\NormalTok{ env}\OperatorTok{[}\NormalTok{i}\OperatorTok{]);}
    \OperatorTok{\}}
\OperatorTok{\}}
\end{Highlighting}
\end{Shaded}

\end{tcolorbox}

Persis seperti memakai \hl{\mbox{\texttt{environ}}} tapi \emph{lebih tidak portabel lagi}. Bagus punya target.

\index{env parameter@\texttt{env} parameter|)} \index{Environment variables|)}

\chapter{C Preprocessor}\label{c-preprocessor}

\index{Preprocessor|(}

Sebelum programmu dikompilasi, ia sebenarnya melewati tahap yang disebut \emph{preprocessing}. Hampir seperti ada bahasa \emph{di atas} bahasa C yang berjalan lebih dulu. Dan ia menghasilkan kode C, yang kemudian dikompilasi.

Kita sudah melihatnya sedikit lewat \hl{\mbox{\texttt{\#include}}}! Itulah C Preprocessor! Di tempat ia menemukan directive itu, ia menyertakan file yang disebut tepat di situ, seolah-olah kamu mengetiknya sendiri. Dan \emph{baru kemudian} compiler membangun semuanya.

Tapi ternyata kemampuannya jauh lebih besar daripada sekadar menyertakan sesuatu. Kamu bisa mendefinisikan \emph{makro} yang disubstitusi\ldots{} bahkan makro yang menerima argumen!

\section{\texorpdfstring{\hl{\mbox{\texttt{\#include}}}}{\#include}}\label{include}

\index{\#include directive@\texttt{\#include} directive|(}

Kita mulai dari yang sudah sering kita lihat. Ini tentu saja cara untuk menyertakan sumber lain ke dalam sumbermu. Sangat umum dipakai dengan header file.

Walaupun spesifikasi mengizinkan berbagai perilaku untuk \hl{\mbox{\texttt{\#include}}}, kita akan mengambil pendekatan yang lebih pragmatis dan membahas cara kerjanya di setiap sistem yang pernah kulihat.

Header file bisa dibagi dua kategori: sistem dan lokal. Yang bawaan, seperti \hl{\mbox{\texttt{stdio.h}}}, \hl{\mbox{\texttt{stdlib.h}}}, \hl{\mbox{\texttt{math.h}}}, dan sebagainya, di-include dengan kurung siku:

\begin{tcolorbox}[enhanced jigsaw,left=3pt,breakable,colback=white!97.2549!black,boxrule=0.2mm]

\begin{Shaded}
\begin{Highlighting}[]
\PreprocessorTok{\#include }\ImportTok{\textless{}stdio.h\textgreater{}}
\PreprocessorTok{\#include }\ImportTok{\textless{}stdlib.h\textgreater{}}
\end{Highlighting}
\end{Shaded}

\end{tcolorbox}

Kurung siku memberi tahu C, ``Hei, jangan cari header file ini di direktori saat ini---cari di direktori include level sistem.''

\index{\#include directive@\texttt{\#include} directive!local files|(}

Itu tentu saja menyiratkan bahwa harus ada cara untuk meng-include file lokal dari direktori saat ini. Dan memang ada: dengan tanda kutip ganda:

\begin{tcolorbox}[enhanced jigsaw,left=3pt,breakable,colback=white!97.2549!black,boxrule=0.2mm]

\begin{Shaded}
\begin{Highlighting}[]
\PreprocessorTok{\#include }\ImportTok{"myheader.h"}
\end{Highlighting}
\end{Shaded}

\end{tcolorbox}

Atau kamu juga hampir pasti bisa mencari di direktori relatif memakai garis miring dan titik, seperti ini:

\begin{tcolorbox}[enhanced jigsaw,left=3pt,breakable,colback=white!97.2549!black,boxrule=0.2mm]

\begin{Shaded}
\begin{Highlighting}[]
\PreprocessorTok{\#include }\ImportTok{"mydir/myheader.h"}
\PreprocessorTok{\#include }\ImportTok{"../someheader.py"}
\end{Highlighting}
\end{Shaded}

\end{tcolorbox}

Jangan pakai backslash (\hl{\mbox{\texttt{\textbackslash{}}}}) sebagai pemisah path di \hl{\mbox{\texttt{\#include}}}-mu! Itu undefined behavior! Pakai garis miring biasa (\hl{\mbox{\texttt{/}}}) saja, bahkan di Windows.

Ringkasnya, pakai kurung siku (\hl{\mbox{\texttt{\textless{}}}} dan \hl{\mbox{\texttt{\textgreater{}}}}) untuk include sistem, dan pakai tanda kutip ganda (\hl{\mbox{\texttt{"}}}) untuk include pribadimu.

\index{\#include directive@\texttt{\#include} directive!local files|)} \index{\#include directive@\texttt{\#include} directive|)}

\section{Makro Sederhana}\label{makro-sederhana}

\index{Preprocessor!macros|(}

\emph{Makro} adalah identifier yang \emph{diperluas} menjadi potongan kode lain sebelum compiler sempat melihatnya. Anggap seperti placeholder---ketika preprocessor melihat salah satu identifier itu, ia menggantinya dengan nilai lain yang sudah kamu definisikan.

\index{\#define directive@\texttt{\#define} directive|(}

Kita melakukannya dengan \hl{\mbox{\texttt{\#define}}} (sering dibaca ``pound define'', atau mungkin ``hash define'', tapi jarang, kalaupun ada, ``octothorpe define''). Berikut contohnya:

\begin{tcolorbox}[enhanced jigsaw,left=3pt,breakable,colback=white!97.2549!black,boxrule=0.2mm]

\begin{Shaded}
\begin{Highlighting}[numbers=left,,]
\PreprocessorTok{\#include }\ImportTok{\textless{}stdio.h\textgreater{}}

\PreprocessorTok{\#define HELLO }\StringTok{"Hello, world"}
\PreprocessorTok{\#define PI }\FloatTok{3.14159}

\DataTypeTok{int}\NormalTok{ main}\OperatorTok{(}\DataTypeTok{void}\OperatorTok{)}
\OperatorTok{\{}
\NormalTok{    printf}\OperatorTok{(}\StringTok{"}\SpecialCharTok{\%s}\StringTok{, }\SpecialCharTok{\%f\textbackslash{}n}\StringTok{"}\OperatorTok{,}\NormalTok{ HELLO}\OperatorTok{,}\NormalTok{ PI}\OperatorTok{);}
\OperatorTok{\}}
\end{Highlighting}
\end{Shaded}

\end{tcolorbox}

Di baris 3 dan 4 kita mendefinisikan beberapa makro. Di mana pun makro itu muncul di kode (baris 8), mereka akan disubstitusi dengan nilai yang didefinisikan.

Dari sudut pandang C compiler, itu persis seperti kalau kita menulis ini:

\begin{tcolorbox}[enhanced jigsaw,left=3pt,breakable,colback=white!97.2549!black,boxrule=0.2mm]

\begin{Shaded}
\begin{Highlighting}[numbers=left,,]
\PreprocessorTok{\#include }\ImportTok{\textless{}stdio.h\textgreater{}}

\DataTypeTok{int}\NormalTok{ main}\OperatorTok{(}\DataTypeTok{void}\OperatorTok{)}
\OperatorTok{\{}
\NormalTok{    printf}\OperatorTok{(}\StringTok{"}\SpecialCharTok{\%s}\StringTok{, }\SpecialCharTok{\%f\textbackslash{}n}\StringTok{"}\OperatorTok{,} \StringTok{"Hello, world"}\OperatorTok{,} \FloatTok{3.14159}\OperatorTok{);}
\OperatorTok{\}}
\end{Highlighting}
\end{Shaded}

\end{tcolorbox}

Lihat bagaimana \hl{\mbox{\texttt{HELLO}}} diganti dengan \hl{\mbox{\texttt{"Hello,\ world"}}} dan \hl{\mbox{\texttt{PI}}} diganti dengan \hl{\mbox{\texttt{3.14159}}}? Dari sudut pandang compiler, itu persis seperti nilai-nilai itu muncul langsung di kode.

Perhatikan bahwa makro tidak punya tipe tertentu, \emph{per se}. Yang terjadi hanyalah mereka diganti utuh dengan apa pun yang di-\hl{\mbox{\texttt{\#define}}}-kan. Kalau kode C hasilnya tidak valid, compiler akan muntah.

Kamu juga bisa mendefinisikan makro tanpa nilai:

\begin{tcolorbox}[enhanced jigsaw,left=3pt,breakable,colback=white!97.2549!black,boxrule=0.2mm]

\begin{Shaded}
\begin{Highlighting}[]
\PreprocessorTok{\#define EXTRA\_HAPPY}
\end{Highlighting}
\end{Shaded}

\end{tcolorbox}

dalam kasus itu, makro-nya ada dan terdefinisi, tapi didefinisikan sebagai kosong. Jadi di mana pun ia muncul di teks, akan diganti dengan kosong. Kita akan melihat kegunaannya nanti.

Secara konvensi, nama makro ditulis \hl{\mbox{\texttt{HURUF\_BESAR\_SEMUA}}} walaupun secara teknis itu tidak wajib.

\index{\#define directive@\texttt{\#define} directive!versus const@versus \texttt{const}|(}

Secara keseluruhan, ini memberimu cara untuk mendefinisikan nilai konstan yang efektif bersifat global dan bisa dipakai di tempat \emph{mana saja}. Bahkan di tempat yang tidak bisa memakai variabel \hl{\mbox{\texttt{const}}}, misalnya pada \hl{\mbox{\texttt{case}}} di \hl{\mbox{\texttt{switch}}} dan panjang array tetap.

Meski begitu, perdebatan masih ramai di internet soal apakah variabel \hl{\mbox{\texttt{const}}} bertipe lebih baik daripada makro \hl{\mbox{\texttt{\#define}}} secara umum.

Ia juga bisa dipakai untuk mengganti atau memodifikasi kata kunci, konsep yang sama sekali asing bagi \hl{\mbox{\texttt{const}}}, walaupun praktik ini sebaiknya jarang dilakukan.

\index{\#define directive@\texttt{\#define} directive!versus const@versus \texttt{const}|)} \index{\#define directive@\texttt{\#define} directive|)} \index{Preprocessor!macros|)}

\section{Kompilasi Kondisional}\label{kompilasi-kondisional}

\index{Conditional compilation|(}

Preprocessor bisa diminta menentukan apakah blok kode tertentu ditampilkan ke compiler atau dibuang sepenuhnya sebelum kompilasi.

Caranya dengan membungkus kode dalam blok kondisional, mirip pernyataan \hl{\mbox{\texttt{if}}}-\hl{\mbox{\texttt{else}}}.

\subsection{\texorpdfstring{Jika Terdefinisi, \hl{\mbox{\texttt{\#ifdef}}} dan \hl{\mbox{\texttt{\#endif}}}}{Jika Terdefinisi, \#ifdef dan \#endif}}\label{jika-terdefinisi-ifdef-dan-endif}

Pertama, mari kita coba mengompilasi kode tertentu bergantung pada apakah sebuah makro terdefinisi atau tidak.

\index{\#ifdef directive@\texttt{\#ifdef} directive|(} \index{\#endif directive@\texttt{\#endif} directive|(}

\begin{tcolorbox}[enhanced jigsaw,left=3pt,breakable,colback=white!97.2549!black,boxrule=0.2mm]

\begin{Shaded}
\begin{Highlighting}[numbers=left,,]
\PreprocessorTok{\#include }\ImportTok{\textless{}stdio.h\textgreater{}}

\PreprocessorTok{\#define EXTRA\_HAPPY}

\DataTypeTok{int}\NormalTok{ main}\OperatorTok{(}\DataTypeTok{void}\OperatorTok{)}
\OperatorTok{\{}

\PreprocessorTok{\#ifdef EXTRA\_HAPPY}
\NormalTok{    printf}\OperatorTok{(}\StringTok{"I\textquotesingle{}m extra happy!}\SpecialCharTok{\textbackslash{}n}\StringTok{"}\OperatorTok{);}
\PreprocessorTok{\#endif}

\NormalTok{    printf}\OperatorTok{(}\StringTok{"OK!}\SpecialCharTok{\textbackslash{}n}\StringTok{"}\OperatorTok{);}
\OperatorTok{\}}
\end{Highlighting}
\end{Shaded}

\end{tcolorbox}

Pada contoh itu, kita mendefinisikan \hl{\mbox{\texttt{EXTRA\_HAPPY}}} (sebagai kosong, tapi ia \emph{terdefinisi}), lalu di baris 8 kita memeriksa apakah ia terdefinisi dengan directive \hl{\mbox{\texttt{\#ifdef}}}. Kalau terdefinisi, kode setelahnya disertakan sampai \hl{\mbox{\texttt{\#endif}}}.

\index{\#ifdef directive@\texttt{\#ifdef} directive|)} \index{\#endif directive@\texttt{\#endif} directive|)}

Jadi karena terdefinisi, kode itu disertakan untuk kompilasi dan outputnya:

\begin{tcolorbox}[enhanced jigsaw,left=3pt,breakable,colback=white!97.2549!black,boxrule=0.2mm]

\begin{Shaded}
\begin{Highlighting}[]
\NormalTok{I\textquotesingle{}m extra happy!}
\NormalTok{OK!}
\end{Highlighting}
\end{Shaded}

\end{tcolorbox}

Kalau \hl{\mbox{\texttt{\#define}}}-nya kita jadikan komentar, seperti ini:

\begin{tcolorbox}[enhanced jigsaw,left=3pt,breakable,colback=white!97.2549!black,boxrule=0.2mm]

\begin{Shaded}
\begin{Highlighting}[]
\CommentTok{//\#define EXTRA\_HAPPY}
\end{Highlighting}
\end{Shaded}

\end{tcolorbox}

maka ia tidak terdefinisi, dan kode itu tidak disertakan dalam kompilasi. Dan outputnya hanya:

\begin{tcolorbox}[enhanced jigsaw,left=3pt,breakable,colback=white!97.2549!black,boxrule=0.2mm]

\begin{Shaded}
\begin{Highlighting}[]
\NormalTok{OK!}
\end{Highlighting}
\end{Shaded}

\end{tcolorbox}

Penting diingat bahwa keputusan ini terjadi saat compile time! Kodenya benar-benar dikompilasi atau dibuang tergantung kondisi. Ini berbeda dengan pernyataan \hl{\mbox{\texttt{if}}} biasa yang dievaluasi saat program berjalan.

\subsection{\texorpdfstring{Jika Tidak Terdefinisi, \hl{\mbox{\texttt{\#ifndef}}}}{Jika Tidak Terdefinisi, \#ifndef}}\label{jika-tidak-terdefinisi-ifndef}

Ada juga bentuk negatif dari ``jika terdefinisi'': ``jika tidak terdefinisi'', atau \hl{\mbox{\texttt{\#ifndef}}}. Kita bisa mengubah contoh sebelumnya supaya mengeluarkan hal berbeda tergantung apakah sesuatu terdefinisi:

\index{\#ifndef directive@\texttt{\#ifndef} directive|(} \index{\#endif directive@\texttt{\#endif} directive|(}

\begin{tcolorbox}[enhanced jigsaw,left=3pt,breakable,colback=white!97.2549!black,boxrule=0.2mm]

\begin{Shaded}
\begin{Highlighting}[numbers=left,,firstnumber=8,]
\PreprocessorTok{\#ifdef EXTRA\_HAPPY}
\NormalTok{    printf}\OperatorTok{(}\StringTok{"I\textquotesingle{}m extra happy!}\SpecialCharTok{\textbackslash{}n}\StringTok{"}\OperatorTok{);}
\PreprocessorTok{\#endif}

\PreprocessorTok{\#ifndef EXTRA\_HAPPY}
\NormalTok{    printf}\OperatorTok{(}\StringTok{"I\textquotesingle{}m just regular}\SpecialCharTok{\textbackslash{}n}\StringTok{"}\OperatorTok{);}
\PreprocessorTok{\#endif}
\end{Highlighting}
\end{Shaded}

\end{tcolorbox}

Kita akan lihat cara yang lebih rapi untuk itu di bagian berikutnya.

Kalau dikaitkan kembali ke header file, kita sudah melihat bagaimana header file bisa dibuat hanya di-include satu kali dengan membungkusnya dalam directive preprocessor seperti ini:

\begin{tcolorbox}[enhanced jigsaw,left=3pt,breakable,colback=white!97.2549!black,boxrule=0.2mm]

\begin{Shaded}
\begin{Highlighting}[]
\PreprocessorTok{\#ifndef MYHEADER\_H  }\CommentTok{// Baris pertama myheader.h}
\PreprocessorTok{\#define MYHEADER\_H}

\DataTypeTok{int}\NormalTok{ x }\OperatorTok{=} \DecValTok{12}\OperatorTok{;}

\PreprocessorTok{\#endif  }\CommentTok{// Baris terakhir myheader.h}
\end{Highlighting}
\end{Shaded}

\end{tcolorbox}

\index{\#ifndef directive@\texttt{\#ifndef} directive|)} \index{\#endif directive@\texttt{\#endif} directive|)}

Ini menunjukkan bagaimana sebuah makro bertahan lintas file dan beberapa \hl{\mbox{\texttt{\#include}}}. Kalau belum terdefinisi, kita definisikan lalu kompilasi seluruh header file.

Tapi saat di-include lagi, kita lihat bahwa \hl{\mbox{\texttt{MYHEADER\_H}}} \emph{sudah} terdefinisi, jadi kita tidak mengirim header file itu ke compiler---ia secara efektif dibuang.

\subsection{\texorpdfstring{\hl{\mbox{\texttt{\#else}}}}{\#else}}\label{else}

\index{\#else directive@\texttt{\#else} directive|(}

Tapi itu belum semua yang bisa kita lakukan! Ada juga \hl{\mbox{\texttt{\#else}}} yang bisa kita campurkan.

\index{\#ifdef directive@\texttt{\#ifdef} directive|(} \index{\#endif directive@\texttt{\#endif} directive|(}

Mari kita modifikasi contoh sebelumnya:

\begin{tcolorbox}[enhanced jigsaw,left=3pt,breakable,colback=white!97.2549!black,boxrule=0.2mm]

\begin{Shaded}
\begin{Highlighting}[numbers=left,,firstnumber=8,]
\PreprocessorTok{\#ifdef EXTRA\_HAPPY}
\NormalTok{    printf}\OperatorTok{(}\StringTok{"I\textquotesingle{}m extra happy!}\SpecialCharTok{\textbackslash{}n}\StringTok{"}\OperatorTok{);}
\PreprocessorTok{\#else}
\NormalTok{    printf}\OperatorTok{(}\StringTok{"I\textquotesingle{}m just regular}\SpecialCharTok{\textbackslash{}n}\StringTok{"}\OperatorTok{);}
\PreprocessorTok{\#endif}
\end{Highlighting}
\end{Shaded}

\end{tcolorbox}

\index{\#ifdef directive@\texttt{\#ifdef} directive|)} \index{\#else directive@\texttt{\#else} directive|)} \index{\#endif directive@\texttt{\#endif} directive|)}

Sekarang kalau \hl{\mbox{\texttt{EXTRA\_HAPPY}}} tidak terdefinisi, ia akan masuk ke klausa \hl{\mbox{\texttt{\#else}}} dan mencetak:

\begin{tcolorbox}[enhanced jigsaw,left=3pt,breakable,colback=white!97.2549!black,boxrule=0.2mm]

\begin{Shaded}
\begin{Highlighting}[]
\NormalTok{I\textquotesingle{}m just regular}
\end{Highlighting}
\end{Shaded}

\end{tcolorbox}

\subsection{\texorpdfstring{Else-If: \hl{\mbox{\texttt{\#elifdef}}}, \hl{\mbox{\texttt{\#elifndef}}}}{Else-If: \#elifdef, \#elifndef}}\label{else-if-elifdef-elifndef}

\index{\#elifdef directive@\texttt{\#elifdef} directive|(} \index{\#elifndef directive@\texttt{\#elifndef} directive|(}

Fitur ini baru di C23!

Bagaimana kalau kamu ingin sesuatu yang lebih kompleks? Mungkin kamu butuh struktur rangkaian if-else supaya kodemu terbangun dengan benar?

Untungnya directive-directive ini tersedia. Kita bisa memakai \hl{\mbox{\texttt{\#elifdef}}} untuk ``else if defined'':

\begin{tcolorbox}[enhanced jigsaw,left=3pt,breakable,colback=white!97.2549!black,boxrule=0.2mm]

\begin{Shaded}
\begin{Highlighting}[]
\PreprocessorTok{\#ifdef MODE\_1}
\NormalTok{    printf}\OperatorTok{(}\StringTok{"This is mode 1}\SpecialCharTok{\textbackslash{}n}\StringTok{"}\OperatorTok{);}
\PreprocessorTok{\#elifdef MODE\_2}
\NormalTok{    printf}\OperatorTok{(}\StringTok{"This is mode 2}\SpecialCharTok{\textbackslash{}n}\StringTok{"}\OperatorTok{);}
\PreprocessorTok{\#elifdef MODE\_3}
\NormalTok{    printf}\OperatorTok{(}\StringTok{"This is mode 3}\SpecialCharTok{\textbackslash{}n}\StringTok{"}\OperatorTok{);}
\PreprocessorTok{\#else}
\NormalTok{    printf}\OperatorTok{(}\StringTok{"This is some other mode}\SpecialCharTok{\textbackslash{}n}\StringTok{"}\OperatorTok{);}
\PreprocessorTok{\#endif}
\end{Highlighting}
\end{Shaded}

\end{tcolorbox}

\index{\#elifdef directive@\texttt{\#elifdef} directive|)}

Sebaliknya, kamu bisa memakai \hl{\mbox{\texttt{\#elifndef}}} untuk ``else if not defined''.

\index{\#elifndef directive@\texttt{\#elifndef} directive|)}

\subsection{\texorpdfstring{Kondisional Umum: \hl{\mbox{\texttt{\#if}}}, \hl{\mbox{\texttt{\#elif}}}}{Kondisional Umum: \#if, \#elif}}\label{kondisional-umum-if-elif}

\index{\#if directive@\texttt{\#if} directive|(} \index{\#elif directive@\texttt{\#elif} directive|(}

Ini bekerja sangat mirip directive \hl{\mbox{\texttt{\#ifdef}}} dan \hl{\mbox{\texttt{\#ifndef}}}, yaitu kamu juga bisa punya \hl{\mbox{\texttt{\#else}}} dan semuanya ditutup dengan \hl{\mbox{\texttt{\#endif}}}.

Satu-satunya perbedaan: ekspresi konstan setelah \hl{\mbox{\texttt{\#if}}} harus bernilai true (bukan nol) agar kode di dalam \hl{\mbox{\texttt{\#if}}} dikompilasi. Jadi alih-alih apakah sesuatu terdefinisi atau tidak, kita ingin ekspresi yang bernilai true.

\begin{tcolorbox}[enhanced jigsaw,left=3pt,breakable,colback=white!97.2549!black,boxrule=0.2mm]

\begin{Shaded}
\begin{Highlighting}[numbers=left,,]
\PreprocessorTok{\#include }\ImportTok{\textless{}stdio.h\textgreater{}}

\PreprocessorTok{\#define HAPPY\_FACTOR }\DecValTok{1}

\DataTypeTok{int}\NormalTok{ main}\OperatorTok{(}\DataTypeTok{void}\OperatorTok{)}
\OperatorTok{\{}

\PreprocessorTok{\#if HAPPY\_FACTOR == 0}
\NormalTok{    printf}\OperatorTok{(}\StringTok{"I\textquotesingle{}m not happy!}\SpecialCharTok{\textbackslash{}n}\StringTok{"}\OperatorTok{);}
\PreprocessorTok{\#elif HAPPY\_FACTOR == 1}
\NormalTok{    printf}\OperatorTok{(}\StringTok{"I\textquotesingle{}m just regular}\SpecialCharTok{\textbackslash{}n}\StringTok{"}\OperatorTok{);}
\PreprocessorTok{\#else}
\NormalTok{    printf}\OperatorTok{(}\StringTok{"I\textquotesingle{}m extra happy!}\SpecialCharTok{\textbackslash{}n}\StringTok{"}\OperatorTok{);}
\PreprocessorTok{\#endif}

\NormalTok{    printf}\OperatorTok{(}\StringTok{"OK!}\SpecialCharTok{\textbackslash{}n}\StringTok{"}\OperatorTok{);}
\OperatorTok{\}}
\end{Highlighting}
\end{Shaded}

\end{tcolorbox}

\index{\#elif directive@\texttt{\#elif} directive|)}

Sekali lagi, untuk klausa \hl{\mbox{\texttt{\#if}}} yang tidak cocok, compiler bahkan tidak akan melihat baris-baris itu. Untuk kode di atas, setelah preprocessor selesai, yang dilihat compiler hanyalah:

\begin{tcolorbox}[enhanced jigsaw,left=3pt,breakable,colback=white!97.2549!black,boxrule=0.2mm]

\begin{Shaded}
\begin{Highlighting}[numbers=left,,]
\PreprocessorTok{\#include }\ImportTok{\textless{}stdio.h\textgreater{}}

\DataTypeTok{int}\NormalTok{ main}\OperatorTok{(}\DataTypeTok{void}\OperatorTok{)}
\OperatorTok{\{}

\NormalTok{    printf}\OperatorTok{(}\StringTok{"I\textquotesingle{}m just regular}\SpecialCharTok{\textbackslash{}n}\StringTok{"}\OperatorTok{);}

\NormalTok{    printf}\OperatorTok{(}\StringTok{"OK!}\SpecialCharTok{\textbackslash{}n}\StringTok{"}\OperatorTok{);}
\OperatorTok{\}}
\end{Highlighting}
\end{Shaded}

\end{tcolorbox}

Salah satu pemakaian akal-akalan dari hal ini adalah mengomentari banyak baris dengan cepat\footnote{Kamu tidak selalu bisa langsung membungkus kode dengan komentar \hl{\mbox{\texttt{/*}}} \hl{\mbox{\texttt{*/}}} karena komentar itu tidak bisa bersarang.}.

\index{\#if 0 directive@\texttt{\#if 0} directive|(}

Kalau kamu menaruh \hl{\mbox{\texttt{\#if\ 0}}} (``if false'') di awal blok yang ingin dikomentari dan \hl{\mbox{\texttt{\#endif}}} di akhir, kamu mendapat efek ini:

\begin{tcolorbox}[enhanced jigsaw,left=3pt,breakable,colback=white!97.2549!black,boxrule=0.2mm]

\begin{Shaded}
\begin{Highlighting}[]
\PreprocessorTok{\#if 0}
\CommentTok{    printf(}\StringTok{"All this code"}\CommentTok{); /* is effectively */}
\CommentTok{    printf(}\StringTok{"commented out"}\CommentTok{); // by the \#if 0}
\PreprocessorTok{\#endif}
\end{Highlighting}
\end{Shaded}

\end{tcolorbox}

\index{\#if 0 directive@\texttt{\#if 0} directive|)}

Bagaimana kalau kamu memakai compiler pra-C23 dan tidak ada dukungan directive \hl{\mbox{\texttt{\#elifdef}}} atau \hl{\mbox{\texttt{\#elifndef}}}? Bagaimana mendapat efek yang sama dengan \hl{\mbox{\texttt{\#if}}}? Artinya, bagaimana kalau kita menginginkan ini:

\begin{tcolorbox}[enhanced jigsaw,left=3pt,breakable,colback=white!97.2549!black,boxrule=0.2mm]

\begin{Shaded}
\begin{Highlighting}[]
\PreprocessorTok{\#ifdef FOO}
\NormalTok{    x }\OperatorTok{=} \DecValTok{2}\OperatorTok{;}
\PreprocessorTok{\#elifdef BAR  }\CommentTok{// KEMUNGKINAN ERROR: Tidak didukung sebelum C23}
\NormalTok{    x }\OperatorTok{=} \DecValTok{3}\OperatorTok{;}
\PreprocessorTok{\#endif}
\end{Highlighting}
\end{Shaded}

\end{tcolorbox}

Bagaimana caranya?

Ternyata ada operator preprocessor bernama \hl{\mbox{\texttt{defined}}} yang bisa dipakai dengan pernyataan \hl{\mbox{\texttt{\#if}}}.

Ketiganya setara:

\index{\#if defined directive@\texttt{\#if defined} directive|(}

\begin{tcolorbox}[enhanced jigsaw,left=3pt,breakable,colback=white!97.2549!black,boxrule=0.2mm]

\begin{Shaded}
\begin{Highlighting}[]
\PreprocessorTok{\#ifdef FOO}
\PreprocessorTok{\#if defined FOO}
\PreprocessorTok{\#if defined(FOO)   }\CommentTok{// Tanda kurung opsional}
\end{Highlighting}
\end{Shaded}

\end{tcolorbox}

Begitu juga ketiga ini:

\begin{tcolorbox}[enhanced jigsaw,left=3pt,breakable,colback=white!97.2549!black,boxrule=0.2mm]

\begin{Shaded}
\begin{Highlighting}[]
\PreprocessorTok{\#ifndef FOO}
\PreprocessorTok{\#if !defined FOO}
\PreprocessorTok{\#if !defined(FOO)   }\CommentTok{// Tanda kurung opsional}
\end{Highlighting}
\end{Shaded}

\end{tcolorbox}

Perhatikan bahwa kita bisa memakai operator logika NOT standar (\hl{\mbox{\texttt{!}}}) untuk ``tidak terdefinisi''.

Jadi sekarang kita kembali ke dunia \hl{\mbox{\texttt{\#if}}} dan bisa memakai \hl{\mbox{\texttt{\#elif}}} sesuka hati!

Kode yang rusak ini:

\begin{tcolorbox}[enhanced jigsaw,left=3pt,breakable,colback=white!97.2549!black,boxrule=0.2mm]

\begin{Shaded}
\begin{Highlighting}[]
\PreprocessorTok{\#ifdef FOO}
\NormalTok{    x }\OperatorTok{=} \DecValTok{2}\OperatorTok{;}
\PreprocessorTok{\#elifdef BAR  }\CommentTok{// KEMUNGKINAN ERROR: Tidak didukung sebelum C23}
\NormalTok{    x }\OperatorTok{=} \DecValTok{3}\OperatorTok{;}
\PreprocessorTok{\#endif}
\end{Highlighting}
\end{Shaded}

\end{tcolorbox}

bisa diganti dengan:

\begin{tcolorbox}[enhanced jigsaw,left=3pt,breakable,colback=white!97.2549!black,boxrule=0.2mm]

\begin{Shaded}
\begin{Highlighting}[]
\PreprocessorTok{\#if defined FOO}
\NormalTok{    x }\OperatorTok{=} \DecValTok{2}\OperatorTok{;}
\PreprocessorTok{\#elif defined BAR}
\NormalTok{    x }\OperatorTok{=} \DecValTok{3}\OperatorTok{;}
\PreprocessorTok{\#endif}
\end{Highlighting}
\end{Shaded}

\end{tcolorbox}

\index{\#if defined directive@\texttt{\#if defined} directive|)} \index{\#if directive@\texttt{\#if} directive|)} \index{Conditional compilation|)}

\subsection{\texorpdfstring{Menghapus Makro: \hl{\mbox{\texttt{\#undef}}}}{Menghapus Makro: \#undef}}\label{menghapus-makro-undef}

\index{\#undef directive@\texttt{\#undef} directive|(}

Kalau kamu sudah mendefinisikan sesuatu tapi tidak membutuhkannya lagi, kamu bisa membatalkan definisinya dengan \hl{\mbox{\texttt{\#undef}}}.

\begin{tcolorbox}[enhanced jigsaw,left=3pt,breakable,colback=white!97.2549!black,boxrule=0.2mm]

\begin{Shaded}
\begin{Highlighting}[numbers=left,,]
\PreprocessorTok{\#include }\ImportTok{\textless{}stdio.h\textgreater{}}

\DataTypeTok{int}\NormalTok{ main}\OperatorTok{(}\DataTypeTok{void}\OperatorTok{)}
\OperatorTok{\{}
\PreprocessorTok{\#define GOATS}

\PreprocessorTok{\#ifdef GOATS}
\NormalTok{    printf}\OperatorTok{(}\StringTok{"Goats detected!}\SpecialCharTok{\textbackslash{}n}\StringTok{"}\OperatorTok{);}  \CommentTok{// tercetak}
\PreprocessorTok{\#endif}

\PreprocessorTok{\#undef GOATS  }\CommentTok{// Buat GOATS tidak terdefinisi lagi}

\PreprocessorTok{\#ifdef GOATS}
\NormalTok{    printf}\OperatorTok{(}\StringTok{"Goats detected, again!}\SpecialCharTok{\textbackslash{}n}\StringTok{"}\OperatorTok{);} \CommentTok{// tidak tercetak}
\PreprocessorTok{\#endif}
\OperatorTok{\}}
\end{Highlighting}
\end{Shaded}

\end{tcolorbox}

\index{\#undef directive@\texttt{\#undef} directive|)}

\section{Makro Bawaan}\label{makro-bawaan}

\index{Preprocessor!predefined macros|(}

Standar mendefinisikan banyak makro bawaan yang bisa diuji dan dipakai untuk kompilasi kondisional. Mari kita lihat di sini.

\subsection{Makro Wajib}\label{makro-wajib}

Semua ini terdefinisi:

\index{\_\_DATE\_\_ macro@\texttt{\_\_DATE\_\_} macro|(} \index{\_\_TIME\_\_ macro@\texttt{\_\_TIME\_\_} macro|(} \index{\_\_FILE\_\_ macro@\texttt{\_\_FILE\_\_} macro|(} \index{\_\_LINE\_\_ macro@\texttt{\_\_LINE\_\_} macro|(} \index{\_\_func\_\_ identifier@\texttt{\_\_func\_\_} identifier|(} \index{\_\_STDC\_VERSION\_\_ macro@\texttt{\_\_STDC\_VERSION\_\_} macro|(}

\index{\_\_STDC\_\_ macro@\texttt{\_\_STDC\_\_} macro} \index{\_\_STDC\_HOSTED\_\_ macro@\texttt{\_\_STDC\_HOSTED\_\_} macro}

{\def\LTcaptype{none} % do not increment counter
\begin{longtable}[]{@{}
  >{\raggedright\arraybackslash}p{(\linewidth - 2\tabcolsep) * \real{0.1746}}
  >{\raggedright\arraybackslash}p{(\linewidth - 2\tabcolsep) * \real{0.8254}}@{}}
\toprule\noalign{}
\begin{minipage}[b]{\linewidth}\raggedright
Makro
\end{minipage} & \begin{minipage}[b]{\linewidth}\raggedright
Deskripsi
\end{minipage} \\
\midrule\noalign{}
\endhead
\bottomrule\noalign{}
\endlastfoot
\hl{\mbox{\texttt{\_\_DATE\_\_}}} & Tanggal kompilasi---yaitu saat file ini dikompilasi---dengan format \hl{\mbox{\texttt{Mmm\ dd\ yyyy}}} \\
\hl{\mbox{\texttt{\_\_TIME\_\_}}} & Waktu kompilasi dengan format \hl{\mbox{\texttt{hh:mm:ss}}} \\
\hl{\mbox{\texttt{\_\_FILE\_\_}}} & String berisi nama file ini \\
\hl{\mbox{\texttt{\_\_LINE\_\_}}} & Nomor baris di file tempat makro ini muncul \\
\hl{\mbox{\texttt{\_\_func\_\_}}} & Nama fungsi tempat ini muncul, sebagai string\footnote{Ini sebenarnya bukan makro---secara teknis ia identifier. Tapi ia satu-satunya identifier bawaan dan terasa sangat seperti makro, jadi kumasukkan di sini. Seperti pemberontak.} \\
\hl{\mbox{\texttt{\_\_STDC\_\_}}} & Bernilai \hl{\mbox{\texttt{1}}} kalau ini C compiler standar \\
\hl{\mbox{\texttt{\_\_STDC\_HOSTED\_\_}}} & Bernilai \hl{\mbox{\texttt{1}}} kalau compiler adalah \emph{hosted implementation}\footnote{Hosted implementation pada dasarnya berarti kamu menjalankan standar C yang lengkap, kemungkinan di atas semacam sistem operasi. Dan kemungkinan besar memang begitu keadaanmu. Kalau kamu berjalan di bare metal pada semacam embedded system, kemungkinan kamu memakai \emph{standalone implementation}.}, selain itu \hl{\mbox{\texttt{0}}} \\
\hl{\mbox{\texttt{\_\_STDC\_VERSION\_\_}}} & Versi C ini, berupa konstanta \hl{\mbox{\texttt{long\ int}}} dengan bentuk \hl{\mbox{\texttt{yyyymmL}}}, mis. \hl{\mbox{\texttt{201710L}}} \\
\end{longtable}
}

Mari kita gabungkan semuanya.

\begin{tcolorbox}[enhanced jigsaw,left=3pt,breakable,colback=white!97.2549!black,boxrule=0.2mm]

\begin{Shaded}
\begin{Highlighting}[numbers=left,,]
\PreprocessorTok{\#include }\ImportTok{\textless{}stdio.h\textgreater{}}

\DataTypeTok{int}\NormalTok{ main}\OperatorTok{(}\DataTypeTok{void}\OperatorTok{)}
\OperatorTok{\{}
\NormalTok{    printf}\OperatorTok{(}\StringTok{"This function: }\SpecialCharTok{\%s\textbackslash{}n}\StringTok{"}\OperatorTok{,}\NormalTok{ \_\_func\_\_}\OperatorTok{);}
\NormalTok{    printf}\OperatorTok{(}\StringTok{"This file: }\SpecialCharTok{\%s\textbackslash{}n}\StringTok{"}\OperatorTok{,}\NormalTok{ \_\_FILE\_\_}\OperatorTok{);}
\NormalTok{    printf}\OperatorTok{(}\StringTok{"This line: }\SpecialCharTok{\%d\textbackslash{}n}\StringTok{"}\OperatorTok{,}\NormalTok{ \_\_LINE\_\_}\OperatorTok{);}
\NormalTok{    printf}\OperatorTok{(}\StringTok{"Compiled on: }\SpecialCharTok{\%s}\StringTok{ }\SpecialCharTok{\%s\textbackslash{}n}\StringTok{"}\OperatorTok{,}\NormalTok{ \_\_DATE\_\_}\OperatorTok{,}\NormalTok{ \_\_TIME\_\_}\OperatorTok{);}
\NormalTok{    printf}\OperatorTok{(}\StringTok{"C Version: }\SpecialCharTok{\%ld\textbackslash{}n}\StringTok{"}\OperatorTok{,}\NormalTok{ \_\_STDC\_VERSION\_\_}\OperatorTok{);}
\OperatorTok{\}}
\end{Highlighting}
\end{Shaded}

\end{tcolorbox}

\index{\_\_DATE\_\_ macro@\texttt{\_\_DATE\_\_} macro|)} \index{\_\_TIME\_\_ macro@\texttt{\_\_TIME\_\_} macro|)}

Output di sistemku:

\begin{tcolorbox}[enhanced jigsaw,left=3pt,breakable,colback=white!97.2549!black,boxrule=0.2mm]

\begin{Shaded}
\begin{Highlighting}[]
\NormalTok{This function: main}
\NormalTok{This file: foo.c}
\NormalTok{This line: 7}
\NormalTok{Compiled on: Nov 23 2020 17:16:27}
\NormalTok{C Version: 201710}
\end{Highlighting}
\end{Shaded}

\end{tcolorbox}

\hl{\mbox{\texttt{\_\_FILE\_\_}}}, \hl{\mbox{\texttt{\_\_func\_\_}}} dan \hl{\mbox{\texttt{\_\_LINE\_\_}}} sangat berguna untuk melaporkan kondisi error dalam pesan kepada developer. Makro \hl{\mbox{\texttt{assert()}}} di \hl{\mbox{\texttt{\textless{}assert.h\textgreater{}}}} memakai ketiganya untuk menunjukkan di bagian kode mana assertion gagal.

\index{\_\_FILE\_\_ macro@\texttt{\_\_FILE\_\_} macro|)} \index{\_\_LINE\_\_ macro@\texttt{\_\_LINE\_\_} macro|)} \index{\_\_func\_\_ identifier@\texttt{\_\_func\_\_} identifier|)}

\subsubsection{\texorpdfstring{\hl{\mbox{\texttt{\_\_STDC\_VERSION\_\_}}}}{\_\_STDC\_VERSION\_\_}}\label{stdc_version__}

\index{Language versions|(}

Kalau kamu penasaran, berikut nomor versi untuk rilis-rilis utama Spesifikasi Bahasa C:

{\def\LTcaptype{none} % do not increment counter
\begin{longtable}[]{@{}lll@{}}
\toprule\noalign{}
Rilis & Versi ISO/IEC & \hl{\mbox{\texttt{\_\_STDC\_VERSION\_\_}}} \\
\midrule\noalign{}
\endhead
\bottomrule\noalign{}
\endlastfoot
C89 & ISO/IEC 9899:1990 & tidak terdefinisi \\
\textbf{C89} & ISO/IEC 9899:1990/Amd.1:1995 & \hl{\mbox{\texttt{199409L}}} \\
\textbf{C99} & ISO/IEC 9899:1999 & \hl{\mbox{\texttt{199901L}}} \\
\textbf{C11} & ISO/IEC 9899:2011/Amd.1:2012 & \hl{\mbox{\texttt{201112L}}} \\
\end{longtable}
}

Perhatikan bahwa makro ini awalnya tidak ada di C89.

Perhatikan juga bahwa rencananya nomor versi akan selalu naik, jadi kamu selalu bisa memeriksa, misalnya, ``setidaknya C99'' dengan:

\begin{tcolorbox}[enhanced jigsaw,left=3pt,breakable,colback=white!97.2549!black,boxrule=0.2mm]

\begin{Shaded}
\begin{Highlighting}[]
\PreprocessorTok{\#if \_\_STDC\_VERSION\_\_ \textgreater{}= 1999901L}
\end{Highlighting}
\end{Shaded}

\end{tcolorbox}

\index{Language versions|)} \index{\_\_STDC\_VERSION\_\_ macro@\texttt{\_\_STDC\_VERSION\_\_} macro|)}

\subsection{Makro Opsional}\label{makro-opsional}

Implementasimu mungkin mendefinisikan ini juga. Atau mungkin tidak.

\index{\_\_STDC\_ISO\_10646\_\_ macro@\texttt{\_\_STDC\_ISO\_10646\_\_} macro} \index{\_\_STDC\_MB\_MIGHT\_NEQ\_WC\_\_ macro@\texttt{\_\_STDC\_MB\_MIGHT\_NEQ\_WC\_\_} macro} \index{\_\_STDC\_UTF\_16\_\_ macro@\texttt{\_\_STDC\_UTF\_16\_\_} macro} \index{\_\_STDC\_UTF\_32\_\_ macro@\texttt{\_\_STDC\_UTF\_32\_\_} macro} \index{\_\_STDC\_ANALYZABLE\_\_ macro@\texttt{\_\_STDC\_ANALYZABLE\_\_} macro} \index{\_\_STDC\_IEC\_559\_\_ macro@\texttt{\_\_STDC\_IEC\_559\_\_} macro} \index{\_\_STDC\_IEC\_559\_COMPLEX\_\_ macro@\texttt{\_\_STDC\_IEC\_559\_COMPLEX\_\_} macro} \index{\_\_STDC\_LIB\_EXT1\_\_ macro@\texttt{\_\_STDC\_LIB\_EXT1\_\_} macro} \index{\_\_STDC\_NO\_ATOMICS\_\_ macro@\texttt{\_\_STDC\_NO\_ATOMICS\_\_} macro} \index{\_\_STDC\_NO\_COMPLEX\_\_ macro@\texttt{\_\_STDC\_NO\_COMPLEX\_\_} macro} \index{\_\_STDC\_NO\_THREADS\_\_ macro@\texttt{\_\_STDC\_NO\_THREADS\_\_} macro} \index{\_\_STDC\_NO\_VLA\_\_ macro@\texttt{\_\_STDC\_NO\_VLA\_\_} macro}

{\def\LTcaptype{none} % do not increment counter
\begin{longtable}[]{@{}
  >{\raggedright\arraybackslash}p{(\linewidth - 2\tabcolsep) * \real{0.2424}}
  >{\raggedright\arraybackslash}p{(\linewidth - 2\tabcolsep) * \real{0.7576}}@{}}
\toprule\noalign{}
\begin{minipage}[b]{\linewidth}\raggedright
Makro
\end{minipage} & \begin{minipage}[b]{\linewidth}\raggedright
Deskripsi
\end{minipage} \\
\midrule\noalign{}
\endhead
\bottomrule\noalign{}
\endlastfoot
\hl{\mbox{\texttt{\_\_STDC\_ISO\_10646\_\_}}} & Jika terdefinisi, \hl{\mbox{\texttt{wchar\_t}}} menyimpan nilai Unicode, selain itu sesuatu yang lain \\
\hl{\mbox{\texttt{\_\_STDC\_MB\_MIGHT\_NEQ\_WC\_\_}}} & Nilai \hl{\mbox{\texttt{1}}} menandakan bahwa nilai pada karakter multibyte mungkin tidak sama dengan nilai pada wide character \\
\hl{\mbox{\texttt{\_\_STDC\_UTF\_16\_\_}}} & Nilai \hl{\mbox{\texttt{1}}} menandakan bahwa sistem memakai encoding UTF-16 pada tipe \hl{\mbox{\texttt{char16\_t}}} \\
\hl{\mbox{\texttt{\_\_STDC\_UTF\_32\_\_}}} & Nilai \hl{\mbox{\texttt{1}}} menandakan bahwa sistem memakai encoding UTF-32 pada tipe \hl{\mbox{\texttt{char32\_t}}} \\
\hl{\mbox{\texttt{\_\_STDC\_ANALYZABLE\_\_}}} & Nilai \hl{\mbox{\texttt{1}}} menandakan bahwa kode itu analyzable\footnote{Oke, aku tahu itu jawaban yang mengelak. Pada dasarnya ada ekstensi opsional yang bisa diimplementasikan compiler, yaitu mereka setuju membatasi jenis undefined behavior tertentu supaya kode C lebih mudah dianalisis secara statis. Kecil kemungkinan kamu perlu memakainya.} \\
\hl{\mbox{\texttt{\_\_STDC\_IEC\_559\_\_}}} & \hl{\mbox{\texttt{1}}} jika floating point IEEE-754 (alias IEC 60559) didukung \\
\hl{\mbox{\texttt{\_\_STDC\_IEC\_559\_COMPLEX\_\_}}} & \hl{\mbox{\texttt{1}}} jika floating point kompleks IEC 60559 didukung \\
\hl{\mbox{\texttt{\_\_STDC\_LIB\_EXT1\_\_}}} & \hl{\mbox{\texttt{1}}} jika implementasi ini mendukung berbagai fungsi standard library alternatif yang ``aman'' (nama mereka berakhiran \hl{\mbox{\texttt{\_s}}}) \\
\hl{\mbox{\texttt{\_\_STDC\_NO\_ATOMICS\_\_}}} & \hl{\mbox{\texttt{1}}} jika implementasi ini \textbf{tidak} mendukung \hl{\mbox{\texttt{\_Atomic}}} atau \hl{\mbox{\texttt{\textless{}stdatomic.h\textgreater{}}}} \\
\hl{\mbox{\texttt{\_\_STDC\_NO\_COMPLEX\_\_}}} & \hl{\mbox{\texttt{1}}} jika implementasi ini \textbf{tidak} mendukung tipe kompleks atau \hl{\mbox{\texttt{\textless{}complex.h\textgreater{}}}} \\
\hl{\mbox{\texttt{\_\_STDC\_NO\_THREADS\_\_}}} & \hl{\mbox{\texttt{1}}} jika implementasi ini \textbf{tidak} mendukung \hl{\mbox{\texttt{\textless{}threads.h\textgreater{}}}} \\
\hl{\mbox{\texttt{\_\_STDC\_NO\_VLA\_\_}}} & \hl{\mbox{\texttt{1}}} jika implementasi ini \textbf{tidak} mendukung variable-length array \\
\end{longtable}
}

\index{Preprocessor!predefined macros|)}

\section{Makro dengan Argumen}\label{makro-dengan-argumen}

\index{Preprocessor!macros with arguments|(}

Makro lebih kuat daripada sekadar substitusi. Kamu bisa mengaturnya agar menerima argumen yang ikut disubstitusi.

Pertanyaan yang sering muncul: kapan memakai makro berparameter versus fungsi. Jawaban singkat: pakai fungsi. Tapi kamu akan melihat banyak makro di alam liar dan di standard library. Orang cenderung memakainya untuk hal-hal singkat yang bersifat matematis, dan juga untuk fitur yang mungkin berubah dari satu platform ke platform lain. Kamu bisa mendefinisikan kata kunci berbeda untuk satu platform atau lainnya.

\subsection{Makro dengan Satu Argumen}\label{makro-dengan-satu-argumen}

Kita mulai dengan yang sederhana, yang mengkuadratkan sebuah angka:

\index{\#define directive@\texttt{\#define} directive|(}

\begin{tcolorbox}[enhanced jigsaw,left=3pt,breakable,colback=white!97.2549!black,boxrule=0.2mm]

\begin{Shaded}
\begin{Highlighting}[numbers=left,,]
\PreprocessorTok{\#include }\ImportTok{\textless{}stdio.h\textgreater{}}

\PreprocessorTok{\#define SQR}\OperatorTok{(}\PreprocessorTok{x}\OperatorTok{)}\PreprocessorTok{ x }\OperatorTok{*}\PreprocessorTok{ x  }\CommentTok{// Belum tepat, tapi ikuti saja dulu}

\DataTypeTok{int}\NormalTok{ main}\OperatorTok{(}\DataTypeTok{void}\OperatorTok{)}
\OperatorTok{\{}
\NormalTok{    printf}\OperatorTok{(}\StringTok{"}\SpecialCharTok{\%d\textbackslash{}n}\StringTok{"}\OperatorTok{,}\NormalTok{ SQR}\OperatorTok{(}\DecValTok{12}\OperatorTok{));}  \CommentTok{// 144}
\OperatorTok{\}}
\end{Highlighting}
\end{Shaded}

\end{tcolorbox}

Maksudnya: ``di mana pun kamu melihat \hl{\mbox{\texttt{SQR}}} dengan suatu nilai, ganti dengan nilai itu dikali dirinya sendiri''.

Jadi baris 7 akan diubah menjadi:

\begin{tcolorbox}[enhanced jigsaw,left=3pt,breakable,colback=white!97.2549!black,boxrule=0.2mm]

\begin{Shaded}
\begin{Highlighting}[numbers=left,,firstnumber=7,]
\NormalTok{    printf}\OperatorTok{(}\StringTok{"}\SpecialCharTok{\%d\textbackslash{}n}\StringTok{"}\OperatorTok{,} \DecValTok{12} \OperatorTok{*} \DecValTok{12}\OperatorTok{);}  \CommentTok{// 144}
\end{Highlighting}
\end{Shaded}

\end{tcolorbox}

yang dengan nyaman dikonversi C menjadi 144.

Tapi kita membuat kesalahan elementer pada makro itu, kesalahan yang harus dihindari.

Mari kita periksa. Bagaimana kalau kita ingin menghitung \hl{\mbox{\texttt{SQR(3\ +\ 4)}}}? Nah, \(3+4=7\), jadi tentu kita ingin menghitung \(7^2=49\). Itu dia; \hl{\mbox{\texttt{49}}}--- jawaban akhir.

Mari kita masukkan ke kode dan lihat bahwa kita mendapat\ldots{} 19?

\begin{tcolorbox}[enhanced jigsaw,left=3pt,breakable,colback=white!97.2549!black,boxrule=0.2mm]

\begin{Shaded}
\begin{Highlighting}[numbers=left,,firstnumber=7,]
\NormalTok{    printf}\OperatorTok{(}\StringTok{"}\SpecialCharTok{\%d\textbackslash{}n}\StringTok{"}\OperatorTok{,}\NormalTok{ SQR}\OperatorTok{(}\DecValTok{3} \OperatorTok{+} \DecValTok{4}\OperatorTok{));}  \CommentTok{// 19!!??}
\end{Highlighting}
\end{Shaded}

\end{tcolorbox}

Apa yang terjadi?

Kalau kita ikuti perluasan makronya, kita mendapat

\begin{tcolorbox}[enhanced jigsaw,left=3pt,breakable,colback=white!97.2549!black,boxrule=0.2mm]

\begin{Shaded}
\begin{Highlighting}[numbers=left,,firstnumber=7,]
\NormalTok{    printf}\OperatorTok{(}\StringTok{"}\SpecialCharTok{\%d\textbackslash{}n}\StringTok{"}\OperatorTok{,} \DecValTok{3} \OperatorTok{+} \DecValTok{4} \OperatorTok{*} \DecValTok{3} \OperatorTok{+} \DecValTok{4}\OperatorTok{);}  \CommentTok{// 19!}
\end{Highlighting}
\end{Shaded}

\end{tcolorbox}

Waduh! Karena perkalian lebih diutamakan, kita menghitung \(4\times3=12\) lebih dulu, dan mendapat \(3+12+4=19\). Bukan yang kita inginkan.

Jadi kita harus memperbaikinya agar benar.

\textbf{Hal ini sangat umum, jadi lakukan saja secara otomatis setiap kali membuat makro matematika berparameter!}

Perbaikannya mudah: tambahkan tanda kurung!

\begin{tcolorbox}[enhanced jigsaw,left=3pt,breakable,colback=white!97.2549!black,boxrule=0.2mm]

\begin{Shaded}
\begin{Highlighting}[numbers=left,,firstnumber=3,]
\PreprocessorTok{\#define SQR}\OperatorTok{(}\PreprocessorTok{x}\OperatorTok{)}\PreprocessorTok{ }\OperatorTok{(}\PreprocessorTok{x}\OperatorTok{)}\PreprocessorTok{ }\OperatorTok{*}\PreprocessorTok{ }\OperatorTok{(}\PreprocessorTok{x}\OperatorTok{)}\PreprocessorTok{   }\CommentTok{// Lebih baik... tapi masih belum cukup bagus!}
\end{Highlighting}
\end{Shaded}

\end{tcolorbox}

Dan sekarang makro kita diperluas menjadi:

\begin{tcolorbox}[enhanced jigsaw,left=3pt,breakable,colback=white!97.2549!black,boxrule=0.2mm]

\begin{Shaded}
\begin{Highlighting}[numbers=left,,firstnumber=7,]
\NormalTok{    printf}\OperatorTok{(}\StringTok{"}\SpecialCharTok{\%d\textbackslash{}n}\StringTok{"}\OperatorTok{,} \OperatorTok{(}\DecValTok{3} \OperatorTok{+} \DecValTok{4}\OperatorTok{)} \OperatorTok{*} \OperatorTok{(}\DecValTok{3} \OperatorTok{+} \DecValTok{4}\OperatorTok{));}  \CommentTok{// 49! Horeee!}
\end{Highlighting}
\end{Shaded}

\end{tcolorbox}

Tapi sebenarnya kita masih punya masalah yang sama, yang bisa muncul kalau ada operator berprioritas lebih tinggi daripada perkalian (\hl{\mbox{\texttt{*}}}) di dekatnya.

Jadi cara yang aman dan benar untuk menyusun makro itu adalah dengan membungkus keseluruhannya dalam tanda kurung tambahan, seperti ini:

\begin{tcolorbox}[enhanced jigsaw,left=3pt,breakable,colback=white!97.2549!black,boxrule=0.2mm]

\begin{Shaded}
\begin{Highlighting}[numbers=left,,firstnumber=3,]
\PreprocessorTok{\#define SQR}\OperatorTok{(}\PreprocessorTok{x}\OperatorTok{)}\PreprocessorTok{ }\OperatorTok{((}\PreprocessorTok{x}\OperatorTok{)}\PreprocessorTok{ }\OperatorTok{*}\PreprocessorTok{ }\OperatorTok{(}\PreprocessorTok{x}\OperatorTok{))}\PreprocessorTok{   }\CommentTok{// Bagus!}
\end{Highlighting}
\end{Shaded}

\end{tcolorbox}

Biasakan melakukan itu setiap kali membuat makro matematika dan kamu tidak akan salah.

\subsection{Makro dengan Lebih dari Satu Argumen}\label{makro-dengan-lebih-dari-satu-argumen}

Kamu bisa menumpuk ini sebanyak yang kamu mau:

\begin{tcolorbox}[enhanced jigsaw,left=3pt,breakable,colback=white!97.2549!black,boxrule=0.2mm]

\begin{Shaded}
\begin{Highlighting}[]
\PreprocessorTok{\#define TRIANGLE\_AREA}\OperatorTok{(}\PreprocessorTok{w}\OperatorTok{,}\PreprocessorTok{ h}\OperatorTok{)}\PreprocessorTok{ }\OperatorTok{(}\FloatTok{0.5}\PreprocessorTok{ }\OperatorTok{*}\PreprocessorTok{ }\OperatorTok{(}\PreprocessorTok{w}\OperatorTok{)}\PreprocessorTok{ }\OperatorTok{*}\PreprocessorTok{ }\OperatorTok{(}\PreprocessorTok{h}\OperatorTok{))}
\end{Highlighting}
\end{Shaded}

\end{tcolorbox}

Mari kita buat beberapa makro yang menyelesaikan \(x\) dengan rumus kuadrat. Kalau-kalau kamu belum hafal, untuk persamaan berbentuk:

\(ax^2+bx+c=0\)

kamu bisa mencari \(x\) dengan rumus kuadrat:

\(x=\displaystyle\frac{-b\pm\sqrt{b^2-4ac}}{2a}\)

Rumus yang gila. Perhatikan juga plus-minus (\(\pm\)) di sana, yang menandakan bahwa sebenarnya ada dua solusi.

Jadi mari kita buat makro untuk keduanya:

\begin{tcolorbox}[enhanced jigsaw,left=3pt,breakable,colback=white!97.2549!black,boxrule=0.2mm]

\begin{Shaded}
\begin{Highlighting}[]
\PreprocessorTok{\#define QUADP}\OperatorTok{(}\PreprocessorTok{a}\OperatorTok{,}\PreprocessorTok{ b}\OperatorTok{,}\PreprocessorTok{ c}\OperatorTok{)}\PreprocessorTok{ }\OperatorTok{(({-}(}\PreprocessorTok{b}\OperatorTok{)}\PreprocessorTok{ }\OperatorTok{+}\PreprocessorTok{ sqrt}\OperatorTok{((}\PreprocessorTok{b}\OperatorTok{)*(}\PreprocessorTok{b}\OperatorTok{)}\PreprocessorTok{ }\OperatorTok{{-}}\PreprocessorTok{ }\DecValTok{4}\PreprocessorTok{ }\OperatorTok{*}\PreprocessorTok{ }\OperatorTok{(}\PreprocessorTok{a}\OperatorTok{)}\PreprocessorTok{ }\OperatorTok{*}\PreprocessorTok{ }\OperatorTok{(}\PreprocessorTok{c}\OperatorTok{)))}\PreprocessorTok{ }\OperatorTok{/}\PreprocessorTok{ }\OperatorTok{(}\DecValTok{2}\PreprocessorTok{ }\OperatorTok{*}\PreprocessorTok{ }\OperatorTok{(}\PreprocessorTok{a}\OperatorTok{)))}
\PreprocessorTok{\#define QUADM}\OperatorTok{(}\PreprocessorTok{a}\OperatorTok{,}\PreprocessorTok{ b}\OperatorTok{,}\PreprocessorTok{ c}\OperatorTok{)}\PreprocessorTok{ }\OperatorTok{(({-}(}\PreprocessorTok{b}\OperatorTok{)}\PreprocessorTok{ }\OperatorTok{{-}}\PreprocessorTok{ sqrt}\OperatorTok{((}\PreprocessorTok{b}\OperatorTok{)*(}\PreprocessorTok{b}\OperatorTok{)}\PreprocessorTok{ }\OperatorTok{{-}}\PreprocessorTok{ }\DecValTok{4}\PreprocessorTok{ }\OperatorTok{*}\PreprocessorTok{ }\OperatorTok{(}\PreprocessorTok{a}\OperatorTok{)}\PreprocessorTok{ }\OperatorTok{*}\PreprocessorTok{ }\OperatorTok{(}\PreprocessorTok{c}\OperatorTok{)))}\PreprocessorTok{ }\OperatorTok{/}\PreprocessorTok{ }\OperatorTok{(}\DecValTok{2}\PreprocessorTok{ }\OperatorTok{*}\PreprocessorTok{ }\OperatorTok{(}\PreprocessorTok{a}\OperatorTok{)))}
\end{Highlighting}
\end{Shaded}

\end{tcolorbox}

Itu sudah memberi kita sedikit matematika. Tapi mari kita definisikan satu lagi yang bisa dipakai sebagai argumen \hl{\mbox{\texttt{printf()}}} untuk mencetak kedua jawaban.

\begin{tcolorbox}[enhanced jigsaw,left=3pt,breakable,colback=white!97.2549!black,boxrule=0.2mm]

\begin{Shaded}
\begin{Highlighting}[]
\CommentTok{//          makro              penggantian}
\CommentTok{//      |{-}{-}{-}{-}{-}{-}{-}{-}{-}{-}{-}| |{-}{-}{-}{-}{-}{-}{-}{-}{-}{-}{-}{-}{-}{-}{-}{-}{-}{-}{-}{-}{-}{-}{-}{-}{-}{-}{-}{-}|}
\PreprocessorTok{\#define QUAD}\OperatorTok{(}\PreprocessorTok{a}\OperatorTok{,}\PreprocessorTok{ b}\OperatorTok{,}\PreprocessorTok{ c}\OperatorTok{)}\PreprocessorTok{ QUADP}\OperatorTok{(}\PreprocessorTok{a}\OperatorTok{,}\PreprocessorTok{ b}\OperatorTok{,}\PreprocessorTok{ c}\OperatorTok{),}\PreprocessorTok{ QUADM}\OperatorTok{(}\PreprocessorTok{a}\OperatorTok{,}\PreprocessorTok{ b}\OperatorTok{,}\PreprocessorTok{ c}\OperatorTok{)}
\end{Highlighting}
\end{Shaded}

\end{tcolorbox}

Itu hanya dua nilai yang dipisahkan koma---dan kita bisa memakainya sebagai argumen ``gabungan'' untuk \hl{\mbox{\texttt{printf()}}} seperti ini:

\begin{tcolorbox}[enhanced jigsaw,left=3pt,breakable,colback=white!97.2549!black,boxrule=0.2mm]

\begin{Shaded}
\begin{Highlighting}[]
\NormalTok{printf}\OperatorTok{(}\StringTok{"x = }\SpecialCharTok{\%f}\StringTok{ or x = }\SpecialCharTok{\%f\textbackslash{}n}\StringTok{"}\OperatorTok{,}\NormalTok{ QUAD}\OperatorTok{(}\DecValTok{2}\OperatorTok{,} \DecValTok{10}\OperatorTok{,} \DecValTok{5}\OperatorTok{));}
\end{Highlighting}
\end{Shaded}

\end{tcolorbox}

Mari kita satukan dalam kode:

\begin{tcolorbox}[enhanced jigsaw,left=3pt,breakable,colback=white!97.2549!black,boxrule=0.2mm]

\begin{Shaded}
\begin{Highlighting}[numbers=left,,]
\PreprocessorTok{\#include }\ImportTok{\textless{}stdio.h\textgreater{}}
\PreprocessorTok{\#include }\ImportTok{\textless{}math.h\textgreater{}}\PreprocessorTok{  }\CommentTok{// Untuk sqrt()}

\PreprocessorTok{\#define QUADP}\OperatorTok{(}\PreprocessorTok{a}\OperatorTok{,}\PreprocessorTok{ b}\OperatorTok{,}\PreprocessorTok{ c}\OperatorTok{)}\PreprocessorTok{ }\OperatorTok{(({-}(}\PreprocessorTok{b}\OperatorTok{)}\PreprocessorTok{ }\OperatorTok{+}\PreprocessorTok{ sqrt}\OperatorTok{((}\PreprocessorTok{b}\OperatorTok{)}\PreprocessorTok{ }\OperatorTok{*}\PreprocessorTok{ }\OperatorTok{(}\PreprocessorTok{b}\OperatorTok{)}\PreprocessorTok{ }\OperatorTok{{-}}\PreprocessorTok{ }\DecValTok{4}\PreprocessorTok{ }\OperatorTok{*}\PreprocessorTok{ }\OperatorTok{(}\PreprocessorTok{a}\OperatorTok{)}\PreprocessorTok{ }\OperatorTok{*}\PreprocessorTok{ }\OperatorTok{(}\PreprocessorTok{c}\OperatorTok{)))}\PreprocessorTok{ }\OperatorTok{/}\PreprocessorTok{ }\OperatorTok{(}\DecValTok{2}\PreprocessorTok{ }\OperatorTok{*}\PreprocessorTok{ }\OperatorTok{(}\PreprocessorTok{a}\OperatorTok{)))}
\PreprocessorTok{\#define QUADM}\OperatorTok{(}\PreprocessorTok{a}\OperatorTok{,}\PreprocessorTok{ b}\OperatorTok{,}\PreprocessorTok{ c}\OperatorTok{)}\PreprocessorTok{ }\OperatorTok{(({-}(}\PreprocessorTok{b}\OperatorTok{)}\PreprocessorTok{ }\OperatorTok{{-}}\PreprocessorTok{ sqrt}\OperatorTok{((}\PreprocessorTok{b}\OperatorTok{)}\PreprocessorTok{ }\OperatorTok{*}\PreprocessorTok{ }\OperatorTok{(}\PreprocessorTok{b}\OperatorTok{)}\PreprocessorTok{ }\OperatorTok{{-}}\PreprocessorTok{ }\DecValTok{4}\PreprocessorTok{ }\OperatorTok{*}\PreprocessorTok{ }\OperatorTok{(}\PreprocessorTok{a}\OperatorTok{)}\PreprocessorTok{ }\OperatorTok{*}\PreprocessorTok{ }\OperatorTok{(}\PreprocessorTok{c}\OperatorTok{)))}\PreprocessorTok{ }\OperatorTok{/}\PreprocessorTok{ }\OperatorTok{(}\DecValTok{2}\PreprocessorTok{ }\OperatorTok{*}\PreprocessorTok{ }\OperatorTok{(}\PreprocessorTok{a}\OperatorTok{)))}
\PreprocessorTok{\#define QUAD}\OperatorTok{(}\PreprocessorTok{a}\OperatorTok{,}\PreprocessorTok{ b}\OperatorTok{,}\PreprocessorTok{ c}\OperatorTok{)}\PreprocessorTok{ QUADP}\OperatorTok{(}\PreprocessorTok{a}\OperatorTok{,}\PreprocessorTok{ b}\OperatorTok{,}\PreprocessorTok{ c}\OperatorTok{),}\PreprocessorTok{ QUADM}\OperatorTok{(}\PreprocessorTok{a}\OperatorTok{,}\PreprocessorTok{ b}\OperatorTok{,}\PreprocessorTok{ c}\OperatorTok{)}

\DataTypeTok{int}\NormalTok{ main}\OperatorTok{(}\DataTypeTok{void}\OperatorTok{)}
\OperatorTok{\{}
\NormalTok{    printf}\OperatorTok{(}\StringTok{"2*x\^{}2 + 10*x + 5 = 0}\SpecialCharTok{\textbackslash{}n}\StringTok{"}\OperatorTok{);}
\NormalTok{    printf}\OperatorTok{(}\StringTok{"x = }\SpecialCharTok{\%f}\StringTok{ or x = }\SpecialCharTok{\%f\textbackslash{}n}\StringTok{"}\OperatorTok{,}\NormalTok{ QUAD}\OperatorTok{(}\DecValTok{2}\OperatorTok{,} \DecValTok{10}\OperatorTok{,} \DecValTok{5}\OperatorTok{));}
\OperatorTok{\}}
\end{Highlighting}
\end{Shaded}

\end{tcolorbox}

Dan ini memberi kita output:

\begin{tcolorbox}[enhanced jigsaw,left=3pt,breakable,colback=white!97.2549!black,boxrule=0.2mm]

\begin{Shaded}
\begin{Highlighting}[]
\NormalTok{2*x\^{}2 + 10*x + 5 = 0}
\NormalTok{x = {-}0.563508 or x = {-}4.436492}
\end{Highlighting}
\end{Shaded}

\end{tcolorbox}

Memasukkan salah satu nilai itu memberi hasil kira-kira nol (agak meleset karena angkanya tidak persis):

\(2\times-0.563508^2+10\times-0.563508+5\approx0.000003\)

\subsection{Makro dengan Argumen Variabel}\label{makro-dengan-argumen-variabel}

\index{Preprocessor!macros with variable arguments|(}

Ada juga cara untuk mengoper jumlah argumen yang bervariasi ke makro, dengan elipsis (\hl{\mbox{\texttt{...}}}) setelah argumen bernama yang sudah diketahui. Saat makro diperluas, semua argumen tambahan akan berada dalam daftar yang dipisahkan koma di makro \hl{\mbox{\texttt{\_\_VA\_ARGS\_\_}}}, dan bisa diganti dari sana:

\begin{tcolorbox}[enhanced jigsaw,left=3pt,breakable,colback=white!97.2549!black,boxrule=0.2mm]

\begin{Shaded}
\begin{Highlighting}[numbers=left,,]
\PreprocessorTok{\#include }\ImportTok{\textless{}stdio.h\textgreater{}}

\CommentTok{// Gabungkan dua argumen pertama menjadi satu angka,}
\CommentTok{// lalu jadikan sisanya daftar yang dipisahkan koma:}

\PreprocessorTok{\#define X}\OperatorTok{(}\PreprocessorTok{a}\OperatorTok{,}\PreprocessorTok{ b}\OperatorTok{,}\PreprocessorTok{ }\OperatorTok{...)}\PreprocessorTok{ }\OperatorTok{(}\DecValTok{10}\OperatorTok{*(}\PreprocessorTok{a}\OperatorTok{)}\PreprocessorTok{ }\OperatorTok{+}\PreprocessorTok{ }\DecValTok{20}\OperatorTok{*(}\PreprocessorTok{b}\OperatorTok{)),}\PreprocessorTok{ \_\_VA\_ARGS\_\_}

\DataTypeTok{int}\NormalTok{ main}\OperatorTok{(}\DataTypeTok{void}\OperatorTok{)}
\OperatorTok{\{}
\NormalTok{    printf}\OperatorTok{(}\StringTok{"}\SpecialCharTok{\%d}\StringTok{ }\SpecialCharTok{\%f}\StringTok{ }\SpecialCharTok{\%s}\StringTok{ }\SpecialCharTok{\%d\textbackslash{}n}\StringTok{"}\OperatorTok{,}\NormalTok{ X}\OperatorTok{(}\DecValTok{5}\OperatorTok{,} \DecValTok{4}\OperatorTok{,} \FloatTok{3.14}\OperatorTok{,} \StringTok{"Hi!"}\OperatorTok{,} \DecValTok{12}\OperatorTok{));}
\OperatorTok{\}}
\end{Highlighting}
\end{Shaded}

\end{tcolorbox}

Substitusi yang terjadi di baris 10 adalah:

\begin{tcolorbox}[enhanced jigsaw,left=3pt,breakable,colback=white!97.2549!black,boxrule=0.2mm]

\begin{Shaded}
\begin{Highlighting}[numbers=left,,firstnumber=10,]
\NormalTok{    printf}\OperatorTok{(}\StringTok{"}\SpecialCharTok{\%d}\StringTok{ }\SpecialCharTok{\%f}\StringTok{ }\SpecialCharTok{\%s}\StringTok{ }\SpecialCharTok{\%d\textbackslash{}n}\StringTok{"}\OperatorTok{,} \OperatorTok{(}\DecValTok{10}\OperatorTok{*(}\DecValTok{5}\OperatorTok{)} \OperatorTok{+} \DecValTok{20}\OperatorTok{*(}\DecValTok{4}\OperatorTok{)),} \FloatTok{3.14}\OperatorTok{,} \StringTok{"Hi!"}\OperatorTok{,} \DecValTok{12}\OperatorTok{);}
\end{Highlighting}
\end{Shaded}

\end{tcolorbox}

dengan output:

\begin{tcolorbox}[enhanced jigsaw,left=3pt,breakable,colback=white!97.2549!black,boxrule=0.2mm]

\begin{Shaded}
\begin{Highlighting}[]
\NormalTok{130 3.140000 Hi! 12}
\end{Highlighting}
\end{Shaded}

\end{tcolorbox}

Kamu juga bisa men-``stringify'' \hl{\mbox{\texttt{\_\_VA\_ARGS\_\_}}} dengan menaruh \hl{\mbox{\texttt{\#}}} di depannya:

\begin{tcolorbox}[enhanced jigsaw,left=3pt,breakable,colback=white!97.2549!black,boxrule=0.2mm]

\begin{Shaded}
\begin{Highlighting}[]
\PreprocessorTok{\#define X}\OperatorTok{(...)}\PreprocessorTok{ }\OperatorTok{\#}\PreprocessorTok{\_\_VA\_ARGS\_\_}

\NormalTok{printf}\OperatorTok{(}\StringTok{"}\SpecialCharTok{\%s\textbackslash{}n}\StringTok{"}\OperatorTok{,}\NormalTok{ X}\OperatorTok{(}\DecValTok{1}\OperatorTok{,}\DecValTok{2}\OperatorTok{,}\DecValTok{3}\OperatorTok{));}  \CommentTok{// Mencetak "1, 2, 3"}
\end{Highlighting}
\end{Shaded}

\end{tcolorbox}

\index{\#define directive@\texttt{\#define} directive|)} \index{Preprocessor!macros with variable arguments|)} \index{Preprocessor!macros with arguments|)}

\subsection{Stringification}\label{stringification}

\index{\# stringification@\texttt{\#} stringification|(}

Sudah disebut di atas, kamu bisa mengubah argumen apa pun menjadi string dengan menaruh \hl{\mbox{\texttt{\#}}} di depannya pada teks penggantian.

Misalnya, kita bisa mencetak apa pun sebagai string dengan makro ini dan \hl{\mbox{\texttt{printf()}}}:

\begin{tcolorbox}[enhanced jigsaw,left=3pt,breakable,colback=white!97.2549!black,boxrule=0.2mm]

\begin{Shaded}
\begin{Highlighting}[]
\PreprocessorTok{\#define STR}\OperatorTok{(}\PreprocessorTok{x}\OperatorTok{)}\PreprocessorTok{ }\OperatorTok{\#}\PreprocessorTok{x}

\NormalTok{printf}\OperatorTok{(}\StringTok{"}\SpecialCharTok{\%s\textbackslash{}n}\StringTok{"}\OperatorTok{,}\NormalTok{ STR}\OperatorTok{(}\FloatTok{3.14159}\OperatorTok{));}
\end{Highlighting}
\end{Shaded}

\end{tcolorbox}

Dalam kasus itu, substitusinya menghasilkan:

\begin{tcolorbox}[enhanced jigsaw,left=3pt,breakable,colback=white!97.2549!black,boxrule=0.2mm]

\begin{Shaded}
\begin{Highlighting}[]
\NormalTok{printf}\OperatorTok{(}\StringTok{"}\SpecialCharTok{\%s\textbackslash{}n}\StringTok{"}\OperatorTok{,} \StringTok{"3.14159"}\OperatorTok{);}
\end{Highlighting}
\end{Shaded}

\end{tcolorbox}

Mari kita lihat apakah ini bisa dipakai lebih efektif sehingga kita bisa mengoper nama variabel \hl{\mbox{\texttt{int}}} mana pun ke makro, dan makro itu mencetak nama dan nilainya.

\begin{tcolorbox}[enhanced jigsaw,left=3pt,breakable,colback=white!97.2549!black,boxrule=0.2mm]

\begin{Shaded}
\begin{Highlighting}[numbers=left,,]
\PreprocessorTok{\#include }\ImportTok{\textless{}stdio.h\textgreater{}}

\PreprocessorTok{\#define PRINT\_INT\_VAL}\OperatorTok{(}\PreprocessorTok{x}\OperatorTok{)}\PreprocessorTok{ printf}\OperatorTok{(}\StringTok{"}\SpecialCharTok{\%s}\StringTok{ = }\SpecialCharTok{\%d\textbackslash{}n}\StringTok{"}\OperatorTok{,}\PreprocessorTok{ }\OperatorTok{\#}\PreprocessorTok{x}\OperatorTok{,}\PreprocessorTok{ x}\OperatorTok{)}

\DataTypeTok{int}\NormalTok{ main}\OperatorTok{(}\DataTypeTok{void}\OperatorTok{)}
\OperatorTok{\{}
    \DataTypeTok{int}\NormalTok{ a }\OperatorTok{=} \DecValTok{5}\OperatorTok{;}

\NormalTok{    PRINT\_INT\_VAL}\OperatorTok{(}\NormalTok{a}\OperatorTok{);}  \CommentTok{// mencetak "a = 5"}
\OperatorTok{\}}
\end{Highlighting}
\end{Shaded}

\end{tcolorbox}

Di baris 9, kita mendapat penggantian makro berikut:

\begin{tcolorbox}[enhanced jigsaw,left=3pt,breakable,colback=white!97.2549!black,boxrule=0.2mm]

\begin{Shaded}
\begin{Highlighting}[numbers=left,,firstnumber=9,]
\NormalTok{    printf}\OperatorTok{(}\StringTok{"}\SpecialCharTok{\%s}\StringTok{ = }\SpecialCharTok{\%d\textbackslash{}n}\StringTok{"}\OperatorTok{,} \StringTok{"a"}\OperatorTok{,}\NormalTok{ a}\OperatorTok{);}
\end{Highlighting}
\end{Shaded}

\end{tcolorbox}

\index{\# stringification@\texttt{\#} stringification|)}

\subsection{Concatenation}\label{concatenation}

\index{\#\# concatenation@\texttt{\#\#} concatenation|(}

Kita juga bisa menyambung dua argumen dengan \hl{\mbox{\texttt{\#\#}}}. Seru!

\begin{tcolorbox}[enhanced jigsaw,left=3pt,breakable,colback=white!97.2549!black,boxrule=0.2mm]

\begin{Shaded}
\begin{Highlighting}[]
\PreprocessorTok{\#define CAT}\OperatorTok{(}\PreprocessorTok{a}\OperatorTok{,}\PreprocessorTok{ b}\OperatorTok{)}\PreprocessorTok{ a }\OperatorTok{\#\#}\PreprocessorTok{ b}

\NormalTok{printf}\OperatorTok{(}\StringTok{"}\SpecialCharTok{\%f\textbackslash{}n}\StringTok{"}\OperatorTok{,}\NormalTok{ CAT}\OperatorTok{(}\FloatTok{3.14}\OperatorTok{,} \DecValTok{1592}\OperatorTok{));}   \CommentTok{// 3.141592}
\end{Highlighting}
\end{Shaded}

\end{tcolorbox}

\index{\#\# concatenation@\texttt{\#\#} concatenation|)}

\section{Makro Multibaris}\label{makro-multibaris}

\index{Preprocessor!multiline macros|(}

Makro bisa dilanjutkan ke beberapa baris kalau newline di-escape dengan backslash (\hl{\mbox{\texttt{\textbackslash{}}}}).

Mari kita tulis makro multibaris yang mencetak angka dari \hl{\mbox{\texttt{0}}} sampai hasil kali dua argumen yang dioper.

\index{do-while statement@\texttt{do}-\texttt{while} statement!in multiline macros|(}

\begin{tcolorbox}[enhanced jigsaw,left=3pt,breakable,colback=white!97.2549!black,boxrule=0.2mm]

\begin{Shaded}
\begin{Highlighting}[numbers=left,,]
\PreprocessorTok{\#include }\ImportTok{\textless{}stdio.h\textgreater{}}

\PreprocessorTok{\#define PRINT\_NUMS\_TO\_PRODUCT}\OperatorTok{(}\PreprocessorTok{a}\OperatorTok{,}\PreprocessorTok{ b}\OperatorTok{)}\PreprocessorTok{ }\ControlFlowTok{do}\PreprocessorTok{ }\OperatorTok{\{}\PreprocessorTok{ }\OperatorTok{\textbackslash{}}
\PreprocessorTok{    }\DataTypeTok{int}\PreprocessorTok{ product }\OperatorTok{=}\PreprocessorTok{ }\OperatorTok{(}\PreprocessorTok{a}\OperatorTok{)}\PreprocessorTok{ }\OperatorTok{*}\PreprocessorTok{ }\OperatorTok{(}\PreprocessorTok{b}\OperatorTok{);}\PreprocessorTok{ }\OperatorTok{\textbackslash{}}
\PreprocessorTok{    }\ControlFlowTok{for}\PreprocessorTok{ }\OperatorTok{(}\DataTypeTok{int}\PreprocessorTok{ i }\OperatorTok{=}\PreprocessorTok{ }\DecValTok{0}\OperatorTok{;}\PreprocessorTok{ i }\OperatorTok{\textless{}}\PreprocessorTok{ product}\OperatorTok{;}\PreprocessorTok{ i}\OperatorTok{++)}\PreprocessorTok{ }\OperatorTok{\{}\PreprocessorTok{ }\OperatorTok{\textbackslash{}}
\PreprocessorTok{        printf}\OperatorTok{(}\StringTok{"}\SpecialCharTok{\%d\textbackslash{}n}\StringTok{"}\OperatorTok{,}\PreprocessorTok{ i}\OperatorTok{);}\PreprocessorTok{ }\OperatorTok{\textbackslash{}}
\PreprocessorTok{    }\OperatorTok{\}}\PreprocessorTok{ }\OperatorTok{\textbackslash{}}
\OperatorTok{\}}\PreprocessorTok{ }\ControlFlowTok{while}\OperatorTok{(}\DecValTok{0}\OperatorTok{)}

\DataTypeTok{int}\NormalTok{ main}\OperatorTok{(}\DataTypeTok{void}\OperatorTok{)}
\OperatorTok{\{}
\NormalTok{    PRINT\_NUMS\_TO\_PRODUCT}\OperatorTok{(}\DecValTok{2}\OperatorTok{,} \DecValTok{4}\OperatorTok{);}  \CommentTok{// Mengeluarkan angka dari 0 sampai 7}
\OperatorTok{\}}
\end{Highlighting}
\end{Shaded}

\end{tcolorbox}

Beberapa hal yang perlu dicatat:

\begin{itemize}
\tightlist
\item
  Escape di akhir setiap baris kecuali yang terakhir untuk menandakan bahwa makro berlanjut.
\item
  Seluruhnya dibungkus dalam loop \hl{\mbox{\texttt{do}}}-\hl{\mbox{\texttt{while(0)}}} dengan kurung kurawal.
\end{itemize}

Poin kedua mungkin agak aneh, tapi semuanya soal menyerap \hl{\mbox{\texttt{;}}} di belakang yang dijatuhkan programmer setelah makro.

Awalnya aku mengira kurung kurawal saja sudah cukup, tapi ada kasus yang gagal kalau programmer menaruh titik koma setelah makro. Inilah kasusnya:

\begin{tcolorbox}[enhanced jigsaw,left=3pt,breakable,colback=white!97.2549!black,boxrule=0.2mm]

\begin{Shaded}
\begin{Highlighting}[numbers=left,,]
\PreprocessorTok{\#include }\ImportTok{\textless{}stdio.h\textgreater{}}

\PreprocessorTok{\#define FOO}\OperatorTok{(}\PreprocessorTok{x}\OperatorTok{)}\PreprocessorTok{ }\OperatorTok{\{}\PreprocessorTok{ }\OperatorTok{(}\PreprocessorTok{x}\OperatorTok{)++;}\PreprocessorTok{ }\OperatorTok{\}}

\DataTypeTok{int}\NormalTok{ main}\OperatorTok{(}\DataTypeTok{void}\OperatorTok{)}
\OperatorTok{\{}
    \DataTypeTok{int}\NormalTok{ i }\OperatorTok{=} \DecValTok{0}\OperatorTok{;}

    \ControlFlowTok{if} \OperatorTok{(}\NormalTok{i }\OperatorTok{==} \DecValTok{0}\OperatorTok{)}
\NormalTok{        FOO}\OperatorTok{(}\NormalTok{i}\OperatorTok{);}
    \ControlFlowTok{else}
\NormalTok{        printf}\OperatorTok{(}\StringTok{":{-}(}\SpecialCharTok{\textbackslash{}n}\StringTok{"}\OperatorTok{);}

\NormalTok{    printf}\OperatorTok{(}\StringTok{"}\SpecialCharTok{\%d\textbackslash{}n}\StringTok{"}\OperatorTok{,}\NormalTok{ i}\OperatorTok{);}
\OperatorTok{\}}
\end{Highlighting}
\end{Shaded}

\end{tcolorbox}

Kelihatannya sederhana, tapi tidak akan terbangun tanpa syntax error:

\begin{tcolorbox}[enhanced jigsaw,left=3pt,breakable,colback=white!97.2549!black,boxrule=0.2mm]

\begin{Shaded}
\begin{Highlighting}[]
\NormalTok{foo.c:11:5: error: ‘else’ without a previous ‘if’  }
\end{Highlighting}
\end{Shaded}

\end{tcolorbox}

Kamu bisa lihat masalahnya?

Mari kita lihat perluasannya:

\begin{tcolorbox}[enhanced jigsaw,left=3pt,breakable,colback=white!97.2549!black,boxrule=0.2mm]

\begin{Shaded}
\begin{Highlighting}[]

    \ControlFlowTok{if} \OperatorTok{(}\NormalTok{i }\OperatorTok{==} \DecValTok{0}\OperatorTok{)} \OperatorTok{\{}
        \OperatorTok{(}\NormalTok{i}\OperatorTok{)++;}
    \OperatorTok{\};}             \CommentTok{// \textless{}{-}{-} Masalah dengan M huruf besar!}

    \ControlFlowTok{else}
\NormalTok{        printf}\OperatorTok{(}\StringTok{":{-}(}\SpecialCharTok{\textbackslash{}n}\StringTok{"}\OperatorTok{);}
\end{Highlighting}
\end{Shaded}

\end{tcolorbox}

\hl{\mbox{\texttt{;}}} itu mengakhiri pernyataan \hl{\mbox{\texttt{if}}}, jadi \hl{\mbox{\texttt{else}}} melayang-layang di sana secara ilegal\footnote{\emph{Breakin' the law\ldots{} breakin' the law\ldots{}}}.

Jadi bungkus makro multibaris itu dengan \hl{\mbox{\texttt{do}}}-\hl{\mbox{\texttt{while(0)}}}.

\index{do-while statement@\texttt{do}-\texttt{while} statement!in multiline macros|)} \index{Preprocessor!multiline macros|)}

\section{Contoh: Makro Assert}\label{my-assert}

Menambahkan assert ke kodemu adalah cara bagus untuk menangkap kondisi yang menurutmu seharusnya tidak terjadi. C menyediakan fungsionalitas \hl{\mbox{\texttt{assert()}}}. Ia memeriksa sebuah kondisi, dan kalau false, program berhenti sambil memberi tahu file dan nomor baris tempat assertion gagal.

Tapi ini kurang memadai.

\begin{enumerate}
\def\labelenumi{\arabic{enumi}.}
\tightlist
\item
  Pertama, kamu tidak bisa menambahkan pesan tambahan pada assert.
\item
  Kedua, tidak ada saklar on-off yang mudah untuk semua assert.
\end{enumerate}

Yang pertama bisa kita atasi dengan makro.

Pada dasarnya, kalau aku punya kode ini:

\begin{tcolorbox}[enhanced jigsaw,left=3pt,breakable,colback=white!97.2549!black,boxrule=0.2mm]

\begin{Shaded}
\begin{Highlighting}[]
\NormalTok{ASSERT}\OperatorTok{(}\NormalTok{x }\OperatorTok{\textless{}} \DecValTok{20}\OperatorTok{,} \StringTok{"x must be under 20"}\OperatorTok{);}
\end{Highlighting}
\end{Shaded}

\end{tcolorbox}

aku ingin sesuatu seperti ini terjadi (dengan asumsi \hl{\mbox{\texttt{ASSERT()}}} ada di baris 220 \hl{\mbox{\texttt{foo.c}}}):

\begin{tcolorbox}[enhanced jigsaw,left=3pt,breakable,colback=white!97.2549!black,boxrule=0.2mm]

\begin{Shaded}
\begin{Highlighting}[]
\ControlFlowTok{if} \OperatorTok{(!(}\NormalTok{x }\OperatorTok{\textless{}} \DecValTok{20}\OperatorTok{))} \OperatorTok{\{}
\NormalTok{    fprintf}\OperatorTok{(}\NormalTok{stderr}\OperatorTok{,} \StringTok{"foo.c:220: assertion x \textless{} 20 failed: "}\OperatorTok{);}
\NormalTok{    fprintf}\OperatorTok{(}\NormalTok{stderr}\OperatorTok{,} \StringTok{"x must be under 20}\SpecialCharTok{\textbackslash{}n}\StringTok{"}\OperatorTok{);}
\NormalTok{    exit}\OperatorTok{(}\DecValTok{1}\OperatorTok{);}
\OperatorTok{\}}
\end{Highlighting}
\end{Shaded}

\end{tcolorbox}

Kita bisa mengambil nama file dari makro \hl{\mbox{\texttt{\_\_FILE\_\_}}}, dan nomor baris dari \hl{\mbox{\texttt{\_\_LINE\_\_}}}. Pesannya sudah berupa string, tapi \hl{\mbox{\texttt{x\ \textless{}\ 20}}} bukan, jadi kita harus men-stringify-nya dengan \hl{\mbox{\texttt{\#}}}. Kita bisa membuat makro multibaris dengan escape backslash di akhir baris.

\begin{tcolorbox}[enhanced jigsaw,left=3pt,breakable,colback=white!97.2549!black,boxrule=0.2mm]

\begin{Shaded}
\begin{Highlighting}[]
\PreprocessorTok{\#define ASSERT}\OperatorTok{(}\PreprocessorTok{c}\OperatorTok{,}\PreprocessorTok{ m}\OperatorTok{)}\PreprocessorTok{ }\OperatorTok{\textbackslash{}}
\ControlFlowTok{do}\PreprocessorTok{ }\OperatorTok{\{}\PreprocessorTok{ }\OperatorTok{\textbackslash{}}
\PreprocessorTok{    }\ControlFlowTok{if}\PreprocessorTok{ }\OperatorTok{(!(}\PreprocessorTok{c}\OperatorTok{))}\PreprocessorTok{ }\OperatorTok{\{}\PreprocessorTok{ }\OperatorTok{\textbackslash{}}
\PreprocessorTok{        fprintf}\OperatorTok{(}\PreprocessorTok{stderr}\OperatorTok{,}\PreprocessorTok{ \_\_FILE\_\_ }\StringTok{":}\SpecialCharTok{\%d}\StringTok{: assertion }\SpecialCharTok{\%s}\StringTok{ failed: }\SpecialCharTok{\%s\textbackslash{}n}\StringTok{"}\OperatorTok{,}\PreprocessorTok{ }\OperatorTok{\textbackslash{}}
\PreprocessorTok{                        \_\_LINE\_\_}\OperatorTok{,}\PreprocessorTok{ }\OperatorTok{\#}\PreprocessorTok{c}\OperatorTok{,}\PreprocessorTok{ m}\OperatorTok{);}\PreprocessorTok{ }\OperatorTok{\textbackslash{}}
\PreprocessorTok{        exit}\OperatorTok{(}\DecValTok{1}\OperatorTok{);}\PreprocessorTok{ }\OperatorTok{\textbackslash{}}
\PreprocessorTok{    }\OperatorTok{\}}\PreprocessorTok{ }\OperatorTok{\textbackslash{}}
\OperatorTok{\}}\PreprocessorTok{ }\ControlFlowTok{while}\OperatorTok{(}\DecValTok{0}\OperatorTok{)}
\end{Highlighting}
\end{Shaded}

\end{tcolorbox}

(Kelihatannya agak aneh dengan \hl{\mbox{\texttt{\_\_FILE\_\_}}} di depan seperti itu, tapi ingat ia adalah string literal, dan string literal yang bersebelahan otomatis disambung. Sedangkan \hl{\mbox{\texttt{\_\_LINE\_\_}}}, ia hanyalah \hl{\mbox{\texttt{int}}}.)

Dan itu berfungsi! Kalau kujalankan ini:

\begin{tcolorbox}[enhanced jigsaw,left=3pt,breakable,colback=white!97.2549!black,boxrule=0.2mm]

\begin{Shaded}
\begin{Highlighting}[]
\DataTypeTok{int}\NormalTok{ x }\OperatorTok{=} \DecValTok{30}\OperatorTok{;}

\NormalTok{ASSERT}\OperatorTok{(}\NormalTok{x }\OperatorTok{\textless{}} \DecValTok{20}\OperatorTok{,} \StringTok{"x must be under 20"}\OperatorTok{);}
\end{Highlighting}
\end{Shaded}

\end{tcolorbox}

aku mendapat output ini:

\begin{tcolorbox}[enhanced jigsaw,left=3pt,breakable,colback=white!97.2549!black,boxrule=0.2mm]

\begin{verbatim}
foo.c:23: assertion x < 20 failed: x must be under 20
\end{verbatim}

\end{tcolorbox}

Bagus sekali!

Tinggal satu hal lagi, yaitu cara menyalakan dan mematikannya, dan itu bisa dilakukan dengan kompilasi kondisional.

Berikut contoh lengkapnya:

\begin{tcolorbox}[enhanced jigsaw,left=3pt,breakable,colback=white!97.2549!black,boxrule=0.2mm]

\begin{Shaded}
\begin{Highlighting}[numbers=left,,]
\PreprocessorTok{\#include }\ImportTok{\textless{}stdio.h\textgreater{}}
\PreprocessorTok{\#include }\ImportTok{\textless{}stdlib.h\textgreater{}}

\PreprocessorTok{\#define ASSERT\_ENABLED }\DecValTok{1}

\PreprocessorTok{\#if ASSERT\_ENABLED}
\PreprocessorTok{\#define ASSERT}\OperatorTok{(}\PreprocessorTok{c}\OperatorTok{,}\PreprocessorTok{ m}\OperatorTok{)}\PreprocessorTok{ }\OperatorTok{\textbackslash{}}
\ControlFlowTok{do}\PreprocessorTok{ }\OperatorTok{\{}\PreprocessorTok{ }\OperatorTok{\textbackslash{}}
\PreprocessorTok{    }\ControlFlowTok{if}\PreprocessorTok{ }\OperatorTok{(!(}\PreprocessorTok{c}\OperatorTok{))}\PreprocessorTok{ }\OperatorTok{\{}\PreprocessorTok{ }\OperatorTok{\textbackslash{}}
\PreprocessorTok{        fprintf}\OperatorTok{(}\PreprocessorTok{stderr}\OperatorTok{,}\PreprocessorTok{ \_\_FILE\_\_ }\StringTok{":}\SpecialCharTok{\%d}\StringTok{: assertion }\SpecialCharTok{\%s}\StringTok{ failed: }\SpecialCharTok{\%s\textbackslash{}n}\StringTok{"}\OperatorTok{,}\PreprocessorTok{ }\OperatorTok{\textbackslash{}}
\PreprocessorTok{                        \_\_LINE\_\_}\OperatorTok{,}\PreprocessorTok{ }\OperatorTok{\#}\PreprocessorTok{c}\OperatorTok{,}\PreprocessorTok{ m}\OperatorTok{);}\PreprocessorTok{ }\OperatorTok{\textbackslash{}}
\PreprocessorTok{        exit}\OperatorTok{(}\DecValTok{1}\OperatorTok{);}\PreprocessorTok{ }\OperatorTok{\textbackslash{}}
\PreprocessorTok{    }\OperatorTok{\}}\PreprocessorTok{ }\OperatorTok{\textbackslash{}}
\OperatorTok{\}}\PreprocessorTok{ }\ControlFlowTok{while}\OperatorTok{(}\DecValTok{0}\OperatorTok{)}
\PreprocessorTok{\#else}
\PreprocessorTok{\#define ASSERT}\OperatorTok{(}\PreprocessorTok{c}\OperatorTok{,}\PreprocessorTok{ m}\OperatorTok{)}\PreprocessorTok{  }\CommentTok{// Makro kosong kalau tidak diaktifkan}
\PreprocessorTok{\#endif}

\DataTypeTok{int}\NormalTok{ main}\OperatorTok{(}\DataTypeTok{void}\OperatorTok{)}
\OperatorTok{\{}
    \DataTypeTok{int}\NormalTok{ x }\OperatorTok{=} \DecValTok{30}\OperatorTok{;}

\NormalTok{    ASSERT}\OperatorTok{(}\NormalTok{x }\OperatorTok{\textless{}} \DecValTok{20}\OperatorTok{,} \StringTok{"x must be under 20"}\OperatorTok{);}
\OperatorTok{\}}
\end{Highlighting}
\end{Shaded}

\end{tcolorbox}

Ini menghasilkan output:

\begin{tcolorbox}[enhanced jigsaw,left=3pt,breakable,colback=white!97.2549!black,boxrule=0.2mm]

\begin{Shaded}
\begin{Highlighting}[]
\NormalTok{foo.c:23: assertion x \textless{} 20 failed: x must be under 20}
\end{Highlighting}
\end{Shaded}

\end{tcolorbox}

\section{\texorpdfstring{Directive \hl{\mbox{\texttt{\#error}}}}{Directive \#error}}\label{directive-error}

\index{\#error directive@\texttt{\#error} directive|(}

Directive ini membuat compiler langsung error begitu menemukannya.

Biasanya, ini dipakai di dalam kondisional untuk mencegah kompilasi kecuali prasyarat tertentu terpenuhi:

\begin{tcolorbox}[enhanced jigsaw,left=3pt,breakable,colback=white!97.2549!black,boxrule=0.2mm]

\begin{Shaded}
\begin{Highlighting}[]
\PreprocessorTok{\#ifndef \_\_STDC\_IEC\_559\_\_}
    \PreprocessorTok{\#error I really need IEEE{-}754 floating point to compile. Sorry!}
\PreprocessorTok{\#endif}
\end{Highlighting}
\end{Shaded}

\end{tcolorbox}

\index{\#error directive@\texttt{\#error} directive|)} \index{\#warning directive@\texttt{\#warning} directive|(}

Beberapa compiler punya directive \hl{\mbox{\texttt{\#warning}}} pelengkap yang non-standar, yang mengeluarkan peringatan tapi tidak menghentikan kompilasi, namun ini tidak ada di spesifikasi C11.

\index{\#warning directive@\texttt{\#warning} directive|)}

\section{\texorpdfstring{Directive \hl{\mbox{\texttt{\#embed}}}}{Directive \#embed}}\label{directive-embed}

\index{\#embed directive@\texttt{\#embed} directive|(}

Baru di C23!

Dan saat ini belum berfungsi di compiler mana pun milikku, jadi jangan telan bagian ini bulat-bulat!

Intinya, kamu bisa menyertakan byte-byte dari sebuah file sebagai konstanta integer seolah-olah kamu mengetiknya sendiri.

Misalnya, kalau kamu punya file biner bernama \hl{\mbox{\texttt{foo.bin}}} yang berisi empat byte bernilai desimal 11, 22, 33, dan 44, lalu kamu menulis ini:

\begin{tcolorbox}[enhanced jigsaw,left=3pt,breakable,colback=white!97.2549!black,boxrule=0.2mm]

\begin{Shaded}
\begin{Highlighting}[]
\DataTypeTok{int}\NormalTok{ a}\OperatorTok{[]} \OperatorTok{=} \OperatorTok{\{}
\PreprocessorTok{\#embed "foo.bin"}
\OperatorTok{\};}
\end{Highlighting}
\end{Shaded}

\end{tcolorbox}

Hasilnya persis seperti kamu mengetik ini:

\begin{tcolorbox}[enhanced jigsaw,left=3pt,breakable,colback=white!97.2549!black,boxrule=0.2mm]

\begin{Shaded}
\begin{Highlighting}[]
\DataTypeTok{int}\NormalTok{ a}\OperatorTok{[]} \OperatorTok{=} \OperatorTok{\{}\DecValTok{11}\OperatorTok{,}\DecValTok{22}\OperatorTok{,}\DecValTok{33}\OperatorTok{,}\DecValTok{44}\OperatorTok{\};}
\end{Highlighting}
\end{Shaded}

\end{tcolorbox}

Ini cara yang sangat kuat untuk menginisialisasi array dengan data biner tanpa perlu mengonversinya menjadi kode dulu---preprocessor yang mengerjakannya untukmu!

Kasus pemakaian yang lebih khas mungkin file berisi gambar kecil yang akan ditampilkan dan tidak ingin kamu muat saat runtime.

Berikut contoh lain:

\begin{tcolorbox}[enhanced jigsaw,left=3pt,breakable,colback=white!97.2549!black,boxrule=0.2mm]

\begin{Shaded}
\begin{Highlighting}[]
\DataTypeTok{int}\NormalTok{ a}\OperatorTok{[]} \OperatorTok{=} \OperatorTok{\{}
\PreprocessorTok{\#embed \textless{}foo.bin\textgreater{}}
\OperatorTok{\};}
\end{Highlighting}
\end{Shaded}

\end{tcolorbox}

Kalau kamu memakai kurung siku, preprocessor mencari di serangkaian tempat yang ditentukan implementasi untuk menemukan file, persis seperti \hl{\mbox{\texttt{\#include}}}. Kalau kamu memakai tanda kutip ganda dan resource tidak ditemukan, compiler akan mencobanya seolah-olah kamu memakai kurung siku sebagai upaya putus asa terakhir untuk menemukan file itu.

\hl{\mbox{\texttt{\#embed}}} bekerja seperti \hl{\mbox{\texttt{\#include}}} dalam arti ia secara efektif menempelkan nilai sebelum compiler melihatnya. Artinya kamu bisa memakainya di berbagai tempat:

\begin{tcolorbox}[enhanced jigsaw,left=3pt,breakable,colback=white!97.2549!black,boxrule=0.2mm]

\begin{verbatim}
return
#embed "somevalue.dat"
;
\end{verbatim}

\end{tcolorbox}

atau

\begin{tcolorbox}[enhanced jigsaw,left=3pt,breakable,colback=white!97.2549!black,boxrule=0.2mm]

\begin{verbatim}
int x =
#embed "xvalue.dat"
;
\end{verbatim}

\end{tcolorbox}

Nah---apakah ini selalu byte? Maksudnya bernilai 0 sampai 255, inklusif? Jawabannya, secara bawaan jelas ``ya'', kecuali saat jawabannya ``tidak''.

Secara teknis, elemennya selebar \hl{\mbox{\texttt{CHAR\_BIT}}} bit. Dan kemungkinan besar itu 8 di sistemmu, jadi kamu mendapat rentang 0-255 pada nilainya. (Nilainya selalu non-negatif.)

Mungkin juga implementasi membolehkan ini diubah dengan cara tertentu, misalnya lewat command line atau parameter.

Ukuran file dalam bit harus kelipatan ukuran elemen. Artinya, kalau setiap elemen 8 bit, ukuran file (dalam bit) harus kelipatan 8. Dalam pemakaian sehari-hari, ini cara membingungkan untuk mengatakan bahwa setiap file harus berupa jumlah byte yang bulat\ldots{} yang tentu saja memang begitu. Jujur saja, aku sendiri tidak yakin kenapa aku repot-repot menulis paragraf ini. Baca spesifikasi kalau kamu benar-benar penasaran.

\subsection{\texorpdfstring{Parameter \hl{\mbox{\texttt{\#embed}}}}{Parameter \#embed}}\label{parameter-embed}

Ada berbagai parameter yang bisa disebutkan pada directive \hl{\mbox{\texttt{\#embed}}}. Berikut contoh dengan parameter \hl{\mbox{\texttt{limit()}}} yang belum diperkenalkan:

\begin{tcolorbox}[enhanced jigsaw,left=3pt,breakable,colback=white!97.2549!black,boxrule=0.2mm]

\begin{Shaded}
\begin{Highlighting}[]
\DataTypeTok{int}\NormalTok{ a}\OperatorTok{[]} \OperatorTok{=} \OperatorTok{\{}
\PreprocessorTok{\#embed "/dev/random" limit(5)}
\OperatorTok{\};}
\end{Highlighting}
\end{Shaded}

\end{tcolorbox}

Tapi bagaimana kalau \hl{\mbox{\texttt{limit}}} sudah didefinisikan di tempat lain? Untungnya kamu bisa menaruh \hl{\mbox{\texttt{\_\_}}} di sekeliling kata kuncinya dan ia bekerja dengan cara yang sama:

\begin{tcolorbox}[enhanced jigsaw,left=3pt,breakable,colback=white!97.2549!black,boxrule=0.2mm]

\begin{Shaded}
\begin{Highlighting}[]
\DataTypeTok{int}\NormalTok{ a}\OperatorTok{[]} \OperatorTok{=} \OperatorTok{\{}
\PreprocessorTok{\#embed "/dev/random" \_\_limit\_\_(5)}
\OperatorTok{\};}
\end{Highlighting}
\end{Shaded}

\end{tcolorbox}

Nah\ldots{} apa itu \hl{\mbox{\texttt{limit}}}?

\subsection{\texorpdfstring{Parameter \hl{\mbox{\texttt{limit()}}}}{Parameter limit()}}\label{parameter-limit}

Kamu bisa menentukan batas jumlah elemen yang disematkan dengan parameter ini.

Ini nilai maksimum, bukan nilai mutlak. Kalau file yang disematkan lebih pendek dari batas yang ditentukan, hanya sebanyak itu byte yang diimpor.

Contoh \hl{\mbox{\texttt{/dev/random}}} di atas adalah contoh motivasi untuk ini---di Unix, itu \emph{character device file} yang mengembalikan aliran tak terbatas bilangan yang cukup acak.

Menyematkan byte dalam jumlah tak terbatas memberatkan RAM, jadi parameter \hl{\mbox{\texttt{limit}}} memberi cara untuk berhenti setelah jumlah tertentu.

Terakhir, kamu boleh memakai makro \hl{\mbox{\texttt{\#define}}} di dalam \hl{\mbox{\texttt{limit}}}, kalau kamu penasaran.

\subsection{\texorpdfstring{Parameter \hl{\mbox{\texttt{if\_empty}}}}{Parameter if\_empty}}\label{parameter-if_empty}

\index{if\_empty() embed parameter@\texttt{if\_empty()} embed parameter|(}

Parameter ini menentukan apa hasil embed kalau file ada tapi tidak berisi data. Misalkan file \hl{\mbox{\texttt{foo.dat}}} berisi satu byte bernilai 123. Kalau kita menulis ini:

\begin{tcolorbox}[enhanced jigsaw,left=3pt,breakable,colback=white!97.2549!black,boxrule=0.2mm]

\begin{Shaded}
\begin{Highlighting}[]
\DataTypeTok{int}\NormalTok{ x }\OperatorTok{=} 
\PreprocessorTok{\#embed "foo.dat" if\_empty(999)}
\OperatorTok{;}
\end{Highlighting}
\end{Shaded}

\end{tcolorbox}

kita mendapat:

\begin{tcolorbox}[enhanced jigsaw,left=3pt,breakable,colback=white!97.2549!black,boxrule=0.2mm]

\begin{Shaded}
\begin{Highlighting}[]
\DataTypeTok{int}\NormalTok{ x }\OperatorTok{=} \DecValTok{123}\OperatorTok{;}   \CommentTok{// Ketika foo.dat berisi byte 123}
\end{Highlighting}
\end{Shaded}

\end{tcolorbox}

Tapi bagaimana kalau file \hl{\mbox{\texttt{foo.dat}}} berukuran nol byte (artinya tidak berisi data dan kosong)? Kalau begitu, hasilnya diperluas menjadi:

\begin{tcolorbox}[enhanced jigsaw,left=3pt,breakable,colback=white!97.2549!black,boxrule=0.2mm]

\begin{Shaded}
\begin{Highlighting}[]
\DataTypeTok{int}\NormalTok{ x }\OperatorTok{=} \DecValTok{999}\OperatorTok{;}   \CommentTok{// Ketika foo.dat kosong}
\end{Highlighting}
\end{Shaded}

\end{tcolorbox}

Perlu dicatat, kalau \hl{\mbox{\texttt{limit}}} diatur ke \hl{\mbox{\texttt{0}}}, maka \hl{\mbox{\texttt{if\_empty}}} akan selalu dipakai. Artinya, batas nol secara efektif berarti file itu kosong.

Ini akan selalu menghasilkan \hl{\mbox{\texttt{x\ =\ 999}}} apa pun isi \hl{\mbox{\texttt{foo.dat}}}:

\begin{tcolorbox}[enhanced jigsaw,left=3pt,breakable,colback=white!97.2549!black,boxrule=0.2mm]

\begin{Shaded}
\begin{Highlighting}[]
\DataTypeTok{int}\NormalTok{ x }\OperatorTok{=} 
\PreprocessorTok{\#embed "foo.dat" limit(0) if\_empty(999)}
\OperatorTok{;}
\end{Highlighting}
\end{Shaded}

\end{tcolorbox}

\index{if\_empty() embed parameter@\texttt{if\_empty()} embed parameter|)}

\subsection{\texorpdfstring{Parameter \hl{\mbox{\texttt{prefix()}}} dan \hl{\mbox{\texttt{suffix()}}}}{Parameter prefix() dan suffix()}}\label{parameter-prefix-dan-suffix}

\index{prefix() embed parameter@\texttt{prefix()} embed parameter|(} \index{suffix() embed parameter@\texttt{suffix()} embed parameter|(}

Ini cara untuk menambahkan sesuatu di depan embed.

Perhatikan bahwa ini hanya berpengaruh pada data yang tidak kosong! Kalau file kosong, baik \hl{\mbox{\texttt{prefix}}} maupun \hl{\mbox{\texttt{suffix}}} tidak berpengaruh.

Berikut contoh ketika kita menyematkan tiga bilangan acak, tapi mengawali angka-angka itu dengan \hl{\mbox{\texttt{11,}}} dan mengakhirinya dengan \hl{\mbox{\texttt{,99}}}:

\begin{tcolorbox}[enhanced jigsaw,left=3pt,breakable,colback=white!97.2549!black,boxrule=0.2mm]

\begin{Shaded}
\begin{Highlighting}[]
\DataTypeTok{int}\NormalTok{ x}\OperatorTok{[]} \OperatorTok{=} \OperatorTok{\{}
\PreprocessorTok{\#embed "/dev/urandom" limit(3) prefix(11,) suffix(,99)}
\OperatorTok{\};}
\end{Highlighting}
\end{Shaded}

\end{tcolorbox}

Contoh hasil:

\begin{tcolorbox}[enhanced jigsaw,left=3pt,breakable,colback=white!97.2549!black,boxrule=0.2mm]

\begin{Shaded}
\begin{Highlighting}[]
\DataTypeTok{int}\NormalTok{ x}\OperatorTok{[]} \OperatorTok{=} \OperatorTok{\{}\DecValTok{11}\OperatorTok{,}\DecValTok{135}\OperatorTok{,}\DecValTok{116}\OperatorTok{,}\DecValTok{220}\OperatorTok{,}\DecValTok{99}\OperatorTok{\};}
\end{Highlighting}
\end{Shaded}

\end{tcolorbox}

Tidak ada kewajiban memakai \hl{\mbox{\texttt{prefix}}} dan \hl{\mbox{\texttt{suffix}}} sekaligus. Kamu bisa memakai keduanya, salah satu, atau tidak sama sekali.

Kita bisa memanfaatkan sifat bahwa keduanya hanya berlaku pada file yang tidak kosong untuk efek yang rapi, seperti pada contoh berikut yang dicuri tanpa malu dari spesifikasi.

Katakanlah kita punya file \hl{\mbox{\texttt{foo.dat}}} yang berisi sejumlah data. Dan kita ingin memakainya untuk menginisialisasi array, lalu kita ingin sufiks pada array berupa elemen nol.

Tidak masalah, kan?

\begin{tcolorbox}[enhanced jigsaw,left=3pt,breakable,colback=white!97.2549!black,boxrule=0.2mm]

\begin{Shaded}
\begin{Highlighting}[]
\DataTypeTok{int}\NormalTok{ x}\OperatorTok{[]} \OperatorTok{=} \OperatorTok{\{}
\PreprocessorTok{\#embed "foo.dat" suffix(,0)}
\OperatorTok{\};}
\end{Highlighting}
\end{Shaded}

\end{tcolorbox}

Kalau \hl{\mbox{\texttt{foo.dat}}} berisi 11, 22, dan 33, kita mendapat:

\begin{tcolorbox}[enhanced jigsaw,left=3pt,breakable,colback=white!97.2549!black,boxrule=0.2mm]

\begin{Shaded}
\begin{Highlighting}[]
\DataTypeTok{int}\NormalTok{ x}\OperatorTok{[]} \OperatorTok{=} \OperatorTok{\{}\DecValTok{11}\OperatorTok{,}\DecValTok{22}\OperatorTok{,}\DecValTok{33}\OperatorTok{,}\DecValTok{0}\OperatorTok{\};}
\end{Highlighting}
\end{Shaded}

\end{tcolorbox}

Tapi tunggu! Bagaimana kalau \hl{\mbox{\texttt{foo.dat}}} kosong? Kita mendapat:

\begin{tcolorbox}[enhanced jigsaw,left=3pt,breakable,colback=white!97.2549!black,boxrule=0.2mm]

\begin{Shaded}
\begin{Highlighting}[]
\DataTypeTok{int}\NormalTok{ x}\OperatorTok{[]} \OperatorTok{=} \OperatorTok{\{\};}
\end{Highlighting}
\end{Shaded}

\end{tcolorbox}

dan itu tidak bagus.

Tapi kita bisa memperbaikinya seperti ini:

\begin{tcolorbox}[enhanced jigsaw,left=3pt,breakable,colback=white!97.2549!black,boxrule=0.2mm]

\begin{Shaded}
\begin{Highlighting}[]
\DataTypeTok{int}\NormalTok{ x}\OperatorTok{[]} \OperatorTok{=} \OperatorTok{\{}
\PreprocessorTok{\#embed "foo.dat" suffix(,)}
    \DecValTok{0}
\OperatorTok{\};}
\end{Highlighting}
\end{Shaded}

\end{tcolorbox}

Karena parameter \hl{\mbox{\texttt{suffix}}} dihilangkan kalau file kosong, ini hanya akan menjadi:

\begin{tcolorbox}[enhanced jigsaw,left=3pt,breakable,colback=white!97.2549!black,boxrule=0.2mm]

\begin{Shaded}
\begin{Highlighting}[]
\DataTypeTok{int}\NormalTok{ x}\OperatorTok{[]} \OperatorTok{=} \OperatorTok{\{}\DecValTok{0}\OperatorTok{\};}
\end{Highlighting}
\end{Shaded}

\end{tcolorbox}

yang tidak masalah.

\index{prefix() embed parameter@\texttt{prefix()} embed parameter|)} \index{suffix() embed parameter@\texttt{suffix()} embed parameter|)}

\subsection{\texorpdfstring{Identifier \hl{\mbox{\texttt{\_\_has\_embed()}}}}{Identifier \_\_has\_embed()}}\label{identifier-__has_embed}

\index{\_\_has\_embed() identifier@\texttt{\_\_has\_embed()} identifier|(}

Ini cara yang bagus untuk menguji apakah file tertentu tersedia untuk disematkan, dan juga apakah file itu kosong atau tidak.

Kamu memakainya dengan directive \hl{\mbox{\texttt{\#if}}}.

Berikut potongan kode yang mengambil 5 bilangan acak dari character device generator bilangan acak. Kalau itu tidak ada, ia mencoba mengambilnya dari file \hl{\mbox{\texttt{myrandoms.dat}}}. Kalau itu juga tidak ada, ia memakai beberapa nilai yang ditulis langsung:

\begin{tcolorbox}[enhanced jigsaw,left=3pt,breakable,colback=white!97.2549!black,boxrule=0.2mm]

\begin{Shaded}
\begin{Highlighting}[]
    \DataTypeTok{int}\NormalTok{ random\_nums}\OperatorTok{[]} \OperatorTok{=} \OperatorTok{\{}
\PreprocessorTok{\#if \_\_has\_embed("/dev/urandom")}
    \PreprocessorTok{\#embed "/dev/urandom" limit(5)}
\PreprocessorTok{\#elif \_\_has\_embed("myrandoms.dat")}
    \PreprocessorTok{\#embed "myrandoms.dat" limit(5)}
\PreprocessorTok{\#else}
    \DecValTok{140}\OperatorTok{,}\DecValTok{178}\OperatorTok{,}\DecValTok{92}\OperatorTok{,}\DecValTok{167}\OperatorTok{,}\DecValTok{120}
\PreprocessorTok{\#endif}
    \OperatorTok{\};}
\end{Highlighting}
\end{Shaded}

\end{tcolorbox}

Secara teknis, identifier \hl{\mbox{\texttt{\_\_has\_embed()}}} bernilai salah satu dari tiga nilai:

{\def\LTcaptype{none} % do not increment counter
\begin{longtable}[]{@{}ll@{}}
\toprule\noalign{}
Hasil \hl{\mbox{\texttt{\_\_has\_embed()}}} & Deskripsi \\
\midrule\noalign{}
\endhead
\bottomrule\noalign{}
\endlastfoot
\hl{\mbox{\texttt{\_\_STDC\_EMBED\_NOT\_FOUND\_\_}}} & Jika file tidak ditemukan. \\
\hl{\mbox{\texttt{\_\_STDC\_EMBED\_FOUND\_\_}}} & Jika file ditemukan dan tidak kosong. \\
\hl{\mbox{\texttt{\_\_STDC\_EMBED\_EMPTY\_\_}}} & Jika file ditemukan dan kosong. \\
\end{longtable}
}

Aku punya alasan kuat untuk percaya bahwa \hl{\mbox{\texttt{\_\_STDC\_EMBED\_NOT\_FOUND\_\_}}} adalah \hl{\mbox{\texttt{0}}} dan yang lain bukan nol (karena itu tersirat dalam proposal dan masuk akal secara logika), tapi aku kesulitan menemukannya di versi draf spesifikasi ini.

\index{\_\_has\_embed() identifier@\texttt{\_\_has\_embed()} identifier|)}

TODO

\subsection{Parameter Lainnya}\label{parameter-lainnya}

Implementasi compiler boleh mendefinisikan parameter embed lain sesukanya ---cari parameter non-standar ini di dokumentasi compiler-mu.

Misalnya:

\begin{tcolorbox}[enhanced jigsaw,left=3pt,breakable,colback=white!97.2549!black,boxrule=0.2mm]

\begin{Shaded}
\begin{Highlighting}[]
\PreprocessorTok{\#embed "foo.bin" limit(12) frotz(lamp)}
\end{Highlighting}
\end{Shaded}

\end{tcolorbox}

Parameter ini biasanya diberi awalan untuk membantu namespacing:

\begin{tcolorbox}[enhanced jigsaw,left=3pt,breakable,colback=white!97.2549!black,boxrule=0.2mm]

\begin{Shaded}
\begin{Highlighting}[]
\PreprocessorTok{\#embed "foo.bin" limit(12) fmc::frotz(lamp)}
\end{Highlighting}
\end{Shaded}

\end{tcolorbox}

Mungkin masuk akal mencoba mendeteksi apakah parameter ini tersedia sebelum dipakai, dan untungnya kita bisa memakai \hl{\mbox{\texttt{\_\_has\_embed}}} untuk itu.

Biasanya, \hl{\mbox{\texttt{\_\_has\_embed()}}} hanya memberi tahu apakah file ada atau tidak. Tapi---dan inilah bagian serunya---ia juga mengembalikan false kalau parameter tambahan tidak didukung!

Jadi kalau kita memberinya file yang \emph{kita tahu} ada serta parameter yang ingin kita uji keberadaannya, ia secara efektif memberi tahu apakah parameter itu didukung.

Tapi file apa yang \emph{selalu} ada? Ternyata kita bisa memakai makro \hl{\mbox{\texttt{\_\_FILE\_\_}}}, yang diperluas menjadi nama file sumber yang merujuknya! File itu \emph{pasti} ada, atau ada yang sangat salah di departemen ayam-dan-telur.

Mari kita uji parameter \hl{\mbox{\texttt{frotz}}} itu untuk melihat apakah bisa dipakai:

\begin{tcolorbox}[enhanced jigsaw,left=3pt,breakable,colback=white!97.2549!black,boxrule=0.2mm]

\begin{Shaded}
\begin{Highlighting}[]
\PreprocessorTok{\#if \_\_has\_embed(\_\_FILE\_\_ fmc::frotz(lamp))}
\NormalTok{    puts}\OperatorTok{(}\StringTok{"fmc::frotz(lamp) is supported!"}\OperatorTok{);}
\PreprocessorTok{\#else}
\NormalTok{    puts}\OperatorTok{(}\StringTok{"fmc::frotz(lamp) is NOT supported!"}\OperatorTok{);}
\PreprocessorTok{\#endif}
\end{Highlighting}
\end{Shaded}

\end{tcolorbox}

\subsection{Menyematkan Nilai Multi-Byte}\label{menyematkan-nilai-multi-byte}

Bagaimana kalau yang dimasukkan \hl{\mbox{\texttt{int}}}, bukan byte individual? Bagaimana dengan nilai multi-byte di file yang disematkan?

Ini tidak didukung standar C23, tapi mungkin nanti ada ekstensi implementasi yang mendefinisikannya.

\index{\#embed directive@\texttt{\#embed} directive|)}

\section{\texorpdfstring{Directive \hl{\mbox{\texttt{\#pragma}}}}{Directive \#pragma}}\label{pragma}

\index{\#pragma directive@\texttt{\#pragma} directive|(}

Ini directive yang aneh, kependekan dari ``pragmatic''. Kamu bisa memakainya untuk\ldots{} ya, apa pun yang didukung compiler-mu.

Pada dasarnya satu-satunya saat kamu menambahkan ini ke kodemu adalah kalau ada dokumentasi yang menyuruhmu melakukannya.

\subsection{Pragma Non-Standar}\label{pragma-non-standar}

\index{\#pragma directive@\texttt{\#pragma} directive!nonstandard pragmas|(}

Berikut satu contoh non-standar memakai \hl{\mbox{\texttt{\#pragma}}} untuk membuat compiler menjalankan loop \hl{\mbox{\texttt{for}}} secara paralel dengan beberapa thread (kalau compiler mendukung ekstensi \href{https://www.openmp.org/}{OpenMP}\footnote{https://www.openmp.org/}):

\begin{tcolorbox}[enhanced jigsaw,left=3pt,breakable,colback=white!97.2549!black,boxrule=0.2mm]

\begin{Shaded}
\begin{Highlighting}[]
\PreprocessorTok{\#pragma omp parallel for}
\ControlFlowTok{for} \OperatorTok{(}\DataTypeTok{int}\NormalTok{ i }\OperatorTok{=} \DecValTok{0}\OperatorTok{;}\NormalTok{ i }\OperatorTok{\textless{}} \DecValTok{10}\OperatorTok{;}\NormalTok{ i}\OperatorTok{++)} \OperatorTok{\{} \OperatorTok{...} \OperatorTok{\}}
\end{Highlighting}
\end{Shaded}

\end{tcolorbox}

Ada berbagai macam directive \hl{\mbox{\texttt{\#pragma}}} yang terdokumentasi di seluruh penjuru dunia.

Semua \hl{\mbox{\texttt{\#pragma}}} yang tidak dikenali diabaikan oleh compiler.

\index{\#pragma directive@\texttt{\#pragma} directive!nonstandard pragmas|)}

\subsection{Pragma Standar}\label{pragma-standar}

Ada juga beberapa yang standar, dan ini diawali \hl{\mbox{\texttt{STDC}}}, dengan bentuk yang sama:

\begin{tcolorbox}[enhanced jigsaw,left=3pt,breakable,colback=white!97.2549!black,boxrule=0.2mm]

\begin{Shaded}
\begin{Highlighting}[]
\PreprocessorTok{\#pragma STDC pragma\_name on{-}off}
\end{Highlighting}
\end{Shaded}

\end{tcolorbox}

Bagian \hl{\mbox{\texttt{on-off}}} bisa berupa \hl{\mbox{\texttt{ON}}}, \hl{\mbox{\texttt{OFF}}}, atau \hl{\mbox{\texttt{DEFAULT}}}.

Dan \hl{\mbox{\texttt{pragma\_name}}} bisa salah satu dari ini:

\index{FP\_CONTRACT pragma@\texttt{FP\_CONTRACT} pragma|(} \index{CX\_LIMITED\_RANGE pragma@\texttt{CX\_LIMITED\_RANGE} pragma|(}

\index{FENV\_ACCESS pragma@\texttt{FENV\_ACCESS} pragma}

{\def\LTcaptype{none} % do not increment counter
\begin{longtable}[]{@{}
  >{\raggedright\arraybackslash}p{(\linewidth - 2\tabcolsep) * \real{0.2063}}
  >{\raggedright\arraybackslash}p{(\linewidth - 2\tabcolsep) * \real{0.7937}}@{}}
\toprule\noalign{}
\begin{minipage}[b]{\linewidth}\raggedright
Nama Pragma
\end{minipage} & \begin{minipage}[b]{\linewidth}\raggedright
Deskripsi
\end{minipage} \\
\midrule\noalign{}
\endhead
\bottomrule\noalign{}
\endlastfoot
\hl{\mbox{\texttt{FP\_CONTRACT}}} & Mengizinkan ekspresi floating point diringkas menjadi satu operasi untuk menghindari galat pembulatan yang mungkin terjadi dari beberapa operasi. \\
\hl{\mbox{\texttt{FENV\_ACCESS}}} & Atur ke \hl{\mbox{\texttt{ON}}} jika kamu berencana mengakses flag status floating point. Jika \hl{\mbox{\texttt{OFF}}}, compiler mungkin melakukan optimasi yang membuat nilai di flag tidak konsisten atau tidak valid. \\
\hl{\mbox{\texttt{CX\_LIMITED\_RANGE}}} & Atur ke \hl{\mbox{\texttt{ON}}} untuk mengizinkan compiler melewatkan pengecekan overflow saat melakukan aritmetika kompleks. Bawaannya \hl{\mbox{\texttt{OFF}}}. \\
\end{longtable}
}

Misalnya:

\begin{tcolorbox}[enhanced jigsaw,left=3pt,breakable,colback=white!97.2549!black,boxrule=0.2mm]

\begin{Shaded}
\begin{Highlighting}[]
\PreprocessorTok{\#pragma STDC FP\_CONTRACT OFF}
\PreprocessorTok{\#pragma STDC CX\_LIMITED\_RANGE ON}
\end{Highlighting}
\end{Shaded}

\end{tcolorbox}

\index{FP\_CONTRACT pragma@\texttt{FP\_CONTRACT} pragma|)}

Soal \hl{\mbox{\texttt{CX\_LIMITED\_RANGE}}}, spesifikasi menunjukkan:

\begin{tcolorbox}[blanker,colback=white,enhanced jigsaw,breakable,arc=0pt,outer arc=0pt,grow to left by=-1cm,enlarge top by=0.1cm,enlarge bottom by=0.1cm,borderline west={1.5pt}{-.3cm}{black}]

\begin{quote}
The purpose of the pragma is to allow the implementation to use the formulas:

\((x+iy)\times(u+iv) = (xu-yv)+i(yu+xv)\)

\((x+iy)/(u+iv) = [(xu+yv)+i(yu-xv)]/(u^2+v^2)\)

\(|x+iy|=\sqrt{x^2+y^2}\)

where the programmer can determine they are safe.
\end{quote}

\end{tcolorbox}

\index{CX\_LIMITED\_RANGE pragma@\texttt{CX\_LIMITED\_RANGE} pragma|)}

\subsection{\texorpdfstring{Operator \hl{\mbox{\texttt{\_Pragma}}}}{Operator \_Pragma}}\label{operator-_pragma}

\index{\_Pragma operator@\texttt{\_Pragma} operator|(}

Ini cara lain untuk mendeklarasikan pragma yang bisa dipakai di makro.

Keduanya setara:

\begin{tcolorbox}[enhanced jigsaw,left=3pt,breakable,colback=white!97.2549!black,boxrule=0.2mm]

\begin{Shaded}
\begin{Highlighting}[]
\PreprocessorTok{\#pragma "Unnecessary" quotes}
\NormalTok{\_Pragma}\OperatorTok{(}\StringTok{"}\SpecialCharTok{\textbackslash{}"}\StringTok{Unnecessary}\SpecialCharTok{\textbackslash{}"}\StringTok{ quotes"}\OperatorTok{)}
\end{Highlighting}
\end{Shaded}

\end{tcolorbox}

\index{\_Pragma operator@\texttt{\_Pragma} operator!in a macro|(}

Ini bisa dipakai di makro, kalau perlu:

\begin{tcolorbox}[enhanced jigsaw,left=3pt,breakable,colback=white!97.2549!black,boxrule=0.2mm]

\begin{Shaded}
\begin{Highlighting}[]
\PreprocessorTok{\#define PRAGMA}\OperatorTok{(}\PreprocessorTok{x}\OperatorTok{)}\PreprocessorTok{ \_Pragma}\OperatorTok{(\#}\PreprocessorTok{x}\OperatorTok{)}
\end{Highlighting}
\end{Shaded}

\end{tcolorbox}

\index{\_Pragma operator@\texttt{\_Pragma} operator!in a macro|)} \index{\_Pragma operator@\texttt{\_Pragma} operator|)} \index{\#pragma directive@\texttt{\#pragma} directive|)}

\section{\texorpdfstring{Directive \hl{\mbox{\texttt{\#line}}}}{Directive \#line}}\label{directive-line}

\index{\#line directive@\texttt{\#line} directive|(} \index{\_\_LINE\_\_ macro@\texttt{\_\_LINE\_\_} macro|(}

Ini memungkinkan kamu menimpa nilai \hl{\mbox{\texttt{\_\_LINE\_\_}}} dan \hl{\mbox{\texttt{\_\_FILE\_\_}}}. Kalau mau.

Aku belum pernah ingin melakukannya, tapi di K\&R2 mereka menulis:

\begin{tcolorbox}[blanker,colback=white,enhanced jigsaw,breakable,arc=0pt,outer arc=0pt,grow to left by=-1cm,enlarge top by=0.1cm,enlarge bottom by=0.1cm,borderline west={1.5pt}{-.3cm}{black}]

\begin{quote}
For the benefit of other preprocessors that generate C programs {[}\ldots{]}
\end{quote}

\end{tcolorbox}

Jadi mungkin itulah alasannya.

Untuk menimpa nomor baris menjadi, misalnya, 300:

\begin{tcolorbox}[enhanced jigsaw,left=3pt,breakable,colback=white!97.2549!black,boxrule=0.2mm]

\begin{Shaded}
\begin{Highlighting}[]
\PreprocessorTok{\#line 300}
\end{Highlighting}
\end{Shaded}

\end{tcolorbox}

dan \hl{\mbox{\texttt{\_\_LINE\_\_}}} akan terus menghitung naik dari sana.

\index{\_\_LINE\_\_ macro@\texttt{\_\_LINE\_\_} macro|)}

Untuk menimpa nomor baris dan nama file:

\begin{tcolorbox}[enhanced jigsaw,left=3pt,breakable,colback=white!97.2549!black,boxrule=0.2mm]

\begin{Shaded}
\begin{Highlighting}[]
\PreprocessorTok{\#line 300 "newfilename"}
\end{Highlighting}
\end{Shaded}

\end{tcolorbox}

\index{\#line directive@\texttt{\#line} directive|)}

\section{Null Directive}\label{null-directive}

\index{\# null directive@\texttt{\#} null directive|(}

Tanda \hl{\mbox{\texttt{\#}}} sendirian di satu baris diabaikan oleh preprocessor. Nah, jujur saja, aku tidak tahu apa kasus pemakaiannya.

Aku pernah melihat contoh seperti ini:

\begin{tcolorbox}[enhanced jigsaw,left=3pt,breakable,colback=white!97.2549!black,boxrule=0.2mm]

\begin{Shaded}
\begin{Highlighting}[]
\PreprocessorTok{\#ifdef FOO}
    \ErrorTok{\#}
\PreprocessorTok{\#else}
\NormalTok{    printf}\OperatorTok{(}\StringTok{"Something"}\OperatorTok{);}
\PreprocessorTok{\#endif}
\end{Highlighting}
\end{Shaded}

\end{tcolorbox}

yang hanya kosmetik; baris dengan \hl{\mbox{\texttt{\#}}} sendirian itu bisa dihapus tanpa efek buruk.

Atau mungkin demi konsistensi kosmetik, seperti ini:

\begin{tcolorbox}[enhanced jigsaw,left=3pt,breakable,colback=white!97.2549!black,boxrule=0.2mm]

\begin{Shaded}
\begin{Highlighting}[]
\ErrorTok{\#}
\PreprocessorTok{\#ifdef FOO}
\NormalTok{    x }\OperatorTok{=} \DecValTok{2}\OperatorTok{;}
\PreprocessorTok{\#endif}
\ErrorTok{\#}
\PreprocessorTok{\#if BAR == 17}
\NormalTok{    x }\OperatorTok{=} \DecValTok{12}\OperatorTok{;}
\PreprocessorTok{\#endif}
\ErrorTok{\#}
\end{Highlighting}
\end{Shaded}

\end{tcolorbox}

Tapi, soal kosmetik, itu jelek.

Satu postingan lain menyebut penghapusan komentar---bahwa di GCC, komentar setelah \hl{\mbox{\texttt{\#}}} tidak akan dilihat compiler. Itu tidak kuragukan, tapi spesifikasi tampaknya tidak menyebut ini sebagai perilaku standar.

Pencarianku akan alasannya tidak membuahkan banyak hasil. Jadi aku akan bilang saja bahwa ini adalah esoterika C kuno yang sesungguhnya.

\index{\# null directive@\texttt{\#} null directive|)} \index{Preprocessor|)}

\chapter{\texorpdfstring{\hl{\mbox{\texttt{struct}}} II: Lebih Seru dengan \hl{\mbox{\texttt{struct}}}}{struct II: Lebih Seru dengan struct}}\label{struct-ii-lebih-seru-dengan-struct}

\index{struct keyword@\texttt{struct} keyword|(}

Ternyata ada jauh lebih banyak yang bisa dilakukan dengan \hl{\mbox{\texttt{struct}}} daripada yang sudah kita bahas, tapi isinya hanya tumpukan besar hal-hal serba-serbi. Jadi kita kumpulkan saja di bab ini.

Kalau kamu sudah nyaman dengan dasar-dasar \hl{\mbox{\texttt{struct}}}, kamu bisa melengkapi pengetahuanmu di sini.

\section{\texorpdfstring{Initializer untuk \hl{\mbox{\texttt{struct}}} Bersarang dan Array}{Initializer untuk struct Bersarang dan Array}}\label{initializer-untuk-struct-bersarang-dan-array}

\index{struct keyword@\texttt{struct} keyword!initializers|(}

Ingat bagaimana kamu bisa \hyperref[struct-initializers]{menginisialisasi anggota struktur dengan cara seperti ini}?

\begin{tcolorbox}[enhanced jigsaw,left=3pt,breakable,colback=white!97.2549!black,boxrule=0.2mm]

\begin{Shaded}
\begin{Highlighting}[]
\KeywordTok{struct}\NormalTok{ foo x }\OperatorTok{=} \OperatorTok{\{.}\NormalTok{a}\OperatorTok{=}\DecValTok{12}\OperatorTok{,} \OperatorTok{.}\NormalTok{b}\OperatorTok{=}\FloatTok{3.14}\OperatorTok{\};}
\end{Highlighting}
\end{Shaded}

\end{tcolorbox}

Ternyata initializer ini punya kemampuan lebih daripada yang awalnya kita sampaikan. Seru!

Untuk satu hal, kalau ada substruktur bersarang seperti berikut, kamu bisa menginisialisasi anggota substruktur itu dengan mengikuti nama variabel berurutan ke bawah:

\begin{tcolorbox}[enhanced jigsaw,left=3pt,breakable,colback=white!97.2549!black,boxrule=0.2mm]

\begin{Shaded}
\begin{Highlighting}[]
\KeywordTok{struct}\NormalTok{ foo x }\OperatorTok{=} \OperatorTok{\{.}\NormalTok{a}\OperatorTok{.}\NormalTok{b}\OperatorTok{.}\NormalTok{c}\OperatorTok{=}\DecValTok{12}\OperatorTok{\};}
\end{Highlighting}
\end{Shaded}

\end{tcolorbox}

Kita lihat contohnya:

\begin{tcolorbox}[enhanced jigsaw,left=3pt,breakable,colback=white!97.2549!black,boxrule=0.2mm]

\begin{Shaded}
\begin{Highlighting}[numbers=left,,]
\PreprocessorTok{\#include }\ImportTok{\textless{}stdio.h\textgreater{}}

\KeywordTok{struct}\NormalTok{ cabin\_information }\OperatorTok{\{}
    \DataTypeTok{int}\NormalTok{ window\_count}\OperatorTok{;}
    \DataTypeTok{int}\NormalTok{ o2level}\OperatorTok{;}
\OperatorTok{\};}

\KeywordTok{struct}\NormalTok{ spaceship }\OperatorTok{\{}
    \DataTypeTok{char} \OperatorTok{*}\NormalTok{manufacturer}\OperatorTok{;}
    \KeywordTok{struct}\NormalTok{ cabin\_information ci}\OperatorTok{;}
\OperatorTok{\};}

\DataTypeTok{int}\NormalTok{ main}\OperatorTok{(}\DataTypeTok{void}\OperatorTok{)}
\OperatorTok{\{}
    \KeywordTok{struct}\NormalTok{ spaceship s }\OperatorTok{=} \OperatorTok{\{}
        \OperatorTok{.}\NormalTok{manufacturer}\OperatorTok{=}\StringTok{"General Products"}\OperatorTok{,}
        \OperatorTok{.}\NormalTok{ci}\OperatorTok{.}\NormalTok{window\_count }\OperatorTok{=} \DecValTok{8}\OperatorTok{,}   \CommentTok{// \textless{}{-}{-} INITIALIZER BERSARANG!}
        \OperatorTok{.}\NormalTok{ci}\OperatorTok{.}\NormalTok{o2level }\OperatorTok{=} \DecValTok{21}
    \OperatorTok{\};}

\NormalTok{    printf}\OperatorTok{(}\StringTok{"}\SpecialCharTok{\%s}\StringTok{: }\SpecialCharTok{\%d}\StringTok{ seats, }\SpecialCharTok{\%d\%\%}\StringTok{ oxygen}\SpecialCharTok{\textbackslash{}n}\StringTok{"}\OperatorTok{,}
\NormalTok{        s}\OperatorTok{.}\NormalTok{manufacturer}\OperatorTok{,}\NormalTok{ s}\OperatorTok{.}\NormalTok{ci}\OperatorTok{.}\NormalTok{window\_count}\OperatorTok{,}\NormalTok{ s}\OperatorTok{.}\NormalTok{ci}\OperatorTok{.}\NormalTok{o2level}\OperatorTok{);}
\OperatorTok{\}}
\end{Highlighting}
\end{Shaded}

\end{tcolorbox}

Lihat baris 16-17! Di situlah kita menginisialisasi anggota \hl{\mbox{\texttt{struct\ cabin\_information}}} dalam definisi \hl{\mbox{\texttt{s}}}, \hl{\mbox{\texttt{struct\ spaceship}}} kita.

Dan ini opsi lain untuk initializer yang sama---kali ini kita buat yang tampilannya lebih standar, tapi kedua cara berfungsi:

\begin{tcolorbox}[enhanced jigsaw,left=3pt,breakable,colback=white!97.2549!black,boxrule=0.2mm]

\begin{Shaded}
\begin{Highlighting}[numbers=left,,firstnumber=15,]
    \KeywordTok{struct}\NormalTok{ spaceship s }\OperatorTok{=} \OperatorTok{\{}
        \OperatorTok{.}\NormalTok{manufacturer}\OperatorTok{=}\StringTok{"General Products"}\OperatorTok{,}
        \OperatorTok{.}\NormalTok{ci}\OperatorTok{=\{}
            \OperatorTok{.}\NormalTok{window\_count }\OperatorTok{=} \DecValTok{8}\OperatorTok{,}
            \OperatorTok{.}\NormalTok{o2level }\OperatorTok{=} \DecValTok{21}
        \OperatorTok{\}}
    \OperatorTok{\};}
\end{Highlighting}
\end{Shaded}

\end{tcolorbox}

Nah, seolah-olah informasi di atas belum cukup spektakuler, kita juga bisa mencampurkan array initializer di dalamnya.

Mari kita ubah contohnya supaya ada array informasi penumpang di dalamnya, dan kita lihat bagaimana initializer bekerja di sana juga.

\begin{tcolorbox}[enhanced jigsaw,left=3pt,breakable,colback=white!97.2549!black,boxrule=0.2mm]

\begin{Shaded}
\begin{Highlighting}[numbers=left,,]
\PreprocessorTok{\#include }\ImportTok{\textless{}stdio.h\textgreater{}}

\KeywordTok{struct}\NormalTok{ passenger }\OperatorTok{\{}
    \DataTypeTok{char} \OperatorTok{*}\NormalTok{name}\OperatorTok{;}
    \DataTypeTok{int}\NormalTok{ covid\_vaccinated}\OperatorTok{;} \CommentTok{// Boolean}
\OperatorTok{\};}

\PreprocessorTok{\#define MAX\_PASSENGERS }\DecValTok{8}

\KeywordTok{struct}\NormalTok{ spaceship }\OperatorTok{\{}
    \DataTypeTok{char} \OperatorTok{*}\NormalTok{manufacturer}\OperatorTok{;}
    \KeywordTok{struct}\NormalTok{ passenger passenger}\OperatorTok{[}\NormalTok{MAX\_PASSENGERS}\OperatorTok{];}
\OperatorTok{\};}

\DataTypeTok{int}\NormalTok{ main}\OperatorTok{(}\DataTypeTok{void}\OperatorTok{)}
\OperatorTok{\{}
    \KeywordTok{struct}\NormalTok{ spaceship s }\OperatorTok{=} \OperatorTok{\{}
        \OperatorTok{.}\NormalTok{manufacturer}\OperatorTok{=}\StringTok{"General Products"}\OperatorTok{,}
        \OperatorTok{.}\NormalTok{passenger }\OperatorTok{=} \OperatorTok{\{}
            \CommentTok{// Inisialisasi satu field per satu}
            \OperatorTok{[}\DecValTok{0}\OperatorTok{].}\NormalTok{name }\OperatorTok{=} \StringTok{"Gridley, Lewis"}\OperatorTok{,}
            \OperatorTok{[}\DecValTok{0}\OperatorTok{].}\NormalTok{covid\_vaccinated }\OperatorTok{=} \DecValTok{0}\OperatorTok{,}

            \CommentTok{// Atau sekaligus}
            \OperatorTok{[}\DecValTok{7}\OperatorTok{]} \OperatorTok{=} \OperatorTok{\{.}\NormalTok{name}\OperatorTok{=}\StringTok{"Brown, Teela"}\OperatorTok{,} \OperatorTok{.}\NormalTok{covid\_vaccinated}\OperatorTok{=}\DecValTok{1}\OperatorTok{\},}
        \OperatorTok{\}}
    \OperatorTok{\};}

\NormalTok{    printf}\OperatorTok{(}\StringTok{"Passengers for }\SpecialCharTok{\%s}\StringTok{ ship:}\SpecialCharTok{\textbackslash{}n}\StringTok{"}\OperatorTok{,}\NormalTok{ s}\OperatorTok{.}\NormalTok{manufacturer}\OperatorTok{);}

    \ControlFlowTok{for} \OperatorTok{(}\DataTypeTok{int}\NormalTok{ i }\OperatorTok{=} \DecValTok{0}\OperatorTok{;}\NormalTok{ i }\OperatorTok{\textless{}}\NormalTok{ MAX\_PASSENGERS}\OperatorTok{;}\NormalTok{ i}\OperatorTok{++)}
        \ControlFlowTok{if} \OperatorTok{(}\NormalTok{s}\OperatorTok{.}\NormalTok{passenger}\OperatorTok{[}\NormalTok{i}\OperatorTok{].}\NormalTok{name }\OperatorTok{!=}\NormalTok{ NULL}\OperatorTok{)}
\NormalTok{            printf}\OperatorTok{(}\StringTok{"    }\SpecialCharTok{\%s}\StringTok{ (}\SpecialCharTok{\%s}\StringTok{vaccinated)}\SpecialCharTok{\textbackslash{}n}\StringTok{"}\OperatorTok{,}
\NormalTok{                s}\OperatorTok{.}\NormalTok{passenger}\OperatorTok{[}\NormalTok{i}\OperatorTok{].}\NormalTok{name}\OperatorTok{,}
\NormalTok{                s}\OperatorTok{.}\NormalTok{passenger}\OperatorTok{[}\NormalTok{i}\OperatorTok{].}\NormalTok{covid\_vaccinated}\OperatorTok{?} \StringTok{""}\OperatorTok{:} \StringTok{"not "}\OperatorTok{);}
\OperatorTok{\}}
\end{Highlighting}
\end{Shaded}

\end{tcolorbox}

\index{struct keyword@\texttt{struct} keyword!initializers|)}

\section{\texorpdfstring{\hl{\mbox{\texttt{struct}}} Anonim}{struct Anonim}}\label{struct-anonim}

\index{struct keyword@\texttt{struct} keyword!anonymous|(}

Inilah ``\hl{\mbox{\texttt{struct}}} tanpa nama''. Kita juga menyinggungnya di bagian \hyperref[typedef-struct]{\hl{\mbox{\texttt{typedef}}}}, tapi kita segarkan lagi di sini.

Ini \hl{\mbox{\texttt{struct}}} biasa:

\begin{tcolorbox}[enhanced jigsaw,left=3pt,breakable,colback=white!97.2549!black,boxrule=0.2mm]

\begin{Shaded}
\begin{Highlighting}[]
\KeywordTok{struct}\NormalTok{ animal }\OperatorTok{\{}
    \DataTypeTok{char} \OperatorTok{*}\NormalTok{name}\OperatorTok{;}
    \DataTypeTok{int}\NormalTok{ leg\_count}\OperatorTok{,}\NormalTok{ speed}\OperatorTok{;}
\OperatorTok{\};}
\end{Highlighting}
\end{Shaded}

\end{tcolorbox}

Dan ini padanannya yang anonim:

\begin{tcolorbox}[enhanced jigsaw,left=3pt,breakable,colback=white!97.2549!black,boxrule=0.2mm]

\begin{Shaded}
\begin{Highlighting}[]
\KeywordTok{struct} \OperatorTok{\{}              \CommentTok{// \textless{}{-}{-} Tanpa nama!}
    \DataTypeTok{char} \OperatorTok{*}\NormalTok{name}\OperatorTok{;}
    \DataTypeTok{int}\NormalTok{ leg\_count}\OperatorTok{,}\NormalTok{ speed}\OperatorTok{;}
\OperatorTok{\};}
\end{Highlighting}
\end{Shaded}

\end{tcolorbox}

Okeeee. Jadi kita punya \hl{\mbox{\texttt{struct}}}, tapi tanpa nama, sehingga tidak ada cara memakainya nanti? Kelihatannya sia-sia.

Harus diakui, pada contoh itu memang begitu. Tapi kita masih bisa memanfaatkannya dengan beberapa cara.

Satu cara yang jarang: karena \hl{\mbox{\texttt{struct}}} anonim merepresentasikan sebuah tipe, kita bisa menaruh nama variabel setelahnya dan memakainya.

\begin{tcolorbox}[enhanced jigsaw,left=3pt,breakable,colback=white!97.2549!black,boxrule=0.2mm]

\begin{Shaded}
\begin{Highlighting}[]
\KeywordTok{struct} \OperatorTok{\{}                   \CommentTok{// \textless{}{-}{-} Tanpa nama!}
    \DataTypeTok{char} \OperatorTok{*}\NormalTok{name}\OperatorTok{;}
    \DataTypeTok{int}\NormalTok{ leg\_count}\OperatorTok{,}\NormalTok{ speed}\OperatorTok{;}
\OperatorTok{\}}\NormalTok{ a}\OperatorTok{,}\NormalTok{ b}\OperatorTok{,}\NormalTok{ c}\OperatorTok{;}                 \CommentTok{// 3 variabel bertipe struct ini}

\NormalTok{a}\OperatorTok{.}\NormalTok{name }\OperatorTok{=} \StringTok{"antelope"}\OperatorTok{;}
\NormalTok{c}\OperatorTok{.}\NormalTok{leg\_count }\OperatorTok{=} \DecValTok{4}\OperatorTok{;}           \CommentTok{// misalnya}
\end{Highlighting}
\end{Shaded}

\end{tcolorbox}

Tapi itu pun belum terlalu berguna.

Jauh lebih umum adalah memakai \hl{\mbox{\texttt{struct}}} anonim bersama \hl{\mbox{\texttt{typedef}}} supaya bisa dipakai nanti (misalnya untuk mengoper variabel ke fungsi).

\begin{tcolorbox}[enhanced jigsaw,left=3pt,breakable,colback=white!97.2549!black,boxrule=0.2mm]

\begin{Shaded}
\begin{Highlighting}[]
\KeywordTok{typedef} \KeywordTok{struct} \OperatorTok{\{}                   \CommentTok{// \textless{}{-}{-} Tanpa nama!}
    \DataTypeTok{char} \OperatorTok{*}\NormalTok{name}\OperatorTok{;}
    \DataTypeTok{int}\NormalTok{ leg\_count}\OperatorTok{,}\NormalTok{ speed}\OperatorTok{;}
\OperatorTok{\}}\NormalTok{ animal}\OperatorTok{;}                          \CommentTok{// Tipe baru: animal}

\NormalTok{animal a}\OperatorTok{,}\NormalTok{ b}\OperatorTok{,}\NormalTok{ c}\OperatorTok{;}

\NormalTok{a}\OperatorTok{.}\NormalTok{name }\OperatorTok{=} \StringTok{"antelope"}\OperatorTok{;}
\NormalTok{c}\OperatorTok{.}\NormalTok{leg\_count }\OperatorTok{=} \DecValTok{4}\OperatorTok{;}           \CommentTok{// misalnya}
\end{Highlighting}
\end{Shaded}

\end{tcolorbox}

Secara pribadi, aku tidak banyak memakai \hl{\mbox{\texttt{struct}}} anonim. Menurutku lebih enak melihat seluruh \hl{\mbox{\texttt{struct\ animal}}} sebelum nama variabel pada deklarasi.

Tapi itu cuma, ya, pendapatku, kawan.

\index{struct keyword@\texttt{struct} keyword!anonymous|)}

\section{\texorpdfstring{\hl{\mbox{\texttt{struct}}} yang Merujuk Diri Sendiri}{struct yang Merujuk Diri Sendiri}}\label{struct-yang-merujuk-diri-sendiri}

\index{struct keyword@\texttt{struct} keyword!self-referential|(}

Untuk struktur data apa pun yang mirip graf, berguna sekali punya pointer ke node/vertex yang terhubung. Tapi ini berarti pada definisi sebuah node, kamu perlu punya pointer ke node. Seperti ayam dan telur!

Tapi ternyata di C ini bisa dilakukan tanpa masalah sama sekali.

Misalnya, ini node linked list:

\begin{tcolorbox}[enhanced jigsaw,left=3pt,breakable,colback=white!97.2549!black,boxrule=0.2mm]

\begin{Shaded}
\begin{Highlighting}[]
\KeywordTok{struct}\NormalTok{ node }\OperatorTok{\{}
    \DataTypeTok{int}\NormalTok{ data}\OperatorTok{;}
    \KeywordTok{struct}\NormalTok{ node }\OperatorTok{*}\NormalTok{next}\OperatorTok{;}
\OperatorTok{\};}
\end{Highlighting}
\end{Shaded}

\end{tcolorbox}

Penting dicatat bahwa \hl{\mbox{\texttt{next}}} adalah pointer. Inilah yang membuat semuanya bisa dibangun. Walaupun compiler belum tahu seperti apa seluruh \hl{\mbox{\texttt{struct\ node}}}, semua pointer berukuran sama.

Berikut program linked list ala kadarnya untuk mengujinya:

\begin{tcolorbox}[enhanced jigsaw,left=3pt,breakable,colback=white!97.2549!black,boxrule=0.2mm]

\begin{Shaded}
\begin{Highlighting}[numbers=left,,]
\PreprocessorTok{\#include }\ImportTok{\textless{}stdio.h\textgreater{}}
\PreprocessorTok{\#include }\ImportTok{\textless{}stdlib.h\textgreater{}}

\KeywordTok{struct}\NormalTok{ node }\OperatorTok{\{}
    \DataTypeTok{int}\NormalTok{ data}\OperatorTok{;}
    \KeywordTok{struct}\NormalTok{ node }\OperatorTok{*}\NormalTok{next}\OperatorTok{;}
\OperatorTok{\};}

\DataTypeTok{int}\NormalTok{ main}\OperatorTok{(}\DataTypeTok{void}\OperatorTok{)}
\OperatorTok{\{}
    \KeywordTok{struct}\NormalTok{ node }\OperatorTok{*}\NormalTok{head}\OperatorTok{;}

    \CommentTok{// Susun linked list secara asal{-}asalan (11){-}\textgreater{}(22){-}\textgreater{}(33)}
\NormalTok{    head }\OperatorTok{=}\NormalTok{ malloc}\OperatorTok{(}\KeywordTok{sizeof}\OperatorTok{(}\KeywordTok{struct}\NormalTok{ node}\OperatorTok{));}
\NormalTok{    head}\OperatorTok{{-}\textgreater{}}\NormalTok{data }\OperatorTok{=} \DecValTok{11}\OperatorTok{;}
\NormalTok{    head}\OperatorTok{{-}\textgreater{}}\NormalTok{next }\OperatorTok{=}\NormalTok{ malloc}\OperatorTok{(}\KeywordTok{sizeof}\OperatorTok{(}\KeywordTok{struct}\NormalTok{ node}\OperatorTok{));}
\NormalTok{    head}\OperatorTok{{-}\textgreater{}}\NormalTok{next}\OperatorTok{{-}\textgreater{}}\NormalTok{data }\OperatorTok{=} \DecValTok{22}\OperatorTok{;}
\NormalTok{    head}\OperatorTok{{-}\textgreater{}}\NormalTok{next}\OperatorTok{{-}\textgreater{}}\NormalTok{next }\OperatorTok{=}\NormalTok{ malloc}\OperatorTok{(}\KeywordTok{sizeof}\OperatorTok{(}\KeywordTok{struct}\NormalTok{ node}\OperatorTok{));}
\NormalTok{    head}\OperatorTok{{-}\textgreater{}}\NormalTok{next}\OperatorTok{{-}\textgreater{}}\NormalTok{next}\OperatorTok{{-}\textgreater{}}\NormalTok{data }\OperatorTok{=} \DecValTok{33}\OperatorTok{;}
\NormalTok{    head}\OperatorTok{{-}\textgreater{}}\NormalTok{next}\OperatorTok{{-}\textgreater{}}\NormalTok{next}\OperatorTok{{-}\textgreater{}}\NormalTok{next }\OperatorTok{=}\NormalTok{ NULL}\OperatorTok{;}

    \CommentTok{// Telusuri}
    \ControlFlowTok{for} \OperatorTok{(}\KeywordTok{struct}\NormalTok{ node }\OperatorTok{*}\NormalTok{cur }\OperatorTok{=}\NormalTok{ head}\OperatorTok{;}\NormalTok{ cur }\OperatorTok{!=}\NormalTok{ NULL}\OperatorTok{;}\NormalTok{ cur }\OperatorTok{=}\NormalTok{ cur}\OperatorTok{{-}\textgreater{}}\NormalTok{next}\OperatorTok{)} \OperatorTok{\{}
\NormalTok{        printf}\OperatorTok{(}\StringTok{"}\SpecialCharTok{\%d\textbackslash{}n}\StringTok{"}\OperatorTok{,}\NormalTok{ cur}\OperatorTok{{-}\textgreater{}}\NormalTok{data}\OperatorTok{);}
    \OperatorTok{\}}
\OperatorTok{\}}
\end{Highlighting}
\end{Shaded}

\end{tcolorbox}

Menjalankannya mencetak:

\begin{tcolorbox}[enhanced jigsaw,left=3pt,breakable,colback=white!97.2549!black,boxrule=0.2mm]

\begin{Shaded}
\begin{Highlighting}[]
\NormalTok{11}
\NormalTok{22}
\NormalTok{33}
\end{Highlighting}
\end{Shaded}

\end{tcolorbox}

\index{struct keyword@\texttt{struct} keyword!self-referential|)}

\section{Flexible Array Member}\label{flexible-array-member}

\index{struct keyword@\texttt{struct} keyword!flexible array members|(}

Di zaman dahulu kala, ketika orang masih mengukir kode C dari kayu, sebagian orang berpikir alangkah kerennya kalau bisa mengalokasikan \hl{\mbox{\texttt{struct}}} yang punya array dengan panjang bervariasi di ujungnya.

Aku ingin menegaskan bahwa paruh pertama bagian ini adalah cara lama, dan setelah itu kita akan memakai cara baru.

Misalnya, mungkin kamu bisa mendefinisikan \hl{\mbox{\texttt{struct}}} untuk menyimpan string dan panjang string itu. Isinya sebuah panjang dan sebuah array untuk menampung data. Mungkin seperti ini:

\begin{tcolorbox}[enhanced jigsaw,left=3pt,breakable,colback=white!97.2549!black,boxrule=0.2mm]

\begin{Shaded}
\begin{Highlighting}[]
\KeywordTok{struct}\NormalTok{ len\_string }\OperatorTok{\{}
    \DataTypeTok{int}\NormalTok{ length}\OperatorTok{;}
    \DataTypeTok{char}\NormalTok{ data}\OperatorTok{[}\DecValTok{8}\OperatorTok{];}
\OperatorTok{\};}
\end{Highlighting}
\end{Shaded}

\end{tcolorbox}

Tapi itu menetapkan \hl{\mbox{\texttt{8}}} sebagai panjang maksimum string, dan itu tidak banyak. Bagaimana kalau kita melakukan sesuatu yang \emph{cerdik} dan \hl{\mbox{\texttt{malloc()}}} sedikit ruang ekstra di ujung setelah struct, lalu membiarkan data meluap ke ruang itu?

Mari kita lakukan, lalu alokasikan 40 byte lagi di atasnya:

\begin{tcolorbox}[enhanced jigsaw,left=3pt,breakable,colback=white!97.2549!black,boxrule=0.2mm]

\begin{Shaded}
\begin{Highlighting}[]
\KeywordTok{struct}\NormalTok{ len\_string }\OperatorTok{*}\NormalTok{s }\OperatorTok{=}\NormalTok{ malloc}\OperatorTok{(}\KeywordTok{sizeof} \OperatorTok{*}\NormalTok{s }\OperatorTok{+} \DecValTok{40}\OperatorTok{);}
\end{Highlighting}
\end{Shaded}

\end{tcolorbox}

Karena \hl{\mbox{\texttt{data}}} adalah field terakhir dari \hl{\mbox{\texttt{struct}}}, kalau kita meluapi field itu, ia mengalir ke ruang yang sudah kita alokasikan! Karena alasan ini, trik ini hanya bekerja kalau array pendek itu adalah field \emph{terakhir} dalam \hl{\mbox{\texttt{struct}}}.

\begin{tcolorbox}[enhanced jigsaw,left=3pt,breakable,colback=white!97.2549!black,boxrule=0.2mm]

\begin{Shaded}
\begin{Highlighting}[]
\CommentTok{// Salin lebih dari 8 byte!}

\NormalTok{strcpy}\OperatorTok{(}\NormalTok{s}\OperatorTok{{-}\textgreater{}}\NormalTok{data}\OperatorTok{,} \StringTok{"Hello, world!"}\OperatorTok{);}  \CommentTok{// Tidak akan crash. Mungkin.}
\end{Highlighting}
\end{Shaded}

\end{tcolorbox}

\index{Arrays!zero length|(}

Bahkan, dulu ada solusi umum dari compiler untuk ini, yaitu mengalokasikan array berukuran nol di ujung:

\begin{tcolorbox}[enhanced jigsaw,left=3pt,breakable,colback=white!97.2549!black,boxrule=0.2mm]

\begin{Shaded}
\begin{Highlighting}[]
\KeywordTok{struct}\NormalTok{ len\_string }\OperatorTok{\{}
    \DataTypeTok{int}\NormalTok{ length}\OperatorTok{;}
    \DataTypeTok{char}\NormalTok{ data}\OperatorTok{[}\DecValTok{0}\OperatorTok{];}
\OperatorTok{\};}
\end{Highlighting}
\end{Shaded}

\end{tcolorbox}

Lalu setiap byte ekstra yang kamu alokasikan siap dipakai untuk string itu.

Karena \hl{\mbox{\texttt{data}}} adalah field terakhir dari \hl{\mbox{\texttt{struct}}}, kalau kita meluapi field itu, ia mengalir ke ruang yang sudah kita alokasikan!

\begin{tcolorbox}[enhanced jigsaw,left=3pt,breakable,colback=white!97.2549!black,boxrule=0.2mm]

\begin{Shaded}
\begin{Highlighting}[]
\CommentTok{// Salin lebih dari 8 byte!}

\NormalTok{strcpy}\OperatorTok{(}\NormalTok{s}\OperatorTok{{-}\textgreater{}}\NormalTok{data}\OperatorTok{,} \StringTok{"Hello, world!"}\OperatorTok{);}  \CommentTok{// Tidak akan crash. Mungkin.}
\end{Highlighting}
\end{Shaded}

\end{tcolorbox}

Tapi, tentu saja, mengakses data di luar ujung array itu adalah undefined behavior! Di zaman modern ini, kita tidak lagi sudi memakai kebiadaban seperti itu.

\index{Arrays!zero length|)}

Untungnya, kita masih bisa mendapat efek yang sama dengan C99 dan sesudahnya, tapi sekarang legal.

Kita cukup mengubah definisi di atas agar array tidak punya ukuran\footnote{Secara teknis kita bilang ia punya \emph{incomplete type}.}:

\begin{tcolorbox}[enhanced jigsaw,left=3pt,breakable,colback=white!97.2549!black,boxrule=0.2mm]

\begin{Shaded}
\begin{Highlighting}[]
\KeywordTok{struct}\NormalTok{ len\_string }\OperatorTok{\{}
    \DataTypeTok{int}\NormalTok{ length}\OperatorTok{;}
    \DataTypeTok{char}\NormalTok{ data}\OperatorTok{[];}
\OperatorTok{\};}
\end{Highlighting}
\end{Shaded}

\end{tcolorbox}

Sekali lagi, ini hanya bekerja kalau flexible array member adalah field \emph{terakhir} dalam \hl{\mbox{\texttt{struct}}}.

Lalu kita bisa mengalokasikan ruang sebanyak yang kita mau untuk string-string itu dengan \hl{\mbox{\texttt{malloc()}}} lebih besar dari \hl{\mbox{\texttt{struct\ len\_string}}}, seperti pada contoh ini yang membuat \hl{\mbox{\texttt{struct\ len\_string}}} baru dari string C:

\begin{tcolorbox}[enhanced jigsaw,left=3pt,breakable,colback=white!97.2549!black,boxrule=0.2mm]

\begin{Shaded}
\begin{Highlighting}[]
\KeywordTok{struct}\NormalTok{ len\_string }\OperatorTok{*}\NormalTok{len\_string\_from\_c\_string}\OperatorTok{(}\DataTypeTok{char} \OperatorTok{*}\NormalTok{s}\OperatorTok{)}
\OperatorTok{\{}
    \DataTypeTok{int}\NormalTok{ len }\OperatorTok{=}\NormalTok{ strlen}\OperatorTok{(}\NormalTok{s}\OperatorTok{);}

    \CommentTok{// Alokasikan "len" byte lebih banyak dari yang biasanya dibutuhkan}
    \KeywordTok{struct}\NormalTok{ len\_string }\OperatorTok{*}\NormalTok{ls }\OperatorTok{=}\NormalTok{ malloc}\OperatorTok{(}\KeywordTok{sizeof} \OperatorTok{*}\NormalTok{ls }\OperatorTok{+}\NormalTok{ len}\OperatorTok{);}

\NormalTok{    ls}\OperatorTok{{-}\textgreater{}}\NormalTok{length }\OperatorTok{=}\NormalTok{ len}\OperatorTok{;}

    \CommentTok{// Salin string ke byte{-}byte ekstra itu}
\NormalTok{    memcpy}\OperatorTok{(}\NormalTok{ls}\OperatorTok{{-}\textgreater{}}\NormalTok{data}\OperatorTok{,}\NormalTok{ s}\OperatorTok{,}\NormalTok{ len}\OperatorTok{);}

    \ControlFlowTok{return}\NormalTok{ ls}\OperatorTok{;}
\OperatorTok{\}}
\end{Highlighting}
\end{Shaded}

\end{tcolorbox}

\index{struct keyword@\texttt{struct} keyword!flexible array members|)}

\section{Padding Byte}\label{struct-padding-bytes}

\index{struct keyword@\texttt{struct} keyword!padding bytes|(}

Waspadai bahwa C boleh menambahkan padding byte di dalam atau setelah \hl{\mbox{\texttt{struct}}} sesuai kehendaknya. Kamu tidak bisa memastikan bahwa field-field itu berdekatan langsung di memori\footnote{Walaupun beberapa compiler punya opsi untuk memaksa hal itu terjadi---cari \hl{\mbox{\texttt{\_\_attribute\_\_((packed))}}} untuk melihat caranya di GCC.}.

Mari lihat program ini. Kita mengeluarkan dua angka. Satu adalah jumlah \hl{\mbox{\texttt{sizeof}}} tipe-tipe field individual. Satunya lagi \hl{\mbox{\texttt{sizeof}}} seluruh \hl{\mbox{\texttt{struct}}}.

Orang akan mengira keduanya sama. Ukuran total adalah ukuran jumlah bagian-bagiannya, kan?

\begin{tcolorbox}[enhanced jigsaw,left=3pt,breakable,colback=white!97.2549!black,boxrule=0.2mm]

\begin{Shaded}
\begin{Highlighting}[numbers=left,,]
\PreprocessorTok{\#include }\ImportTok{\textless{}stdio.h\textgreater{}}

\KeywordTok{struct}\NormalTok{ foo }\OperatorTok{\{}
    \DataTypeTok{int}\NormalTok{ a}\OperatorTok{;}
    \DataTypeTok{char}\NormalTok{ b}\OperatorTok{;}
    \DataTypeTok{int}\NormalTok{ c}\OperatorTok{;}
    \DataTypeTok{char}\NormalTok{ d}\OperatorTok{;}
\OperatorTok{\};}

\DataTypeTok{int}\NormalTok{ main}\OperatorTok{(}\DataTypeTok{void}\OperatorTok{)}
\OperatorTok{\{}
\NormalTok{    printf}\OperatorTok{(}\StringTok{"}\SpecialCharTok{\%zu\textbackslash{}n}\StringTok{"}\OperatorTok{,} \KeywordTok{sizeof}\OperatorTok{(}\DataTypeTok{int}\OperatorTok{)} \OperatorTok{+} \KeywordTok{sizeof}\OperatorTok{(}\DataTypeTok{char}\OperatorTok{)} \OperatorTok{+} \KeywordTok{sizeof}\OperatorTok{(}\DataTypeTok{int}\OperatorTok{)} \OperatorTok{+} \KeywordTok{sizeof}\OperatorTok{(}\DataTypeTok{char}\OperatorTok{));}
\NormalTok{    printf}\OperatorTok{(}\StringTok{"}\SpecialCharTok{\%zu\textbackslash{}n}\StringTok{"}\OperatorTok{,} \KeywordTok{sizeof}\OperatorTok{(}\KeywordTok{struct}\NormalTok{ foo}\OperatorTok{));}
\OperatorTok{\}}
\end{Highlighting}
\end{Shaded}

\end{tcolorbox}

Tapi di sistemku, ini mengeluarkan:

\begin{tcolorbox}[enhanced jigsaw,left=3pt,breakable,colback=white!97.2549!black,boxrule=0.2mm]

\begin{Shaded}
\begin{Highlighting}[]
\NormalTok{10}
\NormalTok{16}
\end{Highlighting}
\end{Shaded}

\end{tcolorbox}

Ternyata tidak sama! Compiler menambahkan 6 byte padding agar performanya lebih baik. Mungkin di compiler-mu hasilnya berbeda, tapi kecuali kamu memaksanya, kamu tidak bisa memastikan tidak ada padding.

\index{struct keyword@\texttt{struct} keyword!padding bytes|)}

\section{\texorpdfstring{\hl{\mbox{\texttt{offsetof}}}}{offsetof}}\label{offsetof}

\index{offsetof() macro@\texttt{offsetof()} macro|(}

Di bagian sebelumnya, kita melihat bahwa compiler bisa menyisipkan padding byte sesukanya di dalam struktur.

Bagaimana kalau kita perlu tahu di mana letaknya? Kita bisa mengukurnya dengan \hl{\mbox{\texttt{offsetof}}}, yang didefinisikan di \hl{\mbox{\texttt{\textless{}stddef.h\textgreater{}}}}.

Mari kita ubah kode di atas untuk mencetak offset field-field individual dalam \hl{\mbox{\texttt{struct}}}:

\begin{tcolorbox}[enhanced jigsaw,left=3pt,breakable,colback=white!97.2549!black,boxrule=0.2mm]

\begin{Shaded}
\begin{Highlighting}[numbers=left,,]
\PreprocessorTok{\#include }\ImportTok{\textless{}stdio.h\textgreater{}}
\PreprocessorTok{\#include }\ImportTok{\textless{}stddef.h\textgreater{}}

\KeywordTok{struct}\NormalTok{ foo }\OperatorTok{\{}
    \DataTypeTok{int}\NormalTok{ a}\OperatorTok{;}
    \DataTypeTok{char}\NormalTok{ b}\OperatorTok{;}
    \DataTypeTok{int}\NormalTok{ c}\OperatorTok{;}
    \DataTypeTok{char}\NormalTok{ d}\OperatorTok{;}
\OperatorTok{\};}

\DataTypeTok{int}\NormalTok{ main}\OperatorTok{(}\DataTypeTok{void}\OperatorTok{)}
\OperatorTok{\{}
\NormalTok{    printf}\OperatorTok{(}\StringTok{"}\SpecialCharTok{\%zu\textbackslash{}n}\StringTok{"}\OperatorTok{,}\NormalTok{ offsetof}\OperatorTok{(}\KeywordTok{struct}\NormalTok{ foo}\OperatorTok{,}\NormalTok{ a}\OperatorTok{));}
\NormalTok{    printf}\OperatorTok{(}\StringTok{"}\SpecialCharTok{\%zu\textbackslash{}n}\StringTok{"}\OperatorTok{,}\NormalTok{ offsetof}\OperatorTok{(}\KeywordTok{struct}\NormalTok{ foo}\OperatorTok{,}\NormalTok{ b}\OperatorTok{));}
\NormalTok{    printf}\OperatorTok{(}\StringTok{"}\SpecialCharTok{\%zu\textbackslash{}n}\StringTok{"}\OperatorTok{,}\NormalTok{ offsetof}\OperatorTok{(}\KeywordTok{struct}\NormalTok{ foo}\OperatorTok{,}\NormalTok{ c}\OperatorTok{));}
\NormalTok{    printf}\OperatorTok{(}\StringTok{"}\SpecialCharTok{\%zu\textbackslash{}n}\StringTok{"}\OperatorTok{,}\NormalTok{ offsetof}\OperatorTok{(}\KeywordTok{struct}\NormalTok{ foo}\OperatorTok{,}\NormalTok{ d}\OperatorTok{));}
\OperatorTok{\}}
\end{Highlighting}
\end{Shaded}

\end{tcolorbox}

Di komputerku, ini mengeluarkan:

\begin{tcolorbox}[enhanced jigsaw,left=3pt,breakable,colback=white!97.2549!black,boxrule=0.2mm]

\begin{Shaded}
\begin{Highlighting}[]
\NormalTok{0}
\NormalTok{4}
\NormalTok{8}
\NormalTok{12}
\end{Highlighting}
\end{Shaded}

\end{tcolorbox}

yang menunjukkan bahwa kita memakai 4 byte untuk setiap field. Agak aneh, karena \hl{\mbox{\texttt{char}}} hanya 1 byte, kan? Compiler menaruh 3 padding byte setelah setiap \hl{\mbox{\texttt{char}}} agar semua field berukuran 4 byte. Mungkin ini akan berjalan lebih cepat di CPU-ku.

\index{offsetof() macro@\texttt{offsetof()} macro|)}

\section{Fake OOP}\label{fake-oop}

Ada hal yang agak kurang ajar dan mirip OOP yang bisa dilakukan dengan \hl{\mbox{\texttt{struct}}}.

Karena pointer ke \hl{\mbox{\texttt{struct}}} sama dengan pointer ke elemen pertama \hl{\mbox{\texttt{struct}}}, kamu bebas meng-cast pointer ke \hl{\mbox{\texttt{struct}}} menjadi pointer ke elemen pertama.

Artinya kita bisa menyusun situasi seperti ini:

\begin{tcolorbox}[enhanced jigsaw,left=3pt,breakable,colback=white!97.2549!black,boxrule=0.2mm]

\begin{Shaded}
\begin{Highlighting}[]
\KeywordTok{struct}\NormalTok{ parent }\OperatorTok{\{}
    \DataTypeTok{int}\NormalTok{ a}\OperatorTok{,}\NormalTok{ b}\OperatorTok{;}
\OperatorTok{\};}

\KeywordTok{struct}\NormalTok{ child }\OperatorTok{\{}
    \KeywordTok{struct}\NormalTok{ parent super}\OperatorTok{;}  \CommentTok{// HARUS yang pertama}
    \DataTypeTok{int}\NormalTok{ c}\OperatorTok{,}\NormalTok{ d}\OperatorTok{;}
\OperatorTok{\};}
\end{Highlighting}
\end{Shaded}

\end{tcolorbox}

Lalu kita bisa mengoper pointer ke \hl{\mbox{\texttt{struct\ child}}} ke fungsi yang mengharapkan itu \emph{atau} pointer ke \hl{\mbox{\texttt{struct\ parent}}}!

Karena \hl{\mbox{\texttt{struct\ parent\ super}}} adalah item pertama dalam \hl{\mbox{\texttt{struct\ child}}}, pointer ke \hl{\mbox{\texttt{struct\ child}}} mana pun sama dengan pointer ke field \hl{\mbox{\texttt{super}}} itu\footnote{Omong-omong, \hl{\mbox{\texttt{super}}} bukan kata kunci. Aku cuma meminjam sedikit terminologi OOP.}.

Mari kita buat contohnya. Kita buat \hl{\mbox{\texttt{struct}}} seperti di atas, lalu kita oper pointer ke \hl{\mbox{\texttt{struct\ child}}} ke fungsi yang butuh pointer ke \hl{\mbox{\texttt{struct\ parent}}}\ldots{} dan itu tetap berfungsi.

\begin{tcolorbox}[enhanced jigsaw,left=3pt,breakable,colback=white!97.2549!black,boxrule=0.2mm]

\begin{Shaded}
\begin{Highlighting}[numbers=left,,]
\PreprocessorTok{\#include }\ImportTok{\textless{}stdio.h\textgreater{}}

\KeywordTok{struct}\NormalTok{ parent }\OperatorTok{\{}
    \DataTypeTok{int}\NormalTok{ a}\OperatorTok{,}\NormalTok{ b}\OperatorTok{;}
\OperatorTok{\};}

\KeywordTok{struct}\NormalTok{ child }\OperatorTok{\{}
    \KeywordTok{struct}\NormalTok{ parent super}\OperatorTok{;}  \CommentTok{// HARUS yang pertama}
    \DataTypeTok{int}\NormalTok{ c}\OperatorTok{,}\NormalTok{ d}\OperatorTok{;}
\OperatorTok{\};}

\CommentTok{// Menjadikan argumen \textasciigrave{}void*\textasciigrave{} supaya kita bisa mengoper tipe apa pun}
\CommentTok{// (yaitu struct parent atau struct child)}
\DataTypeTok{void}\NormalTok{ print\_parent}\OperatorTok{(}\DataTypeTok{void} \OperatorTok{*}\NormalTok{p}\OperatorTok{)}
\OperatorTok{\{}
    \CommentTok{// Mengharapkan struct parent{-}{-}tapi struct child juga berfungsi}
    \CommentTok{// karena pointer{-}nya menunjuk ke struct parent di field}
    \CommentTok{// pertama:}
    \KeywordTok{struct}\NormalTok{ parent }\OperatorTok{*}\NormalTok{self }\OperatorTok{=}\NormalTok{ p}\OperatorTok{;}

\NormalTok{    printf}\OperatorTok{(}\StringTok{"Parent: }\SpecialCharTok{\%d}\StringTok{, }\SpecialCharTok{\%d\textbackslash{}n}\StringTok{"}\OperatorTok{,}\NormalTok{ self}\OperatorTok{{-}\textgreater{}}\NormalTok{a}\OperatorTok{,}\NormalTok{ self}\OperatorTok{{-}\textgreater{}}\NormalTok{b}\OperatorTok{);}
\OperatorTok{\}}

\DataTypeTok{void}\NormalTok{ print\_child}\OperatorTok{(}\KeywordTok{struct}\NormalTok{ child }\OperatorTok{*}\NormalTok{self}\OperatorTok{)}
\OperatorTok{\{}
\NormalTok{    printf}\OperatorTok{(}\StringTok{"Child: }\SpecialCharTok{\%d}\StringTok{, }\SpecialCharTok{\%d\textbackslash{}n}\StringTok{"}\OperatorTok{,}\NormalTok{ self}\OperatorTok{{-}\textgreater{}}\NormalTok{c}\OperatorTok{,}\NormalTok{ self}\OperatorTok{{-}\textgreater{}}\NormalTok{d}\OperatorTok{);}
\OperatorTok{\}}

\DataTypeTok{int}\NormalTok{ main}\OperatorTok{(}\DataTypeTok{void}\OperatorTok{)}
\OperatorTok{\{}
    \KeywordTok{struct}\NormalTok{ child c }\OperatorTok{=} \OperatorTok{\{.}\NormalTok{super}\OperatorTok{.}\NormalTok{a}\OperatorTok{=}\DecValTok{1}\OperatorTok{,} \OperatorTok{.}\NormalTok{super}\OperatorTok{.}\NormalTok{b}\OperatorTok{=}\DecValTok{2}\OperatorTok{,} \OperatorTok{.}\NormalTok{c}\OperatorTok{=}\DecValTok{3}\OperatorTok{,} \OperatorTok{.}\NormalTok{d}\OperatorTok{=}\DecValTok{4}\OperatorTok{\};}

\NormalTok{    print\_child}\OperatorTok{(\&}\NormalTok{c}\OperatorTok{);}
\NormalTok{    print\_parent}\OperatorTok{(\&}\NormalTok{c}\OperatorTok{);}  \CommentTok{// Juga berfungsi meskipun ini struct child!}
\OperatorTok{\}}
\end{Highlighting}
\end{Shaded}

\end{tcolorbox}

Lihat apa yang kita lakukan di baris terakhir \hl{\mbox{\texttt{main()}}}? Kita memanggil \hl{\mbox{\texttt{print\_parent()}}} tapi mengoper \hl{\mbox{\texttt{struct\ child*}}} sebagai argumen! Walaupun \hl{\mbox{\texttt{print\_parent()}}} butuh argumen yang menunjuk ke \hl{\mbox{\texttt{struct\ parent}}}, kita \emph{lolos} karena field pertama dalam \hl{\mbox{\texttt{struct\ child}}} adalah \hl{\mbox{\texttt{struct\ parent}}}.

Sekali lagi, ini berfungsi karena pointer ke \hl{\mbox{\texttt{struct}}} bernilai sama dengan pointer ke field pertama dalam \hl{\mbox{\texttt{struct}}} itu.

Semua ini bergantung pada bagian spesifikasi ini:

\begin{tcolorbox}[blanker,colback=white,enhanced jigsaw,breakable,arc=0pt,outer arc=0pt,grow to left by=-1cm,enlarge top by=0.1cm,enlarge bottom by=0.1cm,borderline west={1.5pt}{-.3cm}{black}]

\begin{quote}
\textbf{§6.7.2.1¶15} {[}\ldots{]} A pointer to a structure object, suitably converted, points to its initial member {[}\ldots{]}, and vice versa.
\end{quote}

\end{tcolorbox}

dan

\begin{tcolorbox}[blanker,colback=white,enhanced jigsaw,breakable,arc=0pt,outer arc=0pt,grow to left by=-1cm,enlarge top by=0.1cm,enlarge bottom by=0.1cm,borderline west={1.5pt}{-.3cm}{black}]

\begin{quote}
\textbf{§6.5¶7} An object shall have its stored value accessed only by an lvalue expression that has one of the following types:

\begin{itemize}
\tightlist
\item
  a type compatible with the effective type of the object
\item
  {[}\ldots{]}
\end{itemize}
\end{quote}

\end{tcolorbox}

serta asumsiku bahwa ``suitably converted'' berarti ``di-cast ke effective type dari anggota pertama''.

\section{Bit-Field}\label{bit-field}

\index{struct keyword@\texttt{struct} keyword!bit fields|(}

Menurut pengalamanku, ini jarang dipakai, tapi kamu mungkin melihatnya sesekali, terutama pada aplikasi level rendah yang mengemas bit-bit bersama ke dalam ruang yang lebih besar.

Mari kita lihat kode untuk mendemonstrasikan satu kasus pemakaian:

\begin{tcolorbox}[enhanced jigsaw,left=3pt,breakable,colback=white!97.2549!black,boxrule=0.2mm]

\begin{Shaded}
\begin{Highlighting}[numbers=left,,]
\PreprocessorTok{\#include }\ImportTok{\textless{}stdio.h\textgreater{}}

\KeywordTok{struct}\NormalTok{ foo }\OperatorTok{\{}
    \DataTypeTok{unsigned} \DataTypeTok{int}\NormalTok{ a}\OperatorTok{;}
    \DataTypeTok{unsigned} \DataTypeTok{int}\NormalTok{ b}\OperatorTok{;}
    \DataTypeTok{unsigned} \DataTypeTok{int}\NormalTok{ c}\OperatorTok{;}
    \DataTypeTok{unsigned} \DataTypeTok{int}\NormalTok{ d}\OperatorTok{;}
\OperatorTok{\};}

\DataTypeTok{int}\NormalTok{ main}\OperatorTok{(}\DataTypeTok{void}\OperatorTok{)}
\OperatorTok{\{}
\NormalTok{    printf}\OperatorTok{(}\StringTok{"}\SpecialCharTok{\%zu\textbackslash{}n}\StringTok{"}\OperatorTok{,} \KeywordTok{sizeof}\OperatorTok{(}\KeywordTok{struct}\NormalTok{ foo}\OperatorTok{));}
\OperatorTok{\}}
\end{Highlighting}
\end{Shaded}

\end{tcolorbox}

Di sistemku, ini mencetak \hl{\mbox{\texttt{16}}}. Itu masuk akal, karena \hl{\mbox{\texttt{unsigned}}} berukuran 4 byte di sistemku.

Tapi bagaimana kalau kita tahu bahwa semua nilai yang akan disimpan di \hl{\mbox{\texttt{a}}} dan \hl{\mbox{\texttt{b}}} bisa disimpan dalam 5 bit, dan nilai di \hl{\mbox{\texttt{c}}} dan \hl{\mbox{\texttt{d}}} bisa disimpan dalam 3 bit? Totalnya hanya 16 bit. Kenapa mencadangkan 128 bit untuk mereka kalau kita hanya akan memakai 16?

Nah, kita bisa meminta C dengan sangat sopan untuk mencoba mengemas nilai-nilai ini. Kita bisa menentukan jumlah bit maksimum yang bisa dipakai nilai (dari 1 sampai ukuran tipe penampungnya).

Caranya dengan menaruh titik dua setelah nama field, diikuti lebar field dalam bit.

\begin{tcolorbox}[enhanced jigsaw,left=3pt,breakable,colback=white!97.2549!black,boxrule=0.2mm]

\begin{Shaded}
\begin{Highlighting}[numbers=left,,firstnumber=3,]
\KeywordTok{struct}\NormalTok{ foo }\OperatorTok{\{}
    \DataTypeTok{unsigned} \DataTypeTok{int}\NormalTok{ a}\OperatorTok{:}\DecValTok{5}\OperatorTok{;}
    \DataTypeTok{unsigned} \DataTypeTok{int}\NormalTok{ b}\OperatorTok{:}\DecValTok{5}\OperatorTok{;}
    \DataTypeTok{unsigned} \DataTypeTok{int}\NormalTok{ c}\OperatorTok{:}\DecValTok{3}\OperatorTok{;}
    \DataTypeTok{unsigned} \DataTypeTok{int}\NormalTok{ d}\OperatorTok{:}\DecValTok{3}\OperatorTok{;}
\OperatorTok{\};}
\end{Highlighting}
\end{Shaded}

\end{tcolorbox}

Sekarang kalau aku bertanya ke C seberapa besar \hl{\mbox{\texttt{struct\ foo}}}-ku, jawabannya 4! Tadinya 16 byte, sekarang hanya 4. Keempat nilai itu ``dikemas'' menjadi 4 byte, penghematan memori empat kali lipat.

Konsekuensinya, tentu saja, field 5-bit hanya bisa menampung nilai 0-31 dan field 3-bit hanya bisa menampung nilai 0-7. Tapi toh hidup memang soal kompromi.

\subsection{Bit-Field yang Tidak Berdekatan}\label{bit-field-yang-tidak-berdekatan}

Satu jebakan: C hanya menggabungkan bit-field yang \textbf{berdekatan}. Kalau mereka disela oleh non-bit-field, tidak ada penghematan:

\begin{tcolorbox}[enhanced jigsaw,left=3pt,breakable,colback=white!97.2549!black,boxrule=0.2mm]

\begin{Shaded}
\begin{Highlighting}[]
\KeywordTok{struct}\NormalTok{ foo }\OperatorTok{\{}            \CommentTok{// sizeof(struct foo) == 16 (di sistemku)}
    \DataTypeTok{unsigned} \DataTypeTok{int}\NormalTok{ a}\OperatorTok{:}\DecValTok{1}\OperatorTok{;}   \CommentTok{// karena a tidak berdekatan dengan c.}
    \DataTypeTok{unsigned} \DataTypeTok{int}\NormalTok{ b}\OperatorTok{;}
    \DataTypeTok{unsigned} \DataTypeTok{int}\NormalTok{ c}\OperatorTok{:}\DecValTok{1}\OperatorTok{;}
    \DataTypeTok{unsigned} \DataTypeTok{int}\NormalTok{ d}\OperatorTok{;}
\OperatorTok{\};}
\end{Highlighting}
\end{Shaded}

\end{tcolorbox}

Pada contoh itu, karena \hl{\mbox{\texttt{a}}} tidak berdekatan dengan \hl{\mbox{\texttt{c}}}, keduanya ``dikemas'' dalam \hl{\mbox{\texttt{int}}}-nya masing-masing.

Jadi kita punya satu \hl{\mbox{\texttt{int}}} masing-masing untuk \hl{\mbox{\texttt{a}}}, \hl{\mbox{\texttt{b}}}, \hl{\mbox{\texttt{c}}}, dan \hl{\mbox{\texttt{d}}}. Karena \hl{\mbox{\texttt{int}}} di sistemku 4 byte, totalnya 16 byte.

Penataan ulang singkat menghemat ruang dari 16 byte menjadi 12 byte (di sistemku):

\begin{tcolorbox}[enhanced jigsaw,left=3pt,breakable,colback=white!97.2549!black,boxrule=0.2mm]

\begin{Shaded}
\begin{Highlighting}[]
\KeywordTok{struct}\NormalTok{ foo }\OperatorTok{\{}            \CommentTok{// sizeof(struct foo) == 12 (di sistemku)}
    \DataTypeTok{unsigned} \DataTypeTok{int}\NormalTok{ a}\OperatorTok{:}\DecValTok{1}\OperatorTok{;}
    \DataTypeTok{unsigned} \DataTypeTok{int}\NormalTok{ c}\OperatorTok{:}\DecValTok{1}\OperatorTok{;}
    \DataTypeTok{unsigned} \DataTypeTok{int}\NormalTok{ b}\OperatorTok{;}
    \DataTypeTok{unsigned} \DataTypeTok{int}\NormalTok{ d}\OperatorTok{;}
\OperatorTok{\};}
\end{Highlighting}
\end{Shaded}

\end{tcolorbox}

Dan sekarang, karena \hl{\mbox{\texttt{a}}} bersebelahan dengan \hl{\mbox{\texttt{c}}}, compiler menaruh keduanya dalam satu \hl{\mbox{\texttt{int}}}.

Jadi kita punya satu \hl{\mbox{\texttt{int}}} untuk gabungan \hl{\mbox{\texttt{a}}} dan \hl{\mbox{\texttt{c}}}, dan satu \hl{\mbox{\texttt{int}}} masing-masing untuk \hl{\mbox{\texttt{b}}} dan \hl{\mbox{\texttt{d}}}. Total 3 \hl{\mbox{\texttt{int}}}, atau 12 byte.

Taruh semua bitfield-mu berdekatan supaya compiler menggabungkannya.

\subsection{\texorpdfstring{\hl{\mbox{\texttt{int}}} Signed atau Unsigned}{int Signed atau Unsigned}}\label{int-signed-atau-unsigned}

Kalau kamu hanya mendeklarasikan bit-field sebagai \hl{\mbox{\texttt{int}}}, compiler yang berbeda akan memperlakukannya sebagai \hl{\mbox{\texttt{signed}}} atau \hl{\mbox{\texttt{unsigned}}}. Sama seperti situasi pada \hl{\mbox{\texttt{char}}}.

Jelaskan tanda (signedness)-nya secara spesifik saat memakai bit-field.

\subsection{Bit-Field Tanpa Nama}\label{bit-field-tanpa-nama}

Pada keadaan tertentu, kamu mungkin perlu mencadangkan beberapa bit untuk alasan hardware, tapi tidak perlu memakainya di kode.

Misalnya, katakanlah ada satu byte yang 2 bit teratasnya punya makna, 1 bit terbawahnya punya makna, tapi 5 bit di tengah tidak kamu pakai\footnote{Dengan asumsi \hl{\mbox{\texttt{char}}} 8-bit, yaitu \hl{\mbox{\texttt{CHAR\_BIT\ ==\ 8}}}.}.

Kita \emph{bisa} melakukan sesuatu seperti ini:

\begin{tcolorbox}[enhanced jigsaw,left=3pt,breakable,colback=white!97.2549!black,boxrule=0.2mm]

\begin{Shaded}
\begin{Highlighting}[]
\KeywordTok{struct}\NormalTok{ foo }\OperatorTok{\{}
    \DataTypeTok{unsigned} \DataTypeTok{char}\NormalTok{ a}\OperatorTok{:}\DecValTok{2}\OperatorTok{;}
    \DataTypeTok{unsigned} \DataTypeTok{char}\NormalTok{ dummy}\OperatorTok{:}\DecValTok{5}\OperatorTok{;}
    \DataTypeTok{unsigned} \DataTypeTok{char}\NormalTok{ b}\OperatorTok{:}\DecValTok{1}\OperatorTok{;}
\OperatorTok{\};}
\end{Highlighting}
\end{Shaded}

\end{tcolorbox}

Dan itu berfungsi---di kode kita memakai \hl{\mbox{\texttt{a}}} dan \hl{\mbox{\texttt{b}}}, tapi tidak pernah \hl{\mbox{\texttt{dummy}}}. Ia hanya ada untuk menghabiskan 5 bit supaya \hl{\mbox{\texttt{a}}} dan \hl{\mbox{\texttt{b}}} berada di posisi ``yang dibutuhkan'' (oleh contoh yang dibuat-buat ini) di dalam byte.

C menyediakan cara untuk merapikannya: \emph{bit-field tanpa nama}. Kamu cukup menghilangkan nama (\hl{\mbox{\texttt{dummy}}}) dalam kasus ini, dan C tetap senang dengan efek yang sama:

\begin{tcolorbox}[enhanced jigsaw,left=3pt,breakable,colback=white!97.2549!black,boxrule=0.2mm]

\begin{Shaded}
\begin{Highlighting}[]
\KeywordTok{struct}\NormalTok{ foo }\OperatorTok{\{}
    \DataTypeTok{unsigned} \DataTypeTok{char}\NormalTok{ a}\OperatorTok{:}\DecValTok{2}\OperatorTok{;}
    \DataTypeTok{unsigned} \DataTypeTok{char} \OperatorTok{:}\DecValTok{5}\OperatorTok{;}   \CommentTok{// \textless{}{-}{-} bit{-}field tanpa nama!}
    \DataTypeTok{unsigned} \DataTypeTok{char}\NormalTok{ b}\OperatorTok{:}\DecValTok{1}\OperatorTok{;}
\OperatorTok{\};}
\end{Highlighting}
\end{Shaded}

\end{tcolorbox}

\subsection{Bit-Field Tanpa Nama dengan Lebar Nol}\label{bit-field-tanpa-nama-dengan-lebar-nol}

Esoterika lagi di sini\ldots{} Katakanlah kamu mengemas bit ke dalam \hl{\mbox{\texttt{unsigned\ int}}}, dan kamu butuh beberapa bit-field berdekatan untuk dikemas ke \hl{\mbox{\texttt{unsigned\ int}}} \emph{berikutnya}.

Artinya, kalau kamu menulis ini:

\begin{tcolorbox}[enhanced jigsaw,left=3pt,breakable,colback=white!97.2549!black,boxrule=0.2mm]

\begin{Shaded}
\begin{Highlighting}[]
\KeywordTok{struct}\NormalTok{ foo }\OperatorTok{\{}
    \DataTypeTok{unsigned} \DataTypeTok{int}\NormalTok{ a}\OperatorTok{:}\DecValTok{1}\OperatorTok{;}
    \DataTypeTok{unsigned} \DataTypeTok{int}\NormalTok{ b}\OperatorTok{:}\DecValTok{2}\OperatorTok{;}
    \DataTypeTok{unsigned} \DataTypeTok{int}\NormalTok{ c}\OperatorTok{:}\DecValTok{3}\OperatorTok{;}
    \DataTypeTok{unsigned} \DataTypeTok{int}\NormalTok{ d}\OperatorTok{:}\DecValTok{4}\OperatorTok{;}
\OperatorTok{\};}
\end{Highlighting}
\end{Shaded}

\end{tcolorbox}

compiler mengemas semuanya ke dalam satu \hl{\mbox{\texttt{unsigned\ int}}}. Tapi bagaimana kalau kamu butuh \hl{\mbox{\texttt{a}}} dan \hl{\mbox{\texttt{b}}} dalam satu \hl{\mbox{\texttt{int}}}, dan \hl{\mbox{\texttt{c}}} dan \hl{\mbox{\texttt{d}}} dalam \hl{\mbox{\texttt{int}}} yang lain?

Ada solusinya: taruh bit-field tanpa nama berlebar \hl{\mbox{\texttt{0}}} di tempat kamu ingin compiler mulai mengemas bit dari awal pada \hl{\mbox{\texttt{int}}} yang berbeda:

\begin{tcolorbox}[enhanced jigsaw,left=3pt,breakable,colback=white!97.2549!black,boxrule=0.2mm]

\begin{Shaded}
\begin{Highlighting}[]
\KeywordTok{struct}\NormalTok{ foo }\OperatorTok{\{}
    \DataTypeTok{unsigned} \DataTypeTok{int}\NormalTok{ a}\OperatorTok{:}\DecValTok{1}\OperatorTok{;}
    \DataTypeTok{unsigned} \DataTypeTok{int}\NormalTok{ b}\OperatorTok{:}\DecValTok{2}\OperatorTok{;}
    \DataTypeTok{unsigned} \DataTypeTok{int} \OperatorTok{:}\DecValTok{0}\OperatorTok{;}   \CommentTok{// \textless{}{-}{-}Bit{-}field tanpa nama berlebar nol}
    \DataTypeTok{unsigned} \DataTypeTok{int}\NormalTok{ c}\OperatorTok{:}\DecValTok{3}\OperatorTok{;}
    \DataTypeTok{unsigned} \DataTypeTok{int}\NormalTok{ d}\OperatorTok{:}\DecValTok{4}\OperatorTok{;}
\OperatorTok{\};}
\end{Highlighting}
\end{Shaded}

\end{tcolorbox}

Ini seperti page break eksplisit di word processor. Kamu memberi tahu compiler, ``Berhenti mengemas bit di \hl{\mbox{\texttt{unsigned}}} ini, dan mulai mengemas di yang berikutnya.''

Dengan menambahkan bit-field tanpa nama berlebar nol di tempat itu, compiler menaruh \hl{\mbox{\texttt{a}}} dan \hl{\mbox{\texttt{b}}} dalam satu \hl{\mbox{\texttt{unsigned\ int}}}, dan \hl{\mbox{\texttt{c}}} dan \hl{\mbox{\texttt{d}}} dalam \hl{\mbox{\texttt{unsigned\ int}}} lain. Total dua, dengan ukuran 8 byte di sistemku (\hl{\mbox{\texttt{unsigned\ int}}} masing-masing 4 byte).

\index{struct keyword@\texttt{struct} keyword!bit fields|)}

\section{Union}\label{union}

\index{union keyword@\texttt{union} keyword|(}

Pada dasarnya ini sama seperti \hl{\mbox{\texttt{struct}}}, hanya saja field-nya saling tumpang tindih di memori. \hl{\mbox{\texttt{union}}} hanya akan sebesar field terbesarnya, dan kamu hanya bisa memakai satu field dalam satu waktu.

Ini cara untuk memakai ulang ruang memori yang sama untuk berbagai jenis data.

Deklarasinya sama seperti \hl{\mbox{\texttt{struct}}}, hanya saja memakai \hl{\mbox{\texttt{union}}}. Lihat ini:

\begin{tcolorbox}[enhanced jigsaw,left=3pt,breakable,colback=white!97.2549!black,boxrule=0.2mm]

\begin{Shaded}
\begin{Highlighting}[]
\KeywordTok{union}\NormalTok{ foo }\OperatorTok{\{}
    \DataTypeTok{int}\NormalTok{ a}\OperatorTok{,}\NormalTok{ b}\OperatorTok{,}\NormalTok{ c}\OperatorTok{,}\NormalTok{ d}\OperatorTok{,}\NormalTok{ e}\OperatorTok{,}\NormalTok{ f}\OperatorTok{;}
    \DataTypeTok{float}\NormalTok{ g}\OperatorTok{,}\NormalTok{ h}\OperatorTok{;}
    \DataTypeTok{char}\NormalTok{ i}\OperatorTok{,}\NormalTok{ j}\OperatorTok{,}\NormalTok{ k}\OperatorTok{,}\NormalTok{ l}\OperatorTok{;}
\OperatorTok{\};}
\end{Highlighting}
\end{Shaded}

\end{tcolorbox}

Nah, itu banyak sekali field. Kalau ini \hl{\mbox{\texttt{struct}}}, sistemku akan bilang dibutuhkan 36 byte untuk menampung semuanya.

Tapi ini \hl{\mbox{\texttt{union}}}, jadi semua field itu saling tumpang tindih dalam bentangan memori yang sama. Field terbesarnya \hl{\mbox{\texttt{int}}} (atau \hl{\mbox{\texttt{float}}}), yang memakai 4 byte di sistemku. Dan memang, kalau kutanya \hl{\mbox{\texttt{sizeof}}} dari \hl{\mbox{\texttt{union\ foo}}}, jawabannya 4!

Konsekuensinya, hanya satu field yang bisa dipakai secara portabel dalam satu waktu. Namun\ldots{}

\subsection{Union dan Type Punning}\label{union-type-punning}

\index{union keyword@\texttt{union} keyword!type punning|(}

Kamu bisa menulis ke satu field \hl{\mbox{\texttt{union}}} dan membaca dari field lain secara non-portabel!

Melakukannya disebut \href{https://en.wikipedia.org/wiki/Type_punning}{type punning}\footnote{https://en.wikipedia.org/wiki/Type\_punning}, dan kamu memakainya kalau benar-benar tahu apa yang kamu lakukan, biasanya pada pemrograman level rendah.

Karena anggota union berbagi memori yang sama, menulis ke satu anggota pasti memengaruhi yang lain. Dan kalau kamu membaca dari satu anggota sesuatu yang ditulis ke anggota lain, kamu mendapat efek yang aneh.

\begin{tcolorbox}[enhanced jigsaw,left=3pt,breakable,colback=white!97.2549!black,boxrule=0.2mm]

\begin{Shaded}
\begin{Highlighting}[numbers=left,,]
\PreprocessorTok{\#include }\ImportTok{\textless{}stdio.h\textgreater{}}

\KeywordTok{union}\NormalTok{ foo }\OperatorTok{\{}
    \DataTypeTok{float}\NormalTok{ b}\OperatorTok{;}
    \DataTypeTok{short}\NormalTok{ a}\OperatorTok{;}
\OperatorTok{\};}

\DataTypeTok{int}\NormalTok{ main}\OperatorTok{(}\DataTypeTok{void}\OperatorTok{)}
\OperatorTok{\{}
    \KeywordTok{union}\NormalTok{ foo x}\OperatorTok{;}

\NormalTok{    x}\OperatorTok{.}\NormalTok{b }\OperatorTok{=} \FloatTok{3.14159}\OperatorTok{;}

\NormalTok{    printf}\OperatorTok{(}\StringTok{"}\SpecialCharTok{\%f\textbackslash{}n}\StringTok{"}\OperatorTok{,}\NormalTok{ x}\OperatorTok{.}\NormalTok{b}\OperatorTok{);}  \CommentTok{// 3.14159, wajar}

\NormalTok{    printf}\OperatorTok{(}\StringTok{"}\SpecialCharTok{\%d\textbackslash{}n}\StringTok{"}\OperatorTok{,}\NormalTok{ x}\OperatorTok{.}\NormalTok{a}\OperatorTok{);}  \CommentTok{// Tapi bagaimana dengan ini?}
\OperatorTok{\}}
\end{Highlighting}
\end{Shaded}

\end{tcolorbox}

Di sistemku, ini mencetak:

\begin{tcolorbox}[enhanced jigsaw,left=3pt,breakable,colback=white!97.2549!black,boxrule=0.2mm]

\begin{verbatim}
3.141590
4048
\end{verbatim}

\end{tcolorbox}

karena di balik layar, representasi object dari float \hl{\mbox{\texttt{3.14159}}} sama dengan representasi object dari short \hl{\mbox{\texttt{4048}}}. Di sistemku. Hasilmu mungkin berbeda.

\index{union keyword@\texttt{union} keyword!type punning|)}

\subsection{\texorpdfstring{Pointer ke \hl{\mbox{\texttt{union}}}}{Pointer ke union}}\label{pointer-ke-union}

\index{union keyword@\texttt{union} keyword!pointers to|(}

Kalau kamu punya pointer ke \hl{\mbox{\texttt{union}}}, kamu bisa meng-cast pointer itu ke tipe field mana pun dalam \hl{\mbox{\texttt{union}}} itu dan mengambil nilainya dengan cara itu.

Pada contoh ini, kita lihat \hl{\mbox{\texttt{union}}} berisi \hl{\mbox{\texttt{int}}} dan \hl{\mbox{\texttt{float}}}. Kita ambil pointer ke \hl{\mbox{\texttt{union}}}, tapi kita cast ke tipe \hl{\mbox{\texttt{int*}}} dan \hl{\mbox{\texttt{float*}}} (cast membungkam peringatan compiler). Lalu kalau kita men-dereference pointer-pointer itu, kita lihat bahwa nilainya sama dengan yang kita simpan langsung di \hl{\mbox{\texttt{union}}}.

\begin{tcolorbox}[enhanced jigsaw,left=3pt,breakable,colback=white!97.2549!black,boxrule=0.2mm]

\begin{Shaded}
\begin{Highlighting}[numbers=left,,]
\PreprocessorTok{\#include }\ImportTok{\textless{}stdio.h\textgreater{}}

\KeywordTok{union}\NormalTok{ foo }\OperatorTok{\{}
    \DataTypeTok{int}\NormalTok{ a}\OperatorTok{,}\NormalTok{ b}\OperatorTok{,}\NormalTok{ c}\OperatorTok{,}\NormalTok{ d}\OperatorTok{,}\NormalTok{ e}\OperatorTok{,}\NormalTok{ f}\OperatorTok{;}
    \DataTypeTok{float}\NormalTok{ g}\OperatorTok{,}\NormalTok{ h}\OperatorTok{;}
    \DataTypeTok{char}\NormalTok{ i}\OperatorTok{,}\NormalTok{ j}\OperatorTok{,}\NormalTok{ k}\OperatorTok{,}\NormalTok{ l}\OperatorTok{;}
\OperatorTok{\};}

\DataTypeTok{int}\NormalTok{ main}\OperatorTok{(}\DataTypeTok{void}\OperatorTok{)}
\OperatorTok{\{}
    \KeywordTok{union}\NormalTok{ foo x}\OperatorTok{;}

    \DataTypeTok{int} \OperatorTok{*}\NormalTok{foo\_int\_p }\OperatorTok{=} \OperatorTok{(}\DataTypeTok{int} \OperatorTok{*)\&}\NormalTok{x}\OperatorTok{;}
    \DataTypeTok{float} \OperatorTok{*}\NormalTok{foo\_float\_p }\OperatorTok{=} \OperatorTok{(}\DataTypeTok{float} \OperatorTok{*)\&}\NormalTok{x}\OperatorTok{;}

\NormalTok{    x}\OperatorTok{.}\NormalTok{a }\OperatorTok{=} \DecValTok{12}\OperatorTok{;}
\NormalTok{    printf}\OperatorTok{(}\StringTok{"}\SpecialCharTok{\%d\textbackslash{}n}\StringTok{"}\OperatorTok{,}\NormalTok{ x}\OperatorTok{.}\NormalTok{a}\OperatorTok{);}           \CommentTok{// 12}
\NormalTok{    printf}\OperatorTok{(}\StringTok{"}\SpecialCharTok{\%d\textbackslash{}n}\StringTok{"}\OperatorTok{,} \OperatorTok{*}\NormalTok{foo\_int\_p}\OperatorTok{);}    \CommentTok{// 12, lagi}

\NormalTok{    x}\OperatorTok{.}\NormalTok{g }\OperatorTok{=} \FloatTok{3.141592}\OperatorTok{;}
\NormalTok{    printf}\OperatorTok{(}\StringTok{"}\SpecialCharTok{\%f\textbackslash{}n}\StringTok{"}\OperatorTok{,}\NormalTok{ x}\OperatorTok{.}\NormalTok{g}\OperatorTok{);}           \CommentTok{// 3.141592}
\NormalTok{    printf}\OperatorTok{(}\StringTok{"}\SpecialCharTok{\%f\textbackslash{}n}\StringTok{"}\OperatorTok{,} \OperatorTok{*}\NormalTok{foo\_float\_p}\OperatorTok{);}  \CommentTok{// 3.141592, lagi}
\OperatorTok{\}}
\end{Highlighting}
\end{Shaded}

\end{tcolorbox}

Kebalikannya juga benar. Kalau kita punya pointer ke tipe di dalam \hl{\mbox{\texttt{union}}}, kita bisa meng-cast-nya ke pointer ke \hl{\mbox{\texttt{union}}} dan mengakses anggotanya.

\begin{tcolorbox}[enhanced jigsaw,left=3pt,breakable,colback=white!97.2549!black,boxrule=0.2mm]

\begin{Shaded}
\begin{Highlighting}[]
\KeywordTok{union}\NormalTok{ foo x}\OperatorTok{;}
\DataTypeTok{int} \OperatorTok{*}\NormalTok{foo\_int\_p }\OperatorTok{=} \OperatorTok{(}\DataTypeTok{int} \OperatorTok{*)\&}\NormalTok{x}\OperatorTok{;}             \CommentTok{// Pointer ke field int}
\KeywordTok{union}\NormalTok{ foo }\OperatorTok{*}\NormalTok{p }\OperatorTok{=} \OperatorTok{(}\KeywordTok{union}\NormalTok{ foo }\OperatorTok{*)}\NormalTok{foo\_int\_p}\OperatorTok{;}  \CommentTok{// Kembali ke pointer ke union}

\NormalTok{p}\OperatorTok{{-}\textgreater{}}\NormalTok{a }\OperatorTok{=} \DecValTok{12}\OperatorTok{;}  \CommentTok{// Baris ini sama dengan...}
\NormalTok{x}\OperatorTok{.}\NormalTok{a }\OperatorTok{=} \DecValTok{12}\OperatorTok{;}   \CommentTok{// yang ini.}
\end{Highlighting}
\end{Shaded}

\end{tcolorbox}

Semua ini hanya memberi tahu bahwa, di balik layar, semua nilai dalam \hl{\mbox{\texttt{union}}} dimulai di tempat yang sama di memori, dan itu sama dengan lokasi seluruh \hl{\mbox{\texttt{union}}}.

\index{union keyword@\texttt{union} keyword!pointers to|)}

\subsection{Common Initial Sequence pada Union}\label{common-initial-sequence-pada-union}

\index{union keyword@\texttt{union} keyword!common initial sequences|(}

Kalau kamu punya \hl{\mbox{\texttt{union}}} berisi \hl{\mbox{\texttt{struct}}}, dan semua \hl{\mbox{\texttt{struct}}} itu diawali dengan \emph{common initial sequence}, boleh mengakses anggota urutan itu dari anggota \hl{\mbox{\texttt{union}}} mana pun.

Apa?

Berikut dua \hl{\mbox{\texttt{struct}}} dengan common initial sequence:

\begin{tcolorbox}[enhanced jigsaw,left=3pt,breakable,colback=white!97.2549!black,boxrule=0.2mm]

\begin{Shaded}
\begin{Highlighting}[]
\KeywordTok{struct}\NormalTok{ a }\OperatorTok{\{}
    \DataTypeTok{int}\NormalTok{ x}\OperatorTok{;}     \CommentTok{//}
    \DataTypeTok{float}\NormalTok{ y}\OperatorTok{;}   \CommentTok{// Common initial sequence}

    \DataTypeTok{char} \OperatorTok{*}\NormalTok{p}\OperatorTok{;}
\OperatorTok{\};}

\KeywordTok{struct}\NormalTok{ b }\OperatorTok{\{}
    \DataTypeTok{int}\NormalTok{ x}\OperatorTok{;}     \CommentTok{//}
    \DataTypeTok{float}\NormalTok{ y}\OperatorTok{;}   \CommentTok{// Common initial sequence}

    \DataTypeTok{double} \OperatorTok{*}\NormalTok{p}\OperatorTok{;}
    \DataTypeTok{short}\NormalTok{ z}\OperatorTok{;}
\OperatorTok{\};}
\end{Highlighting}
\end{Shaded}

\end{tcolorbox}

Kamu lihat? Keduanya diawali \hl{\mbox{\texttt{int}}} diikuti \hl{\mbox{\texttt{float}}}---itulah common initial sequence. Anggota-anggota dalam urutan pada \hl{\mbox{\texttt{struct}}} harus bertipe kompatibel. Dan kita lihat itu pada \hl{\mbox{\texttt{x}}} dan \hl{\mbox{\texttt{y}}}, yaitu \hl{\mbox{\texttt{int}}} dan \hl{\mbox{\texttt{float}}} masing-masing.

Sekarang kita buat union dari keduanya:

\begin{tcolorbox}[enhanced jigsaw,left=3pt,breakable,colback=white!97.2549!black,boxrule=0.2mm]

\begin{Shaded}
\begin{Highlighting}[]
\KeywordTok{union}\NormalTok{ foo }\OperatorTok{\{}
    \KeywordTok{struct}\NormalTok{ a sa}\OperatorTok{;}
    \KeywordTok{struct}\NormalTok{ b sb}\OperatorTok{;}
\OperatorTok{\};}
\end{Highlighting}
\end{Shaded}

\end{tcolorbox}

Aturan ini menjamin bahwa anggota common initial sequence bisa saling menggantikan di kode. Artinya:

\begin{itemize}
\tightlist
\item
  \hl{\mbox{\texttt{f.sa.x}}} sama dengan \hl{\mbox{\texttt{f.sb.x}}}.
\end{itemize}

dan

\begin{itemize}
\tightlist
\item
  \hl{\mbox{\texttt{f.sa.y}}} sama dengan \hl{\mbox{\texttt{f.sb.y}}}.
\end{itemize}

Karena field \hl{\mbox{\texttt{x}}} dan \hl{\mbox{\texttt{y}}} sama-sama ada di common initial sequence.

Selain itu, nama anggota dalam common initial sequence tidak penting---yang penting hanyalah tipenya sama.

Secara keseluruhan, ini memberi kita cara untuk menambahkan informasi bersama di antara \hl{\mbox{\texttt{struct}}} dalam \hl{\mbox{\texttt{union}}} secara aman. Contoh terbaiknya mungkin memakai sebuah field untuk menentukan \hl{\mbox{\texttt{struct}}} mana, dari semua \hl{\mbox{\texttt{struct}}} dalam \hl{\mbox{\texttt{union}}}, yang sedang ``dipakai''.

Artinya, kalau ini tidak diizinkan dan kita mengoper \hl{\mbox{\texttt{union}}} ke suatu fungsi, bagaimana fungsi itu tahu anggota \hl{\mbox{\texttt{union}}} mana yang harus dilihatnya?

Lihat \hl{\mbox{\texttt{struct}}}-\hl{\mbox{\texttt{struct}}} ini. Perhatikan common initial sequence-nya:

\begin{tcolorbox}[enhanced jigsaw,left=3pt,breakable,colback=white!97.2549!black,boxrule=0.2mm]

\begin{Shaded}
\begin{Highlighting}[numbers=left,,]
\PreprocessorTok{\#include }\ImportTok{\textless{}stdio.h\textgreater{}}

\KeywordTok{struct}\NormalTok{ common }\OperatorTok{\{}
    \DataTypeTok{int}\NormalTok{ type}\OperatorTok{;}   \CommentTok{// common initial sequence}
\OperatorTok{\};}

\KeywordTok{struct}\NormalTok{ antelope }\OperatorTok{\{}
    \DataTypeTok{int}\NormalTok{ type}\OperatorTok{;}   \CommentTok{// common initial sequence}

    \DataTypeTok{int}\NormalTok{ loudness}\OperatorTok{;}
\OperatorTok{\};}

\KeywordTok{struct}\NormalTok{ octopus }\OperatorTok{\{}
    \DataTypeTok{int}\NormalTok{ type}\OperatorTok{;}   \CommentTok{// common initial sequence}

    \DataTypeTok{int}\NormalTok{ sea\_creature}\OperatorTok{;}
    \DataTypeTok{float}\NormalTok{ intelligence}\OperatorTok{;}
\OperatorTok{\};}
\end{Highlighting}
\end{Shaded}

\end{tcolorbox}

Sekarang kita masukkan ke dalam \hl{\mbox{\texttt{union}}}:

\begin{tcolorbox}[enhanced jigsaw,left=3pt,breakable,colback=white!97.2549!black,boxrule=0.2mm]

\begin{Shaded}
\begin{Highlighting}[numbers=left,,firstnumber=20,]
\KeywordTok{union}\NormalTok{ animal }\OperatorTok{\{}
    \KeywordTok{struct}\NormalTok{ common common}\OperatorTok{;}
    \KeywordTok{struct}\NormalTok{ antelope antelope}\OperatorTok{;}
    \KeywordTok{struct}\NormalTok{ octopus octopus}\OperatorTok{;}
\OperatorTok{\};}
\end{Highlighting}
\end{Shaded}

\end{tcolorbox}

Mohon dimaklumi dua \hl{\mbox{\texttt{\#define}}} untuk demo ini:

\begin{tcolorbox}[enhanced jigsaw,left=3pt,breakable,colback=white!97.2549!black,boxrule=0.2mm]

\begin{Shaded}
\begin{Highlighting}[numbers=left,,firstnumber=26,]
\PreprocessorTok{\#define ANTELOPE }\DecValTok{1}
\PreprocessorTok{\#define OCTOPUS  }\DecValTok{2}
\end{Highlighting}
\end{Shaded}

\end{tcolorbox}

Sejauh ini, belum ada yang istimewa. Tampaknya field \hl{\mbox{\texttt{type}}} sama sekali tidak berguna.

Tapi sekarang mari kita buat fungsi generik yang mencetak \hl{\mbox{\texttt{union\ animal}}}. Ia harus bisa membedakan apakah yang dilihatnya \hl{\mbox{\texttt{struct\ antelope}}} atau \hl{\mbox{\texttt{struct\ octopus}}}.

Berkat keajaiban common initial sequence, ia bisa mencari tipe hewan di salah satu dari tempat-tempat ini untuk \hl{\mbox{\texttt{union\ animal\ x}}} tertentu:

\begin{tcolorbox}[enhanced jigsaw,left=3pt,breakable,colback=white!97.2549!black,boxrule=0.2mm]

\begin{Shaded}
\begin{Highlighting}[]
\DataTypeTok{int}\NormalTok{ type }\OperatorTok{=}\NormalTok{ x}\OperatorTok{.}\NormalTok{common}\OperatorTok{.}\NormalTok{type}\OperatorTok{;}    \CommentTok{// atau...}
\DataTypeTok{int}\NormalTok{ type }\OperatorTok{=}\NormalTok{ x}\OperatorTok{.}\NormalTok{antelope}\OperatorTok{.}\NormalTok{type}\OperatorTok{;}  \CommentTok{// atau...}
\DataTypeTok{int}\NormalTok{ type }\OperatorTok{=}\NormalTok{ x}\OperatorTok{.}\NormalTok{octopus}\OperatorTok{.}\NormalTok{type}\OperatorTok{;}
\end{Highlighting}
\end{Shaded}

\end{tcolorbox}

Semuanya merujuk ke nilai yang sama di memori.

Dan, seperti yang mungkin sudah kamu duga, \hl{\mbox{\texttt{struct\ common}}} ada di sana supaya kode bisa melihat tipe secara agnostik tanpa menyebut hewan tertentu.

Mari kita lihat kode untuk mencetak \hl{\mbox{\texttt{union\ animal}}}:

\begin{tcolorbox}[enhanced jigsaw,left=3pt,breakable,colback=white!97.2549!black,boxrule=0.2mm]

\begin{Shaded}
\begin{Highlighting}[numbers=left,,firstnumber=29,]
\DataTypeTok{void}\NormalTok{ print\_animal}\OperatorTok{(}\KeywordTok{union}\NormalTok{ animal }\OperatorTok{*}\NormalTok{x}\OperatorTok{)}
\OperatorTok{\{}
    \ControlFlowTok{switch} \OperatorTok{(}\NormalTok{x}\OperatorTok{{-}\textgreater{}}\NormalTok{common}\OperatorTok{.}\NormalTok{type}\OperatorTok{)} \OperatorTok{\{}
        \ControlFlowTok{case}\NormalTok{ ANTELOPE}\OperatorTok{:}
\NormalTok{            printf}\OperatorTok{(}\StringTok{"Antelope: loudness=}\SpecialCharTok{\%d\textbackslash{}n}\StringTok{"}\OperatorTok{,}\NormalTok{ x}\OperatorTok{{-}\textgreater{}}\NormalTok{antelope}\OperatorTok{.}\NormalTok{loudness}\OperatorTok{);}
            \ControlFlowTok{break}\OperatorTok{;}

        \ControlFlowTok{case}\NormalTok{ OCTOPUS}\OperatorTok{:}
\NormalTok{            printf}\OperatorTok{(}\StringTok{"Octopus : sea\_creature=}\SpecialCharTok{\%d\textbackslash{}n}\StringTok{"}\OperatorTok{,}\NormalTok{ x}\OperatorTok{{-}\textgreater{}}\NormalTok{octopus}\OperatorTok{.}\NormalTok{sea\_creature}\OperatorTok{);}
\NormalTok{            printf}\OperatorTok{(}\StringTok{"          intelligence=}\SpecialCharTok{\%f\textbackslash{}n}\StringTok{"}\OperatorTok{,}\NormalTok{ x}\OperatorTok{{-}\textgreater{}}\NormalTok{octopus}\OperatorTok{.}\NormalTok{intelligence}\OperatorTok{);}
            \ControlFlowTok{break}\OperatorTok{;}
        
        \ControlFlowTok{default}\OperatorTok{:}
\NormalTok{            printf}\OperatorTok{(}\StringTok{"Unknown animal type}\SpecialCharTok{\textbackslash{}n}\StringTok{"}\OperatorTok{);}
    \OperatorTok{\}}

\OperatorTok{\}}

\DataTypeTok{int}\NormalTok{ main}\OperatorTok{(}\DataTypeTok{void}\OperatorTok{)}
\OperatorTok{\{}
    \KeywordTok{union}\NormalTok{ animal a }\OperatorTok{=} \OperatorTok{\{.}\NormalTok{antelope}\OperatorTok{.}\NormalTok{type}\OperatorTok{=}\NormalTok{ANTELOPE}\OperatorTok{,} \OperatorTok{.}\NormalTok{antelope}\OperatorTok{.}\NormalTok{loudness}\OperatorTok{=}\DecValTok{12}\OperatorTok{\};}
    \KeywordTok{union}\NormalTok{ animal b }\OperatorTok{=} \OperatorTok{\{.}\NormalTok{octopus}\OperatorTok{.}\NormalTok{type}\OperatorTok{=}\NormalTok{OCTOPUS}\OperatorTok{,} \OperatorTok{.}\NormalTok{octopus}\OperatorTok{.}\NormalTok{sea\_creature}\OperatorTok{=}\DecValTok{1}\OperatorTok{,}
                                       \OperatorTok{.}\NormalTok{octopus}\OperatorTok{.}\NormalTok{intelligence}\OperatorTok{=}\FloatTok{12.8}\OperatorTok{\};}

\NormalTok{    print\_animal}\OperatorTok{(\&}\NormalTok{a}\OperatorTok{);}
\NormalTok{    print\_animal}\OperatorTok{(\&}\NormalTok{b}\OperatorTok{);}
\OperatorTok{\}}
\end{Highlighting}
\end{Shaded}

\end{tcolorbox}

Lihat bagaimana di baris 29 kita hanya mengoper \hl{\mbox{\texttt{union}}}-nya---kita tidak tahu \hl{\mbox{\texttt{struct}}} hewan apa yang sedang dipakai di dalamnya.

Tapi tidak apa-apa! Karena di baris 31 kita memeriksa tipenya untuk melihat apakah itu antelope atau octopus. Lalu kita bisa melihat \hl{\mbox{\texttt{struct}}} yang tepat untuk mengambil anggotanya.

Efek yang sama pasti bisa didapat hanya dengan \hl{\mbox{\texttt{struct}}}, tapi kamu bisa memakai cara ini kalau menginginkan efek hemat memori dari \hl{\mbox{\texttt{union}}}.

\index{union keyword@\texttt{union} keyword!common initial sequences|)}

\section{Union dan Struct Tanpa Nama}\label{union-dan-struct-tanpa-nama}

\index{union keyword@\texttt{union} keyword!and unnamed structs@and unnamed \texttt{struct}s|(}

Kamu tahu bahwa kamu bisa punya struct tanpa nama, seperti ini:

\begin{tcolorbox}[enhanced jigsaw,left=3pt,breakable,colback=white!97.2549!black,boxrule=0.2mm]

\begin{Shaded}
\begin{Highlighting}[]
\KeywordTok{struct} \OperatorTok{\{}
    \DataTypeTok{int}\NormalTok{ x}\OperatorTok{,}\NormalTok{ y}\OperatorTok{;}
\OperatorTok{\}}\NormalTok{ s}\OperatorTok{;}
\end{Highlighting}
\end{Shaded}

\end{tcolorbox}

Itu mendefinisikan variabel \hl{\mbox{\texttt{s}}} bertipe \hl{\mbox{\texttt{struct}}} anonim (karena \hl{\mbox{\texttt{struct}}} itu tidak punya name tag), dengan anggota \hl{\mbox{\texttt{x}}} dan \hl{\mbox{\texttt{y}}}.

Jadi hal seperti ini valid:

\begin{tcolorbox}[enhanced jigsaw,left=3pt,breakable,colback=white!97.2549!black,boxrule=0.2mm]

\begin{Shaded}
\begin{Highlighting}[]
\NormalTok{s}\OperatorTok{.}\NormalTok{x }\OperatorTok{=} \DecValTok{34}\OperatorTok{;}
\NormalTok{s}\OperatorTok{.}\NormalTok{y }\OperatorTok{=} \DecValTok{90}\OperatorTok{;}

\NormalTok{printf}\OperatorTok{(}\StringTok{"}\SpecialCharTok{\%d}\StringTok{ }\SpecialCharTok{\%d\textbackslash{}n}\StringTok{"}\OperatorTok{,}\NormalTok{ s}\OperatorTok{.}\NormalTok{x}\OperatorTok{,}\NormalTok{ s}\OperatorTok{.}\NormalTok{y}\OperatorTok{);}
\end{Highlighting}
\end{Shaded}

\end{tcolorbox}

Ternyata \hl{\mbox{\texttt{struct}}} tanpa nama itu bisa ditaruh dalam \hl{\mbox{\texttt{union}}} seperti yang mungkin kamu duga:

\begin{tcolorbox}[enhanced jigsaw,left=3pt,breakable,colback=white!97.2549!black,boxrule=0.2mm]

\begin{Shaded}
\begin{Highlighting}[]
\KeywordTok{union}\NormalTok{ foo }\OperatorTok{\{}
    \KeywordTok{struct} \OperatorTok{\{}       \CommentTok{// tanpa nama!}
        \DataTypeTok{int}\NormalTok{ x}\OperatorTok{,}\NormalTok{ y}\OperatorTok{;}
    \OperatorTok{\}}\NormalTok{ a}\OperatorTok{;}

    \KeywordTok{struct} \OperatorTok{\{}       \CommentTok{// tanpa nama!}
        \DataTypeTok{int}\NormalTok{ z}\OperatorTok{,}\NormalTok{ w}\OperatorTok{;}
    \OperatorTok{\}}\NormalTok{ b}\OperatorTok{;}
\OperatorTok{\};}
\end{Highlighting}
\end{Shaded}

\end{tcolorbox}

Lalu diakses seperti biasa:

\begin{tcolorbox}[enhanced jigsaw,left=3pt,breakable,colback=white!97.2549!black,boxrule=0.2mm]

\begin{Shaded}
\begin{Highlighting}[]
\KeywordTok{union}\NormalTok{ foo f}\OperatorTok{;}

\NormalTok{f}\OperatorTok{.}\NormalTok{a}\OperatorTok{.}\NormalTok{x }\OperatorTok{=} \DecValTok{1}\OperatorTok{;}
\NormalTok{f}\OperatorTok{.}\NormalTok{a}\OperatorTok{.}\NormalTok{y }\OperatorTok{=} \DecValTok{2}\OperatorTok{;}
\NormalTok{f}\OperatorTok{.}\NormalTok{b}\OperatorTok{.}\NormalTok{z }\OperatorTok{=} \DecValTok{3}\OperatorTok{;}
\NormalTok{f}\OperatorTok{.}\NormalTok{b}\OperatorTok{.}\NormalTok{w }\OperatorTok{=} \DecValTok{4}\OperatorTok{;}
\end{Highlighting}
\end{Shaded}

\end{tcolorbox}

Tidak ada masalah!

\index{union keyword@\texttt{union} keyword!and unnamed structs@and unnamed \texttt{struct}s|)}

\section{\texorpdfstring{Mengoper dan Mengembalikan \hl{\mbox{\texttt{struct}}} dan \hl{\mbox{\texttt{union}}}}{Mengoper dan Mengembalikan struct dan union}}\label{mengoper-dan-mengembalikan-struct-dan-union}

\index{union keyword@\texttt{union} keyword!passing and returning|(} \index{struct keyword@\texttt{struct} keyword!passing and returning|(}

Kamu bisa mengoper \hl{\mbox{\texttt{struct}}} atau \hl{\mbox{\texttt{union}}} ke fungsi secara by value (bukan pointer ke objeknya)---salinan objek itu ke parameter dibuat seperti pada penugasan biasa.

Kamu juga bisa mengembalikan \hl{\mbox{\texttt{struct}}} atau \hl{\mbox{\texttt{union}}} dari fungsi dan ia juga dioper balik secara by value.

\begin{tcolorbox}[enhanced jigsaw,left=3pt,breakable,colback=white!97.2549!black,boxrule=0.2mm]

\begin{Shaded}
\begin{Highlighting}[numbers=left,,]
\PreprocessorTok{\#include }\ImportTok{\textless{}stdio.h\textgreater{}}

\KeywordTok{struct}\NormalTok{ foo }\OperatorTok{\{}
    \DataTypeTok{int}\NormalTok{ x}\OperatorTok{,}\NormalTok{ y}\OperatorTok{;}
\OperatorTok{\};}

\KeywordTok{struct}\NormalTok{ foo f}\OperatorTok{(}\DataTypeTok{void}\OperatorTok{)}
\OperatorTok{\{}
    \ControlFlowTok{return} \OperatorTok{(}\KeywordTok{struct}\NormalTok{ foo}\OperatorTok{)\{.}\NormalTok{x}\OperatorTok{=}\DecValTok{34}\OperatorTok{,} \OperatorTok{.}\NormalTok{y}\OperatorTok{=}\DecValTok{90}\OperatorTok{\};}
\OperatorTok{\}}

\DataTypeTok{int}\NormalTok{ main}\OperatorTok{(}\DataTypeTok{void}\OperatorTok{)}
\OperatorTok{\{}
    \KeywordTok{struct}\NormalTok{ foo a }\OperatorTok{=}\NormalTok{ f}\OperatorTok{();}  \CommentTok{// Salinan dibuat}

\NormalTok{    printf}\OperatorTok{(}\StringTok{"}\SpecialCharTok{\%d}\StringTok{ }\SpecialCharTok{\%d\textbackslash{}n}\StringTok{"}\OperatorTok{,}\NormalTok{ a}\OperatorTok{.}\NormalTok{x}\OperatorTok{,}\NormalTok{ a}\OperatorTok{.}\NormalTok{y}\OperatorTok{);}
\OperatorTok{\}}
\end{Highlighting}
\end{Shaded}

\end{tcolorbox}

Fakta menarik: kalau kamu melakukan ini, kamu bisa memakai operator \hl{\mbox{\texttt{.}}} langsung pada pemanggilan fungsi:

\begin{tcolorbox}[enhanced jigsaw,left=3pt,breakable,colback=white!97.2549!black,boxrule=0.2mm]

\begin{Shaded}
\begin{Highlighting}[numbers=left,,firstnumber=16,]
\NormalTok{    printf}\OperatorTok{(}\StringTok{"}\SpecialCharTok{\%d}\StringTok{ }\SpecialCharTok{\%d\textbackslash{}n}\StringTok{"}\OperatorTok{,}\NormalTok{ f}\OperatorTok{().}\NormalTok{x}\OperatorTok{,}\NormalTok{ f}\OperatorTok{().}\NormalTok{y}\OperatorTok{);}
\end{Highlighting}
\end{Shaded}

\end{tcolorbox}

(Tentu saja contoh itu memanggil fungsi dua kali, secara tidak efisien.)

Dan hal yang sama berlaku untuk mengembalikan pointer ke \hl{\mbox{\texttt{struct}}} dan \hl{\mbox{\texttt{union}}}---pastikan saja memakai arrow operator \hl{\mbox{\texttt{-\textgreater{}}}} dalam kasus itu.

\index{union keyword@\texttt{union} keyword!passing and returning|)} \index{struct keyword@\texttt{struct} keyword!passing and returning|)} \index{union keyword@\texttt{union} keyword|)} \index{struct keyword@\texttt{struct} keyword|)}

\chapter{Karakter dan String II}\label{karakter-dan-string-ii}

Kita sudah membahas bahwa tipe \hl{\mbox{\texttt{char}}} sebenarnya hanyalah tipe integer kecil\ldots{} dan hal yang sama berlaku untuk karakter dalam tanda kutip tunggal.

Tapi string dalam tanda kutip ganda bertipe \hl{\mbox{\texttt{const\ char\ *}}}.

Ternyata ada beberapa jenis string dan karakter lagi, dan ini membawa kita ke salah satu lubang kelinci paling terkenal di bahasa ini: urusan multibyte/wide/Unicode/lokalisasi itu.

Kita akan mengintip ke lubang kelinci itu, tapi belum masuk. \ldots Belum!

\section{Escape Sequence}\label{escape-sequence}

\index{Escape sequences|(}

Kita terbiasa dengan string dan karakter berisi huruf biasa, tanda baca, dan angka:

\begin{tcolorbox}[enhanced jigsaw,left=3pt,breakable,colback=white!97.2549!black,boxrule=0.2mm]

\begin{Shaded}
\begin{Highlighting}[]
\DataTypeTok{char} \OperatorTok{*}\NormalTok{s }\OperatorTok{=} \StringTok{"Hello!"}\OperatorTok{;}
\DataTypeTok{char}\NormalTok{ t }\OperatorTok{=} \CharTok{\textquotesingle{}c\textquotesingle{}}\OperatorTok{;}
\end{Highlighting}
\end{Shaded}

\end{tcolorbox}

Tapi bagaimana kalau kita ingin memasukkan karakter khusus yang tidak bisa diketik di keyboard karena memang tidak ada (misalnya ``€''), atau bahkan kalau kita ingin karakter berupa tanda kutip tunggal? Jelas kita tidak bisa menulis ini:

\begin{tcolorbox}[enhanced jigsaw,left=3pt,breakable,colback=white!97.2549!black,boxrule=0.2mm]

\begin{Shaded}
\begin{Highlighting}[]
\DataTypeTok{char}\NormalTok{ t }\OperatorTok{=} \CharTok{\textquotesingle{}}\ErrorTok{\textquotesingle{}}\CharTok{\textquotesingle{};}
\end{Highlighting}
\end{Shaded}

\end{tcolorbox}

\index{{\backslash} backslash escape@\texttt{{\backslash}} backslash escape|(}

Untuk melakukan hal-hal ini, kita memakai yang disebut \emph{escape sequence}. Bentuknya karakter backslash (\hl{\mbox{\texttt{\textbackslash{}}}}) diikuti karakter lain. Kedua (atau lebih) karakter itu bersama-sama punya makna khusus.

Untuk contoh karakter kutip tunggal kita, kita bisa menaruh escape (yaitu \hl{\mbox{\texttt{\textbackslash{}}}}) di depan tanda kutip tunggal di tengah untuk mengatasinya:

\index{{\backslash}' single quote@\texttt{{\backslash}'} single quote|(}

\begin{tcolorbox}[enhanced jigsaw,left=3pt,breakable,colback=white!97.2549!black,boxrule=0.2mm]

\begin{Shaded}
\begin{Highlighting}[]
\DataTypeTok{char}\NormalTok{ t }\OperatorTok{=} \CharTok{\textquotesingle{}}\SpecialCharTok{\textbackslash{}\textquotesingle{}}\CharTok{\textquotesingle{}}\OperatorTok{;}
\end{Highlighting}
\end{Shaded}

\end{tcolorbox}

Sekarang C tahu bahwa \hl{\mbox{\texttt{\textbackslash{}\textquotesingle{}}}} hanya berarti tanda kutip biasa yang ingin kita cetak, bukan akhir dari deretan karakter.

\index{{\backslash}' single quote@\texttt{{\backslash}'} single quote|)}

Dalam konteks ini kamu bisa mengatakan ``backslash'' atau ``escape'' (``escape tanda kutip itu'') dan developer C akan paham maksudmu. Selain itu, ``escape'' dalam konteks ini berbeda dari tombol \hl{\mbox{\texttt{Esc}}} di keyboardmu atau kode ASCII \hl{\mbox{\texttt{ESC}}}.

\subsection{Escape yang Sering Dipakai}\label{escape-yang-sering-dipakai}

Menurut pendapatku yang rendah hati, karakter escape ini mencakup 99,2\%\footnote{Angka itu kukarang sendiri, tapi kemungkinan tidak jauh meleset} dari semua escape.

\index{{\backslash}n newline@\texttt{{\backslash}n} newline} \index{{\backslash}' single quote@\texttt{{\backslash}'} single quote} \index{{\backslash}" double quote@\texttt{{\backslash}"} double quote} \index{{\backslash}{\backslash} backslash@\texttt{{\backslash}{\backslash}} backslash}

{\def\LTcaptype{none} % do not increment counter
\begin{longtable}[]{@{}
  >{\raggedright\arraybackslash}p{(\linewidth - 2\tabcolsep) * \real{0.1429}}
  >{\raggedright\arraybackslash}p{(\linewidth - 2\tabcolsep) * \real{0.8571}}@{}}
\toprule\noalign{}
\begin{minipage}[b]{\linewidth}\raggedright
Kode
\end{minipage} & \begin{minipage}[b]{\linewidth}\raggedright
Deskripsi
\end{minipage} \\
\midrule\noalign{}
\endhead
\bottomrule\noalign{}
\endlastfoot
\hl{\mbox{\texttt{\textbackslash{}n}}} & Karakter newline---saat mencetak, lanjutkan output berikutnya di baris selanjutnya \\
\hl{\mbox{\texttt{\textbackslash{}\textquotesingle{}}}} & Tanda kutip tunggal---dipakai untuk konstanta karakter kutip tunggal \\
\hl{\mbox{\texttt{\textbackslash{}"}}} & Tanda kutip ganda---dipakai untuk tanda kutip ganda dalam string literal \\
\hl{\mbox{\texttt{\textbackslash{}\textbackslash{}}}} & Backslash---dipakai untuk \hl{\mbox{\texttt{\textbackslash{}}}} literal dalam string atau karakter \\
\end{longtable}
}

Berikut beberapa contoh escape dan apa yang dikeluarkannya saat dicetak.

\begin{tcolorbox}[enhanced jigsaw,left=3pt,breakable,colback=white!97.2549!black,boxrule=0.2mm]

\begin{Shaded}
\begin{Highlighting}[]
\NormalTok{printf}\OperatorTok{(}\StringTok{"Use }\SpecialCharTok{\textbackslash{}\textbackslash{}}\StringTok{n for newline}\SpecialCharTok{\textbackslash{}n}\StringTok{"}\OperatorTok{);}  \CommentTok{// Use \textbackslash{}n for newline}
\NormalTok{printf}\OperatorTok{(}\StringTok{"Say }\SpecialCharTok{\textbackslash{}"}\StringTok{hello}\SpecialCharTok{\textbackslash{}"}\StringTok{!}\SpecialCharTok{\textbackslash{}n}\StringTok{"}\OperatorTok{);}       \CommentTok{// Say "hello"!}
\NormalTok{printf}\OperatorTok{(}\StringTok{"}\SpecialCharTok{\%c\textbackslash{}n}\StringTok{"}\OperatorTok{,} \CharTok{\textquotesingle{}}\SpecialCharTok{\textbackslash{}\textquotesingle{}}\CharTok{\textquotesingle{}}\OperatorTok{);}             \CommentTok{// \textquotesingle{}}
\end{Highlighting}
\end{Shaded}

\end{tcolorbox}

\subsection{Escape yang Jarang Dipakai}\label{escape-yang-jarang-dipakai}

Tapi masih ada escape lain! Hanya saja jarang terlihat.

\index{{\backslash}a alert@\texttt{{\backslash}a} alert} \index{{\backslash}b backspace@\texttt{{\backslash}b} backspace} \index{{\backslash}f formfeed@\texttt{{\backslash}f} formfeed} \index{{\backslash}r carriage return@\texttt{{\backslash}r} carriage return} \index{{\backslash}t tab@\texttt{{\backslash}t} tab} \index{{\backslash}v vertical tab@\texttt{{\backslash}v} vertical tab} \index{{\backslash}? question mark@\texttt{{\backslash}?} question mark}

{\def\LTcaptype{none} % do not increment counter
\begin{longtable}[]{@{}
  >{\raggedright\arraybackslash}p{(\linewidth - 2\tabcolsep) * \real{0.0909}}
  >{\raggedright\arraybackslash}p{(\linewidth - 2\tabcolsep) * \real{0.9091}}@{}}
\toprule\noalign{}
\begin{minipage}[b]{\linewidth}\raggedright
Kode
\end{minipage} & \begin{minipage}[b]{\linewidth}\raggedright
Deskripsi
\end{minipage} \\
\midrule\noalign{}
\endhead
\bottomrule\noalign{}
\endlastfoot
\hl{\mbox{\texttt{\textbackslash{}a}}} & Alert. Membuat terminal mengeluarkan bunyi atau berkedip, atau keduanya! \\
\hl{\mbox{\texttt{\textbackslash{}b}}} & Backspace. Memundurkan kursor satu karakter. Tidak menghapus karakternya. \\
\hl{\mbox{\texttt{\textbackslash{}f}}} & Formfeed. Berpindah ke ``halaman'' berikutnya, tapi itu tidak punya banyak makna di zaman modern. Di sistemku, ini berperilaku seperti \hl{\mbox{\texttt{\textbackslash{}v}}}. \\
\hl{\mbox{\texttt{\textbackslash{}r}}} & Return. Pindah ke awal baris yang sama. \\
\hl{\mbox{\texttt{\textbackslash{}t}}} & Tab horizontal. Pindah ke tab stop horizontal berikutnya. Di mesinku, ini sejajar dengan kolom kelipatan 8, tapi hasilnya bisa berbeda di tempatmu. \\
\hl{\mbox{\texttt{\textbackslash{}v}}} & Tab vertikal. Pindah ke tab stop vertikal berikutnya. Di mesinku, ini pindah ke kolom yang sama di baris berikutnya. \\
\hl{\mbox{\texttt{\textbackslash{}?}}} & Tanda tanya literal. Kadang kamu butuh ini untuk menghindari trigraph, seperti ditunjukkan di bawah. \\
\end{longtable}
}

\subsubsection{Pembaruan Status Satu Baris}\label{pembaruan-status-satu-baris}

\index{{\backslash}b backspace@\texttt{{\backslash}b} backspace|(} \index{{\backslash}r carriage return@\texttt{{\backslash}r} carriage return|(}

Kasus pemakaian \hl{\mbox{\texttt{\textbackslash{}b}}} atau \hl{\mbox{\texttt{\textbackslash{}r}}} adalah menampilkan pembaruan status yang muncul di baris yang sama di layar dan tidak membuat tampilan bergulir. Berikut contoh yang melakukan hitung mundur dari 10. (Kalau compiler-mu tidak mendukung threading, kamu bisa memakai fungsi POSIX non-standar \hl{\mbox{\texttt{sleep()}}} dari \hl{\mbox{\texttt{\textless{}unistd.h\textgreater{}}}}---kalau bukan di sistem mirip Unix, cari platformmu dan \hl{\mbox{\texttt{sleep}}} untuk padanannya.)

\begin{tcolorbox}[enhanced jigsaw,left=3pt,breakable,colback=white!97.2549!black,boxrule=0.2mm]

\begin{Shaded}
\begin{Highlighting}[numbers=left,,]
\PreprocessorTok{\#include }\ImportTok{\textless{}stdio.h\textgreater{}}
\PreprocessorTok{\#include }\ImportTok{\textless{}threads.h\textgreater{}}

\DataTypeTok{int}\NormalTok{ main}\OperatorTok{(}\DataTypeTok{void}\OperatorTok{)}
\OperatorTok{\{}
    \ControlFlowTok{for} \OperatorTok{(}\DataTypeTok{int}\NormalTok{ i }\OperatorTok{=} \DecValTok{10}\OperatorTok{;}\NormalTok{ i }\OperatorTok{\textgreater{}=} \DecValTok{0}\OperatorTok{;}\NormalTok{ i}\OperatorTok{{-}{-})} \OperatorTok{\{}
\NormalTok{        printf}\OperatorTok{(}\StringTok{"}\SpecialCharTok{\textbackslash{}r}\StringTok{T minus }\SpecialCharTok{\%d}\StringTok{ second}\SpecialCharTok{\%s}\StringTok{... }\SpecialCharTok{\textbackslash{}b}\StringTok{"}\OperatorTok{,}\NormalTok{ i}\OperatorTok{,}\NormalTok{ i }\OperatorTok{!=} \DecValTok{1}\OperatorTok{?} \StringTok{"s"}\OperatorTok{:} \StringTok{""}\OperatorTok{);}

\NormalTok{        fflush}\OperatorTok{(}\NormalTok{stdout}\OperatorTok{);}  \CommentTok{// Paksa output untuk diperbarui}

        \CommentTok{// Tidur selama 1 detik}
\NormalTok{        thrd\_sleep}\OperatorTok{(\&(}\KeywordTok{struct}\NormalTok{ timespec}\OperatorTok{)\{.}\NormalTok{tv\_sec}\OperatorTok{=}\DecValTok{1}\OperatorTok{\},}\NormalTok{ NULL}\OperatorTok{);}
    \OperatorTok{\}}

\NormalTok{    printf}\OperatorTok{(}\StringTok{"}\SpecialCharTok{\textbackslash{}r}\StringTok{Liftoff!             }\SpecialCharTok{\textbackslash{}n}\StringTok{"}\OperatorTok{);}
\OperatorTok{\}}
\end{Highlighting}
\end{Shaded}

\end{tcolorbox}

Cukup banyak hal terjadi di baris 7. Pertama, kita mengawali dengan \hl{\mbox{\texttt{\textbackslash{}r}}} untuk membawa kita ke awal baris saat ini, lalu menimpa apa pun yang ada di sana dengan hitung mundur saat ini. (Ada operator ternary di situ untuk memastikan kita mencetak \hl{\mbox{\texttt{1\ second}}} dan bukan \hl{\mbox{\texttt{1\ seconds}}}.)

Perhatikan juga ada spasi setelah \hl{\mbox{\texttt{...}}} Itu supaya kita menimpa \hl{\mbox{\texttt{.}}} terakhir dengan benar ketika \hl{\mbox{\texttt{i}}} turun dari \hl{\mbox{\texttt{10}}} ke \hl{\mbox{\texttt{9}}} dan kita jadi satu kolom lebih sempit. Coba tanpa spasi itu untuk melihat maksudnya.

Dan kita tutup dengan \hl{\mbox{\texttt{\textbackslash{}b}}} untuk mundur melewati spasi itu supaya kursor berada tepat di ujung baris dengan cara yang enak dipandang.

\index{{\backslash}b backspace@\texttt{{\backslash}b} backspace|)}

Perhatikan bahwa baris 15 juga punya banyak spasi di akhir untuk menimpa karakter yang sudah ada dari hitung mundur.

Terakhir, ada \hl{\mbox{\texttt{fflush(stdout)}}} yang aneh di situ, apa pun artinya. Jawaban singkatnya, kebanyakan terminal secara bawaan \emph{line buffered}, artinya mereka tidak benar-benar menampilkan apa pun sampai ada karakter newline. Karena kita tidak punya newline (kita hanya punya \hl{\mbox{\texttt{\textbackslash{}r}}}), tanpa baris ini program akan diam saja sampai \hl{\mbox{\texttt{Liftoff!}}} lalu mencetak semuanya sekaligus dalam sekejap. \hl{\mbox{\texttt{fflush()}}} menimpa perilaku ini dan memaksa output terjadi \emph{sekarang juga}.

\index{{\backslash}r carriage return@\texttt{{\backslash}r} carriage return|)}

\subsubsection{Escape Tanda Tanya}\label{escape-tanda-tanya}

\index{{\backslash}? question mark@\texttt{{\backslash}?} question mark|(}

Untuk apa repot-repot dengan ini? Toh ini berfungsi dengan baik:

\begin{tcolorbox}[enhanced jigsaw,left=3pt,breakable,colback=white!97.2549!black,boxrule=0.2mm]

\begin{Shaded}
\begin{Highlighting}[]
\NormalTok{printf}\OperatorTok{(}\StringTok{"Doesn\textquotesingle{}t it?}\SpecialCharTok{\textbackslash{}n}\StringTok{"}\OperatorTok{);}
\end{Highlighting}
\end{Shaded}

\end{tcolorbox}

Dan berfungsi baik juga dengan escape:

\begin{tcolorbox}[enhanced jigsaw,left=3pt,breakable,colback=white!97.2549!black,boxrule=0.2mm]

\begin{Shaded}
\begin{Highlighting}[]
\NormalTok{printf}\OperatorTok{(}\StringTok{"Doesn\textquotesingle{}t it}\SpecialCharTok{\textbackslash{}?\textbackslash{}n}\StringTok{"}\OperatorTok{);}   \CommentTok{// Perhatikan \textbackslash{}?}
\end{Highlighting}
\end{Shaded}

\end{tcolorbox}

Jadi apa gunanya??!

\index{Trigraphs|(}

Mari kita lebih tegas dengan satu tanda tanya lagi dan tanda seru:

\begin{tcolorbox}[enhanced jigsaw,left=3pt,breakable,colback=white!97.2549!black,boxrule=0.2mm]

\begin{Shaded}
\begin{Highlighting}[]
\NormalTok{printf}\OperatorTok{(}\StringTok{"Doesn\textquotesingle{}t it??!}\SpecialCharTok{\textbackslash{}n}\StringTok{"}\OperatorTok{);}
\end{Highlighting}
\end{Shaded}

\end{tcolorbox}

Waktu kukompilasi, aku mendapat peringatan ini:

\begin{tcolorbox}[enhanced jigsaw,left=3pt,breakable,colback=white!97.2549!black,boxrule=0.2mm]

\begin{Shaded}
\begin{Highlighting}[]
\ExtensionTok{foo.c:}\NormalTok{ In function ‘main’:}
\ExtensionTok{foo.c:5:23:}\NormalTok{ warning: trigraph }\PreprocessorTok{??}\NormalTok{! converted to }\KeywordTok{|} \ExtensionTok{[{-}Wtrigraphs]}
    \ExtensionTok{5} \KeywordTok{|}     \BuiltInTok{printf}\OperatorTok{(}\StringTok{"Doesn\textquotesingle{}t it??!\textbackslash{}n"}\OperatorTok{)}\KeywordTok{;}
      \KeywordTok{|}    
\end{Highlighting}
\end{Shaded}

\end{tcolorbox}

Dan menjalankannya memberi hasil yang tidak disangka:

\begin{tcolorbox}[enhanced jigsaw,left=3pt,breakable,colback=white!97.2549!black,boxrule=0.2mm]

\begin{Shaded}
\begin{Highlighting}[]
\NormalTok{Doesn\textquotesingle{}t it|}
\end{Highlighting}
\end{Shaded}

\end{tcolorbox}

Jadi \emph{trigraph}? Apa-apaan ini??!

Aku yakin kita akan mengunjungi lagi sudut berdebu bahasa ini nanti, tapi singkatnya compiler mencari tiga serangkai karakter tertentu yang diawali \hl{\mbox{\texttt{??}}} dan menggantinya dengan karakter lain. Jadi kalau kamu di terminal kuno tanpa simbol pipa (\hl{\mbox{\texttt{\textbar{}}}}) di keyboard, kamu bisa mengetik \hl{\mbox{\texttt{??!}}} sebagai gantinya.

Kamu bisa memperbaikinya dengan meng-escape tanda tanya kedua, seperti ini:

\begin{tcolorbox}[enhanced jigsaw,left=3pt,breakable,colback=white!97.2549!black,boxrule=0.2mm]

\begin{Shaded}
\begin{Highlighting}[]
\NormalTok{printf}\OperatorTok{(}\StringTok{"Doesn\textquotesingle{}t it?}\SpecialCharTok{\textbackslash{}?}\StringTok{!}\SpecialCharTok{\textbackslash{}n}\StringTok{"}\OperatorTok{);}
\end{Highlighting}
\end{Shaded}

\end{tcolorbox}

Lalu ia dikompilasi dan berfungsi seperti yang diharapkan.

Zaman sekarang, tentu saja, tidak ada yang memakai trigraph. Tapi \hl{\mbox{\texttt{??!}}} itu kadang muncul kalau kamu memutuskan memakainya dalam string untuk penekanan.

\index{Trigraphs|)} \index{{\backslash}? question mark@\texttt{{\backslash}?} question mark|)}

\subsection{Escape Numerik}\label{escape-numerik}

Selain itu, ada cara untuk menentukan konstanta numerik atau nilai karakter lain di dalam string atau konstanta karakter.

Kalau kamu tahu representasi oktal atau heksadesimal sebuah byte, kamu bisa menyertakannya dalam string atau konstanta karakter.

Tabel berikut memakai angka contoh, tapi angka hex atau oktal apa pun bisa dipakai. Beri nol di depan kalau perlu agar jumlah digitnya sesuai.

\index{{\backslash}123 octal value@\texttt{{\backslash}123} octal value} \index{{\backslash}x12 hexadecimal value@\texttt{{\backslash}x12} hexadecimal value} \index{{\backslash}u Unicode escape@\texttt{{\backslash}u} Unicode escape} \index{{\backslash}U Unicode escape@\texttt{{\backslash}U} Unicode escape}

{\def\LTcaptype{none} % do not increment counter
\begin{longtable}[]{@{}
  >{\raggedright\arraybackslash}p{(\linewidth - 2\tabcolsep) * \real{0.1429}}
  >{\raggedright\arraybackslash}p{(\linewidth - 2\tabcolsep) * \real{0.8571}}@{}}
\toprule\noalign{}
\begin{minipage}[b]{\linewidth}\raggedright
Kode
\end{minipage} & \begin{minipage}[b]{\linewidth}\raggedright
Deskripsi
\end{minipage} \\
\midrule\noalign{}
\endhead
\bottomrule\noalign{}
\endlastfoot
\hl{\mbox{\texttt{\textbackslash{}123}}} & Sematkan byte bernilai oktal \hl{\mbox{\texttt{123}}}, tepat 3 digit. \\
\hl{\mbox{\texttt{\textbackslash{}x4D}}} & Sematkan byte bernilai hex \hl{\mbox{\texttt{4D}}}, 2 digit. \\
\hl{\mbox{\texttt{\textbackslash{}u2620}}} & Sematkan karakter Unicode pada code point bernilai hex \hl{\mbox{\texttt{2620}}}, 4 digit. \\
\hl{\mbox{\texttt{\textbackslash{}U0001243F}}} & Sematkan karakter Unicode pada code point bernilai hex \hl{\mbox{\texttt{1243F}}}, 8 digit. \\
\end{longtable}
}

Berikut contoh notasi oktal yang jarang dipakai untuk merepresentasikan huruf \hl{\mbox{\texttt{B}}} di antara \hl{\mbox{\texttt{A}}} dan \hl{\mbox{\texttt{C}}}. Biasanya ini dipakai untuk semacam karakter khusus yang tidak bisa dicetak, tapi kita punya cara lain untuk itu, di bawah, dan ini hanya demo oktal:

\index{{\backslash}123 octal value@\texttt{{\backslash}123} octal value|(}

\begin{tcolorbox}[enhanced jigsaw,left=3pt,breakable,colback=white!97.2549!black,boxrule=0.2mm]

\begin{Shaded}
\begin{Highlighting}[]
\NormalTok{printf}\OperatorTok{(}\StringTok{"A}\ErrorTok{\textbackslash{}1}\StringTok{02C}\SpecialCharTok{\textbackslash{}n}\StringTok{"}\OperatorTok{);}  \CommentTok{// 102 adalah \textasciigrave{}B\textasciigrave{} di ASCII/UTF{-}8}
\end{Highlighting}
\end{Shaded}

\end{tcolorbox}

Perhatikan tidak ada nol di depan pada angka oktal saat disertakan dengan cara ini. Tapi harus tiga karakter, jadi beri nol di depan kalau perlu.

\index{{\backslash}123 octal value@\texttt{{\backslash}123} octal value|)}

\index{{\backslash}x12 hexadecimal value@\texttt{{\backslash}x12} hexadecimal value|(}

Tapi zaman sekarang jauh lebih umum memakai konstanta hex. Berikut demo yang sebaiknya tidak kamu pakai, tapi ia menunjukkan penyematan byte UTF-8 0xE2, 0x80, dan 0xA2 dalam string, yang sesuai dengan karakter ``bullet'' Unicode (•).

\begin{tcolorbox}[enhanced jigsaw,left=3pt,breakable,colback=white!97.2549!black,boxrule=0.2mm]

\begin{Shaded}
\begin{Highlighting}[]
\NormalTok{printf}\OperatorTok{(}\StringTok{"}\SpecialCharTok{\textbackslash{}xE2\textbackslash{}x80\textbackslash{}xA2}\StringTok{ Bullet 1}\SpecialCharTok{\textbackslash{}n}\StringTok{"}\OperatorTok{);}
\NormalTok{printf}\OperatorTok{(}\StringTok{"}\SpecialCharTok{\textbackslash{}xE2\textbackslash{}x80\textbackslash{}xA2}\StringTok{ Bullet 2}\SpecialCharTok{\textbackslash{}n}\StringTok{"}\OperatorTok{);}
\NormalTok{printf}\OperatorTok{(}\StringTok{"}\SpecialCharTok{\textbackslash{}xE2\textbackslash{}x80\textbackslash{}xA2}\StringTok{ Bullet 3}\SpecialCharTok{\textbackslash{}n}\StringTok{"}\OperatorTok{);}
\end{Highlighting}
\end{Shaded}

\end{tcolorbox}

Menghasilkan output berikut kalau kamu di konsol UTF-8 (atau mungkin sampah kalau bukan):

\begin{tcolorbox}[enhanced jigsaw,left=3pt,breakable,colback=white!97.2549!black,boxrule=0.2mm]

\begin{Shaded}
\begin{Highlighting}[]
\NormalTok{• Bullet 1}
\NormalTok{• Bullet 2}
\NormalTok{• Bullet 3}
\end{Highlighting}
\end{Shaded}

\end{tcolorbox}

\index{{\backslash}x12 hexadecimal value@\texttt{{\backslash}x12} hexadecimal value|)}

\index{{\backslash}u Unicode escape@\texttt{{\backslash}u} Unicode escape|(} \index{{\backslash}U Unicode escape@\texttt{{\backslash}U} Unicode escape|(}

Tapi itu cara yang jelek untuk Unicode. Kamu bisa memakai escape \hl{\mbox{\texttt{\textbackslash{}u}}} (16-bit) atau \hl{\mbox{\texttt{\textbackslash{}U}}} (32-bit) untuk merujuk ke Unicode lewat nomor code point. Bullet adalah \hl{\mbox{\texttt{2022}}} (hex) di Unicode, jadi kamu bisa melakukan ini dan mendapat hasil yang lebih portabel:

\begin{tcolorbox}[enhanced jigsaw,left=3pt,breakable,colback=white!97.2549!black,boxrule=0.2mm]

\begin{Shaded}
\begin{Highlighting}[]
\NormalTok{printf}\OperatorTok{(}\StringTok{"}\SpecialCharTok{\textbackslash{}u2022}\StringTok{ Bullet 1}\SpecialCharTok{\textbackslash{}n}\StringTok{"}\OperatorTok{);}
\NormalTok{printf}\OperatorTok{(}\StringTok{"}\SpecialCharTok{\textbackslash{}u2022}\StringTok{ Bullet 2}\SpecialCharTok{\textbackslash{}n}\StringTok{"}\OperatorTok{);}
\NormalTok{printf}\OperatorTok{(}\StringTok{"}\SpecialCharTok{\textbackslash{}u2022}\StringTok{ Bullet 3}\SpecialCharTok{\textbackslash{}n}\StringTok{"}\OperatorTok{);}
\end{Highlighting}
\end{Shaded}

\end{tcolorbox}

Pastikan memberi \hl{\mbox{\texttt{\textbackslash{}u}}} nol di depan secukupnya untuk mencapai empat karakter, dan \hl{\mbox{\texttt{\textbackslash{}U}}} secukupnya untuk mencapai delapan.

\index{{\backslash}u Unicode escape@\texttt{{\backslash}u} Unicode escape|)}

Misalnya, bullet itu bisa ditulis dengan \hl{\mbox{\texttt{\textbackslash{}U}}} dan empat nol di depan:

\begin{tcolorbox}[enhanced jigsaw,left=3pt,breakable,colback=white!97.2549!black,boxrule=0.2mm]

\begin{Shaded}
\begin{Highlighting}[]
\NormalTok{printf}\OperatorTok{(}\StringTok{"}\SpecialCharTok{\textbackslash{}U00002022}\StringTok{ Bullet 1}\SpecialCharTok{\textbackslash{}n}\StringTok{"}\OperatorTok{);}
\end{Highlighting}
\end{Shaded}

\end{tcolorbox}

\index{{\backslash}U Unicode escape@\texttt{{\backslash}U} Unicode escape|)}

Tapi siapa yang punya waktu untuk bertele-tele begitu?

\index{{\backslash} backslash escape@\texttt{{\backslash}} backslash escape|)} \index{Escape sequences|)}

\chapter{\texorpdfstring{Enumerated Type: \hl{\mbox{\texttt{enum}}}}{Enumerated Type: enum}}\label{enumerated-type-enum}

\index{enum enumerated types@\texttt{enum} enumerated types|(}

C menawarkan cara lain untuk memiliki nilai integer konstan dengan nama: \hl{\mbox{\texttt{enum}}}.

Misalnya:

\begin{tcolorbox}[enhanced jigsaw,left=3pt,breakable,colback=white!97.2549!black,boxrule=0.2mm]

\begin{Shaded}
\begin{Highlighting}[]
\KeywordTok{enum} \OperatorTok{\{}
\NormalTok{  ONE}\OperatorTok{=}\DecValTok{1}\OperatorTok{,}
\NormalTok{  TWO}\OperatorTok{=}\DecValTok{2}
\OperatorTok{\};}

\NormalTok{printf}\OperatorTok{(}\StringTok{"}\SpecialCharTok{\%d}\StringTok{ }\SpecialCharTok{\%d}\StringTok{"}\OperatorTok{,}\NormalTok{ ONE}\OperatorTok{,}\NormalTok{ TWO}\OperatorTok{);}  \CommentTok{// 1 2}
\end{Highlighting}
\end{Shaded}

\end{tcolorbox}

Dalam beberapa hal, ini bisa lebih baik---atau berbeda---dibanding memakai \hl{\mbox{\texttt{\#define}}}. Perbedaan utamanya:

\begin{itemize}
\tightlist
\item
  \hl{\mbox{\texttt{enum}}} hanya bisa berupa tipe integer.
\item
  \hl{\mbox{\texttt{\#define}}} bisa mendefinisikan apa saja.
\item
  \hl{\mbox{\texttt{enum}}} sering ditampilkan dengan nama identifier simbolisnya di debugger.
\item
  Angka hasil \hl{\mbox{\texttt{\#define}}} hanya tampil sebagai angka mentah yang lebih sulit diketahui maknanya saat debugging.
\end{itemize}

Karena berupa tipe integer, mereka bisa dipakai di tempat mana pun integer bisa dipakai, termasuk pada dimensi array dan pernyataan \hl{\mbox{\texttt{case}}}.

Mari kita bedah lebih jauh.

\section{\texorpdfstring{Perilaku \hl{\mbox{\texttt{enum}}}}{Perilaku enum}}\label{perilaku-enum}

\subsection{Penomoran}\label{penomoran}

\index{enum enumerated types@\texttt{enum} enumerated types!numbering order|(}

\hl{\mbox{\texttt{enum}}} diberi nomor otomatis kecuali kamu menimpanya.

Penomoran dimulai dari \hl{\mbox{\texttt{0}}}, dan bertambah otomatis dari sana, secara bawaan:

\begin{tcolorbox}[enhanced jigsaw,left=3pt,breakable,colback=white!97.2549!black,boxrule=0.2mm]

\begin{Shaded}
\begin{Highlighting}[]
\KeywordTok{enum} \OperatorTok{\{}
\NormalTok{    SHEEP}\OperatorTok{,}  \CommentTok{// Nilainya 0}
\NormalTok{    WHEAT}\OperatorTok{,}  \CommentTok{// Nilainya 1}
\NormalTok{    WOOD}\OperatorTok{,}   \CommentTok{// Nilainya 2}
\NormalTok{    BRICK}\OperatorTok{,}  \CommentTok{// Nilainya 3}
\NormalTok{    ORE     }\CommentTok{// Nilainya 4}
\OperatorTok{\};}

\NormalTok{printf}\OperatorTok{(}\StringTok{"}\SpecialCharTok{\%d}\StringTok{ }\SpecialCharTok{\%d\textbackslash{}n}\StringTok{"}\OperatorTok{,}\NormalTok{ SHEEP}\OperatorTok{,}\NormalTok{ BRICK}\OperatorTok{);}  \CommentTok{// 0 3}
\end{Highlighting}
\end{Shaded}

\end{tcolorbox}

Kamu bisa memaksakan nilai integer tertentu, seperti yang sudah kita lihat tadi:

\begin{tcolorbox}[enhanced jigsaw,left=3pt,breakable,colback=white!97.2549!black,boxrule=0.2mm]

\begin{Shaded}
\begin{Highlighting}[]
\KeywordTok{enum} \OperatorTok{\{}
\NormalTok{  X}\OperatorTok{=}\DecValTok{2}\OperatorTok{,}
\NormalTok{  Y}\OperatorTok{=}\DecValTok{18}\OperatorTok{,}
\NormalTok{  Z}\OperatorTok{={-}}\DecValTok{2}
\OperatorTok{\};}
\end{Highlighting}
\end{Shaded}

\end{tcolorbox}

Duplikat tidak masalah:

\begin{tcolorbox}[enhanced jigsaw,left=3pt,breakable,colback=white!97.2549!black,boxrule=0.2mm]

\begin{Shaded}
\begin{Highlighting}[]
\KeywordTok{enum} \OperatorTok{\{}
\NormalTok{  X}\OperatorTok{=}\DecValTok{2}\OperatorTok{,}
\NormalTok{  Y}\OperatorTok{=}\DecValTok{2}\OperatorTok{,}
\NormalTok{  Z}\OperatorTok{=}\DecValTok{2}
\OperatorTok{\};}
\end{Highlighting}
\end{Shaded}

\end{tcolorbox}

kalau nilai dihilangkan, penomoran melanjutkan hitungan ke arah positif dari nilai yang terakhir ditentukan. Misalnya:

\begin{tcolorbox}[enhanced jigsaw,left=3pt,breakable,colback=white!97.2549!black,boxrule=0.2mm]

\begin{Shaded}
\begin{Highlighting}[]
\KeywordTok{enum} \OperatorTok{\{}
\NormalTok{  A}\OperatorTok{,}    \CommentTok{// 0, nilai awal bawaan}
\NormalTok{  B}\OperatorTok{,}    \CommentTok{// 1}
\NormalTok{  C}\OperatorTok{=}\DecValTok{4}\OperatorTok{,}  \CommentTok{// 4, diatur manual}
\NormalTok{  D}\OperatorTok{,}    \CommentTok{// 5}
\NormalTok{  E}\OperatorTok{,}    \CommentTok{// 6}
\NormalTok{  F}\OperatorTok{=}\DecValTok{3}\OperatorTok{,}  \CommentTok{// 3, diatur manual}
\NormalTok{  G}\OperatorTok{,}    \CommentTok{// 4}
\NormalTok{  H     }\CommentTok{// 5}
\OperatorTok{\}}
\end{Highlighting}
\end{Shaded}

\end{tcolorbox}

\index{enum enumerated types@\texttt{enum} enumerated types!numbering order|)}

\subsection{Koma di Akhir}\label{koma-di-akhir}

Ini sepenuhnya boleh, kalau itu gayamu:

\begin{tcolorbox}[enhanced jigsaw,left=3pt,breakable,colback=white!97.2549!black,boxrule=0.2mm]

\begin{Shaded}
\begin{Highlighting}[]
\KeywordTok{enum} \OperatorTok{\{}
\NormalTok{  X}\OperatorTok{=}\DecValTok{2}\OperatorTok{,}
\NormalTok{  Y}\OperatorTok{=}\DecValTok{18}\OperatorTok{,}
\NormalTok{  Z}\OperatorTok{={-}}\DecValTok{2}\OperatorTok{,}   \CommentTok{// \textless{}{-}{-} Koma di akhir}
\OperatorTok{\};}
\end{Highlighting}
\end{Shaded}

\end{tcolorbox}

Gaya ini makin populer di bahasa-bahasa dekade terakhir, jadi kamu mungkin senang melihatnya.

\subsection{Scope}\label{scope-1}

\index{enum enumerated types@\texttt{enum} enumerated types!scope|(}

\hl{\mbox{\texttt{enum}}} mengikuti scope seperti yang kamu harapkan. Kalau di file scope, seluruh file bisa melihatnya. Kalau di dalam blok, ia lokal untuk blok itu.

Sangat umum \hl{\mbox{\texttt{enum}}} didefinisikan di header file sehingga bisa di-\hl{\mbox{\texttt{\#include}}} di file scope.

\index{enum enumerated types@\texttt{enum} enumerated types!scope|)}

\subsection{Gaya Penulisan}\label{gaya-penulisan}

Seperti yang kamu perhatikan, umumnya simbol \hl{\mbox{\texttt{enum}}} dideklarasikan dengan huruf besar (dengan underscore).

Ini bukan keharusan, tapi merupakan idiom yang sangat, sangat umum.

\section{\texorpdfstring{\hl{\mbox{\texttt{enum}}}-mu adalah Sebuah Tipe}{enum-mu adalah Sebuah Tipe}}\label{enum-mu-adalah-sebuah-tipe}

Ini hal penting yang perlu diketahui tentang \hl{\mbox{\texttt{enum}}}: ia adalah sebuah tipe, seperti halnya \hl{\mbox{\texttt{struct}}} adalah sebuah tipe.

Kamu bisa memberinya nama tag sehingga bisa merujuk ke tipe itu nanti dan mendeklarasikan variabel bertipe itu.

Nah, karena \hl{\mbox{\texttt{enum}}} adalah tipe integer, kenapa tidak pakai \hl{\mbox{\texttt{int}}} saja?

Di C, alasan terbaiknya adalah kejelasan kode--ini cara bertipe yang rapi untuk menjelaskan pemikiranmu dalam kode. C (tidak seperti C++) sebenarnya tidak memaksa nilai apa pun berada dalam rentang untuk \hl{\mbox{\texttt{enum}}} tertentu.

Mari kita coba contoh ketika kita mendeklarasikan variabel \hl{\mbox{\texttt{r}}} bertipe \hl{\mbox{\texttt{enum\ resource}}} yang bisa menampung nilai-nilai itu:

\begin{tcolorbox}[enhanced jigsaw,left=3pt,breakable,colback=white!97.2549!black,boxrule=0.2mm]

\begin{Shaded}
\begin{Highlighting}[]
\CommentTok{// enum bernama, tipenya "enum resource"}

\KeywordTok{enum}\NormalTok{ resource }\OperatorTok{\{}
\NormalTok{    SHEEP}\OperatorTok{,}
\NormalTok{    WHEAT}\OperatorTok{,}
\NormalTok{    WOOD}\OperatorTok{,}
\NormalTok{    BRICK}\OperatorTok{,}
\NormalTok{    ORE}
\OperatorTok{\};}

\CommentTok{// Deklarasikan variabel "r" bertipe "enum resource"}

\KeywordTok{enum}\NormalTok{ resource r }\OperatorTok{=}\NormalTok{ BRICK}\OperatorTok{;}

\ControlFlowTok{if} \OperatorTok{(}\NormalTok{r }\OperatorTok{==}\NormalTok{ BRICK}\OperatorTok{)} \OperatorTok{\{}
\NormalTok{    printf}\OperatorTok{(}\StringTok{"I\textquotesingle{}ll trade you a brick for two sheep.}\SpecialCharTok{\textbackslash{}n}\StringTok{"}\OperatorTok{);}
\OperatorTok{\}}
\end{Highlighting}
\end{Shaded}

\end{tcolorbox}

Kamu juga bisa me-\hl{\mbox{\texttt{typedef}}}-nya, tentu saja, walaupun secara pribadi aku tidak suka.

\begin{tcolorbox}[enhanced jigsaw,left=3pt,breakable,colback=white!97.2549!black,boxrule=0.2mm]

\begin{Shaded}
\begin{Highlighting}[]
\KeywordTok{typedef} \KeywordTok{enum} \OperatorTok{\{}
\NormalTok{    SHEEP}\OperatorTok{,}
\NormalTok{    WHEAT}\OperatorTok{,}
\NormalTok{    WOOD}\OperatorTok{,}
\NormalTok{    BRICK}\OperatorTok{,}
\NormalTok{    ORE}
\OperatorTok{\}}\NormalTok{ RESOURCE}\OperatorTok{;}

\NormalTok{RESOURCE r }\OperatorTok{=}\NormalTok{ BRICK}\OperatorTok{;}
\end{Highlighting}
\end{Shaded}

\end{tcolorbox}

Jalan pintas lain yang legal tapi jarang adalah mendeklarasikan variabel saat mendeklarasikan \hl{\mbox{\texttt{enum}}}:

\begin{tcolorbox}[enhanced jigsaw,left=3pt,breakable,colback=white!97.2549!black,boxrule=0.2mm]

\begin{Shaded}
\begin{Highlighting}[]
\CommentTok{// Deklarasikan enum dan beberapa variabel terinisialisasi bertipe itu:}

\KeywordTok{enum} \OperatorTok{\{}
\NormalTok{    SHEEP}\OperatorTok{,}
\NormalTok{    WHEAT}\OperatorTok{,}
\NormalTok{    WOOD}\OperatorTok{,}
\NormalTok{    BRICK}\OperatorTok{,}
\NormalTok{    ORE}
\OperatorTok{\}}\NormalTok{ r }\OperatorTok{=}\NormalTok{ BRICK}\OperatorTok{,}\NormalTok{ s }\OperatorTok{=}\NormalTok{ WOOD}\OperatorTok{;}
\end{Highlighting}
\end{Shaded}

\end{tcolorbox}

Kamu juga bisa memberi nama pada \hl{\mbox{\texttt{enum}}} agar bisa dipakai nanti, yang mungkin ingin kamu lakukan pada kebanyakan kasus:

\begin{tcolorbox}[enhanced jigsaw,left=3pt,breakable,colback=white!97.2549!black,boxrule=0.2mm]

\begin{Shaded}
\begin{Highlighting}[]
\CommentTok{// Deklarasikan enum dan beberapa variabel terinisialisasi bertipe itu:}

\KeywordTok{enum}\NormalTok{ resource }\OperatorTok{\{}   \CommentTok{// \textless{}{-}{-} tipenya sekarang "enum resource"}
\NormalTok{    SHEEP}\OperatorTok{,}
\NormalTok{    WHEAT}\OperatorTok{,}
\NormalTok{    WOOD}\OperatorTok{,}
\NormalTok{    BRICK}\OperatorTok{,}
\NormalTok{    ORE}
\OperatorTok{\}}\NormalTok{ r }\OperatorTok{=}\NormalTok{ BRICK}\OperatorTok{,}\NormalTok{ s }\OperatorTok{=}\NormalTok{ WOOD}\OperatorTok{;}
\end{Highlighting}
\end{Shaded}

\end{tcolorbox}

Singkatnya, \hl{\mbox{\texttt{enum}}} adalah cara yang bagus untuk menulis kode yang rapi, ber-scope, dan bertipe.

\index{enum enumerated types@\texttt{enum} enumerated types|)}

\chapter{Pointer III: Pointer ke Pointer dan Lainnya}\label{pointer-iii-pointer-ke-pointer-dan-lainnya}

Di sinilah kita membahas pemakaian pointer tingkat menengah dan lanjut. Kalau kamu belum menguasai pointer, ulas dulu bab-bab sebelumnya tentang \hyperref[pointers]{pointer} dan \hyperref[pointers2]{aritmetika pointer} sebelum mulai dengan materi ini.

\section{Pointer ke Pointer}\label{pointer-ke-pointer}

\index{Pointers!to pointers|(}

Kalau kamu bisa punya pointer ke sebuah variabel, dan variabel bisa berupa pointer, bisakah kamu punya pointer ke variabel yang itu sendiri adalah pointer?

Bisa! Itulah pointer ke pointer, yang disimpan dalam variabel bertipe pointer-pointer.

Sebelum membedahnya, aku ingin mencoba membangun \emph{gut feel} tentang cara kerja pointer ke pointer.

Ingat bahwa pointer hanyalah sebuah angka. Angka yang mewakili indeks dalam memori komputer, biasanya indeks yang menyimpan nilai yang menarik minat kita karena suatu alasan.

Pointer itu, yang berupa angka, harus disimpan di suatu tempat. Dan tempat itu adalah memori, sama seperti yang lainnya\footnote{Ada setan di detailnya untuk nilai yang hanya disimpan di register, tapi untuk keperluan kita di sini, hal itu aman untuk diabaikan. Lagi pula spesifikasi C tidak mengambil sikap soal ``register'' ini di luar kata kunci \hl{\mbox{\texttt{register}}}, yang deskripsinya tidak menyebut register.}.

Tapi karena disimpan di memori, ia pasti punya indeks tempat ia disimpan, kan? Pointer itu pasti punya indeks di memori tempat ia disimpan. Dan indeks itu adalah sebuah angka. Itulah alamat pointer tersebut. Itulah pointer ke pointer.

Kita mulai dengan pointer biasa ke sebuah \hl{\mbox{\texttt{int}}}, dari bab-bab sebelumnya:

\begin{tcolorbox}[enhanced jigsaw,left=3pt,breakable,colback=white!97.2549!black,boxrule=0.2mm]

\begin{Shaded}
\begin{Highlighting}[numbers=left,,]
\PreprocessorTok{\#include }\ImportTok{\textless{}stdio.h\textgreater{}}

\DataTypeTok{int}\NormalTok{ main}\OperatorTok{(}\DataTypeTok{void}\OperatorTok{)}
\OperatorTok{\{}
    \DataTypeTok{int}\NormalTok{ x }\OperatorTok{=} \DecValTok{3490}\OperatorTok{;}  \CommentTok{// Tipe: int}
    \DataTypeTok{int} \OperatorTok{*}\NormalTok{p }\OperatorTok{=} \OperatorTok{\&}\NormalTok{x}\OperatorTok{;}   \CommentTok{// Tipe: pointer ke int}

\NormalTok{    printf}\OperatorTok{(}\StringTok{"}\SpecialCharTok{\%d\textbackslash{}n}\StringTok{"}\OperatorTok{,} \OperatorTok{*}\NormalTok{p}\OperatorTok{);}  \CommentTok{// 3490}
\OperatorTok{\}}
\end{Highlighting}
\end{Shaded}

\end{tcolorbox}

Cukup lugas, kan? Ada dua tipe yang terlibat: \hl{\mbox{\texttt{int}}} dan \hl{\mbox{\texttt{int*}}}, dan kita mengatur \hl{\mbox{\texttt{p}}} agar menunjuk ke \hl{\mbox{\texttt{x}}}. Lalu kita bisa men-dereference \hl{\mbox{\texttt{p}}} di baris 8 dan mencetak nilai \hl{\mbox{\texttt{3490}}}.

Tapi, seperti yang kita katakan, kita bisa punya pointer ke variabel apa pun\ldots{} jadi apakah itu berarti kita bisa punya pointer ke \hl{\mbox{\texttt{p}}}?

Dengan kata lain, ekspresi ini bertipe apa?

\begin{tcolorbox}[enhanced jigsaw,left=3pt,breakable,colback=white!97.2549!black,boxrule=0.2mm]

\begin{Shaded}
\begin{Highlighting}[]
\DataTypeTok{int}\NormalTok{ x }\OperatorTok{=} \DecValTok{3490}\OperatorTok{;}  \CommentTok{// Tipe: int}
\DataTypeTok{int} \OperatorTok{*}\NormalTok{p }\OperatorTok{=} \OperatorTok{\&}\NormalTok{x}\OperatorTok{;}   \CommentTok{// Tipe: pointer ke int}

\OperatorTok{\&}\NormalTok{p  }\CommentTok{// \textless{}{-}{-} Apa tipe alamat p? Alias pointer ke p?}
\end{Highlighting}
\end{Shaded}

\end{tcolorbox}

Kalau \hl{\mbox{\texttt{x}}} adalah \hl{\mbox{\texttt{int}}}, maka \hl{\mbox{\texttt{\&x}}} adalah pointer ke \hl{\mbox{\texttt{int}}} yang kita simpan di \hl{\mbox{\texttt{p}}} yang bertipe \hl{\mbox{\texttt{int*}}}. Ikut? (Ulangi paragraf ini sampai paham!)

Dan karena itu \hl{\mbox{\texttt{\&p}}} adalah pointer ke \hl{\mbox{\texttt{int*}}}, alias ``pointer ke pointer ke \hl{\mbox{\texttt{int}}}''. Alias ``\hl{\mbox{\texttt{int}}}-pointer-pointer''.

Paham? (Ulangi paragraf sebelumnya sampai paham!)

Kita menulis tipe ini dengan dua asterisk: \hl{\mbox{\texttt{int\ **}}}. Kita lihat pemakaiannya.

\begin{tcolorbox}[enhanced jigsaw,left=3pt,breakable,colback=white!97.2549!black,boxrule=0.2mm]

\begin{Shaded}
\begin{Highlighting}[numbers=left,,]
\PreprocessorTok{\#include }\ImportTok{\textless{}stdio.h\textgreater{}}

\DataTypeTok{int}\NormalTok{ main}\OperatorTok{(}\DataTypeTok{void}\OperatorTok{)}
\OperatorTok{\{}
    \DataTypeTok{int}\NormalTok{ x }\OperatorTok{=} \DecValTok{3490}\OperatorTok{;}  \CommentTok{// Tipe: int}
    \DataTypeTok{int} \OperatorTok{*}\NormalTok{p }\OperatorTok{=} \OperatorTok{\&}\NormalTok{x}\OperatorTok{;}   \CommentTok{// Tipe: pointer ke int}
    \DataTypeTok{int} \OperatorTok{**}\NormalTok{q }\OperatorTok{=} \OperatorTok{\&}\NormalTok{p}\OperatorTok{;}  \CommentTok{// Tipe: pointer ke pointer ke int}

\NormalTok{    printf}\OperatorTok{(}\StringTok{"}\SpecialCharTok{\%d}\StringTok{ }\SpecialCharTok{\%d\textbackslash{}n}\StringTok{"}\OperatorTok{,} \OperatorTok{*}\NormalTok{p}\OperatorTok{,} \OperatorTok{**}\NormalTok{q}\OperatorTok{);}  \CommentTok{// 3490 3490}
\OperatorTok{\}}
\end{Highlighting}
\end{Shaded}

\end{tcolorbox}

Mari kita karang beberapa alamat pura-pura untuk nilai-nilai di atas sebagai contoh dan lihat seperti apa ketiga variabel ini di memori. Nilai alamat di bawah ini hanya karanganku untuk keperluan contoh:

{\def\LTcaptype{none} % do not increment counter
\begin{longtable}[]{@{}lll@{}}
\toprule\noalign{}
Variabel & Disimpan di Alamat & Nilai yang Tersimpan di Sana \\
\midrule\noalign{}
\endhead
\bottomrule\noalign{}
\endlastfoot
\hl{\mbox{\texttt{x}}} & \hl{\mbox{\texttt{28350}}} & \hl{\mbox{\texttt{3490}}}---nilai dari kode \\
\hl{\mbox{\texttt{p}}} & \hl{\mbox{\texttt{29122}}} & \hl{\mbox{\texttt{28350}}}---alamat \hl{\mbox{\texttt{x}}}! \\
\hl{\mbox{\texttt{q}}} & \hl{\mbox{\texttt{30840}}} & \hl{\mbox{\texttt{29122}}}---alamat \hl{\mbox{\texttt{p}}}! \\
\end{longtable}
}

Sekarang mari kita coba sungguhan di komputerku\footnote{Kemungkinan besar angka di komputermu berbeda.} dan mencetak nilai pointer dengan \hl{\mbox{\texttt{\%p}}}, lalu aku ulangi tabel yang sama dengan referensi sebenarnya (dicetak dalam hex).

{\def\LTcaptype{none} % do not increment counter
\begin{longtable}[]{@{}lll@{}}
\toprule\noalign{}
Variabel & Disimpan di Alamat & Nilai yang Tersimpan di Sana \\
\midrule\noalign{}
\endhead
\bottomrule\noalign{}
\endlastfoot
\hl{\mbox{\texttt{x}}} & \hl{\mbox{\texttt{0x7ffd96a07b94}}} & \hl{\mbox{\texttt{3490}}}---nilai dari kode \\
\hl{\mbox{\texttt{p}}} & \hl{\mbox{\texttt{0x7ffd96a07b98}}} & \hl{\mbox{\texttt{0x7ffd96a07b94}}}---alamat \hl{\mbox{\texttt{x}}}! \\
\hl{\mbox{\texttt{q}}} & \hl{\mbox{\texttt{0x7ffd96a07ba0}}} & \hl{\mbox{\texttt{0x7ffd96a07b98}}}---alamat \hl{\mbox{\texttt{p}}}! \\
\end{longtable}
}

Kamu bisa lihat alamat-alamat itu sama kecuali byte terakhir, jadi fokus saja ke sana.

Di sistemku, \hl{\mbox{\texttt{int}}} berukuran 4 byte, itulah kenapa kita melihat alamat naik 4 dari \hl{\mbox{\texttt{x}}} ke \hl{\mbox{\texttt{p}}}\footnote{Sama sekali tidak ada di spesifikasi yang mengatakan ini akan selalu begini, tapi kebetulan di sistemku begitu.} lalu naik 8 dari \hl{\mbox{\texttt{p}}} ke \hl{\mbox{\texttt{q}}}. Di sistemku, semua pointer berukuran 8 byte.

Apakah penting kalau itu \hl{\mbox{\texttt{int*}}} atau \hl{\mbox{\texttt{int**}}}? Apakah salah satunya lebih banyak byte-nya? Tidak! Ingat bahwa semua pointer adalah alamat, yaitu indeks ke dalam memori. Dan di mesinku, sebuah indeks bisa direpresentasikan dengan 8 byte\ldots{} tidak peduli apa yang disimpan di indeks itu.

Sekarang lihat apa yang kita lakukan di baris 9 pada contoh sebelumnya: kita \emph{men-dereference ganda} \hl{\mbox{\texttt{q}}} untuk kembali ke \hl{\mbox{\texttt{3490}}}.

Ini bagian penting tentang pointer dan pointer ke pointer:

\begin{itemize}
\tightlist
\item
  Kamu bisa mendapat pointer ke apa pun dengan \hl{\mbox{\texttt{\&}}} (termasuk ke pointer!)
\item
  Kamu bisa mendapat yang ditunjuk sebuah pointer dengan \hl{\mbox{\texttt{*}}} (termasuk pointer!)
\end{itemize}

Jadi kamu bisa menganggap \hl{\mbox{\texttt{\&}}} dipakai untuk membuat pointer, dan \hl{\mbox{\texttt{*}}} adalah kebalikannya---ia berjalan ke arah yang berlawanan dari \hl{\mbox{\texttt{\&}}}---untuk sampai ke yang ditunjuk.

Dari sisi tipe, setiap kali kamu memakai \hl{\mbox{\texttt{\&}}}, itu menambah satu tingkat pointer pada tipenya.

{\def\LTcaptype{none} % do not increment counter
\begin{longtable}[]{@{}lcl@{}}
\toprule\noalign{}
Jika kamu punya & Lalu kamu jalankan & Tipe hasilnya \\
\midrule\noalign{}
\endhead
\bottomrule\noalign{}
\endlastfoot
\hl{\mbox{\texttt{int\ x}}} & \hl{\mbox{\texttt{\&x}}} & \hl{\mbox{\texttt{int\ *}}} \\
\hl{\mbox{\texttt{int\ *x}}} & \hl{\mbox{\texttt{\&x}}} & \hl{\mbox{\texttt{int\ **}}} \\
\hl{\mbox{\texttt{int\ **x}}} & \hl{\mbox{\texttt{\&x}}} & \hl{\mbox{\texttt{int\ ***}}} \\
\hl{\mbox{\texttt{int\ ***x}}} & \hl{\mbox{\texttt{\&x}}} & \hl{\mbox{\texttt{int\ ****}}} \\
\end{longtable}
}

Dan setiap kali kamu memakai dereference (\hl{\mbox{\texttt{*}}}), ia melakukan kebalikannya:

{\def\LTcaptype{none} % do not increment counter
\begin{longtable}[]{@{}lcl@{}}
\toprule\noalign{}
Jika kamu punya & Lalu kamu jalankan & Tipe hasilnya \\
\midrule\noalign{}
\endhead
\bottomrule\noalign{}
\endlastfoot
\hl{\mbox{\texttt{int\ ****x}}} & \hl{\mbox{\texttt{*x}}} & \hl{\mbox{\texttt{int\ ***}}} \\
\hl{\mbox{\texttt{int\ ***x}}} & \hl{\mbox{\texttt{*x}}} & \hl{\mbox{\texttt{int\ **}}} \\
\hl{\mbox{\texttt{int\ **x}}} & \hl{\mbox{\texttt{*x}}} & \hl{\mbox{\texttt{int\ *}}} \\
\hl{\mbox{\texttt{int\ *x}}} & \hl{\mbox{\texttt{*x}}} & \hl{\mbox{\texttt{int}}} \\
\end{longtable}
}

Perhatikan bahwa kamu bisa memakai beberapa \hl{\mbox{\texttt{*}}} berturut-turut untuk men-dereference dengan cepat, seperti yang kita lihat pada kode contoh dengan \hl{\mbox{\texttt{**q}}} di atas. Masing-masing mengupas satu tingkat indirection.

{\def\LTcaptype{none} % do not increment counter
\begin{longtable}[]{@{}lcl@{}}
\toprule\noalign{}
Jika kamu punya & Lalu kamu jalankan & Tipe hasilnya \\
\midrule\noalign{}
\endhead
\bottomrule\noalign{}
\endlastfoot
\hl{\mbox{\texttt{int\ ****x}}} & \hl{\mbox{\texttt{***x}}} & \hl{\mbox{\texttt{int\ *}}} \\
\hl{\mbox{\texttt{int\ ***x}}} & \hl{\mbox{\texttt{**x}}} & \hl{\mbox{\texttt{int\ *}}} \\
\hl{\mbox{\texttt{int\ **x}}} & \hl{\mbox{\texttt{**x}}} & \hl{\mbox{\texttt{int}}} \\
\end{longtable}
}

Secara umum, \hl{\mbox{\texttt{\&*E\ ==\ E}}}\footnote{Bahkan kalau \hl{\mbox{\texttt{E}}} adalah \hl{\mbox{\texttt{NULL}}}, ternyata, anehnya.}. Dereference ``membatalkan'' address-of.

Tapi \hl{\mbox{\texttt{\&}}} tidak bekerja dengan cara yang sama---kamu hanya bisa melakukannya satu per satu, dan harus menyimpan hasilnya di variabel perantara:

\begin{tcolorbox}[enhanced jigsaw,left=3pt,breakable,colback=white!97.2549!black,boxrule=0.2mm]

\begin{Shaded}
\begin{Highlighting}[]
\DataTypeTok{int}\NormalTok{ x }\OperatorTok{=} \DecValTok{3490}\OperatorTok{;}     \CommentTok{// Tipe: int}
\DataTypeTok{int} \OperatorTok{*}\NormalTok{p }\OperatorTok{=} \OperatorTok{\&}\NormalTok{x}\OperatorTok{;}      \CommentTok{// Tipe: int *, pointer ke int}
\DataTypeTok{int} \OperatorTok{**}\NormalTok{q }\OperatorTok{=} \OperatorTok{\&}\NormalTok{p}\OperatorTok{;}     \CommentTok{// Tipe: int **, pointer ke pointer ke int}
\DataTypeTok{int} \OperatorTok{***}\NormalTok{r }\OperatorTok{=} \OperatorTok{\&}\NormalTok{q}\OperatorTok{;}    \CommentTok{// Tipe: int ***, pointer ke pointer ke pointer ke int}
\DataTypeTok{int} \OperatorTok{****}\NormalTok{s }\OperatorTok{=} \OperatorTok{\&}\NormalTok{r}\OperatorTok{;}   \CommentTok{// Tipe: int ****, kamu paham maksudnya}
\DataTypeTok{int} \OperatorTok{*****}\NormalTok{t }\OperatorTok{=} \OperatorTok{\&}\NormalTok{s}\OperatorTok{;}  \CommentTok{// Tipe: int *****}
\end{Highlighting}
\end{Shaded}

\end{tcolorbox}

\index{Pointers!to pointers|)}

\subsection{\texorpdfstring{Pointer Pointer dan \hl{\mbox{\texttt{const}}}}{Pointer Pointer dan const}}\label{pointer-pointer-dan-const}

\index{Pointers!to pointers, const@to pointers, \texttt{const}|(}

Kalau kamu ingat, mendeklarasikan pointer seperti ini:

\begin{tcolorbox}[enhanced jigsaw,left=3pt,breakable,colback=white!97.2549!black,boxrule=0.2mm]

\begin{Shaded}
\begin{Highlighting}[]
\DataTypeTok{int} \OperatorTok{*}\DataTypeTok{const}\NormalTok{ p}\OperatorTok{;}
\end{Highlighting}
\end{Shaded}

\end{tcolorbox}

berarti kamu tidak bisa mengubah \hl{\mbox{\texttt{p}}}. Mencoba \hl{\mbox{\texttt{p++}}} akan memberi error saat kompilasi.

Tapi bagaimana ini bekerja dengan \hl{\mbox{\texttt{int\ **}}} atau \hl{\mbox{\texttt{int\ ***}}}? Di mana \hl{\mbox{\texttt{const}}} ditaruh, dan apa artinya?

Kita mulai dari bagian sederhana. \hl{\mbox{\texttt{const}}} yang tepat di sebelah nama variabel merujuk ke variabel itu. Jadi kalau kamu ingin \hl{\mbox{\texttt{int***}}} yang tidak bisa diubah, kamu bisa menulis:

\begin{tcolorbox}[enhanced jigsaw,left=3pt,breakable,colback=white!97.2549!black,boxrule=0.2mm]

\begin{Shaded}
\begin{Highlighting}[]
\DataTypeTok{int} \OperatorTok{***}\DataTypeTok{const}\NormalTok{ p}\OperatorTok{;}

\NormalTok{p}\OperatorTok{++;}  \CommentTok{// Tidak diperbolehkan}
\end{Highlighting}
\end{Shaded}

\end{tcolorbox}

Tapi di sinilah semuanya jadi agak aneh.

Bagaimana kalau kita punya situasi ini:

\begin{tcolorbox}[enhanced jigsaw,left=3pt,breakable,colback=white!97.2549!black,boxrule=0.2mm]

\begin{Shaded}
\begin{Highlighting}[numbers=left,,]
\DataTypeTok{int}\NormalTok{ main}\OperatorTok{(}\DataTypeTok{void}\OperatorTok{)}
\OperatorTok{\{}
    \DataTypeTok{int}\NormalTok{ x }\OperatorTok{=} \DecValTok{3490}\OperatorTok{;}
    \DataTypeTok{int} \OperatorTok{*}\DataTypeTok{const}\NormalTok{ p }\OperatorTok{=} \OperatorTok{\&}\NormalTok{x}\OperatorTok{;}
    \DataTypeTok{int} \OperatorTok{**}\NormalTok{q }\OperatorTok{=} \OperatorTok{\&}\NormalTok{p}\OperatorTok{;}
\OperatorTok{\}}
\end{Highlighting}
\end{Shaded}

\end{tcolorbox}

Waktu kubangun, aku mendapat peringatan ini:

\begin{tcolorbox}[enhanced jigsaw,left=3pt,breakable,colback=white!97.2549!black,boxrule=0.2mm]

\begin{Shaded}
\begin{Highlighting}[]
\NormalTok{warning: initialization discards ‘const’ qualifier from pointer target type}
\NormalTok{    7 |     int **q = \&p;}
\NormalTok{      |               \^{}}
\end{Highlighting}
\end{Shaded}

\end{tcolorbox}

Apa yang terjadi? Compiler memberi tahu kita di sini bahwa ada variabel yang \hl{\mbox{\texttt{const}}}, dan kita menugaskan nilainya ke variabel lain yang tidak \hl{\mbox{\texttt{const}}} dengan cara yang sama. Sifat ``\hl{\mbox{\texttt{const}}}'' dibuang, yang mungkin bukan yang kita inginkan.

Tipe \hl{\mbox{\texttt{p}}} adalah \hl{\mbox{\texttt{int\ *const\ p}}}, jadi \hl{\mbox{\texttt{\&p}}} bertipe \hl{\mbox{\texttt{int\ *const\ *}}}. Dan kita mencoba menugaskannya ke \hl{\mbox{\texttt{q}}}.

Tapi \hl{\mbox{\texttt{q}}} adalah \hl{\mbox{\texttt{int\ **}}}! Tipe dengan sifat \hl{\mbox{\texttt{const}}} yang berbeda pada \hl{\mbox{\texttt{*}}} pertama! Jadi kita mendapat peringatan bahwa \hl{\mbox{\texttt{const}}} pada \hl{\mbox{\texttt{int\ *const\ *}}} milik \hl{\mbox{\texttt{p}}} diabaikan dan dibuang.

Kita bisa memperbaikinya dengan memastikan tipe \hl{\mbox{\texttt{q}}} setidaknya sama \hl{\mbox{\texttt{const}}}-nya dengan \hl{\mbox{\texttt{p}}}.

\begin{tcolorbox}[enhanced jigsaw,left=3pt,breakable,colback=white!97.2549!black,boxrule=0.2mm]

\begin{Shaded}
\begin{Highlighting}[]
\DataTypeTok{int}\NormalTok{ x }\OperatorTok{=} \DecValTok{3490}\OperatorTok{;}
\DataTypeTok{int} \OperatorTok{*}\DataTypeTok{const}\NormalTok{ p }\OperatorTok{=} \OperatorTok{\&}\NormalTok{x}\OperatorTok{;}
\DataTypeTok{int} \OperatorTok{*}\DataTypeTok{const} \OperatorTok{*}\NormalTok{q }\OperatorTok{=} \OperatorTok{\&}\NormalTok{p}\OperatorTok{;}
\end{Highlighting}
\end{Shaded}

\end{tcolorbox}

Dan sekarang kita senang.

Kita bahkan bisa membuat \hl{\mbox{\texttt{q}}} lebih \hl{\mbox{\texttt{const}}} lagi. Seperti di atas, kita berkata, ``\hl{\mbox{\texttt{q}}} itu sendiri tidak \hl{\mbox{\texttt{const}}}, tapi yang ditunjuknya \hl{\mbox{\texttt{const}}}.'' Tapi kita bisa membuat keduanya \hl{\mbox{\texttt{const}}}:

\begin{tcolorbox}[enhanced jigsaw,left=3pt,breakable,colback=white!97.2549!black,boxrule=0.2mm]

\begin{Shaded}
\begin{Highlighting}[]
\DataTypeTok{int}\NormalTok{ x }\OperatorTok{=} \DecValTok{3490}\OperatorTok{;}
\DataTypeTok{int} \OperatorTok{*}\DataTypeTok{const}\NormalTok{ p }\OperatorTok{=} \OperatorTok{\&}\NormalTok{x}\OperatorTok{;}
\DataTypeTok{int} \OperatorTok{*}\DataTypeTok{const} \OperatorTok{*}\DataTypeTok{const}\NormalTok{ q }\OperatorTok{=} \OperatorTok{\&}\NormalTok{p}\OperatorTok{;}  \CommentTok{// Lebih const!}
\end{Highlighting}
\end{Shaded}

\end{tcolorbox}

Dan itu juga berfungsi. Sekarang kita tidak bisa mengubah \hl{\mbox{\texttt{q}}}, maupun pointer yang ditunjuk \hl{\mbox{\texttt{q}}}.

\index{Pointers!to pointers, const@to pointers, \texttt{const}|)}

\section{Nilai Multibyte}\label{multibyte-values}

\index{Pointers!to multibyte values|(}

Kita sudah menyiratkan ini di berbagai tempat sebelumnya, tapi jelas tidak setiap nilai bisa disimpan dalam satu byte memori. Banyak hal memakai beberapa byte memori (dengan asumsi bukan \hl{\mbox{\texttt{char}}}). Kamu bisa mengetahui jumlah byte dengan \hl{\mbox{\texttt{sizeof}}}. Dan kamu bisa mengetahui alamat memori mana yang merupakan byte \emph{pertama} dari object dengan operator \hl{\mbox{\texttt{\&}}} standar, yang selalu mengembalikan alamat byte pertama.

Dan ini fakta menarik lainnya! Kalau kamu menelusuri byte-byte dari object mana pun, kamu mendapat \emph{object representation}-nya. Dua hal dengan object representation yang sama di memori adalah sama.

Kalau kamu ingin menelusuri object representation, lakukan dengan pointer ke \hl{\mbox{\texttt{unsigned\ char}}}.

Mari kita buat versi kita sendiri dari \href{https://beej.us/guide/bgclr/html/split/stringref.html\#man-memcpy}{\hl{\mbox{\texttt{memcpy()}}}}\footnote{https://beej.us/guide/bgclr/html/split/stringref.html\#man-memcpy} yang melakukan persis itu:

\begin{tcolorbox}[enhanced jigsaw,left=3pt,breakable,colback=white!97.2549!black,boxrule=0.2mm]

\begin{Shaded}
\begin{Highlighting}[]
\DataTypeTok{void} \OperatorTok{*}\NormalTok{my\_memcpy}\OperatorTok{(}\DataTypeTok{void} \OperatorTok{*}\NormalTok{dest}\OperatorTok{,} \DataTypeTok{const} \DataTypeTok{void} \OperatorTok{*}\NormalTok{src}\OperatorTok{,} \DataTypeTok{size\_t}\NormalTok{ n}\OperatorTok{)}
\OperatorTok{\{}
    \CommentTok{// Buat variabel lokal untuk src dan dest, tapi bertipe unsigned char}

    \DataTypeTok{const} \DataTypeTok{unsigned} \DataTypeTok{char} \OperatorTok{*}\NormalTok{s }\OperatorTok{=}\NormalTok{ src}\OperatorTok{;}
    \DataTypeTok{unsigned} \DataTypeTok{char} \OperatorTok{*}\NormalTok{d }\OperatorTok{=}\NormalTok{ dest}\OperatorTok{;}

    \ControlFlowTok{while} \OperatorTok{(}\NormalTok{n}\OperatorTok{{-}{-}} \OperatorTok{\textgreater{}} \DecValTok{0}\OperatorTok{)}   \CommentTok{// Untuk jumlah byte yang diberikan}
        \OperatorTok{*}\NormalTok{d}\OperatorTok{++} \OperatorTok{=} \OperatorTok{*}\NormalTok{s}\OperatorTok{++;}  \CommentTok{// Salin byte sumber ke byte tujuan}

    \CommentTok{// Kebanyakan fungsi salin mengembalikan pointer ke dest demi}
    \CommentTok{// kemudahan caller}

    \ControlFlowTok{return}\NormalTok{ dest}\OperatorTok{;}
\OperatorTok{\}}
\end{Highlighting}
\end{Shaded}

\end{tcolorbox}

(Ada beberapa contoh bagus post-increment dan post-decrement di situ yang bisa kamu pelajari juga.)

Penting dicatat bahwa versi di atas kemungkinan kurang efisien dibanding yang disertakan di sistemmu.

Tapi kamu bisa mengoper pointer ke apa pun ke dalamnya, dan ia akan menyalin object itu. Bisa \hl{\mbox{\texttt{int*}}}, \hl{\mbox{\texttt{struct\ animal*}}}, atau apa saja.

Mari kita coba contoh lain yang mencetak byte object representation dari sebuah \hl{\mbox{\texttt{struct}}} supaya kita bisa melihat apakah ada padding di dalamnya dan nilainya apa\footnote{C compiler-mu tidak wajib menambahkan padding byte, dan nilai padding byte yang ditambahkan itu tak tentu.}.

\begin{tcolorbox}[enhanced jigsaw,left=3pt,breakable,colback=white!97.2549!black,boxrule=0.2mm]

\begin{Shaded}
\begin{Highlighting}[numbers=left,,]
\PreprocessorTok{\#include }\ImportTok{\textless{}stdio.h\textgreater{}}

\KeywordTok{struct}\NormalTok{ foo }\OperatorTok{\{}
    \DataTypeTok{char}\NormalTok{ a}\OperatorTok{;}
    \DataTypeTok{int}\NormalTok{ b}\OperatorTok{;}
\OperatorTok{\};}

\DataTypeTok{int}\NormalTok{ main}\OperatorTok{(}\DataTypeTok{void}\OperatorTok{)}
\OperatorTok{\{}
    \KeywordTok{struct}\NormalTok{ foo x }\OperatorTok{=} \OperatorTok{\{}\BaseNTok{0x12}\OperatorTok{,} \BaseNTok{0x12345678}\OperatorTok{\};}
    \DataTypeTok{unsigned} \DataTypeTok{char} \OperatorTok{*}\NormalTok{p }\OperatorTok{=} \OperatorTok{(}\DataTypeTok{unsigned} \DataTypeTok{char} \OperatorTok{*)\&}\NormalTok{x}\OperatorTok{;}

    \ControlFlowTok{for} \OperatorTok{(}\DataTypeTok{size\_t}\NormalTok{ i }\OperatorTok{=} \DecValTok{0}\OperatorTok{;}\NormalTok{ i }\OperatorTok{\textless{}} \KeywordTok{sizeof}\NormalTok{ x}\OperatorTok{;}\NormalTok{ i}\OperatorTok{++)} \OperatorTok{\{}
\NormalTok{        printf}\OperatorTok{(}\StringTok{"}\SpecialCharTok{\%02X\textbackslash{}n}\StringTok{"}\OperatorTok{,}\NormalTok{ p}\OperatorTok{[}\NormalTok{i}\OperatorTok{]);}
    \OperatorTok{\}}
\OperatorTok{\}}
\end{Highlighting}
\end{Shaded}

\end{tcolorbox}

Di situ ada \hl{\mbox{\texttt{struct\ foo}}} yang disusun sedemikian rupa agar compiler terdorong menyisipkan padding byte (walaupun tidak wajib). Lalu kita ambil \hl{\mbox{\texttt{unsigned\ char\ *}}} ke byte pertama variabel \hl{\mbox{\texttt{struct\ foo}}} \hl{\mbox{\texttt{x}}}.

Dari sana, yang kita perlu tahu hanya \hl{\mbox{\texttt{sizeof\ x}}} dan kita bisa menelusuri sebanyak itu byte, mencetak nilainya (dalam hex agar mudah).

Menjalankannya memberi banyak angka sebagai output. Aku sudah memberinya anotasi di bawah untuk menunjukkan di mana nilai-nilai itu disimpan:

\begin{tcolorbox}[enhanced jigsaw,left=3pt,breakable,colback=white!97.2549!black,boxrule=0.2mm]

\begin{Shaded}
\begin{Highlighting}[]
\NormalTok{12  | x.a == 0x12}

\NormalTok{AB  |}
\NormalTok{BF  | padding bytes with "random" value}
\NormalTok{26  |}

\NormalTok{78  | x.b == 0x12345678}
\NormalTok{56  |}
\NormalTok{34  |}
\NormalTok{12  |}
\end{Highlighting}
\end{Shaded}

\end{tcolorbox}

Di semua sistem, \hl{\mbox{\texttt{sizeof(char)}}} adalah 1, dan kita lihat byte pertama di bagian atas output menyimpan nilai \hl{\mbox{\texttt{0x12}}} yang kita simpan di sana.

Lalu ada beberapa padding byte---di sistemku, nilainya berubah dari satu eksekusi ke eksekusi berikutnya.

Terakhir, di sistemku, \hl{\mbox{\texttt{sizeof(int)}}} adalah 4, dan kita bisa melihat 4 byte itu di bagian akhir. Perhatikan bahwa byte-nya sama dengan nilai hex \hl{\mbox{\texttt{0x12345678}}}, tapi anehnya dalam urutan terbalik\footnote{Ini bervariasi tergantung arsitektur, tapi sistemku \emph{little endian}, artinya byte paling tidak signifikan dari angka disimpan lebih dulu. Sistem \emph{big endian} akan punya \hl{\mbox{\texttt{12}}} di depan dan \hl{\mbox{\texttt{78}}} di belakang. Tapi spesifikasi tidak menentukan apa pun soal representasi ini.}.

Jadi itu sedikit intipan ke bawah kap tentang byte-byte dari entitas yang lebih kompleks di memori.

\index{Pointers!to multibyte values|)}

\section{\texorpdfstring{Pointer \hl{\mbox{\texttt{NULL}}} dan Nol}{Pointer NULL dan Nol}}\label{pointer-null-dan-nol}

\index{NULL pointer@\texttt{NULL} pointer!zero equivalence|(}

Hal-hal berikut bisa dipakai secara bergantian:

\begin{itemize}
\tightlist
\item
  \hl{\mbox{\texttt{NULL}}}
\item
  \hl{\mbox{\texttt{0}}}
\item
  \hl{\mbox{\texttt{\textquotesingle{}\textbackslash{}0\textquotesingle{}}}}
\item
  \hl{\mbox{\texttt{(void\ *)0}}}
\end{itemize}

Secara pribadi, aku selalu memakai \hl{\mbox{\texttt{NULL}}} kalau maksudnya \hl{\mbox{\texttt{NULL}}}, tapi kamu mungkin sesekali melihat varian lain. Walaupun \hl{\mbox{\texttt{\textquotesingle{}\textbackslash{}0\textquotesingle{}}}} (byte dengan semua bit bernilai nol) juga akan dibandingkan sebagai sama, membandingkannya dengan pointer itu \emph{aneh}; kamu seharusnya membandingkan \hl{\mbox{\texttt{NULL}}} dengan pointer. (Tentu saja, sering kali dalam pemrosesan string, kamu membandingkan \emph{yang ditunjuk pointer} dengan \hl{\mbox{\texttt{\textquotesingle{}\textbackslash{}0\textquotesingle{}}}}, dan itu benar.)

\hl{\mbox{\texttt{0}}} disebut \emph{null pointer constant}, dan, ketika dibandingkan atau ditugaskan ke pointer lain, ia dikonversi menjadi null pointer dengan tipe yang sama.

\index{NULL pointer@\texttt{NULL} pointer!zero equivalence|)}

\section{Pointer sebagai Integer}\label{pointer-sebagai-integer}

\index{Pointers!as integers|(}

Kamu bisa meng-cast pointer ke integer dan sebaliknya (karena pointer hanyalah indeks ke memori), tapi kemungkinan kamu hanya perlu melakukannya kalau mengerjakan hal-hal hardware level rendah. Hasil dari akal-akalan semacam itu ditentukan oleh implementasi, jadi tidak portabel. Dan \emph{hal aneh} bisa terjadi.

Meski begitu, C memberi satu jaminan: kamu bisa mengonversi pointer ke tipe \hl{\mbox{\texttt{uintptr\_t}}} dan kamu bisa mengonversinya kembali ke pointer tanpa kehilangan data.

\hl{\mbox{\texttt{uintptr\_t}}} didefinisikan di \hl{\mbox{\texttt{\textless{}stdint.h\textgreater{}}}}\footnote{Ini fitur opsional, jadi mungkin tidak ada---tapi kemungkinan besar ada.}.

Selain itu, kalau kamu ingin yang signed, kamu bisa memakai \hl{\mbox{\texttt{intptr\_t}}} dengan efek yang sama.

\index{Pointers!as integers|)}

\section{Meng-cast Pointer ke Pointer Lain}\label{meng-cast-pointer-ke-pointer-lain}

\index{Pointers!casting|(}

Hanya ada satu konversi pointer yang aman:

\begin{enumerate}
\def\labelenumi{\arabic{enumi}.}
\tightlist
\item
  Mengonversi ke \hl{\mbox{\texttt{intptr\_t}}} atau \hl{\mbox{\texttt{uintptr\_t}}}.
\item
  Mengonversi ke dan dari \hl{\mbox{\texttt{void*}}}.
\end{enumerate}

DUA! Dua konversi pointer yang aman.

\begin{enumerate}
\def\labelenumi{\arabic{enumi}.}
\setcounter{enumi}{2}
\tightlist
\item
  Mengonversi ke dan dari \hl{\mbox{\texttt{char*}}} (atau \hl{\mbox{\texttt{signed\ char*}}}/\hl{\mbox{\texttt{unsigned\ char*}}}).
\end{enumerate}

TIGA! Tiga konversi yang aman!

\begin{enumerate}
\def\labelenumi{\arabic{enumi}.}
\setcounter{enumi}{3}
\tightlist
\item
  Mengonversi ke dan dari pointer ke sebuah \hl{\mbox{\texttt{struct}}} dan pointer ke anggota pertamanya, dan sebaliknya.
\end{enumerate}

EMPAT! Empat konversi yang aman!

Kalau kamu meng-cast ke pointer bertipe lain lalu mengakses object yang ditunjuknya, perilakunya undefined karena sesuatu yang disebut \emph{strict aliasing}.

\emph{Aliasing} biasa merujuk ke kemampuan untuk punya lebih dari satu cara mengakses object yang sama. Titik-titik akses itu adalah alias satu sama lain.

\emph{Strict aliasing} menyatakan kamu hanya boleh mengakses sebuah object lewat pointer ke \emph{tipe yang kompatibel} dengan object itu.

Misalnya, ini jelas diperbolehkan:

\begin{tcolorbox}[enhanced jigsaw,left=3pt,breakable,colback=white!97.2549!black,boxrule=0.2mm]

\begin{Shaded}
\begin{Highlighting}[]
\DataTypeTok{int}\NormalTok{ a }\OperatorTok{=} \DecValTok{1}\OperatorTok{;}
\DataTypeTok{int} \OperatorTok{*}\NormalTok{p }\OperatorTok{=} \OperatorTok{\&}\NormalTok{a}\OperatorTok{;}
\end{Highlighting}
\end{Shaded}

\end{tcolorbox}

\hl{\mbox{\texttt{p}}} adalah pointer ke \hl{\mbox{\texttt{int}}}, dan ia menunjuk ke tipe yang kompatibel---yaitu \hl{\mbox{\texttt{int}}}---jadi aman.

Tapi yang berikut tidak bagus karena \hl{\mbox{\texttt{int}}} dan \hl{\mbox{\texttt{float}}} bukan tipe yang kompatibel:

\begin{tcolorbox}[enhanced jigsaw,left=3pt,breakable,colback=white!97.2549!black,boxrule=0.2mm]

\begin{Shaded}
\begin{Highlighting}[]
\DataTypeTok{int}\NormalTok{ a }\OperatorTok{=} \DecValTok{1}\OperatorTok{;}
\DataTypeTok{float} \OperatorTok{*}\NormalTok{p }\OperatorTok{=} \OperatorTok{(}\DataTypeTok{float} \OperatorTok{*)\&}\NormalTok{a}\OperatorTok{;}
\end{Highlighting}
\end{Shaded}

\end{tcolorbox}

Berikut program demo yang melakukan aliasing. Ia mengambil variabel \hl{\mbox{\texttt{v}}} bertipe \hl{\mbox{\texttt{int32\_t}}} dan meng-alias-kannya ke pointer ke \hl{\mbox{\texttt{struct\ words}}}. \hl{\mbox{\texttt{struct}}} itu berisi dua \hl{\mbox{\texttt{int16\_t}}}. Tipe-tipe ini tidak kompatibel, jadi kita melanggar aturan strict aliasing. Compiler akan menganggap kedua pointer ini tidak pernah menunjuk ke object yang sama\ldots{} tapi kita membuatnya menunjuk. Nakal sekali kita ini.

Mari kita lihat apakah kita bisa merusak sesuatu.

\begin{tcolorbox}[enhanced jigsaw,left=3pt,breakable,colback=white!97.2549!black,boxrule=0.2mm]

\begin{Shaded}
\begin{Highlighting}[numbers=left,,]
\PreprocessorTok{\#include }\ImportTok{\textless{}stdio.h\textgreater{}}
\PreprocessorTok{\#include }\ImportTok{\textless{}stdint.h\textgreater{}}

\KeywordTok{struct}\NormalTok{ words }\OperatorTok{\{}
    \DataTypeTok{int16\_t}\NormalTok{ v}\OperatorTok{[}\DecValTok{2}\OperatorTok{];}
\OperatorTok{\};}

\DataTypeTok{void}\NormalTok{ fun}\OperatorTok{(}\DataTypeTok{int32\_t} \OperatorTok{*}\NormalTok{pv}\OperatorTok{,} \KeywordTok{struct}\NormalTok{ words }\OperatorTok{*}\NormalTok{pw}\OperatorTok{)}
\OperatorTok{\{}
    \ControlFlowTok{for} \OperatorTok{(}\DataTypeTok{int}\NormalTok{ i }\OperatorTok{=} \DecValTok{0}\OperatorTok{;}\NormalTok{ i }\OperatorTok{\textless{}} \DecValTok{5}\OperatorTok{;}\NormalTok{ i}\OperatorTok{++)} \OperatorTok{\{}
        \OperatorTok{(*}\NormalTok{pv}\OperatorTok{)++;}

        \CommentTok{// Cetak nilai 32{-}bit dan nilai{-}nilai 16{-}bit:}

\NormalTok{        printf}\OperatorTok{(}\StringTok{"}\SpecialCharTok{\%x}\StringTok{, }\SpecialCharTok{\%x}\StringTok{{-}}\SpecialCharTok{\%x\textbackslash{}n}\StringTok{"}\OperatorTok{,} \OperatorTok{*}\NormalTok{pv}\OperatorTok{,}\NormalTok{ pw}\OperatorTok{{-}\textgreater{}}\NormalTok{v}\OperatorTok{[}\DecValTok{1}\OperatorTok{],}\NormalTok{ pw}\OperatorTok{{-}\textgreater{}}\NormalTok{v}\OperatorTok{[}\DecValTok{0}\OperatorTok{]);}
    \OperatorTok{\}}
\OperatorTok{\}}

\DataTypeTok{int}\NormalTok{ main}\OperatorTok{(}\DataTypeTok{void}\OperatorTok{)}
\OperatorTok{\{}
    \DataTypeTok{int32\_t}\NormalTok{ v }\OperatorTok{=} \BaseNTok{0x12345678}\OperatorTok{;}

    \KeywordTok{struct}\NormalTok{ words }\OperatorTok{*}\NormalTok{pw }\OperatorTok{=} \OperatorTok{(}\KeywordTok{struct}\NormalTok{ words }\OperatorTok{*)\&}\NormalTok{v}\OperatorTok{;}  \CommentTok{// Melanggar strict aliasing}

\NormalTok{    fun}\OperatorTok{(\&}\NormalTok{v}\OperatorTok{,}\NormalTok{ pw}\OperatorTok{);}
\OperatorTok{\}}
\end{Highlighting}
\end{Shaded}

\end{tcolorbox}

Lihat bagaimana aku mengoper kedua pointer yang tidak kompatibel ke \hl{\mbox{\texttt{fun()}}}? Satu bertipe \hl{\mbox{\texttt{int32\_t*}}} dan yang lain \hl{\mbox{\texttt{struct\ words*}}}.

Tapi keduanya menunjuk ke object yang sama: nilai 32-bit yang diinisialisasi dengan \hl{\mbox{\texttt{0x12345678}}}.

Jadi kalau kita lihat field-field di \hl{\mbox{\texttt{struct\ words}}}, kita seharusnya melihat dua setengah bagian 16-bit dari angka itu. Benar?

Dan di loop \hl{\mbox{\texttt{fun()}}}, kita meng-increment pointer ke \hl{\mbox{\texttt{int32\_t}}}. Itu saja. Tapi karena \hl{\mbox{\texttt{struct}}} menunjuk ke memori yang sama, ia juga seharusnya diperbarui ke nilai yang sama.

Mari kita jalankan dan dapatkan ini, dengan nilai 32-bit di kiri dan dua bagian 16-bit di kanan. Seharusnya cocok\footnote{Aku mencetak nilai 16-bit secara terbalik karena aku memakai mesin little-endian dan ini membuatnya lebih mudah dibaca di sini.}:

\begin{tcolorbox}[enhanced jigsaw,left=3pt,breakable,colback=white!97.2549!black,boxrule=0.2mm]

\begin{Shaded}
\begin{Highlighting}[]
\NormalTok{12345679, 1234{-}5679}
\NormalTok{1234567a, 1234{-}567a}
\NormalTok{1234567b, 1234{-}567b}
\NormalTok{1234567c, 1234{-}567c}
\NormalTok{1234567d, 1234{-}567d}
\end{Highlighting}
\end{Shaded}

\end{tcolorbox}

dan memang cocok\ldots{} \emph{SAMPAI BESOK!}

Mari kita coba mengompilasinya dengan GCC memakai \hl{\mbox{\texttt{-O3}}} dan \hl{\mbox{\texttt{-fstrict-aliasing}}}:

\begin{tcolorbox}[enhanced jigsaw,left=3pt,breakable,colback=white!97.2549!black,boxrule=0.2mm]

\begin{Shaded}
\begin{Highlighting}[]
\NormalTok{12345679, 1234{-}5678}
\NormalTok{1234567a, 1234{-}5679}
\NormalTok{1234567b, 1234{-}567a}
\NormalTok{1234567c, 1234{-}567b}
\NormalTok{1234567d, 1234{-}567c}
\end{Highlighting}
\end{Shaded}

\end{tcolorbox}

Selisih satu! Tapi mereka menunjuk ke memori yang sama! Bagaimana bisa? Jawabannya: melakukan aliasing memori seperti itu adalah undefined behavior. \emph{Apa pun mungkin terjadi}, hanya saja bukan dengan cara yang baik.

Kalau kodemu melanggar aturan strict aliasing, apakah ia berfungsi atau tidak bergantung pada bagaimana seseorang memutuskan mengompilasinya. Dan itu menyebalkan karena di luar kendalimu. Kecuali kamu semacam dewa mahakuasa.

Kemungkinan tidak, maaf.

GCC bisa dipaksa untuk tidak memakai aturan strict aliasing dengan \hl{\mbox{\texttt{-fno-strict-aliasing}}}. Mengompilasi program demo di atas dengan \hl{\mbox{\texttt{-O3}}} dan flag ini membuat outputnya sesuai harapan.

Terakhir, \emph{type punning} adalah memakai pointer bertipe berbeda untuk melihat data yang sama. Sebelum ada strict aliasing, hal seperti ini cukup umum:

\begin{tcolorbox}[enhanced jigsaw,left=3pt,breakable,colback=white!97.2549!black,boxrule=0.2mm]

\begin{Shaded}
\begin{Highlighting}[]
\DataTypeTok{int}\NormalTok{ a }\OperatorTok{=} \BaseNTok{0x12345678}\OperatorTok{;}
\DataTypeTok{short}\NormalTok{ b }\OperatorTok{=} \OperatorTok{*((}\DataTypeTok{short} \OperatorTok{*)\&}\NormalTok{a}\OperatorTok{);}   \CommentTok{// Melanggar strict aliasing}
\end{Highlighting}
\end{Shaded}

\end{tcolorbox}

Kalau kamu ingin melakukan type punning dengan (relatif) aman, lihat bagian tentang \hyperref[union-type-punning]{Union dan Type Punning}.

\index{Pointers!casting|)}

\section{Selisih Pointer}\label{ptr_differences}

\index{Pointers!subtracting|(}

Seperti yang kamu tahu dari bagian tentang aritmetika pointer, kamu bisa mengurangkan satu pointer dari pointer lain\footnote{Dengan asumsi keduanya menunjuk ke object array yang sama.} untuk mendapatkan selisihnya dalam jumlah elemen array.

Nah, \emph{tipe selisih itu} ditentukan oleh implementasi, jadi bisa berbeda dari satu sistem ke sistem lain.

\index{ptrdiff\_t type@\texttt{ptrdiff\_t} type|(}

Agar lebih portabel, kamu bisa menyimpan hasilnya dalam variabel bertipe \hl{\mbox{\texttt{ptrdiff\_t}}} yang didefinisikan di \hl{\mbox{\texttt{\textless{}stddef.h\textgreater{}}}}.

\begin{tcolorbox}[enhanced jigsaw,left=3pt,breakable,colback=white!97.2549!black,boxrule=0.2mm]

\begin{Shaded}
\begin{Highlighting}[]
\DataTypeTok{int}\NormalTok{ cats}\OperatorTok{[}\DecValTok{100}\OperatorTok{];}

\DataTypeTok{int} \OperatorTok{*}\NormalTok{f }\OperatorTok{=}\NormalTok{ cats }\OperatorTok{+} \DecValTok{20}\OperatorTok{;}
\DataTypeTok{int} \OperatorTok{*}\NormalTok{g }\OperatorTok{=}\NormalTok{ cats }\OperatorTok{+} \DecValTok{60}\OperatorTok{;}

\DataTypeTok{ptrdiff\_t}\NormalTok{ d }\OperatorTok{=}\NormalTok{ g }\OperatorTok{{-}}\NormalTok{ f}\OperatorTok{;}  \CommentTok{// selisihnya 40}
\end{Highlighting}
\end{Shaded}

\end{tcolorbox}

\index{ptrdiff\_t type@\texttt{ptrdiff\_t} type!printing|(} Dan kamu bisa mencetaknya dengan menambahkan awalan \hl{\mbox{\texttt{t}}} pada format specifier integer:

\begin{tcolorbox}[enhanced jigsaw,left=3pt,breakable,colback=white!97.2549!black,boxrule=0.2mm]

\begin{Shaded}
\begin{Highlighting}[]
\NormalTok{printf}\OperatorTok{(}\StringTok{"}\SpecialCharTok{\%td\textbackslash{}n}\StringTok{"}\OperatorTok{,}\NormalTok{ d}\OperatorTok{);}  \CommentTok{// Cetak desimal: 40}
\NormalTok{printf}\OperatorTok{(}\StringTok{"}\SpecialCharTok{\%tX\textbackslash{}n}\StringTok{"}\OperatorTok{,}\NormalTok{ d}\OperatorTok{);}  \CommentTok{// Cetak hex:     28}
\end{Highlighting}
\end{Shaded}

\end{tcolorbox}

\index{ptrdiff\_t type@\texttt{ptrdiff\_t} type!printing|)} \index{ptrdiff\_t type@\texttt{ptrdiff\_t} type|)} \index{Pointers!subtracting|)}

\section{Pointer ke Fungsi}\label{pointer-ke-fungsi}

\index{Pointers!to functions|(}

Fungsi hanyalah kumpulan instruksi mesin di memori, jadi tidak ada alasan kita tidak bisa mengambil pointer ke instruksi pertama fungsi itu.

Lalu memanggilnya.

Ini bisa berguna untuk mengoper pointer ke sebuah fungsi ke fungsi lain sebagai argumen. Lalu fungsi kedua itu bisa memanggil apa pun yang dioper.

Bagian yang rumit, meski begitu, adalah C perlu tahu tipe variabel yang berupa pointer ke fungsi itu.

Dan ia benar-benar ingin tahu semua detailnya.

Seperti ``ini adalah pointer ke fungsi yang menerima dua argumen \hl{\mbox{\texttt{int}}} dan mengembalikan \hl{\mbox{\texttt{void}}}''.

Bagaimana menuliskan semua itu agar bisa mendeklarasikan variabel?

Ternyata bentuknya sangat mirip function prototype, hanya dengan beberapa tanda kurung tambahan:

\begin{tcolorbox}[enhanced jigsaw,left=3pt,breakable,colback=white!97.2549!black,boxrule=0.2mm]

\begin{Shaded}
\begin{Highlighting}[]
\CommentTok{// Deklarasikan p sebagai pointer ke fungsi.}
\CommentTok{// Fungsi ini mengembalikan float, dan menerima dua int sebagai argumen.}

\DataTypeTok{float} \OperatorTok{(*}\NormalTok{p}\OperatorTok{)(}\DataTypeTok{int}\OperatorTok{,} \DataTypeTok{int}\OperatorTok{);}
\end{Highlighting}
\end{Shaded}

\end{tcolorbox}

Perhatikan juga bahwa kamu tidak harus memberi nama pada parameter. Tapi boleh saja kalau mau; namanya hanya diabaikan.

\begin{tcolorbox}[enhanced jigsaw,left=3pt,breakable,colback=white!97.2549!black,boxrule=0.2mm]

\begin{Shaded}
\begin{Highlighting}[]
\CommentTok{// Deklarasikan p sebagai pointer ke fungsi.}
\CommentTok{// Fungsi ini mengembalikan float, dan menerima dua int sebagai argumen.}

\DataTypeTok{float} \OperatorTok{(*}\NormalTok{p}\OperatorTok{)(}\DataTypeTok{int}\NormalTok{ a}\OperatorTok{,} \DataTypeTok{int}\NormalTok{ b}\OperatorTok{);}
\end{Highlighting}
\end{Shaded}

\end{tcolorbox}

Jadi sekarang setelah kita tahu cara mendeklarasikan variabel, bagaimana kita tahu apa yang ditugaskan ke dalamnya? Bagaimana mendapatkan alamat sebuah fungsi?

Ternyata ada jalan pintas seperti saat mendapatkan pointer ke array: kamu cukup menyebut nama fungsi tanpa kurung. (Kamu boleh menaruh \hl{\mbox{\texttt{\&}}} di depannya kalau mau, tapi itu tidak perlu dan tidak idiomatis.)

Setelah punya pointer ke fungsi, kamu bisa memanggilnya hanya dengan menambahkan tanda kurung dan daftar argumen.

Kita coba contoh sederhana yang secara efektif membuat alias untuk sebuah fungsi dengan mengarahkan pointer ke sana. Lalu kita memanggilnya.

Kode ini mencetak \hl{\mbox{\texttt{3490}}}:

\begin{tcolorbox}[enhanced jigsaw,left=3pt,breakable,colback=white!97.2549!black,boxrule=0.2mm]

\begin{Shaded}
\begin{Highlighting}[numbers=left,,]
\PreprocessorTok{\#include }\ImportTok{\textless{}stdio.h\textgreater{}}

\DataTypeTok{void}\NormalTok{ print\_int}\OperatorTok{(}\DataTypeTok{int}\NormalTok{ n}\OperatorTok{)}
\OperatorTok{\{}
\NormalTok{    printf}\OperatorTok{(}\StringTok{"}\SpecialCharTok{\%d\textbackslash{}n}\StringTok{"}\OperatorTok{,}\NormalTok{ n}\OperatorTok{);}
\OperatorTok{\}}

\DataTypeTok{int}\NormalTok{ main}\OperatorTok{(}\DataTypeTok{void}\OperatorTok{)}
\OperatorTok{\{}
    \CommentTok{// Tugaskan p agar menunjuk ke print\_int:}

    \DataTypeTok{void} \OperatorTok{(*}\NormalTok{p}\OperatorTok{)(}\DataTypeTok{int}\OperatorTok{)} \OperatorTok{=}\NormalTok{ print\_int}\OperatorTok{;}

\NormalTok{    p}\OperatorTok{(}\DecValTok{3490}\OperatorTok{);}          \CommentTok{// Panggil print\_int lewat pointer}
\OperatorTok{\}}
\end{Highlighting}
\end{Shaded}

\end{tcolorbox}

Perhatikan bagaimana tipe \hl{\mbox{\texttt{p}}} merepresentasikan tipe nilai kembalian dan parameter \hl{\mbox{\texttt{print\_int}}}. Harus begitu, kalau tidak C akan protes soal tipe pointer yang tidak kompatibel.

Satu contoh lagi di sini menunjukkan bagaimana kita bisa mengoper pointer ke fungsi sebagai argumen ke fungsi lain.

Kita akan menulis fungsi yang menerima beberapa argumen integer, ditambah pointer ke fungsi yang beroperasi pada kedua argumen itu. Lalu ia mencetak hasilnya.

\begin{tcolorbox}[enhanced jigsaw,left=3pt,breakable,colback=white!97.2549!black,boxrule=0.2mm]

\begin{Shaded}
\begin{Highlighting}[numbers=left,,]
\PreprocessorTok{\#include }\ImportTok{\textless{}stdio.h\textgreater{}}

\DataTypeTok{int}\NormalTok{ add}\OperatorTok{(}\DataTypeTok{int}\NormalTok{ a}\OperatorTok{,} \DataTypeTok{int}\NormalTok{ b}\OperatorTok{)}
\OperatorTok{\{}
    \ControlFlowTok{return}\NormalTok{ a }\OperatorTok{+}\NormalTok{ b}\OperatorTok{;}
\OperatorTok{\}}

\DataTypeTok{int}\NormalTok{ mult}\OperatorTok{(}\DataTypeTok{int}\NormalTok{ a}\OperatorTok{,} \DataTypeTok{int}\NormalTok{ b}\OperatorTok{)}
\OperatorTok{\{}
    \ControlFlowTok{return}\NormalTok{ a }\OperatorTok{*}\NormalTok{ b}\OperatorTok{;}
\OperatorTok{\}}

\DataTypeTok{void}\NormalTok{ print\_math}\OperatorTok{(}\DataTypeTok{int} \OperatorTok{(*}\NormalTok{op}\OperatorTok{)(}\DataTypeTok{int}\OperatorTok{,} \DataTypeTok{int}\OperatorTok{),} \DataTypeTok{int}\NormalTok{ x}\OperatorTok{,} \DataTypeTok{int}\NormalTok{ y}\OperatorTok{)}
\OperatorTok{\{}
    \DataTypeTok{int}\NormalTok{ result }\OperatorTok{=}\NormalTok{ op}\OperatorTok{(}\NormalTok{x}\OperatorTok{,}\NormalTok{ y}\OperatorTok{);}

\NormalTok{    printf}\OperatorTok{(}\StringTok{"}\SpecialCharTok{\%d\textbackslash{}n}\StringTok{"}\OperatorTok{,}\NormalTok{ result}\OperatorTok{);}
\OperatorTok{\}}

\DataTypeTok{int}\NormalTok{ main}\OperatorTok{(}\DataTypeTok{void}\OperatorTok{)}
\OperatorTok{\{}
\NormalTok{    print\_math}\OperatorTok{(}\NormalTok{add}\OperatorTok{,} \DecValTok{5}\OperatorTok{,} \DecValTok{7}\OperatorTok{);}   \CommentTok{// 12}
\NormalTok{    print\_math}\OperatorTok{(}\NormalTok{mult}\OperatorTok{,} \DecValTok{5}\OperatorTok{,} \DecValTok{7}\OperatorTok{);}  \CommentTok{// 35}
\OperatorTok{\}}
\end{Highlighting}
\end{Shaded}

\end{tcolorbox}

Luangkan waktu untuk mencernanya. Idenya, kita akan mengoper pointer ke sebuah fungsi ke \hl{\mbox{\texttt{print\_math()}}}, dan ia akan memanggil fungsi itu untuk melakukan perhitungan.

Dengan begitu kita bisa mengubah perilaku \hl{\mbox{\texttt{print\_math()}}} dengan mengoper fungsi lain ke dalamnya. Kamu bisa lihat kita melakukannya di baris 22-23 saat mengoper pointer ke fungsi \hl{\mbox{\texttt{add}}} dan \hl{\mbox{\texttt{mult}}}.

Nah, di baris 13, kita semua pasti setuju signature fungsi \hl{\mbox{\texttt{print\_math()}}} adalah pemandangan yang patut dikagumi. Dan, percaya atau tidak, yang ini sebenarnya cukup lugas dibanding beberapa hal yang bisa kamu susun\footnote{Bahasa Pemrograman Go mengambil inspirasi sintaks deklarasi tipenya dari kebalikan dari yang dilakukan C.}.

Tapi mari kita cerna. Ternyata hanya ada tiga parameter, tapi agak sulit dilihat:

\begin{tcolorbox}[enhanced jigsaw,left=3pt,breakable,colback=white!97.2549!black,boxrule=0.2mm]

\begin{Shaded}
\begin{Highlighting}[]
\CommentTok{//                      op             x      y}
\CommentTok{//              |{-}{-}{-}{-}{-}{-}{-}{-}{-}{-}{-}{-}{-}{-}{-}{-}{-}|  |{-}{-}{-}|  |{-}{-}{-}|}
\DataTypeTok{void}\NormalTok{ print\_math}\OperatorTok{(}\DataTypeTok{int} \OperatorTok{(*}\NormalTok{op}\OperatorTok{)(}\DataTypeTok{int}\OperatorTok{,} \DataTypeTok{int}\OperatorTok{),} \DataTypeTok{int}\NormalTok{ x}\OperatorTok{,} \DataTypeTok{int}\NormalTok{ y}\OperatorTok{)}
\end{Highlighting}
\end{Shaded}

\end{tcolorbox}

Parameter pertama, \hl{\mbox{\texttt{op}}}, berupa pointer ke fungsi yang menerima dua \hl{\mbox{\texttt{int}}} sebagai argumen dan mengembalikan \hl{\mbox{\texttt{int}}}. Ini cocok dengan signature \hl{\mbox{\texttt{add()}}} maupun \hl{\mbox{\texttt{mult()}}}.

Yang kedua dan ketiga, \hl{\mbox{\texttt{x}}} dan \hl{\mbox{\texttt{y}}}, hanyalah parameter \hl{\mbox{\texttt{int}}} biasa.

Pelan-pelan dan sengaja, biarkan matamu menelusuri signature itu sambil mengenali bagian-bagian yang bekerja. Satu hal yang selalu menonjol bagiku adalah urutan \hl{\mbox{\texttt{(*op)(}}}, tanda kurung dan asterisk itu. Itulah tanda bahwa ia pointer ke fungsi.

Terakhir, lompat kembali ke bab \emph{Pointer II} untuk melihat \hyperref[qsort-example]{contoh pointer-ke-fungsi memakai \hl{\mbox{\texttt{qsort()}}} bawaan}.

\index{Pointers!to functions|)}

\chapter{Operasi Bitwise}\label{operasi-bitwise}

\index{Bitwise operations|(}

Operasi numerik ini memungkinkan kamu memanipulasi bit individual dalam variabel, cocok karena C adalah bahasa yang begitu low-level\footnote{Bukan berarti bahasa lain tidak melakukannya---mereka melakukannya. Menarik betapa banyak bahasa modern memakai operator yang sama dengan C untuk bitwise.}.

Kalau kamu belum akrab dengan operasi bitwise, \href{https://en.wikipedia.org/wiki/Bitwise_operation}{Wikipedia punya artikel bitwise yang bagus}\footnote{https://en.wikipedia.org/wiki/Bitwise\_operation}.

\section{Bitwise AND, OR, XOR, dan NOT}\label{bitwise-and-or-xor-dan-not}

Untuk masing-masing, \hyperref[usual-arithmetic-conversions]{usual arithmetic conversion} terjadi pada operand (yang dalam kasus ini harus bertipe integer), lalu operasi bitwise yang sesuai dilakukan.

\index{\& bitwise AND@\texttt{\&} bitwise AND} \index{\textbar{} bitwise OR@\texttt{\textbar{}} bitwise OR} \index{\^{} bitwise XOR@\texttt{\^{}} bitwise XOR} \index{~ bitwise NOT@\texttt{~} bitwise NOT}

{\def\LTcaptype{none} % do not increment counter
\begin{longtable}[]{@{}lcl@{}}
\toprule\noalign{}
Operasi & Operator & Contoh \\
\midrule\noalign{}
\endhead
\bottomrule\noalign{}
\endlastfoot
AND & \hl{\mbox{\texttt{\&}}} & \hl{\mbox{\texttt{a\ =\ b\ \&\ c}}} \\
OR & \hl{\mbox{\texttt{\textbar{}}}} & \hl{\mbox{\texttt{a\ =\ b\ \textbar{}\ c}}} \\
XOR & \hl{\mbox{\texttt{\^{}}}} & \hl{\mbox{\texttt{a\ =\ b\ \^{}\ c}}} \\
NOT & \hl{\mbox{\texttt{\textasciitilde{}}}} & \hl{\mbox{\texttt{a\ =\ \textasciitilde{}c}}} \\
\end{longtable}
}

Perhatikan betapa miripnya mereka dengan operator Boolean \hl{\mbox{\texttt{\&\&}}} dan \hl{\mbox{\texttt{\textbar{}\textbar{}}}}.

Operator-operator ini punya varian singkat penugasan seperti \hl{\mbox{\texttt{+=}}} dan \hl{\mbox{\texttt{-=}}}:

\index{\&= assignment@\texttt{\&=} assignment} \index{\textbar{}= assignment@\texttt{\textbar{}=} assignment} \index{\^{}= assignment@\texttt{\^{}=} assignment}

{\def\LTcaptype{none} % do not increment counter
\begin{longtable}[]{@{}lll@{}}
\toprule\noalign{}
Operator & Contoh & Bentuk panjang yang setara \\
\midrule\noalign{}
\endhead
\bottomrule\noalign{}
\endlastfoot
\hl{\mbox{\texttt{\&=}}} & \hl{\mbox{\texttt{a\ \&=\ c}}} & \hl{\mbox{\texttt{a\ =\ a\ \&\ c}}} \\
\hl{\mbox{\texttt{\textbar{}=}}} & \hl{\mbox{\texttt{a\ \textbar{}=\ c}}} & \hl{\mbox{\texttt{a\ =\ a\ \textbar{}\ c}}} \\
\hl{\mbox{\texttt{\^{}=}}} & \hl{\mbox{\texttt{a\ \^{}=\ c}}} & \hl{\mbox{\texttt{a\ =\ a\ \^{}\ c}}} \\
\end{longtable}
}

\section{Bitwise Shift}\label{bitwise-shift}

Untuk ini, \index{Integer promotions} \hyperref[integer-promotions]{integer promotion} dilakukan pada setiap operand (yang harus bertipe integer) lalu bitwise shift dieksekusi. Tipe hasilnya adalah tipe operand kiri yang sudah dipromosikan.

Bit baru diisi dengan nol, dengan kemungkinan pengecualian yang disebut pada perilaku yang ditentukan implementasi di bawah.

\index{<< shift left@\texttt{<<} shift left} \index{>> shift right@\texttt{>>} shift right}

{\def\LTcaptype{none} % do not increment counter
\begin{longtable}[]{@{}lcl@{}}
\toprule\noalign{}
Operasi & Operator & Contoh \\
\midrule\noalign{}
\endhead
\bottomrule\noalign{}
\endlastfoot
Shift kiri & \hl{\mbox{\texttt{\textless{}\textless{}}}} & \hl{\mbox{\texttt{a\ =\ b\ \textless{}\textless{}\ c}}} \\
Shift kanan & \hl{\mbox{\texttt{\textgreater{}\textgreater{}}}} & \hl{\mbox{\texttt{a\ =\ b\ \textgreater{}\textgreater{}\ c}}} \\
\end{longtable}
}

Ada juga bentuk singkat yang serupa untuk shift:

\index{>>= assignment@\texttt{>>=} assignment} \index{<<= assignment@\texttt{<<=} assignment}

{\def\LTcaptype{none} % do not increment counter
\begin{longtable}[]{@{}lll@{}}
\toprule\noalign{}
Operator & Contoh & Bentuk panjang yang setara \\
\midrule\noalign{}
\endhead
\bottomrule\noalign{}
\endlastfoot
\hl{\mbox{\texttt{\textgreater{}\textgreater{}=}}} & \hl{\mbox{\texttt{a\ \textgreater{}\textgreater{}=\ c}}} & \hl{\mbox{\texttt{a\ =\ a\ \textgreater{}\textgreater{}\ c}}} \\
\hl{\mbox{\texttt{\textless{}\textless{}=}}} & \hl{\mbox{\texttt{a\ \textless{}\textless{}=\ c}}} & \hl{\mbox{\texttt{a\ =\ a\ \textless{}\textless{}\ c}}} \\
\end{longtable}
}

Waspadai undefined behavior: tidak boleh shift negatif, dan tidak boleh shift yang lebih besar dari ukuran operand kiri yang sudah dipromosikan.

Waspadai juga perilaku yang ditentukan implementasi: kalau kamu melakukan shift kanan pada angka negatif, hasilnya ditentukan implementasi. (Aman saja melakukan shift kanan pada \hl{\mbox{\texttt{int}}} signed, asalkan nilainya positif.)

\index{Bitwise operations|)}

\chapter{Fungsi Variadic}\label{fungsi-variadic}

\index{Variadic functions|(}

\emph{Variadic} adalah kata keren untuk fungsi yang menerima jumlah argumen sembarang.

Fungsi biasa menerima jumlah argumen yang tertentu, misalnya:

\begin{tcolorbox}[enhanced jigsaw,left=3pt,breakable,colback=white!97.2549!black,boxrule=0.2mm]

\begin{Shaded}
\begin{Highlighting}[]
\DataTypeTok{int}\NormalTok{ add}\OperatorTok{(}\DataTypeTok{int}\NormalTok{ x}\OperatorTok{,} \DataTypeTok{int}\NormalTok{ y}\OperatorTok{)}
\OperatorTok{\{}
    \ControlFlowTok{return}\NormalTok{ x }\OperatorTok{+}\NormalTok{ y}\OperatorTok{;}
\OperatorTok{\}}
\end{Highlighting}
\end{Shaded}

\end{tcolorbox}

Kamu hanya bisa memanggilnya dengan tepat dua argumen yang sesuai dengan parameter \hl{\mbox{\texttt{x}}} dan \hl{\mbox{\texttt{y}}}.

\begin{tcolorbox}[enhanced jigsaw,left=3pt,breakable,colback=white!97.2549!black,boxrule=0.2mm]

\begin{Shaded}
\begin{Highlighting}[]
\NormalTok{add}\OperatorTok{(}\DecValTok{2}\OperatorTok{,} \DecValTok{3}\OperatorTok{);}
\NormalTok{add}\OperatorTok{(}\DecValTok{5}\OperatorTok{,} \DecValTok{12}\OperatorTok{);}
\end{Highlighting}
\end{Shaded}

\end{tcolorbox}

Tapi kalau kamu mencoba dengan lebih banyak, compiler tidak akan mengizinkan:

\begin{tcolorbox}[enhanced jigsaw,left=3pt,breakable,colback=white!97.2549!black,boxrule=0.2mm]

\begin{Shaded}
\begin{Highlighting}[]
\NormalTok{add}\OperatorTok{(}\DecValTok{2}\OperatorTok{,} \DecValTok{3}\OperatorTok{,} \DecValTok{4}\OperatorTok{);}  \CommentTok{// ERROR}
\NormalTok{add}\OperatorTok{(}\DecValTok{5}\OperatorTok{);}        \CommentTok{// ERROR}
\end{Highlighting}
\end{Shaded}

\end{tcolorbox}

Fungsi variadic mengatasi batasan ini sampai tingkat tertentu.

Kita sudah melihat contoh terkenalnya di \hl{\mbox{\texttt{printf()}}}! Kamu bisa mengoper berbagai macam hal kepadanya.

\begin{tcolorbox}[enhanced jigsaw,left=3pt,breakable,colback=white!97.2549!black,boxrule=0.2mm]

\begin{Shaded}
\begin{Highlighting}[]
\NormalTok{printf}\OperatorTok{(}\StringTok{"Hello, world!}\SpecialCharTok{\textbackslash{}n}\StringTok{"}\OperatorTok{);}
\NormalTok{printf}\OperatorTok{(}\StringTok{"The number is }\SpecialCharTok{\%d\textbackslash{}n}\StringTok{"}\OperatorTok{,} \DecValTok{2}\OperatorTok{);}
\NormalTok{printf}\OperatorTok{(}\StringTok{"The number is }\SpecialCharTok{\%d}\StringTok{ and pi is }\SpecialCharTok{\%f\textbackslash{}n}\StringTok{"}\OperatorTok{,} \DecValTok{2}\OperatorTok{,} \FloatTok{3.14159}\OperatorTok{);}
\end{Highlighting}
\end{Shaded}

\end{tcolorbox}

Kelihatannya ia tidak peduli berapa banyak argumen yang diberikan!

Nah, itu tidak sepenuhnya benar. Nol argumen akan memberi error:

\begin{tcolorbox}[enhanced jigsaw,left=3pt,breakable,colback=white!97.2549!black,boxrule=0.2mm]

\begin{Shaded}
\begin{Highlighting}[]
\NormalTok{printf}\OperatorTok{();}  \CommentTok{// ERROR}
\end{Highlighting}
\end{Shaded}

\end{tcolorbox}

Ini membawa kita ke salah satu batasan fungsi variadic di C: mereka harus punya setidaknya satu argumen.

Tapi selain itu, mereka cukup fleksibel, bahkan membolehkan argumen bertipe berbeda-beda seperti \hl{\mbox{\texttt{printf()}}}.

Mari kita lihat cara kerjanya!

\section{Elipsis pada Signature Fungsi}\label{elipsis-pada-signature-fungsi}

Jadi bagaimana cara kerjanya, secara sintaksis?

\index{... variadic arguments@\texttt{...} variadic arguments|(}

Taruh semua argumen yang \emph{harus} dioper lebih dulu (dan ingat harus ada setidaknya satu), lalu setelah itu taruh \hl{\mbox{\texttt{...}}}. Seperti ini:

\begin{tcolorbox}[enhanced jigsaw,left=3pt,breakable,colback=white!97.2549!black,boxrule=0.2mm]

\begin{Shaded}
\begin{Highlighting}[]
\DataTypeTok{void}\NormalTok{ func}\OperatorTok{(}\DataTypeTok{int}\NormalTok{ a}\OperatorTok{,} \OperatorTok{...)}   \CommentTok{// Secara harfiah 3 titik di sini}
\end{Highlighting}
\end{Shaded}

\end{tcolorbox}

Berikut kode untuk mendemonstrasikannya:

\begin{tcolorbox}[enhanced jigsaw,left=3pt,breakable,colback=white!97.2549!black,boxrule=0.2mm]

\begin{Shaded}
\begin{Highlighting}[]
\PreprocessorTok{\#include }\ImportTok{\textless{}stdio.h\textgreater{}}

\DataTypeTok{void}\NormalTok{ func}\OperatorTok{(}\DataTypeTok{int}\NormalTok{ a}\OperatorTok{,} \OperatorTok{...)}
\OperatorTok{\{}
\NormalTok{    printf}\OperatorTok{(}\StringTok{"a is }\SpecialCharTok{\%d\textbackslash{}n}\StringTok{"}\OperatorTok{,}\NormalTok{ a}\OperatorTok{);}  \CommentTok{// Mencetak "a is 2"}
\OperatorTok{\}}

\DataTypeTok{int}\NormalTok{ main}\OperatorTok{(}\DataTypeTok{void}\OperatorTok{)}
\OperatorTok{\{}
\NormalTok{    func}\OperatorTok{(}\DecValTok{2}\OperatorTok{,} \DecValTok{3}\OperatorTok{,} \DecValTok{4}\OperatorTok{,} \DecValTok{5}\OperatorTok{,} \DecValTok{6}\OperatorTok{);}
\OperatorTok{\}}
\end{Highlighting}
\end{Shaded}

\end{tcolorbox}

\index{... variadic arguments@\texttt{...} variadic arguments|)}

Jadi, bagus, kita bisa mendapatkan argumen pertama yang ada di variabel \hl{\mbox{\texttt{a}}}, tapi bagaimana dengan argumen sisanya? Bagaimana cara mengaksesnya?

Di sinilah keseruan dimulai!

\section{Mengambil Argumen Tambahan}\label{mengambil-argumen-tambahan}

Kamu perlu meng-include \index{stdarg.h header file@\texttt{stdarg.h} header file} \hl{\mbox{\texttt{\textless{}stdarg.h\textgreater{}}}} supaya semua ini berfungsi.

\index{va\_list type@\texttt{va\_list} type|(}

Pertama-tama, kita akan memakai variabel khusus bertipe \hl{\mbox{\texttt{va\_list}}} (variable argument list) untuk melacak variabel mana yang sedang kita akses.

\index{va\_start() macro@\texttt{va\_start()} macro|(} \index{va\_arg() macro@\texttt{va\_arg()} macro|(} \index{va\_end() macro@\texttt{va\_end()} macro|(}

Idenya: kita mulai memproses argumen dengan memanggil \hl{\mbox{\texttt{va\_start()}}}, memproses tiap argumen secara berurutan dengan \hl{\mbox{\texttt{va\_arg()}}}, lalu setelah selesai, menutupnya dengan \hl{\mbox{\texttt{va\_end()}}}.

Saat memanggil \hl{\mbox{\texttt{va\_start()}}}, kamu perlu mengoper \emph{parameter bernama terakhir} (yang tepat sebelum \hl{\mbox{\texttt{...}}}) supaya ia tahu dari mana mulai mencari argumen tambahan.

Dan saat memanggil \hl{\mbox{\texttt{va\_arg()}}} untuk mengambil argumen berikutnya, kamu harus memberi tahu tipe argumen yang akan diambil.

Berikut demo yang menjumlahkan sejumlah integer sembarang. Argumen pertama adalah jumlah integer yang akan dijumlahkan. Kita memakainya untuk mengetahui berapa kali harus memanggil \hl{\mbox{\texttt{va\_arg()}}}.

\begin{tcolorbox}[enhanced jigsaw,left=3pt,breakable,colback=white!97.2549!black,boxrule=0.2mm]

\begin{Shaded}
\begin{Highlighting}[numbers=left,,]
\PreprocessorTok{\#include }\ImportTok{\textless{}stdio.h\textgreater{}}
\PreprocessorTok{\#include }\ImportTok{\textless{}stdarg.h\textgreater{}}

\DataTypeTok{int}\NormalTok{ add}\OperatorTok{(}\DataTypeTok{int}\NormalTok{ count}\OperatorTok{,} \OperatorTok{...)}
\OperatorTok{\{}
    \DataTypeTok{int}\NormalTok{ total }\OperatorTok{=} \DecValTok{0}\OperatorTok{;}
    \DataTypeTok{va\_list}\NormalTok{ va}\OperatorTok{;}

\NormalTok{    va\_start}\OperatorTok{(}\NormalTok{va}\OperatorTok{,}\NormalTok{ count}\OperatorTok{);}   \CommentTok{// Mulai dengan argumen setelah "count"}

    \ControlFlowTok{for} \OperatorTok{(}\DataTypeTok{int}\NormalTok{ i }\OperatorTok{=} \DecValTok{0}\OperatorTok{;}\NormalTok{ i }\OperatorTok{\textless{}}\NormalTok{ count}\OperatorTok{;}\NormalTok{ i}\OperatorTok{++)} \OperatorTok{\{}
        \DataTypeTok{int}\NormalTok{ n }\OperatorTok{=}\NormalTok{ va\_arg}\OperatorTok{(}\NormalTok{va}\OperatorTok{,} \DataTypeTok{int}\OperatorTok{);}   \CommentTok{// Ambil int berikutnya}

\NormalTok{        total }\OperatorTok{+=}\NormalTok{ n}\OperatorTok{;}
    \OperatorTok{\}}

\NormalTok{    va\_end}\OperatorTok{(}\NormalTok{va}\OperatorTok{);}  \CommentTok{// Selesai}

    \ControlFlowTok{return}\NormalTok{ total}\OperatorTok{;}
\OperatorTok{\}}

\DataTypeTok{int}\NormalTok{ main}\OperatorTok{(}\DataTypeTok{void}\OperatorTok{)}
\OperatorTok{\{}
\NormalTok{    printf}\OperatorTok{(}\StringTok{"}\SpecialCharTok{\%d\textbackslash{}n}\StringTok{"}\OperatorTok{,}\NormalTok{ add}\OperatorTok{(}\DecValTok{4}\OperatorTok{,} \DecValTok{6}\OperatorTok{,} \DecValTok{2}\OperatorTok{,} \OperatorTok{{-}}\DecValTok{4}\OperatorTok{,} \DecValTok{17}\OperatorTok{));}  \CommentTok{// 6 + 2 {-} 4 + 17 = 21}
\NormalTok{    printf}\OperatorTok{(}\StringTok{"}\SpecialCharTok{\%d\textbackslash{}n}\StringTok{"}\OperatorTok{,}\NormalTok{ add}\OperatorTok{(}\DecValTok{2}\OperatorTok{,} \DecValTok{22}\OperatorTok{,} \DecValTok{44}\OperatorTok{));}        \CommentTok{// 22 + 44 = 66}
\OperatorTok{\}}
\end{Highlighting}
\end{Shaded}

\end{tcolorbox}

\index{va\_start() macro@\texttt{va\_start()} macro|)} \index{va\_end() macro@\texttt{va\_end()} macro|)}

(Perhatikan bahwa saat \hl{\mbox{\texttt{printf()}}} dipanggil, ia memakai jumlah \hl{\mbox{\texttt{\%d}}} (atau sejenisnya) di format string untuk mengetahui berapa banyak argumen tambahan!)

Kalau sintaks \hl{\mbox{\texttt{va\_arg()}}} terlihat aneh bagimu (karena ada nama tipe yang melayang bebas di dalamnya), kamu tidak sendirian. Ini diimplementasikan dengan makro preprocessor supaya semua sihirnya bisa masuk.

\section{\texorpdfstring{Fungsionalitas \hl{\mbox{\texttt{va\_list}}}}{Fungsionalitas va\_list}}\label{fungsionalitas-va_list}

Apa itu variabel \hl{\mbox{\texttt{va\_list}}} yang kita pakai di atas? Ia adalah variabel opaque\footnote{Maksudnya, kita para developer rendahan tidak seharusnya tahu apa isinya atau artinya. Spesifikasi tidak menentukan secara rinci apa itu.} yang menyimpan informasi tentang argumen mana yang akan kita ambil berikutnya dengan \hl{\mbox{\texttt{va\_arg()}}}. Lihat bagaimana kita memanggil \hl{\mbox{\texttt{va\_arg()}}} berulang kali? Variabel \hl{\mbox{\texttt{va\_list}}} adalah placeholder yang melacak kemajuan sejauh ini.

\index{va\_start() macro@\texttt{va\_start()} macro|(}

Tapi kita harus menginisialisasi variabel itu ke nilai yang masuk akal. Di situlah \hl{\mbox{\texttt{va\_start()}}} berperan.

Saat memanggil \hl{\mbox{\texttt{va\_start(va,\ count)}}} di atas, kita berkata, ``Inisialisasi variabel \hl{\mbox{\texttt{va}}} agar menunjuk ke argumen variabel \emph{tepat setelah} \hl{\mbox{\texttt{count}}}.''

\index{va\_end() macro@\texttt{va\_end()} macro|(}

Dan itulah \emph{alasan} kita perlu punya setidaknya satu variabel bernama di daftar argumen\footnote{Jujur saja, batasan itu sebenarnya bisa dihapus dari bahasa, tapi idenya adalah makro \hl{\mbox{\texttt{va\_start()}}}, \hl{\mbox{\texttt{va\_arg()}}}, dan \hl{\mbox{\texttt{va\_end()}}} harus bisa ditulis dalam C. Dan supaya itu terwujud, kita butuh cara untuk menginisialisasi pointer ke lokasi parameter pertama. Dan untuk itu, kita butuh \emph{nama} parameter pertama. Itu akan butuh ekstensi bahasa untuk dimungkinkan, dan sejauh ini komite belum menemukan alasan untuk melakukannya.}.

Setelah punya pointer ke parameter awal, kamu bisa dengan mudah mengambil nilai argumen berikutnya dengan memanggil \hl{\mbox{\texttt{va\_arg()}}} berulang kali. Saat melakukannya, kamu harus mengoper variabel \hl{\mbox{\texttt{va\_list}}} (supaya ia bisa terus melacak posisimu), serta tipe argumen yang akan kamu salin.

Kamu sebagai programmer yang harus menentukan tipe apa yang dioper ke \hl{\mbox{\texttt{va\_arg()}}}. Pada contoh di atas, kita hanya memakai \hl{\mbox{\texttt{int}}}. Tapi pada \hl{\mbox{\texttt{printf()}}}, ia memakai format specifier untuk menentukan tipe apa yang diambil berikutnya.

Dan saat selesai, panggil \hl{\mbox{\texttt{va\_end()}}} untuk menutupnya. Kamu \textbf{harus} (kata spesifikasi) memanggilnya pada variabel \hl{\mbox{\texttt{va\_list}}} tertentu sebelum memutuskan memanggil \hl{\mbox{\texttt{va\_start()}}} atau \hl{\mbox{\texttt{va\_copy()}}} lagi padanya. Aku tahu kita belum membahas \hl{\mbox{\texttt{va\_copy()}}}.

Jadi urutan standarnya adalah:

\begin{itemize}
\tightlist
\item
  \hl{\mbox{\texttt{va\_start()}}} untuk menginisialisasi variabel \hl{\mbox{\texttt{va\_list}}}
\item
  \hl{\mbox{\texttt{va\_arg()}}} berulang kali untuk mengambil nilai-nilainya
\item
  \hl{\mbox{\texttt{va\_end()}}} untuk men-deinisialisasi variabel \hl{\mbox{\texttt{va\_list}}}
\end{itemize}

\index{va\_start() macro@\texttt{va\_start()} macro|)} \index{va\_end() macro@\texttt{va\_end()} macro|)} \index{va\_arg() macro@\texttt{va\_arg()} macro|)}

\index{va\_copy() macro@\texttt{va\_copy()} macro|(}

Aku juga sudah menyebut \hl{\mbox{\texttt{va\_copy()}}} di atas; ia membuat salinan variabel \hl{\mbox{\texttt{va\_list}}}-mu dalam keadaan yang persis sama. Artinya, kalau kamu belum mulai dengan \hl{\mbox{\texttt{va\_arg()}}} pada variabel sumber, yang baru juga belum dimulai. Kalau kamu sudah mengambil 5 variabel dengan \hl{\mbox{\texttt{va\_arg()}}} sejauh ini, salinannya juga mencerminkan itu.

\hl{\mbox{\texttt{va\_copy()}}} bisa berguna kalau kamu perlu memindai argumen lebih ke depan tapi juga perlu mengingat posisimu saat ini.

\index{va\_copy() macro@\texttt{va\_copy()} macro|)}

\section{\texorpdfstring{Fungsi Library yang Memakai \hl{\mbox{\texttt{va\_list}}}}{Fungsi Library yang Memakai va\_list}}\label{fungsi-library-yang-memakai-va_list}

\index{va\_list type@\texttt{va\_list} type!passing to functions|(}

Salah satu kegunaan lain cukup keren: menulis varian \hl{\mbox{\texttt{printf()}}} kustommu sendiri. Repot sekali kalau harus menangani semua format specifier itu, kan? Semiliar jumlahnya?

Untungnya, ada varian \hl{\mbox{\texttt{printf()}}} yang menerima \hl{\mbox{\texttt{va\_list}}} yang sudah jalan sebagai argumen. Kamu bisa memakainya untuk membungkus dan membuat \hl{\mbox{\texttt{printf()}}} kustommu sendiri!

\index{vprintf() function@\texttt{vprintf()} function|(}

Fungsi-fungsi ini diawali huruf \hl{\mbox{\texttt{v}}}, seperti \hl{\mbox{\texttt{vprintf()}}}, \hl{\mbox{\texttt{vfprintf()}}}, \hl{\mbox{\texttt{vsprintf()}}}, dan \hl{\mbox{\texttt{vsnprintf()}}}. Pada dasarnya semua \hl{\mbox{\texttt{printf()}}} lawas favoritmu, hanya dengan \hl{\mbox{\texttt{v}}} di depan.

Mari kita buat fungsi \hl{\mbox{\texttt{my\_printf()}}} yang bekerja persis seperti \hl{\mbox{\texttt{printf()}}} kecuali ia menerima satu argumen ekstra di depan.

\begin{tcolorbox}[enhanced jigsaw,left=3pt,breakable,colback=white!97.2549!black,boxrule=0.2mm]

\begin{Shaded}
\begin{Highlighting}[numbers=left,,]
\PreprocessorTok{\#include }\ImportTok{\textless{}stdio.h\textgreater{}}
\PreprocessorTok{\#include }\ImportTok{\textless{}stdarg.h\textgreater{}}

\DataTypeTok{int}\NormalTok{ my\_printf}\OperatorTok{(}\DataTypeTok{int}\NormalTok{ serial}\OperatorTok{,} \DataTypeTok{const} \DataTypeTok{char} \OperatorTok{*}\NormalTok{format}\OperatorTok{,} \OperatorTok{...)}
\OperatorTok{\{}
    \DataTypeTok{va\_list}\NormalTok{ va}\OperatorTok{;}
    \DataTypeTok{int}\NormalTok{ rv}\OperatorTok{;}

    \CommentTok{// Kerjakan bagian kustom}
\NormalTok{    printf}\OperatorTok{(}\StringTok{"The serial number is: }\SpecialCharTok{\%d\textbackslash{}n}\StringTok{"}\OperatorTok{,}\NormalTok{ serial}\OperatorTok{);}

    \CommentTok{// Lalu serahkan sisanya ke vprintf()}
\NormalTok{    va\_start}\OperatorTok{(}\NormalTok{va}\OperatorTok{,}\NormalTok{ format}\OperatorTok{);}
\NormalTok{    rv }\OperatorTok{=}\NormalTok{ vprintf}\OperatorTok{(}\NormalTok{format}\OperatorTok{,}\NormalTok{ va}\OperatorTok{);}
\NormalTok{    va\_end}\OperatorTok{(}\NormalTok{va}\OperatorTok{);}

    \ControlFlowTok{return}\NormalTok{ rv}\OperatorTok{;}
\OperatorTok{\}}

\DataTypeTok{int}\NormalTok{ main}\OperatorTok{(}\DataTypeTok{void}\OperatorTok{)}
\OperatorTok{\{}
    \DataTypeTok{int}\NormalTok{ x }\OperatorTok{=} \DecValTok{10}\OperatorTok{;}
    \DataTypeTok{float}\NormalTok{ y }\OperatorTok{=} \FloatTok{3.2}\OperatorTok{;}

\NormalTok{    my\_printf}\OperatorTok{(}\DecValTok{3490}\OperatorTok{,} \StringTok{"x is }\SpecialCharTok{\%d}\StringTok{, y is }\SpecialCharTok{\%f\textbackslash{}n}\StringTok{"}\OperatorTok{,}\NormalTok{ x}\OperatorTok{,}\NormalTok{ y}\OperatorTok{);}
\OperatorTok{\}}
\end{Highlighting}
\end{Shaded}

\end{tcolorbox}

Lihat apa yang kita lakukan di situ? Di baris 12-14 kita memulai variabel \hl{\mbox{\texttt{va\_list}}} baru, lalu langsung mengopernya ke \hl{\mbox{\texttt{vprintf()}}}. Dan \hl{\mbox{\texttt{vprintf()}}} tahu persis apa yang harus dilakukan dengannya, karena ia punya semua kecerdasan \hl{\mbox{\texttt{printf()}}} bawaan.

\index{vprintf() function@\texttt{vprintf()} function|)}

Kita tetap harus memanggil \hl{\mbox{\texttt{va\_end()}}} saat selesai, jadi jangan lupa!

\index{va\_list type@\texttt{va\_list} type!passing to functions|)} \index{va\_list type@\texttt{va\_list} type|)}

\section{Jebakan Makro Variadic}\label{jebakan-makro-variadic}

Seperti sudah kusebut, \hl{\mbox{\texttt{va\_start()}}} dan \hl{\mbox{\texttt{va\_end()}}} bisa berupa makro. Salah satu konsekuensinya, mereka mungkin memperkenalkan scope lokal baru. (Yaitu, kalau \hl{\mbox{\texttt{va\_start()}}} berisi \hl{\mbox{\texttt{\{}}} dan \hl{\mbox{\texttt{va\_end()}}} berisi \hl{\mbox{\texttt{\}}}}.)

\index{va\_start() macro@\texttt{va\_start()} macro!scoping issues|)}

Jadi kita perlu waspada terhadap kemungkinan masalah scope. Lihat contoh berikut:

\begin{tcolorbox}[enhanced jigsaw,left=3pt,breakable,colback=white!97.2549!black,boxrule=0.2mm]

\begin{Shaded}
\begin{Highlighting}[]
\NormalTok{va\_start}\OperatorTok{(}\NormalTok{va}\OperatorTok{,}\NormalTok{ format}\OperatorTok{);}          \CommentTok{// mungkin berisi \{}
\DataTypeTok{int}\NormalTok{ rv }\OperatorTok{=}\NormalTok{ vprintf}\OperatorTok{(}\NormalTok{format}\OperatorTok{,}\NormalTok{ va}\OperatorTok{);}
\NormalTok{va\_end}\OperatorTok{(}\NormalTok{va}\OperatorTok{);}                    \CommentTok{// mungkin berisi \}}

\ControlFlowTok{return}\NormalTok{ rv}\OperatorTok{;}
\end{Highlighting}
\end{Shaded}

\end{tcolorbox}

Kalau \hl{\mbox{\texttt{va\_start()}}} membuka scope baru, \hl{\mbox{\texttt{rv}}} akan lokal di scope itu dan pernyataan \hl{\mbox{\texttt{return}}} akan gagal. Tapi ini diam-diam hanya akan terjadi di compiler yang kebetulan melakukan itu dengan makro \hl{\mbox{\texttt{va}}}.

\index{Variadic functions|)}

\chapter{Locale dan Internasionalisasi}\label{locale-dan-internasionalisasi}

\index{Locale|(}

\emph{Lokalisasi} adalah proses membuat aplikasimu siap bekerja dengan baik di locale (atau negara) yang berbeda.

Seperti yang mungkin kamu tahu, tidak semua orang memakai karakter yang sama untuk titik desimal atau pemisah ribuan\ldots{} atau untuk mata uang.

Locale-locale ini punya nama, dan kamu bisa memilih salah satu untuk dipakai. Misalnya, locale AS mungkin menulis angka seperti:

100,000.00

Sedangkan di Brasil, angka yang sama mungkin ditulis dengan koma dan titik desimal bertukar tempat:

100.000,00

Jadi lebih mudah menulis kodemu supaya bisa dibawa ke negara lain dengan mulus!

Ya, kurang lebih begitu. Ternyata C hanya punya satu locale bawaan, dan itu pun terbatas. Spesifikasi benar-benar menyisakan banyak ambiguitas di sini; sulit untuk sepenuhnya portabel.

Tapi kita akan berusaha sebaik mungkin!

\section{Mengatur Lokalisasi, Cara Cepat dan Kotor}\label{mengatur-lokalisasi-cara-cepat-dan-kotor}

Untuk pemanggilan ini, include \index{locale.h header file@\texttt{locale.h} header file} \hl{\mbox{\texttt{\textless{}locale.h\textgreater{}}}}.

Pada dasarnya hanya ada satu hal yang bisa kamu lakukan secara portabel dalam hal mendeklarasikan locale tertentu. Ini kemungkinan besar yang ingin kamu lakukan kalau kamu akan mengerjakan apa pun terkait locale:

\index{setlocale() function@\texttt{setlocale()} function|(}

\begin{tcolorbox}[enhanced jigsaw,left=3pt,breakable,colback=white!97.2549!black,boxrule=0.2mm]

\begin{Shaded}
\begin{Highlighting}[]
\NormalTok{setlocale}\OperatorTok{(}\NormalTok{LC\_ALL}\OperatorTok{,} \StringTok{""}\OperatorTok{);}  \CommentTok{// Pakai locale environment ini untuk semuanya}
\end{Highlighting}
\end{Shaded}

\end{tcolorbox}

Kamu perlu memanggilnya supaya program diinisialisasi dengan locale yang sedang berlaku.

Lebih rinci, ada satu hal lagi yang bisa dilakukan dan tetap portabel:

\begin{tcolorbox}[enhanced jigsaw,left=3pt,breakable,colback=white!97.2549!black,boxrule=0.2mm]

\begin{Shaded}
\begin{Highlighting}[]
\NormalTok{setlocale}\OperatorTok{(}\NormalTok{LC\_ALL}\OperatorTok{,} \StringTok{"C"}\OperatorTok{);}  \CommentTok{// Pakai locale C bawaan}
\end{Highlighting}
\end{Shaded}

\end{tcolorbox}

tapi itu dipanggil secara bawaan setiap kali programmu mulai, jadi tidak banyak alasan untuk melakukannya sendiri.

Pada string kedua itu, kamu bisa menentukan locale apa pun yang didukung sistemmu. Ini sepenuhnya bergantung pada sistem, jadi bervariasi. Di sistemku, aku bisa menentukan ini:

\begin{tcolorbox}[enhanced jigsaw,left=3pt,breakable,colback=white!97.2549!black,boxrule=0.2mm]

\begin{Shaded}
\begin{Highlighting}[]
\NormalTok{setlocale}\OperatorTok{(}\NormalTok{LC\_ALL}\OperatorTok{,} \StringTok{"en\_US.UTF{-}8"}\OperatorTok{);}  \CommentTok{// Tidak portabel!}
\end{Highlighting}
\end{Shaded}

\end{tcolorbox}

Dan itu akan berfungsi. Tapi hanya portabel ke sistem yang punya nama yang persis sama untuk locale yang persis sama, dan itu tidak bisa dijamin.

Dengan mengoper string kosong (\hl{\mbox{\texttt{""}}}) sebagai argumen kedua, kamu berkata kepada C, ``Hei, cari tahu locale sistem ini saat ini apa supaya aku tidak perlu memberi tahu.''

\index{setlocale() function@\texttt{setlocale()} function|)}

\section{Mengambil Pengaturan Locale Moneter}\label{mengambil-pengaturan-locale-moneter}

\index{Locale!money|(}

Karena memindah-mindahkan lembaran kertas hijau menjanjikan kunci kebahagiaan\footnote{``This planet has---or rather had---a problem, which was this: most of the people living on it were unhappy for pretty much of the time. Many solutions were suggested for this problem, but most of these were largely concerned with the movement of small green pieces of paper, which was odd because on the whole it wasn't the small green pieces of paper that were unhappy.'' ---The Hitchhiker's Guide to the Galaxy, Douglas Adams}, mari bicara tentang locale moneter. Saat menulis kode portabel, kamu harus tahu apa yang diketik untuk uang tunai, kan? Entah itu ``\$'', ``€'', ``¥'', atau ``£''.

\index{localeconv() function@\texttt{localeconv()} function|(}

Bagaimana menulis kode itu tanpa jadi gila? Untungnya, setelah kamu memanggil \hl{\mbox{\texttt{setlocale(LC\_ALL,\ "")}}}, kamu bisa mencari semuanya dengan memanggil \hl{\mbox{\texttt{localeconv()}}}:

\begin{tcolorbox}[enhanced jigsaw,left=3pt,breakable,colback=white!97.2549!black,boxrule=0.2mm]

\begin{Shaded}
\begin{Highlighting}[]
\KeywordTok{struct}\NormalTok{ lconv }\OperatorTok{*}\NormalTok{x }\OperatorTok{=}\NormalTok{ localeconv}\OperatorTok{();}
\end{Highlighting}
\end{Shaded}

\end{tcolorbox}

Fungsi ini mengembalikan pointer ke \hl{\mbox{\texttt{struct\ lconv}}} yang dialokasikan secara statis dan berisi semua informasi menarik yang kamu cari.

Berikut field-field \hl{\mbox{\texttt{struct\ lconv}}} beserta artinya.

Pertama, beberapa konvensi. \hl{\mbox{\texttt{\_p\_}}} berarti ``positif'', dan \hl{\mbox{\texttt{\_n\_}}} berarti ``negatif'', dan \hl{\mbox{\texttt{int\_}}} berarti ``internasional''. Walaupun banyak di antaranya bertipe \hl{\mbox{\texttt{char}}} atau \hl{\mbox{\texttt{char*}}}, kebanyakan (atau string yang ditunjuknya) sebenarnya diperlakukan sebagai integer\footnote{Ingat bahwa \hl{\mbox{\texttt{char}}} hanyalah integer seukuran byte.}.

Sebelum lanjut, ketahuilah bahwa \hl{\mbox{\texttt{CHAR\_MAX}}} (dari \hl{\mbox{\texttt{\textless{}limits.h\textgreater{}}}}) adalah nilai maksimum yang bisa ditampung \hl{\mbox{\texttt{char}}}. Dan banyak nilai \hl{\mbox{\texttt{char}}} berikut memakainya untuk menandakan bahwa nilai itu tidak tersedia di locale tersebut.

{\def\LTcaptype{none} % do not increment counter
\begin{longtable}[]{@{}
  >{\raggedright\arraybackslash}p{(\linewidth - 2\tabcolsep) * \real{0.3125}}
  >{\raggedright\arraybackslash}p{(\linewidth - 2\tabcolsep) * \real{0.6875}}@{}}
\toprule\noalign{}
\begin{minipage}[b]{\linewidth}\raggedright
Field
\end{minipage} & \begin{minipage}[b]{\linewidth}\raggedright
Deskripsi
\end{minipage} \\
\midrule\noalign{}
\endhead
\bottomrule\noalign{}
\endlastfoot
\hl{\mbox{\texttt{char\ *mon\_decimal\_point}}} & Karakter titik desimal untuk uang, mis. \hl{\mbox{\texttt{"."}}}. \\
\hl{\mbox{\texttt{char\ *mon\_thousands\_sep}}} & Karakter pemisah ribuan untuk uang, mis. \hl{\mbox{\texttt{","}}}. \\
\hl{\mbox{\texttt{char\ *mon\_grouping}}} & Deskripsi pengelompokan untuk uang (lihat di bawah). \\
\hl{\mbox{\texttt{char\ *positive\_sign}}} & Tanda positif untuk uang, mis. \hl{\mbox{\texttt{"+"}}} atau \hl{\mbox{\texttt{""}}}. \\
\hl{\mbox{\texttt{char\ *negative\_sign}}} & Tanda negatif untuk uang, mis. \hl{\mbox{\texttt{"-"}}}. \\
\hl{\mbox{\texttt{char\ *currency\_symbol}}} & Simbol mata uang, mis. \hl{\mbox{\texttt{"\$"}}}. \\
\hl{\mbox{\texttt{char\ frac\_digits}}} & Saat mencetak jumlah uang, berapa digit yang dicetak setelah titik desimal, mis. \hl{\mbox{\texttt{2}}}. \\
\hl{\mbox{\texttt{char\ p\_cs\_precedes}}} & \hl{\mbox{\texttt{1}}} jika \hl{\mbox{\texttt{currency\_symbol}}} muncul sebelum nilai untuk jumlah uang non-negatif, \hl{\mbox{\texttt{0}}} jika sesudah. \\
\hl{\mbox{\texttt{char\ n\_cs\_precedes}}} & \hl{\mbox{\texttt{1}}} jika \hl{\mbox{\texttt{currency\_symbol}}} muncul sebelum nilai untuk jumlah uang negatif, \hl{\mbox{\texttt{0}}} jika sesudah. \\
\hl{\mbox{\texttt{char\ p\_sep\_by\_space}}} & Menentukan pemisahan \hl{\mbox{\texttt{currency\ symbol}}} dari nilai untuk jumlah non-negatif (lihat di bawah). \\
\hl{\mbox{\texttt{char\ n\_sep\_by\_space}}} & Menentukan pemisahan \hl{\mbox{\texttt{currency\ symbol}}} dari nilai untuk jumlah negatif (lihat di bawah). \\
\hl{\mbox{\texttt{char\ p\_sign\_posn}}} & Menentukan posisi \hl{\mbox{\texttt{positive\_sign}}} untuk nilai non-negatif. \\
\hl{\mbox{\texttt{char\ n\_sign\_posn}}} & Menentukan posisi \hl{\mbox{\texttt{positive\_sign}}} untuk nilai negatif. \\
\hl{\mbox{\texttt{char\ *int\_curr\_symbol}}} & Simbol mata uang internasional, mis. \hl{\mbox{\texttt{"USD\ "}}}. \\
\hl{\mbox{\texttt{char\ int\_frac\_digits}}} & Nilai internasional untuk \hl{\mbox{\texttt{frac\_digits}}}. \\
\hl{\mbox{\texttt{char\ int\_p\_cs\_precedes}}} & Nilai internasional untuk \hl{\mbox{\texttt{p\_cs\_precedes}}}. \\
\hl{\mbox{\texttt{char\ int\_n\_cs\_precedes}}} & Nilai internasional untuk \hl{\mbox{\texttt{n\_cs\_precedes}}}. \\
\hl{\mbox{\texttt{char\ int\_p\_sep\_by\_space}}} & Nilai internasional untuk \hl{\mbox{\texttt{p\_sep\_by\_space}}}. \\
\hl{\mbox{\texttt{char\ int\_n\_sep\_by\_space}}} & Nilai internasional untuk \hl{\mbox{\texttt{n\_sep\_by\_space}}}. \\
\hl{\mbox{\texttt{char\ int\_p\_sign\_posn}}} & Nilai internasional untuk \hl{\mbox{\texttt{p\_sign\_posn}}}. \\
\hl{\mbox{\texttt{char\ int\_n\_sign\_posn}}} & Nilai internasional untuk \hl{\mbox{\texttt{n\_sign\_posn}}}. \\
\end{longtable}
}

\index{localeconv() function@\texttt{localeconv()} function|)}

\subsection{Pengelompokan Digit Moneter}\label{monetary-digit-grouping}

\index{localeconv() function@\texttt{localeconv()} function!mon\_grouping@\texttt{mon\_grouping}|(}

Oke, yang ini membingungkan. \hl{\mbox{\texttt{mon\_grouping}}} adalah \hl{\mbox{\texttt{char*}}}, jadi mungkin kamu mengira itu string. Tapi dalam kasus ini, bukan, ia sebenarnya array berisi \hl{\mbox{\texttt{char}}}. Ia harus selalu berakhir dengan \hl{\mbox{\texttt{0}}} atau \hl{\mbox{\texttt{CHAR\_MAX}}}.

Nilai-nilai ini menjelaskan cara mengelompokkan deretan angka pada mata uang di sebelah \emph{kiri} desimal (bagian bilangan bulat).

Misalnya, kita mungkin punya:

\begin{tcolorbox}[enhanced jigsaw,left=3pt,breakable,colback=white!97.2549!black,boxrule=0.2mm]

\begin{Shaded}
\begin{Highlighting}[]
\NormalTok{  2   1   0}
\NormalTok{ {-}{-}{-} {-}{-}{-} {-}{-}{-}}
\NormalTok{$100,000,000.00}
\end{Highlighting}
\end{Shaded}

\end{tcolorbox}

Ini kelompok-kelompok tiga. Grup 0 (tepat di kiri desimal) punya 3 digit. Grup 1 (grup berikutnya ke kiri) punya 3 digit, dan yang terakhir juga 3.

Jadi kita bisa menggambarkan grup-grup ini, dari kanan (desimal) ke kiri dengan serangkaian nilai integer yang mewakili ukuran grup:

\begin{tcolorbox}[enhanced jigsaw,left=3pt,breakable,colback=white!97.2549!black,boxrule=0.2mm]

\begin{Shaded}
\begin{Highlighting}[]
\NormalTok{3 3 3}
\end{Highlighting}
\end{Shaded}

\end{tcolorbox}

Dan itu akan berfungsi untuk nilai sampai \$100,000,000.

Tapi bagaimana kalau lebih banyak? Kita bisa terus menambah \hl{\mbox{\texttt{3}}}\ldots{}

\begin{tcolorbox}[enhanced jigsaw,left=3pt,breakable,colback=white!97.2549!black,boxrule=0.2mm]

\begin{Shaded}
\begin{Highlighting}[]
\NormalTok{3 3 3 3 3 3 3 3 3 3 3 3 3 3 3 3 3}
\end{Highlighting}
\end{Shaded}

\end{tcolorbox}

tapi itu gila. Untungnya, kita bisa menentukan \hl{\mbox{\texttt{0}}} untuk menandakan bahwa ukuran grup sebelumnya berulang:

\begin{tcolorbox}[enhanced jigsaw,left=3pt,breakable,colback=white!97.2549!black,boxrule=0.2mm]

\begin{Shaded}
\begin{Highlighting}[]
\NormalTok{3 0}
\end{Highlighting}
\end{Shaded}

\end{tcolorbox}

Yang berarti mengulang setiap 3. Itu akan menangani \$100, \$1,000, \$10,000, \$10,000,000, \$100,000,000,000, dan seterusnya.

Kamu bisa benar-benar menggila dengan ini untuk menandakan pengelompokan yang aneh.

Misalnya:

\begin{tcolorbox}[enhanced jigsaw,left=3pt,breakable,colback=white!97.2549!black,boxrule=0.2mm]

\begin{Shaded}
\begin{Highlighting}[]
\NormalTok{4 3 2 1 0}
\end{Highlighting}
\end{Shaded}

\end{tcolorbox}

akan menandakan:

\begin{tcolorbox}[enhanced jigsaw,left=3pt,breakable,colback=white!97.2549!black,boxrule=0.2mm]

\begin{Shaded}
\begin{Highlighting}[]
\NormalTok{$1,0,0,0,0,00,000,0000.00}
\end{Highlighting}
\end{Shaded}

\end{tcolorbox}

Satu nilai lagi yang bisa muncul adalah \hl{\mbox{\texttt{CHAR\_MAX}}}. Ini menandakan bahwa tidak ada pengelompokan lagi, dan bisa muncul di mana saja dalam array, termasuk di nilai pertama.

\begin{tcolorbox}[enhanced jigsaw,left=3pt,breakable,colback=white!97.2549!black,boxrule=0.2mm]

\begin{Shaded}
\begin{Highlighting}[]
\NormalTok{3 2 CHAR\_MAX}
\end{Highlighting}
\end{Shaded}

\end{tcolorbox}

akan menandakan:

\begin{tcolorbox}[enhanced jigsaw,left=3pt,breakable,colback=white!97.2549!black,boxrule=0.2mm]

\begin{Shaded}
\begin{Highlighting}[]
\NormalTok{100000000,00,000.00}
\end{Highlighting}
\end{Shaded}

\end{tcolorbox}

misalnya.

Dan hanya dengan adanya \hl{\mbox{\texttt{CHAR\_MAX}}} di posisi pertama array, itu memberi tahu bahwa tidak ada pengelompokan sama sekali.

\index{localeconv() function@\texttt{localeconv()} function!mon\_grouping@\texttt{mon\_grouping}|)}

\subsection{Pemisah dan Posisi Tanda}\label{pemisah-dan-posisi-tanda}

\index{localeconv() function@\texttt{localeconv()} function!sep\_by\_space@\texttt{sep\_by\_space}|(}

Semua varian \hl{\mbox{\texttt{sep\_by\_space}}} berkaitan dengan spasi di sekitar tanda mata uang. Nilai yang valid:

{\def\LTcaptype{none} % do not increment counter
\begin{longtable}[]{@{}
  >{\centering\arraybackslash}p{(\linewidth - 2\tabcolsep) * \real{0.2500}}
  >{\raggedright\arraybackslash}p{(\linewidth - 2\tabcolsep) * \real{0.7500}}@{}}
\toprule\noalign{}
\begin{minipage}[b]{\linewidth}\centering
Nilai
\end{minipage} & \begin{minipage}[b]{\linewidth}\raggedright
Deskripsi
\end{minipage} \\
\midrule\noalign{}
\endhead
\bottomrule\noalign{}
\endlastfoot
\hl{\mbox{\texttt{0}}} & Tidak ada spasi antara simbol mata uang dan nilai. \\
\hl{\mbox{\texttt{1}}} & Pisahkan simbol mata uang (dan tanda, jika ada) dari nilai dengan spasi. \\
\hl{\mbox{\texttt{2}}} & Pisahkan simbol tanda dari simbol mata uang (jika bersebelahan) dengan spasi, selain itu pisahkan simbol tanda dari nilai dengan spasi. \\
\end{longtable}
}

Varian \hl{\mbox{\texttt{sign\_posn}}} ditentukan oleh nilai-nilai berikut:

{\def\LTcaptype{none} % do not increment counter
\begin{longtable}[]{@{}cl@{}}
\toprule\noalign{}
Nilai & Deskripsi \\
\midrule\noalign{}
\endhead
\bottomrule\noalign{}
\endlastfoot
\hl{\mbox{\texttt{0}}} & Taruh tanda kurung di sekitar nilai dan simbol mata uang. \\
\hl{\mbox{\texttt{1}}} & Taruh string tanda di depan simbol mata uang dan nilai. \\
\hl{\mbox{\texttt{2}}} & Taruh string tanda setelah simbol mata uang dan nilai. \\
\hl{\mbox{\texttt{3}}} & Taruh string tanda tepat di depan simbol mata uang. \\
\hl{\mbox{\texttt{4}}} & Taruh string tanda tepat di belakang simbol mata uang. \\
\end{longtable}
}

\index{localeconv() function@\texttt{localeconv()} function!sep\_by\_space@\texttt{sep\_by\_space}|)} \index{Locale!money|)}

\subsection{Contoh Nilai}\label{contoh-nilai}

Waktu kuambil nilai-nilainya di sistemku, inilah yang kulihat (string pengelompokan ditampilkan sebagai nilai byte individual):

\begin{tcolorbox}[enhanced jigsaw,left=3pt,breakable,colback=white!97.2549!black,boxrule=0.2mm]

\begin{Shaded}
\begin{Highlighting}[]
\NormalTok{mon\_decimal\_point  }\OperatorTok{=} \StringTok{"."}
\NormalTok{mon\_thousands\_sep  }\OperatorTok{=} \StringTok{","}
\NormalTok{mon\_grouping       }\OperatorTok{=} \DecValTok{3} \DecValTok{3} \DecValTok{0}
\NormalTok{positive\_sign      }\OperatorTok{=} \StringTok{""}
\NormalTok{negative\_sign      }\OperatorTok{=} \StringTok{"{-}"}
\NormalTok{currency\_symbol    }\OperatorTok{=} \StringTok{"$"}
\NormalTok{frac\_digits        }\OperatorTok{=} \DecValTok{2}
\NormalTok{p\_cs\_precedes      }\OperatorTok{=} \DecValTok{1}
\NormalTok{n\_cs\_precedes      }\OperatorTok{=} \DecValTok{1}
\NormalTok{p\_sep\_by\_space     }\OperatorTok{=} \DecValTok{0}
\NormalTok{n\_sep\_by\_space     }\OperatorTok{=} \DecValTok{0}
\NormalTok{p\_sign\_posn        }\OperatorTok{=} \DecValTok{1}
\NormalTok{n\_sign\_posn        }\OperatorTok{=} \DecValTok{1}
\NormalTok{int\_curr\_symbol    }\OperatorTok{=} \StringTok{"USD "}
\NormalTok{int\_frac\_digits    }\OperatorTok{=} \DecValTok{2}
\NormalTok{int\_p\_cs\_precedes  }\OperatorTok{=} \DecValTok{1}
\NormalTok{int\_n\_cs\_precedes  }\OperatorTok{=} \DecValTok{1}
\NormalTok{int\_p\_sep\_by\_space }\OperatorTok{=} \DecValTok{1}
\NormalTok{int\_n\_sep\_by\_space }\OperatorTok{=} \DecValTok{1}
\NormalTok{int\_p\_sign\_posn    }\OperatorTok{=} \DecValTok{1}
\NormalTok{int\_n\_sign\_posn    }\OperatorTok{=} \DecValTok{1}
\end{Highlighting}
\end{Shaded}

\end{tcolorbox}

\section{Spesifik Lokalisasi}\label{spesifik-lokalisasi}

Perhatikan bagaimana tadi kita mengoper makro \hl{\mbox{\texttt{LC\_ALL}}} ke \hl{\mbox{\texttt{setlocale()}}}\ldots{} ini mengisyaratkan bahwa mungkin ada varian yang memungkinkan kamu lebih presisi soal \emph{bagian} locale mana yang diatur.

Mari kita lihat nilai-nilai yang bisa dipakai:

\index{setlocale() function@\texttt{setlocale()} function!LC\_ALL macro@\texttt{LC\_ALL} macro} \index{setlocale() function@\texttt{setlocale()} function!LC\_COLLATE macro@\texttt{LC\_COLLATE} macro} \index{setlocale() function@\texttt{setlocale()} function!LC\_CTYPE macro@\texttt{LC\_CTYPE} macro} \index{setlocale() function@\texttt{setlocale()} function!LC\_MONETARY macro@\texttt{LC\_MONETARY} macro} \index{setlocale() function@\texttt{setlocale()} function!LC\_NUMERIC macro@\texttt{LC\_NUMERIC} macro} \index{setlocale() function@\texttt{setlocale()} function!LC\_TIME macro@\texttt{LC\_TIME} macro}

{\def\LTcaptype{none} % do not increment counter
\begin{longtable}[]{@{}
  >{\raggedright\arraybackslash}p{(\linewidth - 2\tabcolsep) * \real{0.2222}}
  >{\raggedright\arraybackslash}p{(\linewidth - 2\tabcolsep) * \real{0.7778}}@{}}
\toprule\noalign{}
\begin{minipage}[b]{\linewidth}\raggedright
Makro
\end{minipage} & \begin{minipage}[b]{\linewidth}\raggedright
Deskripsi
\end{minipage} \\
\midrule\noalign{}
\endhead
\bottomrule\noalign{}
\endlastfoot
\hl{\mbox{\texttt{LC\_ALL}}} & Atur semua yang berikut ke locale yang diberikan. \\
\hl{\mbox{\texttt{LC\_COLLATE}}} & Mengatur perilaku fungsi \hl{\mbox{\texttt{strcoll()}}} dan \hl{\mbox{\texttt{strxfrm()}}}. \\
\hl{\mbox{\texttt{LC\_CTYPE}}} & Mengatur perilaku fungsi penanganan karakter\footnote{Kecuali \hl{\mbox{\texttt{isdigit()}}} dan \hl{\mbox{\texttt{isxdigit()}}}.}. \\
\hl{\mbox{\texttt{LC\_MONETARY}}} & Mengatur nilai yang dikembalikan \hl{\mbox{\texttt{localeconv()}}}. \\
\hl{\mbox{\texttt{LC\_NUMERIC}}} & Mengatur titik desimal untuk keluarga fungsi \hl{\mbox{\texttt{printf()}}}. \\
\hl{\mbox{\texttt{LC\_TIME}}} & Mengatur format waktu pada fungsi pencetak waktu dan tanggal \hl{\mbox{\texttt{strftime()}}} dan \hl{\mbox{\texttt{wcsftime()}}}. \\
\end{longtable}
}

Cukup umum melihat \index{setlocale() function@\texttt{setlocale()} function!LC\_ALL macro@\texttt{LC\_ALL} macro} \hl{\mbox{\texttt{LC\_ALL}}} yang diatur, tapi, setidaknya kamu punya pilihan.

Aku juga perlu menunjukkan bahwa \index{setlocale() function@\texttt{setlocale()} function!LC\_CTYPE macro@\texttt{LC\_CTYPE} macro} \hl{\mbox{\texttt{LC\_CTYPE}}} adalah salah satu yang besar karena terkait dengan wide character, kaleng cacing yang signifikan yang akan kita bahas nanti.

\index{Locale|)}

\chapter{Unicode, Wide Character, dan Semacamnya}\label{unicode-wide-character-dan-semacamnya}

\index{Unicode|(}

Sebelum mulai, perlu diketahui bahwa ini area yang masih aktif berkembang di C, karena bahasa ini sedang berusaha melewati, ehm, \emph{sakit pertumbuhan}. Sekarang C23 sudah terbit, pembaruan di sini kemungkinan akan ada.

Kebanyakan orang pada dasarnya hanya tertarik pada pertanyaan yang tampak sederhana, ``Bagaimana cara memakai character set ini-itu di C?'' Kita akan sampai ke sana. Tapi seperti yang akan kita lihat, mungkin itu sudah berfungsi di sistemmu. Atau kamu mungkin harus menyerahkannya ke library pihak ketiga.

Kita akan membahas banyak hal di bab ini---sebagian tidak bergantung platform, dan sebagian khusus C.

Mari kita buat garis besarnya dulu:

\begin{itemize}
\tightlist
\item
  Latar belakang Unicode
\item
  Latar belakang character encoding
\item
  Source dan Execution Character Set
\item
  Memakai Unicode dan UTF-8
\item
  Memakai tipe karakter lain seperti \hl{\mbox{\texttt{wchar\_t}}}, \hl{\mbox{\texttt{char16\_t}}}, dan \hl{\mbox{\texttt{char32\_t}}}
\end{itemize}

Kita mulai menyelam!

\section{Apa Itu Unicode?}\label{apa-itu-unicode}

Dulu, di AS dan sebagian besar dunia, populer memakai encoding 7-bit atau 8-bit untuk karakter di memori. Artinya kita hanya bisa punya total 128 atau 256 karakter (termasuk karakter yang tidak bisa dicetak). Itu cukup untuk dunia yang berpusat pada AS, tapi ternyata ada alfabet lain di luar sana---siapa sangka? Bahasa Mandarin punya lebih dari 50.000 karakter, dan itu tidak muat dalam satu byte.

Maka orang-orang menciptakan berbagai cara alternatif untuk merepresentasikan character set kustom mereka sendiri. Dan itu tidak masalah, tapi berubah menjadi mimpi buruk kompatibilitas.

Untuk lolos darinya, Unicode diciptakan. Satu character set untuk menguasai semuanya. Ia memanjang sampai tak terhingga (secara efektif) sehingga kita tidak akan pernah kehabisan ruang untuk karakter baru. Ia punya huruf Mandarin, Latin, Yunani, aksara paku, simbol catur, emoji\ldots{} hampir semuanya, sungguh! Dan terus bertambah dari waktu ke waktu!

\section{Code Point}\label{code-point}

\index{Unicode!code points|(}

Aku ingin membahas dua konsep di sini. Membingungkan karena keduanya berupa angka\ldots{} angka berbeda untuk hal yang sama. Tapi sabar dulu ya.

Mari kita definisikan \emph{code point} secara longgar sebagai nilai numerik yang mewakili sebuah karakter. (Code point juga bisa mewakili karakter kontrol yang tidak bisa dicetak, tapi anggap saja yang kumaksud adalah sesuatu seperti huruf ``B'' atau karakter ``π''.)

Setiap code point mewakili satu karakter unik. Dan setiap karakter punya code point numerik unik yang terkait dengannya.

Misalnya, di Unicode, nilai numerik 66 mewakili ``B'', dan 960 mewakili ``π''. Pemetaan karakter lain yang bukan Unicode mungkin memakai nilai berbeda, tapi lupakan saja dan fokus pada Unicode, masa depan!

Jadi itu satu hal: ada angka yang mewakili setiap karakter. Di Unicode, angka-angka ini berkisar dari 0 sampai lebih dari 1 juta.

\index{Unicode!code points|)}

Paham?

Karena kita akan sedikit membalik meja.

\section{Encoding}\label{encoding}

\index{Unicode!encoding|(}

Kalau kamu ingat, byte 8-bit bisa menampung nilai 0-255, inklusif. Itu bagus untuk ``B'' yang bernilai 66---muat dalam satu byte. Tapi ``π'' bernilai 960, dan itu tidak muat dalam satu byte! Kita butuh byte lain. Bagaimana menyimpan semua itu di memori? Atau bagaimana dengan angka yang lebih besar, seperti 195.024? Itu butuh beberapa byte untuk menampungnya.

Pertanyaan Besarnya: bagaimana angka-angka ini direpresentasikan di memori? Inilah yang kita sebut \emph{encoding} karakter.

Jadi kita punya dua hal: yang satu adalah code point yang secara efektif memberi tahu nomor seri sebuah karakter. Dan satu lagi encoding yang memberi tahu bagaimana kita merepresentasikan angka itu di memori.

Ada banyak encoding. Kamu bisa membuat sendiri sekarang juga, kalau mau\footnote{Misalnya, kita bisa menyimpan code point dalam integer 32-bit big-endian. Mudah! Kita baru saja menciptakan sebuah encoding! Sebenarnya tidak; itulah yang disebut encoding UTF-32BE. Ya sudahlah---kembali bekerja!}. Tapi kita akan melihat beberapa encoding yang benar-benar umum dan dipakai bersama Unicode.

\index{Unicode!UTF-8|(} \index{Unicode!UTF-16|(} \index{Unicode!UTF-32|(}

{\def\LTcaptype{none} % do not increment counter
\begin{longtable}[]{@{}
  >{\centering\arraybackslash}p{(\linewidth - 2\tabcolsep) * \real{0.2317}}
  >{\raggedright\arraybackslash}p{(\linewidth - 2\tabcolsep) * \real{0.7683}}@{}}
\toprule\noalign{}
\begin{minipage}[b]{\linewidth}\centering
Encoding
\end{minipage} & \begin{minipage}[b]{\linewidth}\raggedright
Deskripsi
\end{minipage} \\
\midrule\noalign{}
\endhead
\bottomrule\noalign{}
\endlastfoot
UTF-8 & Encoding berorientasi byte yang memakai jumlah byte bervariasi per karakter. Inilah yang sebaiknya dipakai. \\
UTF-16 & Encoding 16-bit per karakter\footnote{Kurang lebih. Secara teknis, lebarnya bervariasi---ada cara untuk merepresentasikan code point yang lebih tinggi dari \(2^{16}\) dengan menggabungkan dua karakter UTF-16.}. \\
UTF-32 & Encoding 32-bit per karakter. \\
\end{longtable}
}

Pada UTF-16 dan UTF-32, urutan byte penting, jadi kamu mungkin melihat UTF-16BE untuk big-endian dan UTF-16LE untuk little-endian. Sama untuk UTF-32. Secara teknis, kalau tidak disebutkan, kamu harus mengasumsikan big-endian. Tapi karena Windows memakai UTF-16 secara luas dan little-endian, kadang itulah yang diasumsikan\footnote{Ada karakter khusus bernama \emph{Byte Order Mark} (BOM), code point 0xFEFF, yang boleh mendahului aliran data dan menandakan endianess-nya. Tapi ini tidak wajib.}.

Mari kita lihat beberapa contoh. Aku akan menulis nilainya dalam hex karena itu tepat dua digit per byte 8-bit, dan memudahkan kita melihat bagaimana susunannya di memori.

{\def\LTcaptype{none} % do not increment counter
\begin{longtable}[]{@{}cccccccc@{}}
\toprule\noalign{}
Karakter & Code Point & UTF-16BE & UTF-32BE & UTF-16LE & UTF-32LE & UTF-8 & \\
\midrule\noalign{}
\endhead
\bottomrule\noalign{}
\endlastfoot
\hl{\mbox{\texttt{A}}} & 41 & 0041 & 00000041 & 4100 & 41000000 & 41 & \\
\hl{\mbox{\texttt{B}}} & 42 & 0042 & 00000042 & 4200 & 42000000 & 42 & \\
\hl{\mbox{\texttt{\textasciitilde{}}}} & 7E & 007E & 0000007E & 7E00 & 7E000000 & 7E & \\
\hl{\mbox{\texttt{π}}} & 3C0 & 03C0 & 000003C0 & C003 & C0030000 & CF80 & \\
\hl{\mbox{\texttt{€}}} & 20AC & 20AC & 000020AC & AC20 & AC200000 & E282AC & \\
\end{longtable}
}

\index{Unicode!endianess|(}

Cari pola di dalamnya. Perhatikan bahwa UTF-16BE dan UTF-32BE hanyalah code point yang direpresentasikan langsung sebagai nilai 16- dan 32-bit\footnote{Sekali lagi, ini hanya benar di UTF-16 untuk karakter yang muat dalam dua byte.}.

\index{Unicode!UTF-16|)} \index{Unicode!UTF-32|)}

Little-endian sama saja, hanya saja urutan byte-nya little-endian.

\index{Unicode!endianess|)}

Lalu ada UTF-8 di bagian akhir. Pertama, mungkin kamu memperhatikan bahwa code point satu byte direpresentasikan sebagai satu byte. Itu bagus. Kamu mungkin juga memperhatikan bahwa code point yang berbeda memakai jumlah byte yang berbeda. Inilah encoding dengan lebar bervariasi.

Jadi begitu kita melewati nilai tertentu, UTF-8 mulai memakai byte tambahan untuk menyimpan nilainya. Dan byte-byte itu tampaknya tidak berkorelasi dengan nilai code point.

\href{https://en.wikipedia.org/wiki/UTF-8}{Detail encoding UTF-8}\footnote{https://en.wikipedia.org/wiki/UTF-8} di luar cakupan panduan ini, tapi cukup diketahui bahwa ia memakai jumlah byte yang bervariasi per code point, dan nilai byte itu tidak cocok dengan code point \emph{kecuali untuk 128 code point pertama}. Kalau kamu benar-benar ingin belajar lebih jauh, \href{https://www.youtube.com/watch?v=MijmeoH9LT4}{Computerphile punya video UTF-8 yang bagus bersama Tom Scott}\footnote{https://www.youtube.com/watch?v=MijmeoH9LT4}.

Soal terakhir itu hal yang keren tentang Unicode dan UTF-8 dari sudut pandang Amerika Utara: ia kompatibel mundur dengan encoding ASCII 7-bit! Jadi kalau kamu terbiasa dengan ASCII, UTF-8 sama saja! Setiap dokumen ber-encoding ASCII juga ber-encoding UTF-8! (Tapi tidak sebaliknya, tentu saja.)

Kemungkinan poin terakhir itulah, lebih dari yang lain, yang mendorong UTF-8 menguasai dunia.

\index{Unicode!UTF-8|)} \index{Unicode!encoding|)}

\section{Source dan Execution Character Set}\label{src-exec-charset}

\index{Character sets|(}

Saat memprogram di C, ada (setidaknya) tiga character set yang berperan:

\begin{itemize}
\tightlist
\item
  Yang dipakai kodemu saat tersimpan di disk.
\item
  \index{Character sets!source}Yang diterjemahkan compiler dari itu tepat saat kompilasi dimulai (\emph{source character set}). Ini mungkin sama dengan yang di disk, mungkin juga tidak.
\item
  \index{Character sets!execution}Yang menjadi tujuan terjemahan source character set oleh compiler untuk eksekusi (\emph{execution character set}). Ini mungkin sama dengan source character set, mungkin juga tidak.
\end{itemize}

Compiler-mu kemungkinan punya opsi untuk memilih character set ini saat build.

\index{Character sets!basic|(}

Basic character set untuk source maupun execution berisi karakter-karakter berikut:

\begin{tcolorbox}[enhanced jigsaw,left=3pt,breakable,colback=white!97.2549!black,boxrule=0.2mm]

\begin{Shaded}
\begin{Highlighting}[]
\NormalTok{A B C D E F G H I J K L M}
\NormalTok{N O P Q R S T U V W X Y Z}
\NormalTok{a b c d e f g h i j k l m}
\NormalTok{n o p q r s t u v w x y z}
\NormalTok{0 1 2 3 4 5 6 7 8 9}
\NormalTok{! " \# \% \& \textquotesingle{} ( ) * + , {-} . / :}
\NormalTok{; \textless{} = \textgreater{} ? [ \textbackslash{} ] \^{} \_ \{ | \} \textasciitilde{}}
\NormalTok{space tab vertical{-}tab}
\NormalTok{form{-}feed end{-}of{-}line}
\end{Highlighting}
\end{Shaded}

\end{tcolorbox}

Itulah karakter yang bisa kamu pakai di source dan tetap 100\% portabel.

Execution character set akan punya tambahan karakter untuk alert (bell/flash), backspace, carriage return, dan newline.

Tapi kebanyakan orang tidak sampai seekstrem itu dan dengan bebas memakai extended character set di source dan executable, terutama sekarang ketika Unicode dan UTF-8 makin umum. Maksudnya, basic character set bahkan tidak mengizinkan \hl{\mbox{\texttt{@}}}, \hl{\mbox{\texttt{\$}}}, atau \hl{\mbox{\texttt{\textasciigrave{}}}}!

Perlu dicatat, menulis karakter Unicode hanya dengan basic character set itu merepotkan (walaupun bisa dengan escape sequence).

\index{Character sets!basic|)} \index{Character sets|)}

\section{Unicode di C}\label{unicode-in-c}

Sebelum membahas encoding di C, mari kita bahas Unicode dari sudut pandang code point dulu. Ada cara di C untuk menentukan karakter Unicode dan karakter-karakter ini akan diterjemahkan compiler ke execution character set\footnote{Mungkin compiler berusaha sebaik mungkin menerjemahkan code point ke encoding output apa pun itu, tapi aku tidak menemukan jaminannya di spesifikasi.}.

Jadi bagaimana caranya?

Bagaimana dengan simbol euro, code point 0x20AC. (Ditulis dalam hex karena kedua cara merepresentasikannya di C membutuhkan hex.) Bagaimana kita bisa memasukkannya ke kode C kita?

\index{{\backslash}u Unicode escape@\texttt{{\backslash}u} Unicode escape|(} Pakai escape \hl{\mbox{\texttt{\textbackslash{}u}}} untuk memasukkannya ke string, mis. \hl{\mbox{\texttt{"\textbackslash{}u20AC"}}} (huruf besar/kecil pada hex tidak berpengaruh). Kamu harus menaruh \textbf{tepat empat} digit hex setelah \hl{\mbox{\texttt{\textbackslash{}u}}}, dengan nol di depan kalau perlu.

Berikut contohnya:

\begin{tcolorbox}[enhanced jigsaw,left=3pt,breakable,colback=white!97.2549!black,boxrule=0.2mm]

\begin{Shaded}
\begin{Highlighting}[]
\DataTypeTok{char} \OperatorTok{*}\NormalTok{s }\OperatorTok{=} \StringTok{"}\SpecialCharTok{\textbackslash{}u20AC}\StringTok{1.23"}\OperatorTok{;}

\NormalTok{printf}\OperatorTok{(}\StringTok{"}\SpecialCharTok{\%s\textbackslash{}n}\StringTok{"}\OperatorTok{,}\NormalTok{ s}\OperatorTok{);}  \CommentTok{// €1.23}
\end{Highlighting}
\end{Shaded}

\end{tcolorbox}

\index{{\backslash}U Unicode escape@\texttt{{\backslash}U} Unicode escape|(}

Jadi \hl{\mbox{\texttt{\textbackslash{}u}}} berfungsi untuk code point Unicode 16-bit, tapi bagaimana dengan yang lebih besar dari 16 bit? Untuk itu kita butuh huruf kapital: \hl{\mbox{\texttt{\textbackslash{}U}}}.

Misalnya:

\begin{tcolorbox}[enhanced jigsaw,left=3pt,breakable,colback=white!97.2549!black,boxrule=0.2mm]

\begin{Shaded}
\begin{Highlighting}[]
\DataTypeTok{char} \OperatorTok{*}\NormalTok{s }\OperatorTok{=} \StringTok{"}\SpecialCharTok{\textbackslash{}U0001D4D1}\StringTok{"}\OperatorTok{;}

\NormalTok{printf}\OperatorTok{(}\StringTok{"}\SpecialCharTok{\%s\textbackslash{}n}\StringTok{"}\OperatorTok{,}\NormalTok{ s}\OperatorTok{);}  \CommentTok{// Mencetak huruf matematika "B"}
\end{Highlighting}
\end{Shaded}

\end{tcolorbox}

Sama seperti \hl{\mbox{\texttt{\textbackslash{}u}}}, hanya dengan 32 bit alih-alih 16. Keduanya setara:

\begin{tcolorbox}[enhanced jigsaw,left=3pt,breakable,colback=white!97.2549!black,boxrule=0.2mm]

\begin{Shaded}
\begin{Highlighting}[]
\NormalTok{\textbackslash{}u03C0}
\NormalTok{\textbackslash{}U000003C0}
\end{Highlighting}
\end{Shaded}

\end{tcolorbox}

\index{{\backslash}u Unicode escape@\texttt{{\backslash}u} Unicode escape|)} \index{{\backslash}U Unicode escape@\texttt{{\backslash}U} Unicode escape|)}

Sekali lagi, ini diterjemahkan ke \index{Character sets!execution}execution character set saat kompilasi. Mereka merepresentasikan code point Unicode, bukan encoding tertentu. Selain itu, kalau sebuah code point Unicode tidak bisa direpresentasikan di execution character set, compiler boleh melakukan apa saja terhadapnya.

Nah, mungkin kamu bertanya kenapa tidak langsung begini:

\begin{tcolorbox}[enhanced jigsaw,left=3pt,breakable,colback=white!97.2549!black,boxrule=0.2mm]

\begin{Shaded}
\begin{Highlighting}[]
\DataTypeTok{char} \OperatorTok{*}\NormalTok{s }\OperatorTok{=} \StringTok{"€1.23"}\OperatorTok{;}

\NormalTok{printf}\OperatorTok{(}\StringTok{"}\SpecialCharTok{\%s\textbackslash{}n}\StringTok{"}\OperatorTok{,}\NormalTok{ s}\OperatorTok{);}  \CommentTok{// €1.23}
\end{Highlighting}
\end{Shaded}

\end{tcolorbox}

Dan kemungkinan besar bisa, dengan compiler modern. \index{Character sets!source}Source character set akan diterjemahkan oleh compiler ke \index{Character sets!execution}execution character set untukmu. Tapi compiler bebas muntah kalau menemukan karakter yang tidak termasuk dalam extended character set-nya, dan simbol € jelas tidak ada di \index{Character sets!basic}basic character set.

\index{{\backslash}u Unicode escape@\texttt{{\backslash}u} Unicode escape|(} \index{{\backslash}U Unicode escape@\texttt{{\backslash}U} Unicode escape|(}

Peringatan dari spesifikasi: kamu tidak boleh memakai \hl{\mbox{\texttt{\textbackslash{}u}}} atau \hl{\mbox{\texttt{\textbackslash{}U}}} untuk meng-encode code point di bawah 0xA0 kecuali 0x24 (\hl{\mbox{\texttt{\$}}}), 0x40 (\hl{\mbox{\texttt{@}}}), dan 0x60 (\hl{\mbox{\texttt{\textasciigrave{}}}})---ya, itulah tiga serangkai tanda baca umum yang tidak ada di basic character set. Tampaknya pembatasan ini dilonggarkan di versi spesifikasi yang akan datang.

\index{{\backslash}u Unicode escape@\texttt{{\backslash}u} Unicode escape|)} \index{{\backslash}U Unicode escape@\texttt{{\backslash}U} Unicode escape|)}

Terakhir, kamu juga bisa memakai ini di identifier dalam kodemu, dengan beberapa batasan. Tapi aku tidak ingin membahasnya di sini. Bab ini fokus pada penanganan string.

Dan itu saja soal Unicode di C (kecuali encoding).

\section{Catatan Singkat tentang UTF-8 Sebelum Kita Terjun ke Semak-Semak}\label{utf8-quick}

\index{Unicode!UTF-8|(}

Bisa jadi file source di disk, extended source character, dan extended execution character-mu semuanya berformat UTF-8. Dan library yang kamu pakai mengharapkan UTF-8. Inilah masa depan gemilang dengan UTF-8 di mana-mana.

Kalau begitu, dan kamu tidak keberatan tidak portabel ke sistem yang tidak seperti itu, jalankan saja. Taruh karakter Unicode di source dan datamu sesukamu. Pakai string C biasa dan berbahagialah.

Banyak hal akan langsung berfungsi (walaupun tidak portabel) karena string UTF-8 aman diakhiri NUL seperti string C lainnya. Tapi mungkin kehilangan portabilitas demi penanganan karakter yang lebih mudah adalah pertukaran yang sepadan bagimu.

Meski begitu, ada beberapa peringatan:

\begin{itemize}
\item
  Fungsi seperti \hl{\mbox{\texttt{strlen()}}} melaporkan jumlah byte dalam string, belum tentu jumlah karakternya. (Fungsi \index{mbstowcs() function@\texttt{mbstowcs()} function!with UTF-8}\hl{\mbox{\texttt{mbstowcs()}}} mengembalikan jumlah karakter dalam string ketika kamu mengonversinya menjadi wide character. POSIX memperluas ini sehingga kamu bisa mengoper \hl{\mbox{\texttt{NULL}}} sebagai argumen pertama kalau hanya ingin jumlah karakter.)
\item
  Yang berikut tidak akan bekerja dengan benar pada karakter lebih dari satu byte: \index{strtok() function@\texttt{strtok()} function!with UTF-8}\hl{\mbox{\texttt{strtok()}}}, \index{strchr() function@\texttt{strchr()} function!with UTF-8}\hl{\mbox{\texttt{strchr()}}} (pakai \index{strstr() function@\texttt{strstr()} function!with UTF-8}\hl{\mbox{\texttt{strstr()}}} sebagai gantinya), fungsi tipe \hl{\mbox{\texttt{strspn()}}}, \index{toupper() function@\texttt{toupper()} function!with UTF-8}\hl{\mbox{\texttt{toupper()}}}, \index{tolower() function@\texttt{tolower()} function!with UTF-8}\hl{\mbox{\texttt{tolower()}}}, fungsi tipe \index{isalpha() function@\texttt{isalpha()} function!with UTF-8}\hl{\mbox{\texttt{isalpha()}}}, dan mungkin lebih banyak lagi. Waspadai apa pun yang beroperasi pada byte.
\item
  Varian \index{printf() function@\texttt{printf()} function!with UTF-8}\hl{\mbox{\texttt{printf()}}} menyediakan cara untuk hanya mencetak sebagian byte dari sebuah string\footnote{Dengan format specifier seperti \hl{\mbox{\texttt{"\%.12s"}}}, misalnya.}. Pastikan kamu mencetak jumlah byte yang benar agar berakhir di batas karakter.
\item
  \index{malloc() function@\texttt{malloc()} function!with UTF-8}Kalau kamu ingin \hl{\mbox{\texttt{malloc()}}} ruang untuk sebuah string, atau mendeklarasikan array \hl{\mbox{\texttt{char}}} untuknya, ketahuilah bahwa ukuran maksimumnya bisa lebih besar dari yang kamu duga. Setiap karakter bisa memakai sampai \index{MB\_LEN\_MAX macro@\texttt{MB\_LEN\_MAX} macro}\hl{\mbox{\texttt{MB\_LEN\_MAX}}} byte (dari \hl{\mbox{\texttt{\textless{}limits.h\textgreater{}}}})---kecuali karakter di basic character set yang dijamin satu byte.
\end{itemize}

Dan mungkin masih ada yang lain yang belum kutemukan. Kabari aku kalau ada jebakan lain di luar sana\ldots{}

\index{Unicode!UTF-8|)}

\section{Berbagai Tipe Karakter}\label{berbagai-tipe-karakter}

Aku ingin memperkenalkan lebih banyak tipe karakter. Kita sudah terbiasa dengan \hl{\mbox{\texttt{char}}}, kan?

Tapi itu terlalu mudah. Mari kita buat semuanya jauh lebih sulit! Hore!

\subsection{Karakter Multibyte}\label{karakter-multibyte}

\index{Multibyte characters|(}

Pertama-tama, aku ingin mengubah cara berpikirmu tentang apa itu string (array \hl{\mbox{\texttt{char}}}). Ini adalah \emph{string multibyte} yang terdiri dari \emph{karakter multibyte}.

Betul---string karakter biasa yang kamu kenal itu multibyte. Ketika seseorang bilang ``string C'', maksudnya ``string multibyte C''.

Bahkan kalau sebuah karakter dalam string hanya satu byte, atau kalau sebuah string hanya terdiri dari karakter tunggal, ia tetap dikenal sebagai string multibyte.

Misalnya:

\begin{tcolorbox}[enhanced jigsaw,left=3pt,breakable,colback=white!97.2549!black,boxrule=0.2mm]

\begin{Shaded}
\begin{Highlighting}[]
\DataTypeTok{char}\NormalTok{ c}\OperatorTok{[}\DecValTok{128}\OperatorTok{]} \OperatorTok{=} \StringTok{"Hello, world!"}\OperatorTok{;}  \CommentTok{// String multibyte}
\end{Highlighting}
\end{Shaded}

\end{tcolorbox}

Maksudnya, sebuah karakter yang tidak ada di basic character set bisa tersusun dari beberapa byte. Sampai sebanyak \index{MB\_LEN\_MAX macro@\texttt{MB\_LEN\_MAX} macro}\hl{\mbox{\texttt{MB\_LEN\_MAX}}} (dari \hl{\mbox{\texttt{\textless{}limits.h\textgreater{}}}}). Memang di layar hanya tampak seperti satu karakter, tapi ia bisa berupa beberapa byte.

Nilai Unicode juga bisa dimasukkan ke sana, seperti yang kita lihat tadi:

\begin{tcolorbox}[enhanced jigsaw,left=3pt,breakable,colback=white!97.2549!black,boxrule=0.2mm]

\begin{Shaded}
\begin{Highlighting}[]
\DataTypeTok{char} \OperatorTok{*}\NormalTok{s }\OperatorTok{=} \StringTok{"}\SpecialCharTok{\textbackslash{}u20AC}\StringTok{1.23"}\OperatorTok{;}

\NormalTok{printf}\OperatorTok{(}\StringTok{"}\SpecialCharTok{\%s\textbackslash{}n}\StringTok{"}\OperatorTok{,}\NormalTok{ s}\OperatorTok{);}  \CommentTok{// €1.23}
\end{Highlighting}
\end{Shaded}

\end{tcolorbox}

Tapi di sini kita mulai masuk ke keanehan, karena lihat ini:

\index{strlen() function@\texttt{strlen()} function!with UTF-8|(}

\begin{tcolorbox}[enhanced jigsaw,left=3pt,breakable,colback=white!97.2549!black,boxrule=0.2mm]

\begin{Shaded}
\begin{Highlighting}[]
\DataTypeTok{char} \OperatorTok{*}\NormalTok{s }\OperatorTok{=} \StringTok{"}\SpecialCharTok{\textbackslash{}u20AC}\StringTok{1.23"}\OperatorTok{;}  \CommentTok{// €1.23}

\NormalTok{printf}\OperatorTok{(}\StringTok{"}\SpecialCharTok{\%zu\textbackslash{}n}\StringTok{"}\OperatorTok{,}\NormalTok{ strlen}\OperatorTok{(}\NormalTok{s}\OperatorTok{));}  \CommentTok{// 7!}
\end{Highlighting}
\end{Shaded}

\end{tcolorbox}

Panjang string \hl{\mbox{\texttt{"€1.23"}}} adalah \hl{\mbox{\texttt{7}}}?! Ya! Setidaknya di sistemku, ya! Ingat bahwa \hl{\mbox{\texttt{strlen()}}} mengembalikan jumlah byte dalam string, bukan jumlah karakter. (Saat membahas ``wide character'' nanti, kita akan melihat cara mendapat jumlah karakter dalam string.)

\index{strlen() function@\texttt{strlen()} function!with UTF-8|)}

Perhatikan bahwa walaupun C mengizinkan konstanta \hl{\mbox{\texttt{char}}} multibyte individual (berbeda dari \hl{\mbox{\texttt{char*}}}), perilakunya bervariasi menurut implementasi dan compiler-mu mungkin memperingatkannya.

GCC, misalnya, memperingatkan soal multi-character character constant pada dua baris berikut (dan, di sistemku, mencetak encoding UTF-8-nya):

\begin{tcolorbox}[enhanced jigsaw,left=3pt,breakable,colback=white!97.2549!black,boxrule=0.2mm]

\begin{Shaded}
\begin{Highlighting}[]
\NormalTok{printf}\OperatorTok{(}\StringTok{"}\SpecialCharTok{\%x\textbackslash{}n}\StringTok{"}\OperatorTok{,} \CharTok{\textquotesingle{}€\textquotesingle{}}\OperatorTok{);}
\NormalTok{printf}\OperatorTok{(}\StringTok{"}\SpecialCharTok{\%x\textbackslash{}n}\StringTok{"}\OperatorTok{,} \CharTok{\textquotesingle{}\textbackslash{}}\ErrorTok{u20ac}\CharTok{\textquotesingle{}}\OperatorTok{);}
\end{Highlighting}
\end{Shaded}

\end{tcolorbox}

\index{Multibyte characters|)}

\subsection{Wide Character}\label{wide-characters}

\index{Wide characters|(}

Kalau bukan karakter multibyte, berarti kamu \emph{wide character}.

Wide character adalah satu nilai yang bisa mewakili secara unik karakter apa pun di locale saat ini. Analog dengan code point Unicode. Tapi bisa jadi bukan. Atau bisa jadi iya.

Pada dasarnya, kalau string karakter multibyte adalah array byte, string wide character adalah array \emph{karakter}. Jadi kamu bisa mulai berpikir per karakter, bukan per byte (yang terakhir jadi berantakan ketika karakter mulai memakai jumlah byte yang bervariasi).

\index{wchar\_t type@\texttt{wchar\_t} type|(}

Wide character bisa direpresentasikan oleh sejumlah tipe, tapi yang paling menonjol adalah \hl{\mbox{\texttt{wchar\_t}}}. Ia yang utama. Mirip \hl{\mbox{\texttt{char}}}, hanya saja lebar.

Mungkin kamu bertanya-tanya, kalau kita tidak bisa tahu apakah itu Unicode atau bukan, bagaimana itu memberi banyak fleksibilitas dalam menulis kode? \hl{\mbox{\texttt{wchar\_t}}} membuka sebagian pintu itu, karena ada sekumpulan fungsi lengkap yang bisa dipakai untuk menangani string \hl{\mbox{\texttt{wchar\_t}}} (seperti mengambil panjang, dan sebagainya) tanpa peduli encoding-nya.

\section{\texorpdfstring{Memakai Wide Character dan \hl{\mbox{\texttt{wchar\_t}}}}{Memakai Wide Character dan wchar\_t}}\label{memakai-wide-character-dan-wchar_t}

Saatnya tipe baru: \hl{\mbox{\texttt{wchar\_t}}}. Ini tipe wide character utama. Ingat bahwa \hl{\mbox{\texttt{char}}} hanya satu byte? Dan satu byte mungkin tidak cukup untuk merepresentasikan semua karakter? Nah, yang ini cukup.

Untuk memakai \hl{\mbox{\texttt{wchar\_t}}}, include \hl{\mbox{\texttt{\textless{}wchar.h\textgreater{}}}}.

Berapa byte ukurannya? Tidak sepenuhnya jelas. Bisa 16 bit. Bisa 32 bit.

Tapi tunggu, katamu---kalau hanya 16 bit, ia tidak cukup besar untuk menampung semua code point Unicode, kan? Benar---memang tidak. Spesifikasi tidak mewajibkannya. Ia hanya harus bisa merepresentasikan semua karakter di locale saat ini.

Ini bisa menyebabkan masalah dengan Unicode pada platform yang \hl{\mbox{\texttt{wchar\_t}}}-nya 16-bit (ehem---Windows). Tapi itu di luar cakupan panduan ini.

\index{L wide character prefix@\texttt{L} wide character prefix|(}

Kamu bisa mendeklarasikan string atau karakter bertipe ini dengan awalan \hl{\mbox{\texttt{L}}}, dan mencetaknya dengan format specifier \hl{\mbox{\texttt{\%ls}}} (``ell ess''). Atau mencetak satu \hl{\mbox{\texttt{wchar\_t}}} dengan \hl{\mbox{\texttt{\%lc}}}.

\begin{tcolorbox}[enhanced jigsaw,left=3pt,breakable,colback=white!97.2549!black,boxrule=0.2mm]

\begin{Shaded}
\begin{Highlighting}[]
\DataTypeTok{wchar\_t} \OperatorTok{*}\NormalTok{s }\OperatorTok{=} \StringTok{L"Hello, world!"}\OperatorTok{;}
\DataTypeTok{wchar\_t}\NormalTok{ c }\OperatorTok{=} \CharTok{L\textquotesingle{}B\textquotesingle{}}\OperatorTok{;}

\NormalTok{printf}\OperatorTok{(}\StringTok{"}\SpecialCharTok{\%ls}\StringTok{ }\SpecialCharTok{\%lc\textbackslash{}n}\StringTok{"}\OperatorTok{,}\NormalTok{ s}\OperatorTok{,}\NormalTok{ c}\OperatorTok{);}
\end{Highlighting}
\end{Shaded}

\end{tcolorbox}

\index{L wide character prefix@\texttt{L} wide character prefix|)}

Nah---apakah karakter-karakter itu disimpan sebagai code point Unicode atau tidak? Tergantung implementasi. Tapi kamu bisa mengujinya dengan makro \index{\_\_STDC\_ISO\_10646\_\_ macro@\texttt{\_\_STDC\_ISO\_10646\_\_} macro} \hl{\mbox{\texttt{\_\_STDC\_ISO\_10646\_\_}}}. Kalau makro ini terdefinisi, jawabannya, ``Itu Unicode!''

Lebih rinci, nilai di makro itu adalah integer berbentuk \hl{\mbox{\texttt{yyyymm}}} yang memberi tahu standar Unicode mana yang bisa kamu andalkan---yang berlaku pada tanggal itu.

Tapi bagaimana cara memakainya?

\subsection{\texorpdfstring{Konversi Multibyte ke \hl{\mbox{\texttt{wchar\_t}}}}{Konversi Multibyte ke wchar\_t}}\label{konversi-multibyte-ke-wchar_t}

Jadi bagaimana kita berpindah dari string standar berorientasi byte ke string lebar berorientasi karakter dan sebaliknya?

Kita bisa memakai beberapa fungsi konversi string untuk mewujudkannya.

Pertama, beberapa konvensi penamaan yang akan kamu lihat pada fungsi-fungsi ini:

\begin{itemize}
\tightlist
\item
  \hl{\mbox{\texttt{mb}}}: multibyte
\item
  \hl{\mbox{\texttt{wc}}}: wide character
\item
  \hl{\mbox{\texttt{mbs}}}: string multibyte
\item
  \hl{\mbox{\texttt{wcs}}}: string wide character
\end{itemize}

Jadi kalau kita ingin mengonversi string multibyte menjadi string wide character, kita bisa memanggil \hl{\mbox{\texttt{mbstowcs()}}}. Dan sebaliknya: \hl{\mbox{\texttt{wcstombs()}}}.

\index{mbtowc() function@\texttt{mbtowc()} function} \index{wctomb() function@\texttt{wctomb()} function} \index{mbstowcs() function@\texttt{mbstowcs()} function} \index{wcstombs() function@\texttt{wcstombs()} function}

{\def\LTcaptype{none} % do not increment counter
\begin{longtable}[]{@{}ll@{}}
\toprule\noalign{}
Fungsi Konversi & Deskripsi \\
\midrule\noalign{}
\endhead
\bottomrule\noalign{}
\endlastfoot
\hl{\mbox{\texttt{mbtowc()}}} & Mengonversi karakter multibyte menjadi wide character. \\
\hl{\mbox{\texttt{wctomb()}}} & Mengonversi wide character menjadi karakter multibyte. \\
\hl{\mbox{\texttt{mbstowcs()}}} & Mengonversi string multibyte menjadi wide string. \\
\hl{\mbox{\texttt{wcstombs()}}} & Mengonversi wide string menjadi string multibyte. \\
\end{longtable}
}

Mari kita buat demo singkat yang mengonversi string multibyte menjadi string wide character, lalu membandingkan panjang kedua string itu dengan fungsi masing-masing.

\index{mbstowcs() function@\texttt{mbstowcs()} function|(}

\begin{tcolorbox}[enhanced jigsaw,left=3pt,breakable,colback=white!97.2549!black,boxrule=0.2mm]

\begin{Shaded}
\begin{Highlighting}[numbers=left,,]
\PreprocessorTok{\#include }\ImportTok{\textless{}stdio.h\textgreater{}}
\PreprocessorTok{\#include }\ImportTok{\textless{}stdlib.h\textgreater{}}
\PreprocessorTok{\#include }\ImportTok{\textless{}wchar.h\textgreater{}}
\PreprocessorTok{\#include }\ImportTok{\textless{}string.h\textgreater{}}
\PreprocessorTok{\#include }\ImportTok{\textless{}locale.h\textgreater{}}

\DataTypeTok{int}\NormalTok{ main}\OperatorTok{(}\DataTypeTok{void}\OperatorTok{)}
\OperatorTok{\{}
    \CommentTok{// Keluar dari locale C ke locale yang kemungkinan punya simbol euro}
\NormalTok{    setlocale}\OperatorTok{(}\NormalTok{LC\_ALL}\OperatorTok{,} \StringTok{""}\OperatorTok{);}

    \CommentTok{// String multibyte asli dengan simbol euro (code point Unicode 20ac)}
    \DataTypeTok{char} \OperatorTok{*}\NormalTok{mb\_string }\OperatorTok{=} \StringTok{"The cost is }\SpecialCharTok{\textbackslash{}u20ac}\StringTok{1.23"}\OperatorTok{;}  \CommentTok{// €1.23}
    \DataTypeTok{size\_t}\NormalTok{ mb\_len }\OperatorTok{=}\NormalTok{ strlen}\OperatorTok{(}\NormalTok{mb\_string}\OperatorTok{);}

    \CommentTok{// Array wide character yang akan menampung string hasil konversi}
    \DataTypeTok{wchar\_t}\NormalTok{ wc\_string}\OperatorTok{[}\DecValTok{128}\OperatorTok{];}  \CommentTok{// Menampung sampai 128 wide character}

    \CommentTok{// Konversi string MB ke WC; ini mengembalikan jumlah wide char}
    \DataTypeTok{size\_t}\NormalTok{ wc\_len }\OperatorTok{=}\NormalTok{ mbstowcs}\OperatorTok{(}\NormalTok{wc\_string}\OperatorTok{,}\NormalTok{ mb\_string}\OperatorTok{,} \DecValTok{128}\OperatorTok{);}

    \CommentTok{// Cetak hasil{-}{-}perhatikan \%ls untuk string wide char}
\NormalTok{    printf}\OperatorTok{(}\StringTok{"multibyte: }\SpecialCharTok{\textbackslash{}"\%s\textbackslash{}"}\StringTok{ (}\SpecialCharTok{\%zu}\StringTok{ bytes)}\SpecialCharTok{\textbackslash{}n}\StringTok{"}\OperatorTok{,}\NormalTok{ mb\_string}\OperatorTok{,}\NormalTok{ mb\_len}\OperatorTok{);}
\NormalTok{    printf}\OperatorTok{(}\StringTok{"wide char: }\SpecialCharTok{\textbackslash{}"\%ls\textbackslash{}"}\StringTok{ (}\SpecialCharTok{\%zu}\StringTok{ characters)}\SpecialCharTok{\textbackslash{}n}\StringTok{"}\OperatorTok{,}\NormalTok{ wc\_string}\OperatorTok{,}\NormalTok{ wc\_len}\OperatorTok{);}
\OperatorTok{\}}
\end{Highlighting}
\end{Shaded}

\end{tcolorbox}

\index{wchar\_t type@\texttt{wchar\_t} type|)}

Di sistemku, ini mengeluarkan:

\begin{tcolorbox}[enhanced jigsaw,left=3pt,breakable,colback=white!97.2549!black,boxrule=0.2mm]

\begin{Shaded}
\begin{Highlighting}[]
\NormalTok{multibyte: "The cost is €1.23" (19 bytes)}
\NormalTok{wide char: "The cost is €1.23" (17 characters)}
\end{Highlighting}
\end{Shaded}

\end{tcolorbox}

(Jumlah byte di sistemmu mungkin berbeda tergantung locale.)

Satu hal menarik: \hl{\mbox{\texttt{mbstowcs()}}}, selain mengonversi string multibyte menjadi wide, mengembalikan panjang (dalam karakter) string wide character itu. Pada sistem yang patuh POSIX, kamu bisa memanfaatkan mode khusus ketika ia \emph{hanya} mengembalikan panjang-dalam-karakter dari string multibyte tertentu: cukup oper \hl{\mbox{\texttt{NULL}}} untuk tujuan, dan \hl{\mbox{\texttt{0}}} untuk jumlah karakter maksimum yang dikonversi (nilai ini diabaikan).

(Pada kode di bawah, aku memakai extended source character set milikku---kamu mungkin harus menggantinya dengan escape \hl{\mbox{\texttt{\textbackslash{}u}}}.)

\begin{tcolorbox}[enhanced jigsaw,left=3pt,breakable,colback=white!97.2549!black,boxrule=0.2mm]

\begin{Shaded}
\begin{Highlighting}[]
\NormalTok{setlocale}\OperatorTok{(}\NormalTok{LC\_ALL}\OperatorTok{,} \StringTok{""}\OperatorTok{);}

\CommentTok{// String berikut punya 7 karakter}
\DataTypeTok{size\_t}\NormalTok{ len\_in\_chars }\OperatorTok{=}\NormalTok{ mbstowcs}\OperatorTok{(}\NormalTok{NULL}\OperatorTok{,} \StringTok{"§¶°±π€•"}\OperatorTok{,} \DecValTok{0}\OperatorTok{);}

\NormalTok{printf}\OperatorTok{(}\StringTok{"}\SpecialCharTok{\%zu}\StringTok{"}\OperatorTok{,}\NormalTok{ len\_in\_chars}\OperatorTok{);}  \CommentTok{// 7}
\end{Highlighting}
\end{Shaded}

\end{tcolorbox}

\index{mbstowcs() function@\texttt{mbstowcs()} function|)}

Sekali lagi, itu ekstensi POSIX yang tidak portabel.

Dan, tentu saja, kalau ingin mengonversi ke arah sebaliknya, pakai \index{wcstombs() function@\texttt{wcstombs()} function} \hl{\mbox{\texttt{wcstombs()}}}.

\section{Fungsionalitas Wide Character}\label{fungsionalitas-wide-character}

Setelah masuk ke wilayah wide character, kita punya berbagai fungsionalitas yang bisa dipakai. Aku hanya akan meringkas sejumlah fungsi di sini, tapi pada dasarnya yang ada adalah versi wide character dari fungsi string multibyte yang sudah kita kenal. (Misalnya, kita kenal \hl{\mbox{\texttt{strlen()}}} untuk string multibyte; ada \index{wcslen() function@\texttt{wcslen()} function} \hl{\mbox{\texttt{wcslen()}}} untuk string wide character.)

\subsection{\texorpdfstring{\hl{\mbox{\texttt{wint\_t}}}}{wint\_t}}\label{wint_t}

\index{wint\_t type@\texttt{wint\_t} type|(}

Banyak fungsi ini memakai \hl{\mbox{\texttt{wint\_t}}} untuk menyimpan karakter tunggal, baik yang dioper maupun yang dikembalikan.

Secara sifat ia terkait dengan \hl{\mbox{\texttt{wchar\_t}}}. \hl{\mbox{\texttt{wint\_t}}} adalah integer yang bisa merepresentasikan semua nilai di extended character set, dan juga karakter end-of-file khusus, \hl{\mbox{\texttt{WEOF}}}.

\index{wint\_t type@\texttt{wint\_t} type|)}

Ini dipakai oleh sejumlah fungsi wide character yang berorientasi satu karakter.

\subsection{Orientasi I/O Stream}\label{io-stream-orientation}

\index{I/O stream orientation|(}

Ringkasnya: jangan mencampur fungsi berorientasi byte (seperti \hl{\mbox{\texttt{fprintf()}}}) dengan fungsi berorientasi wide (seperti \hl{\mbox{\texttt{fwprintf()}}}). Putuskan apakah sebuah stream akan berorientasi byte atau wide dan tetaplah memakai jenis fungsi I/O itu.

Lebih rinci: stream bisa berorientasi byte atau wide. Saat sebuah stream pertama kali dibuat, ia belum punya orientasi, tapi pembacaan atau penulisan pertama akan menetapkan orientasinya.

Kalau kamu pertama kali memakai operasi wide (seperti \hl{\mbox{\texttt{fwprintf()}}}), stream itu berorientasi wide.

Kalau kamu pertama kali memakai operasi byte (seperti \hl{\mbox{\texttt{fprintf()}}}), stream itu berorientasi byte.

Kamu bisa mengatur secara manual stream yang belum berorientasi ke satu arah atau yang lain dengan memanggil \index{fwide() function@\texttt{fwide()} function} \hl{\mbox{\texttt{fwide()}}}. Fungsi yang sama bisa dipakai untuk mengetahui orientasi sebuah stream.

Kalau perlu mengubah orientasi di tengah jalan, bisa dilakukan dengan \hl{\mbox{\texttt{freopen()}}}.

\index{I/O stream orientation|)}

\subsection{Fungsi I/O}\label{fungsi-io}

Biasanya include \hl{\mbox{\texttt{\textless{}stdio.h\textgreater{}}}} dan \hl{\mbox{\texttt{\textless{}wchar.h\textgreater{}}}} untuk ini.

\index{wprintf() function@\texttt{wprintf()} function} \index{wscanf() function@\texttt{wscanf()} function} \index{getwchar() function@\texttt{getwchar()} function} \index{putwchar() function@\texttt{putwchar()} function} \index{fwprintf() function@\texttt{fwprintf()} function} \index{fwscanf() function@\texttt{fwscanf()} function} \index{fgetwc() function@\texttt{fgetwc()} function} \index{fputwc() function@\texttt{fputwc()} function} \index{fgetws() function@\texttt{fgetws()} function} \index{fputws() function@\texttt{fputws()} function} \index{swprintf() function@\texttt{swprintf()} function} \index{swscanf() function@\texttt{swscanf()} function} \index{vfwprintf() function@\texttt{vfwprintf()} function} \index{vfwscanf() function@\texttt{vfwscanf()} function} \index{vswprintf() function@\texttt{vswprintf()} function} \index{vswscanf() function@\texttt{vswscanf()} function} \index{vwprintf() function@\texttt{vwprintf()} function} \index{vwscanf() function@\texttt{vwscanf()} function} \index{ungetwc() function@\texttt{ungetwc()} function} \index{fwide() function@\texttt{fwide()} function}

{\def\LTcaptype{none} % do not increment counter
\begin{longtable}[]{@{}
  >{\raggedright\arraybackslash}p{(\linewidth - 2\tabcolsep) * \real{0.1905}}
  >{\raggedright\arraybackslash}p{(\linewidth - 2\tabcolsep) * \real{0.8095}}@{}}
\toprule\noalign{}
\begin{minipage}[b]{\linewidth}\raggedright
Fungsi I/O
\end{minipage} & \begin{minipage}[b]{\linewidth}\raggedright
Deskripsi
\end{minipage} \\
\midrule\noalign{}
\endhead
\bottomrule\noalign{}
\endlastfoot
\hl{\mbox{\texttt{wprintf()}}} & Output terformat ke konsol. \\
\hl{\mbox{\texttt{wscanf()}}} & Input terformat dari konsol. \\
\hl{\mbox{\texttt{getwchar()}}} & Input konsol berbasis karakter. \\
\hl{\mbox{\texttt{putwchar()}}} & Output konsol berbasis karakter. \\
\hl{\mbox{\texttt{fwprintf()}}} & Output terformat ke file. \\
\hl{\mbox{\texttt{fwscanf()}}} & Input terformat dari file. \\
\hl{\mbox{\texttt{fgetwc()}}} & Input file berbasis karakter. \\
\hl{\mbox{\texttt{fputwc()}}} & Output file berbasis karakter. \\
\hl{\mbox{\texttt{fgetws()}}} & Input file berbasis string. \\
\hl{\mbox{\texttt{fputws()}}} & Output file berbasis string. \\
\hl{\mbox{\texttt{swprintf()}}} & Output terformat ke string. \\
\hl{\mbox{\texttt{swscanf()}}} & Input terformat dari string. \\
\hl{\mbox{\texttt{vfwprintf()}}} & Output terformat variadic ke file. \\
\hl{\mbox{\texttt{vfwscanf()}}} & Input terformat variadic dari file. \\
\hl{\mbox{\texttt{vswprintf()}}} & Output terformat variadic ke string. \\
\hl{\mbox{\texttt{vswscanf()}}} & Input terformat variadic dari string. \\
\hl{\mbox{\texttt{vwprintf()}}} & Output terformat variadic ke konsol. \\
\hl{\mbox{\texttt{vwscanf()}}} & Input terformat variadic dari konsol. \\
\hl{\mbox{\texttt{ungetwc()}}} & Mengembalikan sebuah wide character ke stream. \\
\hl{\mbox{\texttt{fwide()}}} & Mengambil atau mengatur orientasi multibyte/wide sebuah stream. \\
\end{longtable}
}

\subsection{Fungsi Konversi Tipe}\label{fungsi-konversi-tipe}

Biasanya include \hl{\mbox{\texttt{\textless{}wchar.h\textgreater{}}}} untuk ini.

\index{wcstod() function@\texttt{wcstod()} function} \index{wcstof() function@\texttt{wcstof()} function} \index{wcstold() function@\texttt{wcstold()} function} \index{wcstol() function@\texttt{wcstol()} function} \index{wcstoll() function@\texttt{wcstoll()} function} \index{wcstoul() function@\texttt{wcstoul()} function} \index{wcstoull() function@\texttt{wcstoull()} function}

{\def\LTcaptype{none} % do not increment counter
\begin{longtable}[]{@{}ll@{}}
\toprule\noalign{}
Fungsi Konversi & Deskripsi \\
\midrule\noalign{}
\endhead
\bottomrule\noalign{}
\endlastfoot
\hl{\mbox{\texttt{wcstod()}}} & Mengonversi string menjadi \hl{\mbox{\texttt{double}}}. \\
\hl{\mbox{\texttt{wcstof()}}} & Mengonversi string menjadi \hl{\mbox{\texttt{float}}}. \\
\hl{\mbox{\texttt{wcstold()}}} & Mengonversi string menjadi \hl{\mbox{\texttt{long\ double}}}. \\
\hl{\mbox{\texttt{wcstol()}}} & Mengonversi string menjadi \hl{\mbox{\texttt{long}}}. \\
\hl{\mbox{\texttt{wcstoll()}}} & Mengonversi string menjadi \hl{\mbox{\texttt{long\ long}}}. \\
\hl{\mbox{\texttt{wcstoul()}}} & Mengonversi string menjadi \hl{\mbox{\texttt{unsigned\ long}}}. \\
\hl{\mbox{\texttt{wcstoull()}}} & Mengonversi string menjadi \hl{\mbox{\texttt{unsigned\ long\ long}}}. \\
\end{longtable}
}

\subsection{Fungsi Penyalinan String dan Memori}\label{fungsi-penyalinan-string-dan-memori}

Biasanya include \hl{\mbox{\texttt{\textless{}wchar.h\textgreater{}}}} untuk ini.

\index{wcscpy() function@\texttt{wcscpy()} function} \index{wcsncpy() function@\texttt{wcsncpy()} function} \index{wmemcpy() function@\texttt{wmemcpy()} function} \index{wmemmove() function@\texttt{wmemmove()} function} \index{wcscat() function@\texttt{wcscat()} function} \index{wcsncat() function@\texttt{wcsncat()} function}

{\def\LTcaptype{none} % do not increment counter
\begin{longtable}[]{@{}ll@{}}
\toprule\noalign{}
Fungsi Penyalinan & Deskripsi \\
\midrule\noalign{}
\endhead
\bottomrule\noalign{}
\endlastfoot
\hl{\mbox{\texttt{wcscpy()}}} & Menyalin string. \\
\hl{\mbox{\texttt{wcsncpy()}}} & Menyalin string, dibatasi panjang. \\
\hl{\mbox{\texttt{wmemcpy()}}} & Menyalin memori. \\
\hl{\mbox{\texttt{wmemmove()}}} & Menyalin memori yang mungkin tumpang tindih. \\
\hl{\mbox{\texttt{wcscat()}}} & Menyambung string. \\
\hl{\mbox{\texttt{wcsncat()}}} & Menyambung string, dibatasi panjang. \\
\end{longtable}
}

\subsection{Fungsi Perbandingan String dan Memori}\label{fungsi-perbandingan-string-dan-memori}

Biasanya include \hl{\mbox{\texttt{\textless{}wchar.h\textgreater{}}}} untuk ini.

\index{wcscmp() function@\texttt{wcscmp()} function} \index{wcsncmp() function@\texttt{wcsncmp()} function} \index{wcscoll() function@\texttt{wcscoll()} function} \index{wmemcmp() function@\texttt{wmemcmp()} function} \index{wcsxfrm() function@\texttt{wcsxfrm()} function}

{\def\LTcaptype{none} % do not increment counter
\begin{longtable}[]{@{}
  >{\raggedright\arraybackslash}p{(\linewidth - 2\tabcolsep) * \real{0.2317}}
  >{\raggedright\arraybackslash}p{(\linewidth - 2\tabcolsep) * \real{0.7683}}@{}}
\toprule\noalign{}
\begin{minipage}[b]{\linewidth}\raggedright
Fungsi Perbandingan
\end{minipage} & \begin{minipage}[b]{\linewidth}\raggedright
Deskripsi
\end{minipage} \\
\midrule\noalign{}
\endhead
\bottomrule\noalign{}
\endlastfoot
\hl{\mbox{\texttt{wcscmp()}}} & Membandingkan string secara leksikografis. \\
\hl{\mbox{\texttt{wcsncmp()}}} & Membandingkan string secara leksikografis, dibatasi panjang. \\
\hl{\mbox{\texttt{wcscoll()}}} & Membandingkan string menurut urutan kamus berdasarkan locale. \\
\hl{\mbox{\texttt{wmemcmp()}}} & Membandingkan memori secara leksikografis. \\
\hl{\mbox{\texttt{wcsxfrm()}}} & Mengubah string menjadi versi yang membuat \hl{\mbox{\texttt{wcscmp()}}} berperilaku seperti \hl{\mbox{\texttt{wcscoll()}}}\footnote{\hl{\mbox{\texttt{wcscoll()}}} sama dengan \hl{\mbox{\texttt{wcsxfrm()}}} diikuti \hl{\mbox{\texttt{wcscmp()}}}.}. \\
\end{longtable}
}

\subsection{Fungsi Pencarian String}\label{fungsi-pencarian-string}

Biasanya include \hl{\mbox{\texttt{\textless{}wchar.h\textgreater{}}}} untuk ini.

\index{wcschr() function@\texttt{wcschr()} function} \index{wcsrchr() function@\texttt{wcsrchr()} function} \index{wmemchr() function@\texttt{wmemchr()} function} \index{wcsstr() function@\texttt{wcsstr()} function} \index{wcspbrk() function@\texttt{wcspbrk()} function} \index{wcsspn() function@\texttt{wcsspn()} function} \index{wcscspn() function@\texttt{wcscspn()} function} \index{wcstok() function@\texttt{wcstok()} function}

{\def\LTcaptype{none} % do not increment counter
\begin{longtable}[]{@{}
  >{\raggedright\arraybackslash}p{(\linewidth - 2\tabcolsep) * \real{0.5000}}
  >{\raggedright\arraybackslash}p{(\linewidth - 2\tabcolsep) * \real{0.5000}}@{}}
\toprule\noalign{}
\begin{minipage}[b]{\linewidth}\raggedright
Fungsi Pencarian
\end{minipage} & \begin{minipage}[b]{\linewidth}\raggedright
Deskripsi
\end{minipage} \\
\midrule\noalign{}
\endhead
\bottomrule\noalign{}
\endlastfoot
\hl{\mbox{\texttt{wcschr()}}} & Mencari karakter dalam string. \\
\hl{\mbox{\texttt{wcsrchr()}}} & Mencari karakter dalam string dari belakang. \\
\hl{\mbox{\texttt{wmemchr()}}} & Mencari karakter di memori. \\
\hl{\mbox{\texttt{wcsstr()}}} & Mencari substring dalam string. \\
\hl{\mbox{\texttt{wcspbrk()}}} & Mencari salah satu dari sekumpulan karakter dalam string. \\
\hl{\mbox{\texttt{wcsspn()}}} & Mencari panjang substring yang berisi salah satu dari sekumpulan karakter. \\
\hl{\mbox{\texttt{wcscspn()}}} & Mencari panjang substring sebelum salah satu dari sekumpulan karakter. \\
\hl{\mbox{\texttt{wcstok()}}} & Mencari token dalam string. \\
\end{longtable}
}

\subsection{Fungsi Panjang/Lain-lain}\label{fungsi-panjanglain-lain}

Biasanya include \hl{\mbox{\texttt{\textless{}wchar.h\textgreater{}}}} untuk ini.

\index{wcslen() function@\texttt{wcslen()} function} \index{wmemset() function@\texttt{wmemset()} function} \index{wcsftime() function@\texttt{wcsftime()} function}

{\def\LTcaptype{none} % do not increment counter
\begin{longtable}[]{@{}ll@{}}
\toprule\noalign{}
Fungsi Panjang/Lain-lain & Deskripsi \\
\midrule\noalign{}
\endhead
\bottomrule\noalign{}
\endlastfoot
\hl{\mbox{\texttt{wcslen()}}} & Mengembalikan panjang string. \\
\hl{\mbox{\texttt{wmemset()}}} & Mengisi karakter di memori. \\
\hl{\mbox{\texttt{wcsftime()}}} & Output tanggal dan waktu terformat. \\
\end{longtable}
}

\subsection{Fungsi Klasifikasi Karakter}\label{fungsi-klasifikasi-karakter}

Include \hl{\mbox{\texttt{\textless{}wctype.h\textgreater{}}}} untuk ini.

\index{iswalnum() function@\texttt{iswalnum()} function} \index{iswalpha() function@\texttt{iswalpha()} function} \index{iswblank() function@\texttt{iswblank()} function} \index{iswcntrl() function@\texttt{iswcntrl()} function} \index{iswdigit() function@\texttt{iswdigit()} function} \index{iswgraph() function@\texttt{iswgraph()} function} \index{iswlower() function@\texttt{iswlower()} function} \index{iswprint() function@\texttt{iswprint()} function} \index{iswpunct() function@\texttt{iswpunct()} function} \index{iswspace() function@\texttt{iswspace()} function} \index{iswupper() function@\texttt{iswupper()} function} \index{iswxdigit() function@\texttt{iswxdigit()} function} \index{towlower() function@\texttt{towlower()} function} \index{towupper() function@\texttt{towupper()} function}

{\def\LTcaptype{none} % do not increment counter
\begin{longtable}[]{@{}
  >{\raggedright\arraybackslash}p{(\linewidth - 2\tabcolsep) * \real{0.1905}}
  >{\raggedright\arraybackslash}p{(\linewidth - 2\tabcolsep) * \real{0.8095}}@{}}
\toprule\noalign{}
\begin{minipage}[b]{\linewidth}\raggedright
Fungsi Panjang/Lain-lain
\end{minipage} & \begin{minipage}[b]{\linewidth}\raggedright
Deskripsi
\end{minipage} \\
\midrule\noalign{}
\endhead
\bottomrule\noalign{}
\endlastfoot
\hl{\mbox{\texttt{iswalnum()}}} & True jika karakter alfanumerik. \\
\hl{\mbox{\texttt{iswalpha()}}} & True jika karakter berupa huruf. \\
\hl{\mbox{\texttt{iswblank()}}} & True jika karakter kosong (mirip spasi, tapi bukan newline). \\
\hl{\mbox{\texttt{iswcntrl()}}} & True jika karakter adalah karakter kontrol. \\
\hl{\mbox{\texttt{iswdigit()}}} & True jika karakter berupa digit. \\
\hl{\mbox{\texttt{iswgraph()}}} & True jika karakter bisa dicetak (kecuali spasi). \\
\hl{\mbox{\texttt{iswlower()}}} & True jika karakter huruf kecil. \\
\hl{\mbox{\texttt{iswprint()}}} & True jika karakter bisa dicetak (termasuk spasi). \\
\hl{\mbox{\texttt{iswpunct()}}} & True jika karakter berupa tanda baca. \\
\hl{\mbox{\texttt{iswspace()}}} & True jika karakter berupa whitespace. \\
\hl{\mbox{\texttt{iswupper()}}} & True jika karakter huruf besar. \\
\hl{\mbox{\texttt{iswxdigit()}}} & True jika karakter berupa digit hex. \\
\hl{\mbox{\texttt{towlower()}}} & Mengonversi karakter menjadi huruf kecil. \\
\hl{\mbox{\texttt{towupper()}}} & Mengonversi karakter menjadi huruf besar. \\
\end{longtable}
}

\section{Parse State, Fungsi Restartable}\label{parse-state-fungsi-restartable}

\index{Multibyte characters!parse state|(}

Kita akan masuk sedikit ke jeroan konversi multibyte, tapi ini hal yang bagus untuk dipahami secara konseptual.

Bayangkan bagaimana programmu mengambil deretan karakter multibyte dan mengubahnya menjadi wide character, atau sebaliknya. Pada suatu titik, ia mungkin sedang di tengah mem-parse sebuah karakter, atau mungkin harus menunggu byte tambahan sebelum menentukan nilai akhirnya.

Parse state ini disimpan dalam variabel opaque bertipe \hl{\mbox{\texttt{mbstate\_t}}} dan dipakai setiap kali konversi dilakukan. Begitulah cara fungsi konversi melacak posisinya di tengah pekerjaan.

Dan kalau kamu berpindah ke deretan karakter lain di tengah stream, atau mencoba seek ke tempat lain di deretan inputmu, ia bisa bingung karenanya.

Mungkin kamu ingin menegurku soal ini: tadi kita sudah melakukan beberapa konversi, dan aku tidak pernah menyebut \hl{\mbox{\texttt{mbstate\_t}}} di mana pun.

Itu karena fungsi konversi seperti \hl{\mbox{\texttt{mbstowcs()}}}, \hl{\mbox{\texttt{wctomb()}}}, dan sebagainya masing-masing punya variabel \hl{\mbox{\texttt{mbstate\_t}}} sendiri yang mereka pakai. Tapi hanya ada satu per fungsi, jadi kalau kamu menulis kode multithread, fungsi-fungsi itu tidak aman dipakai.

Untungnya, C mendefinisikan versi \emph{restartable} dari fungsi-fungsi ini yang memungkinkan kamu mengoper \hl{\mbox{\texttt{mbstate\_t}}}-mu sendiri per thread kalau perlu. Kalau kamu bekerja dengan multithread, pakai ini!

Catatan singkat tentang menginisialisasi variabel \hl{\mbox{\texttt{mbstate\_t}}}: cukup \hl{\mbox{\texttt{memset()}}} ke nol. Tidak ada fungsi bawaan untuk memaksa inisialisasi.

\begin{tcolorbox}[enhanced jigsaw,left=3pt,breakable,colback=white!97.2549!black,boxrule=0.2mm]

\begin{Shaded}
\begin{Highlighting}[]
\NormalTok{mbstate\_t mbs}\OperatorTok{;}

\CommentTok{// Atur state ke state awal}
\NormalTok{memset}\OperatorTok{(\&}\NormalTok{mbs}\OperatorTok{,} \DecValTok{0}\OperatorTok{,} \KeywordTok{sizeof}\NormalTok{ mbs}\OperatorTok{);}
\end{Highlighting}
\end{Shaded}

\end{tcolorbox}

Berikut daftar fungsi konversi restartable---perhatikan konvensi penamaan dengan menaruh ``\hl{\mbox{\texttt{r}}}'' setelah tipe ``asal'':

\begin{itemize}
\tightlist
\item
  \hl{\mbox{\texttt{mbrtowc()}}}---multibyte ke wide character
\item
  \hl{\mbox{\texttt{wcrtomb()}}}---wide character ke multibyte
\item
  \hl{\mbox{\texttt{mbsrtowcs()}}}---string multibyte ke string wide character
\item
  \hl{\mbox{\texttt{wcsrtombs()}}}---string wide character ke string multibyte
\end{itemize}

Fungsi-fungsi ini sangat mirip dengan padanannya yang non-restartable, hanya saja mereka mewajibkan kamu mengoper pointer ke variabel \hl{\mbox{\texttt{mbstate\_t}}} milikmu sendiri. Mereka juga mengubah pointer string sumber (untuk membantu kalau ditemukan byte yang tidak valid), jadi mungkin berguna menyimpan salinan aslinya.

Berikut contoh dari awal bab yang disusun ulang agar mengoper \hl{\mbox{\texttt{mbstate\_t}}} kita sendiri.

\begin{tcolorbox}[enhanced jigsaw,left=3pt,breakable,colback=white!97.2549!black,boxrule=0.2mm]

\begin{Shaded}
\begin{Highlighting}[numbers=left,,]
\PreprocessorTok{\#include }\ImportTok{\textless{}stdio.h\textgreater{}}
\PreprocessorTok{\#include }\ImportTok{\textless{}stdlib.h\textgreater{}}
\PreprocessorTok{\#include }\ImportTok{\textless{}stddef.h\textgreater{}}
\PreprocessorTok{\#include }\ImportTok{\textless{}wchar.h\textgreater{}}
\PreprocessorTok{\#include }\ImportTok{\textless{}string.h\textgreater{}}
\PreprocessorTok{\#include }\ImportTok{\textless{}locale.h\textgreater{}}

\DataTypeTok{int}\NormalTok{ main}\OperatorTok{(}\DataTypeTok{void}\OperatorTok{)}
\OperatorTok{\{}
    \CommentTok{// Keluar dari locale C ke locale yang kemungkinan punya simbol euro}
\NormalTok{    setlocale}\OperatorTok{(}\NormalTok{LC\_ALL}\OperatorTok{,} \StringTok{""}\OperatorTok{);}

    \CommentTok{// String multibyte asli dengan simbol euro (code point Unicode 20ac)}
    \DataTypeTok{char} \OperatorTok{*}\NormalTok{mb\_string }\OperatorTok{=} \StringTok{"The cost is }\SpecialCharTok{\textbackslash{}u20ac}\StringTok{1.23"}\OperatorTok{;}  \CommentTok{// €1.23}
    \DataTypeTok{size\_t}\NormalTok{ mb\_len }\OperatorTok{=}\NormalTok{ strlen}\OperatorTok{(}\NormalTok{mb\_string}\OperatorTok{);}

    \CommentTok{// Array wide character yang akan menampung string hasil konversi}
    \DataTypeTok{wchar\_t}\NormalTok{ wc\_string}\OperatorTok{[}\DecValTok{128}\OperatorTok{];}  \CommentTok{// Menampung sampai 128 wide character}

    \CommentTok{// Siapkan state konversi}
\NormalTok{    mbstate\_t mbs}\OperatorTok{;}
\NormalTok{    memset}\OperatorTok{(\&}\NormalTok{mbs}\OperatorTok{,} \DecValTok{0}\OperatorTok{,} \KeywordTok{sizeof}\NormalTok{ mbs}\OperatorTok{);}  \CommentTok{// State awal}

    \CommentTok{// mbsrtowcs() mengubah pointer input agar menunjuk ke karakter tidak}
    \CommentTok{// valid pertama, atau NULL jika berhasil. Kita buat salinan pointer}
    \CommentTok{// untuk diutak{-}atik mbsrtowcs() sehingga aslinya tidak berubah.}
    \CommentTok{//}
    \CommentTok{// Contoh ini kemungkinan berhasil, tapi kita periksa lebih ke bawah}
    \CommentTok{// untuk memastikan.}
    \DataTypeTok{const} \DataTypeTok{char} \OperatorTok{*}\NormalTok{invalid }\OperatorTok{=}\NormalTok{ mb\_string}\OperatorTok{;}

    \CommentTok{// Konversi string MB ke WC; ini mengembalikan jumlah wide char}
    \DataTypeTok{size\_t}\NormalTok{ wc\_len }\OperatorTok{=}\NormalTok{ mbsrtowcs}\OperatorTok{(}\NormalTok{wc\_string}\OperatorTok{,} \OperatorTok{\&}\NormalTok{invalid}\OperatorTok{,} \DecValTok{128}\OperatorTok{,} \OperatorTok{\&}\NormalTok{mbs}\OperatorTok{);}

    \ControlFlowTok{if} \OperatorTok{(}\NormalTok{invalid }\OperatorTok{==}\NormalTok{ NULL}\OperatorTok{)} \OperatorTok{\{}
\NormalTok{        printf}\OperatorTok{(}\StringTok{"No invalid characters found}\SpecialCharTok{\textbackslash{}n}\StringTok{"}\OperatorTok{);}

        \CommentTok{// Cetak hasil{-}{-}perhatikan \%ls untuk string wide char}
\NormalTok{        printf}\OperatorTok{(}\StringTok{"multibyte: }\SpecialCharTok{\textbackslash{}"\%s\textbackslash{}"}\StringTok{ (}\SpecialCharTok{\%zu}\StringTok{ bytes)}\SpecialCharTok{\textbackslash{}n}\StringTok{"}\OperatorTok{,}\NormalTok{ mb\_string}\OperatorTok{,}\NormalTok{ mb\_len}\OperatorTok{);}
\NormalTok{        printf}\OperatorTok{(}\StringTok{"wide char: }\SpecialCharTok{\textbackslash{}"\%ls\textbackslash{}"}\StringTok{ (}\SpecialCharTok{\%zu}\StringTok{ characters)}\SpecialCharTok{\textbackslash{}n}\StringTok{"}\OperatorTok{,}\NormalTok{ wc\_string}\OperatorTok{,}\NormalTok{ wc\_len}\OperatorTok{);}
    \OperatorTok{\}} \ControlFlowTok{else} \OperatorTok{\{}
        \DataTypeTok{ptrdiff\_t}\NormalTok{ offset }\OperatorTok{=}\NormalTok{ invalid }\OperatorTok{{-}}\NormalTok{ mb\_string}\OperatorTok{;}
\NormalTok{        printf}\OperatorTok{(}\StringTok{"Invalid character at offset }\SpecialCharTok{\%td\textbackslash{}n}\StringTok{"}\OperatorTok{,}\NormalTok{ offset}\OperatorTok{);}
    \OperatorTok{\}}
\OperatorTok{\}}
\end{Highlighting}
\end{Shaded}

\end{tcolorbox}

Untuk fungsi konversi yang mengelola state-nya sendiri, kamu bisa mereset state internalnya ke state awal dengan mengoper \hl{\mbox{\texttt{NULL}}} untuk argumen \hl{\mbox{\texttt{char*}}}-nya, misalnya:

\begin{tcolorbox}[enhanced jigsaw,left=3pt,breakable,colback=white!97.2549!black,boxrule=0.2mm]

\begin{Shaded}
\begin{Highlighting}[]
\NormalTok{mbstowcs}\OperatorTok{(}\NormalTok{NULL}\OperatorTok{,}\NormalTok{ NULL}\OperatorTok{,} \DecValTok{0}\OperatorTok{);}   \CommentTok{// Reset parse state untuk mbstowcs()}
\NormalTok{mbstowcs}\OperatorTok{(}\NormalTok{dest}\OperatorTok{,}\NormalTok{ src}\OperatorTok{,} \DecValTok{100}\OperatorTok{);}  \CommentTok{// Parse beberapa hal}
\end{Highlighting}
\end{Shaded}

\end{tcolorbox}

Untuk I/O, setiap wide stream mengelola \hl{\mbox{\texttt{mbstate\_t}}}-nya sendiri dan memakainya untuk konversi input dan output selama berjalan.

Dan beberapa fungsi I/O berorientasi byte seperti \hl{\mbox{\texttt{printf()}}} dan \hl{\mbox{\texttt{scanf()}}} menyimpan state internalnya sendiri saat bekerja.

Terakhir, fungsi konversi restartable ini sebenarnya punya state internal sendiri kalau kamu mengoper \hl{\mbox{\texttt{NULL}}} sebagai parameter \hl{\mbox{\texttt{mbstate\_t}}}. Ini membuat mereka berperilaku lebih mirip padanannya yang non-restartable.

\index{Multibyte characters!parse state|)} \index{Wide characters|)}

\section{Encoding Unicode dan C}\label{encoding-unicode-dan-c}

Di bagian ini, kita akan melihat apa yang bisa (dan tidak bisa) dilakukan C pada tiga encoding Unicode tertentu: UTF-8, UTF-16, dan UTF-32.

\subsection{UTF-8}\label{utf-8}

\index{Unicode!UTF-8|(}

Untuk menyegarkan ingatan sebelum bagian ini, baca \hyperref[utf8-quick]{catatan singkat UTF-8, di atas}.

Selain itu, apa kemampuan UTF-8 di C?

Sayangnya, tidak banyak.

\index{u8 UTF-8 prefix@\texttt{u8} UTF-8 prefix|(}

Kamu bisa memberi tahu C bahwa kamu secara spesifik menginginkan string literal ber-encoding UTF-8, dan ia akan melakukannya. Caranya dengan awalan \hl{\mbox{\texttt{u8}}} pada string:

\begin{tcolorbox}[enhanced jigsaw,left=3pt,breakable,colback=white!97.2549!black,boxrule=0.2mm]

\begin{Shaded}
\begin{Highlighting}[]
\DataTypeTok{char} \OperatorTok{*}\NormalTok{s }\OperatorTok{=} \StringTok{u8"Hello, world!"}\OperatorTok{;}

\NormalTok{printf}\OperatorTok{(}\StringTok{"}\SpecialCharTok{\%s\textbackslash{}n}\StringTok{"}\OperatorTok{,}\NormalTok{ s}\OperatorTok{);}   \CommentTok{// Hello, world!{-}{-}kalau kamu bisa mengeluarkan UTF{-}8}
\end{Highlighting}
\end{Shaded}

\end{tcolorbox}

Nah, bisakah karakter Unicode dimasukkan ke sana?

\begin{tcolorbox}[enhanced jigsaw,left=3pt,breakable,colback=white!97.2549!black,boxrule=0.2mm]

\begin{Shaded}
\begin{Highlighting}[]
\DataTypeTok{char} \OperatorTok{*}\NormalTok{s }\OperatorTok{=} \StringTok{u8"€123"}\OperatorTok{;}
\end{Highlighting}
\end{Shaded}

\end{tcolorbox}

\index{u8 UTF-8 prefix@\texttt{u8} UTF-8 prefix|)}

Tentu! Kalau extended source character set mendukungnya. (gcc mendukung.)

Bagaimana kalau tidak? Kamu bisa menentukan code point Unicode dengan \hl{\mbox{\texttt{\textbackslash{}u}}} dan \hl{\mbox{\texttt{\textbackslash{}U}}} yang akrab di sekitar kita, \hyperref[unicode-in-c]{seperti dicatat di atas}.

Tapi hanya itu. Tidak ada cara portabel di standard library untuk mengambil input sembarang dan mengubahnya menjadi UTF-8 kecuali locale-mu UTF-8. Atau untuk mem-parse UTF-8 kecuali locale-mu UTF-8.

Jadi kalau ingin melakukannya, berada di locale UTF-8 dan:

\index{setlocale() function@\texttt{setlocale()} function|(}

\begin{tcolorbox}[enhanced jigsaw,left=3pt,breakable,colback=white!97.2549!black,boxrule=0.2mm]

\begin{Shaded}
\begin{Highlighting}[]
\NormalTok{setlocale}\OperatorTok{(}\NormalTok{LC\_ALL}\OperatorTok{,} \StringTok{""}\OperatorTok{);}
\end{Highlighting}
\end{Shaded}

\end{tcolorbox}

atau cari tahu nama locale UTF-8 di mesin lokalmu dan atur secara eksplisit seperti ini:

\begin{tcolorbox}[enhanced jigsaw,left=3pt,breakable,colback=white!97.2549!black,boxrule=0.2mm]

\begin{Shaded}
\begin{Highlighting}[]
\NormalTok{setlocale}\OperatorTok{(}\NormalTok{LC\_ALL}\OperatorTok{,} \StringTok{"en\_US.UTF{-}8"}\OperatorTok{);}  \CommentTok{// Nama non{-}portabel}
\end{Highlighting}
\end{Shaded}

\end{tcolorbox}

\index{setlocale() function@\texttt{setlocale()} function|)}

Atau pakai \hyperref[utf-3rd-party]{library pihak ketiga}.

\index{Unicode!UTF-8|)}

\subsection{\texorpdfstring{UTF-16, UTF-32, \hl{\mbox{\texttt{char16\_t}}}, dan \hl{\mbox{\texttt{char32\_t}}}}{UTF-16, UTF-32, char16\_t, dan char32\_t}}\label{utf-16-utf-32-char16_t-dan-char32_t}

\index{Unicode!UTF-16|(} \index{Unicode!UTF-32|(} \index{char16\_t type@\texttt{char16\_t} type|(} \index{char32\_t type@\texttt{char32\_t} type|(}

\hl{\mbox{\texttt{char16\_t}}} dan \hl{\mbox{\texttt{char32\_t}}} adalah dua tipe karakter yang berpotensi lebar, masing-masing berukuran 16 bit dan 32 bit. Belum tentu lebar, karena kalau tidak bisa merepresentasikan setiap karakter di locale saat ini, mereka kehilangan sifat wide character-nya. Tapi spesifikasi menyebut mereka tipe ``wide character'' di mana-mana, jadi begitulah.

Tipe-tipe ini ada untuk membuat segalanya sedikit lebih ramah Unicode, mudah-mudahan.

Untuk memakainya, include \hl{\mbox{\texttt{\textless{}uchar.h\textgreater{}}}}. (Itu ``u'', bukan ``w''.)

Header file ini tidak ada di OS X---sayang sekali. Kalau kamu hanya ingin tipenya, kamu bisa:

\begin{tcolorbox}[enhanced jigsaw,left=3pt,breakable,colback=white!97.2549!black,boxrule=0.2mm]

\begin{Shaded}
\begin{Highlighting}[]
\PreprocessorTok{\#include }\ImportTok{\textless{}stdint.h\textgreater{}}

\KeywordTok{typedef} \DataTypeTok{int\_least16\_t} \DataTypeTok{char16\_t}\OperatorTok{;}
\KeywordTok{typedef} \DataTypeTok{int\_least32\_t} \DataTypeTok{char32\_t}\OperatorTok{;}
\end{Highlighting}
\end{Shaded}

\end{tcolorbox}

Tapi kalau kamu juga ingin fungsinya, itu urusanmu sendiri.

\index{u Unicode prefix@\texttt{u} Unicode prefix|(} \index{U Unicode prefix@\texttt{U} Unicode prefix|(}

Dengan asumsi kamu masih bisa lanjut, kamu bisa mendeklarasikan string atau karakter bertipe ini dengan awalan \hl{\mbox{\texttt{u}}} dan \hl{\mbox{\texttt{U}}}:

\begin{tcolorbox}[enhanced jigsaw,left=3pt,breakable,colback=white!97.2549!black,boxrule=0.2mm]

\begin{Shaded}
\begin{Highlighting}[]
\DataTypeTok{char16\_t} \OperatorTok{*}\NormalTok{s }\OperatorTok{=} \StringTok{u"Hello, world!"}\OperatorTok{;}
\DataTypeTok{char16\_t}\NormalTok{ c }\OperatorTok{=} \CharTok{u\textquotesingle{}B\textquotesingle{}}\OperatorTok{;}

\DataTypeTok{char32\_t} \OperatorTok{*}\NormalTok{t }\OperatorTok{=} \StringTok{U"Hello, world!"}\OperatorTok{;}
\DataTypeTok{char32\_t}\NormalTok{ d }\OperatorTok{=} \CharTok{U\textquotesingle{}B\textquotesingle{}}\OperatorTok{;}
\end{Highlighting}
\end{Shaded}

\end{tcolorbox}

\index{char32\_t type@\texttt{char32\_t} type|)} \index{u Unicode prefix@\texttt{u} Unicode prefix|)} \index{U Unicode prefix@\texttt{U} Unicode prefix|)}

Nah---apakah nilai di dalamnya disimpan dalam UTF-16 atau UTF-32? Tergantung implementasi.

Tapi kamu bisa mengujinya. Kalau makro \index{\_\_STDC\_UTF\_16\_\_ macro@\texttt{\_\_STDC\_UTF\_16\_\_} macro} \hl{\mbox{\texttt{\_\_STDC\_UTF\_16\_\_}}} atau \index{\_\_STDC\_UTF\_32\_\_ macro@\texttt{\_\_STDC\_UTF\_32\_\_ macro}} \hl{\mbox{\texttt{\_\_STDC\_UTF\_32\_\_}}} terdefinisi (bernilai \hl{\mbox{\texttt{1}}}), artinya tipe-tipe itu masing-masing menyimpan UTF-16 atau UTF-32.

Kalau kamu penasaran, dan aku tahu kamu penasaran, nilainya, kalau UTF-16 atau UTF-32, disimpan dalam endianess native. Artinya, kamu seharusnya bisa langsung membandingkannya dengan nilai code point Unicode:

\begin{tcolorbox}[enhanced jigsaw,left=3pt,breakable,colback=white!97.2549!black,boxrule=0.2mm]

\begin{Shaded}
\begin{Highlighting}[]
\DataTypeTok{char16\_t}\NormalTok{ pi }\OperatorTok{=} \StringTok{u"}\SpecialCharTok{\textbackslash{}u03C0}\StringTok{"}\OperatorTok{;}  \CommentTok{// simbol pi}

\PreprocessorTok{\#if \_\_STDC\_UTF\_16\_\_}
\NormalTok{pi }\OperatorTok{==} \BaseNTok{0x3C0}\OperatorTok{;}  \CommentTok{// Selalu true}
\PreprocessorTok{\#else}
\NormalTok{pi }\OperatorTok{==} \BaseNTok{0x3C0}\OperatorTok{;}  \CommentTok{// Kemungkinan tidak true}
\PreprocessorTok{\#endif}
\end{Highlighting}
\end{Shaded}

\end{tcolorbox}

\index{char16\_t type@\texttt{char16\_t} type|)} \index{Unicode!UTF-16|)} \index{Unicode!UTF-32|)}

\subsection{Konversi Multibyte}\label{konversi-multibyte}

Kamu bisa mengonversi dari encoding multibyte-mu ke \hl{\mbox{\texttt{char16\_t}}} atau \hl{\mbox{\texttt{char32\_t}}} dengan sejumlah fungsi bantu.

(Seperti kubilang tadi, hasilnya mungkin bukan UTF-16 atau UTF-32 kecuali makro yang sesuai bernilai \hl{\mbox{\texttt{1}}}.)

Semua fungsi ini restartable (yaitu kamu mengoper \hl{\mbox{\texttt{mbstate\_t}}}-mu sendiri), dan semuanya beroperasi karakter demi karakter\footnote{Kurang lebih---urusannya jadi aneh pada encoding UTF-16 dengan beberapa \hl{\mbox{\texttt{char16\_t}}}.}.

\index{mbrtoc16() function@\texttt{mbrtoc16()} function} \index{mbrtoc32() function@\texttt{mbrtoc32()} function} \index{c16rtomb() function@\texttt{c16rtomb()} function} \index{c32rtomb() function@\texttt{c32rtomb()} function}

{\def\LTcaptype{none} % do not increment counter
\begin{longtable}[]{@{}
  >{\raggedright\arraybackslash}p{(\linewidth - 2\tabcolsep) * \real{0.5000}}
  >{\raggedright\arraybackslash}p{(\linewidth - 2\tabcolsep) * \real{0.5000}}@{}}
\toprule\noalign{}
\begin{minipage}[b]{\linewidth}\raggedright
Fungsi Konversi
\end{minipage} & \begin{minipage}[b]{\linewidth}\raggedright
Deskripsi
\end{minipage} \\
\midrule\noalign{}
\endhead
\bottomrule\noalign{}
\endlastfoot
\hl{\mbox{\texttt{mbrtoc16()}}} & Mengonversi karakter multibyte menjadi karakter \hl{\mbox{\texttt{char16\_t}}}. \\
\hl{\mbox{\texttt{mbrtoc32()}}} & Mengonversi karakter multibyte menjadi karakter \hl{\mbox{\texttt{char32\_t}}}. \\
\hl{\mbox{\texttt{c16rtomb()}}} & Mengonversi karakter \hl{\mbox{\texttt{char16\_t}}} menjadi karakter multibyte. \\
\hl{\mbox{\texttt{c32rtomb()}}} & Mengonversi karakter \hl{\mbox{\texttt{char32\_t}}} menjadi karakter multibyte. \\
\end{longtable}
}

\subsection{Library Pihak Ketiga}\label{utf-3rd-party}

Untuk konversi berat antara berbagai encoding spesifik, ada beberapa library matang yang layak dicoba. Perhatikan bahwa aku belum pernah memakai keduanya.

\begin{itemize}
\tightlist
\item
  \href{https://en.wikipedia.org/wiki/Iconv}{iconv}\footnote{https://en.wikipedia.org/wiki/Iconv}---Internationalization Conversion, API standar POSIX yang umum dan tersedia di platform-platform besar.
\item
  \href{http://site.icu-project.org/}{ICU}\footnote{http://site.icu-project.org/}---International Components for Unicode. Setidaknya satu blogger menganggapnya mudah dipakai.
\end{itemize}

Kalau kamu punya library lain yang patut dicatat, kabari aku.

\index{Unicode|)}

\chapter{Keluar dari Program}\label{keluar-dari-program}

\index{Exiting|(}

Ternyata ada banyak cara untuk melakukannya, bahkan ada cara untuk memasang ``hook'' supaya sebuah fungsi dijalankan ketika program keluar.

Di bab ini kita akan menyelam dan melihatnya.

Kita sudah membahas makna kode exit status di bagian \hyperref[exit-status]{Exit Status}, jadi lompat kembali ke sana dan ulas kalau perlu.

Semua fungsi di bagian ini ada di \hl{\mbox{\texttt{\textless{}stdlib.h\textgreater{}}}}.

\section{Keluar Normal}\label{keluar-normal}

Kita mulai dengan cara-cara biasa keluar dari program, lalu melompat ke beberapa yang lebih jarang dan esoterik.

Ketika kamu keluar dari program secara normal, semua I/O stream yang terbuka di-flush dan file sementara dihapus. Pada dasarnya ini keluar yang rapi ketika semuanya dibereskan dan ditangani. Inilah yang kamu inginkan hampir sepanjang waktu kecuali kamu punya alasan lain.

\subsection{\texorpdfstring{Return dari \hl{\mbox{\texttt{main()}}}}{Return dari main()}}\label{return-dari-main}

\index{Exiting!return from main()@return from \texttt{main()}|(}

Kalau kamu perhatikan, \hl{\mbox{\texttt{main()}}} punya tipe kembalian \hl{\mbox{\texttt{int}}}\ldots{} namun aku jarang sekali, kalaupun ada, me-\hl{\mbox{\texttt{return}}} sesuatu dari \hl{\mbox{\texttt{main()}}}.

Ini karena hanya untuk \hl{\mbox{\texttt{main()}}} (dan aku tidak bisa cukup menekankan bahwa kasus khusus ini \emph{hanya} berlaku untuk \hl{\mbox{\texttt{main()}}} dan tidak ada fungsi lain di mana pun) ada \hl{\mbox{\texttt{return\ 0}}} \emph{implisit} kalau eksekusi jatuh melewati ujung.

Kamu bisa \hl{\mbox{\texttt{return}}} secara eksplisit dari \hl{\mbox{\texttt{main()}}} kapan saja, dan sebagian programmer merasa lebih \emph{Benar} untuk selalu menaruh \hl{\mbox{\texttt{return}}} di akhir \hl{\mbox{\texttt{main()}}}. Tapi kalau tidak ditulis, C akan menaruhnya untukmu.

Jadi\ldots{} berikut aturan \hl{\mbox{\texttt{return}}} untuk \hl{\mbox{\texttt{main()}}}:

\begin{itemize}
\tightlist
\item
  Kamu bisa mengembalikan exit status dari \hl{\mbox{\texttt{main()}}} dengan pernyataan \hl{\mbox{\texttt{return}}}. \hl{\mbox{\texttt{main()}}} adalah satu-satunya fungsi dengan perilaku khusus ini. Memakai \hl{\mbox{\texttt{return}}} di fungsi lain hanya kembali dari fungsi itu ke caller-nya.
\item
  Kalau kamu tidak \hl{\mbox{\texttt{return}}} secara eksplisit dan hanya jatuh melewati ujung \hl{\mbox{\texttt{main()}}}, itu sama saja seperti kamu mengembalikan \hl{\mbox{\texttt{0}}} atau \hl{\mbox{\texttt{EXIT\_SUCCESS}}}.
\end{itemize}

\index{Exiting!return from main()@return from \texttt{main()}|)}

\subsection{\texorpdfstring{\hl{\mbox{\texttt{exit()}}}}{exit()}}\label{exit}

\index{Exiting!return from main()@return from \texttt{main()}|)}

Yang ini juga sudah muncul beberapa kali. Kalau kamu memanggil \hl{\mbox{\texttt{exit()}}} dari mana pun di programmu, program akan keluar di titik itu.

Argumen yang kamu oper ke \hl{\mbox{\texttt{exit()}}} adalah exit status.

\subsection{\texorpdfstring{Memasang Exit Handler dengan \hl{\mbox{\texttt{atexit()}}}}{Memasang Exit Handler dengan atexit()}}\label{memasang-exit-handler-dengan-atexit}

\index{atexit() function@\texttt{atexit()} function|(}

Kamu bisa mendaftarkan fungsi yang dipanggil saat program keluar, baik lewat return dari \hl{\mbox{\texttt{main()}}} maupun lewat pemanggilan fungsi \hl{\mbox{\texttt{exit()}}}.

Pemanggilan \hl{\mbox{\texttt{atexit()}}} dengan nama fungsi handler akan melakukannya. Kamu bisa mendaftarkan beberapa exit handler, dan mereka akan dipanggil dalam urutan terbalik dari pendaftaran.

Berikut contohnya:

\begin{tcolorbox}[enhanced jigsaw,left=3pt,breakable,colback=white!97.2549!black,boxrule=0.2mm]

\begin{Shaded}
\begin{Highlighting}[numbers=left,,]
\PreprocessorTok{\#include }\ImportTok{\textless{}stdio.h\textgreater{}}
\PreprocessorTok{\#include }\ImportTok{\textless{}stdlib.h\textgreater{}}

\DataTypeTok{void}\NormalTok{ on\_exit\_1}\OperatorTok{(}\DataTypeTok{void}\OperatorTok{)}
\OperatorTok{\{}
\NormalTok{    printf}\OperatorTok{(}\StringTok{"Exit handler 1 called!}\SpecialCharTok{\textbackslash{}n}\StringTok{"}\OperatorTok{);}
\OperatorTok{\}}

\DataTypeTok{void}\NormalTok{ on\_exit\_2}\OperatorTok{(}\DataTypeTok{void}\OperatorTok{)}
\OperatorTok{\{}
\NormalTok{    printf}\OperatorTok{(}\StringTok{"Exit handler 2 called!}\SpecialCharTok{\textbackslash{}n}\StringTok{"}\OperatorTok{);}
\OperatorTok{\}}

\DataTypeTok{int}\NormalTok{ main}\OperatorTok{(}\DataTypeTok{void}\OperatorTok{)}
\OperatorTok{\{}
\NormalTok{    atexit}\OperatorTok{(}\NormalTok{on\_exit\_1}\OperatorTok{);}
\NormalTok{    atexit}\OperatorTok{(}\NormalTok{on\_exit\_2}\OperatorTok{);}
    
\NormalTok{    printf}\OperatorTok{(}\StringTok{"About to exit...}\SpecialCharTok{\textbackslash{}n}\StringTok{"}\OperatorTok{);}
\OperatorTok{\}}
\end{Highlighting}
\end{Shaded}

\end{tcolorbox}

Dan outputnya:

\begin{tcolorbox}[enhanced jigsaw,left=3pt,breakable,colback=white!97.2549!black,boxrule=0.2mm]

\begin{Shaded}
\begin{Highlighting}[]
\NormalTok{About to exit...}
\NormalTok{Exit handler 2 called!}
\NormalTok{Exit handler 1 called!}
\end{Highlighting}
\end{Shaded}

\end{tcolorbox}

\index{atexit() function@\texttt{atexit()} function|)}

\section{\texorpdfstring{Keluar Lebih Cepat dengan \hl{\mbox{\texttt{quick\_exit()}}}}{Keluar Lebih Cepat dengan quick\_exit()}}\label{keluar-lebih-cepat-dengan-quick_exit}

\index{quick\_exit() function@\texttt{quick\_exit()} function|(}

Ini mirip keluar normal, kecuali:

\begin{itemize}
\tightlist
\item
  File yang terbuka mungkin tidak di-flush.
\item
  File sementara mungkin tidak dihapus.
\item
  Handler \hl{\mbox{\texttt{atexit()}}} tidak akan dipanggil.
\end{itemize}

Tapi ada cara untuk mendaftarkan exit handler: panggil \hl{\mbox{\texttt{at\_quick\_exit()}}} dengan cara yang analog dengan memanggil \hl{\mbox{\texttt{atexit()}}}.

\begin{tcolorbox}[enhanced jigsaw,left=3pt,breakable,colback=white!97.2549!black,boxrule=0.2mm]

\begin{Shaded}
\begin{Highlighting}[numbers=left,,]
\PreprocessorTok{\#include }\ImportTok{\textless{}stdio.h\textgreater{}}
\PreprocessorTok{\#include }\ImportTok{\textless{}stdlib.h\textgreater{}}

\DataTypeTok{void}\NormalTok{ on\_quick\_exit\_1}\OperatorTok{(}\DataTypeTok{void}\OperatorTok{)}
\OperatorTok{\{}
\NormalTok{    printf}\OperatorTok{(}\StringTok{"Quick exit handler 1 called!}\SpecialCharTok{\textbackslash{}n}\StringTok{"}\OperatorTok{);}
\OperatorTok{\}}

\DataTypeTok{void}\NormalTok{ on\_quick\_exit\_2}\OperatorTok{(}\DataTypeTok{void}\OperatorTok{)}
\OperatorTok{\{}
\NormalTok{    printf}\OperatorTok{(}\StringTok{"Quick exit handler 2 called!}\SpecialCharTok{\textbackslash{}n}\StringTok{"}\OperatorTok{);}
\OperatorTok{\}}

\DataTypeTok{void}\NormalTok{ on\_exit}\OperatorTok{(}\DataTypeTok{void}\OperatorTok{)}
\OperatorTok{\{}
\NormalTok{    printf}\OperatorTok{(}\StringTok{"Normal exit{-}{-}I won\textquotesingle{}t be called!}\SpecialCharTok{\textbackslash{}n}\StringTok{"}\OperatorTok{);}
\OperatorTok{\}}

\DataTypeTok{int}\NormalTok{ main}\OperatorTok{(}\DataTypeTok{void}\OperatorTok{)}
\OperatorTok{\{}
\NormalTok{    at\_quick\_exit}\OperatorTok{(}\NormalTok{on\_quick\_exit\_1}\OperatorTok{);}
\NormalTok{    at\_quick\_exit}\OperatorTok{(}\NormalTok{on\_quick\_exit\_2}\OperatorTok{);}

\NormalTok{    atexit}\OperatorTok{(}\NormalTok{on\_exit}\OperatorTok{);}  \CommentTok{// Ini tidak akan dipanggil}

\NormalTok{    printf}\OperatorTok{(}\StringTok{"About to quick exit...}\SpecialCharTok{\textbackslash{}n}\StringTok{"}\OperatorTok{);}

\NormalTok{    quick\_exit}\OperatorTok{(}\DecValTok{0}\OperatorTok{);}
\OperatorTok{\}}
\end{Highlighting}
\end{Shaded}

\end{tcolorbox}

Yang memberi output ini:

\begin{tcolorbox}[enhanced jigsaw,left=3pt,breakable,colback=white!97.2549!black,boxrule=0.2mm]

\begin{Shaded}
\begin{Highlighting}[]
\NormalTok{About to quick exit...}
\NormalTok{Quick exit handler 2 called!}
\NormalTok{Quick exit handler 1 called!}
\end{Highlighting}
\end{Shaded}

\end{tcolorbox}

Cara kerjanya persis seperti \hl{\mbox{\texttt{exit()}}}/\hl{\mbox{\texttt{atexit()}}}, kecuali bahwa pem-flush-an file dan pembersihan mungkin tidak dilakukan.

\index{quick\_exit() function@\texttt{quick\_exit()} function|)}

\section{\texorpdfstring{Nuke it from Orbit: \hl{\mbox{\texttt{\_Exit()}}}}{Nuke it from Orbit: \_Exit()}}\label{nuke-it-from-orbit-_exit}

\index{\_Exit() function@\texttt{\_Exit()} function|(}

Memanggil \hl{\mbox{\texttt{\_Exit()}}} langsung keluar, titik. Tidak ada fungsi callback on-exit yang dijalankan. File tidak di-flush. File sementara tidak dihapus.

Pakai ini kalau kamu harus keluar \emph{detik ini juga, demi apa pun}.

\section{\texorpdfstring{Keluar Kadang-kadang: \hl{\mbox{\texttt{assert()}}}}{Keluar Kadang-kadang: assert()}}\label{keluar-kadang-kadang-assert}

Pernyataan \hl{\mbox{\texttt{assert()}}} dipakai untuk menegaskan bahwa sesuatu harus benar, kalau tidak program akan keluar.

Developer sering memakai assert untuk menangkap error jenis Seharusnya-Tidak-Mungkin-Terjadi.

\begin{tcolorbox}[enhanced jigsaw,left=3pt,breakable,colback=white!97.2549!black,boxrule=0.2mm]

\begin{Shaded}
\begin{Highlighting}[]
\PreprocessorTok{\#define PI }\FloatTok{3.14159}

\NormalTok{assert}\OperatorTok{(}\NormalTok{PI }\OperatorTok{\textgreater{}} \DecValTok{3}\OperatorTok{);}   \CommentTok{// Benar saja, jadi lanjut}
\end{Highlighting}
\end{Shaded}

\end{tcolorbox}

versus:

\begin{tcolorbox}[enhanced jigsaw,left=3pt,breakable,colback=white!97.2549!black,boxrule=0.2mm]

\begin{Shaded}
\begin{Highlighting}[]
\NormalTok{goats }\OperatorTok{{-}=} \DecValTok{100}\OperatorTok{;}

\NormalTok{assert}\OperatorTok{(}\NormalTok{goats }\OperatorTok{\textgreater{}=} \DecValTok{0}\OperatorTok{);}  \CommentTok{// Tidak boleh ada kambing negatif}
\end{Highlighting}
\end{Shaded}

\end{tcolorbox}

Dalam kasus itu, kalau kujalankan dan \hl{\mbox{\texttt{goats}}} turun di bawah \hl{\mbox{\texttt{0}}}, ini yang terjadi:

\begin{tcolorbox}[enhanced jigsaw,left=3pt,breakable,colback=white!97.2549!black,boxrule=0.2mm]

\begin{Shaded}
\begin{Highlighting}[]
\NormalTok{goat\_counter: goat\_counter.c:8: main: Assertion \textasciigrave{}goats \textgreater{}= 0\textquotesingle{} failed.}
\NormalTok{Aborted}
\end{Highlighting}
\end{Shaded}

\end{tcolorbox}

dan aku dilempar kembali ke command line.

Ini kurang ramah pengguna, jadi hanya dipakai untuk hal-hal yang tidak akan pernah dilihat pengguna. Dan sering orang \hyperref[my-assert]{menulis makro assert sendiri yang lebih mudah dimatikan}.

\index{\_Exit() function@\texttt{\_Exit()} function|)}

\section{\texorpdfstring{Keluar Tidak Normal: \hl{\mbox{\texttt{abort()}}}}{Keluar Tidak Normal: abort()}}\label{keluar-tidak-normal-abort}

\index{abort() function@\texttt{abort()} function|(}

Kamu bisa memakai ini kalau ada yang salah sekali dan ingin menandakannya kepada lingkungan luar. Ini juga tidak selalu membersihkan file yang terbuka, dan sebagainya.

Aku jarang melihat ini dipakai.

Sedikit bocoran tentang \emph{signal}: ini sebenarnya bekerja dengan memunculkan \index{SIGABRT signal@\texttt{SIGABRT} signal} \hl{\mbox{\texttt{SIGABRT}}} yang akan mengakhiri proses.

Apa yang terjadi setelah itu tergantung sistem, tapi di sistem mirip Unix, umumnya \href{https://en.wikipedia.org/wiki/Core_dump}{dump core}\footnote{https://en.wikipedia.org/wiki/Core\_dump} saat program berakhir.

\index{abort() function@\texttt{abort()} function|)} \index{Exiting|)}

\chapter{Penanganan Signal}\label{penanganan-signal}

\index{Signal handling|(}

Sebelum mulai, aku hanya ingin menyarankan agar kamu pada umumnya mengabaikan seluruh bab ini dan memakai fungsi penanganan signal bawaan OS-mu yang (sangat mungkin) lebih unggul. Sistem mirip Unix punya keluarga fungsi \index{sigaction() function@\texttt{sigaction()} function} \hl{\mbox{\texttt{sigaction()}}}, dan Windows punya\ldots{} entah apa yang dilakukannya\footnote{Tampaknya jauh di dalamnya Windows sama sekali tidak mengenal signal ala Unix, dan untuk aplikasi konsol signal itu disimulasikan.}.

Setelah itu beres, apa itu signal?

\section{Apa Itu Signal?}\label{apa-itu-signal}

Sebuah \emph{signal} di-\emph{raise} pada berbagai kejadian eksternal. Programmu bisa diatur agar diinterupsi untuk \emph{menangani} signal itu, dan, kalau mau, melanjutkan dari tempat ia berhenti setelah signal selesai ditangani.

Bayangkan seperti fungsi yang otomatis dipanggil ketika salah satu kejadian eksternal ini terjadi.

Apa saja kejadian itu? Di sistemmu kemungkinan ada banyak sekali, tapi di spesifikasi C hanya ada beberapa:

\index{SIGABRT signal@\texttt{SIGABRT} signal} \index{SIGFPE signal@\texttt{SIGFPE} signal} \index{SIGILL signal@\texttt{SIGILL} signal} \index{SIGINT signal@\texttt{SIGINT} signal} \index{SIGSEGV signal@\texttt{SIGSEGV} signal} \index{SIGTERM signal@\texttt{SIGTERM} signal}

{\def\LTcaptype{none} % do not increment counter
\begin{longtable}[]{@{}
  >{\raggedright\arraybackslash}p{(\linewidth - 2\tabcolsep) * \real{0.1111}}
  >{\raggedright\arraybackslash}p{(\linewidth - 2\tabcolsep) * \real{0.8889}}@{}}
\toprule\noalign{}
\begin{minipage}[b]{\linewidth}\raggedright
Signal
\end{minipage} & \begin{minipage}[b]{\linewidth}\raggedright
Deskripsi
\end{minipage} \\
\midrule\noalign{}
\endhead
\bottomrule\noalign{}
\endlastfoot
\hl{\mbox{\texttt{SIGABRT}}} & Terminasi abnormal---yang terjadi ketika \hl{\mbox{\texttt{abort()}}} dipanggil. \\
\hl{\mbox{\texttt{SIGFPE}}} & Exception floating point. \\
\hl{\mbox{\texttt{SIGILL}}} & Instruksi ilegal. \\
\hl{\mbox{\texttt{SIGINT}}} & Interupsi---biasanya akibat \hl{\mbox{\texttt{CTRL-C}}} ditekan. \\
\hl{\mbox{\texttt{SIGSEGV}}} & ``Segmentation Violation'': akses memori yang tidak valid. \\
\hl{\mbox{\texttt{SIGTERM}}} & Permintaan terminasi. \\
\end{longtable}
}

Kamu bisa mengatur programmu untuk mengabaikan, menangani, atau membiarkan aksi bawaan untuk masing-masing signal itu dengan fungsi \hl{\mbox{\texttt{signal()}}}.

\section{\texorpdfstring{Menangani Signal dengan \hl{\mbox{\texttt{signal()}}}}{Menangani Signal dengan signal()}}\label{menangani-signal-dengan-signal}

\index{signal() function@\texttt{signal()} function|(}

Pemanggilan \hl{\mbox{\texttt{signal()}}} menerima dua parameter: signal yang dimaksud, dan aksi yang diambil ketika signal itu di-raise.

Aksinya bisa berupa salah satu dari tiga hal:

\begin{itemize}
\tightlist
\item
  Pointer ke fungsi handler.
\item
  \index{SIG\_IGN macro@\texttt{SIG\_IGN} macro|(}\hl{\mbox{\texttt{SIG\_IGN}}} untuk mengabaikan signal.
\item
  \index{SIG\_DFL macro@\texttt{SIG\_DFL} macro}\hl{\mbox{\texttt{SIG\_DFL}}} untuk mengembalikan handler bawaan signal itu.
\end{itemize}

Mari kita tulis program yang tidak bisa di-\hl{\mbox{\texttt{CTRL-C}}}. (Tenang saja---pada program berikut, kamu juga bisa menekan \hl{\mbox{\texttt{RETURN}}} dan ia akan keluar.)

\index{SIGINT signal@\texttt{SIGINT} signal|(}

\begin{tcolorbox}[enhanced jigsaw,left=3pt,breakable,colback=white!97.2549!black,boxrule=0.2mm]

\begin{Shaded}
\begin{Highlighting}[numbers=left,,]
\PreprocessorTok{\#include }\ImportTok{\textless{}stdio.h\textgreater{}}
\PreprocessorTok{\#include }\ImportTok{\textless{}signal.h\textgreater{}}

\DataTypeTok{int}\NormalTok{ main}\OperatorTok{(}\DataTypeTok{void}\OperatorTok{)}
\OperatorTok{\{}
    \DataTypeTok{char}\NormalTok{ s}\OperatorTok{[}\DecValTok{1024}\OperatorTok{];}

\NormalTok{    signal}\OperatorTok{(}\NormalTok{SIGINT}\OperatorTok{,}\NormalTok{ SIG\_IGN}\OperatorTok{);}    \CommentTok{// Abaikan SIGINT, yang disebabkan oleh \^{}C}

\NormalTok{    printf}\OperatorTok{(}\StringTok{"Try hitting \^{}C... (hit RETURN to exit)}\SpecialCharTok{\textbackslash{}n}\StringTok{"}\OperatorTok{);}

    \CommentTok{// Tunggu satu baris input supaya program tidak langsung keluar}
\NormalTok{    fgets}\OperatorTok{(}\NormalTok{s}\OperatorTok{,} \KeywordTok{sizeof}\NormalTok{ s}\OperatorTok{,}\NormalTok{ stdin}\OperatorTok{);}
\OperatorTok{\}}
\end{Highlighting}
\end{Shaded}

\end{tcolorbox}

Lihat baris 8---kita menyuruh program mengabaikan \hl{\mbox{\texttt{SIGINT}}}, signal interupsi yang di-raise ketika \hl{\mbox{\texttt{CTRL-C}}} ditekan. Sebanyak apa pun kamu menekannya, signal itu tetap diabaikan. Kalau baris 8 dijadikan komentar, kamu akan melihat bahwa kamu bisa \hl{\mbox{\texttt{CTRL-C}}} sesuka hati dan langsung keluar dari program.

\index{SIGINT signal@\texttt{SIGINT} signal|)} \index{SIG\_IGN macro@\texttt{SIG\_IGN} macro|)}

\section{Menulis Signal Handler}\label{menulis-signal-handler}

Tadi kusebut bahwa kamu juga bisa menulis fungsi handler yang dipanggil ketika signal di-raise.

Ini cukup lugas, tapi kemampuannya juga sangat terbatas menurut spesifikasi.

\index{signal() function@\texttt{signal()} function|(}

Sebelum mulai, mari kita lihat function prototype untuk pemanggilan \hl{\mbox{\texttt{signal()}}}:

\begin{tcolorbox}[enhanced jigsaw,left=3pt,breakable,colback=white!97.2549!black,boxrule=0.2mm]

\begin{Shaded}
\begin{Highlighting}[]
\DataTypeTok{void} \OperatorTok{(*}\NormalTok{signal}\OperatorTok{(}\DataTypeTok{int}\NormalTok{ sig}\OperatorTok{,} \DataTypeTok{void} \OperatorTok{(*}\NormalTok{func}\OperatorTok{)(}\DataTypeTok{int}\OperatorTok{)))(}\DataTypeTok{int}\OperatorTok{);}
\end{Highlighting}
\end{Shaded}

\end{tcolorbox}

Mudah dibaca, kan?

\emph{SALAH!} \hl{\mbox{\texttt{:)}}}

Mari kita bongkar sebentar untuk latihan.

\hl{\mbox{\texttt{signal()}}} menerima dua argumen: integer \hl{\mbox{\texttt{sig}}} yang mewakili signal, dan pointer \hl{\mbox{\texttt{func}}} ke handler (handler mengembalikan \hl{\mbox{\texttt{void}}} dan menerima \hl{\mbox{\texttt{int}}} sebagai argumen), disorot di bawah:

\begin{tcolorbox}[enhanced jigsaw,left=3pt,breakable,colback=white!97.2549!black,boxrule=0.2mm]

\begin{Shaded}
\begin{Highlighting}[]
\NormalTok{                sig          func}
              \OperatorTok{|{-}{-}{-}{-}{-}|}  \OperatorTok{|{-}{-}{-}{-}{-}{-}{-}{-}{-}{-}{-}{-}{-}{-}{-}|}
\DataTypeTok{void} \OperatorTok{(*}\NormalTok{signal}\OperatorTok{(}\DataTypeTok{int}\NormalTok{ sig}\OperatorTok{,} \DataTypeTok{void} \OperatorTok{(*}\NormalTok{func}\OperatorTok{)(}\DataTypeTok{int}\OperatorTok{)))(}\DataTypeTok{int}\OperatorTok{);}
\end{Highlighting}
\end{Shaded}

\end{tcolorbox}

\index{signal() function@\texttt{signal()} function|)}

Pada dasarnya, kita akan mengoper nomor signal yang ingin kita tangkap, dan mengoper pointer ke fungsi berbentuk:

\begin{tcolorbox}[enhanced jigsaw,left=3pt,breakable,colback=white!97.2549!black,boxrule=0.2mm]

\begin{Shaded}
\begin{Highlighting}[]
\DataTypeTok{void}\NormalTok{ f}\OperatorTok{(}\DataTypeTok{int}\NormalTok{ x}\OperatorTok{);}
\end{Highlighting}
\end{Shaded}

\end{tcolorbox}

yang akan melakukan penangkapan sebenarnya.

Nah---bagaimana dengan sisa prototype itu? Isinya pada dasarnya semua tipe kembalian. Begini, \hl{\mbox{\texttt{signal()}}} akan mengembalikan apa pun yang kamu oper sebagai \hl{\mbox{\texttt{func}}} kalau berhasil\ldots{} artinya ia mengembalikan pointer ke fungsi yang mengembalikan \hl{\mbox{\texttt{void}}} dan menerima \hl{\mbox{\texttt{int}}} sebagai argumen.

\begin{tcolorbox}[enhanced jigsaw,left=3pt,breakable,colback=white!97.2549!black,boxrule=0.2mm]

\begin{Shaded}
\begin{Highlighting}[]
\NormalTok{returned}
\NormalTok{function    indicates we}\CharTok{\textquotesingle{}r}\ErrorTok{e              and}
\NormalTok{returns     returning a                  that function}
\DataTypeTok{void}\NormalTok{        pointer to function          takes an }\DataTypeTok{int}
\OperatorTok{|{-}{-}|}        \OperatorTok{|}                                   \OperatorTok{|{-}{-}{-}|}
\DataTypeTok{void}       \OperatorTok{(*}\NormalTok{signal}\OperatorTok{(}\DataTypeTok{int}\NormalTok{ sig}\OperatorTok{,} \DataTypeTok{void} \OperatorTok{(*}\NormalTok{func}\OperatorTok{)(}\DataTypeTok{int}\OperatorTok{)))(}\DataTypeTok{int}\OperatorTok{);}
\end{Highlighting}
\end{Shaded}

\end{tcolorbox}

Ia juga bisa mengembalikan \index{SIG\_ERR macro@\texttt{SIG\_ERR} macro} \hl{\mbox{\texttt{SIG\_ERR}}} kalau terjadi error.

Mari kita buat contoh yang membuatmu harus menekan \hl{\mbox{\texttt{CTRL-C}}} dua kali untuk keluar.

Aku ingin menegaskan bahwa program ini melibatkan undefined behavior dalam beberapa hal. Tapi kemungkinan besar program ini berfungsi di sistemmu, dan sulit membuat demo non-trivial yang portabel.

\index{signal() function@\texttt{signal()} function|(} \index{SIG\_INT signal@\texttt{SIG\_INT} signal|(}

\begin{tcolorbox}[enhanced jigsaw,left=3pt,breakable,colback=white!97.2549!black,boxrule=0.2mm]

\begin{Shaded}
\begin{Highlighting}[numbers=left,,]
\PreprocessorTok{\#include }\ImportTok{\textless{}stdio.h\textgreater{}}
\PreprocessorTok{\#include }\ImportTok{\textless{}stdlib.h\textgreater{}}
\PreprocessorTok{\#include }\ImportTok{\textless{}signal.h\textgreater{}}

\DataTypeTok{int}\NormalTok{ count }\OperatorTok{=} \DecValTok{0}\OperatorTok{;}

\DataTypeTok{void}\NormalTok{ sigint\_handler}\OperatorTok{(}\DataTypeTok{int}\NormalTok{ signum}\OperatorTok{)}
\OperatorTok{\{}
    \CommentTok{// Compiler boleh menjalankan:}
    \CommentTok{//}
    \CommentTok{//   signal(signum, SIG\_DFL)}
    \CommentTok{//}
    \CommentTok{// saat handler dipanggil. Jadi kita reset handler{-}nya di sini:}
\NormalTok{    signal}\OperatorTok{(}\NormalTok{SIGINT}\OperatorTok{,}\NormalTok{ sigint\_handler}\OperatorTok{);}

    \OperatorTok{(}\DataTypeTok{void}\OperatorTok{)}\NormalTok{signum}\OperatorTok{;}   \CommentTok{// Hilangkan peringatan variabel tak terpakai}

\NormalTok{    count}\OperatorTok{++;}                       \CommentTok{// Undefined behavior}
\NormalTok{    printf}\OperatorTok{(}\StringTok{"Count: }\SpecialCharTok{\%d\textbackslash{}n}\StringTok{"}\OperatorTok{,}\NormalTok{ count}\OperatorTok{);}  \CommentTok{// Undefined behavior}

    \ControlFlowTok{if} \OperatorTok{(}\NormalTok{count }\OperatorTok{==} \DecValTok{2}\OperatorTok{)} \OperatorTok{\{}
\NormalTok{        printf}\OperatorTok{(}\StringTok{"Exiting!}\SpecialCharTok{\textbackslash{}n}\StringTok{"}\OperatorTok{);}      \CommentTok{// Undefined behavior}
\NormalTok{        exit}\OperatorTok{(}\DecValTok{0}\OperatorTok{);}
    \OperatorTok{\}}
\OperatorTok{\}}

\DataTypeTok{int}\NormalTok{ main}\OperatorTok{(}\DataTypeTok{void}\OperatorTok{)}
\OperatorTok{\{}
\NormalTok{    signal}\OperatorTok{(}\NormalTok{SIGINT}\OperatorTok{,}\NormalTok{ sigint\_handler}\OperatorTok{);}

\NormalTok{    printf}\OperatorTok{(}\StringTok{"Try hitting \^{}C...}\SpecialCharTok{\textbackslash{}n}\StringTok{"}\OperatorTok{);}

    \ControlFlowTok{for}\OperatorTok{(;;);}  \CommentTok{// Tunggu di sini selamanya}
\OperatorTok{\}}
\end{Highlighting}
\end{Shaded}

\end{tcolorbox}

\index{SIG\_INT signal@\texttt{SIG\_INT} signal|)}

Salah satu hal yang akan kamu lihat: di baris 14 kita mereset signal handler. Ini karena C punya opsi untuk mereset signal handler ke perilaku \index{SIG\_DFL macro@\texttt{SIG\_DFL} macro} \hl{\mbox{\texttt{SIG\_DFL}}}-nya sebelum menjalankan handler kustommu. Dengan kata lain, handler itu bisa jadi sekali pakai. Jadi hal pertama yang kita lakukan adalah mereset-nya agar signal berikutnya juga tertangani.

Dalam kasus ini kita mengabaikan nilai kembalian \hl{\mbox{\texttt{signal()}}}. Kalau sebelumnya kita sudah mengatur handler lain, ia akan mengembalikan pointer ke handler itu, yang bisa kita ambil seperti ini:

\begin{tcolorbox}[enhanced jigsaw,left=3pt,breakable,colback=white!97.2549!black,boxrule=0.2mm]

\begin{Shaded}
\begin{Highlighting}[]
\CommentTok{// old\_handler bertipe "pointer ke fungsi yang menerima satu parameter}
\CommentTok{// int dan mengembalikan void":}

\DataTypeTok{void} \OperatorTok{(*}\NormalTok{old\_handler}\OperatorTok{)(}\DataTypeTok{int}\OperatorTok{);}

\NormalTok{old\_handler }\OperatorTok{=}\NormalTok{ signal}\OperatorTok{(}\NormalTok{SIGINT}\OperatorTok{,}\NormalTok{ sigint\_handler}\OperatorTok{);}
\end{Highlighting}
\end{Shaded}

\end{tcolorbox}

\index{signal() function@\texttt{signal()} function|)}

Meski begitu, aku tidak yakin ada kasus pemakaian umum untuk ini. Tapi kalau kamu butuh handler lama karena suatu alasan, kamu bisa mendapatkannya dengan cara itu.

Catatan singkat soal baris 16---itu hanya untuk memberi tahu compiler agar tidak memperingatkan bahwa variabel ini tidak dipakai. Seperti berkata, ``Aku tahu aku tidak memakainya; tidak perlu memperingatkanku.''

Dan terakhir, kamu akan melihat bahwa aku menandai undefined behavior di beberapa tempat. Soal itu dibahas di bagian berikutnya.

\section{Apa yang Sebenarnya Bisa Kita Lakukan?}\label{apa-yang-sebenarnya-bisa-kita-lakukan}

Ternyata kita cukup terbatas dalam hal apa yang bisa dan tidak bisa dilakukan di signal handler. Ini salah satu alasan aku bilang kamu tidak perlu repot dengan ini dan sebaiknya memakai penanganan signal OS sendiri (misalnya \index{sigaction() function@\texttt{sigaction()} function} \hl{\mbox{\texttt{sigaction()}}} untuk sistem mirip Unix).

Wikipedia bahkan menyebut satu-satunya hal yang benar-benar portabel adalah memanggil \hl{\mbox{\texttt{signal()}}} dengan \hl{\mbox{\texttt{SIG\_IGN}}} atau \hl{\mbox{\texttt{SIG\_DFL}}}, itu saja.

Berikut yang \textbf{tidak bisa} kita lakukan secara portabel:

\index{Signal handling-!limitations|(}

\begin{itemize}
\tightlist
\item
  Memanggil fungsi standard library apa pun.

  \begin{itemize}
  \tightlist
  \item
    Seperti \hl{\mbox{\texttt{printf()}}}, misalnya.
  \item
    Kurasa mungkin aman memanggil fungsi restartable/reentrant, tapi spesifikasi tidak memberikan kebebasan itu.
  \end{itemize}
\item
  Mengambil atau mengisi nilai dari variabel lokal \hl{\mbox{\texttt{static}}}, file scope, atau thread-local.

  \begin{itemize}
  \tightlist
  \item
    Kecuali kalau itu atomic object yang lock-free atau\ldots{}
  \item
    Kamu menugaskan ke variabel bertipe \hl{\mbox{\texttt{volatile\ sig\_atomic\_t}}}.
  \end{itemize}
\end{itemize}

\index{sig\_atomic\_t type@\texttt{sig\_atomic\_t} type|(}

Bagian terakhir itu--\hl{\mbox{\texttt{sig\_atomic\_t}}}--adalah tiketmu untuk mengeluarkan data dari signal handler. (Kecuali kamu ingin memakai atomic object lock-free, yang di luar cakupan bagian ini\footnote{Membingungkan, \hl{\mbox{\texttt{sig\_atomic\_t}}} lebih tua daripada atomic lock-free dan tidak sama dengannya.}.) Ia adalah tipe integer yang mungkin signed, mungkin tidak. Dan ia dibatasi oleh apa yang bisa dimasukkan ke dalamnya.

Kamu bisa melihat nilai minimum dan maksimum yang diizinkan pada makro \hl{\mbox{\texttt{SIG\_ATOMIC\_MIN}}} dan \hl{\mbox{\texttt{SIG\_ATOMIC\_MAX}}}\footnote{Kalau \hl{\mbox{\texttt{sig\_action\_t}}} signed, rentangnya setidaknya \hl{\mbox{\texttt{-127}}} sampai \hl{\mbox{\texttt{127}}}. Kalau unsigned, setidaknya \hl{\mbox{\texttt{0}}} sampai \hl{\mbox{\texttt{255}}}.}.

Membingungkannya, spesifikasi juga menyatakan kamu tidak boleh merujuk ``to any object with static or thread storage duration that is not a lock-free atomic object other than by assigning a value to an object declared as \hl{\mbox{\texttt{volatile\ sig\_atomic\_t}}} {[}\ldots{]}''

Menurut pemahamanku, kamu tidak boleh membaca atau menulis apa pun yang bukan atomic object lock-free. Kamu juga boleh menugaskan ke object \hl{\mbox{\texttt{volatile\ sig\_atomic\_t}}}.

Tapi bolehkah membacanya? Jujur, aku tidak melihat kenapa tidak, hanya saja spesifikasi sangat tegas menyebut menugaskan ke. Tapi kalau kamu harus membacanya dan membuat keputusan berdasarkan nilainya, kamu mungkin membuka peluang terjadinya semacam race condition.

\index{Signal handling-!limitations|)}

Dengan itu dalam pikiran, kita bisa menulis ulang kode ``tekan \hl{\mbox{\texttt{CTRL-C}}} dua kali untuk keluar'' agar sedikit lebih portabel, walaupun outputnya jadi lebih singkat.

Mari kita ubah handler \hl{\mbox{\texttt{SIGINT}}} kita agar tidak melakukan apa-apa selain menaikkan nilai bertipe \hl{\mbox{\texttt{volatile\ sig\_atomic\_t}}}. Jadi ia akan menghitung berapa kali \hl{\mbox{\texttt{CTRL-C}}} ditekan.

Lalu di loop utama, kita periksa apakah penghitung itu melebihi \hl{\mbox{\texttt{2}}}, dan keluar kalau ya.

\begin{tcolorbox}[enhanced jigsaw,left=3pt,breakable,colback=white!97.2549!black,boxrule=0.2mm]

\begin{Shaded}
\begin{Highlighting}[numbers=left,,]
\PreprocessorTok{\#include }\ImportTok{\textless{}stdio.h\textgreater{}}
\PreprocessorTok{\#include }\ImportTok{\textless{}signal.h\textgreater{}}

\DataTypeTok{volatile} \DataTypeTok{sig\_atomic\_t}\NormalTok{ count }\OperatorTok{=} \DecValTok{0}\OperatorTok{;}

\DataTypeTok{void}\NormalTok{ sigint\_handler}\OperatorTok{(}\DataTypeTok{int}\NormalTok{ signum}\OperatorTok{)}
\OperatorTok{\{}
    \OperatorTok{(}\DataTypeTok{void}\OperatorTok{)}\NormalTok{signum}\OperatorTok{;}                    \CommentTok{// Peringatan variabel tak terpakai}

\NormalTok{    signal}\OperatorTok{(}\NormalTok{SIGINT}\OperatorTok{,}\NormalTok{ sigint\_handler}\OperatorTok{);}  \CommentTok{// Reset signal handler}

\NormalTok{    count}\OperatorTok{++;}                         \CommentTok{// Undefined behavior}
\OperatorTok{\}}

\DataTypeTok{int}\NormalTok{ main}\OperatorTok{(}\DataTypeTok{void}\OperatorTok{)}
\OperatorTok{\{}
\NormalTok{    signal}\OperatorTok{(}\NormalTok{SIGINT}\OperatorTok{,}\NormalTok{ sigint\_handler}\OperatorTok{);}

\NormalTok{    printf}\OperatorTok{(}\StringTok{"Hit \^{}C twice to exit.}\SpecialCharTok{\textbackslash{}n}\StringTok{"}\OperatorTok{);}

    \ControlFlowTok{while}\OperatorTok{(}\NormalTok{count }\OperatorTok{\textless{}} \DecValTok{2}\OperatorTok{);}
\OperatorTok{\}}
\end{Highlighting}
\end{Shaded}

\end{tcolorbox}

\index{sig\_atomic\_t type@\texttt{sig\_atomic\_t} type|)}

Undefined behavior lagi? Menurutku iya, karena kita harus membaca nilainya untuk menaikkan dan menyimpannya. Thread lain bisa saja mengutak-atik \hl{\mbox{\texttt{count}}} dan bikin kesal. Tapi pada kasus sederhana ini, tidak ada thread lain yang melakukannya, jadi pelanggaran ini bisa dimaafkan dan kita nikmati saja demonya.

Kalau kita hanya ingin menunda keluar satu kali penekanan \hl{\mbox{\texttt{CTRL-C}}}, itu bisa dilakukan tanpa terlalu repot. Tapi penundaan lebih banyak membutuhkan perangkaian fungsi yang konyol.

Yang akan kita lakukan: tangani sekali, dan handler akan mereset signal ke perilaku bawaannya (yaitu keluar):

\index{SIG\_DFL macro@\texttt{SIG\_DFL} macro|(}

\begin{tcolorbox}[enhanced jigsaw,left=3pt,breakable,colback=white!97.2549!black,boxrule=0.2mm]

\begin{Shaded}
\begin{Highlighting}[numbers=left,,]
\PreprocessorTok{\#include }\ImportTok{\textless{}stdio.h\textgreater{}}
\PreprocessorTok{\#include }\ImportTok{\textless{}signal.h\textgreater{}}

\DataTypeTok{void}\NormalTok{ sigint\_handler}\OperatorTok{(}\DataTypeTok{int}\NormalTok{ signum}\OperatorTok{)}
\OperatorTok{\{}
    \OperatorTok{(}\DataTypeTok{void}\OperatorTok{)}\NormalTok{signum}\OperatorTok{;}                      \CommentTok{// Peringatan variabel tak terpakai}
\NormalTok{    signal}\OperatorTok{(}\NormalTok{SIGINT}\OperatorTok{,}\NormalTok{ SIG\_DFL}\OperatorTok{);}           \CommentTok{// Reset signal handler}
\OperatorTok{\}}

\DataTypeTok{int}\NormalTok{ main}\OperatorTok{(}\DataTypeTok{void}\OperatorTok{)}
\OperatorTok{\{}
\NormalTok{    signal}\OperatorTok{(}\NormalTok{SIGINT}\OperatorTok{,}\NormalTok{ sigint\_handler}\OperatorTok{);}

\NormalTok{    printf}\OperatorTok{(}\StringTok{"Hit \^{}C twice to exit.}\SpecialCharTok{\textbackslash{}n}\StringTok{"}\OperatorTok{);}

    \ControlFlowTok{while}\OperatorTok{(}\DecValTok{1}\OperatorTok{);}
\OperatorTok{\}}
\end{Highlighting}
\end{Shaded}

\end{tcolorbox}

\index{SIG\_DFL macro@\texttt{SIG\_DFL} macro|)}

Nanti saat membahas variabel atomic lock-free, kita akan melihat cara memperbaiki versi \hl{\mbox{\texttt{count}}} (dengan asumsi variabel atomic lock-free tersedia di sistemmu).

Inilah alasan di awal tadi aku menyarankan melihat sistem signal bawaan OS-mu sebagai alternatif yang kemungkinan lebih unggul.

\section{\texorpdfstring{Teman Tidak Membiarkan Temannya \hl{\mbox{\texttt{signal()}}}}{Teman Tidak Membiarkan Temannya signal()}}\label{teman-tidak-membiarkan-temannya-signal}

Sekali lagi, pakai penanganan signal bawaan OS-mu atau padanannya. Itu tidak ada di spesifikasi, kurang portabel, tapi kemungkinan jauh lebih mumpuni. Lagi pula OS-mu mungkin punya sejumlah signal yang tidak ada di spesifikasi C. Dan menulis kode portabel dengan \hl{\mbox{\texttt{signal()}}} memang sulit.

\index{signal() function@\texttt{signal()} function|)} \index{Signal handling|)}

\chapter{Variable-Length Array (VLA)}\label{variable-length-array-vla}

\index{Variable-length array|(}

C menyediakan cara untuk mendeklarasikan array yang ukurannya ditentukan saat runtime. Dengan begitu kamu mendapat manfaat ukuran dinamis saat runtime seperti pada \hl{\mbox{\texttt{malloc()}}}, tapi tanpa perlu repot men-\hl{\mbox{\texttt{free()}}} memorinya sesudahnya.

Nah, banyak orang tidak menyukai VLA. Misalnya, VLA dilarang di kernel Linux. Alasannya akan dibahas lebih jauh \hyperref[vla-general-issues]{nanti}.

\index{\_\_STDC\_NO\_VLA\_\_ macro@\texttt{\_\_STDC\_NO\_VLA\_\_} macro|(}

Ini fitur bahasa yang opsional. Makro \hl{\mbox{\texttt{\_\_STDC\_NO\_VLA\_\_}}} bernilai \hl{\mbox{\texttt{1}}} kalau VLA \emph{tidak} ada. (VLA wajib di C99, lalu menjadi opsional di C11.)

\begin{tcolorbox}[enhanced jigsaw,left=3pt,breakable,colback=white!97.2549!black,boxrule=0.2mm]

\begin{Shaded}
\begin{Highlighting}[]
\PreprocessorTok{\#if \_\_STDC\_NO\_VLA\_\_ == 1}
   \PreprocessorTok{\#error Sorry, need VLAs for this program!}
\PreprocessorTok{\#endif}
\end{Highlighting}
\end{Shaded}

\end{tcolorbox}

\index{\_\_STDC\_NO\_VLA\_\_ macro@\texttt{\_\_STDC\_NO\_VLA\_\_} macro|)}

Tapi karena GCC maupun Clang tidak repot-repot mendefinisikan makro ini, manfaatnya mungkin terbatas.

Kita mulai dulu dengan contoh, lalu kita cari setan di detailnya.

\section{Dasar-Dasarnya}\label{dasar-dasarnya}

Array biasa dideklarasikan dengan ukuran konstan, seperti ini:

\begin{tcolorbox}[enhanced jigsaw,left=3pt,breakable,colback=white!97.2549!black,boxrule=0.2mm]

\begin{Shaded}
\begin{Highlighting}[]
\DataTypeTok{int}\NormalTok{ v}\OperatorTok{[}\DecValTok{10}\OperatorTok{];}
\end{Highlighting}
\end{Shaded}

\end{tcolorbox}

\index{Variable-length array!defining|(}

Tapi dengan VLA, kita bisa memakai ukuran yang ditentukan saat runtime untuk membuat array, seperti ini:

\begin{tcolorbox}[enhanced jigsaw,left=3pt,breakable,colback=white!97.2549!black,boxrule=0.2mm]

\begin{Shaded}
\begin{Highlighting}[]
\DataTypeTok{int}\NormalTok{ n }\OperatorTok{=} \DecValTok{10}\OperatorTok{;}
\DataTypeTok{int}\NormalTok{ v}\OperatorTok{[}\NormalTok{n}\OperatorTok{];}
\end{Highlighting}
\end{Shaded}

\end{tcolorbox}

Kelihatannya sama, dan dalam banyak hal memang sama, tapi cara ini memberi keleluasaan untuk menghitung ukuran yang dibutuhkan, lalu mendapatkan array dengan ukuran persis itu.

Mari kita minta pengguna memasukkan ukuran array, lalu menyimpan indeks-kali-10 di setiap elemen array itu:

\begin{tcolorbox}[enhanced jigsaw,left=3pt,breakable,colback=white!97.2549!black,boxrule=0.2mm]

\begin{Shaded}
\begin{Highlighting}[numbers=left,,]
\PreprocessorTok{\#include }\ImportTok{\textless{}stdio.h\textgreater{}}

\DataTypeTok{int}\NormalTok{ main}\OperatorTok{(}\DataTypeTok{void}\OperatorTok{)}
\OperatorTok{\{}
    \DataTypeTok{int}\NormalTok{ n}\OperatorTok{;}
    \DataTypeTok{char}\NormalTok{ buf}\OperatorTok{[}\DecValTok{32}\OperatorTok{];}

\NormalTok{    printf}\OperatorTok{(}\StringTok{"Enter a number: "}\OperatorTok{);}\NormalTok{ fflush}\OperatorTok{(}\NormalTok{stdout}\OperatorTok{);}
\NormalTok{    fgets}\OperatorTok{(}\NormalTok{buf}\OperatorTok{,} \KeywordTok{sizeof}\NormalTok{ buf}\OperatorTok{,}\NormalTok{ stdin}\OperatorTok{);}
\NormalTok{    n }\OperatorTok{=}\NormalTok{ strtoul}\OperatorTok{(}\NormalTok{buf}\OperatorTok{,}\NormalTok{ NULL}\OperatorTok{,} \DecValTok{10}\OperatorTok{);}

    \DataTypeTok{int}\NormalTok{ v}\OperatorTok{[}\NormalTok{n}\OperatorTok{];}

    \ControlFlowTok{for} \OperatorTok{(}\DataTypeTok{int}\NormalTok{ i }\OperatorTok{=} \DecValTok{0}\OperatorTok{;}\NormalTok{ i }\OperatorTok{\textless{}}\NormalTok{ n}\OperatorTok{;}\NormalTok{ i}\OperatorTok{++)}
\NormalTok{        v}\OperatorTok{[}\NormalTok{i}\OperatorTok{]} \OperatorTok{=}\NormalTok{ i }\OperatorTok{*} \DecValTok{10}\OperatorTok{;}

    \ControlFlowTok{for} \OperatorTok{(}\DataTypeTok{int}\NormalTok{ i }\OperatorTok{=} \DecValTok{0}\OperatorTok{;}\NormalTok{ i }\OperatorTok{\textless{}}\NormalTok{ n}\OperatorTok{;}\NormalTok{ i}\OperatorTok{++)}
\NormalTok{        printf}\OperatorTok{(}\StringTok{"v[}\SpecialCharTok{\%d}\StringTok{] = }\SpecialCharTok{\%d\textbackslash{}n}\StringTok{"}\OperatorTok{,}\NormalTok{ i}\OperatorTok{,}\NormalTok{ v}\OperatorTok{[}\NormalTok{i}\OperatorTok{]);}
\OperatorTok{\}}
\end{Highlighting}
\end{Shaded}

\end{tcolorbox}

(Di baris 8 aku menaruh \hl{\mbox{\texttt{fflush()}}} yang seharusnya memaksa baris itu keluar meskipun aku tidak menaruh newline di akhirnya.)

Baris 12 adalah tempat VLA dideklarasikan---begitu eksekusi melewati baris itu, ukuran array ditetapkan sesuai nilai \hl{\mbox{\texttt{n}}} pada saat itu. Panjang array tidak bisa diubah belakangan.

Kamu juga bisa menaruh ekspresi di dalam kurung siku:

\begin{tcolorbox}[enhanced jigsaw,left=3pt,breakable,colback=white!97.2549!black,boxrule=0.2mm]

\begin{Shaded}
\begin{Highlighting}[]
\DataTypeTok{int}\NormalTok{ v}\OperatorTok{[}\NormalTok{x }\OperatorTok{*} \DecValTok{100}\OperatorTok{];}
\end{Highlighting}
\end{Shaded}

\end{tcolorbox}

Beberapa batasan:

\begin{itemize}
\tightlist
\item
  VLA tidak bisa dideklarasikan di file scope, dan tidak bisa dibuat \hl{\mbox{\texttt{static}}} di block scope\footnote{Ini karena VLA biasanya dialokasikan di stack, sedangkan variabel \hl{\mbox{\texttt{static}}} ada di heap. Dan seluruh gagasan VLA adalah mereka otomatis didealokasikan saat stack frame di-pop di akhir fungsi.}.
\item
  Kamu tidak bisa memakai initializer list untuk menginisialisasi array.
\end{itemize}

Selain itu, memasukkan nilai negatif untuk ukuran array memicu undefined behavior---setidaknya di alam semesta ini.

\index{Variable-length array!defining|)}

\section{\texorpdfstring{\hl{\mbox{\texttt{sizeof}}} dan VLA}{sizeof dan VLA}}\label{sizeof-dan-vla}

\index{Variable-length array!and sizeof()@and \texttt{sizeof()}|(}

Kita terbiasa dengan \hl{\mbox{\texttt{sizeof}}} yang memberi ukuran dalam byte untuk object tertentu, termasuk array. VLA pun tidak terkecuali.

Perbedaan utamanya: \hl{\mbox{\texttt{sizeof}}} pada VLA dieksekusi saat \emph{runtime}, sedangkan pada variabel berukuran tetap ia dihitung saat \emph{compile time}.

Tapi pemakaiannya sama.

Kamu bahkan bisa menghitung jumlah elemen dalam VLA dengan trik array biasa:

\begin{tcolorbox}[enhanced jigsaw,left=3pt,breakable,colback=white!97.2549!black,boxrule=0.2mm]

\begin{Shaded}
\begin{Highlighting}[]
\DataTypeTok{size\_t}\NormalTok{ num\_elems }\OperatorTok{=} \KeywordTok{sizeof}\NormalTok{ v }\OperatorTok{/} \KeywordTok{sizeof}\NormalTok{ v}\OperatorTok{[}\DecValTok{0}\OperatorTok{];}
\end{Highlighting}
\end{Shaded}

\end{tcolorbox}

Baris di atas menyiratkan hal yang halus dan benar: aritmetika pointer bekerja persis seperti yang kamu harapkan pada array biasa. Jadi pakai saja sepuasmu:

\begin{tcolorbox}[enhanced jigsaw,left=3pt,breakable,colback=white!97.2549!black,boxrule=0.2mm]

\begin{Shaded}
\begin{Highlighting}[numbers=left,,]
\PreprocessorTok{\#include }\ImportTok{\textless{}stdio.h\textgreater{}}

\DataTypeTok{int}\NormalTok{ main}\OperatorTok{(}\DataTypeTok{void}\OperatorTok{)}
\OperatorTok{\{}
    \DataTypeTok{int}\NormalTok{ n }\OperatorTok{=} \DecValTok{5}\OperatorTok{;}
    \DataTypeTok{int}\NormalTok{ v}\OperatorTok{[}\NormalTok{n}\OperatorTok{];}

    \DataTypeTok{int} \OperatorTok{*}\NormalTok{p }\OperatorTok{=}\NormalTok{ v}\OperatorTok{;}

    \OperatorTok{*(}\NormalTok{p}\OperatorTok{+}\DecValTok{2}\OperatorTok{)} \OperatorTok{=} \DecValTok{12}\OperatorTok{;}
\NormalTok{    printf}\OperatorTok{(}\StringTok{"}\SpecialCharTok{\%d\textbackslash{}n}\StringTok{"}\OperatorTok{,}\NormalTok{ v}\OperatorTok{[}\DecValTok{2}\OperatorTok{]);}  \CommentTok{// 12}

\NormalTok{    p}\OperatorTok{[}\DecValTok{3}\OperatorTok{]} \OperatorTok{=} \DecValTok{34}\OperatorTok{;}
\NormalTok{    printf}\OperatorTok{(}\StringTok{"}\SpecialCharTok{\%d\textbackslash{}n}\StringTok{"}\OperatorTok{,}\NormalTok{ v}\OperatorTok{[}\DecValTok{3}\OperatorTok{]);}  \CommentTok{// 34}
\OperatorTok{\}}
\end{Highlighting}
\end{Shaded}

\end{tcolorbox}

Seperti pada array biasa, kamu bisa memakai tanda kurung dengan \hl{\mbox{\texttt{sizeof()}}} untuk mengetahui ukuran calon VLA tanpa benar-benar mendeklarasikannya:

\begin{tcolorbox}[enhanced jigsaw,left=3pt,breakable,colback=white!97.2549!black,boxrule=0.2mm]

\begin{Shaded}
\begin{Highlighting}[]
\DataTypeTok{int}\NormalTok{ x }\OperatorTok{=} \DecValTok{12}\OperatorTok{;}

\NormalTok{printf}\OperatorTok{(}\StringTok{"}\SpecialCharTok{\%zu\textbackslash{}n}\StringTok{"}\OperatorTok{,} \KeywordTok{sizeof}\OperatorTok{(}\DataTypeTok{int} \OperatorTok{[}\NormalTok{x}\OperatorTok{]));}  \CommentTok{// Mencetak 48 di sistemku}
\end{Highlighting}
\end{Shaded}

\end{tcolorbox}

\index{Variable-length array!and sizeof()@and \texttt{sizeof()}|)}

\section{VLA Multidimensi}\label{vla-multidimensi}

\index{Variable-length array!multidimensional|(}

Kamu boleh membuat berbagai macam VLA dengan satu dimensi atau lebih berupa variabel

\begin{tcolorbox}[enhanced jigsaw,left=3pt,breakable,colback=white!97.2549!black,boxrule=0.2mm]

\begin{Shaded}
\begin{Highlighting}[]
\DataTypeTok{int}\NormalTok{ w }\OperatorTok{=} \DecValTok{10}\OperatorTok{;}
\DataTypeTok{int}\NormalTok{ h }\OperatorTok{=} \DecValTok{20}\OperatorTok{;}

\DataTypeTok{int}\NormalTok{ x}\OperatorTok{[}\NormalTok{h}\OperatorTok{][}\NormalTok{w}\OperatorTok{];}
\DataTypeTok{int}\NormalTok{ y}\OperatorTok{[}\DecValTok{5}\OperatorTok{][}\NormalTok{w}\OperatorTok{];}
\DataTypeTok{int}\NormalTok{ z}\OperatorTok{[}\DecValTok{10}\OperatorTok{][}\NormalTok{w}\OperatorTok{][}\DecValTok{20}\OperatorTok{];}
\end{Highlighting}
\end{Shaded}

\end{tcolorbox}

Sekali lagi, kamu bisa menjelajahinya persis seperti array biasa.

\index{Variable-length array!multidimensional|)}

\section{Mengoper VLA Satu Dimensi ke Fungsi}\label{mengoper-vla-satu-dimensi-ke-fungsi}

\index{Variable-length array!passing to functions|(}

Mengoper VLA satu dimensi ke fungsi tidak berbeda dari mengoper array biasa. Langsung saja.

\begin{tcolorbox}[enhanced jigsaw,left=3pt,breakable,colback=white!97.2549!black,boxrule=0.2mm]

\begin{Shaded}
\begin{Highlighting}[numbers=left,,]
\PreprocessorTok{\#include }\ImportTok{\textless{}stdio.h\textgreater{}}

\DataTypeTok{int}\NormalTok{ sum}\OperatorTok{(}\DataTypeTok{int}\NormalTok{ count}\OperatorTok{,} \DataTypeTok{int} \OperatorTok{*}\NormalTok{v}\OperatorTok{)}
\OperatorTok{\{}
    \DataTypeTok{int}\NormalTok{ total }\OperatorTok{=} \DecValTok{0}\OperatorTok{;}

    \ControlFlowTok{for} \OperatorTok{(}\DataTypeTok{int}\NormalTok{ i }\OperatorTok{=} \DecValTok{0}\OperatorTok{;}\NormalTok{ i }\OperatorTok{\textless{}}\NormalTok{ count}\OperatorTok{;}\NormalTok{ i}\OperatorTok{++)}
\NormalTok{        total }\OperatorTok{+=}\NormalTok{ v}\OperatorTok{[}\NormalTok{i}\OperatorTok{];}

    \ControlFlowTok{return}\NormalTok{ total}\OperatorTok{;}
\OperatorTok{\}}

\DataTypeTok{int}\NormalTok{ main}\OperatorTok{(}\DataTypeTok{void}\OperatorTok{)}
\OperatorTok{\{}
    \DataTypeTok{int}\NormalTok{ x}\OperatorTok{[}\DecValTok{5}\OperatorTok{];}   \CommentTok{// Array standar}

    \DataTypeTok{int}\NormalTok{ a }\OperatorTok{=} \DecValTok{5}\OperatorTok{;}
    \DataTypeTok{int}\NormalTok{ y}\OperatorTok{[}\NormalTok{a}\OperatorTok{];}   \CommentTok{// VLA}

    \ControlFlowTok{for} \OperatorTok{(}\DataTypeTok{int}\NormalTok{ i }\OperatorTok{=} \DecValTok{0}\OperatorTok{;}\NormalTok{ i }\OperatorTok{\textless{}}\NormalTok{ a}\OperatorTok{;}\NormalTok{ i}\OperatorTok{++)}
\NormalTok{        x}\OperatorTok{[}\NormalTok{i}\OperatorTok{]} \OperatorTok{=}\NormalTok{ y}\OperatorTok{[}\NormalTok{i}\OperatorTok{]} \OperatorTok{=}\NormalTok{ i }\OperatorTok{+} \DecValTok{1}\OperatorTok{;}

\NormalTok{    printf}\OperatorTok{(}\StringTok{"}\SpecialCharTok{\%d\textbackslash{}n}\StringTok{"}\OperatorTok{,}\NormalTok{ sum}\OperatorTok{(}\DecValTok{5}\OperatorTok{,}\NormalTok{ x}\OperatorTok{));}
\NormalTok{    printf}\OperatorTok{(}\StringTok{"}\SpecialCharTok{\%d\textbackslash{}n}\StringTok{"}\OperatorTok{,}\NormalTok{ sum}\OperatorTok{(}\NormalTok{a}\OperatorTok{,}\NormalTok{ y}\OperatorTok{));}
\OperatorTok{\}}
\end{Highlighting}
\end{Shaded}

\end{tcolorbox}

Tapi ada sedikit lebih banyak daripada itu. Kamu juga bisa memberi tahu C bahwa array itu berukuran VLA tertentu dengan mengopernya lebih dulu lalu menyebut dimensi itu di daftar parameter:

\begin{tcolorbox}[enhanced jigsaw,left=3pt,breakable,colback=white!97.2549!black,boxrule=0.2mm]

\begin{Shaded}
\begin{Highlighting}[]
\DataTypeTok{int}\NormalTok{ sum}\OperatorTok{(}\DataTypeTok{int}\NormalTok{ count}\OperatorTok{,} \DataTypeTok{int}\NormalTok{ v}\OperatorTok{[}\NormalTok{count}\OperatorTok{])}
\OperatorTok{\{}
    \CommentTok{// ...}
\OperatorTok{\}}
\end{Highlighting}
\end{Shaded}

\end{tcolorbox}

\index{Variable-length array!in function prototypes|(} \index{* for VLA function prototypes@\texttt{*} for VLA function prototypes|(}

Omong-omong, ada beberapa cara menulis prototype untuk fungsi di atas; salah satunya memakai \hl{\mbox{\texttt{*}}} kalau kamu tidak ingin menyebut nama nilai di VLA secara spesifik. Ia hanya menandakan bahwa tipenya VLA, bukan pointer biasa.

Prototype VLA:

\begin{tcolorbox}[enhanced jigsaw,left=3pt,breakable,colback=white!97.2549!black,boxrule=0.2mm]

\begin{Shaded}
\begin{Highlighting}[]
\DataTypeTok{void}\NormalTok{ do\_something}\OperatorTok{(}\DataTypeTok{int}\NormalTok{ count}\OperatorTok{,} \DataTypeTok{int}\NormalTok{ v}\OperatorTok{[}\NormalTok{count}\OperatorTok{]);}  \CommentTok{// Dengan nama}
\DataTypeTok{void}\NormalTok{ do\_something}\OperatorTok{(}\DataTypeTok{int}\OperatorTok{,} \DataTypeTok{int}\NormalTok{ v}\OperatorTok{[*]);}            \CommentTok{// Tanpa nama}
\end{Highlighting}
\end{Shaded}

\end{tcolorbox}

Sekali lagi, \hl{\mbox{\texttt{*}}} itu hanya berfungsi di prototype---di fungsinya sendiri, kamu harus menulis ukuran secara eksplisit.

\index{* for VLA function prototypes@\texttt{*} for VLA function prototypes|)} \index{Variable-length array!in function prototypes|)}

Nah---\emph{saatnya multidimensi}! Di sinilah keseruan dimulai.

\section{Mengoper VLA Multidimensi ke Fungsi}\label{mengoper-vla-multidimensi-ke-fungsi}

Sama seperti yang kita lakukan pada bentuk kedua VLA satu dimensi di atas, tapi kali ini kita mengoper dua dimensi dan memakainya.

Pada contoh berikut, kita membuat matriks tabel perkalian dengan lebar dan tinggi yang bervariasi, lalu mengopernya ke fungsi untuk dicetak.

\begin{tcolorbox}[enhanced jigsaw,left=3pt,breakable,colback=white!97.2549!black,boxrule=0.2mm]

\begin{Shaded}
\begin{Highlighting}[numbers=left,,]
\PreprocessorTok{\#include }\ImportTok{\textless{}stdio.h\textgreater{}}

\DataTypeTok{void}\NormalTok{ print\_matrix}\OperatorTok{(}\DataTypeTok{int}\NormalTok{ h}\OperatorTok{,} \DataTypeTok{int}\NormalTok{ w}\OperatorTok{,} \DataTypeTok{int}\NormalTok{ m}\OperatorTok{[}\NormalTok{h}\OperatorTok{][}\NormalTok{w}\OperatorTok{])}
\OperatorTok{\{}
    \ControlFlowTok{for} \OperatorTok{(}\DataTypeTok{int}\NormalTok{ row }\OperatorTok{=} \DecValTok{0}\OperatorTok{;}\NormalTok{ row }\OperatorTok{\textless{}}\NormalTok{ h}\OperatorTok{;}\NormalTok{ row}\OperatorTok{++)} \OperatorTok{\{}
        \ControlFlowTok{for} \OperatorTok{(}\DataTypeTok{int}\NormalTok{ col }\OperatorTok{=} \DecValTok{0}\OperatorTok{;}\NormalTok{ col }\OperatorTok{\textless{}}\NormalTok{ w}\OperatorTok{;}\NormalTok{ col}\OperatorTok{++)}
\NormalTok{            printf}\OperatorTok{(}\StringTok{"}\SpecialCharTok{\%2d}\StringTok{ "}\OperatorTok{,}\NormalTok{ m}\OperatorTok{[}\NormalTok{row}\OperatorTok{][}\NormalTok{col}\OperatorTok{]);}
\NormalTok{        printf}\OperatorTok{(}\StringTok{"}\SpecialCharTok{\textbackslash{}n}\StringTok{"}\OperatorTok{);}
    \OperatorTok{\}}
\OperatorTok{\}}

\DataTypeTok{int}\NormalTok{ main}\OperatorTok{(}\DataTypeTok{void}\OperatorTok{)}
\OperatorTok{\{}
    \DataTypeTok{int}\NormalTok{ rows }\OperatorTok{=} \DecValTok{4}\OperatorTok{;}
    \DataTypeTok{int}\NormalTok{ cols }\OperatorTok{=} \DecValTok{7}\OperatorTok{;}

    \DataTypeTok{int}\NormalTok{ matrix}\OperatorTok{[}\NormalTok{rows}\OperatorTok{][}\NormalTok{cols}\OperatorTok{];}

    \ControlFlowTok{for} \OperatorTok{(}\DataTypeTok{int}\NormalTok{ row }\OperatorTok{=} \DecValTok{0}\OperatorTok{;}\NormalTok{ row }\OperatorTok{\textless{}}\NormalTok{ rows}\OperatorTok{;}\NormalTok{ row}\OperatorTok{++)}
        \ControlFlowTok{for} \OperatorTok{(}\DataTypeTok{int}\NormalTok{ col }\OperatorTok{=} \DecValTok{0}\OperatorTok{;}\NormalTok{ col }\OperatorTok{\textless{}}\NormalTok{ cols}\OperatorTok{;}\NormalTok{ col}\OperatorTok{++)}
\NormalTok{            matrix}\OperatorTok{[}\NormalTok{row}\OperatorTok{][}\NormalTok{col}\OperatorTok{]} \OperatorTok{=}\NormalTok{ row }\OperatorTok{*}\NormalTok{ col}\OperatorTok{;}

\NormalTok{    print\_matrix}\OperatorTok{(}\NormalTok{rows}\OperatorTok{,}\NormalTok{ cols}\OperatorTok{,}\NormalTok{ matrix}\OperatorTok{);}
\OperatorTok{\}}
\end{Highlighting}
\end{Shaded}

\end{tcolorbox}

\subsection{VLA Multidimensi Parsial}\label{vla-multidimensi-parsial}

Sebagian dimensi bisa tetap dan sebagian variabel. Misalkan kita punya panjang record yang tetap 5 elemen, tapi kita tidak tahu ada berapa record.

\begin{tcolorbox}[enhanced jigsaw,left=3pt,breakable,colback=white!97.2549!black,boxrule=0.2mm]

\begin{Shaded}
\begin{Highlighting}[numbers=left,,]
\PreprocessorTok{\#include }\ImportTok{\textless{}stdio.h\textgreater{}}

\DataTypeTok{void}\NormalTok{ print\_records}\OperatorTok{(}\DataTypeTok{int}\NormalTok{ count}\OperatorTok{,} \DataTypeTok{int}\NormalTok{ record}\OperatorTok{[}\NormalTok{count}\OperatorTok{][}\DecValTok{5}\OperatorTok{])}
\OperatorTok{\{}
    \ControlFlowTok{for} \OperatorTok{(}\DataTypeTok{int}\NormalTok{ i }\OperatorTok{=} \DecValTok{0}\OperatorTok{;}\NormalTok{ i }\OperatorTok{\textless{}}\NormalTok{ count}\OperatorTok{;}\NormalTok{ i}\OperatorTok{++)} \OperatorTok{\{}
        \ControlFlowTok{for} \OperatorTok{(}\DataTypeTok{int}\NormalTok{ j }\OperatorTok{=} \DecValTok{0}\OperatorTok{;}\NormalTok{ j }\OperatorTok{\textless{}} \DecValTok{5}\OperatorTok{;}\NormalTok{ j}\OperatorTok{++)}
\NormalTok{            printf}\OperatorTok{(}\StringTok{"}\SpecialCharTok{\%2d}\StringTok{ "}\OperatorTok{,}\NormalTok{ record}\OperatorTok{[}\NormalTok{i}\OperatorTok{][}\NormalTok{j}\OperatorTok{]);}
\NormalTok{        printf}\OperatorTok{(}\StringTok{"}\SpecialCharTok{\textbackslash{}n}\StringTok{"}\OperatorTok{);}
    \OperatorTok{\}}
\OperatorTok{\}}

\DataTypeTok{int}\NormalTok{ main}\OperatorTok{(}\DataTypeTok{void}\OperatorTok{)}
\OperatorTok{\{}
    \DataTypeTok{int}\NormalTok{ rec\_count }\OperatorTok{=} \DecValTok{3}\OperatorTok{;}
    \DataTypeTok{int}\NormalTok{ records}\OperatorTok{[}\NormalTok{rec\_count}\OperatorTok{][}\DecValTok{5}\OperatorTok{];}

    \CommentTok{// Isi dengan beberapa data tiruan}
    \ControlFlowTok{for} \OperatorTok{(}\DataTypeTok{int}\NormalTok{ i }\OperatorTok{=} \DecValTok{0}\OperatorTok{;}\NormalTok{ i }\OperatorTok{\textless{}}\NormalTok{ rec\_count}\OperatorTok{;}\NormalTok{ i}\OperatorTok{++)}
        \ControlFlowTok{for} \OperatorTok{(}\DataTypeTok{int}\NormalTok{ j }\OperatorTok{=} \DecValTok{0}\OperatorTok{;}\NormalTok{ j }\OperatorTok{\textless{}} \DecValTok{5}\OperatorTok{;}\NormalTok{ j}\OperatorTok{++)}
\NormalTok{            records}\OperatorTok{[}\NormalTok{i}\OperatorTok{][}\NormalTok{j}\OperatorTok{]} \OperatorTok{=} \OperatorTok{(}\NormalTok{i}\OperatorTok{+}\DecValTok{1}\OperatorTok{)*(}\NormalTok{j}\OperatorTok{+}\DecValTok{2}\OperatorTok{);}

\NormalTok{    print\_records}\OperatorTok{(}\NormalTok{rec\_count}\OperatorTok{,}\NormalTok{ records}\OperatorTok{);}
\OperatorTok{\}}
\end{Highlighting}
\end{Shaded}

\end{tcolorbox}

\index{Variable-length array!passing to functions|)}

\section{Kompatibilitas dengan Array Biasa}\label{kompatibilitas-dengan-array-biasa}

\index{Variable-length array!with regular arrays|(}

Karena VLA di memori sama saja dengan array biasa, boleh sekali mengoper keduanya secara bergantian\ldots{} selama dimensinya cocok.

Misalnya, kalau ada fungsi yang secara spesifik menginginkan array \(3\times5\), kita tetap bisa mengoper VLA ke dalamnya.

\begin{tcolorbox}[enhanced jigsaw,left=3pt,breakable,colback=white!97.2549!black,boxrule=0.2mm]

\begin{Shaded}
\begin{Highlighting}[]
\DataTypeTok{int}\NormalTok{ foo}\OperatorTok{(}\DataTypeTok{int}\NormalTok{ m}\OperatorTok{[}\DecValTok{5}\OperatorTok{][}\DecValTok{3}\OperatorTok{])} \OperatorTok{\{...\}}

\NormalTok{\textbackslash{}\textbackslash{} }\OperatorTok{...}

\DataTypeTok{int}\NormalTok{ w }\OperatorTok{=} \DecValTok{3}\OperatorTok{,}\NormalTok{ h }\OperatorTok{=} \DecValTok{5}\OperatorTok{;}
\DataTypeTok{int}\NormalTok{ matrix}\OperatorTok{[}\NormalTok{h}\OperatorTok{][}\NormalTok{w}\OperatorTok{];}

\NormalTok{foo}\OperatorTok{(}\NormalTok{matrix}\OperatorTok{);}   \CommentTok{// OK!}
\end{Highlighting}
\end{Shaded}

\end{tcolorbox}

Begitu juga, kalau kamu punya fungsi VLA, kamu bisa mengoper array biasa ke dalamnya:

\begin{tcolorbox}[enhanced jigsaw,left=3pt,breakable,colback=white!97.2549!black,boxrule=0.2mm]

\begin{Shaded}
\begin{Highlighting}[]
\DataTypeTok{int}\NormalTok{ foo}\OperatorTok{(}\DataTypeTok{int}\NormalTok{ h}\OperatorTok{,} \DataTypeTok{int}\NormalTok{ w}\OperatorTok{,} \DataTypeTok{int}\NormalTok{ m}\OperatorTok{[}\NormalTok{h}\OperatorTok{][}\NormalTok{w}\OperatorTok{])} \OperatorTok{\{...\}}

\NormalTok{\textbackslash{}\textbackslash{} }\OperatorTok{...}

\DataTypeTok{int}\NormalTok{ matrix}\OperatorTok{[}\DecValTok{3}\OperatorTok{][}\DecValTok{5}\OperatorTok{];}

\NormalTok{foo}\OperatorTok{(}\DecValTok{3}\OperatorTok{,} \DecValTok{5}\OperatorTok{,}\NormalTok{ matrix}\OperatorTok{);}   \CommentTok{// OK!}
\end{Highlighting}
\end{Shaded}

\end{tcolorbox}

Tapi waspada: kalau dimensinya tidak cocok, kemungkinan besar kamu akan mengalami undefined behavior.

\index{Variable-length array!with regular arrays|)}

\section{\texorpdfstring{\hl{\mbox{\texttt{typedef}}} dan VLA}{typedef dan VLA}}\label{typedef-dan-vla}

\index{Variable-length array!with typedef@with \texttt{typedef}|(}

Kamu bisa me-\hl{\mbox{\texttt{typedef}}} VLA, tapi perilakunya mungkin tidak seperti yang kamu harapkan.

Pada dasarnya, \hl{\mbox{\texttt{typedef}}} membuat tipe baru dengan nilai-nilai seperti yang ada pada saat \hl{\mbox{\texttt{typedef}}} dieksekusi.

Jadi ini bukan \hl{\mbox{\texttt{typedef}}} dari VLA, melainkan tipe array berukuran tetap baru dengan dimensi pada saat itu.

\begin{tcolorbox}[enhanced jigsaw,left=3pt,breakable,colback=white!97.2549!black,boxrule=0.2mm]

\begin{Shaded}
\begin{Highlighting}[numbers=left,,]
\PreprocessorTok{\#include }\ImportTok{\textless{}stdio.h\textgreater{}}

\DataTypeTok{int}\NormalTok{ main}\OperatorTok{(}\DataTypeTok{void}\OperatorTok{)}
\OperatorTok{\{}
    \DataTypeTok{int}\NormalTok{ w }\OperatorTok{=} \DecValTok{10}\OperatorTok{;}

    \KeywordTok{typedef} \DataTypeTok{int}\NormalTok{ goat}\OperatorTok{[}\NormalTok{w}\OperatorTok{];}

    \CommentTok{// goat adalah array berisi 10 int}
\NormalTok{    goat x}\OperatorTok{;}

    \CommentTok{// Inisialisasi dengan kuadrat angka}
    \ControlFlowTok{for} \OperatorTok{(}\DataTypeTok{int}\NormalTok{ i }\OperatorTok{=} \DecValTok{0}\OperatorTok{;}\NormalTok{ i }\OperatorTok{\textless{}}\NormalTok{ w}\OperatorTok{;}\NormalTok{ i}\OperatorTok{++)}
\NormalTok{        x}\OperatorTok{[}\NormalTok{i}\OperatorTok{]} \OperatorTok{=}\NormalTok{ i}\OperatorTok{*}\NormalTok{i}\OperatorTok{;}

    \CommentTok{// Cetak semuanya}
    \ControlFlowTok{for} \OperatorTok{(}\DataTypeTok{int}\NormalTok{ i }\OperatorTok{=} \DecValTok{0}\OperatorTok{;}\NormalTok{ i }\OperatorTok{\textless{}}\NormalTok{ w}\OperatorTok{;}\NormalTok{ i}\OperatorTok{++)}
\NormalTok{        printf}\OperatorTok{(}\StringTok{"}\SpecialCharTok{\%d\textbackslash{}n}\StringTok{"}\OperatorTok{,}\NormalTok{ x}\OperatorTok{[}\NormalTok{i}\OperatorTok{]);}

    \CommentTok{// Sekarang kita ubah w...}

\NormalTok{    w }\OperatorTok{=} \DecValTok{20}\OperatorTok{;}

    \CommentTok{// Tapi goat TETAP array berisi 10 int, karena itulah nilai w}
    \CommentTok{// saat typedef dieksekusi.}
\OperatorTok{\}}
\end{Highlighting}
\end{Shaded}

\end{tcolorbox}

Jadi ia bertingkah seperti array berukuran tetap.

Tapi kamu tetap tidak bisa memakai initializer list padanya.

\index{Variable-length array!with typedef@with \texttt{typedef}|)}

\section{Jebakan Lompatan}\label{jebakan-lompatan}

\index{Variable-length array!with goto@with \texttt{goto}|(}

Kamu harus berhati-hati saat memakai \hl{\mbox{\texttt{goto}}} di dekat VLA karena banyak hal yang tidak legal.

\index{Variable-length array!with goto@with \texttt{goto}|)} \index{Variable-length array!with longjmp()@with \texttt{longjmp()}|(}

Dan saat memakai \hl{\mbox{\texttt{longjmp()}}} ada kasus ketika kamu bisa membocorkan memori dengan VLA.

\index{Variable-length array!with longjmp()@with \texttt{longjmp()}|)}

Tapi kedua hal ini akan kita bahas di bab masing-masing.

\section{Masalah Umum}\label{vla-general-issues}

\index{Variable-length array!controversy|(}

VLA dilarang di kernel Linux karena beberapa alasan:

\begin{itemize}
\tightlist
\item
  Banyak tempat yang memakainya seharusnya cukup dengan ukuran tetap.
\item
  Kode di balik VLA lebih lambat (sampai tingkat yang tidak disadari kebanyakan orang, tapi berpengaruh pada sistem operasi).
\item
  VLA tidak didukung sama baiknya oleh semua C compiler.
\item
  Ukuran stack terbatas, dan VLA masuk ke stack. Kalau ada kode yang tidak sengaja (atau sengaja) mengoper nilai besar ke fungsi kernel yang mengalokasikan VLA, \emph{Hal-Hal Buruk}™ bisa terjadi.
\end{itemize}

Orang lain di internet menunjukkan bahwa tidak ada cara mendeteksi kegagalan alokasi VLA, dan program yang mengalami masalah seperti itu kemungkinan besar langsung crash. Array berukuran tetap juga punya masalah yang sama, tapi jauh lebih mungkin seseorang tanpa sengaja membuat \emph{VLA Berukuran Tidak Lazim} daripada entah bagaimana tanpa sengaja mendeklarasikan array berukuran tetap, katakanlah, 30 megabyte.

\index{Variable-length array!controversy|)} \index{Variable-length array|)}

\chapter{\texorpdfstring{\hl{\mbox{\texttt{goto}}}}{goto}}\label{goto}

\index{goto statement@\texttt{goto} statement|(}

Pernyataan \hl{\mbox{\texttt{goto}}} dipuja secara universal dan bisa disajikan di sini tanpa perdebatan.

Bercanda! Bertahun-tahun lamanya, orang bolak-balik berdebat apakah \hl{\mbox{\texttt{goto}}} \href{https://en.wikipedia.org/wiki/Goto\#Criticism}{dianggap berbahaya}\footnote{https://en.wikipedia.org/wiki/Goto\#Criticism} atau tidak (sering kali dianggap tidak boleh).

Menurut pendapat programmer yang satu ini, kamu sebaiknya memakai konstruksi apa pun yang menghasilkan kode \emph{terbaik}, dengan mempertimbangkan kemudahan perawatan dan kecepatan. Dan kadang itu mungkin \hl{\mbox{\texttt{goto}}}!

Di bab ini, kita akan melihat cara kerja \hl{\mbox{\texttt{goto}}} di C, lalu mengintip beberapa kasus umum tempat ia dipakai\footnote{Aku ingin menunjukkan bahwa memakai \hl{\mbox{\texttt{goto}}} pada semua kasus ini bisa dihindari. Kamu bisa memakai variabel dan loop sebagai gantinya. Hanya saja sebagian orang menganggap \hl{\mbox{\texttt{goto}}} menghasilkan kode \emph{terbaik} pada keadaan-keadaan itu.}.

\section{Contoh Sederhana}\label{contoh-sederhana}

\index{Labels|(}

Pada contoh ini, kita memakai \hl{\mbox{\texttt{goto}}} untuk melewati satu baris kode dan melompat ke sebuah \emph{label}. Label adalah identifier yang bisa menjadi target \hl{\mbox{\texttt{goto}}}---ia diakhiri titik dua (\hl{\mbox{\texttt{:}}}).

\begin{tcolorbox}[enhanced jigsaw,left=3pt,breakable,colback=white!97.2549!black,boxrule=0.2mm]

\begin{Shaded}
\begin{Highlighting}[numbers=left,,]
\PreprocessorTok{\#include }\ImportTok{\textless{}stdio.h\textgreater{}}

\DataTypeTok{int}\NormalTok{ main}\OperatorTok{(}\DataTypeTok{void}\OperatorTok{)}
\OperatorTok{\{}
\NormalTok{    printf}\OperatorTok{(}\StringTok{"One}\SpecialCharTok{\textbackslash{}n}\StringTok{"}\OperatorTok{);}
\NormalTok{    printf}\OperatorTok{(}\StringTok{"Two}\SpecialCharTok{\textbackslash{}n}\StringTok{"}\OperatorTok{);}

    \ControlFlowTok{goto}\NormalTok{ skip\_3}\OperatorTok{;}

\NormalTok{    printf}\OperatorTok{(}\StringTok{"Three}\SpecialCharTok{\textbackslash{}n}\StringTok{"}\OperatorTok{);}

\NormalTok{skip\_3}\OperatorTok{:}

\NormalTok{    printf}\OperatorTok{(}\StringTok{"Five!}\SpecialCharTok{\textbackslash{}n}\StringTok{"}\OperatorTok{);}
\OperatorTok{\}}
\end{Highlighting}
\end{Shaded}

\end{tcolorbox}

Outputnya:

\begin{tcolorbox}[enhanced jigsaw,left=3pt,breakable,colback=white!97.2549!black,boxrule=0.2mm]

\begin{Shaded}
\begin{Highlighting}[]
\NormalTok{One}
\NormalTok{Two}
\NormalTok{Five!}
\end{Highlighting}
\end{Shaded}

\end{tcolorbox}

\hl{\mbox{\texttt{goto}}} membuat eksekusi melompat ke label yang disebutkan, melewati semua yang ada di antaranya.

Kamu bisa melompat maju maupun mundur dengan \hl{\mbox{\texttt{goto}}}.

\begin{tcolorbox}[enhanced jigsaw,left=3pt,breakable,colback=white!97.2549!black,boxrule=0.2mm]

\begin{Shaded}
\begin{Highlighting}[]
\NormalTok{infinite\_loop}\OperatorTok{:}
\NormalTok{    print}\OperatorTok{(}\StringTok{"Hello, world!}\SpecialCharTok{\textbackslash{}n}\StringTok{"}\OperatorTok{);}
    \ControlFlowTok{goto}\NormalTok{ infinite\_loop}\OperatorTok{;}
\end{Highlighting}
\end{Shaded}

\end{tcolorbox}

Label dilewati saat eksekusi. Yang berikut akan mencetak ketiga angka secara berurutan seolah-olah label itu tidak ada:

\begin{tcolorbox}[enhanced jigsaw,left=3pt,breakable,colback=white!97.2549!black,boxrule=0.2mm]

\begin{Shaded}
\begin{Highlighting}[]
\NormalTok{    printf}\OperatorTok{(}\StringTok{"Zero}\SpecialCharTok{\textbackslash{}n}\StringTok{"}\OperatorTok{);}
\NormalTok{label\_1}\OperatorTok{:}
\NormalTok{label\_2}\OperatorTok{:}
\NormalTok{    printf}\OperatorTok{(}\StringTok{"One}\SpecialCharTok{\textbackslash{}n}\StringTok{"}\OperatorTok{);}
\NormalTok{label\_3}\OperatorTok{:}
\NormalTok{    printf}\OperatorTok{(}\StringTok{"Two}\SpecialCharTok{\textbackslash{}n}\StringTok{"}\OperatorTok{);}
\NormalTok{label\_4}\OperatorTok{:}
\NormalTok{    printf}\OperatorTok{(}\StringTok{"Three}\SpecialCharTok{\textbackslash{}n}\StringTok{"}\OperatorTok{);}
\end{Highlighting}
\end{Shaded}

\end{tcolorbox}

Seperti yang kamu perhatikan, konvensi umumnya adalah merapatkan label sepenuhnya ke kiri. Ini meningkatkan keterbacaan karena pembaca bisa dengan cepat memindai untuk menemukan tujuannya.

Label punya \emph{function scope}. Artinya, sedalam apa pun label itu berada di dalam blok, kamu tetap bisa \hl{\mbox{\texttt{goto}}} ke sana dari mana saja di dalam fungsi.

Artinya juga kamu hanya bisa \hl{\mbox{\texttt{goto}}} ke label yang berada di fungsi yang sama dengan \hl{\mbox{\texttt{goto}}}-nya. Label di fungsi lain berada di luar scope dari sudut pandang \hl{\mbox{\texttt{goto}}}. Dan artinya kamu bisa memakai nama label yang sama di dua fungsi---hanya tidak boleh nama label yang sama di fungsi yang sama.

\index{Labels|)}

\section{\texorpdfstring{\hl{\mbox{\texttt{continue}}} Berlabel}{continue Berlabel}}\label{continue-berlabel}

\index{goto statement@\texttt{goto} statement!as labeled continue@as labeled \texttt{continue}|(}

Di beberapa bahasa, kamu sebenarnya bisa menentukan label untuk pernyataan \hl{\mbox{\texttt{continue}}}. C tidak mengizinkannya, tapi kamu dengan mudah bisa memakai \hl{\mbox{\texttt{goto}}} sebagai gantinya.

Untuk menunjukkan masalahnya, lihat \hl{\mbox{\texttt{continue}}} pada loop bersarang ini:

\begin{tcolorbox}[enhanced jigsaw,left=3pt,breakable,colback=white!97.2549!black,boxrule=0.2mm]

\begin{Shaded}
\begin{Highlighting}[]
\ControlFlowTok{for} \OperatorTok{(}\DataTypeTok{int}\NormalTok{ i }\OperatorTok{=} \DecValTok{0}\OperatorTok{;}\NormalTok{ i }\OperatorTok{\textless{}} \DecValTok{3}\OperatorTok{;}\NormalTok{ i}\OperatorTok{++)} \OperatorTok{\{}
    \ControlFlowTok{for} \OperatorTok{(}\DataTypeTok{int}\NormalTok{ j }\OperatorTok{=} \DecValTok{0}\OperatorTok{;}\NormalTok{ j }\OperatorTok{\textless{}} \DecValTok{3}\OperatorTok{;}\NormalTok{ j}\OperatorTok{++)} \OperatorTok{\{}
\NormalTok{        printf}\OperatorTok{(}\StringTok{"}\SpecialCharTok{\%d}\StringTok{, }\SpecialCharTok{\%d\textbackslash{}n}\StringTok{"}\OperatorTok{,}\NormalTok{ i}\OperatorTok{,}\NormalTok{ j}\OperatorTok{);}
        \ControlFlowTok{continue}\OperatorTok{;}   \CommentTok{// Selalu menuju j berikutnya}
    \OperatorTok{\}}
\OperatorTok{\}}
\end{Highlighting}
\end{Shaded}

\end{tcolorbox}

Seperti yang kita lihat, \hl{\mbox{\texttt{continue}}} itu, seperti semua \hl{\mbox{\texttt{continue}}}, menuju iterasi berikutnya dari loop terdekat yang melingkupinya. Bagaimana kalau kita ingin \hl{\mbox{\texttt{continue}}} di loop yang lebih luar, yaitu loop \hl{\mbox{\texttt{i}}}?

Nah, kita bisa \hl{\mbox{\texttt{break}}} untuk kembali ke loop luar, kan?

\begin{tcolorbox}[enhanced jigsaw,left=3pt,breakable,colback=white!97.2549!black,boxrule=0.2mm]

\begin{Shaded}
\begin{Highlighting}[]
\ControlFlowTok{for} \OperatorTok{(}\DataTypeTok{int}\NormalTok{ i }\OperatorTok{=} \DecValTok{0}\OperatorTok{;}\NormalTok{ i }\OperatorTok{\textless{}} \DecValTok{3}\OperatorTok{;}\NormalTok{ i}\OperatorTok{++)} \OperatorTok{\{}
    \ControlFlowTok{for} \OperatorTok{(}\DataTypeTok{int}\NormalTok{ j }\OperatorTok{=} \DecValTok{0}\OperatorTok{;}\NormalTok{ j }\OperatorTok{\textless{}} \DecValTok{3}\OperatorTok{;}\NormalTok{ j}\OperatorTok{++)} \OperatorTok{\{}
\NormalTok{        printf}\OperatorTok{(}\StringTok{"}\SpecialCharTok{\%d}\StringTok{, }\SpecialCharTok{\%d\textbackslash{}n}\StringTok{"}\OperatorTok{,}\NormalTok{ i}\OperatorTok{,}\NormalTok{ j}\OperatorTok{);}
        \ControlFlowTok{break}\OperatorTok{;}     \CommentTok{// Membawa kita ke iterasi i berikutnya}
    \OperatorTok{\}}
\OperatorTok{\}}
\end{Highlighting}
\end{Shaded}

\end{tcolorbox}

Itu membawa kita melewati dua tingkat loop bersarang. Tapi kalau kita menyarangkan loop lagi, pilihan kita habis. Bagaimana dengan yang ini, ketika tidak ada pernyataan yang bisa membawa kita keluar ke iterasi \hl{\mbox{\texttt{i}}} berikutnya?

\begin{tcolorbox}[enhanced jigsaw,left=3pt,breakable,colback=white!97.2549!black,boxrule=0.2mm]

\begin{Shaded}
\begin{Highlighting}[]
\ControlFlowTok{for} \OperatorTok{(}\DataTypeTok{int}\NormalTok{ i }\OperatorTok{=} \DecValTok{0}\OperatorTok{;}\NormalTok{ i }\OperatorTok{\textless{}} \DecValTok{3}\OperatorTok{;}\NormalTok{ i}\OperatorTok{++)} \OperatorTok{\{}
    \ControlFlowTok{for} \OperatorTok{(}\DataTypeTok{int}\NormalTok{ j }\OperatorTok{=} \DecValTok{0}\OperatorTok{;}\NormalTok{ j }\OperatorTok{\textless{}} \DecValTok{3}\OperatorTok{;}\NormalTok{ j}\OperatorTok{++)} \OperatorTok{\{}
        \ControlFlowTok{for} \OperatorTok{(}\DataTypeTok{int}\NormalTok{ k }\OperatorTok{=} \DecValTok{0}\OperatorTok{;}\NormalTok{ k }\OperatorTok{\textless{}} \DecValTok{3}\OperatorTok{;}\NormalTok{ k}\OperatorTok{++)} \OperatorTok{\{}
\NormalTok{            printf}\OperatorTok{(}\StringTok{"}\SpecialCharTok{\%d}\StringTok{, }\SpecialCharTok{\%d}\StringTok{, }\SpecialCharTok{\%d\textbackslash{}n}\StringTok{"}\OperatorTok{,}\NormalTok{ i}\OperatorTok{,}\NormalTok{ j}\OperatorTok{,}\NormalTok{ k}\OperatorTok{);}

            \ControlFlowTok{continue}\OperatorTok{;}  \CommentTok{// Membawa kita ke iterasi k berikutnya}
            \ControlFlowTok{break}\OperatorTok{;}     \CommentTok{// Membawa kita ke iterasi j berikutnya}
            \OperatorTok{????;}      \CommentTok{// Membawa kita ke iterasi i berikutnya???}

        \OperatorTok{\}}
    \OperatorTok{\}}
\OperatorTok{\}}
\end{Highlighting}
\end{Shaded}

\end{tcolorbox}

Pernyataan \hl{\mbox{\texttt{goto}}} menawarkan jalannya!

\begin{tcolorbox}[enhanced jigsaw,left=3pt,breakable,colback=white!97.2549!black,boxrule=0.2mm]

\begin{Shaded}
\begin{Highlighting}[]
    \ControlFlowTok{for} \OperatorTok{(}\DataTypeTok{int}\NormalTok{ i }\OperatorTok{=} \DecValTok{0}\OperatorTok{;}\NormalTok{ i }\OperatorTok{\textless{}} \DecValTok{3}\OperatorTok{;}\NormalTok{ i}\OperatorTok{++)} \OperatorTok{\{}
        \ControlFlowTok{for} \OperatorTok{(}\DataTypeTok{int}\NormalTok{ j }\OperatorTok{=} \DecValTok{0}\OperatorTok{;}\NormalTok{ j }\OperatorTok{\textless{}} \DecValTok{3}\OperatorTok{;}\NormalTok{ j}\OperatorTok{++)} \OperatorTok{\{}
            \ControlFlowTok{for} \OperatorTok{(}\DataTypeTok{int}\NormalTok{ k }\OperatorTok{=} \DecValTok{0}\OperatorTok{;}\NormalTok{ k }\OperatorTok{\textless{}} \DecValTok{3}\OperatorTok{;}\NormalTok{ k}\OperatorTok{++)} \OperatorTok{\{}
\NormalTok{                printf}\OperatorTok{(}\StringTok{"}\SpecialCharTok{\%d}\StringTok{, }\SpecialCharTok{\%d}\StringTok{, }\SpecialCharTok{\%d\textbackslash{}n}\StringTok{"}\OperatorTok{,}\NormalTok{ i}\OperatorTok{,}\NormalTok{ j}\OperatorTok{,}\NormalTok{ k}\OperatorTok{);}

                \ControlFlowTok{goto}\NormalTok{ continue\_i}\OperatorTok{;}   \CommentTok{// Sekarang melanjutkan loop i!!}
            \OperatorTok{\}}
        \OperatorTok{\}}
\NormalTok{continue\_i}\OperatorTok{:} \OperatorTok{;}
    \OperatorTok{\}}
\end{Highlighting}
\end{Shaded}

\end{tcolorbox}

Ada \hl{\mbox{\texttt{;}}} di bagian akhir di sana---itu karena label tidak boleh menunjuk ke ujung polos dari sebuah compound statement (atau sebelum deklarasi variabel).

\index{goto statement@\texttt{goto} statement!as labeled continue@as labeled \texttt{continue}|)}

\section{Keluar Darurat}\label{keluar-darurat}

\index{goto statement@\texttt{goto} statement!for bailing out|(}

Ketika kamu berada di dalam kode yang sangat bersarang, kamu bisa memakai \hl{\mbox{\texttt{goto}}} untuk keluar dengan cara yang sering kali lebih bersih daripada menyarangkan lebih banyak \hl{\mbox{\texttt{if}}} dan memakai variabel flag.

\begin{tcolorbox}[enhanced jigsaw,left=3pt,breakable,colback=white!97.2549!black,boxrule=0.2mm]

\begin{Shaded}
\begin{Highlighting}[]
    \CommentTok{// Pseudocode}

    \ControlFlowTok{for}\OperatorTok{(...)} \OperatorTok{\{}
        \ControlFlowTok{for} \OperatorTok{(...)} \OperatorTok{\{}
            \ControlFlowTok{while} \OperatorTok{(...)} \OperatorTok{\{}
                \ControlFlowTok{do} \OperatorTok{\{}
                    \ControlFlowTok{if} \OperatorTok{(}\NormalTok{some\_error\_condition}\OperatorTok{)}
                        \ControlFlowTok{goto}\NormalTok{ bail}\OperatorTok{;}

                \OperatorTok{\}} \ControlFlowTok{while}\OperatorTok{(...);}
            \OperatorTok{\}}
        \OperatorTok{\}}
    \OperatorTok{\}}

\NormalTok{bail}\OperatorTok{:}
    \CommentTok{// Pembersihan di sini}
\end{Highlighting}
\end{Shaded}

\end{tcolorbox}

Tanpa \hl{\mbox{\texttt{goto}}}, kamu harus memeriksa flag kondisi error di semua loop untuk bisa keluar sepenuhnya.

\index{goto statement@\texttt{goto} statement!for bailing out|)}

\section{\texorpdfstring{\hl{\mbox{\texttt{break}}} Berlabel}{break Berlabel}}\label{break-berlabel}

\index{goto statement@\texttt{goto} statement!as labeled break@as labeled \texttt{break}|(}

Situasinya sangat mirip dengan \hl{\mbox{\texttt{continue}}} yang hanya melanjutkan loop paling dalam. \hl{\mbox{\texttt{break}}} juga hanya keluar dari loop paling dalam.

\begin{tcolorbox}[enhanced jigsaw,left=3pt,breakable,colback=white!97.2549!black,boxrule=0.2mm]

\begin{Shaded}
\begin{Highlighting}[]
    \ControlFlowTok{for} \OperatorTok{(}\DataTypeTok{int}\NormalTok{ i }\OperatorTok{=} \DecValTok{0}\OperatorTok{;}\NormalTok{ i }\OperatorTok{\textless{}} \DecValTok{3}\OperatorTok{;}\NormalTok{ i}\OperatorTok{++)} \OperatorTok{\{}
        \ControlFlowTok{for} \OperatorTok{(}\DataTypeTok{int}\NormalTok{ j }\OperatorTok{=} \DecValTok{0}\OperatorTok{;}\NormalTok{ j }\OperatorTok{\textless{}} \DecValTok{3}\OperatorTok{;}\NormalTok{ j}\OperatorTok{++)} \OperatorTok{\{}
\NormalTok{            printf}\OperatorTok{(}\StringTok{"}\SpecialCharTok{\%d}\StringTok{, }\SpecialCharTok{\%d\textbackslash{}n}\StringTok{"}\OperatorTok{,}\NormalTok{ i}\OperatorTok{,}\NormalTok{ j}\OperatorTok{);}
            \ControlFlowTok{break}\OperatorTok{;}   \CommentTok{// Hanya keluar dari loop j}
        \OperatorTok{\}}
    \OperatorTok{\}}

\NormalTok{    printf}\OperatorTok{(}\StringTok{"Done!}\SpecialCharTok{\textbackslash{}n}\StringTok{"}\OperatorTok{);}
\end{Highlighting}
\end{Shaded}

\end{tcolorbox}

Tapi kita bisa memakai \hl{\mbox{\texttt{goto}}} untuk keluar lebih jauh:

\begin{tcolorbox}[enhanced jigsaw,left=3pt,breakable,colback=white!97.2549!black,boxrule=0.2mm]

\begin{Shaded}
\begin{Highlighting}[]
    \ControlFlowTok{for} \OperatorTok{(}\DataTypeTok{int}\NormalTok{ i }\OperatorTok{=} \DecValTok{0}\OperatorTok{;}\NormalTok{ i }\OperatorTok{\textless{}} \DecValTok{3}\OperatorTok{;}\NormalTok{ i}\OperatorTok{++)} \OperatorTok{\{}
        \ControlFlowTok{for} \OperatorTok{(}\DataTypeTok{int}\NormalTok{ j }\OperatorTok{=} \DecValTok{0}\OperatorTok{;}\NormalTok{ j }\OperatorTok{\textless{}} \DecValTok{3}\OperatorTok{;}\NormalTok{ j}\OperatorTok{++)} \OperatorTok{\{}
\NormalTok{            printf}\OperatorTok{(}\StringTok{"}\SpecialCharTok{\%d}\StringTok{, }\SpecialCharTok{\%d\textbackslash{}n}\StringTok{"}\OperatorTok{,}\NormalTok{ i}\OperatorTok{,}\NormalTok{ j}\OperatorTok{);}
            \ControlFlowTok{goto}\NormalTok{ break\_i}\OperatorTok{;}   \CommentTok{// Sekarang keluar dari loop i!}
        \OperatorTok{\}}
    \OperatorTok{\}}

\NormalTok{break\_i}\OperatorTok{:}

\NormalTok{    printf}\OperatorTok{(}\StringTok{"Done!}\SpecialCharTok{\textbackslash{}n}\StringTok{"}\OperatorTok{);}
\end{Highlighting}
\end{Shaded}

\end{tcolorbox}

\index{goto statement@\texttt{goto} statement!as labeled break@as labeled \texttt{break}|)}

\section{Pembersihan Bertingkat}\label{pembersihan-bertingkat}

\index{goto statement@\texttt{goto} statement!multilevel cleanup|(}

Kalau kamu memanggil beberapa fungsi untuk menginisialisasi beberapa sistem dan salah satunya gagal, kamu hanya boleh men-deinisialisasi yang sudah kamu capai sejauh ini.

Mari kita buat contoh palsu ketika kita mulai menginisialisasi sistem dan memeriksa apakah ada yang mengembalikan error (kita pakai \hl{\mbox{\texttt{-1}}} untuk menandakan error). Kalau ada yang begitu, kita hanya perlu mematikan sistem yang sudah kita inisialisasi sejauh ini.

\begin{tcolorbox}[enhanced jigsaw,left=3pt,breakable,colback=white!97.2549!black,boxrule=0.2mm]

\begin{Shaded}
\begin{Highlighting}[]
    \ControlFlowTok{if} \OperatorTok{(}\NormalTok{init\_system\_1}\OperatorTok{()} \OperatorTok{==} \OperatorTok{{-}}\DecValTok{1}\OperatorTok{)}
        \ControlFlowTok{goto}\NormalTok{ shutdown}\OperatorTok{;}

    \ControlFlowTok{if} \OperatorTok{(}\NormalTok{init\_system\_2}\OperatorTok{()} \OperatorTok{==} \OperatorTok{{-}}\DecValTok{1}\OperatorTok{)}
        \ControlFlowTok{goto}\NormalTok{ shutdown\_1}\OperatorTok{;}

    \ControlFlowTok{if} \OperatorTok{(}\NormalTok{init\_system\_3}\OperatorTok{()} \OperatorTok{==} \OperatorTok{{-}}\DecValTok{1}\OperatorTok{)}
        \ControlFlowTok{goto}\NormalTok{ shutdown\_2}\OperatorTok{;}

    \ControlFlowTok{if} \OperatorTok{(}\NormalTok{init\_system\_4}\OperatorTok{()} \OperatorTok{==} \OperatorTok{{-}}\DecValTok{1}\OperatorTok{)}
        \ControlFlowTok{goto}\NormalTok{ shutdown\_3}\OperatorTok{;}

\NormalTok{    do\_main\_thing}\OperatorTok{();}   \CommentTok{// Jalankan program kita}

\NormalTok{    shutdown\_system4}\OperatorTok{();}

\NormalTok{shutdown\_3}\OperatorTok{:}
\NormalTok{    shutdown\_system3}\OperatorTok{();}

\NormalTok{shutdown\_2}\OperatorTok{:}
\NormalTok{    shutdown\_system2}\OperatorTok{();}

\NormalTok{shutdown\_1}\OperatorTok{:}
\NormalTok{    shutdown\_system1}\OperatorTok{();}

\NormalTok{shutdown}\OperatorTok{:}
\NormalTok{    print}\OperatorTok{(}\StringTok{"All subsystems shut down.}\SpecialCharTok{\textbackslash{}n}\StringTok{"}\OperatorTok{);}
\end{Highlighting}
\end{Shaded}

\end{tcolorbox}

Perhatikan bahwa kita mematikan dalam urutan terbalik dari urutan inisialisasi subsistem. Jadi kalau subsistem 4 gagal mulai, ia akan mematikan 3, 2, lalu 1 dalam urutan itu.

\index{goto statement@\texttt{goto} statement!multilevel cleanup|)}

\section{Tail Call Optimization}\label{tail-call-optimization}

\index{goto statement@\texttt{goto} statement!tail call optimzation|(} \index{Tail call optimzation!with goto@with \texttt{goto}|(}

Kurang lebih. Hanya untuk fungsi rekursif.

Kalau kamu belum akrab, \href{https://en.wikipedia.org/wiki/Tail_call}{Tail Call Optimization (TCO)}\footnote{https://en.wikipedia.org/wiki/Tail\_call} adalah cara agar tidak menyia-nyiakan ruang stack saat memanggil fungsi lain pada keadaan yang sangat spesifik. Sayangnya detailnya di luar cakupan panduan ini.

Tapi kalau kamu punya fungsi rekursif yang kamu tahu bisa dioptimasi dengan cara ini, kamu bisa memakai teknik ini. (Perhatikan bahwa kamu tidak bisa melakukan tail call ke fungsi lain karena label punya function scope.)

Kita buat contoh yang lugas, faktorial.

Berikut versi rekursif yang bukan TCO, tapi bisa dijadikan TCO!

\begin{tcolorbox}[enhanced jigsaw,left=3pt,breakable,colback=white!97.2549!black,boxrule=0.2mm]

\begin{Shaded}
\begin{Highlighting}[numbers=left,,]
\PreprocessorTok{\#include }\ImportTok{\textless{}stdio.h\textgreater{}}
\PreprocessorTok{\#include }\ImportTok{\textless{}complex.h\textgreater{}}

\DataTypeTok{int}\NormalTok{ factorial}\OperatorTok{(}\DataTypeTok{int}\NormalTok{ n}\OperatorTok{,} \DataTypeTok{int}\NormalTok{ a}\OperatorTok{)}
\OperatorTok{\{}
    \ControlFlowTok{if} \OperatorTok{(}\NormalTok{n }\OperatorTok{==} \DecValTok{0}\OperatorTok{)}
        \ControlFlowTok{return}\NormalTok{ a}\OperatorTok{;}

    \ControlFlowTok{return}\NormalTok{ factorial}\OperatorTok{(}\NormalTok{n }\OperatorTok{{-}} \DecValTok{1}\OperatorTok{,}\NormalTok{ a }\OperatorTok{*}\NormalTok{ n}\OperatorTok{);}
\OperatorTok{\}}

\DataTypeTok{int}\NormalTok{ main}\OperatorTok{(}\DataTypeTok{void}\OperatorTok{)}
\OperatorTok{\{}
    \ControlFlowTok{for} \OperatorTok{(}\DataTypeTok{int}\NormalTok{ i }\OperatorTok{=} \DecValTok{0}\OperatorTok{;}\NormalTok{ i }\OperatorTok{\textless{}} \DecValTok{8}\OperatorTok{;}\NormalTok{ i}\OperatorTok{++)}
\NormalTok{        printf}\OperatorTok{(}\StringTok{"}\SpecialCharTok{\%d}\StringTok{! == }\SpecialCharTok{\%ld\textbackslash{}n}\StringTok{"}\OperatorTok{,}\NormalTok{ i}\OperatorTok{,}\NormalTok{ factorial}\OperatorTok{(}\NormalTok{i}\OperatorTok{,} \DecValTok{1}\OperatorTok{));}
\OperatorTok{\}}
\end{Highlighting}
\end{Shaded}

\end{tcolorbox}

Untuk mewujudkannya, kamu bisa mengganti pemanggilan itu dengan dua langkah:

\begin{enumerate}
\def\labelenumi{\arabic{enumi}.}
\tightlist
\item
  Atur nilai parameter menjadi nilai yang akan dipakai pada pemanggilan berikutnya.
\item
  \hl{\mbox{\texttt{goto}}} ke label di baris pertama fungsi.
\end{enumerate}

Kita coba:

\begin{tcolorbox}[enhanced jigsaw,left=3pt,breakable,colback=white!97.2549!black,boxrule=0.2mm]

\begin{Shaded}
\begin{Highlighting}[numbers=left,,]
\PreprocessorTok{\#include }\ImportTok{\textless{}stdio.h\textgreater{}}

\DataTypeTok{int}\NormalTok{ factorial}\OperatorTok{(}\DataTypeTok{int}\NormalTok{ n}\OperatorTok{,} \DataTypeTok{int}\NormalTok{ a}\OperatorTok{)}
\OperatorTok{\{}
\NormalTok{tco}\OperatorTok{:}  \CommentTok{// tambahkan ini}

    \ControlFlowTok{if} \OperatorTok{(}\NormalTok{n }\OperatorTok{==} \DecValTok{0}\OperatorTok{)}
        \ControlFlowTok{return}\NormalTok{ a}\OperatorTok{;}

    \CommentTok{// ganti return dengan mengatur nilai parameter baru dan}
    \CommentTok{// goto ke awal fungsi}

    \CommentTok{//return factorial(n {-} 1, a * n);}

    \DataTypeTok{int}\NormalTok{ next\_n }\OperatorTok{=}\NormalTok{ n }\OperatorTok{{-}} \DecValTok{1}\OperatorTok{;}  \CommentTok{// Lihat bagaimana ini cocok dengan}
    \DataTypeTok{int}\NormalTok{ next\_a }\OperatorTok{=}\NormalTok{ a }\OperatorTok{*}\NormalTok{ n}\OperatorTok{;}  \CommentTok{// argumen rekursif di atas?}

\NormalTok{    n }\OperatorTok{=}\NormalTok{ next\_n}\OperatorTok{;}   \CommentTok{// Atur parameter ke nilai baru}
\NormalTok{    a }\OperatorTok{=}\NormalTok{ next\_a}\OperatorTok{;}

    \ControlFlowTok{goto}\NormalTok{ tco}\OperatorTok{;}   \CommentTok{// Dan ulangi!}
\OperatorTok{\}}

\DataTypeTok{int}\NormalTok{ main}\OperatorTok{(}\DataTypeTok{void}\OperatorTok{)}
\OperatorTok{\{}
    \ControlFlowTok{for} \OperatorTok{(}\DataTypeTok{int}\NormalTok{ i }\OperatorTok{=} \DecValTok{0}\OperatorTok{;}\NormalTok{ i }\OperatorTok{\textless{}} \DecValTok{8}\OperatorTok{;}\NormalTok{ i}\OperatorTok{++)}
\NormalTok{        printf}\OperatorTok{(}\StringTok{"}\SpecialCharTok{\%d}\StringTok{! == }\SpecialCharTok{\%d\textbackslash{}n}\StringTok{"}\OperatorTok{,}\NormalTok{ i}\OperatorTok{,}\NormalTok{ factorial}\OperatorTok{(}\NormalTok{i}\OperatorTok{,} \DecValTok{1}\OperatorTok{));}
\OperatorTok{\}}
\end{Highlighting}
\end{Shaded}

\end{tcolorbox}

Di atas, aku memakai variabel sementara untuk mengatur nilai parameter berikutnya sebelum melompat ke awal fungsi. Lihat bagaimana mereka cocok dengan argumen rekursif pada pemanggilan rekursifnya?

Nah, kenapa memakai variabel sementara? Aku bisa saja melakukan ini:

\begin{tcolorbox}[enhanced jigsaw,left=3pt,breakable,colback=white!97.2549!black,boxrule=0.2mm]

\begin{Shaded}
\begin{Highlighting}[]
\NormalTok{    a }\OperatorTok{*=}\NormalTok{ n}\OperatorTok{;}
\NormalTok{    n }\OperatorTok{{-}=} \DecValTok{1}\OperatorTok{;}

    \ControlFlowTok{goto}\NormalTok{ tco}\OperatorTok{;}
\end{Highlighting}
\end{Shaded}

\end{tcolorbox}

dan itu sebenarnya berfungsi dengan baik. Tapi kalau aku ceroboh menukar dua baris kode itu:

\begin{tcolorbox}[enhanced jigsaw,left=3pt,breakable,colback=white!97.2549!black,boxrule=0.2mm]

\begin{Shaded}
\begin{Highlighting}[]
\NormalTok{    n }\OperatorTok{{-}=} \DecValTok{1}\OperatorTok{;}  \CommentTok{// KABAR BURUK}
\NormalTok{    a }\OperatorTok{*=}\NormalTok{ n}\OperatorTok{;}
\end{Highlighting}
\end{Shaded}

\end{tcolorbox}

---sekarang kita dalam masalah. Kita mengubah \hl{\mbox{\texttt{n}}} sebelum memakainya untuk mengubah \hl{\mbox{\texttt{a}}}. Itu Buruk karena memang bukan begitu cara kerjanya saat memanggil secara rekursif. Memakai variabel sementara menghindari masalah ini bahkan kalau kamu tidak waspada. Dan compiler kemungkinan mengoptimasi mereka hingga hilang, bagaimanapun juga.

\index{goto statement@\texttt{goto} statement!tail call optimzation|)} \index{Tail call optimzation!with goto@with \texttt{goto}|)}

\section{Memulai Ulang System Call yang Terinterupsi}\label{memulai-ulang-system-call-yang-terinterupsi}

\index{goto statement@\texttt{goto} statement!restarting system calls|(}

Ini di luar spesifikasi, tapi umum terlihat di sistem mirip Unix.

Beberapa system call yang berjalan lama mungkin mengembalikan error kalau diinterupsi oleh signal, dan \hl{\mbox{\texttt{errno}}} akan diatur ke \hl{\mbox{\texttt{EINTR}}} untuk menandakan bahwa syscall itu baik-baik saja; ia hanya terinterupsi.

Dalam kasus itu, sangat umum programmer ingin memulai ulang pemanggilan itu dan mencobanya lagi.

\begin{tcolorbox}[enhanced jigsaw,left=3pt,breakable,colback=white!97.2549!black,boxrule=0.2mm]

\begin{Shaded}
\begin{Highlighting}[]
\NormalTok{retry}\OperatorTok{:}
\NormalTok{    byte\_count }\OperatorTok{=}\NormalTok{ read}\OperatorTok{(}\DecValTok{0}\OperatorTok{,}\NormalTok{ buf}\OperatorTok{,} \KeywordTok{sizeof}\OperatorTok{(}\NormalTok{buf}\OperatorTok{)} \OperatorTok{{-}} \DecValTok{1}\OperatorTok{);}  \CommentTok{// syscall read() Unix}

    \ControlFlowTok{if} \OperatorTok{(}\NormalTok{byte\_count }\OperatorTok{==} \OperatorTok{{-}}\DecValTok{1}\OperatorTok{)} \OperatorTok{\{}            \CommentTok{// Terjadi error...}
        \ControlFlowTok{if} \OperatorTok{(}\NormalTok{errno }\OperatorTok{==}\NormalTok{ EINTR}\OperatorTok{)} \OperatorTok{\{}          \CommentTok{// Tapi hanya karena terinterupsi}
\NormalTok{            printf}\OperatorTok{(}\StringTok{"Restarting...}\SpecialCharTok{\textbackslash{}n}\StringTok{"}\OperatorTok{);}
            \ControlFlowTok{goto}\NormalTok{ retry}\OperatorTok{;}
        \OperatorTok{\}}
\end{Highlighting}
\end{Shaded}

\end{tcolorbox}

Banyak sistem mirip Unix punya flag \hl{\mbox{\texttt{SA\_RESTART}}} yang bisa dioper ke \hl{\mbox{\texttt{sigaction()}}} untuk meminta OS memulai ulang secara otomatis syscall lambat apa pun alih-alih gagal dengan \hl{\mbox{\texttt{EINTR}}}.

Sekali lagi, ini khusus Unix dan di luar standar C.

Meski begitu, teknik serupa bisa dipakai kapan pun sebuah fungsi perlu dimulai ulang.

\index{goto statement@\texttt{goto} statement!restarting system calls|)}

\section{\texorpdfstring{\hl{\mbox{\texttt{goto}}} dan Thread Preemption}{goto dan Thread Preemption}}\label{goto-dan-thread-preemption}

\index{goto statement@\texttt{goto} statement!thread preemption|(}

Contoh ini diambil langsung dari \href{http://www.ostep.org/}{\emph{Operating Systems: Three Easy Pieces}}, buku luar biasa lain dari penulis yang sepaham dan juga merasa buku bermutu seharusnya bisa diunduh gratis. Bukan berarti aku keras kepala atau semacamnya.

\begin{tcolorbox}[enhanced jigsaw,left=3pt,breakable,colback=white!97.2549!black,boxrule=0.2mm]

\begin{Shaded}
\begin{Highlighting}[]
\NormalTok{retry}\OperatorTok{:}

\NormalTok{    pthread\_mutex\_lock}\OperatorTok{(}\NormalTok{L1}\OperatorTok{);}

    \ControlFlowTok{if} \OperatorTok{(}\NormalTok{pthread\_mutex\_trylock}\OperatorTok{(}\NormalTok{L2}\OperatorTok{)} \OperatorTok{!=} \DecValTok{0}\OperatorTok{)} \OperatorTok{\{}
\NormalTok{        pthread\_mutex\_unlock}\OperatorTok{(}\NormalTok{L1}\OperatorTok{);}
        \ControlFlowTok{goto}\NormalTok{ retry}\OperatorTok{;}
    \OperatorTok{\}}

\NormalTok{    save\_the\_day}\OperatorTok{();}

\NormalTok{    pthread\_mutex\_unlock}\OperatorTok{(}\NormalTok{L2}\OperatorTok{);}
\NormalTok{    pthread\_mutex\_unlock}\OperatorTok{(}\NormalTok{L1}\OperatorTok{);}
\end{Highlighting}
\end{Shaded}

\end{tcolorbox}

Di situ thread dengan senang hati mendapatkan mutex \hl{\mbox{\texttt{L1}}}, tapi kemudian berpotensi gagal mendapatkan sumber daya kedua yang dijaga mutex \hl{\mbox{\texttt{L2}}} (misalnya kalau thread lain yang tidak kooperatif memegangnya). Kalau thread kita tidak bisa mendapatkan lock \hl{\mbox{\texttt{L2}}}, ia melepas \hl{\mbox{\texttt{L1}}} lalu memakai \hl{\mbox{\texttt{goto}}} untuk mencoba ulang dengan rapi.

Kita berharap thread pahlawan kita akhirnya berhasil mendapatkan kedua mutex dan menyelamatkan hari, sambil menghindari deadlock jahat.

\index{goto statement@\texttt{goto} statement!thread preemption|)}

\section{\texorpdfstring{\hl{\mbox{\texttt{goto}}} dan Scope Variabel}{goto dan Scope Variabel}}\label{goto-dan-scope-variabel}

\index{goto statement@\texttt{goto} statement!variable scope|(}

Kita sudah melihat bahwa label punya function scope, tapi hal-hal aneh bisa terjadi kalau kita melompati inisialisasi variabel.

Lihat contoh ini ketika kita melompat dari tempat variabel \hl{\mbox{\texttt{x}}} berada di luar scope ke tengah scope-nya (di dalam blok).

\begin{tcolorbox}[enhanced jigsaw,left=3pt,breakable,colback=white!97.2549!black,boxrule=0.2mm]

\begin{Shaded}
\begin{Highlighting}[]
    \ControlFlowTok{goto}\NormalTok{ label}\OperatorTok{;}

    \OperatorTok{\{}
        \DataTypeTok{int}\NormalTok{ x }\OperatorTok{=} \DecValTok{12345}\OperatorTok{;}

\NormalTok{label}\OperatorTok{:}
\NormalTok{        printf}\OperatorTok{(}\StringTok{"}\SpecialCharTok{\%d\textbackslash{}n}\StringTok{"}\OperatorTok{,}\NormalTok{ x}\OperatorTok{);}
    \OperatorTok{\}}
\end{Highlighting}
\end{Shaded}

\end{tcolorbox}

Ini akan dikompilasi dan dijalankan, tapi aku mendapat peringatan:

\begin{tcolorbox}[enhanced jigsaw,left=3pt,breakable,colback=white!97.2549!black,boxrule=0.2mm]

\begin{Shaded}
\begin{Highlighting}[]
\NormalTok{warning: ‘x’ is used uninitialized in this function}
\end{Highlighting}
\end{Shaded}

\end{tcolorbox}

Lalu mencetak \hl{\mbox{\texttt{0}}} saat kujalankan (hasilmu mungkin berbeda).

Pada dasarnya yang terjadi: kita melompat ke scope \hl{\mbox{\texttt{x}}} (jadi boleh merujuknya di \hl{\mbox{\texttt{printf()}}}) tapi melompati baris yang benar-benar menginisialisasinya ke \hl{\mbox{\texttt{12345}}}. Jadi nilainya tak tentu.

Perbaikannya, tentu saja, menempatkan inisialisasi \emph{setelah} label dengan satu atau lain cara.

\begin{tcolorbox}[enhanced jigsaw,left=3pt,breakable,colback=white!97.2549!black,boxrule=0.2mm]

\begin{Shaded}
\begin{Highlighting}[]
    \ControlFlowTok{goto}\NormalTok{ label}\OperatorTok{;}

    \OperatorTok{\{}
        \DataTypeTok{int}\NormalTok{ x}\OperatorTok{;}

\NormalTok{label}\OperatorTok{:}
\NormalTok{        x }\OperatorTok{=} \DecValTok{12345}\OperatorTok{;}
\NormalTok{        printf}\OperatorTok{(}\StringTok{"}\SpecialCharTok{\%d\textbackslash{}n}\StringTok{"}\OperatorTok{,}\NormalTok{ x}\OperatorTok{);}
    \OperatorTok{\}}
\end{Highlighting}
\end{Shaded}

\end{tcolorbox}

Mari kita lihat satu contoh lagi.

\begin{tcolorbox}[enhanced jigsaw,left=3pt,breakable,colback=white!97.2549!black,boxrule=0.2mm]

\begin{Shaded}
\begin{Highlighting}[]
    \OperatorTok{\{}
        \DataTypeTok{int}\NormalTok{ x }\OperatorTok{=} \DecValTok{10}\OperatorTok{;}

\NormalTok{label}\OperatorTok{:}

\NormalTok{        printf}\OperatorTok{(}\StringTok{"}\SpecialCharTok{\%d\textbackslash{}n}\StringTok{"}\OperatorTok{,}\NormalTok{ x}\OperatorTok{);}
    \OperatorTok{\}}

    \ControlFlowTok{goto}\NormalTok{ label}\OperatorTok{;}
\end{Highlighting}
\end{Shaded}

\end{tcolorbox}

Apa yang terjadi di sini?

Pada kali pertama melewati blok, semuanya baik. \hl{\mbox{\texttt{x}}} bernilai \hl{\mbox{\texttt{10}}} dan itulah yang tercetak.

Tapi setelah \hl{\mbox{\texttt{goto}}}, kita melompat ke scope \hl{\mbox{\texttt{x}}}, tapi melewati inisialisasinya. Artinya kita masih bisa mencetaknya, tapi nilainya tak tentu (karena belum diinisialisasi ulang).

Di mesinku, ia mencetak \hl{\mbox{\texttt{10}}} lagi (sampai tak terhingga), tapi itu hanya keberuntungan. Ia bisa mencetak nilai apa pun setelah \hl{\mbox{\texttt{goto}}} karena \hl{\mbox{\texttt{x}}} belum diinisialisasi.

\index{goto statement@\texttt{goto} statement!variable scope|)}

\section{\texorpdfstring{\hl{\mbox{\texttt{goto}}} dan Variable-Length Array}{goto dan Variable-Length Array}}\label{goto-dan-variable-length-array}

\index{goto statement@\texttt{goto} statement!with variable-length arrays|(}

Soal VLA dan \hl{\mbox{\texttt{goto}}}, hanya ada satu aturan: kamu tidak boleh melompat dari luar scope VLA ke dalam scope VLA itu.

Kalau aku mencoba ini:

\begin{tcolorbox}[enhanced jigsaw,left=3pt,breakable,colback=white!97.2549!black,boxrule=0.2mm]

\begin{Shaded}
\begin{Highlighting}[]
    \DataTypeTok{int}\NormalTok{ x }\OperatorTok{=} \DecValTok{10}\OperatorTok{;}

    \ControlFlowTok{goto}\NormalTok{ label}\OperatorTok{;}

    \OperatorTok{\{}
        \DataTypeTok{int}\NormalTok{ v}\OperatorTok{[}\NormalTok{x}\OperatorTok{];}

\NormalTok{label}\OperatorTok{:}

\NormalTok{        printf}\OperatorTok{(}\StringTok{"Hi!}\SpecialCharTok{\textbackslash{}n}\StringTok{"}\OperatorTok{);}
    \OperatorTok{\}}
\end{Highlighting}
\end{Shaded}

\end{tcolorbox}

aku mendapat error:

\begin{tcolorbox}[enhanced jigsaw,left=3pt,breakable,colback=white!97.2549!black,boxrule=0.2mm]

\begin{Shaded}
\begin{Highlighting}[]
\NormalTok{error: jump into scope of identifier with variably modified type}
\end{Highlighting}
\end{Shaded}

\end{tcolorbox}

Kamu bisa melompat masuk sebelum deklarasi VLA, seperti ini:

\begin{tcolorbox}[enhanced jigsaw,left=3pt,breakable,colback=white!97.2549!black,boxrule=0.2mm]

\begin{Shaded}
\begin{Highlighting}[]
    \DataTypeTok{int}\NormalTok{ x }\OperatorTok{=} \DecValTok{10}\OperatorTok{;}

    \ControlFlowTok{goto}\NormalTok{ label}\OperatorTok{;}

    \OperatorTok{\{}
\NormalTok{label}\OperatorTok{:}  \OperatorTok{;}
        \DataTypeTok{int}\NormalTok{ v}\OperatorTok{[}\NormalTok{x}\OperatorTok{];}

\NormalTok{        printf}\OperatorTok{(}\StringTok{"Hi!}\SpecialCharTok{\textbackslash{}n}\StringTok{"}\OperatorTok{);}
    \OperatorTok{\}}
\end{Highlighting}
\end{Shaded}

\end{tcolorbox}

Karena dengan begitu VLA dialokasikan dengan benar sebelum dealokasinya yang tak terelakkan saat keluar dari scope.

\index{goto statement@\texttt{goto} statement!with variable-length arrays|)} \index{goto statement@\texttt{goto} statement|)}

\chapter{Tipe Bagian V: Compound Literal dan Generic Selection}\label{tipe-bagian-v-compound-literal-dan-generic-selection}

Ini bab terakhir tentang tipe! Kita akan membahas dua hal:

\begin{itemize}
\tightlist
\item
  Cara memiliki object ``anonim'' tanpa nama dan apa gunanya.
\item
  Cara menghasilkan kode yang bergantung pada tipe.
\end{itemize}

Keduanya tidak terlalu berkaitan, tapi masing-masing tidak cukup layak dijadikan satu bab sendiri. Jadi aku menjejalkannya di sini seperti pemberontak!

\section{Compound Literal}\label{compound-literal}

\index{Compound literals|(}

Ini fitur bahasa yang keren, yang memungkinkan kamu membuat object bertipe tertentu secara langsung tanpa pernah menugaskannya ke variabel. Kamu bisa membuat tipe sederhana, array, \hl{\mbox{\texttt{struct}}}, apa saja.

Salah satu kegunaan utamanya adalah mengoper argumen kompleks ke fungsi ketika kamu tidak ingin membuat variabel sementara untuk menampung nilainya.

Cara membuat compound literal adalah menaruh nama tipe dalam tanda kurung, lalu menaruh initializer list setelahnya. Misalnya, array \hl{\mbox{\texttt{int}}} tanpa nama mungkin tampak seperti ini:

\begin{tcolorbox}[enhanced jigsaw,left=3pt,breakable,colback=white!97.2549!black,boxrule=0.2mm]

\begin{Shaded}
\begin{Highlighting}[]
\OperatorTok{(}\DataTypeTok{int} \OperatorTok{[])\{}\DecValTok{1}\OperatorTok{,}\DecValTok{2}\OperatorTok{,}\DecValTok{3}\OperatorTok{,}\DecValTok{4}\OperatorTok{\}}
\end{Highlighting}
\end{Shaded}

\end{tcolorbox}

Nah, baris kode itu sendiri tidak melakukan apa-apa. Ia membuat array tanpa nama berisi 4 \hl{\mbox{\texttt{int}}}, lalu membuangnya tanpa dipakai.

Kita bisa memakai pointer untuk menyimpan referensi ke array itu\ldots{}

\begin{tcolorbox}[enhanced jigsaw,left=3pt,breakable,colback=white!97.2549!black,boxrule=0.2mm]

\begin{Shaded}
\begin{Highlighting}[]
\DataTypeTok{int} \OperatorTok{*}\NormalTok{p }\OperatorTok{=} \OperatorTok{(}\DataTypeTok{int} \OperatorTok{[])\{}\DecValTok{1} \OperatorTok{,}\DecValTok{2} \OperatorTok{,}\DecValTok{3} \OperatorTok{,}\DecValTok{4}\OperatorTok{\};}

\NormalTok{printf}\OperatorTok{(}\StringTok{"}\SpecialCharTok{\%d\textbackslash{}n}\StringTok{"}\OperatorTok{,}\NormalTok{ p}\OperatorTok{[}\DecValTok{1}\OperatorTok{]);}  \CommentTok{// 2}
\end{Highlighting}
\end{Shaded}

\end{tcolorbox}

Tapi itu terasa seperti cara bertele-tele untuk punya array. Maksudnya, kita bisa saja melakukan ini\footnote{Yang tidak persis sama, karena ia array, bukan pointer ke \hl{\mbox{\texttt{int}}}.}:

\begin{tcolorbox}[enhanced jigsaw,left=3pt,breakable,colback=white!97.2549!black,boxrule=0.2mm]

\begin{Shaded}
\begin{Highlighting}[]
\DataTypeTok{int}\NormalTok{ p}\OperatorTok{[]} \OperatorTok{=} \OperatorTok{\{}\DecValTok{1}\OperatorTok{,} \DecValTok{2}\OperatorTok{,} \DecValTok{3}\OperatorTok{,} \DecValTok{4}\OperatorTok{\};}

\NormalTok{printf}\OperatorTok{(}\StringTok{"}\SpecialCharTok{\%d\textbackslash{}n}\StringTok{"}\OperatorTok{,}\NormalTok{ p}\OperatorTok{[}\DecValTok{1}\OperatorTok{]);}  \CommentTok{// 2}
\end{Highlighting}
\end{Shaded}

\end{tcolorbox}

Jadi mari kita lihat contoh yang lebih berguna.

\subsection{Mengoper Object Tanpa Nama ke Fungsi}\label{mengoper-object-tanpa-nama-ke-fungsi}

\index{Compound literals!passing to functions|(}

Katakanlah kita punya fungsi untuk menjumlahkan array \hl{\mbox{\texttt{int}}}:

\begin{tcolorbox}[enhanced jigsaw,left=3pt,breakable,colback=white!97.2549!black,boxrule=0.2mm]

\begin{Shaded}
\begin{Highlighting}[]
\DataTypeTok{int}\NormalTok{ sum}\OperatorTok{(}\DataTypeTok{int}\NormalTok{ p}\OperatorTok{[],} \DataTypeTok{int}\NormalTok{ count}\OperatorTok{)}
\OperatorTok{\{}
    \DataTypeTok{int}\NormalTok{ total }\OperatorTok{=} \DecValTok{0}\OperatorTok{;}

    \ControlFlowTok{for} \OperatorTok{(}\DataTypeTok{int}\NormalTok{ i }\OperatorTok{=} \DecValTok{0}\OperatorTok{;}\NormalTok{ i }\OperatorTok{\textless{}}\NormalTok{ count}\OperatorTok{;}\NormalTok{ i}\OperatorTok{++)}
\NormalTok{        total }\OperatorTok{+=}\NormalTok{ p}\OperatorTok{[}\NormalTok{i}\OperatorTok{];}

    \ControlFlowTok{return}\NormalTok{ total}\OperatorTok{;}
\OperatorTok{\}}
\end{Highlighting}
\end{Shaded}

\end{tcolorbox}

Kalau ingin memanggilnya, biasanya kita harus melakukan sesuatu seperti ini, mendeklarasikan array dan mengisinya dengan nilai untuk dioper ke fungsi:

\begin{tcolorbox}[enhanced jigsaw,left=3pt,breakable,colback=white!97.2549!black,boxrule=0.2mm]

\begin{Shaded}
\begin{Highlighting}[]
\DataTypeTok{int}\NormalTok{ a}\OperatorTok{[]} \OperatorTok{=} \OperatorTok{\{}\DecValTok{1}\OperatorTok{,} \DecValTok{2}\OperatorTok{,} \DecValTok{3}\OperatorTok{,} \DecValTok{4}\OperatorTok{\};}

\DataTypeTok{int}\NormalTok{ s }\OperatorTok{=}\NormalTok{ sum}\OperatorTok{(}\NormalTok{a}\OperatorTok{,} \DecValTok{4}\OperatorTok{);}
\end{Highlighting}
\end{Shaded}

\end{tcolorbox}

Tapi object tanpa nama memberi kita cara untuk melewati variabel itu dengan mengopernya langsung (nama parameter tercantum di atas). Lihat ini---kita akan mengganti variabel \hl{\mbox{\texttt{a}}} dengan array tanpa nama yang dioper sebagai argumen pertama:

\begin{tcolorbox}[enhanced jigsaw,left=3pt,breakable,colback=white!97.2549!black,boxrule=0.2mm]

\begin{Shaded}
\begin{Highlighting}[]
\CommentTok{//                   p[]         count}
\CommentTok{//           |{-}{-}{-}{-}{-}{-}{-}{-}{-}{-}{-}{-}{-}{-}{-}{-}{-}|  |}
\DataTypeTok{int}\NormalTok{ s }\OperatorTok{=}\NormalTok{ sum}\OperatorTok{((}\DataTypeTok{int} \OperatorTok{[])\{}\DecValTok{1}\OperatorTok{,} \DecValTok{2}\OperatorTok{,} \DecValTok{3}\OperatorTok{,} \DecValTok{4}\OperatorTok{\},} \DecValTok{4}\OperatorTok{);}
\end{Highlighting}
\end{Shaded}

\end{tcolorbox}

Keren sekali!

\index{Compound literals!passing to functions|)}

\subsection{\texorpdfstring{\hl{\mbox{\texttt{struct}}} Tanpa Nama}{struct Tanpa Nama}}\label{struct-tanpa-nama}

\index{Compound literals!with struct@with \texttt{struct}|(} \index{struct keyword@\texttt{struct} keyword!compound literals|(}

Kita bisa melakukan hal serupa dengan \hl{\mbox{\texttt{struct}}}.

Pertama, kita kerjakan tanpa object tanpa nama. Kita definisikan \hl{\mbox{\texttt{struct}}} untuk menyimpan koordinat \hl{\mbox{\texttt{x}}}/\hl{\mbox{\texttt{y}}}. Lalu kita definisikan satu, dengan mengoper nilai ke initializer-nya. Terakhir, kita oper ke fungsi untuk mencetak nilainya:

\begin{tcolorbox}[enhanced jigsaw,left=3pt,breakable,colback=white!97.2549!black,boxrule=0.2mm]

\begin{Shaded}
\begin{Highlighting}[numbers=left,,]
\PreprocessorTok{\#include }\ImportTok{\textless{}stdio.h\textgreater{}}

\KeywordTok{struct}\NormalTok{ coord }\OperatorTok{\{}
    \DataTypeTok{int}\NormalTok{ x}\OperatorTok{,}\NormalTok{ y}\OperatorTok{;}
\OperatorTok{\};}

\DataTypeTok{void}\NormalTok{ print\_coord}\OperatorTok{(}\KeywordTok{struct}\NormalTok{ coord c}\OperatorTok{)}
\OperatorTok{\{}
\NormalTok{    printf}\OperatorTok{(}\StringTok{"}\SpecialCharTok{\%d}\StringTok{, }\SpecialCharTok{\%d\textbackslash{}n}\StringTok{"}\OperatorTok{,}\NormalTok{ c}\OperatorTok{.}\NormalTok{x}\OperatorTok{,}\NormalTok{ c}\OperatorTok{.}\NormalTok{y}\OperatorTok{);}
\OperatorTok{\}}

\DataTypeTok{int}\NormalTok{ main}\OperatorTok{(}\DataTypeTok{void}\OperatorTok{)}
\OperatorTok{\{}
    \KeywordTok{struct}\NormalTok{ coord t }\OperatorTok{=} \OperatorTok{\{.}\NormalTok{x}\OperatorTok{=}\DecValTok{10}\OperatorTok{,} \OperatorTok{.}\NormalTok{y}\OperatorTok{=}\DecValTok{20}\OperatorTok{\};}

\NormalTok{    print\_coord}\OperatorTok{(}\NormalTok{t}\OperatorTok{);}   \CommentTok{// mencetak "10, 20"}
\OperatorTok{\}}
\end{Highlighting}
\end{Shaded}

\end{tcolorbox}

Cukup lugas?

Mari kita ubah agar memakai object tanpa nama sebagai ganti variabel \hl{\mbox{\texttt{t}}} yang kita oper ke \hl{\mbox{\texttt{print\_coord()}}}.

Kita cukup mengambil \hl{\mbox{\texttt{t}}} dari sana dan menggantinya dengan \hl{\mbox{\texttt{struct}}} tanpa nama:

\begin{tcolorbox}[enhanced jigsaw,left=3pt,breakable,colback=white!97.2549!black,boxrule=0.2mm]

\begin{Shaded}
\begin{Highlighting}[numbers=left,,firstnumber=7,]
    \CommentTok{//struct coord t = \{.x=10, .y=20\};}

\NormalTok{    print\_coord}\OperatorTok{((}\KeywordTok{struct}\NormalTok{ coord}\OperatorTok{)\{.}\NormalTok{x}\OperatorTok{=}\DecValTok{10}\OperatorTok{,} \OperatorTok{.}\NormalTok{y}\OperatorTok{=}\DecValTok{20}\OperatorTok{\});}   \CommentTok{// mencetak "10, 20"}
\end{Highlighting}
\end{Shaded}

\end{tcolorbox}

Tetap berfungsi!

\index{struct keyword@\texttt{struct} keyword!compound literals|)} \index{Compound literals!with struct@with \texttt{struct}|)}

\subsection{Pointer ke Object Tanpa Nama}\label{pointer-ke-object-tanpa-nama}

\index{Compound literals!pointers to|(}

Mungkin kamu memperhatikan pada contoh terakhir bahwa walaupun kita memakai \hl{\mbox{\texttt{struct}}}, yang dioper ke \hl{\mbox{\texttt{print\_coord()}}} adalah salinan \hl{\mbox{\texttt{struct}}}, bukan pointer ke \hl{\mbox{\texttt{struct}}}.

Ternyata kita bisa mengambil alamat object tanpa nama dengan \hl{\mbox{\texttt{\&}}} seperti biasa.

Ini karena, secara umum, kalau suatu operator bisa bekerja pada variabel bertipe tertentu, kamu bisa memakai operator itu pada object tanpa nama bertipe sama.

Mari kita ubah kode di atas agar mengoper pointer ke object tanpa nama

\begin{tcolorbox}[enhanced jigsaw,left=3pt,breakable,colback=white!97.2549!black,boxrule=0.2mm]

\begin{Shaded}
\begin{Highlighting}[numbers=left,,]
\PreprocessorTok{\#include }\ImportTok{\textless{}stdio.h\textgreater{}}

\KeywordTok{struct}\NormalTok{ coord }\OperatorTok{\{}
    \DataTypeTok{int}\NormalTok{ x}\OperatorTok{,}\NormalTok{ y}\OperatorTok{;}
\OperatorTok{\};}

\DataTypeTok{void}\NormalTok{ print\_coord}\OperatorTok{(}\KeywordTok{struct}\NormalTok{ coord }\OperatorTok{*}\NormalTok{c}\OperatorTok{)}
\OperatorTok{\{}
\NormalTok{    printf}\OperatorTok{(}\StringTok{"}\SpecialCharTok{\%d}\StringTok{, }\SpecialCharTok{\%d\textbackslash{}n}\StringTok{"}\OperatorTok{,}\NormalTok{ c}\OperatorTok{{-}\textgreater{}}\NormalTok{x}\OperatorTok{,}\NormalTok{ c}\OperatorTok{{-}\textgreater{}}\NormalTok{y}\OperatorTok{);}
\OperatorTok{\}}

\DataTypeTok{int}\NormalTok{ main}\OperatorTok{(}\DataTypeTok{void}\OperatorTok{)}
\OperatorTok{\{}
    \CommentTok{//     Perhatikan \&}
    \CommentTok{//          |}
\NormalTok{    print\_coord}\OperatorTok{(\&(}\KeywordTok{struct}\NormalTok{ coord}\OperatorTok{)\{.}\NormalTok{x}\OperatorTok{=}\DecValTok{10}\OperatorTok{,} \OperatorTok{.}\NormalTok{y}\OperatorTok{=}\DecValTok{20}\OperatorTok{\});}   \CommentTok{// mencetak "10, 20"}
\OperatorTok{\}}
\end{Highlighting}
\end{Shaded}

\end{tcolorbox}

Selain itu, ini bisa menjadi cara yang bagus untuk mengoper pointer ke object sederhana sekalipun:

\begin{tcolorbox}[enhanced jigsaw,left=3pt,breakable,colback=white!97.2549!black,boxrule=0.2mm]

\begin{Shaded}
\begin{Highlighting}[]
\CommentTok{// Oper pointer ke int bernilai 3490}
\NormalTok{foo}\OperatorTok{(\&(}\DataTypeTok{int}\OperatorTok{)\{}\DecValTok{3490}\OperatorTok{\});}
\end{Highlighting}
\end{Shaded}

\end{tcolorbox}

Semudah itu.

\index{Compound literals!pointers to|)}

\subsection{Object Tanpa Nama dan Scope}\label{object-tanpa-nama-dan-scope}

\index{Compound literals!scope|(}

Masa hidup object tanpa nama berakhir di akhir scope-nya. Jebakan terbesarnya: kamu membuat object tanpa nama baru, mengambil pointer ke sana, lalu meninggalkan scope object itu. Dalam kasus itu, pointer akan merujuk ke object yang sudah mati.

Jadi ini undefined behavior:

\begin{tcolorbox}[enhanced jigsaw,left=3pt,breakable,colback=white!97.2549!black,boxrule=0.2mm]

\begin{Shaded}
\begin{Highlighting}[]
\DataTypeTok{int} \OperatorTok{*}\NormalTok{p}\OperatorTok{;}

\OperatorTok{\{}
\NormalTok{    p }\OperatorTok{=} \OperatorTok{\&(}\DataTypeTok{int}\OperatorTok{)\{}\DecValTok{10}\OperatorTok{\};}
\OperatorTok{\}}

\NormalTok{printf}\OperatorTok{(}\StringTok{"}\SpecialCharTok{\%d\textbackslash{}n}\StringTok{"}\OperatorTok{,} \OperatorTok{*}\NormalTok{p}\OperatorTok{);}  \CommentTok{// TIDAK VALID: (int)\{10\} sudah keluar dari scope}
\end{Highlighting}
\end{Shaded}

\end{tcolorbox}

Begitu juga, kamu tidak bisa mengembalikan pointer ke object tanpa nama dari sebuah fungsi. Object itu didealokasikan saat keluar dari scope:

\begin{tcolorbox}[enhanced jigsaw,left=3pt,breakable,colback=white!97.2549!black,boxrule=0.2mm]

\begin{Shaded}
\begin{Highlighting}[numbers=left,,]
\PreprocessorTok{\#include }\ImportTok{\textless{}stdio.h\textgreater{}}

\DataTypeTok{int} \OperatorTok{*}\NormalTok{get3490}\OperatorTok{(}\DataTypeTok{void}\OperatorTok{)}
\OperatorTok{\{}
    \CommentTok{// Jangan lakukan ini}
    \ControlFlowTok{return} \OperatorTok{\&(}\DataTypeTok{int}\OperatorTok{)\{}\DecValTok{3490}\OperatorTok{\};}
\OperatorTok{\}}

\DataTypeTok{int}\NormalTok{ main}\OperatorTok{(}\DataTypeTok{void}\OperatorTok{)}
\OperatorTok{\{}
\NormalTok{    printf}\OperatorTok{(}\StringTok{"}\SpecialCharTok{\%d\textbackslash{}n}\StringTok{"}\OperatorTok{,} \OperatorTok{*}\NormalTok{get3490}\OperatorTok{());}  \CommentTok{// TIDAK VALID: (int)\{3490\} sudah keluar dari scope}
\OperatorTok{\}}
\end{Highlighting}
\end{Shaded}

\end{tcolorbox}

Bayangkan saja scope-nya seperti scope variabel lokal biasa. Kamu juga tidak bisa mengembalikan pointer ke variabel lokal.

\index{Compound literals!scope|)}

\subsection{Contoh Konyol Object Tanpa Nama}\label{contoh-konyol-object-tanpa-nama}

Tipe apa pun bisa ditaruh di sana untuk membuat object tanpa nama.

Misalnya, ketiga ini secara efektif setara:

\begin{tcolorbox}[enhanced jigsaw,left=3pt,breakable,colback=white!97.2549!black,boxrule=0.2mm]

\begin{Shaded}
\begin{Highlighting}[]
\DataTypeTok{int}\NormalTok{ x }\OperatorTok{=} \DecValTok{3490}\OperatorTok{;}

\NormalTok{printf}\OperatorTok{(}\StringTok{"}\SpecialCharTok{\%d\textbackslash{}n}\StringTok{"}\OperatorTok{,}\NormalTok{ x}\OperatorTok{);}               \CommentTok{// 3490 (variabel)}
\NormalTok{printf}\OperatorTok{(}\StringTok{"}\SpecialCharTok{\%d\textbackslash{}n}\StringTok{"}\OperatorTok{,} \DecValTok{3490}\OperatorTok{);}            \CommentTok{// 3490 (konstanta)}
\NormalTok{printf}\OperatorTok{(}\StringTok{"}\SpecialCharTok{\%d\textbackslash{}n}\StringTok{"}\OperatorTok{,} \OperatorTok{(}\DataTypeTok{int}\OperatorTok{)\{}\DecValTok{3490}\OperatorTok{\});}     \CommentTok{// 3490 (object tanpa nama)}
\end{Highlighting}
\end{Shaded}

\end{tcolorbox}

Yang terakhir itu tanpa nama, tapi konyol. Lebih baik pakai yang sederhana di baris sebelumnya.

Tapi semoga ini memberi sedikit kejelasan tentang sintaksnya.

\index{Compound literals|)}

\section{Generic Selection}\label{type-generics}

\index{Generic selections|(}

Ini ekspresi yang memungkinkan kamu memilih potongan kode berbeda tergantung pada \emph{tipe} argumen pertama ekspresi itu.

Sebentar lagi kita lihat contohnya, tapi penting diketahui bahwa ini diproses saat compile time, \emph{bukan saat runtime}. Tidak ada analisis runtime di sini.

\index{\_Generic keyword@\texttt{\_Generic} keyword|(}

Ekspresi ini diawali \hl{\mbox{\texttt{\_Generic}}}, bekerja mirip \hl{\mbox{\texttt{switch}}}, dan menerima setidaknya dua argumen.

Argumen pertama adalah ekspresi (atau variabel\footnote{Variabel yang dipakai di sini \emph{adalah} ekspresi.}) yang punya \emph{tipe}. Semua ekspresi punya tipe. Argumen sisanya untuk \hl{\mbox{\texttt{\_Generic}}} adalah kasus-kasus tentang apa yang disubstitusikan sebagai hasil ekspresi kalau argumen pertama bertipe itu.

Wat?

Mari kita coba dan lihat.

\begin{tcolorbox}[enhanced jigsaw,left=3pt,breakable,colback=white!97.2549!black,boxrule=0.2mm]

\begin{Shaded}
\begin{Highlighting}[numbers=left,,]
\PreprocessorTok{\#include }\ImportTok{\textless{}stdio.h\textgreater{}}

\DataTypeTok{int}\NormalTok{ main}\OperatorTok{(}\DataTypeTok{void}\OperatorTok{)}
\OperatorTok{\{}
    \DataTypeTok{int}\NormalTok{ i}\OperatorTok{;}
    \DataTypeTok{float}\NormalTok{ f}\OperatorTok{;}
    \DataTypeTok{char}\NormalTok{ c}\OperatorTok{;}

    \DataTypeTok{char} \OperatorTok{*}\NormalTok{s }\OperatorTok{=} \DataTypeTok{\_Generic}\OperatorTok{(}\NormalTok{i}\OperatorTok{,}
                    \DataTypeTok{int}\OperatorTok{:} \StringTok{"that variable is an int"}\OperatorTok{,}
                    \DataTypeTok{float}\OperatorTok{:} \StringTok{"that variable is a float"}\OperatorTok{,}
                    \ControlFlowTok{default}\OperatorTok{:} \StringTok{"that variable is some type"}
                \OperatorTok{);}

\NormalTok{    printf}\OperatorTok{(}\StringTok{"}\SpecialCharTok{\%s\textbackslash{}n}\StringTok{"}\OperatorTok{,}\NormalTok{ s}\OperatorTok{);}
\OperatorTok{\}}
\end{Highlighting}
\end{Shaded}

\end{tcolorbox}

Lihat ekspresi \hl{\mbox{\texttt{\_Generic}}} yang dimulai di baris 9.

Saat compiler melihatnya, ia melihat tipe argumen pertama. (Pada contoh ini, tipe variabel \hl{\mbox{\texttt{i}}}.) Lalu ia menelusuri kasus-kasus untuk mencari yang bertipe itu. Setelah itu ia menyubstitusikan argumen tersebut menggantikan seluruh ekspresi \hl{\mbox{\texttt{\_Generic}}}.

Dalam kasus ini, \hl{\mbox{\texttt{i}}} adalah \hl{\mbox{\texttt{int}}}, jadi cocok dengan kasus itu. Lalu string itu disubstitusikan ke ekspresi. Jadi baris itu berubah menjadi ini saat dilihat compiler:

\begin{tcolorbox}[enhanced jigsaw,left=3pt,breakable,colback=white!97.2549!black,boxrule=0.2mm]

\begin{Shaded}
\begin{Highlighting}[]
    \DataTypeTok{char} \OperatorTok{*}\NormalTok{s }\OperatorTok{=} \StringTok{"that variable is an int"}\OperatorTok{;}
\end{Highlighting}
\end{Shaded}

\end{tcolorbox}

Kalau compiler tidak menemukan kecocokan tipe di \hl{\mbox{\texttt{\_Generic}}}, ia mencari kasus \hl{\mbox{\texttt{default}}} yang opsional dan memakainya.

Kalau tidak ada kecocokan tipe dan tidak ada \hl{\mbox{\texttt{default}}}, kamu akan mendapat error kompilasi. Ekspresi pertama \textbf{harus} cocok dengan salah satu tipe atau \hl{\mbox{\texttt{default}}}.

Karena merepotkan menulis \hl{\mbox{\texttt{\_Generic}}} berulang-ulang, ia sering dipakai untuk membuat badan makro yang bisa dengan mudah dipakai ulang.

Mari kita buat makro \hl{\mbox{\texttt{TYPESTR(x)}}} yang menerima satu argumen dan mengembalikan string berisi tipe argumen itu.

Jadi \hl{\mbox{\texttt{TYPESTR(1)}}} akan mengembalikan string \hl{\mbox{\texttt{"int"}}}, misalnya.

Begini caranya:

\begin{tcolorbox}[enhanced jigsaw,left=3pt,breakable,colback=white!97.2549!black,boxrule=0.2mm]

\begin{Shaded}
\begin{Highlighting}[]
\PreprocessorTok{\#include }\ImportTok{\textless{}stdio.h\textgreater{}}

\PreprocessorTok{\#define TYPESTR}\OperatorTok{(}\PreprocessorTok{x}\OperatorTok{)}\PreprocessorTok{ }\DataTypeTok{\_Generic}\OperatorTok{((}\PreprocessorTok{x}\OperatorTok{),}\PreprocessorTok{ }\OperatorTok{\textbackslash{}}
\PreprocessorTok{                        }\DataTypeTok{int}\OperatorTok{:}\PreprocessorTok{ }\StringTok{"int"}\OperatorTok{,}\PreprocessorTok{ }\OperatorTok{\textbackslash{}}
\PreprocessorTok{                        }\DataTypeTok{long}\OperatorTok{:}\PreprocessorTok{ }\StringTok{"long"}\OperatorTok{,}\PreprocessorTok{ }\OperatorTok{\textbackslash{}}
\PreprocessorTok{                        }\DataTypeTok{float}\OperatorTok{:}\PreprocessorTok{ }\StringTok{"float"}\OperatorTok{,}\PreprocessorTok{ }\OperatorTok{\textbackslash{}}
\PreprocessorTok{                        }\DataTypeTok{double}\OperatorTok{:}\PreprocessorTok{ }\StringTok{"double"}\OperatorTok{,}\PreprocessorTok{ }\OperatorTok{\textbackslash{}}
\PreprocessorTok{                        }\ControlFlowTok{default}\OperatorTok{:}\PreprocessorTok{ }\StringTok{"something else"}\OperatorTok{)}

\DataTypeTok{int}\NormalTok{ main}\OperatorTok{(}\DataTypeTok{void}\OperatorTok{)}
\OperatorTok{\{}
    \DataTypeTok{int}\NormalTok{ i}\OperatorTok{;}
    \DataTypeTok{long}\NormalTok{ l}\OperatorTok{;}
    \DataTypeTok{float}\NormalTok{ f}\OperatorTok{;}
    \DataTypeTok{double}\NormalTok{ d}\OperatorTok{;}
    \DataTypeTok{char}\NormalTok{ c}\OperatorTok{;}

\NormalTok{    printf}\OperatorTok{(}\StringTok{"i is type }\SpecialCharTok{\%s\textbackslash{}n}\StringTok{"}\OperatorTok{,}\NormalTok{ TYPESTR}\OperatorTok{(}\NormalTok{i}\OperatorTok{));}
\NormalTok{    printf}\OperatorTok{(}\StringTok{"l is type }\SpecialCharTok{\%s\textbackslash{}n}\StringTok{"}\OperatorTok{,}\NormalTok{ TYPESTR}\OperatorTok{(}\NormalTok{l}\OperatorTok{));}
\NormalTok{    printf}\OperatorTok{(}\StringTok{"f is type }\SpecialCharTok{\%s\textbackslash{}n}\StringTok{"}\OperatorTok{,}\NormalTok{ TYPESTR}\OperatorTok{(}\NormalTok{f}\OperatorTok{));}
\NormalTok{    printf}\OperatorTok{(}\StringTok{"d is type }\SpecialCharTok{\%s\textbackslash{}n}\StringTok{"}\OperatorTok{,}\NormalTok{ TYPESTR}\OperatorTok{(}\NormalTok{d}\OperatorTok{));}
\NormalTok{    printf}\OperatorTok{(}\StringTok{"c is type }\SpecialCharTok{\%s\textbackslash{}n}\StringTok{"}\OperatorTok{,}\NormalTok{ TYPESTR}\OperatorTok{(}\NormalTok{c}\OperatorTok{));}
\OperatorTok{\}}
\end{Highlighting}
\end{Shaded}

\end{tcolorbox}

Ini menghasilkan output:

\begin{tcolorbox}[enhanced jigsaw,left=3pt,breakable,colback=white!97.2549!black,boxrule=0.2mm]

\begin{Shaded}
\begin{Highlighting}[]
\NormalTok{i is type int}
\NormalTok{l is type long}
\NormalTok{f is type float}
\NormalTok{d is type double}
\NormalTok{c is type something else}
\end{Highlighting}
\end{Shaded}

\end{tcolorbox}

Yang seharusnya tidak mengejutkan, karena, seperti yang kita katakan, kode di \hl{\mbox{\texttt{main()}}} itu diganti dengan yang berikut saat dikompilasi:

\begin{tcolorbox}[enhanced jigsaw,left=3pt,breakable,colback=white!97.2549!black,boxrule=0.2mm]

\begin{Shaded}
\begin{Highlighting}[]
\NormalTok{    printf}\OperatorTok{(}\StringTok{"i is type }\SpecialCharTok{\%s\textbackslash{}n}\StringTok{"}\OperatorTok{,} \StringTok{"int"}\OperatorTok{);}
\NormalTok{    printf}\OperatorTok{(}\StringTok{"l is type }\SpecialCharTok{\%s\textbackslash{}n}\StringTok{"}\OperatorTok{,} \StringTok{"long"}\OperatorTok{);}
\NormalTok{    printf}\OperatorTok{(}\StringTok{"f is type }\SpecialCharTok{\%s\textbackslash{}n}\StringTok{"}\OperatorTok{,} \StringTok{"float"}\OperatorTok{);}
\NormalTok{    printf}\OperatorTok{(}\StringTok{"d is type }\SpecialCharTok{\%s\textbackslash{}n}\StringTok{"}\OperatorTok{,} \StringTok{"double"}\OperatorTok{);}
\NormalTok{    printf}\OperatorTok{(}\StringTok{"c is type }\SpecialCharTok{\%s\textbackslash{}n}\StringTok{"}\OperatorTok{,} \StringTok{"something else"}\OperatorTok{);}
\end{Highlighting}
\end{Shaded}

\end{tcolorbox}

Dan itulah persis output yang kita lihat.

Satu contoh lagi. Di sini aku menyertakan beberapa makro agar ketika kamu menjalankan:

\begin{tcolorbox}[enhanced jigsaw,left=3pt,breakable,colback=white!97.2549!black,boxrule=0.2mm]

\begin{Shaded}
\begin{Highlighting}[]
\DataTypeTok{int}\NormalTok{ i }\OperatorTok{=} \DecValTok{10}\OperatorTok{;}
\DataTypeTok{char} \OperatorTok{*}\NormalTok{s }\OperatorTok{=} \StringTok{"Foo!"}\OperatorTok{;}

\NormalTok{PRINT\_VAL}\OperatorTok{(}\NormalTok{i}\OperatorTok{);}
\NormalTok{PRINT\_VAL}\OperatorTok{(}\NormalTok{s}\OperatorTok{);}
\end{Highlighting}
\end{Shaded}

\end{tcolorbox}

kamu mendapat output:

\begin{tcolorbox}[enhanced jigsaw,left=3pt,breakable,colback=white!97.2549!black,boxrule=0.2mm]

\begin{Shaded}
\begin{Highlighting}[]
\NormalTok{i = 10}
\NormalTok{s = Foo!}
\end{Highlighting}
\end{Shaded}

\end{tcolorbox}

Kita perlu memakai sedikit sihir makro untuk itu.

\begin{tcolorbox}[enhanced jigsaw,left=3pt,breakable,colback=white!97.2549!black,boxrule=0.2mm]

\begin{Shaded}
\begin{Highlighting}[numbers=left,,]
\PreprocessorTok{\#include }\ImportTok{\textless{}stdio.h\textgreater{}}
\PreprocessorTok{\#include }\ImportTok{\textless{}string.h\textgreater{}}

\CommentTok{// Makro yang mengembalikan format specifier untuk sebuah tipe}
\PreprocessorTok{\#define FMTSPEC}\OperatorTok{(}\PreprocessorTok{x}\OperatorTok{)}\PreprocessorTok{ }\DataTypeTok{\_Generic}\OperatorTok{((}\PreprocessorTok{x}\OperatorTok{),}\PreprocessorTok{ }\OperatorTok{\textbackslash{}}
\PreprocessorTok{                        }\DataTypeTok{int}\OperatorTok{:}\PreprocessorTok{ }\StringTok{"}\SpecialCharTok{\%d}\StringTok{"}\OperatorTok{,}\PreprocessorTok{ }\OperatorTok{\textbackslash{}}
\PreprocessorTok{                        }\DataTypeTok{long}\OperatorTok{:}\PreprocessorTok{ }\StringTok{"}\SpecialCharTok{\%ld}\StringTok{"}\OperatorTok{,}\PreprocessorTok{ }\OperatorTok{\textbackslash{}}
\PreprocessorTok{                        }\DataTypeTok{float}\OperatorTok{:}\PreprocessorTok{ }\StringTok{"}\SpecialCharTok{\%f}\StringTok{"}\OperatorTok{,}\PreprocessorTok{ }\OperatorTok{\textbackslash{}}
\PreprocessorTok{                        }\DataTypeTok{double}\OperatorTok{:}\PreprocessorTok{ }\StringTok{"}\SpecialCharTok{\%f}\StringTok{"}\OperatorTok{,}\PreprocessorTok{ }\OperatorTok{\textbackslash{}}
\PreprocessorTok{                        }\DataTypeTok{char}\PreprocessorTok{ }\OperatorTok{*:}\PreprocessorTok{ }\StringTok{"}\SpecialCharTok{\%s}\StringTok{"}\OperatorTok{)}
                        \CommentTok{// }\AlertTok{TODO}\CommentTok{: tambahkan tipe lain}
                        
\CommentTok{// Makro yang mencetak variabel dalam bentuk "nama = nilai"}
\PreprocessorTok{\#define PRINT\_VAL}\OperatorTok{(}\PreprocessorTok{x}\OperatorTok{)}\PreprocessorTok{ }\ControlFlowTok{do}\PreprocessorTok{ }\OperatorTok{\{}\PreprocessorTok{ }\OperatorTok{\textbackslash{}}
\PreprocessorTok{    }\DataTypeTok{char}\PreprocessorTok{ fmt}\OperatorTok{[}\DecValTok{512}\OperatorTok{];}\PreprocessorTok{ }\OperatorTok{\textbackslash{}}
\PreprocessorTok{    snprintf}\OperatorTok{(}\PreprocessorTok{fmt}\OperatorTok{,}\PreprocessorTok{ }\KeywordTok{sizeof}\PreprocessorTok{ fmt}\OperatorTok{,}\PreprocessorTok{ }\OperatorTok{\#}\PreprocessorTok{x }\StringTok{" = }\SpecialCharTok{\%s\textbackslash{}n}\StringTok{"}\OperatorTok{,}\PreprocessorTok{ FMTSPEC}\OperatorTok{(}\PreprocessorTok{x}\OperatorTok{));}\PreprocessorTok{ }\OperatorTok{\textbackslash{}}
\PreprocessorTok{    printf}\OperatorTok{(}\PreprocessorTok{fmt}\OperatorTok{,}\PreprocessorTok{ }\OperatorTok{(}\PreprocessorTok{x}\OperatorTok{));}\PreprocessorTok{ }\OperatorTok{\textbackslash{}}
\OperatorTok{\}}\PreprocessorTok{ }\ControlFlowTok{while}\OperatorTok{(}\DecValTok{0}\OperatorTok{)}

\DataTypeTok{int}\NormalTok{ main}\OperatorTok{(}\DataTypeTok{void}\OperatorTok{)}
\OperatorTok{\{}
    \DataTypeTok{int}\NormalTok{ i }\OperatorTok{=} \DecValTok{10}\OperatorTok{;}
    \DataTypeTok{float}\NormalTok{ f }\OperatorTok{=} \FloatTok{3.14159}\OperatorTok{;}
    \DataTypeTok{char} \OperatorTok{*}\NormalTok{s }\OperatorTok{=} \StringTok{"Hello, world!"}\OperatorTok{;}

\NormalTok{    PRINT\_VAL}\OperatorTok{(}\NormalTok{i}\OperatorTok{);}
\NormalTok{    PRINT\_VAL}\OperatorTok{(}\NormalTok{f}\OperatorTok{);}
\NormalTok{    PRINT\_VAL}\OperatorTok{(}\NormalTok{s}\OperatorTok{);}
\OperatorTok{\}}
\end{Highlighting}
\end{Shaded}

\end{tcolorbox}

\index{\_Generic keyword@\texttt{\_Generic} keyword|)}

untuk output:

\begin{tcolorbox}[enhanced jigsaw,left=3pt,breakable,colback=white!97.2549!black,boxrule=0.2mm]

\begin{Shaded}
\begin{Highlighting}[]
\NormalTok{i = 10}
\NormalTok{f = 3.141590}
\NormalTok{s = Hello, world!}
\end{Highlighting}
\end{Shaded}

\end{tcolorbox}

Kita bisa saja menjejalkan semuanya dalam satu makro besar, tapi aku memecahnya jadi dua supaya mata tidak berdarah.

\index{Generic selections|)}

\chapter{Array Bagian II}\label{array-bagian-ii}

Di bab ini kita akan membahas beberapa hal serba-serbi tambahan tentang array.

\begin{itemize}
\tightlist
\item
  Type qualifier pada parameter array
\item
  Kata kunci \hl{\mbox{\texttt{static}}} pada parameter array
\item
  Initializer array multidimensi parsial
\end{itemize}

Hal-hal ini tidak terlalu sering terlihat, tapi kita intip saja karena termasuk bagian spesifikasi yang lebih baru.

\section{Type Qualifier untuk Array dalam Daftar Parameter}\label{type-qualifier-untuk-array-dalam-daftar-parameter}

\index{Type qualifiers!arrays in parameter lists|(} \index{Arrays!type qualifiers in parameter lists|(}

Kalau kamu ingat dari bab sebelumnya, dua hal ini setara dalam daftar parameter fungsi:

\begin{tcolorbox}[enhanced jigsaw,left=3pt,breakable,colback=white!97.2549!black,boxrule=0.2mm]

\begin{Shaded}
\begin{Highlighting}[]
\DataTypeTok{int}\NormalTok{ func}\OperatorTok{(}\DataTypeTok{int} \OperatorTok{*}\NormalTok{p}\OperatorTok{)} \OperatorTok{\{...\}}
\DataTypeTok{int}\NormalTok{ func}\OperatorTok{(}\DataTypeTok{int}\NormalTok{ p}\OperatorTok{[])} \OperatorTok{\{...\}}
\end{Highlighting}
\end{Shaded}

\end{tcolorbox}

Dan mungkin kamu juga ingat bahwa kamu bisa menambahkan type qualifier ke variabel pointer seperti ini:

\begin{tcolorbox}[enhanced jigsaw,left=3pt,breakable,colback=white!97.2549!black,boxrule=0.2mm]

\begin{Shaded}
\begin{Highlighting}[]
\DataTypeTok{int} \OperatorTok{*}\DataTypeTok{const}\NormalTok{ p}\OperatorTok{;}
\DataTypeTok{int} \OperatorTok{*}\DataTypeTok{volatile}\NormalTok{ p}\OperatorTok{;}
\DataTypeTok{int} \OperatorTok{*}\DataTypeTok{const} \DataTypeTok{volatile}\NormalTok{ p}\OperatorTok{;}
\CommentTok{// dan seterusnya.}
\end{Highlighting}
\end{Shaded}

\end{tcolorbox}

Tapi bagaimana melakukannya kalau kita memakai notasi array di daftar parameter?

Ternyata qualifier itu ditaruh di dalam kurung siku. Dan kamu bisa menaruh jumlah opsional setelahnya. Ketiga baris berikut setara:

\begin{tcolorbox}[enhanced jigsaw,left=3pt,breakable,colback=white!97.2549!black,boxrule=0.2mm]

\begin{Shaded}
\begin{Highlighting}[]
\DataTypeTok{int}\NormalTok{ func}\OperatorTok{(}\DataTypeTok{int} \OperatorTok{*}\DataTypeTok{const} \DataTypeTok{volatile}\NormalTok{ p}\OperatorTok{)} \OperatorTok{\{...\}}
\DataTypeTok{int}\NormalTok{ func}\OperatorTok{(}\DataTypeTok{int}\NormalTok{ p}\OperatorTok{[}\DataTypeTok{const} \DataTypeTok{volatile}\OperatorTok{])} \OperatorTok{\{...\}}
\DataTypeTok{int}\NormalTok{ func}\OperatorTok{(}\DataTypeTok{int}\NormalTok{ p}\OperatorTok{[}\DataTypeTok{const} \DataTypeTok{volatile} \DecValTok{10}\OperatorTok{])} \OperatorTok{\{...\}}
\end{Highlighting}
\end{Shaded}

\end{tcolorbox}

Kalau kamu punya array multidimensi, type qualifier harus ditaruh di pasangan kurung siku pertama.

\index{Type qualifiers!arrays in parameter lists|)} \index{Arrays!type qualifiers in parameter lists|)}

\section{\texorpdfstring{\hl{\mbox{\texttt{static}}} untuk Array dalam Daftar Parameter}{static untuk Array dalam Daftar Parameter}}\label{static-untuk-array-dalam-daftar-parameter}

\index{Arrays!static in parameter lists@\texttt{static} in parameter lists|(}

Demikian pula, kata kunci static bisa dipakai pada array di daftar parameter.

Aku belum pernah melihat ini di alam liar. Ia \textbf{selalu} diikuti dimensi:

\begin{tcolorbox}[enhanced jigsaw,left=3pt,breakable,colback=white!97.2549!black,boxrule=0.2mm]

\begin{Shaded}
\begin{Highlighting}[]
\DataTypeTok{int}\NormalTok{ func}\OperatorTok{(}\DataTypeTok{int}\NormalTok{ p}\OperatorTok{[}\DataTypeTok{static} \DecValTok{4}\OperatorTok{])} \OperatorTok{\{...\}}
\end{Highlighting}
\end{Shaded}

\end{tcolorbox}

Artinya, pada contoh di atas, compiler akan menganggap bahwa array apa pun yang dioper ke fungsi berisi \emph{setidaknya} 4 elemen.

Selain itu adalah undefined behavior.

\begin{tcolorbox}[enhanced jigsaw,left=3pt,breakable,colback=white!97.2549!black,boxrule=0.2mm]

\begin{Shaded}
\begin{Highlighting}[]
\DataTypeTok{int}\NormalTok{ func}\OperatorTok{(}\DataTypeTok{int}\NormalTok{ p}\OperatorTok{[}\DataTypeTok{static} \DecValTok{4}\OperatorTok{])} \OperatorTok{\{...\}}

\DataTypeTok{int}\NormalTok{ main}\OperatorTok{(}\DataTypeTok{void}\OperatorTok{)}
\OperatorTok{\{}
    \DataTypeTok{int}\NormalTok{ a}\OperatorTok{[]} \OperatorTok{=} \OperatorTok{\{}\DecValTok{11}\OperatorTok{,} \DecValTok{22}\OperatorTok{,} \DecValTok{33}\OperatorTok{,} \DecValTok{44}\OperatorTok{\};}
    \DataTypeTok{int}\NormalTok{ b}\OperatorTok{[]} \OperatorTok{=} \OperatorTok{\{}\DecValTok{11}\OperatorTok{,} \DecValTok{22}\OperatorTok{,} \DecValTok{33}\OperatorTok{,} \DecValTok{44}\OperatorTok{,} \DecValTok{55}\OperatorTok{\};}
    \DataTypeTok{int}\NormalTok{ c}\OperatorTok{[]} \OperatorTok{=} \OperatorTok{\{}\DecValTok{11}\OperatorTok{,} \DecValTok{22}\OperatorTok{\};}

\NormalTok{    func}\OperatorTok{(}\NormalTok{a}\OperatorTok{);}  \CommentTok{// OK! a berisi 4 elemen, batas minimum}
\NormalTok{    func}\OperatorTok{(}\NormalTok{b}\OperatorTok{);}  \CommentTok{// OK! b berisi setidaknya 4 elemen}
\NormalTok{    func}\OperatorTok{(}\NormalTok{c}\OperatorTok{);}  \CommentTok{// Undefined behavior! c berisi kurang dari 4 elemen!}
\OperatorTok{\}}
\end{Highlighting}
\end{Shaded}

\end{tcolorbox}

Ini pada dasarnya menetapkan ukuran minimum array yang boleh dioper.

Catatan penting: tidak ada apa pun di compiler yang melarangmu mengoper array yang lebih kecil. Compiler kemungkinan tidak akan memperingatkan, dan tidak akan mendeteksinya saat runtime.

Dengan menaruh \hl{\mbox{\texttt{static}}} di sana, kamu berkata, ``Aku JANJI dengan sumpah rahasia ganda bahwa aku tidak akan pernah mengoper array yang lebih kecil dari ini.'' Dan compiler menjawab, ``Ya, baiklah,'' dan memercayaimu untuk tidak melakukannya.

Lalu compiler bisa melakukan optimasi kode tertentu, dengan tenang karena tahu bahwa kamu, sang programmer, akan selalu melakukan hal yang benar.

\index{Arrays!static in parameter lists@\texttt{static} in parameter lists|)}

\section{Initializer yang Setara}\label{initializer-yang-setara}

\index{Arrays!multidimensional initializers|(}

C itu, bolehlah dibilang, agak \emph{fleksibel} soal array initializer.

Kita sudah melihat sebagian, yaitu nilai yang tidak disebutkan diganti dengan nol.

Misalnya, kita bisa menginisialisasi array 5 elemen menjadi \hl{\mbox{\texttt{1,2,0,0,0}}} dengan ini:

\begin{tcolorbox}[enhanced jigsaw,left=3pt,breakable,colback=white!97.2549!black,boxrule=0.2mm]

\begin{Shaded}
\begin{Highlighting}[]
\DataTypeTok{int}\NormalTok{ a}\OperatorTok{[}\DecValTok{5}\OperatorTok{]} \OperatorTok{=} \OperatorTok{\{}\DecValTok{1}\OperatorTok{,} \DecValTok{2}\OperatorTok{\};}
\end{Highlighting}
\end{Shaded}

\end{tcolorbox}

Atau mengisi seluruh array dengan nol memakai:

\begin{tcolorbox}[enhanced jigsaw,left=3pt,breakable,colback=white!97.2549!black,boxrule=0.2mm]

\begin{Shaded}
\begin{Highlighting}[]
\DataTypeTok{int}\NormalTok{ a}\OperatorTok{[}\DecValTok{5}\OperatorTok{]} \OperatorTok{=} \OperatorTok{\{}\DecValTok{0}\OperatorTok{\};}
\end{Highlighting}
\end{Shaded}

\end{tcolorbox}

Tapi segalanya jadi menarik saat menginisialisasi array multidimensi.

Kita buat array dengan 3 baris dan 2 kolom:

\begin{tcolorbox}[enhanced jigsaw,left=3pt,breakable,colback=white!97.2549!black,boxrule=0.2mm]

\begin{Shaded}
\begin{Highlighting}[]
\DataTypeTok{int}\NormalTok{ a}\OperatorTok{[}\DecValTok{3}\OperatorTok{][}\DecValTok{2}\OperatorTok{];}
\end{Highlighting}
\end{Shaded}

\end{tcolorbox}

Kita tulis kode untuk menginisialisasinya dan mencetak hasilnya:

\begin{tcolorbox}[enhanced jigsaw,left=3pt,breakable,colback=white!97.2549!black,boxrule=0.2mm]

\begin{Shaded}
\begin{Highlighting}[]
\PreprocessorTok{\#include }\ImportTok{\textless{}stdio.h\textgreater{}}

\DataTypeTok{int}\NormalTok{ main}\OperatorTok{(}\DataTypeTok{void}\OperatorTok{)}
\OperatorTok{\{}
    \DataTypeTok{int}\NormalTok{ a}\OperatorTok{[}\DecValTok{3}\OperatorTok{][}\DecValTok{2}\OperatorTok{]} \OperatorTok{=} \OperatorTok{\{}
        \OperatorTok{\{}\DecValTok{1}\OperatorTok{,} \DecValTok{2}\OperatorTok{\},}
        \OperatorTok{\{}\DecValTok{3}\OperatorTok{,} \DecValTok{4}\OperatorTok{\},}
        \OperatorTok{\{}\DecValTok{5}\OperatorTok{,} \DecValTok{6}\OperatorTok{\}}
    \OperatorTok{\};}

    \ControlFlowTok{for} \OperatorTok{(}\DataTypeTok{int}\NormalTok{ row }\OperatorTok{=} \DecValTok{0}\OperatorTok{;}\NormalTok{ row }\OperatorTok{\textless{}} \DecValTok{3}\OperatorTok{;}\NormalTok{ row}\OperatorTok{++)} \OperatorTok{\{}
        \ControlFlowTok{for} \OperatorTok{(}\DataTypeTok{int}\NormalTok{ col }\OperatorTok{=} \DecValTok{0}\OperatorTok{;}\NormalTok{ col }\OperatorTok{\textless{}} \DecValTok{2}\OperatorTok{;}\NormalTok{ col}\OperatorTok{++)}
\NormalTok{            printf}\OperatorTok{(}\StringTok{"}\SpecialCharTok{\%d}\StringTok{ "}\OperatorTok{,}\NormalTok{ a}\OperatorTok{[}\NormalTok{row}\OperatorTok{][}\NormalTok{col}\OperatorTok{]);}
\NormalTok{        printf}\OperatorTok{(}\StringTok{"}\SpecialCharTok{\textbackslash{}n}\StringTok{"}\OperatorTok{);}
    \OperatorTok{\}}
\OperatorTok{\}}
\end{Highlighting}
\end{Shaded}

\end{tcolorbox}

Dan saat dijalankan, kita mendapat hasil yang diharapkan:

\begin{tcolorbox}[enhanced jigsaw,left=3pt,breakable,colback=white!97.2549!black,boxrule=0.2mm]

\begin{Shaded}
\begin{Highlighting}[]
\NormalTok{1 2}
\NormalTok{3 4}
\NormalTok{5 6}
\end{Highlighting}
\end{Shaded}

\end{tcolorbox}

Mari kita hilangkan beberapa elemen initializer dan lihat bahwa mereka diisi nol:

\begin{tcolorbox}[enhanced jigsaw,left=3pt,breakable,colback=white!97.2549!black,boxrule=0.2mm]

\begin{Shaded}
\begin{Highlighting}[]
    \DataTypeTok{int}\NormalTok{ a}\OperatorTok{[}\DecValTok{3}\OperatorTok{][}\DecValTok{2}\OperatorTok{]} \OperatorTok{=} \OperatorTok{\{}
        \OperatorTok{\{}\DecValTok{1}\OperatorTok{,} \DecValTok{2}\OperatorTok{\},}
        \OperatorTok{\{}\DecValTok{3}\OperatorTok{\},}    \CommentTok{// Angka 4 dihilangkan!}
        \OperatorTok{\{}\DecValTok{5}\OperatorTok{,} \DecValTok{6}\OperatorTok{\}}
    \OperatorTok{\};}
\end{Highlighting}
\end{Shaded}

\end{tcolorbox}

yang menghasilkan:

\begin{tcolorbox}[enhanced jigsaw,left=3pt,breakable,colback=white!97.2549!black,boxrule=0.2mm]

\begin{Shaded}
\begin{Highlighting}[]
\NormalTok{1 2}
\NormalTok{3 0}
\NormalTok{5 6}
\end{Highlighting}
\end{Shaded}

\end{tcolorbox}

Sekarang kita hilangkan seluruh elemen tengah:

\begin{tcolorbox}[enhanced jigsaw,left=3pt,breakable,colback=white!97.2549!black,boxrule=0.2mm]

\begin{Shaded}
\begin{Highlighting}[]
    \DataTypeTok{int}\NormalTok{ a}\OperatorTok{[}\DecValTok{3}\OperatorTok{][}\DecValTok{2}\OperatorTok{]} \OperatorTok{=} \OperatorTok{\{}
        \OperatorTok{\{}\DecValTok{1}\OperatorTok{,} \DecValTok{2}\OperatorTok{\},}
        \CommentTok{// \{3, 4\},   // Potong saja seluruh bagian ini}
        \OperatorTok{\{}\DecValTok{5}\OperatorTok{,} \DecValTok{6}\OperatorTok{\}}
    \OperatorTok{\};}
\end{Highlighting}
\end{Shaded}

\end{tcolorbox}

Dan sekarang kita mendapat ini, yang mungkin bukan yang kamu harapkan:

\begin{tcolorbox}[enhanced jigsaw,left=3pt,breakable,colback=white!97.2549!black,boxrule=0.2mm]

\begin{Shaded}
\begin{Highlighting}[]
\NormalTok{1 2}
\NormalTok{5 6}
\NormalTok{0 0}
\end{Highlighting}
\end{Shaded}

\end{tcolorbox}

Tapi kalau dipikir-pikir, kita hanya memberi initializer yang cukup untuk dua baris, jadi initializer itu dipakai untuk dua baris pertama. Dan elemen sisanya diinisialisasi dengan nol.

Sejauh ini bagus. Secara umum, kalau kita menghilangkan sebagian initializer, compiler mengisi elemen yang bersesuaian dengan \hl{\mbox{\texttt{0}}}.

Tapi sekarang mari kita menggila.

\begin{tcolorbox}[enhanced jigsaw,left=3pt,breakable,colback=white!97.2549!black,boxrule=0.2mm]

\begin{Shaded}
\begin{Highlighting}[]
    \DataTypeTok{int}\NormalTok{ a}\OperatorTok{[}\DecValTok{3}\OperatorTok{][}\DecValTok{2}\OperatorTok{]} \OperatorTok{=} \OperatorTok{\{} \DecValTok{1}\OperatorTok{,} \DecValTok{2}\OperatorTok{,} \DecValTok{3}\OperatorTok{,} \DecValTok{4}\OperatorTok{,} \DecValTok{5}\OperatorTok{,} \DecValTok{6} \OperatorTok{\};}
\end{Highlighting}
\end{Shaded}

\end{tcolorbox}

Apa---? Itu array 2D, tapi initializer-nya hanya 1D!

Ternyata itu legal (walaupun GCC akan memperingatkannya kalau peringatan yang tepat diaktifkan).

Pada dasarnya, yang terjadi: pengisian elemen dimulai di baris 0, lalu baris 1, lalu baris 2, dari kiri ke kanan.

Jadi saat dicetak, urutannya:

\begin{tcolorbox}[enhanced jigsaw,left=3pt,breakable,colback=white!97.2549!black,boxrule=0.2mm]

\begin{Shaded}
\begin{Highlighting}[]
\NormalTok{1 2}
\NormalTok{3 4}
\NormalTok{5 6}
\end{Highlighting}
\end{Shaded}

\end{tcolorbox}

Kalau sebagian dihilangkan:

\begin{tcolorbox}[enhanced jigsaw,left=3pt,breakable,colback=white!97.2549!black,boxrule=0.2mm]

\begin{Shaded}
\begin{Highlighting}[]
    \DataTypeTok{int}\NormalTok{ a}\OperatorTok{[}\DecValTok{3}\OperatorTok{][}\DecValTok{2}\OperatorTok{]} \OperatorTok{=} \OperatorTok{\{} \DecValTok{1}\OperatorTok{,} \DecValTok{2}\OperatorTok{,} \DecValTok{3} \OperatorTok{\};}
\end{Highlighting}
\end{Shaded}

\end{tcolorbox}

sisanya diisi \hl{\mbox{\texttt{0}}}:

\begin{tcolorbox}[enhanced jigsaw,left=3pt,breakable,colback=white!97.2549!black,boxrule=0.2mm]

\begin{Shaded}
\begin{Highlighting}[]
\NormalTok{1 2}
\NormalTok{3 0}
\NormalTok{0 0}
\end{Highlighting}
\end{Shaded}

\end{tcolorbox}

Jadi kalau ingin mengisi seluruh array dengan \hl{\mbox{\texttt{0}}}, silakan:

\begin{tcolorbox}[enhanced jigsaw,left=3pt,breakable,colback=white!97.2549!black,boxrule=0.2mm]

\begin{Shaded}
\begin{Highlighting}[]
    \DataTypeTok{int}\NormalTok{ a}\OperatorTok{[}\DecValTok{3}\OperatorTok{][}\DecValTok{2}\OperatorTok{]} \OperatorTok{=} \OperatorTok{\{}\DecValTok{0}\OperatorTok{\};}
\end{Highlighting}
\end{Shaded}

\end{tcolorbox}

Tapi rekomendasiku: kalau punya array 2D, pakai initializer 2D. Kode jadi lebih mudah dibaca. (Kecuali untuk menginisialisasi seluruh array dengan \hl{\mbox{\texttt{0}}}, yang idiomatis memakai \hl{\mbox{\texttt{\{0\}}}} berapa pun dimensi arraynya.)

\index{Arrays!multidimensional initializers|)}

\chapter{\texorpdfstring{Long Jump dengan \hl{\mbox{\texttt{setjmp}}}, \hl{\mbox{\texttt{longjmp}}}}{Long Jump dengan setjmp, longjmp}}\label{setjmp-longjmp}

\index{Long jumps|(}

Kita sudah melihat \hl{\mbox{\texttt{goto}}}, yang melompat dalam function scope. Tapi \hl{\mbox{\texttt{longjmp()}}} memungkinkanmu melompat kembali ke titik eksekusi yang lebih awal, kembali ke fungsi yang memanggil fungsi ini.

Ada banyak batasan dan peringatan, tapi ini bisa menjadi fungsi yang berguna untuk keluar dari kedalaman call stack kembali ke keadaan yang lebih awal.

Menurut pengalamanku, fungsionalitas ini sangat jarang dipakai.

\section{\texorpdfstring{Memakai \hl{\mbox{\texttt{setjmp}}} dan \hl{\mbox{\texttt{longjmp}}}}{Memakai setjmp dan longjmp}}\label{memakai-setjmp-dan-longjmp}

\index{setjmp() function@\texttt{setjmp()} function|(} \index{longjmp() function@\texttt{longjmp()} function|(}

Tarian yang akan kita lakukan di sini pada dasarnya adalah menaruh penanda buku dalam eksekusi dengan \hl{\mbox{\texttt{setjmp()}}}. Nanti, kita memanggil \hl{\mbox{\texttt{longjmp()}}} dan ia akan melompat kembali ke titik eksekusi sebelumnya tempat kita menaruh penanda itu dengan \hl{\mbox{\texttt{setjmp()}}}.

Dan ini bisa dilakukan bahkan kalau kamu sudah memanggil subfungsi.

Berikut demo singkat ketika kita memanggil fungsi beberapa tingkat ke dalam lalu keluar darurat.

Kita akan memakai variabel file scope \hl{\mbox{\texttt{env}}} untuk menyimpan \emph{state} saat kita memanggil \hl{\mbox{\texttt{setjmp()}}} sehingga bisa memulihkannya ketika memanggil \hl{\mbox{\texttt{longjmp()}}} nanti. Inilah variabel tempat kita mengingat ``posisi'' kita.

Variabel \hl{\mbox{\texttt{env}}} bertipe \hl{\mbox{\texttt{jmp\_buf}}}, tipe opaque yang dideklarasikan di \hl{\mbox{\texttt{\textless{}setjmp.h\textgreater{}}}}.

\begin{tcolorbox}[enhanced jigsaw,left=3pt,breakable,colback=white!97.2549!black,boxrule=0.2mm]

\begin{Shaded}
\begin{Highlighting}[numbers=left,,]
\PreprocessorTok{\#include }\ImportTok{\textless{}stdio.h\textgreater{}}
\PreprocessorTok{\#include }\ImportTok{\textless{}setjmp.h\textgreater{}}

\NormalTok{jmp\_buf env}\OperatorTok{;}

\DataTypeTok{void}\NormalTok{ depth2}\OperatorTok{(}\DataTypeTok{void}\OperatorTok{)}
\OperatorTok{\{}
\NormalTok{    printf}\OperatorTok{(}\StringTok{"Entering depth 2}\SpecialCharTok{\textbackslash{}n}\StringTok{"}\OperatorTok{);}
\NormalTok{    longjmp}\OperatorTok{(}\NormalTok{env}\OperatorTok{,} \DecValTok{3490}\OperatorTok{);}           \CommentTok{// Keluar darurat}
\NormalTok{    printf}\OperatorTok{(}\StringTok{"Leaving depth 2}\SpecialCharTok{\textbackslash{}n}\StringTok{"}\OperatorTok{);}  \CommentTok{// Ini tidak akan terjadi}
\OperatorTok{\}}

\DataTypeTok{void}\NormalTok{ depth1}\OperatorTok{(}\DataTypeTok{void}\OperatorTok{)}
\OperatorTok{\{}
\NormalTok{    printf}\OperatorTok{(}\StringTok{"Entering depth 1}\SpecialCharTok{\textbackslash{}n}\StringTok{"}\OperatorTok{);}
\NormalTok{    depth2}\OperatorTok{();}
\NormalTok{    printf}\OperatorTok{(}\StringTok{"Leaving depth 1}\SpecialCharTok{\textbackslash{}n}\StringTok{"}\OperatorTok{);}  \CommentTok{// Ini tidak akan terjadi}
\OperatorTok{\}}

\DataTypeTok{int}\NormalTok{ main}\OperatorTok{(}\DataTypeTok{void}\OperatorTok{)}
\OperatorTok{\{}
    \ControlFlowTok{switch} \OperatorTok{(}\NormalTok{setjmp}\OperatorTok{(}\NormalTok{env}\OperatorTok{))} \OperatorTok{\{}
      \ControlFlowTok{case} \DecValTok{0}\OperatorTok{:}
\NormalTok{          printf}\OperatorTok{(}\StringTok{"Calling into functions, setjmp() returned 0}\SpecialCharTok{\textbackslash{}n}\StringTok{"}\OperatorTok{);}
\NormalTok{          depth1}\OperatorTok{();}
\NormalTok{          printf}\OperatorTok{(}\StringTok{"Returned from functions}\SpecialCharTok{\textbackslash{}n}\StringTok{"}\OperatorTok{);}  \CommentTok{// Ini tidak akan terjadi}
          \ControlFlowTok{break}\OperatorTok{;}

      \ControlFlowTok{case} \DecValTok{3490}\OperatorTok{:}
\NormalTok{          printf}\OperatorTok{(}\StringTok{"Bailed back to main, setjmp() returned 3490}\SpecialCharTok{\textbackslash{}n}\StringTok{"}\OperatorTok{);}
          \ControlFlowTok{break}\OperatorTok{;}
    \OperatorTok{\}}
\OperatorTok{\}}
\end{Highlighting}
\end{Shaded}

\end{tcolorbox}

Saat dijalankan, ini mengeluarkan:

\begin{tcolorbox}[enhanced jigsaw,left=3pt,breakable,colback=white!97.2549!black,boxrule=0.2mm]

\begin{Shaded}
\begin{Highlighting}[]
\NormalTok{Calling into functions, setjmp() returned 0}
\NormalTok{Entering depth 1}
\NormalTok{Entering depth 2}
\NormalTok{Bailed back to main, setjmp() returned 3490}
\end{Highlighting}
\end{Shaded}

\end{tcolorbox}

Kalau kamu mencoba mencocokkan output itu dengan kodenya, jelas ada hal yang benar-benar \emph{aneh} sedang terjadi.

Salah satu yang paling mencolok: \hl{\mbox{\texttt{setjmp()}}} mengembalikan nilai \emph{dua kali}. Apa-apaan ini? Sihir apa ini?!

Begini ceritanya: kalau \hl{\mbox{\texttt{setjmp()}}} mengembalikan \hl{\mbox{\texttt{0}}}, artinya kamu berhasil menaruh ``penanda buku'' di titik itu.

Kalau ia mengembalikan nilai bukan nol, artinya kamu baru saja kembali ke ``penanda buku'' yang ditaruh sebelumnya. (Dan nilai yang dikembalikan adalah nilai yang kamu oper ke \hl{\mbox{\texttt{longjmp()}}}.)

Dengan cara ini kamu bisa membedakan antara menaruh penanda dan kembali ke penanda itu nanti.

Jadi ketika kode di atas memanggil \hl{\mbox{\texttt{setjmp()}}} untuk pertama kali, \hl{\mbox{\texttt{setjmp()}}} \emph{menyimpan} state di variabel \hl{\mbox{\texttt{env}}} dan mengembalikan \hl{\mbox{\texttt{0}}}. Nanti saat kita memanggil \hl{\mbox{\texttt{longjmp()}}} dengan \hl{\mbox{\texttt{env}}} yang sama, ia memulihkan state itu dan \hl{\mbox{\texttt{setjmp()}}} mengembalikan nilai yang dioper ke \hl{\mbox{\texttt{longjmp()}}}.

\index{setjmp() function@\texttt{setjmp()} function|)} \index{longjmp() function@\texttt{longjmp()} function|)}

\section{Jebakan}\label{jebakan}

Di balik layar, ini cukup lugas. Biasanya \emph{stack pointer} melacak lokasi di memori tempat variabel lokal disimpan, dan \emph{program counter} melacak alamat instruksi yang sedang dieksekusi\footnote{Baik ``stack pointer'' maupun ``program counter'' terkait dengan arsitektur dan implementasi C di bawahnya, dan bukan bagian dari spesifikasi.}.

Jadi kalau kita ingin melompat kembali ke fungsi yang lebih awal, pada dasarnya tinggal memulihkan stack pointer dan program counter ke nilai yang disimpan di variabel \index{jmp\_buf type@\texttt{jmp\_buf} type} \hl{\mbox{\texttt{jmp\_buf}}}, dan memastikan nilai kembaliannya diatur dengan benar. Lalu eksekusi akan dilanjutkan di sana.

Tapi berbagai faktor mengacaukan ini, menciptakan banyak jebakan undefined behavior.

\subsection{Nilai Variabel Lokal}\label{nilai-variabel-lokal}

\index{setjmp() function@\texttt{setjmp()} function|(}

Kalau kamu ingin nilai variabel lokal otomatis (bukan \hl{\mbox{\texttt{static}}} dan bukan \hl{\mbox{\texttt{extern}}}) tetap bertahan di fungsi yang memanggil \hl{\mbox{\texttt{setjmp()}}} setelah \index{longjmp()@\texttt{longjmp()}} \hl{\mbox{\texttt{longjmp()}}} terjadi, kamu harus mendeklarasikan variabel itu sebagai \index{volatile type qualifier@\texttt{volatile} type qualifier!with setjmp()@with \texttt{setjmp()}|(} \hl{\mbox{\texttt{volatile}}}.

Secara teknis, mereka hanya perlu \hl{\mbox{\texttt{volatile}}} kalau nilainya berubah antara saat \hl{\mbox{\texttt{setjmp()}}} dipanggil dan \hl{\mbox{\texttt{longjmp()}}} dipanggil\footnote{Alasannya, program mungkin menyimpan nilai sementara di \emph{register CPU} saat mengerjakannya. Pada rentang waktu itu, register menyimpan nilai yang benar, dan nilai di stack mungkin sudah usang. Lalu nanti nilai register akan tertimpa dan perubahan pada variabel itu hilang.}.

Misalnya, kalau kita menjalankan kode ini:

\begin{tcolorbox}[enhanced jigsaw,left=3pt,breakable,colback=white!97.2549!black,boxrule=0.2mm]

\begin{Shaded}
\begin{Highlighting}[]
\DataTypeTok{int}\NormalTok{ x }\OperatorTok{=} \DecValTok{20}\OperatorTok{;}

\ControlFlowTok{if} \OperatorTok{(}\NormalTok{setjmp}\OperatorTok{(}\NormalTok{env}\OperatorTok{)} \OperatorTok{==} \DecValTok{0}\OperatorTok{)} \OperatorTok{\{}
\NormalTok{    x }\OperatorTok{=} \DecValTok{30}\OperatorTok{;}
\OperatorTok{\}}
\end{Highlighting}
\end{Shaded}

\end{tcolorbox}

dan nanti \hl{\mbox{\texttt{longjmp()}}} kembali, nilai \hl{\mbox{\texttt{x}}} akan tak tentu.

Kalau ingin memperbaikinya, \hl{\mbox{\texttt{x}}} harus \hl{\mbox{\texttt{volatile}}}:

\begin{tcolorbox}[enhanced jigsaw,left=3pt,breakable,colback=white!97.2549!black,boxrule=0.2mm]

\begin{Shaded}
\begin{Highlighting}[]
\DataTypeTok{volatile} \DataTypeTok{int}\NormalTok{ x }\OperatorTok{=} \DecValTok{20}\OperatorTok{;}

\ControlFlowTok{if} \OperatorTok{(}\NormalTok{setjmp}\OperatorTok{(}\NormalTok{env}\OperatorTok{)} \OperatorTok{==} \DecValTok{0}\OperatorTok{)} \OperatorTok{\{}
\NormalTok{    x }\OperatorTok{=} \DecValTok{30}\OperatorTok{;}
\OperatorTok{\}}
\end{Highlighting}
\end{Shaded}

\end{tcolorbox}

\index{setjmp() function@\texttt{setjmp()} function|)}

\index{volatile type qualifier@\texttt{volatile} type qualifier!with setjmp()@with \texttt{setjmp()}|)}

Sekarang nilainya akan benar, yaitu \hl{\mbox{\texttt{30}}}, setelah \index{longjmp()@\texttt{longjmp()}} \hl{\mbox{\texttt{longjmp()}}} membawa kita kembali ke titik ini.

\subsection{Berapa Banyak State yang Disimpan?}\label{berapa-banyak-state-yang-disimpan}

\index{setjmp() function@\texttt{setjmp()} function|(} \index{longjmp() function@\texttt{longjmp()} function|(}

Saat kamu \hl{\mbox{\texttt{longjmp()}}}, eksekusi dilanjutkan di titik \hl{\mbox{\texttt{setjmp()}}} yang bersesuaian. Dan hanya itu.

\index{setjmp() function@\texttt{setjmp()} function|)}

Spesifikasi menunjukkan bahwa itu sama saja seolah-olah kamu melompat kembali ke fungsi pada titik itu dengan variabel lokal bernilai apa pun yang mereka miliki saat pemanggilan \hl{\mbox{\texttt{longjmp()}}} dilakukan.

\index{longjmp() function@\texttt{longjmp()} function|)}

Hal-hal yang tidak dipulihkan, memparafrasekan spesifikasi, mencakup:

\begin{itemize}
\tightlist
\item
  Flag status floating point
\item
  File yang terbuka
\item
  Komponen lain apa pun dari abstract machine
\end{itemize}

\subsection{\texorpdfstring{Kamu Tidak Boleh Menamai Apa Pun \hl{\mbox{\texttt{setjmp}}}}{Kamu Tidak Boleh Menamai Apa Pun setjmp}}\label{kamu-tidak-boleh-menamai-apa-pun-setjmp}

Kamu tidak boleh punya identifier \hl{\mbox{\texttt{extern}}} bernama \hl{\mbox{\texttt{setjmp}}}. Atau, kalau \hl{\mbox{\texttt{setjmp}}} adalah makro, kamu tidak boleh meng-undefine-nya.

Keduanya undefined behavior.

\subsection{\texorpdfstring{\hl{\mbox{\texttt{setjmp()}}} Tidak Boleh Berada dalam Ekspresi yang Lebih Besar}{setjmp() Tidak Boleh Berada dalam Ekspresi yang Lebih Besar}}\label{setjmp-tidak-boleh-berada-dalam-ekspresi-yang-lebih-besar}

\index{setjmp()@\texttt{setjmp()}!in an expression|(}

Artinya, kamu tidak bisa melakukan sesuatu seperti ini:

\begin{tcolorbox}[enhanced jigsaw,left=3pt,breakable,colback=white!97.2549!black,boxrule=0.2mm]

\begin{Shaded}
\begin{Highlighting}[]
\ControlFlowTok{if} \OperatorTok{(}\NormalTok{x }\OperatorTok{==} \DecValTok{12} \OperatorTok{\&\&}\NormalTok{ setjmp}\OperatorTok{(}\NormalTok{env}\OperatorTok{)} \OperatorTok{==} \DecValTok{0}\OperatorTok{)} \OperatorTok{\{} \OperatorTok{...} \OperatorTok{\}}
\end{Highlighting}
\end{Shaded}

\end{tcolorbox}

Itu terlalu kompleks untuk diizinkan spesifikasi karena rekayasa yang harus terjadi saat membongkar stack dan sebagainya. Kita tidak bisa \hl{\mbox{\texttt{longjmp()}}} kembali ke ekspresi kompleks yang baru dieksekusi sebagian.

Jadi ada batas kompleksitas pada ekspresi itu.

\begin{itemize}
\item
  Ia boleh menjadi seluruh ekspresi pengendali dari kondisional.

  \begin{tcolorbox}[enhanced jigsaw,left=3pt,breakable,colback=white!97.2549!black,boxrule=0.2mm]

\begin{Shaded}
\begin{Highlighting}[]
\ControlFlowTok{if} \OperatorTok{(}\NormalTok{setjmp}\OperatorTok{(}\NormalTok{env}\OperatorTok{))} \OperatorTok{\{...\}}
\end{Highlighting}
\end{Shaded}

  \end{tcolorbox}

  \begin{tcolorbox}[enhanced jigsaw,left=3pt,breakable,colback=white!97.2549!black,boxrule=0.2mm]

\begin{Shaded}
\begin{Highlighting}[]
\ControlFlowTok{switch} \OperatorTok{(}\NormalTok{setjmp}\OperatorTok{(}\NormalTok{env}\OperatorTok{))} \OperatorTok{\{...\}}
\end{Highlighting}
\end{Shaded}

  \end{tcolorbox}
\item
  Ia boleh menjadi bagian dari ekspresi relasional atau kesamaan, selama operand lainnya adalah konstanta integer. Dan keseluruhannya adalah ekspresi pengendali dari kondisional.

  \begin{tcolorbox}[enhanced jigsaw,left=3pt,breakable,colback=white!97.2549!black,boxrule=0.2mm]

\begin{Shaded}
\begin{Highlighting}[]
\ControlFlowTok{if} \OperatorTok{(}\NormalTok{setjmp}\OperatorTok{(}\NormalTok{env}\OperatorTok{)} \OperatorTok{==} \DecValTok{0}\OperatorTok{)} \OperatorTok{\{...\}}
\end{Highlighting}
\end{Shaded}

  \end{tcolorbox}
\item
  Operand dari operasi logical NOT (\hl{\mbox{\texttt{!}}}), yang menjadi seluruh ekspresi pengendali.

  \begin{tcolorbox}[enhanced jigsaw,left=3pt,breakable,colback=white!97.2549!black,boxrule=0.2mm]

\begin{Shaded}
\begin{Highlighting}[]
\ControlFlowTok{if} \OperatorTok{(!}\NormalTok{setjmp}\OperatorTok{(}\NormalTok{env}\OperatorTok{))} \OperatorTok{\{...\}}
\end{Highlighting}
\end{Shaded}

  \end{tcolorbox}
\item
  Ekspresi mandiri, mungkin di-cast ke \hl{\mbox{\texttt{void}}}.

  \begin{tcolorbox}[enhanced jigsaw,left=3pt,breakable,colback=white!97.2549!black,boxrule=0.2mm]

\begin{Shaded}
\begin{Highlighting}[]
\NormalTok{setjmp}\OperatorTok{(}\NormalTok{env}\OperatorTok{);}
\end{Highlighting}
\end{Shaded}

  \end{tcolorbox}

  \begin{tcolorbox}[enhanced jigsaw,left=3pt,breakable,colback=white!97.2549!black,boxrule=0.2mm]

\begin{Shaded}
\begin{Highlighting}[]
\OperatorTok{(}\DataTypeTok{void}\OperatorTok{)}\NormalTok{setjmp}\OperatorTok{(}\NormalTok{env}\OperatorTok{);}
\end{Highlighting}
\end{Shaded}

  \end{tcolorbox}
\end{itemize}

\index{setjmp()@\texttt{setjmp()}!in an expression|)}

\subsection{\texorpdfstring{Kapan \hl{\mbox{\texttt{longjmp()}}} Tidak Boleh Dilakukan?}{Kapan longjmp() Tidak Boleh Dilakukan?}}\label{kapan-longjmp-tidak-boleh-dilakukan}

\index{lonjmp()@\texttt{lonjmp()}|(}

Ini undefined behavior kalau:

\begin{itemize}
\tightlist
\item
  Kamu belum memanggil \hl{\mbox{\texttt{setjmp()}}} sebelumnya
\item
  Kamu memanggil \hl{\mbox{\texttt{setjmp()}}} dari thread lain
\item
  Kamu memanggil \hl{\mbox{\texttt{setjmp()}}} di scope sebuah variable length array (VLA), dan eksekusi sudah meninggalkan scope VLA itu sebelum \hl{\mbox{\texttt{longjmp()}}} dipanggil.
\item
  Fungsi yang berisi \hl{\mbox{\texttt{setjmp()}}} sudah keluar sebelum \hl{\mbox{\texttt{longjmp()}}} dipanggil.
\end{itemize}

Untuk yang terakhir, ``keluar'' mencakup return normal dari fungsi itu, serta kasus ketika \hl{\mbox{\texttt{longjmp()}}} lain sudah melompat kembali ke titik yang ``lebih awal'' di call stack daripada fungsi yang dimaksud.

\subsection{\texorpdfstring{Kamu Tidak Boleh Mengoper \hl{\mbox{\texttt{0}}} ke \hl{\mbox{\texttt{longjmp()}}}}{Kamu Tidak Boleh Mengoper 0 ke longjmp()}}\label{kamu-tidak-boleh-mengoper-0-ke-longjmp}

Kalau kamu mencoba mengoper nilai \hl{\mbox{\texttt{0}}} ke \hl{\mbox{\texttt{longjmp()}}}, nilai itu akan diam-diam diubah menjadi \hl{\mbox{\texttt{1}}}.

Karena \hl{\mbox{\texttt{setjmp()}}} pada akhirnya mengembalikan nilai ini, dan \hl{\mbox{\texttt{setjmp()}}} yang mengembalikan \hl{\mbox{\texttt{0}}} punya makna khusus, mengembalikan \hl{\mbox{\texttt{0}}} dilarang.

\subsection{\texorpdfstring{\hl{\mbox{\texttt{longjmp()}}} dan Variable Length Array}{longjmp() dan Variable Length Array}}\label{longjmp-dan-variable-length-array}

Kalau kamu berada di scope sebuah VLA dan \hl{\mbox{\texttt{longjmp()}}} keluar dari sana, memori yang dialokasikan untuk VLA itu bisa bocor\footnote{Artinya, tetap teralokasi sampai program berakhir tanpa cara untuk membebaskannya.}.

Hal yang sama terjadi kalau kamu \hl{\mbox{\texttt{longjmp()}}} melewati fungsi-fungsi sebelumnya yang masih punya VLA dalam scope.

Ini salah satu hal yang benar-benar mengganggu aku soal VLA---kamu bisa menulis kode C yang sepenuhnya sah tapi menyia-nyiakan memori. Tapi ya, aku bukan pemegang spesifikasi.

\index{lonjmp()@\texttt{lonjmp()}|)} \index{Long jumps|)}

\chapter{Incomplete Type}\label{incomplete-type}

\index{Incomplete types|(}

Mungkin mengejutkan bahwa kode ini dibangun tanpa error:

\begin{tcolorbox}[enhanced jigsaw,left=3pt,breakable,colback=white!97.2549!black,boxrule=0.2mm]

\begin{Shaded}
\begin{Highlighting}[]
\KeywordTok{extern} \DataTypeTok{int}\NormalTok{ a}\OperatorTok{[];}

\DataTypeTok{int}\NormalTok{ main}\OperatorTok{(}\DataTypeTok{void}\OperatorTok{)}
\OperatorTok{\{}
    \KeywordTok{struct}\NormalTok{ foo }\OperatorTok{*}\NormalTok{x}\OperatorTok{;}
    \KeywordTok{union}\NormalTok{ bar }\OperatorTok{*}\NormalTok{y}\OperatorTok{;}
    \KeywordTok{enum}\NormalTok{ baz }\OperatorTok{*}\NormalTok{z}\OperatorTok{;}
\OperatorTok{\}}
\end{Highlighting}
\end{Shaded}

\end{tcolorbox}

Kita tidak pernah memberi ukuran untuk \hl{\mbox{\texttt{a}}}. Dan kita punya pointer ke \hl{\mbox{\texttt{struct}}} \hl{\mbox{\texttt{foo}}}, \hl{\mbox{\texttt{bar}}}, dan \hl{\mbox{\texttt{baz}}} yang tampaknya tidak dideklarasikan di mana pun.

Dan satu-satunya peringatan yang kudapat hanyalah bahwa \hl{\mbox{\texttt{x}}}, \hl{\mbox{\texttt{y}}}, dan \hl{\mbox{\texttt{z}}} tidak dipakai.

Inilah contoh \emph{incomplete type}.

Incomplete type adalah tipe yang ukurannya (yaitu ukuran yang akan kamu dapat dari \hl{\mbox{\texttt{sizeof}}}) tidak diketahui. Cara lain memikirkannya: tipe yang belum selesai kamu deklarasikan.

Kamu bisa punya pointer ke incomplete type, tapi tidak bisa men-dereference-nya atau memakai aritmetika pointer padanya. Dan kamu tidak bisa meng-\hl{\mbox{\texttt{sizeof}}}-nya.

Jadi apa yang bisa dilakukan dengannya?

\section{Kasus Pemakaian: Struktur yang Merujuk Diri Sendiri}\label{kasus-pemakaian-struktur-yang-merujuk-diri-sendiri}

\index{Incomplete types!self-referential structs@self-referential \texttt{struct}s|(}

Aku hanya tahu satu kasus pemakaian yang nyata: forward reference ke \hl{\mbox{\texttt{struct}}} atau \hl{\mbox{\texttt{union}}} dengan struktur yang merujuk diri sendiri atau saling bergantung. (Untuk sisa contoh, aku akan memakai \hl{\mbox{\texttt{struct}}}, tapi semuanya berlaku sama untuk \hl{\mbox{\texttt{union}}}.)

Kita kerjakan contoh klasiknya dulu.

Tapi sebelum aku mulai, ketahuilah ini! Saat kamu mendeklarasikan sebuah \hl{\mbox{\texttt{struct}}}, \hl{\mbox{\texttt{struct}}} itu incomplete sampai kurung kurawal penutup tercapai!

\begin{tcolorbox}[enhanced jigsaw,left=3pt,breakable,colback=white!97.2549!black,boxrule=0.2mm]

\begin{Shaded}
\begin{Highlighting}[]
\KeywordTok{struct}\NormalTok{ antelope }\OperatorTok{\{}              \CommentTok{// struct antelope incomplete di sini}
    \DataTypeTok{int}\NormalTok{ leg\_count}\OperatorTok{;}             \CommentTok{// Masih incomplete}
    \DataTypeTok{float}\NormalTok{ stomach\_fullness}\OperatorTok{;}    \CommentTok{// Masih incomplete}
    \DataTypeTok{float}\NormalTok{ top\_speed}\OperatorTok{;}           \CommentTok{// Masih incomplete}
    \DataTypeTok{char} \OperatorTok{*}\NormalTok{nickname}\OperatorTok{;}            \CommentTok{// Masih incomplete}
\OperatorTok{\};}                             \CommentTok{// SEKARANG sudah complete.}
\end{Highlighting}
\end{Shaded}

\end{tcolorbox}

Lalu kenapa? Kelihatannya cukup waras.

Tapi bagaimana kalau kita membuat linked list? Setiap node linked list perlu punya referensi ke node lain. Tapi bagaimana kita membuat referensi ke node lain kalau deklarasi node saja belum selesai?

Kelonggaran C untuk incomplete type memungkinkannya. Kita tidak bisa mendeklarasikan sebuah node, tapi \emph{bisa} mendeklarasikan pointer ke node, walaupun ia incomplete!

\begin{tcolorbox}[enhanced jigsaw,left=3pt,breakable,colback=white!97.2549!black,boxrule=0.2mm]

\begin{Shaded}
\begin{Highlighting}[]
\KeywordTok{struct}\NormalTok{ node }\OperatorTok{\{}
    \DataTypeTok{int}\NormalTok{ val}\OperatorTok{;}
    \KeywordTok{struct}\NormalTok{ node }\OperatorTok{*}\NormalTok{next}\OperatorTok{;}  \CommentTok{// struct node incomplete, tapi tidak masalah!}
\OperatorTok{\};}
\end{Highlighting}
\end{Shaded}

\end{tcolorbox}

Meskipun \hl{\mbox{\texttt{struct\ node}}} incomplete di baris 3, kita tetap bisa mendeklarasikan pointer ke sana\footnote{Ini bisa karena di C, pointer berukuran sama apa pun tipe data yang ditunjuknya. Jadi compiler tidak perlu tahu ukuran \hl{\mbox{\texttt{struct\ node}}} pada titik ini; ia hanya perlu tahu ukuran pointer.}.

Hal yang sama bisa dilakukan kalau kita punya dua \hl{\mbox{\texttt{struct}}} berbeda yang saling merujuk:

\begin{tcolorbox}[enhanced jigsaw,left=3pt,breakable,colback=white!97.2549!black,boxrule=0.2mm]

\begin{Shaded}
\begin{Highlighting}[]
\KeywordTok{struct}\NormalTok{ a }\OperatorTok{\{}
    \KeywordTok{struct}\NormalTok{ b }\OperatorTok{*}\NormalTok{x}\OperatorTok{;}  \CommentTok{// Merujuk ke \textasciigrave{}struct b\textasciigrave{}}
\OperatorTok{\};}

\KeywordTok{struct}\NormalTok{ b }\OperatorTok{\{}
    \KeywordTok{struct}\NormalTok{ a }\OperatorTok{*}\NormalTok{x}\OperatorTok{;}  \CommentTok{// Merujuk ke \textasciigrave{}struct a\textasciigrave{}}
\OperatorTok{\};}
\end{Highlighting}
\end{Shaded}

\end{tcolorbox}

Kita tidak akan pernah bisa membuat pasangan struktur itu tanpa aturan yang dilonggarkan untuk incomplete type.

\index{Incomplete types!self-referential structs@self-referential \texttt{struct}s|)}

\section{Pesan Error Incomplete Type}\label{pesan-error-incomplete-type}

Apakah kamu mendapat error seperti ini?

\begin{tcolorbox}[enhanced jigsaw,left=3pt,breakable,colback=white!97.2549!black,boxrule=0.2mm]

\begin{Shaded}
\begin{Highlighting}[]
\NormalTok{invalid application of ‘sizeof’ to incomplete type}

\NormalTok{invalid use of undefined type}

\NormalTok{dereferencing pointer to incomplete type}
\end{Highlighting}
\end{Shaded}

\end{tcolorbox}

Tersangka paling mungkin: kamu kemungkinan lupa meng-\hl{\mbox{\texttt{\#include}}} header file yang mendeklarasikan tipe itu.

\section{Incomplete Type Lainnya}\label{incomplete-type-lainnya}

Mendeklarasikan \hl{\mbox{\texttt{struct}}} atau \hl{\mbox{\texttt{union}}} tanpa badan membuat incomplete type, mis. \hl{\mbox{\texttt{struct\ foo;}}}.

\hl{\mbox{\texttt{enum}}} incomplete sampai kurung kurawal penutup.

\hl{\mbox{\texttt{void}}} adalah incomplete type.

Array yang dideklarasikan \hl{\mbox{\texttt{extern}}} tanpa ukuran adalah incomplete, mis.:

\begin{tcolorbox}[enhanced jigsaw,left=3pt,breakable,colback=white!97.2549!black,boxrule=0.2mm]

\begin{Shaded}
\begin{Highlighting}[]
\KeywordTok{extern} \DataTypeTok{int}\NormalTok{ a}\OperatorTok{[];}
\end{Highlighting}
\end{Shaded}

\end{tcolorbox}

Kalau ia array non-\hl{\mbox{\texttt{extern}}} tanpa ukuran yang diikuti initializer, ia incomplete sampai kurung kurawal penutup initializer.

\section{Kasus Pemakaian: Array di Header File}\label{kasus-pemakaian-array-di-header-file}

Berguna mendeklarasikan tipe array incomplete di header file. Dalam kasus itu, penyimpanan sebenarnya (tempat array lengkap dideklarasikan) harus ada di satu file \hl{\mbox{\texttt{.c}}}. Kalau ditaruh di file \hl{\mbox{\texttt{.h}}}, ia akan diduplikasi setiap kali header file di-include.

Jadi yang bisa dilakukan adalah membuat header file dengan incomplete type yang merujuk ke array, seperti ini:

\begin{tcolorbox}[enhanced jigsaw,left=3pt,breakable,colback=white!97.2549!black,boxrule=0.2mm]

\begin{Shaded}
\begin{Highlighting}[numbers=left,,]
\CommentTok{// File: bar.h}

\PreprocessorTok{\#ifndef BAR\_H}
\PreprocessorTok{\#define BAR\_H}

\KeywordTok{extern} \DataTypeTok{int}\NormalTok{ my\_array}\OperatorTok{[];}  \CommentTok{// Incomplete type}

\PreprocessorTok{\#endif}
\end{Highlighting}
\end{Shaded}

\end{tcolorbox}

Lalu di file \hl{\mbox{\texttt{.c}}}, definisikan array-nya yang sebenarnya:

\begin{tcolorbox}[enhanced jigsaw,left=3pt,breakable,colback=white!97.2549!black,boxrule=0.2mm]

\begin{Shaded}
\begin{Highlighting}[numbers=left,,]
\CommentTok{// File: bar.c}

\DataTypeTok{int}\NormalTok{ my\_array}\OperatorTok{[}\DecValTok{1024}\OperatorTok{];}     \CommentTok{// Complete type!}
\end{Highlighting}
\end{Shaded}

\end{tcolorbox}

Setelah itu kamu bisa meng-include header itu dari sebanyak mungkin tempat yang kamu mau, dan setiap tempat itu akan merujuk ke \hl{\mbox{\texttt{my\_array}}} yang sama.

\begin{tcolorbox}[enhanced jigsaw,left=3pt,breakable,colback=white!97.2549!black,boxrule=0.2mm]

\begin{Shaded}
\begin{Highlighting}[numbers=left,,]
\CommentTok{// File: foo.c}

\PreprocessorTok{\#include }\ImportTok{\textless{}stdio.h\textgreater{}}
\PreprocessorTok{\#include }\ImportTok{"bar.h"}\PreprocessorTok{    }\CommentTok{// menyertakan incomplete type untuk my\_array}

\DataTypeTok{int}\NormalTok{ main}\OperatorTok{(}\DataTypeTok{void}\OperatorTok{)}
\OperatorTok{\{}
\NormalTok{    my\_array}\OperatorTok{[}\DecValTok{0}\OperatorTok{]} \OperatorTok{=} \DecValTok{10}\OperatorTok{;}

\NormalTok{    printf}\OperatorTok{(}\StringTok{"}\SpecialCharTok{\%d\textbackslash{}n}\StringTok{"}\OperatorTok{,}\NormalTok{ my\_array}\OperatorTok{[}\DecValTok{0}\OperatorTok{]);}
\OperatorTok{\}}
\end{Highlighting}
\end{Shaded}

\end{tcolorbox}

Saat mengompilasi banyak file, ingat untuk menyebutkan semua file \hl{\mbox{\texttt{.c}}} ke compiler, tapi bukan file \hl{\mbox{\texttt{.h}}}, mis.:

\begin{tcolorbox}[enhanced jigsaw,left=3pt,breakable,colback=white!97.2549!black,boxrule=0.2mm]

\begin{Shaded}
\begin{Highlighting}[]
\FunctionTok{gcc} \AttributeTok{{-}o}\NormalTok{ foo foo.c bar.c}
\end{Highlighting}
\end{Shaded}

\end{tcolorbox}

\section{Melengkapi Incomplete Type}\label{melengkapi-incomplete-type}

Kalau kamu punya incomplete type, kamu bisa melengkapinya dengan mendefinisikan \hl{\mbox{\texttt{struct}}}, \hl{\mbox{\texttt{union}}}, \hl{\mbox{\texttt{enum}}}, atau array yang lengkap di scope yang sama.

\begin{tcolorbox}[enhanced jigsaw,left=3pt,breakable,colback=white!97.2549!black,boxrule=0.2mm]

\begin{Shaded}
\begin{Highlighting}[]
\KeywordTok{struct}\NormalTok{ foo}\OperatorTok{;}        \CommentTok{// incomplete type}

\KeywordTok{struct}\NormalTok{ foo }\OperatorTok{*}\NormalTok{p}\OperatorTok{;}     \CommentTok{// pointer, tidak masalah}

\CommentTok{// struct foo f;   // Error: incomplete type!}

\KeywordTok{struct}\NormalTok{ foo }\OperatorTok{\{}
    \DataTypeTok{int}\NormalTok{ x}\OperatorTok{,}\NormalTok{ y}\OperatorTok{,}\NormalTok{ z}\OperatorTok{;}
\OperatorTok{\};}                 \CommentTok{// Sekarang struct foo sudah complete!}

\KeywordTok{struct}\NormalTok{ foo f}\OperatorTok{;}      \CommentTok{// Berhasil!}
\end{Highlighting}
\end{Shaded}

\end{tcolorbox}

Perhatikan bahwa walaupun \hl{\mbox{\texttt{void}}} adalah incomplete type, tidak ada cara untuk melengkapinya. Bukan berarti ada yang pernah terpikir melakukan hal aneh itu. Tapi ini menjelaskan kenapa kamu bisa melakukan ini:

\begin{tcolorbox}[enhanced jigsaw,left=3pt,breakable,colback=white!97.2549!black,boxrule=0.2mm]

\begin{Shaded}
\begin{Highlighting}[]
\DataTypeTok{void} \OperatorTok{*}\NormalTok{p}\OperatorTok{;}             \CommentTok{// OK: pointer ke incomplete type}
\end{Highlighting}
\end{Shaded}

\end{tcolorbox}

dan tidak salah satu dari dua ini:

\begin{tcolorbox}[enhanced jigsaw,left=3pt,breakable,colback=white!97.2549!black,boxrule=0.2mm]

\begin{Shaded}
\begin{Highlighting}[]
\DataTypeTok{void}\NormalTok{ v}\OperatorTok{;}              \CommentTok{// Error: mendeklarasikan variabel bertipe incomplete}

\NormalTok{printf}\OperatorTok{(}\StringTok{"}\SpecialCharTok{\%d\textbackslash{}n}\StringTok{"}\OperatorTok{,} \OperatorTok{*}\NormalTok{p}\OperatorTok{);}  \CommentTok{// Error: men{-}dereference incomplete type}
\end{Highlighting}
\end{Shaded}

\end{tcolorbox}

Makin tahu aja\ldots{}

\index{Incomplete types|(}

\chapter{Bilangan Kompleks}\label{bilangan-kompleks}

\index{Complex numbers|(}

Primer singkat tentang \href{https://en.wikipedia.org/wiki/Complex_number}{bilangan kompleks}\footnote{https://en.wikipedia.org/wiki/Complex\_number} yang dicuri langsung dari Wikipedia:

\begin{tcolorbox}[blanker,colback=white,enhanced jigsaw,breakable,arc=0pt,outer arc=0pt,grow to left by=-1cm,enlarge top by=0.1cm,enlarge bottom by=0.1cm,borderline west={1.5pt}{-.3cm}{black}]

\begin{quote}
A \textbf{complex number} is a number that can be expressed in the form \(a+bi\), where \(a\) and \(b\) are real numbers {[}i.e.~floating point types in C{]}, and \(i\) represents the imaginary unit, satisfying the equation \(i^2=−1\). Because no real number satisfies this equation, \(i\) is called an imaginary number. For the complex number \(a+bi\), \(a\) is called the \textbf{real part}, and \(b\) is called the \textbf{imaginary part}.
\end{quote}

\end{tcolorbox}

Tapi hanya sampai di situ yang akan kujelaskan. Kita anggap kalau kamu membaca bab ini, kamu sudah tahu apa itu bilangan kompleks dan apa yang ingin kamu lakukan dengannya.

Yang perlu kita bahas hanyalah kemampuan C untuk itu.

Ternyata, dukungan bilangan kompleks di compiler adalah fitur \emph{opsional}. Tidak semua compiler yang patuh standar bisa melakukannya. Dan yang bisa pun mungkin mendukungnya dengan tingkat kelengkapan yang berbeda-beda.

Kamu bisa menguji apakah sistemmu mendukung bilangan kompleks dengan:

\index{\_\_STDC\_NO\_COMPLEX\_\_ macro@\texttt{\_\_STDC\_NO\_COMPLEX\_\_} macro|(}

\begin{tcolorbox}[enhanced jigsaw,left=3pt,breakable,colback=white!97.2549!black,boxrule=0.2mm]

\begin{Shaded}
\begin{Highlighting}[]
\PreprocessorTok{\#ifdef \_\_STDC\_NO\_COMPLEX\_\_}
\PreprocessorTok{\#error Complex numbers not supported!}
\PreprocessorTok{\#endif}
\end{Highlighting}
\end{Shaded}

\end{tcolorbox}

\index{\_\_STDC\_NO\_COMPLEX\_\_ macro@\texttt{\_\_STDC\_NO\_COMPLEX\_\_} macro|)}

\index{\_\_STDC\_IEC\_559\_COMPLEX\_\_ macro@\texttt{\_\_STDC\_IEC\_559\_COMPLEX\_\_} macro|(}

Selain itu, ada makro yang menandakan kepatuhan pada standar ISO 60559 (IEEE 754) untuk matematika floating point dengan bilangan kompleks, serta adanya tipe \hl{\mbox{\texttt{\_Imaginary}}}.

\begin{tcolorbox}[enhanced jigsaw,left=3pt,breakable,colback=white!97.2549!black,boxrule=0.2mm]

\begin{Shaded}
\begin{Highlighting}[]
\PreprocessorTok{\#if \_\_STDC\_IEC\_559\_COMPLEX\_\_ != 1}
\PreprocessorTok{\#error Need IEC 60559 complex support!}
\PreprocessorTok{\#endif}
\end{Highlighting}
\end{Shaded}

\end{tcolorbox}

\index{\_\_STDC\_IEC\_559\_COMPLEX\_\_ macro@\texttt{\_\_STDC\_IEC\_559\_COMPLEX\_\_} macro|)}

Detail lebih lanjut dijabarkan di Annex G spesifikasi C11.

\section{Tipe Kompleks}\label{tipe-kompleks}

Untuk memakai bilangan kompleks, \index{complex.h header file@\texttt{complex.h} header file} \hl{\mbox{\texttt{\#include\ \textless{}complex.h\textgreater{}}}}.

Dengan itu, kamu mendapat setidaknya dua tipe:

\index{\_Complex type@\texttt{\_Complex} type} \index{complex type@\texttt{complex} type|(}

\begin{tcolorbox}[enhanced jigsaw,left=3pt,breakable,colback=white!97.2549!black,boxrule=0.2mm]

\begin{Shaded}
\begin{Highlighting}[]
\DataTypeTok{\_Complex}
\DataTypeTok{complex}
\end{Highlighting}
\end{Shaded}

\end{tcolorbox}

Keduanya berarti sama, jadi sebaiknya pakai \hl{\mbox{\texttt{complex}}} yang lebih cantik.

\index{complex type@\texttt{complex} type|)}

Kamu juga mendapat beberapa tipe untuk bilangan imajiner kalau implementasimu patuh IEC 60559:

\index{\_Imaginary type@\texttt{\_Imaginary} type} \index{imaginary type@\texttt{imaginary} type|(}

\begin{tcolorbox}[enhanced jigsaw,left=3pt,breakable,colback=white!97.2549!black,boxrule=0.2mm]

\begin{Shaded}
\begin{Highlighting}[]
\DataTypeTok{\_Imaginary}
\DataTypeTok{imaginary}
\end{Highlighting}
\end{Shaded}

\end{tcolorbox}

Ini juga keduanya berarti sama, jadi sebaiknya pakai \hl{\mbox{\texttt{imaginary}}} yang lebih cantik.

\index{imaginary type@\texttt{imaginary} type|)}

Kamu juga mendapat nilai untuk bilangan imajiner \(i\) itu sendiri:

\index{I macro@\texttt{I} macro|(} \index{\_Complex\_I macro@\texttt{\_Complex\_I} macro|(} \index{\_Imaginary\_I macro@\texttt{\_Imaginary\_I} macro|(}

\begin{tcolorbox}[enhanced jigsaw,left=3pt,breakable,colback=white!97.2549!black,boxrule=0.2mm]

\begin{Shaded}
\begin{Highlighting}[]
\NormalTok{I}
\NormalTok{\_Complex\_I}
\NormalTok{\_Imaginary\_I}
\end{Highlighting}
\end{Shaded}

\end{tcolorbox}

Makro \hl{\mbox{\texttt{I}}} diatur menjadi \hl{\mbox{\texttt{\_Imaginary\_I}}} (kalau tersedia), atau \hl{\mbox{\texttt{\_Complex\_I}}}. Jadi pakai saja \hl{\mbox{\texttt{I}}} untuk bilangan imajiner.

\index{I macro@\texttt{I} macro|)} \index{\_Complex\_I macro@\texttt{\_Complex\_I} macro|)}

\index{\_\_STDC\_IEC\_559\_COMPLEX\_\_ macro@\texttt{\_\_STDC\_IEC\_559\_COMPLEX\_\_} macro|(}

Satu catatan sampingan: tadi kubilang bahwa kalau compiler punya \hl{\mbox{\texttt{\_\_STDC\_IEC\_559\_COMPLEX\_\_}}} bernilai \hl{\mbox{\texttt{1}}}, ia harus mendukung tipe \hl{\mbox{\texttt{\_Imaginary}}} agar patuh. Itu bacaanku atas spesifikasi. Namun, aku tidak tahu satu pun compiler yang benar-benar mendukung \hl{\mbox{\texttt{\_Imaginary}}} meskipun \hl{\mbox{\texttt{\_\_STDC\_IEC\_559\_COMPLEX\_\_}}}-nya diatur. Jadi aku akan menulis sebagian kode dengan tipe itu di sini tanpa bisa mengujinya. Maaf!

\index{\_\_STDC\_IEC\_559\_COMPLEX\_\_ macro@\texttt{\_\_STDC\_IEC\_559\_COMPLEX\_\_} macro|)} \index{\_Imaginary\_I macro@\texttt{\_Imaginary\_I} macro|)}

Oke, sekarang kita tahu ada tipe \hl{\mbox{\texttt{complex}}}, bagaimana cara memakainya?

\section{Menugaskan Bilangan Kompleks}\label{menugaskan-bilangan-kompleks}

\index{Complex numbers!declaring|(}

Karena bilangan kompleks punya bagian real dan imajiner, dan keduanya bergantung pada bilangan floating point untuk menyimpan nilai, kita juga perlu memberi tahu C presisi apa yang dipakai untuk bagian-bagian bilangan kompleks itu.

\index{complex float type@\texttt{complex float} type|(} \index{complex double type@\texttt{complex double} type|(} \index{complex long double type@\texttt{complex long double} type|(}

Caranya, tinggal menyematkan \hl{\mbox{\texttt{float}}}, \hl{\mbox{\texttt{double}}}, atau \hl{\mbox{\texttt{long\ double}}} pada \hl{\mbox{\texttt{complex}}}, baik sebelum maupun sesudahnya.

\index{complex long double type@\texttt{complex long double} type|)}

Mari kita definisikan bilangan kompleks yang memakai \hl{\mbox{\texttt{float}}} untuk komponennya:

\begin{tcolorbox}[enhanced jigsaw,left=3pt,breakable,colback=white!97.2549!black,boxrule=0.2mm]

\begin{Shaded}
\begin{Highlighting}[]
\DataTypeTok{float} \DataTypeTok{complex}\NormalTok{ c}\OperatorTok{;}   \CommentTok{// Spesifikasi lebih suka cara ini}
\DataTypeTok{complex} \DataTypeTok{float}\NormalTok{ c}\OperatorTok{;}   \CommentTok{// Sama saja{-}{-}urutan tidak berpengaruh}
\end{Highlighting}
\end{Shaded}

\end{tcolorbox}

Itu bagus untuk deklarasi, tapi bagaimana cara menginisialisasi atau menugaskannya?

Ternyata kita bisa memakai notasi yang cukup alami. Contoh!

\index{I macro@\texttt{I} macro|(}

\begin{tcolorbox}[enhanced jigsaw,left=3pt,breakable,colback=white!97.2549!black,boxrule=0.2mm]

\begin{Shaded}
\begin{Highlighting}[]
\DataTypeTok{double} \DataTypeTok{complex}\NormalTok{ x }\OperatorTok{=} \DecValTok{5} \OperatorTok{+} \DecValTok{2}\OperatorTok{*}\NormalTok{I}\OperatorTok{;}
\DataTypeTok{double} \DataTypeTok{complex}\NormalTok{ y }\OperatorTok{=} \DecValTok{10} \OperatorTok{+} \DecValTok{3}\OperatorTok{*}\NormalTok{I}\OperatorTok{;}
\end{Highlighting}
\end{Shaded}

\end{tcolorbox}

\index{I macro@\texttt{I} macro|)}

Masing-masing untuk \(5+2i\) dan \(10+3i\).

\index{complex float type@\texttt{complex float} type|)} \index{complex double type@\texttt{complex double} type|)} \index{Complex numbers!declaring|)}

\section{Membentuk, Membongkar, dan Mencetak}\label{membentuk-membongkar-dan-mencetak}

Kita hampir sampai\ldots{}

Kita sudah melihat satu cara menulis bilangan kompleks:

\begin{tcolorbox}[enhanced jigsaw,left=3pt,breakable,colback=white!97.2549!black,boxrule=0.2mm]

\begin{Shaded}
\begin{Highlighting}[]
\DataTypeTok{double} \DataTypeTok{complex}\NormalTok{ x }\OperatorTok{=} \DecValTok{5} \OperatorTok{+} \DecValTok{2}\OperatorTok{*}\NormalTok{I}\OperatorTok{;}
\end{Highlighting}
\end{Shaded}

\end{tcolorbox}

Tidak masalah juga memakai bilangan floating point lain untuk membentuknya:

\begin{tcolorbox}[enhanced jigsaw,left=3pt,breakable,colback=white!97.2549!black,boxrule=0.2mm]

\begin{Shaded}
\begin{Highlighting}[]
\DataTypeTok{double}\NormalTok{ a }\OperatorTok{=} \DecValTok{5}\OperatorTok{;}
\DataTypeTok{double}\NormalTok{ b }\OperatorTok{=} \DecValTok{2}\OperatorTok{;}
\DataTypeTok{double} \DataTypeTok{complex}\NormalTok{ x }\OperatorTok{=}\NormalTok{ a }\OperatorTok{+}\NormalTok{ b}\OperatorTok{*}\NormalTok{I}\OperatorTok{;}
\end{Highlighting}
\end{Shaded}

\end{tcolorbox}

\index{CMPLX() macro@\texttt{CMPLX()} macro|(}

Ada juga sekumpulan makro untuk membantu membentuknya. Kode di atas bisa ditulis memakai makro \hl{\mbox{\texttt{CMPLX()}}}, seperti ini:

\begin{tcolorbox}[enhanced jigsaw,left=3pt,breakable,colback=white!97.2549!black,boxrule=0.2mm]

\begin{Shaded}
\begin{Highlighting}[]
\DataTypeTok{double} \DataTypeTok{complex}\NormalTok{ x }\OperatorTok{=}\NormalTok{ CMPLX}\OperatorTok{(}\DecValTok{5}\OperatorTok{,} \DecValTok{2}\OperatorTok{);}
\end{Highlighting}
\end{Shaded}

\end{tcolorbox}

Sejauh yang kutemukan dalam risetku, keduanya \emph{hampir} setara:

\begin{tcolorbox}[enhanced jigsaw,left=3pt,breakable,colback=white!97.2549!black,boxrule=0.2mm]

\begin{Shaded}
\begin{Highlighting}[]
\DataTypeTok{double} \DataTypeTok{complex}\NormalTok{ x }\OperatorTok{=} \DecValTok{5} \OperatorTok{+} \DecValTok{2}\OperatorTok{*}\NormalTok{I}\OperatorTok{;}
\DataTypeTok{double} \DataTypeTok{complex}\NormalTok{ x }\OperatorTok{=}\NormalTok{ CMPLX}\OperatorTok{(}\DecValTok{5}\OperatorTok{,} \DecValTok{2}\OperatorTok{);}
\end{Highlighting}
\end{Shaded}

\end{tcolorbox}

Tapi makro \hl{\mbox{\texttt{CMPLX()}}} akan selalu menangani nol negatif pada bagian imajiner dengan benar, sedangkan cara yang lain mungkin mengonversinya menjadi nol positif. Aku \emph{kira}\footnote{Yang satu ini lebih sulit diriset, dan aku menerima informasi tambahan dari siapa pun. \hl{\mbox{\texttt{I}}} bisa didefinisikan sebagai \hl{\mbox{\texttt{\_Complex\_I}}} atau \hl{\mbox{\texttt{\_Imaginary\_I}}}, kalau yang terakhir ada. \hl{\mbox{\texttt{\_Imaginary\_I}}} akan menangani nol bertanda, tapi \hl{\mbox{\texttt{\_Complex\_I}}} \emph{mungkin} tidak. Ini punya implikasi pada branch cut dan hal-hal lain yang berbau matematika-kompleks. Mungkin. Kelihatan kan aku mulai keluar dari zona nyamanku? Bagaimanapun, makro \hl{\mbox{\texttt{CMPLX()}}} berperilaku seolah-olah \hl{\mbox{\texttt{I}}} didefinisikan sebagai \hl{\mbox{\texttt{\_Imaginary\_I}}}, dengan nol bertanda, bahkan kalau \hl{\mbox{\texttt{\_Imaginary\_I}}} tidak ada di sistem.} ini menyiratkan bahwa kalau ada kemungkinan bagian imajinernya nol, sebaiknya pakai makro\ldots{} tapi seseorang tolong koreksi kalau aku keliru!

Makro \hl{\mbox{\texttt{CMPLX()}}} bekerja pada tipe \hl{\mbox{\texttt{double}}}. Ada dua makro lain untuk \hl{\mbox{\texttt{float}}} dan \hl{\mbox{\texttt{long\ double}}}: \index{CMPLXF() macro@\texttt{CMPLXF()} macro} \hl{\mbox{\texttt{CMPLXF()}}} dan \index{CMPLXL() macro@\texttt{CMPLXL()} macro} \hl{\mbox{\texttt{CMPLXL()}}}. (Sufiks ``f'' dan ``l'' ini muncul di hampir semua fungsi terkait bilangan kompleks.)

\index{CMPLX() macro@\texttt{CMPLX()} macro|)}

Sekarang kita coba kebalikannya: kalau kita punya bilangan kompleks, bagaimana memecahnya menjadi bagian real dan imajinernya?

\index{creal() function@\texttt{creal()} function|(} \index{cimag() function@\texttt{cimag()} function|(}

Ada beberapa fungsi yang mengambil bagian real dan imajiner dari angka itu: \hl{\mbox{\texttt{creal()}}} dan \hl{\mbox{\texttt{cimag()}}}:

\begin{tcolorbox}[enhanced jigsaw,left=3pt,breakable,colback=white!97.2549!black,boxrule=0.2mm]

\begin{Shaded}
\begin{Highlighting}[]
\DataTypeTok{double} \DataTypeTok{complex}\NormalTok{ x }\OperatorTok{=} \DecValTok{5} \OperatorTok{+} \DecValTok{2}\OperatorTok{*}\NormalTok{I}\OperatorTok{;}
\DataTypeTok{double} \DataTypeTok{complex}\NormalTok{ y }\OperatorTok{=} \DecValTok{10} \OperatorTok{+} \DecValTok{3}\OperatorTok{*}\NormalTok{I}\OperatorTok{;}

\NormalTok{printf}\OperatorTok{(}\StringTok{"x = }\SpecialCharTok{\%f}\StringTok{ + }\SpecialCharTok{\%f}\StringTok{i}\SpecialCharTok{\textbackslash{}n}\StringTok{"}\OperatorTok{,}\NormalTok{ creal}\OperatorTok{(}\NormalTok{x}\OperatorTok{),}\NormalTok{ cimag}\OperatorTok{(}\NormalTok{x}\OperatorTok{));}
\NormalTok{printf}\OperatorTok{(}\StringTok{"y = }\SpecialCharTok{\%f}\StringTok{ + }\SpecialCharTok{\%f}\StringTok{i}\SpecialCharTok{\textbackslash{}n}\StringTok{"}\OperatorTok{,}\NormalTok{ creal}\OperatorTok{(}\NormalTok{y}\OperatorTok{),}\NormalTok{ cimag}\OperatorTok{(}\NormalTok{y}\OperatorTok{));}
\end{Highlighting}
\end{Shaded}

\end{tcolorbox}

dengan output:

\begin{tcolorbox}[enhanced jigsaw,left=3pt,breakable,colback=white!97.2549!black,boxrule=0.2mm]

\begin{Shaded}
\begin{Highlighting}[]
\NormalTok{x = 5.000000 + 2.000000i}
\NormalTok{y = 10.000000 + 3.000000i}
\end{Highlighting}
\end{Shaded}

\end{tcolorbox}

Perhatikan bahwa \hl{\mbox{\texttt{i}}} yang kupakai pada format string \hl{\mbox{\texttt{printf()}}} adalah \hl{\mbox{\texttt{i}}} literal yang dicetak---ia bukan bagian dari format specifier. Kedua nilai kembalian \hl{\mbox{\texttt{creal()}}} dan \hl{\mbox{\texttt{cimag()}}} bertipe \hl{\mbox{\texttt{double}}}.

Dan seperti biasa, ada varian \hl{\mbox{\texttt{float}}} dan \hl{\mbox{\texttt{long\ double}}} dari fungsi-fungsi ini: \index{crealf() function@\texttt{crealf()} function} \hl{\mbox{\texttt{crealf()}}}, \index{cimagf() function@\texttt{cimagf()} function} \hl{\mbox{\texttt{cimagf()}}}, \index{creall() function@\texttt{creall()} function} \hl{\mbox{\texttt{creall()}}}, dan \index{cimagl() function@\texttt{cimagl()} function} \hl{\mbox{\texttt{cimagl()}}}.

\index{creal() function@\texttt{creal()} function|)} \index{cimag() function@\texttt{cimag()} function|)}

\section{Aritmetika dan Perbandingan Kompleks}\label{aritmetika-dan-perbandingan-kompleks}

\index{Complex numbers!arithmetic|(}

Aritmetika bisa dilakukan pada bilangan kompleks, walaupun cara kerjanya secara matematis di luar cakupan panduan ini.

\begin{tcolorbox}[enhanced jigsaw,left=3pt,breakable,colback=white!97.2549!black,boxrule=0.2mm]

\begin{Shaded}
\begin{Highlighting}[numbers=left,,]
\PreprocessorTok{\#include }\ImportTok{\textless{}stdio.h\textgreater{}}
\PreprocessorTok{\#include }\ImportTok{\textless{}complex.h\textgreater{}}

\DataTypeTok{int}\NormalTok{ main}\OperatorTok{(}\DataTypeTok{void}\OperatorTok{)}
\OperatorTok{\{}
    \DataTypeTok{double} \DataTypeTok{complex}\NormalTok{ x }\OperatorTok{=} \DecValTok{1} \OperatorTok{+} \DecValTok{2}\OperatorTok{*}\NormalTok{I}\OperatorTok{;}
    \DataTypeTok{double} \DataTypeTok{complex}\NormalTok{ y }\OperatorTok{=} \DecValTok{3} \OperatorTok{+} \DecValTok{4}\OperatorTok{*}\NormalTok{I}\OperatorTok{;}
    \DataTypeTok{double} \DataTypeTok{complex}\NormalTok{ z}\OperatorTok{;}

\NormalTok{    z }\OperatorTok{=}\NormalTok{ x }\OperatorTok{+}\NormalTok{ y}\OperatorTok{;}
\NormalTok{    printf}\OperatorTok{(}\StringTok{"x + y = }\SpecialCharTok{\%f}\StringTok{ + }\SpecialCharTok{\%f}\StringTok{i}\SpecialCharTok{\textbackslash{}n}\StringTok{"}\OperatorTok{,}\NormalTok{ creal}\OperatorTok{(}\NormalTok{z}\OperatorTok{),}\NormalTok{ cimag}\OperatorTok{(}\NormalTok{z}\OperatorTok{));}

\NormalTok{    z }\OperatorTok{=}\NormalTok{ x }\OperatorTok{{-}}\NormalTok{ y}\OperatorTok{;}
\NormalTok{    printf}\OperatorTok{(}\StringTok{"x {-} y = }\SpecialCharTok{\%f}\StringTok{ + }\SpecialCharTok{\%f}\StringTok{i}\SpecialCharTok{\textbackslash{}n}\StringTok{"}\OperatorTok{,}\NormalTok{ creal}\OperatorTok{(}\NormalTok{z}\OperatorTok{),}\NormalTok{ cimag}\OperatorTok{(}\NormalTok{z}\OperatorTok{));}

\NormalTok{    z }\OperatorTok{=}\NormalTok{ x }\OperatorTok{*}\NormalTok{ y}\OperatorTok{;}
\NormalTok{    printf}\OperatorTok{(}\StringTok{"x * y = }\SpecialCharTok{\%f}\StringTok{ + }\SpecialCharTok{\%f}\StringTok{i}\SpecialCharTok{\textbackslash{}n}\StringTok{"}\OperatorTok{,}\NormalTok{ creal}\OperatorTok{(}\NormalTok{z}\OperatorTok{),}\NormalTok{ cimag}\OperatorTok{(}\NormalTok{z}\OperatorTok{));}

\NormalTok{    z }\OperatorTok{=}\NormalTok{ x }\OperatorTok{/}\NormalTok{ y}\OperatorTok{;}
\NormalTok{    printf}\OperatorTok{(}\StringTok{"x / y = }\SpecialCharTok{\%f}\StringTok{ + }\SpecialCharTok{\%f}\StringTok{i}\SpecialCharTok{\textbackslash{}n}\StringTok{"}\OperatorTok{,}\NormalTok{ creal}\OperatorTok{(}\NormalTok{z}\OperatorTok{),}\NormalTok{ cimag}\OperatorTok{(}\NormalTok{z}\OperatorTok{));}
\OperatorTok{\}}
\end{Highlighting}
\end{Shaded}

\end{tcolorbox}

dengan hasil:

\begin{tcolorbox}[enhanced jigsaw,left=3pt,breakable,colback=white!97.2549!black,boxrule=0.2mm]

\begin{Shaded}
\begin{Highlighting}[]
\NormalTok{x + y = 4.000000 + 6.000000i}
\NormalTok{x {-} y = {-}2.000000 + {-}2.000000i}
\NormalTok{x * y = {-}5.000000 + 10.000000i}
\NormalTok{x / y = 0.440000 + 0.080000i}
\end{Highlighting}
\end{Shaded}

\end{tcolorbox}

Kamu juga bisa membandingkan dua bilangan kompleks untuk kesamaan (atau ketidaksamaan):

\begin{tcolorbox}[enhanced jigsaw,left=3pt,breakable,colback=white!97.2549!black,boxrule=0.2mm]

\begin{Shaded}
\begin{Highlighting}[numbers=left,,]
\PreprocessorTok{\#include }\ImportTok{\textless{}stdio.h\textgreater{}}
\PreprocessorTok{\#include }\ImportTok{\textless{}complex.h\textgreater{}}

\DataTypeTok{int}\NormalTok{ main}\OperatorTok{(}\DataTypeTok{void}\OperatorTok{)}
\OperatorTok{\{}
    \DataTypeTok{double} \DataTypeTok{complex}\NormalTok{ x }\OperatorTok{=} \DecValTok{1} \OperatorTok{+} \DecValTok{2}\OperatorTok{*}\NormalTok{I}\OperatorTok{;}
    \DataTypeTok{double} \DataTypeTok{complex}\NormalTok{ y }\OperatorTok{=} \DecValTok{3} \OperatorTok{+} \DecValTok{4}\OperatorTok{*}\NormalTok{I}\OperatorTok{;}

\NormalTok{    printf}\OperatorTok{(}\StringTok{"x == y = }\SpecialCharTok{\%d\textbackslash{}n}\StringTok{"}\OperatorTok{,}\NormalTok{ x }\OperatorTok{==}\NormalTok{ y}\OperatorTok{);}  \CommentTok{// 0}
\NormalTok{    printf}\OperatorTok{(}\StringTok{"x != y = }\SpecialCharTok{\%d\textbackslash{}n}\StringTok{"}\OperatorTok{,}\NormalTok{ x }\OperatorTok{!=}\NormalTok{ y}\OperatorTok{);}  \CommentTok{// 1}
\OperatorTok{\}}
\end{Highlighting}
\end{Shaded}

\end{tcolorbox}

dengan output:

\begin{tcolorbox}[enhanced jigsaw,left=3pt,breakable,colback=white!97.2549!black,boxrule=0.2mm]

\begin{Shaded}
\begin{Highlighting}[]
\NormalTok{x == y = 0}
\NormalTok{x != y = 1}
\end{Highlighting}
\end{Shaded}

\end{tcolorbox}

Keduanya sama kalau kedua komponennya sama. Perhatikan bahwa seperti pada semua floating point, keduanya bisa dianggap sama kalau cukup dekat karena galat pembulatan\footnote{Kesederhanaan pernyataan ini tidak mencerminkan besarnya upaya yang dibutuhkan hanya untuk memahami bagaimana floating point sebenarnya bekerja. https://randomascii.wordpress.com/2012/02/25/comparing-floating-point-numbers-2012-edition/}.

\index{Complex numbers!arithmetic|)}

\section{Matematika Kompleks}\label{matematika-kompleks}

Tapi tunggu! Ada lebih dari sekadar aritmetika kompleks sederhana!

Berikut tabel ringkasan semua fungsi matematika yang tersedia untuk bilangan kompleks.

Aku hanya mencantumkan versi \hl{\mbox{\texttt{double}}} dari setiap fungsi, tapi untuk semuanya ada versi \hl{\mbox{\texttt{float}}} yang didapat dengan menambahkan \hl{\mbox{\texttt{f}}} pada nama fungsi, dan versi \hl{\mbox{\texttt{long\ double}}} dengan menambahkan \hl{\mbox{\texttt{l}}}.

Misalnya, fungsi \hl{\mbox{\texttt{cabs()}}} untuk menghitung nilai mutlak bilangan kompleks juga punya varian \hl{\mbox{\texttt{cabsf()}}} dan \hl{\mbox{\texttt{cabsl()}}}. Aku menghilangkannya demi keringkasan.

\subsection{Fungsi Trigonometri}\label{fungsi-trigonometri}

\index{ccos() function@\texttt{ccos()} function} \index{csin() function@\texttt{csin()} function} \index{ctan() function@\texttt{ctan()} function} \index{cacos() function@\texttt{cacos()} function} \index{casin() function@\texttt{casin()} function} \index{catan() function@\texttt{catan()} function} \index{ccosh() function@\texttt{ccosh()} function} \index{csinh() function@\texttt{csinh()} function} \index{ctanh() function@\texttt{ctanh()} function} \index{cacosh() function@\texttt{cacosh()} function} \index{casinh() function@\texttt{casinh()} function} \index{catanh() function@\texttt{catanh()} function}

{\def\LTcaptype{none} % do not increment counter
\begin{longtable}[]{@{}ll@{}}
\toprule\noalign{}
Fungsi & Deskripsi \\
\midrule\noalign{}
\endhead
\bottomrule\noalign{}
\endlastfoot
\hl{\mbox{\texttt{ccos()}}} & Kosinus \\
\hl{\mbox{\texttt{csin()}}} & Sinus \\
\hl{\mbox{\texttt{ctan()}}} & Tangen \\
\hl{\mbox{\texttt{cacos()}}} & Arc kosinus \\
\hl{\mbox{\texttt{casin()}}} & Arc sinus \\
\hl{\mbox{\texttt{catan()}}} & Bermain \emph{Settlers of Catan} \\
\hl{\mbox{\texttt{ccosh()}}} & Kosinus hiperbolik \\
\hl{\mbox{\texttt{csinh()}}} & Sinus hiperbolik \\
\hl{\mbox{\texttt{ctanh()}}} & Tangen hiperbolik \\
\hl{\mbox{\texttt{cacosh()}}} & Arc kosinus hiperbolik \\
\hl{\mbox{\texttt{casinh()}}} & Arc sinus hiperbolik \\
\hl{\mbox{\texttt{catanh()}}} & Arc tangen hiperbolik \\
\end{longtable}
}

\subsection{Fungsi Eksponensial dan Logaritmik}\label{fungsi-eksponensial-dan-logaritmik}

\index{cexp() function@\texttt{cexp()} function} \index{clog() function@\texttt{clog()} function}

{\def\LTcaptype{none} % do not increment counter
\begin{longtable}[]{@{}ll@{}}
\toprule\noalign{}
Fungsi & Deskripsi \\
\midrule\noalign{}
\endhead
\bottomrule\noalign{}
\endlastfoot
\hl{\mbox{\texttt{cexp()}}} & Eksponensial basis-\(e\) \\
\hl{\mbox{\texttt{clog()}}} & Logaritma natural (basis-\(e\)) \\
\end{longtable}
}

\subsection{Fungsi Pangkat dan Nilai Mutlak}\label{fungsi-pangkat-dan-nilai-mutlak}

\index{cabs() function@\texttt{cabs()} function} \index{cpow() function@\texttt{cpow()} function} \index{csqrt() function@\texttt{csqrt()} function}

{\def\LTcaptype{none} % do not increment counter
\begin{longtable}[]{@{}ll@{}}
\toprule\noalign{}
Fungsi & Deskripsi \\
\midrule\noalign{}
\endhead
\bottomrule\noalign{}
\endlastfoot
\hl{\mbox{\texttt{cabs()}}} & Nilai mutlak \\
\hl{\mbox{\texttt{cpow()}}} & Pangkat \\
\hl{\mbox{\texttt{csqrt()}}} & Akar kuadrat \\
\end{longtable}
}

\subsection{Fungsi Manipulasi}\label{fungsi-manipulasi}

\index{creal() function@\texttt{creal()} function} \index{cimag() function@\texttt{cimag()} function} \index{CMPLX() macro@\texttt{CMPLX()} macro} \index{carg() function@\texttt{carg()} function} \index{conj() function@\texttt{conj()} function} \index{cproj() function@\texttt{cproj()} function}

{\def\LTcaptype{none} % do not increment counter
\begin{longtable}[]{@{}ll@{}}
\toprule\noalign{}
Fungsi & Deskripsi \\
\midrule\noalign{}
\endhead
\bottomrule\noalign{}
\endlastfoot
\hl{\mbox{\texttt{creal()}}} & Mengembalikan bagian real \\
\hl{\mbox{\texttt{cimag()}}} & Mengembalikan bagian imajiner \\
\hl{\mbox{\texttt{CMPLX()}}} & Membentuk bilangan kompleks \\
\hl{\mbox{\texttt{carg()}}} & Argumen/sudut fase \\
\hl{\mbox{\texttt{conj()}}} & Konjugat\footnote{Ini satu-satunya yang tidak diawali \hl{\mbox{\texttt{c}}} tambahan, anehnya.} \\
\hl{\mbox{\texttt{cproj()}}} & Proyeksi pada bola Riemann \\
\end{longtable}
}

\index{Complex numbers|)}

\chapter{Tipe Integer Berlebar Tetap}\label{tipe-integer-berlebar-tetap}

\index{Fixed width integers|(}

C punya semua tipe integer kecil, lebih besar, dan terbesar seperti \hl{\mbox{\texttt{int}}} dan \hl{\mbox{\texttt{long}}} dan sebagainya. Dan kamu bisa melihat \hyperref[limits-macros]{bagian tentang batas} untuk mengetahui \hl{\mbox{\texttt{int}}} terbesar dengan \hl{\mbox{\texttt{INT\_MAX}}} dan seterusnya.

Seberapa besar tipe-tipe itu? Artinya, berapa byte yang mereka pakai? Kita bisa memakai \hl{\mbox{\texttt{sizeof}}} untuk mendapat jawabannya.

Tapi bagaimana kalau aku ingin kebalikannya? Bagaimana kalau aku butuh tipe yang persis 32 bit (4 byte) atau setidaknya 16 bit atau semacamnya?

Bagaimana kita mendeklarasikan tipe dengan ukuran tertentu?

Header \index{stdint.h header file@\texttt{stdint.h} header file} \hl{\mbox{\texttt{\textless{}stdint.h\textgreater{}}}} memberi kita caranya.

\section{Tipe Berukuran Bit}\label{tipe-berukuran-bit}

Untuk integer signed maupun unsigned, kita bisa menentukan tipe yang berukuran sejumlah bit tertentu, dengan beberapa syarat, tentu saja.

Dan ada tiga kelas utama tipe ini (pada contoh-contoh ini, \hl{\mbox{\texttt{N}}} diganti dengan jumlah bit tertentu):

\begin{itemize}
\tightlist
\item
  Integer dengan ukuran persis tertentu (\hl{\mbox{\texttt{intN\_t}}})
\item
  Integer dengan ukuran setidaknya tertentu (\hl{\mbox{\texttt{int\_leastN\_t}}})
\item
  Integer dengan ukuran setidaknya tertentu dan secepat mungkin (\hl{\mbox{\texttt{int\_fastN\_t}}})\footnote{Beberapa arsitektur punya data berukuran berbeda yang bisa diolah CPU dan RAM dengan laju lebih cepat daripada yang lain. Dalam kasus itu, kalau kamu butuh angka 8-bit tercepat, mungkin yang diberikan adalah tipe 16- atau 32-bit karena itu memang lebih cepat. Jadi dengan ini, kamu tidak akan tahu seberapa besar tipenya, tapi ukurannya setidaknya sebesar yang kamu sebutkan.}
\end{itemize}

Seberapa lebih cepat \hl{\mbox{\texttt{fast}}} itu? Pasti mungkin agak lebih cepat. Kira-kira. Spesifikasi tidak menyebut seberapa lebih cepat, hanya bahwa mereka akan menjadi yang tercepat di arsitektur ini. Kebanyakan C compiler cukup pandai, jadi kemungkinan ini hanya dipakai di tempat yang butuh jaminan kecepatan setinggi mungkin (bukan sekadar berharap compiler menghasilkan kode yang cepat banget, yang memang begitu).

Terakhir, tipe angka unsigned ini punya awalan \hl{\mbox{\texttt{u}}} untuk membedakannya.

Misalnya, tipe-tipe ini punya makna yang tercantum:

\index{int\_leastN\_t types@\texttt{int\_leastN\_t} types|(} \index{uint\_leastN\_t types@\texttt{uint\_leastN\_t} types|(} \index{int\_fastN\_t types@\texttt{int\_fastN\_t} types|(} \index{uint\_fastN\_t types@\texttt{uint\_fastN\_t} types|(} \index{intN\_t types@\texttt{intN\_t} types|(} \index{uintN\_t types@\texttt{uintN\_t} types|(}

\begin{tcolorbox}[enhanced jigsaw,left=3pt,breakable,colback=white!97.2549!black,boxrule=0.2mm]

\begin{Shaded}
\begin{Highlighting}[]
\DataTypeTok{int32\_t}\NormalTok{ w}\OperatorTok{;}        \CommentTok{// w persis 32 bit, signed}
\DataTypeTok{uint16\_t}\NormalTok{ x}\OperatorTok{;}       \CommentTok{// x persis 16 bit, unsigned}

\DataTypeTok{int\_least8\_t}\NormalTok{ y}\OperatorTok{;}   \CommentTok{// y setidaknya 8 bit, signed}

\CommentTok{// z adalah representasi tercepat dengan setidaknya 64 bit, unsigned:}
\DataTypeTok{uint\_fast64\_t}\NormalTok{ z}\OperatorTok{;}
\end{Highlighting}
\end{Shaded}

\end{tcolorbox}

Tipe-tipe berikut dijamin terdefinisi:

\begin{tcolorbox}[enhanced jigsaw,left=3pt,breakable,colback=white!97.2549!black,boxrule=0.2mm]

\begin{Shaded}
\begin{Highlighting}[]
\DataTypeTok{int\_least8\_t}      \DataTypeTok{uint\_least8\_t}
\DataTypeTok{int\_least16\_t}     \DataTypeTok{uint\_least16\_t}
\DataTypeTok{int\_least32\_t}     \DataTypeTok{uint\_least32\_t}
\DataTypeTok{int\_least64\_t}     \DataTypeTok{uint\_least64\_t}

\DataTypeTok{int\_fast8\_t}       \DataTypeTok{uint\_fast8\_t}
\DataTypeTok{int\_fast16\_t}      \DataTypeTok{uint\_fast16\_t}
\DataTypeTok{int\_fast32\_t}      \DataTypeTok{uint\_fast32\_t}
\DataTypeTok{int\_fast64\_t}      \DataTypeTok{uint\_fast64\_t}
\end{Highlighting}
\end{Shaded}

\end{tcolorbox}

Mungkin ada yang lain dengan lebar berbeda, tapi itu opsional.

Hei! Di mana tipe berlebar tetap seperti \hl{\mbox{\texttt{int16\_t}}}? Ternyata itu sepenuhnya opsional\ldots kecuali kondisi tertentu terpenuhi\footnote{Yaitu, sistem punya integer 8, 16, 32, atau 64 bit tanpa padding yang memakai representasi two's complement, dan dalam kasus itu varian \hl{\mbox{\texttt{intN\_t}}} untuk jumlah bit tersebut \emph{harus} didefinisikan.}. Dan kalau kamu memakai sistem komputer modern biasa, kondisi itu kemungkinan terpenuhi. Dan kalau terpenuhi, kamu akan punya tipe-tipe ini:

\begin{tcolorbox}[enhanced jigsaw,left=3pt,breakable,colback=white!97.2549!black,boxrule=0.2mm]

\begin{Shaded}
\begin{Highlighting}[]
\DataTypeTok{int8\_t}      \DataTypeTok{uint8\_t}
\DataTypeTok{int16\_t}     \DataTypeTok{uint16\_t}
\DataTypeTok{int32\_t}     \DataTypeTok{uint32\_t}
\DataTypeTok{int64\_t}     \DataTypeTok{uint64\_t}
\end{Highlighting}
\end{Shaded}

\end{tcolorbox}

Varian lain dengan lebar berbeda mungkin didefinisikan, tapi itu opsional.

\index{int\_leastN\_t types@\texttt{int\_leastN\_t} types|)} \index{uint\_leastN\_t types@\texttt{uint\_leastN\_t} types|)} \index{int\_fastN\_t types@\texttt{int\_fastN\_t} types|)} \index{uint\_fastN\_t types@\texttt{uint\_fastN\_t} types|)} \index{intN\_t types@\texttt{intN\_t} types|)} \index{uintN\_t types@\texttt{uintN\_t} types|)}

\section{Tipe Ukuran Integer Maksimum}\label{tipe-ukuran-integer-maksimum}

Ada tipe yang bisa dipakai untuk menampung integer terbesar yang bisa direpresentasikan di sistem, baik signed maupun unsigned:

\index{intmax\_t type@\texttt{intmax\_t} type|(} \index{uintmax\_t type@\texttt{uintmax\_t} type|(}

\begin{tcolorbox}[enhanced jigsaw,left=3pt,breakable,colback=white!97.2549!black,boxrule=0.2mm]

\begin{Shaded}
\begin{Highlighting}[]
\DataTypeTok{intmax\_t}
\DataTypeTok{uintmax\_t}
\end{Highlighting}
\end{Shaded}

\end{tcolorbox}

\index{intmax\_t type@\texttt{intmax\_t} type|)} \index{uintmax\_t type@\texttt{uintmax\_t} type|)}

Pakai tipe ini kalau kamu ingin sebesar mungkin.

Jelas nilai dari tipe integer lain dengan tanda yang sama pasti muat di tipe ini.

\section{Memakai Konstanta Berukuran Tetap}\label{memakai-konstanta-berukuran-tetap}

Kalau kamu punya konstanta yang ingin dimuat dalam sejumlah bit tertentu, kamu bisa memakai makro berikut untuk menambahkan sufiks yang tepat secara otomatis pada angka itu (mis. \hl{\mbox{\texttt{22L}}} atau \hl{\mbox{\texttt{3490ULL}}}).

\index{INTn\_C() macros@\texttt{INTn\_C()} macros|(} \index{UINTn\_C() macros@\texttt{UINTn\_C()} macros|(} \index{INTMAX\_C() macro@\texttt{INTMAX\_C()} macro|(} \index{UINTMAX\_C() macro@\texttt{UINTMAX\_C()} macro|(}

\begin{tcolorbox}[enhanced jigsaw,left=3pt,breakable,colback=white!97.2549!black,boxrule=0.2mm]

\begin{Shaded}
\begin{Highlighting}[]
\NormalTok{INT8\_C}\OperatorTok{(}\NormalTok{x}\OperatorTok{)}\NormalTok{     UINT8\_C}\OperatorTok{(}\NormalTok{x}\OperatorTok{)}
\NormalTok{INT16\_C}\OperatorTok{(}\NormalTok{x}\OperatorTok{)}\NormalTok{    UINT16\_C}\OperatorTok{(}\NormalTok{x}\OperatorTok{)}
\NormalTok{INT32\_C}\OperatorTok{(}\NormalTok{x}\OperatorTok{)}\NormalTok{    UINT32\_C}\OperatorTok{(}\NormalTok{x}\OperatorTok{)}
\NormalTok{INT64\_C}\OperatorTok{(}\NormalTok{x}\OperatorTok{)}\NormalTok{    UINT64\_C}\OperatorTok{(}\NormalTok{x}\OperatorTok{)}
\NormalTok{INTMAX\_C}\OperatorTok{(}\NormalTok{x}\OperatorTok{)}\NormalTok{   UINTMAX\_C}\OperatorTok{(}\NormalTok{x}\OperatorTok{)}
\end{Highlighting}
\end{Shaded}

\end{tcolorbox}

\index{INTn\_C() macros@\texttt{INTn\_C()} macros|)} \index{UINTMAX\_C() macro@\texttt{UINTMAX\_C()} macro|)}

Sekali lagi, ini hanya bekerja dengan nilai integer konstan.

Misalnya, kita bisa memakai salah satunya untuk menugaskan nilai konstan seperti ini:

\begin{tcolorbox}[enhanced jigsaw,left=3pt,breakable,colback=white!97.2549!black,boxrule=0.2mm]

\begin{Shaded}
\begin{Highlighting}[]
\DataTypeTok{uint16\_t}\NormalTok{ x }\OperatorTok{=}\NormalTok{ UINT16\_C}\OperatorTok{(}\DecValTok{12}\OperatorTok{);}
\DataTypeTok{intmax\_t}\NormalTok{ y }\OperatorTok{=}\NormalTok{ INTMAX\_C}\OperatorTok{(}\DecValTok{3490}\OperatorTok{);}
\end{Highlighting}
\end{Shaded}

\end{tcolorbox}

\index{UINTn\_C() macros@\texttt{UINTn\_C()} macros|)} \index{INTMAX\_C() macro@\texttt{INTMAX\_C()} macro|)}

\section{Batas Integer Berukuran Tetap}\label{batas-integer-berukuran-tetap}

Kita juga punya beberapa batas yang didefinisikan agar kamu bisa mendapat nilai maksimum dan minimum untuk tipe-tipe ini:

\index{INTn\_MAX macros@\texttt{INTn\_MAX} macros|(} \index{INTn\_MIN macros@\texttt{INTn\_MIN} macros|(} \index{UINTn\_MAX macros@\texttt{UINTn\_MAX} macros|(} \index{INT\_LEASTn\_MAX macros@\texttt{INT\_LEASTn\_MAX} macros|(} \index{INT\_LEASTn\_MIN macros@\texttt{INT\_LEASTn\_MIN} macros|(} \index{UINT\_LEASTn\_MAX macros@\texttt{UINT\_LEASTn\_MAX} macros|(} \index{INT\_FASTn\_MAX macros@\texttt{INT\_FASTn\_MAX} macros|(} \index{INT\_FASTn\_MIN macros@\texttt{INT\_FASTn\_MIN} macros|(} \index{UINT\_FASTn\_MAX macros@\texttt{UINT\_FASTn\_MAX} macros|(} \index{INTMAX\_MAX macro@\texttt{INTMAX\_MAX} macro|(} \index{INTMAX\_MIN macro@\texttt{INTMAX\_MIN} macro|(} \index{UINTMAX\_MAX macro@\texttt{UINTMAX\_MAX} macro|(}

\begin{tcolorbox}[enhanced jigsaw,left=3pt,breakable,colback=white!97.2549!black,boxrule=0.2mm]

\begin{Shaded}
\begin{Highlighting}[]
\NormalTok{INT8\_MAX           INT8\_MIN           UINT8\_MAX}
\NormalTok{INT16\_MAX          INT16\_MIN          UINT16\_MAX}
\NormalTok{INT32\_MAX          INT32\_MIN          UINT32\_MAX}
\NormalTok{INT64\_MAX          INT64\_MIN          UINT64\_MAX}

\NormalTok{INT\_LEAST8\_MAX     INT\_LEAST8\_MIN     UINT\_LEAST8\_MAX}
\NormalTok{INT\_LEAST16\_MAX    INT\_LEAST16\_MIN    UINT\_LEAST16\_MAX}
\NormalTok{INT\_LEAST32\_MAX    INT\_LEAST32\_MIN    UINT\_LEAST32\_MAX}
\NormalTok{INT\_LEAST64\_MAX    INT\_LEAST64\_MIN    UINT\_LEAST64\_MAX}

\NormalTok{INT\_FAST8\_MAX      INT\_FAST8\_MIN      UINT\_FAST8\_MAX}
\NormalTok{INT\_FAST16\_MAX     INT\_FAST16\_MIN     UINT\_FAST16\_MAX}
\NormalTok{INT\_FAST32\_MAX     INT\_FAST32\_MIN     UINT\_FAST32\_MAX}
\NormalTok{INT\_FAST64\_MAX     INT\_FAST64\_MIN     UINT\_FAST64\_MAX}

\NormalTok{INTMAX\_MAX         INTMAX\_MIN         UINTMAX\_MAX}
\end{Highlighting}
\end{Shaded}

\end{tcolorbox}

\index{INTn\_MAX macros@\texttt{INTn\_MAX} macros|)} \index{UINTn\_MAX macros@\texttt{UINTn\_MAX} macros|)} \index{INT\_LEASTn\_MAX macros@\texttt{INT\_LEASTn\_MAX} macros|)} \index{UINT\_LEASTn\_MAX macros@\texttt{UINT\_LEASTn\_MAX} macros|)} \index{INT\_FASTn\_MAX macros@\texttt{INT\_FASTn\_MAX} macros|)} \index{UINT\_FASTn\_MAX macros@\texttt{UINT\_FASTn\_MAX} macros|)} \index{INTMAX\_MAX macro@\texttt{INTMAX\_MAX} macro|)} \index{UINTMAX\_MAX macro@\texttt{UINTMAX\_MAX} macro|)}

Perhatikan bahwa \hl{\mbox{\texttt{MIN}}} untuk semua tipe unsigned adalah \hl{\mbox{\texttt{0}}}, sehingga tidak ada makronya.

\index{INTn\_MIN macros@\texttt{INTn\_MIN} macros|)} \index{INT\_LEASTn\_MIN macros@\texttt{INT\_LEASTn\_MIN} macros|)} \index{INT\_FASTn\_MIN macros@\texttt{INT\_FASTn\_MIN} macros|)} \index{INTMAX\_MIN macro@\texttt{INTMAX\_MIN} macro|)}

\section{Format Specifier}\label{format-specifier}

Untuk mencetak tipe-tipe ini, kamu perlu mengirim format specifier yang tepat ke \index{printf() function@\texttt{printf()} function} \hl{\mbox{\texttt{printf()}}}. (Dan masalah yang sama untuk mengambil input dengan \index{scanf() function@\texttt{scanf()} function} \hl{\mbox{\texttt{scanf()}}}.)

Tapi bagaimana kamu akan tahu ukuran tipe di baliknya? Untungnya, sekali lagi, C menyediakan beberapa makro untuk membantu.

Semuanya ada di \hl{\mbox{\texttt{\textless{}inttypes.h\textgreater{}}}}.

Nah, ada banyak sekali makro. Seperti ledakan kompleksitas makro. Jadi aku akan berhenti mendaftar satu per satu dan cukup menaruh huruf kecil \hl{\mbox{\texttt{n}}} di tempat kamu seharusnya menaruh \hl{\mbox{\texttt{8}}}, \hl{\mbox{\texttt{16}}}, \hl{\mbox{\texttt{32}}}, atau \hl{\mbox{\texttt{64}}} sesuai kebutuhan.

Mari kita lihat makro untuk mencetak integer signed:

\index{PRIdn macros@\texttt{PRIdn} macros|(} \index{PRIin macros@\texttt{PRIin} macros|(} \index{PRIdLEASTn macros@\texttt{PRIdLEASTn} macros|(} \index{PRIiLEASTn macros@\texttt{PRIiLEASTn} macros|(} \index{PRIdFASTn macros@\texttt{PRIdFASTn} macros|(} \index{PRIiFASTn macros@\texttt{PRIiFASTn} macros|(} \index{PRIdMAX macro@\texttt{PRIdMAX} macro|(} \index{PRIiMAX macro@\texttt{PRIiMAX} macro|(}

\begin{tcolorbox}[enhanced jigsaw,left=3pt,breakable,colback=white!97.2549!black,boxrule=0.2mm]

\begin{Shaded}
\begin{Highlighting}[]
\NormalTok{PRIdn    PRIdLEASTn    PRIdFASTn    PRIdMAX}
\NormalTok{PRIin    PRIiLEASTn    PRIiFASTn    PRIiMAX}
\end{Highlighting}
\end{Shaded}

\end{tcolorbox}

Cari polanya. Kamu bisa melihat ada varian untuk tipe fixed, least, fast, dan max.

Dan ada juga \hl{\mbox{\texttt{d}}} huruf kecil dan \hl{\mbox{\texttt{i}}} huruf kecil. Keduanya sesuai dengan format specifier \hl{\mbox{\texttt{printf()}}} \hl{\mbox{\texttt{\%d}}} dan \hl{\mbox{\texttt{\%i}}}.

Jadi kalau aku punya sesuatu bertipe:

\begin{tcolorbox}[enhanced jigsaw,left=3pt,breakable,colback=white!97.2549!black,boxrule=0.2mm]

\begin{Shaded}
\begin{Highlighting}[]
\DataTypeTok{int\_least16\_t}\NormalTok{ x }\OperatorTok{=} \DecValTok{3490}\OperatorTok{;}
\end{Highlighting}
\end{Shaded}

\end{tcolorbox}

Aku bisa mencetaknya dengan padanan format specifier \hl{\mbox{\texttt{\%d}}} dengan memakai \hl{\mbox{\texttt{PRIdLEAST16}}}.

Tapi bagaimana caranya? Bagaimana memakai makro itu?

Pertama, makro itu menentukan string berisi huruf atau huruf-huruf yang dibutuhkan \hl{\mbox{\texttt{printf()}}} untuk mencetak tipe itu. Misalnya, bisa berupa \hl{\mbox{\texttt{"d"}}} atau \hl{\mbox{\texttt{"ld"}}}.

Jadi yang perlu kita lakukan hanyalah menyematkannya ke format string pada pemanggilan \hl{\mbox{\texttt{printf()}}}.

Untuk itu, kita bisa memanfaatkan fakta tentang C yang mungkin sudah kamu lupakan: string literal yang bersebelahan otomatis digabung menjadi satu string. Misalnya:

\begin{tcolorbox}[enhanced jigsaw,left=3pt,breakable,colback=white!97.2549!black,boxrule=0.2mm]

\begin{Shaded}
\begin{Highlighting}[]
\NormalTok{printf}\OperatorTok{(}\StringTok{"Hello, "} \StringTok{"world!}\SpecialCharTok{\textbackslash{}n}\StringTok{"}\OperatorTok{);}   \CommentTok{// Mencetak "Hello, world!"}
\end{Highlighting}
\end{Shaded}

\end{tcolorbox}

Dan karena makro-makro ini adalah string literal, kita bisa memakainya seperti ini:

\begin{tcolorbox}[enhanced jigsaw,left=3pt,breakable,colback=white!97.2549!black,boxrule=0.2mm]

\begin{Shaded}
\begin{Highlighting}[numbers=left,,]
\PreprocessorTok{\#include }\ImportTok{\textless{}stdio.h\textgreater{}}
\PreprocessorTok{\#include }\ImportTok{\textless{}stdint.h\textgreater{}}
\PreprocessorTok{\#include }\ImportTok{\textless{}inttypes.h\textgreater{}}

\DataTypeTok{int}\NormalTok{ main}\OperatorTok{(}\DataTypeTok{void}\OperatorTok{)}
\OperatorTok{\{}
    \DataTypeTok{int\_least16\_t}\NormalTok{ x }\OperatorTok{=} \DecValTok{3490}\OperatorTok{;}

\NormalTok{    printf}\OperatorTok{(}\StringTok{"The value is \%"}\NormalTok{ PRIdLEAST16 }\StringTok{"!}\SpecialCharTok{\textbackslash{}n}\StringTok{"}\OperatorTok{,}\NormalTok{ x}\OperatorTok{);}
\OperatorTok{\}}
\end{Highlighting}
\end{Shaded}

\end{tcolorbox}

\index{PRIdn macros@\texttt{PRIdn} macros|)} \index{PRIin macros@\texttt{PRIin} macros|)} \index{PRIdLEASTn macros@\texttt{PRIdLEASTn} macros|)} \index{PRIiLEASTn macros@\texttt{PRIiLEASTn} macros|)} \index{PRIdFASTn macros@\texttt{PRIdFASTn} macros|)} \index{PRIiFASTn macros@\texttt{PRIiFASTn} macros|)} \index{PRIdMAX macro@\texttt{PRIdMAX} macro|)} \index{PRIiMAX macro@\texttt{PRIiMAX} macro|)}

Kita juga punya setumpuk makro untuk mencetak tipe unsigned:

\index{PRIon macros@\texttt{PRIon} macros|(} \index{PRIun macros@\texttt{PRIun} macros|(} \index{PRIxn macros@\texttt{PRIxn} macros|(} \index{PRIXn macros@\texttt{PRIXn} macros|(} \index{PRIoLEASTn macros@\texttt{PRIoLEASTn} macros|(} \index{PRIuLEASTn macros@\texttt{PRIuLEASTn} macros|(} \index{PRIxLEASTn macros@\texttt{PRIxLEASTn} macros|(} \index{PRIXLEASTn macros@\texttt{PRIXLEASTn} macros|(} \index{PRIoFASTn macros@\texttt{PRIoFASTn} macros|(} \index{PRIuFASTn macros@\texttt{PRIuFASTn} macros|(} \index{PRIxFASTn macros@\texttt{PRIxFASTn} macros|(} \index{PRIXFASTn macros@\texttt{PRIXFASTn} macros|(} \index{PRIoMAX macros@\texttt{PRIoMAX} macros|(} \index{PRIuMAX macros@\texttt{PRIuMAX} macros|(} \index{PRIxMAX macros@\texttt{PRIxMAX} macros|(} \index{PRIXMAX macros@\texttt{PRIXMAX} macros|(}

\begin{tcolorbox}[enhanced jigsaw,left=3pt,breakable,colback=white!97.2549!black,boxrule=0.2mm]

\begin{Shaded}
\begin{Highlighting}[]
\NormalTok{PRIon    PRIoLEASTn    PRIoFASTn    PRIoMAX}
\NormalTok{PRIun    PRIuLEASTn    PRIuFASTn    PRIuMAX}
\NormalTok{PRIxn    PRIxLEASTn    PRIxFASTn    PRIxMAX}
\NormalTok{PRIXn    PRIXLEASTn    PRIXFASTn    PRIXMAX}
\end{Highlighting}
\end{Shaded}

\end{tcolorbox}

Dalam kasus ini, \hl{\mbox{\texttt{o}}}, \hl{\mbox{\texttt{u}}}, \hl{\mbox{\texttt{x}}}, dan \hl{\mbox{\texttt{X}}} sesuai dengan format specifier yang terdokumentasi di \hl{\mbox{\texttt{printf()}}}.

Dan, seperti tadi, huruf kecil \hl{\mbox{\texttt{n}}} harus diganti dengan \hl{\mbox{\texttt{8}}}, \hl{\mbox{\texttt{16}}}, \hl{\mbox{\texttt{32}}}, atau \hl{\mbox{\texttt{64}}}.

\index{PRIon macros@\texttt{PRIon} macros|)} \index{PRIun macros@\texttt{PRIun} macros|)} \index{PRIxn macros@\texttt{PRIxn} macros|)} \index{PRIXn macros@\texttt{PRIXn} macros|)} \index{PRIoLEASTn macros@\texttt{PRIoLEASTn} macros|)} \index{PRIuLEASTn macros@\texttt{PRIuLEASTn} macros|)} \index{PRIxLEASTn macros@\texttt{PRIxLEASTn} macros|)} \index{PRIXLEASTn macros@\texttt{PRIXLEASTn} macros|)} \index{PRIoFASTn macros@\texttt{PRIoFASTn} macros|)} \index{PRIuFASTn macros@\texttt{PRIuFASTn} macros|)} \index{PRIxFASTn macros@\texttt{PRIxFASTn} macros|)} \index{PRIXFASTn macros@\texttt{PRIXFASTn} macros|)} \index{PRIoMAX macros@\texttt{PRIoMAX} macros|)} \index{PRIuMAX macros@\texttt{PRIuMAX} macros|)} \index{PRIxMAX macros@\texttt{PRIxMAX} macros|)} \index{PRIXMAX macros@\texttt{PRIXMAX} macros|)}

Tapi tepat ketika kamu mengira sudah cukup dengan makro, ternyata kita punya satu set pelengkap lengkap untuk \index{scanf() function@\texttt{scanf()} function} \hl{\mbox{\texttt{scanf()}}}!

\index{SCNdn macros@\texttt{SCNdn} macros|(} \index{SCNin macros@\texttt{SCNin} macros|(} \index{SCNon macros@\texttt{SCNon} macros|(} \index{SCNun macros@\texttt{SCNun} macros|(} \index{SCNxn macros@\texttt{SCNxn} macros|(} \index{SCNdLEASTn macros@\texttt{SCNdLEASTn} macros|(} \index{SCNiLEASTn macros@\texttt{SCNiLEASTn} macros|(} \index{SCNoLEASTn macros@\texttt{SCNoLEASTn} macros|(} \index{SCNuLEASTn macros@\texttt{SCNuLEASTn} macros|(} \index{SCNxLEASTn macros@\texttt{SCNxLEASTn} macros|(} \index{SCNdFASTn macros@\texttt{SCNdFASTn} macros|(} \index{SCNiFASTn macros@\texttt{SCNiFASTn} macros|(} \index{SCNoFASTn macros@\texttt{SCNoFASTn} macros|(} \index{SCNuFASTn macros@\texttt{SCNuFASTn} macros|(} \index{SCNxFASTn macros@\texttt{SCNxFASTn} macros|(} \index{SCNdMAX macros@\texttt{SCNdMAX} macros|(} \index{SCNiMAX macros@\texttt{SCNiMAX} macros|(} \index{SCNoMAX macros@\texttt{SCNoMAX} macros|(} \index{SCNuMAX macros@\texttt{SCNuMAX} macros|(} \index{SCNxMAX macros@\texttt{SCNxMAX} macros|(}

\begin{tcolorbox}[enhanced jigsaw,left=3pt,breakable,colback=white!97.2549!black,boxrule=0.2mm]

\begin{Shaded}
\begin{Highlighting}[]
\NormalTok{SCNdn    SCNdLEASTn    SCNdFASTn    SCNdMAX}
\NormalTok{SCNin    SCNiLEASTn    SCNiFASTn    SCNiMAX}
\NormalTok{SCNon    SCNoLEASTn    SCNoFASTn    SCNoMAX}
\NormalTok{SCNun    SCNuLEASTn    SCNuFASTn    SCNuMAX}
\NormalTok{SCNxn    SCNxLEASTn    SCNxFASTn    SCNxMAX}
\end{Highlighting}
\end{Shaded}

\end{tcolorbox}

\index{SCNdn macros@\texttt{SCNdn} macros|)} \index{SCNin macros@\texttt{SCNin} macros|)} \index{SCNon macros@\texttt{SCNon} macros|)} \index{SCNun macros@\texttt{SCNun} macros|)} \index{SCNxn macros@\texttt{SCNxn} macros|)} \index{SCNdLEASTn macros@\texttt{SCNdLEASTn} macros|)} \index{SCNiLEASTn macros@\texttt{SCNiLEASTn} macros|)} \index{SCNoLEASTn macros@\texttt{SCNoLEASTn} macros|)} \index{SCNuLEASTn macros@\texttt{SCNuLEASTn} macros|)} \index{SCNxLEASTn macros@\texttt{SCNxLEASTn} macros|)} \index{SCNdFASTn macros@\texttt{SCNdFASTn} macros|)} \index{SCNiFASTn macros@\texttt{SCNiFASTn} macros|)} \index{SCNoFASTn macros@\texttt{SCNoFASTn} macros|)} \index{SCNuFASTn macros@\texttt{SCNuFASTn} macros|)} \index{SCNxFASTn macros@\texttt{SCNxFASTn} macros|)} \index{SCNdMAX macros@\texttt{SCNdMAX} macros|)} \index{SCNiMAX macros@\texttt{SCNiMAX} macros|)} \index{SCNoMAX macros@\texttt{SCNoMAX} macros|)} \index{SCNuMAX macros@\texttt{SCNuMAX} macros|)} \index{SCNxMAX macros@\texttt{SCNxMAX} macros|)}

Ingat: saat ingin mencetak tipe integer berukuran tetap dengan \hl{\mbox{\texttt{printf()}}} atau \hl{\mbox{\texttt{scanf()}}}, ambil format specifier yang sesuai dari \hl{\mbox{\texttt{\textless{}inttypes.h\textgreater{}}}}.

\index{Fixed width integers|)}

\chapter{Fungsionalitas Tanggal dan Waktu}\label{fungsionalitas-tanggal-dan-waktu}

\begin{tcolorbox}[blanker,colback=white,enhanced jigsaw,breakable,arc=0pt,outer arc=0pt,grow to left by=-1cm,enlarge top by=0.1cm,enlarge bottom by=0.1cm,borderline west={1.5pt}{-.3cm}{black}]

\begin{quote}
``Time is an illusion. Lunchtime doubly so.''\\
---Ford Prefect, The Hitchhikers Guide to the Galaxy
\end{quote}

\end{tcolorbox}

\index{Date and time|(}

Ini tidak terlalu rumit, tapi bisa sedikit mengintimidasi pada awalnya, baik karena banyaknya tipe yang tersedia maupun cara kita mengonversi antar tipe itu.

Campurkan GMT (UTC) dan waktu lokal, maka kita dapat semua \emph{Keseruan Biasa}™ yang selalu ada pada waktu dan tanggal.

Dan tentu saja jangan pernah lupa aturan emas tanggal dan waktu: \emph{Jangan pernah mencoba menulis fungsionalitas tanggal dan waktu sendiri. Pakai saja yang disediakan library.}

Waktu terlalu rumit untuk ditangani dengan benar oleh programmer manusia biasa. Serius, kita semua berutang satu poin kepada setiap orang yang mengerjakan library tanggal dan waktu mana pun, jadi masukkan itu ke anggaranmu.

\section{Terminologi dan Informasi Singkat}\label{terminologi-dan-informasi-singkat}

Beberapa istilah singkat kalau-kalau kamu belum hafal.

\begin{itemize}
\item
  \index{Universal Coordinated Time} \textbf{UTC}: Coordinated Universal Time adalah waktu absolut yang disepakati secara universal\footnote{Di Bumi, setidaknya. Siapa tahu sistem gila apa yang dipakai \emph{di luar sana}\ldots{}}. Semua orang di planet ini menganggap sekarang waktunya sama di UTC\ldots{} meskipun waktu lokal mereka berbeda.
\item
  \index{Greenwich Mean time} \textbf{GMT}: Greenwich Mean Time, secara efektif sama dengan UTC\footnote{Oke, jangan hajar aku! GMT secara teknis adalah zona waktu, sedangkan UTC adalah sistem waktu global. Selain itu beberapa negara mungkin menyesuaikan GMT untuk daylight saving time, sedangkan UTC tidak pernah disesuaikan untuk daylight saving time.}. Kemungkinan kamu ingin mengatakan UTC, atau ``universal time''. Kalau kamu membicarakan zona waktu GMT secara khusus, sebut GMT. Membingungkan, banyak fungsi UTC di C lebih tua daripada UTC dan masih merujuk ke Greenwich Mean Time. Saat melihat itu, ketahuilah bahwa maksud C adalah UTC.
\item
  \index{Local time} \textbf{Waktu lokal}: jam berapa sekarang di tempat komputer yang menjalankan program berada. Ini dinyatakan sebagai selisih dari UTC. Walaupun ada banyak zona waktu di dunia, kebanyakan komputer bekerja dalam waktu lokal atau UTC.
\end{itemize}

\index{Universal Coordinated Time|(}

Sebagai aturan umum, kalau kamu menggambarkan kejadian yang terjadi satu kali, seperti entri log, atau peluncuran roket, atau saat pointer akhirnya klik bagimu, pakai UTC.

\index{Local time|(}

Sebaliknya, kalau itu sesuatu yang terjadi pada waktu yang sama \emph{di setiap zona waktu}, seperti malam Tahun Baru atau jam makan malam, pakai waktu lokal.

Karena banyak bahasa hanya pandai mengonversi antara UTC dan waktu lokal, kamu bisa menyakiti dirimu sendiri dengan memilih menyimpan tanggalmu dalam bentuk yang salah. (Tanya aku bagaimana aku tahu.)

\index{Local time|)} \index{Universal Coordinated Time|)}

\section{Tipe Tanggal}\label{tipe-tanggal}

Ada dua\footnote{Harus diakui, ada lebih dari dua.} tipe utama di C untuk tanggal: \hl{\mbox{\texttt{time\_t}}} dan \hl{\mbox{\texttt{struct\ tm}}}.

Spesifikasi sebenarnya tidak banyak bicara soal keduanya:

\begin{itemize}
\item
  \index{time\_t type@\texttt{time\_t} type} \hl{\mbox{\texttt{time\_t}}}: tipe real yang mampu menampung waktu. Jadi menurut spesifikasi, ini bisa berupa tipe floating atau integer. Di POSIX (sistem mirip Unix), ia integer. Ia menyimpan \emph{calendar time}. Yang bisa dianggap sebagai waktu UTC.
\item
  \index{struct tm type@\texttt{struct tm} type} \hl{\mbox{\texttt{struct\ tm}}}: menyimpan komponen sebuah calendar time. Ini \emph{broken-down time}, yaitu komponen-komponen waktu, seperti jam, menit, detik, hari, bulan, tahun, dan sebagainya.
\end{itemize}

\index{time\_t type@\texttt{time\_t} type|(}

Di banyak sistem, \hl{\mbox{\texttt{time\_t}}} mewakili jumlah detik sejak \href{https://en.wikipedia.org/wiki/Unix_time}{\emph{Epoch}}\footnote{https://en.wikipedia.org/wiki/Unix\_time}. Epoch dalam beberapa hal adalah awal waktu dari sudut pandang komputer, yang umumnya 1 Januari 1970 UTC. \hl{\mbox{\texttt{time\_t}}} bisa negatif untuk mewakili waktu sebelum Epoch. Windows berperilaku sama dengan Unix sejauh yang kutahu.

\index{time\_t type@\texttt{time\_t} type|)} \index{struct tm type@\texttt{struct tm} type|(}

Dan apa isi \hl{\mbox{\texttt{struct\ tm}}}? Field-field berikut:

\begin{tcolorbox}[enhanced jigsaw,left=3pt,breakable,colback=white!97.2549!black,boxrule=0.2mm]

\begin{Shaded}
\begin{Highlighting}[]
\KeywordTok{struct}\NormalTok{ tm }\OperatorTok{\{}
    \DataTypeTok{int}\NormalTok{ tm\_sec}\OperatorTok{;}    \CommentTok{// detik setelah menit {-}{-} [0, 60]}
    \DataTypeTok{int}\NormalTok{ tm\_min}\OperatorTok{;}    \CommentTok{// menit setelah jam {-}{-} [0, 59]}
    \DataTypeTok{int}\NormalTok{ tm\_hour}\OperatorTok{;}   \CommentTok{// jam sejak tengah malam {-}{-} [0, 23]}
    \DataTypeTok{int}\NormalTok{ tm\_mday}\OperatorTok{;}   \CommentTok{// hari dalam bulan {-}{-} [1, 31]}
    \DataTypeTok{int}\NormalTok{ tm\_mon}\OperatorTok{;}    \CommentTok{// bulan sejak Januari {-}{-} [0, 11]}
    \DataTypeTok{int}\NormalTok{ tm\_year}\OperatorTok{;}   \CommentTok{// tahun sejak 1900}
    \DataTypeTok{int}\NormalTok{ tm\_wday}\OperatorTok{;}   \CommentTok{// hari sejak Minggu {-}{-} [0, 6]}
    \DataTypeTok{int}\NormalTok{ tm\_yday}\OperatorTok{;}   \CommentTok{// hari sejak 1 Januari {-}{-} [0, 365]}
    \DataTypeTok{int}\NormalTok{ tm\_isdst}\OperatorTok{;}  \CommentTok{// Flag Daylight Saving Time}
\OperatorTok{\};}
\end{Highlighting}
\end{Shaded}

\end{tcolorbox}

Perhatikan bahwa semuanya berbasis nol kecuali hari dalam bulan.

Penting diketahui bahwa kamu bisa memasukkan nilai apa pun yang kamu mau ke tipe-tipe ini. Ada fungsi untuk membantu mengambil waktu \emph{sekarang}, tapi tipe-tipe itu menyimpan \emph{suatu} waktu, bukan \emph{satu-satunya} waktu.

\index{struct tm type@\texttt{struct tm} type|)}

Jadi pertanyaannya menjadi: ``Bagaimana menginisialisasi data bertipe ini, dan bagaimana mengonversi di antara mereka?''

\section{Inisialisasi dan Konversi Antartipe}\label{inisialisasi-dan-konversi-antartipe}

\index{time() function@\texttt{time()} function|(}

Pertama, kamu bisa mengambil waktu saat ini dan menyimpannya dalam \hl{\mbox{\texttt{time\_t}}} dengan fungsi \hl{\mbox{\texttt{time()}}}.

\begin{tcolorbox}[enhanced jigsaw,left=3pt,breakable,colback=white!97.2549!black,boxrule=0.2mm]

\begin{Shaded}
\begin{Highlighting}[]
\DataTypeTok{time\_t}\NormalTok{ now}\OperatorTok{;}  \CommentTok{// Variabel untuk menampung waktu sekarang}

\NormalTok{now }\OperatorTok{=}\NormalTok{ time}\OperatorTok{(}\NormalTok{NULL}\OperatorTok{);}  \CommentTok{// Kamu bisa mengambilnya seperti ini...}

\NormalTok{time}\OperatorTok{(\&}\NormalTok{now}\OperatorTok{);}        \CommentTok{// ...atau ini. Sama saja dengan baris sebelumnya.}
\end{Highlighting}
\end{Shaded}

\end{tcolorbox}

Bagus! Kamu punya variabel yang menyimpan waktu sekarang.

\index{time() function@\texttt{time()} function|)} \index{ctime() function@\texttt{ctime()} function|(}

Lucunya, hanya ada satu cara portabel untuk mencetak isi \hl{\mbox{\texttt{time\_t}}}, yaitu fungsi \hl{\mbox{\texttt{ctime()}}} yang jarang dipakai, yang mencetak nilai dalam waktu lokal:

\begin{tcolorbox}[enhanced jigsaw,left=3pt,breakable,colback=white!97.2549!black,boxrule=0.2mm]

\begin{Shaded}
\begin{Highlighting}[]
\NormalTok{now }\OperatorTok{=}\NormalTok{ time}\OperatorTok{(}\NormalTok{NULL}\OperatorTok{);}
\NormalTok{printf}\OperatorTok{(}\StringTok{"}\SpecialCharTok{\%s}\StringTok{"}\OperatorTok{,}\NormalTok{ ctime}\OperatorTok{(\&}\NormalTok{now}\OperatorTok{));}
\end{Highlighting}
\end{Shaded}

\end{tcolorbox}

Ini mengembalikan string dengan bentuk yang sangat spesifik, termasuk newline di akhir:

\begin{tcolorbox}[enhanced jigsaw,left=3pt,breakable,colback=white!97.2549!black,boxrule=0.2mm]

\begin{Shaded}
\begin{Highlighting}[]
\NormalTok{Sun Feb 28 18:47:25 2021}
\end{Highlighting}
\end{Shaded}

\end{tcolorbox}

\index{ctime() function@\texttt{ctime()} function|)}

Jadi itu agak kaku. Kalau ingin kendali lebih besar, kamu harus mengonversi \hl{\mbox{\texttt{time\_t}}} itu menjadi \hl{\mbox{\texttt{struct\ tm}}}.

\subsection{\texorpdfstring{Mengonversi \hl{\mbox{\texttt{time\_t}}} ke \hl{\mbox{\texttt{struct\ tm}}}}{Mengonversi time\_t ke struct tm}}\label{mengonversi-time_t-ke-struct-tm}

\index{time\_t type@\texttt{time\_t} type!conversion to struct tm@conversion to \texttt{struct tm}|(}

Ada dua cara luar biasa untuk melakukan konversi ini:

\index{localtime() function@\texttt{localtime()} function|(}

\begin{itemize}
\tightlist
\item
  \hl{\mbox{\texttt{localtime()}}}: fungsi ini mengonversi \hl{\mbox{\texttt{time\_t}}} menjadi \hl{\mbox{\texttt{struct\ tm}}} dalam waktu lokal.
\end{itemize}

\index{gmtime() function@\texttt{gmtime()} function|(}

\begin{itemize}
\tightlist
\item
  \hl{\mbox{\texttt{gmtime()}}}: fungsi ini mengonversi \hl{\mbox{\texttt{time\_t}}} menjadi \hl{\mbox{\texttt{struct\ tm}}} dalam UTC. (Lihat GMT kuno menyelinap ke nama fungsi ini?)
\end{itemize}

\index{asctime() function@\texttt{asctime()} function|(}

Mari kita lihat jam berapa sekarang dengan mencetak \hl{\mbox{\texttt{struct\ tm}}} memakai fungsi \hl{\mbox{\texttt{asctime()}}}:

\begin{tcolorbox}[enhanced jigsaw,left=3pt,breakable,colback=white!97.2549!black,boxrule=0.2mm]

\begin{Shaded}
\begin{Highlighting}[]
\NormalTok{printf}\OperatorTok{(}\StringTok{"Local: }\SpecialCharTok{\%s}\StringTok{"}\OperatorTok{,}\NormalTok{ asctime}\OperatorTok{(}\NormalTok{localtime}\OperatorTok{(\&}\NormalTok{now}\OperatorTok{)));}
\NormalTok{printf}\OperatorTok{(}\StringTok{"  UTC: }\SpecialCharTok{\%s}\StringTok{"}\OperatorTok{,}\NormalTok{ asctime}\OperatorTok{(}\NormalTok{gmtime}\OperatorTok{(\&}\NormalTok{now}\OperatorTok{)));}
\end{Highlighting}
\end{Shaded}

\end{tcolorbox}

\index{asctime() function@\texttt{asctime()} function|)} \index{localtime() function@\texttt{localtime()} function|)} \index{gmtime() function@\texttt{gmtime()} function|)}

Output (aku berada di zona waktu Pacific Standard Time):

\begin{tcolorbox}[enhanced jigsaw,left=3pt,breakable,colback=white!97.2549!black,boxrule=0.2mm]

\begin{Shaded}
\begin{Highlighting}[]
\NormalTok{Local: Sun Feb 28 20:15:27 2021}
\NormalTok{  UTC: Mon Mar  1 04:15:27 2021}
\end{Highlighting}
\end{Shaded}

\end{tcolorbox}

Begitu \hl{\mbox{\texttt{time\_t}}}-mu ada di \hl{\mbox{\texttt{struct\ tm}}}, banyak pintu terbuka. Kamu bisa mencetak waktu dengan berbagai cara, mencari tahu hari apa dalam seminggu sebuah tanggal, dan sebagainya. Atau mengonversinya kembali menjadi \hl{\mbox{\texttt{time\_t}}}.

Lebih lanjut soal itu sebentar lagi!

\index{time\_t type@\texttt{time\_t} type!conversion to struct tm@conversion to \texttt{struct tm}|)}

\subsection{\texorpdfstring{Mengonversi \hl{\mbox{\texttt{struct\ tm}}} ke \hl{\mbox{\texttt{time\_t}}}}{Mengonversi struct tm ke time\_t}}\label{mengonversi-struct-tm-ke-time_t}

\index{struct tm type@\texttt{struct tm} type!conversion to time\_t@conversion to \texttt{time\_t}|(} \index{mktime() function@\texttt{mktime()} function|(}

Kalau ingin ke arah sebaliknya, kamu bisa memakai \hl{\mbox{\texttt{mktime()}}} untuk mendapatkan informasi itu.

\hl{\mbox{\texttt{mktime()}}} mengatur nilai \hl{\mbox{\texttt{tm\_wday}}} dan \hl{\mbox{\texttt{tm\_yday}}} untukmu, jadi tidak perlu repot mengisinya karena akan tertimpa.

Selain itu, kamu bisa mengatur \hl{\mbox{\texttt{tm\_isdst}}} menjadi \hl{\mbox{\texttt{-1}}} agar ia yang menentukannya. Atau mengaturnya secara manual menjadi true atau false.

\begin{tcolorbox}[enhanced jigsaw,left=3pt,breakable,colback=white!97.2549!black,boxrule=0.2mm]

\begin{Shaded}
\begin{Highlighting}[]
\CommentTok{// Jangan tergoda menaruh nol di depan angka{-}angka ini (kecuali kamu}
\CommentTok{// memang bermaksud memakai oktal)!}

\KeywordTok{struct}\NormalTok{ tm some\_time }\OperatorTok{=} \OperatorTok{\{}
    \OperatorTok{.}\NormalTok{tm\_year}\OperatorTok{=}\DecValTok{82}\OperatorTok{,}   \CommentTok{// tahun sejak 1900}
    \OperatorTok{.}\NormalTok{tm\_mon}\OperatorTok{=}\DecValTok{3}\OperatorTok{,}     \CommentTok{// bulan sejak Januari {-}{-} [0, 11]}
    \OperatorTok{.}\NormalTok{tm\_mday}\OperatorTok{=}\DecValTok{12}\OperatorTok{,}   \CommentTok{// hari dalam bulan {-}{-} [1, 31]}
    \OperatorTok{.}\NormalTok{tm\_hour}\OperatorTok{=}\DecValTok{12}\OperatorTok{,}   \CommentTok{// jam sejak tengah malam {-}{-} [0, 23]}
    \OperatorTok{.}\NormalTok{tm\_min}\OperatorTok{=}\DecValTok{0}\OperatorTok{,}     \CommentTok{// menit setelah jam {-}{-} [0, 59]}
    \OperatorTok{.}\NormalTok{tm\_sec}\OperatorTok{=}\DecValTok{4}\OperatorTok{,}     \CommentTok{// detik setelah menit {-}{-} [0, 60]}
    \OperatorTok{.}\NormalTok{tm\_isdst}\OperatorTok{={-}}\DecValTok{1}\OperatorTok{,}  \CommentTok{// Flag Daylight Saving Time}
\OperatorTok{\};}

\DataTypeTok{time\_t}\NormalTok{ some\_time\_epoch}\OperatorTok{;}

\NormalTok{some\_time\_epoch }\OperatorTok{=}\NormalTok{ mktime}\OperatorTok{(\&}\NormalTok{some\_time}\OperatorTok{);}

\NormalTok{printf}\OperatorTok{(}\StringTok{"}\SpecialCharTok{\%s}\StringTok{"}\OperatorTok{,}\NormalTok{ ctime}\OperatorTok{(\&}\NormalTok{some\_time\_epoch}\OperatorTok{));}
\NormalTok{printf}\OperatorTok{(}\StringTok{"Is DST: }\SpecialCharTok{\%d\textbackslash{}n}\StringTok{"}\OperatorTok{,}\NormalTok{ some\_time}\OperatorTok{.}\NormalTok{tm\_isdst}\OperatorTok{);}
\end{Highlighting}
\end{Shaded}

\end{tcolorbox}

Output:

\begin{tcolorbox}[enhanced jigsaw,left=3pt,breakable,colback=white!97.2549!black,boxrule=0.2mm]

\begin{Shaded}
\begin{Highlighting}[]
\NormalTok{Mon Apr 12 12:00:04 1982}
\NormalTok{Is DST: 0}
\end{Highlighting}
\end{Shaded}

\end{tcolorbox}

Saat memuat \hl{\mbox{\texttt{struct\ tm}}} secara manual seperti itu, nilainya seharusnya dalam waktu lokal. \hl{\mbox{\texttt{mktime()}}} akan mengonversi waktu lokal itu menjadi calendar time \hl{\mbox{\texttt{time\_t}}}.

\index{mktime() function@\texttt{mktime()} function|)}

Anehnya, standar tidak memberi kita cara untuk memuat \hl{\mbox{\texttt{struct\ tm}}} dengan waktu UTC lalu mengonversinya menjadi \hl{\mbox{\texttt{time\_t}}}. Kalau ingin melakukannya di sistem mirip Unix, coba \index{timegm() Unix function@\texttt{timegm()} Unix function} \hl{\mbox{\texttt{timegm()}}} yang non-standar. Di Windows, \index{\_mkgmtime() Windows function@\texttt{\_mkgmtime()} Windows function} \hl{\mbox{\texttt{\_mkgmtime()}}}.

\index{struct tm type@\texttt{struct tm} type!conversion to time\_t@conversion to \texttt{time\_t}|)}

\section{Output Tanggal Terformat}\label{output-tanggal-terformat}

Kita sudah melihat beberapa cara mencetak output tanggal terformat ke layar. Dengan \hl{\mbox{\texttt{time\_t}}} kita bisa memakai \hl{\mbox{\texttt{ctime()}}}, dan dengan \hl{\mbox{\texttt{struct\ tm}}} kita bisa memakai \hl{\mbox{\texttt{asctime()}}}.

\begin{tcolorbox}[enhanced jigsaw,left=3pt,breakable,colback=white!97.2549!black,boxrule=0.2mm]

\begin{Shaded}
\begin{Highlighting}[]
\DataTypeTok{time\_t}\NormalTok{ now }\OperatorTok{=}\NormalTok{ time}\OperatorTok{(}\NormalTok{NULL}\OperatorTok{);}
\KeywordTok{struct}\NormalTok{ tm }\OperatorTok{*}\NormalTok{local }\OperatorTok{=}\NormalTok{ localtime}\OperatorTok{(\&}\NormalTok{now}\OperatorTok{);}
\KeywordTok{struct}\NormalTok{ tm }\OperatorTok{*}\NormalTok{utc }\OperatorTok{=}\NormalTok{ gmtime}\OperatorTok{(\&}\NormalTok{now}\OperatorTok{);}

\NormalTok{printf}\OperatorTok{(}\StringTok{"Local time: }\SpecialCharTok{\%s}\StringTok{"}\OperatorTok{,}\NormalTok{ ctime}\OperatorTok{(\&}\NormalTok{now}\OperatorTok{));}     \CommentTok{// Waktu lokal dengan time\_t}
\NormalTok{printf}\OperatorTok{(}\StringTok{"Local time: }\SpecialCharTok{\%s}\StringTok{"}\OperatorTok{,}\NormalTok{ asctime}\OperatorTok{(}\NormalTok{local}\OperatorTok{));}  \CommentTok{// Waktu lokal dengan struct tm}
\NormalTok{printf}\OperatorTok{(}\StringTok{"UTC       : }\SpecialCharTok{\%s}\StringTok{"}\OperatorTok{,}\NormalTok{ asctime}\OperatorTok{(}\NormalTok{utc}\OperatorTok{));}    \CommentTok{// UTC dengan struct tm}
\end{Highlighting}
\end{Shaded}

\end{tcolorbox}

Tapi bagaimana kalau kukatakan kepadamu, pembaca yang budiman, bahwa ada cara untuk punya kendali yang jauh lebih besar atas cara tanggal dicetak?

\index{strftime() function@\texttt{strftime()} function|(}

Memang kita bisa memancing field individual dari \hl{\mbox{\texttt{struct\ tm}}}, tapi ada fungsi hebat bernama \hl{\mbox{\texttt{strftime()}}} yang mengerjakan banyak pekerjaan berat untukmu. Ia seperti \hl{\mbox{\texttt{printf()}}}, hanya untuk tanggal!

Mari kita lihat beberapa contoh. Pada masing-masing, kita mengoper buffer tujuan, jumlah karakter maksimum yang ditulis, lalu format string (bergaya ---tapi tidak sama dengan---\hl{\mbox{\texttt{printf()}}}) yang memberi tahu \hl{\mbox{\texttt{strftime()}}} komponen \hl{\mbox{\texttt{struct\ tm}}} mana yang dicetak dan bagaimana.

Kamu juga bisa menambahkan karakter konstan lain untuk disertakan dalam output pada format string, seperti pada \hl{\mbox{\texttt{printf()}}}.

Pada kasus ini kita mendapat \hl{\mbox{\texttt{struct\ tm}}} dari \hl{\mbox{\texttt{localtime()}}}, tapi sumber mana pun bisa dipakai.

\begin{tcolorbox}[enhanced jigsaw,left=3pt,breakable,colback=white!97.2549!black,boxrule=0.2mm]

\begin{Shaded}
\begin{Highlighting}[numbers=left,,]
\PreprocessorTok{\#include }\ImportTok{\textless{}stdio.h\textgreater{}}
\PreprocessorTok{\#include }\ImportTok{\textless{}time.h\textgreater{}}

\DataTypeTok{int}\NormalTok{ main}\OperatorTok{(}\DataTypeTok{void}\OperatorTok{)}
\OperatorTok{\{}
    \DataTypeTok{char}\NormalTok{ s}\OperatorTok{[}\DecValTok{128}\OperatorTok{];}
    \DataTypeTok{time\_t}\NormalTok{ now }\OperatorTok{=}\NormalTok{ time}\OperatorTok{(}\NormalTok{NULL}\OperatorTok{);}

    \CommentTok{// \%c: cetak tanggal sesuai locale saat ini}
\NormalTok{    strftime}\OperatorTok{(}\NormalTok{s}\OperatorTok{,} \KeywordTok{sizeof}\NormalTok{ s}\OperatorTok{,} \StringTok{"}\SpecialCharTok{\%c}\StringTok{"}\OperatorTok{,}\NormalTok{ localtime}\OperatorTok{(\&}\NormalTok{now}\OperatorTok{));}
\NormalTok{    puts}\OperatorTok{(}\NormalTok{s}\OperatorTok{);}   \CommentTok{// Sun Feb 28 22:29:00 2021}

    \CommentTok{// \%A: nama lengkap hari}
    \CommentTok{// \%B: nama lengkap bulan}
    \CommentTok{// \%d: hari dalam bulan}
\NormalTok{    strftime}\OperatorTok{(}\NormalTok{s}\OperatorTok{,} \KeywordTok{sizeof}\NormalTok{ s}\OperatorTok{,} \StringTok{"}\SpecialCharTok{\%A}\StringTok{, \%B }\SpecialCharTok{\%d}\StringTok{"}\OperatorTok{,}\NormalTok{ localtime}\OperatorTok{(\&}\NormalTok{now}\OperatorTok{));}
\NormalTok{    puts}\OperatorTok{(}\NormalTok{s}\OperatorTok{);}   \CommentTok{// Sunday, February 28}

    \CommentTok{// \%I: jam (jam 12{-}an)}
    \CommentTok{// \%M: menit}
    \CommentTok{// \%S: detik}
    \CommentTok{// \%p: AM atau PM}
\NormalTok{    strftime}\OperatorTok{(}\NormalTok{s}\OperatorTok{,} \KeywordTok{sizeof}\NormalTok{ s}\OperatorTok{,} \StringTok{"It\textquotesingle{}s \%I:\%M:\%S }\SpecialCharTok{\%p}\StringTok{"}\OperatorTok{,}\NormalTok{ localtime}\OperatorTok{(\&}\NormalTok{now}\OperatorTok{));}
\NormalTok{    puts}\OperatorTok{(}\NormalTok{s}\OperatorTok{);}   \CommentTok{// It\textquotesingle{}s 10:29:00 PM}

    \CommentTok{// \%F: ISO 8601 yyyy{-}mm{-}dd}
    \CommentTok{// \%T: ISO 8601 hh:mm:ss}
    \CommentTok{// \%z: selisih zona waktu ISO 8601}
\NormalTok{    strftime}\OperatorTok{(}\NormalTok{s}\OperatorTok{,} \KeywordTok{sizeof}\NormalTok{ s}\OperatorTok{,} \StringTok{"ISO 8601: }\SpecialCharTok{\%F}\StringTok{T\%T\%z"}\OperatorTok{,}\NormalTok{ localtime}\OperatorTok{(\&}\NormalTok{now}\OperatorTok{));}
\NormalTok{    puts}\OperatorTok{(}\NormalTok{s}\OperatorTok{);}   \CommentTok{// ISO 8601: 2021{-}02{-}28T22:29:00{-}0800}
\OperatorTok{\}}
\end{Highlighting}
\end{Shaded}

\end{tcolorbox}

Ada \emph{banyak sekali} format specifier pencetakan tanggal untuk \hl{\mbox{\texttt{strftime()}}}, jadi pastikan melihatnya di \href{https://beej.us/guide/bgclr/html/split/time.html\#man-strftime}{halaman referensi \hl{\mbox{\texttt{strftime()}}}}\footnote{https://beej.us/guide/bgclr/html/split/time.html\#man-strftime}.

\index{strftime() function@\texttt{strftime()} function|)}

\section{\texorpdfstring{Resolusi Lebih Tinggi dengan \hl{\mbox{\texttt{timespec\_get()}}}}{Resolusi Lebih Tinggi dengan timespec\_get()}}\label{resolusi-lebih-tinggi-dengan-timespec_get}

\index{timespec\_get() function@\texttt{timespec\_get()} function|(}

Kamu bisa mendapat jumlah detik dan nanodetik sejak Epoch dengan \hl{\mbox{\texttt{timespec\_get()}}}.

Mungkin.

Implementasi mungkin tidak punya resolusi nanodetik (itu sepersemiliar detik) jadi tidak ada yang tahu berapa banyak tempat signifikan yang kamu dapat, tapi coba saja dan lihat.

\index{struct timespec type@\texttt{struct timespec} type|(}

\hl{\mbox{\texttt{timespec\_get()}}} menerima dua argumen. Satu adalah pointer ke \hl{\mbox{\texttt{struct\ timespec}}} untuk menampung informasi waktu. Yang lain adalah \hl{\mbox{\texttt{base}}}, yang oleh spesifikasi boleh diatur ke \hl{\mbox{\texttt{TIME\_UTC}}} untuk menandakan bahwa kamu tertarik pada detik sejak Epoch. (Implementasi lain mungkin memberi lebih banyak pilihan untuk \hl{\mbox{\texttt{base}}}.)

Dan struktur itu sendiri punya dua field:

\begin{tcolorbox}[enhanced jigsaw,left=3pt,breakable,colback=white!97.2549!black,boxrule=0.2mm]

\begin{Shaded}
\begin{Highlighting}[]
\KeywordTok{struct}\NormalTok{ timespec }\OperatorTok{\{}
    \DataTypeTok{time\_t}\NormalTok{ tv\_sec}\OperatorTok{;}   \CommentTok{// Detik}
    \DataTypeTok{long}\NormalTok{   tv\_nsec}\OperatorTok{;}  \CommentTok{// Nanodetik (sepersemiliar detik)}
\OperatorTok{\};}
\end{Highlighting}
\end{Shaded}

\end{tcolorbox}

Berikut contoh ketika kita mengambil waktu dan mencetaknya sebagai nilai integer dan juga nilai floating:

\begin{tcolorbox}[enhanced jigsaw,left=3pt,breakable,colback=white!97.2549!black,boxrule=0.2mm]

\begin{Shaded}
\begin{Highlighting}[]
\KeywordTok{struct}\NormalTok{ timespec ts}\OperatorTok{;}

\NormalTok{timespec\_get}\OperatorTok{(\&}\NormalTok{ts}\OperatorTok{,}\NormalTok{ TIME\_UTC}\OperatorTok{);}

\NormalTok{printf}\OperatorTok{(}\StringTok{"}\SpecialCharTok{\%ld}\StringTok{ s, }\SpecialCharTok{\%ld}\StringTok{ ns}\SpecialCharTok{\textbackslash{}n}\StringTok{"}\OperatorTok{,}\NormalTok{ ts}\OperatorTok{.}\NormalTok{tv\_sec}\OperatorTok{,}\NormalTok{ ts}\OperatorTok{.}\NormalTok{tv\_nsec}\OperatorTok{);}

\DataTypeTok{double}\NormalTok{ float\_time }\OperatorTok{=}\NormalTok{ ts}\OperatorTok{.}\NormalTok{tv\_sec }\OperatorTok{+}\NormalTok{ ts}\OperatorTok{.}\NormalTok{tv\_nsec}\OperatorTok{/}\FloatTok{1000000000.0}\OperatorTok{;}
\NormalTok{printf}\OperatorTok{(}\StringTok{"}\SpecialCharTok{\%f}\StringTok{ seconds since epoch}\SpecialCharTok{\textbackslash{}n}\StringTok{"}\OperatorTok{,}\NormalTok{ float\_time}\OperatorTok{);}
\end{Highlighting}
\end{Shaded}

\end{tcolorbox}

Contoh output:

\begin{tcolorbox}[enhanced jigsaw,left=3pt,breakable,colback=white!97.2549!black,boxrule=0.2mm]

\begin{Shaded}
\begin{Highlighting}[]
\NormalTok{1614581530 s, 806325800 ns}
\NormalTok{1614581530.806326 seconds since epoch}
\end{Highlighting}
\end{Shaded}

\end{tcolorbox}

\hl{\mbox{\texttt{struct\ timespec}}} juga muncul di sejumlah fungsi threading yang perlu bisa menentukan waktu dengan resolusi itu.

\index{struct timespec type@\texttt{struct timespec} type|)} \index{timespec\_get() function@\texttt{timespec\_get()} function|)}

\section{Selisih Antarwaktu}\label{selisih-antarwaktu}

\index{Date and time!differences|(}

Satu catatan singkat soal mengambil selisih antara dua \hl{\mbox{\texttt{time\_t}}}: karena spesifikasi tidak menentukan bagaimana tipe itu merepresentasikan waktu, kamu mungkin tidak bisa begitu saja mengurangkan dua \hl{\mbox{\texttt{time\_t}}} dan mendapat sesuatu yang masuk akal\footnote{Di POSIX bisa, karena \hl{\mbox{\texttt{time\_t}}} jelas integer. Sayangnya seluruh dunia bukan POSIX, jadi begitulah.}.

\index{difftime() function@\texttt{difftime()} function|(}

Untungnya kamu bisa memakai \hl{\mbox{\texttt{difftime()}}} untuk menghitung selisih dalam detik antara dua tanggal.

Pada contoh berikut, kita punya dua kejadian yang terjadi pada waktu yang berjauhan, dan kita memakai \hl{\mbox{\texttt{difftime()}}} untuk menghitung selisihnya.

\begin{tcolorbox}[enhanced jigsaw,left=3pt,breakable,colback=white!97.2549!black,boxrule=0.2mm]

\begin{Shaded}
\begin{Highlighting}[numbers=left,,]
\PreprocessorTok{\#include }\ImportTok{\textless{}stdio.h\textgreater{}}
\PreprocessorTok{\#include }\ImportTok{\textless{}time.h\textgreater{}}

\DataTypeTok{int}\NormalTok{ main}\OperatorTok{(}\DataTypeTok{void}\OperatorTok{)}
\OperatorTok{\{}
    \KeywordTok{struct}\NormalTok{ tm time\_a }\OperatorTok{=} \OperatorTok{\{}
        \OperatorTok{.}\NormalTok{tm\_year}\OperatorTok{=}\DecValTok{82}\OperatorTok{,}   \CommentTok{// tahun sejak 1900}
        \OperatorTok{.}\NormalTok{tm\_mon}\OperatorTok{=}\DecValTok{3}\OperatorTok{,}     \CommentTok{// bulan sejak Januari {-}{-} [0, 11]}
        \OperatorTok{.}\NormalTok{tm\_mday}\OperatorTok{=}\DecValTok{12}\OperatorTok{,}   \CommentTok{// hari dalam bulan {-}{-} [1, 31]}
        \OperatorTok{.}\NormalTok{tm\_hour}\OperatorTok{=}\DecValTok{4}\OperatorTok{,}    \CommentTok{// jam sejak tengah malam {-}{-} [0, 23]}
        \OperatorTok{.}\NormalTok{tm\_min}\OperatorTok{=}\BaseNTok{00}\OperatorTok{,}    \CommentTok{// menit setelah jam {-}{-} [0, 59]}
        \OperatorTok{.}\NormalTok{tm\_sec}\OperatorTok{=}\BaseNTok{04}\OperatorTok{,}    \CommentTok{// detik setelah menit {-}{-} [0, 60]}
        \OperatorTok{.}\NormalTok{tm\_isdst}\OperatorTok{={-}}\DecValTok{1}\OperatorTok{,}  \CommentTok{// Flag Daylight Saving Time}
    \OperatorTok{\};}

    \KeywordTok{struct}\NormalTok{ tm time\_b }\OperatorTok{=} \OperatorTok{\{}
        \OperatorTok{.}\NormalTok{tm\_year}\OperatorTok{=}\DecValTok{120}\OperatorTok{,}  \CommentTok{// tahun sejak 1900}
        \OperatorTok{.}\NormalTok{tm\_mon}\OperatorTok{=}\DecValTok{10}\OperatorTok{,}    \CommentTok{// bulan sejak Januari {-}{-} [0, 11]}
        \OperatorTok{.}\NormalTok{tm\_mday}\OperatorTok{=}\DecValTok{15}\OperatorTok{,}   \CommentTok{// hari dalam bulan {-}{-} [1, 31]}
        \OperatorTok{.}\NormalTok{tm\_hour}\OperatorTok{=}\DecValTok{16}\OperatorTok{,}   \CommentTok{// jam sejak tengah malam {-}{-} [0, 23]}
        \OperatorTok{.}\NormalTok{tm\_min}\OperatorTok{=}\DecValTok{27}\OperatorTok{,}    \CommentTok{// menit setelah jam {-}{-} [0, 59]}
        \OperatorTok{.}\NormalTok{tm\_sec}\OperatorTok{=}\BaseNTok{00}\OperatorTok{,}    \CommentTok{// detik setelah menit {-}{-} [0, 60]}
        \OperatorTok{.}\NormalTok{tm\_isdst}\OperatorTok{={-}}\DecValTok{1}\OperatorTok{,}  \CommentTok{// Flag Daylight Saving Time}
    \OperatorTok{\};}

    \DataTypeTok{time\_t}\NormalTok{ cal\_a }\OperatorTok{=}\NormalTok{ mktime}\OperatorTok{(\&}\NormalTok{time\_a}\OperatorTok{);}
    \DataTypeTok{time\_t}\NormalTok{ cal\_b }\OperatorTok{=}\NormalTok{ mktime}\OperatorTok{(\&}\NormalTok{time\_b}\OperatorTok{);}

    \DataTypeTok{double}\NormalTok{ diff }\OperatorTok{=}\NormalTok{ difftime}\OperatorTok{(}\NormalTok{cal\_b}\OperatorTok{,}\NormalTok{ cal\_a}\OperatorTok{);}

    \DataTypeTok{double}\NormalTok{ years }\OperatorTok{=}\NormalTok{ diff }\OperatorTok{/} \DecValTok{60} \OperatorTok{/} \DecValTok{60} \OperatorTok{/} \DecValTok{24} \OperatorTok{/} \FloatTok{365.2425}\OperatorTok{;}  \CommentTok{// cukup dekat}

\NormalTok{    printf}\OperatorTok{(}\StringTok{"}\SpecialCharTok{\%f}\StringTok{ seconds (}\SpecialCharTok{\%f}\StringTok{ years) between events}\SpecialCharTok{\textbackslash{}n}\StringTok{"}\OperatorTok{,}\NormalTok{ diff}\OperatorTok{,}\NormalTok{ years}\OperatorTok{);}
\OperatorTok{\}}
\end{Highlighting}
\end{Shaded}

\end{tcolorbox}

Output:

\begin{tcolorbox}[enhanced jigsaw,left=3pt,breakable,colback=white!97.2549!black,boxrule=0.2mm]

\begin{Shaded}
\begin{Highlighting}[]
\NormalTok{1217996816.000000 seconds (38.596783 years) between events}
\end{Highlighting}
\end{Shaded}

\end{tcolorbox}

Nah, begitulah! Ingat untuk memakai \hl{\mbox{\texttt{difftime()}}} guna mengambil selisih waktu. Walaupun di sistem POSIX kamu bisa langsung mengurangkan, sekalian saja tetap portabel.

\index{difftime() function@\texttt{difftime()} function|)} \index{Date and time!differences|)} \index{Date and time|)}

\chapter{Multithreading}\label{multithreading}

\index{Multithreading|(}

C11 secara resmi memperkenalkan multithreading ke bahasa C. Rasanya sangat mirip dengan \href{https://en.wikipedia.org/wiki/POSIX_Threads}{POSIX threads}\footnote{https://en.wikipedia.org/wiki/POSIX\_Threads}, kalau kamu pernah memakainya.

Dan kalau belum, tidak apa-apa. Kita akan membahasnya.

Namun perlu dicatat, aku tidak bermaksud membuat panduan multithreading klasik yang lengkap\footnote{Aku sendiri lebih suka shared-nothing, dan kemampuanku dengan konstruksi multithreading klasik berkarat, paling tidak begitu.}; untuk itu kamu harus mencari buku lain yang sangat tebal. Maaf!

\index{\_\_STDC\_NO\_THREADS\_\_ macro@\texttt{\_\_STDC\_NO\_THREADS\_\_} macro|(}

Threading adalah fitur opsional. Kalau compiler C11+ mendefinisikan \hl{\mbox{\texttt{\_\_STDC\_NO\_THREADS\_\_}}}, thread \textbf{tidak} ada di library. Kenapa mereka memilih makna negatif pada makro itu, aku tidak mengerti, tapi begitulah adanya.

Kamu bisa mengujinya seperti ini:

\begin{tcolorbox}[enhanced jigsaw,left=3pt,breakable,colback=white!97.2549!black,boxrule=0.2mm]

\begin{Shaded}
\begin{Highlighting}[]
\PreprocessorTok{\#ifdef \_\_STDC\_NO\_THREADS\_\_}
\PreprocessorTok{\#error I need threads to build this program!}
\PreprocessorTok{\#endif}
\end{Highlighting}
\end{Shaded}

\end{tcolorbox}

\index{\_\_STDC\_NO\_THREADS\_\_ macro@\texttt{\_\_STDC\_NO\_THREADS\_\_} macro|)}

\index{gcc compiler@\texttt{gcc} compiler!with threads|(}

Selain itu, kamu mungkin perlu menentukan opsi linker tertentu saat membangun. Di sistem mirip Unix, coba tambahkan \hl{\mbox{\texttt{-lpthreads}}} di akhir command line untuk me-link library \hl{\mbox{\texttt{pthreads}}}\footnote{Ya, \hl{\mbox{\texttt{pthreads}}} dengan ``\hl{\mbox{\texttt{p}}}''. Itu singkatan POSIX threads, library yang banyak dipinjam C11 untuk implementasi thread-nya.}:

\begin{tcolorbox}[enhanced jigsaw,left=3pt,breakable,colback=white!97.2549!black,boxrule=0.2mm]

\begin{Shaded}
\begin{Highlighting}[]
\FunctionTok{gcc} \AttributeTok{{-}std}\OperatorTok{=}\NormalTok{c11 }\AttributeTok{{-}o}\NormalTok{ foo foo.c }\AttributeTok{{-}lpthreads}
\end{Highlighting}
\end{Shaded}

\end{tcolorbox}

\index{gcc compiler@\texttt{gcc} compiler!with threads|)}

Kalau kamu mendapat error linker di sistemmu, bisa jadi karena library yang sesuai tidak disertakan.

\section{Latar Belakang}\label{latar-belakang}

Thread adalah cara agar semua core CPU mengilap yang sudah kamu bayar bekerja untukmu dalam program yang sama.

Biasanya, program C hanya berjalan di satu core CPU. Tapi kalau kamu tahu cara membagi pekerjaan, kamu bisa memberikan potongannya ke sejumlah thread dan menyuruh mereka mengerjakannya secara bersamaan.

Walaupun spesifikasi tidak menyebutnya, di sistemmu kemungkinan besar C (atau OS atas perintahnya) akan berusaha menyeimbangkan thread di semua core CPU-mu.

Dan kalau thread-nya lebih banyak daripada core, tidak apa-apa. Hanya saja kamu tidak akan merasakan semua keuntungannya kalau semuanya berebut waktu CPU.

\section{Hal-Hal yang Bisa Kamu Lakukan}\label{hal-hal-yang-bisa-kamu-lakukan}

Kamu bisa membuat thread. Ia akan mulai menjalankan fungsi yang kamu tentukan. Parent thread yang melahirkannya juga tetap berjalan.

Dan kamu bisa menunggu thread selesai. Ini disebut \emph{joining}.

Atau kalau kamu tidak peduli kapan thread selesai dan tidak mau menunggu, kamu bisa \emph{melepasnya (detach)}.

Sebuah thread bisa secara eksplisit \emph{keluar}, atau secara implisit berhenti dengan return dari fungsi utamanya.

Sebuah thread juga bisa \emph{tidur} selama jangka waktu tertentu, tidak melakukan apa-apa selagi thread lain berjalan.

Program \hl{\mbox{\texttt{main()}}} juga sebuah thread.

Selain itu, kita punya thread local storage, mutex, dan conditional variable. Tapi soal itu nanti. Untuk sekarang kita lihat dasar-dasarnya saja.

\section{Data Race dan Standard Library}\label{data-race-dan-standard-library}

\index{Multithreading!and the standard library|(}

Sebagian fungsi di standard library (mis. \hl{\mbox{\texttt{asctime()}}} dan \hl{\mbox{\texttt{strtok()}}}) mengembalikan atau memakai elemen data \hl{\mbox{\texttt{static}}} yang tidak threadsafe. Tapi secara umum, kecuali dinyatakan lain, standard library berusaha agar threadsafe\footnote{Menurut §7.1.4¶5.}.

Tapi tetap waspada. Kalau sebuah fungsi standard library menyimpan state antarpemanggilan di variabel yang bukan milikmu, atau kalau sebuah fungsi mengembalikan pointer ke sesuatu yang tidak kamu oper, ia tidak threadsafe.

\index{Multithreading!and the standard library|)}

\section{Membuat dan Menunggu Thread}\label{membuat-dan-menunggu-thread}

Mari kita buat sesuatu!

Kita akan membuat beberapa thread (create) dan menunggu mereka selesai (join).

Tapi ada sedikit hal yang perlu dipahami dulu.

Setiap thread diidentifikasi oleh variabel opaque bertipe \index{thrd\_t type@\texttt{thrd\_t} type} \hl{\mbox{\texttt{thrd\_t}}}. Ini identifier unik per thread di programmu. Saat kamu membuat thread, ia diberi ID baru.

Saat membuat thread, kamu juga harus memberinya pointer ke fungsi yang akan dijalankan, dan pointer ke argumen yang dioper kepadanya (atau \hl{\mbox{\texttt{NULL}}} kalau tidak ada yang dioper).

Thread akan mulai mengeksekusi fungsi yang kamu tentukan.

Saat ingin menunggu sebuah thread selesai, kamu harus menyebutkan ID thread-nya supaya C tahu yang mana yang ditunggu.

Jadi gagasan dasarnya:

\begin{enumerate}
\def\labelenumi{\arabic{enumi}.}
\tightlist
\item
  Tulis fungsi yang berperan sebagai ``\hl{\mbox{\texttt{main}}}'' thread. Ini bukan \hl{\mbox{\texttt{main()}}} yang sebenarnya, tapi analog dengannya. Thread akan mulai berjalan di sana.
\item
  Dari thread utama, jalankan thread baru dengan \index{thrd\_create() function@\texttt{thrd\_create()} function}\hl{\mbox{\texttt{thrd\_create()}}}, dan oper pointer ke fungsi yang akan dijalankan.
\item
  Di fungsi itu, biarkan thread mengerjakan apa pun yang harus dikerjakannya.
\item
  Sementara itu, thread utama bisa melanjutkan apa pun yang harus dikerjakan \emph{olehnya}.
\item
  Ketika thread utama memutuskan, ia bisa menunggu child thread selesai dengan memanggil \index{thrd\_join() function@\texttt{thrd\_join()} function} \hl{\mbox{\texttt{thrd\_join()}}}. Pada umumnya kamu \textbf{harus} me-\hl{\mbox{\texttt{thrd\_join()}}} thread itu untuk membersihkan sisanya, kalau tidak memori akan bocor\footnote{Kecuali kamu \hl{\mbox{\texttt{thrd\_detach()}}}. Soal ini dibahas nanti.}
\end{enumerate}

\index{thrd\_create() function@\texttt{thrd\_create()} function|(} \index{thrd\_start\_t type@\texttt{thrd\_start\_t} type|(}

\hl{\mbox{\texttt{thrd\_create()}}} menerima pointer ke fungsi yang akan dijalankan, dan bertipe \hl{\mbox{\texttt{thrd\_start\_t}}}, yaitu \hl{\mbox{\texttt{int\ (*)(void\ *)}}}. Itu bahasa Yunani untuk ``pointer ke fungsi yang menerima \hl{\mbox{\texttt{void*}}} sebagai argumen, dan mengembalikan \hl{\mbox{\texttt{int}}}.''

\index{thrd\_join() function@\texttt{thrd\_join()} function|(}

Mari kita buat thread! Kita akan menjalankannya dari thread utama dengan \hl{\mbox{\texttt{thrd\_create()}}} untuk menjalankan sebuah fungsi, mengerjakan hal lain, lalu menunggunya selesai dengan \hl{\mbox{\texttt{thrd\_join()}}}. Aku menamai fungsi utama thread ini \hl{\mbox{\texttt{run()}}}, tapi kamu bisa menamainya apa saja asalkan tipenya cocok dengan \hl{\mbox{\texttt{thrd\_start\_t}}}.

\begin{tcolorbox}[enhanced jigsaw,left=3pt,breakable,colback=white!97.2549!black,boxrule=0.2mm]

\begin{Shaded}
\begin{Highlighting}[numbers=left,,]
\PreprocessorTok{\#include }\ImportTok{\textless{}stdio.h\textgreater{}}
\PreprocessorTok{\#include }\ImportTok{\textless{}threads.h\textgreater{}}

\CommentTok{// Ini fungsi yang akan dijalankan thread. Boleh diberi nama apa saja.}
\CommentTok{//}
\CommentTok{// arg adalah pointer argumen yang dioper ke \textasciigrave{}thrd\_create()\textasciigrave{}.}
\CommentTok{//}
\CommentTok{// Parent thread akan mendapat nilai kembaliannya dari \textasciigrave{}thrd\_join()\textasciigrave{}\textquotesingle{}}
\CommentTok{// nanti.}

\DataTypeTok{int}\NormalTok{ run}\OperatorTok{(}\DataTypeTok{void} \OperatorTok{*}\NormalTok{arg}\OperatorTok{)}
\OperatorTok{\{}
    \DataTypeTok{int} \OperatorTok{*}\NormalTok{a }\OperatorTok{=}\NormalTok{ arg}\OperatorTok{;}  \CommentTok{// Kita akan mengoper int* dari thrd\_create()}

\NormalTok{    printf}\OperatorTok{(}\StringTok{"THREAD: Running thread with arg }\SpecialCharTok{\%d\textbackslash{}n}\StringTok{"}\OperatorTok{,} \OperatorTok{*}\NormalTok{a}\OperatorTok{);}

    \ControlFlowTok{return} \DecValTok{12}\OperatorTok{;}  \CommentTok{// Nilai yang akan diambil thrd\_join() (12 dipilih acak)}
\OperatorTok{\}}

\DataTypeTok{int}\NormalTok{ main}\OperatorTok{(}\DataTypeTok{void}\OperatorTok{)}
\OperatorTok{\{}
\NormalTok{    thrd\_t t}\OperatorTok{;}  \CommentTok{// t akan menampung ID thread}
    \DataTypeTok{int}\NormalTok{ arg }\OperatorTok{=} \DecValTok{3490}\OperatorTok{;}

\NormalTok{    printf}\OperatorTok{(}\StringTok{"Launching a thread}\SpecialCharTok{\textbackslash{}n}\StringTok{"}\OperatorTok{);}

    \CommentTok{// Jalankan thread ke fungsi run(), dengan mengoper pointer ke 3490}
    \CommentTok{// sebagai argumen. ID thread juga disimpan di t:}

\NormalTok{    thrd\_create}\OperatorTok{(\&}\NormalTok{t}\OperatorTok{,}\NormalTok{ run}\OperatorTok{,} \OperatorTok{\&}\NormalTok{arg}\OperatorTok{);}

\NormalTok{    printf}\OperatorTok{(}\StringTok{"Doing other things while the thread runs}\SpecialCharTok{\textbackslash{}n}\StringTok{"}\OperatorTok{);}

\NormalTok{    printf}\OperatorTok{(}\StringTok{"Waiting for thread to complete...}\SpecialCharTok{\textbackslash{}n}\StringTok{"}\OperatorTok{);}

    \DataTypeTok{int}\NormalTok{ res}\OperatorTok{;}  \CommentTok{// Menampung nilai kembalian dari keluarnya thread}

    \CommentTok{// Tunggu di sini sampai thread selesai; simpan nilai kembalian}
    \CommentTok{// di res:}

\NormalTok{    thrd\_join}\OperatorTok{(}\NormalTok{t}\OperatorTok{,} \OperatorTok{\&}\NormalTok{res}\OperatorTok{);}

\NormalTok{    printf}\OperatorTok{(}\StringTok{"Thread exited with return value }\SpecialCharTok{\%d\textbackslash{}n}\StringTok{"}\OperatorTok{,}\NormalTok{ res}\OperatorTok{);}
\OperatorTok{\}}
\end{Highlighting}
\end{Shaded}

\end{tcolorbox}

\index{thrd\_start\_t type@\texttt{thrd\_start\_t} type|)}

Lihat bagaimana kita memakai \hl{\mbox{\texttt{thrd\_create()}}} di situ untuk memanggil fungsi \hl{\mbox{\texttt{run()}}}? Lalu kita mengerjakan hal lain di \hl{\mbox{\texttt{main()}}} dan kemudian berhenti dan menunggu thread selesai dengan \hl{\mbox{\texttt{thrd\_join()}}}.

\index{thrd\_create() function@\texttt{thrd\_create()} function|)} \index{thrd\_join() function@\texttt{thrd\_join()} function|)}

Contoh output (hasilmu mungkin berbeda):

\begin{tcolorbox}[enhanced jigsaw,left=3pt,breakable,colback=white!97.2549!black,boxrule=0.2mm]

\begin{Shaded}
\begin{Highlighting}[]
\NormalTok{Launching a thread}
\NormalTok{Doing other things while the thread runs}
\NormalTok{Waiting for thread to complete...}
\NormalTok{THREAD: Running thread with arg 3490}
\NormalTok{Thread exited with return value 12}
\end{Highlighting}
\end{Shaded}

\end{tcolorbox}

\hl{\mbox{\texttt{arg}}} yang dioper ke fungsi harus punya masa hidup yang cukup panjang agar thread bisa mengambilnya sebelum ia hilang. Selain itu, ia tidak boleh tertimpa oleh thread utama sebelum thread baru sempat memakainya.

Mari kita lihat contoh yang menjalankan 5 thread. Satu hal yang perlu dicatat di sini adalah cara kita memakai array \hl{\mbox{\texttt{thrd\_t}}} untuk melacak semua ID thread.

\index{thrd\_create() function@\texttt{thrd\_create()} function|(} \index{thrd\_join() function@\texttt{thrd\_join()} function|(}

\begin{tcolorbox}[enhanced jigsaw,left=3pt,breakable,colback=white!97.2549!black,boxrule=0.2mm]

\begin{Shaded}
\begin{Highlighting}[numbers=left,,]
\PreprocessorTok{\#include }\ImportTok{\textless{}stdio.h\textgreater{}}
\PreprocessorTok{\#include }\ImportTok{\textless{}threads.h\textgreater{}}

\DataTypeTok{int}\NormalTok{ run}\OperatorTok{(}\DataTypeTok{void} \OperatorTok{*}\NormalTok{arg}\OperatorTok{)}
\OperatorTok{\{}
    \DataTypeTok{int}\NormalTok{ i }\OperatorTok{=} \OperatorTok{*(}\DataTypeTok{int}\OperatorTok{*)}\NormalTok{arg}\OperatorTok{;}

\NormalTok{    printf}\OperatorTok{(}\StringTok{"THREAD }\SpecialCharTok{\%d}\StringTok{: running!}\SpecialCharTok{\textbackslash{}n}\StringTok{"}\OperatorTok{,}\NormalTok{ i}\OperatorTok{);}

    \ControlFlowTok{return}\NormalTok{ i}\OperatorTok{;}
\OperatorTok{\}}

\PreprocessorTok{\#define THREAD\_COUNT }\DecValTok{5}

\DataTypeTok{int}\NormalTok{ main}\OperatorTok{(}\DataTypeTok{void}\OperatorTok{)}
\OperatorTok{\{}
\NormalTok{    thrd\_t t}\OperatorTok{[}\NormalTok{THREAD\_COUNT}\OperatorTok{];}

    \DataTypeTok{int}\NormalTok{ i}\OperatorTok{;}

\NormalTok{    printf}\OperatorTok{(}\StringTok{"Launching threads...}\SpecialCharTok{\textbackslash{}n}\StringTok{"}\OperatorTok{);}
    \ControlFlowTok{for} \OperatorTok{(}\NormalTok{i }\OperatorTok{=} \DecValTok{0}\OperatorTok{;}\NormalTok{ i }\OperatorTok{\textless{}}\NormalTok{ THREAD\_COUNT}\OperatorTok{;}\NormalTok{ i}\OperatorTok{++)}

        \CommentTok{// PERHATIAN! Pada baris berikut, kita mengoper pointer ke i,}
        \CommentTok{// tapi setiap thread melihat pointer yang sama. Jadi mereka akan}
        \CommentTok{// mencetak hal{-}hal aneh saat i berubah nilai di sini di}
        \CommentTok{// thread utama! (Lebih lanjut di teks, di bawah.)}

\NormalTok{        thrd\_create}\OperatorTok{(}\NormalTok{t }\OperatorTok{+}\NormalTok{ i}\OperatorTok{,}\NormalTok{ run}\OperatorTok{,} \OperatorTok{\&}\NormalTok{i}\OperatorTok{);}

\NormalTok{    printf}\OperatorTok{(}\StringTok{"Doing other things while the thread runs...}\SpecialCharTok{\textbackslash{}n}\StringTok{"}\OperatorTok{);}
\NormalTok{    printf}\OperatorTok{(}\StringTok{"Waiting for thread to complete...}\SpecialCharTok{\textbackslash{}n}\StringTok{"}\OperatorTok{);}

    \ControlFlowTok{for} \OperatorTok{(}\DataTypeTok{int}\NormalTok{ i }\OperatorTok{=} \DecValTok{0}\OperatorTok{;}\NormalTok{ i }\OperatorTok{\textless{}}\NormalTok{ THREAD\_COUNT}\OperatorTok{;}\NormalTok{ i}\OperatorTok{++)} \OperatorTok{\{}
        \DataTypeTok{int}\NormalTok{ res}\OperatorTok{;}
\NormalTok{        thrd\_join}\OperatorTok{(}\NormalTok{t}\OperatorTok{[}\NormalTok{i}\OperatorTok{],} \OperatorTok{\&}\NormalTok{res}\OperatorTok{);}

\NormalTok{        printf}\OperatorTok{(}\StringTok{"Thread }\SpecialCharTok{\%d}\StringTok{ complete!}\SpecialCharTok{\textbackslash{}n}\StringTok{"}\OperatorTok{,}\NormalTok{ res}\OperatorTok{);}
    \OperatorTok{\}}

\NormalTok{    printf}\OperatorTok{(}\StringTok{"All threads complete!}\SpecialCharTok{\textbackslash{}n}\StringTok{"}\OperatorTok{);}
\OperatorTok{\}}
\end{Highlighting}
\end{Shaded}

\end{tcolorbox}

\index{thrd\_join() function@\texttt{thrd\_join()} function|)}

Saat aku menjalankan thread, aku menghitung \hl{\mbox{\texttt{i}}} naik dari 0 sampai 4. Lalu aku mengoper pointer ke \hl{\mbox{\texttt{i}}} ke \hl{\mbox{\texttt{thrd\_create()}}}. Pointer ini berakhir di rutin \hl{\mbox{\texttt{run()}}} tempat kita membuat salinannya.

\index{thrd\_create() function@\texttt{thrd\_create()} function|)}

Cukup sederhana? Berikut outputnya:

\begin{tcolorbox}[enhanced jigsaw,left=3pt,breakable,colback=white!97.2549!black,boxrule=0.2mm]

\begin{Shaded}
\begin{Highlighting}[]
\NormalTok{Launching threads...}
\NormalTok{THREAD 2: running!}
\NormalTok{THREAD 3: running!}
\NormalTok{THREAD 4: running!}
\NormalTok{THREAD 2: running!}
\NormalTok{Doing other things while the thread runs...}
\NormalTok{Waiting for thread to complete...}
\NormalTok{Thread 2 complete!}
\NormalTok{Thread 2 complete!}
\NormalTok{THREAD 5: running!}
\NormalTok{Thread 3 complete!}
\NormalTok{Thread 4 complete!}
\NormalTok{Thread 5 complete!}
\NormalTok{All threads complete!}
\end{Highlighting}
\end{Shaded}

\end{tcolorbox}

Haaah---? Di mana \hl{\mbox{\texttt{THREAD\ 0}}}? Dan kenapa ada \hl{\mbox{\texttt{THREAD\ 5}}} padahal jelas \hl{\mbox{\texttt{i}}} tidak pernah lebih dari \hl{\mbox{\texttt{4}}} saat kita memanggil \hl{\mbox{\texttt{thrd\_create()}}}? Dan ada dua \hl{\mbox{\texttt{THREAD\ 2}}}? Gila!

Ini sudah masuk ke wilayah seru \index{Multithreading!race conditions} \emph{race condition}. Thread utama mengubah \hl{\mbox{\texttt{i}}} sebelum thread sempat menyalinnya. Bahkan, \hl{\mbox{\texttt{i}}} sampai ke \hl{\mbox{\texttt{5}}} dan loop berakhir sebelum thread terakhir sempat menyalinnya.

Kita harus punya variabel per-thread yang bisa dirujuk agar bisa dioper sebagai \hl{\mbox{\texttt{arg}}}.

Kita bisa membuat array besar berisi variabel itu. Atau kita bisa \hl{\mbox{\texttt{malloc()}}} ruangnya (dan membebaskannya di suatu tempat---mungkin di dalam thread itu sendiri).

Kita coba:

\begin{tcolorbox}[enhanced jigsaw,left=3pt,breakable,colback=white!97.2549!black,boxrule=0.2mm]

\begin{Shaded}
\begin{Highlighting}[numbers=left,,]
\PreprocessorTok{\#include }\ImportTok{\textless{}stdio.h\textgreater{}}
\PreprocessorTok{\#include }\ImportTok{\textless{}stdlib.h\textgreater{}}
\PreprocessorTok{\#include }\ImportTok{\textless{}threads.h\textgreater{}}

\DataTypeTok{int}\NormalTok{ run}\OperatorTok{(}\DataTypeTok{void} \OperatorTok{*}\NormalTok{arg}\OperatorTok{)}
\OperatorTok{\{}
    \DataTypeTok{int}\NormalTok{ i }\OperatorTok{=} \OperatorTok{*(}\DataTypeTok{int}\OperatorTok{*)}\NormalTok{arg}\OperatorTok{;}  \CommentTok{// Salin arg}

\NormalTok{    free}\OperatorTok{(}\NormalTok{arg}\OperatorTok{);}  \CommentTok{// Sudah selesai dengan ini}

\NormalTok{    printf}\OperatorTok{(}\StringTok{"THREAD }\SpecialCharTok{\%d}\StringTok{: running!}\SpecialCharTok{\textbackslash{}n}\StringTok{"}\OperatorTok{,}\NormalTok{ i}\OperatorTok{);}

    \ControlFlowTok{return}\NormalTok{ i}\OperatorTok{;}
\OperatorTok{\}}

\PreprocessorTok{\#define THREAD\_COUNT }\DecValTok{5}

\DataTypeTok{int}\NormalTok{ main}\OperatorTok{(}\DataTypeTok{void}\OperatorTok{)}
\OperatorTok{\{}
\NormalTok{    thrd\_t t}\OperatorTok{[}\NormalTok{THREAD\_COUNT}\OperatorTok{];}

    \DataTypeTok{int}\NormalTok{ i}\OperatorTok{;}

\NormalTok{    printf}\OperatorTok{(}\StringTok{"Launching threads...}\SpecialCharTok{\textbackslash{}n}\StringTok{"}\OperatorTok{);}
    \ControlFlowTok{for} \OperatorTok{(}\NormalTok{i }\OperatorTok{=} \DecValTok{0}\OperatorTok{;}\NormalTok{ i }\OperatorTok{\textless{}}\NormalTok{ THREAD\_COUNT}\OperatorTok{;}\NormalTok{ i}\OperatorTok{++)} \OperatorTok{\{}

        \CommentTok{// Siapkan ruang untuk argumen per{-}thread:}

        \DataTypeTok{int} \OperatorTok{*}\NormalTok{arg }\OperatorTok{=}\NormalTok{ malloc}\OperatorTok{(}\KeywordTok{sizeof} \OperatorTok{*}\NormalTok{arg}\OperatorTok{);}
        \OperatorTok{*}\NormalTok{arg }\OperatorTok{=}\NormalTok{ i}\OperatorTok{;}

\NormalTok{        thrd\_create}\OperatorTok{(}\NormalTok{t }\OperatorTok{+}\NormalTok{ i}\OperatorTok{,}\NormalTok{ run}\OperatorTok{,}\NormalTok{ arg}\OperatorTok{);}
    \OperatorTok{\}}

    \CommentTok{// ...}
\end{Highlighting}
\end{Shaded}

\end{tcolorbox}

Perhatikan di baris 27-30 kita \hl{\mbox{\texttt{malloc()}}} ruang untuk \hl{\mbox{\texttt{int}}} dan menyalin nilai \hl{\mbox{\texttt{i}}} ke dalamnya. Setiap thread baru mendapat variabel hasil \hl{\mbox{\texttt{malloc()}}} yang masih baru untuk dirinya sendiri, dan kita mengoper pointer ke variabel itu ke fungsi \hl{\mbox{\texttt{run()}}}.

Setelah \hl{\mbox{\texttt{run()}}} membuat salinannya sendiri dari \hl{\mbox{\texttt{arg}}} di baris 7, ia men-\hl{\mbox{\texttt{free()}}} \hl{\mbox{\texttt{int}}} hasil \hl{\mbox{\texttt{malloc()}}} itu. Dan karena sudah punya salinan sendiri, ia bebas melakukan apa pun padanya.

Dan saat dijalankan, hasilnya:

\begin{tcolorbox}[enhanced jigsaw,left=3pt,breakable,colback=white!97.2549!black,boxrule=0.2mm]

\begin{Shaded}
\begin{Highlighting}[]
\NormalTok{Launching threads...}
\NormalTok{THREAD 0: running!}
\NormalTok{THREAD 1: running!}
\NormalTok{THREAD 2: running!}
\NormalTok{THREAD 3: running!}
\NormalTok{Doing other things while the thread runs...}
\NormalTok{Waiting for thread to complete...}
\NormalTok{Thread 0 complete!}
\NormalTok{Thread 1 complete!}
\NormalTok{Thread 2 complete!}
\NormalTok{Thread 3 complete!}
\NormalTok{THREAD 4: running!}
\NormalTok{Thread 4 complete!}
\NormalTok{All threads complete!}
\end{Highlighting}
\end{Shaded}

\end{tcolorbox}

Nah begitu! Thread 0-4 semuanya berjalan!

Hasil jalanmu mungkin berbeda---bagaimana thread dijadwalkan berada di luar spesifikasi C. Pada contoh di atas kita lihat thread 4 bahkan belum mulai sampai thread 0-1 selesai. Bahkan, kalau kujalankan lagi, kemungkinan besar aku mendapat output berbeda. Kita tidak bisa menjamin urutan eksekusi thread.

\section{Melepas Thread (Detach)}\label{melepas-thread-detach}

Kalau kamu ingin melepas thread begitu saja (fire-and-forget), yaitu supaya tidak perlu me-\hl{\mbox{\texttt{thrd\_join()}}}-nya nanti, kamu bisa melakukannya dengan \hl{\mbox{\texttt{thrd\_detach()}}}.

Ini menghilangkan kemampuan parent thread untuk mengambil nilai kembalian dari child thread, tapi kalau kamu tidak peduli soal itu dan hanya ingin thread membereskan dirinya sendiri dengan rapi, inilah caranya.

Pada dasarnya kita akan melakukan ini:

\index{thrd\_detach() function@\texttt{thrd\_detach()} function|(}

\begin{tcolorbox}[enhanced jigsaw,left=3pt,breakable,colback=white!97.2549!black,boxrule=0.2mm]

\begin{Shaded}
\begin{Highlighting}[]
\NormalTok{thrd\_create}\OperatorTok{(\&}\NormalTok{t}\OperatorTok{,}\NormalTok{ run}\OperatorTok{,}\NormalTok{ NULL}\OperatorTok{);}
\NormalTok{thrd\_detach}\OperatorTok{(}\NormalTok{t}\OperatorTok{);}
\end{Highlighting}
\end{Shaded}

\end{tcolorbox}

Pemanggilan \hl{\mbox{\texttt{thrd\_detach()}}} adalah parent thread yang berkata, ``Hei, aku tidak akan menunggu child thread ini selesai dengan \hl{\mbox{\texttt{thrd\_join()}}}. Jadi silakan bereskan sendiri saat selesai.''

\begin{tcolorbox}[enhanced jigsaw,left=3pt,breakable,colback=white!97.2549!black,boxrule=0.2mm]

\begin{Shaded}
\begin{Highlighting}[numbers=left,,]
\PreprocessorTok{\#include }\ImportTok{\textless{}stdio.h\textgreater{}}
\PreprocessorTok{\#include }\ImportTok{\textless{}threads.h\textgreater{}}

\DataTypeTok{int}\NormalTok{ run}\OperatorTok{(}\DataTypeTok{void} \OperatorTok{*}\NormalTok{arg}\OperatorTok{)}
\OperatorTok{\{}
    \OperatorTok{(}\DataTypeTok{void}\OperatorTok{)}\NormalTok{arg}\OperatorTok{;}

    \CommentTok{//printf("Thread running! \%lu\textbackslash{}n", thrd\_current()); // tidak portabel!}
\NormalTok{    printf}\OperatorTok{(}\StringTok{"Thread running!}\SpecialCharTok{\textbackslash{}n}\StringTok{"}\OperatorTok{);}

    \ControlFlowTok{return} \DecValTok{0}\OperatorTok{;}
\OperatorTok{\}}

\PreprocessorTok{\#define THREAD\_COUNT }\DecValTok{10}

\DataTypeTok{int}\NormalTok{ main}\OperatorTok{(}\DataTypeTok{void}\OperatorTok{)}
\OperatorTok{\{}
\NormalTok{    thrd\_t t}\OperatorTok{;}

    \ControlFlowTok{for} \OperatorTok{(}\DataTypeTok{int}\NormalTok{ i }\OperatorTok{=} \DecValTok{0}\OperatorTok{;}\NormalTok{ i }\OperatorTok{\textless{}}\NormalTok{ THREAD\_COUNT}\OperatorTok{;}\NormalTok{ i}\OperatorTok{++)} \OperatorTok{\{}
\NormalTok{        thrd\_create}\OperatorTok{(\&}\NormalTok{t}\OperatorTok{,}\NormalTok{ run}\OperatorTok{,}\NormalTok{ NULL}\OperatorTok{);}
\NormalTok{        thrd\_detach}\OperatorTok{(}\NormalTok{t}\OperatorTok{);}               \CommentTok{// \textless{}{-}{-} DETACH!}
    \OperatorTok{\}}

    \CommentTok{// Tidur selama satu detik supaya semua thread sempat selesai}
\NormalTok{    thrd\_sleep}\OperatorTok{(\&(}\KeywordTok{struct}\NormalTok{ timespec}\OperatorTok{)\{.}\NormalTok{tv\_sec}\OperatorTok{=}\DecValTok{1}\OperatorTok{\},}\NormalTok{ NULL}\OperatorTok{);}
\OperatorTok{\}}
\end{Highlighting}
\end{Shaded}

\end{tcolorbox}

\index{thrd\_detach() function@\texttt{thrd\_detach()} function|)}

Perhatikan bahwa pada kode ini, thread utama ditidurkan selama 1 detik dengan \hl{\mbox{\texttt{thrd\_sleep()}}}---soal itu dibahas nanti.

Di fungsi \hl{\mbox{\texttt{run()}}} aku juga menaruh baris yang dikomentari untuk mencetak ID thread sebagai \hl{\mbox{\texttt{unsigned\ long}}}. Ini tidak portabel, karena spesifikasi tidak menyebut tipe \hl{\mbox{\texttt{thrd\_t}}} sebenarnya apa---bisa saja ia \hl{\mbox{\texttt{struct}}}. Tapi baris itu berfungsi di sistemku.

Hal menarik yang kulihat saat menjalankan kode di atas dan mencetak ID thread: sebagian thread punya ID kembar! Ini seharusnya mustahil, tapi C boleh \emph{memakai ulang} ID thread setelah thread yang bersangkutan keluar. Jadi yang kulihat adalah sebagian thread selesai berjalan sebelum thread lain diluncurkan.

\section{Thread Local Data}\label{thread-local-data}

\index{Thread local data|(}

Thread itu menarik karena mereka tidak punya memori sendiri di luar variabel lokal. Kalau ada variabel \hl{\mbox{\texttt{static}}} atau variabel file scope, semua thread akan melihat variabel yang sama.

Ini bisa menyebabkan race condition, ketika \emph{Hal-Hal Aneh}™ terjadi.

Lihat contoh ini. Ada variabel \hl{\mbox{\texttt{static}}} \hl{\mbox{\texttt{foo}}} di block scope pada \hl{\mbox{\texttt{run()}}}. Variabel ini terlihat oleh semua thread yang melewati fungsi \hl{\mbox{\texttt{run()}}}. Dan berbagai thread secara efektif bisa saling menginjak kaki.

Setiap thread menyalin \hl{\mbox{\texttt{foo}}} ke variabel lokal \hl{\mbox{\texttt{x}}} (yang tidak dibagi antar thread---semua thread punya call stack sendiri). Jadi mereka \emph{seharusnya} sama, kan?

Dan pada pencetakan pertama, memang sama\footnote{Walaupun kurasa tidak harus begitu. Hanya saja thread tampaknya tidak dijadwalkan ulang sampai ada semacam system call, seperti yang mungkin terjadi pada \hl{\mbox{\texttt{printf()}}}\ldots{} itulah kenapa ada \hl{\mbox{\texttt{printf()}}} di situ.}. Tapi tepat setelah itu, kita periksa untuk memastikan keduanya masih sama.

Dan \emph{biasanya} memang sama. Tapi tidak selalu!

\begin{tcolorbox}[enhanced jigsaw,left=3pt,breakable,colback=white!97.2549!black,boxrule=0.2mm]

\begin{Shaded}
\begin{Highlighting}[numbers=left,,]
\PreprocessorTok{\#include }\ImportTok{\textless{}stdio.h\textgreater{}}
\PreprocessorTok{\#include }\ImportTok{\textless{}stdlib.h\textgreater{}}
\PreprocessorTok{\#include }\ImportTok{\textless{}threads.h\textgreater{}}

\DataTypeTok{int}\NormalTok{ run}\OperatorTok{(}\DataTypeTok{void} \OperatorTok{*}\NormalTok{arg}\OperatorTok{)}
\OperatorTok{\{}
    \DataTypeTok{int}\NormalTok{ n }\OperatorTok{=} \OperatorTok{*(}\DataTypeTok{int}\OperatorTok{*)}\NormalTok{arg}\OperatorTok{;}  \CommentTok{// Nomor thread agar manusia bisa membedakan}

\NormalTok{    free}\OperatorTok{(}\NormalTok{arg}\OperatorTok{);}

    \DataTypeTok{static} \DataTypeTok{int}\NormalTok{ foo }\OperatorTok{=} \DecValTok{10}\OperatorTok{;}  \CommentTok{// Nilai static yang dibagi antar thread}

    \DataTypeTok{int}\NormalTok{ x }\OperatorTok{=}\NormalTok{ foo}\OperatorTok{;}  \CommentTok{// Variabel lokal otomatis{-}{-}tiap thread punya sendiri}

    \CommentTok{// Kita baru saja menugaskan x dari foo, jadi mereka harus sama di sini.}
    \CommentTok{// (Di semua percobaanku memang sama, tapi bahkan ini tidak}
    \CommentTok{// dijamin!)}

\NormalTok{    printf}\OperatorTok{(}\StringTok{"Thread }\SpecialCharTok{\%d}\StringTok{: x = }\SpecialCharTok{\%d}\StringTok{, foo = }\SpecialCharTok{\%d\textbackslash{}n}\StringTok{"}\OperatorTok{,}\NormalTok{ n}\OperatorTok{,}\NormalTok{ x}\OperatorTok{,}\NormalTok{ foo}\OperatorTok{);}

    \CommentTok{// Dan mereka seharusnya sama di sini, tapi tidak selalu!}
    \CommentTok{// (Kadang sama, kadang tidak!)}

    \CommentTok{// Yang terjadi: thread lain masuk dan meng{-}increment foo}
    \CommentTok{// saat ini juga, tapi x milik thread ini tetap seperti sebelumnya!}

    \ControlFlowTok{if} \OperatorTok{(}\NormalTok{x }\OperatorTok{!=}\NormalTok{ foo}\OperatorTok{)} \OperatorTok{\{}
\NormalTok{        printf}\OperatorTok{(}\StringTok{"Thread }\SpecialCharTok{\%d}\StringTok{: Craziness! x != foo! }\SpecialCharTok{\%d}\StringTok{ != }\SpecialCharTok{\%d\textbackslash{}n}\StringTok{"}\OperatorTok{,}\NormalTok{ n}\OperatorTok{,}\NormalTok{ x}\OperatorTok{,}\NormalTok{ foo}\OperatorTok{);}
    \OperatorTok{\}}

\NormalTok{    foo}\OperatorTok{++;}  \CommentTok{// Increment nilai yang dibagi}

    \ControlFlowTok{return} \DecValTok{0}\OperatorTok{;}
\OperatorTok{\}}

\PreprocessorTok{\#define THREAD\_COUNT }\DecValTok{5}

\DataTypeTok{int}\NormalTok{ main}\OperatorTok{(}\DataTypeTok{void}\OperatorTok{)}
\OperatorTok{\{}
\NormalTok{    thrd\_t t}\OperatorTok{[}\NormalTok{THREAD\_COUNT}\OperatorTok{];}

    \ControlFlowTok{for} \OperatorTok{(}\DataTypeTok{int}\NormalTok{ i }\OperatorTok{=} \DecValTok{0}\OperatorTok{;}\NormalTok{ i }\OperatorTok{\textless{}}\NormalTok{ THREAD\_COUNT}\OperatorTok{;}\NormalTok{ i}\OperatorTok{++)} \OperatorTok{\{}
        \DataTypeTok{int} \OperatorTok{*}\NormalTok{n }\OperatorTok{=}\NormalTok{ malloc}\OperatorTok{(}\KeywordTok{sizeof} \OperatorTok{*}\NormalTok{n}\OperatorTok{);}  \CommentTok{// Menampung nomor seri thread}
        \OperatorTok{*}\NormalTok{n }\OperatorTok{=}\NormalTok{ i}\OperatorTok{;}
\NormalTok{        thrd\_create}\OperatorTok{(}\NormalTok{t }\OperatorTok{+}\NormalTok{ i}\OperatorTok{,}\NormalTok{ run}\OperatorTok{,}\NormalTok{ n}\OperatorTok{);}
    \OperatorTok{\}}

    \ControlFlowTok{for} \OperatorTok{(}\DataTypeTok{int}\NormalTok{ i }\OperatorTok{=} \DecValTok{0}\OperatorTok{;}\NormalTok{ i }\OperatorTok{\textless{}}\NormalTok{ THREAD\_COUNT}\OperatorTok{;}\NormalTok{ i}\OperatorTok{++)} \OperatorTok{\{}
\NormalTok{        thrd\_join}\OperatorTok{(}\NormalTok{t}\OperatorTok{[}\NormalTok{i}\OperatorTok{],}\NormalTok{ NULL}\OperatorTok{);}
    \OperatorTok{\}}
\OperatorTok{\}}
\end{Highlighting}
\end{Shaded}

\end{tcolorbox}

Berikut contoh outputnya (walau ini berbeda dari satu eksekusi ke eksekusi lain):

\begin{tcolorbox}[enhanced jigsaw,left=3pt,breakable,colback=white!97.2549!black,boxrule=0.2mm]

\begin{Shaded}
\begin{Highlighting}[]
\NormalTok{Thread 0: x = 10, foo = 10}
\NormalTok{Thread 1: x = 10, foo = 10}
\NormalTok{Thread 1: Craziness! x != foo! 10 != 11}
\NormalTok{Thread 2: x = 12, foo = 12}
\NormalTok{Thread 4: x = 13, foo = 13}
\NormalTok{Thread 3: x = 14, foo = 14}
\end{Highlighting}
\end{Shaded}

\end{tcolorbox}

Pada thread 1, di antara dua \hl{\mbox{\texttt{printf()}}}, nilai \hl{\mbox{\texttt{foo}}} entah bagaimana berubah dari \hl{\mbox{\texttt{10}}} menjadi \hl{\mbox{\texttt{11}}}, padahal jelas tidak ada increment di antara kedua \hl{\mbox{\texttt{printf()}}}!

Ada thread lain yang masuk ke sana (kemungkinan thread 0, kalau dilihat dari hasilnya) dan meng-increment nilai \hl{\mbox{\texttt{foo}}} di belakang punggung thread 1!

Mari kita selesaikan masalah ini dengan dua cara berbeda. (Kalau kamu ingin semua thread berbagi variabel \emph{dan} tidak saling menginjak kaki, kamu harus lanjut membaca bagian \hyperref[mutex]{mutex}.)

\subsection{\texorpdfstring{Storage-Class \hl{\mbox{\texttt{\_Thread\_local}}}}{Storage-Class \_Thread\_local}}\label{thread-local}

\index{\_Thread\_local storage class@\texttt{\_Thread\_local} storage class|(}

Pertama-tama, mari kita lihat cara mudah mengatasinya: storage-class \hl{\mbox{\texttt{\_Thread\_local}}}.

Pada dasarnya kita tinggal menempelkannya di depan variabel \hl{\mbox{\texttt{static}}} block scope kita dan semuanya berfungsi! Ia memberi tahu C bahwa setiap thread harus punya versinya sendiri dari variabel ini, sehingga tidak ada yang saling menginjak kaki.

\index{thread\_local storage class@\texttt{thread\_local} storage class|(}

Header \index{threads.h header file@\texttt{threads.h} header file} \hl{\mbox{\texttt{\textless{}threads.h\textgreater{}}}} mendefinisikan \hl{\mbox{\texttt{thread\_local}}} sebagai alias untuk \hl{\mbox{\texttt{\_Thread\_local}}} supaya kodemu tidak perlu tampak begitu jelek.

Mari kita ambil contoh sebelumnya dan jadikan \hl{\mbox{\texttt{foo}}} variabel \hl{\mbox{\texttt{thread\_local}}} agar data itu tidak dibagi.

\begin{tcolorbox}[enhanced jigsaw,left=3pt,breakable,colback=white!97.2549!black,boxrule=0.2mm]

\begin{Shaded}
\begin{Highlighting}[numbers=left,,firstnumber=5,]
\DataTypeTok{int}\NormalTok{ run}\OperatorTok{(}\DataTypeTok{void} \OperatorTok{*}\NormalTok{arg}\OperatorTok{)}
\OperatorTok{\{}
    \DataTypeTok{int}\NormalTok{ n }\OperatorTok{=} \OperatorTok{*(}\DataTypeTok{int}\OperatorTok{*)}\NormalTok{arg}\OperatorTok{;}  \CommentTok{// Nomor thread agar manusia bisa membedakan}

\NormalTok{    free}\OperatorTok{(}\NormalTok{arg}\OperatorTok{);}

    \DataTypeTok{thread\_local} \DataTypeTok{static} \DataTypeTok{int}\NormalTok{ foo }\OperatorTok{=} \DecValTok{10}\OperatorTok{;}  \CommentTok{// \textless{}{-}{-} Tidak lagi dibagi!!}
\end{Highlighting}
\end{Shaded}

\end{tcolorbox}

Dan saat dijalankan, kita mendapat:

\begin{tcolorbox}[enhanced jigsaw,left=3pt,breakable,colback=white!97.2549!black,boxrule=0.2mm]

\begin{Shaded}
\begin{Highlighting}[]
\NormalTok{Thread 0: x = 10, foo = 10}
\NormalTok{Thread 1: x = 10, foo = 10}
\NormalTok{Thread 2: x = 10, foo = 10}
\NormalTok{Thread 4: x = 10, foo = 10}
\NormalTok{Thread 3: x = 10, foo = 10}
\end{Highlighting}
\end{Shaded}

\end{tcolorbox}

Tidak ada lagi masalah aneh!

Satu hal: kalau variabel \hl{\mbox{\texttt{thread\_local}}} berada di block scope, ia \textbf{harus} \hl{\mbox{\texttt{static}}}. Begitulah aturannya. (Tapi ini tidak masalah karena variabel non-\hl{\mbox{\texttt{static}}} memang sudah per-thread, sebab tiap thread punya variabel non-\hl{\mbox{\texttt{static}}}-nya sendiri.)

Sedikit kebohongan di situ: variabel \hl{\mbox{\texttt{thread\_local}}} di block scope juga boleh \hl{\mbox{\texttt{extern}}}.

\index{thread\_local storage class@\texttt{thread\_local} storage class|)} \index{\_Thread\_local storage class@\texttt{\_Thread\_local} storage class|)} \index{Thread local data|)}

\subsection{Opsi Lain: Thread-Specific Storage}\label{opsi-lain-thread-specific-storage}

\index{Thread-specific storage|(}

Thread-specific storage (TSS) adalah cara lain untuk mendapatkan data per-thread.

Satu fitur tambahan: fungsi-fungsi ini memungkinkanmu menentukan destructor yang akan dipanggil pada data saat variabel TSS dihapus. Biasanya destructor ini adalah \hl{\mbox{\texttt{free()}}} untuk membersihkan data per-thread hasil \hl{\mbox{\texttt{malloc()}}} secara otomatis. Atau \hl{\mbox{\texttt{NULL}}} kalau tidak ada yang perlu dihancurkan.

Destructor bertipe \index{tss\_dtor\_t type@\texttt{tss\_dtor\_t} type} \hl{\mbox{\texttt{tss\_dtor\_t}}}, yaitu pointer ke fungsi yang mengembalikan \hl{\mbox{\texttt{void}}} dan menerima \hl{\mbox{\texttt{void*}}} sebagai argumen (\hl{\mbox{\texttt{void*}}} itu menunjuk ke data yang disimpan di variabel). Dengan kata lain, ia \hl{\mbox{\texttt{void\ (*)(void*)}}}, kalau itu membantu menjelaskan. Dan aku akui, kemungkinan itu tidak membantu. Lihat contoh di bawah.

Secara umum, \hl{\mbox{\texttt{thread\_local}}} kemungkinan pilihan utamamu, tapi kalau kamu menyukai gagasan destructor, kamu bisa memanfaatkannya.

Pemakaiannya agak aneh karena kita butuh variabel bertipe \index{tss\_t type@\texttt{tss\_t} type} \hl{\mbox{\texttt{tss\_t}}} yang tetap hidup untuk mewakili nilai per thread. Lalu kita menginisialisasinya dengan \index{tss\_create() function@\texttt{tss\_create()} function} \hl{\mbox{\texttt{tss\_create()}}}. Akhirnya kita membuangnya dengan \index{tss\_delete() function@\texttt{tss\_delete()} function|(} \hl{\mbox{\texttt{tss\_delete()}}}. Perhatikan bahwa memanggil \hl{\mbox{\texttt{tss\_delete()}}} tidak menjalankan semua destructor---yang melakukannya adalah \hl{\mbox{\texttt{thrd\_exit()}}} (atau return dari fungsi run). \hl{\mbox{\texttt{tss\_delete()}}} hanya melepas memori yang dialokasikan oleh \hl{\mbox{\texttt{tss\_create()}}}. \index{tss\_delete() function@\texttt{tss\_delete()} function|)}

Di tengah-tengahnya, thread bisa memanggil \index{tss\_set() function@\texttt{tss\_set()} function} \hl{\mbox{\texttt{tss\_set()}}} dan \index{tss\_get() function@\texttt{tss\_get()} function} \hl{\mbox{\texttt{tss\_get()}}} untuk mengisi dan mengambil nilai.

Pada kode berikut, kita menyiapkan variabel TSS sebelum membuat thread, lalu membersihkan setelah thread selesai.

Pada fungsi \hl{\mbox{\texttt{run()}}}, thread melakukan \hl{\mbox{\texttt{malloc()}}} ruang untuk string dan menyimpan pointer itu di variabel TSS.

Saat thread keluar, fungsi destructor (\hl{\mbox{\texttt{free()}}} dalam kasus ini) dipanggil untuk \emph{semua} thread.

\index{tss\_t type@\texttt{tss\_t} type|(} \index{tss\_get() function@\texttt{tss\_get()} function|(} \index{tss\_set() function@\texttt{tss\_set()} function|(} \index{tss\_create() function@\texttt{tss\_create()} function|(} \index{tss\_delete() function@\texttt{tss\_delete()} function|(}

\begin{tcolorbox}[enhanced jigsaw,left=3pt,breakable,colback=white!97.2549!black,boxrule=0.2mm]

\begin{Shaded}
\begin{Highlighting}[numbers=left,,]
\PreprocessorTok{\#include }\ImportTok{\textless{}stdio.h\textgreater{}}
\PreprocessorTok{\#include }\ImportTok{\textless{}stdlib.h\textgreater{}}
\PreprocessorTok{\#include }\ImportTok{\textless{}threads.h\textgreater{}}

\NormalTok{tss\_t str}\OperatorTok{;}

\DataTypeTok{void}\NormalTok{ some\_function}\OperatorTok{(}\DataTypeTok{void}\OperatorTok{)}
\OperatorTok{\{}
    \CommentTok{// Ambil nilai per{-}thread dari string ini}
    \DataTypeTok{char} \OperatorTok{*}\NormalTok{tss\_string }\OperatorTok{=}\NormalTok{ tss\_get}\OperatorTok{(}\NormalTok{str}\OperatorTok{);}

    \CommentTok{// Lalu cetak}
\NormalTok{    printf}\OperatorTok{(}\StringTok{"TSS string: }\SpecialCharTok{\%s\textbackslash{}n}\StringTok{"}\OperatorTok{,}\NormalTok{ tss\_string}\OperatorTok{);}
\OperatorTok{\}}

\DataTypeTok{int}\NormalTok{ run}\OperatorTok{(}\DataTypeTok{void} \OperatorTok{*}\NormalTok{arg}\OperatorTok{)}
\OperatorTok{\{}
    \DataTypeTok{int}\NormalTok{ serial }\OperatorTok{=} \OperatorTok{*(}\DataTypeTok{int}\OperatorTok{*)}\NormalTok{arg}\OperatorTok{;}  \CommentTok{// Ambil nomor seri thread ini}
\NormalTok{    free}\OperatorTok{(}\NormalTok{arg}\OperatorTok{);}

    \CommentTok{// malloc() ruang untuk menampung data thread ini}
    \DataTypeTok{char} \OperatorTok{*}\NormalTok{s }\OperatorTok{=}\NormalTok{ malloc}\OperatorTok{(}\DecValTok{64}\OperatorTok{);}
\NormalTok{    sprintf}\OperatorTok{(}\NormalTok{s}\OperatorTok{,} \StringTok{"thread }\SpecialCharTok{\%d}\StringTok{! :)"}\OperatorTok{,}\NormalTok{ serial}\OperatorTok{);}  \CommentTok{// String kecil yang riang}

    \CommentTok{// Atur variabel TSS ini agar menunjuk ke string}
\NormalTok{    tss\_set}\OperatorTok{(}\NormalTok{str}\OperatorTok{,}\NormalTok{ s}\OperatorTok{);}

    \CommentTok{// Panggil fungsi yang akan mengambil variabelnya}
\NormalTok{    some\_function}\OperatorTok{();}

    \ControlFlowTok{return} \DecValTok{0}\OperatorTok{;}   \CommentTok{// Setara dengan thrd\_exit(0)}
\OperatorTok{\}}

\PreprocessorTok{\#define THREAD\_COUNT }\DecValTok{15}

\DataTypeTok{int}\NormalTok{ main}\OperatorTok{(}\DataTypeTok{void}\OperatorTok{)}
\OperatorTok{\{}
\NormalTok{    thrd\_t t}\OperatorTok{[}\NormalTok{THREAD\_COUNT}\OperatorTok{];}

    \CommentTok{// Buat variabel TSS baru, fungsi free() adalah destructor{-}nya}
\NormalTok{    tss\_create}\OperatorTok{(\&}\NormalTok{str}\OperatorTok{,}\NormalTok{ free}\OperatorTok{);}

    \ControlFlowTok{for} \OperatorTok{(}\DataTypeTok{int}\NormalTok{ i }\OperatorTok{=} \DecValTok{0}\OperatorTok{;}\NormalTok{ i }\OperatorTok{\textless{}}\NormalTok{ THREAD\_COUNT}\OperatorTok{;}\NormalTok{ i}\OperatorTok{++)} \OperatorTok{\{}
        \DataTypeTok{int} \OperatorTok{*}\NormalTok{n }\OperatorTok{=}\NormalTok{ malloc}\OperatorTok{(}\KeywordTok{sizeof} \OperatorTok{*}\NormalTok{n}\OperatorTok{);}  \CommentTok{// Menampung nomor seri thread}
        \OperatorTok{*}\NormalTok{n }\OperatorTok{=}\NormalTok{ i}\OperatorTok{;}
\NormalTok{        thrd\_create}\OperatorTok{(}\NormalTok{t }\OperatorTok{+}\NormalTok{ i}\OperatorTok{,}\NormalTok{ run}\OperatorTok{,}\NormalTok{ n}\OperatorTok{);}
    \OperatorTok{\}}

    \ControlFlowTok{for} \OperatorTok{(}\DataTypeTok{int}\NormalTok{ i }\OperatorTok{=} \DecValTok{0}\OperatorTok{;}\NormalTok{ i }\OperatorTok{\textless{}}\NormalTok{ THREAD\_COUNT}\OperatorTok{;}\NormalTok{ i}\OperatorTok{++)} \OperatorTok{\{}
\NormalTok{        thrd\_join}\OperatorTok{(}\NormalTok{t}\OperatorTok{[}\NormalTok{i}\OperatorTok{],}\NormalTok{ NULL}\OperatorTok{);}
    \OperatorTok{\}}

    \CommentTok{// Semua thread sudah selesai, jadi kita selesai dengan ini}
\NormalTok{    tss\_delete}\OperatorTok{(}\NormalTok{str}\OperatorTok{);}
\OperatorTok{\}}
\end{Highlighting}
\end{Shaded}

\end{tcolorbox}

\index{tss\_t type@\texttt{tss\_t} type|)} \index{tss\_get() function@\texttt{tss\_get()} function|)} \index{tss\_set() function@\texttt{tss\_set()} function|)} \index{tss\_create() function@\texttt{tss\_create()} function|)} \index{tss\_delete() function@\texttt{tss\_delete()} function|)}

Sekali lagi, ini cara yang agak menyakitkan dibanding \hl{\mbox{\texttt{thread\_local}}}, jadi kecuali kamu benar-benar membutuhkan fungsionalitas destructor itu, aku akan memakai \hl{\mbox{\texttt{thread\_local}}} saja.

\index{Thread-specific storage|)}

\section{Mutex}\label{mutex}

\index{Mutexes|(}

Kalau kamu hanya ingin mengizinkan satu thread masuk ke bagian kritis kode pada satu waktu, kamu bisa melindungi bagian itu dengan mutex\footnote{Singkatan ``mutual exclusion'', alias ``lock'' pada bagian kode yang hanya boleh dieksekusi satu thread.}.

Misalnya, kalau ada variabel \hl{\mbox{\texttt{static}}} dan kita ingin bisa mengambil dan mengisinya dalam dua operasi tanpa thread lain melompat ke tengah dan merusaknya, kita bisa memakai mutex untuk itu.

Kamu bisa memperoleh (acquire) mutex atau melepasnya. Kalau kamu mencoba memperoleh mutex dan berhasil, kamu boleh melanjutkan eksekusi. Kalau mencoba dan gagal (karena orang lain memegangnya), kamu akan \emph{block}\footnote{Artinya, prosesmu akan tidur.} sampai mutex dilepas.

Kalau beberapa thread diblokir menunggu mutex dilepas, salah satunya akan dipilih untuk berjalan (secara acak, dari sudut pandang kita), dan yang lain tetap tidur.

Rencana permainannya: pertama kita inisialisasi variabel mutex agar siap dipakai dengan \index{mtx\_init() function@\texttt{mtx\_init()} function} \hl{\mbox{\texttt{mtx\_init()}}}.

Lalu thread-thread berikutnya bisa memanggil \index{mtx\_lock() function@\texttt{mtx\_lock()} function} \hl{\mbox{\texttt{mtx\_lock()}}} dan \index{mtx\_unlock function@\texttt{mtx\_unlock} function} \hl{\mbox{\texttt{mtx\_unlock()}}} untuk memperoleh dan melepas mutex.

Saat kita benar-benar selesai dengan mutex, kita bisa menghancurkannya dengan \index{mtx\_destroy() function@\texttt{mtx\_destroy()} function} \hl{\mbox{\texttt{mtx\_destroy()}}}, kebalikan logis dari \index{mtx\_init() function@\texttt{mtx\_init()} function} \hl{\mbox{\texttt{mtx\_init()}}}.

\index{Multithreading!race conditions|(}

Pertama, mari kita lihat kode yang \emph{tidak} memakai mutex, dan berusaha mencetak nomor seri yang dibagi (\hl{\mbox{\texttt{static}}}) lalu meng-increment-nya. Karena kita tidak memakai mutex pada pengambilan nilai (untuk dicetak) dan pengisiannya (untuk di-increment), thread bisa saling menghalangi di bagian kritis itu.

\begin{tcolorbox}[enhanced jigsaw,left=3pt,breakable,colback=white!97.2549!black,boxrule=0.2mm]

\begin{Shaded}
\begin{Highlighting}[numbers=left,,]
\PreprocessorTok{\#include }\ImportTok{\textless{}stdio.h\textgreater{}}
\PreprocessorTok{\#include }\ImportTok{\textless{}threads.h\textgreater{}}

\DataTypeTok{int}\NormalTok{ run}\OperatorTok{(}\DataTypeTok{void} \OperatorTok{*}\NormalTok{arg}\OperatorTok{)}
\OperatorTok{\{}
    \OperatorTok{(}\DataTypeTok{void}\OperatorTok{)}\NormalTok{arg}\OperatorTok{;}

    \DataTypeTok{static} \DataTypeTok{int}\NormalTok{ serial }\OperatorTok{=} \DecValTok{0}\OperatorTok{;}   \CommentTok{// Variabel static yang dibagi!}

\NormalTok{    printf}\OperatorTok{(}\StringTok{"Thread running! }\SpecialCharTok{\%d\textbackslash{}n}\StringTok{"}\OperatorTok{,}\NormalTok{ serial}\OperatorTok{);}

\NormalTok{    serial}\OperatorTok{++;}

    \ControlFlowTok{return} \DecValTok{0}\OperatorTok{;}
\OperatorTok{\}}

\PreprocessorTok{\#define THREAD\_COUNT }\DecValTok{10}

\DataTypeTok{int}\NormalTok{ main}\OperatorTok{(}\DataTypeTok{void}\OperatorTok{)}
\OperatorTok{\{}
\NormalTok{    thrd\_t t}\OperatorTok{[}\NormalTok{THREAD\_COUNT}\OperatorTok{];}

    \ControlFlowTok{for} \OperatorTok{(}\DataTypeTok{int}\NormalTok{ i }\OperatorTok{=} \DecValTok{0}\OperatorTok{;}\NormalTok{ i }\OperatorTok{\textless{}}\NormalTok{ THREAD\_COUNT}\OperatorTok{;}\NormalTok{ i}\OperatorTok{++)} \OperatorTok{\{}
\NormalTok{        thrd\_create}\OperatorTok{(}\NormalTok{t }\OperatorTok{+}\NormalTok{ i}\OperatorTok{,}\NormalTok{ run}\OperatorTok{,}\NormalTok{ NULL}\OperatorTok{);}
    \OperatorTok{\}}

    \ControlFlowTok{for} \OperatorTok{(}\DataTypeTok{int}\NormalTok{ i }\OperatorTok{=} \DecValTok{0}\OperatorTok{;}\NormalTok{ i }\OperatorTok{\textless{}}\NormalTok{ THREAD\_COUNT}\OperatorTok{;}\NormalTok{ i}\OperatorTok{++)} \OperatorTok{\{}
\NormalTok{        thrd\_join}\OperatorTok{(}\NormalTok{t}\OperatorTok{[}\NormalTok{i}\OperatorTok{],}\NormalTok{ NULL}\OperatorTok{);}
    \OperatorTok{\}}
\OperatorTok{\}}
\end{Highlighting}
\end{Shaded}

\end{tcolorbox}

Waktu kujalankan, hasilnya kira-kira seperti ini:

\begin{tcolorbox}[enhanced jigsaw,left=3pt,breakable,colback=white!97.2549!black,boxrule=0.2mm]

\begin{Shaded}
\begin{Highlighting}[]
\NormalTok{Thread running! 0}
\NormalTok{Thread running! 0}
\NormalTok{Thread running! 0}
\NormalTok{Thread running! 3}
\NormalTok{Thread running! 4}
\NormalTok{Thread running! 5}
\NormalTok{Thread running! 6}
\NormalTok{Thread running! 7}
\NormalTok{Thread running! 8}
\NormalTok{Thread running! 9}
\end{Highlighting}
\end{Shaded}

\end{tcolorbox}

Jelas beberapa thread masuk ke sana dan menjalankan \hl{\mbox{\texttt{printf()}}} sebelum ada yang sempat memperbarui variabel \hl{\mbox{\texttt{serial}}}.

\index{Multithreading!race conditions|)}

Yang ingin kita lakukan adalah membungkus pengambilan variabel dan pengisiannya ke dalam satu bentangan kode yang dilindungi mutex.

Kita tambahkan variabel baru untuk mewakili mutex bertipe \index{mtx\_t type@\texttt{mtx\_t} type} \hl{\mbox{\texttt{mtx\_t}}} di file scope, menginisialisasinya, lalu thread bisa mengunci dan membukanya di fungsi \hl{\mbox{\texttt{run()}}}.

\index{mtx\_lock() function@\texttt{mtx\_lock()} function|(} \index{mtx\_unlock() function@\texttt{mtx\_unlock()} function|(} \index{mtx\_init() function@\texttt{mtx\_init()} function|(} \index{mtx\_destroy() function@\texttt{mtx\_destroy()} function|(}

\begin{tcolorbox}[enhanced jigsaw,left=3pt,breakable,colback=white!97.2549!black,boxrule=0.2mm]

\begin{Shaded}
\begin{Highlighting}[numbers=left,,]
\PreprocessorTok{\#include }\ImportTok{\textless{}stdio.h\textgreater{}}
\PreprocessorTok{\#include }\ImportTok{\textless{}threads.h\textgreater{}}

\NormalTok{mtx\_t serial\_mtx}\OperatorTok{;}     \CommentTok{// \textless{}{-}{-} VARIABEL MUTEX}

\DataTypeTok{int}\NormalTok{ run}\OperatorTok{(}\DataTypeTok{void} \OperatorTok{*}\NormalTok{arg}\OperatorTok{)}
\OperatorTok{\{}
    \OperatorTok{(}\DataTypeTok{void}\OperatorTok{)}\NormalTok{arg}\OperatorTok{;}

    \DataTypeTok{static} \DataTypeTok{int}\NormalTok{ serial }\OperatorTok{=} \DecValTok{0}\OperatorTok{;}   \CommentTok{// Variabel static yang dibagi!}

    \CommentTok{// Peroleh mutex{-}{-}semua thread akan terblokir di pemanggilan ini}
    \CommentTok{// sampai mereka mendapat lock:}

\NormalTok{    mtx\_lock}\OperatorTok{(\&}\NormalTok{serial\_mtx}\OperatorTok{);}           \CommentTok{// \textless{}{-}{-} PEROLEH MUTEX}

\NormalTok{    printf}\OperatorTok{(}\StringTok{"Thread running! }\SpecialCharTok{\%d\textbackslash{}n}\StringTok{"}\OperatorTok{,}\NormalTok{ serial}\OperatorTok{);}

\NormalTok{    serial}\OperatorTok{++;}

    \CommentTok{// Selesai mengambil dan mengisi data, jadi lepas lock. Ini akan}
    \CommentTok{// membuka blokir thread pada pemanggilan mtx\_lock():}

\NormalTok{    mtx\_unlock}\OperatorTok{(\&}\NormalTok{serial\_mtx}\OperatorTok{);}         \CommentTok{// \textless{}{-}{-} LEPAS MUTEX}

    \ControlFlowTok{return} \DecValTok{0}\OperatorTok{;}
\OperatorTok{\}}

\PreprocessorTok{\#define THREAD\_COUNT }\DecValTok{10}

\DataTypeTok{int}\NormalTok{ main}\OperatorTok{(}\DataTypeTok{void}\OperatorTok{)}
\OperatorTok{\{}
\NormalTok{    thrd\_t t}\OperatorTok{[}\NormalTok{THREAD\_COUNT}\OperatorTok{];}

    \CommentTok{// Inisialisasi variabel mutex, menandakan ini mutex biasa tanpa}
    \CommentTok{// embel{-}embel:}

\NormalTok{    mtx\_init}\OperatorTok{(\&}\NormalTok{serial\_mtx}\OperatorTok{,}\NormalTok{ mtx\_plain}\OperatorTok{);}        \CommentTok{// \textless{}{-}{-} BUAT MUTEX}

    \ControlFlowTok{for} \OperatorTok{(}\DataTypeTok{int}\NormalTok{ i }\OperatorTok{=} \DecValTok{0}\OperatorTok{;}\NormalTok{ i }\OperatorTok{\textless{}}\NormalTok{ THREAD\_COUNT}\OperatorTok{;}\NormalTok{ i}\OperatorTok{++)} \OperatorTok{\{}
\NormalTok{        thrd\_create}\OperatorTok{(}\NormalTok{t }\OperatorTok{+}\NormalTok{ i}\OperatorTok{,}\NormalTok{ run}\OperatorTok{,}\NormalTok{ NULL}\OperatorTok{);}
    \OperatorTok{\}}

    \ControlFlowTok{for} \OperatorTok{(}\DataTypeTok{int}\NormalTok{ i }\OperatorTok{=} \DecValTok{0}\OperatorTok{;}\NormalTok{ i }\OperatorTok{\textless{}}\NormalTok{ THREAD\_COUNT}\OperatorTok{;}\NormalTok{ i}\OperatorTok{++)} \OperatorTok{\{}
\NormalTok{        thrd\_join}\OperatorTok{(}\NormalTok{t}\OperatorTok{[}\NormalTok{i}\OperatorTok{],}\NormalTok{ NULL}\OperatorTok{);}
    \OperatorTok{\}}

    \CommentTok{// Selesai dengan mutex, hancurkan:}

\NormalTok{    mtx\_destroy}\OperatorTok{(\&}\NormalTok{serial\_mtx}\OperatorTok{);}                \CommentTok{// \textless{}{-}{-} HANCURKAN MUTEX}
\OperatorTok{\}}
\end{Highlighting}
\end{Shaded}

\end{tcolorbox}

Lihat bagaimana di baris 38 dan 50 pada \hl{\mbox{\texttt{main()}}} kita menginisialisasi dan menghancurkan mutex.

\index{mtx\_init() function@\texttt{mtx\_init()} function|)} \index{mtx\_destroy() function@\texttt{mtx\_destroy()} function|)}

Tapi setiap thread memperoleh mutex di baris 15 dan melepasnya di baris 24.

Di antara \hl{\mbox{\texttt{mtx\_lock()}}} dan \hl{\mbox{\texttt{mtx\_unlock()}}} adalah \emph{critical section}, bagian kode yang tidak ingin kita biarkan dikacau beberapa thread sekaligus.

\index{mtx\_lock() function@\texttt{mtx\_lock()} function|)} \index{mtx\_unlock() function@\texttt{mtx\_unlock()} function|)}

Dan sekarang kita mendapat output yang benar!

\begin{tcolorbox}[enhanced jigsaw,left=3pt,breakable,colback=white!97.2549!black,boxrule=0.2mm]

\begin{Shaded}
\begin{Highlighting}[]
\NormalTok{Thread running! 0}
\NormalTok{Thread running! 1}
\NormalTok{Thread running! 2}
\NormalTok{Thread running! 3}
\NormalTok{Thread running! 4}
\NormalTok{Thread running! 5}
\NormalTok{Thread running! 6}
\NormalTok{Thread running! 7}
\NormalTok{Thread running! 8}
\NormalTok{Thread running! 9}
\end{Highlighting}
\end{Shaded}

\end{tcolorbox}

Kalau butuh beberapa mutex, tidak masalah: punya saja beberapa variabel mutex.

Dan selalu ingat Aturan Nomor Satu untuk Banyak Mutex: \emph{Buka kunci mutex dalam urutan terbalik dari urutan saat mengunci!}

\subsection{Berbagai Tipe Mutex}\label{berbagai-tipe-mutex}

\index{Mutexes!types|(}

Seperti yang diisyaratkan tadi, ada beberapa tipe mutex yang bisa dibuat dengan \hl{\mbox{\texttt{mtx\_init()}}}. (Sebagian tipe ini adalah hasil operasi bitwise-OR, seperti tercatat di tabel.)

\index{Mutexes!timeouts|(}

\index{mtx\_plain macro@\texttt{mtx\_plain} macro} \index{mtx\_timed macro@\texttt{mtx\_timed} macro} \index{mtx\_plain macro@\texttt{mtx\_plain} macro} \index{mtx\_timed macro@\texttt{mtx\_timed} macro} \index{mtx\_recursive macro@\texttt{mtx\_recursive} macro}

{\def\LTcaptype{none} % do not increment counter
\begin{longtable}[]{@{}ll@{}}
\toprule\noalign{}
Tipe & Deskripsi \\
\midrule\noalign{}
\endhead
\bottomrule\noalign{}
\endlastfoot
\hl{\mbox{\texttt{mtx\_plain}}} & Mutex biasa \\
\hl{\mbox{\texttt{mtx\_timed}}} & Mutex yang mendukung timeout \\
\hl{\mbox{\texttt{mtx\_plain\textbar{}mtx\_recursive}}} & Mutex rekursif \\
\hl{\mbox{\texttt{mtx\_timed\textbar{}mtx\_recursive}}} & Mutex rekursif yang mendukung timeout \\
\end{longtable}
}

``Rekursif'' berarti pemegang lock boleh memanggil \hl{\mbox{\texttt{mtx\_lock()}}} berkali-kali pada lock yang sama. (Ia harus membukanya sebanyak itu juga sebelum orang lain bisa mengambil mutex.) Ini kadang memudahkan penulisan kode, terutama kalau kamu memanggil fungsi yang perlu mengunci mutex padahal kamu sudah memegang mutex itu.

Dan timeout memberi thread kesempatan untuk \emph{mencoba} mendapat lock selama beberapa waktu, tapi menyerah kalau tidak berhasil dalam jangka waktu itu.

\index{mtx\_timed macro@\texttt{mtx\_timed} macro|(}

Untuk mutex dengan timeout, pastikan membuatnya dengan \hl{\mbox{\texttt{mtx\_timed}}}:

\begin{tcolorbox}[enhanced jigsaw,left=3pt,breakable,colback=white!97.2549!black,boxrule=0.2mm]

\begin{Shaded}
\begin{Highlighting}[]
\NormalTok{mtx\_init}\OperatorTok{(\&}\NormalTok{serial\_mtx}\OperatorTok{,}\NormalTok{ mtx\_timed}\OperatorTok{);}
\end{Highlighting}
\end{Shaded}

\end{tcolorbox}

\index{mtx\_timed macro@\texttt{mtx\_timed} macro|)}

Lalu saat menunggunya, kamu harus menentukan waktu dalam UTC kapan ia akan dibuka\footnote{Mungkin kamu mengira itu ``waktu dari sekarang'', tapi kamu pasti ingin berpikir begitu, kan!}.

\index{timespec\_get() function@\texttt{timespec\_get()} function|(}

Fungsi \hl{\mbox{\texttt{timespec\_get()}}} dari \hl{\mbox{\texttt{\textless{}time.h\textgreater{}}}} bisa membantu di sini. Ia memberimu waktu saat ini dalam UTC di \hl{\mbox{\texttt{struct\ timespec}}}, yang persis kita butuhkan. Bahkan, ia tampaknya ada hanya untuk tujuan ini.

Ia punya dua field: \hl{\mbox{\texttt{tv\_sec}}} berisi waktu saat ini dalam detik sejak epoch, dan \hl{\mbox{\texttt{tv\_nsec}}} berisi nanodetik (sepersemiliar detik) sebagai bagian ``pecahan''.

Jadi kamu bisa mengisinya dengan waktu saat ini, lalu menambahkannya untuk mendapat timeout tertentu.

\index{mtx\_timedlock() function@\texttt{mtx\_timedlock()} function|(}

Lalu panggil \hl{\mbox{\texttt{mtx\_timedlock()}}} alih-alih \hl{\mbox{\texttt{mtx\_lock()}}}. Kalau ia mengembalikan nilai \hl{\mbox{\texttt{thrd\_timedout}}}, berarti timeout.

\begin{tcolorbox}[enhanced jigsaw,left=3pt,breakable,colback=white!97.2549!black,boxrule=0.2mm]

\begin{Shaded}
\begin{Highlighting}[]
\KeywordTok{struct}\NormalTok{ timespec timeout}\OperatorTok{;}

\NormalTok{timespec\_get}\OperatorTok{(\&}\NormalTok{timeout}\OperatorTok{,}\NormalTok{ TIME\_UTC}\OperatorTok{);}  \CommentTok{// Ambil waktu saat ini}
\NormalTok{timeout}\OperatorTok{.}\NormalTok{tv\_sec }\OperatorTok{+=} \DecValTok{1}\OperatorTok{;}               \CommentTok{// Timeout 1 detik setelah sekarang}

\DataTypeTok{int}\NormalTok{ result }\OperatorTok{=}\NormalTok{ mtx\_timedlock}\OperatorTok{(\&}\NormalTok{serial\_mtx}\OperatorTok{,} \OperatorTok{\&}\NormalTok{timeout}\OperatorTok{));}

\ControlFlowTok{if} \OperatorTok{(}\NormalTok{result }\OperatorTok{==}\NormalTok{ thrd\_timedout}\OperatorTok{)} \OperatorTok{\{}
\NormalTok{    printf}\OperatorTok{(}\StringTok{"Mutex lock timed out!}\SpecialCharTok{\textbackslash{}n}\StringTok{"}\OperatorTok{);}
\OperatorTok{\}}
\end{Highlighting}
\end{Shaded}

\end{tcolorbox}

\index{mtx\_timedlock() function@\texttt{mtx\_timedlock()} function|)} \index{timespec\_get() function@\texttt{timespec\_get()} function|)}

Selain itu, timed lock sama saja dengan lock biasa.

\index{Mutexes!timeouts|)} \index{Mutexes!types|)} \index{Mutexes|)}

\section{Condition Variable}\label{condition-variable}

\index{Condition variables|(}

Condition Variable adalah kepingan terakhir yang kita butuhkan untuk membuat aplikasi multithread yang berkinerja baik dan menyusun struktur multithread yang lebih kompleks.

Condition variable menyediakan cara agar thread tidur sampai suatu kejadian di thread lain terjadi.

Dengan kata lain, mungkin ada sejumlah thread yang sudah tidak sabar berangkat, tapi mereka harus menunggu sampai suatu kejadian benar sebelum melanjutkan. Pada dasarnya mereka disuruh ``tunggu dulu!'' sampai diberi tahu.

Dan ini bekerja beriringan dengan mutex karena yang kita tunggu biasanya bergantung pada nilai suatu data, dan data itu biasanya perlu dilindungi mutex.

Penting dicatat bahwa condition variable itu sendiri bukan penyimpan data tertentu dari sudut pandang kita. Ia hanyalah variabel yang dipakai C untuk melacak status menunggu/tidak menunggu dari sebuah thread atau kelompok thread.

Mari kita tulis program rekaan yang membaca kelompok berisi 5 angka dari thread utama satu per satu. Lalu, ketika 5 angka sudah dimasukkan, child thread bangun, menjumlahkan kelima angka itu, dan mencetak hasilnya.

Angka-angka itu akan disimpan di array global yang dibagi, begitu juga indeks ke array untuk angka yang akan dimasukkan.

Karena ini nilai yang dibagi, setidaknya kita harus menyembunyikannya di balik mutex untuk thread utama maupun child thread. (Thread utama akan menulis data ke sana dan child thread akan membaca data darinya.)

Tapi itu belum cukup. Child thread perlu diblokir (``tidur'') sampai 5 angka terbaca ke array. Lalu parent thread perlu membangunkan child thread agar bisa bekerja.

Dan saat bangun, ia perlu memegang mutex itu. Dan memang akan! Ketika thread menunggu pada condition variable, ia juga memperoleh mutex saat bangun.

Semua ini terjadi seputar variabel tambahan bertipe \index{cnd\_t type@\texttt{cnd\_t} type} \hl{\mbox{\texttt{cnd\_t}}} yang adalah \emph{condition variable}. Kita membuat variabel ini dengan fungsi \index{cnd\_init() function@\texttt{cnd\_init()} function} \hl{\mbox{\texttt{cnd\_init()}}} dan menghancurkannya saat selesai dengan fungsi \index{cnd\_destroy() function@\texttt{cnd\_destroy()} function} \hl{\mbox{\texttt{cnd\_destroy()}}}.

Tapi bagaimana semuanya bekerja? Mari kita lihat garis besar apa yang akan dilakukan child thread:

\index{mtx\_lock() function@\texttt{mtx\_lock()} function|(} \index{mtx\_unlock() function@\texttt{mtx\_unlock()} function|(} \index{cnd\_wait() function@\texttt{cnd\_wait()} function|(} \index{cnd\_signal() function@\texttt{cnd\_signal()} function|(}

\begin{enumerate}
\def\labelenumi{\arabic{enumi}.}
\tightlist
\item
  Kunci mutex dengan \hl{\mbox{\texttt{mtx\_lock()}}}
\item
  Kalau kita belum memasukkan semua angka, tunggu pada condition variable dengan \hl{\mbox{\texttt{cnd\_wait()}}}
\item
  Kerjakan pekerjaan yang harus dikerjakan
\item
  Buka kunci mutex dengan \hl{\mbox{\texttt{mtx\_unlock()}}}
\end{enumerate}

Sementara itu thread utama akan melakukan ini:

\begin{enumerate}
\def\labelenumi{\arabic{enumi}.}
\tightlist
\item
  Kunci mutex dengan \hl{\mbox{\texttt{mtx\_lock()}}}
\item
  Simpan angka yang baru dibaca ke array
\item
  Kalau array penuh, beri sinyal agar child bangun dengan \hl{\mbox{\texttt{cnd\_signal()}}}
\item
  Buka kunci mutex dengan \hl{\mbox{\texttt{mtx\_unlock()}}}
\end{enumerate}

\index{mtx\_lock() function@\texttt{mtx\_lock()} function|)} \index{mtx\_unlock() function@\texttt{mtx\_unlock()} function|)}

Kalau kamu tidak terlalu cepat membacanya sekilas (tidak apa-apa---aku tidak tersinggung), mungkin kamu memperhatikan sesuatu yang aneh: bagaimana thread utama bisa memegang kunci mutex dan memberi sinyal ke child, kalau child harus memegang kunci mutex untuk menunggu sinyal? Keduanya tidak bisa memegang kunci!

Dan memang tidak! Ada sihir di balik layar pada condition variable: saat kamu \hl{\mbox{\texttt{cnd\_wait()}}}, ia melepas mutex yang kamu tentukan dan thread tidur. Dan ketika seseorang memberi sinyal agar thread itu bangun, ia memperoleh kembali kunci itu seolah-olah tidak terjadi apa-apa.

Di sisi \hl{\mbox{\texttt{cnd\_signal()}}} sedikit berbeda. Ia tidak melakukan apa-apa terhadap mutex. Thread yang memberi sinyal masih harus melepas mutex secara manual sebelum thread yang menunggu bisa bangun.

\index{cnd\_signal() function@\texttt{cnd\_signal()} function|)}

Satu hal lagi soal \hl{\mbox{\texttt{cnd\_wait()}}}. Kamu kemungkinan akan memanggil \hl{\mbox{\texttt{cnd\_wait()}}} kalau suatu kondisi\footnote{Dan itulah kenapa mereka disebut \emph{condition variable}!} belum terpenuhi (mis. dalam kasus ini, kalau belum semua angka dimasukkan). Intinya: kondisi ini harus berada dalam loop \hl{\mbox{\texttt{while}}}, bukan pernyataan \hl{\mbox{\texttt{if}}}. Kenapa?

Itu karena fenomena misterius bernama \index{Condition variables!spurious wakeup} \emph{spurious wakeup}. Kadang, pada beberapa implementasi, sebuah thread bisa terbangun dari tidur \hl{\mbox{\texttt{cnd\_wait()}}} tanpa \emph{alasan} yang jelas. \emph{{[}Musik X-Files{]}}\footnote{Aku tidak bilang ini ulah alien\ldots{} tapi ini ulah alien. Oke, sebenarnya lebih mungkin thread lain sudah terbangun dan sampai ke pekerjaannya lebih dulu.}. Karena itu kita harus memeriksa bahwa kondisi yang kita butuhkan masih benar-benar terpenuhi saat bangun. Dan kalau tidak, kembali tidur!

\index{cnd\_wait() function@\texttt{cnd\_wait()} function|)}

Jadi mari kita kerjakan! Mulai dari thread utama:

\begin{itemize}
\item
  Thread utama akan menyiapkan mutex dan condition variable, dan meluncurkan child thread.
\item
  Lalu ia akan, dalam loop tak terbatas, mengambil angka sebagai input dari konsol.
\item
  Ia juga akan memperoleh mutex untuk menyimpan angka yang dimasukkan ke array global.
\item
  Ketika array berisi 5 angka, thread utama akan memberi sinyal ke child thread bahwa saatnya bangun dan bekerja.
\item
  Lalu thread utama akan membuka kunci mutex dan kembali membaca angka berikutnya dari konsol.
\end{itemize}

Sementara itu, child thread sibuk dengan kenakalannya sendiri:

\begin{itemize}
\item
  Child thread mengambil mutex
\item
  Selama kondisi belum terpenuhi (yaitu selama array bersama belum berisi 5 angka), child thread tidur dengan menunggu pada condition variable. Saat menunggu, ia secara implisit membuka kunci mutex.
\item
  Begitu thread utama memberi sinyal agar child thread bangun, ia bangun untuk bekerja dan mendapatkan kembali kunci mutex.
\item
  Child thread menjumlahkan angka-angka itu dan mereset variabel yang menjadi indeks ke array.
\item
  Lalu ia melepas mutex dan berjalan lagi dalam loop tak terbatas.
\end{itemize}

Dan inilah kodenya! Pelajari sebentar supaya kamu bisa melihat di mana semua bagian di atas ditangani:

\index{mtx\_lock() function@\texttt{mtx\_lock()} function|(} \index{mtx\_unlock() function@\texttt{mtx\_unlock()} function|(} \index{mtx\_init() function@\texttt{mtx\_init()} function|(} \index{mtx\_destroy() function@\texttt{mtx\_destroy()} function|(} \index{cnd\_init() function@\texttt{cnd\_init()} function|(} \index{cnd\_destroy() function@\texttt{cnd\_destroy()} function|(} \index{cnd\_wait() function@\texttt{cnd\_wait()} function|(} \index{cnd\_signal() function@\texttt{cnd\_signal()} function|(} \index{cnd\_t type@\texttt{cnd\_t} type|(}

\begin{tcolorbox}[enhanced jigsaw,left=3pt,breakable,colback=white!97.2549!black,boxrule=0.2mm]

\begin{Shaded}
\begin{Highlighting}[numbers=left,,]
\PreprocessorTok{\#include }\ImportTok{\textless{}stdio.h\textgreater{}}
\PreprocessorTok{\#include }\ImportTok{\textless{}threads.h\textgreater{}}

\PreprocessorTok{\#define VALUE\_COUNT\_MAX }\DecValTok{5}

\DataTypeTok{int}\NormalTok{ value}\OperatorTok{[}\NormalTok{VALUE\_COUNT\_MAX}\OperatorTok{];}  \CommentTok{// Global yang dibagi}
\DataTypeTok{int}\NormalTok{ value\_count }\OperatorTok{=} \DecValTok{0}\OperatorTok{;}   \CommentTok{// Global yang dibagi juga}

\NormalTok{mtx\_t value\_mtx}\OperatorTok{;}   \CommentTok{// Mutex untuk value}
\NormalTok{cnd\_t value\_cnd}\OperatorTok{;}   \CommentTok{// Condition variable untuk value}

\DataTypeTok{int}\NormalTok{ run}\OperatorTok{(}\DataTypeTok{void} \OperatorTok{*}\NormalTok{arg}\OperatorTok{)}
\OperatorTok{\{}
    \OperatorTok{(}\DataTypeTok{void}\OperatorTok{)}\NormalTok{arg}\OperatorTok{;}

    \ControlFlowTok{for} \OperatorTok{(;;)} \OperatorTok{\{}
\NormalTok{        mtx\_lock}\OperatorTok{(\&}\NormalTok{value\_mtx}\OperatorTok{);}      \CommentTok{// \textless{}{-}{-} AMBIL MUTEX}

        \ControlFlowTok{while} \OperatorTok{(}\NormalTok{value\_count }\OperatorTok{\textless{}}\NormalTok{ VALUE\_COUNT\_MAX}\OperatorTok{)} \OperatorTok{\{}
\NormalTok{            printf}\OperatorTok{(}\StringTok{"Thread: is waiting}\SpecialCharTok{\textbackslash{}n}\StringTok{"}\OperatorTok{);}
\NormalTok{            cnd\_wait}\OperatorTok{(\&}\NormalTok{value\_cnd}\OperatorTok{,} \OperatorTok{\&}\NormalTok{value\_mtx}\OperatorTok{);}  \CommentTok{// \textless{}{-}{-} TUNGGU KONDISI}
        \OperatorTok{\}}

\NormalTok{        printf}\OperatorTok{(}\StringTok{"Thread: is awake!}\SpecialCharTok{\textbackslash{}n}\StringTok{"}\OperatorTok{);}

        \DataTypeTok{int}\NormalTok{ t }\OperatorTok{=} \DecValTok{0}\OperatorTok{;}

        \CommentTok{// Jumlahkan semuanya}
        \ControlFlowTok{for} \OperatorTok{(}\DataTypeTok{int}\NormalTok{ i }\OperatorTok{=} \DecValTok{0}\OperatorTok{;}\NormalTok{ i }\OperatorTok{\textless{}}\NormalTok{ VALUE\_COUNT\_MAX}\OperatorTok{;}\NormalTok{ i}\OperatorTok{++)}
\NormalTok{            t }\OperatorTok{+=}\NormalTok{ value}\OperatorTok{[}\NormalTok{i}\OperatorTok{];}

\NormalTok{        printf}\OperatorTok{(}\StringTok{"Thread: total is }\SpecialCharTok{\%d\textbackslash{}n}\StringTok{"}\OperatorTok{,}\NormalTok{ t}\OperatorTok{);}

        \CommentTok{// Reset indeks input untuk thread utama}
\NormalTok{        value\_count }\OperatorTok{=} \DecValTok{0}\OperatorTok{;}

\NormalTok{        mtx\_unlock}\OperatorTok{(\&}\NormalTok{value\_mtx}\OperatorTok{);}   \CommentTok{// \textless{}{-}{-} BUKA KUNCI MUTEX}
    \OperatorTok{\}}

    \ControlFlowTok{return} \DecValTok{0}\OperatorTok{;}
\OperatorTok{\}}

\DataTypeTok{int}\NormalTok{ main}\OperatorTok{(}\DataTypeTok{void}\OperatorTok{)}
\OperatorTok{\{}
\NormalTok{    thrd\_t t}\OperatorTok{;}

    \CommentTok{// Luncurkan thread baru}

\NormalTok{    thrd\_create}\OperatorTok{(\&}\NormalTok{t}\OperatorTok{,}\NormalTok{ run}\OperatorTok{,}\NormalTok{ NULL}\OperatorTok{);}
\NormalTok{    thrd\_detach}\OperatorTok{(}\NormalTok{t}\OperatorTok{);}

    \CommentTok{// Siapkan mutex dan condition variable}

\NormalTok{    mtx\_init}\OperatorTok{(\&}\NormalTok{value\_mtx}\OperatorTok{,}\NormalTok{ mtx\_plain}\OperatorTok{);}
\NormalTok{    cnd\_init}\OperatorTok{(\&}\NormalTok{value\_cnd}\OperatorTok{);}

    \ControlFlowTok{for} \OperatorTok{(;;)} \OperatorTok{\{}
        \DataTypeTok{int}\NormalTok{ n}\OperatorTok{;}

\NormalTok{        scanf}\OperatorTok{(}\StringTok{"}\SpecialCharTok{\%d}\StringTok{"}\OperatorTok{,} \OperatorTok{\&}\NormalTok{n}\OperatorTok{);}

\NormalTok{        mtx\_lock}\OperatorTok{(\&}\NormalTok{value\_mtx}\OperatorTok{);}    \CommentTok{// \textless{}{-}{-} KUNCI MUTEX}

\NormalTok{        value}\OperatorTok{[}\NormalTok{value\_count}\OperatorTok{++]} \OperatorTok{=}\NormalTok{ n}\OperatorTok{;}

        \ControlFlowTok{if} \OperatorTok{(}\NormalTok{value\_count }\OperatorTok{==}\NormalTok{ VALUE\_COUNT\_MAX}\OperatorTok{)} \OperatorTok{\{}
\NormalTok{            printf}\OperatorTok{(}\StringTok{"Main: signaling thread}\SpecialCharTok{\textbackslash{}n}\StringTok{"}\OperatorTok{);}
\NormalTok{            cnd\_signal}\OperatorTok{(\&}\NormalTok{value\_cnd}\OperatorTok{);}  \CommentTok{// \textless{}{-}{-} BERI SINYAL KONDISI}
        \OperatorTok{\}}

\NormalTok{        mtx\_unlock}\OperatorTok{(\&}\NormalTok{value\_mtx}\OperatorTok{);}  \CommentTok{// \textless{}{-}{-} BUKA KUNCI MUTEX}
    \OperatorTok{\}}

    \CommentTok{// Bersih{-}bersih (loop di atas ini tak terbatas, tapi kita ingin}
    \CommentTok{// setidaknya berpura{-}pura rapi):}

\NormalTok{    mtx\_destroy}\OperatorTok{(\&}\NormalTok{value\_mtx}\OperatorTok{);}
\NormalTok{    cnd\_destroy}\OperatorTok{(\&}\NormalTok{value\_cnd}\OperatorTok{);}
\OperatorTok{\}}
\end{Highlighting}
\end{Shaded}

\end{tcolorbox}

\index{mtx\_lock() function@\texttt{mtx\_lock()} function|)} \index{mtx\_unlock() function@\texttt{mtx\_unlock()} function|)} \index{mtx\_init() function@\texttt{mtx\_init()} function|)} \index{mtx\_destroy() function@\texttt{mtx\_destroy()} function|)} \index{cnd\_init() function@\texttt{cnd\_init()} function|)} \index{cnd\_destroy() function@\texttt{cnd\_destroy()} function|)} \index{cnd\_wait() function@\texttt{cnd\_wait()} function|)} \index{cnd\_signal() function@\texttt{cnd\_signal()} function|)} \index{cnd\_t type@\texttt{cnd\_t} type|)}

Dan berikut contoh output (angka-angka tunggal di baris adalah inputku):

\begin{tcolorbox}[enhanced jigsaw,left=3pt,breakable,colback=white!97.2549!black,boxrule=0.2mm]

\begin{Shaded}
\begin{Highlighting}[]
\NormalTok{Thread: is waiting}
\NormalTok{1}
\NormalTok{1}
\NormalTok{1}
\NormalTok{1}
\NormalTok{1}
\NormalTok{Main: signaling thread}
\NormalTok{Thread: is awake!}
\NormalTok{Thread: total is 5}
\NormalTok{Thread: is waiting}
\NormalTok{2}
\NormalTok{8}
\NormalTok{5}
\NormalTok{9}
\NormalTok{0}
\NormalTok{Main: signaling thread}
\NormalTok{Thread: is awake!}
\NormalTok{Thread: total is 24}
\NormalTok{Thread: is waiting}
\end{Highlighting}
\end{Shaded}

\end{tcolorbox}

Ini pemakaian condition variable yang umum pada situasi producer-consumer seperti ini. Kalau kita tidak punya cara untuk menidurkan child thread selagi ia menunggu suatu kondisi terpenuhi, ia akan terpaksa melakukan polling, yang sangat memboroskan CPU.

\subsection{Condition Wait dengan Timeout}\label{condition-wait-dengan-timeout}

\index{Condition variables!timeouts|(}

Ada varian \hl{\mbox{\texttt{cnd\_wait()}}} yang memungkinkanmu menentukan timeout agar bisa berhenti menunggu.

Karena child thread harus mengunci ulang mutex, ini belum tentu berarti kamu langsung hidup kembali tepat saat timeout terjadi; kamu masih harus menunggu thread lain melepas mutex.

Tapi artinya kamu tidak akan menunggu sampai \hl{\mbox{\texttt{cnd\_signal()}}} terjadi.

Agar ini berfungsi, panggil \index{cnd\_timedwait() function@\texttt{cnd\_timedwait()} function} \hl{\mbox{\texttt{cnd\_timedwait()}}} alih-alih \hl{\mbox{\texttt{cnd\_wait()}}}. Kalau ia mengembalikan nilai \index{thrd\_timedout macro@\texttt{thrd\_timedout} macro} \hl{\mbox{\texttt{thrd\_timedout}}}, berarti timeout.

Timestamp-nya adalah waktu absolut dalam UTC, bukan waktu-dari-sekarang. Untungnya fungsi \index{timespec\_get() function@\texttt{timespec\_get()} function} \hl{\mbox{\texttt{timespec\_get()}}} di \hl{\mbox{\texttt{\textless{}time.h\textgreater{}}}} tampaknya dibuat khusus untuk kasus ini.

\index{timespec\_get() function@\texttt{timespec\_get()} function|(} \index{cnd\_timedwait() function@\texttt{cnd\_timedwait()} function|(} \index{thrd\_timedout() macro@\texttt{thrd\_timedout()} macro|(}

\begin{tcolorbox}[enhanced jigsaw,left=3pt,breakable,colback=white!97.2549!black,boxrule=0.2mm]

\begin{Shaded}
\begin{Highlighting}[]
\KeywordTok{struct}\NormalTok{ timespec timeout}\OperatorTok{;}

\NormalTok{timespec\_get}\OperatorTok{(\&}\NormalTok{timeout}\OperatorTok{,}\NormalTok{ TIME\_UTC}\OperatorTok{);}  \CommentTok{// Ambil waktu saat ini}
\NormalTok{timeout}\OperatorTok{.}\NormalTok{tv\_sec }\OperatorTok{+=} \DecValTok{1}\OperatorTok{;}               \CommentTok{// Timeout 1 detik setelah sekarang}

\DataTypeTok{int}\NormalTok{ result }\OperatorTok{=}\NormalTok{ cnd\_timedwait}\OperatorTok{(\&}\NormalTok{condition}\OperatorTok{,} \OperatorTok{\&}\NormalTok{mutex}\OperatorTok{,} \OperatorTok{\&}\NormalTok{timeout}\OperatorTok{));}

\ControlFlowTok{if} \OperatorTok{(}\NormalTok{result }\OperatorTok{==}\NormalTok{ thrd\_timedout}\OperatorTok{)} \OperatorTok{\{}
\NormalTok{    printf}\OperatorTok{(}\StringTok{"Condition variable timed out!}\SpecialCharTok{\textbackslash{}n}\StringTok{"}\OperatorTok{);}
\OperatorTok{\}}
\end{Highlighting}
\end{Shaded}

\end{tcolorbox}

\index{timespec\_get() function@\texttt{timespec\_get()} function|)} \index{cnd\_timedwait() function@\texttt{cnd\_timedwait()} function|)} \index{thrd\_timedout() macro@\texttt{thrd\_timedout()} macro|)} \index{Condition variables!timeouts|)}

\subsection{Broadcast: Membangunkan Semua Thread yang Menunggu}\label{broadcast-membangunkan-semua-thread-yang-menunggu}

\index{Condition variables!broadcasting|(}

\index{cnd\_signal() function@\texttt{cnd\_signal()} function|)} \hl{\mbox{\texttt{cnd\_signal()}}} hanya membangunkan satu thread untuk melanjutkan pekerjaan. Tergantung bagaimana logikamu disusun, mungkin masuk akal membangunkan lebih dari satu thread untuk melanjutkan begitu kondisi terpenuhi.

Tentu saja hanya satu dari mereka yang bisa mengambil mutex, tapi kalau kamu punya situasi ketika:

\begin{itemize}
\item
  Thread yang baru bangun bertanggung jawab membangunkan thread berikutnya, dan---
\item
  Ada kemungkinan kondisi loop spurious-wakeup mencegahnya melakukan itu, maka---
\end{itemize}

kamu perlu men-broadcast pembangunan itu agar yakin setidaknya satu thread keluar dari loop itu untuk meluncurkan thread berikutnya.

Bagaimana caranya, tanyamu?

\index{cnd\_broadcast() function@\texttt{cnd\_broadcast()} function|(}

Cukup pakai \hl{\mbox{\texttt{cnd\_broadcast()}}} alih-alih \hl{\mbox{\texttt{cnd\_signal()}}}. Pemakaiannya persis sama, hanya saja \hl{\mbox{\texttt{cnd\_broadcast()}}} membangunkan \textbf{semua} thread yang tidur menunggu condition variable itu.

\index{cnd\_broadcast() function@\texttt{cnd\_broadcast()} function|)} \index{Condition variables!broadcasting|)} \index{Condition variables|)}

\section{Menjalankan Fungsi Satu Kali}\label{menjalankan-fungsi-satu-kali}

\index{Multithreading!one-time functions|(}

Katakanlah kamu punya fungsi yang \emph{bisa} dijalankan banyak thread, tapi kamu tidak tahu kapan, dan tidak sepadan repot menulis semua logikanya.

Ada jalan keluarnya: pakai \index{call\_once() function@\texttt{call\_once()} function} \hl{\mbox{\texttt{call\_once()}}}. Banyak sekali thread bisa mencoba menjalankan fungsi itu, tapi hanya yang pertama yang dihitung\footnote{Yang kuat yang bertahan! Benar? Harus diakui ini sebenarnya sama sekali tidak seperti itu.}

Untuk memakainya, kamu butuh variabel flag khusus yang dideklarasikan untuk melacak apakah fungsi itu sudah dijalankan atau belum. Dan kamu butuh fungsi yang akan dijalankan, yang tidak menerima parameter dan tidak mengembalikan nilai.

\index{once\_flag type@\texttt{once\_flag} type|(} \index{ONCE\_FLAG\_INIT macro@\texttt{ONCE\_FLAG\_INIT} macro|(} \index{call\_once() function@\texttt{call\_once()} function|(}

\begin{tcolorbox}[enhanced jigsaw,left=3pt,breakable,colback=white!97.2549!black,boxrule=0.2mm]

\begin{Shaded}
\begin{Highlighting}[]
\NormalTok{once\_flag of }\OperatorTok{=}\NormalTok{ ONCE\_FLAG\_INIT}\OperatorTok{;}  \CommentTok{// Inisialisasi seperti ini}

\DataTypeTok{void}\NormalTok{ run\_once\_function}\OperatorTok{(}\DataTypeTok{void}\OperatorTok{)}
\OperatorTok{\{}
\NormalTok{    printf}\OperatorTok{(}\StringTok{"I\textquotesingle{}ll only run once!}\SpecialCharTok{\textbackslash{}n}\StringTok{"}\OperatorTok{);}
\OperatorTok{\}}

\DataTypeTok{int}\NormalTok{ run}\OperatorTok{(}\DataTypeTok{void} \OperatorTok{*}\NormalTok{arg}\OperatorTok{)}
\OperatorTok{\{}
    \OperatorTok{(}\DataTypeTok{void}\OperatorTok{)}\NormalTok{arg}\OperatorTok{;}

\NormalTok{    call\_once}\OperatorTok{(\&}\NormalTok{of}\OperatorTok{,}\NormalTok{ run\_once\_function}\OperatorTok{);}

    \CommentTok{// ...}
\end{Highlighting}
\end{Shaded}

\end{tcolorbox}

\index{once\_flag type@\texttt{once\_flag} type|)} \index{ONCE\_FLAG\_INIT macro@\texttt{ONCE\_FLAG\_INIT} macro|)} \index{call\_once() function@\texttt{call\_once()} function|)}

Pada contoh ini, berapa pun banyaknya thread yang sampai ke fungsi \hl{\mbox{\texttt{run()}}}, \hl{\mbox{\texttt{run\_once\_function()}}} hanya akan dipanggil satu kali.

\index{Multithreading!one-time functions|)} \index{Multithreading|)}

\chapter{Atomics}\label{chapter-atomics}

\begin{tcolorbox}[blanker,colback=white,enhanced jigsaw,breakable,arc=0pt,outer arc=0pt,grow to left by=-1cm,enlarge top by=0.1cm,enlarge bottom by=0.1cm,borderline west={1.5pt}{-.3cm}{black}]

\begin{quote}
\emph{``They tried and failed, all of them?''}\\
\emph{``Oh, no.'' She shook her head. ``They tried and died.''}

---Paul Atreides and The Reverend Mother Gaius Helen Mohiam, \emph{Dune}
\end{quote}

\end{tcolorbox}

\index{Atomic variables|(}

Ini salah satu aspek multithreading di C yang paling menantang. Tapi kita akan berusaha santai.

Pada dasarnya, aku akan membahas pemakaian variabel atomic yang lebih lugas, apa itu, bagaimana cara kerjanya, dan sebagainya. Dan aku akan menyebut beberapa jalur yang jauh lebih rumit yang tersedia bagimu.

Tapi aku tidak akan menempuh jalur-jalur itu. Bukan hanya karena aku nyaris tidak punya kualifikasi untuk menulis tentangnya, tapi kupikir kalau kamu tahu kamu membutuhkannya, kamu sudah tahu lebih banyak daripada aku.

Tapi ada hal-hal aneh bahkan di dasar-dasarnya. Jadi kencangkan sabuk pengaman, semuanya, karena Kansas akan pamit pergi.

\section{Menguji Dukungan Atomic}\label{menguji-dukungan-atomic}

\index{\_\_STDC\_NO\_ATOMICS\_\_ macro@\texttt{\_\_STDC\_NO\_ATOMICS\_\_} macro|(}

Atomic adalah fitur opsional. Ada makro \hl{\mbox{\texttt{\_\_STDC\_NO\_ATOMICS\_\_}}} yang bernilai \hl{\mbox{\texttt{1}}} kalau kamu \emph{tidak} punya atomic.

Makro itu mungkin tidak ada sebelum C11, jadi kita sebaiknya menguji versi bahasanya dengan \hl{\mbox{\texttt{\_\_STDC\_VERSION\_\_}}}\footnote{Makro \hl{\mbox{\texttt{\_\_STDC\_VERSION\_\_}}} tidak ada di C89 awal, jadi kalau kamu khawatir soal itu, periksa dengan \hl{\mbox{\texttt{\#ifdef}}}.}.

\begin{tcolorbox}[enhanced jigsaw,left=3pt,breakable,colback=white!97.2549!black,boxrule=0.2mm]

\begin{Shaded}
\begin{Highlighting}[]
\PreprocessorTok{\#if \_\_STDC\_VERSION\_\_ \textless{} 201112L || \_\_STDC\_NO\_ATOMICS\_\_ == 1}
\PreprocessorTok{\#define HAS\_ATOMICS }\DecValTok{0}
\PreprocessorTok{\#else}
\PreprocessorTok{\#define HAS\_ATOMICS }\DecValTok{1}
\PreprocessorTok{\#endif}
\end{Highlighting}
\end{Shaded}

\end{tcolorbox}

\index{\_\_STDC\_NO\_ATOMICS\_\_ macro@\texttt{\_\_STDC\_NO\_ATOMICS\_\_} macro|)}

\index{Atomic variables!compiling with|(}

Kalau pengujian itu lolos, kamu aman meng-include \hl{\mbox{\texttt{\textless{}stdatomic.h\textgreater{}}}}, header yang menjadi dasar sisa bab ini. Tapi kalau tidak ada dukungan atomic, header itu mungkin bahkan tidak ada.

Di beberapa sistem, kamu mungkin perlu menambahkan \hl{\mbox{\texttt{-latomic}}} di akhir command line kompilasi untuk memakai fungsi apa pun di header file itu.

\index{Atomic variables!compiling with|)}

\section{Variabel Atomic}\label{variabel-atomic}

Berikut \emph{sebagian} dari cara kerja variabel atomic:

Kalau kamu punya variabel atomic yang dibagi dan kamu menulis ke sana dari satu thread, penulisan itu akan \emph{semua-atau-tidak-sama-sekali} di thread yang berbeda.

Artinya, thread lain akan melihat seluruh penulisan, misalnya nilai 32-bit. Bukan setengahnya. Tidak ada cara bagi satu thread untuk menginterupsi thread lain yang sedang berada di \emph{tengah} penulisan atomic multibyte.

Hampir seperti ada kunci kecil di sekitar pengambilan dan pengisian variabel itu. (Dan \emph{mungkin} memang ada! Lihat \hyperref[lock-free-atomic]{Variabel Atomic Lock-Free}, di bawah.)

Dan soal itu, kamu bisa saja tidak pernah memakai atomic kalau memakai mutex untuk mengunci critical section-mu. Hanya saja ada kelas \emph{struktur data lock-free} yang selalu mengizinkan thread lain membuat kemajuan alih-alih terblokir oleh mutex\ldots{} tapi ini sulit dibuat dengan benar dari nol, dan merupakan salah satu hal yang di luar cakupan panduan ini, sayangnya.

Itu baru sebagian ceritanya. Tapi itulah bagian yang kita mulai.

Sebelum lanjut, bagaimana cara mendeklarasikan variabel sebagai atomic?

Pertama, include \index{stdatomic.h header@\texttt{stdatomic.h} header} \hl{\mbox{\texttt{\textless{}stdatomic.h\textgreater{}}}}.

\index{atomic\_int type@\texttt{atomic\_int} type|(}

Ini memberi kita tipe seperti \hl{\mbox{\texttt{atomic\_int}}}.

Lalu kita tinggal mendeklarasikan variabel bertipe itu.

Tapi mari kita buat demo dengan dua thread. Yang pertama berjalan sebentar lalu mengisi variabel dengan nilai tertentu, lalu keluar. Yang lain berjalan sampai ia melihat nilai itu terisi, lalu keluar.

\begin{tcolorbox}[enhanced jigsaw,left=3pt,breakable,colback=white!97.2549!black,boxrule=0.2mm]

\begin{Shaded}
\begin{Highlighting}[numbers=left,,]
\PreprocessorTok{\#include }\ImportTok{\textless{}stdio.h\textgreater{}}
\PreprocessorTok{\#include }\ImportTok{\textless{}threads.h\textgreater{}}
\PreprocessorTok{\#include }\ImportTok{\textless{}stdatomic.h\textgreater{}}

\NormalTok{atomic\_int x}\OperatorTok{;}   \CommentTok{// KEKUATAN ATOMIC! BWAHAHAHA!}

\DataTypeTok{int}\NormalTok{ thread1}\OperatorTok{(}\DataTypeTok{void} \OperatorTok{*}\NormalTok{arg}\OperatorTok{)}
\OperatorTok{\{}
    \OperatorTok{(}\DataTypeTok{void}\OperatorTok{)}\NormalTok{arg}\OperatorTok{;}

\NormalTok{    printf}\OperatorTok{(}\StringTok{"Thread 1: Sleeping for 1.5 seconds}\SpecialCharTok{\textbackslash{}n}\StringTok{"}\OperatorTok{);}
\NormalTok{    thrd\_sleep}\OperatorTok{(\&(}\KeywordTok{struct}\NormalTok{ timespec}\OperatorTok{)\{.}\NormalTok{tv\_sec}\OperatorTok{=}\DecValTok{1}\OperatorTok{,} \OperatorTok{.}\NormalTok{tv\_nsec}\OperatorTok{=}\DecValTok{500000000}\OperatorTok{\},}\NormalTok{ NULL}\OperatorTok{);}

\NormalTok{    printf}\OperatorTok{(}\StringTok{"Thread 1: Setting x to 3490}\SpecialCharTok{\textbackslash{}n}\StringTok{"}\OperatorTok{);}
\NormalTok{    x }\OperatorTok{=} \DecValTok{3490}\OperatorTok{;}

\NormalTok{    printf}\OperatorTok{(}\StringTok{"Thread 1: Exiting}\SpecialCharTok{\textbackslash{}n}\StringTok{"}\OperatorTok{);}
    \ControlFlowTok{return} \DecValTok{0}\OperatorTok{;}
\OperatorTok{\}}

\DataTypeTok{int}\NormalTok{ thread2}\OperatorTok{(}\DataTypeTok{void} \OperatorTok{*}\NormalTok{arg}\OperatorTok{)}
\OperatorTok{\{}
    \OperatorTok{(}\DataTypeTok{void}\OperatorTok{)}\NormalTok{arg}\OperatorTok{;}

\NormalTok{    printf}\OperatorTok{(}\StringTok{"Thread 2: Waiting for 3490}\SpecialCharTok{\textbackslash{}n}\StringTok{"}\OperatorTok{);}
    \ControlFlowTok{while} \OperatorTok{(}\NormalTok{x }\OperatorTok{!=} \DecValTok{3490}\OperatorTok{)} \OperatorTok{\{\}}  \CommentTok{// berputar di sini}

\NormalTok{    printf}\OperatorTok{(}\StringTok{"Thread 2: Got 3490{-}{-}exiting!}\SpecialCharTok{\textbackslash{}n}\StringTok{"}\OperatorTok{);}
    \ControlFlowTok{return} \DecValTok{0}\OperatorTok{;}
\OperatorTok{\}}

\DataTypeTok{int}\NormalTok{ main}\OperatorTok{(}\DataTypeTok{void}\OperatorTok{)}
\OperatorTok{\{}
\NormalTok{    x }\OperatorTok{=} \DecValTok{0}\OperatorTok{;}

\NormalTok{    thrd\_t t1}\OperatorTok{,}\NormalTok{ t2}\OperatorTok{;}

\NormalTok{    thrd\_create}\OperatorTok{(\&}\NormalTok{t1}\OperatorTok{,}\NormalTok{ thread1}\OperatorTok{,}\NormalTok{ NULL}\OperatorTok{);}
\NormalTok{    thrd\_create}\OperatorTok{(\&}\NormalTok{t2}\OperatorTok{,}\NormalTok{ thread2}\OperatorTok{,}\NormalTok{ NULL}\OperatorTok{);}

\NormalTok{    thrd\_join}\OperatorTok{(}\NormalTok{t1}\OperatorTok{,}\NormalTok{ NULL}\OperatorTok{);}
\NormalTok{    thrd\_join}\OperatorTok{(}\NormalTok{t2}\OperatorTok{,}\NormalTok{ NULL}\OperatorTok{);}

\NormalTok{    printf}\OperatorTok{(}\StringTok{"Main    : Threads are done, so x better be 3490}\SpecialCharTok{\textbackslash{}n}\StringTok{"}\OperatorTok{);}
\NormalTok{    printf}\OperatorTok{(}\StringTok{"Main    : And indeed, x == }\SpecialCharTok{\%d\textbackslash{}n}\StringTok{"}\OperatorTok{,}\NormalTok{ x}\OperatorTok{);}
\OperatorTok{\}}
\end{Highlighting}
\end{Shaded}

\end{tcolorbox}

\index{atomic\_int type@\texttt{atomic\_int} type|)}

Thread kedua berputar di tempat, memandangi flag dan menunggunya diisi dengan nilai \hl{\mbox{\texttt{3490}}}. Dan thread pertama melakukannya.

Dan outputnya:

\begin{tcolorbox}[enhanced jigsaw,left=3pt,breakable,colback=white!97.2549!black,boxrule=0.2mm]

\begin{Shaded}
\begin{Highlighting}[]
\NormalTok{Thread 1: Sleeping for 1.5 seconds}
\NormalTok{Thread 2: Waiting for 3490}
\NormalTok{Thread 1: Setting x to 3490}
\NormalTok{Thread 1: Exiting}
\NormalTok{Thread 2: Got 3490{-}{-}exiting!}
\NormalTok{Main    : Threads are done, so x better be 3490}
\NormalTok{Main    : And indeed, x == 3490}
\end{Highlighting}
\end{Shaded}

\end{tcolorbox}

Lihat, Ma! Kita mengakses variabel dari thread berbeda dan tidak memakai mutex! Dan itu akan berhasil setiap kali berkat sifat atomic dari variabel atomic.

Mungkin kamu bertanya-tanya apa yang terjadi kalau itu \hl{\mbox{\texttt{int}}} biasa yang non-atomic. Nah, di sistemku tetap berfungsi\ldots{} kecuali kalau aku melakukan build dengan optimasi, dan dalam kasus itu ia menggantung di thread 2 menunggu melihat 3490 terisi\footnote{Alasannya, saat dioptimasi, compilerku menaruh nilai \hl{\mbox{\texttt{x}}} di register agar loop \hl{\mbox{\texttt{while}}} cepat. Tapi register tidak punya cara untuk tahu bahwa variabel diperbarui di thread lain, jadi ia tidak pernah melihat \hl{\mbox{\texttt{3490}}}. Ini sebenarnya tidak terkait dengan bagian \emph{semua-atau-tidak-sama-sekali} dari atomicity, tapi lebih terkait dengan aspek sinkronisasi di bagian berikutnya.}.

Tapi itu baru awal ceritanya. Bagian berikutnya akan membutuhkan lebih banyak tenaga otak dan berkaitan dengan sesuatu yang disebut \emph{sinkronisasi}.

\section{Sinkronisasi}\label{sinkronisasi}

\index{Atomic variables!synchronization|(}

Bagian cerita berikutnya sepenuhnya tentang kapan penulisan memori tertentu di satu thread menjadi terlihat oleh thread lain.

Mungkin kamu mengira, langsung terlihat, kan? Tapi tidak. Banyak hal bisa salah. Salah secara aneh.

Compiler mungkin sudah menata ulang akses memori sehingga ketika kamu mengira sebuah nilai diisi lebih dulu dibanding yang lain, itu belum tentu benar. Dan kalaupun compiler tidak melakukannya, CPU-mu bisa melakukannya secara langsung. Atau mungkin ada hal lain di arsitektur ini yang membuat penulisan di satu CPU tertunda sebelum terlihat di CPU lain.

Kabar baiknya, semua potensi masalah ini bisa kita ringkas menjadi satu: akses memori yang tidak tersinkronisasi bisa tampak tidak berurutan tergantung thread mana yang mengamati, seolah-olah baris kodenya sendiri telah ditata ulang.

Sebagai contoh, mana yang terjadi lebih dulu pada kode berikut, penulisan ke \hl{\mbox{\texttt{x}}} atau penulisan ke \hl{\mbox{\texttt{y}}}?

\begin{tcolorbox}[enhanced jigsaw,left=3pt,breakable,colback=white!97.2549!black,boxrule=0.2mm]

\begin{Shaded}
\begin{Highlighting}[numbers=left,,]
\DataTypeTok{int}\NormalTok{ x}\OperatorTok{,}\NormalTok{ y}\OperatorTok{;}  \CommentTok{// global}

\CommentTok{// ...}

\NormalTok{x }\OperatorTok{=} \DecValTok{2}\OperatorTok{;}
\NormalTok{y }\OperatorTok{=} \DecValTok{3}\OperatorTok{;}

\NormalTok{printf}\OperatorTok{(}\StringTok{"}\SpecialCharTok{\%d}\StringTok{ }\SpecialCharTok{\%d\textbackslash{}n}\StringTok{"}\OperatorTok{,}\NormalTok{ x}\OperatorTok{,}\NormalTok{ y}\OperatorTok{);}
\end{Highlighting}
\end{Shaded}

\end{tcolorbox}

Jawabannya: kita tidak tahu. Compiler atau CPU bisa diam-diam menukar baris 5 dan 6 dan kita tidak akan tahu. Kode berjalan single-thread \emph{seolah-olah} dieksekusi sesuai urutan kode.

Dalam skenario multithread, kita mungkin punya sesuatu seperti pseudocode ini:

\begin{tcolorbox}[enhanced jigsaw,left=3pt,breakable,colback=white!97.2549!black,boxrule=0.2mm]

\begin{Shaded}
\begin{Highlighting}[numbers=left,,]
\DataTypeTok{int}\NormalTok{ x }\OperatorTok{=} \DecValTok{0}\OperatorTok{,}\NormalTok{ y }\OperatorTok{=} \DecValTok{0}\OperatorTok{;}

\NormalTok{thread1}\OperatorTok{()} \OperatorTok{\{}
\NormalTok{    x }\OperatorTok{=} \DecValTok{2}\OperatorTok{;}
\NormalTok{    y }\OperatorTok{=} \DecValTok{3}\OperatorTok{;}
\OperatorTok{\}}

\NormalTok{thread2}\OperatorTok{()} \OperatorTok{\{}
    \ControlFlowTok{while} \OperatorTok{(}\NormalTok{y }\OperatorTok{!=} \DecValTok{3}\OperatorTok{)} \OperatorTok{\{\}}  \CommentTok{// berputar}
\NormalTok{    printf}\OperatorTok{(}\StringTok{"x is now }\SpecialCharTok{\%d\textbackslash{}n}\StringTok{"}\OperatorTok{,}\NormalTok{ x}\OperatorTok{);}  \CommentTok{// 2? ...atau 0?}
\OperatorTok{\}}
\end{Highlighting}
\end{Shaded}

\end{tcolorbox}

Apa output dari thread 2?

Nah, kalau \hl{\mbox{\texttt{x}}} diisi \hl{\mbox{\texttt{2}}} \emph{sebelum} \hl{\mbox{\texttt{y}}} diisi \hl{\mbox{\texttt{3}}}, maka aku akan mengharapkan output yang sangat masuk akal:

\begin{tcolorbox}[enhanced jigsaw,left=3pt,breakable,colback=white!97.2549!black,boxrule=0.2mm]

\begin{Shaded}
\begin{Highlighting}[]
\NormalTok{x is now 2 }
\end{Highlighting}
\end{Shaded}

\end{tcolorbox}

Tapi sesuatu yang licik bisa menata ulang baris 4 dan 5 sehingga kita melihat nilai \hl{\mbox{\texttt{0}}} untuk \hl{\mbox{\texttt{x}}} saat mencetaknya.

Dengan kata lain, semua taruhan batal kecuali kita bisa entah bagaimana berkata, ``Mulai titik ini, aku mengharapkan semua penulisan sebelumnya di thread lain terlihat di thread ini.''

Dua thread \emph{tersinkronisasi} ketika mereka sepakat tentang keadaan memori bersama. Seperti yang sudah kita lihat, mereka tidak selalu sepakat dengan kode. Jadi bagaimana mereka sepakat?

Memakai variabel atomic bisa memaksa kesepakatan itu\footnote{Sampai kukatakan lain, aku berbicara secara umum tentang operasi \emph{sequentially consistent}. Soal artinya dibahas sebentar lagi.}. Kalau sebuah thread menulis ke variabel atomic, ia berkata ``siapa pun yang membaca variabel atomic ini di masa depan juga akan melihat semua perubahan yang kubuat pada memori (atomic atau bukan) sampai dan termasuk variabel atomic itu''.

Atau, dalam bahasa yang lebih manusiawi: mari duduk mengelilingi meja konferensi dan pastikan kita sepakat soal bagian memori bersama mana yang berisi nilai apa. Kamu setuju bahwa perubahan memori yang kamu buat sampai-dan-termasuk atomic store akan terlihat olehku setelah aku melakukan load pada variabel atomic yang sama.

Jadi kita bisa dengan mudah memperbaiki contoh kita:

\begin{tcolorbox}[enhanced jigsaw,left=3pt,breakable,colback=white!97.2549!black,boxrule=0.2mm]

\begin{Shaded}
\begin{Highlighting}[numbers=left,,]
\DataTypeTok{int}\NormalTok{ x }\OperatorTok{=} \DecValTok{0}\OperatorTok{;}
\NormalTok{atomic }\DataTypeTok{int}\NormalTok{ y }\OperatorTok{=} \DecValTok{0}\OperatorTok{;}  \CommentTok{// Jadikan y atomic}

\NormalTok{thread1}\OperatorTok{()} \OperatorTok{\{}
\NormalTok{    x }\OperatorTok{=} \DecValTok{2}\OperatorTok{;}
\NormalTok{    y }\OperatorTok{=} \DecValTok{3}\OperatorTok{;}             \CommentTok{// Sinkronisasi saat menulis}
\OperatorTok{\}}

\NormalTok{thread2}\OperatorTok{()} \OperatorTok{\{}
    \ControlFlowTok{while} \OperatorTok{(}\NormalTok{y }\OperatorTok{!=} \DecValTok{3}\OperatorTok{)} \OperatorTok{\{\}}  \CommentTok{// Sinkronisasi saat membaca}
\NormalTok{    printf}\OperatorTok{(}\StringTok{"x is now }\SpecialCharTok{\%d\textbackslash{}n}\StringTok{"}\OperatorTok{,}\NormalTok{ x}\OperatorTok{);}  \CommentTok{// 2, titik.}
\OperatorTok{\}}
\end{Highlighting}
\end{Shaded}

\end{tcolorbox}

Karena thread-thread itu tersinkronisasi melalui \hl{\mbox{\texttt{y}}}, semua penulisan di thread 1 yang terjadi \emph{sebelum} penulisan ke \hl{\mbox{\texttt{y}}} terlihat di thread 2 \emph{setelah} pembacaan dari \hl{\mbox{\texttt{y}}} (di loop \hl{\mbox{\texttt{while}}}).

Penting dicatat beberapa hal di sini:

\begin{enumerate}
\def\labelenumi{\arabic{enumi}.}
\item
  Tidak ada yang tidur. Sinkronisasi bukan operasi yang memblokir. Kedua thread berjalan penuh sampai keluar. Bahkan yang terjebak di spin loop pun tidak menghalangi yang lain berjalan.
\item
  Sinkronisasi terjadi ketika satu thread membaca variabel atomic yang ditulis thread lain. Jadi ketika thread 2 membaca \hl{\mbox{\texttt{y}}}, semua penulisan memori sebelumnya di thread 1 (yaitu mengisi \hl{\mbox{\texttt{x}}}) akan terlihat di thread 2.
\item
  Perhatikan bahwa \hl{\mbox{\texttt{x}}} tidak atomic. Itu tidak masalah karena kita tidak menyinkronkan lewat \hl{\mbox{\texttt{x}}}, dan sinkronisasi lewat \hl{\mbox{\texttt{y}}} saat kita menulisnya di thread 1 berarti semua penulisan sebelumnya---termasuk \hl{\mbox{\texttt{x}}}---di thread 1 akan terlihat oleh thread lain\ldots{} kalau thread lain itu membaca \hl{\mbox{\texttt{y}}} untuk sinkronisasi.
\end{enumerate}

Memaksa sinkronisasi ini tidak efisien dan bisa jauh lebih lambat daripada sekadar memakai variabel biasa. Itulah kenapa kita tidak memakai atomic kecuali terpaksa untuk aplikasi tertentu.

Jadi itulah dasar-dasarnya. Mari kita lihat lebih dalam.

\index{Atomic variables!synchronization|(}

\section{Acquire dan Release}\label{acquire-dan-release}

\index{Atomic variables!acquire|(} \index{Atomic variables!release|(}

Terminologi lagi! Akan bermanfaat mempelajarinya sekarang.

Ketika sebuah thread membaca variabel atomic, itu disebut operasi \emph{acquire}.

Ketika sebuah thread menulis variabel atomic, itu disebut operasi \emph{release}.

Apa maksudnya? Mari kita sejajarkan dengan istilah yang sudah kamu kenal terkait variabel atomic:

\textbf{Read = Load = Acquire}. Seperti saat kamu membandingkan variabel atomic atau membacanya untuk disalin ke nilai lain.

\textbf{Write = Store = Release}. Seperti saat kamu menugaskan nilai ke variabel atomic.

Saat memakai variabel atomic dengan semantik acquire/release ini, C menjabarkan apa yang bisa terjadi kapan.

Acquire/release membentuk dasar sinkronisasi yang baru saja kita bahas.

Ketika sebuah thread meng-acquire variabel atomic, ia bisa melihat nilai yang diisi di thread lain yang me-release variabel yang sama.

Dengan kata lain:

Ketika sebuah thread membaca variabel atomic, ia bisa melihat nilai yang diisi di thread lain yang menulis ke variabel yang sama.

Sinkronisasi terjadi di sepasang acquire/release.

Detail lebih lanjut:

Dengan read/load/acquire pada variabel atomic tertentu:

\begin{itemize}
\item
  Semua penulisan (atomic maupun non-atomic) di thread lain yang terjadi sebelum thread lain itu write/store/release variabel atomic ini kini terlihat di thread ini.
\item
  Nilai baru variabel atomic yang diisi thread lain juga terlihat di thread ini.
\item
  Tidak ada pembacaan atau penulisan variabel/memori apa pun di thread saat ini yang boleh ditata ulang untuk terjadi sebelum acquire ini.
\item
  Acquire berperan sebagai penghalang satu arah dalam penataan ulang kode; pembacaan dan penulisan di thread saat ini bisa dipindah turun dari \emph{sebelum} acquire ke \emph{sesudah}-nya. Tapi, yang lebih penting untuk sinkronisasi, tidak ada yang bisa pindah naik dari \emph{sesudah} acquire ke \emph{sebelum}-nya.
\end{itemize}

Dengan write/store/release pada variabel atomic tertentu:

\begin{itemize}
\item
  Semua penulisan (atomic maupun non-atomic) di thread saat ini yang terjadi sebelum release ini menjadi terlihat oleh thread lain yang telah read/load/acquire variabel atomic yang sama.
\item
  Nilai yang ditulis ke variabel atomic ini oleh thread ini juga terlihat oleh thread lain.
\item
  Tidak ada pembacaan atau penulisan variabel/memori apa pun di thread saat ini yang boleh ditata ulang untuk terjadi setelah release ini.
\item
  Release berperan sebagai penghalang satu arah dalam penataan ulang kode: pembacaan dan penulisan di thread saat ini bisa dipindah naik dari \emph{sesudah} release ke \emph{sebelum}-nya. Tapi, yang lebih penting untuk sinkronisasi, tidak ada yang bisa pindah turun dari \emph{sebelum} release ke \emph{sesudah}-nya.
\end{itemize}

Sekali lagi, intinya adalah sinkronisasi memori dari satu thread ke thread lain. Thread kedua bisa yakin bahwa variabel dan memori ditulis dalam urutan yang dimaksudkan programmer.

\begin{tcolorbox}[enhanced jigsaw,left=3pt,breakable,colback=white!97.2549!black,boxrule=0.2mm]

\begin{verbatim}
int x, y, z = 0;
atomic_int a = 0;

thread1() {
    x = 10;
    y = 20;
    a = 999;  // Release
    z = 30;
}

thread2()
{
    while (a != 999) { } // Acquire

    assert(x == 10);  // tidak pernah gagal, x selalu 10
    assert(y == 20);  // tidak pernah gagal, y selalu 20

    assert(z == 0);  // mungkin gagal!!
}
\end{verbatim}

\end{tcolorbox}

Pada contoh di atas, \hl{\mbox{\texttt{thread2}}} bisa yakin akan nilai di \hl{\mbox{\texttt{x}}} dan \hl{\mbox{\texttt{y}}} setelah meng-acquire \hl{\mbox{\texttt{a}}} karena keduanya diisi sebelum \hl{\mbox{\texttt{thread1}}} me-release atomic \hl{\mbox{\texttt{a}}}.

Tapi \hl{\mbox{\texttt{thread2}}} tidak bisa yakin akan nilai \hl{\mbox{\texttt{z}}} karena itu terjadi setelah release. Mungkin penugasan ke \hl{\mbox{\texttt{z}}} dipindah ke sebelum penugasan ke \hl{\mbox{\texttt{a}}}.

Catatan penting: me-release satu variabel atomic tidak berpengaruh pada acquire variabel atomic yang berbeda. Setiap variabel terisolasi dari yang lain.

\index{Atomic variables!acquire|)} \index{Atomic variables!release|)}

\section{Sequential Consistency}\label{sequential-consistency}

\index{Atomic variables!sequential consistency|(}

Masih bertahan? Kita sudah melewati inti pemakaian atomic yang lebih sederhana. Dan karena kita bahkan tidak akan membahas pemakaian yang lebih kompleks di sini, kamu bisa sedikit santai.

\emph{Sequential consistency} adalah yang disebut \emph{memory ordering}. Ada banyak memory ordering, tapi sequential consistency adalah yang paling waras\footnote{Paling waras dari sudut pandang programmer.} yang ditawarkan C. Ia juga yang bawaan. Kamu harus berusaha keras untuk memakai memory ordering lain.

Semua yang sudah kita bahas sejauh ini terjadi dalam ranah sequential consistency.

Kita sudah membahas bagaimana compiler atau CPU bisa menata ulang pembacaan dan penulisan memori dalam satu thread selama mengikuti aturan \emph{as-if}.

Dan kita sudah melihat bagaimana kita bisa mengerem perilaku ini dengan menyinkronkan lewat variabel atomic.

Mari kita formalkan sedikit lagi.

Kalau operasi-operasi \emph{sequentially consistent}, artinya pada akhirnya, ketika semuanya selesai, semua thread bisa menyelonjorkan kaki, membuka minuman pilihan masing-masing, dan semuanya sepakat tentang urutan terjadinya perubahan memori selama eksekusi. Dan urutan itu adalah yang ditentukan oleh kode.

Tidak akan ada yang berkata, ``Bukankah \emph{B} terjadi sebelum \emph{A}?'' kalau yang lain berkata, ``\emph{A} pasti terjadi sebelum \emph{B}''. Mereka semua berteman di sini.

Khususnya, dalam satu thread, tidak ada acquire dan release yang boleh ditata ulang terhadap satu sama lain. Ini tambahan dari aturan tentang akses memori lain apa yang boleh ditata ulang di sekitarnya.

Aturan ini memberi tingkat kewarasan tambahan pada perkembangan load/acquire dan store/release atomic.

Setiap memory order lain di C melibatkan pelonggaran aturan penataan ulang, baik untuk acquire/release maupun akses memori lain, atomic atau bukan. Kamu melakukannya kalau kamu \emph{benar-benar} tahu apa yang kamu lakukan dan butuh dorongan kecepatan. \emph{Di sini ada pasukan naga\ldots{}}

Soal itu dibahas nanti, tapi untuk sekarang, kita bertahan di yang aman dan praktis.

\index{Atomic variables!sequential consistency|)}

\section{Penugasan dan Operator Atomic}\label{penugasan-dan-operator-atomic}

\index{Atomic variables!assignments and operators|(}

Operator tertentu pada variabel atomic bersifat atomic. Dan yang lain tidak.

Kita mulai dengan contoh tandingan:

\begin{tcolorbox}[enhanced jigsaw,left=3pt,breakable,colback=white!97.2549!black,boxrule=0.2mm]

\begin{Shaded}
\begin{Highlighting}[]
\NormalTok{atomic\_int x }\OperatorTok{=} \DecValTok{0}\OperatorTok{;}

\NormalTok{thread1}\OperatorTok{()} \OperatorTok{\{}
\NormalTok{    x }\OperatorTok{=}\NormalTok{ x }\OperatorTok{+} \DecValTok{3}\OperatorTok{;}  \CommentTok{// TIDAK atomic!}
\OperatorTok{\}}
\end{Highlighting}
\end{Shaded}

\end{tcolorbox}

Karena ada pembacaan \hl{\mbox{\texttt{x}}} di sisi kanan penugasan dan penulisan di sisi kiri, ini dua operasi. Thread lain bisa menyelinap di tengah dan membuatmu kesal.

Tapi kamu \emph{bisa} memakai bentuk singkat \hl{\mbox{\texttt{+=}}} untuk mendapat operasi atomic:

\begin{tcolorbox}[enhanced jigsaw,left=3pt,breakable,colback=white!97.2549!black,boxrule=0.2mm]

\begin{Shaded}
\begin{Highlighting}[]
\NormalTok{atomic\_int x }\OperatorTok{=} \DecValTok{0}\OperatorTok{;}

\NormalTok{thread1}\OperatorTok{()} \OperatorTok{\{}
\NormalTok{    x }\OperatorTok{+=} \DecValTok{3}\OperatorTok{;}   \CommentTok{// ATOMIC!}
\OperatorTok{\}}
\end{Highlighting}
\end{Shaded}

\end{tcolorbox}

Dalam kasus itu, \hl{\mbox{\texttt{x}}} akan di-increment sebesar \hl{\mbox{\texttt{3}}} secara atomic---tidak ada thread lain yang bisa melompat ke tengah.

Khususnya, operator-operator berikut adalah operasi atomic read-modify-write dengan sequential consistency, jadi pakailah dengan riang gembira. (Pada contoh, \hl{\mbox{\texttt{a}}} bersifat atomic.)

\begin{tcolorbox}[enhanced jigsaw,left=3pt,breakable,colback=white!97.2549!black,boxrule=0.2mm]

\begin{Shaded}
\begin{Highlighting}[]
\NormalTok{a}\OperatorTok{++}\NormalTok{       a}\OperatorTok{{-}{-}}       \OperatorTok{{-}{-}}\NormalTok{a       }\OperatorTok{++}\NormalTok{a}
\NormalTok{a }\OperatorTok{+=}\NormalTok{ b    a }\OperatorTok{{-}=}\NormalTok{ b    a }\OperatorTok{*=}\NormalTok{ b    a }\OperatorTok{/=}\NormalTok{ b    a }\OperatorTok{\%=}\NormalTok{ b}
\NormalTok{a }\OperatorTok{\&=}\NormalTok{ b    a }\OperatorTok{|=}\NormalTok{ b    a }\OperatorTok{\^{}=}\NormalTok{ b    a }\OperatorTok{\textgreater{}\textgreater{}=}\NormalTok{ b   a }\OperatorTok{\textless{}\textless{}=}\NormalTok{ b}
\end{Highlighting}
\end{Shaded}

\end{tcolorbox}

\index{Atomic variables!assignments and operators|)}

\section{Fungsi Library yang Otomatis Menyinkronkan}\label{fungsi-library-yang-otomatis-menyinkronkan}

\index{Atomic variables!synchronized library functions|(}

Sejauh ini kita sudah membahas bagaimana kamu bisa menyinkronkan dengan variabel atomic, tapi ternyata ada beberapa fungsi library yang melakukan sinkronisasi terbatas di balik layar.

\begin{tcolorbox}[enhanced jigsaw,left=3pt,breakable,colback=white!97.2549!black,boxrule=0.2mm]

\begin{Shaded}
\begin{Highlighting}[]
\NormalTok{call\_once}\OperatorTok{()}\NormalTok{      thrd\_create}\OperatorTok{()}\NormalTok{       thrd\_join}\OperatorTok{()}
\NormalTok{mtx\_lock}\OperatorTok{()}\NormalTok{       mtx\_timedlock}\OperatorTok{()}\NormalTok{     mtx\_trylock}\OperatorTok{()}
\NormalTok{malloc}\OperatorTok{()}\NormalTok{         calloc}\OperatorTok{()}\NormalTok{            realloc}\OperatorTok{()}
\NormalTok{aligned\_alloc}\OperatorTok{()}
\end{Highlighting}
\end{Shaded}

\end{tcolorbox}

\textbf{\hl{\mbox{\texttt{call\_once()}}}}---Menyinkronkan dengan semua pemanggilan \hl{\mbox{\texttt{call\_once()}}} berikutnya untuk flag tertentu. Dengan begitu pemanggilan berikutnya bisa tenang bahwa kalau thread lain mengisi flag, mereka akan melihatnya.

\textbf{\hl{\mbox{\texttt{thrd\_create()}}}}---Menyinkronkan dengan awal thread baru. Thread baru bisa yakin akan melihat semua penulisan memori bersama dari parent thread sebelum pemanggilan \hl{\mbox{\texttt{thrd\_create()}}}.

\textbf{\hl{\mbox{\texttt{thrd\_join()}}}}---Ketika sebuah thread mati, ia menyinkronkan dengan fungsi ini. Thread yang memanggil \hl{\mbox{\texttt{thrd\_join()}}} bisa yakin bahwa ia melihat semua penulisan bersama dari thread yang telah berpulang itu.

\textbf{\hl{\mbox{\texttt{mtx\_lock()}}}}---Pemanggilan \hl{\mbox{\texttt{mtx\_unlock()}}} sebelumnya pada mutex yang sama menyinkronkan dengan pemanggilan ini. Inilah kasus yang paling mencerminkan proses acquire/release yang sudah kita bahas. \hl{\mbox{\texttt{mtx\_unlock()}}} melakukan release pada variabel mutex, menjamin thread berikutnya yang melakukan acquire dengan \hl{\mbox{\texttt{mtx\_lock()}}} bisa melihat semua perubahan memori bersama di critical section.

\textbf{\hl{\mbox{\texttt{mtx\_timedlock()}}}} dan \textbf{\hl{\mbox{\texttt{mtx\_trylock()}}}}---Mirip situasi \hl{\mbox{\texttt{mtx\_lock()}}}, kalau pemanggilan ini berhasil, pemanggilan \hl{\mbox{\texttt{mtx\_unlock()}}} sebelumnya menyinkronkan dengannya.

\textbf{Fungsi Memori Dinamis}: kalau kamu mengalokasikan memori, ia menyinkronkan dengan dealokasi sebelumnya dari memori yang sama. Dan alokasi serta dealokasi wilayah memori tertentu itu terjadi dalam satu urutan total yang disepakati semua thread. Aku \emph{kira} idenya, dealokasi boleh menghapus wilayah itu kalau mau, dan kita ingin yakin bahwa alokasi berikutnya tidak melihat data yang belum dihapus. Kabari aku kalau ada yang lebih dari itu.

\index{Atomic variables!synchronized library functions|)}

\section{Type Specifier dan Qualifier Atomic}\label{type-specifier-dan-qualifier-atomic}

Mari kita turunkan sedikit tingkatnya dan lihat tipe apa yang tersedia, dan bagaimana kita bahkan bisa membuat tipe atomic baru.

\index{\_Atomic type qualifier@\texttt{\_Atomic} type qualifier|(}

Pertama, mari kita lihat tipe atomic bawaan dan di-\hl{\mbox{\texttt{typedef}}} ke apa mereka. (Bocoran: \hl{\mbox{\texttt{\_Atomic}}} adalah type qualifier!)

\index{atomic\_bool type@\texttt{atomic\_bool} type} \index{atomic\_char type@\texttt{atomic\_char} type} \index{atomic\_schar type@\texttt{atomic\_schar} type} \index{atomic\_uchar type@\texttt{atomic\_uchar} type} \index{atomic\_short type@\texttt{atomic\_short} type} \index{atomic\_ushort type@\texttt{atomic\_ushort} type} \index{atomic\_int type@\texttt{atomic\_int} type} \index{atomic\_uint type@\texttt{atomic\_uint} type} \index{atomic\_long type@\texttt{atomic\_long} type} \index{atomic\_ulong type@\texttt{atomic\_ulong} type} \index{atomic\_llong type@\texttt{atomic\_llong} type} \index{atomic\_ullong type@\texttt{atomic\_ullong} type} \index{atomic\_char16\_t type@\texttt{atomic\_char16\_t} type} \index{atomic\_char32\_t type@\texttt{atomic\_char32\_t} type} \index{atomic\_wchar\_t type@\texttt{atomic\_wchar\_t} type} \index{atomic\_int\_least8\_t type@\texttt{atomic\_int\_least8\_t} type} \index{atomic\_uint\_least8\_t type@\texttt{atomic\_uint\_least8\_t} type} \index{atomic\_int\_least16\_t type@\texttt{atomic\_int\_least16\_t} type} \index{atomic\_uint\_least16\_t type@\texttt{atomic\_uint\_least16\_t} type} \index{atomic\_int\_least32\_t type@\texttt{atomic\_int\_least32\_t} type} \index{atomic\_uint\_least32\_t type@\texttt{atomic\_uint\_least32\_t} type} \index{atomic\_int\_least64\_t type@\texttt{atomic\_int\_least64\_t} type} \index{atomic\_uint\_least64\_t type@\texttt{atomic\_uint\_least64\_t} type} \index{atomic\_int\_fast8\_t type@\texttt{atomic\_int\_fast8\_t} type} \index{atomic\_uint\_fast8\_t type@\texttt{atomic\_uint\_fast8\_t} type} \index{atomic\_int\_fast16\_t type@\texttt{atomic\_int\_fast16\_t} type} \index{atomic\_uint\_fast16\_t type@\texttt{atomic\_uint\_fast16\_t} type} \index{atomic\_int\_fast32\_t type@\texttt{atomic\_int\_fast32\_t} type} \index{atomic\_uint\_fast32\_t type@\texttt{atomic\_uint\_fast32\_t} type} \index{atomic\_int\_fast64\_t type@\texttt{atomic\_int\_fast64\_t} type} \index{atomic\_uint\_fast64\_t type@\texttt{atomic\_uint\_fast64\_t} type} \index{atomic\_intptr\_t type@\texttt{atomic\_intptr\_t} type} \index{atomic\_uintptr\_t type@\texttt{atomic\_uintptr\_t} type} \index{atomic\_size\_t type@\texttt{atomic\_size\_t} type} \index{atomic\_ptrdiff\_t type@\texttt{atomic\_ptrdiff\_t} type} \index{atomic\_intmax\_t type@\texttt{atomic\_intmax\_t} type} \index{atomic\_uintmax\_t type@\texttt{atomic\_uintmax\_t} type}

{\def\LTcaptype{none} % do not increment counter
\begin{longtable}[]{@{}ll@{}}
\toprule\noalign{}
Tipe atomic & Padanan bentuk panjang \\
\midrule\noalign{}
\endhead
\bottomrule\noalign{}
\endlastfoot
\hl{\mbox{\texttt{atomic\_bool}}} & \hl{\mbox{\texttt{\_Atomic\ \_Bool}}} \\
\hl{\mbox{\texttt{atomic\_char}}} & \hl{\mbox{\texttt{\_Atomic\ char}}} \\
\hl{\mbox{\texttt{atomic\_schar}}} & \hl{\mbox{\texttt{\_Atomic\ signed\ char}}} \\
\hl{\mbox{\texttt{atomic\_uchar}}} & \hl{\mbox{\texttt{\_Atomic\ unsigned\ char}}} \\
\hl{\mbox{\texttt{atomic\_short}}} & \hl{\mbox{\texttt{\_Atomic\ short}}} \\
\hl{\mbox{\texttt{atomic\_ushort}}} & \hl{\mbox{\texttt{\_Atomic\ unsigned\ short}}} \\
\hl{\mbox{\texttt{atomic\_int}}} & \hl{\mbox{\texttt{\_Atomic\ int}}} \\
\hl{\mbox{\texttt{atomic\_uint}}} & \hl{\mbox{\texttt{\_Atomic\ unsigned\ int}}} \\
\hl{\mbox{\texttt{atomic\_long}}} & \hl{\mbox{\texttt{\_Atomic\ long}}} \\
\hl{\mbox{\texttt{atomic\_ulong}}} & \hl{\mbox{\texttt{\_Atomic\ unsigned\ long}}} \\
\hl{\mbox{\texttt{atomic\_llong}}} & \hl{\mbox{\texttt{\_Atomic\ long\ long}}} \\
\hl{\mbox{\texttt{atomic\_ullong}}} & \hl{\mbox{\texttt{\_Atomic\ unsigned\ long\ long}}} \\
\hl{\mbox{\texttt{atomic\_char16\_t}}} & \hl{\mbox{\texttt{\_Atomic\ char16\_t}}} \\
\hl{\mbox{\texttt{atomic\_char32\_t}}} & \hl{\mbox{\texttt{\_Atomic\ char32\_t}}} \\
\hl{\mbox{\texttt{atomic\_wchar\_t}}} & \hl{\mbox{\texttt{\_Atomic\ wchar\_t}}} \\
\hl{\mbox{\texttt{atomic\_int\_least8\_t}}} & \hl{\mbox{\texttt{\_Atomic\ int\_least8\_t}}} \\
\hl{\mbox{\texttt{atomic\_uint\_least8\_t}}} & \hl{\mbox{\texttt{\_Atomic\ uint\_least8\_t}}} \\
\hl{\mbox{\texttt{atomic\_int\_least16\_t}}} & \hl{\mbox{\texttt{\_Atomic\ int\_least16\_t}}} \\
\hl{\mbox{\texttt{atomic\_uint\_least16\_t}}} & \hl{\mbox{\texttt{\_Atomic\ uint\_least16\_t}}} \\
\hl{\mbox{\texttt{atomic\_int\_least32\_t}}} & \hl{\mbox{\texttt{\_Atomic\ int\_least32\_t}}} \\
\hl{\mbox{\texttt{atomic\_uint\_least32\_t}}} & \hl{\mbox{\texttt{\_Atomic\ uint\_least32\_t}}} \\
\hl{\mbox{\texttt{atomic\_int\_least64\_t}}} & \hl{\mbox{\texttt{\_Atomic\ int\_least64\_t}}} \\
\hl{\mbox{\texttt{atomic\_uint\_least64\_t}}} & \hl{\mbox{\texttt{\_Atomic\ uint\_least64\_t}}} \\
\hl{\mbox{\texttt{atomic\_int\_fast8\_t}}} & \hl{\mbox{\texttt{\_Atomic\ int\_fast8\_t}}} \\
\hl{\mbox{\texttt{atomic\_uint\_fast8\_t}}} & \hl{\mbox{\texttt{\_Atomic\ uint\_fast8\_t}}} \\
\hl{\mbox{\texttt{atomic\_int\_fast16\_t}}} & \hl{\mbox{\texttt{\_Atomic\ int\_fast16\_t}}} \\
\hl{\mbox{\texttt{atomic\_uint\_fast16\_t}}} & \hl{\mbox{\texttt{\_Atomic\ uint\_fast16\_t}}} \\
\hl{\mbox{\texttt{atomic\_int\_fast32\_t}}} & \hl{\mbox{\texttt{\_Atomic\ int\_fast32\_t}}} \\
\hl{\mbox{\texttt{atomic\_uint\_fast32\_t}}} & \hl{\mbox{\texttt{\_Atomic\ uint\_fast32\_t}}} \\
\hl{\mbox{\texttt{atomic\_int\_fast64\_t}}} & \hl{\mbox{\texttt{\_Atomic\ int\_fast64\_t}}} \\
\hl{\mbox{\texttt{atomic\_uint\_fast64\_t}}} & \hl{\mbox{\texttt{\_Atomic\ uint\_fast64\_t}}} \\
\hl{\mbox{\texttt{atomic\_intptr\_t}}} & \hl{\mbox{\texttt{\_Atomic\ intptr\_t}}} \\
\hl{\mbox{\texttt{atomic\_uintptr\_t}}} & \hl{\mbox{\texttt{\_Atomic\ uintptr\_t}}} \\
\hl{\mbox{\texttt{atomic\_size\_t}}} & \hl{\mbox{\texttt{\_Atomic\ size\_t}}} \\
\hl{\mbox{\texttt{atomic\_ptrdiff\_t}}} & \hl{\mbox{\texttt{\_Atomic\ ptrdiff\_t}}} \\
\hl{\mbox{\texttt{atomic\_intmax\_t}}} & \hl{\mbox{\texttt{\_Atomic\ intmax\_t}}} \\
\hl{\mbox{\texttt{atomic\_uintmax\_t}}} & \hl{\mbox{\texttt{\_Atomic\ uintmax\_t}}} \\
\end{longtable}
}

\index{\_Atomic type qualifier@\texttt{\_Atomic} type qualifier|)}

Pakai saja sesukamu! Tipe-tipe ini konsisten dengan alias atomic di C++, kalau itu membantu.

Tapi bagaimana kalau kamu ingin lebih?

Kamu bisa melakukannya dengan type qualifier atau type specifier.

\index{\_Atomic type specifier@\texttt{\_Atomic} type specifier|(}

Pertama, specifier! Ia adalah kata kunci \hl{\mbox{\texttt{\_Atomic}}} dengan tipe dalam kurung setelahnya\footnote{Tampaknya C++23 menambahkan ini sebagai makro.}---cocok dipakai dengan \hl{\mbox{\texttt{typedef}}}:

\begin{tcolorbox}[enhanced jigsaw,left=3pt,breakable,colback=white!97.2549!black,boxrule=0.2mm]

\begin{Shaded}
\begin{Highlighting}[]
\KeywordTok{typedef} \KeywordTok{\_Atomic}\OperatorTok{(}\DataTypeTok{double}\OperatorTok{)}\NormalTok{ atomic\_double}\OperatorTok{;}

\NormalTok{atomic\_double f}\OperatorTok{;}
\end{Highlighting}
\end{Shaded}

\end{tcolorbox}

Batasan pada specifier: tipe yang dijadikan atomic tidak boleh bertipe array atau fungsi, dan tidak boleh atomic atau sudah ber-qualifier lain.

\index{\_Atomic type specifier@\texttt{\_Atomic} type specifier|)} \index{\_Atomic type qualifier@\texttt{\_Atomic} type qualifier|(}

Berikutnya, qualifier! Ia adalah kata kunci \hl{\mbox{\texttt{\_Atomic}}} \emph{tanpa} tipe dalam kurung.

Jadi keduanya melakukan hal serupa\footnote{Spesifikasi mencatat bahwa keduanya mungkin berbeda dalam ukuran, representasi, dan alignment.}:

\begin{tcolorbox}[enhanced jigsaw,left=3pt,breakable,colback=white!97.2549!black,boxrule=0.2mm]

\begin{Shaded}
\begin{Highlighting}[]
\KeywordTok{\_Atomic}\OperatorTok{(}\DataTypeTok{int}\OperatorTok{)}\NormalTok{ i}\OperatorTok{;}   \CommentTok{// type specifier}
\KeywordTok{\_Atomic} \DataTypeTok{int}\NormalTok{  j}\OperatorTok{;}   \CommentTok{// type qualifier}
\end{Highlighting}
\end{Shaded}

\end{tcolorbox}

Bedanya, kamu bisa menyertakan type qualifier lain pada yang terakhir:

\begin{tcolorbox}[enhanced jigsaw,left=3pt,breakable,colback=white!97.2549!black,boxrule=0.2mm]

\begin{Shaded}
\begin{Highlighting}[]
\KeywordTok{\_Atomic} \DataTypeTok{volatile} \DataTypeTok{int}\NormalTok{ k}\OperatorTok{;}   \CommentTok{// variabel atomic ber{-}qualifier}
\end{Highlighting}
\end{Shaded}

\end{tcolorbox}

Batasan pada qualifier: tipe yang dijadikan atomic tidak boleh bertipe array atau fungsi.

\index{\_Atomic type qualifier@\texttt{\_Atomic} type qualifier|)}

\section{Variabel Atomic Lock-Free}\label{lock-free-atomic}

\index{Atomic variables!lock-free|(}

Arsitektur hardware terbatas dalam jumlah data yang bisa dibaca dan ditulis secara atomic. Itu tergantung bagaimana ia dirakit. Dan bervariasi.

Kalau kamu memakai tipe atomic, kamu bisa yakin bahwa akses ke tipe itu akan atomic\ldots{} tapi ada syaratnya: kalau hardware tidak bisa melakukannya, ia dilakukan dengan lock sebagai gantinya.

Jadi akses atomic menjadi lock-akses-unlock, yang cukup lebih lambat dan punya implikasi untuk signal handler.

\hyperref[atomic-flags]{Atomic flag}, di bawah, adalah satu-satunya tipe atomic yang dijamin lock-free di semua implementasi yang patuh. Di dunia komputer desktop/laptop biasa, tipe lain yang lebih besar kemungkinan lock-free.

Untungnya, ada beberapa cara untuk menentukan apakah tipe tertentu adalah atomic lock-free atau bukan.

Pertama, beberapa makro---ini bisa dipakai saat compile time dengan \hl{\mbox{\texttt{\#if}}}. Berlaku untuk tipe signed maupun unsigned.

\index{ATOMIC\_BOOL\_LOCK\_FREE macro@\texttt{ATOMIC\_BOOL\_LOCK\_FREE} macro} \index{ATOMIC\_CHAR\_LOCK\_FREE macro@\texttt{ATOMIC\_CHAR\_LOCK\_FREE} macro} \index{ATOMIC\_CHAR16\_T\_LOCK\_FREE macro@\texttt{ATOMIC\_CHAR16\_T\_LOCK\_FREE} macro} \index{ATOMIC\_CHAR32\_T\_LOCK\_FREE macro@\texttt{ATOMIC\_CHAR32\_T\_LOCK\_FREE} macro} \index{ATOMIC\_WCHAR\_T\_LOCK\_FREE macro@\texttt{ATOMIC\_WCHAR\_T\_LOCK\_FREE} macro} \index{ATOMIC\_SHORT\_LOCK\_FREE macro@\texttt{ATOMIC\_SHORT\_LOCK\_FREE} macro} \index{ATOMIC\_INT\_LOCK\_FREE macro@\texttt{ATOMIC\_INT\_LOCK\_FREE} macro} \index{ATOMIC\_LONG\_LOCK\_FREE macro@\texttt{ATOMIC\_LONG\_LOCK\_FREE} macro} \index{ATOMIC\_LLONG\_LOCK\_FREE macro@\texttt{ATOMIC\_LLONG\_LOCK\_FREE} macro} \index{ATOMIC\_POINTER\_LOCK\_FREE macro@\texttt{ATOMIC\_POINTER\_LOCK\_FREE} macro}

{\def\LTcaptype{none} % do not increment counter
\begin{longtable}[]{@{}ll@{}}
\toprule\noalign{}
Tipe Atomic & Makro Lock Free \\
\midrule\noalign{}
\endhead
\bottomrule\noalign{}
\endlastfoot
\hl{\mbox{\texttt{atomic\_bool}}} & \hl{\mbox{\texttt{ATOMIC\_BOOL\_LOCK\_FREE}}} \\
\hl{\mbox{\texttt{atomic\_char}}} & \hl{\mbox{\texttt{ATOMIC\_CHAR\_LOCK\_FREE}}} \\
\hl{\mbox{\texttt{atomic\_char16\_t}}} & \hl{\mbox{\texttt{ATOMIC\_CHAR16\_T\_LOCK\_FREE}}} \\
\hl{\mbox{\texttt{atomic\_char32\_t}}} & \hl{\mbox{\texttt{ATOMIC\_CHAR32\_T\_LOCK\_FREE}}} \\
\hl{\mbox{\texttt{atomic\_wchar\_t}}} & \hl{\mbox{\texttt{ATOMIC\_WCHAR\_T\_LOCK\_FREE}}} \\
\hl{\mbox{\texttt{atomic\_short}}} & \hl{\mbox{\texttt{ATOMIC\_SHORT\_LOCK\_FREE}}} \\
\hl{\mbox{\texttt{atomic\_int}}} & \hl{\mbox{\texttt{ATOMIC\_INT\_LOCK\_FREE}}} \\
\hl{\mbox{\texttt{atomic\_long}}} & \hl{\mbox{\texttt{ATOMIC\_LONG\_LOCK\_FREE}}} \\
\hl{\mbox{\texttt{atomic\_llong}}} & \hl{\mbox{\texttt{ATOMIC\_LLONG\_LOCK\_FREE}}} \\
\hl{\mbox{\texttt{atomic\_intptr\_t}}} & \hl{\mbox{\texttt{ATOMIC\_POINTER\_LOCK\_FREE}}} \\
\end{longtable}
}

Makro-makro ini menariknya bisa punya \emph{tiga} nilai berbeda:

{\def\LTcaptype{none} % do not increment counter
\begin{longtable}[]{@{}ll@{}}
\toprule\noalign{}
Nilai & Arti \\
\midrule\noalign{}
\endhead
\bottomrule\noalign{}
\endlastfoot
\hl{\mbox{\texttt{0}}} & Tidak pernah lock-free. \\
\hl{\mbox{\texttt{1}}} & \emph{Kadang-kadang} lock-free. \\
\hl{\mbox{\texttt{2}}} & Selalu lock-free. \\
\end{longtable}
}

Tunggu---bagaimana sesuatu bisa \emph{kadang-kadang} lock-free? Ini hanya berarti jawabannya tidak diketahui saat compile time, tapi bisa diketahui nanti saat runtime. Mungkin jawabannya bervariasi tergantung apakah kamu menjalankan kode ini di Genuine Intel atau AMD, atau semacam itu\footnote{Aku mengarang contoh itu begitu saja. Mungkin tidak berpengaruh di Intel/AMD, tapi bisa berpengaruh di suatu tempat, astaga!}.

Tapi kamu selalu bisa menguji saat runtime dengan fungsi \index{atomic\_is\_lock\_free() function@\texttt{atomic\_is\_lock\_free()} function} \hl{\mbox{\texttt{atomic\_is\_lock\_free()}}}. Fungsi ini mengembalikan true atau false apakah tipe tertentu atomic saat ini.

Jadi kenapa kita peduli?

Lock-free lebih cepat, jadi mungkin ada kekhawatiran kecepatan yang akan kamu atasi dengan cara lain. Atau mungkin kamu perlu memakai variabel atomic di signal handler.

\index{Atomic variables!lock-free|)}

\subsection{Signal Handler dan Atomic Lock-Free}\label{signal-handler-dan-atomic-lock-free}

\index{Signal handlers!with lock-free atomics|(} \index{Atomic variables!with signal handlers|(}

Kalau kamu membaca atau menulis variabel bersama (static storage duration atau \hl{\mbox{\texttt{\_Thread\_Local}}}) di signal handler, itu undefined behavior {[}astaga!{]}\ldots{} Kecuali kamu melakukan salah satu dari berikut:

\begin{enumerate}
\def\labelenumi{\arabic{enumi}.}
\item
  Menulis ke variabel bertipe \hl{\mbox{\texttt{volatile\ sig\_atomic\_t}}}.
\item
  Membaca atau menulis variabel atomic lock-free.
\end{enumerate}

Sejauh yang kutahu, variabel atomic lock-free adalah salah satu dari sedikit cara untuk mengambil informasi dari signal handler secara portabel.

Menurut bacaanku, spesifikasi agak samar soal memory order saat meng-acquire atau me-release variabel atomic di signal handler. C++ menyatakan, dan itu masuk akal, bahwa akses semacam itu unsequenced terhadap sisa program\footnote{C++ menjelaskan bahwa kalau signal adalah hasil pemanggilan \index{raise() function@\texttt{raise()} function} \hl{\mbox{\texttt{raise()}}}, ia disekuensikan \emph{setelah} \hl{\mbox{\texttt{raise()}}}.}. Bagaimanapun, signal bisa di-raise kapan saja. Jadi aku menganggap perilaku C serupa.

\index{Signal handlers!with lock-free atomics|)} \index{Atomic variables!with signal handlers|)}

\section{Atomic Flag}\label{atomic-flags}

\index{Atomic variables!atomic flags|(} \index{atomic\_flag type@\texttt{atomic\_flag} type|(}

Hanya ada satu tipe yang dijamin standar sebagai atomic lock-free: \hl{\mbox{\texttt{atomic\_flag}}}. Ini tipe opaque untuk operasi \href{https://en.wikipedia.org/wiki/Test-and-set}{test-and-set}\footnote{https://en.wikipedia.org/wiki/Test-and-set}.

Ia bisa \emph{set} atau \emph{clear}. Kamu bisa menginisialisasinya menjadi clear dengan:

\index{ATOMIC\_FLAG\_INIT macro@\texttt{ATOMIC\_FLAG\_INIT} macro|(}

\begin{tcolorbox}[enhanced jigsaw,left=3pt,breakable,colback=white!97.2549!black,boxrule=0.2mm]

\begin{Shaded}
\begin{Highlighting}[]
\NormalTok{atomic\_flag f }\OperatorTok{=}\NormalTok{ ATOMIC\_FLAG\_INIT}\OperatorTok{;}
\end{Highlighting}
\end{Shaded}

\end{tcolorbox}

\index{ATOMIC\_FLAG\_INIT macro@\texttt{ATOMIC\_FLAG\_INIT} macro|)}

\index{atomic\_flag\_test\_and\_set() function@\texttt{atomic\_flag\_test\_and\_set()} function|(}

Kamu bisa men-set flag secara atomic dengan \hl{\mbox{\texttt{atomic\_flag\_test\_and\_set()}}}, yang akan men-set flag dan mengembalikan status sebelumnya sebagai \hl{\mbox{\texttt{\_Bool}}} (true untuk set).

\index{atomic\_flag\_clear() function@\texttt{atomic\_flag\_clear()} function|(}

Kamu bisa meng-clear flag secara atomic dengan \hl{\mbox{\texttt{atomic\_flag\_clear()}}}.

Berikut contoh ketika kita menginisialisasi flag menjadi clear, men-set-nya dua kali, lalu meng-clear-nya lagi.

\begin{tcolorbox}[enhanced jigsaw,left=3pt,breakable,colback=white!97.2549!black,boxrule=0.2mm]

\begin{Shaded}
\begin{Highlighting}[]
\PreprocessorTok{\#include }\ImportTok{\textless{}stdio.h\textgreater{}}
\PreprocessorTok{\#include }\ImportTok{\textless{}stdbool.h\textgreater{}}\PreprocessorTok{   }\CommentTok{// Tidak diperlukan di C23}
\PreprocessorTok{\#include }\ImportTok{\textless{}stdatomic.h\textgreater{}}

\NormalTok{atomic\_flag f }\OperatorTok{=}\NormalTok{ ATOMIC\_FLAG\_INIT}\OperatorTok{;}

\DataTypeTok{int}\NormalTok{ main}\OperatorTok{(}\DataTypeTok{void}\OperatorTok{)}
\OperatorTok{\{}
    \DataTypeTok{bool}\NormalTok{ r }\OperatorTok{=}\NormalTok{ atomic\_flag\_test\_and\_set}\OperatorTok{(\&}\NormalTok{f}\OperatorTok{);}
\NormalTok{    printf}\OperatorTok{(}\StringTok{"Value was: }\SpecialCharTok{\%d\textbackslash{}n}\StringTok{"}\OperatorTok{,}\NormalTok{ r}\OperatorTok{);}           \CommentTok{// 0}

\NormalTok{    r }\OperatorTok{=}\NormalTok{ atomic\_flag\_test\_and\_set}\OperatorTok{(\&}\NormalTok{f}\OperatorTok{);}
\NormalTok{    printf}\OperatorTok{(}\StringTok{"Value was: }\SpecialCharTok{\%d\textbackslash{}n}\StringTok{"}\OperatorTok{,}\NormalTok{ r}\OperatorTok{);}           \CommentTok{// 1}

\NormalTok{    atomic\_flag\_clear}\OperatorTok{(\&}\NormalTok{f}\OperatorTok{);}
\NormalTok{    r }\OperatorTok{=}\NormalTok{ atomic\_flag\_test\_and\_set}\OperatorTok{(\&}\NormalTok{f}\OperatorTok{);}
\NormalTok{    printf}\OperatorTok{(}\StringTok{"Value was: }\SpecialCharTok{\%d\textbackslash{}n}\StringTok{"}\OperatorTok{,}\NormalTok{ r}\OperatorTok{);}           \CommentTok{// 0}
\OperatorTok{\}}
\end{Highlighting}
\end{Shaded}

\end{tcolorbox}

\index{atomic\_flag\_clear() function@\texttt{atomic\_flag\_clear()} function|)} \index{atomic\_flag\_test\_and\_set() function@\texttt{atomic\_flag\_test\_and\_set()} function|)} \index{Atomic variables!atomic flags|)} \index{atomic\_flag type@\texttt{atomic\_flag} type|)}

\section{\texorpdfstring{\hl{\mbox{\texttt{struct}}} dan \hl{\mbox{\texttt{union}}} Atomic}{struct dan union Atomic}}\label{struct-dan-union-atomic}

\index{Atomic variables!struct and union@\texttt{struct} and \texttt{union}|(}

Dengan qualifier atau specifier \hl{\mbox{\texttt{\_Atomic}}}, kamu bisa membuat \hl{\mbox{\texttt{struct}}} atau \hl{\mbox{\texttt{union}}} atomic! Cukup mencengangkan.

Kalau tidak banyak data di dalamnya (yaitu segenggam byte), tipe atomic yang dihasilkan mungkin lock-free. Ujilah dengan \hl{\mbox{\texttt{atomic\_is\_lock\_free()}}}.

\begin{tcolorbox}[enhanced jigsaw,left=3pt,breakable,colback=white!97.2549!black,boxrule=0.2mm]

\begin{Shaded}
\begin{Highlighting}[numbers=left,,]
\PreprocessorTok{\#include }\ImportTok{\textless{}stdio.h\textgreater{}}
\PreprocessorTok{\#include }\ImportTok{\textless{}stdatomic.h\textgreater{}}

\DataTypeTok{int}\NormalTok{ main}\OperatorTok{(}\DataTypeTok{void}\OperatorTok{)}
\OperatorTok{\{}
    \KeywordTok{struct}\NormalTok{ point }\OperatorTok{\{}
        \DataTypeTok{float}\NormalTok{ x}\OperatorTok{,}\NormalTok{ y}\OperatorTok{;}
    \OperatorTok{\};}

    \KeywordTok{\_Atomic}\OperatorTok{(}\KeywordTok{struct}\NormalTok{ point}\OperatorTok{)}\NormalTok{ p}\OperatorTok{;}

\NormalTok{    printf}\OperatorTok{(}\StringTok{"Is lock free: }\SpecialCharTok{\%d\textbackslash{}n}\StringTok{"}\OperatorTok{,}\NormalTok{ atomic\_is\_lock\_free}\OperatorTok{(\&}\NormalTok{p}\OperatorTok{));}
\OperatorTok{\}}
\end{Highlighting}
\end{Shaded}

\end{tcolorbox}

Ini syaratnya: kamu tidak bisa mengakses field dari \hl{\mbox{\texttt{struct}}} atau \hl{\mbox{\texttt{union}}} atomic\ldots{} lalu apa gunanya? Nah, kamu bisa \emph{menyalin} seluruh \hl{\mbox{\texttt{struct}}} secara atomic ke variabel non-atomic lalu memakainya. Kamu juga bisa menyalin secara atomic ke arah sebaliknya.

\begin{tcolorbox}[enhanced jigsaw,left=3pt,breakable,colback=white!97.2549!black,boxrule=0.2mm]

\begin{Shaded}
\begin{Highlighting}[numbers=left,,]
\PreprocessorTok{\#include }\ImportTok{\textless{}stdio.h\textgreater{}}
\PreprocessorTok{\#include }\ImportTok{\textless{}stdatomic.h\textgreater{}}

\DataTypeTok{int}\NormalTok{ main}\OperatorTok{(}\DataTypeTok{void}\OperatorTok{)}
\OperatorTok{\{}
    \KeywordTok{struct}\NormalTok{ point }\OperatorTok{\{}
        \DataTypeTok{float}\NormalTok{ x}\OperatorTok{,}\NormalTok{ y}\OperatorTok{;}
    \OperatorTok{\};}

    \KeywordTok{\_Atomic}\OperatorTok{(}\KeywordTok{struct}\NormalTok{ point}\OperatorTok{)}\NormalTok{ p}\OperatorTok{;}
    \KeywordTok{struct}\NormalTok{ point t}\OperatorTok{;}

\NormalTok{    p }\OperatorTok{=} \OperatorTok{(}\KeywordTok{struct}\NormalTok{ point}\OperatorTok{)\{}\DecValTok{1}\OperatorTok{,} \DecValTok{2}\OperatorTok{\};}  \CommentTok{// Salinan atomic}

    \CommentTok{//printf("\%f\textbackslash{}n", p.x);  // Error}

\NormalTok{    t }\OperatorTok{=}\NormalTok{ p}\OperatorTok{;}   \CommentTok{// Salinan atomic}

\NormalTok{    printf}\OperatorTok{(}\StringTok{"}\SpecialCharTok{\%f\textbackslash{}n}\StringTok{"}\OperatorTok{,}\NormalTok{ t}\OperatorTok{.}\NormalTok{x}\OperatorTok{);}  \CommentTok{// OK!}
\OperatorTok{\}}
\end{Highlighting}
\end{Shaded}

\end{tcolorbox}

Kamu juga bisa mendeklarasikan \hl{\mbox{\texttt{struct}}} yang field-field individualnya atomic. Apakah tipe atomic diizinkan pada bitfield ditentukan oleh implementasi.

\index{Atomic variables!struct and union@\texttt{struct} and \texttt{union}|)}

\section{Pointer Atomic}\label{pointer-atomic}

\index{Atomic variables!pointers|(}

Sekadar catatan di sini tentang penempatan \hl{\mbox{\texttt{\_Atomic}}} untuk pointer.

Pertama, pointer ke atomic (yaitu nilai pointer-nya tidak atomic, tapi yang ditunjuknya atomic):

\begin{tcolorbox}[enhanced jigsaw,left=3pt,breakable,colback=white!97.2549!black,boxrule=0.2mm]

\begin{Shaded}
\begin{Highlighting}[]
\KeywordTok{\_Atomic} \DataTypeTok{int}\NormalTok{ x}\OperatorTok{;}
\KeywordTok{\_Atomic} \DataTypeTok{int} \OperatorTok{*}\NormalTok{p}\OperatorTok{;}  \CommentTok{// p adalah pointer ke int atomic}

\NormalTok{p }\OperatorTok{=} \OperatorTok{\&}\NormalTok{x}\OperatorTok{;}  \CommentTok{// OK!}
\end{Highlighting}
\end{Shaded}

\end{tcolorbox}

Kedua, pointer atomic ke nilai non-atomic (yaitu nilai pointer itu sendiri atomic, tapi yang ditunjuknya tidak):

\begin{tcolorbox}[enhanced jigsaw,left=3pt,breakable,colback=white!97.2549!black,boxrule=0.2mm]

\begin{Shaded}
\begin{Highlighting}[]
\DataTypeTok{int}\NormalTok{ x}\OperatorTok{;}
\DataTypeTok{int} \OperatorTok{*} \KeywordTok{\_Atomic}\NormalTok{ p}\OperatorTok{;}  \CommentTok{// p adalah pointer atomic ke int}

\NormalTok{p }\OperatorTok{=} \OperatorTok{\&}\NormalTok{x}\OperatorTok{;}  \CommentTok{// OK!}
\end{Highlighting}
\end{Shaded}

\end{tcolorbox}

Terakhir, pointer atomic ke nilai atomic (yaitu pointer dan yang ditunjuknya sama-sama atomic):

\begin{tcolorbox}[enhanced jigsaw,left=3pt,breakable,colback=white!97.2549!black,boxrule=0.2mm]

\begin{Shaded}
\begin{Highlighting}[]
\KeywordTok{\_Atomic} \DataTypeTok{int}\NormalTok{ x}\OperatorTok{;}
\KeywordTok{\_Atomic} \DataTypeTok{int} \OperatorTok{*} \KeywordTok{\_Atomic}\NormalTok{ p}\OperatorTok{;}  \CommentTok{// p adalah pointer atomic ke int atomic}

\NormalTok{p }\OperatorTok{=} \OperatorTok{\&}\NormalTok{x}\OperatorTok{;}  \CommentTok{// OK!}
\end{Highlighting}
\end{Shaded}

\end{tcolorbox}

\index{Atomic variables!pointers|)}

\section{Memory Order}\label{memory-order}

\index{Atomic variables!memory order|(} \index{Memory order|(}

Kita sudah membahas sequential consistency, yang paling masuk akal di antara semuanya. Tapi ada beberapa yang lain:

\index{memory\_order\_seq\_cst enumerated type@\texttt{memory\_order\_seq\_cst} enumerated type} \index{memory\_order\_acq\_rel enumerated type@\texttt{memory\_order\_acq\_rel} enumerated type} \index{memory\_order\_release enumerated type@\texttt{memory\_order\_release} enumerated type} \index{memory\_order\_acquire enumerated type@\texttt{memory\_order\_acquire} enumerated type} \index{memory\_order\_consume enumerated type@\texttt{memory\_order\_consume} enumerated type} \index{memory\_order\_relaxed enumerated type@\texttt{memory\_order\_relaxed} enumerated type}

{\def\LTcaptype{none} % do not increment counter
\begin{longtable}[]{@{}ll@{}}
\toprule\noalign{}
\hl{\mbox{\texttt{memory\_order}}} & Deskripsi \\
\midrule\noalign{}
\endhead
\bottomrule\noalign{}
\endlastfoot
\hl{\mbox{\texttt{memory\_order\_seq\_cst}}} & Sequential Consistency \\
\hl{\mbox{\texttt{memory\_order\_acq\_rel}}} & Acquire/Release \\
\hl{\mbox{\texttt{memory\_order\_release}}} & Release \\
\hl{\mbox{\texttt{memory\_order\_acquire}}} & Acquire \\
\hl{\mbox{\texttt{memory\_order\_consume}}} & Consume \\
\hl{\mbox{\texttt{memory\_order\_relaxed}}} & Relaxed \\
\end{longtable}
}

Kamu bisa menentukan yang lain dengan fungsi library tertentu. Misalnya, kamu bisa menambahkan nilai ke variabel atomic seperti ini:

\begin{tcolorbox}[enhanced jigsaw,left=3pt,breakable,colback=white!97.2549!black,boxrule=0.2mm]

\begin{Shaded}
\begin{Highlighting}[]
\NormalTok{atomic\_int x }\OperatorTok{=} \DecValTok{0}\OperatorTok{;}

\NormalTok{x }\OperatorTok{+=} \DecValTok{5}\OperatorTok{;}  \CommentTok{// Sequential consistency, yang bawaan}
\end{Highlighting}
\end{Shaded}

\end{tcolorbox}

Atau kamu bisa melakukan hal yang sama dengan fungsi library ini:

\index{atomic\_fetch\_add() function@\texttt{atomic\_fetch\_add()} function|(}

\begin{tcolorbox}[enhanced jigsaw,left=3pt,breakable,colback=white!97.2549!black,boxrule=0.2mm]

\begin{Shaded}
\begin{Highlighting}[]
\NormalTok{atomic\_int x }\OperatorTok{=} \DecValTok{0}\OperatorTok{;}

\NormalTok{atomic\_fetch\_add}\OperatorTok{(\&}\NormalTok{x}\OperatorTok{,} \DecValTok{5}\OperatorTok{);}  \CommentTok{// Sequential consistency, yang bawaan}
\end{Highlighting}
\end{Shaded}

\end{tcolorbox}

\index{atomic\_fetch\_add() function@\texttt{atomic\_fetch\_add()} function|)}

Atau kamu bisa melakukan hal yang sama dengan memory ordering eksplisit:

\index{atomic\_fetch\_add\_explicit() function@\texttt{atomic\_fetch\_add\_explicit()} function|(}

\begin{tcolorbox}[enhanced jigsaw,left=3pt,breakable,colback=white!97.2549!black,boxrule=0.2mm]

\begin{Shaded}
\begin{Highlighting}[]
\NormalTok{atomic\_int x }\OperatorTok{=} \DecValTok{0}\OperatorTok{;}

\NormalTok{atomic\_fetch\_add\_explicit}\OperatorTok{(\&}\NormalTok{x}\OperatorTok{,} \DecValTok{5}\OperatorTok{,}\NormalTok{ memory\_order\_seq\_cst}\OperatorTok{);}
\end{Highlighting}
\end{Shaded}

\end{tcolorbox}

Tapi bagaimana kalau kita tidak menginginkan sequential consistency? Dan kamu menginginkan acquire/release karena suatu alasan? Sebutkan saja:

\begin{tcolorbox}[enhanced jigsaw,left=3pt,breakable,colback=white!97.2549!black,boxrule=0.2mm]

\begin{Shaded}
\begin{Highlighting}[]
\NormalTok{atomic\_int x }\OperatorTok{=} \DecValTok{0}\OperatorTok{;}

\NormalTok{atomic\_fetch\_add\_explicit}\OperatorTok{(\&}\NormalTok{x}\OperatorTok{,} \DecValTok{5}\OperatorTok{,}\NormalTok{ memory\_order\_acq\_rel}\OperatorTok{);}
\end{Highlighting}
\end{Shaded}

\end{tcolorbox}

\index{atomic\_fetch\_add\_explicit() function@\texttt{atomic\_fetch\_add\_explicit()} function|)}

Kita akan membahas rincian memory order yang berbeda di bawah. Jangan mengutak-atik apa pun selain sequential consistency kecuali kamu tahu apa yang kamu lakukan. Sangat mudah membuat kesalahan yang menyebabkan kegagalan langka yang sulit direproduksi.

\subsection{Sequential Consistency}\label{sequential-consistency-1}

\index{Atomic variables!sequential consistency|(} \index{Memory order!sequential consistency|(}

\begin{itemize}
\tightlist
\item
  Operasi load melakukan acquire (lihat di bawah).
\item
  Operasi store melakukan release (lihat di bawah).
\item
  Operasi read-modify-write melakukan acquire lalu release.
\end{itemize}

Selain itu, untuk menjaga urutan total acquire dan release, tidak ada acquire atau release yang akan ditata ulang terhadap satu sama lain. (Aturan acquire/release tidak melarang penataan ulang release yang diikuti acquire. Tapi aturan sequentially consistent melarangnya.)

\index{Memory order!sequential consistency|)} \index{Atomic variables!sequential consistency|)}

\subsection{Acquire}\label{acquire}

\index{Atomic variables!acquire|(} \index{Memory order!acquire|(}

Inilah yang terjadi pada operasi load/read variabel atomic.

\begin{itemize}
\item
  Kalau thread lain me-release variabel atomic ini, semua penulisan yang dilakukan thread itu kini terlihat di thread ini.
\item
  Akses memori di thread ini yang terjadi setelah load ini tidak bisa ditata ulang ke sebelumnya.
\end{itemize}

\index{Memory order!acquire|)} \index{Atomic variables!acquire|)}

\subsection{Release}\label{release}

\index{Atomic variables!release|(} \index{Memory order!acquire|(}

Inilah yang terjadi pada store/write variabel atomic.

\begin{itemize}
\item
  Kalau thread lain nanti meng-acquire variabel atomic ini, semua penulisan memori di thread ini sebelum penulisan atomic-nya menjadi terlihat oleh thread lain itu.
\item
  Akses memori di thread ini yang terjadi sebelum release tidak bisa ditata ulang ke setelahnya.
\end{itemize}

\index{Atomic variables!release|)} \index{Memory order!release|)}

\subsection{Consume}\label{consume}

\index{Atomic variables!consume|(} \index{Memory order!consume|(}

Yang ini aneh, mirip versi acquire yang kurang ketat. Ia memengaruhi akses memori yang \emph{bergantung data} pada variabel atomic.

Menjadi ``bergantung data'' secara samar berarti variabel atomic dipakai dalam suatu perhitungan.

Artinya, kalau sebuah thread meng-consume variabel atomic, maka semua operasi di thread itu yang kemudian memakai variabel atomic tersebut akan bisa melihat penulisan memori di thread yang melakukan release.

Bandingkan dengan acquire, ketika penulisan memori di thread yang melakukan release akan terlihat oleh \emph{semua} operasi di thread saat ini, bukan hanya yang bergantung data.

Seperti acquire juga, ada batasan operasi mana yang boleh ditata ulang \emph{sebelum} consume. Dengan acquire, tidak ada yang boleh ditata ulang ke sebelumnya. Dengan consume, tidak ada yang bergantung pada nilai atomic yang dimuat yang boleh ditata ulang ke sebelumnya.

\index{Atomic variables!consume|)} \index{Memory order!consume|)}

\subsection{Acquire/Release}\label{acquirerelease}

\index{Atomic variables!acquire/release|(} \index{Memory order!acquire/release|(}

Ini hanya berlaku untuk operasi read-modify-write. Ia acquire dan release yang digabung menjadi satu.

\begin{itemize}
\tightlist
\item
  Acquire terjadi untuk pembacaan.
\item
  Release terjadi untuk penulisan.
\end{itemize}

\index{Atomic variables!acquire/release|)} \index{Memory order!acquire/release|)}

\subsection{Relaxed}\label{relaxed}

\index{Atomic variables!relaxed|(} \index{Memory order!relaxed|(}

Tidak ada aturan; ini anarki! Semua orang boleh menata ulang semuanya di mana saja! Anjing dan kucing hidup bersama---histeria massal!

Sebenarnya ada satu aturan. Pembacaan dan penulisan atomic tetap semua-atau-tidak-sama-sekali. Tapi operasi bisa ditata ulang sesuka hati dan sama sekali tidak ada sinkronisasi antar thread.

Ada beberapa kasus pemakaian untuk memory order ini, yang bisa kamu temukan dengan sedikit pencarian, mis. counter sederhana.

Dan kamu bisa memakai fence untuk memaksa sinkronisasi setelah sekumpulan penulisan relaxed.

\index{Atomic variables!relaxed|)} \index{Memory order!relaxed|)} \index{Memory order|)} \index{Atomic variables!memory order|)}

\section{Fence}\label{fence}

\index{Atomic variables!fences|(}

Kamu tahu bagaimana release dan acquire variabel atomic terjadi saat kamu membaca dan menulisnya?

Nah, mungkin juga melakukan release atau acquire \emph{tanpa} variabel atomic.

Ini disebut \emph{fence}. Jadi kalau kamu ingin semua penulisan di sebuah thread terlihat di tempat lain, kamu bisa memasang release fence di satu thread dan acquire fence di thread lain, persis seperti cara kerja variabel atomic.

Karena operasi consume tidak benar-benar masuk akal pada fence\footnote{Karena consume sepenuhnya tentang operasi yang bergantung pada nilai variabel atomic yang di-acquire, dan tidak ada variabel atomic pada fence.}, \hl{\mbox{\texttt{memory\_order\_consume}}} diperlakukan sebagai acquire.

Kamu bisa memasang fence dengan order apa pun yang ditentukan:

\index{atomic\_thread\_fence() function@\texttt{atomic\_thread\_fence()} function|(}

\begin{tcolorbox}[enhanced jigsaw,left=3pt,breakable,colback=white!97.2549!black,boxrule=0.2mm]

\begin{Shaded}
\begin{Highlighting}[]
\NormalTok{atomic\_thread\_fence}\OperatorTok{(}\NormalTok{memory\_order\_release}\OperatorTok{);}
\end{Highlighting}
\end{Shaded}

\end{tcolorbox}

\index{atomic\_thread\_fence() function@\texttt{atomic\_thread\_fence()} function|)} \index{atomic\_signal\_fence() function@\texttt{atomic\_signal\_fence()} function|(}

Ada juga versi ringan fence untuk dipakai dengan signal handler, bernama \hl{\mbox{\texttt{atomic\_signal\_fence()}}}.

Cara kerjanya sama persis dengan \hl{\mbox{\texttt{atomic\_thread\_fence()}}}, kecuali:

\begin{itemize}
\item
  Ia hanya menangani visibilitas nilai dalam thread yang sama; tidak ada sinkronisasi dengan thread lain.
\item
  Tidak ada instruksi fence hardware yang dikeluarkan.
\end{itemize}

Kalau kamu ingin yakin efek samping operasi non-atomic (dan operasi atomic relaxed) terlihat di signal handler, kamu bisa memakai fence ini.

Idenya: signal handler dieksekusi di thread \emph{ini}, bukan thread lain, jadi ini cara yang lebih ringan untuk memastikan perubahan di luar signal handler terlihat di dalamnya (yaitu belum ditata ulang).

\index{atomic\_signal\_fence() function@\texttt{atomic\_signal\_fence()} function|)} \index{Atomic variables!fences|)}

\section{Referensi}\label{referensi}

Kalau kamu ingin mempelajari hal-hal ini lebih jauh, berikut beberapa bahan yang membantuku menembusnya:

\begin{itemize}
\item
  Ceramah Herb Sutter \emph{\hl{\mbox{\texttt{atomic\textless{}\textgreater{}}}} Weapons}:

  \begin{itemize}
  \tightlist
  \item
    \href{https://www.youtube.com/watch?v=A8eCGOqgvH4}{Bagian 1}\footnote{https://www.youtube.com/watch?v=A8eCGOqgvH4}
  \item
    \href{https://www.youtube.com/watch?v=KeLBd2EJLOU}{Bagian 2}\footnote{https://www.youtube.com/watch?v=KeLBd2EJLOU}
  \end{itemize}
\item
  \href{https://preshing.com/archives/}{Materi Jeff Preshing}\footnote{https://preshing.com/archives/}, khususnya:

  \begin{itemize}
  \tightlist
  \item
    \href{https://preshing.com/20120612/an-introduction-to-lock-free-programming/}{An Introduction to Lock-Free Programming}\footnote{https://preshing.com/20120612/an-introduction-to-lock-free-programming/}
  \item
    \href{https://preshing.com/20120913/acquire-and-release-semantics/}{Acquire and Release Semantics}\footnote{https://preshing.com/20120913/acquire-and-release-semantics/}
  \item
    \href{https://preshing.com/20130702/the-happens-before-relation/}{The \emph{Happens-Before} Relation}\footnote{https://preshing.com/20130702/the-happens-before-relation/}
  \item
    \href{https://preshing.com/20130823/the-synchronizes-with-relation/}{The \emph{Synchronizes-With} Relation}\footnote{https://preshing.com/20130823/the-synchronizes-with-relation/}
  \item
    \href{https://preshing.com/20140709/the-purpose-of-memory_order_consume-in-cpp11/}{The Purpose of \hl{\mbox{\texttt{memory\_order\_consume}}} in C++11}\footnote{https://preshing.com/20140709/the-purpose-of-memory\_order\_consume-in-cpp11/}
  \item
    \href{https://preshing.com/20150402/you-can-do-any-kind-of-atomic-read-modify-write-operation/}{You Can Do Any Kind of Atomic Read-Modify-Write Operation}\footnote{https://preshing.com/20150402/you-can-do-any-kind-of-atomic-read-modify-write-operation/}
  \end{itemize}
\item
  CPPReference:

  \begin{itemize}
  \tightlist
  \item
    \href{https://en.cppreference.com/w/c/atomic/memory_order}{Memory Order}\footnote{https://en.cppreference.com/w/c/atomic/memory\_order}
  \item
    \href{https://en.cppreference.com/w/c/language/atomic}{Atomic Types}\footnote{https://en.cppreference.com/w/c/language/atomic}
  \end{itemize}
\item
  \href{https://docs.microsoft.com/en-us/windows/win32/dxtecharts/lockless-programming}{Lockless Programming Considerations}\footnote{https://docs.microsoft.com/en-us/windows/win32/dxtecharts/lockless-programming} oleh Bruce Dawson
\item
  Orang-orang yang suka membantu dan berpengetahuan di \href{https://www.reddit.com/r/C_Programming/}{r/C\_Programming}\footnote{https://www.reddit.com/r/C\_Programming/}
\end{itemize}

\index{Atomic variables|)}

\chapter{Function Specifier, Alignment Specifier dan Operator}\label{function-specifier-alignment-specifier-dan-operator}

Menurut pengalamanku, ini tidak terlalu sering dipakai, tapi kita bahas di sini demi kelengkapan.

\section{Function Specifier}\label{function-specifier}

\index{Function specifiers|(}

Saat mendeklarasikan fungsi, kamu bisa memberi compiler beberapa petunjuk tentang bagaimana fungsi itu mungkin atau akan dipakai. Ini memungkinkan atau mendorong compiler melakukan optimasi tertentu.

\subsection{\texorpdfstring{\hl{\mbox{\texttt{inline}}} untuk Kecepatan---Mungkin}{inline untuk Kecepatan---Mungkin}}\label{inline-untuk-kecepatanmungkin}

\index{inline function specifier@\texttt{inline} function specifier|(}

Kamu bisa mendeklarasikan fungsi sebagai inline seperti ini:

\begin{tcolorbox}[enhanced jigsaw,left=3pt,breakable,colback=white!97.2549!black,boxrule=0.2mm]

\begin{Shaded}
\begin{Highlighting}[]
\DataTypeTok{static} \KeywordTok{inline} \DataTypeTok{int}\NormalTok{ add}\OperatorTok{(}\DataTypeTok{int}\NormalTok{ x}\OperatorTok{,} \DataTypeTok{int}\NormalTok{ y}\OperatorTok{)} \OperatorTok{\{}
    \ControlFlowTok{return}\NormalTok{ x }\OperatorTok{+}\NormalTok{ y}\OperatorTok{;}
\OperatorTok{\}}
\end{Highlighting}
\end{Shaded}

\end{tcolorbox}

Ini dimaksudkan untuk mendorong compiler membuat pemanggilan fungsi ini secepat mungkin. Dan, secara historis, salah satu caranya adalah \emph{inlining}, artinya badan fungsi disematkan seluruhnya di tempat pemanggilan dilakukan. Ini menghindari semua overhead menyiapkan dan membongkar pemanggilan fungsi, dengan konsekuensi ukuran kode lebih besar karena fungsi disalin ke mana-mana, bukan dipakai ulang.

Hal-hal cepat dan kasar yang perlu diingat:

\begin{enumerate}
\def\labelenumi{\arabic{enumi}.}
\item
  Kamu kemungkinan tidak perlu memakai \hl{\mbox{\texttt{inline}}} untuk kecepatan. Compiler modern tahu mana yang terbaik.
\item
  Kalau kamu memakainya untuk kecepatan, pakai dengan file scope, yaitu \hl{\mbox{\texttt{static\ inline}}}. Ini menghindari aturan rumit external linkage dan fungsi inline.
\end{enumerate}

Berhenti membaca bagian ini sekarang.

Suka menyiksa diri, ya?

Mari kita coba menghilangkan \hl{\mbox{\texttt{static}}}.

\begin{tcolorbox}[enhanced jigsaw,left=3pt,breakable,colback=white!97.2549!black,boxrule=0.2mm]

\begin{Shaded}
\begin{Highlighting}[numbers=left,,]
\PreprocessorTok{\#include }\ImportTok{\textless{}stdio.h\textgreater{}}

\KeywordTok{inline} \DataTypeTok{int}\NormalTok{ add}\OperatorTok{(}\DataTypeTok{int}\NormalTok{ x}\OperatorTok{,} \DataTypeTok{int}\NormalTok{ y}\OperatorTok{)}
\OperatorTok{\{}
    \ControlFlowTok{return}\NormalTok{ x }\OperatorTok{+}\NormalTok{ y}\OperatorTok{;}
\OperatorTok{\}}

\DataTypeTok{int}\NormalTok{ main}\OperatorTok{(}\DataTypeTok{void}\OperatorTok{)}
\OperatorTok{\{}
\NormalTok{    printf}\OperatorTok{(}\StringTok{"}\SpecialCharTok{\%d\textbackslash{}n}\StringTok{"}\OperatorTok{,}\NormalTok{ add}\OperatorTok{(}\DecValTok{1}\OperatorTok{,} \DecValTok{2}\OperatorTok{));}
\OperatorTok{\}}
\end{Highlighting}
\end{Shaded}

\end{tcolorbox}

\hl{\mbox{\texttt{gcc}}} memberi error linker pada \hl{\mbox{\texttt{add()}}}\footnote{Kecuali kamu mengompilasi dengan optimasi aktif (kemungkinan)! Tapi kurasa, ketika itu terjadi, perilakunya tidak sesuai spesifikasi.}. Spesifikasi mewajibkan bahwa kalau kamu punya fungsi inline non-\hl{\mbox{\texttt{extern}}}, kamu juga harus menyediakan versi dengan external linkage.

Jadi kamu harus punya versi \hl{\mbox{\texttt{extern}}} di tempat lain agar ini berfungsi. Kalau compiler punya fungsi \hl{\mbox{\texttt{inline}}} di file saat ini dan versi eksternal dari fungsi yang sama di tempat lain, ia boleh memilih mana yang dipanggil. Jadi aku sangat menyarankan keduanya sama.

Hal lain yang bisa dilakukan adalah mendeklarasikan fungsi sebagai \hl{\mbox{\texttt{extern\ inline}}}. Ini akan mencoba meng-inline di file yang sama (demi kecepatan), tapi juga membuat versi dengan external linkage.

\index{inline function specifier@\texttt{inline} function specifier|)}

\subsection{\texorpdfstring{\hl{\mbox{\texttt{noreturn}}} dan \hl{\mbox{\texttt{\_Noreturn}}}}{noreturn dan \_Noreturn}}\label{noreturn}

\index{noreturn function specifier@\texttt{noreturn} function specifier|(} \index{\_Noreturn function specifier@\texttt{\_Noreturn} function specifier|(}

Ini menandakan kepada compiler bahwa suatu fungsi tidak akan pernah kembali ke pemanggilnya, yaitu program akan keluar lewat suatu mekanisme sebelum fungsi itu kembali.

Ini memungkinkan compiler melakukan beberapa optimasi di sekitar pemanggilan fungsi tersebut.

Ini juga memungkinkanmu memberi tahu developer lain bahwa sebagian logika program bergantung pada fungsi yang \emph{tidak} kembali.

Kemungkinan kamu tidak akan pernah memakainya, tapi kamu akan melihatnya pada beberapa pemanggilan library seperti \href{https://beej.us/guide/bgclr/html/split/stdlib.html\#man-exit}{\hl{\mbox{\texttt{exit()}}}}\footnote{https://beej.us/guide/bgclr/html/split/stdlib.html\#man-exit} dan \href{https://beej.us/guide/bgclr/html/split/stdlib.html\#man-abort}{\hl{\mbox{\texttt{abort()}}}}\footnote{https://beej.us/guide/bgclr/html/split/stdlib.html\#man-abort}.

Kata kunci bawaannya adalah \hl{\mbox{\texttt{\_Noreturn}}}, tapi kalau tidak merusak kode yang sudah ada, semua orang akan menyarankan meng-include \hl{\mbox{\texttt{\textless{}stdnoreturn.h\textgreater{}}}} dan memakai \hl{\mbox{\texttt{noreturn}}} yang lebih mudah dibaca.

Ini undefined behavior kalau fungsi yang dinyatakan \hl{\mbox{\texttt{noreturn}}} ternyata kembali. Itu tidak jujur secara komputasi, tahu.

Berikut contoh pemakaian \hl{\mbox{\texttt{noreturn}}} dengan benar:

\begin{tcolorbox}[enhanced jigsaw,left=3pt,breakable,colback=white!97.2549!black,boxrule=0.2mm]

\begin{Shaded}
\begin{Highlighting}[numbers=left,,]
\PreprocessorTok{\#include }\ImportTok{\textless{}stdio.h\textgreater{}}
\PreprocessorTok{\#include }\ImportTok{\textless{}stdlib.h\textgreater{}}
\PreprocessorTok{\#include }\ImportTok{\textless{}stdnoreturn.h\textgreater{}}

\NormalTok{noreturn }\DataTypeTok{void}\NormalTok{ foo}\OperatorTok{(}\DataTypeTok{void}\OperatorTok{)} \CommentTok{// Fungsi ini tidak boleh pernah kembali!}
\OperatorTok{\{}
\NormalTok{    printf}\OperatorTok{(}\StringTok{"Happy days}\SpecialCharTok{\textbackslash{}n}\StringTok{"}\OperatorTok{);}

\NormalTok{    exit}\OperatorTok{(}\DecValTok{1}\OperatorTok{);}            \CommentTok{// Dan ia tidak kembali{-}{-}ia keluar di sini!}
\OperatorTok{\}}

\DataTypeTok{int}\NormalTok{ main}\OperatorTok{(}\DataTypeTok{void}\OperatorTok{)}
\OperatorTok{\{}
\NormalTok{    foo}\OperatorTok{();}
\OperatorTok{\}}
\end{Highlighting}
\end{Shaded}

\end{tcolorbox}

Kalau compiler mendeteksi bahwa fungsi \hl{\mbox{\texttt{noreturn}}} bisa kembali, ia mungkin akan memperingatkanmu, dengan baik hati.

Mengganti fungsi \hl{\mbox{\texttt{foo()}}} dengan ini:

\begin{tcolorbox}[enhanced jigsaw,left=3pt,breakable,colback=white!97.2549!black,boxrule=0.2mm]

\begin{Shaded}
\begin{Highlighting}[]
\NormalTok{noreturn }\DataTypeTok{void}\NormalTok{ foo}\OperatorTok{(}\DataTypeTok{void}\OperatorTok{)}
\OperatorTok{\{}
\NormalTok{    printf}\OperatorTok{(}\StringTok{"Breakin\textquotesingle{} the law}\SpecialCharTok{\textbackslash{}n}\StringTok{"}\OperatorTok{);}
\OperatorTok{\}}
\end{Highlighting}
\end{Shaded}

\end{tcolorbox}

menghasilkan peringatan:

\begin{tcolorbox}[enhanced jigsaw,left=3pt,breakable,colback=white!97.2549!black,boxrule=0.2mm]

\begin{Shaded}
\begin{Highlighting}[]
\NormalTok{foo.c:7:1: warning: function declared \textquotesingle{}noreturn\textquotesingle{} should not return}
\end{Highlighting}
\end{Shaded}

\end{tcolorbox}

\index{noreturn function specifier@\texttt{noreturn} function specifier|)} \index{\_Noreturn function specifier@\texttt{\_Noreturn} function specifier|)} \index{Function specifiers|)}

\section{Alignment Specifier dan Operator}\label{alignment-specifier-dan-operator}

\index{Alignment|(}

\href{https://en.wikipedia.org/wiki/Data_structure_alignment}{\emph{Alignment}}\footnote{https://en.wikipedia.org/wiki/Data\_structure\_alignment} sepenuhnya tentang kelipatan alamat tempat object bisa disimpan. Bisakah ini disimpan di alamat mana pun? Atau harus alamat awal yang habis dibagi 2? Atau 8? Atau 16?

Kalau kamu membuat sesuatu yang low-level seperti memory allocator yang berinteraksi dengan OS-mu, kamu mungkin perlu mempertimbangkan ini. Kebanyakan developer melewati seluruh kariernya tanpa memakai fungsionalitas ini di C.

\subsection{\texorpdfstring{\hl{\mbox{\texttt{alignas}}} dan \hl{\mbox{\texttt{\_Alignas}}}}{alignas dan \_Alignas}}\label{alignas-dan-_alignas}

\index{alignas alignment specifier@\texttt{alignas} alignment specifier|(} \index{\_Alignas alignment specifier@\texttt{\_Alignas} alignment specifier|(}

Ini bukan fungsi. Melainkan \emph{alignment specifier} yang bisa dipakai pada deklarasi variabel.

Specifier bawaannya adalah \hl{\mbox{\texttt{\_Alignas}}}, tapi header \hl{\mbox{\texttt{\textless{}stdalign.h\textgreater{}}}} mendefinisikannya sebagai \hl{\mbox{\texttt{alignas}}} agar lebih sedap dipandang.

Kalau kamu butuh \hl{\mbox{\texttt{char}}}-mu ter-align seperti \hl{\mbox{\texttt{int}}}, kamu bisa memaksanya seperti ini saat mendeklarasikannya:

\begin{tcolorbox}[enhanced jigsaw,left=3pt,breakable,colback=white!97.2549!black,boxrule=0.2mm]

\begin{Shaded}
\begin{Highlighting}[]
\DataTypeTok{char} \KeywordTok{alignas}\OperatorTok{(}\DataTypeTok{int}\OperatorTok{)}\NormalTok{ c}\OperatorTok{;}
\end{Highlighting}
\end{Shaded}

\end{tcolorbox}

Kamu juga bisa mengoper nilai konstan atau ekspresi sebagai alignment. Ini harus sesuatu yang didukung sistem, tapi spesifikasi tidak sampai menentukan nilai apa yang boleh dimasukkan. Pangkat 2 kecil (1, 2, 4, 8, dan 16) umumnya taruhan yang aman.

\begin{tcolorbox}[enhanced jigsaw,left=3pt,breakable,colback=white!97.2549!black,boxrule=0.2mm]

\begin{Shaded}
\begin{Highlighting}[]
\DataTypeTok{char} \KeywordTok{alignas}\OperatorTok{(}\DecValTok{8}\OperatorTok{)}\NormalTok{ c}\OperatorTok{;}   \CommentTok{// align pada batas 8 byte}
\end{Highlighting}
\end{Shaded}

\end{tcolorbox}

Kalau ingin meng-align pada alignment maksimum yang dipakai sistemmu, include \hl{\mbox{\texttt{\textless{}stddef.h\textgreater{}}}} dan pakai tipe \hl{\mbox{\texttt{max\_align\_t}}}, seperti ini:

\begin{tcolorbox}[enhanced jigsaw,left=3pt,breakable,colback=white!97.2549!black,boxrule=0.2mm]

\begin{Shaded}
\begin{Highlighting}[]
\DataTypeTok{char} \KeywordTok{alignas}\OperatorTok{(}\DataTypeTok{max\_align\_t}\OperatorTok{)}\NormalTok{ c}\OperatorTok{;}
\end{Highlighting}
\end{Shaded}

\end{tcolorbox}

Kamu berpotensi melakukan \emph{over-align} dengan menentukan alignment lebih besar dari \hl{\mbox{\texttt{max\_align\_t}}}, tapi apakah hal seperti itu diizinkan tergantung sistem.

\index{alignas alignment specifier@\texttt{alignas} alignment specifier|)} \index{\_Alignas alignment specifier@\texttt{\_Alignas} alignment specifier|)}

\subsection{\texorpdfstring{\hl{\mbox{\texttt{alignof}}} dan \hl{\mbox{\texttt{\_Alignof}}}}{alignof dan \_Alignof}}\label{alignof-dan-_alignof}

\index{alignof operator@\texttt{alignof} operator|(} \index{\_Alignof operator@\texttt{\_Alignof} operator|(}

Operator ini mengembalikan kelipatan alamat yang dipakai tipe tertentu untuk alignment di sistem ini. Misalnya, mungkin \hl{\mbox{\texttt{char}}} di-align setiap 1 alamat, dan \hl{\mbox{\texttt{int}}} di-align setiap 4 alamat.

Operator bawaannya adalah \hl{\mbox{\texttt{\_Alignof}}}, tapi header \hl{\mbox{\texttt{\textless{}stdalign.h\textgreater{}}}} mendefinisikannya sebagai \hl{\mbox{\texttt{alignof}}} kalau kamu ingin tampak lebih keren.

Berikut program yang mencetak alignment berbagai tipe. Sekali lagi, ini akan bervariasi dari satu sistem ke sistem lain. Perhatikan bahwa tipe \hl{\mbox{\texttt{max\_align\_t}}} akan memberi alignment maksimum yang dipakai sistem.

\begin{tcolorbox}[enhanced jigsaw,left=3pt,breakable,colback=white!97.2549!black,boxrule=0.2mm]

\begin{Shaded}
\begin{Highlighting}[numbers=left,,]
\PreprocessorTok{\#include }\ImportTok{\textless{}stdalign.h\textgreater{}}
\PreprocessorTok{\#include }\ImportTok{\textless{}stdio.h\textgreater{}}\PreprocessorTok{     }\CommentTok{// untuk printf()}
\PreprocessorTok{\#include }\ImportTok{\textless{}stddef.h\textgreater{}}\PreprocessorTok{    }\CommentTok{// untuk max\_align\_t}

\KeywordTok{struct}\NormalTok{ t }\OperatorTok{\{}
    \DataTypeTok{int}\NormalTok{ a}\OperatorTok{;}
    \DataTypeTok{char}\NormalTok{ b}\OperatorTok{;}
    \DataTypeTok{float}\NormalTok{ c}\OperatorTok{;}
\OperatorTok{\};}

\DataTypeTok{int}\NormalTok{ main}\OperatorTok{(}\DataTypeTok{void}\OperatorTok{)}
\OperatorTok{\{}
\NormalTok{    printf}\OperatorTok{(}\StringTok{"char       : }\SpecialCharTok{\%zu\textbackslash{}n}\StringTok{"}\OperatorTok{,} \KeywordTok{alignof}\OperatorTok{(}\DataTypeTok{char}\OperatorTok{));}
\NormalTok{    printf}\OperatorTok{(}\StringTok{"short      : }\SpecialCharTok{\%zu\textbackslash{}n}\StringTok{"}\OperatorTok{,} \KeywordTok{alignof}\OperatorTok{(}\DataTypeTok{short}\OperatorTok{));}
\NormalTok{    printf}\OperatorTok{(}\StringTok{"int        : }\SpecialCharTok{\%zu\textbackslash{}n}\StringTok{"}\OperatorTok{,} \KeywordTok{alignof}\OperatorTok{(}\DataTypeTok{int}\OperatorTok{));}
\NormalTok{    printf}\OperatorTok{(}\StringTok{"long       : }\SpecialCharTok{\%zu\textbackslash{}n}\StringTok{"}\OperatorTok{,} \KeywordTok{alignof}\OperatorTok{(}\DataTypeTok{long}\OperatorTok{));}
\NormalTok{    printf}\OperatorTok{(}\StringTok{"long long  : }\SpecialCharTok{\%zu\textbackslash{}n}\StringTok{"}\OperatorTok{,} \KeywordTok{alignof}\OperatorTok{(}\DataTypeTok{long} \DataTypeTok{long}\OperatorTok{));}
\NormalTok{    printf}\OperatorTok{(}\StringTok{"double     : }\SpecialCharTok{\%zu\textbackslash{}n}\StringTok{"}\OperatorTok{,} \KeywordTok{alignof}\OperatorTok{(}\DataTypeTok{double}\OperatorTok{));}
\NormalTok{    printf}\OperatorTok{(}\StringTok{"long double: }\SpecialCharTok{\%zu\textbackslash{}n}\StringTok{"}\OperatorTok{,} \KeywordTok{alignof}\OperatorTok{(}\DataTypeTok{long} \DataTypeTok{double}\OperatorTok{));}
\NormalTok{    printf}\OperatorTok{(}\StringTok{"struct t   : }\SpecialCharTok{\%zu\textbackslash{}n}\StringTok{"}\OperatorTok{,} \KeywordTok{alignof}\OperatorTok{(}\KeywordTok{struct}\NormalTok{ t}\OperatorTok{));}
\NormalTok{    printf}\OperatorTok{(}\StringTok{"max\_align\_t: }\SpecialCharTok{\%zu\textbackslash{}n}\StringTok{"}\OperatorTok{,} \KeywordTok{alignof}\OperatorTok{(}\DataTypeTok{max\_align\_t}\OperatorTok{));}
\OperatorTok{\}}
\end{Highlighting}
\end{Shaded}

\end{tcolorbox}

Output di sistemku:

\begin{tcolorbox}[enhanced jigsaw,left=3pt,breakable,colback=white!97.2549!black,boxrule=0.2mm]

\begin{Shaded}
\begin{Highlighting}[]
\NormalTok{char       : 1}
\NormalTok{short      : 2}
\NormalTok{int        : 4}
\NormalTok{long       : 8}
\NormalTok{long long  : 8}
\NormalTok{double     : 8}
\NormalTok{long double: 16}
\NormalTok{struct t   : 16}
\NormalTok{max\_align\_t: 16}
\end{Highlighting}
\end{Shaded}

\end{tcolorbox}

\index{alignof operator@\texttt{alignof} operator|)} \index{\_Alignof operator@\texttt{\_Alignof} operator|)}

\section{\texorpdfstring{Fungsi \hl{\mbox{\texttt{memalignment()}}}}{Fungsi memalignment()}}\label{fungsi-memalignment}

\index{memalignment() function@\texttt{memalignment()} function|(}

Baru di C23!

(Peringatan: tidak ada compilerku yang mendukung fungsi ini, jadi kodenya sebagian besar belum teruji.)

\hl{\mbox{\texttt{alignof}}} bagus kalau kamu tahu tipe datamu. Tapi bagaimana kalau kamu \emph{sangat tidak tahu} tipenya, dan hanya punya pointer ke data itu?

Bagaimana itu bisa terjadi?

Tentu saja dengan teman baik kita \hl{\mbox{\texttt{void*}}}. Kita tidak bisa mengopernya ke \hl{\mbox{\texttt{alignof}}}, tapi bagaimana kalau kita perlu tahu alignment sesuatu yang ditunjuknya?

Mungkin kita ingin tahu ini kalau akan memakai memori itu untuk sesuatu yang punya kebutuhan alignment yang signifikan. Misalnya, tipe atomic dan floating sering berperilaku buruk kalau tidak ter-align.

Jadi dengan fungsi ini kita bisa memeriksa alignment sebuah data selama kita punya pointer ke data itu, bahkan kalau berupa \hl{\mbox{\texttt{void*}}}.

Mari kita buat tes cepat untuk melihat apakah void pointer ter-align dengan baik untuk dipakai sebagai tipe atomic, dan, kalau ya, kita ambil variabel untuk memakainya sebagai tipe itu:

\begin{tcolorbox}[enhanced jigsaw,left=3pt,breakable,colback=white!97.2549!black,boxrule=0.2mm]

\begin{Shaded}
\begin{Highlighting}[]
\DataTypeTok{void}\NormalTok{ foo}\OperatorTok{(}\DataTypeTok{void} \OperatorTok{*}\NormalTok{p}\OperatorTok{)}
\OperatorTok{\{}
    \ControlFlowTok{if} \OperatorTok{(}\NormalTok{memalignment}\OperatorTok{(}\NormalTok{p}\OperatorTok{)} \OperatorTok{\textgreater{}=} \KeywordTok{alignof}\OperatorTok{(}\NormalTok{atomic }\DataTypeTok{int}\OperatorTok{))} \OperatorTok{\{}
\NormalTok{        atomic }\DataTypeTok{int} \OperatorTok{*}\NormalTok{i }\OperatorTok{=}\NormalTok{ p}\OperatorTok{;}
\NormalTok{        do\_things}\OperatorTok{(}\NormalTok{i}\OperatorTok{);}
    \OperatorTok{\}} \ControlFlowTok{else}
\NormalTok{        puts}\OperatorTok{(}\StringTok{"This pointer is no good as an atomic int}\SpecialCharTok{\textbackslash{}n}\StringTok{"}\OperatorTok{);}

\OperatorTok{...}
\end{Highlighting}
\end{Shaded}

\end{tcolorbox}

Kurasa kamu akan jarang (sampai kemungkinan tidak pernah) memerlukan fungsi ini kecuali kamu mengerjakan hal-hal low-level.

\index{memalignment() function@\texttt{memalignment()} function|)}

Dan begitulah. Alignment!

\index{Alignment|)}

\backmatter
% Add the index
\printindex

\end{document}
